\documentclass[12pt]{book}%

%\usepackage{fontspec}[greek]
%\usepackage{fourier-otf} % pour les arcs (double brève qui compte comme une longue)
%\usepackage{arcs} % pour les arcs (double brève qui compte comme une longue)
\usepackage{graphicx} % pour l'image de la coupe d'Onésime
\usepackage{fontspec}%
\usepackage{calc} % \ratio
\usepackage{fancyhdr}% en-têtes personnalisés
\usepackage[strict]{changepage}%pour modifier les marges des footnotes
\pagestyle{fancy}
\fancyfoot[LE,RO]{\thepage}
\fancyfoot[C]{}
\fancyhead[LE,RO]{}
%\fancyhead[RE]{\leftmark}
%\fancyhead[LO]{\rightmark}
\newcommand{\enTeteG}[1]{\fancyhead[LE]{#1}}% En-tête principale (pages paires de gauche)
\newcommand{\enTeteD}[1]{\fancyhead[RO]{#1}}% En-tête secondaire (pages impaires de droite)
\renewcommand{\headrulewidth}[0]{0pt}%filet de en-tête haut supprimé
\renewcommand{\footrulewidth}[0]{0.05mm}%filet de bas page
%\usepackage{totpages}
%\cfoot{{\thepage} / \ref{TotPages}}
%\usepackage{xunicode}%
%\usepackage{yfonts}%
\usepackage{hyperref}% pour le courriel
\usepackage{amsmath}% pour \overset
\usepackage{cancel}% pour \inRas
\usepackage{quoting}%
%\usepackage{csquotes} % Pour remplacer les guillemets anglais par des français
%\MakeOuterQuote{"}
\usepackage{ulem}% pour souligner en pointillé
\usepackage{comment}%
\usepackage{xcolor}%
\usepackage{fourier}% pour \widearc de la synizèze
%\usepackage{ragged2e}% pour la justification
%scale = {hscale, vscale}
%\usepackage[scale={0.75,0.8}]{geometry}%\usepackage[total={12cm,20cm}%Harmattan
%\usepackage[total={12cm,22cm}, top=3.5cm, right=4.5cm, footnotesep=1cm, includefoot]{geometry}% pour une page impaire
\usepackage[twoside, bindingoffset=1cm,total={12cm,22.5cm}, top=3.5cm, right=4.5cm, footnotesep=1cm, includefoot]{geometry}
%\usepackage{xpatch}
%\setlength{\hoffset}{-18pt}        
%\setlength{\oddsidemargin}{20pt} % Marge gauche sur pages impaires ATTENTION écrit 2 fois
%\setlength{\evensidemargin}{9pt} % Marge gauche sur pages paires
%\setlength{\marginparwidth}{54pt} % Largeur de note dans la marge
%\setlength{\textwidth}{481pt} % Largeur de la zone de texte (17cm)
%\setlength{\voffset}{-18pt} % Bon pour DOS
%\setlength{\marginparsep}{7pt} % Séparation de la marge
%\setlength{\topmargin}{0pt} % Pas de marge en haut
%\setlength{\headheight}{13pt} % Haut de page
%\setlength{\headsep}{10pt} % Entre le haut de page et le texte
%\setlength{\footskip}{27pt} % Bas de page + séparation
%\setlength{\textheight}{708pt} % Hauteur de la zone de texte (25cm)


%\newfontfamily\policeGr{Liberation Serif}[
%  Script=Greek,
%  NFSSFamily=liberationserif,
%
%]
%\newfontfamily\policeGr{Liberation Serif}%

%\newfontfamily{\policeFr}{Garamond Premier Pro}%
%\newfontfamily{\policeLat}{Book Antiqua}%
%\newfontfamily{\policeAn}{Comic Sans MS}%
%\newfontfamily{\policeAll}{Fraktur}%
%\newfontfamily{\policeIt}{BodoniFLF}%
%\newfontfamily{\policeEsp}{Great Moravia}%
% ...
%\setmainfont{Liberation Serif}%
%\setmainfont{Liberation Serif}%
%\newfontfamily{\policeFr}{Garamond Premier Pro}%
%\newfontfamily{\policeFr}{EB Garamond}%
\IfFontExistsTF{GFS Bodoni}{\newfontfamily{\policeFr}{GFS Bodoni}\newfontfamily{\policeGr}{GFS Bodoni}\setmainfont{GFS Bodoni}}{\newfontfamily{\policeFr}{GFS Didot}\newfontfamily{\policeGr}{GFS Didot}\setmainfont{GFS Didot}}%
%\newfontfamily{\policeLat}{Book Antiqua}%
\newfontfamily{\policeLat}{GFS Neohellenic}%
%\newfontfamily{\policeAn}{Comic Sans MS}%
\IfFontExistsTF{Courier 10 Pitch}{\newfontfamily{\policeAn}{Courier 10 Pitch}}{\newfontfamily{\policeAn}{FreeMono}}%
\newfontfamily{\policeEsp}{Abyssinica SIL}%
%\newfontfamily{\policeAll}{Fraktur}%
%\newfontfamily{\policeAll}{fraktur-bt.ttf}%
\IfFontExistsTF{Bertholdr Mainzer Fraktur}{\newfontfamily{\policeAll}{Bertholdr Mainzer Fraktur}}{\usepackage{yfonts}}%
%\newfontfamily{\policeIt}{BodoniFLF}%
\newfontfamily{\policeIt}{GFS Artemisia}%
%\newfontfamily{\policeEsp}{Great Moravia}%
%\newfontfamily\policeGr{GFSBodoni.otf}[BoldFont = GFSBodoniBold.otf, ItalicFont = GFSBodoniIt.otf, BoldItalicFont = GFSBodoniBoldIt.otf]%
%\setmainfont{GFSBodoni.otf}[BoldFont = GFSBodoniBold.otf, ItalicFont = GFSBodoniIt.otf, BoldItalicFont = GFSBodoniBoldIt.otf]

%\setdefaultlanguage[variant=ancient]{greek}%
\usepackage{polyglossia}%
\setdefaultlanguage{french}%
%\setotherlanguage[variant=ancient]{greek}%
%\setmainlanguage[variant=ancient]{greek} 
%\setotherlanguage[variant=ancient]{greek} 
%\setmainfont{Libertinus Serif} 
%\newfontfamily{\greekfont}{Liberation Serif}[
%  Script=Greek,
%  NFSSFamily=liberationserif,
%]

\usepackage[maxlevel=3]{csquotes}%
%\MakeOuterQuote{"}% pour les guillemets français
%\usepackage{verse}%
%\usepackage}[squarebr]{gmverse}% pour le rejet avec un crochet droit
\usepackage[parapparatus]{reledmac}%
\usepackage[parledgroup, nomaxlines]{reledpar}% 
%\usepackage[parledgroup, nomaxlines, nopbinverse]{reledpar}% nopbinverse empêche de casser un groupe
%\usepackage[parledgroup]{reledpar}% nomaxlines permettait d'avoir un nombre de lignes différents sur des pages qui se font face.
%\parledgrouptrue
%\setgoalfraction{0.9}% sept 2022
%\makeatletter
%\xpatchcmd{\mp@append@Xnotes}{\@parboxrestore}{\@parboxrestore\leftskip-1in}{}{}
%\makeatother
% Pour que les commentaires apparaissent sur les deux pages
\AtBeginDocument{%

%\Xbeforeinserting[B]{\espace{-2cm}}% code exécuté avant l'appel à la note
%\Xafternote[B]{\espace{20cm}}
%\Xbhookgroup[A,C]{\espace{-2cm}}
%\Xprenotes{2cm}
%\setlength{\skip\footins}{-2cm}
%\Xhangindent[B]{3cm}
%\sethangingsymbol{\protect\hfill}
%\versehangrightsquare% le rejet d'un vers se fait avec un crochet droit
\rightnoteupfalse 
\Xmaxhnotes[A,B]{0.50\textheight}%
\Xmaxhnotes[C]{0.22\textheight}%
%\Xmaxhnotes{0.9\textheight}%
%\maxhnotesX{0.45\textheight}%
\setRlineflag{}%
%\newcommand{\decoupeVers}[1]{\linenumannotation{#1}}%
\newcommand{\decoupeVers}[1]{}%
%\setlength{\oddsidemargin}{100pt} % Marge gauche sur pages impaires ATTENTION écrit 2 fois
% La Largeur du commentaire
%\Xwidth{\columnwidth+\marginparsep+\ledrsnotewidth}
% \Xgroupbyline[A,C]
\Xgroupbyline[A,C]
% \Xonlyside[A]{R}% imprime les testimonia page de droite
%\Xafterrule[A]{-1em}% supprime un peu d'esp vertical
%\Xwidth{\columnwidth+\marginparsep+\ledrsnotewidth}
\Xwidth{14cm}
%\Xsethangindent{20pt}%décale les footnotes
%\Xhangindent{-3cm}
%\Xparindent{-10cm}
%\sethangindentX{20pt}
%\Xhangindent = 20pt%
}%


\newcommand\Petit{\fontsize{9pt}{11pt}\selectfont} % x 1,2 pour e second argument qui est l'espacement interligne
\newcommand{\Taille}[1]{\fontsize{#1pt}{12pt}\selectfont}
\newcommand{\TaillePt}[1]{\fontsize{#1}{15pt}\selectfont}
\newlength{\LargAuto}%
\newlength{\LargTheor}%
\newlength{\Vrai}%
\newlength{\TP}%
\setlength{\TP}{12pt}%
\newcommand{\Automatique}[1]{%
\setlength{\LargTheor}{-72pt}%
\addtolength{\LargTheor}{\textwidth}%
\settowidth{\LargAuto}{#1}%
%\divide \LargTheor by \LargAuto%  \real{1.} *
\setlength{\Vrai}{\TP * \ratio{\LargTheor}{\LargAuto}}%
%\edef\TP{12* \ratio{\LargTheor}{\LargAuto}}%
{\TaillePt{\Vrai}#1}%
}%
\newlength{\largtext}%
\newcommand{\Strophe}[1]{\ledleftnote{%
\settowidth{\textelargeur}{\footnotesize{\francais{#1}}}%
%\espace{-\textelargeur}
\begin{minipage}{\textelargeur}\footnotesize{\francais{#1}}\end{minipage}}}%
\newcommand{\renverse}[1]{\rotatebox[origin=c]{180}{#1}}
\newcommand{\italique}[1]{\textit{#1}}%
\newcommand{\citComm}[1]{\begin{quoting}[vskip=0pt] #1 \end{quoting}}% vskip = \topsep + \partopsep par défaut
\newcommand{\passeLigne}[0]{\skipnumbering}%
\newcommand{\fin}[0]{\linebreak\passeLigne\protect\hfill $[$ }% Pour le rejet 
\newcommand{\ligne}[1]{\setline{#1}\hidenumbering}%
\newcommand{\espace}[1]{\hspace{#1}}%
\newcommand{\cesure}[0]{ [...] }%
\newcommand{\nieme}[1]{#1\italique{e}}%
\newcommand{\lettreBreve}[1]{$\breve{\mbox{#1}}$}
\newcommand{\synizese}[1]{$\widearc{\mbox{#1}}$}
\newcommand{\inRas}[2]{\if\relax\detokenize{#2}\relax\text{\sout{#1}}\else$\overset{\text{#2}}{\text{\sout{#1}}}$\fi}
\newcommand{\auDessus}[2]{$\overset{\text{#2}}{\text{#1}}$}
%\newcommand{\etc}[0]{κτλ}%
%\newcommand{\cit}[0]{\espace{1cm}} % indent. pour une citation
\newcommand{\expo}[1]{\textsuperscript{#1}}%
\newcommand{\alinea}[0]{ \ \ }%
\newcommand{\Fr}[0]{ Ar. \italique{fragm.} (Kassel-Austin) }%
\newcommand{\Syn}[0]{ Συναγωγή }%
\newcommand{\SynB}[0]{ Συναγωγή\expo{B} }% la deuxième version pour la lettre A
\newcommand{\ClemAlex}{ Clemens Alexandrinus }
\newcommand{\Eustathe}{ Eustathius }
\newcommand{\AelDion}{ Aelius Dionysius }
\newcommand{\EtymMagnum}{ Etymologicum magnum }
\newcommand{\EtymGen}{ Etymologicum genuinum }
\newcommand{\EtymSym}{ Etymologicum Symeonis }
\newcommand{\Apostol}{ Michael Apostolius }
\newcommand{\Macaire}{ Macarius } 
\newcommand{\Phryn}{ Phrynichus \italique{Praep. soph.} }
\newcommand{\Stobee}{ Joannes Stobaeus }
\newcommand{\EpimAlphab}{ Epimerismi alphabetici in Homerum (ed. Cramer) }
\newcommand{\GregCor}{ Gregorius Corinthius }
\newcommand{\Antiatt}{ Antiatticista (ed. Bekker) }
\newcommand{\Harp}{ Harpocration }% Autre édition de référence : Dindorf
\newcommand{\Rpc}[0]{R\expo{pc}}%
\newcommand{\Rac}[0]{R\expo{ac}}%
\newcommand{\Rdeux}[0]{R\expo{2}}%
\newcommand{\Mon}[0]{ M }%
\newcommand{\PsiUn}[0]{PSI-1194}
\newcommand{\PsiDeux}[0]{POxy-3839}
\newcommand{\PsiTrois}[0]{POxy-3840}
\newcommand{\PsiQuatre}[0]{POxy-4935}
\newcommand{\PsiCinq}[0]{POxy-1176}%P50

\newcommand{\francais}[1]{{\policeFr #1}}%
\newcommand{\latin}[1]{{\policeLat #1}}%
\IfFontExistsTF{Bertholdr Mainzer Fraktur}{\newcommand{\allemand}[1]{{\policeAll #1}}}{\newcommand{\allemand}[1]{{\textfrak{#1}}}}%
%\newcommand{\allemand}[1]{{\textfrak{#1}}}%
%\newcommand{\grec}[1]{{\policeGr #1}}%
\newcommand{\grec}[1]{#1}%
\newcommand{\anglais}[1]{{\policeAn #1}}%
\newcommand{\italien}[1]{{\policeIt #1}}%
\newcommand{\espagnol}[1]{{\policeEsp #1}}%

\newcommand{\Nauck}[0]{N\expo{2}}%
\newcommand*{\auteur}[1]{\textsc{#1}*}%
\newcommand{\cf}[0]{\italique{cf.}\xspace}%
\newcommand{\chap}[1]{\^{#1}}%
%\newcommand{\chap}[1]{$\hat{#1}$}


%\marginparwidth = 1.2\marginparwidth
%\ledrsnotewidth = 1.2\marginparwidth
\newlength{\larg}%
\newcommand{\decalage}[1]{%
\settowidth{\larg}{#1}%
\espace{\larg}%
}
\newcommand{\decalageN}[1]{%
\settowidth{\larg}{#1}%
\espace{-\larg}%
}
%\newcommand{\apparat}[2]{\ledrightnote{\espace{-1.2\marginparwidth}\begin{minipage}{1.4\marginparwidth}\footnotesize{\textbf{#1} #2}\end{minipage}}}%
\newcommand{\apparat}[2]{\ledrightnote{\espace{-2cm}\begin{minipage}{4.5cm}\footnotesize{\textbf{#1} #2}\end{minipage}}}%
%\newcommand{\lecon}[1]{\dotuline{#1}} % la leçon retenue de l'apparat critiqe
\newcommand{\lecon}[1]{\textbf{#1}}
%\newcommand{\douteux}[1]{\dotuline{#1}} % pour les lettres douteuses : un point en-dessous
\newcommand{\douteux}[1]{$\underset{\cdot}{\text{#1}}$}
\newcommand{\difference}[1]{\dotuline{#1}}


%\setsidenotesep{ sep }%
\newcommand{\espApp}[0]{ || }%\espace{10pt}%
\newcommand{\espTest}[0]{ || }%\espace{10pt}%
\newcommand{\ligneApp}[1]{ \ \ \textbf{#1} }%
%% 3 arguments : les numéros de ligne, le lemme et la note elle-même



%\newcommand{\commentaire}[3]{\edtext{}{\Bfootnote[nonum, nosep]{\espace{-2cm}\begin{minipage}{16cm}\espace{-1cm}\textbf{#1~~#2} #3\end{minipage}}}}%
\newcommand{\commentaire}[3]{\edtext{}{\Bfootnote[nonum,nosep]{\begin{adjustwidth}{-1.5cm}{-2cm}\espace{-1.5cm}\textbf{#1}\decalageN{\textbf{#1}}\espace{0pt}\espace{1.5cm}\textbf{#2} #3\end{adjustwidth}}}}% Avec une * : adjustwidth* pour la symétrie des pages
%\renewcommand{\footfudgefiddle}{68}%Note that you must call it before 
% \Xarrangement[C]{paragraph}% entrâine trop d'erreurs, cf. documentation
\newcommand{\debutNote}[1]{\edtext{}{\Cfootnote[nonum, nosep]{\espace{-1cm}\espace{-15pt}\textbf{\textbf{#1}}}}}%
%\newcommand{\debutNote}[1]{\edtext{}{\Cfootnote[nonum, nosep]{\espace{-1cm}\espace{-15pt}\textbf{\dotuline{\textbf{#1}}}}}}%
\newcommand{\note}[2]{\edtext{}{\Cfootnote[nonum, nosep]{\espace{-15pt}\textbf{#1} #2}}}%
\newcommand{\testimonia}[2]{\edtext{}{\Afootnote[nonum, nosep]{\espace{-15pt}\textbf{#1} #2}}}%
%\newcommand{\testimonia}[2]{}
\newcommand{\marqueur}[1]{\textsuperscript{#1}}%
%\newcommand{\cat}[0]{$""_"^"$}%
\newcommand{\cat}[0]{\texorpdfstring{\textsubscript{\textasciicircum}}{cat}}%
%\fnpos{critical-familiar}%
%\mpfnpos{familiar-critical}%

\newcommand{\eff}[1]{%
{\color{orange}#1}%
}%

\newlength{\textlarg}%
\newcommand{\typeVers}[1]{%
\settowidth{\textlarg}{\footnotesize{#1}}%
\hspace{-0.8cm}{\footnotesize{#1}}\hspace{-\textlarg}\hspace{0.8cm}}%
%\newlength{\largtext}%
%\newcommand{\strophe}[1]{%
%\settowidth{\largtext}{~~\footnotesize{\francais{#1}}}%
%~~\footnotesize{\francais{#1}}%\hspace{-\largtext}%
%}%
\newcommand{\strophe}[1]{\footnotesize{~\francais{#1}}}%
%\newcommand{\strophe}[1]{%
%\settowidth{\largtext}{\footnotesize{\francais{#1}}}%
%	\ledleftnote{%
		%\hspace{-0.5cm}\footnotesize{\francais{#1}}\hspace{\largtext}\hspace{0.5cm%}%
%		\hspace{-0.5cm}\footnotesize{\francais{#1}}\hspace{\largtext}\hspace{0.5cm}%
%	}%
%}%
%\newcommand{\personnageGrec}[1]{\ledleftnote{\textbf{\large{#1}}\espace{-1.5cm}}}%
\newlength{\textelargeur}%
\newcommand{\personnageGrec}[1]{\ledleftnote{%
%\settowidth{\textelargeur}{\large{\textbf{#1}}}%
\settowidth{\textelargeur}{\large{\textbf{Στεφανόπωλις }}}%
%\espace{-\textelargeur}
\begin{minipage}{\textelargeur}\begin{flushright}\large{\textbf{#1}}\end{flushright}\end{minipage}}}%
%\newcommand{\Strophe}[1]{\footnotesize{\francais{#1}}}%
%\newcommand{\Strophe}[1]{\ledleftnote{%
%\settowidth{\textelargeur}{\footnotesize{\francais{#1}}}%
%\espace{-\textelargeur}
%\begin{minipage}{\textelargeur}\footnotesize{\francais{#1}}\end{minipage}}}%
%\newcommand{\personnageFr}[1]{\ledleftnote{\textbf{\large{\francais{#1}}}\espace{-1.5cm}}}%
\newlength{\texto}%
\newcommand{\personnageFr}[1]{\ledleftnote{%
%\settowidth{\texto}{\large{\textbf{#1}}}%
\settowidth{\texto}{\large{\textbf{Marchande d}}}%
%\espace{-\textelargeur}
\begin{minipage}{\texto}\begin{flushright}\large{\textbf{#1}}\end{flushright}\end{minipage}}}%
%% Les différents types de vers
\newcommand{\trimIamb}[1]{%
\hangindent=20pt%
\espace{0pt}% la valeur précédente -20
#1%
}%
\newcommand{\trimIambFr}[1]{%
\hangindent=20pt%
\espace{0pt}% la valeur précédente -20
\francais{#1%
}}%
\newcommand{\trimIambCat}[1]{%
\hangindent=53pt%
\espace{33pt}% la valeur précédente -20
#1%
}%
\newcommand{\trimIambCatFr}[1]{%
\hangindent=53pt%
\espace{33pt}% la valeur précédente -20
\francais{#1%
}}%
\newcommand{\tetramIamb}[1]{%
\hangindent=10pt%
\espace{-10pt}% la valeur précédente -20
#1%
}%
\newcommand{\tetramIambFr}[1]{%
\hangindent=20pt%
\espace{0pt}% la valeur précédente -20
\francais{#1%
}}%
\newcommand{\tetramTroch}[1]{%
\hangindent=20pt%
\espace{0pt}% la valeur précédente -20
#1%
}%
\newcommand{\tetramTrochFr}[1]{%
\hangindent=20pt%
\espace{0pt}% la valeur précédente -20
\francais{#1%
}}%
\newcommand{\tetramIambCat}[1]{%
\hangindent=10pt%
\espace{-10pt}% la valeur précédente -20
#1%
}%
\newcommand{\tetramIambCatFr}[1]{%
\hangindent=20pt%
\espace{0pt}% la valeur précédente -20
\francais{#1%
}}%
\newcommand{\dimAnap}[1]{%
\hangindent=70pt%
\espace{50pt}% la valeur précédente -20
#1%
}%
\newcommand{\dimAnapFr}[1]{%
\hangindent=70pt%
\espace{50pt}% la valeur précédente -20
\francais{#1%
}}%
\newcommand{\anap}[1]{%
\hangindent=70pt%
\espace{50pt}% la valeur précédente -20
#1%
}%
\newcommand{\anapFr}[1]{%
\hangindent=70pt%
\espace{50pt}% la valeur précédente -20
\francais{#1%
}}%
% 2 dochmiaques
\newcommand{\diDoch}[1]{%
\hangindent=60pt%
\espace{40pt}% la valeur précédente -20
#1%
}%
\newcommand{\diDochFr}[1]{%
\hangindent=60pt%
\espace{40pt}% la valeur précédente -20
\francais{#1%
}}%
\newcommand{\tetramAnap}[1]{%
\hangindent=20pt%
\espace{0pt}% la valeur précédente -20
#1%
}%
\newcommand{\tetramAnapFr}[1]{%
\hangindent=20pt%
\espace{0pt}% la valeur précédente -20
\francais{#1%
}}%
\newcommand{\chantAgath}[1]{%
\hangindent=90pt%
\espace{70pt}% la valeur précédente -20
#1%
}%
\newcommand{\chantAgathFr}[1]{%
\hangindent=90pt%
\espace{70pt}% la valeur précédente -20
\francais{#1%
}}%
%% Pour les "vers" 295 à 311
\newcommand{\prose}[1]{%
\hangindent=30pt%
\espace{10pt}% la valeur précédente -20
#1%
}%
\newcommand{\proseFr}[1]{%
\hangindent=30pt%
\espace{10pt}% la valeur précédente -20
\francais{#1%
}}%
%% pour les vers courts de structure changeante dans un choeur
\newcommand{\chantChoeur}[1]{%
%{\centering
%	#1\par
%}%
\hangindent=50pt%
\espace{30pt}% la valeur précédente -20
#1%
}%
\newcommand{\chantChoeurFr}[1]{%
%{\centering
%	\francais{#1}\par
%}%
\hangindent=50pt%
\espace{30pt}% la valeur précédente -20
\francais{#1%
}}%
%% pour les vers 1015 à 1055
\newcommand{\paratragiques}[1]{%
%{\centering
%	#1\par
%}%
\hangindent=50pt%
\espace{30pt}% la valeur précédente -20
#1%
}%
\newcommand{\paratragiquesFr}[1]{%
%{\centering
%	\francais{#1}\par
%}%
\hangindent=30pt%
\espace{10pt}% la valeur précédente -20
\francais{#1%
}}%
% pour indiquer que les notes marginales s'alignent au niveau de la première ligne
%\leftnoteupfalse
%\rightnoteupfalse



\newenvironment{groupe}{%
\pstart%
\ledgroup%
}{%
\pend %
}%

\newcommand{\LSJ}[0]{LSJ}%

%% pour voir tous les numéros de ligne (version brouillonne)
\linenumincrement{1}%
\firstlinenum{1}%
%% Lignes à imprimer
\gdef\linenumberlist{2,6,10,15,20,24,30,36,40,45,50,55,60,65,69,74,80,87,91,94,102,105,110,115,120,125,131,135,140,145,150,155,160,165,170,174,180,185,190,196,201,205,215,220,225,229,234,239,245,250,261,265,270,276,280,285,290,295,300,305,310,315,320,325,330,335,340,345,350,355,360,365,370,375,379,385,390,395,400,405,410,415,420,425,430,435,440,445,450,455,460,465,470,475,480,485,490,495,499,501,505,510,515,521,526,530,535,541,545,550,556,561,565,572,575,580,585,591,596,600,609,614,620,624,629,635,640,650,656,660,665,670,675,680,685,690,695,700,708,713,720,725,729,735,741,745,750,755,761,766,770,775,780,786,790,795,800,805,810,815,820,825,831,835,840,845,850,856,860,865,870,875,880,884,890,894,900,905,915,921,930,935,940,946,950,955,960,965,970,975,980,985,990,996,1000,1010,1016,1020,1025,1030,1035,1040,1045,1050,1055,1060,1066,1071,1076,1100,1106,1116,1120,1124,1130,1141,1147,1150,1156,1161,1166,1171,1175,1180,1185,1190,1195,1200,1205,1211,1215,1221,1224,1230}



\usepackage{adjustbox}
\usepackage{array}
\usepackage{multicol}
\setlength{\columnsep}{9pt}
\setlength{\multicolsep}{1pt}
\newsavebox{\commMeasureBox}
\newcommand{\renderCommBox}[2]{%
  \sbox{\commMeasureBox}{#2}%
  \typeout{COMMBOX_DIM_#1: ht=\the\dimexpr\ht\commMeasureBox+\dp\commMeasureBox\relax}%
  \begin{adjustbox}{max totalheight=9.8cm, max width=16.5cm, keepaspectratio}%
  \usebox{\commMeasureBox}%
  \end{adjustbox}%
}
\setmonofont{FreeMono}
\setlength{\headheight}{16pt}
\renewcommand{\enTeteG}[1]{\fancyhead[LE]{\small #1}}
\renewcommand{\enTeteD}[1]{\fancyhead[RO]{\small #1}}

% In Spread 2, critical apparatus and testimonia are already displayed on Spread 1 Left Page:
\renewcommand{\apparat}[2]{}
\renewcommand{\testimonia}[2]{}
\renewcommand{\fin}[0]{}
\renewcommand{\Automatique}[1]{#1}
\renewcommand{\allemand}[1]{\textit{#1}}
\makeatletter
\renewcommand*{\led@err@LeftOnRightPage}{}
\renewcommand*{\led@err@RightOnLeftPage}{}
\makeatother
\AtBeginDocument{%
  \Xmaxhnotes[A,B]{0.82\textheight}%
  \Xmaxhnotes[C]{0.55\textheight}%
}
\renewcommand{\commentaire}[3]{\edtext{}{\Bfootnote[nonum,nosep]{\begin{adjustwidth}{-1.5cm}{-2cm}\fontsize{8.0pt}{9.4pt}\selectfont\espace{-1.5cm}\textbf{#1}\decalageN{\textbf{#1}}\espace{0pt}\espace{1.5cm}\textbf{#2} #3\end{adjustwidth}}}}%
\renewcommand{\debutNote}[1]{\edtext{}{\Cfootnote[nonum, nosep]{\fontsize{7.8pt}{9.2pt}\selectfont\espace{-1cm}\espace{-15pt}\textbf{#1}}}}%
\renewcommand{\note}[2]{\edtext{}{\Cfootnote[nonum, nosep]{\fontsize{7.8pt}{9.2pt}\selectfont\espace{-15pt}\textbf{#1} #2}}}%

% Row macro for Spread 1 Left Page (Diplomatic-yet-corrected Ravennas + inline apparatus)
\newcommand{\ligneRav}[4]{%
  \noindent\hbox to 16.5cm{%
    \hbox to 0.55cm{\hfill\tiny\texttt{#1}\hspace{0.1cm}}%
    \hbox to 1.25cm{\hfill\scriptsize #2\hspace{0.15cm}}%
    \hbox to 8.1cm{\adjustbox{max width=8.0cm}{#3}\hfill}%
    \hspace{0.2cm}%
    \parbox[t]{6.4cm}{\raggedright\fontsize{6.7pt}{7.9pt}\selectfont #4}%
  }\par\vspace{0.35pt}%
}

% Row macro for Spread 2 in the Strict 4-Page-Per-Folio Edition (Established Greek / French)
\newcommand{\ligneEst}[3]{%
  \noindent\hbox to 16.5cm{%
    \hbox to 2.5cm{\hfill\scriptsize\adjustbox{max width=2.35cm}{#1}\hspace{0.15cm}}%
    \hbox to 13.2cm{\adjustbox{max width=13.1cm}{#2}\hfill}%
    \hbox to 0.8cm{\hfill\scriptsize\texttt{#3}}%
  }\par\vspace{0.3pt}%
}

\begin{document}
\thispagestyle{empty}
\noindent\begin{adjustbox}{max totalheight=19.5cm, max width=\textwidth, keepaspectratio}\begin{minipage}{\textwidth}\begin{center}
\vspace*{0.4cm}
{\Huge\textbf{ΑΡΙΣΤΟΦΑΝΟΥΣ ΘΕΣΜΟΦΟΡΙΑΖΟΥΣΑΙ}}\par\vspace{0.5cm}
{\LARGE\textbf{ARISTOPHANE --- LES THESMOPHORIES}}\par\vspace{0.8cm}
{\Large\textit{Nouvelle édition philologique et paléographique sur quatre pages}}\par\vspace{0.4cm}
{\large Édition synoptique stricte (4 pages par folio : Transcription R + Fac-similé / Grec établi + Français + Commentaire)}\par\vspace{0.5cm}
\hrule height 0.8pt\vspace{0.6cm}
{\normalsize
\textbf{Architecture des doubles pages (30 folios du \textit{Codex Ravennas} 429, \texttt{ff. 162v–177r}, vv. 1–1231) :}\par\vspace{0.3cm}
\begin{itemize}
  \item \textbf{Première double page — Page de gauche (pair) :} Transcription diplomatique corrigée du \textit{Codex Ravennas} 429 (\textbf{R}), avec \textit{notae personarum} tachygraphiques, \textit{paragraphoi}, \textit{dikola}, grattages (\textit{in rasura}), ajouts marginaux, les quatre \textit{parepigraphai} scéniques, soulignement pointillé des émendations (\dotuline{texte}), \textbf{apparat critique dans le prolongement de la ligne}, et en bas de page : (1) \textbf{\textit{loci paralleli}} et (2) \textbf{suite (en commentaire) de l'apparat critique détaillé}.
  \item \textbf{Première double page — Page de droite (impair) :} Fac-similé haute résolution en pleine page du folio correspondant du \textit{Codex Ravennas} 429 (Biblioteca Classense, Ravenne).
  \item \textbf{Seconde double page — Page de gauche (pair) :} Texte grec établi proprement, avec colométrie et attribution complète des répliques.
  \item \textbf{Seconde double page — Page de droite (impair) :} Traduction française en regard.
  \item \textbf{Bas de la seconde double page (gauche et droite) :} Long commentaire philologique, historique, métrique et scénique.
\end{itemize}
}\vspace{0.6cm}
\hrule height 0.8pt\vspace{0.5cm}
{\Large\textbf{Édition, traduction et commentaire par Louis HAUSEUX (\textit{Bibi})}}\par
\end{center}\end{minipage}\end{adjustbox}
\newpage

\clearpage\ifodd\value{page}\hbox{}\newpage\fi
\enTeteG{ΠΡΟΛΟΓΟΣ --- Codex Ravennas 429, f. 162v (Classense 164v) --- vv. 1–42}
\begin{adjustwidth}{-1.5cm}{-3.0cm}
\begin{adjustbox}{max totalheight=21.0cm, max width=16.5cm, keepaspectratio}
\begin{minipage}{16.5cm}
{\fontsize{7.2pt}{8.5pt}\selectfont
\noindent\textbf{Folio 01 --- Codex Ravennas 429, \texttt{f. 162v} (\texttt{Classense 164v}) : ΠΡΟΛΟΓΟΣ, vv. 1–42}\hfill \textit{Témoins :} $\text{R}^1, \Rpc, \Rdeux$ (\Mon)\par\vspace{2pt}
\hrule height 0.4pt\vspace{2.5pt}
\noindent\textbf{\scriptsize [$\Rdeux$ — Notice liminaire (\italique{Hypothesis}) en semi-onciale de glose en tête du \texttt{f. 162v} (sans liste des personnages) :]}\par\vspace{0.5pt}
{\scriptsize\noindent\texttt{1.} ὁ χορὸς ἐκ θεσμοφοριαζουσῶν $\cdot$ καὶ τοῦτο τὸ δρᾶμα τῶν κατ' Εὐριπίδου πεποιημένων $\cdot$ ἀπὸ δὲ τῶν θεσμο\\
\texttt{2.} φορίων, τὰς θεσμοφοριαζούσας γυναῖκας ἐπέγραψεν, ἀφ' ὧν καὶ ὁ χορ(ός) $\cdot$ Εὐριπίδου γυνὴ μὲν Χοιρίνη,\\
\texttt{3.} μή(τ)ηρ δὲ Κλειτώ :$\sim$\hspace{1.5em}προλογίζει Μνησίλοχ(ος) ὁ κηδεστὴς Εὐριπίδου $\div$}\par\vspace{1.5pt}
\centerline{\footnotesize\textbf{$\ast\;\cdot\;\ast$ ἀριστοφάνους θεσμοφοριάζουσαι —}}\vspace{1.5pt}
\hrule height 0.3pt\vspace{2pt}
}
{\fontsize{8.1pt}{9.5pt}\selectfont
\ligneRav{1}{μ\auDessus{ν}{Η}}{ὦ Ζεῦ \difference{Χελιδὼν} \difference{ἆρά ποτε} φανήσεται;}{\textbf{1} ἄρα πότε R : \lecon{ἆρά ποτε} Biset}
\ligneRav{2}{}{ἀπολεῖ μ᾽ ἀλοῶν \difference{ἅνθρωπος} ἐξ ἑωθινοῦ·}{\textbf{2} ἄνθρωπος R : \lecon{ἅνθρωπος} Boissonade}
\ligneRav{3}{}{οἷόν \difference{τε}, πρὶν τὸν σπλῆνα κομιδῇ μ᾽ ἐκβαλεῖν \hfill\textbf{[β \Rpc]}}{\textbf{3} τι R : \lecon{τε} Σ (δυνατόν), Suid. οι 150 \espApp vv. 3–4 inv. \Rac{} : litt. $\beta', \alpha'$ corr. \Rpc{}}
\ligneRav{4}{}{παρὰ σοῦ πυθέσθαι ποῖ μ᾽ ἄγεις ωὖριπίδη; \hfill\textbf{[α \Rpc]}}{}
\ligneRav{5}{\auDessus{εὐ}{ρ}}{ἀλλ᾽ οὐκ ἀκούειν δεῖ σε πάνθ᾽ \difference{ὅσ᾽} αὐτίκα}{\textbf{5} ὅσ R : \lecon{ὅσ᾽} \Mon}
\ligneRav{}{—}{ὄψει παρεστώς. \auDessus{:}{\tiny μν} πῶς λέγεις; αὖθις φράσον·}{}
\ligneRav{}{—}{οὐ δεῖ μ᾽ ἀκούειν; \auDessus{:}{\tiny εὐ} οὐχ ἅ γ᾽ ἂν μέλλῃς ὁρᾶν·}{}
\ligneRav{}{μ\auDessus{ν}{Η}}{οὐδ᾽ ἆρ᾽ ὁρᾶν δεῖ μ᾽; \auDessus{:}{\tiny εὐ} οὐχ ἅ γ᾽ ἂν ἀκούειν δέῃ·}{}
\ligneRav{}{μ\auDessus{ν}{Η}}{πῶς μοι παραινεῖς; δεξιῶς μέντοι λέγεις·}{}
\ligneRav{10}{}{οὐ φῂς σὺ χρῆναί μ᾽ οὔτ᾽ ἀκούειν οὔθ᾽ ὁρᾶν·}{\textbf{10} \lecon{ὁρᾶν.} R : ὁρᾶν; Gelenius}
\ligneRav{}{\auDessus{εὐ}{ρ}}{χωρὶς γὰρ αὐτοῖν ἑκατέρου 'στὶν ἡ φύσις·}{}
\ligneRav{}{μ\auDessus{ν}{Η}}{τοῦ μήτ᾽ ἀκούειν μήθ᾽ ὁρᾶν; \auDessus{:}{\tiny εὐ} εὖ ἴσθ᾽ ὅτι·}{}
\ligneRav{}{μ\auDessus{ν}{Η}}{πῶς χωρίς; \auDessus{:}{\tiny εὐ} οὕτω ταῦτα διεκρίθη τότε·}{}
\ligneRav{}{}{αἰθὴρ γὰρ ὅτε τὰ πρῶτα διεχωρίζετο}{}
\ligneRav{15}{}{καὶ ζᾦ᾽ ἐν αὑτῷ ξυνετέκνου κινούμενα,}{}
\ligneRav{16}{}{ᾧ μὲν βλέπειν χρὴ \difference{πρῶτ᾽} ἐμηχανήσατο}{\textbf{16} πρῶτα R : \lecon{πρῶτ᾽} Grynaeus}
\ligneRav{}{}{ὀφθαλμὸν ἀντίμιμον ἡλίου τροχῷ,}{}
\ligneRav{18}{}{ἀκοὴν δὲ χοάνης ὦτα διετετρήνατο·}{\textbf{18} \lecon{ἀκοὴν δὲ χοάνης} R, Willems : ἀκοὴς δὲ χοάνην Biset : ἀκοῇ δὲ χοάνην Boissonade}
\ligneRav{}{μ\auDessus{ν}{Η}}{διὰ τὴν χοάνην οὖν μήτ᾽ ἀκούω μήθ᾽ ὁρῶ;}{}
\ligneRav{20}{}{νὴ τὸν Δί᾽ ἥδομαί γε τουτὶ προσμαθών·}{}
\ligneRav{21}{}{οἷόν \difference{γέ} που ᾽στιν αἱ σοφαὶ ξυνουσίαι·}{\textbf{21} τε R : \lecon{γε} Σ}
\ligneRav{}{\auDessus{εὐ}{ρ}}{πόλλ᾽ ἂν μάθοις τοιαῦτα παρ᾽ ἐμοῦ. \auDessus{:}{\tiny μν} πῶς ἂν οὖν}{}
\ligneRav{(23)}{}{$[\![\text{\sout{πρὸς τοῖς ἀγαθοῖς τούτοισιν ἐξεύροιμ᾽ ὅπως}}]\!]$}{\textbf{v. 23} πρὸς τοῖς ἀγαθοῖς τούτοισιν ἐξεύροιμ᾽ ὅπως R : \lecon{del.} Willems}
\ligneRav{24}{}{ἔτι \difference{προσμάθοιμι} χωλὸς εἶναι τὼ σκέλει·}{\textbf{24} προσμάθοι· μὴ R : \lecon{προσμάθοιμι} Ellebodius}
\ligneRav{25}{\auDessus{εὐ}{ρ}}{βάδιζε δευρὶ καὶ πρόσεχε τὸν νοῦν. \auDessus{:}{\tiny μν} ἰδού·}{}
\ligneRav{}{\auDessus{εὐ}{ρ}}{ὁρᾷς τὸ θύριον τοῦτο; \auDessus{:}{\tiny μν} νὴ τὸν Ἡρακλέα}{}
\ligneRav{27}{—}{οἶμαί γε \auDessus{:}{\tiny εὐ} σίγα \difference{νυν} \auDessus{:}{\tiny μν} σιωπῶ τὸ θύριον· \hfill\textbf{[α \Rpc]}}{\textbf{27} νῦν R : \lecon{νυν} Boissonade \espApp θύριον R : \lecon{θύριον;} Dindorf}
\ligneRav{28}{\auDessus{εὐ}{ρ}}{ἄκου᾽ \auDessus{:}{\tiny μν} \difference{ἀκούω} καὶ σιωπῶ τὸ θύριον· \hfill\textbf{[β \Rpc]}}{\textbf{28} ἀκούσω R : \lecon{ἀκούω} Brunck \espApp vers. 28 post 30 transp. \Rac : \lecon{corr.} \Rpc}
\ligneRav{29}{\auDessus{εὐ}{ρ}}{ἐνταῦθ᾽ Ἀγάθων ὁ κλεινὸς οἰκῶν τυγχάνει \hfill\textbf{[γ \Rpc]}}{}
\ligneRav{30}{—}{ὁ τραγῳδοποιός \auDessus{:}{\tiny μν} ποῖος οὗτος \difference{Ἁγάθων}· \hfill\textbf{[δ \Rpc]}}{\textbf{30-1} Ἀγάθων (bis) R : \lecon{Ἁγάθων} et \lecon{ἀγαθῶν} Bibi}
\ligneRav{31}{}{ἔστιν τις \difference{ἀγαθῶν}; μῶν ὁ μέλας ὁ καρτερός·}{\textbf{31} ἔστιν τις ἀγαθῶν \lecon{(Μνη.) continuat} R : (Εὐρ.) \Mon \espApp Ἀγάθων R : \lecon{ἀγαθῶν;} Bibi}
\ligneRav{32}{\auDessus{εὐ}{ρ}}{οὔκ, ἀλλ᾽ ἕτερός τις· $[\![\text{\sout{μ\auDessus{ν}{Η}}}]\!]$ οὐχ \difference{ἑόρακας} πώποτε;}{\textbf{32} (Μνη.) οὐχ ... πώποτε \Rac : \lecon{(Εὐρ.) continuat} \Rpc, Σ \espApp ἑόρακα R : \lecon{ἑόρακας} Bentley \ligneApp{v. 33} omm. \Rac, \Mon : in marg. \Rpc}
\ligneRav{33}{μ\auDessus{ν}{Η}}{$[\![\text{in marg. \Rpc{} : }$μῶν ὁ δασυπώγων : οὐχ ἑόρακας πώποτε·$]\!]$}{\textbf{33} omm. $\text{R}^{\text{ac}}, \text{M}$ : \lecon{in marg.} \Rpc{}}
\ligneRav{34}{μ\auDessus{ν}{Η}}{μὰ τὸν Δί᾽ οὔτοι γ᾽ ὥστε \difference{κἀμέ γ᾽} (\inRas{καὶ}{κἀ}μέ γ᾽) εἰδέναι·}{\textbf{34} καὶ μεγ᾽ \Rac \ (ὶ eras. \Rpc) : καὶ μεγ᾽ \Mon : \lecon{κἀμέ γ᾽} edd. omn.}
\ligneRav{35}{\auDessus{εὐ}{ρ}}{καὶ μὴν \difference{†}βε\inRas{α}{β}ίνηκας\difference{†} σύ γ᾽, ἀλλ᾽ οὐκ οἶσθ᾽ ἴσως·}{\textbf{35} βε\inRas{α}{β}ίνηκας \Rpc{} (ex βεαίνηκας \Rac{} ut vid.) : \lecon{†βεβίνηκας†} \lecon{Bibi}}
\ligneRav{}{}{ἀλλ᾽ ἐκποδὼν πτήξωμεν, ὡς ἐξέρχεται}{}
\ligneRav{}{}{θεράπων τις αὐτοῦ πῦρ ἔχων καὶ μυρρίνας}{}
\ligneRav{}{}{προθυσόμενος ἔοικε τῆς ποιήσεως·}{}
\ligneRav{39}{θερ/}{\hspace{1.1em}εὔφημος πᾶς ἔστω \difference{λαός},}{\textbf{39} λέως R, Suid. ε 3792 codd. AF : \lecon{λαός} Σ, Suid. codd. GIVM}
\ligneRav{40}{}{\hspace{1.1em}στόμα συγκλῄσας· ἐπιδημεῖ γὰρ}{}
\ligneRav{}{}{\hspace{1.1em}θίασος Μουσῶν ἔνδον μελάθρων}{}
\ligneRav{}{}{\hspace{1.1em}τῶν δεσποσύνων μελοποιῶν·}{}
}
\vspace{3pt}\noindent\rule{4.5cm}{0.4pt}\par\vspace{1.5pt}
{\fontsize{7.0pt}{8.3pt}\selectfont
\noindent\textbf{Loci paralleli (vv. 1–42) ---} \textbf{1} \Fr 617 πυθοῦ χελιδὼν πηνίκ᾽ ἄττα φαίνεται \espTest \textbf{2} (ἀλοῶν) Suidas α 1330 ; Photius α 1029 $=$ \SynB α 984 \espTest \textbf{18-9} (χοάνη = χώνη) Moeris χ 1 ~ Hesychius χ 576 χοάνη \espTest \textbf{24} \Macaire VI 90 ; \Apostol II 94 \espTest \textbf{33} (δασυπώγων) Pollux II 88 \espTest \textbf{36} (ἀλλ᾽ -- πτήξωμεν) Suidas π 3024 \espTest \textbf{39-41} Suidas ε 3791 (λεὼς codd. AF : λαὸς codd. GIVM, συγκλείσας, ἐνδημεῖ pro ἐπιδημεῖ ) \espTest \textbf{40-2} (ἐπιδημεῖ -- μελοποιῶν) Suidas θ 378\par
}
\vspace{2.5pt}\noindent\rule{4.5cm}{0.4pt}\par\vspace{1.5pt}
{\fontsize{6.9pt}{8.2pt}\selectfont
\noindent\textbf{Suite (en commentaire) de l'apparat critique \& notes paléographiques (\texttt{f. 162v}) ---} En haut de la page, \Rdeux{} (Main~B) a copié en semi-onciale de glose une notice liminaire de trois lignes ($\text{ὁ χορὸς ἐκ θεσμοφοριαζουσῶν... Εὐριπίδου γυνὴ μὲν Χοιρίνη, / μή(τ)ηρ δὲ Κλειτώ :}\sim\ \text{προλογίζει Μνησίλοχ(ος) ὁ κηδεστὴς Εὐριπίδου }\div$, sans liste séparée des \italique{Dramatis personae}), source unique des noms de \textbf{$\text{Μνησίλοχος}$} et de \textbf{$\text{Χοιρίνη / Χοιρίλη}$} défendus par Louis Hauseux (« Le dossier Mnésiloque », \texttt{Introduction.tex}, ll. 325–357). --- Au-dessous, le titre en semi-onciale d'apparat précédé d'un \italique{asteriscus} : $\ast\;\cdot\;\ast$ ἀριστοφάνους θεσμοφοριάζουσαι —. --- \textbf{Inversion des vv. 3–4 et transposition des vv. 27–30 corrigées par lettres-chiffres :} Aux vv.~3–4, \Rac{} avait copié le v.~4 avant le v.~3 (ordre rétabli par $\beta', \alpha'$ de \Rpc{}) ; aux vv.~27–30, \Rac{} avait copié le v. 28 après le v. 30 ; \Rpc{} a rétabli l'ordre exact en inscrivant à l'extrémité droite des quatre lignes les lettres numérales \textbf{\textbf{α, γ, δ, β}}. --- \textbf{Sigle tachygraphique gratté et scholie au v. 32 :} Comme le relève Louis Hauseux (\texttt{vv 1-641.tex}, ll. 777–787), le manuscrit comportait à l'intérieur du vers, devant οὐχ ἑόρακα πώποτε, le sigle tachygraphique μ\auDessus{ν}{Η} ($\overset{\text{Η}}{\text{μν}}$ = $\text{Μνησίλοχος}$), qui a été ensuite \textbf{gratté (\italique{ut videtur})} ; l'erreur de ce sigle est en outre prouvée par la scholie $\Sigma^{\text{R}}$ du v. 32 (ἐπειδὴ οὐ πάλαι ἤρξατο διδάσκειν, ἀλλὰ πρὸ τριῶν ἐτῶν πρὸ τούτου), qui ne se comprend qu'avec la 2\ieme{} personne ἑόρακας prononcée par Euripide. --- \textbf{Omission de vers réparée en marge (v. 33) :} Le \textbf{v. 33} ($\text{μῶν ὁ δασυπώγων; --- οὐχ ἑόρακας πώποτε}$) avait été omis par saut du même au même dans la colonne par \Rac{} (et reste \textbf{omis dans \Mon{}}) ; il est rajouté sur deux lignes dans la marge droite par \Rpc{}. --- \textbf{Grattage au v. 34 et correction \latin{in rasura} au v. 35 :} Au v.~34, \Rac{} avait écrit καὶ μεγ' (copié tel quel par \Mon{}) ; \Rpc{} a gratté l'iota (ὶ eras.) pour amorcer la crase κἀμέ γ'. Au \textbf{v.~35}, le second $\beta$ de \textbf{βεβίνηκας} est corrigé \latin{in rasura} sur un $\alpha$ gratté (\textbf{βεαίνηκας} \Rac{}), ce qui, joint au silence complet de $\Sigma^{\text{R}}$, motive les \latin{cruces} \lecon{†βεβίνηκας†} de \lecon{Bibi}. --- \textbf{v. 1 :} ἄρα πότε R : \lecon{ἆρά ποτε} Biset \espApp χελιδὼν R : \lecon{Χελιδὼν} \lecon{Bibi} ap. Clements (dans l'\italique{Index nominum} et le commentaire : référence au devin \italique{Khelidôn}, πυθοῦ Χελιδόνος). --- \textbf{v. 3 :} τι R : \lecon{τε} $\Sigma$ (δυνατόν), Souda $\text{οι}\ 150$. --- \textbf{v. 10 :} \lecon{ὁρᾶν.} R (ponctuation haute maintenue par L. Hauseux) : ὁρᾶν; Gelenius. --- \textbf{v. 18 :} \lecon{ἀκοὴν δὲ χοάνης} R, défendu par Alphonse Willems et adopté par L. Hauseux : ἀκοὴς δὲ χοάνην Biset : ἀκοῇ δὲ χοάνην Boissonade. --- \textbf{v. 21 :} τε R : \lecon{γε} $\Sigma$. --- \textbf{v. 23 :} πρὸς τοῖς ἀγαθοῖς τούτοισιν ἐξεύροιμ' ὅπως R : \textbf{vers athétisé (\italique{del.}) par Alphonse Willems} et L. Hauseux comme interpolation issue d'une glose syntaxique. --- \textbf{v. 24 :} προσμάθοι· μὴ R (médivision avec \italique{dikolon} fautif) : \lecon{προσμάθοιμι} Ellebodius. --- \textbf{v. 28 :} ἀκούσω R : \lecon{ἀκούω} Brunck \espApp vers. 28 post 30 transp. \Rac{} : \lecon{corr.} \Rpc{} (litt. $\alpha, \gamma, \delta, \beta$ in mg.). --- \textbf{vv. 30–31 :} Ἀγάθων (bis) R : \lecon{Ἁγάθων} (v. 30, avec crase $\text{ὁ Ἀγάθων}$) et \lecon{ἀγαθῶν} (v. 31, génitif pluriel en calembour) \lecon{Bibi} \espApp Au \textbf{v. 31}, ἔστιν τις ἀγαθῶν \lecon{(Μνη.) continuat} R (maintenu par \textbf{Bibi}) : (Εὐρ.) \Mon{}. --- \textbf{v. 32 :} (Μνη.) (μ\auDessus{ν}{Η}) οὐχ ... πώποτε \Rac{} (sigle gratté \italique{ut vid.}) : \lecon{(Εὐρ.) continuat} $\text{R}^{\text{pc}}, \Sigma$ \espApp ἑόρακα R : \lecon{ἑόρακας} Bentley. --- \textbf{v. 35 :} βε\inRas{α}{β}ίνηκας \Rpc{} (ex βεαίνηκας \Rac{} ut vid.) : \lecon{†βεβίνηκας†} \lecon{Bibi}. --- \textbf{v. 39 :} λέως R, Souda $\varepsilon\ 3792$ codd. AF : \lecon{λαός} $\Sigma$, Souda codd. GIVM.\par
}
\end{minipage}
\end{adjustbox}
\end{adjustwidth}
\newpage
\enTeteD{Codex Ravennas 429 (Biblioteca Classense) --- Fac-similé du f. 162v (vv. 1–42)}
\begin{adjustwidth}{-1.5cm}{-3.0cm}
\begin{center}
\includegraphics[width=16.5cm,height=20.5cm,keepaspectratio]{/workspaces/Thesmophories/dossier_ravennas_thesmophories/folios_individuels_ff162v-177r/01_f162v_classense164v_vv0001-0042.jpg}\par
\vspace{2pt}
{\scriptsize\textbf{Codex Ravennas 429, f. 162v (Classense 164v)} --- ΠΡΟΛΟΓΟΣ, \textbf{vv. 1–42} ($\text{R}^1, \Rpc, \Rdeux$ (\Mon)).}
\end{center}
\end{adjustwidth}
\newpage
\clearpage\ifodd\value{page}\hbox{}\newpage\fi% --- FOLIO 01 (f. 162v) : SPREAD 2.1 (Texte grec établi + Traduction française + Commentaire 1/8 & 2/8) ---
\clearpage\ifodd\value{page}\hbox{}\newpage\fi
\enTeteG{ΘΕΣΜΟΦΟΡΙΑΖΟΥΣΑΙ --- Texte grec établi (f. 162v, vv. 1–42 [1/4])}
\enTeteD{LES THESMOPHORIES --- Traduction française (f. 162v, vv. 1–42 [1/4])}
\begin{adjustwidth}{-2cm}{-2cm}
\noindent\textbf{\footnotesize Texte grec établi --- Folio 01 (\texttt{f. 162v}, vv. 1–42 [1/4]) \hfill ΠΡΟΛΟΓΟΣ}\\[-2pt]
\noindent\rule{16.5cm}{0.5pt}\\[2pt]
{\fontsize{8.8pt}{10.2pt}\selectfont
\ligneEst{\textbf{Μνησίλοχος} {\tiny\textit{$3$ ia.}}}{ὦ Ζεῦ χελιδὼν ἆρά ποτε φανήσεται;}{1}
\ligneEst{}{ἀπολεῖ μ᾽ ἀλοῶν ἅνθρωπος ἐξ ἑωθινοῦ.}{}
\ligneEst{}{παρὰ σοῦ πυθέσθαι ποῖ μ᾽ ἄγεις ωὖριπίδη;}{}
\ligneEst{}{οἷόν τε, πρὶν τὸν σπλῆνα κομιδῇ μ᾽ ἐκβαλεῖν.}{}
\ligneEst{\textbf{Εὐριπίδης}}{ἀλλ᾽ οὐκ ἀκούειν δεῖ σε πάνθ᾽ ὅσ᾽ αὐτίκα}{5}
\ligneEst{}{ὄψει παρεστώς.}{}
\ligneEst{\textbf{Μνησίλοχος}}{\decalage{ὄψει παρεστώς.}πῶς λέγεις; αὖθις φράσον.}{}
\ligneEst{}{οὐ δεῖ μ᾽ ἀκούειν;}{}
\ligneEst{\textbf{Εὐριπίδης}}{\decalage{οὐ δεῖ μ᾽ ἀκούειν;}οὐχ ἅ γ᾽ ἂν μέλλῃς ὁρᾶν.}{}
\ligneEst{\textbf{Μνησίλοχος}}{οὐδ᾽ ἆρ᾽ ὁρᾶν δεῖ μ᾽;}{}
\ligneEst{\textbf{Εὐριπίδης}}{\decalage{οὐδ᾽ ἆρ᾽ ὁρᾶν δεῖ μ᾽;}οὐχ ἅ γ᾽ ἂν ἀκούειν δέῃ.}{}
\ligneEst{\textbf{Μνησίλοχος}}{πῶς μοι παραινεῖς; δεξιῶς μέντοι λέγεις.}{}
\ligneEst{}{οὐ φῂς σὺ χρῆναί μ᾽ οὔτ᾽ ἀκούειν οὔθ᾽ ὁρᾶν.}{10}
}
\vspace{2pt}
\noindent\rule{16.5cm}{0.6pt}\\[2pt]
\noindent\textbf{\small Commentaire philologique, métrique et scénique (vv. 1–42 --- partie 1/8)}\\[2pt]
\begin{adjustbox}{max totalheight=15.2cm, max width=16.5cm, keepaspectratio}%
\parbox[t]{16.5cm}{
\setlength{\columnsep}{9pt}%
\begin{multicols}{2}
\fontsize{7.6pt}{8.8pt}\selectfont
\setlength{\parindent}{0pt}%
\setlength{\parskip}{1.2pt}%
\noindent\textbf{1-10} \textbf{ὦ Ζεῦ ... οὐχ ἅ γ᾽ ἂν ἀκούειν δέῃ.} Nous sommes au tout début de la pièce, il fait encore nuit, un personnage entre d'un pas pressé ; il est immédiatement reconnaissable à son masque bien connu des spectateurs : c'est Euripide\marqueur{1}. Il est suivi d'un vieillard qui traîne la patte et crie sa détresse -- \italique{topos} comique. (On apprendra plus loin qu'il est un parent d'Euripide ; la tradition l'identifie à Mnésiloque, le père de la première femme d'Euripide, Chœrile\marqueur{2}.) Réplique fracassante d'Euripide, et bien selon ses façons -- les fameuses γνῶμαι d'Euripide fort goûtées des spectateurs : ἀλλ᾽ οὐκ ἀκούειν δεῖ σε πάνθ᾽ ὅσ᾽ αὐτίκα ὄψει παρεστώς, « il n'est pas besoin que tu entendes tout ce que tu vas bientôt voir ».
\noindent Hors d'haleine, Mnésiloque n'est pas en état de suivre Euripide dans ses aphorismes : il s'est arrêté à la première partie : οὐ δεῖ μ᾽ ἀκούειν, « il ne faut pas que j'entende... » Euripide reprend : οὐχ ἅ γ᾽ ἂν μέλῃς ὁρᾶν, « pas certes ce que tu es sur le point de voir ».\\ C'est là que nous arrivons à l'étrange réponse de Mnésiloque au vers 7b : οὐδ ἆρ ὁρᾶν δεῖ μ᾽, « je ne dois pas non plus voir ? ». Ne comprenant pas pourquoi  Mnésiloque passe d'une défense d'entendre à une défense de voir, l'on est incité à mieux scruter les vers précédents.
\noindent Le scholiaste nous met sur la piste : les constructions tortueuses d'Euripide sont sujettes à l'amphibologie\marqueur{3}.
\noindent Mnésiloque commet une erreur (volontaire) d'interprétation sur la dernière réplique d'Euripide : dans `οὐχ ἅ γ ἂν μέλῃς ὁρᾶν', il feint d'entendre `ἄγαν' à la place de `ἅ γ᾽ ἂν'. Aristophane ayant mis à dessein ces particules dans un ordre inhabituel\marqueur{4}.
\noindent Il est remarquable de noter que, des onze pièces conservées de notre poète, le cas ne se rencontre qu'une seule autre fois... et c'est encore un pastiche d'un Euripide postiche (celui des \italique{Grenouilles}, cette fois, vers 960) :
\citComm{ἐξ \textbf{ὧν γ᾽ ἂν} ἐξηλεγχόμην.}
\noindent Mnésiloque croit comprendre que le voir est lui aussi soumis au μηδὲν ἄγαν grec : `οὐκ ἄγαν μέλῃς ὁρᾶν', « tu ne dois pas trop voir »... d'où sa réaction du vers 7b.\\
\noindent\textbf{1-4} \textbf{ὦ Ζεῦ ... πρὶν τὸν σπλῆνα κομιδῇ μ᾽ ἐκβαλεῖν.} Ce quatrain est analogue à celui des vers 71-74. À deux reprises, Mnésiloque adresse une prière plaintive à Zeus avant de questionner Euripide. Ressemblance saisissante aussi avec les premiers vers des \italique{Nuées} (vv. 1-5) :
\citComm{Στρεψιάδης --- ἰοὺ ἰού·\\
ὦ Ζεῦ βασιλεῦ τὸ χρῆμα τῶν νυκτῶν ὅσον·\\
ἀπέραντον. οὐδέποθ᾽ ἡμέρα γενήσεται;\\
καὶ μὴν πάλαι γ᾽ ἀλεκτρυόνος ἤκουσ᾽ ἐγώ·\\
οἱ δ᾽ οἰκέται ῥέγκουσιν· ἀλλ᾽ οὐκ ἂν πρὸ τοῦ.\\
Strepsiade --- Iou ! iou !\\
Ah ! Zeus roi, quel fardeau que ces nuits !\\
C'est interminable... Il ne fera donc jamais jour ?\\
Pourtant voilà longtemps que j'ai entendu le coq,\\
et mes serviteurs ronflent. Ah ! que non pas autrefois !}
\end{multicols}
}%
\end{adjustbox}%
\vfill
\end{adjustwidth}
\newpage
% --- PAGE DE DROITE : Traduction française + Commentaire (2/8) ---
\begin{adjustwidth}{-2cm}{-2cm}
\noindent\textbf{\footnotesize Traduction française --- Folio 01 (\texttt{f. 162v}, vv. 1–42 [1/4]) \hfill ΠΡΟΛΟΓΟΣ}\\[-2pt]
\noindent\rule{16.5cm}{0.5pt}\\[2pt]
{\fontsize{8.8pt}{10.2pt}\selectfont
\ligneEst{\textbf{Mnésiloque}}{\francais{Ô Zeus, est-ce qu'à la fin une hirondelle paraîtra ?}}{1}
\ligneEst{}{\francais{Cet homme me tuera de me faire tribuler dès le matin.}}{}
\ligneEst{}{\francais{Y a-t-il moyen, avant que perdre tout de bon la rate,}}{}
\ligneEst{}{\francais{de savoir de toi où tu me mènes, Euripide ?}}{}
\ligneEst{\textbf{Euripide}}{\francais{Eh ! il n'est pas besoin que tu entendes tout ce que}}{5}
\ligneEst{}{\francais{tout à l'heure tu verras de tes yeux.}}{}
\ligneEst{\textbf{Mnésiloque}}{\francais{\espace{2.5cm}Comment dis-tu ? Répète.}}{}
\ligneEst{}{\francais{Il n'est pas besoin que j'entende.}}{}
\ligneEst{\textbf{Euripide}}{\francais{\espace{2.7cm}Non, ce que tu vas voir.}}{}
\ligneEst{\textbf{Mnésiloque}}{\francais{Ni donc que je voie...}}{}
\ligneEst{\textbf{Euripide}}{\francais{\espace{3.1cm}Non, ce qu'il faut que tu entendes.}}{}
\ligneEst{\textbf{Mnésiloque}}{\francais{Singulier conseil ! Ma foi, c'est finement dit.}}{}
\ligneEst{}{\francais{Tu prétends que je ne dois ni entendre ni voir.}}{10}
}
\vspace{2pt}
\noindent\rule{16.5cm}{0.6pt}\\[2pt]
\noindent\textbf{\small Commentaire philologique, métrique et scénique (vv. 1–42 --- partie 2/8)}\\[2pt]
\begin{adjustbox}{max totalheight=15.2cm, max width=16.5cm, keepaspectratio}%
\parbox[t]{16.5cm}{
\setlength{\columnsep}{9pt}%
\begin{multicols}{2}
\fontsize{7.6pt}{8.8pt}\selectfont
\setlength{\parindent}{0pt}%
\setlength{\parskip}{1.2pt}%
\noindent Cette analogie permet de saisir le sens imagé de l'hirondelle (χελιδών) de notre passage : elle est mise à la place du coq des \italique{Nuées} pour ce qu'elle a de semblable : en chantant dès avant l'aube, elle annonce le jour bientôt renaissant\marqueur{1}. Que l'hirondelle était réputée chez les Grecs chanter\marqueur{2} au matin, cela est appuyé par l'adjectif ὀρθρογόη, \italique{aux lamentations matinales}, qu'Hésiode (\italique{Opera et Dies}, v. 568) applique à l'hirondelle\marqueur{3} :
\citComm{ὀρθρογόη Πανδιονὶς ὦρτο χελιδὼν\\
La fille matinale et gémissante de Pandion, l'hirondelle, réveille...}
\noindent Ce strident réveille-matin avait donné lieu à un expression proverbiale dont Aristophane fait mention ailleurs\marqueur{4}:
\citComm{πυθοῦ χελιδὼν πηνίκ᾽ ἄττα φαίνεται\\
Renseigne-toi à quelle heure l'hirondelle apparaît.}
\noindent Ici, ce n'est pas une hirondelle mais Euripide qui a jeté hors du lit Mnésiloque bien trop tôt.. Ce sera donc lui qui sera interrogé : \textbf{παρὰ σοῦ} πυθέσθαι (v. 3).
\noindent\textbf{2} \textbf{ἀλοῶν.} Ἀλοᾶν est un terme agricole signifiant \italique{battre le grain}, selon une technique bien définie : sur l'aire à battre (ἡ ἅλως), on fait tourner des bœufs, attachés à un piquet, qui piétinent et qui broient la céréale jonchant le sol. De ce fait, ἁλοᾶν est devenu un synonyme expressif de \grec{τύπτειν} avec le sens de \italique{broyer de coups}. (Note de Taillardat\marqueur{1})\\ Ici, ἁλοᾶν a le sens moral de \italique{tourmenter}. Conformité du grec et du français -- ou du latin, plutôt : l'on possède la même image en avec notre \italique{tribulation}, le latin \latin{tribulare} signifie « battre le blé, presser avec la herse (\latin{tribulum}) ».
\noindent\textbf{[Notes ad 1-10]} \textbf{(1)} À quoi pouvait ressembler ce masque d'Euripide ? On peut supposer qu'il présentait une barbe très touffue (cf. le vers 190) ainsi que de multiples taches sur son visage : Ἐλέγετο δὲ καὶ βαθὺν πώγωνα θρέψαι καὶ ἐπὶ τῆς ὄψεως φακοὺς ἐσχηκέναι. (Tiré d'une des \italique{Vies} d'Euripide, rassemblées par Richard Kannicht pour son édition des fragments d'Euripide, cinquième volume des \italique{Tragicorum Graecorum Fragmenta} (2004). Ce détail se trouve à la page 47.)
\noindent \textbf{(2)} \italique{Confer} l'introduction sur la vie d'Euripide et de Mnésiloque et le vers 291.
\noindent \textbf{(3)} ὁ μὲν τραγικώτερον καὶ ὑψηλότερον φράζει, ὁ δὲ ταπεινότερον ἀκούει, « l'un s'exprime dans un style plus tragique et élevé, quand l'autre comprend de façon plus basse ».
\noindent \textbf{(4)} Le véritable Euripide s'exprimait ainsi : (\italique{Oreste}, vers 81) : Ἑλένη, τί σοι λέγοιμ᾽ \textbf{ἂν ἅ γε} παροῦσ᾽ ὁρᾷς;
\noindent \textbf{(1)} Et non pas le printemps -- qui n'est pas de saison ici --, comme on peut le lire drôlement dans nombre de commentaires.
\noindent \textbf{(2)} Chanter ?.. ou plutôt babiller : l'hirondelle était proverbialement brocardée pour son caquet, cf. notamment Taillardat §§520-522 et les notes d'Alphonse Willems à propos du vers 682 des \italique{Grenouilles}.
\noindent \textbf{(3)} Voir aussi Hésychius (ο 1189) : ὀρθ<ρ>ογόη· ἡ χελιδών.
\end{multicols}
}%
\end{adjustbox}%
\vfill
\end{adjustwidth}
\newpage

% --- FOLIO 01 (f. 162v) : SPREAD 2.2 (Texte grec établi + Traduction française + Commentaire 3/8 & 4/8) ---
\clearpage\ifodd\value{page}\hbox{}\newpage\fi
\enTeteG{ΘΕΣΜΟΦΟΡΙΑΖΟΥΣΑΙ --- Texte grec établi (f. 162v, vv. 1–42 [2/4])}
\enTeteD{LES THESMOPHORIES --- Traduction française (f. 162v, vv. 1–42 [2/4])}
\begin{adjustwidth}{-2cm}{-2cm}
\noindent\textbf{\footnotesize Texte grec établi --- Folio 01 (\texttt{f. 162v}, vv. 1–42 [2/4]) \hfill ΠΡΟΛΟΓΟΣ}\\[-2pt]
\noindent\rule{16.5cm}{0.5pt}\\[2pt]
{\fontsize{8.8pt}{10.2pt}\selectfont
\ligneEst{\textbf{Εὐριπίδης}}{χωρὶς γὰρ αὐτοῖν ἑκατέρου 'στὶν ἡ φύσις.}{}
\ligneEst{\textbf{Μνησίλοχος}}{τοῦ μήτ᾽ ἀκούειν μήθ᾽ ὁρᾶν;}{}
\ligneEst{\textbf{Εὐριπίδης}}{\decalage{τοῦ μήτ᾽ ἀκούειν μήθ᾽ ὁρᾶν;}εὖ ἴσθ᾽ ὅτι.}{}
\ligneEst{\textbf{Μνησίλοχος}}{πῶς χωρίς;}{}
\ligneEst{\textbf{Εὐριπίδης}}{\decalage{πῶς χωρίς;}οὕτω ταῦτα διεκρίθη τότε.}{}
\ligneEst{}{αἰθὴρ γὰρ ὅτε τὰ πρῶτα διεχωρίζετο}{}
\ligneEst{}{καὶ ζᾦ᾽ ἐν αὑτῷ ξυνετέκνου κινούμενα,}{15}
\ligneEst{}{ᾧ μὲν βλέπειν χρὴ πρῶτ᾽ ἐμηχανήσατο}{}
\ligneEst{}{ὀφθαλμὸν ἀντίμιμον ἡλίου τροχῷ,}{}
\ligneEst{}{ἀκοὴν δὲ χοάνης ὦτα διετετρήνατο.}{}
\ligneEst{\textbf{Μνησίλοχος}}{διὰ τὴν χοάνην οὖν μήτ᾽ ἀκούω μήθ᾽ ὁρῶ;}{}
\ligneEst{}{νὴ τὸν Δί᾽ ἥδομαί γε τουτὶ προσμαθών.}{20}
\ligneEst{}{οἷόν γέ που 'στιν αἱ σοφαὶ ξυνουσίαι.}{}
\ligneEst{\textbf{Εὐριπίδης}}{πόλλ᾽ ἂν μάθοις τοιαῦτα παρ᾽ ἐμοῦ.}{}
}
\vspace{2pt}
\noindent\rule{16.5cm}{0.6pt}\\[2pt]
\noindent\textbf{\small Commentaire philologique, métrique et scénique (vv. 1–42 --- partie 3/8)}\\[2pt]
\begin{adjustbox}{max totalheight=15.2cm, max width=16.5cm, keepaspectratio}%
\parbox[t]{16.5cm}{
\setlength{\columnsep}{9pt}%
\begin{multicols}{2}
\fontsize{7.6pt}{8.8pt}\selectfont
\setlength{\parindent}{0pt}%
\setlength{\parskip}{1.2pt}%
\noindent \textbf{(4)} Fragment 617 Kassel-Austin. Diversement recensé par les parémiographes (souvent sous la forme πυθοῦ Χελιδόνος) : Hésychius π 4311 ; Photius π 1519 ; Suidas π 3136 ; \Apostol XV, 14 ; \&c.
\noindent \textbf{(1)} Je me référerai souvent – et avec force citations ! – au magistral ouvrage que Jean Taillardat a consacré à la langue bigarrée et haute en couleurs d'Aristophane : \italique{Les images d'Aristophane} édité aux Belles Lettres (seconde édition de 1965). Je cite ici le §138 ; voir aussi les références du §606 : la métaphore était usuelle et pas de l'invention de notre poète : \italique{Grenouilles} 149, Sophocle fr. 19 \Nauck [$=$ fr. 20 Radt (dub.)] ; Hérondas 2, 34 ; Babrius, \italique{Le lion amoureux}, 98.15 ; καταλοᾶν Eschine 2, 140 ; Xénoph. \italique{Cyrop.} 7, 1, 31 ; Lucien \italique{Icarom.} 15 ; συναλοᾶν Hippon. 14AD\expo{3}.4 [$=$ fr. 92.4 Masson $=$ fr 95.4 Degani] ; Théocr. 22.128.
\noindent\textbf{3} \textbf{πρὶν τὸν σπλῆνα κομιδῇ μ᾽ ἐκβαλεῖν.} Les dérangements (palpitations) dus à la course – d'autant plus aggravés par l'âge avancé de Mnésiloque – étaient imputés par les Anciens à des troubles de la rate (σπλήν, \anglais{spleen} en anglais)\marqueur{1}. En français, celui qui court ventre à terre court aussi comme un « dé-raté ».
\noindent\textbf{[Notes ad 3]} \textbf{(1)} Voir Pline l'Ancien, \italique{Histoire naturelle} 11,205 ; Plaute, \italique{Mercator}, v. 124.
\noindent\textbf{13-8} \textbf{} Aristophane nous peint Euripide comme un de ces faiseurs de cosmogonies (tel le Socrate des \italique{Nuées} avec son Δῖνος). Le véritable Euripide avait effectivement un penchant pour ce genre de nuées\marqueur{1}. En l'occurrence, le pastiche de cosmogonie qu'Aristophane met dans la bouche d'Euripide semble inspiré d'Empédocle\marqueur{2}.
\noindent\textbf{16-8} \textbf{πρῶτα ἐμηχανήσατο ... ἀκοὴν δὲ χοάνης ὦτα διετετρήνατο.} Pour l'explication de ce passage diversement retouché au gré des éditeurs, je ne peux que renvoyer au brillant exercice d'exégèse d'Alphonse Willems. (Ses notes sur les \italique{Thesmophories} sont reproduites en annexe, celle relative à ce passage en est la première. À propos d'Alphonse Willems, \italique{confer} l'Introduction.) Euripide recourt à une image tirée de la fonderie -- certainement plus parlante aux spectateurs de l'époque qu'à nous : l'oreille (ἀκοή) est comparée aux ὦτα (littéralement les \italique{oreilles} et, par analogie, les \italique{anses}, \italique{poignées} ou, comme ici, les \italique{godets} ou \italique{entonnoirs} d'un moule, χοάνη). Willems construit ainsi la phrase, en accord avec la sobre indication du scholiaste (λείπει ὡς) : διετετρήνατο δὲ ἀκοὴν (ὡς) χοάνης ὦτα.
\noindent\textbf{[Notes ad 13-8]} \textbf{(1)} Voir par exemple le fragment 5 de \italique{Mélanippe la sage} dans l'édition de Jouan-van Looy aux Belles Lettres, tome VIII, \nieme{2} partie ($=$ fragment 484 de l'édition Kannicht).
\noindent \textbf{(2)} Cf. notamment ce passage d'Eusèbe de Césarée (\italique{Préparation évangélique}, XV, 37 ; fragment A49 Diels-Kranz) : Ἐμπεδοκλῆς τὸν μὲν αἰθέρα πρῶτον διακριθῆναι et la création de l'œil par Aphrodite (B84.9 D.-K.) reconstitué par Blass (1883) : χοάνῃσι δίαντα τετρήατο θεσπεσίῃσιν.
\end{multicols}
}%
\end{adjustbox}%
\vfill
\end{adjustwidth}
\newpage
% --- PAGE DE DROITE : Traduction française + Commentaire (4/8) ---
\begin{adjustwidth}{-2cm}{-2cm}
\noindent\textbf{\footnotesize Traduction française --- Folio 01 (\texttt{f. 162v}, vv. 1–42 [2/4]) \hfill ΠΡΟΛΟΓΟΣ}\\[-2pt]
\noindent\rule{16.5cm}{0.5pt}\\[2pt]
{\fontsize{8.8pt}{10.2pt}\selectfont
\ligneEst{\textbf{Euripide}}{\francais{C'est que distincte est l'essence de chacune de ces choses.}}{}
\ligneEst{\textbf{Mnésiloque}}{\francais{Ne pas entendre et ne pas voir ?}}{}
\ligneEst{\textbf{Euripide}}{\francais{\espace{4.4cm}Tiens-le pour assuré.}}{}
\ligneEst{\textbf{Mnésiloque}}{\francais{Comment distincte ?}}{}
\ligneEst{\textbf{Euripide}}{\francais{\espace{1.8cm}Voici comment elles furent départies à l'origine.}}{}
\ligneEst{}{\francais{L'Éther, quand d'abord il se dégagea et procréa}}{}
\ligneEst{}{\francais{en soi-même des animaux doués de mouvement,}}{15}
\ligneEst{}{\francais{commença par fabriquer pour qui doit voir}}{}
\ligneEst{}{\francais{un œil analogique de l'orbe solaire}}{}
\ligneEst{}{\francais{et fora une oreille à l'image des godets d'un moule.}}{}
\ligneEst{\textbf{Mnésiloque}}{\francais{Et à cause de ce moule je ne dois ni entendre ni voir ?}}{}
\ligneEst{}{\francais{Par Zeus, je suis ravi d'apprendre encore cela.}}{20}
\ligneEst{}{\francais{Quelle chose tout de même, que les doctes accointances !}}{}
\ligneEst{\textbf{Euripide}}{\francais{Auprès de moi, tu en apprendras bien d'autres du même genre.}}{}
}
\vspace{2pt}
\noindent\rule{16.5cm}{0.6pt}\\[2pt]
\noindent\textbf{\small Commentaire philologique, métrique et scénique (vv. 1–42 --- partie 4/8)}\\[2pt]
\begin{adjustbox}{max totalheight=15.2cm, max width=16.5cm, keepaspectratio}%
\parbox[t]{16.5cm}{
\setlength{\columnsep}{9pt}%
\begin{multicols}{2}
\fontsize{7.6pt}{8.8pt}\selectfont
\setlength{\parindent}{0pt}%
\setlength{\parskip}{1.2pt}%
\noindent\textbf{19} \textbf{διὰ τὴν χοάνην οὖν μήτ᾽ ἀκούω μήθ᾽ ὁρῶ;} Mnésiloque ne goûte guère aux subtilités des poètereaux et tord à plaisir leurs propos. Ses saillies fort désobligeantes se repèrent aisément puisqu'elles sont toujours superbement ignorées. C'en est ici le premier exemple ; Euripide le-faiseur-de-cosmogonies euripidise en lui racontant des histoires de fonderie et de moule... conclusion de Mnésiloque pour expliquer son accablement : « C'est donc à cause de ton moule que je suis moulu ? », διὰ τὴν χοάνην οὖν μήτ᾽ ἀκούω μήθ᾽ ὁρῶ\marqueur{1}. C'est bien le sens qu'il faut donner au tour `μήτ᾽ ἀκούω μήθ᾽ ὁρῶ' et qu'on peut rapprocher du latin \citComm{\latin{Nec aures me credo habere nec tango},\\
Je crois n'avoir plus ni oreilles ni main,} qui signifie justement qu'on ne se sent plus la force de faire quoi que ce soit\marqueur{2}.\\
\noindent La langue d'Aristophane, comme nous l'allons voir en chemin, est truffée d'échos aux proverbes grecs ; on ne peut que regretter le manque d'intérêt criant des commentateurs d'Aristophane pour ce domaine. Nous nous efforcerons au contraire, et selon nos moyens, de souligner chacun des emprunts d'Aristophane à cette langue proverbiale. Et ils sont nombreux ! Le cas présent est malheureusement moins probant puisque je n'ai pas pu trouver d'autres occurrences dans la littérature grecque. (Surtout : le jeu de mots exact en grec, qui tourne autour de χοάνη, nous échappe.)
\noindent\textbf{21} \textbf{οἷόν τέ που 'στιν αἱ σοφαὶ ξυνουσίαι.} Le scholiaste note à ce vers une allusion à un mot célèbre : σοφοὶ τύραννοι τῶν σοφῶν ξυνουσίᾳ, qu'Aristophane dans ses \italique{Héros}\marqueur{1} mais aussi Platon\marqueur{2} et Anthistène\marqueur{3} attribuent à Euripide. Cela n'empêche pas moins notre scholiaste, bien sûr de lui, d'accorder la paternité de cette formule à Sophocle\marqueur{4}. Toujours est-il qu'Euripide n'en passait certainement pas moins aux yeux d'Aristophane et de ses spectateurs pour le père putatif de la citation. Sur le sens de γέ που, on lira avec intérêt la seconde note de Willems reproduite ci-après en annexe.
\noindent\textbf{[Notes ad 19]} \textbf{(1)} Ou bien : « Je ne vois plus Claire ; j'ai perdu Louis. » (traduction plus libre.)
\noindent \textbf{(2)} Je renvoie le lecteur à l'admirable \italique{Dictionnaire des sentences latines et grecques} de Renzo Tosi (traduit de l'italien par Rebecca Lenoir ; proverbe n° 1971, cité par saint Jérôme notamment). Érasme avait recensé le proverbe dans ses \italique{Adages}, n° 2500. (Je cite les \italique{Adages} ainsi que leur traduction française selon l'édition de Jean-Christophe Saladin aux Belles Lettres (2019).) Toujours dans le §1971, Tosi cite aussi la sentence médiévale \latin{Nec audio nec video}, qui se superpose parfaitement au grec μήτ᾽ ἀκούω μήθ᾽ ὁρῶ.
\noindent \textbf{(1)} Fragment 323 de l'édition Kassel-Austin.
\noindent \textbf{(2)} \italique{Théagès}, 125b et la \italique{République}, VIII, 568b. Platon devait avoir un point de vue sur Euripide assez proche de celui d'Aristophane : il chasse joyeusement de sa République idéale Euripide, cet homme qui occupe une place prestigieuse parmi les semi-habiles de sa cité.
\noindent \textbf{(3)} Fragment 59 de l'édition Decleva-Caizzi
\end{multicols}
}%
\end{adjustbox}%
\vfill
\end{adjustwidth}
\newpage

% --- FOLIO 01 (f. 162v) : SPREAD 2.3 (Texte grec établi + Traduction française + Commentaire 5/8 & 6/8) ---
\clearpage\ifodd\value{page}\hbox{}\newpage\fi
\enTeteG{ΘΕΣΜΟΦΟΡΙΑΖΟΥΣΑΙ --- Texte grec établi (f. 162v, vv. 1–42 [3/4])}
\enTeteD{LES THESMOPHORIES --- Traduction française (f. 162v, vv. 1–42 [3/4])}
\begin{adjustwidth}{-2cm}{-2cm}
\noindent\textbf{\footnotesize Texte grec établi --- Folio 01 (\texttt{f. 162v}, vv. 1–42 [3/4]) \hfill ΠΡΟΛΟΓΟΣ}\\[-2pt]
\noindent\rule{16.5cm}{0.5pt}\\[2pt]
{\fontsize{8.8pt}{10.2pt}\selectfont
\ligneEst{\textbf{Μνησίλοχος}}{\decalage{πόλλ᾽ ἂν μάθοις τοιαῦτα παρ᾽ ἐμοῦ.}πῶς ἂν οὖν}{}
\ligneEst{}{ἔτι προσμάθοιμι χωλὸς εἶναι τὼ σκέλει;}{}
\ligneEst{\textbf{Εὐριπίδης}}{βάδιζε δευρὶ καὶ πρόσεχε τὸν νοῦν.}{25}
\ligneEst{\textbf{Μνησίλοχος}}{\decalage{βάδιζε δευρὶ καὶ πρόσεχε τὸν νοῦν.}ἰδού.}{}
\ligneEst{\textbf{Εὐριπίδης}}{ὁρᾷς τὸ θύριον τοῦτο;}{}
\ligneEst{\textbf{Μνησίλοχος}}{\decalage{ὁρᾷς τὸ θύριον τοῦτο;}νὴ τὸν Ἡρακλέα}{}
\ligneEst{}{οἶμαί γε.}{}
\ligneEst{\textbf{Εὐριπίδης}}{\decalage{οἶμαί γε.}σίγα νυν.}{}
\ligneEst{\textbf{Μνησίλοχος}}{\decalage{οἶμαί γε.σίγα νυν.}σιωπῶ τὸ θύριον.}{}
\ligneEst{\textbf{Εὐριπίδης}}{ἄκου᾽.}{}
\ligneEst{\textbf{Μνησίλοχος}}{\decalage{ἄκου᾽.}ἀκούω καὶ σιωπῶ τὸ θύριον.}{}
\ligneEst{\textbf{Εὐριπίδης}}{ἐνταῦθ᾽ Ἀγάθων ὁ κλεινὸς οἰκῶν τυγχάνει}{}
\ligneEst{}{ὁ τραγῳδοποιός.}{30}
}
\vspace{2pt}
\noindent\rule{16.5cm}{0.6pt}\\[2pt]
\noindent\textbf{\small Commentaire philologique, métrique et scénique (vv. 1–42 --- partie 5/8)}\\[2pt]
\begin{adjustbox}{max totalheight=15.2cm, max width=16.5cm, keepaspectratio}%
\parbox[t]{16.5cm}{
\setlength{\columnsep}{9pt}%
\begin{multicols}{2}
\fontsize{7.6pt}{8.8pt}\selectfont
\setlength{\parindent}{0pt}%
\setlength{\parskip}{1.2pt}%
\noindent \textbf{(4)} (Dans son Αἴας Λοκρός.) Fr. 14 de l'édition de Stefan Radt, \italique{Tragicorum Graecorum Fragmenta}, vol. 4, édition corrigée et augmentée de 1999. Il figure aussi parmi les fragments incertains d'Euripide (n° 888b) dans l'édition Kannicht des \italique{Tragicorum Graecorum Fragmenta}, vol 5.
\noindent Le ton de la scholie est d'une telle assurance doublée d'une forte érudition qu'on serait tenté de l'en croire. Preuve que la scholie est de Didyme, selon le critère de Pierre Boudreaux ? (\italique{Le texte d'Aristophane et ses commentateurs} (1920). Boudreaux, brillant helléniste fauché bien jeune par la Grande guerre, réussit le tour de force d'écrire un bouquin -- publié de façon posthume -- passionnant quoique traitant de sujets à priori fort arides.) Après tout, Euripide est coutumier du fait : Aulu-Gelle (\italique{Nuits attiques}, XIII, 18) nous rapporte entre autres exemples que le vers 193 des \italique{Bacchantes}, γέρων γέροντα παιδαγωγήσω σ᾽ ἐγώ, se retrouve aussi à l'identique chez Sophocle (fr. 695 Radt) ; de même, Aristote (Περὶ ποιητικῆς, XIII, 1458b22) cite une reprise d'Eschyle par Euripide.
\noindent À rebours de leur querelle dans les \italique{Grenouilles}, Aristote cite un vers iambique d'Eschyle jugé plat en raison du recours à un mot courant tandis qu'Euripide le reprit en l'agrémentant d'un mot emprunté et rare : οἷον τὸ αὐτὸ ποιήσαντος ἰαμβεῖον Αἰσχύλου καὶ Εὐριπίδου, ἓν δὲ μόνον ὄνομα μεταθέντος, ἀντὶ κυρίου εἰωθότος γλῶτταν, τὸ μὲν φαίνεται καλὸν τὸ δ᾽ εὐτελές. Αἰσχύλος μὲν γὰρ ἐν τῷ Φιλοκτήτῃ ἐποίησε·\textbf{φαγέδαιναν ἥ μου σάρκας ἐσθίει ποδός,} ὁ δὲ ἀντὶ τοῦ `ἐσθίει' τὸ `θοινᾶται' μετέθηκεν.
\noindent\textbf{22-4} \textbf{πῶς ἂν οὖν ... προσμάθοιμι χωλὸς εἶναι τὼ σκέλει;} Aristophane joue sur un détournement de proverbe équivalent à notre «~Qui fréquente le boiteux apprend à boiter~» : Παρὰ χωλὸν οἰκῶν κἂν ἐπισκάζειν μαθήσῃ\marqueur{1}.
\noindent Le boiteux -- ou plutôt le \italique{faiseur de boiteux}, χωλοποιός\marqueur{2} -- c'est bien-sûr Euripide et sa fâcheuse prédilection pour ce genre de personnages\marqueur{3}.
\noindent La pointe de la plaisanterie de Mnésiloque est rejetée à la toute fin ; quand il précise vouloir être boiteux τὼ σκέλει, \italique{des deux jambes}.
\noindent Le scholiaste a raison :
\citComm{χωλὸς εἶναι· ἵνα μὴ περιπατήσῃ μετ᾽ αὐτοῦ.}
\noindent Si Mnésiloque venait à boiter des deux jambes, il ne serait plus forcé de suivre Euripide. Le grec aime à préciser que l'on est rompu \italique{des deux jambes} pour souligner une fatigue intense\marqueur{4}.
\noindent (Voir là-dessus la \nieme{3} note d'Alphonse Willems, reproduite en annexe, qui justifie la suppression du vers 23 -- une glose explicative du suivant.) Mais Euripide ne s'en laisse pas compter et lui ordonne de lui emboîter le pas (βάδιζε δευρὶ).\\
\noindent\textbf{[Notes ad 22-4]} \textbf{(1)} C'est le proverbe n° 1674 du \italique{Dictionnaire des sentences latines \& grecques}. Recensé par Macaire au \nieme{XIV} siècle dans son anthologie Ῥοδωνιά (VI, 90 de l'édition de Leutsch \& Schneidewim du \italique{Corpus Paroemiographorum Graecorum}, vol. II) ainsi que sous d'autres formes proches par Michel Apostolius, Plutarque et un scholiaste de Pindare.
\noindent \textbf{(2)} \italique{Les Grenouilles}, vers 846.
\end{multicols}
}%
\end{adjustbox}%
\vfill
\end{adjustwidth}
\newpage
% --- PAGE DE DROITE : Traduction française + Commentaire (6/8) ---
\begin{adjustwidth}{-2cm}{-2cm}
\noindent\textbf{\footnotesize Traduction française --- Folio 01 (\texttt{f. 162v}, vv. 1–42 [3/4]) \hfill ΠΡΟΛΟΓΟΣ}\\[-2pt]
\noindent\rule{16.5cm}{0.5pt}\\[2pt]
{\fontsize{8.8pt}{10.2pt}\selectfont
\ligneEst{\textbf{Mnésiloque}}{\francais{\espace{5.3cm}Ah ! que ne puis-je}}{}
\ligneEst{}{\francais{en outre apprendre aussi à être perclus des jambes ?}}{}
\ligneEst{\textbf{Euripide}}{\francais{Avance ici et prête-moi ton attention.}}{25}
\ligneEst{\textbf{Mnésiloque}}{\francais{\espace{5.2cm}J'y suis.}}{}
\ligneEst{\textbf{Euripide}}{\francais{Tu vois cette petite porte ?}}{}
\ligneEst{\textbf{Mnésiloque}}{\francais{\espace{3.4cm}Oui, par Héraclès,}}{}
\ligneEst{}{\francais{à ce que je crois, du moins.}}{}
\ligneEst{\textbf{Euripide}}{\francais{\espace{1.4cm}Maintenant, motus.}}{}
\ligneEst{\textbf{Mnésiloque}}{\francais{\espace{3.4cm}Je ne dis mot de la petite porte.}}{}
\ligneEst{\textbf{Euripide}}{\francais{Écoute !}}{}
\ligneEst{\textbf{Mnésiloque}}{\francais{\espace{1.4cm}J'écoute et ne dis mot de la petite porte.}}{}
\ligneEst{\textbf{Euripide}}{\francais{C'est là que se trouve demeurer l'illustre Agathon,}}{}
\ligneEst{}{\francais{le poète tragique.}}{30}
}
\vspace{2pt}
\noindent\rule{16.5cm}{0.6pt}\\[2pt]
\noindent\textbf{\small Commentaire philologique, métrique et scénique (vv. 1–42 --- partie 6/8)}\\[2pt]
\begin{adjustbox}{max totalheight=15.2cm, max width=16.5cm, keepaspectratio}%
\parbox[t]{16.5cm}{
\setlength{\columnsep}{9pt}%
\begin{multicols}{2}
\fontsize{7.6pt}{8.8pt}\selectfont
\setlength{\parindent}{0pt}%
\setlength{\parskip}{1.2pt}%
\noindent \textbf{(3)} Ou autres aussi peu reluisants. Cf. les \italique{Acharniens} v. 411 sq. : Télèphe, Bellérophon, Philoctète le mendiant ; voir aussi la \italique{Paix} v. 147 ; les \italique{Oiseaux} v. 1293 et le fragment 53 de l'édition Kassel-Austin des \italique{Poetae Comici Graeci}, vol. III 2 (1984).
\noindent \textbf{(4)} Voir par exemple la \italique{Paix} et la douloureuse ascension de l'orchestre à la scène de Trygée : ἔγωγέ τοι πεπόνηκα κομιδῇ τὼ σκέλει (v. 820) et encore : ἤλγουν τὼ σκέλει μακρὰν ὁδὸν διεληλυθώς(v. 825).
\noindent Nous ne sommes pas dans la chanson de Nadaud où souhaiter claudiquer des deux côtés, c'est souhaiter ne plus claudiquer :
\citComm{Si je marche plus droit qu'un autre\\
C'est que je boite des deux pieds.}
\noindent\textbf{27-8} \textbf{σίγα νυν --- σιωπῶ τὸ θύριον --- ἄκου᾽ --- ἀκούω καὶ σιωπῶ τὸ θύριον.} Le trait d'esprit de Mnésiloque nous échappe.
\noindent\textbf{29} \textbf{Ἀγάθων.} Voici enfin la réponse à la question de Mnésiloque (v. 4) sur la destination de leur équipée : la demeure d'Agathon, le fameux poète tragique\marqueur{1} – sauf pour Mnésiloque –, rendu célèbre par sa victoire au concours tragique des Lénéennes de 416\marqueur{2}. Platon immortalisa dans son \italique{Banquet} la réception qu'il fit chez lui le lendemain de cette victoire. Le rôle d'Agathon dans les \italique{Thesmophories} est semblable à celui d'Euripide dans les \italique{Acharniens} (v. 410-489) : il fournit le héros en matériel de travestissement. Euripide confirmera cette parenté aux vers 173-174\marqueur{3}. Après la mort d'Agathon, Aristophane lui rendit un hommage posthume dans les \italique{Grenouilles} aux vers 83-85.
\noindent\textbf{[Notes ad 29]} \textbf{(1)} On n'a conservé d'Agathon que des fragments et quelques témoignages. Ils ont été regroupés par Snell \& Kannicht au premier volume des \italique{Tragicorum Graecorum Fragmenta} (1986).
\noindent \textbf{(2)} Athénée, \italique{Deipnosophistes}, V, 216f-217a.
\noindent \textbf{(3)} Aristophane n'est d'ailleurs pas le seul à rapprocher Euripide \& Agathon : ainsi fait Aristote dans sa \italique{Poétique} (1456a25-32) : Agathon introduisit des interludes (ἐμβόλιμα) chantés par le chœur sans rapport avec l'action, suivant la voie tracée par Euripide.
\noindent\textbf{30-3} \textbf{ποῖος οὗτος Ἁγάθων; ... οὐχ ἑόρακας πώποτε.} La répartition des répliques de ce passage fluctue d'une édition à l'autre.
\noindent Tâchons d'y voir plus clair !\\ À l'instar du manuscrit de Ravenne, nous laissons à Mnésiloque le début du vers 31 : ἔστιν τις Ἀγάθων;  La clef de compréhension est la suivante : lisons :
\citComm{ἔστιν τις ἀγαθῶν;}
\noindent Mnésiloque joue sur le nom d'Agathon ; \latin{nomen omen}, demande-t-il : Agathon est-il bien l'\italique{honnête homme} que suggère son nom ? Or, pour une nature rustre comme celle de Mnésiloque, un bon-homme se doit d'être un solide gaillard, à la peau bronzée par le travail en extérieur et non pas au teint fromage-blanc comme ces désœuvrés de fémoïdes !
\end{multicols}
}%
\end{adjustbox}%
\vfill
\end{adjustwidth}
\newpage

% --- FOLIO 01 (f. 162v) : SPREAD 2.4 (Texte grec établi + Traduction française + Commentaire 7/8 & 8/8) ---
\clearpage\ifodd\value{page}\hbox{}\newpage\fi
\enTeteG{ΘΕΣΜΟΦΟΡΙΑΖΟΥΣΑΙ --- Texte grec établi (f. 162v, vv. 1–42 [4/4])}
\enTeteD{LES THESMOPHORIES --- Traduction française (f. 162v, vv. 1–42 [4/4])}
\begin{adjustwidth}{-2cm}{-2cm}
\noindent\textbf{\footnotesize Texte grec établi --- Folio 01 (\texttt{f. 162v}, vv. 1–42 [4/4]) \hfill ΠΡΟΛΟΓΟΣ}\\[-2pt]
\noindent\rule{16.5cm}{0.5pt}\\[2pt]
{\fontsize{8.8pt}{10.2pt}\selectfont
\ligneEst{\textbf{Μνησίλοχος}}{\decalage{ὁ τραγῳδοποιός.}ποῖος οὗτος Ἁγάθων;}{}
\ligneEst{}{ἔστιν τις ἀγαθῶν; μῶν ὁ μέλας ὁ καρτερός;}{}
\ligneEst{\textbf{Εὐριπίδης}}{οὔκ, ἀλλ᾽ ἕτερός τις· οὐχ ἑόρακας πώποτε;}{}
\ligneEst{\textbf{Μνησίλοχος}}{μῶν ὁ δασυπώγων;}{}
\ligneEst{\textbf{Εὐριπίδης}}{\decalage{μῶν ὁ δασυπώγων;}οὐχ ἑόρακας πώποτε.}{}
\ligneEst{\textbf{Μνησίλοχος}}{μὰ τὸν Δί᾽ οὔτοι γ᾽ ὥστε κἀμέ γ᾽ εἰδέναι.}{}
\ligneEst{\textbf{Εὐριπίδης}}{καὶ μὴν †βεβίνηκας† σύ γ᾽, ἀλλ᾽ οὐκ οἶσθ᾽ ἴσως.}{35}
\ligneEst{}{ἀλλ᾽ ἐκποδὼν πτήξωμεν, ὡς ἐξέρχεται}{}
\ligneEst{}{θεράπων τις αὐτοῦ πῦρ ἔχων καὶ μυρρίνας}{}
\ligneEst{}{προθυσόμενος ἔοικε τῆς ποιήσεως.}{}
\ligneEst{\textbf{Θεράπων} {\tiny\textit{dim. anap.}}}{εὔφημος πᾶς ἔστω λαός,}{}
\ligneEst{}{στόμα συγκλῄσας· ἐπιδημεῖ γὰρ}{40}
\ligneEst{}{θίασος Μουσῶν ἔνδον μελάθρων}{}
\ligneEst{{\tiny\textit{(catalect.)}}}{τῶν δεσποσύνων μελοποιῶν.}{42}
}
\vspace{2pt}
\noindent\rule{16.5cm}{0.6pt}\\[2pt]
\noindent\textbf{\small Commentaire philologique, métrique et scénique (vv. 1–42 --- partie 7/8)}\\[2pt]
\begin{adjustbox}{max totalheight=15.2cm, max width=16.5cm, keepaspectratio}%
\parbox[t]{16.5cm}{
\setlength{\columnsep}{9pt}%
\begin{multicols}{2}
\fontsize{7.6pt}{8.8pt}\selectfont
\setlength{\parindent}{0pt}%
\setlength{\parskip}{1.2pt}%
\noindent Il continue sur sa lancée : μῶν ὁ μέλας ὁ καρτερός; « n'est-ce pas le noiraud, le costaud ? » (Agathon est en réalité tout le contraire d'un rustaud ; Aristophane donne ainsi facilement à rire à son public aux dépens de l'original bien connu.)\\ Le vers suivant est écrit ainsi dans le manuscrit de Ravenne :
\citComm{οὔκ ἀλλ᾽ ἕτερός τις --- οὐχ \textbf{ἑόρακα} πώποτε}
\noindent Le signe mis pour signifier un changement d'interlocuteur\marqueur{1} ayant par la suite été gratté (\latin{ut videtur}).
\noindent L'entièreté du vers est effectivement à restituer à Euripide : une scholie écrite juste au-dessus du vers l'atteste :
\citComm{\textbf{ἐπειδὴ} οὐ πάλαι ἤρξατο διδάσκειν, ἀλλὰ τρισὶν πρὸ τούτων ἔτεσιν.\\
\textbf{Puisqu}'il n'a pas commencé récemment de faire jouer des pièces, mais depuis trois ans. [sept, en réalité, et selon la façon grecque de compter]}
\noindent Cette scholie est évidemment là pour expliquer l'étonnement d'Euripide : οὐχ ἑόρακα\textbf{ς} πώποτε; « Tu ne l'as donc jamais vu ?! [alors que cela fait trois année qu'il monte des pièces] » ; le scholiaste ne l'aurait pas ainsi rédigée (en commençant par ἐπειδὴ...) si, au vers 32, le `οὐχ ἑόρακα πώποτε' était prononcé par Mnésiloque, « je ne l'ai jamais vu ».\\ Cette interprétation, outre qu'elle est anciennement corroborée par la scholie, est aussi celle qui permet la lecture la plus drôle : οὐχ ἑόρακας πώποτε; --- μῶν ὁ δασυπώγων; --- οὐχ ἑόρακας πώποτε, « Tu ne l'as jamais vu ?! --- N'est-il pas barbu ? --- Tu ne l'as jamais vu !.. », bien meilleure par le rythme et la répétition déconvenue d'Euripide que l'échange premièrement transcrit sur le manuscrit de Ravenne où Mnésiloque commençait par avouer son ignorance : (Μνη.) οὐχ ἑόρακα πώποτε. μῶν ὁ δασυπώγων; --- (Εὐρ.) οὐχ ἑόρακας πώποτε.
\noindent\textbf{[Notes ad 30-3]} \textbf{(1)} Non un tiret, \italique{paragraphos}, ni un double point ` : ', dont l'usage est courant dans notre manuscrit lorsque le changement d'interlocuteur intervient au milieu d'un vers, mais en l'occurrence une indication de personnage : μ\auDessus{ν}{Η}.
\noindent\textbf{35} \textbf{καὶ μὴν βεβίνηκας σύ γ᾽, ἀλλ᾽ οὐκ οἶσθ᾽ ἴσως.} Le ton est inhabituel -- voire plutôt choquant -- dans la bouche d'Euripide.
\noindent Qu'à cela ne tienne : puisqu'il semble d'autre part qu'Agathon tourne le dos lorsqu'il accorde ses faveurs\marqueur{1}, nos commentateurs en concluent qu'il est logique que Mnésiloque ne le connaisse pas -- sauf éventuellement au sens biblique... de Sodome, s'entend.
\noindent Et c'est le grave et docte Euripide qui colporterait ce genre de ragots.
\noindent Y a-t-il besoin de commenter ce genre de commentaires ?.. je précise : ces commentaires \italique{modernes} ; car le βεβίνηκας du \italique{Ravennas} n'a pas laissé derrière lui la moindre scholie ancienne ; chose étonnante lorsqu'on sait que les anciens commentateurs étaient presqu'autant portés au scabreux que nos modernes.
\noindent Il faut croire que le texte original d'Aristophane, l'\allemand{Urtext} des \italique{Thesmophories}, comportait une tout autre leçon : une banalité sur laquelle aucun commentateur ancien ne se sera arrêté.
\end{multicols}
}%
\end{adjustbox}%
\vfill
\end{adjustwidth}
\newpage
% --- PAGE DE DROITE : Traduction française + Commentaire (8/8) ---
\begin{adjustwidth}{-2cm}{-2cm}
\noindent\textbf{\footnotesize Traduction française --- Folio 01 (\texttt{f. 162v}, vv. 1–42 [4/4]) \hfill ΠΡΟΛΟΓΟΣ}\\[-2pt]
\noindent\rule{16.5cm}{0.5pt}\\[2pt]
{\fontsize{8.8pt}{10.2pt}\selectfont
\ligneEst{\textbf{Mnésiloque}}{\francais{\espace{2.4cm}Qui est-ce, cet Agathon ?}}{}
\ligneEst{}{\francais{Est-ce l'un de ces honnêtes hommes... celui au teint hâlé, plein d'allant ?}}{}
\ligneEst{\textbf{Euripide}}{\francais{Non, c'est un autre. Tu ne l'as jamais vu ?}}{}
\ligneEst{\textbf{Mnésiloque}}{\francais{L'homme à la barbe touffue ?}}{}
\ligneEst{\textbf{Euripide}}{\francais{\espace{3cm}Tu ne l'as jamais vu !..}}{}
\ligneEst{\textbf{Mnésiloque}}{\francais{Non, par Zeus, pas du moins que je sache.}}{}
\ligneEst{\textbf{Euripide}}{\francais{À coup sûr tu l'as enfourché, mais peut-être ne t'en souvient-il pas.}}{35}
\ligneEst{}{\francais{Mais allons nous tapir à l'écart. Car sort de chez lui}}{}
\ligneEst{}{\francais{un serviteur avec du feu et des myrtes.}}{}
\ligneEst{}{\francais{Il va sacrifier, ce semble, en faveur du succès de sa poésie.}}{}
\ligneEst{\textbf{Serviteur d'Agathon}}{\francais{Que le peuple entier se recueille}}{}
\ligneEst{}{\francais{et ferme la bouche, car est venu séjourner}}{40}
\ligneEst{}{\francais{un thiase de muses dans la demeure}}{}
\ligneEst{}{\francais{de mon maître pour composer des chants.}}{42}
}
\vspace{2pt}
\noindent\rule{16.5cm}{0.6pt}\\[2pt]
\noindent\textbf{\small Commentaire philologique, métrique et scénique (vv. 1–42 --- partie 8/8)}\\[2pt]
\begin{adjustbox}{max totalheight=15.2cm, max width=16.5cm, keepaspectratio}%
\parbox[t]{16.5cm}{
\setlength{\columnsep}{9pt}%
\begin{multicols}{2}
\fontsize{7.6pt}{8.8pt}\selectfont
\setlength{\parindent}{0pt}%
\setlength{\parskip}{1.2pt}%
\noindent Las !.. on ne détruit que ce que l'on remplace ; à défaut d'avoir une proposition alternative solidement étayée, il faut nous résoudre à \latin{crux}-désespérer...
\noindent\textbf{37} \textbf{μυρρίνας} désigne les \italique{rameaux de myrte} que l'on brûlait lors des sacrifices (cf. les \italique{Guêpes}, v. 860 ; les \italique{Oiseaux} au v. 43).
\noindent\textbf{[Notes ad 35]} \textbf{(1)} Voir le vers 258 et sa scholie ; Agathon se livrerait au travestissement dans l'interlope \italique{Athens by night} pour y être éprouvé à la manière d'une femme -- mais le commentaire déluré du scholiaste n'est peut-être qu'une mauvaise interprétation du présent vers 35 adjointe à une forte dose d'imagination lubrique. Ce qui est sûr, c'est qu'au vers 205, Agathon dit craindre s'il était reconnu d'être lynché pour avoir chapardé nuitamment les œuvres de la Cypris : les Athéniennes voient en lui un concurrent irrégulier.
\noindent\textbf{39-62} \textbf{Passage en dimètres anapestiques.} (Ceux aux vers 42 et 62 sont catalectiques.) Ce mètre est d'usage courant chez Aristophane pour les prières (cf. les \italique{Guêpes} v. 323-333 et la \italique{Paix} v. 947-1015). Le serviteur d'Agathon est tout empreint de la préciosité du maître, comme, du reste, le serviteur de l'Euripide des \italique{Acharniens} (v. 396-400) l'était.
\noindent\textbf{39-40} \textbf{εὔφημος πᾶς ἔστω λέως / στόμα συγκλῄσας·} Ce genre de formule est classique au début d'une prière ou d'un sacrifice. On peut néanmoins constater ici un peu plus d'emphase avec le tour tragique\marqueur{1} et l'utilisation de l'adjectif εὔφημος plutôt que le classique verbe εὐφημεῖν\marqueur{2} ou substantif εὐφημία\marqueur{3}.
\noindent\textbf{40} \textbf{ἐπιδημεῖ.} \italique{Séjourner}, tel me paraît bien le sens de ἐπιδημεῖ, \latin{peregre advenit et adest}. (Xénophon, \italique{Mémor.}, I, 2, 61.) Cf. ἐπιχωρτάζει, \latin{venit ibique commoratur}. Mελοποιῶν a le sens du futur (voir note [de Willems] sur les \italique{Grenouilles} au v. 148, t. III p. 109-110). (Note de Willems.)
\noindent\textbf{[Notes ad 39-40]} \textbf{(1)} C'est ce que souligne le scholiaste : παρὰ τοῖς τραγικοῖς οὕτω λέγεται, ``<εὔφημος> πᾶς ἔστω λαός''.
\noindent \textbf{(2)} \italique{Acharniens} v. 241 ; \italique{Cavaliers} v. 1316 εὐφημεῖν χρὴ καὶ στόμα κλείεν ; \italique{Nuées} v. 263 ; \italique{Paix} v. 434 et 1316 ; \italique{Grenouilles} v. 354.
\noindent \textbf{(3)} \italique{Guêpes} v. 868 ; \italique{Oiseaux} v. 959.
\noindent\textbf{41-2} \textbf{θίασος Μουσῶν ἔνδον μελάθρων / τῶν δεσποσύνων μελοποιῶν.} Pour bien goûter ces vers, il faut savoir ce qu'est un \italique{thiase} : une bruyante troupe de gens célébrant des rites (bachiques la plupart du temps) et parcourant les rues. Cf. les \italique{Bacchantes} d'Euripide où des femmes munies de tambourins interprètent des danses scandées de cris rituels (évohé) avec une fougue poussée jusqu'à la frénésie. Par extension, ce mot très prisé des tragiques désigne une foule extrêmement bruyante ; Euripide l'emploie à propos d'une bande de centaures (\italique{Iphigénie à Aulis}, v. 1059) ; une troupe endiablée de jeunes danseuses (\italique{Iphigénie en Tauride}, v. 1146) ou encore une meute de guerriers (\italique{Phaéton}, v. 796).
\noindent Pour décrire cette multitude de Muses, au lieu de recourir à la banale image de l'\italique{essaim}\marqueur{1}, à dessein Aristophane préfère évoquer un thiase :
\noindent comme le remarque Sommerstein, l'inversion est totale : un lieu d'un cortège célébrant une divinité, c'est un thiase de Muses qui vient honorer Agathon.
\noindent\textbf{[Notes ad 41-2]} \textbf{(1)} Le Socrate des \italique{Nuées} (vers 297) parle d'un « essaim de déesses », μέγα γάρ τι θεῶν κινεῖται σμῆνος. Voir Jean Taillardat, \italique{Les images d'Aristophane}, §662.
\end{multicols}
}%
\end{adjustbox}%
\vfill
\end{adjustwidth}
\newpage
\clearpage\ifodd\value{page}\hbox{}\newpage\fi
\enTeteG{ΠΡΟΛΟΓΟΣ --- Codex Ravennas 429, f. 163r (Classense 165r) --- vv. 43–79}
\begin{adjustwidth}{-1.5cm}{-3.0cm}
\begin{adjustbox}{max totalheight=21.0cm, max width=16.5cm, keepaspectratio}
\begin{minipage}{16.5cm}
{\fontsize{7.2pt}{8.5pt}\selectfont
\noindent\textbf{Folio 02 --- Codex Ravennas 429, \texttt{f. 163r} (\texttt{Classense 165r}) : ΠΡΟΛΟΓΟΣ, vv. 43–79}\hfill \textit{Témoins :} $\text{R}^1, \Rpc, \Rdeux$ (\Mon)\par\vspace{2pt}
\hrule height 0.4pt\vspace{2.5pt}
}
{\fontsize{8.1pt}{9.5pt}\selectfont
\ligneRav{43}{θερ/}{\hspace{1.1em}\auDessus{ἐχέτω}{\tiny ἐπεχέτω} δὲ πνοὰς νήνεμος αἰθήρ,}{\textbf{43} \lecon{ἐχέτω} R (avec la glose interne ἐπεχέτω)}
\ligneRav{44}{}{\hspace{1.1em}κῦμα δὲ πόντου μὴ κελαδείτω}{\textbf{44} \lecon{δὲ} R : τε Austin, cf. v.46, 47}
\ligneRav{45}{—}{\hspace{1.1em}γλαυκόν — : βομβάξ· : σίγα· τί λέγει\inRas{ς}{};}{\textbf{45} λέγει\inRas{ς}{} \Rpc}
\ligneRav{}{θερ/}{\hspace{1.1em}πτηνῶν τε γένη κατακοιμάσθω,}{}
\ligneRav{}{}{\hspace{1.1em}θηρῶν τ᾽ ἀγρίων πόδες ὑλοδρόμων}{}
\ligneRav{}{—}{\hspace{1.1em}μὴ λυέσθων. : βομβαλοβομβάξ·}{}
\ligneRav{}{θερ/}{\hspace{1.1em}μέλλει γὰρ ὁ καλλιεπὴς Ἀγάθων}{}
\ligneRav{50}{—}{\hspace{1.1em}\difference{πρανοῦν προτέρους} — : μῶν βινεῖσθαι;}{\textbf{50} πράμος R : πρόμος Grynaeus : \lecon{πρανοῦν} Bibi ap. Hesych. π 3188 \espApp ἡμέτερος R : \lecon{προτέρους} Bibi, cf. \italique{Nubes} v.562}
\ligneRav{}{θερ/}{\hspace{1.1em}τίς ὁ φωνήσας; : νήνεμος αἰθήρ·}{}
\ligneRav{}{θερ/}{\hspace{1.1em}δρυόχους, τιθέναι δράματος ἀρχάς·}{}
\ligneRav{53}{}{\hspace{1.1em}κάμπτει δὲ νέας \difference{ἁψῖδας} ἐπῶν,}{\textbf{53} ἀσπίδας R : \lecon{ἁψῖδας} Suid. δ 1547, λ 181}
\ligneRav{}{}{\hspace{1.1em}τὰ δὲ τορνεύει, τὰ δὲ κολλομελεῖ,}{}
\ligneRav{55}{}{\hspace{1.1em}καὶ γνωμοτυπεῖ κἀντονομάζει}{}
\ligneRav{56}{}{\hspace{1.1em}καὶ κηροχυτεῖ καὶ \difference{γογγύλλει}}{\textbf{56} γογγύλίζει R, Suid. γ 361 : \lecon{γογγύλλει} Porson : γογγυλίει Hermann (ad \italique{Nub.} 992)}
\ligneRav{}{—}{\hspace{1.1em}καὶ χοανεύει— : καὶ λαικάζει·}{}
\ligneRav{58}{θερ/}{\hspace{1.1em}τίς \difference{ἀγροιώτας} πελάθει θριγκοῖς;}{\textbf{58} ἀγριώτας R : \lecon{ἀγροιώτας} anonym. Parisinus}
\ligneRav{}{μ\auDessus{ν}{Η}}{\hspace{1.1em}ὃς ἕτοιμος σοῦ τοῦ τε ποιητοῦ}{}
\ligneRav{60}{}{\hspace{1.1em}τοῦ καλλιεποῦς \expo{κατὰ} τοῦ θριγκοῦ}{\textbf{60} κατὰ om. \Rac : add. \Rpc}
\ligneRav{61}{}{\hspace{1.1em}\difference{συγγογγύλας} καὶ συστρέψας}{\textbf{61} γογγυλίσας R : \lecon{συγγογγύλας} Enger}
\ligneRav{}{}{\hspace{1.1em}τουτὶ τὸ πέος χοανεῦσαι·}{}
\ligneRav{}{θερ/}{ἦ που νέος γ᾽ ὢν ἦσθ᾽ ὑβριστὴς ὦ γέρον·}{}
\ligneRav{}{\auDessus{εὐ}{ρ}}{ὦ δαιμόνιε τοῦτον μὲν ἔα χαίρειν, σὺ δὲ}{}
\ligneRav{65}{}{Ἀγάθωνά μοι δεῦρ᾽ ἐκκάλεσον πάσῃ τέχνῃ·}{}
\ligneRav{}{θερ/}{μηδὲν ἱκέτευ᾽· αὐτὸς γὰρ ἔξεισιν τάχα·}{}
\ligneRav{}{}{καὶ γὰρ μελοποιεῖν ἄρχεται· χειμῶνος οὖν}{}
\ligneRav{}{}{ὄντος κατακάμπτειν τὰς στροφὰς οὐ ῥᾴδιον,}{}
\ligneRav{69}{}{ἢν μὴ προίῃ θύρασι πρὸς τὸν ἥλιον·}{\textbf{69} \lecon{θύρασι} R : θύραζε Zanetti}
\ligneRav{70}{μ\auDessus{ν}{Η}}{τί οὖν ἐγὼ δρῶ; : περίμεν᾽, ὡς ἐξερχεται·}{}
\ligneRav{71}{μ\auDessus{ν}{Η}}{ὦ Ζεῦ τί δρᾶσαι διανοεῖ με τήμερον;}{\textbf{71-4} (Μνη.) (71) et (Εὐρ.) (72-4) R : \lecon{(Μνη.) (71-4)} Bibi ap. Σ : alii alia}
\ligneRav{}{}{νὴ τοὺς θεοὺς ἐγὼ πυθέσθαι βούλομαι}{}
\ligneRav{}{}{τί τὸ πρᾶγμα τουτί. τί στένεις; τί δυσφορεῖς;}{}
\ligneRav{74}{}{οὐ χρῆν σε κρύπτειν ὄντα κηδεστὴν ἐμόν·}{\textbf{74} \lecon{χρῆν} R, Suid. κ 1487 codd. FV : χρή Suid. codd. AGM \espApp \lecon{ἐμόν} R : ἐμέ Valckenaer ap. Dobree (p. 237)}
\ligneRav{75}{\auDessus{εὐ}{ρ}}{ἔστιν κακόν μοι μέγα τι προπεφυραμένον·}{}
\ligneRav{}{μ\auDessus{ν}{Η}}{ποῖόν τι; : τῇδε θἠμέρᾳ κριθήσεται}{}
\ligneRav{}{}{εἴτ᾽ ἔστ᾽ ἔτι ζῶν εἴτ᾽ ἀπόλωλ᾽ Εὐριπίδης·}{}
\ligneRav{}{μ\auDessus{ν}{Η}}{καὶ πῶς; ἐπεὶ νῦν γ᾽ οὔτε τὰ δικαστήρια}{}
\ligneRav{}{}{μέλλει δικάζειν οὔτε βουλῆς ἐσθ᾽ ἕδρα,}{}
}
\vspace{3pt}\noindent\rule{4.5cm}{0.4pt}\par\vspace{1.5pt}
{\fontsize{7.0pt}{8.3pt}\selectfont
\noindent\textbf{Loci paralleli (vv. 43–79) ---} \textbf{43-5} (ἐχέτω -- γλαυκόν) Suidas ε 3792 ; Zonaras 941 ἐχέτω \espTest \textbf{44} (κελαδείτω) Suidas κ 1278 \espTest \textbf{45} (βομβάξ) Hesychius β 787 βομβάξ ; Suidas ε 2807 ; \Eustathe 1718.9 \espTest \textbf{48} (βομβαλοβομβάξ) Suidas ε 2807 ; \Eustathe 790.13, 1718.13 \espTest \textbf{49} (καλλιεπὴς) Philostrate \italique{Vit. soph.} I 9.1 [$=$ II 12.8 Kayser] \espTest \textbf{50} (πραῖνοι) Hesychius π 3188 †πραῖνοι· πρηνίζειν, καταστρέφειν \espTest \textbf{52-7} Suidas δ 1547 (sed τορεύει pro τορνεύει, χιανεύει pro χοανεύει) \espTest \textbf{52} Zonaras 569 δρύοχοι \espTest \textbf{53} (κάμπτει) Hesychius κ 617 κάμπτειν \espTest \textbf{54} (τορνεύει) Hesychius τ 1168 τορνεύει \espTest \textbf{56} (γογγύλλει) Hesychius γ 764 γογγύλλειν \espTest Suidas γ 361 $=$ Zonaras 449 γογγυλίζει \espTest \textbf{57} (χοανεύει) \Eustathe 1153.40 \espTest \textbf{58} (θριγκοῖς) Hesychius θ 750 θριγκοῖς· φραγμοῖς \espTest \textbf{62} (χοανεῦσαι $=$ χωνεῦσαι) Moeris χ 21 ; Hesychius χ 575 χοανεῦσαι \espTest \textbf{74} Suidas κ 1487, cf. Moeris κ 38\par
}
\vspace{2.5pt}\noindent\rule{4.5cm}{0.4pt}\par\vspace{1.5pt}
{\fontsize{6.9pt}{8.2pt}\selectfont
\noindent\textbf{Suite (en commentaire) de l'apparat critique \& notes paléographiques (\texttt{f. 163r}) ---} \textbf{v. 43 :} Glose interlinéaire/interne dans R : ἐπεχέτω glosant ἐχέτω. --- \textbf{v. 45 :} Grattage du sigma final dans λέγει\inRas{ς}{} par \Rpc{}. --- \textbf{v. 50 :} Lieu célèbre de corruption paléographique (πράμος ἡμέτερος R), magistralement résolu par Louis Hauseux (\textbf{Bibi}). --- \textbf{v. 60 :} La préposition κατὰ, omise par \Rac{}, a été rajoutée par \Rpc{} au-dessus de la ligne. --- \textbf{v. 43 :} \lecon{ἐχέτω} R (avec la glose interne ἐπεχέτω). --- \textbf{v. 44 :} \lecon{δὲ} R (maintenu par \textbf{Bibi}) : τε Austin (cf. vv. 46, 47). --- \textbf{v. 50 :} πράμος R (πρόμος Grynaeus) : \lecon{πρανοῦν} \lecon{Bibi} ap. Hésychius $\pi\ 3188$ (et πραῖνοι dans l'\italique{Index}) \espApp ἡμέτερος R : \lecon{προτέρους} \lecon{Bibi}, cf. \italique{Nuées} v. 562. \textbf{Explication paléographique (Commentaire ad v. 50) :} Le texte primitif πρανοῦν προτέρους (ou πρανοι ν) a subi : 1° une médivision avec confusion minuscule $\nu \to \mu$ et $\iota \to \varsigma$ donnant πράμος ; 2° la syllabe προ- de προτέρους, écrite au moyen du signe tachygraphique en forme de Ʒ inversé (semblable à un $\varepsilon$, cf. Gardthausen II p. 340), a été rattachée au $\nu$ précédent et lue <η>μ- + -ετερος $\to$ ἡμέτερος ! --- \textbf{v. 53 :} ἀσπίδας R (« boucliers », faute par attraction d'un mot militaire banal) : \lecon{ἁψῖδας} Souda $\delta\ 1547, \lambda\ 181$ (« cintres, voûtes d'hymnes »). --- \textbf{v. 60 :} κατὰ om. \Rac{} : \textbf{add.} \Rpc{}. --- \textbf{v. 69 :} \lecon{θύρασι} R (maintenu par \textbf{Bibi}) : θύραζε Zanetti. --- \textbf{vv. 71–74 :} (Μνη.) (71) et (Εὐρ.) (72–74) R : \lecon{(Μνη.) (71–74)} \lecon{Bibi} d'après $\Sigma^{\text{R}}$ (toute la tirade appartient à Mnésiloque).\par
}
\end{minipage}
\end{adjustbox}
\end{adjustwidth}
\newpage
\enTeteD{Codex Ravennas 429 (Biblioteca Classense) --- Fac-similé du f. 163r (vv. 43–79)}
\begin{adjustwidth}{-1.5cm}{-3.0cm}
\begin{center}
\includegraphics[width=16.5cm,height=20.5cm,keepaspectratio]{/workspaces/Thesmophories/dossier_ravennas_thesmophories/folios_individuels_ff162v-177r/02_f163r_classense165r_vv0043-0079.jpg}\par
\vspace{2pt}
{\scriptsize\textbf{Codex Ravennas 429, f. 163r (Classense 165r)} --- ΠΡΟΛΟΓΟΣ, \textbf{vv. 43–79} ($\text{R}^1, \Rpc, \Rdeux$ (\Mon)).}
\end{center}
\end{adjustwidth}
\newpage
\clearpage\ifodd\value{page}\hbox{}\newpage\fi% --- FOLIO 02 (f. 163r) : SPREAD 2.1 (Texte grec établi + Traduction française + Commentaire 1/8 & 2/8) ---
\clearpage\ifodd\value{page}\hbox{}\newpage\fi
\enTeteG{ΘΕΣΜΟΦΟΡΙΑΖΟΥΣΑΙ --- Texte grec établi (f. 163r, vv. 43–79 [1/4])}
\enTeteD{LES THESMOPHORIES --- Traduction française (f. 163r, vv. 43–79 [1/4])}
\begin{adjustwidth}{-2cm}{-2cm}
\noindent\textbf{\footnotesize Texte grec établi --- Folio 02 (\texttt{f. 163r}, vv. 43–79 [1/4]) \hfill ΠΡΟΛΟΓΟΣ}\\[-2pt]
\noindent\rule{16.5cm}{0.5pt}\\[2pt]
{\fontsize{8.8pt}{10.2pt}\selectfont
\ligneEst{}{ἐχέτω δὲ πνοὰς νήνεμος αἰθήρ,}{43}
\ligneEst{}{κῦμα δὲ πόντου μὴ κελαδείτω}{}
\ligneEst{}{γλαυκόν—}{45}
\ligneEst{\textbf{Μνησίλοχος}}{\decalage{γλαυκόν—}βομβάξ.}{}
\ligneEst{\textbf{Εὐριπίδης}}{\decalage{γλαυκόν—βομβάξ.}σίγα. τί λέγει;}{}
\ligneEst{\textbf{Θεράπων}}{πτηνῶν τε γένη κατακοιμάσθω,}{}
\ligneEst{}{θηρῶν τ᾽ ἀγρίων πόδες ὑλοδρόμων}{}
\ligneEst{}{μὴ λυέσθων.}{}
\ligneEst{\textbf{Μνησίλοχος}}{\decalage{μὴ λυέσθων.}βομβαλοβομβάξ.}{}
\ligneEst{\textbf{Θεράπων}}{μέλλει γὰρ ὁ καλλιεπὴς Ἀγάθων}{}
\ligneEst{}{πρανοῦν προτέρους—}{50}
}
\vspace{2pt}
\noindent\rule{16.5cm}{0.6pt}\\[2pt]
\noindent\textbf{\small Commentaire philologique, métrique et scénique (vv. 43–79 --- partie 1/8)}\\[2pt]
\begin{adjustbox}{max totalheight=15.2cm, max width=16.5cm, keepaspectratio}%
\parbox[t]{16.5cm}{
\setlength{\columnsep}{9pt}%
\begin{multicols}{2}
\fontsize{8.4pt}{9.7pt}\selectfont
\setlength{\parindent}{0pt}%
\setlength{\parskip}{1.2pt}%
\noindent\textbf{43-7} \textbf{ἐχέτω δὲ πνοὰς νήνεμος αἰθήρ ...} De même que dans les \italique{Bacchantes} d'Euripide (v. 1084), à la voix de Dionysos « l'éther se tut, le ravin herbeux étouffa le murmure des feuilles, et l'on n'eût point entendu le cri des bêtes sauvages ».
\citComm{σίγησε δ᾽ αἰθήρ, σῖγα δ᾽ ὕλιμος νάπη\\
φύλλ᾽ εἶχε, θηρῶν δ᾽ οὐκ ἂν ἤκουσας βοήν.}
\noindent Et aussi dans Théocrite (II, 38), tout se tait à l'arrivée d'Hécate. (Note de Willems.)
\noindent\textbf{43} \textbf{ἐχέτω} est mis pour ἐπεχέτω (ici : \italique{retenir}) comme l'indique la scholie intérieure au texte du \italique{Ravennas}.
\noindent\textbf{45} \textbf{βομβάξ.} Forgée sur le modèle d'ὄμπαξ -- interjection pour réclamer le silence « paix ! (πάξ) », « assez ! » --, βομβάξ fait entendre le \italique{bourdonnement}, βόμβος (ὁ βομβυλιός désigne tout type d'insecte bourdonnant). Plaisante manière pour Aristophane de ridiculiser l'image détournée du serviteur : le bruyant thiase et sa pompe de muses venues honorer la demeure d'Agathon sont ramenées à la plus commune métaphore de la \italique{ruche}.
\noindent\textbf{49-50} \textbf{μέλλει γὰρ ὁ καλλιεπὴς Ἀγάθων / πράμος ἡμέτερος.} πράμος n'est ni autrement attesté  ni même relevé par les grammairiens et doit nécessairement être scandé $\cup \cup $ pour rentrer dans le mètre anapestique, empêchant tout rapprochement étymologique pertinent.
\noindent Autant d'indices qui font soupçonner une erreur textuelle plutôt qu'un hapax agathonien.
\noindent πρόμος, \italique{le chef}, fut proposé par Grynaeus et quasi-unanimement repris depuis.
\noindent La correction est très légère mais n'est pas satisfaisante :
\noindent il manque au terme πρόμος une acception \italique{sensu obsceno} pour expliquer l'exclamation de Mnésiloque : μῶν βινεῖσθαι;\marqueur{1}
\noindent De plus,  l'on s'attendrait plutôt à trouver un infinitif réclamé par μέλλει.
\noindent Une glose d'Hésychius va nous permettre de débrouiller l'écheveau.
\noindent Cette glose est elle-même légèrement corrompue (π 3188) :
\citComm{†πραῖνοι†· πρηνίζειν, καταστρέφειν.}
\noindent Sans nous attarder pour l'instant sur les hypothèses de corruptions, constatons que Hésychius atteste l'existence du verbe πρᾳνόω (ou πρανόω\marqueur{2}).
\noindent Pour le sens, πρηνίζειν signifie \italique{précipiter}, \italique{renverser}, \italique{jeter} (en avant).
\end{multicols}
}%
\end{adjustbox}%
\vfill
\end{adjustwidth}
\newpage
% --- PAGE DE DROITE : Traduction française + Commentaire (2/8) ---
\begin{adjustwidth}{-2cm}{-2cm}
\noindent\textbf{\footnotesize Traduction française --- Folio 02 (\texttt{f. 163r}, vv. 43–79 [1/4]) \hfill ΠΡΟΛΟΓΟΣ}\\[-2pt]
\noindent\rule{16.5cm}{0.5pt}\\[2pt]
{\fontsize{8.8pt}{10.2pt}\selectfont
\ligneEst{}{\francais{Que retienne ses haleines l'éther paisible,}}{43}
\ligneEst{}{\francais{que le flot de la mer point ne retentisse}}{}
\ligneEst{}{\francais{glauque...}}{45}
\ligneEst{\textbf{Mnésiloque}}{\francais{\espace{1.5cm}Bombax !}}{}
\ligneEst{\textbf{Euripide}}{\francais{\espace{2.8cm}Chut ! que dit-il ?}}{}
\ligneEst{\textbf{Serviteur}}{\francais{Que la gent ailée s'assoupisse,}}{}
\ligneEst{}{\francais{que les pieds des bêtes sauvages vaguant par les bois}}{}
\ligneEst{}{\francais{point ne se meuvent...}}{}
\ligneEst{\textbf{Mnésiloque}}{\francais{\espace{1.9cm}Bombalo bombax !}}{}
\ligneEst{\textbf{Serviteur}}{\francais{... Car il s'apprête l'homme au beau langage, Agathon,}}{}
\ligneEst{}{\francais{à aller au devant du gaillard...}}{50}
}
\vspace{2pt}
\noindent\rule{16.5cm}{0.6pt}\\[2pt]
\noindent\textbf{\small Commentaire philologique, métrique et scénique (vv. 43–79 --- partie 2/8)}\\[2pt]
\begin{adjustbox}{max totalheight=15.2cm, max width=16.5cm, keepaspectratio}%
\parbox[t]{16.5cm}{
\setlength{\columnsep}{9pt}%
\begin{multicols}{2}
\fontsize{8.4pt}{9.7pt}\selectfont
\setlength{\parindent}{0pt}%
\setlength{\parskip}{1.2pt}%
\noindent Le serviteur d'Agathon construit ainsi sa phrase :
\citComm{μέλλει γὰρ Ἁγάθων πρανοῦν δρυόχους\marqueur{3}.\\
Agathon s'apprête à jeter les fondations\marqueur{3}.}
\noindent Seulement, πρηνίζειν peut désigner aussi l'action de s'accroupir tête baissée\marqueur{4}. (πρανοῦν, comme accroupir, fait entendre la \italique{croupe}, πρανής.) C'est bien-sûr dans ce sens détourné que le prend Mnésiloque -- image d'un Agathon se jetant à son travail dans la position du penseur de Rodin ? L'occasion est trop belle : Mnésiloque s'imagine Agathon en position κύβδα (cf. le vers 489) et tournant le dos à d'éventuels assaillants --- μῶν βινεῖσθαι; demande-t-il. Le véritable Agathon devait effectivement travailler dans une posture qui prêtait à rire aux mauvais esprits : c'est sur son lit et dans cette même posture qu'il apparaît sur scène quelques vers plus loin ; cette position étrange lui vaudra d'être comparé par Mnésiloque à Cyrène \italique{aux douze postures} (δωδεκαμήχανος) (v. 98).\\ Reste à expliquer la présence dans le \italique{Ravennas} de †ἡμέτερος† après †πράμος†.
\noindent Hésychius glose †πραῖνοι† par deux infinitifs, πρηνίζειν et καταστρέφειν, preuve que le mot relevé en vedette devait originellement être lui aussi un infinitif. Le rétablissement de πρανοῦν se fait alors sans difficulté.
\noindent C'est la même mécoupure du verbe πρανοῦν, lu †πρανοι \ ν† qui pourrait expliquer le †πράμος† et le début de †ἡμέτερος† du \italique{Ravennas} selon le \italique{stemma} suivant :
\par\vspace{1pt}\noindent\adjustbox{max width=\columnwidth}{$\displaystyle\begin{aligned}
\overset{\overset{\text{{\allemand{Urtext}}}}{\text{πρηνίζειν, καταστρέφειν}}}{\text{πρανοῦν [...]}}\xrightarrow{\text{mécoup. \ } + \text{\ conf. \  υ/ι}}\overset{\overset{\text{{Archétype}}}{\text{πρηνίζειν, καταστρέφειν}}}{\text{πρανοι \ \ ν[...]}}&\xrightarrow{\text{\ \ \ \ \ }}\overset{\text{{Hésychius π 3188}}}{\text{πραῖνοι· πρηνίζειν, καταστρέφειν}}\\
& \xrightarrow{\text{confus. \ ν/μ et ι/ς}}\overset{\text{{\italique{Ravennas}}}}{\text{πραμος <η>μ[ετερος]}}
\end{aligned}$}\par\vspace{1pt}
\noindent Rien de plus semblable qu'un ι et un ς, ainsi qu'un μ et un ν  -- dans les manuscrits écrits en minuscules à l'instar du \italique{Ravennas}.
\end{multicols}
}%
\end{adjustbox}%
\vfill
\end{adjustwidth}
\newpage

% --- FOLIO 02 (f. 163r) : SPREAD 2.2 (Texte grec établi + Traduction française + Commentaire 3/8 & 4/8) ---
\clearpage\ifodd\value{page}\hbox{}\newpage\fi
\enTeteG{ΘΕΣΜΟΦΟΡΙΑΖΟΥΣΑΙ --- Texte grec établi (f. 163r, vv. 43–79 [2/4])}
\enTeteD{LES THESMOPHORIES --- Traduction française (f. 163r, vv. 43–79 [2/4])}
\begin{adjustwidth}{-2cm}{-2cm}
\noindent\textbf{\footnotesize Texte grec établi --- Folio 02 (\texttt{f. 163r}, vv. 43–79 [2/4]) \hfill ΠΡΟΛΟΓΟΣ}\\[-2pt]
\noindent\rule{16.5cm}{0.5pt}\\[2pt]
{\fontsize{8.8pt}{10.2pt}\selectfont
\ligneEst{\textbf{Μνησίλοχος}}{\decalage{πραῖνοι ἡμέτερος—}μῶν βινεῖσθαι;}{}
\ligneEst{\textbf{Θεράπων}}{τίς ὁ φωνήσας;}{}
\ligneEst{\textbf{Μνησίλοχος}}{\decalage{τίς ὁ φωνήσας;}νήνεμος αἰθήρ.}{}
\ligneEst{\textbf{Θεράπων}}{δρυόχους, τιθέναι δράματος ἀρχάς.}{}
\ligneEst{}{κάμπτει δὲ νέας ἁψῖδας ἐπῶν,}{}
\ligneEst{}{τὰ δὲ τορνεύει, τὰ δὲ κολλομελεῖ,}{}
\ligneEst{}{καὶ γνωμοτυπεῖ κἀντονομάζει}{55}
\ligneEst{}{καὶ κηροχυτεῖ καὶ γογγύλλει}{}
\ligneEst{}{καὶ χοανεύει—}{}
\ligneEst{\textbf{Μνησίλοχος}}{\decalage{καὶ χοανεύει—}καὶ λαικάζει.}{}
\ligneEst{\textbf{Θεράπων}}{τίς ἀγροιώτας πελάθει θριγκοῖς;}{}
}
\vspace{2pt}
\noindent\rule{16.5cm}{0.6pt}\\[2pt]
\noindent\textbf{\small Commentaire philologique, métrique et scénique (vv. 43–79 --- partie 3/8)}\\[2pt]
\begin{adjustbox}{max totalheight=15.2cm, max width=16.5cm, keepaspectratio}%
\parbox[t]{16.5cm}{
\setlength{\columnsep}{9pt}%
\begin{multicols}{2}
\fontsize{8.4pt}{9.7pt}\selectfont
\setlength{\parindent}{0pt}%
\setlength{\parskip}{1.2pt}%
\noindent Il s'agit maintenant de comprendre la présence de †ετερος†. Sans tenter d'explications, je constate qu'au vers 562 des \italique{Nuées}, les deux principaux manuscrits des pièces d'Aristophane, notre \italique{Ravennas} et le \italique{Venetus}, portent deux leçons différentes : ἑτέρας \italique{versus} προτέρας\marqueur{5}.
\noindent Je suppose ici la même corruption :
\citComm{μέλλει γὰρ ὁ καλλιεπὴς Ἀγάθων πρανοῦν προτέρους δρυόχους, τιθέναι δράματος ἀρχάς. κάμπτει δὲ νέας ἁψῖδας ἐπῶν κτλ.\\
Agathon au beau langage s'apprête à jeter les premiers fondements, à disposer les prémices d'un drame. Il tourne de nouvelles roues de mots etc.}
\noindent\textbf{[Notes ad 49-50]} \textbf{(1)} Sommerstein voit en πρόμος une légère déformation suggérant à l'oreille πόρνος... « Notre chef de file --- Celui qu'on enfile ? » Mais que c'est loin en grec !.. Certains ont voulu justifier le rapprochement πρόμος -- πόρνος en arguant de la plaisanterie de \italique{Nuées} (v. 394) sur l'indissociabilité entre βροντή et πορδή -- certes pas au niveau de la prononciation. Nous avons profité des annexes auxquelles nous renvoyons le lecteur pour exégéter ce jeu de mots qui roule sur la langue proverbiale.
\noindent \textbf{(2)} Même genre de variantes orthographiques entre πρᾷος et πρᾶος.
\noindent \textbf{(3)} Les δρύοχοι sont les \italique{chantiers}, ou poutres sur lesquelles s'établit la quille d'un vaisseau en construction. On rapprochera la métaphore de \italique{Guêpes} 30 : τὴν τρόπιν τοῦ πράγματος, « la base de l'affaire ». (Note de Taillardat §758 n. 1 ; voir aussi §673.)
\noindent \textbf{(4)} Cf. Hésychius π 3199, π 3269 πρηνιχθέντα· ἐπὶ πρόσωπον καὶ ἐπὶ στόμα πεσόντα ~ Photius π 1138 et L'Etymologicum genuinum : πρηνίξαι πρηνιχθῆναι καὶ πεπρηνίχθαι· γνὺξ γάρ ἐστιν κυρίως τὸ εἰς γόνυ πεσεῖν ... τὸ δὲ πρηνιχθῆναι ἐπὶ τὸ πρόσωπον [πεσεῖν] et d'autres références dans l'abondant et récent apparat des témoins de l'édition Theodoridis de Photius.
\noindent \textbf{(5)} Les écrits anciens comportaient de nombreuses abréviations. (Cf. par exemple la très ingénieuse reconstruction du vers 80 par Nauck.) Par exemple, προσ était souvent abrégé d'une sorte de Ʒ renversé qui ressemblait à un epsilon : \renverse{Ʒ}.
\noindent Voir G. Bady, \italique{Paléographie grecque : éléments de cours}, ou bien V. Gardthausen, \italique{Griechische palaeographie}\expo{2}, II (1913), p. 340.
\end{multicols}
}%
\end{adjustbox}%
\vfill
\end{adjustwidth}
\newpage
% --- PAGE DE DROITE : Traduction française + Commentaire (4/8) ---
\begin{adjustwidth}{-2cm}{-2cm}
\noindent\textbf{\footnotesize Traduction française --- Folio 02 (\texttt{f. 163r}, vv. 43–79 [2/4]) \hfill ΠΡΟΛΟΓΟΣ}\\[-2pt]
\noindent\rule{16.5cm}{0.5pt}\\[2pt]
{\fontsize{8.8pt}{10.2pt}\selectfont
\ligneEst{\textbf{Mnésiloque}}{\francais{\espace{2.6cm}Et se faire enfiler ?}}{}
\ligneEst{\textbf{Serviteur}}{\francais{Qui a parlé?}}{}
\ligneEst{\textbf{Mnésiloque}}{\francais{\espace{2.3cm}L'éther paisible.}}{}
\ligneEst{\textbf{Serviteur}}{\francais{... du gaillard d'avant ; il va composer les débuts d'un drame.}}{}
\ligneEst{}{\francais{Il cambre, en jantes nouvelles, des vers.}}{}
\ligneEst{}{\francais{Il met les uns sur le tour, il assemble les autres.}}{}
\ligneEst{}{\francais{Il frappe des sentences, fait s'opposer des mots,}}{55}
\ligneEst{}{\francais{malaxe la cire, l'entoure d'un moule,}}{}
\ligneEst{}{\francais{puis introduit le godet...}}{}
\ligneEst{\textbf{Mnésiloque}}{\francais{\espace{2.2cm}Et avale.}}{}
\ligneEst{\textbf{Serviteur}}{\francais{Quel est le manant qui approche cette enceinte ?}}{}
}
\vspace{2pt}
\noindent\rule{16.5cm}{0.6pt}\\[2pt]
\noindent\textbf{\small Commentaire philologique, métrique et scénique (vv. 43–79 --- partie 4/8)}\\[2pt]
\begin{adjustbox}{max totalheight=15.2cm, max width=16.5cm, keepaspectratio}%
\parbox[t]{16.5cm}{
\setlength{\columnsep}{9pt}%
\begin{multicols}{2}
\fontsize{7.6pt}{8.8pt}\selectfont
\setlength{\parindent}{0pt}%
\setlength{\parskip}{1.2pt}%
\noindent\textbf{52-7} \textbf{τιθέναι δράματος ἀρχάς... καὶ χοανεύει.} Agathon le tragique n'est pas un génie bâtisseur, un τέκτων ; il ne crée pas : il fabrique, il bricole. Il est donc tout-à-tour charpentier de marine, menuisier, modeleur : « il va poser des chantiers, fondements d'un drame. etc. » Ces métaphores sont évidemment burlesques ; elles avilissent délibérément l'art d'Agathon, ravalé au rang de métier et ne peuvent que tourner le poète en dérision...\\ Agathon est donc à mille lieues du mystère de l'inspiration, décrit par Platon (dans l'\italique{Ion}) : κοῦφον γὰρ χρῆμα ποιητής ἐστιν καὶ πτηνὸν καὶ ἱερόν, « le poète est chose légère, ailée, sacrée ». De fait, Platon connaît bien les images du menuisier et du modeleur\marqueur{1} mais jamais il ne les applique au poète (ἀποτορνεύειν est dit de Lysias, \italique{Phèdre} 243e : ἀκριβῶς ἕκαστα τῶν ὀνομάτων ἀποτετόρνευται et reste péjoratif. La même observation vaut pour κολλᾶν, en parlant d'un auteur, \italique{Ménexène} 263b ; \italique{Phèdre} 278e). À noter que ces images semblent banales ; ainsi, dans le \italique{Timée} (81b), on trouve : οἷον ἐκ δρυόχων, « comme depuis les longerons de la carcasse », c'est-à-dire \italique{entièrement}. κάμπτειν : \italique{courber}, \italique{plier} désigne en musique les courbes capricieuses d'un air ou les inflexions de la voix. (Notes de Taillardat\marqueur{2}.)
\noindent\textbf{[Notes ad 52-7]} \textbf{(1)} Voir Pierre Louis, \italique{Les Métaphores de Platon}, Belles Lettres, 1945, p. 98, 170, 204-5.
\noindent \textbf{(2)} \italique{Les images d'Aristophane}, §758 et §784.
\noindent\textbf{56-7} \textbf{καὶ κηροχυτεῖ καὶ γογγύλίζει / καὶ χοανεύει — καὶ λαικάζει.} La métaphore du serviteur est tirée de la technique de la \italique{cire perdue}\marqueur{1} ; un procédé de moulage pour obtenir une statue en métal à partir d'un modèle en cire : on commence par modeler à la cire la statue (κηροχυτεῖν), modèle qu'on recouvre d'une gangue d'argile qui fera office de moule. C'est à cette opération que renvoie probablement le verbe γογγυλίζειν (corr. γογγύλλειν), associant la gangue à une \italique{rave ronde}, γογγυλίς.
\noindent Il faut laisser à ce moule d'argile des cheminées, certaines pour insérer des entonnoirs et y couler le métal en fusion (χοανεύειν) qui va prendre la place de la cire, et d'autres cheminées pour permettre à la cire fondue d'être évacuée.
\noindent Cet esprit tordu de Mnésiloque y voit la description d'une scène érotique.
\noindent Il interrompt d'une obscénité: « καὶ λαικάζει ».
\noindent Ce verbe a le sens de \latin{fellare}\marqueur{2}.
\noindent\textbf{[Notes ad 56-7]} \textbf{(1)} Voir sur cette technique l'article de C. C. Mattusch « The Berlin Foundry Cup: The Casting of Greek Bronze Statuary in the Early Fifth Century B. C. », \italique{American Journal of Archaeology}, 84 n° 4 (1980), pp. 435-444, ou son livre \italique{Greek Bronze Statuary: From the Beginnings Through the Fifth Century B.C.} (1988).
\end{multicols}
}%
\end{adjustbox}%
\vfill
\end{adjustwidth}
\newpage

% --- FOLIO 02 (f. 163r) : SPREAD 2.3 (Texte grec établi + Traduction française + Commentaire 5/8 & 6/8) ---
\clearpage\ifodd\value{page}\hbox{}\newpage\fi
\enTeteG{ΘΕΣΜΟΦΟΡΙΑΖΟΥΣΑΙ --- Texte grec établi (f. 163r, vv. 43–79 [3/4])}
\enTeteD{LES THESMOPHORIES --- Traduction française (f. 163r, vv. 43–79 [3/4])}
\begin{adjustwidth}{-2cm}{-2cm}
\noindent\textbf{\footnotesize Texte grec établi --- Folio 02 (\texttt{f. 163r}, vv. 43–79 [3/4]) \hfill ΠΡΟΛΟΓΟΣ}\\[-2pt]
\noindent\rule{16.5cm}{0.5pt}\\[2pt]
{\fontsize{8.8pt}{10.2pt}\selectfont
\ligneEst{\textbf{Μνησίλοχος}}{ὃς ἕτοιμος σοῦ τοῦ τε ποιητοῦ}{}
\ligneEst{}{τοῦ καλλιεποῦς κατὰ τοῦ θριγκοῦ}{60}
\ligneEst{}{συγγογγύλας καὶ συστρέψας}{}
\ligneEst{{\tiny\textit{(catalect.)}}}{τουτὶ τὸ πέος χοανεῦσαι.}{}
\ligneEst{\textbf{Θεράπων} {\tiny\textit{3 ia.}}}{ἦ που νέος γ᾽ ὢν ἦσθ᾽ ὑβριστὴς ὦ γέρον.}{}
\ligneEst{\textbf{Εὐριπίδης}}{ὦ δαιμόνιε τοῦτον μὲν ἔα χαίρειν, σὺ δὲ}{}
\ligneEst{}{Ἀγάθωνά μοι δεῦρ᾽ ἐκκάλεσον πάσῃ τέχνῃ.}{65}
\ligneEst{\textbf{Θεράπων}}{μηδὲν ἱκέτευ᾽· αὐτὸς γὰρ ἔξεισιν τάχα.}{}
\ligneEst{}{καὶ γὰρ μελοποιεῖν ἄρχεται· χειμῶνος οὖν}{}
\ligneEst{}{ὄντος κατακάμπτειν τὰς στροφὰς οὐ ῥᾴδιον,}{}
\ligneEst{}{ἢν μὴ προίῃ θύρασι πρὸς τὸν ἥλιον.}{}
}
\vspace{2pt}
\noindent\rule{16.5cm}{0.6pt}\\[2pt]
\noindent\textbf{\small Commentaire philologique, métrique et scénique (vv. 43–79 --- partie 5/8)}\\[2pt]
\begin{adjustbox}{max totalheight=15.2cm, max width=16.5cm, keepaspectratio}%
\parbox[t]{16.5cm}{
\setlength{\columnsep}{9pt}%
\begin{multicols}{2}
\fontsize{7.6pt}{8.8pt}\selectfont
\setlength{\parindent}{0pt}%
\setlength{\parskip}{1.2pt}%
\noindent \textbf{(2)} \italique{Confer} par exemple Chantraine, \italique{Dictionnaire étymologique de la langue grecque}. γογγύλλειν peut avoir le sens de \italique{rendre compact}, \italique{serré}. Cf. la \nieme{3} note §166 de Taillardat : \italique{ramassé}, \italique{compact}, \italique{serré} sont les sens que Taillardat donne à l'adjectif γογγύλος aux vers 25-28 de la \italique{Paix}.
\noindent\textbf{61} \textbf{γογγυλίσας καὶ συστρέψας.} Vers retouché \italique{gratia metri} d'après un passage de \italique{Lysistrata} (v. 975), jouée probablement l'année précédente (cf. Introduction).
\noindent Allons-y voir d'un peu plus près pour comprendre précisément l'image de Mnésiloque : Kinésias est dans un état lamentable après que Myrrhine l'ait fortement échaudé... et finalement laissé raide. Le demi-chœur des hommes en vient à souhaiter un traitement de faveur pour cette s[célérate] (v. 972-979) :
\citComm{μιαρὰ δῆτ' ὦ Ζεῦ Ζεῦ· |
εἴθ' αὐτὴν ὥσπερ τοὺς θωμοὺς |
μεγάλῳ τυφῷ καὶ πρηστῆρι |
\textbf{ξυστρέψας καὶ ξυγγογγύλας} |
οἴχοιο φέρων, εἶτα μεθείης, |
ἡ δὲ φέροιτ' αὖ πάλιν ἐς τὴν γῆν, |
κᾆτ' ἐξαίφνης |
περὶ τὴν ψωλὴν περιβαίη.\\
Une gueuse ! ô Zeus, Zeus, puisse-t-elle, comme les tas de blé dans un coup de vent et une tourmente, roulée par toi [Kinésias] et ramassée en boule, être portée dans l'air, puis, lâchée, retomber à terre et du coup enfourcher ce priape.\\
\textbf{975} ξυγγογγυλίσας R, Γ, Suid. : \lecon{ξυγγογγύλας} Cobet}
\noindent L'image est tirée de la technique du foulage et du vannage du grain. Voici ce qu'en dit le Daremberg \& Saglio\marqueur{1} :
\citComm{Foulage (ἀλόησις) --- Les Grecs ne battaient les céréales ni en grange, ni en plein air, car ils n'avaient point de fléau pour faire sortir le grain de l'épi. Ils employaient la vieille méthode égyptienne du dépiquage ou foulage, qui n'est plus guère en usage que chez les Provençaux et les Languedociens. En juin, c’est-à-dire, à l'apparition d’Orion\marqueur{2}, on portait les gerbes sur une aire, qualifiée de μεγάλη\marqueur{3} , ἐϋκτίμενη\marqueur{4}, ἐϋτροχάλος\marqueur{5}, ἱερά\marqueur{6}, et construite, plus ou moins, d'après le modèle de l'aire sacrée d'Éleusis\marqueur{7}. Des ouvriers, ἐπαλώστες\marqueur{8}, dressent les gerbes a côté les unes des autres en les inclinant un peu vers le centre et on fait entrer les animaux : bœufs, chevaux, mulets, parfois même des ânes\marqueur{9} ; ils sont accouplés deux à deux et marchent en cercle, d’abord au pas, puis au trot, sous la direction d'un homme, placé au centre de l'aire, qui tient en main les guides, δεσμά\marqueur{10}. D'autres ouvriers, armés de fourches en bois, τρίναξ\marqueur{11}, régularisent la foulaison, ὁμαλιεῖται ὁ ἁλοητός, en retournant la paille et en la ramenant sous les pieds des animaux, pour faciliter le manège, δῖνος.\\
Nettoyage du grain (λίκμησις)  — Les Grecs ne connaissaient point le vannage proprement dit, opération de nos batteurs en grange, qui nettoye simplement le grain des glumes ou glumelles et que l'on exécute près de la porte d’une grange, dans un léger courant d'air, mais à l'abri du vent. Au contraire, les Grecs attendaient que le vent fût assez fort\marqueur{12} pour emporter les fétus de paille et leur permît de pratiquer le \italique{ventage}, qui peut se faire de deux façons :  1°  Avec une corbeille, λικμός, πλόκανον [\latin{Ventilabrum}] [...]    2°   Avec une pelle en bois, πτύον [\latin{Pala}]. C'est notre \italique{nettoyage à la roue} ; on jette circulairement contre le vent et, aussi loin que possible, à la hauteur d’un mètre au moins, une pelletée de grains qui retombent sur le sol, tandis que la paille et les balles sont emportées à l'extrémité de l'aire.}
\end{multicols}
}%
\end{adjustbox}%
\vfill
\end{adjustwidth}
\newpage
% --- PAGE DE DROITE : Traduction française + Commentaire (6/8) ---
\begin{adjustwidth}{-2cm}{-2cm}
\noindent\textbf{\footnotesize Traduction française --- Folio 02 (\texttt{f. 163r}, vv. 43–79 [3/4]) \hfill ΠΡΟΛΟΓΟΣ}\\[-2pt]
\noindent\rule{16.5cm}{0.5pt}\\[2pt]
{\fontsize{8.8pt}{10.2pt}\selectfont
\ligneEst{\textbf{Mnésiloque}}{\francais{Quelqu'un qui s'apprête à te mettre en boule et te retourner,}}{}
\ligneEst{}{\francais{toi et ton poète au beau langage,}}{60}
\ligneEst{}{\francais{et à introduire par votre enceinte}}{}
\ligneEst{}{\francais{cet outil-ci.}}{}
\ligneEst{\textbf{Serviteur}}{\francais{Dans ta jeunesse tu devais être un ribaud, vieillard..}}{}
\ligneEst{\textbf{Euripide}}{\francais{(À Mnésiloque :) Diantre d'homme, laisse-le tranquille. (Au serviteur :) Toi,}}{}
\ligneEst{}{\francais{va m'appeler Agathon bon gré, mal gré.}}{65}
\ligneEst{\textbf{Serviteur}}{\francais{Inutile d'insister. Lui-même ne tardera pas à sortir,}}{}
\ligneEst{}{\francais{car il commence un chant lyrique. Et de fait, en hiver,}}{}
\ligneEst{}{\francais{il n'est pas facile de contourner ses strophes,}}{}
\ligneEst{}{\francais{sans aller se mettre dehors au soleil.}}{}
}
\vspace{2pt}
\noindent\rule{16.5cm}{0.6pt}\\[2pt]
\noindent\textbf{\small Commentaire philologique, métrique et scénique (vv. 43–79 --- partie 6/8)}\\[2pt]
\begin{adjustbox}{max totalheight=15.2cm, max width=16.5cm, keepaspectratio}%
\parbox[t]{16.5cm}{
\setlength{\columnsep}{9pt}%
\begin{multicols}{2}
\fontsize{8.4pt}{9.7pt}\selectfont
\setlength{\parindent}{0pt}%
\setlength{\parskip}{1.2pt}%
\noindent Ces longues mais utiles précisions ayant été apportées, le passage de \italique{Lysistrata} est maintenant clair : Myrrhine doit être traitée comme un fétu de paille qu'on retourne et compactifie en boule (ξυστρέψας καὶ ξυγγογγύλας\marqueur{13}) avant de le jeter (εἶτα μεθείης) pour qu'il jonche l'aire vide ; ici ἡ ψωλή $=$ τὸ κενόν\marqueur{8}, mais avec un double sens obscène : le [phallos] \italique{nu}, c'est-à-dire décalotté.\\ Revenons aux \italique{Thesmophories}. Le serviteur demande qui est le paysan (τίς ἀγροιώτας) qui se permet d'interrompre de la sorte son poème. Mnésiloque lui répond donc par une image agricole certainement populaire : il se prépare à le retourner et le mettre en boule (συγγογγύλας καὶ συστρέψας v. 61) aussi rustiquement que le chœur des vieillards voulait faire avec Myrrhine.
\noindent\textbf{[Notes ad 61]} \textbf{(1)} \italique{Dictionnaire des antiquités grecques et romaines}, IV-2, pp. 906-907. Article « \latin{Rustica Res} ».
\noindent \textbf{(2)} Hésiode, \italique{Travaux et les jours}, v. 596.
\noindent \textbf{(3)} \italique{Iliade}, XIII, v. 588.
\noindent \textbf{(4)} \italique{Iliade}, XX, v. 496.
\noindent \textbf{(5)} Hésiode \italique{Travaux et les jours}, v. 597.
\noindent \textbf{(6)} \italique{Iliade}, V, v. 499.
\noindent \textbf{(7)} Pausanias, I, 38.
\noindent \textbf{(8)} Xénophon, \italique{Économiq.}, XVIII.
\noindent \textbf{(9)} \italique{Anthol. palat.} IX, 301.
\noindent \textbf{(10)} \italique{Anthol. palat.} VI, 104.
\noindent \textbf{(11)} \italique{Anthol. palat.} VI, 95.
\noindent \textbf{(12)} \italique{Iliade}, V, vv. 499 sq. ; XIII, vv. 588 sq. ; cf. \italique{Anth. palat.}, VI, 53.
\noindent \textbf{(13)} Soit dit en passant, c'est certainement l'un de ces deux passages que glose Hésychius (γ 764) : γογγύλλειν· συστρέφειν.
\noindent\textbf{67-8} \textbf{χειμῶνος οὖν / ὄντος κατακάμπτειν τὰς στροφὰς οὐ ῥᾴδιον.} Les Thesmophories avaient lieu du 11 au 13 du mois de Pyanepsion, jours qui tombent en général pendant notre seconde moitié du mois d'octobre.
\noindent Le serviteur continue de filer sa métaphore : la cire est plus facile à modeler lorsqu'elle est exposée au Soleil.
\end{multicols}
}%
\end{adjustbox}%
\vfill
\end{adjustwidth}
\newpage

% --- FOLIO 02 (f. 163r) : SPREAD 2.4 (Texte grec établi + Traduction française + Commentaire 7/8 & 8/8) ---
\clearpage\ifodd\value{page}\hbox{}\newpage\fi
\enTeteG{ΘΕΣΜΟΦΟΡΙΑΖΟΥΣΑΙ --- Texte grec établi (f. 163r, vv. 43–79 [4/4])}
\enTeteD{LES THESMOPHORIES --- Traduction française (f. 163r, vv. 43–79 [4/4])}
\begin{adjustwidth}{-2cm}{-2cm}
\noindent\textbf{\footnotesize Texte grec établi --- Folio 02 (\texttt{f. 163r}, vv. 43–79 [4/4]) \hfill ΠΡΟΛΟΓΟΣ}\\[-2pt]
\noindent\rule{16.5cm}{0.5pt}\\[2pt]
{\fontsize{8.8pt}{10.2pt}\selectfont
\ligneEst{\textbf{Μνησίλοχος}}{τί οὖν ἐγὼ δρῶ;}{70}
\ligneEst{\textbf{Εὐριπίδης}}{\decalage{τί οὖν ἐγὼ δρῶ;}περίμεν᾽, ὡς ἐξερχεται.}{}
\ligneEst{\textbf{Μνησίλοχος}}{ὦ Ζεῦ τί δρᾶσαι διανοεῖ με τήμερον;}{}
\ligneEst{}{νὴ τοὺς θεοὺς ἐγὼ πυθέσθαι βούλομαι}{}
\ligneEst{}{τί τὸ πρᾶγμα τουτί. τί στένεις; τί δυσφορεῖς;}{}
\ligneEst{}{οὐ χρῆν σε κρύπτειν ὄντα κηδεστὴν ἐμόν.}{}
\ligneEst{\textbf{Εὐριπίδης}}{ἔστιν κακόν μοι μέγα τι προπεφυραμένον.}{75}
\ligneEst{\textbf{Μνησίλοχος}}{ποῖόν τι;}{}
\ligneEst{\textbf{Εὐριπίδης}}{\decalage{ποῖόν τι;}τῇδε θἠμέρᾳ κριθήσεται}{}
\ligneEst{}{εἴτ᾽ ἔστ᾽ ἔτι ζῶν εἴτ᾽ ἀπόλωλ᾽ Εὐριπίδης.}{}
\ligneEst{\textbf{Μνησίλοχος}}{καὶ πῶς; ἐπεὶ νῦν γ᾽ οὔτε τὰ δικαστήρια}{}
\ligneEst{}{μέλλει δικάζειν οὔτε βουλῆς ἐσθ᾽ ἕδρα,}{79}
}
\vspace{2pt}
\noindent\rule{16.5cm}{0.6pt}\\[2pt]
\noindent\textbf{\small Commentaire philologique, métrique et scénique (vv. 43–79 --- partie 7/8)}\\[2pt]
\begin{adjustbox}{max totalheight=15.2cm, max width=16.5cm, keepaspectratio}%
\parbox[t]{16.5cm}{
\setlength{\columnsep}{9pt}%
\begin{multicols}{2}
\fontsize{8.4pt}{9.7pt}\selectfont
\setlength{\parindent}{0pt}%
\setlength{\parskip}{1.2pt}%
\noindent\textbf{70-4} \textbf{τί οὖν ἐγὼ δρῶ; ... ὄντα κηδεστὴν ἐμόν.} La répartition des ce passage diffère d'un éditeur à un autre ; chacun voulant apporter son grain de sel. Inscrivons-nous, nous aussi, dans cette grande tradition philologique et apportons notre pierre à l'édifice : une nouvelle répartition des répliques qui, espérons-le, clôturera un débat qui n'a que trop duré. La répartition des deux répliques du vers 70 ne pose aucun problème : il suffit de songer aux caractères des trois personnages en présence : le ton empoté du « τί οὖν ἐγὼ δρῶ; » ne peut convenir qu'au personnage pour l'instant passif et qui subit l'action : Mnésiloque\marqueur{1} (Euripide ruisselle de défauts mais garde toujours une certaine dignité ; jamais il ne s'abaisserait à lancer un mot aussi trivial et décontenancé – ses plaintes ont davantage de lyrisme !) La réponse : « περίμεν᾽, ὡς ἐξερχεται » ne peut être attribuée qu'à Euripide : c'est lui qui a besoin de Mnésiloque ; de plus, le serviteur ne délivre ni conseil ni ordre : il se contente de répondre le plus sèchement possible aux questions (cette répartition est d'ailleurs celle du manuscrit de Ravenne). Les vers 72 à 74 sont évidemment prononcés par Mnésiloque qui cherche à percer la détresse d'Euripide et nous révèle par la même occasion son lien de parenté avec le poète. Reste le vers 71 : il est plus que tentant de l'attribuer à Mnésiloque et faire ainsi des quatre vers 71-74 l'exact pendant des quatre vers 1-4 : une invocation à Zeus suivie d'une question : là pour savoir où Euripide le mène et ici pour savoir ce qui le tourmente. Problème : la présence d'un \italique{signe} (σημεῖον), probablement un \italique{paragraphos}, sur d'anciens manuscrits a été interprétée comme indiquant un changement d'interlocuteurs entre le vers 71 et les vers 72-74. Mais tel n'est pas le cas ici, et c'est ce que nous indiquait une ancienne scholie, très précieuse et incomprise\marqueur{2} :
\citComm{τὸ σημεῖον· οὐ μόνον κηδεστὴς ὁ ἐκδιδούς, ἀλλὰ καὶ ὁ λαμβάνων.\\
\francais{Le signe [typographique] : non seulement le Parent est celui qui abandonne [la parole], mais aussi celui qui la [re]prend.}}
\noindent La scholie ne précise pas à quoi servait ce σημεῖον ; serait-ce une pause ? quoi qu'il en soit, le protagoniste devait certainement faire quelque jeu de théâtre après l'invocation à Zeus du vers 71, identique à celui déjà fait après le premier vers et montrer ainsi aux spectateurs le parallélisme entre les deux quatrains.
\noindent\textbf{[Notes ad 70-4]} \textbf{(1)} Réplique analogue de Mnésiloque au vers 925 : ἐγὼ δ᾽ ὁ κακοδαίμων τί δρῶ;
\end{multicols}
}%
\end{adjustbox}%
\vfill
\end{adjustwidth}
\newpage
% --- PAGE DE DROITE : Traduction française + Commentaire (8/8) ---
\begin{adjustwidth}{-2cm}{-2cm}
\noindent\textbf{\footnotesize Traduction française --- Folio 02 (\texttt{f. 163r}, vv. 43–79 [4/4]) \hfill ΠΡΟΛΟΓΟΣ}\\[-2pt]
\noindent\rule{16.5cm}{0.5pt}\\[2pt]
{\fontsize{8.8pt}{10.2pt}\selectfont
\ligneEst{\textbf{Mnésiloque}}{\francais{Et moi, que vais-je faire ?}}{70}
\ligneEst{\textbf{Euripide}}{\francais{\espace{2.6cm}Demeure, puisqu'il arrive.}}{}
\ligneEst{\textbf{Mnésiloque}}{\francais{Ô Zeus, que prétends-tu faire de moi aujourd'hui ?}}{}
\ligneEst{}{\francais{Par les dieux, je tiens à savoir de quoi il retourne.}}{}
\ligneEst{}{\francais{Qu'as-tu à geindre ? à te contrarier ?}}{}
\ligneEst{}{\francais{Il n'y a pas de mystère à faire avec moi, puisque tu es mon gendre.}}{}
\ligneEst{\textbf{Euripide}}{\francais{Un grand malheur, manigancé d'avance, me menace.}}{75}
\ligneEst{\textbf{Mnésiloque}}{\francais{De quelle sorte ?}}{}
\ligneEst{\textbf{Euripide}}{\francais{\espace{1.7cm}Ce jour décidera}}{}
\ligneEst{}{\francais{si Euripide vit encore ou s'il est mort.}}{}
\ligneEst{\textbf{Mnésiloque}}{\francais{Comment cela ? Car ni les tribunaux}}{}
\ligneEst{}{\francais{ne sont dans le cas de juger, ni la Boulé ne siège,}}{79}
}
\vspace{2pt}
\noindent\rule{16.5cm}{0.6pt}\\[2pt]
\noindent\textbf{\small Commentaire philologique, métrique et scénique (vv. 43–79 --- partie 8/8)}\\[2pt]
\begin{adjustbox}{max totalheight=15.2cm, max width=16.5cm, keepaspectratio}%
\parbox[t]{16.5cm}{
\setlength{\columnsep}{9pt}%
\begin{multicols}{2}
\fontsize{8.4pt}{9.7pt}\selectfont
\setlength{\parindent}{0pt}%
\setlength{\parskip}{1.2pt}%
\noindent \textbf{(2)} Incomprise déjà de Suidas – lexicographe hypothétique du \nieme{IX} ou \nieme{X} siècle, à qui l'on attribue le lexique de la \italique{Souda}. Ce dernier fourmille en références et citations d'Aristophane : la \italique{Souda} est un outil de recherche sur Aristophane indispensable.
\noindent Il en existe une version numérique : \href{https://www.cs.uky.edu/~raphael/sol/sol-html/}{https://www.cs.uky.edu/~raphael/sol/sol-html/} – incomprise donc de Suidas ce qui tend à prouver son ancienneté : Suidas (κ 1487) biffe le lemme `τὸ σημεῖον', pourtant indispensable à la bonne compréhension et égare ainsi les différents éditeurs des scholies qui l'appliquent au vers 74, voire qui la retouchent pour finir de la rendre tout-à-fait inintelligible.
\noindent\textbf{71} \textbf{ὦ Ζεῦ τί δρᾶσαι διανοεῖ με τήμερον;} Formule courante ; se reporter, par exemple, aux \italique{Oiseaux}, vv. 58 \& 62.
\noindent\textbf{75} \textbf{προπεφυραμένον} Pour comprendre cette image, l'on se rapportera à un passage des \italique{Oiseaux} (v. 462-463) et le commentaire qu'en fait Taillardat (§755) :
\citComm{καὶ προπεφύραται λόγος εἷς μοι, / ὃν διαμάττειν οὐ κωλύει\\
\francais{j'ai sous la main un discours déjà bien délayé et rien n'empêche que j'achève de le pétrir.}}
\noindent « Φυρᾶν signifie proprement \italique{délayer} de la farine dans un liquide\marqueur{1}. » Est donc προ-πεφυραμένον ce qui a déjà été délayé à l'avance mais point encore pétri ; autrement dit : le complot est ourdi mais rien n'est encore joué. L'avantage de ce genre de métaphore boulangère est qu'elle se transpose assez aisément en français\marqueur{2}.
\noindent\textbf{[Notes ad 75]} \textbf{(1)} Cf. Hippocrate \italique{Du régime} (\italique{Περί διαίτης}) 2,40 et Thucydide 3,49.
\noindent \textbf{(2)} Comme l'épisode raconté par Démosthène dans les \italique{Cavaliers} (v. 55-57) : il s'est fait voler par Cléon les pains qu'il avait distribués... alors que c'est lui qui avait mis la main à la pâte ! Le verbe \italique{pétrir} et le \italique{pétrin} ont donné nombre d'expressions en français.
\noindent\textbf{76-7} \textbf{τῇδε θἠμέρᾳ κριθήσεται / εἴτ᾽ ἔστ᾽ ἔτι ζῶν εἴτ᾽ ἀπόλωλ᾽ Εὐριπίδης.} Le style est paratragique (avec l'expression poétique `τῇδε θἠμέρᾳ' plutôt que la forme plus courante τήμερον). Apprécier surtout à sa juste valeur l'allitération ! Lyrisme ? emphase ? Les `ttttttttt' saccadent-ils le vers ou lui donnent-ils du rythme ? sans doute un peu des deux. (Euripide est souvent à la frontière du sublime et du ridicule, frontière qu'il passe et repasse allègrement).
\end{multicols}
}%
\end{adjustbox}%
\vfill
\end{adjustwidth}
\newpage
\clearpage\ifodd\value{page}\hbox{}\newpage\fi
\enTeteG{ΠΡΟΛΟΓΟΣ --- Codex Ravennas 429, f. 163v (Classense 165v) --- vv. 80–122}
\begin{adjustwidth}{-1.5cm}{-3.0cm}
\begin{adjustbox}{max totalheight=21.0cm, max width=16.5cm, keepaspectratio}
\begin{minipage}{16.5cm}
{\fontsize{7.2pt}{8.5pt}\selectfont
\noindent\textbf{Folio 03 --- Codex Ravennas 429, \texttt{f. 163v} (\texttt{Classense 165v}) : ΠΡΟΛΟΓΟΣ, vv. 80–122}\hfill \textit{Témoins :} $\text{R}^1, \Rpc, \Rdeux$ (\Mon)\par\vspace{2pt}
\hrule height 0.4pt\vspace{2.5pt}
}
{\fontsize{8.1pt}{9.5pt}\selectfont
\ligneRav{80}{μ\auDessus{ν}{Η}}{\difference{ἐπείπερ ἐστὶ} Θεσμοφορίων ἡ μέση·}{\textbf{80} ἐπεὶ τρίτη 'στὶ R : \lecon{ἐπείπερ ἐστὶ} Nauck (\italique{Mél.} (1866), II pp.468-71)}
\ligneRav{}{\auDessus{εὐ}{ρ}}{τοῦτ᾽ αὐτὸ γὰρ τοι κἀπολεῖν με προσδοκῶ·}{}
\ligneRav{}{}{αἱ γὰρ γυναῖκες ἐπιβεβουλεύκασί μοι}{}
\ligneRav{83}{}{κἀν Θεσμοφόροιν μέλλουσι περί \difference{μου} τήμερον}{\textbf{83} ἐμοῦ R : \lecon{μου} Biset}
\ligneRav{}{—}{ἐκκλησιάζειν ἐπ᾽ ὀλέθρῳ. : τιὴ τί δή;}{}
\ligneRav{85}{\auDessus{εὐ}{ρ}}{ὁτιὴ τραγῳδῶ καὶ κακῶς αὐτὰς λέγω·}{}
\ligneRav{86}{μ\auDessus{ν}{Η}}{νὴ τὸν Ποσειδῶ καὶ \difference{δικαίαν} ἂν πάθοις·}{\textbf{86} δί᾽ και᾽ R : δίκαιά γ᾽ Grynaeus : \lecon{δικαίαν} Wilamowitz}
\ligneRav{}{}{ἀτὰρ τίν᾽ ἐκ ταύτης σὺ μηχανὴν ἔχεις;}{}
\ligneRav{}{\auDessus{εὐ}{ρ}}{Ἀγάθωνα πεῖσαι τὸν τραγῳδοδιδάσκαλον}{}
\ligneRav{}{—}{ἐς Θεσμοφόροιν ἐλθεῖν. : τί δράσοντ᾽; εἰπέ μοι·}{}
\ligneRav{90}{\auDessus{εὐ}{ρ}}{ἐκκλησιάσοντ᾽ ἐν ταῖς \difference{γυναιξὶ} \difference{χἂν} δέῃ}{\textbf{90} γυναιξίν R : \lecon{γυναιξὶ} \Mon \espApp κἂν R : \lecon{χἂν} Markland (1763) ad Eur. \italique{Suppl.} 364}
\ligneRav{91}{—}{λέξονθ᾽ ὑπὲρ ἐμοῦ. : πότερα φανερὸν ἢ λάθρᾳ;}{\textbf{91} \lecon{φανερὸν} R : φανερῶς Meineke}
\ligneRav{}{\auDessus{εὐ}{ρ}}{λάθρᾳ, στολὴν γυναικὸς ἠμφιεσμένον·}{}
\ligneRav{}{μ\auDessus{ν}{Η}}{τὸ πρᾶγμα κομψὸν καὶ σφόδρ᾽ ἐκ τοῦ σοῦ τρόπου·}{}
\ligneRav{}{}{τοῦ γὰρ τεχνάζειν ἡμέτερος ὁ πυραμοῦς·}{}
\ligneRav{95}{\auDessus{εὐ}{ρ}}{σίγα. : τί δ᾽ ἔστιν; : Ἁγάθων ἐξέρχεται·}{\textbf{95} \lecon{τί δ᾽} R : τί Kaibel}
\ligneRav{96}{μ\auDessus{ν}{Η}}{καὶ ποῖός ἐστιν; : οὗτος οὑκκυκλούμενος·}{\textbf{96} ποῖο ἐστὶν οὗτος — οὑκκυκλούμενος R : \lecon{ποῖος ἐστὶν; — οὗτος οὑκκυκλούμενος} Tyrwhitt (p. 65) \espApp ποῖο ἐστὶν οὗτος — οὑκκυκλούμενος R : ποῖος ἐστὶν οὗτος — οὑκκυκλούμενος \Mon : ποῦ ᾽σθ᾽; — ὅπου ᾽στίν; οὗτος οὑκκυκλούμενος Meineke : \lecon{ποῖος ἐστὶν; — οὗτος οὑκκυκλούμενος} Tyrwhitt (p. 65)}
\ligneRav{97}{μ\auDessus{ν}{Η}}{ἀλλ᾽ ἦ τυφλὸς μέν εἰμ᾽· ἐγὼ γὰρ οὐχ ὁρῶ}{\textbf{97} ἢ R : \lecon{ἦ} Biset}
\ligneRav{}{}{ἄνδρ᾽ οὐδέν᾽ ἐνθάδ᾽ ὄντα, Κυρήνην δ᾽ ὁρῶ·}{}
\ligneRav{99}{\auDessus{εὐ}{ρ}}{σίγα· μελῳδεῖν νῦν παρασκευάζεται·}{\textbf{99} ἂν R : \lecon{νῦν} Rogers : alii alia}
\ligneRav{}{}{\hspace{2.2em}\textbf{[παρεπιγραφή R\expo{1} :]} \lecon{μινυρισμός}}{\textbf{post 99} \lecon{μινυρισμός} (παρεπιγραφή) R, $\Sigma^{\text{R}}$}
\ligneRav{100}{μ\auDessus{ν}{Η}}{Μύρμηκος ἀτραπούς, ἢ τί διαμινύρεται;}{\textbf{100} μύρμ- R : \lecon{Μύρμ-} Bibi ap Hesych. μ 1904}
\ligneRav{101}{ἀγάθ/}{\hspace{1.1em}ἱερὰν Χθονίαιν δε-}{\textbf{101} χθονίαις R : \lecon{χθονίαιν} Meineke}
\ligneRav{}{}{\hspace{1.1em}ξάμεναι λαμπάδα κοῦραι}{}
\ligneRav{103}{}{\hspace{1.1em}ξὺν ἐλευθέρᾳ πατρί-}{\textbf{103} \lecon{πατρίδι} R : πραπίδι Wecklein}
\ligneRav{}{}{\hspace{1.1em} δι χορεύσασθε βοάν·}{}
\ligneRav{}{}{\hspace{1.1em}τίνι δαιμόνων ὁ κῶμος;}{}
\ligneRav{105}{}{\hspace{1.1em}λέγε \difference{νιν} εὐπείστως δὲ τοὐμὸν}{\textbf{105} νυν R : \lecon{νιν} Dindorf \espApp εὐπίστως R, Suid. σ 178, 183 : \lecon{εὐπείστως} Reiske}
\ligneRav{106}{}{\hspace{1.1em}δαίμονας \difference{ἔχει} σεβίσαι·}{\textbf{106} ἔχεις R : \lecon{ἔχει} Suid.}
\ligneRav{107}{}{\hspace{1.1em}ἄγε νῦν \difference{ὄλβιζε} Μοῦσα}{\textbf{107} ὅπλιζε R : \lecon{ὄλβιζε} Bentley \espApp \lecon{Μοῦσα} R : μούσᾳ Bergk}
\ligneRav{}{}{\hspace{1.1em}χρυσέων ῥύτορα τόξων}{}
\ligneRav{}{}{\hspace{1.1em}Φοῖβον, ὃς ἱδρύσατο χώρας}{}
\ligneRav{110}{}{\hspace{1.1em}γύαλα Σιμουντίδι γᾷ·}{}
\ligneRav{}{}{\hspace{1.1em}χαῖρε καλλίσταις ἀοιδαῖς}{}
\ligneRav{}{}{\hspace{1.1em}Φοῖβ᾽ ἐν εὐμούσοισι τιμαῖς}{}
\ligneRav{}{}{\hspace{1.1em}γέρας ἱερὸν προφέρων·}{}
\ligneRav{}{}{\hspace{1.1em}τάν τ᾽ ἐν ὄρεσι δρυογόνοισι}{}
\ligneRav{115}{}{\hspace{1.1em}κόραν \difference{ἀείσατ᾽} Ἄρτεμιν ἀγροτέραν·}{\textbf{115} ἀεισαντ᾽ R : \lecon{ἀείσατ᾽} Zanetti (\italique{laudate} Divus)}
\ligneRav{}{}{\hspace{1.1em}ἕπομαι κλῄζουσα σεμνὰν}{}
\ligneRav{117}{}{\hspace{1.1em}γόνον ὀλβίζουσα Λατοῦς}{\textbf{117} σεμνὸν R : \lecon{σεμνὰν} Dindorf}
\ligneRav{}{}{\hspace{1.1em}Ἄρτεμιν ἀπειρολεχῆ·}{}
\ligneRav{}{}{\hspace{1.1em}Λατώ τε κρούματά τ' Ἀσιάδος ποδὶ}{}
\ligneRav{120}{}{\hspace{1.1em}παράρυθμ' εὔρυθμα Φρυγίων}{}
\ligneRav{122}{}{\hspace{1.1em}\difference{διὰ νεύματα} Χαρίτων·}{\textbf{122} διανεύματα R : δινεύματα Bentley : \lecon{διὰ νεύματα} Fritzsche}
}
\vspace{3pt}\noindent\rule{4.5cm}{0.4pt}\par\vspace{1.5pt}
{\fontsize{7.0pt}{8.3pt}\selectfont
\noindent\textbf{Loci paralleli (vv. 80–122) ---} \textbf{80} Hesychius τ 1440 τρίτη Θεσμοφορίων \espTest \textbf{93-4} Suidas τ 439 \espTest \textbf{93} Suidas κ 2025 \espTest \textbf{94} (ἡμέτερος ὁ πυραμοῦς) Suidas π 3192, cf. Tzetzes \italique{Chil.} XIII 281-288 \espTest \textbf{96} (οὑκκυκλούμενος) Zonaras 678 ἐκκυκλεῖ \espTest \textbf{100} Hesychius μ 1904 Μύρμηκος ἀτραπούς ; Photius μ 606 Μύρμηκος ἀτραπός ; Suidas μ 1445 \espTest \textbf{104-6} (εὐπίστως -- σεβίσαι) Suidas σ 183 \espTest \textbf{115-6} (Ἄρτεμιν ἀγροτέραν) \Eustathe 361.37 \espTest \textbf{120} \EtymMagnum 153.21 Ἀσιάδος κρούματα· τῆς κιθάρας. οὔτως Αριστοφάνης εἶπε, παρῳδῶν τὸ ἐξ Ἐρεχθέως Εὐριπίδου ; \SynB α 2226 $=$ Photius α 2957 $=$ Suidas 4147 $=$ \EtymGen (AB) α 1271 ed. Lasserre-Livadaras\par
}
\vspace{2.5pt}\noindent\rule{4.5cm}{0.4pt}\par\vspace{1.5pt}
{\fontsize{6.9pt}{8.2pt}\selectfont
\noindent\textbf{Suite (en commentaire) de l'apparat critique \& notes paléographiques (\texttt{f. 163v}) ---} \textbf{Première \italique{parepigraphè} (entre vv. 99 et 100) :} $\text{R}^1$ copie dans la colonne le mot scénique \lecon{μινυρισμός} (« fredonnement d'Agathon »), accompagné de la précieuse scholie $\Sigma^{\text{R}}$ : μεταξὺ τῶν δυοῖν ἀξιοῦσί τινες γράφειν 'μινυρισμός', ὡς πολλὰ τοιαῦτα παρεπιγράφεται. --- \textbf{Erreurs d'attribution de \Rdeux{} dans la monodie d'Agathon (vv. 101–129) :} Prenant le dialogue lyrique fictif d'Agathon (qui chante successivement le rôle de la coryphée troyenne et celui du chœur) pour un vrai dialogue dramatique, \Rdeux{} inscrit à tort en marge des \italique{notae personarum} faisant alterner αγαθ/ et χορ/. --- \textbf{v. 80 :} ἐπεὶ τρίτη 'στὶ R : \lecon{ἐπείπερ ἐστὶ} Nauck (\italique{Mélanges gréco-romains}, 1866, II, pp. 468–471), adopté et expliqué paléographiquement par L. Hauseux : dans l'archétype en onciale, l'abréviation \textbf{Π' = περ} dans ΕΠΕΙΠ'ΕΣΤΙ a été lue \textbf{Γ' = τρίτη} (ΕΠΕΙΓ'ΕΣΤΙ), sous l'influence d'une glose savante sur le troisième jour des Thesmophories (la \italique{Nesteia}). --- \textbf{post v. 99 :} \italique{parepigraphè} \lecon{μινυρισμός} R ($\Sigma^{\text{R}}$ μεταξὺ τῶν δυοῖν ἀξιοῦσί τινες γράφειν 'μινυρισμός', ὡς πολλὰ τοιαῦτα παρεπιγράφεται). --- \textbf{v. 100 :} μύρμ- R : \lecon{Μύρμ-} (Μύρμηκος ἀτραπούς) \lecon{Bibi} d'après Hésychius $\mu\ 1904$. --- \textbf{v. 107 :} ὅπλιζε R (métathèse de consonnes) : \lecon{ὄλβιζε} Bentley \espApp \lecon{Μοῦσα} R (μούσᾳ Bergk).\par
}
\end{minipage}
\end{adjustbox}
\end{adjustwidth}
\newpage
\enTeteD{Codex Ravennas 429 (Biblioteca Classense) --- Fac-similé du f. 163v (vv. 80–122)}
\begin{adjustwidth}{-1.5cm}{-3.0cm}
\begin{center}
\includegraphics[width=16.5cm,height=20.5cm,keepaspectratio]{/workspaces/Thesmophories/dossier_ravennas_thesmophories/folios_individuels_ff162v-177r/03_f163v_classense165v_vv0080-0122.jpg}\par
\vspace{2pt}
{\scriptsize\textbf{Codex Ravennas 429, f. 163v (Classense 165v)} --- ΠΡΟΛΟΓΟΣ, \textbf{vv. 80–122} ($\text{R}^1, \Rpc, \Rdeux$ (\Mon)).}
\end{center}
\end{adjustwidth}
\newpage
\clearpage\ifodd\value{page}\hbox{}\newpage\fi% --- FOLIO 03 (f. 163v) : SPREAD 2.1 (Texte grec établi + Traduction française + Commentaire 1/8 & 2/8) ---
\clearpage\ifodd\value{page}\hbox{}\newpage\fi
\enTeteG{ΘΕΣΜΟΦΟΡΙΑΖΟΥΣΑΙ --- Texte grec établi (f. 163v, vv. 80–122 [1/4])}
\enTeteD{LES THESMOPHORIES --- Traduction française (f. 163v, vv. 80–122 [1/4])}
\begin{adjustwidth}{-2cm}{-2cm}
\noindent\textbf{\footnotesize Texte grec établi --- Folio 03 (\texttt{f. 163v}, vv. 80–122 [1/4]) \hfill ΠΡΟΛΟΓΟΣ}\\[-2pt]
\noindent\rule{16.5cm}{0.5pt}\\[2pt]
{\fontsize{8.8pt}{10.2pt}\selectfont
\ligneEst{}{ἐπείπερ ἐστὶ Θεσμοφορίων ἡ μέση.}{80}
\ligneEst{\textbf{Εὐριπίδης}}{τοῦτ᾽ αὐτὸ γὰρ τοι κἀπολεῖν με προσδοκῶ.}{}
\ligneEst{}{αἱ γὰρ γυναῖκες ἐπιβεβουλεύκασί μοι}{}
\ligneEst{}{κἀν Θεσμοφόροιν μέλλουσι περί μου τήμερον}{}
\ligneEst{}{ἐκκλησιάζειν ἐπ᾽ ὀλέθρῳ.}{}
\ligneEst{\textbf{Μνησίλοχος}}{\decalage{ἐκκλησιάζειν ἐπ᾽ ὀλέθρῳ.}τιὴ τί δή;}{}
\ligneEst{\textbf{Εὐριπίδης}}{ὁτιὴ τραγῳδῶ καὶ κακῶς αὐτὰς λέγω.}{85}
\ligneEst{\textbf{Μνησίλοχος}}{νὴ τὸν Ποσειδῶ καὶ δικαίαν ἂν πάθοις.}{}
\ligneEst{}{ἀτὰρ τίν᾽ ἐκ ταύτης σὺ μηχανὴν ἔχεις;}{}
\ligneEst{\textbf{Εὐριπίδης}}{Ἀγάθωνα πεῖσαι τὸν τραγῳδοδιδάσκαλον}{}
\ligneEst{}{ἐς Θεσμοφόροιν ἐλθεῖν.}{}
\ligneEst{\textbf{Μνησίλοχος}}{\decalage{ἐς Θεσμοφόροιν ἐλθεῖν.}τί δράσοντ᾽; εἰπέ μοι.}{}
}
\vspace{2pt}
\noindent\rule{16.5cm}{0.6pt}\\[2pt]
\noindent\textbf{\small Commentaire philologique, métrique et scénique (vv. 80–122 --- partie 1/8)}\\[2pt]
\begin{adjustbox}{max totalheight=15.2cm, max width=16.5cm, keepaspectratio}%
\parbox[t]{16.5cm}{
\setlength{\columnsep}{9pt}%
\begin{multicols}{2}
\fontsize{8.4pt}{9.7pt}\selectfont
\setlength{\parindent}{0pt}%
\setlength{\parskip}{1.2pt}%
\noindent\textbf{80} \textbf{ἐπεὶ τρίτη 'στὶ Θεσμοφορίων ἡ μέση.} Le manuscrit de Ravenne donne au vers 80 : « puisque c'est le troisième [jour] des Thesmophories, celui du milieu. » Problème : si les Thesmophories duraient bien trois jours, celui du milieu ne peut en aucune façon être le troisième. Cette contradiction eut la fâcheuse conséquence de déchaîner des torrents de commentaires (à commencer par celui du scholiaste) ; on a inventé, en contradiction avec tous les témoignages\marqueur{1}, des Thesmophories de quatre jours à grand renfort d'érudition et de savantes notes, et ce sous le prétexte qu'il ait pu y avoir avant ces trois jours d'autres festivités ailleurs (surprenant, non ? mais il suffit de prendre un calendrier des fêtes grecques pour constater qu'avec ce genre de raisonnements, on pourrait tout aussi bien \italique{prouver} que les Thesmophories duraient 365 jours par an ; au reste, il faudrait encore montrer que le milieu d'un ensemble de quatre éléments en est le \nieme{3} ; pour finir – \italique{last but not least} –, faire préciser à Aristophane quel jour tombe ἡ μέση serait un détail inepte tant cela était connu de tous les Athéniens). Le texte avait été pourtant brillamment rétabli par August Nauck\marqueur{2} et intégré avec explications dès l'édition aux Belles Lettres de Victor Coulon ; cela n'empêche pas la faute d'apparaître dans nombre d'éditions modernes. Venons-en donc à la bonne correction : l'apparition du τρίτη, \italique{troisième}, est due à la mauvaise lecture d'un Π' (abréviation de περί ou περ) en Γ' (\italique{troisième} lettre de l'alphabet) dans un ancien papyrus\marqueur{3} – cette abréviation est fréquente, et source d'erreurs aussi dans l'autre sens ; ainsi a-t-il fallu en maintes occasions rétablir un `τρίτη' à partir d'un gamma\marqueur{4}. Bref, voici l'histoire de notre début de \nieme{80} vers, très ingénieusement reconstituée : \par\vspace{1pt}\noindent\adjustbox{max width=\columnwidth}{$\displaystyle \overset{\text{{\allemand{Urtext}}}}{\text{ἐπείπερ ἐστὶ}}\xrightarrow{\text{onciales}}\text{ΕΠΕΙΠ'ΕCΤΙ}\xrightarrow{\text{corruption}}\text{ΕΠΕΙΓ'ΕCΤΙ}\xrightarrow{\text{interprétation}}\overset{\text{{\italique{Ravennas}}}}{\text{ἐπεὶ τρίτη ᾽στὶ.}}$}\par\vspace{1pt}
\noindent\textbf{[Notes ad 80]} \textbf{(1)} Outre le scholiaste, qui \italique{sait} que les Thesmophories ne duraient que trois jours – sinon il ne prendrait pas la peine de tant ergoter -- et la glose dubitative d'Hésychius (τ 1440), citons : Alciphron 2,37.2 ; 3,39 ; et la « \anglais{speculative restauration} » du papyrus n° 3540.1-2, \italique{The Oxyrhynchus Papyri}, L (1983), ed. E. W. Handley
\end{multicols}
}%
\end{adjustbox}%
\vfill
\end{adjustwidth}
\newpage
% --- PAGE DE DROITE : Traduction française + Commentaire (2/8) ---
\begin{adjustwidth}{-2cm}{-2cm}
\noindent\textbf{\footnotesize Traduction française --- Folio 03 (\texttt{f. 163v}, vv. 80–122 [1/4]) \hfill ΠΡΟΛΟΓΟΣ}\\[-2pt]
\noindent\rule{16.5cm}{0.5pt}\\[2pt]
{\fontsize{8.8pt}{10.2pt}\selectfont
\ligneEst{}{\francais{puisque nous sommes au jour intermédiaire des Thesmophories.}}{80}
\ligneEst{\textbf{Euripide}}{\francais{C'est précisément ce qui fait que je m'attends à périr.}}{}
\ligneEst{}{\francais{Car les femmes ont monté un coup contre moi.}}{}
\ligneEst{}{\francais{Elles doivent, dans le temple des Thesmophories, à mon sujet aujourd'hui,}}{}
\ligneEst{}{\francais{tenir assemblée afin de me perdre.}}{}
\ligneEst{\textbf{Mnésiloque}}{\francais{\espace{4cm}Pourquoi cela ?}}{}
\ligneEst{\textbf{Euripide}}{\francais{Parce que je les théâtralise et parle mal d'elles.}}{85}
\ligneEst{\textbf{Mnésiloque}}{\francais{Par Poseidôn, tu serais châtié en justice.}}{}
\ligneEst{}{\francais{Mais, pour y échapper, de quelle machine disposes-tu ?}}{}
\ligneEst{\textbf{Euripide}}{\francais{De persuader à Agathon, le poète tragique,}}{}
\ligneEst{}{\francais{de se rendre aux Thesmophories.}}{}
\ligneEst{\textbf{Mnésiloque}}{\francais{\espace{3.5cm}Pour quoi faire, dis-moi ?}}{}
}
\vspace{2pt}
\noindent\rule{16.5cm}{0.6pt}\\[2pt]
\noindent\textbf{\small Commentaire philologique, métrique et scénique (vv. 80–122 --- partie 2/8)}\\[2pt]
\begin{adjustbox}{max totalheight=15.2cm, max width=16.5cm, keepaspectratio}%
\parbox[t]{16.5cm}{
\setlength{\columnsep}{9pt}%
\begin{multicols}{2}
\fontsize{8.4pt}{9.7pt}\selectfont
\setlength{\parindent}{0pt}%
\setlength{\parskip}{1.2pt}%
\noindent \textbf{(2)} in \italique{Mélanges gréco-romains}, St.-Pétersburg (1866), t. II, pp. 468-71 ; voir aussi pp. 344-5 pour des témoignages sur les \italique{Thesmophories}.
\noindent \textbf{(3)} La corruption doit être ancienne puisqu'elle – conjointe au bavardage du scholiaste – rendait déjà perplexe Hésychius (τ 1440): τρίτη Θεσμοφορίων· ζητεῖται πῶς ἅμα μὲν λέγει τρίτην Θεσμοφορίων εἶναι, ἅμα δὲ μέσην, τεττάρων οὐσῶν ἡμερῶν.
\noindent \textbf{(4)} Pour rester avec notre cher Aristophane, un `τρεῖς' a dû être restitué au vers 1507 des \italique{Guêpes} là où tous les manuscrits comportent un `γε'.
\noindent\textbf{83} \textbf{κἀν Θεσμοφόροιν.} Le scholiaste précise que c'est une Ἀττικὴ σύνταξις καὶ δυϊκὴ, ἀντὶ τοῦ `ἐν τῷ τῶν θεσμοφόρων ἱερῷ' ; les \italique{deux} déesses thesmophores étant Déméter et sa fille Perséphone.
\noindent\textbf{86-7} \textbf{νὴ τὸν Ποσειδῶ καὶ δί᾽ και᾽ ἂν πάθοις / ἀτὰρ τίν᾽ ἐκ ταύτης σὺ μηχανὴν ἔχεις;} Ne voyant pas à quel nom féminin renvoyait `ἐκ ταύτης', on a retouché le texte du vers 87 à l'envi (ἐκ τούτων Bergler, εἰς ταύτας Sandbach, \&c.).
\noindent C'est à mon avis aller un peu vite en besogne alors qu'il y a justement un problème textuel au vers précédent : il manque une syllabe, probablement au niveau du fautif `δί᾽ και᾽'. Je fais l'hypothèse que c'est justement à cet endroit que se trouvait ce nom féminin, repris par le `ἐκ ταύτης' du vers suivant. De quoi Euripide peut-il se mettre \italique{hors} en déployant quelque machine ? Réponse : tel le Strepsiade des \italique{Nuées}, c'est à la justice qu'il veut échapper. Au vu du texte conservé, la correction est évidente : il faut substituer au `δί᾽ και᾽' un `δικαίαν' ; la chute de l'`αν' s'expliquant par la courante erreur de «~saut du même au même~» (`δικαίαν' étant directement suivi de la particule `ἂν').
\noindent Lisons donc :
\citComm{νὴ τὸν Ποσειδῶ καὶ δικαίαν ἂν πάθοις.\\
ἀτὰρ τίν᾽ ἐκ ταύτης σὺ μηχανὴν ἔχεις;\\
\francais{Sainte Mer ! ce serait en justice que tu pâtirais.\\
Mais pour y échapper, de quelle machine disposes-tu ?}}
\noindent Cette correction est d'autant plus intéressante que ἡ δικαία, mot rare, est un synonyme poétique de δίκη\marqueur{1} ; Mnésiloque retourne-t-il contre Euripide une de ses propres sentences\marqueur{2} ?\\ Wilamowitz\marqueur{3} avait déjà proposé cette émendation sans qu'elle soit jamais reprise dans les éditions publiées depuis -- à une notable exception : celle de Prato.
\end{multicols}
}%
\end{adjustbox}%
\vfill
\end{adjustwidth}
\newpage

% --- FOLIO 03 (f. 163v) : SPREAD 2.2 (Texte grec établi + Traduction française + Commentaire 3/8 & 4/8) ---
\clearpage\ifodd\value{page}\hbox{}\newpage\fi
\enTeteG{ΘΕΣΜΟΦΟΡΙΑΖΟΥΣΑΙ --- Texte grec établi (f. 163v, vv. 80–122 [2/4])}
\enTeteD{LES THESMOPHORIES --- Traduction française (f. 163v, vv. 80–122 [2/4])}
\begin{adjustwidth}{-2cm}{-2cm}
\noindent\textbf{\footnotesize Texte grec établi --- Folio 03 (\texttt{f. 163v}, vv. 80–122 [2/4]) \hfill ΠΡΟΛΟΓΟΣ}\\[-2pt]
\noindent\rule{16.5cm}{0.5pt}\\[2pt]
{\fontsize{8.8pt}{10.2pt}\selectfont
\ligneEst{\textbf{Εὐριπίδης}}{ἐκκλησιάσοντ᾽ ἐν ταῖς γυναιξὶ χἂν δέῃ}{90}
\ligneEst{}{λέξονθ᾽ ὑπὲρ ἐμοῦ.}{}
\ligneEst{\textbf{Μνησίλοχος}}{\decalage{λέξονθ᾽ ὑπὲρ ἐμοῦ.}πότερα φανερὸν ἢ λάθρᾳ;}{}
\ligneEst{\textbf{Εὐριπίδης}}{λάθρᾳ, στολὴν γυναικὸς ἠμφιεσμένον.}{}
\ligneEst{\textbf{Μνησίλοχος}}{τὸ πρᾶγμα κομψὸν καὶ σφόδρ᾽ ἐκ τοῦ σοῦ τρόπου·}{}
\ligneEst{}{τοῦ γὰρ τεχνάζειν ἡμέτερος ὁ πυραμοῦς.}{}
\ligneEst{\textbf{Εὐριπίδης}}{σίγα.}{95}
\ligneEst{\textbf{Μνησίλοχος}}{\decalage{σίγα.}τί δ᾽ ἔστιν;}{}
\ligneEst{\textbf{Εὐριπίδης}}{\decalage{σίγα.τί δ᾽ ἔστιν;}Ἁγάθων ἐξέρχεται.}{}
\ligneEst{\textbf{Μνησίλοχος}}{καὶ ποῖός ἐστιν;}{}
\ligneEst{\textbf{Εὐριπίδης}}{\decalage{καὶ ποῖός ἐστιν;}οὗτος οὑκκυκλούμενος.}{}
\ligneEst{\textbf{Μνησίλοχος}}{ἀλλ᾽ ἦ τυφλὸς μέν εἰμ᾽· ἐγὼ γὰρ οὐχ ὁρῶ}{}
}
\vspace{2pt}
\noindent\rule{16.5cm}{0.6pt}\\[2pt]
\noindent\textbf{\small Commentaire philologique, métrique et scénique (vv. 80–122 --- partie 3/8)}\\[2pt]
\begin{adjustbox}{max totalheight=15.2cm, max width=16.5cm, keepaspectratio}%
\parbox[t]{16.5cm}{
\setlength{\columnsep}{9pt}%
\begin{multicols}{2}
\fontsize{7.6pt}{8.8pt}\selectfont
\setlength{\parindent}{0pt}%
\setlength{\parskip}{1.2pt}%
\noindent\textbf{87} \textbf{μηχανὴν.} Au sens imagé, une μηχανή est un \italique{expédient}, un \italique{tour}, un \italique{stratagème}. Certes, Euripide est un roué, un filou, un retors... mais Euripide est avant tout un homme de théâtre ! Il dispose de \italique{machineries} dont il use et abuse\marqueur{1}. Ce sont ces deux sens de μηχανή, le propre et le figuré, qui étaient instantanément saisis du spectateur du \nieme{V} siècle av. J.-C. Conformité du langage français avec le grec : le mot « truc », qui désignait au \nieme{XVII} siècle une machine de théâtre, s'est approprié le même sens imagé.
\noindent\textbf{[Notes ad 86-7]} \textbf{(1)} Le mot est attesté dans l'\EtymMagnum à l'entrée Ἀθηνᾶ : le glosateur y recense Ἀθηναία, qui est à Ἀθήνη ce que δικαία est à δίκη (Ἀθήνη Ἀθηναία, ὡς δίκη δικαία).
\noindent \textbf{(2)} Tout au long de la pièce, Euripide est pris au piège de ses propres γνῶμαι (aux vers 194, 275-276, etc). Est-ce aussi le cas ici ? Quoi qu'il en soit, on trouve des expressions très proches aux vers 1379 des \italique{Nuées} :
\citComm{Φειδιππίδης -- νὴ τὸν Δί᾽ ἐν δίκῃ γ᾽ ἄν,}
\noindent ou 403 de \italique{Lysistrata} :
\citComm{Πρόβουλος -- νὴ τὸν Ποσειδῶ τὸν ἁλυκὸν δίκαιά [δικαία ?] γε.}
\noindent \textbf{(3)} Au commentaire du vers 681 de \italique{La Folie d'Héraclès} : second volume de son \italique{Euripides Herakles erklärt von Ulrich von Wilamowitz-Moellendorff}. Wilamowitz cite \italique{en passant} le vers 86 des \italique{Thesmophories} ainsi corrigé sans même préciser la correction -- dont il semble bien être l'auteur -- tant elle devait lui sembler aller de soi.
\noindent \textbf{(1)} Euripide recourut à la μηχανή massivement dans ses œuvres tardives. Aristote, dans sa \italique{Poétique} (XV, 1454a37), le lui reproche pour sa \italique{Médée} : Euripide privilégia un dénouement ἀπὸ μηχανῆς plutôt que par l'histoire elle-même : φανερὸν οὖν ὅτι καὶ τὰς λύσεις τῶν μύθων ἐξ αὐτοῦ δεῖ τοῦ μύθου συμβαίνειν, καὶ μὴ ὥσπερ ἐν τῇ Μηδείᾳ ἀπὸ μηχανῆς. Aristophane parodie en maints endroits cet artifice de théâtre. Pour un dossier plus complet, aller voir l'article «~Machines et \italique{deus ex machina} dans les spectacles comiques~» de Leonardo Fiorentini, in \italique{L'appareil scénique dans les spectacles de l'Antiquité} (2013).
\noindent\textbf{90-1} \textbf{ἐκκλησιάσοντ᾽ ἐν ταῖς γυναιξὶν κἂν δέῃ / λέξονθ᾽ ὑπὲρ ἐμοῦ.} Je reprends la correction de Markland : substituer χἂν ($=$ καὶ ἃ ἂν) à κἂν ($=$ καὶ ἐὰν) qui présente un sens bien meilleur : Euripide ne dit certes pas « qu'au besoin » ou « s'il le faut » (κἂν) Agathon interviendra mais bien plutôt qu'il lui faut intervenir autant qu'il le pourra (χἂν). Euripide a sa vie en jeu – si l'on peut s'exprimer ainsi – et sait bien le courroux de la gent féminine ; l'inaction aurait forcément une issue fatale... heureuse pour Athènes et les Athéniennes mais un peu plus fâcheuse pour sa personne.
\end{multicols}
}%
\end{adjustbox}%
\vfill
\end{adjustwidth}
\newpage
% --- PAGE DE DROITE : Traduction française + Commentaire (4/8) ---
\begin{adjustwidth}{-2cm}{-2cm}
\noindent\textbf{\footnotesize Traduction française --- Folio 03 (\texttt{f. 163v}, vv. 80–122 [2/4]) \hfill ΠΡΟΛΟΓΟΣ}\\[-2pt]
\noindent\rule{16.5cm}{0.5pt}\\[2pt]
{\fontsize{8.8pt}{10.2pt}\selectfont
\ligneEst{\textbf{Euripide}}{\francais{Pour intervenir à l'assemblée devant les femmes et, autant qu'il est nécessaire,}}{90}
\ligneEst{}{\francais{me défendre.}}{}
\ligneEst{\textbf{Mnésiloque}}{\francais{\espace{2.8cm}Ouvertement ou en tapinois ?}}{}
\ligneEst{\textbf{Euripide}}{\francais{En tapinois, revêtu d'un costume féminin.}}{}
\ligneEst{\textbf{Mnésiloque}}{\francais{C'est ingénieux et bien dans tes façons.}}{}
\ligneEst{}{\francais{En fait d'artifices, à nous le gâteau.}}{}
\ligneEst{\textbf{Euripide}}{\francais{Silence.}}{95}
\ligneEst{\textbf{Mnésiloque}}{\francais{\espace{1cm}Qu'y a-t-il ?}}{}
\ligneEst{\textbf{Euripide}}{\francais{\espace{2.7cm}Agathon sort.}}{}
\ligneEst{\textbf{Mnésiloque}}{\francais{Quel homme est-ce ?}}{}
\ligneEst{\textbf{Euripide}}{\francais{\espace{2cm}Celui-là qu'on roule dehors.}}{}
\ligneEst{\textbf{Mnésiloque}}{\francais{Est-ce donc que je suis aveugle ? Je ne vois pas qu'il y ait}}{}
}
\vspace{2pt}
\noindent\rule{16.5cm}{0.6pt}\\[2pt]
\noindent\textbf{\small Commentaire philologique, métrique et scénique (vv. 80–122 --- partie 4/8)}\\[2pt]
\begin{adjustbox}{max totalheight=15.2cm, max width=16.5cm, keepaspectratio}%
\parbox[t]{16.5cm}{
\setlength{\columnsep}{9pt}%
\begin{multicols}{2}
\fontsize{8.4pt}{9.7pt}\selectfont
\setlength{\parindent}{0pt}%
\setlength{\parskip}{1.2pt}%
\noindent Précisons que ἐν a ici le sens technique que donne le Bailly (A, III.10) : « lorsqu'il s'agit d'un orateur, d'un accusé, ἐν correspond au français \italique{devant} » (cf. les \italique{Acharniens} v. 630 ; les \italique{Guêpes} v. 1185 ; le fragment 161 K.-A.) ; il faut donner un sens actif au `ἐκκλησιάσοντ᾽ ἐν ταῖς γυναιξὶν' : « pour intervenir à l'assemblée devant les femmes ».
\noindent\textbf{94} \textbf{ἡμέτερος ὁ πυραμοῦς.} L'expression proverbiale signifie : « à nous la palme de la victoire ». En effet, on offrait un gâteau de miel et de froment (πυραμοῦς), à titre de récompense, à celui qui avait le mieux résisté au sommeil dans une παννυχίς (scholiaste des \italique{Cavaliers}, au v. 277 où l'expression apparaît à l'identique ; Athénée 647b-c, 668c). (Note de Taillardat, §847.)
\noindent\textbf{96} \textbf{οὑκκυκλούμενος} ($=$ ὁ ἐκκυκλούμενος). Agathon apparaît sur l'eccyclème, « machine de théâtre munie de roues qui, en tournant, faisait voir aux spectateurs ce qui était censé se passer dans une maison. » (Chantraine.)  Le scholiaste l'avait noté :
\citComm{ἐπὶ ἐκκυκλήματος γὰρ φαίνεται.}
\noindent\textbf{98} \textbf{Κυρήνην δ᾽ ὁρῶ.} Cyrène était une célèbre courtisane évoquée aussi aux vers 1327-1328 des \italique{Grenouilles} :
\citComm{ἀνὰ τὸ δωδεκαμήχανον / Κυρήνης μελοποιῶν\\
\francais{$[$Euripide$]$ compose ses chants à la mode de Cyrène aux douze postures.}}
\noindent Le scholiaste y est plus loquace que le nôtre et nous apprend le trait distinctif de cette Cyrène :
\citComm{Κυρήνη τις ἑταίρα ἐπίσημος, δωδεκαμήχανος ἐπικαλουμένη διὰ τὸ τοσαῦτα σχήματα ἐν τῇ συνουσίᾳ ποιεῖν.}
\noindent Taillardat (§734) rapproche pertinemment ce sobriquet du titre d'un ``traité'' de Paxamos : Δωδεκάτεχνον, à propos duquel la \italique{Souda} (à l'entrée Πάξαμος) dit : ἔστι δὲ περὶ αἰσχρῶν σχημάτων. Que moque donc Aristophane en rapprochant, non cette fois la muse du poète comme dans les \italique{Grenouilles}, mais Agathon lui-même de Cyrène ? Bien-sûr son air efféminé ! bien-sûr son travestissement !.. mais pourquoi Cyrène en particulier ? Réponse : pour ce qu'elle a de remarquable (ἐπίσημος) qui lui valut son surnom ; il ne me paraît pas douteux qu'Agathon, qui vient de se faire rouler par l'eccyclème, apparaît affalé sur son lit dans une position suggestive, ou du moins qui peut prêter à la confusion de Mnésiloque. (Cf. la raillerie -- μῶν βινεῖσθαι; -- au vers 50.)
\end{multicols}
}%
\end{adjustbox}%
\vfill
\end{adjustwidth}
\newpage

% --- FOLIO 03 (f. 163v) : SPREAD 2.3 (Texte grec établi + Traduction française + Commentaire 5/8 & 6/8) ---
\clearpage\ifodd\value{page}\hbox{}\newpage\fi
\enTeteG{ΘΕΣΜΟΦΟΡΙΑΖΟΥΣΑΙ --- Texte grec établi (f. 163v, vv. 80–122 [3/4])}
\enTeteD{LES THESMOPHORIES --- Traduction française (f. 163v, vv. 80–122 [3/4])}
\begin{adjustwidth}{-2cm}{-2cm}
\noindent\textbf{\footnotesize Texte grec établi --- Folio 03 (\texttt{f. 163v}, vv. 80–122 [3/4]) \hfill ΠΡΟΛΟΓΟΣ}\\[-2pt]
\noindent\rule{16.5cm}{0.5pt}\\[2pt]
{\fontsize{8.8pt}{10.2pt}\selectfont
\ligneEst{}{ἄνδρ᾽ οὐδέν᾽ ἐνθάδ᾽ ὄντα, Κυρήνην δ᾽ ὁρῶ.}{}
\ligneEst{\textbf{Εὐριπίδης}}{σίγα· μελῳδεῖν νῦν παρασκευάζεται.}{}
\ligneEst{\textbf{Μνησίλοχος}}{Μύρμηκος ἀτραπούς, ἢ τί διαμινύρεται;}{100}
\ligneEst{\textbf{Ἀγάθων} {\tiny\textit{io. min\cat + io. min.}}}{ἱερὰν Χθονίαιν δε-}{}
\ligneEst{{\tiny\textit{2 ion. min.}}}{ξάμεναι λαμπάδα κοῦραι}{}
\ligneEst{{\tiny\textit{io. min\cat + bacch.}}}{ξὺν ἐλευθέρᾳ πατρί-}{}
\ligneEst{{\tiny\textit{io. min. + io. min\cat}}}{\alinea δι χορεύσασθε βοάν.}{}
\ligneEst{{\tiny\textit{2 io. min. anacl.}}}{τίνι δαιμόνων ὁ κῶμος;}{}
\ligneEst{{\tiny\textit{2 troch.}}}{λέγε νιν εὐπείστως δὲ τοὐμὸν}{105}
\ligneEst{{\tiny\textit{troch. + io. min\cat}}}{δαίμονας ἔχει σεβίσαι.}{}
\ligneEst{{\tiny\textit{ion. min. + troch.}}}{ἄγε νῦν ὄλβιζε Μοῦσα}{}
\ligneEst{{\tiny\textit{troch. + io. min.}}}{χρυσέων ῥύτορα τόξων}{}
}
\vspace{2pt}
\noindent\rule{16.5cm}{0.6pt}\\[2pt]
\noindent\textbf{\small Commentaire philologique, métrique et scénique (vv. 80–122 --- partie 5/8)}\\[2pt]
\begin{adjustbox}{max totalheight=15.2cm, max width=16.5cm, keepaspectratio}%
\parbox[t]{16.5cm}{
\setlength{\columnsep}{9pt}%
\begin{multicols}{2}
\fontsize{8.4pt}{9.7pt}\selectfont
\setlength{\parindent}{0pt}%
\setlength{\parskip}{1.2pt}%
\noindent\textbf{99} \textbf{μινυρισμός.} Une scholie du \italique{Ravennas} semble avoir gardé la trace d'une ancienne didascalie : μινυρισμός était indiqué marginalement :
\citComm{μεταξὺ τῶν δυοῖν ἀξιοῦσί τινες γράφειν `μινυρισμός', ὡς πολλὰ τοιαῦτα παρεπιγράφεται.}
\noindent\textbf{100} \textbf{μύρμηκος ἀτραπούς.} Willems, en accord avec le scholiaste, se propose d'expliquer ainsi l'expression μύρμηκος ἀτραπούς : « par allusion à la marche saccadée et tortueuse des fourmis. Image familière aux Anciens. Dans l'\italique{Anthologie}, XI, 78, il est question des galeries obliques et droites, λόξα καὶ ὀρθά, d'une fourmilière. Voir aussi Plutarque, \italique{de Solertia animal.}, II, p. 968a. Plaute a dit : \latin{formicinum gradum movere}. » Taillardat (§784) abonde en son sens avec d'autres exemples à l'appui.
\noindent Cependant, pour suggérer l'image de \italique{sentiers de fourmis}, n'aurait-il pas fallu, en français comme en grec, un pluriel (μυρμήκ\textbf{ων} ἀτραπούς) ?\\ Hésychius ne le comprenait pas autrement et voyait dans `μύρμηκος' un nom propre (μ 1904) :
\citComm{Μύρμηκος ἀτραπούς· Ἀθήνησιν ἐν Σκαμβωνιδῶν ἔστι Μύρμηκος ἀτραπός, ἀπὸ ἥρωος Μύρμηκος ὀνομαζομένη.\\
\francais{Dans [le dème] athénien des Scambônides [σκαμβός signifie \italique{courbe}, \italique{tortu}] se trouve un sentier de Myrmex, qui tire son nom du héros Myrmex.}}
\noindent Que l'image soit tirée de la fourmilière ou d'un nom propre, la critique n'en reste pas moins d'ordre musical\marqueur{1}.
\noindent Philoxène de Cythère, un représentant de cette «~nouvelle musique~» dénoncée par Aristophane, reçut d'ailleurs le surnom de Μύρμηξ (voir \italique{Souda} φ 393, à l'entrée Φιλόξενος). Le persiflage de Mnésiloque ne serait-il pas une allusion audit Philoxène, alias Myrmex ?
\noindent\textbf{101} \textbf{Χθονίαιν} est un génitif se rapportant à ἱεράν. Cf. \italique{Ploutos}, 937, et \italique{Caval.}, 301. (Note de Willems.)
\noindent\textbf{[Notes ad 100]} \textbf{(1)} L'image des lignes brisées de la fourmilière, μυρμηκιά, servit en musique pour décrire des trilles ou des arpèges. On trouve l'expression `ἐκτραπέλους μυρμηκιάς' chez le Pseudo-Plutarque, \italique{De Musica}, 30 ; il s'agit du fragment 155(v.23) de Phérécrate dans l'édition Kassel-Austin. À noter que le même Phérécrate a composé une comédie intitulée les μυρμηκάνθρωποι.
\end{multicols}
}%
\end{adjustbox}%
\vfill
\end{adjustwidth}
\newpage
% --- PAGE DE DROITE : Traduction française + Commentaire (6/8) ---
\begin{adjustwidth}{-2cm}{-2cm}
\noindent\textbf{\footnotesize Traduction française --- Folio 03 (\texttt{f. 163v}, vv. 80–122 [3/4]) \hfill ΠΡΟΛΟΓΟΣ}\\[-2pt]
\noindent\rule{16.5cm}{0.5pt}\\[2pt]
{\fontsize{8.8pt}{10.2pt}\selectfont
\ligneEst{}{\francais{un homme ici. C'est Cyrène que je vois.}}{}
\ligneEst{\textbf{Euripide}}{\francais{Chut ! Il se dispose à chanter.}}{}
\ligneEst{\textbf{Mnésiloque}}{\francais{Des sentes de Fourmi, ou que sont ces fredons ?}}{100}
\ligneEst{\textbf{Agathon}}{\francais{1° — Des deux Chtoniennes, recevez}}{}
\ligneEst{}{\francais{la torche consacrée, vierges,}}{}
\ligneEst{}{\francais{et avec votre libre patrie}}{}
\ligneEst{}{\francais{entonnez un chœur de danse.}}{}
\ligneEst{}{\francais{2° — À quelle divinité ce cômos ?}}{}
\ligneEst{}{\francais{Parle. Je me laisse aisément induire}}{105}
\ligneEst{}{\francais{à honorer les dieux.}}{}
\ligneEst{}{\francais{3° — Allons, Muse, glorifie}}{}
\ligneEst{}{\francais{l'archer à l'arc d'or,}}{}
}
\vspace{2pt}
\noindent\rule{16.5cm}{0.6pt}\\[2pt]
\noindent\textbf{\small Commentaire philologique, métrique et scénique (vv. 80–122 --- partie 6/8)}\\[2pt]
\begin{adjustbox}{max totalheight=15.2cm, max width=16.5cm, keepaspectratio}%
\parbox[t]{16.5cm}{
\setlength{\columnsep}{9pt}%
\begin{multicols}{2}
\fontsize{8.4pt}{9.7pt}\selectfont
\setlength{\parindent}{0pt}%
\setlength{\parskip}{1.2pt}%
\noindent\textbf{101-29} \textbf{Agathon chante un fragment de tragédie.} Sujet : un coryphée et un chœur de jeunes filles troyennes célèbrent en strophes alternées Léto et ses deux rejetons. Il est très évident qu'Agathon chante seul. Les arguments fournis par M. P. Mazon\marqueur{1} sont à cet égard sans réplique. (Note de Willems.)\\ Suidas (χ 403) nous confirme qu'Agathon avait l'habitude de chanter de la sorte :
\citComm{Χορικός· ὁ τοῦ χοροῦ ἐξάρχων. τοιοῦτος ἦν Ἀγάθων, ὃς ὑποκριτικὰ μέλη ἐποίει καὶ αὐτὸς ὑποκρίνεται\\
\francais{Chorique : celui qui mène le chœur. Tel Agathon, qui avait l'habitude de créer des chants dramatiques et de les interpréter lui-même.}}
\noindent L'on pas manqué de remarquer en quoi les ruptures de rythmes (παράρυθμ' εὔρυθμα) ou le contexte ionique de cette parodie d'Agathon en font un morceau efféminé et suave, ce qui est tout-à-fait juste quoiqu'il nous manque la musique pour juger sur pièces\marqueur{2} ; mais beaucoup trop de commentateurs se sont arrêté à ces considérations\marqueur{3}.
\noindent (Opinion au demeurant est bien pratique puisqu'elle dispense d'avoir à porter trop d'attention sur le texte.) Ce fut pas le cas du regretté Alain Blanchard avait excellemment analysé la parodie d'Agathon dans son article : \italique{La parodie d'Agathon et la structure des }``Bucoliques''\italique{ de Virgile}\marqueur{4}. Blanchard décela dans ce chant une structure concentrique, en cela analogue aux \italique{Bucoliques} de Virgile ou encore aux cantates de Bach ; le point important étant que « la reconnaissance de ce schéma concentrique » est d'une grande « valeur exégétique ». Laissons-lui la parole : « la coryphée invite le chœur à célébrer Apollon, Artémis et leur mère Léto, et le chœur, docilement, obéit à cette injonction. En réalité, l'étude de composition, une fois admise l'existence d'une structure concentrique conduit à plus de subtilité.
\noindent Le parallélisme entre le premier élément (vv. 101-103) et le neuvième (vv. 126-129) réside dans le thème de la lumière (lumière des torches portées par le chœur, lumière qui brille en réponse dans les yeux divins), thème qui précède un appel à la célébration d'un dieu – l'appel étant plus solennel, plus brusque et plutôt à titre de rappel dans le neuvième. Le parallélisme entre le deuxième élément (104-106) et le huitième (123-125) réside dans le thème du σέβεσθαι, introduit en mineur à la fin du deuxième élément et qui éclate sur le mode majeur au début du huitième.
\end{multicols}
}%
\end{adjustbox}%
\vfill
\end{adjustwidth}
\newpage

% --- FOLIO 03 (f. 163v) : SPREAD 2.4 (Texte grec établi + Traduction française + Commentaire 7/8 & 8/8) ---
\clearpage\ifodd\value{page}\hbox{}\newpage\fi
\enTeteG{ΘΕΣΜΟΦΟΡΙΑΖΟΥΣΑΙ --- Texte grec établi (f. 163v, vv. 80–122 [4/4])}
\enTeteD{LES THESMOPHORIES --- Traduction française (f. 163v, vv. 80–122 [4/4])}
\begin{adjustwidth}{-2cm}{-2cm}
\noindent\textbf{\footnotesize Texte grec établi --- Folio 03 (\texttt{f. 163v}, vv. 80–122 [4/4]) \hfill ΠΡΟΛΟΓΟΣ}\\[-2pt]
\noindent\rule{16.5cm}{0.5pt}\\[2pt]
{\fontsize{8.8pt}{10.2pt}\selectfont
\ligneEst{{\tiny\textit{2 io. min.}}}{Φοῖβον, ὃς ἱδρύσατο χώρας}{}
\ligneEst{{\tiny\textit{io. min. + io. min\cat}}}{γύαλα Σιμουντίδι γᾷ.}{110}
\ligneEst{{\tiny\textit{2 troch.}}}{χαῖρε καλλίσταις ἀοιδαῖς}{}
\ligneEst{}{Φοῖβ᾽ ἐν εὐμούσοισι τιμαῖς}{}
\ligneEst{{\tiny\textit{io. min. + io. min\cat}}}{γέρας ἱερὸν προφέρων.}{}
\ligneEst{{\tiny\textit{2 troch.}}}{τάν τ᾽ ἐν ὄρεσι δρυογόνοισι}{}
\ligneEst{{\tiny\textit{2 io. m. + io. m\cat}}}{κόραν ἀείσατ᾽ Ἄρτεμιν ἀγροτέραν.}{}
\ligneEst{{\tiny\textit{io. min. + troch.}}}{ἕπομαι κλῄζουσα σεμνὰν}{}
\ligneEst{{\tiny\textit{io. min. + troch.}}}{γόνον ὀλβίζουσα Λατοῦς}{}
\ligneEst{{\tiny\textit{io. min. + io. min\cat}}}{Ἄρτεμιν ἀπειρολεχῆ.}{}
\ligneEst{{\tiny\textit{io. m+io. m\cat+io. m}}}{Λατώ τε κρούματά τ' Ἀσιάδος ποδὶ}{}
\ligneEst{{\tiny\textit{2 io. min.}}}{παράρυθμ' εὔρυθμα Φρυγίων}{120}
\ligneEst{{\tiny\textit{io. min. + io. min\cat}}}{διὰ νεύματα Χαρίτων.}{}
}
\vspace{2pt}
\noindent\rule{16.5cm}{0.6pt}\\[2pt]
\noindent\textbf{\small Commentaire philologique, métrique et scénique (vv. 80–122 --- partie 7/8)}\\[2pt]
\begin{adjustbox}{max totalheight=15.2cm, max width=16.5cm, keepaspectratio}%
\parbox[t]{16.5cm}{
\setlength{\columnsep}{9pt}%
\begin{multicols}{2}
\fontsize{7.6pt}{8.8pt}\selectfont
\setlength{\parindent}{0pt}%
\setlength{\parskip}{1.2pt}%
\noindent Le parallélisme entre le troisième élément (107-110) et le septième (120-122) réside dans le caractère national des divinités honorées : caractère national sensible surtout dans le troisième élément (Apollon constructeur des remparts de Troie), mais présent, plus discrètement, dans le septième avec les Charites phrygiennes.
\noindent Enfin, le thème de la majesté divine (majesté d'Apollon comme dispensateur de la gloire) entre le quatrième élément (111-113) et le sixième (117-119). Le cinquième (114-116), central, est un appel à honorer Artémis, divinité féminine. Dans la première partie de la parodie, les trois interventions du coryphée concernent des divinités armées de l'arc qu'il faut honorer avec un esprit libre et d'une voix forte (103), mais le coryphée s'adresse à un chœur de jeunes filles qui, pour ainsi dire, retourne complètement la nature de ces divinités dont on aperçoit alors le caractère ambigu. Déjà, aux appels à un culte viril qu'on leur faisait, elles avaient répondu, avec une soumission apparente bien féminine, par une profession de piété plus neutre (σεβίσαι, 106). Puis, avec elles, Apollon cesse d'être le dieu archer pour être le coryphée des Muses (112). Qu'on leur propose d'honorer Artémis, il n'est plus question avec elle de chasse – occupation virile – mais de maternité : maternité glorieuse de Léto (118), maternité refusée par elle-même par la déesse Vierge (119). Le ton est donné pour la seconde partie du poème où le chœur, à l'appel du coryphée, honorera, après Artémis, sa mère Léto et la cithare. Dans un cadre aussi favorable, le chœur peut même se permettre une concession – nous parlerions plutôt d'un appel spontané et aventureux du cœur – avec l'expression ἄρσενι βοᾷ, \italique{mâle accent}, en 125 qui rappelle en la précisant l'invitation faite en 103b par le coryphée. Mais qu'on invite de nouveau à honorer une divinité masculine comme Apollon : c'est la maternité de Léto qui est chantée (130). »\\ Fera-t-on la même erreur à propos de ce chant que les premiers auditeurs du sonnet d'Oronte en prétendant : « Il est bon » ? C'est en tout cas ce qu'il nous semble à l'analyse de ce petit bijou de malice au subtil renversement : d'un banal hommage au fondateur d'une cité, l'on passe doucement et insidieusement à un morceau d'érotisme sur fond de maternité frustrée, le tout échafaudé sur une structure solidement ancrée\marqueur{5}. Rien d'étonnant donc à ce que Mnésiloque en soit tout émoustillé ; rien d'étonnant non plus à ce que, par la bouche de Dionysos, Aristophane rende hommage à Agathon après sa mort : ἀγαθὸς ποιητὴς καὶ ποθεινὸς τοῖς φίλοις, « un bon poète apprécié de ses amis » (\italique{Grenouilles}, v. 84). Pour finir, l'élaboration complexe du morceau ainsi mise au jour par Blanchard fait voler en éclat la thèse fort répandue d'une improvisation de la part d'Agathon ; d'ailleurs l'on sait par Aristote qu'Agathon avait eu l'ambition dans une de ses tétralogies de traiter intégralement du sac de Troie, ambition qui l'avait mené à l'échec : (\italique{Poétique}, XVIII, 56a16-19) :
\citComm{ὅσοι πέρσιν Ἰλίου ὅλην ἐποίησαν καὶ μὴ κατὰ μέρος ὥσπερ Εὐριπίδης, <ἢ> Νιόβην καὶ μὴ ὥσπερ Αἰσχύλος, ἢ ἐκπίπτουσιν ἢ κακῶς ἀγωνίζονται, ἐπεὶ καὶ Ἀγάθων ἐξέπεσεν ἐν τούτῳ μόνῳ.\\
\francais{Tous les poètes qui veulent traiter la ruine de Troie entièrement et non pas par parties comme Euripide, ou la légende de Niobé tout entière au lieu de faire comme Eschyle, tantôt échouent complètement, tantôt font mauvaise figure dans les concours. Ainsi, Agathon échoua sur ce seul point.}}
\end{multicols}
}%
\end{adjustbox}%
\vfill
\end{adjustwidth}
\newpage
% --- PAGE DE DROITE : Traduction française + Commentaire (8/8) ---
\begin{adjustwidth}{-2cm}{-2cm}
\noindent\textbf{\footnotesize Traduction française --- Folio 03 (\texttt{f. 163v}, vv. 80–122 [4/4]) \hfill ΠΡΟΛΟΓΟΣ}\\[-2pt]
\noindent\rule{16.5cm}{0.5pt}\\[2pt]
{\fontsize{8.8pt}{10.2pt}\selectfont
\ligneEst{}{\francais{Phébus, qui fonda les murs}}{}
\ligneEst{}{\francais{de cette cité dans la contrée du Simois.}}{110}
\ligneEst{}{\francais{4° — Sois salué de nos plus beaux chants,}}{}
\ligneEst{}{\francais{Phébus, qui dans les fêtes des Muses}}{}
\ligneEst{}{\francais{emportes le prix sacré.}}{}
\ligneEst{}{\francais{5° — Et chantez la Vierge qui se plaît dans les}}{}
\ligneEst{}{\francais{montagnes dryogènes, Artémis-chasseresse.}}{}
\ligneEst{}{\francais{6° — À ta suite je célèbre et glorifie}}{}
\ligneEst{}{\francais{l'auguste rejeton de Léto,}}{}
\ligneEst{}{\francais{Artémis ignorant l'hymen.}}{}
\ligneEst{}{\francais{7° — Ainsi que Léto et les percussions de}}{}
\ligneEst{}{\francais{la cithare d'Asie, aux beaux rythmes contrariés marqués par le pied}}{120}
\ligneEst{}{\francais{sous la direction des Charites phrygiennes.}}{}
}
\vspace{2pt}
\noindent\rule{16.5cm}{0.6pt}\\[2pt]
\noindent\textbf{\small Commentaire philologique, métrique et scénique (vv. 80–122 --- partie 8/8)}\\[2pt]
\begin{adjustbox}{max totalheight=15.2cm, max width=16.5cm, keepaspectratio}%
\parbox[t]{16.5cm}{
\setlength{\columnsep}{9pt}%
\begin{multicols}{2}
\fontsize{8.4pt}{9.7pt}\selectfont
\setlength{\parindent}{0pt}%
\setlength{\parskip}{1.2pt}%
\noindent Je me plais à imaginer Aristophane, révolté devant cet échec injuste, remettre sur scène ce qu'il considérait comme la perle (érotique) du cycle troyen d'Agathon (c'est en tout cas ce sens de la justice qui l'avait poussé à vouloir faire rejouer ses \italique{Nuées}, mésestimées par les mauvais jurés athéniens).
\noindent\textbf{[Notes ad 101-29]} \textbf{(1)} in \italique{Essai sur la composition des comédies d'Aristophane} (1904), p. 127-128.
\noindent \textbf{(2)} À ce propos, cf. les \italique{Grenouilles} v. 1299 où Eschyle dit s'être écarté du pré des Muses de Phrynichos – Muses par trop efféminées et ionisantes ; de même, l'\italique{Assemblée des femmes} (v. 883).
\noindent \textbf{(3)} Le style d'Agathon est certes efféminé mais beaucoup moins caricatural qu'on a voulu le faire croire. Cf. le fragment 178 (Γηρυτάδης) de l'édition Kassel-Austin : Aristophane y emploie l'adjectif Ἀγαθώνειος dans un sens loin d'être péjoratif d'après les commentaires anciens (Zénobius I, 2 ; Hésych. α 281 ; Phot. α 83 ; Suid. α 125 ; \&c.) :
\citComm{Ἀγαθώνιος αὔλησις: ἡ μαλακὴ καὶ ἐκλελυμένη: ἢ ἡ μήτε χαλαρὰ, μήτε πικρὰ, ἀλλ' εὔκρατος καὶ ἡδίστη.\\
Un air de flûte agathonien : doux et détendu ; ni mol ni strident, mais tempéré et très doux.}
\noindent \textbf{(4)} \italique{Dans l'ouvroir du poète}, \italique{Hellenica} (2008), pp. 35-45. J'eus la chance de recevoir d'Alain Blanchard peu de temps avant sa mort de vifs encouragements, source d'énergie nécessaire à l'abattage de la somme de travail que représente ce commentaire des \italique{Thesmophories}.
\noindent \textbf{(5)} Cf. pour plus de détails l'article d'A. Blanchard avec les chiffres à la clef.
\end{multicols}
}%
\end{adjustbox}%
\vfill
\end{adjustwidth}
\newpage
\clearpage\ifodd\value{page}\hbox{}\newpage\fi
\enTeteG{ΠΡΟΛΟΓΟΣ --- Codex Ravennas 429, f. 164r (Classense 166r) --- vv. 123–164}
\begin{adjustwidth}{-1.5cm}{-3.0cm}
\begin{adjustbox}{max totalheight=21.0cm, max width=16.5cm, keepaspectratio}
\begin{minipage}{16.5cm}
{\fontsize{7.2pt}{8.5pt}\selectfont
\noindent\textbf{Folio 04 --- Codex Ravennas 429, \texttt{f. 164r} (\texttt{Classense 166r}) : ΠΡΟΛΟΓΟΣ, vv. 123–164}\hfill \textit{Témoins :} $\text{R}^1, \Rpc, \Rdeux$, \PsiUn\par\vspace{2pt}
\hrule height 0.4pt\vspace{2.5pt}
}
{\fontsize{8.1pt}{9.5pt}\selectfont
\ligneRav{123}{ἀγάθ/}{\hspace{1.1em}σέβομαι Λατώ τ' ἄνασσαν}{}
\ligneRav{}{}{\hspace{1.1em}κίθαρίν τε ματέρ᾽ ὕμνων}{}
\ligneRav{125}{}{\hspace{1.1em}ἄρσενι βοᾷ δοκιμῳ·}{\textbf{125} \lecon{δοκιμῳ} R, Σ\expo{1}, Suid. α 4016 : δόκιμον Σ\expo{2} : δοκίμων Schöne}
\ligneRav{126}{}{\hspace{1.1em}τᾷ \difference{φάος} ἔσσυτο δαιμονίοις}{\textbf{126} φῶς R : \lecon{φάος} Burges}
\ligneRav{127}{}{\hspace{1.1em}ὄμμασιν \difference{ἁμετέρας} τε δι᾽ αἰφνι-}{\textbf{127} ἡμετέρας R : ὑμετέρας Nietzsche : \lecon{ἁμετέρας} Meineke}
\ligneRav{127}{}{\hspace{1.1em} δίου ὀπός. ὧν χάριν}{\textbf{127} ἡμετέρας R : ὑμετέρας Nietzsche : \lecon{ἁμετέρας} Meineke}
\ligneRav{128}{}{\hspace{1.1em}ἄνακτ᾽ ἄγαλλε \difference{Φοῖβον}·}{\textbf{128} \lecon{ἄγαλλε} R : ἀγάλλετε Austin (1987) \espApp Φοῖβον τιμᾷ R : \lecon{Φοῖβον} Dindorf}
\ligneRav{}{}{\hspace{1.1em}χαῖρ᾽ ὄλβιε παῖ Λατοῦς·}{}
\ligneRav{}{}{\hspace{2.2em}\textbf{[παρεπιγραφή R\expo{1} :]} \lecon{ὀλολύζει ὁ γέρων}}{\textbf{post 129} \lecon{ὀλολύζει ὁ γέρων} (παρεπιγραφή) R}
\ligneRav{130}{μ\auDessus{ν}{Η}}{ὡς ἡδὺ τὸ μέλος ὦ \difference{πότνιαι} Γενετυλλίδες}{\textbf{130} πότνια R : \lecon{πότνιαι} Suid. γ 141}
\ligneRav{}{}{καὶ θηλυδριῶδες καὶ κατεγλωττισμένον}{}
\ligneRav{}{}{καὶ μανδαλωτόν, ὥστ᾽ ἐμοῦ γ᾽ ἀκροωμένου}{}
\ligneRav{}{}{ὑπὸ τὴν ἕδραν αὐτὴν ὑπῆλθε γάργαλος·}{}
\ligneRav{134}{}{καί σ᾽ ὦ \difference{νεανίσκ᾽ εἴ τις} εἶ, κατ᾽ Αἰσχύλον}{\textbf{134-5} νεανίσχ᾽ εἴ τις R : \lecon{νεανίσκ᾽ εἴ τις} \Mon : νεανίσχ᾽ ἥτις Gelenius \espApp Λυκουργίας R : \lecon{Λυκουργείας} Dindorf}
\ligneRav{135}{}{ἐκ τῆς Λυκουργείας ἐρέσθαι βούλομαι·}{}
\ligneRav{}{}{ποδαπὸς ὁ γύννις; τίς πάτρα; τίς ἡ στολή;}{}
\ligneRav{}{}{τίς ἡ τάραξις τοῦ βίου; τί βάρβιτος}{}
\ligneRav{}{}{λαλεῖ κροκωτῷ; τί δὲ λύρα κεκρυφάλῳ;}{}
\ligneRav{139}{}{τί λήκυθος καὶ στρόφιον; ὡς οὐ ξύμφορον·}{\textbf{139} \lecon{οὐ ξύμ-} R, Suid. β 190 V : ἀξυμ- Suid. M : ὀξύμ- Suid. A,G,I : ξύμ- Suid. F}
\ligneRav{140}{}{τίς δαὶ \difference{κατρόπτου} καὶ ξίφους κοινωνία;}{\textbf{140} \lecon{κατόπτρου} R, Suid. β 110 : \lecon{κατρόπτου} Dover ad \italique{Ran.} 47}
\ligneRav{141}{}{\difference{σύ τ᾽} αὐτὸς ὦ παῖ; πότερον ὡς ἀνὴρ τρέφει;}{\textbf{141} \lecon{σύ τ᾽} \PsiUn, Σ : τίς δ᾽ R}
\ligneRav{142}{}{καὶ ποῦ \difference{πέος}; ποῦ χλαῖνα; ποῦ Λακωνικαί;}{\textbf{142} \lecon{πέος} \PsiUn, R : τὸ πέος Suid. λ 63, π 987}
\ligneRav{}{}{ἀλλ᾽ ὡς γυνὴ δῆτ᾽ εἶτα ποῦ τὰ τιτθία;}{}
\ligneRav{}{}{τί φῄς; τί σιγᾷς; ἀλλὰ δῆτ᾽ ἐκ τοῦ μέλους}{}
\ligneRav{145}{}{ζητῶ σ᾽, ἐπειδή γ᾽ αὐτὸς οὐ βούλει φράσαι;}{}
\ligneRav{}{ἀγάθ/}{ὦ πρέσβυ πρέσβυ, τοῦ φθόνου μὲν τὸν ψόγον}{}
\ligneRav{}{}{ἤκουσα, τὴν δ᾽ ἄλγησιν οὐ παρεσχόμην}{}
\ligneRav{148}{}{ἐγὼ δὲ τὴν ἐσθῆθ᾽ ἅμα \difference{τῇ γνώμῃ} φορῶ·}{\textbf{148} γνώμη R : \lecon{τῇ γνώμῃ} Meineke}
\ligneRav{}{}{χρὴ γὰρ ποιητὴν ἄνδρα πρὸς τὰ δράματα}{}
\ligneRav{150}{}{ἃ δεῖ ποιεῖν πρὸς ταῦτα τοὺς τρόπους ἔχειν·}{}
\ligneRav{}{}{αὐτίκα γυναικεῖ᾽ ἢν ποιῇ τις δράματα,}{}
\ligneRav{}{}{μετουσίαν δεῖ τῶν τρόπων τὸ σῶμ᾽ ἔχειν·}{}
\ligneRav{}{μ\auDessus{ν}{Η}}{οὐκοῦν κελητίζεις, ὅταν Φαίδραν ποιῇς;}{}
\ligneRav{154}{ἀγάθ/}{ἀνδρεῖα δ᾽ \difference{ἄν} ποιῇ τις, ἐν τῷ σώματι}{\textbf{154} \lecon{ἄν} \PsiUn, R : ἢν Dindorf}
\ligneRav{155}{}{ἔνεσθ᾽ ὑπάρχον τοῦθ᾽. ἃ δ᾽ οὐ κεκτήμεθα,}{}
\ligneRav{}{}{μίμησις ἤδη ταῦτα συνθηρεύεται·}{}
\ligneRav{}{μ\auDessus{ν}{Η}}{ὅταν σατύρους τοίνυν ποιῇς, καλεῖν ἐμέ,}{}
\ligneRav{}{}{ἵνα συμποιῶ σοὔπισθεν ἐστυκὼς ἐγώ·}{}
\ligneRav{}{ἀγάθ/}{ἄλλως τ᾽ ἄμουσόν ἐστι ποιητὴν ἰδεῖν}{}
\ligneRav{160}{}{ἀγρεῖον ὄντα καὶ δασύν σκέψαι δ᾽ ὅτι}{}
\ligneRav{}{}{Ἴβυκος ἐκεῖνος κἀνακρέων ὁ Τήιος}{}
\ligneRav{162}{}{κἀλκαῖος, οἳ περὶ ἁρμονίαν \difference{ἐκεχήνεσαν},}{\textbf{162} \lecon{κἀλκαῖος} R, Suid. : (κ)ἀχαιός Σ \espApp ἐχύμησαν R : ἐχύμισαν Suid. : \lecon{ἐκεχήνεσαν} Bibi}
\ligneRav{163}{}{ἐμιτροφόρουν τε καὶ \difference{διέκειντ'} Ἰωνικῶς,}{\textbf{163} διεκίνων R : διεκίνουν Suid. ε 989, ι 495 : διεκλῶντ᾽ Toup : \lecon{διέκειντ'} Bibi}
\ligneRav{}{}{καὶ Φρύνιχος, τοῦτον γὰρ οὖν ἀκήκοας,}{}
}
\vspace{3pt}\noindent\rule{4.5cm}{0.4pt}\par\vspace{1.5pt}
{\fontsize{7.0pt}{8.3pt}\selectfont
\noindent\textbf{Loci paralleli (vv. 123–164) ---} \textbf{125} (ἄρσενι βοᾷ δοκιμῳ) Suidas α 4016 $=$ Scholium \espTest \textbf{128} (ἄγαλλε) Photius α 164 $=$ \SynB α 145 \espTest \textbf{130-2} (ἡδὺ -- μανδαλωτόν) Suidas γ 141 \espTest \textbf{130} (Γενετυλλίδες) Hesychius γ 343 Γενετυλλίς \espTest \textbf{131-2} (θηλυδριῶδες -- μανδαλωτόν) Suidas μ 134 \espTest \textbf{131} (θηλυδριῶδες καὶ κατεγλωττισμένον) Suidas κ 912 \espTest Hesychius θ 491 θηλυδριῶδες \espTest (κατεγλωττισμένον) Pollux II.109 \espTest \textbf{132} Hesychius μ 224 μανδαλωτόν $\simeq$ Photius μ 86, cf. Photius γ 115 \espTest \textbf{132-3} (ὥστ᾽ -- γάργαλος) Suidas γ 83 (sed τοῦ μέλους post ἀκροωμένου addito) \espTest \textbf{133} (γάργαλος $=$ γαργαλισμός) Moeris γ 23, cf. \Phryn 56.9 \espTest \textbf{134} Hesychius η 948 ἥτις· ὅστις (?) \espTest \textbf{136} (ποδαπὸς -- πάτρα) Suidas γ 504 \espTest (γύννις) Hesychius γ 1015 γύννις, cf. \Phryn 60.19 \espTest \textbf{136-40} (τίς ἡ στολή -- κοινωνία) Suidas β 110 (δὲ pro δαὶ) \espTest \textbf{137} (βάρβιτος) Hesychius β 211 βάρβιτος, cf. Photius β 62 \espTest \textbf{141-2} (πότερον -- πέος) Suidas π 987 \espTest \textbf{142} (ποῦ τὸ πέος -- Λακωνικαί) Suidas λ 63, π 987 (τὸ πέος pro πέος) \espTest Moeris χ 19 χλαῖνα \espTest (Λακωνικαί) Hesychius λ 225 Λακωνικαί ; Photius λ 57 Λακωνικαί ; scholium \italique{Vespae} 1158-60 \espTest \textbf{142} (ποῦ τὸ πέος -- Λακωνικαί) Suidas λ 63 \espTest \textbf{146-7} Suidas α 1097 (μὲν omisso et δὲ pro δ') \espTest \textbf{149-50} \Apostol XVIII 41h \espTest \textbf{159-61} Suidas α 1633 (κ' Ἀνακρ.) \espTest \textbf{161-2} Suidas ε 4049 (Ἀνακρέων, Ἀλκαῖος, οἵπερ ἁρμονίαν ἐχύμισαν) \espTest \textbf{162-3} (κ' Ἀλκαῖος -- Ἰωνικῶς) Suidas ε 989 (sed οἵπερ ἁρμονίαν ἐχύμισαν) $\simeq$ Zonaras 707 ἐμιτροφόρουν \espTest \textbf{162} Zonaras 942 οἵπερ ἁρμονίαν ἐχύμισαν \espTest (ἐχύμισαν) Hesychius ε 7681 ἐχύμισαν· ἤρτυσαν \espTest \textbf{162-3} (κ' Ἀλκαῖος -- Ἰωνικῶς) Suidas ε 989 $\simeq$ Zonaras 707 ἐμιτροφόρουν \espTest (ἐχύμισαν) Hesychius ε 7681 ἐχύμισαν \espTest \textbf{163} Suidas ι 495 ἐμιτροφόρουν τε καὶ διεκίνουν Ἰωνικῶς\par
}
\vspace{2.5pt}\noindent\rule{4.5cm}{0.4pt}\par\vspace{1.5pt}
{\fontsize{6.9pt}{8.2pt}\selectfont
\noindent\textbf{Suite (en commentaire) de l'apparat critique \& notes paléographiques (\texttt{f. 164r}) ---} \textbf{Témoin papyrologique :} \textbf{\PsiUn{} (\italique{PSI} XI 1194 + XIV p. XV)} couvre les \textbf{vv. 133–156}. --- \textbf{Deuxième \italique{parepigraphè} (entre vv. 129 et 130) :} $\text{R}^1$ copie sur une ligne propre \lecon{ὀλολύζει ὁ γέρων}. --- \textbf{Confrontation avec \PsiUn{} (v. 141) :} Au v. 141, R porte la leçon fautive τίς δ' αὐτὸς ὦ παῖ· πότερον ὡς ἀνὴρ τρέφει;, tandis que le papyrus \PsiUn{} et la scholie $\Sigma^{\text{R}}$ conservent l'authentique \lecon{σύ τ'}. --- \textbf{Double émendation paléographique de \textbf{Bibi} aux vv. 162–163 :} ἐχύμησαν ... διεκίνων R $\to$ \lecon{ἐκεχήνεσαν ... διέκειντ'} \textbf{Bibi}. --- \textbf{v. 125 :} \lecon{δοκίμῳ} R, $\Sigma^1$, Souda $\alpha\ 4016$ : δόκιμον $\Sigma^2$ : δοκίμων Schöne. --- \textbf{v. 127 :} ἡμετέρας R (ὑμετέρας Nietzsche) : \lecon{ἁμετέρας} Meineke (dorisme lyrique). --- \textbf{v. 128 :} \lecon{ἄγαλλε} R (ἀγάλλετε Austin 1987) \espApp Φοῖβον τιμᾷ R (τιμᾷ est une glose intruse) : \lecon{Φοῖβον} Dindorf. --- \textbf{vv. 134–135 :} νεανίσχ' εἴ τις R (aspiration fautive par anticipation de l'esprit rude) : \lecon{νεανίσκ' εἴ τις} \Mon{} (νεανίσχ' ἥτις Gelenius) \espApp \textbf{v. 135 :} Λυκουργίας R (iotacisme) : \lecon{Λυκουργείας} Dindorf. --- \textbf{v. 141 :} \lecon{σύ τ'} $\Psi^1, \Sigma$ : τίς δ' R. --- \textbf{v. 162 :} \lecon{κἀλκαῖος} R, Souda ((κ)ἀχαιός $\Sigma$) \espApp ἐχύμησαν R (ἐχύμισαν Souda) : \lecon{ἐκεχήνεσαν} \lecon{Bibi}. \textbf{Explication paléographique (Commentaire ad v. 162) :} Dans ἐκεχήνεσαν (plus-que-parfait de $\text{χάσκω}$, « ils restaient bouche bée », cf. \italique{Ach.} 10), le saut visuel εκεχ- $\to$ εχ- (haplographie), joint à l'iotacisme $\eta \to \upsilon$ et à la confusion minuscule constante $\nu \to \mu$, a produit le monstre ἐχύμησαν de R ! --- \textbf{v. 163 :} διεκίνων R (διεκίνουν Souda $\varepsilon\ 989, \iota\ 495$ ; διεκλῶντ' Toup) : \lecon{διέκειντ'} \lecon{Bibi} (ou δίκανον ap. Hésychius $\delta\ 1805$).\par
}
\end{minipage}
\end{adjustbox}
\end{adjustwidth}
\newpage
\enTeteD{Codex Ravennas 429 (Biblioteca Classense) --- Fac-similé du f. 164r (vv. 123–164)}
\begin{adjustwidth}{-1.5cm}{-3.0cm}
\begin{center}
\includegraphics[width=16.5cm,height=20.5cm,keepaspectratio]{/workspaces/Thesmophories/dossier_ravennas_thesmophories/folios_individuels_ff162v-177r/04_f164r_classense166r_vv0123-0164.jpg}\par
\vspace{2pt}
{\scriptsize\textbf{Codex Ravennas 429, f. 164r (Classense 166r)} --- ΠΡΟΛΟΓΟΣ, \textbf{vv. 123–164} ($\text{R}^1, \Rpc, \Rdeux$, \PsiUn).}
\end{center}
\end{adjustwidth}
\newpage
\clearpage\ifodd\value{page}\hbox{}\newpage\fi% --- FOLIO 04 (f. 164r) : SPREAD 2.1 (Texte grec établi + Traduction française + Commentaire 1/8 & 2/8) ---
\clearpage\ifodd\value{page}\hbox{}\newpage\fi
\enTeteG{ΘΕΣΜΟΦΟΡΙΑΖΟΥΣΑΙ --- Texte grec établi (f. 164r, vv. 123–164 [1/4])}
\enTeteD{LES THESMOPHORIES --- Traduction française (f. 164r, vv. 123–164 [1/4])}
\begin{adjustwidth}{-2cm}{-2cm}
\noindent\textbf{\footnotesize Texte grec établi --- Folio 04 (\texttt{f. 164r}, vv. 123–164 [1/4]) \hfill ΠΡΟΛΟΓΟΣ}\\[-2pt]
\noindent\rule{16.5cm}{0.5pt}\\[2pt]
{\fontsize{8.8pt}{10.2pt}\selectfont
\ligneEst{{\tiny\textit{io. min. + troch.}}}{σέβομαι Λατώ τ' ἄνασσαν}{123}
\ligneEst{{\tiny\textit{2 io. min. anacl.}}}{κίθαρίν τε ματέρ᾽ ὕμνων}{}
\ligneEst{{\tiny\textit{troch. + io. min\cat}}}{ἄρσενι βοᾷ δοκιμῳ.}{125}
\ligneEst{{\tiny\textit{4 dact\cat}}}{τᾷ φάος ἔσσυτο δαιμονίοις}{}
\ligneEst{{\tiny\textit{6 dact. acatal.}}}{ὄμμασιν ἁμετέρας τε δι᾽ αἰφνι-}{}
\ligneEst{}{\alinea δίου ὀπός. ὧν χάριν}{}
\ligneEst{{\tiny\textit{2 iamb\cat}}}{ἄνακτ᾽ ἄγαλλε Φοῖβον.}{}
\ligneEst{{\tiny\textit{telesil. scazon}}}{χαῖρ᾽ ὄλβιε παῖ Λατοῦς.}{}
\ligneEst{\textbf{Μνησίλοχος} {\tiny\textit{3 ia.}}}{ὡς ἡδὺ τὸ μέλος ὦ πότνιαι Γενετυλλίδες}{130}
\ligneEst{}{καὶ θηλυδριῶδες καὶ κατεγλωττισμένον}{}
}
\vspace{2pt}
\noindent\rule{16.5cm}{0.6pt}\\[2pt]
\noindent\textbf{\small Commentaire philologique, métrique et scénique (vv. 123–164 --- partie 1/8)}\\[2pt]
\begin{adjustbox}{max totalheight=15.2cm, max width=16.5cm, keepaspectratio}%
\parbox[t]{16.5cm}{
\setlength{\columnsep}{9pt}%
\begin{multicols}{2}
\fontsize{8.4pt}{9.7pt}\selectfont
\setlength{\parindent}{0pt}%
\setlength{\parskip}{1.2pt}%
\noindent\textbf{130} \textbf{ὦ πότνιαι Γενετυλλίδες.} Déesses de la génération, les Génétyllides sont aussi citées par Aristophane au vers 52 des \italique{Nuées} et au second vers de \italique{Lysistrata} ; de ce dernier passage on peut inférer que les femmes étaient réputées célébrer des rites licencieux en leur honneur.
\noindent\textbf{131} \textbf{θηλυδριῶδες.} Semble être un néologisme d'Aristophane. Chantraine\marqueur{1} suppose un relais *θηλύδριον qui expliquerait aussi θηλυδρίας, bien attesté. Aristophane a-t-il porté son choix sur θηλυδριώδης pour faire entendre θηριώδης et toute l'animalité du spasme de Mnésiloque ? (cf. κατεγλωττισμένον et μανδαλωτόν.)
\noindent\textbf{131-2} \textbf{καὶ κατεγλωττισμένον / καὶ μανδαλωτόν} : ces mots renvoient tous deux à ce que les anglophones appellent – ne boudons pas l'hommage qui nous est ainsi rendu ! – \italique{french kiss}. À propos de μανδαλωτόν, je ne résiste pas à citer Taillardat (§195) qui met au jour la bigarrure de la langue aristophanienne : « le substantif, τὸ φίλημα, \italique{le baiser}, est déjà connu d'Eschyle ; il se trouve chez Sophocle, chez Euripide, chez les comiques Platon et Nicophon. En revanche, il n'est pas attesté dans les œuvres conservées d'Aristophane, bien que le poëte emploie φιλεῖν comme synonyme de κυνεῖν [...] Aristophane préfère des mots imagés et comme ``techniques'' : τὸ περιπεταστόν (\italique{Acharniens} 1201), τὸ ἐπιμανδαλωτόν, \italique{ib.}, τὸ μανδαλωτόν (\italique{Thesm.} 132), les deux derniers termes évoquant l'idée d'une porte fermée ; ὁ μάνδαλος, comme βάλανος, désigne la clavette assurant la barre de la porte. La clavette, c'est ici la langue. »
\noindent\textbf{134} \textbf{νεανίσχ᾽ εἴ τις εἶ.} Le manuscrit de Ravenne comporte une coquille : le κ de νεανίσκ᾽ est aspiré alors qu'il est suivi de εἴ. Plusieurs solutions ont été apportées : lire ἥτις (justifiant le féminin par Hésychius (η 948) glosant ἥτις· ὅστις). D'autres se sont contentés de substituer un κ au χ. J'ai une préférence pour cette dernière correction\marqueur{1} qui me semble être une pique contre Agathon : « jeune homme... si tu en es un ! » car Agathon serait âgé au moment de la représentation de 35 ans ou environ ; à croire que son mode de vie le conserve bien... mais tout-de-même ! cette allusion permanente à la jeunesse d'Agathon – comme dans le \italique{Banquet} de Platon censé se passer en 416, cinq ans plus tôt – me paraît plutôt être une allusion au fait qu'Agathon pèche par certain attribut pour être tout-à-fait considéré comme un homme\marqueur{2}.
\end{multicols}
}%
\end{adjustbox}%
\vfill
\end{adjustwidth}
\newpage
% --- PAGE DE DROITE : Traduction française + Commentaire (2/8) ---
\begin{adjustwidth}{-2cm}{-2cm}
\noindent\textbf{\footnotesize Traduction française --- Folio 04 (\texttt{f. 164r}, vv. 123–164 [1/4]) \hfill ΠΡΟΛΟΓΟΣ}\\[-2pt]
\noindent\rule{16.5cm}{0.5pt}\\[2pt]
{\fontsize{8.8pt}{10.2pt}\selectfont
\ligneEst{}{\francais{8° — J'honore Léto la souveraine}}{123}
\ligneEst{}{\francais{et la cithare mère des hymnes,}}{}
\ligneEst{}{\francais{à la voix mâle et éprouvée.}}{125}
\ligneEst{}{\francais{9° — À ces accents, une lumière jaillit}}{}
\ligneEst{}{\francais{dans les yeux des dieux}}{}
\ligneEst{}{\francais{et traverse notre soudain regard.}}{}
\ligneEst{}{\francais{Aussi bien, glorifie le seigneur Phébus.}}{}
\ligneEst{}{\francais{10° — Salut, heureux fils de Léto.}}{}
\ligneEst{\textbf{Mnésiloque}}{\francais{Quel chant suave, ô Génétyllides vénérées}}{130}
\ligneEst{}{\francais{et qui sent la féminie et les coups de langue}}{}
}
\vspace{2pt}
\noindent\rule{16.5cm}{0.6pt}\\[2pt]
\noindent\textbf{\small Commentaire philologique, métrique et scénique (vv. 123–164 --- partie 2/8)}\\[2pt]
\begin{adjustbox}{max totalheight=15.2cm, max width=16.5cm, keepaspectratio}%
\parbox[t]{16.5cm}{
\setlength{\columnsep}{9pt}%
\begin{multicols}{2}
\fontsize{7.6pt}{8.8pt}\selectfont
\setlength{\parindent}{0pt}%
\setlength{\parskip}{1.2pt}%
\noindent\textbf{134-5} \textbf{κατ᾽ Αἰσχύλον / ἐκ τῆς Λυκουργείας.} D'après le scholiaste, la \italique{Lycurgie} est le nom collectif d'une tétralogie d'Eschyle, composée des \italique{Édoniens} (fr.\marqueur{1} 57-67), des \italique{Bassarides} (fr. 23-25), des \italique{Jouvenceaux} (Νεανίσκοι, fr. 146-9) et du drame satyrique \italique{Lycurgue} (fr. 124-6). Elle compte les mésaventures du roi Lycurgue essayant de contrer l'invasion du culte dionysiaque en Thrace. Laissons la parole à François Jouan\marqueur{2} : « La légende de Lycurgue était déjà connue d'Homère. L'\italique{Iliade} rappelait sommairement que Lycurgue avait été aveuglé par Zeus pour avoir pourchassé les Nourrices de Dionysos et forcé le petit dieu à se jeter à la mer\marqueur{3}. La version eschyléenne est très différente, et comporte en particulier une localisation précise : le héros est roi des Édoniens et vit en Thrace, dans la région du Pangée. On admet en général que cette tétralogie [...] datait de la fin de la vie d'Eschyle, entre 466 et 459. Or dans cette période, des forces athéniennes commandées par Cimon guerroyaient précisément contre les Édoniens, autour d'Eion et d'Ennéahodoi. [...] Les \italique{Édoniens} sont la pièce la moins mal connue, en particulier du fait de l'existence de fragments d'un \italique{Lycurgue} du poète latin Naevius, qui semble avoir suivi d'assez près le drame d'Eschyle. Dans le prologue, un messager annonçait au roi thrace l'arrivée d'un personnage entouré d'un bruyant cortège de musiciens et de danseurs – Dionysos et son thiase. Le roi ordonnait de les arrêter et de les faire comparaître devant lui. Dans la parodos, le chœur, sans doute composé de sujets de Lycurgue, décrivait le cortège dionysiaque. Un garde rendait compte de l'arrestation des trublions et de leur docilité. Le roi interrogeait lui-même Dionysos, comme le Penthée des \italique{Bacchantes}. Aussi mécontent des réponses de Dionysos que le prince thébain, il l'emprisonnait avec sa suite dans son palais, où il devait trouver, lui disait-il ironiquement «~un logement gratuit [ἄδασμος, fr. 63]~». Mais la défaite du dieu n'était que provisoire. Bientôt les liens des captifs tombaient d'eux-mêmes, Dionysos apparaissait dans sa gloire sur le toit du palais, qui vacillait sur sa base et se nimbait de flammes : autre motif inclus dans les \italique{Bacchantes}\marqueur{5}. » \\ C'est donc bien à cette première tragédie des \italique{Édoniens} que semble se référer Mnésiloque, interrogeant à la manière de Lycurgue, associant ainsi Agathon à Dionysos (métaphore qui avait commencée d'être filée au vers 41 quand le serviteur d'Agathon parlait du thiase de Muses réunies chez son maître). Rappelons que, pour les Grecs de l'époque, Dionysos était un dieu jeune : « Les particularités de son physique et de son costume sont détaillées aussi bien par Eschyle\marqueur{6} que par Euripide. Il est grand et mince avec de longues jambes, une chevelure blonde, abondante et parfumée, de belles boucles. Il est beau, avec un teint vermeil et des traits de visage plutôt féminins, l'air souriant. Il est vêtu d'une robe asiatique de fin tissu, tombant jusqu'à ses pieds et tient en main le thyrse. »\marqueur{7}
\end{multicols}
}%
\end{adjustbox}%
\vfill
\end{adjustwidth}
\newpage

% --- FOLIO 04 (f. 164r) : SPREAD 2.2 (Texte grec établi + Traduction française + Commentaire 3/8 & 4/8) ---
\clearpage\ifodd\value{page}\hbox{}\newpage\fi
\enTeteG{ΘΕΣΜΟΦΟΡΙΑΖΟΥΣΑΙ --- Texte grec établi (f. 164r, vv. 123–164 [2/4])}
\enTeteD{LES THESMOPHORIES --- Traduction française (f. 164r, vv. 123–164 [2/4])}
\begin{adjustwidth}{-2cm}{-2cm}
\noindent\textbf{\footnotesize Texte grec établi --- Folio 04 (\texttt{f. 164r}, vv. 123–164 [2/4]) \hfill ΠΡΟΛΟΓΟΣ}\\[-2pt]
\noindent\rule{16.5cm}{0.5pt}\\[2pt]
{\fontsize{8.8pt}{10.2pt}\selectfont
\ligneEst{}{καὶ μανδαλωτόν, ὥστ᾽ ἐμοῦ γ᾽ ἀκροωμένου}{}
\ligneEst{}{ὑπὸ τὴν ἕδραν αὐτὴν ὑπῆλθε γάργαλος.}{}
\ligneEst{}{καί σ᾽ ὦ νεανίσκ᾽ εἴ τις εἶ, κατ᾽ Αἰσχύλον}{}
\ligneEst{}{ἐκ τῆς Λυκουργείας ἐρέσθαι βούλομαι.}{135}
\ligneEst{}{ποδαπὸς ὁ γύννις; τίς πάτρα; τίς ἡ στολή;}{}
\ligneEst{}{τίς ἡ τάραξις τοῦ βίου; τί βάρβιτος}{}
\ligneEst{}{λαλεῖ κροκωτῷ; τί δὲ λύρα κεκρυφάλῳ;}{}
\ligneEst{}{τί λήκυθος καὶ στρόφιον; ὡς οὐ ξύμφορον.}{}
\ligneEst{}{τίς δαὶ κατρόπτου καὶ ξίφους κοινωνία;}{140}
\ligneEst{}{σύ τ᾽ αὐτὸς ὦ παῖ; πότερον ὡς ἀνὴρ τρέφει;}{}
\ligneEst{}{καὶ ποῦ πέος; ποῦ χλαῖνα; ποῦ Λακωνικαί;}{}
}
\vspace{2pt}
\noindent\rule{16.5cm}{0.6pt}\\[2pt]
\noindent\textbf{\small Commentaire philologique, métrique et scénique (vv. 123–164 --- partie 3/8)}\\[2pt]
\begin{adjustbox}{max totalheight=15.2cm, max width=16.5cm, keepaspectratio}%
\parbox[t]{16.5cm}{
\setlength{\columnsep}{9pt}%
\begin{multicols}{2}
\fontsize{7.6pt}{8.8pt}\selectfont
\setlength{\parindent}{0pt}%
\setlength{\parskip}{1.2pt}%
\noindent (Ajoutons à propos de cette \italique{crocote} au cœur des péripéties de Mnésiloque qu'elle est aussi associée à Dionysos dans les \italique{Grenouilles} (v.46) ; par Cratinos, fr. 40 de l'édition de Kassel-Austin et, comme on l'a vu, Eschyle, fr. 59.)  Voilà qui explique un peu les associations d'idées de Mnésiloque !
\noindent\textbf{[Notes ad 131]} \textbf{(1)} \italique{Dictionnaire étymologique de la langue grecque}, \italique{sub voce} θηλή.
\noindent \textbf{(1)} Signalons qu'une coquille analogue (οὐχ᾽ εἴρηχ᾽ R en place d'un οὐκ εἴρηχ᾽) se retrouve au vers 556. De plus, les anciens manuscrits sur lesquels on a copié les \italique{Thesophories} n'étaient pas systématiquement élidés ; ce genre de fautes s'explique donc aisément. Cf. par exemple v. 148 ἐσθητ' ἅμα ; v. 279 θραττ' έπου ; v. 1178 ἔρχετ' ὡς.
\noindent \textbf{(2)} καὶ ποῦ πέος; demande Mnésiloque au vers 142.
\noindent \textbf{(1)} Les fragments d'Eschyle sont cités selon l'édition de S. Radt, \italique{Tragicorum Graecorum Fragmenta}, vol. 3 (1985).
\noindent \textbf{(2)} On lira avec intérêt deux de ses articles : « Dionysos chez Eschyle » (in \italique{Kernos} 5, 1992, p. 71-86) et « Une introduction au mythe des \italique{Bacchantes} » (\italique{Kentron} 14, 1998). C'est le premier que nous citons.
\noindent \textbf{(3)} \italique{Iliade}, VI, 130-140.
\noindent \textbf{(4)} Eur. \italique{Bacch}., 586-599 ; 622-624 ; 633.
\noindent \textbf{(5)} Cf. Hérodote, II, 145.
\noindent \textbf{(6)} Voir les fragments 59, 60, 62 des \italique{Édoniens}.
\noindent \textbf{(7)} F. Jouan, \italique{Une introduction au mythe des} Bacchantes, p. 12.
\noindent\textbf{139-44} \textbf{Fragments d'un ancien papyrus.} L'on possède quelques fragments d'un ancien papyrus (\PsiUn \ dans notre apparat) du premier ou second siècle ap. J.-C. pour les vers 133-144 ; 145-156 ; 237-245 ; 273-288 ; 594-596 ; 804-809. Il nous apprend en outre le \latin{numerus versuum} d'une ancienne édition des \italique{Thesmophories} : la pièce comportait alors 1201 vers (contre 1231 pour celle de Brunck sur laquelle se basent tous les éditeurs modernes). La papyrus est reproduit en annexe. Cf. \italique{Papiri della Società Italiana} XI, n° 1194 et  \italique{Papiri della Società Italiana} XIV, p. XV.
\noindent\textbf{139} \textbf{τί λήκυθος καὶ στρόφιον;} Dans sa liste des attributs masculins et féminins qui jurent ensemble chez Agathon, Mnésiloque oppose ici le \italique{lécythe} et le \italique{strophion}. Si le cas du strophion ne pose pas de problème : il s'agit du \latin{strophium} ou \latin{fascia pectoralis}\marqueur{1} : bandelette que les femmes enroulaient au niveau de leur poitrine pour la soutenir -- bref : un soutien-gorge --, en revanche, le lécythe a étonnamment égaré les commentateurs: ils lui ont donné son sens (courant) de \italique{fiole à huile}. D'aucuns ont même déployé force énergie pour tenter de montrer en quoi cet accessoire était typiquement masculin (!?)  Ils eussent plus vite et mieux fait d'ouvrir n'importe quel lexique\marqueur{2} : ils y auraient appris que λήκυθος peut aussi désigner la \italique{pomme d'Adam}. On comprend mieux l'effarement de Mnésiloque face à cette bandelette toute féminine entourant la poitrine désespérément plate d'Agathon, surmontée d'une proéminente et peu ragoûtante pomme d'Adam ; admettez que cela est d'assez mauvais goût (ὡς οὐ ξύμφορον)\marqueur{3}.
\end{multicols}
}%
\end{adjustbox}%
\vfill
\end{adjustwidth}
\newpage
% --- PAGE DE DROITE : Traduction française + Commentaire (4/8) ---
\begin{adjustwidth}{-2cm}{-2cm}
\noindent\textbf{\footnotesize Traduction française --- Folio 04 (\texttt{f. 164r}, vv. 123–164 [2/4]) \hfill ΠΡΟΛΟΓΟΣ}\\[-2pt]
\noindent\rule{16.5cm}{0.5pt}\\[2pt]
{\fontsize{8.8pt}{10.2pt}\selectfont
\ligneEst{}{\francais{et les baisers lascifs. En l'écoutant}}{}
\ligneEst{}{\francais{j'ai ressenti un prurit jusque sous le fondement.}}{}
\ligneEst{}{\francais{Jouvenceau, si tu en es un, à la façon d'Eschyle}}{}
\ligneEst{}{\francais{dans sa \italique{Lycurgie}, je veux t'interroger.}}{135}
\ligneEst{}{\francais{D'où es-tu, ô femmelette ? Quelle est ta patrie ? Qu'est cet accoutrement,}}{}
\ligneEst{}{\francais{cette incohérence dans tes façons ? Un barbitos,}}{}
\ligneEst{}{\francais{qu'a-t-il à dire à une crocote, une lyre à une résille ?}}{}
\ligneEst{}{\francais{Quoi ! une pomme d'Adam avec un strophion ? Quelle  disconvenance !}}{}
\ligneEst{}{\francais{Le rapport d'un miroir à une épée ?}}{140}
\ligneEst{}{\francais{Et toi-même, ô jeunesse, qui es-tu ? Un homme ?}}{}
\ligneEst{}{\francais{Mais où est ton cas ? où ton manteau ? où tes laconiennes ?}}{}
}
\vspace{2pt}
\noindent\rule{16.5cm}{0.6pt}\\[2pt]
\noindent\textbf{\small Commentaire philologique, métrique et scénique (vv. 123–164 --- partie 4/8)}\\[2pt]
\begin{adjustbox}{max totalheight=15.2cm, max width=16.5cm, keepaspectratio}%
\parbox[t]{16.5cm}{
\setlength{\columnsep}{9pt}%
\begin{multicols}{2}
\fontsize{7.6pt}{8.8pt}\selectfont
\setlength{\parindent}{0pt}%
\setlength{\parskip}{1.2pt}%
\noindent\textbf{140} \textbf{κατόπτρου.} La variante κάτροπτον (avec métathèse du ρ) est la seule attestée par l'épigraphie pour les inscriptions contemporaines de la représentation des \italique{Thesmophories}\marqueur{1}. Ce n'est qu'à partir de la fin du \nieme{IV} siècle que la forme κάτοπτρον, plus conforme à l'étymologie, semble prendre le dessus.
\noindent\textbf{140} \textbf{τίς δαὶ κατόπτρου καὶ ξίφους κοινωνία;} Le sel de ce vers ne se peut goûter qu'à condition d'avoir en tête le proverbe grec\marqueur{1} :
\citComm{τίς γὰρ κατόπτρῳ καὶ τυφλῷ κοινωνία;\\
Qu'y a-t-il de commun entre un aveugle et un miroir ?}
\noindent Ce proverbe a été répertorié par les parémiographes\marqueur{2} dans sa version concise (κάτοπτρον ὁ τυφλός) avec d'autres associations tout aussi inutiles. Ainsi : φαλακρὸς κτένα , \italique{un chauve un peigne}.
\noindent L'on trouvait bien évidemment des proverbes équivalents dans notre vieux parler ; \italique{À l'aveugle ne duit} [$=$ convient]\italique{ peinture / Couleur, miroir ne figure}\marqueur{3}.
\noindent Aristophane joue par paronomase entre ξίφους et τυφλοῦ.
\noindent Au vers 138, il devait sans doute se livrer à un détournement du même genre : on recense de nombreux proverbes mettant en relation une lyre avec un âne (ex. : τί γὰρ κοινὸν ... λύρᾳ καὶ ὄνῳ, voir Tosi n° 1934), animal lourd et rétif aux muses ; ou une lyre avec un choucas (\latin{Nihil graculo cum fidibus}, \italique{Adages} d'Érasme n° 337), oiseau criard.
\noindent\textbf{[Notes ad 139]} \textbf{(1)} Je renvoie ici aux articles correspondants de la monumentale somme qu'est le \italique{Dictionnaire des antiquités grecques et romaines} de Daremberg et Saglio
\noindent \textbf{(2)} Les dictionnaires classiques : Bailly, \LSJ, \&c. Hésychius (λ 857): τὸ μεταξὺ τοῦ λαυκανίου καὶ αὐχένος ήχῶδες. Cf. le verbe ληκυθίζειν = λαρυγγίζειν = φαρυγγίζειν, \italique{donner de la voix}.
\noindent \textbf{(3)} Faut-il en conclure que la présence d'une pomme d'Adam était caractéristique du masque comique représentant Agathon ?
\noindent \textbf{(1)} Voir par exemple les \italique{Inscriptiones Graecae} XXII.1471.47, 1544.58.
\noindent Le lecteur pourra aisément vérifier ces dires par lui-même en papillonnant sur les \italique{Inscriptiones Graecae} à l'adresse suivante : \href{https://epigraphy.packhum.org/}{https://epigraphy.packhum.org/}
\noindent \textbf{(1)} \nieme{3} vers d'un fragment comique adespote cité par Stobée (IV, 30, 6a) ; c'est le \nieme{903} fragment des \italique{Poetae Comici Graeci, Adespota}, VIII (1995).
\noindent \textbf{(2)} \italique{Corpus paroemiographorum Graecorum}, II, p. 774 (\italique{Mantissae Proverbiorum} 100).
\end{multicols}
}%
\end{adjustbox}%
\vfill
\end{adjustwidth}
\newpage

% --- FOLIO 04 (f. 164r) : SPREAD 2.3 (Texte grec établi + Traduction française + Commentaire 5/8 & 6/8) ---
\clearpage\ifodd\value{page}\hbox{}\newpage\fi
\enTeteG{ΘΕΣΜΟΦΟΡΙΑΖΟΥΣΑΙ --- Texte grec établi (f. 164r, vv. 123–164 [3/4])}
\enTeteD{LES THESMOPHORIES --- Traduction française (f. 164r, vv. 123–164 [3/4])}
\begin{adjustwidth}{-2cm}{-2cm}
\noindent\textbf{\footnotesize Texte grec établi --- Folio 04 (\texttt{f. 164r}, vv. 123–164 [3/4]) \hfill ΠΡΟΛΟΓΟΣ}\\[-2pt]
\noindent\rule{16.5cm}{0.5pt}\\[2pt]
{\fontsize{8.8pt}{10.2pt}\selectfont
\ligneEst{}{ἀλλ᾽ ὡς γυνὴ δῆτ᾽ εἶτα ποῦ τὰ τιτθία;}{}
\ligneEst{}{τί φῄς; τί σιγᾷς; ἀλλὰ δῆτ᾽ ἐκ τοῦ μέλους}{}
\ligneEst{}{ζητῶ σ᾽, ἐπειδή γ᾽ αὐτὸς οὐ βούλει φράσαι;}{145}
\ligneEst{\textbf{Ἀγάθων}}{ὦ πρέσβυ πρέσβυ, τοῦ φθόνου μὲν τὸν ψόγον}{}
\ligneEst{}{ἤκουσα, τὴν δ᾽ ἄλγησιν οὐ παρεσχόμην}{}
\ligneEst{}{ἐγὼ δὲ τὴν ἐσθῆθ᾽ ἅμα τῇ γνώμῃ φορῶ.}{}
\ligneEst{}{χρὴ γὰρ ποιητὴν ἄνδρα πρὸς τὰ δράματα}{}
\ligneEst{}{ἃ δεῖ ποιεῖν πρὸς ταῦτα τοὺς τρόπους ἔχειν.}{150}
\ligneEst{}{αὐτίκα γυναικεῖ᾽ ἢν ποιῇ τις δράματα,}{}
\ligneEst{}{μετουσίαν δεῖ τῶν τρόπων τὸ σῶμ᾽ ἔχειν.}{}
\ligneEst{\textbf{Μνησίλοχος}}{οὐκοῦν κελητίζεις, ὅταν Φαίδραν ποιῇς;}{}
}
\vspace{2pt}
\noindent\rule{16.5cm}{0.6pt}\\[2pt]
\noindent\textbf{\small Commentaire philologique, métrique et scénique (vv. 123–164 --- partie 5/8)}\\[2pt]
\begin{adjustbox}{max totalheight=15.2cm, max width=16.5cm, keepaspectratio}%
\parbox[t]{16.5cm}{
\setlength{\columnsep}{9pt}%
\begin{multicols}{2}
\fontsize{7.6pt}{8.8pt}\selectfont
\setlength{\parindent}{0pt}%
\setlength{\parskip}{1.2pt}%
\noindent \textbf{(3)} \italique{Le livre des proverbes français} de Le Roux de Lincy, Série n°V, à l'entrée \italique{Aveugle}. À l'initiative du génial Pierre Boutang, le livre a été réédité en 1996 chez Hachette, augmenté d'une sienne préface dans laquelle -- entre autres merveilles -- Boutang questionne avec bonheur l'étymologie du \italique{pro-verbe}.
\noindent\textbf{153} \textbf{κελητίζεις, ὅταν Φαίδραν ποιῇς.} Allusion à la \italique{pendula Venus} des Latins. Cf. \italique{Guêpes}, 500 ; \italique{Lysistrata}, 60 et 773. Le poète Machon disait καθιππάζειν (Athénée, XIII, p. 581d). (Note d'Alphonse Willems. Phèdre était éprise d'Hippolyte, dont le nom signifie littéralement le «~délieur de chevaux~» ; c'est-à-dire celui dont les chevaux courent à bride abattue.)
\noindent\textbf{157} \textbf{ὅταν σατύρους τοίνυν ποιῇς.} Lors des Grandes Dionysies, chaque poète tragique en lice devait présenter une tétralogie composée de trois tragédies et d'un drame satyrique.
\noindent\textbf{158} \textbf{συμποιῶ σοὔπισθεν ἐστυκὼς} ; ἐστυκὼς de στύω, \italique{bander}, alors qu'on attendait ἑστηκώς, \italique{se tenir debout} ; σοὔπισθεν est mis ici pour σοί ὄπισθέν (le datif σοί réclamé par le συμ- de συμποιῶ).
\noindent\textbf{161-3} \textbf{Ἴβυκος ἐκεῖνος κἀνακρέων ὁ Τήιος / κἀλκαῖος, οἳ περὶ ἁρμονίαν ἐχύμησαν / ἐμιτροφόρουν τε καὶ διεκίνων Ἰωνικῶς,} Les trois poètes lyriques Ibycos de Rhégium, Anacréon de Téos et Alcée de Lesbos sont aussi rapprochés par le scholiaste de Pindare (\italique{Isthmiques}, 2.1b\marqueur{1}) pour d'autres raisons que d'ordre harmonique : ce sont trois auteurs pédérastiques ; Cicéron les regroupe aussi sur le critère de la passion amoureuse (\italique{Tusculanes}, IV, 33) :
\citComm{\latin{quid denique homines doctissimi et summi poëtae de se ipsis et carminibus edant et cantibus ? Fortis vir in sua re publica cognitus quae de iuvenum amore scribit Alcaeus ! nam Anacreontis quidem tota poësis est amatoria. Maxume vero omnium flagrasse amore Reginum Ibycum apparet ex scriptis.}\\
\francais{Avec quelle licence de très savants hommes et de grands poètes n'ont-ils pas chanté leurs galanteries ? Alcée, ce fameux guerrier de Mitylène, que n'a-t-il pas écrit de ces inclinations à la Grecque ? Anacréon respire-t-il autre chose que l'amour ? On voit la passion encore poussée bien plus loin dans les poésies d'Ibycus.}}
\noindent Le spectateur du \nieme{V} siècle ne s'y trompait pas et savait à quoi s'en tenir avec ces trois bougres-là et comment interpréter correctement les euphémismes agathoniens.
\noindent Le commentaire ancien le plus pertinent sur notre passage n'est pas écrit en marge du \italique{Ravennas} -- comme c'est hélas bien souvent le cas pour les \italique{Thesmophories} -- mais se trouve dans la \italique{Souda} (ε 989 -- et quasiment à l'identique dans le lexique de Zonaras, p. 707) :
\citComm{ἀντὶ τοῦ μίτραις\marqueur{2} τουτέστι ζώναις\marqueur{2} διαζωσάμενοι ἁβρῶς περιεπάτουν.\\
Signifie : ceints de mitres -- c'est-à-dire de ceintures -- ils vivaient de façon luxurieuse.}
\end{multicols}
}%
\end{adjustbox}%
\vfill
\end{adjustwidth}
\newpage
% --- PAGE DE DROITE : Traduction française + Commentaire (6/8) ---
\begin{adjustwidth}{-2cm}{-2cm}
\noindent\textbf{\footnotesize Traduction française --- Folio 04 (\texttt{f. 164r}, vv. 123–164 [3/4]) \hfill ΠΡΟΛΟΓΟΣ}\\[-2pt]
\noindent\rule{16.5cm}{0.5pt}\\[2pt]
{\fontsize{8.8pt}{10.2pt}\selectfont
\ligneEst{}{\francais{Alors tu es femme ? Mais où sont tes seins ?}}{}
\ligneEst{}{\francais{Que dis-tu ? Pourquoi te taire ? Alors d'après ton chant,}}{}
\ligneEst{}{\francais{je te déchiffre puisque tu ne veux pas t'expliquer toi-même.}}{145}
\ligneEst{\textbf{Agathon}}{\francais{Vieillard, vieillard, ces reproches dictés par l'envie,}}{}
\ligneEst{}{\francais{je les ai entendus, mais n'en ai point senti l'atteinte.}}{}
\ligneEst{}{\francais{Je porte la mise congruente à mes pensées.}}{}
\ligneEst{}{\francais{Car c'est une loi que le poète, suivant les drames}}{}
\ligneEst{}{\francais{qu'il lui faut faire, modèle sur eux ses mœurs.}}{150}
\ligneEst{}{\francais{Et par exemple, si l'on compose des drames féminins,}}{}
\ligneEst{}{\francais{il faut que le corps participe aux mœurs de la femme..}}{}
\ligneEst{\textbf{Mnésiloque}}{\francais{Ainsi tu chevauches quand tu composes une Phèdre ?}}{}
}
\vspace{2pt}
\noindent\rule{16.5cm}{0.6pt}\\[2pt]
\noindent\textbf{\small Commentaire philologique, métrique et scénique (vv. 123–164 --- partie 6/8)}\\[2pt]
\begin{adjustbox}{max totalheight=15.2cm, max width=16.5cm, keepaspectratio}%
\parbox[t]{16.5cm}{
\setlength{\columnsep}{9pt}%
\begin{multicols}{2}
\fontsize{8.4pt}{9.7pt}\selectfont
\setlength{\parindent}{0pt}%
\setlength{\parskip}{1.2pt}%
\noindent Ce commentaire est doublement intéressant : d'abord, il précise que la μίτρα dont il est question est à prendre au sens de \italique{ceinture} et non de bandeau. Ensuite parce qu'il glose trois des mots du vers 163 : ἐμιτροφόρουν par μίτραις διαζωσάμενοι ; Ἰωνικῶς par ἁβρῶς\marqueur{3} ; et -- plus intéressant -- le \italique{contra metrum} †διεκίνων† (†διεκίνουν† dans la \italique{Souda}) par περιεπάτουν (ici avec le sens figuré de \italique{se conduire}, \italique{vivre de telle} ou \italique{telle manière}).
\noindent Cette glose est donc incompatible avec l'émendation de Toup, très largement reprise dans les éditions modernes. Toup proposa la correction suivante : διεκλῶντ' (de διακλάω , \italique{briser en deux}, d'où \italique{briser une vie par la mollesse} (?!)).
\noindent Il est bien préférable -- tant du point de vue paléographique que du sens -- de lire διέκειντ', de διάκειμαι, \italique{être disposé}, synonyme de περιπατέω, et qui se construit fréquemment avec un adverbe (ici : διάκειμαι Ἰωνικῶς, \italique{avoir des dispositions à la vie ionienne}, c'est-à-dire à la mollesse et la luxure\marqueur{3}).\\ Les sous-entendus d'Agathon sont à peine voilés : les trois poètes pédérastes portent la ceinture à la mode ionienne, c'est-à-dire lâche, voire franchement dénouée...
\noindent La métaphore est la suivante : tels qu'ils portaient la ceinture, béante, telles étaient leurs \italique{harmonies}. Le vers 162, οἳ περὶ ἁρμονίαν †ἐχύμησαν†\marqueur{4}, se laisse maintenant comprendre. Il faut lire :
\citComm{οἳ περὶ ἁρμονίαν ἐκεχήνεσαν}
\noindent La correction est très légère : une nouvelle faute de iotacisme (entre η et υ), de confusion μ/ν (quasiment identiques sur notre \italique{Ravennas}, comme sur beaucoup de manuscrits écrits en minuscules) et d'une haplographie bien explicable :  $\text{εκεχ} \xrightarrow{\text{saut du même au m. }} \text{ εχ}$.\\ χάσκειν signifie littéralement \italique{avoir la bouche grande ouverte}\marqueur{5}. En bâillant de la sorte aux corneilles, l'on pouvait émettre de légers sons, comparables aux coquilles grillant sur des chardons\marqueur{6} ou à cette vieille lyre de Cratinos, le rival hors-d'âge d'Aristophane, chevrotant ses derniers et mauvais accords (\italique{Cavaliers}, v. 533)\marqueur{7} :
\citComm{τῶν θ' ἁρμονιῶν διαχασκουσῶν.}
\end{multicols}
}%
\end{adjustbox}%
\vfill
\end{adjustwidth}
\newpage

% --- FOLIO 04 (f. 164r) : SPREAD 2.4 (Texte grec établi + Traduction française + Commentaire 7/8 & 8/8) ---
\clearpage\ifodd\value{page}\hbox{}\newpage\fi
\enTeteG{ΘΕΣΜΟΦΟΡΙΑΖΟΥΣΑΙ --- Texte grec établi (f. 164r, vv. 123–164 [4/4])}
\enTeteD{LES THESMOPHORIES --- Traduction française (f. 164r, vv. 123–164 [4/4])}
\begin{adjustwidth}{-2cm}{-2cm}
\noindent\textbf{\footnotesize Texte grec établi --- Folio 04 (\texttt{f. 164r}, vv. 123–164 [4/4]) \hfill ΠΡΟΛΟΓΟΣ}\\[-2pt]
\noindent\rule{16.5cm}{0.5pt}\\[2pt]
{\fontsize{8.8pt}{10.2pt}\selectfont
\ligneEst{\textbf{Ἀγάθων}}{ἀνδρεῖα δ᾽ ἄν ποιῇ τις, ἐν τῷ σώματι}{}
\ligneEst{}{ἔνεσθ᾽ ὑπάρχον τοῦθ᾽. ἃ δ᾽ οὐ κεκτήμεθα,}{155}
\ligneEst{}{μίμησις ἤδη ταῦτα συνθηρεύεται.}{}
\ligneEst{\textbf{Μνησίλοχος}}{ὅταν σατύρους τοίνυν ποιῇς, καλεῖν ἐμέ,}{}
\ligneEst{}{ἵνα συμποιῶ σοὔπισθεν ἐστυκὼς ἐγώ.}{}
\ligneEst{\textbf{Ἀγάθων}}{ἄλλως τ᾽ ἄμουσόν ἐστι ποιητὴν ἰδεῖν}{}
\ligneEst{}{ἀγρεῖον ὄντα καὶ δασύν σκέψαι δ᾽ ὅτι}{160}
\ligneEst{}{Ἴβυκος ἐκεῖνος κἀνακρέων ὁ Τήιος}{}
\ligneEst{}{κἀλκαῖος, οἳ περὶ ἁρμονίαν ἐκεχήνεσαν,}{}
\ligneEst{}{ἐμιτροφόρουν τε καὶ διέκειντ' Ἰωνικῶς,}{}
\ligneEst{}{καὶ Φρύνιχος, τοῦτον γὰρ οὖν ἀκήκοας,}{164}
}
\vspace{2pt}
\noindent\rule{16.5cm}{0.6pt}\\[2pt]
\noindent\textbf{\small Commentaire philologique, métrique et scénique (vv. 123–164 --- partie 7/8)}\\[2pt]
\begin{adjustbox}{max totalheight=15.2cm, max width=16.5cm, keepaspectratio}%
\parbox[t]{16.5cm}{
\setlength{\columnsep}{9pt}%
\begin{multicols}{2}
\fontsize{7.6pt}{8.8pt}\selectfont
\setlength{\parindent}{0pt}%
\setlength{\parskip}{1.2pt}%
\noindent (Noter la collocation entre (δια-)χάσκω et ἁρμονία.)  C'est la même expression qu'on retrouve dans la bouche d'Agathon : περὶ ἁρμονίαν ἐκεχήνεσαν, « béer une harmonie », c'est-à-dire bêler -- ou chevroter, telle est l'image en français, point très éloignée du grec -- quelques paroles en musique\marqueur{8}. Le sel de la métaphore réside dans le double sens de περιχάσκειν. Le verbe est justement attesté au sens d'\italique{ouvrir une ceinture}\marqueur{9}.\\ Il est remarquable de noter que la tradition millénaire du texte d'Aristophane ait conservé la leçon ἐχύμησαν avec un η jusqu'au scribe du \italique{Ravennas} ; et ce en dépit de manuscrits anciens fautifs portant ἐχύμισαν et de commentaires afférents (ἔγχυμον ἐποίησαν) -- ayant pénétré jusqu'aux marges du \italique{Ravennas} -- suggérant la présence d'un verbe *χυμίζω. Preuve encore une fois que l'œuvre conservée d'Aristophane est, comme l'avait déjà noté Willems, l'une des plus fidèlement transmises avec l'\italique{Iliade} et l'\italique{Odyssée}\marqueur{10}.
\noindent\textbf{164} \textbf{Φρύνιχος, τοῦτον γὰρ οὖν ἀκήκοας.} Le poète Phrynicos, qui vivait au temps des guerres médiques portait naturellement le riche costume de son temps. Cf. les \italique{Cavaliers}, v. 1331-1332 :
\citComm{Ἀλλαντοπώλης — ὅδ᾽ ἐκεῖνος ὁρᾶν τεττιγοφόρας, ἀρχαίῳ σχήματι λαμπρός,\\
οὐ χοιρινῶν ὄζων ἀλλὰ σπονδῶν, σμύρνῃ κατάλειπτος.\\
\francais{Le charcutier — Le voilà qui apparait, la cigale à la chevelure, dans tout l'éclat du costume antique, fleurant, non les coquilles mais les libations de paix, et frotté avec de la myrrhe.}}
\noindent Et Thucydide (I, 6) :
\citComm{καὶ οἱ πρεσβύτεροι αὐτοῖς τῶν εὐδαιμόνων \{διὰ τὸ ἁβροδίαιτον\} οὐ πολὺς χρόνος ἐπειδὴ χιτῶνάς τε λινοῦς ἐπαύσαντο φοροῦντες καὶ χρυσῶν τεττίγων ἐνέρσει κρωβύλον ἀναδούμενοι τῶν ἐν τῇ κεφαλῇ τριχῶν.\\
Les vieillards athéniens de la classe aisée portaient encore, il n'y a pas longtemps, des tuniques (longues) de lin, et piquaient des cigales d'or dans leurs cheveux, ramassés et noués sur le sommet de la tête.}
\noindent (Note de Willems.)\\ Le scholiaste explique l'incise d'Agathon (τοῦτον γὰρ οὖν ἀκήκοας) ainsi : ἐπεὶ Ἀθηναῖος ἦν.
\noindent\textbf{[Notes ad 161-3]} \textbf{(1)} Au tout début de la célèbre Isthmique où Pindare affirme qu'autrefois la Muse ne se faisait pas mercenaire mais chantait la beauté. Le scholiaste de noter que cela s'applique aux trois poètes lyriques : ταῦτα δὲ τείνει καὶ εἰς τοὺς περὶ Ἀλκαῖον καὶ Ἴβυκον καὶ Ἀνακρέοντα... ainsi qu'à tous les poètes apparemment fort portés sur la pédérastie : καὶ εἴ τινες τῶν πρὸ αὐτοῦ δοκοῦσι περὶ τὰ παιδικὰ ἠσχολῆσθαι· οὗτοι γὰρ παλαιότεροι Πινδάρου. Et d'appuyer le trait avec cette anecdote : on interrogeait Anacréon sur la raison qu'il avait d'écrire des hymnes à de jeunes adolescents plutôt qu'aux dieux. Anacréon répondit que c'était justement ces derniers leurs dieux : Ἀνακρέοντα γοῦν ἐρωτηθέντα, φασί, διατί οὐκ εἰς θεοὺς ἀλλ εἰς παῖδας γράφεις τοὺς ὕμνους; εἰπεῖν, ὅτι οὗτοι ἡμῶν θεοί εἰσιν.
\end{multicols}
}%
\end{adjustbox}%
\vfill
\end{adjustwidth}
\newpage
% --- PAGE DE DROITE : Traduction française + Commentaire (8/8) ---
\begin{adjustwidth}{-2cm}{-2cm}
\noindent\textbf{\footnotesize Traduction française --- Folio 04 (\texttt{f. 164r}, vv. 123–164 [4/4]) \hfill ΠΡΟΛΟΓΟΣ}\\[-2pt]
\noindent\rule{16.5cm}{0.5pt}\\[2pt]
{\fontsize{8.8pt}{10.2pt}\selectfont
\ligneEst{\textbf{Agathon}}{\francais{Compose-t-on des drames mâles, dans le corps}}{}
\ligneEst{}{\francais{réside ce qui est requis. Mais ce que nous ne possédons pas,}}{155}
\ligneEst{}{\francais{l'imitation est là qui l'attrape.}}{}
\ligneEst{\textbf{Mnésiloque}}{\francais{Eh bien, quand tu feras des drames satyriques, appelle-moi,}}{}
\ligneEst{}{\francais{et je coopérerai avec toi, de derrière, avec raideur.}}{}
\ligneEst{\textbf{Agathon}}{\francais{D'ailleurs, c'est chose choquante à voir}}{}
\ligneEst{}{\francais{qu'un poète grossier et hirsute. Avise-toi que}}{160}
\ligneEst{}{\francais{le fameux Ibycos, et Anacréon de Téos}}{}
\ligneEst{}{\francais{et Alcée, accordaient leur voix,}}{}
\ligneEst{}{\francais{leur mitre et leur vit à l'ionienne. Et Phrynichos,}}{}
\ligneEst{}{\francais{— celui-là, en effet, tu l'as entendu, —}}{164}
}
\vspace{2pt}
\noindent\rule{16.5cm}{0.6pt}\\[2pt]
\noindent\textbf{\small Commentaire philologique, métrique et scénique (vv. 123–164 --- partie 8/8)}\\[2pt]
\begin{adjustbox}{max totalheight=15.2cm, max width=16.5cm, keepaspectratio}%
\parbox[t]{16.5cm}{
\setlength{\columnsep}{9pt}%
\begin{multicols}{2}
\fontsize{7.6pt}{8.8pt}\selectfont
\setlength{\parindent}{0pt}%
\setlength{\parskip}{1.2pt}%
\noindent \textbf{(2)} Zonaras -- ou le \italique{pseudo}-Zonaras pour les esprits chagrins -- comporte un accusatif en place de l'ablatif : ἀντὶ τοῦ μίτρας ἤτοι ζώνας διαζωσάμενοι.
\noindent \textbf{(3)} Voir Érasme, \italique{Adages} (n° 469, \latin{Risus Ionicus}).
\noindent \textbf{(4)} La rêverie du scholiaste reprise par Suidas, ἐχύμισαν· ἀντὶ τοῦ ἔγχυμον ἐποίησαν, n'est évidemment pas à prendre au sérieux. Pas plus que l'invention du verbe χυμίζω... à moins qu'on ne m'explique à quelle sauce (χυμός) tourner le passage pour qu'il ait un rapport avec les inflexions de voix (\italique{harmonies}). Ces divagations auront probablement troublé davantage encore les scribes, déjà en butte à un texte corrompu.
\noindent \textbf{(5)} Bien souvent : \italique{être béat d'admiration}. Voir Taillardat §§472 et 578.
\noindent \textbf{(6)} ἀνέχασκον εἷς ἕκαστος ἐμφερέστατα ὀπτωμέναις κόγχαισιν ἐπὶ τῶν ἀνθράκων. Aristophane, \italique{fr.} 67 K.-A.
\noindent \textbf{(7)} C'est bien ainsi qu'il faut interpréter τῶν θ' ἁρμονιῶν διαχασκουσῶν : « émettre des sons, ou accords » et non « aux joints béants » (en parlant de la lyre), comme on peut le lire.
\noindent \textbf{(8)} ἁρμονία a le sens d'arrangement de paroles pour un morceau de musique dans Platon, \italique{Theæt.} 175e-176a.
\noindent \textbf{(9)} Voir LSJ, I. b. Attesté par Héliodore le médecin (I\italique{er} ou \nieme{II} s. ap. J.-C). De longs fragments nous en ont été conservés par Oribase (\italique{Collectiones medicae} XLVIII, 58.5) : cf. les \italique{Œuvres d'Oribase}, t. 4, p. 320 (1862), de Bussemaker et Daremberg. À propos d'un patient à qui l'on applique un bandage, appelé ceinture : τῇ ζώνῃ \textbf{περιχάσκων} ζώννυται ὁ πάσχων, καὶ τότε ὁ εὐθὺς τελαμὼν ἀπὸ τῆς ὀσφύος ἄγεται ἐπὶ τὴν ἕδραν καὶ τὸν περίναιον, « On entoure le malade avec la ceinture, en l'\italique{ouvrant} largement ; puis on amène la bande droite des lombes au siège et au périnée » (trad. Daremberg).
\noindent \textbf{(10)} Voir l'introduction à ses premières notes sur les \italique{Cavaliers}. L'une des raisons en est que le Comique, au contraire des Tragiques, ne fut que peu rejoué. Les premières éditions critiques de ses œuvres furent sans doute beaucoup plus aisées à constituer que celles d'un Euripide, dont beaucoup plus de versions différentes des pièces circulaient. Victor Coulon, l'éditeur aux Belles Lettres d'Aristophane, abonda dans le sens de Willems : à l’opposé du courant majoritaire de son temps, il s'est astreint à conserver le plus possible les leçons des manuscrits contre les conjectures des philologues. Ce « conservatisme », qu'il défend pour le cas d'Aristophane dans son \italique{Essai sur la méthode de la critique conjecturale appliquée au texte d’Aristophane} (1933) lui valut d’être encore considéré à la fin des années 90 par Sommerstein comme le meilleur éditeur d’Aristophane.
\end{multicols}
}%
\end{adjustbox}%
\vfill
\end{adjustwidth}
\newpage
\clearpage\ifodd\value{page}\hbox{}\newpage\fi
\enTeteG{ΠΡΟΛΟΓΟΣ --- Codex Ravennas 429, f. 164v (Classense 166v) --- vv. 165–207}
\begin{adjustwidth}{-1.5cm}{-3.0cm}
\begin{adjustbox}{max totalheight=21.0cm, max width=16.5cm, keepaspectratio}
\begin{minipage}{16.5cm}
{\fontsize{7.2pt}{8.5pt}\selectfont
\noindent\textbf{Folio 05 --- Codex Ravennas 429, \texttt{f. 164v} (\texttt{Classense 166v}) : ΠΡΟΛΟΓΟΣ, vv. 165–207}\hfill \textit{Témoins :} $\text{R}^1, \Rpc, \Rdeux$ (\Mon)\par\vspace{2pt}
\hrule height 0.4pt\vspace{2.5pt}
}
{\fontsize{8.1pt}{9.5pt}\selectfont
\ligneRav{165}{ἀγάθ/}{αὐτός τε καλὸς ἦν καὶ καλῶς ἠμπέσχετο}{\textbf{165} \lecon{ἠμπέσχετο} R : ἠμπισχετο Elmsley ad Eur. \italique{Med.} 1128 (\italique{induebatur} Divus)}
\ligneRav{}{}{διὰ τοῦτ᾽ ἄρ᾽ αὐτοῦ καὶ κάλ᾽ ἦν τὰ δράματα·}{}
\ligneRav{}{}{ὅμοια γὰρ ποιεῖν ἀνάγκη τῇ φύσει·}{}
\ligneRav{168}{μ\auDessus{ν}{Η}}{ταῦτ᾽ ἄρ᾽ ὁ \difference{Φιλοκλέης} αἰσχρὸς ὢν αἰσχρῶς ποιεῖ,}{\textbf{168-9} Φιλοκλῆς, Ξενοκλῆς R : \lecon{Φιλοκλέης}, \lecon{Ξενοκλέης} Bentley \espApp \lecon{ὁ δ᾽ αὖ} R : ὁ δὲ Brunck}
\ligneRav{}{}{ὁ δ᾽ αὖ Ξενοκλέης ὢν κακὸς κακῶς ποιεῖ,}{}
\ligneRav{170}{}{ὁ δ᾽ αὖ Θέογνις ψυχρὸς ὢν ψυχρῶς ποιεῖ·}{}
\ligneRav{}{ἀγάθ/}{ἅπασ᾽ ἀνάγκη ταῦτα γάρ τοι γνοὺς ἐγὼ}{}
\ligneRav{}{—}{ἐμαυτὸν ἐθεράπευσα. : πῶς πρὸς τῶν θεῶν;}{}
\ligneRav{173}{\auDessus{εὐ}{ρ}}{παῦσαι βαΰζων καὶ γὰρ ἐγὼ τοιοῦτος ἦν}{\textbf{173} ἦν R, Suid. β 192 : \lecon{ἦ} Dindorf (1825)}
\ligneRav{}{}{ὢν τηλικοῦτος, ἡνίκ᾽ ἠρχόμην ποιεῖν·}{}
\ligneRav{175}{μ\auDessus{ν}{Η}}{μὰ τὸν Δί᾽ οὐ ζηλῶ σε τῆς παιδεύσεως·}{}
\ligneRav{176}{\auDessus{εὐ}{ρ}}{ἀλλ᾽ ὧνπερ οὕνεκ᾽ ἦλθον, ἔα μ᾽ εἰπεῖν. : λέγε·}{\textbf{176} \lecon{(Ἀγάθ.)} λέγε R, Fritzsche : (Μνη.) Küster}
\ligneRav{}{\auDessus{εὐ}{ρ}}{Ἀγάθων, σοφοῦ πρὸς ἀνδρός, ὅστις ἐν βραχεῖ}{}
\ligneRav{178}{}{πολλοὺς καλῶς \difference{οἷός} τε συντέμνειν λόγους·}{\textbf{178} οἷόν R : \lecon{οἷός} Suid. σ 792, οι 154, Stob. III,35.3 \espApp \lecon{συντέμνειν} R, Suid. οι 154, Stob. : συντεμεῖν Suid. σ 792 \espApp κοινῇ R, Suid. κ 2556, σ 792 : \lecon{καινῇ} Biset ap. Σ}
\ligneRav{179}{}{ἐγὼ δὲ \difference{καινῇ} ξυμφορᾷ πεπληγμένος}{\textbf{179} κοινῇ R, Souda $\kappa\ 2556, \sigma\ 792$ (confusion minuscule $\alpha \to \text{ο}$) : \lecon{καινῇ} Biset d'après $\Sigma^{\text{R}}$}
\ligneRav{180}{—}{ἱκέτης ἀφῖγμαι πρὸς σέ. : τοῦ χρείαν ἔχων;}{}
\ligneRav{}{\auDessus{εὐ}{ρ}}{μέλλουσί μ᾽ αἱ γυναῖκες ἀπολεῖν τήμερον}{}
\ligneRav{}{}{τοῖς Θεσμοφορίοις, ὅτι κακῶς αὐτὰς λέγω·}{}
\ligneRav{}{ἀγάθ/}{τίς οὖν παρ᾽ ἡμῶν ἐστιν ὠφέλειά σοι;}{}
\ligneRav{184}{\auDessus{εὐ}{ρ}}{\difference{ἡ πᾶσ᾽·} ἐὰν γὰρ ἐγκαθεζόμενος λάθρᾳ}{\textbf{184} ἡ πᾶσα R : \lecon{ἡ πᾶσ᾽·} Küster : ἅπασ' Scaliger}
\ligneRav{185}{}{ἐν ταῖς γυναιξίν, ὡς δοκῶν εἶναι γυνή,}{}
\ligneRav{}{}{ὑπεραποκρίνῃ μου, σαφῶς σώσεις ἐμέ·}{}
\ligneRav{}{}{μόνος γὰρ ἂν λέξειας ἀξίως ἐμοῦ·}{}
\ligneRav{}{ἀγάθ/}{ἔπειτα πῶς οὐκ αὐτὸς ἀπολογεῖ παρών;}{}
\ligneRav{}{\auDessus{εὐ}{ρ}}{ἐγὼ φράσω σοι. πρῶτα μὲν γιγνώσκομαι}{}
\ligneRav{190}{}{ἔπειτα πολιός εἰμι καὶ πώγων᾽ ἔχω,}{}
\ligneRav{}{}{σὺ δ᾽ εὐπρόσωπος λευκὸς ἐξυρημένος}{}
\ligneRav{}{}{γυναικόφωνος ἁπαλὸς εὐπρεπὴς ἰδεῖν·}{}
\ligneRav{}{ἀγάθ/}{Εὐριπίδη— : τί ἔστιν; : ἐποίησάς ποτε,}{}
\ligneRav{}{}{`χαίρεις ὁρῶν φῶς, πατέρα δ᾽ οὐ χαίρειν δοκεῖς;'}{}
\ligneRav{195}{\auDessus{εὐ}{ρ}}{ἔγωγε. : μή νυν ἐλπίσῃς τὸ σὸν κακὸν}{}
\ligneRav{196}{}{ἡμᾶς ὑφέξειν. καὶ \difference{γὰρ ἂν} μαινοίμεθ᾽ ἂν·}{\textbf{196} γὰρ R : \lecon{γὰρ ἂν} Suid. υ 709}
\ligneRav{}{}{ἀλλ᾽ αὐτὸς ὅ γε σόν ἐστιν οἰκείως φέρε·}{}
\ligneRav{}{}{τὰς συμφορὰς γὰρ οὐχὶ τοῖς τεχνάσμασιν}{}
\ligneRav{}{}{φέρειν δίκαιον ἀλλὰ τοῖς παθήμασιν·}{}
\ligneRav{200}{μ\auDessus{ν}{Η}}{καὶ μὴν σύ γ᾽ ὦ κατάπυγον εὐρύπρωκτος εἶ}{}
\ligneRav{}{}{οὐ τοῖς λόγοισιν ἀλλὰ τοῖς παθήμασιν·}{}
\ligneRav{}{\auDessus{εὐ}{ρ}}{τί δ᾽ ἔστιν ὅτι δέδοικας ἐλθεῖν αὐτόσε;}{}
\ligneRav{}{ἀγάθ/}{κάκιον ἀπολοίμην ἂν ἢ σύ. : πῶς; : ὅπως;}{}
\ligneRav{204}{}{δοκῶν γυναικῶν ἔργα νυκτερείσια}{\textbf{204} \lecon{νυκτερείσια} R : νυκτερήσια Bothe}
\ligneRav{205}{}{κλέπτειν ὑφαρπάζειν τε θήλειαν Κύπριν·}{}
\ligneRav{206}{μ\auDessus{ν}{Η}}{ἰδού γε κλέπτειν· νὴ Δία βινεῖσθαι μὲν οὖν·}{\textbf{206} \lecon{νὴ Δία} R : μà Δία Fritzsche}
\ligneRav{}{}{ἀτὰρ ἡ πρόφασίς γε νὴ Δί᾽ εἰκότως ἔχει·}{}
}
\vspace{3pt}\noindent\rule{4.5cm}{0.4pt}\par\vspace{1.5pt}
{\fontsize{7.0pt}{8.3pt}\selectfont
\noindent\textbf{Loci paralleli (vv. 165–207) ---} \textbf{168} (ὁ Φιλοκλέης αἰσχρὸς $=$ κορυδός) scholium \italique{Aves} 1295a.α \espTest \textbf{173} (βαύ̈ζων -- ἦν) Suidas β 192 $\simeq$ Zonaras 379 (sed γὰρ καὶ pro καὶ γὰρ) \espTest \textbf{175} (οὐ -- παιδεύσεως) Suidas ζ 63 \espTest \textbf{176} Scholium v. 500 (sed ἀλλ' εἰπεῖν βούλομαι pro ἔα μ' εἰπεῖν) \espTest \textbf{177-80} (σοφοῦ -- πρὸς σέ) Suidas σ 792 \espTest \textbf{177-8} (σοφοῦ -- λόγους) Suidas οι 154 \espTest \textbf{179} Suidas κ 2556 \espTest \textbf{192} (γυναικόφωνος) Pollux II 111 \espTest \textbf{193} (ἐποίησάς ποτε) \Eustathe 4.9 \espTest \textbf{195-7} (μή νυν -- φέρε) Suidas υ 709 (νῦν pro νυν et ἀλλ' omisso) \espTest \textbf{198-201} (τὰς συμφορὰς -- παθήμασιν) Suidas τ 440 (-σι pro -σιν) %(τεχνάσμασι et παθήμασι pro -σιν)\par
}
\vspace{2.5pt}\noindent\rule{4.5cm}{0.4pt}\par\vspace{1.5pt}
{\fontsize{6.9pt}{8.2pt}\selectfont
\noindent\textbf{Suite (en commentaire) de l'apparat critique \& notes paléographiques (\texttt{f. 164v}) ---} \textbf{Contexte :} Fin de la tirade d'Agathon (vv. 165–172), requête d'Euripide priant Agathon d'aller plaider sa cause aux Thesmophories et refus prudent d'Agathon (vv. 173–207). --- Page très régulière de 43 lignes. On relève les formes contractes non métriques de noms propres tragiques (Φιλοκλῆς, Ξενοκλῆς R aux vv. 168–169 pour -κλέης), la confusion minuscule $\alpha/\text{ο}$ dans κοινῇ (v. 179 R et Souda) pour καινῇ (Biset d'après $\Sigma^{\text{R}}$), et la graphie iotacisante νυκτερείσια (v. 204 R, maintenue par L. Hauseux). --- \textbf{v. 179 :} κοινῇ R, Souda $\kappa\ 2556, \sigma\ 792$ (confusion minuscule $\alpha \to \text{ο}$) : \lecon{καινῇ} Biset d'après $\Sigma^{\text{R}}$.\par
}
\end{minipage}
\end{adjustbox}
\end{adjustwidth}
\newpage
\enTeteD{Codex Ravennas 429 (Biblioteca Classense) --- Fac-similé du f. 164v (vv. 165–207)}
\begin{adjustwidth}{-1.5cm}{-3.0cm}
\begin{center}
\includegraphics[width=16.5cm,height=20.5cm,keepaspectratio]{/workspaces/Thesmophories/dossier_ravennas_thesmophories/folios_individuels_ff162v-177r/05_f164v_classense166v_vv0165-0207.jpg}\par
\vspace{2pt}
{\scriptsize\textbf{Codex Ravennas 429, f. 164v (Classense 166v)} --- ΠΡΟΛΟΓΟΣ, \textbf{vv. 165–207} ($\text{R}^1, \Rpc, \Rdeux$ (\Mon)).}
\end{center}
\end{adjustwidth}
\newpage
\clearpage\ifodd\value{page}\hbox{}\newpage\fi% --- FOLIO 05 (f. 164v) : SPREAD 2.1 (Texte grec établi + Traduction française + Commentaire 1/6 & 2/6) ---
\clearpage\ifodd\value{page}\hbox{}\newpage\fi
\enTeteG{ΘΕΣΜΟΦΟΡΙΑΖΟΥΣΑΙ --- Texte grec établi (f. 164v, vv. 165–207 [1/3])}
\enTeteD{LES THESMOPHORIES --- Traduction française (f. 164v, vv. 165–207 [1/3])}
\begin{adjustwidth}{-2cm}{-2cm}
\noindent\textbf{\footnotesize Texte grec établi --- Folio 05 (\texttt{f. 164v}, vv. 165–207 [1/3]) \hfill ΠΡΟΛΟΓΟΣ}\\[-2pt]
\noindent\rule{16.5cm}{0.5pt}\\[2pt]
{\fontsize{8.4pt}{9.7pt}\selectfont
\ligneEst{}{αὐτός τε καλὸς ἦν καὶ καλῶς ἠμπέσχετο}{165}
\ligneEst{}{διὰ τοῦτ᾽ ἄρ᾽ αὐτοῦ καὶ κάλ᾽ ἦν τὰ δράματα.}{}
\ligneEst{}{ὅμοια γὰρ ποιεῖν ἀνάγκη τῇ φύσει.}{}
\ligneEst{\textbf{Μνησίλοχος}}{ταῦτ᾽ ἄρ᾽ ὁ Φιλοκλέης αἰσχρὸς ὢν αἰσχρῶς ποιεῖ,}{}
\ligneEst{}{ὁ δ᾽ αὖ Ξενοκλέης ὢν κακὸς κακῶς ποιεῖ,}{}
\ligneEst{}{ὁ δ᾽ αὖ Θέογνις ψυχρὸς ὢν ψυχρῶς ποιεῖ.}{170}
\ligneEst{\textbf{Ἀγάθων}}{ἅπασ᾽ ἀνάγκη ταῦτα γάρ τοι γνοὺς ἐγὼ}{}
\ligneEst{}{ἐμαυτὸν ἐθεράπευσα.}{}
\ligneEst{\textbf{Μνησίλοχος}}{\decalage{ἐμαυτὸν ἐθεράπευσα.}πῶς πρὸς τῶν θεῶν;}{}
\ligneEst{\textbf{Εὐριπίδης}}{παῦσαι βαΰζων καὶ γὰρ ἐγὼ τοιοῦτος ἦν}{}
\ligneEst{}{ὢν τηλικοῦτος, ἡνίκ᾽ ἠρχόμην ποιεῖν.}{}
\ligneEst{\textbf{Μνησίλοχος}}{μὰ τὸν Δί᾽ οὐ ζηλῶ σε τῆς παιδεύσεως.}{175}
\ligneEst{\textbf{Εὐριπίδης}}{ἀλλ᾽ ὧνπερ οὕνεκ᾽ ἦλθον, ἔα μ᾽ εἰπεῖν.}{}
\ligneEst{\textbf{Ἀγάθων}}{\decalage{ἀλλ᾽ ὧνπερ οὕνεκ᾽ ἦλθον, ἔα μ᾽ εἰπεῖν.}λέγε.}{}
\ligneEst{\textbf{Εὐριπίδης}}{Ἀγάθων, σοφοῦ πρὸς ἀνδρός, ὅστις ἐν βραχεῖ}{}
\ligneEst{}{πολλοὺς καλῶς οἷός τε συντέμνειν λόγους.}{}
\ligneEst{}{ἐγὼ δὲ καινῇ ξυμφορᾷ πεπληγμένος}{}
}
\vspace{2pt}
\noindent\rule{16.5cm}{0.6pt}\\[2pt]
\noindent\textbf{\small Commentaire philologique, métrique et scénique (vv. 165–207 --- partie 1/6)}\\[2pt]
\begin{adjustbox}{max totalheight=13.0cm, max width=16.5cm, keepaspectratio}%
\parbox[t]{16.5cm}{
\setlength{\columnsep}{9pt}%
\begin{multicols}{2}
\fontsize{8.4pt}{9.7pt}\selectfont
\setlength{\parindent}{0pt}%
\setlength{\parskip}{1.2pt}%
\noindent\textbf{168-70} \textbf{Φιλοκλῆς, Ξενοκλῆς, Θέογνις.} Trois poètereaux tragiques contemporains d'Aristophane.\\ À propos du vilain Philoclès, deux scholies nous ont été transmises dans le manuscrit de Ravenne\marqueur{1} :
\citComm{Φιλοκλέης: τραγῳδίας ποιητὴς, αἰσχρὸς τὸ πρόσωπον καὶ σιμὸς καὶ μικρὸς τὸ σῶμα\marqueur{2}.\\
ἐκ τούτου δηλονότι καὶ ἐν τοῖς Ὄρνισιν (v. 1295) οὐκ ἄλλως κορυδῷ αὐτὸν ὡ[μοίου] ἀλλ᾽ ἢ ὅτι μικρὸς καὶ αἰσχρός. καὶ Τηλεκλείδης Ἡσιόδοις\marqueur{3}·
`ἀλλ᾽ ἡ τάλαινα Φιλοκλέα βδε\douteux{λ}\douteux{λ}...\douteux{λ}.οθεν οὖν·
εἰ δ᾽ ἐστὶν Αἰσχύλου φρόνημ᾽ ἔχων'.}
\noindent Le scholiaste ne se trompe pas lorsqu'il prétend que la comparaison avec l'\italique{alouette huppée}, κορυδός, était d'ordre physique\marqueur{4} : on connaît par les parémiographes\marqueur{5} le proverbe grec «~être plus hideux que Corydée~» :
\citComm{Κορυδέως εἰδεχθέστερος· ἐπὶ δυσμορφίᾳ ἐκωμῳδεῖτο καὶ οὗτος καὶ οἱ παῖδες αὐτοῦ ὧν ἕνα τὸν Ἀρχέστρατον φασιν.}
\noindent La pauvre alouette huppée n'est pas bien lotie par la tradition populaire : en plus d'être laide, même dans la force de l'âge elle est chétive (plus que l'aigle vieillissant : ἀετοῦ γῆρας, κορύδου νεότης), et n'a pas pour autant le ramage du rossignol (ἐν ἀμούσοις καὶ κόρυδος φθέγγεται\marqueur{6})... « aux pays des aveugles, les borgnes sont rois ».\\ (Voir le commentaire au vers 441 pour le cas Xénoclès.)\\ Jean Taillardat a consacré la note que voici à Théognis\marqueur{7} :\\ Théognis le Tragique est ennuyeux et \italique{froid}, ψυχρός, au point qu'on le surnomme la \italique{Neige} (Χιών, Σ \italique{Acharniens} 11). La représentation d'une de ses pièces suffit à déchaîner les rigueurs de l'hiver, même sous le ciel déjà peu clément de la Thrace (\italique{Ach.} 138-140). Pris de pitié pour un froid rimailleur, Pisthétaire lui offre un vêtement de peau (\italique{Oiseaux} 935) :
\citComm{ἔχε τὴν σπολάδα· πάντως δέ μοι ῥιγῶν δοκεῖς.\\
\francais{Prends ce manteau ; tu me parais complètement gelé.}}
\end{multicols}
}%
\end{adjustbox}%
\vfill
\end{adjustwidth}
\newpage
% --- PAGE DE DROITE : Traduction française + Commentaire (2/6) ---
\begin{adjustwidth}{-2cm}{-2cm}
\noindent\textbf{\footnotesize Traduction française --- Folio 05 (\texttt{f. 164v}, vv. 165–207 [1/3]) \hfill ΠΡΟΛΟΓΟΣ}\\[-2pt]
\noindent\rule{16.5cm}{0.5pt}\\[2pt]
{\fontsize{8.4pt}{9.7pt}\selectfont
\ligneEst{}{\francais{beau lui-même, il s'habillait de beaux vêtements}}{165}
\ligneEst{}{\francais{et c'est aussi ce qui fait que ses drames étaient beaux.}}{}
\ligneEst{}{\francais{Car c'est une nécessité que l'on compose suivant sa nature.}}{}
\ligneEst{\textbf{Mnésiloque}}{\francais{Voilà aussi pourquoi Philoclès, qui est laid, compose de laides choses ;}}{}
\ligneEst{}{\francais{pourquoi Xénoclès, qui est mauvais, en compose de  mauvaises ;}}{}
\ligneEst{}{\francais{pourquoi Théognis, qui est froid, compose froidement.}}{170}
\ligneEst{\textbf{Agathon}}{\francais{De toute rigueur. C'est, en effet, pour avoir su cela}}{}
\ligneEst{}{\francais{que j'ai soigné ma personne.}}{}
\ligneEst{\textbf{Mnésiloque}}{\francais{\espace{3.2cm}Comment, au nom des dieux ?}}{}
\ligneEst{\textbf{Euripide}}{\francais{Cesse de japper. Car tel j'étais aussi moi-même}}{}
\ligneEst{}{\francais{à son âge, quand je commençais à produire.}}{}
\ligneEst{\textbf{Mnésiloque}}{\francais{Par Zeus ! je ne t'envie pas ton éducation.}}{175}
\ligneEst{\textbf{Euripide}}{\francais{Mais laisse-moi dire ce qui m'amène.}}{}
\ligneEst{\textbf{Agathon}}{\francais{\espace{5.5cm}Parle.}}{}
\ligneEst{\textbf{Euripide}}{\francais{Agathon, « c'est le propre d'un sage de savoir résumer en bref}}{}
\ligneEst{}{\francais{de longs discours ».}}{}
\ligneEst{}{\francais{Frappé d'une calamité inouïe,}}{}
}
\vspace{2pt}
\noindent\rule{16.5cm}{0.6pt}\\[2pt]
\noindent\textbf{\small Commentaire philologique, métrique et scénique (vv. 165–207 --- partie 2/6)}\\[2pt]
\begin{adjustbox}{max totalheight=13.0cm, max width=16.5cm, keepaspectratio}%
\parbox[t]{16.5cm}{
\setlength{\columnsep}{9pt}%
\begin{multicols}{2}
\fontsize{8.4pt}{9.7pt}\selectfont
\setlength{\parindent}{0pt}%
\setlength{\parskip}{1.2pt}%
\noindent Cette plaisanterie est continuée aux vers 946 et 955 (ταυταγὶ | τὰ κρυερὰ, \italique{ces vers glacés} ; cf. νιφοβόλους ἀναβολάς (\italique{ib.} 1385), « des préludes battus de neige ». Cet emploi de ψυχρός est banal\marqueur{8}.)
\noindent\textbf{[Notes ad 168-70]} \textbf{(1)} Manuscrit qui possède l'unique exemplaire des \italique{Thesmophories} qui nous soit parvenu -- \italique{unicodex} dit-on. Le manuscrit de Munich (\nieme{XV} siècle) n'étant qu'une recopie du premier (voir l'édition de von Velsen des \italique{Aristophanis Thesmophoriazusae} (1883) pour une comparaison complète du \italique{Ravennas} (R) et du \italique{Monacensis} (G)).
\noindent \textbf{(2)} Suidas (φ 379) reprend cette scholie qu'il joint à d'autres éléments tirés du scholiaste aux \italique{Oiseaux} (v. 281).
\noindent \textbf{(3)} fr. 15 in \italique{Poetae Comici Graeci}, VII, edd. Kassel et Austin.
\noindent \textbf{(4)} Ce qui explique le surnom du \italique{Coryde} dont il est question dans les \italique{Deipnosophistes} d'Athénée (VI, 241a \&c.) «~qui aimait à se faire persifler~».
\noindent \textbf{(5)} Zénobius IV, 59 ; Diogenianus V, 57 ; Michel Apostolius IX, 92 ; \&c. à retrouver dans le \italique{Corpus paroemiographorum graecorum}, ed. Leutsch. Cf. aussi Hésychius (σ 3681). Kassel \& Austin l'ont inclus (n° 827) aux \italique{Comica Adespota} de leur édition des \italique{Poetae Comici Graeci}, VIII.
\noindent \textbf{(6)} Voir ce qu'en dit Érasme, qui a recensé ce proverbe dans ses \italique{Adages} (n° 1192).
\noindent \textbf{(7)} \italique{Les images d'Aristophane} (1965), §759.
\noindent \textbf{(8)} Ψυχρός en parlant d'un \italique{poète} : Alexis 179 [$=$ 184 Kassel-Austin $=$ Athénée, III, 123e], Machon (\italique{ap.} Athénée 580a), d'un \italique{orateur} : Plat. \italique{Euthyd.} 284e ; Théophil. 4, 4, \&c. Ψυχρός se dit aussi d'une tragédie manquée (\italique{Thesm.} 848), d'une plaisanterie sans sel (Eupol. 244 [$=$ 261 Kassel-Austin]).
\end{multicols}
}%
\end{adjustbox}%
\vfill
\end{adjustwidth}
\newpage

% --- FOLIO 05 (f. 164v) : SPREAD 2.2 (Texte grec établi + Traduction française + Commentaire 3/6 & 4/6) ---
\clearpage\ifodd\value{page}\hbox{}\newpage\fi
\enTeteG{ΘΕΣΜΟΦΟΡΙΑΖΟΥΣΑΙ --- Texte grec établi (f. 164v, vv. 165–207 [2/3])}
\enTeteD{LES THESMOPHORIES --- Traduction française (f. 164v, vv. 165–207 [2/3])}
\begin{adjustwidth}{-2cm}{-2cm}
\noindent\textbf{\footnotesize Texte grec établi --- Folio 05 (\texttt{f. 164v}, vv. 165–207 [2/3]) \hfill ΠΡΟΛΟΓΟΣ}\\[-2pt]
\noindent\rule{16.5cm}{0.5pt}\\[2pt]
{\fontsize{8.4pt}{9.7pt}\selectfont
\ligneEst{}{ἱκέτης ἀφῖγμαι πρὸς σέ.}{180}
\ligneEst{\textbf{Ἀγάθων}}{\decalage{ἱκέτης ἀφῖγμαι πρὸς σέ.}τοῦ χρείαν ἔχων;}{}
\ligneEst{\textbf{Εὐριπίδης}}{μέλλουσί μ᾽ αἱ γυναῖκες ἀπολεῖν τήμερον}{}
\ligneEst{}{τοῖς Θεσμοφορίοις, ὅτι κακῶς αὐτὰς λέγω.}{}
\ligneEst{\textbf{Ἀγάθων}}{τίς οὖν παρ᾽ ἡμῶν ἐστιν ὠφέλειά σοι;}{}
\ligneEst{\textbf{Εὐριπίδης}}{ἡ πᾶσ᾽· ἐὰν γὰρ ἐγκαθεζόμενος λάθρᾳ}{}
\ligneEst{}{ἐν ταῖς γυναιξίν, ὡς δοκῶν εἶναι γυνή,}{185}
\ligneEst{}{ὑπεραποκρίνῃ μου, σαφῶς σώσεις ἐμέ.}{}
\ligneEst{}{μόνος γὰρ ἂν λέξειας ἀξίως ἐμοῦ.}{}
\ligneEst{\textbf{Ἀγάθων}}{ἔπειτα πῶς οὐκ αὐτὸς ἀπολογεῖ παρών;}{}
\ligneEst{\textbf{Εὐριπίδης}}{ἐγὼ φράσω σοι. πρῶτα μὲν γιγνώσκομαι}{}
\ligneEst{}{ἔπειτα πολιός εἰμι καὶ πώγων᾽ ἔχω,}{190}
\ligneEst{}{σὺ δ᾽ εὐπρόσωπος λευκὸς ἐξυρημένος}{}
\ligneEst{}{γυναικόφωνος ἁπαλὸς εὐπρεπὴς ἰδεῖν.}{}
\ligneEst{\textbf{Ἀγάθων}}{Εὐριπίδη—}{}
\ligneEst{\textbf{Εὐριπίδης}}{\decalage{Εὐριπίδη—}τί ἔστιν;}{}
\ligneEst{\textbf{Ἀγάθων}}{\decalage{Εὐριπίδη—τί ἔστιν;}ἐποίησάς ποτε,}{}
}
\vspace{2pt}
\noindent\rule{16.5cm}{0.6pt}\\[2pt]
\noindent\textbf{\small Commentaire philologique, métrique et scénique (vv. 165–207 --- partie 3/6)}\\[2pt]
\begin{adjustbox}{max totalheight=13.0cm, max width=16.5cm, keepaspectratio}%
\parbox[t]{16.5cm}{
\setlength{\columnsep}{9pt}%
\begin{multicols}{2}
\fontsize{8.4pt}{9.7pt}\selectfont
\setlength{\parindent}{0pt}%
\setlength{\parskip}{1.2pt}%
\noindent\textbf{172-3} \textbf{πῶς πρὸς τῶν θεῶν; — παῦσαι βαΰζων.} Goûtons à l'allitération en `t', `p' et surtout en `o' aboyée par Mnésiloque : π\textbf{ῶ}ς πρ\textbf{ὸ}ς τ\textbf{ῶ}ν θε\textbf{ῶ}ν ; c'est cette dissonance que lui reproche Euripide – pour la forme, car la question de Mnésiloque est aussi tout-à-fait désobligeante et ironique... il est certainement à l'affût de quelque double sens au «~ἐμαυτὸν ἐθεράπευσα~» d'Agathon.
\noindent Euripide interrompt donc : παῦσαι βαΰζων, « cesse d'aboyer ! » -- quoique le verbe «~aboyer~» ne rende pas le registre tragique du verbe βαΰζω, formé par onomatopée sur le cri βαῢ βαΰ du chien (αὗ αὗ au vers 903 des \italique{Guêpes}).
\noindent\textbf{173-5} \textbf{καὶ γὰρ ἐγὼ τοιοῦτος ἦν ... — οὐ ζηλῶ σε τῆς παιδεύσεως.} Euripide se reconnaît – non sans quelque condescendance – une parenté artistique\marqueur{1} quand Mnésiloque répond comme si Agathon était héritier des mœurs d'Euripide.
\noindent\textbf{[Notes ad 173-5]} \textbf{(1)} Confirmée par Aristote. (Voir le commentaire \italique{infra} au vers 29, n. 3.)
\noindent\textbf{176} \textbf{λέγε.} \latin{Nam neque Euripides neque Agatho Mnesilochi jocos solent curare. 
Apparet vero, orationem sic cohaerere},
\citComm{\grec{παῦσαι βαΰζων — άλλ' ... ἔα μ᾽ εἰπεῖν}\\
\latin{desine me maledictis insequi, sed potius ea me patere loqui.}}
\noindent \latin{Jam necessario is, quis rogatus est, respondere debet, estque certo verbum λέγε dic igitur Agathonis.} (Note de Fritzsche.)
\noindent\textbf{177-8} \textbf{Ἀγάθων σοφοῦ πρὸς ἀνδρός ὅστις ἐν βραχεῖ / πολλοὺς καλῶς οἷός τε συντέμνειν λόγους.} Citation quasi-textuelle de l'\italique{Éole} (Αἰόλος) d'Euripide rapportée par Stobée (3, 35, 3)\marqueur{1} :\citComm{παῖδες, σοφοῦ πρὸς ἀνδρός ὅστις ἐν βραχεῖ,\\
πολλοὺς λόγους οἷός τε συντέμνειν καλῶς}
\end{multicols}
}%
\end{adjustbox}%
\vfill
\end{adjustwidth}
\newpage
% --- PAGE DE DROITE : Traduction française + Commentaire (4/6) ---
\begin{adjustwidth}{-2cm}{-2cm}
\noindent\textbf{\footnotesize Traduction française --- Folio 05 (\texttt{f. 164v}, vv. 165–207 [2/3]) \hfill ΠΡΟΛΟΓΟΣ}\\[-2pt]
\noindent\rule{16.5cm}{0.5pt}\\[2pt]
{\fontsize{8.4pt}{9.7pt}\selectfont
\ligneEst{}{\francais{je viens à toi en suppliant.}}{180}
\ligneEst{\textbf{Agathon}}{\francais{\espace{3.5cm}Que te faut-il ?}}{}
\ligneEst{\textbf{Euripide}}{\francais{Les femmes pensent consommer ma perte, aujourd'hui,}}{}
\ligneEst{}{\francais{aux Thesmophories, parce que je dis du mal d'elles.}}{}
\ligneEst{\textbf{Agathon}}{\francais{Et que pouvons-nous pour toi ?}}{}
\ligneEst{\textbf{Euripide}}{\francais{Tout. Si siégeant en tapinois parmi les femmes,}}{}
\ligneEst{}{\francais{en feignant d'être l'une d'elles, tu prends ma défense,}}{185}
\ligneEst{}{\francais{tu ne manqueras pas à me sauver.}}{}
\ligneEst{}{\francais{Car seul tu parleras d'une manière digne de moi.}}{}
\ligneEst{\textbf{Agathon}}{\francais{Alors, que ne vas-tu te défendre en personne ?}}{}
\ligneEst{\textbf{Euripide}}{\francais{Je te le dirai. D'abord on me connaît.}}{}
\ligneEst{}{\francais{Ensuite je suis tout chenu, et j'ai de la barbe.}}{190}
\ligneEst{}{\francais{Toi, tu es avenant de figure, blanc, rasé,}}{}
\ligneEst{}{\francais{doué d'une voix de femme, mignon, plaisant à voir.}}{}
\ligneEst{\textbf{Agathon}}{\francais{Euripide...}}{}
\ligneEst{\textbf{Euripide}}{\francais{\espace{1.5cm}Eh bien ?}}{}
\ligneEst{\textbf{Agathon}}{\francais{\espace{3cm}Tu écrivis un jour : « Tu prends plaisir}}{}
}
\vspace{2pt}
\noindent\rule{16.5cm}{0.6pt}\\[2pt]
\noindent\textbf{\small Commentaire philologique, métrique et scénique (vv. 165–207 --- partie 4/6)}\\[2pt]
\begin{adjustbox}{max totalheight=13.0cm, max width=16.5cm, keepaspectratio}%
\parbox[t]{16.5cm}{
\setlength{\columnsep}{9pt}%
\begin{multicols}{2}
\fontsize{8.4pt}{9.7pt}\selectfont
\setlength{\parindent}{0pt}%
\setlength{\parskip}{1.2pt}%
\noindent P.-L. Courier, \italique{Pamphlet des pamphlets} :
\citComm{\francais{C'est un mérite non commun, ni facile, de clore en peu de mots beaucoup de sens.}
(Note de Willems.)}
\noindent\textbf{179} \textbf{ἐγὼ δὲ κοινῇ ξυμφορᾷ πεπληγμένος.} L'émendation de Biset : substituer καινῇ à κοινῇ est reprise par quasiment tous les éditeurs. Elle a l'avantage de concorder avec la scholie : μεγάλῃ καὶ θαυμαστῇ (καινός désigne au propre ce qui est \italique{nouveau}, différent de ce qui s'est produit jusqu'à présent, et par suite, ce qui n'était pas attendu, \italique{étrange}, \italique{extraordinaire} – avec une évolution toute identique à notre \italique{in-ouï}).
\noindent Nous sommes devant l'un des cas\marqueur{1} où la scholie rend compte d'un texte en meilleur état que celui auquel elle est attenante. (À noter qu'Euripide emploie l'expression καινάς συμφοράς aux vers 1226-7 de l'\italique{Oreste}.)\\ Cette glose s'est retrouvée dans la \italique{Souda} avec la mauvaise vedette (à l'article Κοινολογία). Pis : Suidas -- ou l'un des rédacteurs de la \italique{Souda} si l'on préfère --, voyant le problème aura ajouté aux deux termes dont le sens est incompatible avec κοινῇ un tiers, qui, lui, l'est, ἀνθρωπίνῃ.
\noindent Il écrit : κοινῇ δὲ ἀντὶ τοῦ ἀνθρωπίνῃ, μεγάλῃ καὶ θαυμαστῇ... Ce sont là des méthodes bien peu rigoureuses qui doivent nous le faire tenir en suspicion\marqueur{2}.
\noindent\textbf{[Notes ad 177-8]} \textbf{(1)} À retrouver dans l'édition de C. Wachsmuth \& O. Hense, \italique{Ioannis Stobaei Anthologium}, vol. 3, 1904, p. 688 l. 5-7. C'est le fragment 21 de l'édition Jouan-van Looy aux Belles Lettres, t. VIII, 1\italique{ère} partie ; il correspond respectivement aux fragments 28 et 33 des éditions Kannicht et Mette.
\noindent \textbf{(1)} Voir pour un autre exemple le commentaire au vers 1174.
\end{multicols}
}%
\end{adjustbox}%
\vfill
\end{adjustwidth}
\newpage

% --- FOLIO 05 (f. 164v) : SPREAD 2.3 (Texte grec établi + Traduction française + Commentaire 5/6 & 6/6) ---
\clearpage\ifodd\value{page}\hbox{}\newpage\fi
\enTeteG{ΘΕΣΜΟΦΟΡΙΑΖΟΥΣΑΙ --- Texte grec établi (f. 164v, vv. 165–207 [3/3])}
\enTeteD{LES THESMOPHORIES --- Traduction française (f. 164v, vv. 165–207 [3/3])}
\begin{adjustwidth}{-2cm}{-2cm}
\noindent\textbf{\footnotesize Texte grec établi --- Folio 05 (\texttt{f. 164v}, vv. 165–207 [3/3]) \hfill ΠΡΟΛΟΓΟΣ}\\[-2pt]
\noindent\rule{16.5cm}{0.5pt}\\[2pt]
{\fontsize{8.4pt}{9.7pt}\selectfont
\ligneEst{}{`χαίρεις ὁρῶν φῶς, πατέρα δ᾽ οὐ χαίρειν δοκεῖς;'}{}
\ligneEst{\textbf{Εὐριπίδης}}{ἔγωγε.}{195}
\ligneEst{\textbf{Ἀγάθων}}{\decalage{ἔγωγε.}μή νυν ἐλπίσῃς τὸ σὸν κακὸν}{}
\ligneEst{}{ἡμᾶς ὑφέξειν. καὶ γὰρ ἂν μαινοίμεθ᾽ ἂν.}{}
\ligneEst{}{ἀλλ᾽ αὐτὸς ὅ γε σόν ἐστιν οἰκείως φέρε.}{}
\ligneEst{}{τὰς συμφορὰς γὰρ οὐχὶ τοῖς τεχνάσμασιν}{}
\ligneEst{}{φέρειν δίκαιον ἀλλὰ τοῖς παθήμασιν.}{}
\ligneEst{\textbf{Μνησίλοχος}}{καὶ μὴν σύ γ᾽ ὦ κατάπυγον εὐρύπρωκτος εἶ}{200}
\ligneEst{}{οὐ τοῖς λόγοισιν ἀλλὰ τοῖς παθήμασιν.}{}
\ligneEst{\textbf{Εὐριπίδης}}{τί δ᾽ ἔστιν ὅτι δέδοικας ἐλθεῖν αὐτόσε;}{}
\ligneEst{\textbf{Ἀγάθων}}{κάκιον ἀπολοίμην ἂν ἢ σύ.}{}
\ligneEst{\textbf{Εὐριπίδης}}{\decalage{κάκιον ἀπολοίμην ἂν ἢ σύ.}πῶς;}{}
\ligneEst{\textbf{Ἀγάθων}}{\decalage{κάκιον ἀπολοίμην ἂν ἢ σύ.πῶς;}ὅπως;}{}
\ligneEst{}{δοκῶν γυναικῶν ἔργα νυκτερείσια}{}
\ligneEst{}{κλέπτειν ὑφαρπάζειν τε θήλειαν Κύπριν.}{205}
\ligneEst{\textbf{Μνησίλοχος}}{ἰδού γε κλέπτειν· νὴ Δία βινεῖσθαι μὲν οὖν.}{}
\ligneEst{}{ἀτὰρ ἡ πρόφασίς γε νὴ Δί᾽ εἰκότως ἔχει.}{207}
}
\vspace{2pt}
\noindent\rule{16.5cm}{0.6pt}\\[2pt]
\noindent\textbf{\small Commentaire philologique, métrique et scénique (vv. 165–207 --- partie 5/6)}\\[2pt]
\begin{adjustbox}{max totalheight=13.0cm, max width=16.5cm, keepaspectratio}%
\parbox[t]{16.5cm}{
\setlength{\columnsep}{9pt}%
\begin{multicols}{2}
\fontsize{8.4pt}{9.7pt}\selectfont
\setlength{\parindent}{0pt}%
\setlength{\parskip}{1.2pt}%
\noindent \textbf{(2)} Voir \italique{infra} le commentaire au vers 735 pour une erreur analogue de Suidas.
\noindent\textbf{180} \textbf{τοῦ χρείαν ἔχων;} Euripide, \italique{Suppliantes}, 115 : Τί χρῆμα θηρῶν καὶ τίνος χρείαν ἔχων;
\noindent\textbf{186} \textbf{σαφῶς σώσεις ἐμέ.} Le sigmatisme tant reproché à Euripide ! toujours avec le verbe σῴζειν, \italique{sauver}, le véritable Euripide avait sifflé au vers 476 de sa \italique{Médée} : ἔσῳσά σ' ὡς ἴσασιν Ἑλλήνων ὅσοι. «~Grand merci de nous avoir sauvé des sigmas d'Euripide~» avait persiflé Platon le Comique\marqueur{1} en reprenant le début de vers :
\citComm{εὖ γέ σοι γενοιθ᾽ ἡμᾶς ὅτι\\
ἔσωσας ἐκ τῶν σίγμα τῶν Εὐριπίδου.}
\noindent\textbf{[Notes ad 186]} \textbf{(1)} Fragment 29(30) de l'édition Kassel-Austin, \italique{Poetae Comici Graeci}, vol. VII (1989) ; chez Meineke, \italique{Fragmenta Comicorum Graecorum}, il se trouve au tome II, p. 626.
\noindent\textbf{193} \textbf{χαίρεις ὁρῶν φῶς, πατέρα δ᾽ οὐ χαίρειν δοκεῖς;} Citation tirée de l'\italique{Alceste} d'Euripide, v. 691 : Phérès envoie paître son fils Admète qui lui demande de mourir à sa place. C'est finalement Alceste qui se sacrifiera pour son mari ; tout comme Mnésiloque, parent par \italique{alliance} d'Euripide, ne va pas tarder à accepter ce qui ressemble fort à un suicide.
\noindent\textbf{200-1} \textbf{εὐρύπρωκτος εἶ / οὐ τοῖς λόγοισιν ἀλλὰ τοῖς παθήμασιν.} Mnésiloque reprend l'opposition τοῖς τεχνάσμασιν / τοῖς παθήμασιν d'Agathon en substituant τοῖς λόγοισιν à τοῖς τεχνάσμασιν.\\ Le spectateur attend donc l'opposition classique τοῖς λόγοις / τοῖς ἔργοις : \italique{en théorie} / \italique{en pratique}... sauf qu'ici les \italique{pratiques} d'Agathon sont jugées par trop passives.
\end{multicols}
}%
\end{adjustbox}%
\vfill
\end{adjustwidth}
\newpage
% --- PAGE DE DROITE : Traduction française + Commentaire (6/6) ---
\begin{adjustwidth}{-2cm}{-2cm}
\noindent\textbf{\footnotesize Traduction française --- Folio 05 (\texttt{f. 164v}, vv. 165–207 [3/3]) \hfill ΠΡΟΛΟΓΟΣ}\\[-2pt]
\noindent\rule{16.5cm}{0.5pt}\\[2pt]
{\fontsize{8.4pt}{9.7pt}\selectfont
\ligneEst{}{\francais{à voir cette lumière, penses-tu que ton père y prenne moins plaisir ? »}}{}
\ligneEst{\textbf{Euripide}}{\francais{Sans doute.}}{195}
\ligneEst{\textbf{Agathon}}{\francais{\espace{1.2cm}N'espère donc pas que je prenne sur mon compte ton infortune.}}{}
\ligneEst{}{\francais{Il faudrait que je fusse fou.}}{}
\ligneEst{}{\francais{Charge-toi personnellement de tes propres affaires.}}{}
\ligneEst{}{\francais{Car les revers, il convient de les subir à force, non d'artifices,}}{}
\ligneEst{}{\francais{mais d'endurance.}}{}
\ligneEst{\textbf{Mnésiloque}}{\francais{Au fait toi-même, bougre de cul, tu as l'anus béant à force}}{200}
\ligneEst{}{\francais{non de discours mais d'endurance.}}{}
\ligneEst{\textbf{Euripide}}{\francais{Mais qu'est-ce qui te fait craindre d'aller là ?}}{}
\ligneEst{\textbf{Agathon}}{\francais{Je finirais plus misérablement que toi.}}{}
\ligneEst{\textbf{Euripide}}{\francais{\espace{4cm}Pourquoi ?}}{}
\ligneEst{\textbf{Agathon}}{\francais{\espace{5cm}Pourquoi ?}}{}
\ligneEst{}{\francais{Je passe pour ravir les travaux nocturnes des femmes,}}{}
\ligneEst{}{\francais{et dérober la Cypris féminine.}}{205}
\ligneEst{\textbf{Mnésiloque}}{\francais{Mon cul, ravir ! Eh par Zeus! dis de te faire enfiler.}}{}
\ligneEst{}{\francais{Mais, par Zeus, la défaite est bonne.}}{207}
}
\vspace{2pt}
\noindent\rule{16.5cm}{0.6pt}\\[2pt]
\noindent\textbf{\small Commentaire philologique, métrique et scénique (vv. 165–207 --- partie 6/6)}\\[2pt]
\begin{adjustbox}{max totalheight=13.0cm, max width=16.5cm, keepaspectratio}%
\parbox[t]{16.5cm}{
\setlength{\columnsep}{9pt}%
\begin{multicols}{2}
\fontsize{8.4pt}{9.7pt}\selectfont
\setlength{\parindent}{0pt}%
\setlength{\parskip}{1.2pt}%
\noindent\textbf{204} \textbf{νυκτερείσια.} La forme canonique est νυκτερήσια, \italique{nocturne} (attestée tardivement par Lucien et Sextus Empiricus), forgée en miroir de ἡμερήσιος.
\noindent Gardons-nous pour autant de retoucher la graphie du codex de Ravenne ; la réplique perdrait de son sel aristophanesque :
\noindent νυκτ-ερείσια fait entendre le substantif du verbe ἐρείδειν, « \latin{quod in comoedia est, subagitare mulierem} » (Fritzsche).
\noindent\textbf{205} \textbf{ὑφαρπάζειν τε θήλειαν Κύπριν.} L'expression ὑφαρπάζειν Κύπριν se retrouve au \nieme{722} vers de l'\italique{Assemblée des femmes}.
\noindent Les Athéniennes ont toujours une bonne raison d'en vouloir aux poètes tragiques : lorsqu'ils ne médisent pas d'elles, ils se livrent à une concurrence déloyale.
\citComm{Mais les jeunes femmes d'Athènes\\
Privées d'leurs époux, d'leurs galants,\\
Accumulèrent la rancune patiemment...\\
Puis un jour, ivres de colère,\\
Elles s'armèrent de bâtons,\\
Et, farouches, elles immolèrent Agathon !}
\noindent\textbf{207} \textbf{ἀτὰρ ἡ πρόφασίς γε νὴ Δί᾽ εἰκότως ἔχει.} Molière, \italique{Misanthrope}, IV, 3 :
\citComm{\francais{Ah! le détour est bon et l'excuse admirable.}}
\noindent \textbf{La défaite}: prétexte fallacieux pour se défaire de ce qui est importun, embarrassant. (J'ai laissé le mot soigneusement choisi par Willems.)
\end{multicols}
}%
\end{adjustbox}%
\vfill
\end{adjustwidth}
\newpage
\clearpage\ifodd\value{page}\hbox{}\newpage\fi
\enTeteG{ΠΡΟΛΟΓΟΣ --- Codex Ravennas 429, f. 165r (Classense 167r) --- vv. 208–251}
\begin{adjustwidth}{-1.5cm}{-3.0cm}
\begin{adjustbox}{max totalheight=21.0cm, max width=16.5cm, keepaspectratio}
\begin{minipage}{16.5cm}
{\fontsize{7.2pt}{8.5pt}\selectfont
\noindent\textbf{Folio 06 --- Codex Ravennas 429, \texttt{f. 165r} (\texttt{Classense 167r}) : ΠΡΟΛΟΓΟΣ, vv. 208–251}\hfill \textit{Témoins :} $\text{R}^1, \Rpc, \Rdeux$, \PsiUn\par\vspace{2pt}
\hrule height 0.4pt\vspace{2.5pt}
}
{\fontsize{8.1pt}{9.5pt}\selectfont
\ligneRav{208}{\auDessus{εὐ}{ρ}}{τί \difference{οὖν}; ποιήσεις \difference{ταῦτα}; : μὴ δόκει γε σύ·}{\textbf{208} οὖν et ταῦτα R : \lecon{οὖν;} et \lecon{ταῦτα;} Ellebodius}
\ligneRav{209}{\auDessus{εὐ}{ρ}}{ὦ τρισκακοδαίμον, ὡς ἀπόλωλ᾽ Εὐριπίδης·}{\textbf{209} \lecon{τρισκακόδαιμον} R : τρισκακοδαίμων \Mon \espApp \lecon{(Εὐρ.)} R : (Μνη.) Elmsley ad \italique{Ach}. 475}
\ligneRav{210}{μ\auDessus{ν}{Η}}{ὦ φίλτατ᾽ ὦ κηδεστὰ μὴ σαυτὸν προδῷς·}{}
\ligneRav{}{\auDessus{εὐ}{ρ}}{πῶς οὖν ποιήσω δῆτα; : τοῦτον μὲν μακρὰ}{}
\ligneRav{}{}{κλάειν κέλευ᾽, ἐμοὶ δ᾽ ὅ τι βούλει χρῶ λαβών·}{}
\ligneRav{}{\auDessus{εὐ}{ρ}}{ἄγε νυν ἐπειδὴ σαυτὸν ἐπιδίδως ἐμοί,}{}
\ligneRav{}{—}{ἀπόδυθι τουτὶ θοἰμάτιον. : καὶ δὴ χαμαί·}{}
\ligneRav{215}{—}{ἀτὰρ τί μέλλεις δρᾶν \difference{μ᾽}; : ἀποξυρεῖν ταδί,}{\textbf{215} με R : \lecon{μ᾽} edd. omnes}
\ligneRav{216}{—}{τὰ κάτω δ᾽ ἀφεύειν. : ἀλλὰ πρᾶττ᾽, \difference{εἴ σοι} δοκεῖ·}{\textbf{216} \lecon{ἀφεύειν} R : ἀφαύειν Suid. α 4565 \espApp είσω R : \lecon{εἴ σοι} Bonini}
\ligneRav{217}{}{ἢ μὴ \difference{διδόναι γ' ἑαυτὸν} ὤφελόν ποτε·}{\textbf{217} διδόναι γ' αὐτὸν R : διδόναι γ' ἑμαυτὸν Scaliger : \lecon{διδόναι γ' ἑαυτὸν} Bibi auctore Willems}
\ligneRav{}{\auDessus{εὐ}{ρ}}{Ἀγάθων σὺ μέντοι ξυροφορεῖς ἑκάστοτε,}{}
\ligneRav{}{—}{χρῆσόν τι νῦν ἡμῖν ξυρόν. : αὐτὸς λάμβανε}{}
\ligneRav{220}{—}{ἐντεῦθεν ἐκ τῆς \difference{ξυροδόκης}. : γενναῖος εἶ·}{\textbf{220} ξυροδίκης R : \lecon{ξυροδόκης} Σ, Pollux II.32, X.140}
\ligneRav{}{}{κάθιξε· φύσα τὴν γνάθον τὴν δεξιάν·}{}
\ligneRav{222}{μ\auDessus{ν}{Η}}{\difference{οἴμοι}. : τί κέκραγας; ἐμβαλῶ σοι πάτταλον,}{\textbf{222} ὤιμοι R : \lecon{οἴμοι} Dindorf}
\ligneRav{223}{—}{ἢν μὴ σιωπᾷς. : \difference{ἀτταταῖ ἰατταταῖ}·}{\textbf{223} ἀττατα· ατταται R : \lecon{ἀτταταῖ ἰατταταῖ} Biset}
\ligneRav{}{\auDessus{εὐ}{ρ}}{οὗτος σὺ ποῖ θεῖς; : εἰς τὸ τῶν σεμνῶν θεῶν·}{}
\ligneRav{225}{}{οὐ γὰρ μὰ τὴν Δήμητρά γ᾽ ἐνταυθοῖ μενῶ}{\textbf{225} \lecon{Δήμητρά γ᾽} R : Δήμητρ᾽ ἔτ᾽ Porson, cf. ad \italique{Eq}. 698 : Δημήτερ᾽ Fritzsche}
\ligneRav{226}{—}{τεμνόμενος. : \difference{οὔκουν} καταγέλαστος δῆτ᾽ ἔσει}{\textbf{226} οὐκ οὖν R : \lecon{οὔκουν} Brunck}
\ligneRav{}{}{τὴν ἡμίκραιραν τὴν ἑτέραν ψιλὴν ἔχων;}{}
\ligneRav{}{μ\auDessus{ν}{Η}}{ὀλίγον μέλει μοι. : μηδαμῶς πρὸς τῶν θεῶν}{}
\ligneRav{}{—}{προδῷς με· χώρει δεῦρο. : κακοδαίμων ἐγώ·}{}
\ligneRav{230}{\auDessus{εὐ}{ρ}}{ἔχ᾽ ἀτρέμα σαυτὸν κἀνάκυπτε· ποῖ στρέφει;}{}
\ligneRav{}{μ\auDessus{ν}{Η}}{μυμῦ. : τί μύζεις; πάντα πεποίηται καλῶς·}{}
\ligneRav{232}{μ\auDessus{ν}{Η}}{οἴμοι κακοδαίμων, ψιλὸς αὖ στρατεύσομαι·}{\textbf{232} \lecon{αὖ} R, Suid. ψ 105 : οὖν von Velsen}
\ligneRav{}{\auDessus{εὐ}{ρ}}{μὴ φροντίσῃ· ὡς εὐπρεπὴς φανεῖ πάνυ·}{}
\ligneRav{234}{—}{βούλει \difference{θεάσασθ᾽ αὑτόν}; : εἰ δοκεῖ, φέρε·}{\textbf{234} θεάσασθαι σεαυτόν R : \lecon{θεάσασθ᾽ αὑτόν} Fritzsche : θεᾶσθαι σαυτόν Porson}
\ligneRav{235}{\auDessus{εὐ}{ρ}}{ὁρᾷς σεαυτόν; : οὐ μὰ Δί᾽ ἀλλὰ \difference{Κλεισθένη}·}{\textbf{235} Κλεισθένην R, Suid. κ 1756 : \lecon{Κλεισθένη} Brunck}
\ligneRav{}{\auDessus{εὐ}{ρ}}{ἀνίστασ᾽, ἵν᾽ ἀφεύσω σε, κἀγκύψας ἔχε·}{}
\ligneRav{}{μ\auDessus{ν}{Η}}{οἴμοι κακοδαίμων δελφάκιον γενήσομαι·}{}
\ligneRav{}{\auDessus{εὐ}{ρ}}{ἐνεγκάτω τις ἔνδοθεν δᾷδ᾽ ἢ λύχνον·}{}
\ligneRav{}{}{ἐπίκυπτε· τὴν κέρκον φυλάττου νυν ἄκραν·}{}
\ligneRav{240}{μ\auDessus{ν}{Η}}{ἐμοὶ μελήσει νὴ Δία, πλήν γ᾽ ὅτι κάομαι·}{}
\ligneRav{}{}{οἴμοι τάλας. ὕδωρ ὕδωρ ὦ γείτονες}{}
\ligneRav{242}{}{τὸν \difference{[ἕ]τερον} πρὶν ἀντιλαβέσθαι φλόγα·}{\textbf{242} \lecon{[ἓ]τερον} \PsiUn \ (restituit Medaglia) : ἓτερον om. R}
\ligneRav{}{\auDessus{εὐ}{ρ}}{θάρρει. : τί θαρρῶ καταπεπυρπολημένος;}{}
\ligneRav{}{\auDessus{εὐ}{ρ}}{ἀλλ᾽ οὐκ ἔτ᾽ οἰδὲν πρᾶγμά σοι· τὰ πλεῖστα γὰρ}{}
\ligneRav{245}{—}{ἀποπεπόνηκας. : φεῦ ἰοὺ τῆς ἀσβόλου·}{\textbf{245} \lecon{φεῦ} R, Suid. α 4132 : φῦ Dindorf, Wilamowitz \italique{Lysistr.} v.294}
\ligneRav{246}{}{αἰθὸς \difference{γεγένημαι} πάντα τὰ περὶ τὴν τράμιν·}{\textbf{246} γεγέννημαι R : \lecon{γεγένημαι} Suid. τ 903}
\ligneRav{247}{\auDessus{εὐ}{ρ}}{μὴ φροντίσῃς· ἕτερος γὰρ αὐτὰ σπογγιεῖ·}{\textbf{247} \lecon{σπογγιεῖ} R : σφογγιεῖ Elmsley ad \italique{Ach.} 463}
\ligneRav{248}{μ\auDessus{ν}{Η}}{\difference{οἰμώξετ᾽} ἆρ᾽ \difference{εἴ τις} τὸν ἐμὸν πρωκτὸν πλυνεῖ·}{\textbf{248} οἰμώζετ᾽ R : \lecon{οἰμώξετ᾽} Zanetti \espApp εἰς R : \lecon{εἴ τις} Brunck}
\ligneRav{}{\auDessus{εὐ}{ρ}}{Ἀγάθων, ἐπειδὴ σαυτὸν ἐπιδοῦναι φθονεῖς,}{}
\ligneRav{250}{}{ἀλλ᾽ ἱμάτιον γοῦν χρῆσον ἡμῖν \difference{τουτῳὶ}}{\textbf{250} τουτωί R : \lecon{τουτῳὶ} edd. plerique}
\ligneRav{}{}{καὶ στρόφιον· οὐ γὰρ ταῦτά γ᾽ ὡς οὐκ ἔστ᾽ ἐρεῖς·}{}
}
\vspace{3pt}\noindent\rule{4.5cm}{0.4pt}\par\vspace{1.5pt}
{\fontsize{7.0pt}{8.3pt}\selectfont
\noindent\textbf{Loci paralleli (vv. 208–251) ---} \textbf{215-6} (ἀποξυρεῖν -- ἀφεύειν) Suidas α 4565 (ἀφαύειν pro ἀφεύειν), cf. Scholium ταῦτα δὲ ἔλαβεν ἐκ τῶν Ἰδαίων Κρατίνου (fr. 90 Kassel-Austin) ; \ClemAlex, \italique{Στρώματα} VI 26.4 Ἀριστοφάνης δὲ ὁ κωμικὸς ἐν ταῖς πρώταις Θεσμοφοριαζούσαις τὰ ἐκ τῶν Κρατίνου Ἐμπιπραμένων μετήνεγκεν ἔπη. \espTest \textbf{219-20} (χρῆσόν -- ξυροδόκης) Pollux X 140 (νῦν omisso) \espTest \textbf{220} (ξυροδόκης) Pollux II 32 ; \Eustathe 130.43, 1399.38 ; cf. Scholium \espTest \textbf{223} (ἀτταταῖ ἰατταταῖ) Suidas ε 2807 \espTest \textbf{227} Suidas η 334 \espTest (ἡμίκραιραν) Hesychius η 496 ἡμίκραιρα ; \EtymGen (AB) 150 ed. Miller ; \Eustathe 710.50, 1127.35 ; \AelDion η 9 \espTest \textbf{231} \EtymGen (AB) 218 ed. Miller \espTest (μῦ μῦ) Suidas ε 2807 ; cf. \Eustathe 723.9 \espTest (τί -- καλῶς) Suidas μ 1383 ; \Apostol XVI 60c \espTest (μύζεις) \Eustathe 440.23 \espTest \textbf{232} Suidas ψ 105 \espTest \textbf{235} Suidas κ 1756 Κλεισθένην ὁρῶ \espTest \textbf{237} Suidas δ 205 (sed γέγονα pro γενήσομαι) \espTest \textbf{244-5} (τὰ πλεῖστα -- ἀσβόλου) Photius α 2946 $=$ \SynB α 2213 ; \AelDion α 186 \espTest \textbf{245} (φεῦ -- ἀσβόλου) Suidas α 4132 \espTest (ἄσβολος $=$ ἀσβόλη) Moeris α 17 $\simeq$ Hesychius α 7634 ; cf. Phrynicus \italique{Praep. soph} 28.1 (ἄσβολος fem.) \espTest \textbf{246} Suidas τ 903 (sed τὸν pro τὴν) \espTest (αἰθὸς γεγένημαι) \Eustathe 83.6, 992.60, 1314.15 \espTest (αἰθὸς) \Eustathe 1385.59 \espTest (τράμιν) Hesychius τ 1243 ; Photius τ 415 \espTest \textbf{248} (πλυνεῖ) Moeris π 79\par
}
\vspace{2.5pt}\noindent\rule{4.5cm}{0.4pt}\par\vspace{1.5pt}
{\fontsize{6.9pt}{8.2pt}\selectfont
\noindent\textbf{Suite (en commentaire) de l'apparat critique \& notes paléographiques (\texttt{f. 165r}) ---} \textbf{Témoin papyrologique :} \textbf{\PsiUn{} (\italique{PSI} XI 1194)} couvre les \textbf{vv. 237–245} et comble au \textbf{v. 242} l'omission de [ἓ]τερον dans R. --- \textbf{v. 242 :} R avait omis un crétique avant πρὶν ἀντιλαβέσθαι ; le papyrus \PsiUn{} (lu par Medaglia) restitue \lecon{[ἓ]τερον}, tombé par saut visuel dans l'ancêtre de R. --- \textbf{v. 250 :} Exemple typique de \textbf{iota muet adscrit} dans $\text{R}^1$ : τουτωί. --- \textbf{v. 217 :} διδόναι γ' αὐτὸν R (διδόναι γ' ἐμαυτὸν Scaliger) : \lecon{διδόναι γ' ἑαυτὸν} \lecon{Bibi} \italique{auctore Willems}. --- \textbf{v. 220 :} ξυροδίκης R (iotacisme/attraction de $\text{δίκη}$) : \lecon{ξυροδόκης} $\Sigma$, Pollux II.32, X.140 (« étui à rasoir »). --- \textbf{v. 234 :} θεάσασθαι σεαυτόν R (glose pronominale non contracte) : \lecon{θεάσασθ' αὑτόν} Fritzsche (θεᾶσθαι σαυτόν Porson). --- \textbf{v. 242 :} \lecon{[ἓ]τερον} \PsiUn{} (restituit Medaglia) : ἕτερον om. R. --- \textbf{v. 248 :} οἰμώζετ' R : \lecon{οἰμώξετ'} Zanetti \espApp εἰς R (haplographie de ΤΙ dans ΕΙΤΙΣ) : \lecon{εἴ τις} Brunck.\par
}
\end{minipage}
\end{adjustbox}
\end{adjustwidth}
\newpage
\enTeteD{Codex Ravennas 429 (Biblioteca Classense) --- Fac-similé du f. 165r (vv. 208–251)}
\begin{adjustwidth}{-1.5cm}{-3.0cm}
\begin{center}
\includegraphics[width=16.5cm,height=20.5cm,keepaspectratio]{/workspaces/Thesmophories/dossier_ravennas_thesmophories/folios_individuels_ff162v-177r/06_f165r_classense167r_vv0208-0251.jpg}\par
\vspace{2pt}
{\scriptsize\textbf{Codex Ravennas 429, f. 165r (Classense 167r)} --- ΠΡΟΛΟΓΟΣ, \textbf{vv. 208–251} ($\text{R}^1, \Rpc, \Rdeux$, \PsiUn).}
\end{center}
\end{adjustwidth}
\newpage
\clearpage\ifodd\value{page}\hbox{}\newpage\fi% --- FOLIO 06 (f. 165r) : SPREAD 2.1 (Texte grec établi + Traduction française + Commentaire 1/6 & 2/6) ---
\clearpage\ifodd\value{page}\hbox{}\newpage\fi
\enTeteG{ΘΕΣΜΟΦΟΡΙΑΖΟΥΣΑΙ --- Texte grec établi (f. 165r, vv. 208–251 [1/3])}
\enTeteD{LES THESMOPHORIES --- Traduction française (f. 165r, vv. 208–251 [1/3])}
\begin{adjustwidth}{-2cm}{-2cm}
\noindent\textbf{\footnotesize Texte grec établi --- Folio 06 (\texttt{f. 165r}, vv. 208–251 [1/3]) \hfill ΠΡΟΛΟΓΟΣ}\\[-2pt]
\noindent\rule{16.5cm}{0.5pt}\\[2pt]
{\fontsize{8.4pt}{9.7pt}\selectfont
\ligneEst{\textbf{Εὐριπίδης}}{τί οὖν; ποιήσεις ταῦτα;}{208}
\ligneEst{\textbf{Ἀγάθων}}{\decalage{τί οὖν; ποιήσεις ταῦτα;}μὴ δόκει γε σύ.}{}
\ligneEst{\textbf{Εὐριπίδης}}{ὦ τρισκακοδαίμον, ὡς ἀπόλωλ᾽ Εὐριπίδης.}{}
\ligneEst{\textbf{Μνησίλοχος}}{ὦ φίλτατ᾽ ὦ κηδεστὰ μὴ σαυτὸν προδῷς.}{210}
\ligneEst{\textbf{Εὐριπίδης}}{πῶς οὖν ποιήσω δῆτα;}{}
\ligneEst{\textbf{Μνησίλοχος}}{\decalage{πῶς οὖν ποιήσω δῆτα;}τοῦτον μὲν μακρὰ}{}
\ligneEst{}{κλάειν κέλευ᾽, ἐμοὶ δ᾽ ὅ τι βούλει χρῶ λαβών.}{}
\ligneEst{\textbf{Εὐριπίδης}}{ἄγε νυν ἐπειδὴ σαυτὸν ἐπιδίδως ἐμοί,}{}
\ligneEst{}{ἀπόδυθι τουτὶ θοἰμάτιον.}{}
\ligneEst{\textbf{Μνησίλοχος}}{\decalage{ἀπόδυθι τουτὶ θοἰμάτιον.}καὶ δὴ χαμαί.}{}
\ligneEst{}{ἀτὰρ τί μέλλεις δρᾶν μ᾽;}{215}
\ligneEst{\textbf{Εὐριπίδης}}{\decalage{ἀτὰρ τί μέλλεις δρᾶν μ᾽;}ἀποξυρεῖν ταδί,}{}
\ligneEst{}{τὰ κάτω δ᾽ ἀφεύειν.}{}
\ligneEst{\textbf{Μνησίλοχος}}{\decalage{τὰ κάτω δ᾽ ἀφεύειν.}ἀλλὰ πρᾶττ᾽, εἴ σοι δοκεῖ.}{}
\ligneEst{}{ἢ μὴ διδόναι γ' ἑαυτὸν ὤφελόν ποτε.}{}
\ligneEst{\textbf{Εὐριπίδης}}{Ἀγάθων σὺ μέντοι ξυροφορεῖς ἑκάστοτε,}{}
\ligneEst{}{χρῆσόν τι νῦν ἡμῖν ξυρόν.}{}
\ligneEst{\textbf{Ἀγάθων}}{\decalage{χρῆσόν τί νυν ἡμῖν ξυρόν.}αὐτὸς λάμβανε}{}
\ligneEst{}{ἐντεῦθεν ἐκ τῆς ξυροδόκης.}{220}
\ligneEst{\textbf{Εὐριπίδης}}{\decalage{ἐντεῦθεν ἐκ τῆς ξυροδόκης.}γενναῖος εἶ.}{}
}
\vspace{2pt}
\noindent\rule{16.5cm}{0.6pt}\\[2pt]
\noindent\textbf{\small Commentaire philologique, métrique et scénique (vv. 208–251 --- partie 1/6)}\\[2pt]
\begin{adjustbox}{max totalheight=13.0cm, max width=16.5cm, keepaspectratio}%
\parbox[t]{16.5cm}{
\setlength{\columnsep}{9pt}%
\begin{multicols}{2}
\fontsize{8.4pt}{9.7pt}\selectfont
\setlength{\parindent}{0pt}%
\setlength{\parskip}{1.2pt}%
\noindent\textbf{209} \textbf{ὦ τρισκακόδαιμον ὡς ἀπόλωλ᾽ Εὐριπίδης.} La plupart des éditeurs modernes découpent cette réplique, attribuant la première partie à Euripide, puis Mnésiloque de lancer l'apostrophe : Εὐριπίδη κτλ. Cela les oblige à retoucher le texte pour faire de ὤ τρισκακοδαίμων le sujet de ἀπόλωλ(α) et changer Εὐριπίδης en Εὐριπίδη...
\noindent Tout cela paraît bien tarabiscoté alors que le texte de notre unicodex se laisse comprendre facilement : « Ô triple malheureux ! ainsi périt Euripide », apostrophe qu'Euripide s'adresse à lui-même d'autant moins étonnante qu'il a déjà parlé de lui à la \nieme{3} personne (v. 77).
\noindent En outre, si l'attribution des répliques dans le manuscrit de Ravenne est souvent folklorique, la découpe des répliques, elle, est bien plus sûre. Philogues, réfrénez vos pulsions à vouloir tailler ainsi le vers.
\noindent\textbf{217} \textbf{ἢ μὴ διδόναι γ' αὐτὸν ὤφελόν ποτε.} Le vers pèche contre le mètre.
\noindent Willems (voir sa note 4 en annexe), en guerre perpétuelle contre les retoucheurs de texte, rendait justice à ce qu'il croyait être le texte original. (En fait légèrement modifié : γ' ἑμαυτὸν à la place de γ' αὐτὸν. Willems a-t-il été trompé par Dindorf ?)  Il s'en justifiait par des exemples qui ne laissent nul doute de l'epsilonesque correction qu'il aurait lui-même proposée s'il avait eu le texte des manuscrits sous les yeux : substituer γ' ἑαυτὸν à γ' αὐτὸν.\\
\noindent Au reste, comme l'avait souligné Willems, « διδόναι ἑαυτὸν τινι est la formule ordinaire, et qui se lit partout : \italique{Cavaliers}., 740 ; Sophocle, \italique{Philoct}., 84 ; \italique{Trach}., 1117; Hérodote, VI, 108; Thucydide, II, 68; Démosthène, \italique{Sur la couronne}, 503d.~»
\noindent\textbf{224} \textbf{οὗτος σὺ ποῖ θεῖς;} formule identique au vers 854 des \italique{Guêpes}.
\noindent\textbf{224} \textbf{εἰς τὸ τῶν σεμνῶν θεῶν·} Le scholiaste :
\citComm{ἀντὶ τοῦ `εἰς <τὸ> τῶν ᾿Ερινύων'. ἄσυλον γὰρ εἶχον οἵ καταλαμβάνοντες τὰ ἱερὰ τῶν Ἐρινύων.}
\end{multicols}
}%
\end{adjustbox}%
\vfill
\end{adjustwidth}
\newpage
% --- PAGE DE DROITE : Traduction française + Commentaire (2/6) ---
\begin{adjustwidth}{-2cm}{-2cm}
\noindent\textbf{\footnotesize Traduction française --- Folio 06 (\texttt{f. 165r}, vv. 208–251 [1/3]) \hfill ΠΡΟΛΟΓΟΣ}\\[-2pt]
\noindent\rule{16.5cm}{0.5pt}\\[2pt]
{\fontsize{8.4pt}{9.7pt}\selectfont
\ligneEst{\textbf{Euripide}}{\francais{Eh bien ? le feras-tu ?}}{208}
\ligneEst{\textbf{Agathon}}{\francais{\espace{3.5cm}N'y compte pas.}}{}
\ligneEst{\textbf{Euripide}}{\francais{Ô triple malheureux, comment périt Euripide.}}{}
\ligneEst{\textbf{Mnésiloque}}{\francais{Mon ami, mon gendre, n'abandonne point.}}{210}
\ligneEst{\textbf{Euripide}}{\francais{Que puis-je faire ?}}{}
\ligneEst{\textbf{Mnésiloque}}{\francais{\espace{3.4cm}Envoie-le se faire}}{}
\ligneEst{}{\francais{pendre, lui, et dispose de moi comme tu voudras.}}{}
\ligneEst{\textbf{Euripide}}{\francais{Voyons, puisque tu te remets entre mes mains,}}{}
\ligneEst{}{\francais{ôte ce manteau.}}{}
\ligneEst{\textbf{Mnésiloque}}{\francais{\espace{3.7cm}Le voilà par terre.}}{}
\ligneEst{}{\francais{Mais que penses-tu faire de moi ?}}{215}
\ligneEst{\textbf{Euripide}}{\francais{\espace{3.7cm}Te raser cela}}{}
\ligneEst{}{\francais{et te flamber par dessous.}}{}
\ligneEst{\textbf{Mnésiloque}}{\francais{\espace{3cm}Eh bien, fais à ta guise ;}}{}
\ligneEst{}{\francais{autrement je n'aurais jamais dû me mettre à ta merci.}}{}
\ligneEst{\textbf{Euripide}}{\francais{Agathon, pour sûr, tu portes toujours sur toi des rasoirs.}}{}
\ligneEst{}{\francais{Prête-moi donc un peu un rasoir.}}{}
\ligneEst{\textbf{Agathon}}{\francais{\espace{3.8cm}Prends-le toi-même,}}{}
\ligneEst{}{\francais{là, dans l'étui.}}{220}
\ligneEst{\textbf{Euripide}}{\francais{\espace{4cm}Tu es galant homme.}}{}
}
\vspace{2pt}
\noindent\rule{16.5cm}{0.6pt}\\[2pt]
\noindent\textbf{\small Commentaire philologique, métrique et scénique (vv. 208–251 --- partie 2/6)}\\[2pt]
\begin{adjustbox}{max totalheight=13.0cm, max width=16.5cm, keepaspectratio}%
\parbox[t]{16.5cm}{
\setlength{\columnsep}{9pt}%
\begin{multicols}{2}
\fontsize{8.4pt}{9.7pt}\selectfont
\setlength{\parindent}{0pt}%
\setlength{\parskip}{1.2pt}%
\noindent\textbf{230} \textbf{ἔχ᾽ ἀτρέμα σαυτὸν κἀνάκυπτε·} C'est la leçon du \italique{Ravennas}.
\noindent M. van Leeuwen, d'après Dobree, corrige : ἔχ᾽ ἀτρέμας αὐτοῦ, \italique{quia ferri nequit accusativus} ; lequel αὐτοῦ était à sa place en \italique{Oiseaux}, 1200, mais pas du tout ici.
\noindent Hérodote a pourtant dit, en parlant de l'armée athénienne : εἶχον ἀτρέμας σφέας αὐτοὺς ἵνα ἐτάχθησαν\marqueur{1}. Hérodote le pouvait, mais non Aristophane, qui ailleurs s'est borné à ἔχ' ἀτρέμα\marqueur{2}. Quoi ! voilà un comique qui, à quatre reprises (car on s'est donné la peine de compter), a écrit : \italique{reste coi}, et la cinquième fois il se permettrait d'écrire : \italique{tiens-toi coi}.
\noindent À quoi servirait la statistique si l'on tolérait de pareilles licences ?
\noindent (Note vengeresse n°5 de Willems.)
\noindent\textbf{[Notes ad 230]} \textbf{(1)} Hérodote, IX, 54 (aussi III, 79). — De même Platon, \italique{Epist.}, 348 B, écrit : ἐκποδὼν ἑαυτὸν ἔσχεν et Euripide, \italique{Iph. Taur.}, 1129 : ἐκποδὼν τοῦδ' ἔχειν μιάσματος.
\noindent \textbf{(2)} ll est vrai qu'une des quatre fois les bons manuscrits portent ἔχ' ἀτρεμί (\italique{Nuées}, 261). Était-ce assez dégoûtant ? Heureusement que M. van Leeuwen y a mis bon ordre.
\noindent\textbf{232} \textbf{ψιλὸς αὖ στρατεύσομαι.} ψιλός, comme \italique{levis} en latin, signifie à la fois \italique{lisse}, \italique{sans poil}, et \italique{léger}, \italique{armé à la légère}. (Note de Willems.)\\
\noindent Cette note sur ψιλός me sera l'occasion de digresser sur l'un des plus beaux échanges à bâtons rompus de la Tragédie grecque : chez Sophocle (\italique{Ajax}, v. 1120 sq.), entre Ménélas et Teucer :
\citComm{Μενέλαος — Ὁ τοξότης ἔοικεν οὐ σμικρὸν φρονεῖν.\\	
Τεύκρος — Οὐ γὰρ βάναυσον τὴν τέχνην ἐκτησάμην.\\
Μενέλαος — Μέγ᾽ ἄν τι κομπάσειας, ἀσπίδ᾽ εἰ λάβοις.\\
Τεύκρος — Κἂν \textbf{ψιλὸς} ἀρκέσαιμι σοί γ᾽ ὡπλισμένῳ.\\
Μενέλαος — Ἡ γλῶσσά σου τὸν θυμὸν ὡς δεινὸν τρέφει.\\
Τεύκρος — Ξὺν τῷ δικαίῳ γὰρ μέγ᾽ ἔξεστιν φρονεῖν.}
\end{multicols}
}%
\end{adjustbox}%
\vfill
\end{adjustwidth}
\newpage

% --- FOLIO 06 (f. 165r) : SPREAD 2.2 (Texte grec établi + Traduction française + Commentaire 3/6 & 4/6) ---
\clearpage\ifodd\value{page}\hbox{}\newpage\fi
\enTeteG{ΘΕΣΜΟΦΟΡΙΑΖΟΥΣΑΙ --- Texte grec établi (f. 165r, vv. 208–251 [2/3])}
\enTeteD{LES THESMOPHORIES --- Traduction française (f. 165r, vv. 208–251 [2/3])}
\begin{adjustwidth}{-2cm}{-2cm}
\noindent\textbf{\footnotesize Texte grec établi --- Folio 06 (\texttt{f. 165r}, vv. 208–251 [2/3]) \hfill ΠΡΟΛΟΓΟΣ}\\[-2pt]
\noindent\rule{16.5cm}{0.5pt}\\[2pt]
{\fontsize{8.4pt}{9.7pt}\selectfont
\ligneEst{}{κάθιξε· φύσα τὴν γνάθον τὴν δεξιάν.}{}
\ligneEst{\textbf{Μνησίλοχος}}{οἴμοι.}{}
\ligneEst{\textbf{Εὐριπίδης}}{\decalage{οἴμοι.}τί κέκραγας; ἐμβαλῶ σοι πάτταλον,}{}
\ligneEst{}{ἢν μὴ σιωπᾷς.}{}
\ligneEst{\textbf{Μνησίλοχος}}{\decalage{ἢν μὴ σιωπᾷς.}ἀτταταῖ ἰατταταῖ.}{}
\ligneEst{\textbf{Εὐριπίδης}}{οὗτος σὺ ποῖ θεῖς;}{}
\ligneEst{\textbf{Μνησίλοχος}}{\decalage{οὗτος σὺ ποῖ θεῖς;}εἰς τὸ τῶν σεμνῶν θεῶν·}{}
\ligneEst{}{οὐ γὰρ μὰ τὴν Δήμητρά γ᾽ ἐνταυθοῖ μενῶ}{225}
\ligneEst{}{τεμνόμενος.}{}
\ligneEst{\textbf{Εὐριπίδης}}{\decalage{τεμνόμενος.}οὔκουν καταγέλαστος δῆτ᾽ ἔσει}{}
\ligneEst{}{τὴν ἡμίκραιραν τὴν ἑτέραν ψιλὴν ἔχων;}{}
\ligneEst{\textbf{Μνησίλοχος}}{ὀλίγον μέλει μοι.}{}
\ligneEst{\textbf{Εὐριπίδης}}{\decalage{ὀλίγον μέλει μοι.}μηδαμῶς πρὸς τῶν θεῶν}{}
\ligneEst{}{προδῷς με· χώρει δεῦρο.}{}
\ligneEst{\textbf{Μνησίλοχος}}{\decalage{προδῷς με· χώρει δεῦρο.}κακοδαίμων ἐγώ.}{}
\ligneEst{\textbf{Εὐριπίδης}}{ἔχ᾽ ἀτρέμα σαυτὸν κἀνάκυπτε· ποῖ στρέφει;}{230}
\ligneEst{\textbf{Μνησίλοχος}}{μυμῦ.}{}
\ligneEst{\textbf{Εὐριπίδης}}{\decalage{μυμῦ.}τί μύζεις; πάντα πεποίηται καλῶς.}{}
\ligneEst{\textbf{Μνησίλοχος}}{οἴμοι κακοδαίμων, ψιλὸς αὖ στρατεύσομαι.}{}
\ligneEst{\textbf{Εὐριπίδης}}{μὴ φροντίσῃ· ὡς εὐπρεπὴς φανεῖ πάνυ.}{}
\ligneEst{}{βούλει θεάσασθ᾽ αὑτόν;}{}
}
\vspace{2pt}
\noindent\rule{16.5cm}{0.6pt}\\[2pt]
\noindent\textbf{\small Commentaire philologique, métrique et scénique (vv. 208–251 --- partie 3/6)}\\[2pt]
\begin{adjustbox}{max totalheight=13.0cm, max width=16.5cm, keepaspectratio}%
\parbox[t]{16.5cm}{
\setlength{\columnsep}{9pt}%
\begin{multicols}{2}
\fontsize{8.4pt}{9.7pt}\selectfont
\setlength{\parindent}{0pt}%
\setlength{\parskip}{1.2pt}%
\noindent Ménélas s'étonne de ce que Teucer se permette de lui adresser des reproches – à lui le roi de Sparte ! –, vil archer de son état --- Mais l'archerie est loin d'être méprisable lui rétorque Teucer --- Ménélas ironise : Que serait-ce s'il portait un bouclier --- La réplique de Teucer fuse : « Même \italique{ψιλὸς}, je tiendrais bon devant toi et ton [royal] armement ». Comprendre : même \italique{fantassin de l'armée légère} ; c'est-à-dire l'échelon le plus bas chez les militaires – encore plus bas qu'archer –, poste souvent réservé à des esclaves et donc le plus opposé à la royale position de Ménélas et le plus inégal en terme d'armements. C'en est trop pour Ménélas qui s'indigne de cette langue par trop pendue. C'est alors que Teucer lance la fameuse et cinglante réplique : « La fierté est permise, quand on a pour soi la justice ! »
\noindent\textbf{234} \textbf{βούλει θεάσασθαι σεαυτόν; — εἰ δοκεῖ, φέρε.} Le vers péchant contre le mètre, la correction la plus convaincante me semble être celle de Fritzsche qui substitue αὑτόν au σεαυτόν, faisant l'hypothèse que σεαυτόν aura été suggéré à un scribe inattentif par la proximité avec ὁρᾷς σεαυτόν du vers suivant.
\noindent\textbf{235} \textbf{Κλεισθένην.} Clisthène, l'efféminé notoire, personnage qui apparaîtra dans la pièce au vers 571 -- et sera d'ailleurs reconnue comme femme -- avant de prêter allégeance au genre féminin.
\noindent Cette tête de Turc d'Aristophane est raillée comme : une pseudo-femme (\italique{Nuées} 355, \italique{Oiseaux} 829-831, \italique{Thesmophories} 571) ; en tout cas femme en amour (\italique{Lysistrata} 621, 1091-1092, \italique{Thesmophories} 574, \italique{Grenouilles} 48) ; imberbe comme femme (\italique{Acharniens} 117-121, \italique{Grenouilles} 422-424) ; Clisthène ne contrebalance même pas ces innombrables défauts – qui ne sont pas tous du ressort de sa volonté – par une misogynie affectée : il trouve encore le moyen d'être traître à son sexe en s'alliant au genre féminin (\italique{Lysistrata} 621, \italique{Thesmophories} 575-576, 602).
\noindent\textbf{237} \textbf{δελφάκιον} est un diminutif d'ἡ δέλφαξ, le porc \italique{mature}, opposé au χοῖρος, \italique{le porcelet}.
\end{multicols}
}%
\end{adjustbox}%
\vfill
\end{adjustwidth}
\newpage
% --- PAGE DE DROITE : Traduction française + Commentaire (4/6) ---
\begin{adjustwidth}{-2cm}{-2cm}
\noindent\textbf{\footnotesize Traduction française --- Folio 06 (\texttt{f. 165r}, vv. 208–251 [2/3]) \hfill ΠΡΟΛΟΓΟΣ}\\[-2pt]
\noindent\rule{16.5cm}{0.5pt}\\[2pt]
{\fontsize{8.4pt}{9.7pt}\selectfont
\ligneEst{}{\francais{(À Mnésiloque.) Assieds-toi. Enfle la joue droite. (Il le rase.)}}{}
\ligneEst{\textbf{Mnésiloque}}{\francais{Aïe !}}{}
\ligneEst{\textbf{Euripide}}{\francais{\espace{1cm}Qu'as-tu à crier ? Je t'enfonce un bâillon}}{}
\ligneEst{}{\francais{si tu ne te tais.}}{}
\ligneEst{\textbf{Mnésiloque}}{\francais{\espace{2.2cm}Aïe, aïe, aïe ! aie, aïe, aïe !}}{}
\ligneEst{\textbf{Euripide}}{\francais{Holà ! où cours-tu ?}}{}
\ligneEst{\textbf{Mnésiloque}}{\francais{\espace{2.6cm}Au temple des augustes}}{}
\ligneEst{}{\francais{déesses. Car, par Déméter, je ne resterai pas là à me faire}}{225}
\ligneEst{}{\francais{taillader.}}{}
\ligneEst{\textbf{Euripide}}{\francais{\espace{1.7cm}Mais tu seras grotesque}}{}
\ligneEst{}{\francais{avec ta moitié de figure dégarnie de poil.}}{}
\ligneEst{\textbf{Mnésiloque}}{\francais{Peu m'en chaut.}}{}
\ligneEst{\textbf{Euripide}}{\francais{\espace{2.7cm}Au nom des dieux,}}{}
\ligneEst{}{\francais{ne m'abandonne pas. Arrive ici.}}{}
\ligneEst{\textbf{Mnésiloque}}{\francais{\espace{3.6cm}Misère de moi !}}{}
\ligneEst{\textbf{Euripide}}{\francais{Tiens-toi coi et lève la tête. Où te tournes-tu?}}{230}
\ligneEst{\textbf{Mnésiloque}}{\francais{Mumû.}}{}
\ligneEst{\textbf{Euripide}}{\francais{\espace{1.1cm}Pourquoi mumû ? Tout est fini, et selon les règles.}}{}
\ligneEst{\textbf{Mnésiloque}}{\francais{Ah malheur ! Je vais resservir dans l'infanterie légère.}}{}
\ligneEst{\textbf{Euripide}}{\francais{Ne t'inquiète pas. Tu auras un air tout avenant.}}{}
\ligneEst{}{\francais{Veux-tu te mirer ? (Il lui tend un miroir.)}}{}
}
\vspace{2pt}
\noindent\rule{16.5cm}{0.6pt}\\[2pt]
\noindent\textbf{\small Commentaire philologique, métrique et scénique (vv. 208–251 --- partie 4/6)}\\[2pt]
\begin{adjustbox}{max totalheight=13.0cm, max width=16.5cm, keepaspectratio}%
\parbox[t]{16.5cm}{
\setlength{\columnsep}{9pt}%
\begin{multicols}{2}
\fontsize{8.4pt}{9.7pt}\selectfont
\setlength{\parindent}{0pt}%
\setlength{\parskip}{1.2pt}%
\noindent N'en déplaise à certains commentateurs, il n'est fait aucune allusion ici au sexe féminin -- pas plus qu'au vers 1061 de \italique{Lysistrata}. Et d'abord, ce serait absurde : il n'est question que du côté pile ; côté épilé par flambage.
\noindent L'image est ailleurs : c'est celle de de viande de cochon dont on retire les poils par le feu.
\noindent Cf. la \italique{Souda} (δ 205):
\citComm{δελφάκιον : ὁ μικρὸς χοῖρος. Ἀριστοφάνης: οἴμοι κακοδαίμων, δελφάκιον γέγονα. τοῦτο εἴρηκεν ἀφευθεὶς τὸν πρωκτόν. ἀντὶ τοῦ φλογισθείς. μετὰ γὰρ τὸ τυθῆναι τὰ δελφάκια φλογίζονται, ἵνα ψιλωθῶσιν.}
\noindent (On pourra aussi se reporter à Homère\marqueur{1} ou Eschyle\marqueur{2}.)
\noindent\textbf{[Notes ad 237]} \textbf{(1)} \italique{Iliade}, IX, 467-8 et \italique{Odyssée}, II, 300 ; XIV, 73-5
\noindent \textbf{(2)} Fragment 310 de l'édition de S. Radt, \italique{Tragicorum Graecorum Fragmenta}, vol. 3 (1985).
\noindent\textbf{239} \textbf{τὴν κέρκον φυλάττου νυν ἄκραν.} Euripide reprend l'image du δελφάκιον flambé ; il parle bien-sûr de l'\italique{olisbos} de Mnésiloque.
\noindent\textbf{241-2} \textbf{οἴμοι τάλας. ὕδωρ ὕδωρ ὦ γείτονες / πρὶν ἀντιλαβέσθαι πρωκτὸν ἓτερον τῆς φλογός.} Mnésiloque est littéralement en train de prendre feu ; toute sa réplique est un pastiche d'un appel au secours, comme si sa propre maison s'embrasait\marqueur{1} et que celles, contiguës, de ses voisins allaient prendre feu à leur tour (sauf qu'il s'agit ici non d'οἶκός mais de πρωκτός)... il appelle donc à puiser de l'eau « ὕδωρ ! ὕδωρ ! »  Horace prodigue ce conseil dans ses \italique{Épîtres} (1, 18, 84) :
\citComm{\latin{Tua res agitur, paries cum priximus ardet}\\
\francais{C'est ton intérêt qui est en jeu, lorsque la maison de ton voisin brûle\marqueur{2}.}}
\noindent L'image de l'incendie qui se propage de maison en maison était si fréquente dans le monde latin et roman\marqueur{3} qu'on a peine à croire qu'elle ne le fût pas aussi dans le monde grec -- témoin ce passage d'Aristophane.
\end{multicols}
}%
\end{adjustbox}%
\vfill
\end{adjustwidth}
\newpage

% --- FOLIO 06 (f. 165r) : SPREAD 2.3 (Texte grec établi + Traduction française + Commentaire 5/6 & 6/6) ---
\clearpage\ifodd\value{page}\hbox{}\newpage\fi
\enTeteG{ΘΕΣΜΟΦΟΡΙΑΖΟΥΣΑΙ --- Texte grec établi (f. 165r, vv. 208–251 [3/3])}
\enTeteD{LES THESMOPHORIES --- Traduction française (f. 165r, vv. 208–251 [3/3])}
\begin{adjustwidth}{-2cm}{-2cm}
\noindent\textbf{\footnotesize Texte grec établi --- Folio 06 (\texttt{f. 165r}, vv. 208–251 [3/3]) \hfill ΠΡΟΛΟΓΟΣ}\\[-2pt]
\noindent\rule{16.5cm}{0.5pt}\\[2pt]
{\fontsize{8.4pt}{9.7pt}\selectfont
\ligneEst{\textbf{Μνησίλοχος}}{\decalage{βούλει θεάσασθ᾽ αὑτόν;}εἰ δοκεῖ, φέρε.}{}
\ligneEst{\textbf{Εὐριπίδης}}{ὁρᾷς σεαυτόν;}{235}
\ligneEst{\textbf{Μνησίλοχος}}{\decalage{ὁρᾷς σεαυτόν;}οὐ μὰ Δί᾽ ἀλλὰ Κλεισθένη.}{}
\ligneEst{\textbf{Εὐριπίδης}}{ἀνίστασ᾽, ἵν᾽ ἀφεύσω σε, κἀγκύψας ἔχε.}{}
\ligneEst{\textbf{Μνησίλοχος}}{οἴμοι κακοδαίμων δελφάκιον γενήσομαι.}{}
\ligneEst{\textbf{Εὐριπίδης}}{ἐνεγκάτω τις ἔνδοθεν δᾷδ᾽ ἢ λύχνον.}{}
\ligneEst{}{ἐπίκυπτε· τὴν κέρκον φυλάττου νυν ἄκραν.}{}
\ligneEst{\textbf{Μνησίλοχος}}{ἐμοὶ μελήσει νὴ Δία, πλήν γ᾽ ὅτι κάομαι.}{240}
\ligneEst{}{οἴμοι τάλας. ὕδωρ ὕδωρ ὦ γείτονες}{}
\ligneEst{}{πρὶν ἀντιλαβέσθαι πρωκτὸν ἓτερον τῆς φλογός.}{}
\ligneEst{\textbf{Εὐριπίδης}}{θάρρει.}{}
\ligneEst{\textbf{Μνησίλοχος}}{\decalage{θάρρει.}τί θαρρῶ καταπεπυρπολημένος;}{}
\ligneEst{\textbf{Εὐριπίδης}}{ἀλλ᾽ οὐκ ἔτ᾽ οἰδὲν πρᾶγμά σοι· τὰ πλεῖστα γὰρ}{}
\ligneEst{}{ἀποπεπόνηκας.}{245}
\ligneEst{\textbf{Μνησίλοχος}}{\decalage{ἀποπεπόνηκας.}φεῦ ἰοὺ τῆς ἀσβόλου.}{}
\ligneEst{}{αἰθὸς γεγένημαι πάντα τὰ περὶ τὴν τράμιν.}{}
\ligneEst{\textbf{Εὐριπίδης}}{μὴ φροντίσῃς· ἕτερος γὰρ αὐτὰ σπογγιεῖ.}{}
\ligneEst{\textbf{Μνησίλοχος}}{οἰμώξετ᾽ ἆρ᾽ εἴ τις τὸν ἐμὸν πρωκτὸν πλυνεῖ.}{}
\ligneEst{\textbf{Εὐριπίδης}}{Ἀγάθων, ἐπειδὴ σαυτὸν ἐπιδοῦναι φθονεῖς,}{}
\ligneEst{}{ἀλλ᾽ ἱμάτιον γοῦν χρῆσον ἡμῖν τουτῳὶ}{250}
\ligneEst{}{καὶ στρόφιον· οὐ γὰρ ταῦτά γ᾽ ὡς οὐκ ἔστ᾽ ἐρεῖς.}{251}
}
\vspace{2pt}
\noindent\rule{16.5cm}{0.6pt}\\[2pt]
\noindent\textbf{\small Commentaire philologique, métrique et scénique (vv. 208–251 --- partie 5/6)}\\[2pt]
\begin{adjustbox}{max totalheight=13.0cm, max width=16.5cm, keepaspectratio}%
\parbox[t]{16.5cm}{
\setlength{\columnsep}{9pt}%
\begin{multicols}{2}
\fontsize{8.4pt}{9.7pt}\selectfont
\setlength{\parindent}{0pt}%
\setlength{\parskip}{1.2pt}%
\noindent\textbf{[Notes ad 241-2]} \textbf{(1)} La scholie au vers 242 : (πρωκτὸν·) ὡς εἰ ἔλεγεν ``οἰκίαν''.
\noindent \textbf{(2)} J'emprunte la traduction au magistral \italique{Dictionnaire des sentences latines et grecques} de Renzo Tosi (traduit de l'italien par Rebecca Lenoir). La consultation de ce dictionnaire me fut d'une grande aide pour découvrir, en ce passage comme en d'autres, des ponts entre la langue d'Aristophane et l'univers des proverbes.
\noindent \textbf{(3)} Voir le \italique{Dictionnaire} sus-cité (§524).
\noindent\textbf{243} \textbf{θάρρει — τί θαρρῶ;} Plaute, \italique{Curculio}, IV, 3, 21 : \latin{bellator vale — quid valeam ?}
\noindent\textbf{247} \textbf{μὴ φροντίσῃς· ἕτερος γὰρ αὐτὰ σπογγιεῖ.} Les Grecs se servaient d'éponges pour la toilette et, dans un usage détourné, pouvaient même les utiliser en guise de torche-cul comme cela ressort d'une scène des \italique{Grenouilles} (vv. 481 sqq.).
\noindent Qui est cet ἕτερος ? Réponse : le voisin d'à côté : « La situation est bien celle qu'on attend : en cas d'incendie, les voisins, soucieux de leur propre intérêt, commençaient par faire la chaîne pour apporter de l'eau. Puis, le feu une fois éteint, il se trouvait presque toujours un voisin complaisant pour aider le propriétaire (ou le locataire) malchanceux à nettoyer sa maison et à y remettre de l'ordre. » (Taillardat\marqueur{1}.)
\noindent\textbf{[Notes ad 247]} \textbf{(1)} Jean Taillardat, «~Aristophane, \italique{Thesmophories} v. 247, 289, 721 sqq. : essai d'interprétation du texte transmis~», in \italique{Revue des Études Grecques}, t. 85, fasc. 406-408, juillet-décembre 1972, p. XVII.
\noindent Le grand Taillardat avait donc compris le sens de ce passage avant même que le \nieme{242} vers, lacunaire, ne fût rétabli par la découverte d'un papyrus.
\noindent\textbf{248} \textbf{οἰμώξετ᾽ ἄρ᾽ εἴ τις τὸν ἐμὸν πρωκτὸν πλυνεῖ.} πλύνειν signifie au propre \italique{nettoyer} (et fait donc écho au σπογγιεῖ du vers précédent) et par suite \italique{dire du mal} (\italique{passer un savon} en français ; l'on manque toujours de la délicatesse requise lorsqu'on s'occupe de la toilette d'autrui).
\end{multicols}
}%
\end{adjustbox}%
\vfill
\end{adjustwidth}
\newpage
% --- PAGE DE DROITE : Traduction française + Commentaire (6/6) ---
\begin{adjustwidth}{-2cm}{-2cm}
\noindent\textbf{\footnotesize Traduction française --- Folio 06 (\texttt{f. 165r}, vv. 208–251 [3/3]) \hfill ΠΡΟΛΟΓΟΣ}\\[-2pt]
\noindent\rule{16.5cm}{0.5pt}\\[2pt]
{\fontsize{8.4pt}{9.7pt}\selectfont
\ligneEst{\textbf{Mnésiloque}}{\francais{\espace{3.5cm}Si c'est ton plaisir. Passe.}}{}
\ligneEst{\textbf{Euripide}}{\francais{Te vois-tu ?}}{235}
\ligneEst{\textbf{Mnésiloque}}{\francais{\espace{2.4cm}Non, pas Zeus. Je vois Clisthène.}}{}
\ligneEst{\textbf{Euripide}}{\francais{Lève-toi, que je te flambe. Tiens-toi penché.}}{}
\ligneEst{\textbf{Mnésiloque}}{\francais{Ah malheur ! Je vais être cochon de lait.}}{}
\ligneEst{\textbf{Euripide}}{\francais{Qu'on m'apporte de dedans une torche ou une lampe.}}{}
\ligneEst{}{\francais{Courbe-toi, et gare au bout de la queue.}}{}
\ligneEst{\textbf{Mnésiloque}}{\francais{J'y veillerai, par Zeus... N'empêche que je brûle.}}{240}
\ligneEst{}{\francais{Ah misère ! de l'eau, de l'eau, voisins,}}{}
\ligneEst{}{\francais{avant que la fesse d'à côté ne s'enflamme.}}{}
\ligneEst{\textbf{Euripide}}{\francais{Du sang froid.}}{}
\ligneEst{\textbf{Mnésiloque}}{\francais{\espace{1.4cm}Du sang froid ? quand je suis dévoré par les flammes !..}}{}
\ligneEst{\textbf{Euripide}}{\francais{Mais tu n'as plus à t'inquiéter. Car te voilà dépêtré}}{}
\ligneEst{}{\francais{du principal.}}{245}
\ligneEst{\textbf{Mnésiloque}}{\francais{\espace{2.3cm}Phu ! Ah! cette suie !}}{}
\ligneEst{}{\francais{Je suis tout feu dans l'entrefesson.}}{}
\ligneEst{\textbf{Euripide}}{\francais{Ne t'inquiète pas. Le voisin d'à côté t'épongera cela.}}{}
\ligneEst{\textbf{Mnésiloque}}{\francais{Il en coûtera à celui qui s'attaquera à mon derrière.}}{}
\ligneEst{\textbf{Euripide}}{\francais{Agathon, puisque tu refuses de te dévouer à moi,}}{}
\ligneEst{}{\francais{prête-nous du moins un vêtement}}{250}
\ligneEst{}{\francais{et un strophion. Tu ne nieras pas que tu en possèdes..}}{251}
}
\vspace{2pt}
\noindent\rule{16.5cm}{0.6pt}\\[2pt]
\noindent\textbf{\small Commentaire philologique, métrique et scénique (vv. 208–251 --- partie 6/6)}\\[2pt]
\begin{adjustbox}{max totalheight=13.0cm, max width=16.5cm, keepaspectratio}%
\parbox[t]{16.5cm}{
\setlength{\columnsep}{9pt}%
\begin{multicols}{2}
\fontsize{8.4pt}{9.7pt}\selectfont
\setlength{\parindent}{0pt}%
\setlength{\parskip}{1.2pt}%
\noindent Le sel de la boutade de Mnésiloque nous échappe...
\noindent Mnésiloque souhaite-t-il bien du courage à qui voudrait \italique{s'attaquer} -- au prime sens de πλύνειν : \italique{commencer de nettoyer} -- à son fondement ?
\noindent Ingrate tâche qui serait aussitôt à recommencer. Travail de Sisyphe qu'on pourrait rapprocher du proverbe grec πλίνθον πλύνειν, « laver une brique »\marqueur{}. Dans le même genre d'idée, Aristophane emploie au vers 604 des \italique{Guêpes} l'expression qui paraît proverbiale πρωκτὸς λουτροῦ περιγιγνόμενος, « le c.. l'emporte sur l'eau lustrale »\marqueur{2}.
\noindent\textbf{[Notes ad 248]} \textbf{(1)} Voir Photius ο 360 (\italique{s.v} ὄνου πόκοι), ο 363 (ὄνου πόκαι) et l'apparat conséquent des \latin{testimonia} et autres \latin{loci paralleli} de l'édition Theodoridis. Le proverbe fut répertorié entre autres par Plutarque, \italique{Corpus paroemiographorum Graecorum}, I, περὶ τῶν ἀδυνάτων ; Diogenianus (VII, 50), Zenobius (VI, 48). Le scholiaste des \italique{Guêpes} (280b) le cite : <λίθον ἕψεις·> ὡς καὶ τὸ `\textbf{πλίνθον πλύνεις}' καὶ `χύτραν ποικίλλεις' καὶ `εἰς ὕδωρ γράφεις' καὶ `Αἰθίοπα λευκαίνεις' καὶ `κατὰ θαλάττης σπείρεις' καὶ τὰ τοιαῦτα. C'est d'ailleurs en écho à ce proverbe qu'Aristophane rapproche au vers 514 du \italique{Ploutos} les verbes πλινθουργεῖν et πλύνειν.
\noindent \textbf{(2)} Bdélycléon lançait cette sentence à la barbe de Philocléon, saoul de discours et qui, quelque effort qu'il fît pour soigner sa glossolalite, n'en retomberait pas moins un jour dans son travers.
\noindent Le proverbe devait donc peu ou prou signifier qu'un naturel grossier, le πρωκτός, retournait inévitablement à sa fange ; en français : un \italique{indécrottable}. Il est repris tel quel par les parémiographes Diogenianus (III, 53) et Michel Apostolius (XIV, 78), Hésychius (π 4144), Photius (π 1450), Suidas (π 2951). La source principale doit être d'anciens commentaires au vers 604 des \italique{Guêpes}. (Cf. les scholies aux manuscrits d'Aristophane ; l'on y apprend que Καλλίστρατος lui-même avait signalé l'expression comme proverbiale.)
\end{multicols}
}%
\end{adjustbox}%
\vfill
\end{adjustwidth}
\newpage
\clearpage\ifodd\value{page}\hbox{}\newpage\fi
\enTeteG{ΠΡΟΛΟΓΟΣ --- Codex Ravennas 429, f. 165v (Classense 167v) --- vv. 252–294}
\begin{adjustwidth}{-1.5cm}{-3.0cm}
\begin{adjustbox}{max totalheight=21.0cm, max width=16.5cm, keepaspectratio}
\begin{minipage}{16.5cm}
{\fontsize{7.2pt}{8.5pt}\selectfont
\noindent\textbf{Folio 07 --- Codex Ravennas 429, \texttt{f. 165v} (\texttt{Classense 167v}) : ΠΡΟΛΟΓΟΣ, vv. 252–294}\hfill \textit{Témoins :} $\text{R}^1, \Rpc, \Rdeux$, \PsiUn\par\vspace{2pt}
\hrule height 0.4pt\vspace{2.5pt}
}
{\fontsize{8.1pt}{9.5pt}\selectfont
\ligneRav{252}{ἀγάθ/}{λαμβάνετε καὶ χρῆσθ᾽· οὐ φθονῶ. : τί οὖν λάβω;}{\textbf{252} (Εὐρ.) R : \lecon{(Μνη.)} Bentley}
\ligneRav{253}{\auDessus{εὐ}{ρ}}{ὅ τι; τὸν κροκωτὸν πρῶτον ἐνδύου λαβών·}{\textbf{253} (Ἀγάθ.) R : \lecon{(Εὐρ.)} Dobree}
\ligneRav{254}{μ\auDessus{ν}{Η}}{νὴ τὴν Ἀφροδίτην ἡδύ γ᾽ ὄζει ποσθίου·}{\textbf{254} (Εὐρ.) R}
\ligneRav{255}{ἀγάθ/}{σύζωσον ἀνύσας. : αἶρε νῦν στρόφιον. : ἰδού·}{\textbf{255} (Ἀγάθ.) et dicolon ante αἶρε praefixit R \espApp (Εὐρ.) ἰδού — (Μνη.) ἴθι ... σκέλη R}
\ligneRav{256}{μ\auDessus{ν}{Η}}{ἴθι νυν κατάστειλόν με τὰ περὶ τὼ \difference{σκέλει}·}{\textbf{256} σκέλη R : \lecon{σκέλει} van Leeuwen}
\ligneRav{}{\auDessus{εὐ}{ρ}}{κεκρυφάλου δεῖ καὶ μίτρας. : ἡδὶ μὲν οὖν}{}
\ligneRav{}{}{κεφαλὴ περίθετος, ἣν ἐγὼ νύκτωρ φορῶ·}{}
\ligneRav{}{\auDessus{εὐ}{ρ}}{νὴ τὸν Δί᾽ ἀλλὰ κἀπιτηδεία πάνυ·}{}
\ligneRav{260}{μ\auDessus{ν}{Η}}{\difference{ἆρ᾽ ἁρμόσει} μοι; : νὴ Δί᾽ ἀλλ᾽ ἄριστ᾽ ἔχει·}{\textbf{260} ἠραρμοσηι R : \lecon{ἆρ᾽ ἁρμόσει} Küster \espApp (Ἀγάθ.) R : \lecon{(Εὐρ.)} Ellebodius}
\ligneRav{261}{—}{φέρ᾽ ἔγκυκλον· : \difference{τουτὶ} \difference{λάβ᾽} ἀπὸ τῆς κλινίδος·}{\textbf{261} τοῦτὶ R : \lecon{τουτὶ} Suid. ε 135, κ 1809 \espApp λάμβαν᾽ R, Suid. : \lecon{λάβ᾽} Daubuz, Küster}
\ligneRav{}{\auDessus{εὐ}{ρ}}{ὑποδημάτων δεῖ. : τἀμὰ ταυτὶ λάμβανε·}{}
\ligneRav{263}{μ\auDessus{ν}{Η}}{χὤπως μὲν ἀνδρείως· : Ἀγάθωνος γ᾽ ο\auDessus{ῦ}{ν} ἐστὶν·}{\textbf{263} μοι R : \lecon{μοι;} Küster \espApp \lecon{(Ἀγάθ.)} χαλαρὰ ... φορῶν R : (Μνη.) Brunck, edd. post. \espApp γ᾽ο\auDessus{ῦ}{ν} \Rpc \espApp χαίρεις R, Suid. A,G,M : \lecon{χαίροις} Suid. F,V}
\ligneRav{263}{μ\auDessus{ν}{Η}}{χὤπως μὲν ἀνδρείως· : Ἀγάθωνος γ᾽ ο\auDessus{ῦ}{ν} ἐστὶν·}{\textbf{263} μοι R : \lecon{μοι;} Küster \espApp \lecon{(Ἀγάθ.)} χαλαρὰ ... φορῶν R : (Μνη.) Brunck, edd. post. \espApp γ᾽ο\auDessus{ῦ}{ν} \Rpc \espApp χαίρεις R, Suid. A,G,M : \lecon{χαίροις} Suid. F,V}
\ligneRav{264}{}{εἴσω τις ὡς τάχιστά μ᾽ ἐσκυκλησάτω·}{\textbf{264} nulla nota mut. pers. R : \lecon{(Ἀγάθ.)} ἀλλ᾽ ... Bibi}
\ligneRav{266}{\auDessus{εὐ}{ρ}}{\difference{ἁνὴρ} μὲν ἡμῖν οὑτοσὶ καὶ δὴ γυνὴ}{\textbf{266} ἀνὴρ R : \lecon{ἁνὴρ} Fritzsche}
\ligneRav{267}{}{\difference{τό γ᾽} εἶδος· ἢν λαλῇς δ᾽, ὅπως τῷ φθέγματι}{\textbf{267} τόδ᾽ R : \lecon{τό γ᾽} Küster}
\ligneRav{}{—}{γυναικιεῖς εὖ καὶ πιθανῶς. : πειράσομαι·}{}
\ligneRav{}{\auDessus{εὐ}{ρ}}{βάδιζε τοίνυν. : μὰ τὸν Ἀπόλλω οὔκ, ἤν γε μὴ}{}
\ligneRav{270}{—}{ὀμόσῃς ἐμοί— : τί χρῆμα; : συσσώσειν ἐμὲ}{}
\ligneRav{}{}{πάσαις τέχναις, ἤν μοί τι περιπίπτῃ κακόν·}{}
\ligneRav{}{\auDessus{εὐ}{ρ}}{`ὄμνυμι τοίνυν αἰθέρ᾽ οἴκησιν Διός.'}{}
\ligneRav{273}{μ\auDessus{ν}{Η}}{τί μᾶλλον ἢ τὴν Ὑποκράτους ξυνοικίαν;}{\textbf{273} \lecon{Ὑποκράτους} R : Ἱπποκράτους Σ}
\ligneRav{}{\auDessus{εὐ}{ρ}}{ὄμνυμι τοίνυν πάντας ἄρδην τοὺς θεούς·}{}
\ligneRav{275}{μ\auDessus{ν}{Η}}{μέμνησο τοίνυν ταῦθ᾽, ὅτι ἡ φρὴν ὤμοσεν,}{}
\ligneRav{}{}{ἡ γλῶττα δ᾽ οὐκ ὀμώμοκ᾽· οὐδ᾽ ὥρκωσ᾽ ἐγώ·}{}
\ligneRav{}{}{\hspace{2.2em}\textbf{[παρεπιγραφή R\expo{1}, \PsiUn{} :]} \lecon{ὀλολύζουσι τε ἱερὸν ὠθεῖται}}{\textbf{post 276} \lecon{ὀλολύζουσι τε ἱερὸν ὠθεῖται} R, ολολυ[ \PsiUn, $\Sigma^{\text{R}}$}
\ligneRav{277}{\auDessus{εὐ}{ρ}}{ἔκσπευδε ταχέως· ὡς τὸ τῆς ἐκκλησίας}{\textbf{277} \lecon{ἔκσπευδε} R : εἄ· σπεῦδε Maas}
\ligneRav{278}{}{σημεῖον ἐν τῷ Θεσμοφορίῳ φαίνεται·}{\textbf{278} \lecon{Θεσμοφορίῳ} R, cf. \italique{IG} II\expo{2} 1177.24 ; 2498.12 : Θεσμοφόρῳ Suid. σ 325 : Θεσμοφορείῳ Scaliger}
\ligneRav{}{—}{ἐγὼ δ᾽ ἄπειμι. : δεῦρό νυν ὦ Θρᾷτθ᾽ ἕπου·}{}
\ligneRav{280}{}{ὦ Θρᾷττα θέασαι, καομένων τῶν λαμπάδων}{}
\ligneRav{}{}{ὅσον τὸ χρῆμ᾽ ἀνέρχεθ᾽ ὑπὸ τῆς λιγνύος·}{}
\ligneRav{282}{}{ἀλλ᾽, ὦ \difference{περικαλλεῖ} \difference{Θεσμοφόρω}, δέξασθέ με}{\textbf{282} Θεσμοφόρῳ R : \lecon{Θεσμοφόρω} Photius α 1386 \espApp περὶ καλλῆ R : \lecon{περικαλλεῖ} Kaibel}
\ligneRav{283}{}{ἀγαθῇ τύχῃ \difference{καὶ δεῦρο καὶ πάλιν} οἴκαδε·}{\textbf{283} καὶ δεῦρο πάλιν R : \lecon{καὶ δεῦρο καὶ πάλιν} Bentley}
\ligneRav{}{}{ὦ Θρᾷττα τὴν κίστην κάθελε, κᾆτ᾽ ἔξελε}{}
\ligneRav{285}{}{\difference{τὰ πόπαν᾽}, ὅπως λαβοῦσα θύσω ταῖν θεαῖν·}{\textbf{285} \lecon{τὰ πόπαν᾽} \PsiUn : τὸ πόπανον R \espApp \lecon{τὰ πόπαν᾽} \PsiUn, Porson : τὸ πόπανον R}
\ligneRav{}{}{δέσποινα πολυτίμητε Δήμητερ φίλη}{}
\ligneRav{287}{}{καὶ \difference{Φερρέφαττα}, πολλὰ πολλάκις μέ σοι}{\textbf{287} \lecon{Φερρέφαττα} \PsiUn, Thiersch : Φερσέφαττα R}
\ligneRav{}{}{θύειν ἔχουσαν, εἰ δὲ μή, ἀλλὰ νῦν λαθεῖν}{}
\ligneRav{289}{}{καὶ τὴν θυγατέρα \difference{Χοιρίλ'} ἀνδρός μοι τυχεῖν}{\textbf{289} χοῖρον R : χοιρίον Σ : \lecon{Χοιρίλ'} Bibi}
\ligneRav{290}{}{πλουτοῦντος, ἄλλως τ᾽ ἠλιθίου κἀβελτέρου,}{}
\ligneRav{291}{}{καὶ \difference{πρὸς θυλάκιον} νοῦν ἔχειν μοι καὶ φρένας·}{\textbf{291} πρὸς θάληκον R : Ποσθαλίσκον Dindorf : \lecon{πρὸς θυλάκιον} Bibi}
\ligneRav{}{}{ποῦ ποῦ καθίζωμ᾽ ἐν καλῷ, τῶν ῥητόρων}{}
\ligneRav{}{}{ἵν᾽ ἐξακούω; σὺ δ᾽ ἄπιθ᾽ ὦ Θρᾷττ᾽ ἐκποδών·}{}
\ligneRav{}{}{δούλοις γὰρ οὐκ ἔξεστ᾽ ἀκούειν τῶν λόγων·}{}
}
\vspace{3pt}\noindent\rule{4.5cm}{0.4pt}\par\vspace{1.5pt}
{\fontsize{7.0pt}{8.3pt}\selectfont
\noindent\textbf{Loci paralleli (vv. 252–294) ---} \textbf{254} Suidas π 2114 \espTest (ποσθίου) Photius π 1113 πόσθιον· αἰδοῖον· ἢ (sic) τὴν Ἀφροδίτην ; Hesychius π 3096 \espTest \textbf{257-8} (ἡδὶ -- φορῶ) Suidas π 1165 (sed ἥδιον pro ἡδὶ) \espTest \textbf{258} Photius π 690 περίθετος κεφαλή ; \Eustathe 1280.63 \espTest (κεφαλὴ περίθετος) Pollux II 35 \espTest \textbf{261} Suidas ε 135 (sed τουτὶ pro τι) \espTest (ἔγκυκλον -- κλινίδος) Suidas κ 1809 \espTest \textbf{263} Suidas χ 7 χαλαρὰ γοῦν χαίρεις φορῶν ὑποδήματα (sed χαίροις codd. FV) \espTest \textbf{274} (ἄρδην) Hesychius α 7097 \espTest \textbf{275-6} cf. \Eustathe 443.37 \espTest \textbf{277-8} (ὡς -- φαίνεται) Suidas σ 325 (sed Θεσμοφόρῳ pro Θεσμοφορίῳ) \espTest \textbf{282} Photius α 1386 ; \AelDion α 113 \espTest \textbf{285} (πόπαν᾽) \Eustathe 437.4 \espTest \textbf{289} (χοῖρον) Euripides, testimonia 4, 18 Kannicht Χοιρίλη \espTest \textbf{291} (θάληκον) Hesychius θ 50 θάλλικα ; Pollux X 152 θυλάκιον [...] ἐπὶ τοῦ αὐτοῦ δὲ καὶ σακίον, ὅταν φῇ ἐν Θεσμοφοριαζούσαις \espTest \textbf{292} (τῶν ῥητόρων) Photius ρ 106 ῥήτορα· καὶ τὴν γυναῖκα. Ἀριστοφάνης\par
}
\vspace{2.5pt}\noindent\rule{4.5cm}{0.4pt}\par\vspace{1.5pt}
{\fontsize{6.9pt}{8.2pt}\selectfont
\noindent\textbf{Suite (en commentaire) de l'apparat critique \& notes paléographiques (\texttt{f. 165v}) ---} \textbf{Témoin papyrologique :} \textbf{\PsiUn{} (\italique{PSI} XI 1194)} couvre les \textbf{vv. 273–288}. --- \textbf{Troisième \italique{parepigraphè} (post v. 276) :} R écrit dans la colonne \lecon{ὀλολύζουσι τε ἱερὸν ὠθεῖται}, avec une \textbf{\italique{koronis}} dans la marge gauche et la scholie $\Sigma^{\text{R}}$ : παρεπιγραφή. ἐκκυκλεῖται ἐπὶ τὸ ἔξω τὸ θεσμοφόριον. Le papyrus \PsiUn{} confirme ολολυ[. --- \textbf{Abréviation supralinéaire au v. 263 :} γ' ο\auDessus{ῦ}{ν} \Rpc{}. --- \textbf{Deux conjectures majeures de \textbf{Bibi} (vv. 289 et 291) :} Au v. 289, \lecon{Χοιρίλ'} (nom de la fille de Mnésiloque et épouse d'Euripide, corrompu en χοῖρον dans R et χοιρίον dans $\Sigma^{\text{R}}$) ; au v. 291, \lecon{πρὸς θυλάκιον} (ou πρὸς θάλικα ap. Hésychius $\theta\ 50$) pour πρὸς θάληκον de R. --- \textbf{v. 260 :} ἠραρμοσηι R (iotacisme et médivision) : \lecon{ἆρ' ἁρμόσει} Küster \espApp (Ἀγάθ.) R : \lecon{(Εὐρ.)} Ellebodius. --- \textbf{v. 263 :} μοι R : \lecon{μοι;} Küster \espApp \textbf{(Ἀγάθ.) χαλαρὰ ... φορῶν} R (maintenu par \textbf{Bibi} contre (Μνη.) Brunck) \espApp \lecon{γ' ο\auDessus{ῦ}{ν}} \Rpc{} \espApp χαίρεις R, Souda AGM : \lecon{χαίροις} Souda FV. --- \textbf{v. 264 :} nulla nota mut. pers. R : \lecon{(Ἀγάθ.)} ἀλλ' ... \lecon{Bibi}. --- \textbf{v. 273 :} \lecon{Ὑποκράτους} R (maintenu par \textbf{Bibi} : calembour comique sur $\text{ὑπό}$, contre la normalisation Ἱπποκράτους de $\Sigma^{\text{R}}$). --- \textbf{v. 289 :} χοῖρον R (χοιρίον $\Sigma$) : \lecon{Χοιρίλ'} (ou Χοιρίν') \lecon{Bibi} (fille de Mnésiloque et épouse d'Euripide). --- \textbf{v. 291 :} πρὸς θάληκον R (Ποσθαλίσκον Dindorf) : \lecon{πρὸς θυλάκιον} (ou πρὸς θάλικα ap. Hésychius $\theta\ 50$) \lecon{Bibi} (par métathèse).\par
}
\end{minipage}
\end{adjustbox}
\end{adjustwidth}
\newpage
\enTeteD{Codex Ravennas 429 (Biblioteca Classense) --- Fac-similé du f. 165v (vv. 252–294)}
\begin{adjustwidth}{-1.5cm}{-3.0cm}
\begin{center}
\includegraphics[width=16.5cm,height=20.5cm,keepaspectratio]{/workspaces/Thesmophories/dossier_ravennas_thesmophories/folios_individuels_ff162v-177r/07_f165v_classense167v_vv0252-0294.jpg}\par
\vspace{2pt}
{\scriptsize\textbf{Codex Ravennas 429, f. 165v (Classense 167v)} --- ΠΡΟΛΟΓΟΣ, \textbf{vv. 252–294} ($\text{R}^1, \Rpc, \Rdeux$, \PsiUn).}
\end{center}
\end{adjustwidth}
\newpage
\clearpage\ifodd\value{page}\hbox{}\newpage\fi% --- FOLIO 07 (f. 165v) : SPREAD 2.1 (Texte grec établi + Traduction française + Commentaire 1/8 & 2/8) ---
\clearpage\ifodd\value{page}\hbox{}\newpage\fi
\enTeteG{ΘΕΣΜΟΦΟΡΙΑΖΟΥΣΑΙ --- Texte grec établi (f. 165v, vv. 252–294 [1/4])}
\enTeteD{LES THESMOPHORIES --- Traduction française (f. 165v, vv. 252–294 [1/4])}
\begin{adjustwidth}{-2cm}{-2cm}
\noindent\textbf{\footnotesize Texte grec établi --- Folio 07 (\texttt{f. 165v}, vv. 252–294 [1/4]) \hfill ΠΡΟΛΟΓΟΣ}\\[-2pt]
\noindent\rule{16.5cm}{0.5pt}\\[2pt]
{\fontsize{8.8pt}{10.2pt}\selectfont
\ligneEst{\textbf{Ἀγάθων}}{λαμβάνετε καὶ χρῆσθ᾽· οὐ φθονῶ.}{252}
\ligneEst{\textbf{Μνησίλοχος}}{\decalage{λαμβάνετε καὶ χρῆσθ᾽· οὐ φθονῶ.}τί οὖν λάβω;}{}
\ligneEst{\textbf{Εὐριπίδης}}{ὅ τι; τὸν κροκωτὸν πρῶτον ἐνδύου λαβών.}{}
\ligneEst{\textbf{Μνησίλοχος}}{νὴ τὴν Ἀφροδίτην ἡδύ γ᾽ ὄζει ποσθίου.}{}
\ligneEst{\textbf{Ἀγάθων}}{σύζωσον ἀνύσας.}{255}
\ligneEst{\textbf{Εὐριπίδης}}{\decalage{σύζωσον ἀνύσας.}αἶρε νῦν στρόφιον.}{}
\ligneEst{\textbf{Ἀγάθων}}{\decalage{σύζωσον ἀνύσας. αἶρε νῦν στρόφιον.}ἰδού.}{}
\ligneEst{\textbf{Μνησίλοχος}}{ἴθι νυν κατάστειλόν με τὰ περὶ τὼ σκέλει.}{}
\ligneEst{\textbf{Εὐριπίδης}}{κεκρυφάλου δεῖ καὶ μίτρας.}{}
\ligneEst{\textbf{Ἀγάθων}}{\decalage{κεκρυφάλου δεῖ καὶ μίτρας.}ἡδὶ μὲν οὖν}{}
\ligneEst{}{κεφαλὴ περίθετος, ἣν ἐγὼ νύκτωρ φορῶ.}{}
\ligneEst{\textbf{Εὐριπίδης}}{νὴ τὸν Δί᾽ ἀλλὰ κἀπιτηδεία πάνυ.}{}
\ligneEst{\textbf{Μνησίλοχος}}{ἆρ᾽ ἁρμόσει μοι;}{260}
\ligneEst{\textbf{Εὐριπίδης}}{\decalage{ἆρ᾽ ἁρμόσει μοι;}νὴ Δί᾽ ἀλλ᾽ ἄριστ᾽ ἔχει.}{}
}
\vspace{2pt}
\noindent\rule{16.5cm}{0.6pt}\\[2pt]
\noindent\textbf{\small Commentaire philologique, métrique et scénique (vv. 252–294 --- partie 1/8)}\\[2pt]
\begin{adjustbox}{max totalheight=15.2cm, max width=16.5cm, keepaspectratio}%
\parbox[t]{16.5cm}{
\setlength{\columnsep}{9pt}%
\begin{multicols}{2}
\fontsize{7.6pt}{8.8pt}\selectfont
\setlength{\parindent}{0pt}%
\setlength{\parskip}{1.2pt}%
\noindent\textbf{255} \textbf{σύζωσον ἀνύσας --- αἶρε νῦν στρόφιον.} La répartition de ce vers comme de tout ce passage est très changeante d'une édition à une autre.
\noindent J'ai essayé pour ma part de respecter le découpage des vers du manuscrit de Ravenne, sans tailler ni fusionner çà et là des répliques qui ne l'étaient pas originellement.
\noindent J'ai confié à Euripide les conseils d'habillement, Agathon continuant d'adopter une attitude froide : il se contente de tendre à Mnésiloque les vêtements pour ne pas que ce rustre aille puiser directement dans sa garde-robe.
\noindent La principale question qui se soit posée est celle de la première partie du vers 255 : σύζωσον ἀνύσας --- αἶρε νῦν στρόφιον.
\noindent Il comporte un \italique{dicôlon} entre ἀνύσας et αἶρε sans pour autant indiquer de changement de locuteur... Ce \italique{dicôlon} est-il là pour signifier une simple pause le temps que Mnésiloque passe la crocote ? ou y a-t-il véritablement changement de locuteur ?
\noindent `σύζωσον ἀνύσας' me semble être une réplique au ton sec, marquant la pointe d'agacement qu'Agathon aurait pu sentir à la remarque de Mnésiloque du vers 254. Aussi lui ai-je accordé cette réplique ; puis Euripide de demander le \italique{strophion} (αἶρε νῦν στρόφιον) et Agathon de le tendre à Mnésiloque (ἰδού).
\noindent\textbf{258} \textbf{κεφαλὴ περίθετος, ἣν ἐγὼ νύκτωρ φορῶ.} De tout l'attirail nécessaire pour se métamorphoser en femme, il est étonnant que ce soit le plus évident à priori qui ait rencontré le plus de difficultés : c'est Fritzsche qui le premier s'avisa que la \grec{κεφαλὴ περίθετος} désignait une \italique{perruque}. Hé oui ! car une résille et un bandeau ne suffiraient à cacher l'absence volumique d'une belle crinière, d'où la suggestion avisée d'Agathon – en homme habitué à de tels travestissements. Ce sens est clairement attesté par Pollux (X, 170) :
\citComm{\grec{πηνήκην δὲ καὶ προκόμιον, καὶ ὡς Ἀριστφάνης ἐν Γηρυτάδῃ περίθεσιν· καὶ περιθέτην δέ, ὡς ἐν Ἄμφιδος Ἀλκμαίωνι καὶ Μενάνδρου Ὀλυνθίᾳ.}}
\noindent Les synonymes \grec{πηνήκη} et \grec{προκόμιον} ne souffrent aucune ambiguïté.\\ Le tacle d'Aristophane envers les mœurs agathoniennes – qui lui fait préciser : ἣν ἐγὼ νύκτωρ φορῶ – n'est pas clair... le scholiaste à l'imagination débridée -- qui n'était probablement pas plus renseigné que nous sur la question -- y voyait une allusion aux chasses masquées contre-nature de notre poètereau dans l'interlope \italique{Athens by night}. Cette rêverie ne trouve qu'un faible écho au vers 35 (et s'appuie peut-être sur une transmission fautive du texte, comme nous en avons fait l'hypothèse \italique{ad loc.}).
\end{multicols}
}%
\end{adjustbox}%
\vfill
\end{adjustwidth}
\newpage
% --- PAGE DE DROITE : Traduction française + Commentaire (2/8) ---
\begin{adjustwidth}{-2cm}{-2cm}
\noindent\textbf{\footnotesize Traduction française --- Folio 07 (\texttt{f. 165v}, vv. 252–294 [1/4]) \hfill ΠΡΟΛΟΓΟΣ}\\[-2pt]
\noindent\rule{16.5cm}{0.5pt}\\[2pt]
{\fontsize{8.8pt}{10.2pt}\selectfont
\ligneEst{\textbf{Agathon}}{\francais{Prenez et usez-en ; je ne m'y oppose pas.}}{252}
\ligneEst{\textbf{Mnésiloque}}{\francais{\espace{5cm}Que faut-il prendre ?}}{}
\ligneEst{\textbf{Euripide}}{\francais{Quoi ? prends d'abord cette crocote et revêts-la.}}{}
\ligneEst{\textbf{Mnésiloque}}{\francais{Par Aphrodite, elle fleure bon le priape.}}{}
\ligneEst{\textbf{Agathon}}{\francais{Serre-t'en vite la taille.}}{255}
\ligneEst{\textbf{Euripide}}{\francais{\espace{2.5cm}Maintenant passe-nous un strophion.}}{}
\ligneEst{\textbf{Agathon}}{\francais{\espace{5.5cm}Là.}}{}
\ligneEst{\textbf{Mnésiloque}}{\francais{Maintenant arrange-moi les jambes.}}{}
\ligneEst{\textbf{Euripide}}{\francais{Une résille et un bandeau ?}}{}
\ligneEst{\textbf{Agathon}}{\francais{\espace{4.2cm}Plutôt cette}}{}
\ligneEst{}{\francais{perruque, que je porte la nuit.}}{}
\ligneEst{\textbf{Euripide}}{\francais{Eh ! par Zeus, c'est même tout juste l'affaire.}}{}
\ligneEst{\textbf{Mnésiloque}}{\francais{M'ira-t-elle ?}}{260}
\ligneEst{\textbf{Euripide}}{\francais{\espace{2.7cm}Mais c'est parfait, par Zeus.}}{}
}
\vspace{2pt}
\noindent\rule{16.5cm}{0.6pt}\\[2pt]
\noindent\textbf{\small Commentaire philologique, métrique et scénique (vv. 252–294 --- partie 2/8)}\\[2pt]
\begin{adjustbox}{max totalheight=15.2cm, max width=16.5cm, keepaspectratio}%
\parbox[t]{16.5cm}{
\setlength{\columnsep}{9pt}%
\begin{multicols}{2}
\fontsize{7.6pt}{8.8pt}\selectfont
\setlength{\parindent}{0pt}%
\setlength{\parskip}{1.2pt}%
\noindent\textbf{261} \textbf{φέρ᾽ ἔγκυκλον.} Se reporter à la note au vers 643 pour un développement sur l'\italique{encycle}. Ce mot désigne une sorte de \italique{jupe courte}.
\noindent\textbf{263} \textbf{ἆρ᾽ ἁρμόσει μοι Ἀγ. χαλαρὰ γ᾽ο\auDessus{ῦ}{ν} χαίρεις φορῶν.} Mnésiloque, devant les chaussures qu'il trouve étroites --- « Mais m'iront-elles ? » --- Agathon « Tu préférerais (χαίροις Suid.) en chausser de larges... »\\ Rien de plus clair : Agathon « n'aurait eu garde de porter des souliers larges, car c'eût été faire preuve de d'inélégance et même de rusticité [...] Mnésiloque, en vieillard qui aime ses aises, est donc fondé à craindre que les souliers ne soient trop étroits pour lui. »\marqueur{1}.
\noindent Seulement, Brunck eut l'idée saugrenue de confier l'intégralité du vers 263 à Mnésiloque alors que le \italique{Ravennas} indiquait clairement le passage de témoin à Agathon pour la seconde partie, χαλαρὰ γο\auDessus{ῦ}{ν} χαίρεις φορῶν. Ainsi, c'est Mnésiloque qui prononcerait « tu aimes donc à te chausser large »... allusion aux mœurs sodomotiques d'Agathon. En voilà bien assez pour que tous les éditeurs postérieurs suivent en meute Brunck.
\noindent Le scholiaste était le premier à avoir ouvert le bal des affreux, qui glosait : <χαλαρὰ·>χαῦνα. διαβάλλει δὲ πάλιν τὸν Ἀγάθωνα ὡς χαῦνον.
\noindent Personne ne s'est avisé que si Mnésiloque avait effectivement lancé le trait à sous-entendu obscène qu'on lui impute, alors jamais Agathon ne lui aurait répondu avec bienveillance `σὺ τοῦτο γίγνωσκ᾽', « à toi de t'en aviser ».
\noindent Mais cette engeance de fouines à ordures perd tout sens critique dès qu'elle se met à flairer du graveleux.
\noindent Willems l'avait noté : Mnésiloque se trouve désappointé devant l'étroitesse des élégants souliers d'Agathon ; cette paire lui ira-t-elle : ἆρ᾽ ἁρμόσει μοι; se demande-t-il.
\noindent À ces mots, Agathon ne se sent pas de joie : ce rustre en préférerait sans doute une bien large ; c'est bien Agathon qui prononce `χαλαρὰ γοῦν χαίροις φορῶν'.
\noindent Et c'est à ce moment seulement que Mnésiloque saisit la perche ; il taquine une nouvelle fois Agathon sur ses mœurs supposées :
\citComm{Μνησίλοχος --- σὺ τοῦτο γίγνωσκ᾽.\\
Mnésiloque --- Ah ça ! tu en sais quelque chose...}
\noindent C'en est trop pour la patience d'Agathon : excédé, il décide de se retirer superbement avec l'eccyclème (exactement comme autrefois l'Euripide des \italique{Acharniens}).
\noindent\textbf{[Notes ad 263]} \textbf{(1)} Se reporter encore une fois au grand Alphonse Willems ; la citation est tirée de sa \nieme{6} note, reproduite à la fin.
\end{multicols}
}%
\end{adjustbox}%
\vfill
\end{adjustwidth}
\newpage

% --- FOLIO 07 (f. 165v) : SPREAD 2.2 (Texte grec établi + Traduction française + Commentaire 3/8 & 4/8) ---
\clearpage\ifodd\value{page}\hbox{}\newpage\fi
\enTeteG{ΘΕΣΜΟΦΟΡΙΑΖΟΥΣΑΙ --- Texte grec établi (f. 165v, vv. 252–294 [2/4])}
\enTeteD{LES THESMOPHORIES --- Traduction française (f. 165v, vv. 252–294 [2/4])}
\begin{adjustwidth}{-2cm}{-2cm}
\noindent\textbf{\footnotesize Texte grec établi --- Folio 07 (\texttt{f. 165v}, vv. 252–294 [2/4]) \hfill ΠΡΟΛΟΓΟΣ}\\[-2pt]
\noindent\rule{16.5cm}{0.5pt}\\[2pt]
{\fontsize{8.8pt}{10.2pt}\selectfont
\ligneEst{}{φέρ᾽ ἔγκυκλον·}{}
\ligneEst{\textbf{Ἀγάθων}}{\decalage{φέρ᾽ ἔγκυκλον·}τουτὶ λάβ᾽ ἀπὸ τῆς κλινίδος.}{}
\ligneEst{\textbf{Εὐριπίδης}}{ὑποδημάτων δεῖ.}{}
\ligneEst{\textbf{Ἀγάθων}}{\decalage{ὑποδημάτων δεῖ.}τἀμὰ ταυτὶ λάμβανε.}{}
\ligneEst{\textbf{Μνησίλοχος}}{ἆρ᾽ ἁρμόσει μοι;}{}
\ligneEst{\textbf{Ἀγάθων}}{\decalage{ἆρ᾽ ἁρμόσει μοι;}χαλαρὰ γοῦν χαίροις φορῶν.}{}
\ligneEst{\textbf{Μνησίλοχος}}{σὺ τοῦτο γίγνωσκ᾽}{}
\ligneEst{\textbf{Ἀγάθων}}{\decalage{σὺ τοῦτο γίγνωσκ᾽}ἀλλ᾽ ἔχεις γὰρ ὧν δέει,}{}
\ligneEst{}{εἴσω τις ὡς τάχιστά μ᾽ ἐσκυκλησάτω.}{265}
\ligneEst{\textbf{Εὐριπίδης}}{ἁνὴρ μὲν ἡμῖν οὑτοσὶ καὶ δὴ γυνὴ}{}
\ligneEst{}{τό γ᾽ εἶδος· ἢν λαλῇς δ᾽, ὅπως τῷ φθέγματι}{}
\ligneEst{}{γυναικιεῖς εὖ καὶ πιθανῶς.}{}
\ligneEst{\textbf{Μνησίλοχος}}{\decalage{γυναικιεῖς εὖ καὶ πιθανῶς.}πειράσομαι.}{}
\ligneEst{\textbf{Εὐριπίδης}}{βάδιζε τοίνυν.}{}
}
\vspace{2pt}
\noindent\rule{16.5cm}{0.6pt}\\[2pt]
\noindent\textbf{\small Commentaire philologique, métrique et scénique (vv. 252–294 --- partie 3/8)}\\[2pt]
\begin{adjustbox}{max totalheight=15.2cm, max width=16.5cm, keepaspectratio}%
\parbox[t]{16.5cm}{
\setlength{\columnsep}{9pt}%
\begin{multicols}{2}
\fontsize{7.6pt}{8.8pt}\selectfont
\setlength{\parindent}{0pt}%
\setlength{\parskip}{1.2pt}%
\noindent\textbf{265} \textbf{εἴσω τις ὡς τάχιστά μ᾽ ἐσκυκλησάτω.} De même que l'Euripide des \italique{Acharniens} (vers 409) demande à être roulé par l'eccyclème pour sortir de la scène : il ne lui agrée pas de descendre (καταβαίνειν, cf. \italique{infra} v. 1215) :
\citComm{ἀλλ' ἐκκυκλήσομαι· καταβαίνειν δ᾽ οὐ σχολή.}
\noindent\textbf{272} \textbf{ὄμνυμι τοίνυν αἰθέρ᾽ οἴκησιν Διός.} Selon le scholiaste, ce vers est une citation de la \italique{Mélanippe [la sage]} d'Euripide\marqueur{1} (ἐκ τῆς Μελανίππης Εὐριπίδου) ; citation qui devait être fameuse à la façon dont Aristophane l'évoque encore aux vers 100 et 311 des \italique{Grenouilles}.
\noindent\textbf{[Notes ad 272]} \textbf{(1)} \nieme{8} fragment de l'édition Jouan-van Looy aux Belles Lettres ; fr. 487 de l'édition Kannicht ; 670 chez Mette.
\noindent\textbf{273} \textbf{τί μᾶλλον ἢ τὴν Ὑποκράτους ξυνοικίαν;} Au Ciel de qui se moque Euripide ?! Jurer sur l'Éther ? Et pourquoi pas `τὴν Ὑποκράτους ξυνοικίαν' (c'est-à-dire « par une porcherie ») ?\\ « Les fils d'Hippocratès étaient connus de tout le monde à Athènes dès l'époque des \italique{Nuées} (v. 1001). Car leur père étant fils d'Ariphron, propre frère de Périclès, ils étaient les petits neveux du fameux homme d'État. On leur avait donné le surnom de porcs à cause de leur \italique{balourdise} (ὑομουσία, \italique{Cavaliers}, v. 986). Les textes sont unanimes sur ce point. Athénée (III, 96e) :
\citComm{ἑξῆς ἐπεισηνέχθη, μητρόπολίς τις ὡς ἀληθῶς οὖσα καὶ μήτηρ τῶν Ἱπποκράτους υἱῶν, οὓς εἰς ὑωδίαν κωμῳδουμένους οἶδα.\\
On servit ensuite une vulve de truie, mère des fils d'Hippocratès, lesquels, je le sais, ont été raillés par les comiques à cause de leur cochonnerie\marqueur{1}.}
\noindent Cela se trouve confirmé par le scholiaste d'Aristophane sur les \italique{Nuées}, v. 1001, et par Suidas. De même Photius (σ 845), au mot σῦς :
\citComm{ὗς τοὺς Ἱπποκράτους υἱοὺς ἔλεγον καὶ τοὺς Παναιτίου καὶ Μέμνονος εἰς ὑηνίαν κωμῳδοῦντες.\\
On nommait porcs les fils d'Hippocratès ainsi que ceux de Panaitios et de Memnon. Les comiques les raillent à cause de leur cochonnerie.}
\noindent (Autres textes dans Meineke, \italique{Fr. Com. Gr.}, t. II, p. 477. [Pour une bibliographie complète et mise à jour, voir les \italique{testimonia} de l'édition Theodoridis de Photius à l'adresse σ 845.]) Le logis de famille d'Hippocratès est donc une porcherie. Il est à croire que l'auteur ne se faisait pas faute de prononcer συῳκία. » (Note de Willems.) \\ Je conserve la leçon du manuscrit de Ravenne : \grec{Ὑποκράτους}, sans reprendre la correction communément admise \grec{Ἱπποκράτους}. En effet, le préfixe \grec{ὑ}- sert à construire des composés sur le mot \italique{porc} ; ainsi, pour se borner aux textes cités par Willems, on recense \grec{ὑωδία} (littéralement, la \italique{ressemblance à un porc}) et \grec{ὑηνία} ( littéralement \italique{nature de porc}).
\end{multicols}
}%
\end{adjustbox}%
\vfill
\end{adjustwidth}
\newpage
% --- PAGE DE DROITE : Traduction française + Commentaire (4/8) ---
\begin{adjustwidth}{-2cm}{-2cm}
\noindent\textbf{\footnotesize Traduction française --- Folio 07 (\texttt{f. 165v}, vv. 252–294 [2/4]) \hfill ΠΡΟΛΟΓΟΣ}\\[-2pt]
\noindent\rule{16.5cm}{0.5pt}\\[2pt]
{\fontsize{8.8pt}{10.2pt}\selectfont
\ligneEst{}{\francais{Un encycle ?}}{}
\ligneEst{\textbf{Agathon}}{\francais{\espace{2.4cm}Prends celui-là, sur la couchette.}}{}
\ligneEst{\textbf{Euripide}}{\francais{Des chaussures ?}}{}
\ligneEst{\textbf{Agathon}}{\francais{\espace{2.6cm}Voici les miennes, prends-les.}}{}
\ligneEst{\textbf{Mnésiloque}}{\francais{Est-ce qu'elles m'iront ?}}{}
\ligneEst{\textbf{Agathon}}{\francais{\espace{3.6cm}Tu préférerais en chausser de larges...}}{}
\ligneEst{\textbf{Mnésiloque}}{\francais{Ça tu en sais quelque chose.}}{}
\ligneEst{\textbf{Agathon}}{\francais{\espace{3.1cm} Or çà, tu as ce qu'il te faut.}}{}
\ligneEst{}{\francais{Qu'on me roule au plus vite à l'intérieur.}}{265}
\ligneEst{\textbf{Euripide}}{\francais{Voilà notre homme devenu femme,}}{}
\ligneEst{}{\francais{du moins en apparence. Si tu parles, tâche par ton timbre de voix}}{}
\ligneEst{}{\francais{de simuler la femme, avec aisance et naturel.}}{}
\ligneEst{\textbf{Mnésiloque}}{\francais{\espace{4.1cm}(Voix de fausset.) J'y tâcherai.}}{}
\ligneEst{\textbf{Euripide}}{\francais{Va maintenant.}}{}
}
\vspace{2pt}
\noindent\rule{16.5cm}{0.6pt}\\[2pt]
\noindent\textbf{\small Commentaire philologique, métrique et scénique (vv. 252–294 --- partie 4/8)}\\[2pt]
\begin{adjustbox}{max totalheight=15.2cm, max width=16.5cm, keepaspectratio}%
\parbox[t]{16.5cm}{
\setlength{\columnsep}{9pt}%
\begin{multicols}{2}
\fontsize{8.4pt}{9.7pt}\selectfont
\setlength{\parindent}{0pt}%
\setlength{\parskip}{1.2pt}%
\noindent En prononçant : \grec{τὴν ὑ-ποκράτου-ς-(σ)υ-ῳκίαν}, Mnésiloque insiste doublement sur la porcinité de l'hippocratique maison.
\noindent Voilà aussi pourquoi les enfants qui l'habitaient sont plus souvent qualifiés d'« \grec{υ-ἱοί Ἱπποκράτους} » que de « παῖδες Ἱπποκράτους ».
\noindent\textbf{[Notes ad 273]} \textbf{(1)} Comprendre : \italique{grossièreté}, \italique{lourdeur}. ὗς, le \italique{porc} est pour les Grecs le symbole de l'ignorance, partant, de la sottise. Voir à ce sujet le toujours bon Taillardat, §451. L'adjectif Ὑάντιος servait à désigner les Béotiens, peuples à l'éternelle réputation de lourdeur et d'ignorance, cf. l'adage n° 906 d'Érasme : ἡ βοιωτία ὗς.
\noindent\textbf{275-6} \textbf{ὅτι ἡ φρὴν ὤμοσεν / ἡ γλῶττα δ᾽ οὐκ ὀμώμοκ᾽· οὐδ᾽ ὥρκωσ᾽ ἐγώ.} Mnésiloque se méfie avec raison d'Euripide-le-Retors, il a notamment en tête le vers 612 de son \italique{Hippolyte}\marqueur{1} :
\citComm{ἡ γλῶσσ' ὀμώμοχ', ἡ δὲ φρὴν ἀνώμοτος.}
\noindent (La citation d'Euripide comportait déjà 7 `o' (dont 4 d'affilée) ; avec Aristophane, le compteur explose !)
\noindent\textbf{post 276} \textbf{ὀλολύζουσι τε ἱερὸν ὠθεῖται.} Le vers 276 est suivi de l'indication scénique : ὀλολύζουσι τε ἱερὸν ὠθεῖται ainsi que d'une κορωνίς en marge, qui marque le changement de scène.
\noindent Le scholiaste confirme : παρεπιγραφή. ἐκκυκλεῖται ἐπὶ τὸ ἔξω τὸ θεσμοφόριον et l'on trouve déjà ολολυ[ sur un fragment de papyrus datant du second siècle après Jésus-Christ.\\ La maison d'Agathon est donc roulée par l'eccyclème au fond pour se transformer en temple des Thesmophores ; est hissé le \italique{signal}, le σημεῖον du vers 278 tandis que le chœur des femmes fait retentir l'ὀλολυγή ; c'est à-dire l'espèce de hurlement sacré qu'Hérodote (IV, 189) faisait remonter aux Libyens.
\noindent (Cette ululation dont il est question déjà dans Homère (\italique{Iliade}, VI, 301), et qui retentissait dans les temples en signe tantôt de douleur tantôt de joie.)
\noindent\textbf{[Notes ad 275-6]} \textbf{(1)} Cf. Érasme, \italique{Adages} n° 1441, \latin{Lingua juravit} et Tosi, \italique{Dictionnaire des sentences latines et grecques}, §2252.
\noindent\textbf{280} \textbf{Θρᾷττα} est un nom courant pour les esclaves Thraces (comme dans les \italique{Acharniens}).
\noindent\textbf{281} \textbf{ὑπὸ τῆς λιγνύος.} Sur le sens de ὑπὸ : \italique{avec}, \italique{accompagné de} (souvent en parlant d'un accompagnement musical. Cf. Photius, \italique{Bibliothèque} 320b : οἱ παλαιοὶ τὴν ὑπό ἀντὶ τῆς μετά πολλάκις ἐλάμβανον), voir la note d'Alphonse Willems n°7, écrite d'une juste colère (et reproduite à la fin).
\end{multicols}
}%
\end{adjustbox}%
\vfill
\end{adjustwidth}
\newpage

% --- FOLIO 07 (f. 165v) : SPREAD 2.3 (Texte grec établi + Traduction française + Commentaire 5/8 & 6/8) ---
\clearpage\ifodd\value{page}\hbox{}\newpage\fi
\enTeteG{ΘΕΣΜΟΦΟΡΙΑΖΟΥΣΑΙ --- Texte grec établi (f. 165v, vv. 252–294 [3/4])}
\enTeteD{LES THESMOPHORIES --- Traduction française (f. 165v, vv. 252–294 [3/4])}
\begin{adjustwidth}{-2cm}{-2cm}
\noindent\textbf{\footnotesize Texte grec établi --- Folio 07 (\texttt{f. 165v}, vv. 252–294 [3/4]) \hfill ΠΡΟΛΟΓΟΣ}\\[-2pt]
\noindent\rule{16.5cm}{0.5pt}\\[2pt]
{\fontsize{8.8pt}{10.2pt}\selectfont
\ligneEst{\textbf{Μνησίλοχος}}{\decalage{βάδιζε τοίνυν.}μὰ τὸν Ἀπόλλω οὔκ, ἤν γε μὴ}{}
\ligneEst{}{ὀμόσῃς ἐμοί—}{270}
\ligneEst{\textbf{Εὐριπίδης}}{\decalage{ὀμόσῃς ἐμοί—}τί χρῆμα;}{}
\ligneEst{\textbf{Μνησίλοχος}}{\decalage{ὀμόσῃς ἐμοί—τί χρῆμα;}συσσώσειν ἐμὲ}{}
\ligneEst{}{πάσαις τέχναις, ἤν μοί τι περιπίπτῃ κακόν.}{}
\ligneEst{\textbf{Εὐριπίδης}}{`ὄμνυμι τοίνυν αἰθέρ᾽ οἴκησιν Διός.'}{}
\ligneEst{\textbf{Μνησίλοχος}}{τί μᾶλλον ἢ τὴν Ὑποκράτους ξυνοικίαν;}{}
\ligneEst{\textbf{Εὐριπίδης}}{ὄμνυμι τοίνυν πάντας ἄρδην τοὺς θεούς.}{}
\ligneEst{\textbf{Μνησίλοχος}}{μέμνησο τοίνυν ταῦθ᾽, ὅτι ἡ φρὴν ὤμοσεν,}{275}
\ligneEst{}{ἡ γλῶττα δ᾽ οὐκ ὀμώμοκ᾽· οὐδ᾽ ὥρκωσ᾽ ἐγώ.}{}
\ligneEst{\textbf{Εὐριπίδης}}{ἔκσπευδε ταχέως· ὡς τὸ τῆς ἐκκλησίας}{}
\ligneEst{}{σημεῖον ἐν τῷ Θεσμοφορίῳ φαίνεται.}{}
\ligneEst{}{ἐγὼ δ᾽ ἄπειμι.}{}
\ligneEst{\textbf{Μνησίλοχος}}{\decalage{ἐγὼ δ᾽ ἄπειμι.}δεῦρό νυν ὦ Θρᾷτθ᾽ ἕπου.}{}
}
\vspace{2pt}
\noindent\rule{16.5cm}{0.6pt}\\[2pt]
\noindent\textbf{\small Commentaire philologique, métrique et scénique (vv. 252–294 --- partie 5/8)}\\[2pt]
\begin{adjustbox}{max totalheight=15.2cm, max width=16.5cm, keepaspectratio}%
\parbox[t]{16.5cm}{
\setlength{\columnsep}{9pt}%
\begin{multicols}{2}
\fontsize{7.6pt}{8.8pt}\selectfont
\setlength{\parindent}{0pt}%
\setlength{\parskip}{1.2pt}%
\noindent\textbf{288} \textbf{εἰ δὲ μή, ἀλλὰ νῦν λαθεῖν.} Mnésiloque, qui déclame sa prière à haute voix afin de mystifier l'assemblée des femmes, prononce en aparté cette réplique.
\noindent\textbf{289} \textbf{καὶ τὴν θυγατέρα χοῖρον ἀνδρός μοι τυχεῖν.} Les vers 288 à 291 sont corrompus en plusieurs endroits. Occasion trop belle pour les philologues de nous infliger leur saleté morale. Et ils s'en sont donné à cœur joie : Mnésiloque userait d'obscénités pour parler de sa fille (τὴν θυγατέρα χοῖρον ; χοῖρον, \italique{le porc}, allusion selon eux au sexe de la femme), quand il ne lui souhaiterait tout bonnement pas les mœurs les plus infamantes.
\noindent Et ce au beau milieu d’une prière aux deux déesses parmi les plus révérées !.. Juste courroux d'Alphonse Willems\marqueur{1}. (Allez donc chercher dans tout Aristophane la moindre irrévérence à Déméter ou Coré.) On ne le répétera jamais assez : « Mais pour qui prennent-ils Aristophane ? »  Las ! on ne détruit que ce que l'on remplace et le coup d'épée de Willems – qu'il avait pourtant aiguisée – porta dans l'eau car sa note était essentiellement aporétique. Prions les déesses thesmophores de conférer à la nôtre assez de force pour dissiper à tout jamais ces nuées.
\noindent Le nom ou surnom de la fille de Mnésiloque, †χοῖρον†, est évidemment corrompu.
\noindent La preuve en est que scholiaste avait une autre leçon devant les yeux : χοιρίον\marqueur{2}.
\noindent N'allons pas inventer un sobriquet obscène ; reportons-nous plutôt aux quelques éléments biographiques qui nous sont parvenus\marqueur{3}. Mnésiloque est le beau-père d'Euripide, et la fille qu'il lui a cédé en mariage s'appelle justement Χοιρίλη\marqueur{4} ; la correction est toute indiquée : la présence d'un nombre propre (rare) a naturellement tendance à induire des fautes de transmission du texte ; de plus, la coquille consistant à substituer un Ν à un Λ est monnaie courante.
\noindent Mnésiloque profite donc de cette prière fictive pour tacler son gendre Euripide et souhaiter à sa fille Chœrile un mari riche et de savoir gérer à bon escient le porte-monnaie conjugal -- Euripide avait-il la réputation d'être mauvais économe ? (cf. la note ci-dessous.)
\noindent\textbf{291} \textbf{καὶ πρὸς θάληκον νοῦν ἔχειν μοι καὶ φρένας.} « Pour ce qui du †θάληκον†, qu'elle ait esprit et intelligence. » Le vers 291 a fait l'objet de nombreux débats ; et pour cause : le mot θάληκον est inconnu au bataillon des lexiques.
\noindent Une scholie en marge du manuscrit de Ravenne se référant probablement au \nieme{291} vers indique :
\citComm{τὸν παιδαρίσκον. ἴσως δὲ παρὰ τὴν πόσθην αὐτὸ συνέθηκεν,\\
le jeune garçon. Peut-être à rapprocher du pénis (πόσθη).}
\end{multicols}
}%
\end{adjustbox}%
\vfill
\end{adjustwidth}
\newpage
% --- PAGE DE DROITE : Traduction française + Commentaire (6/8) ---
\begin{adjustwidth}{-2cm}{-2cm}
\noindent\textbf{\footnotesize Traduction française --- Folio 07 (\texttt{f. 165v}, vv. 252–294 [3/4]) \hfill ΠΡΟΛΟΓΟΣ}\\[-2pt]
\noindent\rule{16.5cm}{0.5pt}\\[2pt]
{\fontsize{8.8pt}{10.2pt}\selectfont
\ligneEst{\textbf{Mnésiloque}}{\francais{\espace{2.3cm}Non, par Apollon, à moins que}}{}
\ligneEst{}{\francais{tu ne me jures...}}{270}
\ligneEst{\textbf{Euripide}}{\francais{\espace{2.3cm}Quoi ?}}{}
\ligneEst{\textbf{Mnésiloque}}{\francais{\espace{4cm}D'aider à me sauver}}{}
\ligneEst{}{\francais{par toutes tes rubriques, s'il me tombe dessus quelque malheur.}}{}
\ligneEst{\textbf{Euripide}}{\francais{Eh bien, « je jure, par l'Éther, la demeure de Zeus ».}}{}
\ligneEst{\textbf{Mnésiloque}}{\francais{Pourquoi pas aussi bien par le logis de famille d'Hippocratès ?}}{}
\ligneEst{\textbf{Euripide}}{\francais{Eh bien! je jure par tous les dieux en bloc.}}{}
\ligneEst{\textbf{Mnésiloque}}{\francais{Au surplus, souviens-toi de ceci : c'est ton cœur qui a juré,}}{275}
\ligneEst{}{\francais{et non ta langue ; je ne la tiens pas pour liée.}}{}
\ligneEst{\textbf{Euripide}}{\francais{Hâte-toi de partir. Car le signal de l'assemblée}}{}
\ligneEst{}{\francais{paraît sur le Thesmophorion. Moi, je m'en vais.}}{}
\ligneEst{}{\francais{(Tous sortent. Le décor du fond change et représente le Thesmophorion. Entrent Mnésiloque suivi d'une servante portant une ciste.)}}{}
\ligneEst{\textbf{Mnésiloque}}{\francais{\espace{2.3cm}Par ici, Thratta, suis-moi.}}{}
}
\vspace{2pt}
\noindent\rule{16.5cm}{0.6pt}\\[2pt]
\noindent\textbf{\small Commentaire philologique, métrique et scénique (vv. 252–294 --- partie 6/8)}\\[2pt]
\begin{adjustbox}{max totalheight=15.2cm, max width=16.5cm, keepaspectratio}%
\parbox[t]{16.5cm}{
\setlength{\columnsep}{9pt}%
\begin{multicols}{2}
\fontsize{8.4pt}{9.7pt}\selectfont
\setlength{\parindent}{0pt}%
\setlength{\parskip}{1.2pt}%
\noindent Les commentateurs s'en sont remis à cette scholie pour en induire que πρὸς θάληκον était une transmission fautive d'un surnom de garçon, forgé sur le mot πόσθη, \italique{le membre}.
\noindent La reconstruction encore la plus en vogue est celle de Dindorf : πρὸς θάληκον $ \to $ Ποσθαλίσκον, où Ποσθαλίσκον désignerait le fils (fictif) de Mnésiloque.
\noindent À ce ποσθαλίσκον, Willems a opposé deux raisons\marqueur{1}, que je crois fortes et fondées : « il est plus que douteux que ce dernier mot [Ποσθαλίσκον] soit grec, les diminutifs se faisant en ισκος et non en αλισκος. (On cite, il est vrai, κωραλίσκος, παιδαρίσκος et μειρακίσκος. Ce serait concluant s'il n'y avait pas l'intermédiaire de κωράλιον, παιδάριον et μειράκιον\marqueur{2}. Si des diminutifs comme ποσθάλιον et ποσθαλίσκος avaient existé, ils n'auraient pu manquer d'être relevés par d'anciens lexicographes.) » À l'école de Willems, j'ai aussi appris à me méfier du scholiaste\marqueur{3}, cet « ignare byzantin », et faire bien davantage confiance à l'« alexandrin » Hésychius. Autres lieux, autres temps et autres sciences !
\noindent Comme l'a déjà fait remarquer Alphonse Willems, Hésychius n'a jamais compulsé de ποσθαλίσκον ou autre composé de πόσθη pouvant expliquer notre πρὸς θάληκον.
\noindent En revanche, l'on trouve chez Hésychius à l'adresse θ 50 :
\citComm{θάλλικα· σάκκου εἶδος\\
Thallika : sorte de sac.}
\noindent Voilà qui est beaucoup plus proche du texte transmis dans le \italique{Ravennas}, θάληκον.
\noindent Pour le sens, j'ai reçu l'aide précieuse de Jean-Victor Vernhes qui « estime qu'il s'agit du sac où cette femme met des bijoux, la clef de sa maison, ses emplettes, sa carte de crédit (et bientôt, vous allez voir, celle de son mari), et qu'il symbolise d'une façon générale, ses affaires, ses intérêts, au profit desquels elle exploite son nigaud de conjoint. Elle finira par tout fourrer dans son sac ! »\marqueur{4}
\noindent Thème cher à Aristophane : confier aux femmes les clefs de la cité tout autant que celles du \italique{portefeuille du ménage} ; pour ces choses-là, elles ont beaucoup plus d'esprit et d'intelligence, νοῦν ἔχειν καὶ φρένας. Il y a de nombreux passages de l'\italique{Assemblée des Femmes} ou de \italique{Lysistrata} qui feraient écho à ce genre de prière.
\noindent Néanmoins, le mot θάλλικα (de θάλλιξ ?) est rare et circonscrit au lexique d'Hésychius. Il semble être une métathèse d'un mot beaucoup plus courant et qu'on trouve chez Aristophane : θυλάκιον.
\end{multicols}
}%
\end{adjustbox}%
\vfill
\end{adjustwidth}
\newpage

% --- FOLIO 07 (f. 165v) : SPREAD 2.4 (Texte grec établi + Traduction française + Commentaire 7/8 & 8/8) ---
\clearpage\ifodd\value{page}\hbox{}\newpage\fi
\enTeteG{ΘΕΣΜΟΦΟΡΙΑΖΟΥΣΑΙ --- Texte grec établi (f. 165v, vv. 252–294 [4/4])}
\enTeteD{LES THESMOPHORIES --- Traduction française (f. 165v, vv. 252–294 [4/4])}
\begin{adjustwidth}{-2cm}{-2cm}
\noindent\textbf{\footnotesize Texte grec établi --- Folio 07 (\texttt{f. 165v}, vv. 252–294 [4/4]) \hfill ΠΡΟΛΟΓΟΣ}\\[-2pt]
\noindent\rule{16.5cm}{0.5pt}\\[2pt]
{\fontsize{8.4pt}{9.7pt}\selectfont
\ligneEst{}{ὦ Θρᾷττα θέασαι, καομένων τῶν λαμπάδων}{280}
\ligneEst{}{ὅσον τὸ χρῆμ᾽ ἀνέρχεθ᾽ ὑπὸ τῆς λιγνύος.}{}
\ligneEst{}{ἀλλ᾽, ὦ περικαλλεῖ Θεσμοφόρω, δέξασθέ με}{}
\ligneEst{}{ἀγαθῇ τύχῃ καὶ δεῦρο καὶ πάλιν οἴκαδε.}{}
\ligneEst{}{ὦ Θρᾷττα τὴν κίστην κάθελε, κᾆτ᾽ ἔξελε}{}
\ligneEst{}{τὰ πόπαν᾽, ὅπως λαβοῦσα θύσω ταῖν θεαῖν.}{285}
\ligneEst{}{δέσποινα πολυτίμητε Δήμητερ φίλη}{}
\ligneEst{}{καὶ Φερρέφαττα, πολλὰ πολλάκις μέ σοι}{}
\ligneEst{}{θύειν ἔχουσαν, εἰ δὲ μή, ἀλλὰ νῦν λαθεῖν}{}
\ligneEst{}{καὶ τὴν θυγατέρα Χοιρίλ' ἀνδρός μοι τυχεῖν}{}
\ligneEst{}{πλουτοῦντος, ἄλλως τ᾽ ἠλιθίου κἀβελτέρου,}{290}
\ligneEst{}{καὶ πρὸς θυλάκιον νοῦν ἔχειν μοι καὶ φρένας.}{}
\ligneEst{}{ποῦ ποῦ καθίζωμ᾽ ἐν καλῷ, τῶν ῥητόρων}{}
\ligneEst{}{ἵν᾽ ἐξακούω; σὺ δ᾽ ἄπιθ᾽ ὦ Θρᾷττ᾽ ἐκποδών.}{}
\ligneEst{}{δούλοις γὰρ οὐκ ἔξεστ᾽ ἀκούειν τῶν λόγων.}{294}
}
\vspace{2pt}
\noindent\rule{16.5cm}{0.6pt}\\[2pt]
\noindent\textbf{\small Commentaire philologique, métrique et scénique (vv. 252–294 --- partie 7/8)}\\[2pt]
\begin{adjustbox}{max totalheight=13.0cm, max width=16.5cm, keepaspectratio}%
\parbox[t]{16.5cm}{
\setlength{\columnsep}{9pt}%
\begin{multicols}{2}
\fontsize{7.6pt}{8.8pt}\selectfont
\setlength{\parindent}{0pt}%
\setlength{\parskip}{1.2pt}%
\noindent Si J.-V. Verhnes avait dû condenser son propos en grec, il aurait écrit quelque chose comme : σακίον, ἐν οἷσπερ τἀργύριον ταμιεύεται. Hé bien, c'est exactement la glose qu'on retrouve dans un passage de l'\italique{Onomasticon} de Pollux citant... les \italique{Thesmophories} d'Aristophane.
\noindent Pollux (X, 152) donne une liste de divers synonymes pour désigner une \italique{bourse}. Voici ce qu'il dit :
\citComm{καὶ μιξάμενοι δ’ ἂν εἴποιμεν σκεύη, βαλάντια καὶ βαλαντίδια ὡς ἐν Αἰξὶν Εὔπολις. καὶ \textbf{θυλάκιον} δὲ καὶ θυλακίσκον·\\
Ἀριστοφάνης γοῦν ἐν Τριφάλητι τοῦτο ὑποδηλοῖ, ὅταν φῇ
\citComm{ἔπειτ' ἐπὶ τοὖψον ἧκε, τὴν σπυρίδα λαβών\\
καὶ θυλακίσκον καὶ τὸ μέγα βαλάντιον. }
$[$ ... $]$
ἐπὶ τοῦ αὐτοῦ δὲ καὶ σακίον, ὅταν φῇ ἐν Θεσμοφοριαζούσαις 
\citComm{\textbf{σακίον, ἐν οἷσπερ τἀργύριον ταμιεύεται.}}
Et les mêlant nous dirions, comme dans les \italique{Chèvres} d'Eupolis, les ustensiles, les grandes bourses et les petites bourses. Ainsi que θυλάκιον et θυλακίσκον. Aristophane l'illustre dans \italique{Triphalès}, lorsqu'il dit :
\citComm{Ensuite il arriva aux hors-d'œuvre, tenant la corbeille, un petit sac (θυλακίσκον) et une grande bourse.}
Ce mot désigne aussi une petite bourse, comme cela apparaît dans les \italique{Thesmophories} :
\citComm{\italique{une petite bourse, dans lesquelles on dépose l'argent.}}}
\noindent Aristophane a écrit deux pièces intitulées \italique{Les Thesmophories}. Ce passage a été considéré comme un fragment des \italique{Thesmophories} perdues (fr. 343 de l'édition Kassel-Austin). En fait, ce petit bout de phrase explicatif, σακίον· ἐν οἷσπερ τἀργύριον ταμιεύεται, ressemble bien davantage à une glose d'un commentateur. Il arrive fréquemment à Pollux de fusionner des gloses avec le texte d'Aristophane lui-même\marqueur{5}. Ce n'est pas en réalité aux \italique{Thesmophories} perdues mais bien au vers 291 des \italique{Thesmophories} conservées que se rattache cette glose, en tant que définition de θυλάκιον :
\citComm{$<$θυλάκιον·$>$ σακίον, ἐν οἷσπερ τἀργύριον ταμιεύεται.}
\noindent Une métathèse ancienne aura donné la leçon fautive du \italique{Ravennas}, comportant †θάληκον† en place de θυλάκιον, tandis que cette métathèse adjointe à sa glose `σακίον' rejoignait le lexique d'Hésychius sous la forme : θάλλικα· σάκκου εἶδος.
\end{multicols}
}%
\end{adjustbox}%
\vfill
\end{adjustwidth}
\newpage
% --- PAGE DE DROITE : Traduction française + Commentaire (8/8) ---
\begin{adjustwidth}{-2cm}{-2cm}
\noindent\textbf{\footnotesize Traduction française --- Folio 07 (\texttt{f. 165v}, vv. 252–294 [4/4]) \hfill ΠΡΟΛΟΓΟΣ}\\[-2pt]
\noindent\rule{16.5cm}{0.5pt}\\[2pt]
{\fontsize{8.4pt}{9.7pt}\selectfont
\ligneEst{}{\francais{Vois, Thratta, sous la lueur fumeuse des torches ardentes,}}{280}
\ligneEst{}{\francais{vois la foule qui monte.}}{}
\ligneEst{}{\francais{— Thesmophores, merveilles de beauté, accueillez-moi}}{}
\ligneEst{}{\francais{sous d'heureux auspices, ici et à mon retour au logis.}}{}
\ligneEst{}{\francais{Thratta, dépose la ciste, et tires-en}}{}
\ligneEst{}{\francais{le gâteau, pour que je l'offre aux deux déesses.}}{285}
\ligneEst{}{\francais{Maîtresse très honorée, chère Déméter,}}{}
\ligneEst{}{\francais{et toi Perséphone, puissé-je mainte et mainte fois}}{}
\ligneEst{}{\francais{être à même de te sacrifier. (À part :) À tout le moins, n'être pas découvert aujourd'hui.}}{}
\ligneEst{}{\francais{(Fort :) Que ma fille Chœrile trouve un mari riche}}{}
\ligneEst{}{\francais{fortuné, et au surplus benêt et stupide.}}{290}
\ligneEst{}{\francais{Et qu'elle me soit sensée et avisée sur le porte-monnaie.}}{}
\ligneEst{}{\francais{Où donc, où m'asseoir en bonne place pour entendre les orateurs ?}}{}
\ligneEst{}{\francais{Toi, Thratta, retire-toi à l'écart,}}{}
\ligneEst{}{\francais{car il n'est pas permis aux esclaves d'ouïr ce qui se dit ici.}}{294}
}
\vspace{2pt}
\noindent\rule{16.5cm}{0.6pt}\\[2pt]
\noindent\textbf{\small Commentaire philologique, métrique et scénique (vv. 252–294 --- partie 8/8)}\\[2pt]
\begin{adjustbox}{max totalheight=13.0cm, max width=16.5cm, keepaspectratio}%
\parbox[t]{16.5cm}{
\setlength{\columnsep}{9pt}%
\begin{multicols}{2}
\fontsize{7.6pt}{8.8pt}\selectfont
\setlength{\parindent}{0pt}%
\setlength{\parskip}{1.2pt}%
\noindent Pour séduisante que soit cette reconstruction, elle n'en demeure pas moins conjecturale ; le \italique{Ravennas} ne comportant pour les \italique{Thesmophories} que très peu de scholies, et parmi les moins intéressantes.
\noindent Finissons d'ailleurs cette note en montrant que notre conjecture est aussi compatible avec l'unique scholie conservée, τὸν παιδαρίσκον. ἴσως δὲ παρὰ τὴν πόσθην αὐτὸ συνέθηκεν. Le scholiaste ayant devant lui θάληκον -- qu'il lisait indifféremment θάλικον -- devait l'analyser comme un mot formé sur τό θάλος, \italique{la jeune pousse} et, dans un sens métaphorique, l'\italique{enfant}, suivi d'un suffixe en -ικος ; l'on comprend le quiproquo, facilité par l'évocation deux vers plus haut de la fille de Mnésiloque et d'un imaginaire sexe féminin (χοιρίον).
\noindent\textbf{[Notes ad 289]} \textbf{(1)} Tome II de son édition posthume d'Aristophane (1919), pp. 548-549.
\noindent Cette note sur ce passage des \italique{Thesmophories}, la \nieme{8} qu'il écrivit, est reproduite en annexe.
\noindent \textbf{(2)} La scholie : ὡς τῆς θυγατρὸς αὐτοῦ κατ' ἐπίκλησιν οὕτω καλουμένης. οἷον βοίδιον ἢ χρυσίον ἢ μυρτ<ί>ον. ἴσως δὲ τὸ αἰδοῖον τὸ γυναικεῖον αἰνίττεται.
\noindent \textbf{(3)} Voir le « Dossier Mnésiloque » de l'Introduction.
\noindent \textbf{(4)} Certains manuscrits portent Χοιρίνη.
\noindent \textbf{(1)} Note 8, reproduite à la fin.
\noindent \textbf{(2)} Chantraine, dans sa \italique{Formation des noms en grec ancien} aux pages 406 et suivantes, parle longuement du suffixe -ισκος sans même évoquer un possible -αλισκος ou -αρισκος ; Chantraine cite παιδαρίσκος et μειρακίσκος en expliquant que le Grec recourait à ce dernier diminutif parce que l'autre, μειράκιον, était devenu d'usage banal.
\noindent \textbf{(3)} Je parle toujours, par abus de langage et conformément à un usage trop bien établi, \italique{du} scholiaste. Il est bien évident qu'en réalité les commentaires en marge des manuscrits sont de différents auteurs, et ont été compulsés les uns à la suite des autres postérieurement. Certaines scholies font preuve d'une grande érudition et apportent parfois des leçons meilleures que celle du texte transmis, d'autres, par contre, – la plupart, hélas ! – ne sont, au mieux, d'aucun intérêt voire souvent sources d'erreurs pour qui y accorderait trop d'importance. Pour finir de gâter l'affaire, les \italique{Thesmophories} est la pièce d'Aristophane la moins pourvue en scholies.
\noindent \textbf{(4)} La glose d'Hésychius sur θαλικοποιοί laisse à penser que le sac en question était fait de cuir et qu'il ressemblait à un petit coffret -- une \italique{cassette}.
\noindent \textbf{(5)} Voir au commentaire du vers 1175 un autre exemple.
\end{multicols}
}%
\end{adjustbox}%
\vfill
\end{adjustwidth}
\newpage
\clearpage\ifodd\value{page}\hbox{}\newpage\fi
\enTeteG{ΠΑΡΟΔΟΣ --- Codex Ravennas 429, f. 166r (Classense 168r) --- vv. 295–336}
\begin{adjustwidth}{-1.5cm}{-3.0cm}
\begin{adjustbox}{max totalheight=21.0cm, max width=16.5cm, keepaspectratio}
\begin{minipage}{16.5cm}
{\fontsize{7.2pt}{8.5pt}\selectfont
\noindent\textbf{Folio 08 --- Codex Ravennas 429, \texttt{f. 166r} (\texttt{Classense 168r}) : ΠΑΡΟΔΟΣ, vv. 295–336}\hfill \textit{Témoins :} $\text{R}^1, \Rpc, \Rdeux$, \PsiUn\par\vspace{2pt}
\hrule height 0.4pt\vspace{2.5pt}
}
{\fontsize{8.1pt}{9.5pt}\selectfont
\ligneRav{295}{\inRas{Κῆρυξ}{}}{\hspace{1.1em}\difference{εὐφημία} ἔστω,}{\textbf{295} εὐφημί᾽ R : \lecon{εὐφημία} Suid. ε 3797}
\ligneRav{}{}{\hspace{1.1em}εὐφημία ἔστω·}{}
\ligneRav{297}{}{\hspace{1.1em}εὔχεσθε ταῖν Θεσμοφόροιν $[\![\text{\sout{τῇ Δήμητρῃ καὶ τῇ Κόρῃ}}]\!]$}{\textbf{297} Θεσμοφόροιν τῇ Δήμητρῃ καὶ τῇ Κόρῃ R : \lecon{Θ.} Blaydes}
\ligneRav{}{}{\hspace{1.1em}καὶ τῷ Πλούτῳ καὶ τῇ Καλλιγενείᾳ}{}
\ligneRav{300}{}{\hspace{1.1em}καὶ τῇ Κουροτρόφῳ $[\![\text{\sout{τῇ Γῇ}}]\!]$}{}
\ligneRav{301}{}{\hspace{1.1em}καὶ τῷ Ἑρμῇ καὶ \difference{ταῖς Χ}άρισιν}{\textbf{301} Χάρισιν R : \lecon{ταῖς Χ.} Meineke}
\ligneRav{}{}{\hspace{1.1em}ἐκκλησίαν τήνδε καὶ σύνοδον τὴν νῦν}{}
\ligneRav{}{}{\hspace{1.1em}κάλλιστα καὶ ἄριστα ποῆσαι,}{}
\ligneRav{304}{}{\hspace{1.1em}πολυωφελῶς μὲν \difference{τῇ π}όλει τῇ Ἀθηναίων,}{\textbf{304} πόλει R : \lecon{τῇ π.} Meineke}
\ligneRav{304}{}{\hspace{1.1em}}{\textbf{304} πόλει R : \lecon{τῇ π.} Meineke}
\ligneRav{305}{}{\hspace{1.1em}τυχηρῶς δ᾽ ἡμῖν αὐταῖς·}{\textbf{305} \lecon{ἡμῖν} R : ὑμῖν Fritzsche}
\ligneRav{306}{}{\hspace{1.1em}\expo{καὶ}τὴν δρῶσαν καὶ ἀγορεύουσαν}{\textbf{306} \expo{καὶ}τὴν ἀγορεύουσαν \Rdeux}
\ligneRav{306}{}{\hspace{1.1em}\expo{καὶ}τὴν δρῶσαν καὶ ἀγορεύουσαν}{\textbf{306} \expo{καὶ}τὴν ἀγορεύουσαν \Rdeux}
\ligneRav{307}{}{\hspace{1.1em}τὰ βέλτιστα περὶ τὸν δῆμον \difference{τὸν Ἀθηναίων}}{\textbf{307} τῶν Ἀθηναίων R : \lecon{τὸν Ἀθηναίων} Grynaeus}
\ligneRav{307}{}{\hspace{1.1em}}{\textbf{307} τῶν Ἀθηναίων R : \lecon{τὸν Ἀθηναίων} Grynaeus}
\ligneRav{}{}{\hspace{1.1em}καὶ τὸν τῶν γυναικῶν}{}
\ligneRav{}{}{\hspace{1.1em}ταύτην νικᾶν·}{}
\ligneRav{310}{}{\hspace{1.1em}ταῦτ᾽ εὔχεσθε καὶ ὑμῖν αὐταῖς τἀγαθά·}{\textbf{310} \lecon{ὑμῖν} R : ἡμῖν \Mon}
\ligneRav{310}{}{\hspace{1.1em}}{\textbf{310} \lecon{ὑμῖν} R : ἡμῖν \Mon}
\ligneRav{}{}{\hspace{1.1em}ἰὴ παιὼν ἰὴ παιὼν ἰὴ παιών. χαίρωμεν·}{}
\ligneRav{312}{χορ/}{\hspace{1.1em}δεχόμεθα καὶ θεῶν γένος}{\textbf{312} δεχόμεσθα R : \lecon{-μεθα} Dindorf}
\ligneRav{313}{}{\hspace{1.1em}\difference{λιτόμεθα} ταῖσδ᾽ ἐπ᾽ εὐχαῖς}{\textbf{313} λιτόμεσθα R : \lecon{-μεθα} Bourdin \espApp λιτόμεσθα R : \lecon{λιτόμεθα} Bourdin}
\ligneRav{}{}{\hspace{1.1em}φανέντας ἐπιχαρῆναι·}{}
\ligneRav{315}{}{\hspace{1.1em}Ζεῦ μεγαλώνυμε χρυσολύρα τε}{}
\ligneRav{}{}{\hspace{1.1em}Δῆλον ὃς ἔχεις ἱεράν,}{}
\ligneRav{317}{}{\hspace{1.1em}καὶ σὺ \difference{παγκρατὲς} \difference{κόρα}}{\textbf{317} παμκρατὲς R : \lecon{παγκρατὲς} ed. pr. \espApp κόρᾳ R : \lecon{κόρα} \Mon}
\ligneRav{}{}{\hspace{1.1em}γλαυκῶπι χρυσόλογχε πόλιν}{}
\ligneRav{}{}{\hspace{1.1em}οἰκοῦσα περιμάχητον, ἐλθὲ δεῦρο·}{}
\ligneRav{320}{}{\hspace{1.1em}\difference{θηροφόνη} (θηροφ\inRas{ω}{ό}νε) παῖ Λατοῦς,}{\textbf{320} θηροφ\inRas{ω}{ό}νε παῖ-- \Rpc : \lecon{θηροφόνη} Hermann}
\ligneRav{}{}{\hspace{1.1em}Λατοῦς χρυσώπιδος ἔρνος,}{}
\ligneRav{}{}{\hspace{1.1em}σύ τε πόντιε σεμνὲ Πόσειδον}{}
\ligneRav{}{}{\hspace{1.1em}ἁλιμέδον}{}
\ligneRav{324}{}{\hspace{1.1em}προλιπὼν μυχὸν \difference{ἰχθυόεντα},}{\textbf{324} ἰχθυόεντ᾽ R, Suid. οι 181 : \lecon{ἰχθυόεντα} Wilamowitz}
\ligneRav{324}{}{\hspace{1.1em}οἰστροδόνητον}{\textbf{324} ἰχθυόεντ᾽ R, Suid. οι 181 : \lecon{ἰχθυόεντα} Wilamowitz}
\ligneRav{325}{}{\hspace{1.1em}Νηρέος \difference{εἰναλίου} τε κόραι}{\textbf{325} ἐναλίου R : ἐνάλιοί Fritzsche : \lecon{εἰναλίου} Bentley}
\ligneRav{326}{}{\hspace{1.1em}Νύμφαι τ᾽ ὀρείπλαγκτοι·}{\textbf{326} \lecon{ὀρείπλαγκτοι} R : ὀρίπλαγκτοι Austin}
\ligneRav{}{}{\hspace{1.1em}χρυσέα τε φόρμιγξ}{}
\ligneRav{}{}{\hspace{1.1em}ἰαχήσειεν ἐπ᾽ εὐχαῖς}{}
\ligneRav{329}{}{\hspace{1.1em}ἡμετέραις· τελέως δ᾽ ἐκκλησιάσαιμεν Ἀθηνῶν}{\textbf{329} Ἀθηναίων R : \lecon{Ἀθηνῶν} Reisig}
\ligneRav{330}{}{\hspace{1.1em}εὐγενεῖς γυναῖκες·}{}
\ligneRav{331}{κηρ/}{εὔχεσθε τοῖς θεοῖσι τοῖς Ὀλυμπίοις}{\textbf{331} (Κηρ.) R : \lecon{(Κορ.)} van Leeuwen}
\ligneRav{}{}{καὶ ταῖς Ὀλυμπίαισι καὶ τοῖς Πυθίοις}{}
\ligneRav{}{}{καὶ ταῖσι Πυθίαισι καὶ τοῖς Δηλίοις}{}
\ligneRav{334}{}{καὶ ταῖσι Δηλίαισι \difference{τοῖς τ᾽} ἄλλοις θεοῖς,}{\textbf{334} καὶ τοῖς R : \lecon{τοῖς τ᾽} Grynaeus}
\ligneRav{335}{}{εἴ τις ἐπιβουλεύει τι τῷ δήμῳ κακὸν}{}
\ligneRav{}{}{τῷ τῶν γυναικῶν ἢ 'πικηρυκεύεται}{}
}
\vspace{3pt}\noindent\rule{4.5cm}{0.4pt}\par\vspace{1.5pt}
{\fontsize{7.0pt}{8.3pt}\selectfont
\noindent\textbf{Loci paralleli (vv. 295–336) ---} \textbf{299} (Καλλιγενείᾳ $=$ Δημήτερι) Hesychius κ 472 Καλλιγένειαν ; cf. Photius κ 118 ($=$ fr. Aristophanes 331 Kassel-Austin) \espTest \textbf{300} (Κουροτρόφῳ $=$ Δημήτηρ) Hesychius κ 3865 κουροτρόφος ; scholium vet. in Hesiodus \italique{Opera et dies} 227-9 \espTest \textbf{322-4} (σύ τε -- οἰστροδόνητον) Suidas οι 181 (sed σεμνὲ et ἁλιμέδον omissis) \espTest \textbf{335-7} Satyrus \italique{Vit. Euripid.} ($=$ \italique{The Oxyrhynchus Papyri} IX 1176) fr. 39 col. XII 8-16 (Εὐριπιδίῳ pro Εὐριπίδῃ (?)) \espTest \textbf{336} ('πικηρυκεύεται) Photius ε 1580 ἐπικηρυκεύεται $=$ \Syn ε 682 $=$ Suidas ε 2375 $=$ \EtymMagnum 360.21 $=$ \EtymGen (B) ; cf. Hesychius ε 4853\par
}
\vspace{2.5pt}\noindent\rule{4.5cm}{0.4pt}\par\vspace{1.5pt}
{\fontsize{6.9pt}{8.2pt}\selectfont
\noindent\textbf{Suite (en commentaire) de l'apparat critique \& notes paléographiques (\texttt{f. 166r}) ---} \textbf{Témoin papyrologique :} \textbf{\PsiCinq{} (\italique{P.Oxy.} IX 1176, Satyrus)} cite les \textbf{vv. 335–337}. --- \textbf{Grattage de la rubrique d'interlocuteur au v. 295 :} \Rdeux{} a gratté Κῆρυξ (\inRas{Κῆρυξ}{} \Rdeux{} : \auDessus{κη}{ρυ} \Mon{}, tandis que $\Sigma^{\text{R}}$ explique γυνὴ μιμουμένη κήρυκα ; L. Hauseux édite Κορ.-Ἱέρεια). --- \textbf{Intrusions de gloses religieuses (vv. 297–298) :} Dans la prière en prose de la Héraut, R a incorporé dans le texte les deux gloses interlinéaires τῇ Δήμητρῃ καὶ τῇ Κόρῃ (glosant τοῖν Θεσμοφόροιν) et τῇ Γῇ (glosant τῇ Κουροτρόφῳ). --- \textbf{Ajout supralinéaire de \Rdeux{} (v. 306) :} \expo{καὶ}τὴν ἀγορεύουσαν. --- \textbf{Correction \italique{in rasura} au v. 320 :} θηροφ\inRas{ω}{ό}νε παῖ-- \Rpc{} ($\omega$ de \Rac{} gratté et corrigé en ό). --- \textbf{v. 295 :} \inRas{Κῆρυξ}{} \Rdeux{} : \auDessus{κη}{ρυ} \Mon{} : γυνὴ μιμουμένη κήρυκα $\Sigma$ : Κρίτυλλα Sommerstein : \lecon{(Κορ.-Ἱέρεια)} \lecon{Bibi} \espApp εὐφημί' R : \lecon{εὐφημία} Souda $\varepsilon\ 3797$. --- \textbf{v. 297 :} Θεσμοφόροιν τῇ Δήμητρῃ καὶ τῇ Κόρῃ R : \lecon{Θεσμοφόροιν} Blaydes (suppression de la glose). --- \textbf{v. 298 :} τῇ Κουροτρόφῳ τῇ Γῇ R : \lecon{τῇ Κουροτρόφῳ} Dobree ex $\Sigma^{\text{R}}$.\par
}
\end{minipage}
\end{adjustbox}
\end{adjustwidth}
\newpage
\enTeteD{Codex Ravennas 429 (Biblioteca Classense) --- Fac-similé du f. 166r (vv. 295–336)}
\begin{adjustwidth}{-1.5cm}{-3.0cm}
\begin{center}
\includegraphics[width=16.5cm,height=20.5cm,keepaspectratio]{/workspaces/Thesmophories/dossier_ravennas_thesmophories/folios_individuels_ff162v-177r/08_f166r_classense168r_vv0295-0336.jpg}\par
\vspace{2pt}
{\scriptsize\textbf{Codex Ravennas 429, f. 166r (Classense 168r)} --- ΠΑΡΟΔΟΣ, \textbf{vv. 295–336} ($\text{R}^1, \Rpc, \Rdeux$, \PsiUn).}
\end{center}
\end{adjustwidth}
\newpage
\clearpage\ifodd\value{page}\hbox{}\newpage\fi% --- FOLIO 08 (f. 166r) : SPREAD 2.1 (Texte grec établi + Traduction française + Commentaire 1/2 & 2/2) ---
\clearpage\ifodd\value{page}\hbox{}\newpage\fi
\enTeteG{ΘΕΣΜΟΦΟΡΙΑΖΟΥΣΑΙ --- Texte grec établi (f. 166r, vv. 295–336)}
\enTeteD{LES THESMOPHORIES --- Traduction française (f. 166r, vv. 295–336)}
\begin{adjustwidth}{-2cm}{-2cm}
\noindent\textbf{\footnotesize Texte grec établi --- Folio 08 (\texttt{f. 166r}, vv. 295–336) \hfill ΠΑΡΟΔΟΣ}\\[-2pt]
\noindent\rule{16.5cm}{0.5pt}\\[2pt]
{\fontsize{7.6pt}{8.7pt}\selectfont
\ligneEst{\textbf{Κορ.-Ἱέρεια} {\tiny\textit{prose}}}{εὐφημία ἔστω, | εὐφημία ἔστω. |}{}
\ligneEst{}{εὔχεσθε ταῖν Θεσμοφόροιν |}{}
\ligneEst{}{καὶ τῷ Πλούτῳ καὶ τῇ Καλλιγενείᾳ |}{}
\ligneEst{}{καὶ τῇ Κουροτρόφῳ |}{300}
\ligneEst{}{καὶ τῷ Ἑρμῇ καὶ ταῖς Χάρισιν |}{}
\ligneEst{}{ἐκκλησίαν τήνδε καὶ σύνοδον τὴν νῦν |}{}
\ligneEst{}{κάλλιστα καὶ ἄριστα ποῆσαι, |}{}
\ligneEst{}{πολυωφελῶς μὲν τῇ πόλει τῇ Ἀθηναίων, |}{}
\ligneEst{}{τυχηρῶς δ᾽ ἡμῖν αὐταῖς. |}{305}
\ligneEst{}{καὶ τὴν δρῶσαν καὶ ἀγορεύουσαν |}{}
\ligneEst{}{τὰ βέλτιστα περὶ τὸν δῆμον τὸν Ἀθηναίων |}{}
\ligneEst{}{καὶ τὸν τῶν γυναικῶν | ταύτην νικᾶν. |}{}
\ligneEst{}{ταῦτ᾽ εὔχεσθε καὶ ὑμῖν αὐταῖς τἀγαθά. |}{}
\ligneEst{}{ἰὴ παιὼν ἰὴ παιὼν ἰὴ παιών. χαίρωμεν.}{310}
\ligneEst{\textbf{Χορός} {\tiny\textit{$2$ iamb.}}}{δεχόμεθα καὶ θεῶν γένος}{}
\ligneEst{{\tiny\textit{$2$ iamb\cat}}}{λιτόμεθα ταῖσδ᾽ ἐπ᾽ εὐχαῖς}{}
\ligneEst{{\tiny\textit{$2$ iamb\cat}}}{φανέντας ἐπιχαρῆναι.}{}
\ligneEst{{\tiny\textit{$2$ adon.}}}{Ζεῦ μεγαλώνυμε χρυσολύρα τε}{315}
\ligneEst{{\tiny\textit{cr. $+$ chor.}}}{Δῆλον ὃς ἔχεις ἱεράν,}{}
\ligneEst{{\tiny\textit{lecyth.}}}{καὶ σὺ παγκρατὲς κόρα}{}
\ligneEst{{\tiny\textit{$2$ iamb.}}}{γλαυκῶπι χρυσόλογχε πόλιν}{}
\ligneEst{{\tiny\textit{$3$ iamb\cat}}}{οἰκοῦσα περιμάχητον, ἐλθὲ δεῦρο.}{}
\ligneEst{{\tiny\textit{alcm.}}}{καὶ πολυώνυμε θηροφόνη}{320}
\ligneEst{{\tiny\textit{paroem.}}}{Λατοῦς χρυσώπιδος ἔρνος,}{}
\ligneEst{{\tiny\textit{paroem.}}}{σύ τε πόντιε σεμνὲ Πόσειδον}{}
\ligneEst{{\tiny\textit{cret.}}}{ἁλιμέδον}{}
\ligneEst{{\tiny\textit{paroem.}}}{προλιπὼν μυχὸν ἰχθυόεντα,}{}
\ligneEst{{\tiny\textit{adon.}}}{οἰστροδόνητον}{}
\ligneEst{{\tiny\textit{alcm.}}}{Νηρέος εἰναλίου τε κόραι}{325}
\ligneEst{{\tiny\textit{pros.}}}{Νύμφαι τ᾽ ὀρείπλαγκτοι.}{}
\ligneEst{{\tiny\textit{ith.}}}{χρυσέα τε φόρμιγξ}{}
\ligneEst{{\tiny\textit{ba. $+$ ad.}}}{ἰαχήσειεν ἐπ᾽ εὐχαῖς}{}
\ligneEst{{\tiny\textit{2 hem.}}}{ἡμετέραις· τελέως δ᾽ ἐκκλησιάσαιμεν Ἀθηνῶν}{}
\ligneEst{{\tiny\textit{ith.}}}{εὐγενεῖς γυναῖκες.}{330}
\ligneEst{\textbf{Κορυφαῖος} {\tiny\textit{3 ia.}}}{εὔχεσθε τοῖς θεοῖσι τοῖς Ὀλυμπίοις}{}
\ligneEst{}{καὶ ταῖς Ὀλυμπίαισι καὶ τοῖς Πυθίοις}{}
\ligneEst{}{καὶ ταῖσι Πυθίαισι καὶ τοῖς Δηλίοις}{}
\ligneEst{}{καὶ ταῖσι Δηλίαισι τοῖς τ᾽ ἄλλοις θεοῖς,}{}
\ligneEst{}{εἴ τις ἐπιβουλεύει τι τῷ δήμῳ κακὸν}{335}
\ligneEst{}{τῷ τῶν γυναικῶν ἢ 'πικηρυκεύεται}{336}
}
\vspace{2pt}
\noindent\rule{16.5cm}{0.6pt}\\[2pt]
\noindent\textbf{\small Commentaire philologique, métrique et scénique (vv. 295–336 --- partie 1/2)}\\[2pt]
\begin{adjustbox}{max totalheight=9.8cm, max width=16.5cm, keepaspectratio}%
\parbox[t]{16.5cm}{
\setlength{\columnsep}{9pt}%
\begin{multicols}{2}
\fontsize{8.4pt}{9.7pt}\selectfont
\setlength{\parindent}{0pt}%
\setlength{\parskip}{1.2pt}%
\noindent\textbf{295-311} \textbf{} La coryphée-prêtresse commence la prière, en prose selon l'usage : (scholie) :
\citComm{πεζῇ εὐφημίᾳ χρῶντα οἱ κωμικοὶ, ἐπειδὰν εὐχὴν ἢ ψήφισμα εἰσάγωσιν.}
\noindent Cf. en \italique{Oiseaux} (vers 864-888) pour une prière analogue.
\noindent\textbf{307-8} \textbf{τὸν δῆμον τῶν Ἀθηναίων καὶ τὸν τῶν γυναικῶν.} Plaisanterie \italique{contre l'attente} (παρὰ προσδοκίαν) : « Le peuple des Athéniens et celui... des femmes », là où l'on attendait \grec{καὶ τὸν τῶν ξυμμάχων}, « et celui de ses alliés ».
\end{multicols}
}%
\end{adjustbox}%
\vfill
\end{adjustwidth}
\newpage
% --- PAGE DE DROITE : Traduction française + Commentaire (2/2) ---
\begin{adjustwidth}{-2cm}{-2cm}
\noindent\textbf{\footnotesize Traduction française --- Folio 08 (\texttt{f. 166r}, vv. 295–336) \hfill ΠΑΡΟΔΟΣ}\\[-2pt]
\noindent\rule{16.5cm}{0.5pt}\\[2pt]
{\fontsize{7.6pt}{8.7pt}\selectfont
\ligneEst{\textbf{Coryphée}}{\francais{(La prêtresse :) Recueillez-vous, recueillez-vous.}}{}
\ligneEst{}{\francais{Priez les Thesmophores,}}{}
\ligneEst{}{\francais{et Ploutos et Calligénie,}}{}
\ligneEst{}{\francais{et la Nourricière,}}{300}
\ligneEst{}{\francais{et Hermès et les Charites,}}{}
\ligneEst{}{\francais{de rendre cette ecclésie et présente réunion}}{}
\ligneEst{}{\francais{la meilleure et la plus avantageuse qu'il se pourra,}}{}
\ligneEst{}{\francais{pour le plus grand profit de la cité des Athéniens}}{}
\ligneEst{}{\francais{et pour notre propre bien.}}{305}
\ligneEst{}{\francais{Que celle qui aura agi et parlé}}{}
\ligneEst{}{\francais{le mieux en faveur du peuple des Athéniens}}{}
\ligneEst{}{\francais{et du peuple des femmes voie prévaloir ses avis.}}{}
\ligneEst{}{\francais{Demandez aux dieux cette faveur, et pour vous-mêmes toutes les prospérités.}}{}
\ligneEst{}{\francais{Io Péan ! Io Péan ! Soyons en Joie.}}{310}
\ligneEst{\textbf{Chœur}}{\francais{Nous acceptons ! Et nous supplions}}{}
\ligneEst{}{\francais{la gent des dieux d'apparaître}}{}
\ligneEst{}{\francais{et de prendre en gré ces prières.}}{}
\ligneEst{}{\francais{Zeus au grand renom, et toi, dieu à la lyre d'or,}}{315}
\ligneEst{}{\francais{qui tiens Délos la sainte,}}{}
\ligneEst{}{\francais{et toi, vierge toute-puissante,}}{}
\ligneEst{}{\francais{aux yeux glauques, à la lance d'or, qui habites}}{}
\ligneEst{}{\francais{une cité enviée, viens ici. Et toi aussi,}}{}
\ligneEst{}{\francais{vierge aux noms multiples, meurtrière de bêtes,}}{320}
\ligneEst{}{\francais{rejeton de Léto aux yeux d'or.}}{}
\ligneEst{}{\francais{Et toi, auguste Poseidon, thalassien,}}{}
\ligneEst{}{\francais{modérateur de la mer,}}{}
\ligneEst{}{\francais{quitte le gouffre poissonneux}}{}
\ligneEst{}{\francais{fouetté des vents.}}{}
\ligneEst{}{\francais{Et vous, filles de Nérée le marin,}}{325}
\ligneEst{}{\francais{et vous, Nymphes des montagnes.}}{}
\ligneEst{}{\francais{Que la phorminx d'or}}{}
\ligneEst{}{\francais{accompagne nos prières, et puissions-nous}}{}
\ligneEst{}{\francais{mener à bonne fin notre ecclésie,}}{}
\ligneEst{}{\francais{femmes de pure lignée athénienne.}}{330}
\ligneEst{\textbf{Coryphée}}{\francais{Invoquez les dieux et les déesses de l'Olympe,}}{}
\ligneEst{}{\francais{les dieux et les déesses de Pytho,}}{}
\ligneEst{}{\francais{les dieux et les déesses de Délos}}{}
\ligneEst{}{\francais{et les autres dieux.}}{}
\ligneEst{}{\francais{Si quelqu'un forme quelque dessein contre le peuple}}{335}
\ligneEst{}{\francais{des femmes, fait des ouvertures}}{336}
}
\vspace{2pt}
\noindent\rule{16.5cm}{0.6pt}\\[2pt]
\noindent\textbf{\small Commentaire philologique, métrique et scénique (vv. 295–336 --- partie 2/2)}\\[2pt]
\begin{adjustbox}{max totalheight=9.8cm, max width=16.5cm, keepaspectratio}%
\parbox[t]{16.5cm}{
\setlength{\columnsep}{9pt}%
\begin{multicols}{2}
\fontsize{8.4pt}{9.7pt}\selectfont
\setlength{\parindent}{0pt}%
\setlength{\parskip}{1.2pt}%
\noindent\textbf{312-30} \textbf{Chant du chœur.} Pour l'étude métrique de ce chant du chœur, l'on pourra consulter Lætitia Parker et son ouvrage de référence, \italique{The Songs of Aristophanes}\marqueur{1}. Elle parvient à faire apparaître comme une sorte de refrain scandé $- \cup \cup - \cup \cup - \cup \cup -$\marqueur{2} aux vers 315, 320 et 325.\\ Quoique la découverte soit extrêmement intéressante et encore à creuser, j'ai préféré pour ma part reprendre la colométrie de Prato, plus proche du manuscrit de Ravenne : «~\italien{Il Coro accoglie l'invito dell'Aralda e rivolge la sua preghiera alle divinità peculiari dell' ecclesia (maschile) escluse nella precedente invocazione in prosa e qui menzionate, non senza adattamenti, col consueto criterio selettivo, quasi per conferire un crisma di legalità maggiore alle decisioni finali.}~» (Prato.)
\noindent\textbf{[Notes ad 312-30]} \textbf{(1)} Lætitia P. E. Parker, \italique{The songs of Aristophanes}, Clarendon Press Oxford (1997), pp. 406-409.
\noindent \textbf{(2)} Lætitia désigne ce vers comme un \italique{énoplien} (dactylo-épitrite prolongé) et Prato comme un \italique{alcméen}.
\noindent\textbf{330} \textbf{εὐγενεῖς γυναῖκες.} Confirme, s'il était besoin, que seules les citoyennes athéniennes assistaient aux Thesmophories (cf. le vers 294 lorsque Mnésiloque renvoie Thratta).
\end{multicols}
}%
\end{adjustbox}%
\vfill
\end{adjustwidth}
\newpage
\clearpage\ifodd\value{page}\hbox{}\newpage\fi
\enTeteG{ΠΑΡΟΔΟΣ --- Codex Ravennas 429, f. 166v (Classense 168v) --- vv. 337–379}
\begin{adjustwidth}{-1.5cm}{-3.0cm}
\begin{adjustbox}{max totalheight=21.0cm, max width=16.5cm, keepaspectratio}
\begin{minipage}{16.5cm}
{\fontsize{7.2pt}{8.5pt}\selectfont
\noindent\textbf{Folio 09 --- Codex Ravennas 429, \texttt{f. 166v} (\texttt{Classense 168v}) : ΠΑΡΟΔΟΣ, vv. 337–379}\hfill \textit{Témoins :} $\text{R}^1, \Rpc, \Rdeux$, \PsiCinq\par\vspace{2pt}
\hrule height 0.4pt\vspace{2.5pt}
}
{\fontsize{8.1pt}{9.5pt}\selectfont
\ligneRav{337}{χορ/}{†\difference{Εὐριπίδῃ Μήδοις τ᾽}† ἐπὶ βλάβῃ τινὶ}{\textbf{337} ευριπιδιω\douteux{ι}[. .]\douteux{υ}τ \PsiCinq : Εὐριπίδῃ Μήδοις R : \lecon{Εὐριπίδῃ Μήδοις τ᾽} Biset}
\ligneRav{}{}{τῇ τῶν γυναικῶν, ἢ τυραννεῖν ἐπινοεῖ}{}
\ligneRav{}{}{ἢ τὸν τύραννον συγκατάγειν, ἢ παιδίον}{}
\ligneRav{340}{}{ὑποβαλλομένης \difference{κατεῖπέ τις}, ἢ δούλη τινὸς}{\textbf{340} κατεῖπεν τις R : \lecon{κατεῖπέ τις} Brunck}
\ligneRav{341}{}{προαγωγὸς οὖσ᾽ ἐνετρύλλισεν τῷ δεσπότῃ}{\textbf{341} ἐνετρύλλισεν R, Suid. ε 1271 : \lecon{ἐνετρύλισεν} Dindorf}
\ligneRav{}{}{ἢ πεμπομένη τις ἀγγελίας ψευδεῖς φέρει,}{}
\ligneRav{}{}{ἢ μοιχὸς εἴ τις ἐξαπατᾷ ψευδῆ λέγων}{}
\ligneRav{344}{}{καὶ μὴ δίδωσιν \difference{ἃν} ὑπόσχηταί ποτε,}{\textbf{344} ἅ ἂν R : \lecon{ἃν} Faber}
\ligneRav{345}{}{ἢ δῶρά τις δίδωσι μοιχῷ γραῦς γυνή,}{}
\ligneRav{346}{}{ἢ καὶ δέχεται προδιδοῦσ᾽ \difference{ἑταίρᾳ} τὸν φίλον,}{\textbf{346} ἑταίρα R : \lecon{ἑταίρᾳ} Fritzsche ap Σ}
\ligneRav{347}{}{κεἴ τις κάπηλος ἢ καπηλὶς τοῦ χοὸς}{\textbf{347} \lecon{χοὸς} R, Suid. λ 394 : χοῶς Brunck}
\ligneRav{}{}{ἢ τῶν κοτυλῶν τὸ νόμισμα διαλυμαίνεται,}{}
\ligneRav{349}{}{κακῶς ἀπολέσθαι τοῦτον αὐτὸν \difference{κᾠκίαν}}{\textbf{349} κ' οἰκίαν R : \lecon{κᾠκίαν} Brunck}
\ligneRav{350}{}{ἀρᾶσθε, ταῖς δ᾽ ἄλλαισιν ὑμῖν τοὺς θεοὺς}{\textbf{350} \lecon{ὑμῖν} R : ἡμῖν \Mon}
\ligneRav{351}{}{εὔχεσθε πάσαις πολλὰ δοῦναι κἀγαθά·}{\textbf{351} \lecon{κἀγαθά} R : τἀγαθά Zanetti}
\ligneRav{352}{χορ/}{\hspace{1.1em}\difference{ξυνευχόμεσθα} τέλεα μὲν}{\textbf{352} ξυνευχόμεθα R : \lecon{ξυνευχόμεσθα} Daubruz}
\ligneRav{353}{}{\hspace{1.1em}πόλει τέλεα τε δήμῳ}{\textbf{353} \lecon{τε} R : δὲ \Mon}
\ligneRav{354}{}{\hspace{1.1em}τάδ᾽ εὔγματ᾽ \difference{ἀποτελεῖσθαι},}{\textbf{354} γενέσθαι R : ἐκγενέσθαι Dindorf : \lecon{ἀποτελεῖσθαι} Willems}
\ligneRav{355}{}{\hspace{1.1em}τὰ δ᾽ ἄρισθ᾽ ὅσαις προσήκει}{}
\ligneRav{}{}{\hspace{1.1em}νικᾶν λεγούσαις· ὁπόσαι δ᾽}{}
\ligneRav{357}{}{\hspace{1.1em}\difference{ἐξαπατῶσιν} παραβαίνουσί τε τοὺς}{\textbf{357} ἐξαπατῶσι R : \lecon{ἐξαπατῶσιν} Bentley}
\ligneRav{}{}{\hspace{1.1em}ὅρκους τοὺς νενομισμένους}{}
\ligneRav{360}{}{\hspace{1.1em}κερδῶν \difference{οὕνεκ᾽} ἐπὶ βλάβῃ,}{\textbf{360} ἕνεκ᾽ R : \lecon{οὕνεκ᾽} Bentley}
\ligneRav{}{}{\hspace{1.1em}ἢ ψηφίσματα καὶ νόμον}{}
\ligneRav{}{}{\hspace{1.1em}ζητοῦσ᾽ ἀντιμεθιστάναι,}{}
\ligneRav{}{}{\hspace{1.1em}τἀπόρρητά τε τοῖσιν ἐχ-}{}
\ligneRav{364}{}{\hspace{1.1em} θροῖς τοῖς ἡμετέροις \difference{λέγουσ᾽},}{\textbf{364} λέγουσιν R, Suid. α 3501 : \lecon{λέγουσ᾽} Bentley}
\ligneRav{365}{}{\hspace{1.1em}ἢ Μήδους ἐπάγουσι τῆς}{}
\ligneRav{}{}{\hspace{1.1em}χώρας οὕνεκ᾽ ἐπὶ βλάβῃ,}{}
\ligneRav{367}{}{\hspace{1.1em}\difference{ἀσεβοῦσ᾽ ἀδικοῦσί} τε τὴν πόλιν·}{\textbf{367} ἀσεβοῦσιν ἀδικοῦσίν R : \lecon{ἀσεβοῦσ᾽ ἀδικοῦσί} Bothe}
\ligneRav{}{}{\hspace{1.1em}ἀλλ᾽ ὦ παγκρατὲς}{}
\ligneRav{}{}{\hspace{1.1em}Ζεῦ ταῦτα κυρώσειας, ὥσθ᾽}{}
\ligneRav{370}{}{\hspace{1.1em}ἡμῖν θεοὺς παραστατεῖν}{}
\ligneRav{}{}{\hspace{1.1em}καίπερ γυναιξὶν οὔσαις·}{}
\ligneRav{372}{χορ/}{ἄκουε \difference{πᾶσ᾽}. ἔδοξε τῇ βουλῇ τάδε}{\textbf{372} πᾶς R : \lecon{πᾶσ᾽} Voss}
\ligneRav{}{}{τῇ τῶν γυναικῶν· Τιμόκλει᾽ ἐπεστάτει,}{}
\ligneRav{374}{}{Λύσιλλ᾽ ἐγραμμάτευεν, εἶπε Ζωστράτη·}{\textbf{374} Ἀρ[. . . .]\douteux{λ}(vel \douteux{a}, \douteux{δ}, \douteux{μ})ειa \PsiCinq, Ἀρ[χίκλ]ει᾽ suppl. Hunt : \lecon{Τιμόκλει᾽} R, Σ, Suid. ε 2616}
\ligneRav{375}{}{ἐκκλησίαν ποιεῖν ἕωθεν τῇ μέσῃ}{}
\ligneRav{376}{}{τῶν Θεσμοφορίων, ᾗ μάλισθ᾽ ἡμῖν \difference{σχολή},}{\textbf{376} σχολῇ R : \lecon{σχολή} Zanetti}
\ligneRav{}{}{καὶ χρηματίζειν πρῶτα περὶ Εὐριπίδου,}{}
\ligneRav{}{}{ὅ τι χρὴ παθεῖν ἐκεῖνον· ἀδικεῖν γὰρ δοκεῖ}{}
\ligneRav{}{}{ἡμῖν ἁπάσαις. τίς ἀγορεύειν βούλεται;}{}
}
\vspace{3pt}\noindent\rule{4.5cm}{0.4pt}\par\vspace{1.5pt}
{\fontsize{7.0pt}{8.3pt}\selectfont
\noindent\textbf{Loci paralleli (vv. 337–379) ---} \textbf{341} Suidas ε 1271 \espTest (προαγωγὸς) Hesychius π 3308 προαγωγός \espTest (ἐνετρύλλισεν) Zonaras 745 οὔ σ' ἐνετρύλλησεν \espTest \textbf{347-8} Suidas χ 394 (sed εἴ pro κεἴ) \espTest \textbf{363-4} (τἀπόρρητά -- λέγουσ᾽) Suidas α 3501 ἀπόρρητά τε τοῖσιν ἐχθροῖσι τοῖς ἡμετέροις λέγουσιν \espTest \textbf{369-71} (ὥσθ᾽ -- οὔσαις) Suidas π 444 \espTest \textbf{373-4} Satyrus \italique{Vit. Euripid.} ($=$ \italique{The Oxyrhynchus Papyri} IX 1176) fr. 39 col. XII 1 (vid. ad 335-7 ; sed Ἀρ[χίκλ]ει᾽ pro Τιμόκλει᾽, suppl. Hunt) \espTest (Τιμόκλει᾽ -- ἐγραμμάτευεν) Suidas ε 2616 (sed ἐπιστάτει pro ἐπεστάτει et ἐγραμμάτευε pro ἐγραμμάτευεν) \espTest \textbf{377-8} (καὶ χρηματίζειν -- παθεῖν) Suidas χ 481 ; Lexicon Vindobonense 190.10\par
}
\vspace{2.5pt}\noindent\rule{4.5cm}{0.4pt}\par\vspace{1.5pt}
{\fontsize{6.9pt}{8.2pt}\selectfont
\noindent\textbf{Suite (en commentaire) de l'apparat critique \& notes paléographiques (\texttt{f. 166v}) ---} \textbf{Témoin papyrologique :} \textbf{\PsiCinq{} (\italique{P.Oxy.} IX 1176, Satyrus)} cite les \textbf{vv. 335–337} et \textbf{vv. 373–375}. --- \textbf{v. 346 :} ἑταίρα R : \lecon{ἑταίρᾳ} Fritzsche ap. $\Sigma^{\text{R}}$. --- \textbf{v. 354 :} γενέσθαι R (ἐκγενέσθαι Dindorf) : \lecon{ἀποτελεῖσθαι} Alphonse Willems (adopté par L. Hauseux : γενέσθαι est une glose banale substituée au verbe rituel ἀποτελεῖσθαι). --- \textbf{v. 374 :} Ἀρ[....]λεια \PsiCinq{} (Ἀρ[χίκλ]ει' suppl. Hunt) : \lecon{Τιμόκλει'} R, $\Sigma^{\text{R}}$, Souda $\varepsilon\ 2616$ (maintenu par L. Hauseux).\par
}
\end{minipage}
\end{adjustbox}
\end{adjustwidth}
\newpage
\enTeteD{Codex Ravennas 429 (Biblioteca Classense) --- Fac-similé du f. 166v (vv. 337–379)}
\begin{adjustwidth}{-1.5cm}{-3.0cm}
\begin{center}
\includegraphics[width=16.5cm,height=20.5cm,keepaspectratio]{/workspaces/Thesmophories/dossier_ravennas_thesmophories/folios_individuels_ff162v-177r/09_f166v_classense168v_vv0337-0379.jpg}\par
\vspace{2pt}
{\scriptsize\textbf{Codex Ravennas 429, f. 166v (Classense 168v)} --- ΠΑΡΟΔΟΣ, \textbf{vv. 337–379} ($\text{R}^1, \Rpc, \Rdeux$, \PsiCinq).}
\end{center}
\end{adjustwidth}
\newpage
\clearpage\ifodd\value{page}\hbox{}\newpage\fi% --- FOLIO 09 (f. 166v) : SPREAD 2.1 (Texte grec établi + Traduction française + Commentaire 1/8 & 2/8) ---
\clearpage\ifodd\value{page}\hbox{}\newpage\fi
\enTeteG{ΘΕΣΜΟΦΟΡΙΑΖΟΥΣΑΙ --- Texte grec établi (f. 166v, vv. 337–379 [1/4])}
\enTeteD{LES THESMOPHORIES --- Traduction française (f. 166v, vv. 337–379 [1/4])}
\begin{adjustwidth}{-2cm}{-2cm}
\noindent\textbf{\footnotesize Texte grec établi --- Folio 09 (\texttt{f. 166v}, vv. 337–379 [1/4]) \hfill ΠΑΡΟΔΟΣ}\\[-2pt]
\noindent\rule{16.5cm}{0.5pt}\\[2pt]
{\fontsize{8.8pt}{10.2pt}\selectfont
\ligneEst{}{†Εὐριπίδῃ Μήδοις τ᾽† ἐπὶ βλάβῃ τινὶ}{337}
\ligneEst{}{τῇ τῶν γυναικῶν, ἢ τυραννεῖν ἐπινοεῖ}{}
\ligneEst{}{ἢ τὸν τύραννον συγκατάγειν, ἢ παιδίον}{}
\ligneEst{}{ὑποβαλλομένης κατεῖπέ τις, ἢ δούλη τινὸς}{340}
\ligneEst{}{προαγωγὸς οὖσ᾽ ἐνετρύλλισεν τῷ δεσπότῃ}{}
\ligneEst{}{ἢ πεμπομένη τις ἀγγελίας ψευδεῖς φέρει,}{}
\ligneEst{}{ἢ μοιχὸς εἴ τις ἐξαπατᾷ ψευδῆ λέγων}{}
\ligneEst{}{καὶ μὴ δίδωσιν ἃν ὑπόσχηταί ποτε,}{}
\ligneEst{}{ἢ δῶρά τις δίδωσι μοιχῷ γραῦς γυνή,}{345}
\ligneEst{}{ἢ καὶ δέχεται προδιδοῦσ᾽ ἑταίρᾳ τὸν φίλον,}{}
}
\vspace{2pt}
\noindent\rule{16.5cm}{0.6pt}\\[2pt]
\noindent\textbf{\small Commentaire philologique, métrique et scénique (vv. 337–379 --- partie 1/8)}\\[2pt]
\begin{adjustbox}{max totalheight=15.2cm, max width=16.5cm, keepaspectratio}%
\parbox[t]{16.5cm}{
\setlength{\columnsep}{9pt}%
\begin{multicols}{2}
\fontsize{8.4pt}{9.7pt}\selectfont
\setlength{\parindent}{0pt}%
\setlength{\parskip}{1.2pt}%
\noindent\textbf{336-7} \textbf{ἢ 'πικηρυκεύεται / Εὐριπίδῃ Μήδοις ἐπὶ βλάβῃ τινὶ.} Dans sa parodie d'ecclésie, Aristophane ne manque pas de reprendre les imprécations contre les Mèdes, qui étaient de règle comme on le sait par Isocrate (\italique{Panégyrique}, 157) :
\citComm{Καὶ πολλῶν μὲν οἱ πατέρες ἡμῶν μηδισμοῦ θάνατον κατέγνωσαν, ἐν δὲ τοῖς συλλόγοις ἔτι καὶ νῦν ἀρὰς ποιοῦνται, πρὶν ἄλλο τι χρηματίζειν, εἴ τις ἐπικηρυκεύεται Πέρσαις τῶν πολιτῶν.\\
\francais{Nos pères ont condamné à mort plusieurs citoyens pour leur adhésion au parti des Perses (μηδισμός). Encore aujourd'hui, dans nos assemblées, avant de traiter aucune affaire, on prononce des imprécations coutre celui des citoyens qui fait des ouvertures aux Mèdes.}}
\noindent L'ennemi héréditaire, pour la féminine assemblée, est remplacé -- on s'y attendait ! -- par Euripide.
\noindent Remplacé ? pas tout-à-fait ; car au côté d'Εὐριπίδῃ continuent de subsister étrangement les Μήδοις... détail qui fleure bon son interpolation : le spectateur athénien n'avait pas besoin qu'on lui souligne ainsi le trait comique...
\noindent Or nous possédons un papyrus d'Oxyrhynque\marqueur{1} du \nieme{II} siècle après Jésus-Christ\marqueur{2} qui comporte une citation des vers 335-337 des \italique{Thesmophories}. Malheureusement, le vers qui nous intéresse est lacunaire :
\citComm{ευριπιδιω\douteux{ι}[. .]\douteux{υ}\marqueur{3}τ|επιβλαβη\douteux{ι}τι|νι}
\noindent Hunt avait déjà remarqué en son temps que la lacune était bien insuffisante pour qu'ait pu s'y trouver Μήδοις τ᾽\marqueur{4}.
\noindent\textbf{338-9} \textbf{ἢ τυραννεῖν ἐπινοεῖ ἢ τὸν τύραννον συγκατάγειν.} Citation quasi-textuelle d'une ancienne loi, peut-être de Dracon, rapportée dans la \italique{Constitution d'Athènes} (XVI, 10) :
\citComm{ἦσαν δὲ καὶ τοῖς Ἀθηναίοις οἱ περὶ τῶν τυράννων νόμοι πρᾷοι κατ᾽ ἐκείνους τοὺς καιρούς, οἵ τ᾽ ἄλλοι καὶ δὴ καὶ ὁ μάλιστα καθήκων πρὸς τὴν τῆς τυραννίδος <κατάστασιν>. νόμος γὰρ αὐτοῖς ἦν ὅδε. `θέσμια τάδε Ἀθηναίων καὶ πάτρια· ἐάν τινες τυραννεῖν ἐπανιστῶνται ἐπὶ τυραννίδι, ἢ συγκαθιστῇ τὴν τυραννίδα, ἄτιμον εἶναι καὶ αὐτὸν καὶ γένος.'\\
\francais{Les lois des Athéniens sur la tyrannie étaient alors [à l'époque de Pisitrate] peu sévères, celle surtout qui visait les entreprises des tyrans et dont voici le texte : « Les lois athéniennes établies par nos pères portent que : quiconque aspire à la tyrannie ou forme un complot pour l'établir, sera frappé d'atimie dans sa personne et dans sa race. »}}
\end{multicols}
}%
\end{adjustbox}%
\vfill
\end{adjustwidth}
\newpage
% --- PAGE DE DROITE : Traduction française + Commentaire (2/8) ---
\begin{adjustwidth}{-2cm}{-2cm}
\noindent\textbf{\footnotesize Traduction française --- Folio 09 (\texttt{f. 166v}, vv. 337–379 [1/4]) \hfill ΠΑΡΟΔΟΣ}\\[-2pt]
\noindent\rule{16.5cm}{0.5pt}\\[2pt]
{\fontsize{8.8pt}{10.2pt}\selectfont
\ligneEst{}{\francais{à Euripide et aux Mèdes au détriment}}{337}
\ligneEst{}{\francais{des femmes, vise à être tyran}}{}
\ligneEst{}{\francais{ou à ramener le tyran,}}{}
\ligneEst{}{\francais{dénonce une femme qui a supposé un enfant; si une esclave}}{340}
\ligneEst{}{\francais{ayant servi d'entremetteuse l'est allé chuchoter au mari,}}{}
\ligneEst{}{\francais{ou si chargée d'un message, elle fait un rapport mensonger ;}}{}
\ligneEst{}{\francais{si un galant suborne par de belles paroles,}}{}
\ligneEst{}{\francais{et ne tient pas les promesses qu'il a faites ;}}{}
\ligneEst{}{\francais{si une vieille fait des présents à un galant,}}{345}
\ligneEst{}{\francais{ou en reçoit en livrant à une courtisane son galant ;}}{}
}
\vspace{2pt}
\noindent\rule{16.5cm}{0.6pt}\\[2pt]
\noindent\textbf{\small Commentaire philologique, métrique et scénique (vv. 337–379 --- partie 2/8)}\\[2pt]
\begin{adjustbox}{max totalheight=15.2cm, max width=16.5cm, keepaspectratio}%
\parbox[t]{16.5cm}{
\setlength{\columnsep}{9pt}%
\begin{multicols}{2}
\fontsize{8.4pt}{9.7pt}\selectfont
\setlength{\parindent}{0pt}%
\setlength{\parskip}{1.2pt}%
\noindent Étant donnée l'époque troublée de la représentation des \italique{Thesmophories}, cette loi devait être sur toutes les bouches ; pour preuve, la même année que le représentation des \italique{Thesmophories}, en 410/09\marqueur{1}, Démophante proposa un décret inspiré de cette dernière garantissant l'impunité à l'assassin de toute personne complice de crime de lèse-démocratie (cette loi fut gravée sur une stèle fameuse à Athènes qu'Andocide cite dans son \italique{Discours sur les mystères}, 97) :
\citComm{\grec{καὶ ἐάν τις τυραννεῖν ἐπαναστῇ ἢ τὸν τύραννον συγκαταστήσῃ}.}
\noindent\textbf{[Notes ad 336-7]} \textbf{(1)} \italique{The Oxyrhynchus Papyri}, IX (1912), ed. Arthur S. Dunt. Il s'agit du papyrus n° 1176, fragment 39, colonne XII 8-16. Ce papyrus date du \nieme{II} siècle ap. J.-C. et comporte un fragment d'une \italique{Vie} d'Euripide par Satyros de Callatis. Je me réfère aussi à l'édition ainsi qu'au commentaire de ce fragment par Stefan Schorn, \italique{Satyros aus Kallatis} (2002).
\noindent \textbf{(2)} Quand notre manuscrit de Ravenne est beaucoup plus tardif : \nieme{IX} ou \nieme{X} s. ap. J.-C.
\noindent \textbf{(3)} Schorn indique que la présence d'un \douteux{c} est moins probable, et que ce n'est assurément pas un \douteux{ε}.
\noindent \textbf{(4)} Ce qui n'a pas empêché les éditeurs depuis de prétendre dans leur apparat critique que ce papyrus confirmait l'addition du τ᾽ de Biset : Εὐριπίδῃ Μήδοις <τ᾽>... sauf que le τ identifié par Hunt ne pouvait en aucun cas se rattacher aux Μήδοις, absents du papyrus.
\noindent \textbf{(1)} C. Carey, Ifigeneia G., Brenda G.-W., \italique{Use and abuse of law in the Athenian courts}, p. 295.
\noindent\textbf{346} \textbf{ἢ καὶ δέχεται προδιδοῦσ᾽ ἑταίρα τὸν φίλον.} Je reprends la correction de Fritzsche qui substitue avec bonheur \grec{ἑταίρᾳ} à \grec{ἑταίρα}. Le scholiaste précise d'ailleurs qu'il y a différentes versions dans les manuscrits : certaines avec le iota souscrit, d'autres sans :
\citComm{γρ<άφεται> ἔν τισιν ἑταίρα καὶ ἔστιν ἕτερος νοῦς [...] ἐὰν μὲν ἄνευ τοῦ ῑ γρ<άφηται> ἑταίρα [...] ἐὰν δὲ σὺν τῷ ῑ}
\end{multicols}
}%
\end{adjustbox}%
\vfill
\end{adjustwidth}
\newpage

% --- FOLIO 09 (f. 166v) : SPREAD 2.2 (Texte grec établi + Traduction française + Commentaire 3/8 & 4/8) ---
\clearpage\ifodd\value{page}\hbox{}\newpage\fi
\enTeteG{ΘΕΣΜΟΦΟΡΙΑΖΟΥΣΑΙ --- Texte grec établi (f. 166v, vv. 337–379 [2/4])}
\enTeteD{LES THESMOPHORIES --- Traduction française (f. 166v, vv. 337–379 [2/4])}
\begin{adjustwidth}{-2cm}{-2cm}
\noindent\textbf{\footnotesize Texte grec établi --- Folio 09 (\texttt{f. 166v}, vv. 337–379 [2/4]) \hfill ΠΑΡΟΔΟΣ}\\[-2pt]
\noindent\rule{16.5cm}{0.5pt}\\[2pt]
{\fontsize{8.8pt}{10.2pt}\selectfont
\ligneEst{}{κεἴ τις κάπηλος ἢ καπηλὶς τοῦ χοὸς}{}
\ligneEst{}{ἢ τῶν κοτυλῶν τὸ νόμισμα διαλυμαίνεται,}{}
\ligneEst{}{κακῶς ἀπολέσθαι τοῦτον αὐτὸν κᾠκίαν}{}
\ligneEst{}{ἀρᾶσθε, ταῖς δ᾽ ἄλλαισιν ὑμῖν τοὺς θεοὺς}{350}
\ligneEst{}{εὔχεσθε πάσαις πολλὰ δοῦναι κἀγαθά.}{}
\ligneEst{\textbf{Χορός} {\tiny\textit{$2$ iamb.}}}{ξυνευχόμεσθα τέλεα μὲν}{}
\ligneEst{{\tiny\textit{$2$ iamb\cat}}}{πόλει τέλεα τε δήμῳ}{}
\ligneEst{{\tiny\textit{$2$ iamb\cat}}}{τάδ᾽ εὔγματ᾽ ἀποτελεῖσθαι,}{}
\ligneEst{{\tiny\textit{$2$ iamb\cat}}}{τὰ δ᾽ ἄρισθ᾽ ὅσαις προσήκει}{355}
\ligneEst{{\tiny\textit{ia. $+$ chor.}}}{νικᾶν λεγούσαις· ὁπόσαι δ᾽}{}
\ligneEst{{\tiny\textit{$3$ chor.}}}{ἐξαπατῶσιν παραβαίνουσί τε τοὺς}{}
}
\vspace{2pt}
\noindent\rule{16.5cm}{0.6pt}\\[2pt]
\noindent\textbf{\small Commentaire philologique, métrique et scénique (vv. 337–379 --- partie 3/8)}\\[2pt]
\begin{adjustbox}{max totalheight=15.2cm, max width=16.5cm, keepaspectratio}%
\parbox[t]{16.5cm}{
\setlength{\columnsep}{9pt}%
\begin{multicols}{2}
\fontsize{8.4pt}{9.7pt}\selectfont
\setlength{\parindent}{0pt}%
\setlength{\parskip}{1.2pt}%
\noindent La leçon sans iota est incohérente : en quoi une hétaïre, dont c'est le métier de recevoir des dons en échange de faveurs \italique{trahirait-elle} (\grec{προδιδοῦσα}) son amant ce faisant ?
\noindent (Willems, sentant bien la difficulté, s'éloigne en traduisant : « à l'insu de son amant ».)  Seule donc la première version est satisfaisante : il est question au vers précédent d'une vieille faisant des cadeaux pour s'attirer les faveurs d'un galant (thème qu'Aristophane développe dans le \italique{Ploutos}, vv. 959 sqq.) ; après la lubricité, le lucre : cette fois la vieille reçoit des cadeaux pour servir d'entremetteuse en livrant (προδιδοῦσ᾽) le soupirant (τὸν φίλον) à une femme irrégulière (ἑταίρᾳ). (Rôle de vieille procureuse que tiendra d'ailleurs Euripide à la fin de la pièce.)
\noindent\textbf{347-8} \textbf{κεἴ τις κάπηλος ἢ καπηλὶς τοῦ χοὸς ἢ τῶν κοτυλῶν τὸ νόμισμα διαλυμαίνεται.} Les débitantes de mauvais aloi sont brocardées de la même façon par Aristophane dans le \italique{Ploutos}, v. 435-6 :
\citComm{Βλεψίδημος — ἆρ᾽ ἐστὶν ἡ καπηλὶς ἡ ᾽κ τῶν γειτόνων,\\
ἣ ταῖς κοτύλαις ἀεί με διαλυμαίνεται;}
\noindent\textbf{352-71} \textbf{} Je reprends encore une fois la colométrie de Prato pour ce chant du chœur.
\noindent\textbf{354} \textbf{τάδ᾽ εὔγματα γενέσθαι.} J'adopte la correction de Willems qui substitue ἀποτελεῖσθαι à γενέσθαι (\nieme{9} note, reproduite en annexe).
\noindent\textbf{356-67} \textbf{} Sur tout ce passage, vers 356 à 367, pour le texte et l'interprétation, allez voir la \nieme{10} note d'Alphonse Willems qui fut le point de départ à l'excellent article de Taillardat \& Blanc reproduit ci-après.
\end{multicols}
}%
\end{adjustbox}%
\vfill
\end{adjustwidth}
\newpage
% --- PAGE DE DROITE : Traduction française + Commentaire (4/8) ---
\begin{adjustwidth}{-2cm}{-2cm}
\noindent\textbf{\footnotesize Traduction française --- Folio 09 (\texttt{f. 166v}, vv. 337–379 [2/4]) \hfill ΠΑΡΟΔΟΣ}\\[-2pt]
\noindent\rule{16.5cm}{0.5pt}\\[2pt]
{\fontsize{8.8pt}{10.2pt}\selectfont
\ligneEst{}{\francais{si un débitant ou une débitante}}{}
\ligneEst{}{\francais{falsifie la mesure légale des chous ou des cotyles,}}{}
\ligneEst{}{\francais{souhaitez qu'il périsse misérablement avec sa maison ;}}{}
\ligneEst{}{\francais{mais à vous autres toutes priez les dieux}}{350}
\ligneEst{}{\francais{d'accorder des biens en abondance.}}{}
\ligneEst{\textbf{Chœur}}{\francais{Nous unissons nos prières aux tiennes}}{}
\ligneEst{}{\francais{pour le bien de la cité, pour le bien du peuple}}{}
\ligneEst{}{\francais{afin que ces vœux s'accomplissent,}}{}
\ligneEst{}{\francais{et que la victoire demeure, comme de raison,}}{355}
\ligneEst{}{\francais{à celles qui auront le mieux parlé. Mais}}{}
\ligneEst{}{\francais{toutes celles qui en imposent ou transgressent}}{}
}
\vspace{2pt}
\noindent\rule{16.5cm}{0.6pt}\\[2pt]
\noindent\textbf{\small Commentaire philologique, métrique et scénique (vv. 337–379 --- partie 4/8)}\\[2pt]
\begin{adjustbox}{max totalheight=15.2cm, max width=16.5cm, keepaspectratio}%
\parbox[t]{16.5cm}{
\setlength{\columnsep}{9pt}%
\begin{multicols}{2}
\fontsize{8.4pt}{9.7pt}\selectfont
\setlength{\parindent}{0pt}%
\setlength{\parskip}{1.2pt}%
\noindent\textbf{365-6} \textbf{ἢ Μήδους ἐπάγουσι τῆς χώρας / οὕνεκ᾽ ἐπὶ βλάβῃ.} Voilà deux vers bien malmenés et retouchés à l'envi par divers éditeurs d'Aristophane. On doit à Jean Taillardat et Alain Blanc d'avoir définitivement confirmé la leçon du manuscrit de Ravenne, tout en y relevant un jeu de mots jusqu'alors passé à l'as. La lecture de leur article\marqueur{1} est si rapicolante que je me ferais scrupule d'y taillarder, en voici donc un très copieux extrait : « Pendant plusieurs siècles, l'interprétation des vers 365-366 a été controversée. Les savants ont d'abord été d'accord pour chercher dans τῆς χώρας, v. 366, une désignation géographique (“l'Attique”) et pour comprendre, comme Andrea Divo, “\latin{quaecunque ... Persas inducunt in regionem nostram}”. Mais traduire τῆς χώρας οὕνεκα par “\latin{in regionem nostram}”\marqueur{2}, c'est prêter à οὕνεκα un sens qu'il n'a jamais. D'où diverses conjectures\marqueur{3} et, enfin, la \italique{crux philologorum} dont Van Leeuwen a marqué les vers 365-367. C'est à Willems\marqueur{4} que revient le mérite d'avoir fait progresser, d'une façon décisive, l'interprétation de ce texte : il a rendu leur juste sens aux leçons transmises, ἐπάγειν et οὕνεκα, en traduisant “toutes celles ... qui sollicitent les Mèdes en considération de notre pays, mais à notre dam”. Il commente ainsi sa version (t. II, p. 551-553) : “ Ici il s'agit de ceux qui sollicitent les Mèdes « dans l'intérêt du pays », τῆς χώρας οὕνεκα [ἐπὶ βλάβῃ]. « Dans l'intérêt du pays » est tout simplement un euphémisme équivalent à « contre les Péloponnésiens », ce qu'il fallait éviter de dire. On sait, en effet, que les oligarques intriguaient auprès des satrapes d'Asie Mineure pour détacher le Grand Roi de l'alliance de Sparte et le faire entrer dans celle d'Athènes (Thucydide VIII, 53). Mais l'on sait non moins bien que cette politique indignait Aristophane ».
\noindent Que Μήδους ἐπάγουσι τῆς χώρας οὕνεκα puisse signifier « elles sollicitent l'alliance des Mèdes pour l'amour de notre pays », cela est sûr\marqueur{5} ; mais on ne voit pas qu'il y ait là « un euphémisme » lié à une sorte d'interdit. Mieux vaudrait dire que le chœur rapporte ironiquement les propos de toute femme « qui cherche à traiter avec lui (= l'ennemi) dans l'intérêt \italique{soi-disant d'Athènes} »\marqueur{5}. On peut admettre en effet que cette phrase transpose adroitement le prétexte donné par les intrigantes (Μήδους ἐπάγομεν τῆς χώρας οὕνεκα) : l'enjambement τῆς | χώρας οὕνεκα invitait le chœur à faire un sort au mot χώρας ainsi mis en relief, en sorte qu'une intonation particulière (ironique) des acteurs pouvait suffire à faire comprendre “pour l'amour de notre pays, disent-elles”.
\end{multicols}
}%
\end{adjustbox}%
\vfill
\end{adjustwidth}
\newpage

% --- FOLIO 09 (f. 166v) : SPREAD 2.3 (Texte grec établi + Traduction française + Commentaire 5/8 & 6/8) ---
\clearpage\ifodd\value{page}\hbox{}\newpage\fi
\enTeteG{ΘΕΣΜΟΦΟΡΙΑΖΟΥΣΑΙ --- Texte grec établi (f. 166v, vv. 337–379 [3/4])}
\enTeteD{LES THESMOPHORIES --- Traduction française (f. 166v, vv. 337–379 [3/4])}
\begin{adjustwidth}{-2cm}{-2cm}
\noindent\textbf{\footnotesize Texte grec établi --- Folio 09 (\texttt{f. 166v}, vv. 337–379 [3/4]) \hfill ΠΑΡΟΔΟΣ}\\[-2pt]
\noindent\rule{16.5cm}{0.5pt}\\[2pt]
{\fontsize{8.8pt}{10.2pt}\selectfont
\ligneEst{{\tiny\textit{glyc.}}}{ὅρκους τοὺς νενομισμένους}{}
\ligneEst{{\tiny\textit{glyc.}}}{κερδῶν οὕνεκ᾽ ἐπὶ βλάβῃ,}{360}
\ligneEst{{\tiny\textit{glyc.}}}{ἢ ψηφίσματα καὶ νόμον}{}
\ligneEst{{\tiny\textit{glyc.}}}{ζητοῦσ᾽ ἀντιμεθιστάναι,}{}
\ligneEst{{\tiny\textit{glyc.}}}{τἀπόρρητά τε τοῖσιν ἐχ-}{}
\ligneEst{{\tiny\textit{glyc.}}}{\alinea θροῖς τοῖς ἡμετέροις λέγουσ᾽,}{}
\ligneEst{{\tiny\textit{glyc.}}}{ἢ Μήδους ἐπάγουσι τῆς}{365}
\ligneEst{{\tiny\textit{glyc.}}}{χώρας οὕνεκ᾽ ἐπὶ βλάβῃ,}{}
\ligneEst{{\tiny\textit{reiz. $+$ ia.}}}{ἀσεβοῦσ᾽ ἀδικοῦσί τε τὴν πόλιν.}{}
\ligneEst{{\tiny\textit{dochm.}}}{ἀλλ᾽ ὦ παγκρατὲς}{}
}
\vspace{2pt}
\noindent\rule{16.5cm}{0.6pt}\\[2pt]
\noindent\textbf{\small Commentaire philologique, métrique et scénique (vv. 337–379 --- partie 5/8)}\\[2pt]
\begin{adjustbox}{max totalheight=15.2cm, max width=16.5cm, keepaspectratio}%
\parbox[t]{16.5cm}{
\setlength{\columnsep}{9pt}%
\begin{multicols}{2}
\fontsize{7.6pt}{8.8pt}\selectfont
\setlength{\parindent}{0pt}%
\setlength{\parskip}{1.2pt}%
\noindent\textbf{} \textbf{} \ \ \ Il y a cependant autre chose dans ce texte : un jeu de mots sur τῆς χώρας, plaisanterie dont personne ne s'est encore avisé. En effet, le parallélisme éclatant, et sûrement dû à Aristophane lui-même, de κερδῶν οὕνεκ᾽ ἐπὶ βλάβῃ (v. 360) et de τῆς χώρας οὕνεκ᾽ ἐπὶ βλάβῃ (v. 366) incite à chercher dans τῆς χώρας, autant que le nom de l'Attique, une variation de la notion exprimée par κερδῶν ; car le chœur peut faire allusion, ici comme là, à l'intérêt personnel et à l'avidité des gens qu'il maudit. Or le mot χώρα désigne, chez plusieurs auteurs, “la place” qu'on occupe dans la vie, par exemple “le rang”, qu'on a dans la hiérarchie sociale ou militaire\marqueur{7}. Dans les \italique{Thesmophories}, v. 366, c'est de rang dans la hiérarchie politique et sociale qu'il s'agit ; si bien que τῆς χώρας οὕνεκα, signifie à la fois “pour l'amour de notre pays (disent-elles)” et “pour se pousser aux (bonnes) places”\marqueur{8}, littéralement “en vue de leur place personnelle”\marqueur{9}.
\noindent Plus que le cas général des σπουδαρχίδαι, “les sires de Brigue-charges”\marqueur{10}, Aristophane vise ici une circonstance particulière. Il pense, comme Willems l'a dit, aux intrigues ambitieuses d'Alcibiade\marqueur{11} et de Pisandre qui, pendant l'hiver 412-411, prétendaient arracher la Perse à son alliance spartiate et cherchaient à obtenir, pour Athènes, l'appui et l'or du Grand Roi\marqueur{12}, pour eux-mêmes, le pouvoir à Athènes. Et que fait Pisandre, à la fin de 412, quand le peuple le délègue auprès d'Alcibiade et de Tissapherne ? Avant de quitter l'Attique pour l'Asie Mineure, “il rend visite aux sociétés secrètes déjà établies à Athènes pour s'occuper des procès et des magistratures”, ξυνωμοσίας, αἵπερ ἐτύγχανον πρότερον ἐν τῇ πόλει οὖσαι ἐπὶ δίκαις καὶ ἀρχαῖς, ἁπάσας ἐπελθὼν... (Thuc. VIII, 54, 4). C'est dire que les comploteurs avaient ouvert la chasse aux places. Les oligarques se consultent ἐπὶ δίκαις καὶ ἀρχαῖς et tout le monde le sait à Athènes : deux mois avant les \italique{Thesmophories}, Aristophane a déjà dénoncé “ces gens qui s'agglomèrent et font bourre pour se pousser aux charges” τούς γε συνισταμένους τούτους καὶ τοὺς πιλοῦντας ἑαυτοὺς | ἐπὶ ταῖς ἀρχαῖσι (\italique{Lysistrata}, v. 577 sq.). Ces ambitieux sont concussionnaires et vénaux : “c'est pour être en mesure de voler que Pisandre et ceux qui ne pensent qu'aux charges créent à chaque instant grabuge et désordre. Qu'ils fassent donc, pour cela, tout ce qu'ils voudront ; mais cet argent-ci\marqueur{13}, il n'y a plus de danger qu'ils l'emportent”,
\citComm{ἵνα γὰρ Πείσανδρος ἔχοι κλέπτειν χοἰ ταῖς ἀρχαῖς ἐπέχοντες,\\
ἀεί τινα κορκορυγὴν ἐκύκων. οἱ δ᾽ οὖν τοῦδ᾽ οὕνεκα δρώντων\\
ὅ τι βούλονται: τὸ γὰρ ἀργύριον τοῦτ᾽ οὐκέτι μὴ καθέλωσιν.\\
(\italique{Lysistrata}, v. 490-492).}
\end{multicols}
}%
\end{adjustbox}%
\vfill
\end{adjustwidth}
\newpage
% --- PAGE DE DROITE : Traduction française + Commentaire (6/8) ---
\begin{adjustwidth}{-2cm}{-2cm}
\noindent\textbf{\footnotesize Traduction française --- Folio 09 (\texttt{f. 166v}, vv. 337–379 [3/4]) \hfill ΠΑΡΟΔΟΣ}\\[-2pt]
\noindent\rule{16.5cm}{0.5pt}\\[2pt]
{\fontsize{8.8pt}{10.2pt}\selectfont
\ligneEst{}{\francais{les serments établis,}}{}
\ligneEst{}{\francais{dans leur propre intérêt, au détriment commun,}}{360}
\ligneEst{}{\francais{celles qui cherchent à bouleverser}}{}
\ligneEst{}{\francais{les décrets et la loi,}}{}
\ligneEst{}{\francais{révèlent nos secrets à nos ennemis,}}{}
\ligneEst{}{\francais{ou se positionnent}}{}
\ligneEst{}{\francais{vis-à-vis des Mèdes pour le bien de notre empire...}}{365}
\ligneEst{}{\francais{et empirent notre position. Celles-là sont}}{}
\ligneEst{}{\francais{coupables de lèse-majesté divine et de lèse-nation.}}{}
\ligneEst{}{\francais{Ô tout-puissant}}{}
}
\vspace{2pt}
\noindent\rule{16.5cm}{0.6pt}\\[2pt]
\noindent\textbf{\small Commentaire philologique, métrique et scénique (vv. 337–379 --- partie 6/8)}\\[2pt]
\begin{adjustbox}{max totalheight=15.2cm, max width=16.5cm, keepaspectratio}%
\parbox[t]{16.5cm}{
\setlength{\columnsep}{9pt}%
\begin{multicols}{2}
\fontsize{8.4pt}{9.7pt}\selectfont
\setlength{\parindent}{0pt}%
\setlength{\parskip}{1.2pt}%
\noindent Ἐπὶ (δίκαις καὶ) ἀρχαῖς, dit Thucydide ; ἐπὶ ταῖς ἀρχαῖσι, οἱ ταῖς ἀρχαῖς ἐπέχοντες, dit Aristophane. Peut-on rêver plus belles gloses au τῆς χώρας οὕνεκα des \italique{Thesmophories} ? Aux yeux d'Aristophane, tremper dans la grenouillère des combinaisons louches de l'hiver de 412-411, c'est donc prétendre à faire d'une pierre deux coups : devenir un personnage important à Athènes et s'enrichir en profitant de tous les avantages, licites ou non\marqueur{14}, attachés aux ἀρχαί ; c'est tout cela que résume τῆς χώρας οὕνεκ᾽ ἐπὶ βλάβῃ, car de telles ambitions ne réussissent qu'au mépris du bien public.
\noindent \textbf{(1)} «~Aristophane, ``Thesmophories'', v. 366 : τῆς χώρας οὕνεκα et l'ambition des oligarques dans l'hiver de 412-411~», in \italique{Revue de philologie, de littérature et d'histoire ancienne}, t. LX, 1986, fasc. 2, p. 183-186
\noindent \textbf{(2)} Andreas Divus, \italique{Aristophanis comoediae undecim e graeco in latinum ad verbum translatae}, Basiliae, 1552. Brunck, en 1783, est même plus précis : « quaecunque ... Medos inducunt in regionem nostram, ad eam perdendam » (Brunck lit ἐπάγουσι τῆς χώρας ἐπὶ βλάβῃ, en supprimant οὕνεκα).
\noindent \textbf{(3)} Par exemple : ἐπάγουσι τῆς | χώρας ἀσεβοῦσιν (Dobrée, suivi par Dindorf et Blaydes ; οὕνεκ᾽ ἐπὶ βλάβῃ disparaît) ; ἐπάγουσι τῆς | χώρας ἐπὶ βλάβῃ (Brunck) ; ἐπάγουσι χώρας | κερδῶν οὕνεκα ἐπὶ βλάβῃ (Reiske) ; ἐπάγουσι τῆς | ἔχθρας οὕνεκ᾽ ἐπὶ βλάβῃ (Hermann) ; ἐπάγουσι τῆς | ὥρας οὕνεκ᾽ ἐπὶ βλάβῃ (Fritzsche) ; ἐπάγουσι γῇ | κερδῶν οὕνεκ᾽ ἐπὶ βλάβῃ (van Velsen, suivi par Hall-Geldart et Rogers, ce dernier « with some hesitation ») ; τῆς | χώρας οὐ μεμελημένοι (C. Austin, voir ci-dessous, n. 5). Cet acharnement à modifier le texte transmis par le Ravennas (seul manuscrit des \italique{Thesmophories}) est d'autant plus étrange qu'il est métriquement sain : les vers 359-366 sont des glyconiens.
\noindent \textbf{(4)} Il s'agit de la note 10 de Willems, reproduite en annexe.
\noindent \textbf{(5)} Malgré C. Austin (\italique{P. C. Ph. S.}, 20, 1974, p. 1-2) qui refuse cette interprétation et corrige le texte (voir note 3).
\end{multicols}
}%
\end{adjustbox}%
\vfill
\end{adjustwidth}
\newpage

% --- FOLIO 09 (f. 166v) : SPREAD 2.4 (Texte grec établi + Traduction française + Commentaire 7/8 & 8/8) ---
\clearpage\ifodd\value{page}\hbox{}\newpage\fi
\enTeteG{ΘΕΣΜΟΦΟΡΙΑΖΟΥΣΑΙ --- Texte grec établi (f. 166v, vv. 337–379 [4/4])}
\enTeteD{LES THESMOPHORIES --- Traduction française (f. 166v, vv. 337–379 [4/4])}
\begin{adjustwidth}{-2cm}{-2cm}
\noindent\textbf{\footnotesize Texte grec établi --- Folio 09 (\texttt{f. 166v}, vv. 337–379 [4/4]) \hfill ΠΑΡΟΔΟΣ}\\[-2pt]
\noindent\rule{16.5cm}{0.5pt}\\[2pt]
{\fontsize{8.8pt}{10.2pt}\selectfont
\ligneEst{{\tiny\textit{$2$ iamb.}}}{Ζεῦ ταῦτα κυρώσειας, ὥσθ᾽}{}
\ligneEst{{\tiny\textit{$2$ iamb.}}}{ἡμῖν θεοὺς παραστατεῖν}{370}
\ligneEst{{\tiny\textit{$2$ iamb\cat}}}{καίπερ γυναιξὶν οὔσαις.}{}
\ligneEst{\textbf{Κορυφαῖος} {\tiny\textit{trim. iamb.}}}{ἄκουε πᾶσ᾽. ἔδοξε τῇ βουλῇ τάδε}{}
\ligneEst{}{τῇ τῶν γυναικῶν· Τιμόκλει᾽ ἐπεστάτει,}{}
\ligneEst{}{Λύσιλλ᾽ ἐγραμμάτευεν, εἶπε Ζωστράτη·}{}
\ligneEst{}{ἐκκλησίαν ποιεῖν ἕωθεν τῇ μέσῃ}{375}
\ligneEst{}{τῶν Θεσμοφορίων, ᾗ μάλισθ᾽ ἡμῖν σχολή,}{}
\ligneEst{}{καὶ χρηματίζειν πρῶτα περὶ Εὐριπίδου,}{}
\ligneEst{}{ὅ τι χρὴ παθεῖν ἐκεῖνον· ἀδικεῖν γὰρ δοκεῖ}{}
\ligneEst{}{ἡμῖν ἁπάσαις. τίς ἀγορεύειν βούλεται;}{379}
}
\vspace{2pt}
\noindent\rule{16.5cm}{0.6pt}\\[2pt]
\noindent\textbf{\small Commentaire philologique, métrique et scénique (vv. 337–379 --- partie 7/8)}\\[2pt]
\begin{adjustbox}{max totalheight=15.2cm, max width=16.5cm, keepaspectratio}%
\parbox[t]{16.5cm}{
\setlength{\columnsep}{9pt}%
\begin{multicols}{2}
\fontsize{8.4pt}{9.7pt}\selectfont
\setlength{\parindent}{0pt}%
\setlength{\parskip}{1.2pt}%
\noindent \textbf{(6)} Willems, II, p. 552. C'est nous [Taillardat \& Blanc] qui soulignons soi-disant.
\noindent \textbf{(7)} Xén., \italique{Anab.}, 5, 6, 13 (ἐν ἀνδραπόδων χώρᾳ ἐσόμεθα), \italique{Cyr.}, 2, 1, 18 (ἐν μισθοφόρου χώρᾳ εἶναι), Polybe, 1, 43, 1 (τῶν ἡγεμόνων τινὲς τῶν τὰς μεγίστας χώρας ἐχόντων ἐν τοῖς μισθοφόροις), Pausanias, 10, 1, 8 (ὀ δὲ χώραν ἐν τοῖς ἄρχουσιν ἔχων τὴν μεγίστην μάντις ἦν Τελλίας ὁ Ἠλεῖος). Avec le sens de « place dans la société », χώρα paraît déjà exister chez Théognis, v. 822.
\noindent \textbf{(8)} On justifiera ci-dessous la traduction « aux \italique{bonnes} places ».
\noindent \textbf{(9)} Tel est le sens de l'article, cf. Kühner-Gerth, \italique{Satzlehre}, I, p. 593, n°2. C'est l'emploi intentionnel du singulier τῆς χώρας (Kühner-Gerth, I, p. 14, n°4), au lieu du pluriel τῶν χωρῶν (plus fréquent dans ce cas, cf. Kühner-Gerth, I, p. 19, Anm. 3), qui permet le jeu de mots. Pour οὕνεκα, \cf \italique{Lysistrata}, v. 491 : τοῦδ᾽ οὕνεκα « en vue de cela (= en vue des magistratures et des vols qu'elles permettent) ».
\noindent \textbf{(10)} \italique{Ach.}, 595 ; pour le fait, \cf \italique{Equ.}, 157-159, 164-165, et \italique{Aves}, 798-800.
\noindent \textbf{(11)} Envisageant de provoquer une révolution à Athènes, il préparait alors son retour dans sa patrie (Thuc., 8, 47, 1-2).
\noindent \textbf{(12)} Pour le rôle de l'or perse à ce moment, voir Thucydide, 8, 29, 1-2 ; 8, 45, 2-6; 8, 48, 2-3 ; 8, 53, 2.
\noindent \textbf{(13)} Le trésor déposé à l'Acropole dont les Athéniennes barricadées interdisent l'entrée.
\noindent \textbf{(14)} Voir J. Taillardat, \italique{Les images d'Aristophane}, pp. 408-423. Pour les avantages licites, \cf \italique{Acharniens}, v. 65 sqq. (indemnité de mission officielle et vie facile des ambassadeurs athéniens auprès du Grand Roi).
\noindent\textbf{373} \textbf{Ἀρ[χίκλ]ει᾽ / Τιμόκλει᾽.} Comme aux vers 336-337, le \italique{papyrus Oxyrhynchi} n° 1176\marqueur{1} diffère de la leçon du \italique{Ravennas} pour le nom de notre \italique{épistate}\marqueur{2} ; Hunt a suppléé aux quatre ou cinq lettres manquantes Ἀρ[. . . . .]ειa en proposant : Ἀρ[χίκλ]ει᾽. Pourquoi pas ? d'autres noms étaient cependant possibles.
\end{multicols}
}%
\end{adjustbox}%
\vfill
\end{adjustwidth}
\newpage
% --- PAGE DE DROITE : Traduction française + Commentaire (8/8) ---
\begin{adjustwidth}{-2cm}{-2cm}
\noindent\textbf{\footnotesize Traduction française --- Folio 09 (\texttt{f. 166v}, vv. 337–379 [4/4]) \hfill ΠΑΡΟΔΟΣ}\\[-2pt]
\noindent\rule{16.5cm}{0.5pt}\\[2pt]
{\fontsize{8.8pt}{10.2pt}\selectfont
\ligneEst{}{\francais{Zeus, ratifie ces vœux et fais}}{}
\ligneEst{}{\francais{que les dieux nous assistent,}}{370}
\ligneEst{}{\francais{quoique nous soyons femmes.}}{}
\ligneEst{\textbf{Coryphée}}{\francais{Écoutez toutes. La Boulé}}{}
\ligneEst{}{\francais{des femmes, Timocléa épistate,}}{}
\ligneEst{}{\francais{Lysilla greffière, Sostraté auteur de la motion, a décrété ce qui suit :}}{}
\ligneEst{}{\francais{Une assemblée aura lieu le matin du jour intermédiaire}}{375}
\ligneEst{}{\francais{des Thesmophories, celui où nous sommes le plus de loisir ;}}{}
\ligneEst{}{\francais{et on délibérera d'abord sur Euripide}}{}
\ligneEst{}{\francais{et la peine qu'il doit subir. Car il est tenu coupable}}{}
\ligneEst{}{\francais{par nous toutes. Qui veut prendre la parole ?}}{379}
}
\vspace{2pt}
\noindent\rule{16.5cm}{0.6pt}\\[2pt]
\noindent\textbf{\small Commentaire philologique, métrique et scénique (vv. 337–379 --- partie 8/8)}\\[2pt]
\begin{adjustbox}{max totalheight=15.2cm, max width=16.5cm, keepaspectratio}%
\parbox[t]{16.5cm}{
\setlength{\columnsep}{9pt}%
\begin{multicols}{2}
\fontsize{8.4pt}{9.7pt}\selectfont
\setlength{\parindent}{0pt}%
\setlength{\parskip}{1.2pt}%
\noindent Quoique le papyrus soit certainement plus fiable que notre \italique{Ravennas}, en l'absence de certitude sur la leçon originale du papyrus, je me bornerai à suivre celle de l'unicodex.
\noindent\textbf{376-7} \textbf{ἐκκλησίαν ποιεῖν ἕωθεν τῇ μέσῃ / τῶν Θεσμοφορίων, ᾗ μάλισθ᾽ ἡμῖν σχολή.} Remarquons au passage que les assemblées se tenaient à l'aube\marqueur{1}.< Le fait que le \italique{jour du milieu} soit celui du \italique{jeûne} (νηστεία) explique que les femmes aient tout le loisir pour tenir assemblée.
\noindent\textbf{373-9} \textbf{ἄκουε πᾶς ... τίς ἀγορεύειν βούλεται;} L'assemblée à huis-clos féminin se donne les formes pour être en droit de juger Euripide : tout ce passage est un pastiche d'un \grec{προβούλευμα}, une \italique{résolution du Conseil}. En voici un exemple : le début du décret proposé par Démophante évoqué plus haut (Andocide, \italique{Discours sur les mystères}, 96) :
\citComm{ΝΟΜΟΣ\\
Ἔδοξε τῇ βουλῇ καὶ τῷ δήμῳ. Αἰαντὶς ἐπρυτάνευε, Κλειγένης ἐγραμμάτευε, Βοηθὸς ἐπεστάτει. Τάδε Δημόφαντος συνέγραψεν. Ἄρχει χρόνος τοῦδε τοῦ ψηφίσματος ἡ βουλὴ οἱ πεντακόσιοι λαχόντες τῷ κυάμῳ, οἷς Κλειγένης πρῶτος ἐγραμμάτευεν.\marqueur{1}}
\noindent\textbf{[Notes ad 373]} \textbf{(1)} \italique{The Oxyrhynchus Papyri}, IX (1912), ed. Arthur S. Hunt. Il s'agit du papyrus n°1176, fragment 39, colonne XII 8-16 ; papyrus du \nieme{II} siècle ap. J.-C. comportant un fragment d'une \italique{Vie} d'Euripide par Satyros de Callatis. Voir aussi l'édition Schorn de Satyros, \italique{Satyros aus Kallatis} (2002).
\noindent \textbf{(2)} On sait que c'est au niveau des noms propres que les textes s'altèrent le plus lors des copies successives, surtout si, comme probablement ici, il s'agit d'un personnage inconnu ailleurs.
\noindent \textbf{(1)} Coutume confirmée par d'autres passages d'Aristophane : \italique{Acharniens}, 19-20 ; \italique{Assemblée des femmes}, 20-21, 84-85, 739-741.
\noindent \textbf{(1)} Autres exemples : Thucydide, IV, 118 ; \italique{Inscriptiones Graecae} I\expo{3} 5, 9, 10, 12, 17 \&c. \&c.
\end{multicols}
}%
\end{adjustbox}%
\vfill
\end{adjustwidth}
\newpage
\clearpage\ifodd\value{page}\hbox{}\newpage\fi
\enTeteG{ΕΠΕΙΣΟΔΙΟΝ Α' --- Codex Ravennas 429, f. 167r (Classense 169r) --- vv. 380–422}
\begin{adjustwidth}{-1.5cm}{-3.0cm}
\begin{adjustbox}{max totalheight=21.0cm, max width=16.5cm, keepaspectratio}
\begin{minipage}{16.5cm}
{\fontsize{7.2pt}{8.5pt}\selectfont
\noindent\textbf{Folio 10 --- Codex Ravennas 429, \texttt{f. 167r} (\texttt{Classense 169r}) : ΕΠΕΙΣΟΔΙΟΝ Α', vv. 380–422}\hfill \textit{Témoins :} $\text{R}^1, \Rpc, \Rdeux$ (\Mon)\par\vspace{2pt}
\hrule height 0.4pt\vspace{2.5pt}
}
{\fontsize{8.1pt}{9.5pt}\selectfont
\ligneRav{380}{\auDessus{γυ}{ν}}{ἐγώ. : περίθου νυν τόνδε πρῶτον πρὶν λέγειν·}{\textbf{380} (Γυνή τις) R : \lecon{(Μίκα)} Rogers, cf. v. 760 \espApp \lecon{πρῶτον} R : πρότερον Suid. π 1166}
\ligneRav{}{}{σίγα σιώπα, πρόσεχε τὸν νοῦν· χρέμπτεται γὰρ ἤδη}{}
\ligneRav{}{}{ὅπερ ποιοῦσ᾽ οἱ ῥήτορες. μακρὰν ἔοικε λέξειν·}{}
\ligneRav{383}{καλλιλεξία}{φιλοτιμίᾳ μὲν οὐδεμιᾷ μὰ τὼ θεὼ}{\textbf{383} (\auDessus{γυ}{ν}) \Rac : (καλλιλεξία γυνή) \Rpc, \Mon : (\lecon{Μίκα}) Rogers}
\ligneRav{}{}{λέξουσ᾽ ἀνέστην ὦ γυναῖκες· ἀλλὰ γὰρ}{}
\ligneRav{385}{}{βαρέως φέρω τάλαινα πολὺν ἤδη χρόνον}{}
\ligneRav{386}{}{προπηλακιζομένας ὁρῶς \difference{ἡμᾶς} ὑπὸ}{\textbf{386} ὑμᾶς R : \lecon{ἡμᾶς} Cobet}
\ligneRav{}{}{Εὐριπίδου τοῦ τῆς λαχανοπωλητρίας}{}
\ligneRav{}{}{καὶ πολλὰ καὶ παντοῖ᾽ ἀκουούσας κακά·}{}
\ligneRav{389}{}{τί γὰρ οὗτος ἡμᾶς οὐκ ἐπισμῇ τῶν κακῶν;}{\textbf{389} \lecon{ἡμᾶς} R, Suid. ε 2592 : ἡμῖν Brunck}
\ligneRav{390}{}{ποῦ δ᾽ οὐχὶ διαβέβληχ᾽, ὅπουπερ \difference{ἔμβραχυ}}{\textbf{390} ἐν βραχὺ R : \lecon{ἔμβραχυ} Suid., Σ \italique{Gorg}. et \italique{Hipp.m}.}
\ligneRav{391}{}{εἰσὶν θεαταὶ καὶ τραγῳδοὶ καὶ χοροί,}{\textbf{391} \lecon{τραγῳδοὶ καὶ} R : τραγῳδικοὶ Schol. Plat. \italique{Theag} : τραγῳδία καὶ Suid. οι 128}
\ligneRav{392}{}{τὰς μοιχοτρόπους, τὰς ἀνδρεραστρίας καλῶν,}{\textbf{392} μυ- R : \lecon{μοι-} Suid. \espApp ἀνδρεα- R : \lecon{ἀνδρερα-} Suid.}
\ligneRav{393}{}{τὰς οἰνοπίπας, τὰς προδότιδας, τὰς λάλους,}{\textbf{393} οἰνοπότιδας R : \lecon{-πίπας} Suid. : -πίπους Σ}
\ligneRav{394}{}{τὰς οὐδὲν ὑγιές, τὰς μέγ᾽ ἀνδράσιν κακόν·}{\textbf{394} ἀνδράσι R : \lecon{-σιν} Daubuz}
\ligneRav{395}{}{ὥστ᾽ εὐθὺς εἰσιόντες ἀπὸ τῶν ἰκρίων}{}
\ligneRav{}{}{ὑποβλέπουσ᾽ ἡμᾶς σκοποῦνταί τ᾽ εὐθέως}{}
\ligneRav{}{}{μὴ μοιχὸς ἔνδον ᾖ τις ἀποκεκρυμμένος·}{}
\ligneRav{398}{}{δρᾶσαι δ᾽ ἔθ᾽ ἡμῖν οὐδὲν \difference{ὧνπερ} καὶ πρὸ τοῦ}{\textbf{398} ὥσπερ R : \lecon{ὧνπερ} Küster}
\ligneRav{}{}{ἔξεστι· τοιαῦθ᾽ οὗτος ἐδίδαξεν κακὰ}{}
\ligneRav{400}{}{τοὺς ἄνδρας ἡμῶν· ὥστ᾽ ἐάν \difference{τις ξυμπλέκῃ}}{\textbf{400} τις πλέκῃ R : \lecon{τις ξυμπλέκῃ} Willems : δή τις πλέκῃ Austin : alii alia}
\ligneRav{}{}{γυνὴ στέφανον, ἐρᾶν δοκεῖ· κἂν ἐκβάλῃ}{}
\ligneRav{}{}{σκεῦός τι κατὰ τὴν οἰκίαν πλανωμένη,}{}
\ligneRav{403}{}{\difference{ἁνὴρ} ἐρωτᾷ· `τῷ κατέαγεν ἡ χύτρα;}{\textbf{403} ἀνὴρ R : \lecon{ἁνὴρ} Brunck}
\ligneRav{}{}{οὐκ ἔσθ᾽ ὅπως οὐ τῷ Κορινθίῳ ξένῳ.'}{}
\ligneRav{405}{}{κάμνει κόρη τις, εὐθὺς \difference{ἁδελφὸς} λέγει·}{\textbf{405} ἀδελφὸς R : \lecon{ἁδελφὸς} Biset}
\ligneRav{}{}{`τὸ χρῶμα τοῦτό μ᾽ οὐκ ἀρέσκει τῆς κόρης.'}{}
\ligneRav{}{}{εἶἑν, γυνή τις ὑποβαλέσθαι βούλεται}{}
\ligneRav{}{}{ἀποροῦσα παίδων, οὐδὲ τοῦτ᾽ ἔστιν λαθεῖν·}{}
\ligneRav{409}{}{\difference{ἅνδρες} γὰρ ἤδη παρακάθηνται πλησίον·}{\textbf{409} ἄνδρες R : \lecon{ἅνδρες} Brunck.}
\ligneRav{410}{}{πρὸς τοὺς γέροντάς θ᾽ οἳ πρὸ τοῦ τὰς μείρακας}{}
\ligneRav{}{}{ἤγοντο, διαβέβληκεν, ὥστ᾽ οὐδεὶς γέρων}{}
\ligneRav{412}{}{γαμεῖν \difference{ἐθέλει} γυναῖκα διὰ τοὔπος τοδὶ}{\textbf{412} θέλει R : \lecon{ἐθέλει} Kappeyne van de Coppello \italique{Obs. crit.} I.2}
\ligneRav{}{}{`δέσποινα γὰρ γέροντι νυμφίῳ γυνή'·}{}
\ligneRav{}{}{εἶτα διὰ τοῦτον ταῖς γυναικωνίτισιν}{}
\ligneRav{415}{}{σφραγῖδας ἐπιβάλλουσιν ἤδη καὶ μοχλοὺς}{}
\ligneRav{}{}{τηροῦντες ἡμᾶς, καὶ προσέτι Μολοττικοὺς}{}
\ligneRav{417}{}{τρέφουσι μορμολυκεῖα τοῖς \difference{μοιχοῖς} κύνας·}{\textbf{417} μυχοῖς R : \lecon{μοιχοῖς} \Mon}
\ligneRav{}{}{καὶ ταῦτα μὲν ξυγγνώσθ᾽. ἃ δ᾽ ἦν ἡμῖν πρὸ τοῦ}{}
\ligneRav{419}{}{αὐταῖς ταμιεῦσαι καὶ προαιρούσαις λαβεῖν}{\textbf{419} ταμιεύεσθαι R : \lecon{ταμιεῦσαι καὶ} Reiske}
\ligneRav{420}{}{ἄλφιτον ἔλαιον οἶνον, οὐδὲ ταῦτ᾽ ἔτι}{}
\ligneRav{}{}{ἔξεστιν. οἱ γὰρ ἄνδρες ἤδη κλῄδια}{}
\ligneRav{}{}{αὐτοὶ φοροῦσι κρυπτὰ κακοηθέστατα}{}
}
\vspace{3pt}\noindent\rule{4.5cm}{0.4pt}\par\vspace{1.5pt}
{\fontsize{7.0pt}{8.3pt}\selectfont
\noindent\textbf{Loci paralleli (vv. 380–422) ---} \textbf{380} Suid. π 1166 περίθου νῦν τόνδε πρότερον \espTest \textbf{381} (χρέμπτεται) \Phryn 126.1 \espTest \textbf{383-4} (φιλοτιμίᾳ -- γυναῖκες) Suidas φ 431 \espTest \textbf{389} Suidas ε 2592 \espTest (ἐπισμῇ) Hesychius ε 5195 ἐπισμῇ \espTest \textbf{390-4} Suidas οι 128 (οἷς θέα pro εἰσὶν θεαταὶ et τραγῳδία pro τραγῳδοὶ) \espTest \textbf{390-1} scholium Plato \italique{Theages} 127c (πῶς pro ποῦ et τραγῳδικοὶ pro τραγῳδοὶ καὶ) $=$ scholium Plato \italique{Gorgias} 457b $=$ scholium Plato \italique{Hippias min.} 365d \espTest (ποῦ -- θεαταὶ) Suidas ε 955 (sed εἰσὶ pro εἰσὶν) \espTest \textbf{390-1} scholium Plato \italique{Theages} 127c $=$ scholium Plato \italique{Gorgias} 457b $=$ scholium Plato \italique{Hippias min.} 365d \espTest \textbf{392} (ἀνδρεραστρίας) Pollux III 71 ; \Phryn 34.18 ; \Eustathe 380.15 \espTest \textbf{393} (οἰνοπίπας) scholium \italique{Aves} 307 \espTest \textbf{395-6} (ὥστ᾽ εὐθὺς -- ἡμᾶς) Suidas ι 275 \espTest \textbf{395} (ἰκρίων) Hesychius ι 501 ἴκρια \espTest \textbf{404} (τῷ Κορινθίῳ ξένῳ) Hesychius κ 3629 Κορίνθιος ξένος \espTest \textbf{413} Menandros \italique{Monostich.} 191 ($=$ Euripides fr. 804.3 ed. Kannicht) ; \Apostol V 93b \espTest \textbf{417} (μορμολυκεῖα) \EtymMagnum 590.53 \espTest \textbf{421-8} (οἱ γὰρ ἄνδρες -- ἐξαψαμένους) Suidas λ 54 (αὐτοὶ, τρεῖς et ᾡκότριψ Εὐριπίδης omissis, Λακωνίκα pro Λακωνίκ᾽, θριπηδέστατ᾽ pro θριπήδεστ᾽ et σφραγῖδα pro σφραγίδια) \espTest \textbf{421-7} (οἱ γὰρ ἄνδρες -- θριπήδεστ᾽) Suidas κ 159 (αὐτοὶ et ᾡκότριψ Εὐριπίδης omissis, Λακωνίκα pro Λακωνίκ᾽ et θριπηδέστατ᾽ pro θριπήδεστ᾽) \espTest \textbf{422} (κακοηθέστατα) \Eustathe 974.37, 1108.25, 1603.52\par
}
\vspace{2.5pt}\noindent\rule{4.5cm}{0.4pt}\par\vspace{1.5pt}
{\fontsize{6.9pt}{8.2pt}\selectfont
\noindent\textbf{Suite (en commentaire) de l'apparat critique \& notes paléographiques (\texttt{f. 167r}) ---} \textbf{Note de personnage descriptive au v. 383 :} Là où \Rac{} avait simplement abrégé (\auDessus{γυ}{ν}), \Rpc{} (suivi par \Mon{}) a développé en marge \lecon{(καλλιλεξία γυνή)} (« la femme au beau langage », glose descriptive passée dans la rubrique de locuteur ; Rogers l'identifie à \textbf{Mika} d'après le v. 760). --- \textbf{Iotacismes $\text{οι} \to \upsilon$ de $\text{R}^1$ :} μυχοτρόπους (v. 392) et μυχοῖς (v. 417) pour μοιχο-, redressés par la Souda et \Mon{}. --- \textbf{Restitution par \textbf{Bibi} du fragment de Sappho (fr. 70/125) dans $\Sigma^{\text{R}}$ ad v. 401 :} La scholie du \italique{Ravennas} au v. 401 conserve un précieux vers de Sappho sous la forme corrompue \lecon{αυταόρα αἱστεφανηπλόκουν} ; \textbf{Louis Hauseux (\textbf{Bibi}, \texttt{vv 1-641.tex}, ll. 3272–3294) restitue l'éolien αὐτὰ κόραις στεφαναπλόκεο} (ou -πλόκην), montrant que αυταόρα αἱ procède de la chute du $\kappa$ et d'une médivision de αὐτὰ κόραις. --- \textbf{v. 390 :} ἐν βραχὺ R : \lecon{ἔμβραχυ} Souda, $\Sigma$ ad Plat. \italique{Gorg.} et \italique{Hipp. min.} --- \textbf{v. 392 :} μυ- R : \lecon{μοι-} (μοιχοτρόπους) Souda \espApp ἀνδρεα- R (haplographie du second $\rho$) : \lecon{ἀνδρερα-} (ἀνδρεραστρίας) Souda. --- \textbf{v. 393 :} οἰνοπότιδας R (substitution d'un mot banal au composé comique rare) : \lecon{οἰνοπίπας} Souda (-πίπους $\Sigma$). --- \textbf{v. 400 :} τις πλέκῃ R : \lecon{τις ξυμπλέκῃ} Alphonse Willems (adopté par L. Hauseux ; δή τις πλέκῃ Austin). --- \textbf{$\Sigma^{\text{R}}$ ad v. 401 (Sappho fr. 70/125) :} αυταόρα αἱστεφανηπλόκουν $\Sigma^{\text{R}}$ : \lecon{αὐτὰ κόραις στεφαναπλόκεο} \lecon{Bibi}.\par
}
\end{minipage}
\end{adjustbox}
\end{adjustwidth}
\newpage
\enTeteD{Codex Ravennas 429 (Biblioteca Classense) --- Fac-similé du f. 167r (vv. 380–422)}
\begin{adjustwidth}{-1.5cm}{-3.0cm}
\begin{center}
\includegraphics[width=16.5cm,height=20.5cm,keepaspectratio]{/workspaces/Thesmophories/dossier_ravennas_thesmophories/folios_individuels_ff162v-177r/10_f167r_classense169r_vv0380-0422.jpg}\par
\vspace{2pt}
{\scriptsize\textbf{Codex Ravennas 429, f. 167r (Classense 169r)} --- ΕΠΕΙΣΟΔΙΟΝ Α', \textbf{vv. 380–422} ($\text{R}^1, \Rpc, \Rdeux$ (\Mon)).}
\end{center}
\end{adjustwidth}
\newpage
\clearpage\ifodd\value{page}\hbox{}\newpage\fi% --- FOLIO 10 (f. 167r) : SPREAD 2.1 (Texte grec établi + Traduction française + Commentaire 1/4 & 2/4) ---
\clearpage\ifodd\value{page}\hbox{}\newpage\fi
\enTeteG{ΘΕΣΜΟΦΟΡΙΑΖΟΥΣΑΙ --- Texte grec établi (f. 167r, vv. 380–422 [1/2])}
\enTeteD{LES THESMOPHORIES --- Traduction française (f. 167r, vv. 380–422 [1/2])}
\begin{adjustwidth}{-2cm}{-2cm}
\noindent\textbf{\footnotesize Texte grec établi --- Folio 10 (\texttt{f. 167r}, vv. 380–422 [1/2]) \hfill ΕΠΕΙΣΟΔΙΟΝ Α'}\\[-2pt]
\noindent\rule{16.5cm}{0.5pt}\\[2pt]
{\fontsize{8.4pt}{9.7pt}\selectfont
\ligneEst{\textbf{Μίκα}}{ἐγώ.}{380}
\ligneEst{\textbf{Κορυφαῖος}}{\decalage{ἐγώ.}περίθου νυν τόνδε πρῶτον πρὶν λέγειν.}{}
\ligneEst{{\tiny\textit{tetr. iamb\cat}}}{σίγα σιώπα, πρόσεχε τὸν νοῦν· χρέμπτεται γὰρ ἤδη}{}
\ligneEst{}{ὅπερ ποιοῦσ᾽ οἱ ῥήτορες. μακρὰν ἔοικε λέξειν.}{}
\ligneEst{\textbf{Μίκα} {\tiny\textit{trim. iamb.}}}{φιλοτιμίᾳ μὲν οὐδεμιᾷ μὰ τὼ θεὼ}{}
\ligneEst{}{λέξουσ᾽ ἀνέστην ὦ γυναῖκες· ἀλλὰ γὰρ}{}
\ligneEst{}{βαρέως φέρω τάλαινα πολὺν ἤδη χρόνον}{385}
\ligneEst{}{προπηλακιζομένας ὁρῶς ἡμᾶς ὑπὸ}{}
\ligneEst{}{Εὐριπίδου τοῦ τῆς λαχανοπωλητρίας}{}
\ligneEst{}{καὶ πολλὰ καὶ παντοῖ᾽ ἀκουούσας κακά.}{}
\ligneEst{}{τί γὰρ οὗτος ἡμᾶς οὐκ ἐπισμῇ τῶν κακῶν;}{}
\ligneEst{}{ποῦ δ᾽ οὐχὶ διαβέβληχ᾽, ὅπουπερ ἔμβραχυ}{390}
\ligneEst{}{εἰσὶν θεαταὶ καὶ τραγῳδοὶ καὶ χοροί,}{}
\ligneEst{}{τὰς μοιχοτρόπους, τὰς ἀνδρεραστρίας καλῶν,}{}
\ligneEst{}{τὰς οἰνοπίπας, τὰς προδότιδας, τὰς λάλους,}{}
\ligneEst{}{τὰς οὐδὲν ὑγιές, τὰς μέγ᾽ ἀνδράσιν κακόν·}{}
\ligneEst{}{ὥστ᾽ εὐθὺς εἰσιόντες ἀπὸ τῶν ἰκρίων}{395}
\ligneEst{}{ὑποβλέπουσ᾽ ἡμᾶς σκοποῦνταί τ᾽ εὐθέως}{}
\ligneEst{}{μὴ μοιχὸς ἔνδον ᾖ τις ἀποκεκρυμμένος.}{}
\ligneEst{}{δρᾶσαι δ᾽ ἔθ᾽ ἡμῖν οὐδὲν ὧνπερ καὶ πρὸ τοῦ}{}
\ligneEst{}{ἔξεστι· τοιαῦθ᾽ οὗτος ἐδίδαξεν κακὰ}{}
\ligneEst{}{τοὺς ἄνδρας ἡμῶν· ὥστ᾽ ἐάν τις ξυμπλέκῃ}{400}
}
\vspace{2pt}
\noindent\rule{16.5cm}{0.6pt}\\[2pt]
\noindent\textbf{\small Commentaire philologique, métrique et scénique (vv. 380–422 --- partie 1/4)}\\[2pt]
\begin{adjustbox}{max totalheight=13.0cm, max width=16.5cm, keepaspectratio}%
\parbox[t]{16.5cm}{
\setlength{\columnsep}{9pt}%
\begin{multicols}{2}
\fontsize{8.4pt}{9.7pt}\selectfont
\setlength{\parindent}{0pt}%
\setlength{\parskip}{1.2pt}%
\noindent\textbf{381} \textbf{χρέμπτεται.} Elle \italique{tousse}, et non elle crache, comme traduisent les dictionnaires. L'ancien français avait un terme excellent, qui s'est perdu avec tant d'autres : « Il se lève, compose sa robe, sa barbe et ses yeux à la modestie, et ayant craché et toussi clair, il commença ». D'Aubigné, \italique{Avent. du baron de Foeneste}, I. IV, ch. 10. (Note de Willems.)
\noindent\textbf{393} \textbf{τὰς οἰνοπίπας,} Encore une fois, c'est le scholiaste, disposant de meilleurs manuscrits, qui nous permet de revenir à la leçon originale : après avoir évoqué Symmaque, il rapproche non sans raison \grec{οἰνοπίπης} de \grec{παιδοπίπης} (on aurait pu rajouter : \grec{γυναικοπίπης}, \grec{παρθενοπίπης} ou encore le \grec{πυροπίπης} / πυρροπίπης des \italique{Cavaliers}, v. 407) ; c'est-à-dire que le suffixe ne renverrait pas aux verbes \grec{πιπίζειν} ou \grec{πιπίσκειν} mais -\grec{οπίπης} serait de la famille de \grec{ὀπιπεύειν}, \italique{guigner}, \italique{guetter}, \italique{épier} (cf. aussi Hésychius (ο 1006) : \grec{ὀπῖπα· ἐξαπατᾷ, ἀπατεών, ἢ ἀπατῶν}).
\noindent Aristophane ne va pas manquer de surenchérir : la difficile relation des femmes au vin sera allègrement mise en scène dans la suite de la pièce.
\noindent\textbf{394} \textbf{τὰς οὐδὲν ὑγιές.} Érasme a compilé l'expression dans ses \italique{Adages} (n° 738, \latin{nihil sanum}).
\noindent\textbf{400} \textbf{τοὺς ἄνδρας ἡμῶν· ὥστ᾽ ἐάν τις πλέκῃ.} Le vers 400 pèche contre le mètre. Willems y supplée en lisant <ξυμ>πλέκῃ ; il s'en justifie dans la \nieme{11} note, reproduite à la fin.
\noindent\textbf{400-1} \textbf{ὥστ᾽ ἐάν τις πλέκῃ / γυνὴ στέφανον, ἐρᾶν δοκεῖ·} Le scholiaste s'est montré court sur ce travers qu'Euripide impute aux femmes tressandières : l'on ne sait de laquelle de ses pièces (perdues) le reproche est tiré.
\noindent Toutefois, tresser une couronne pour son amant(e) est chose bien attestée ; et le scholiaste de citer un vers de Sappho\marqueur{1} :
\citComm{νεωτέρων καὶ ἐρωτικῶν τὸ στεφανοπλοκεῖν. πρὸς τὸ ἔθος, ὅτι ἐστεφανηπλόκουν αἱ παλαιαί.\\
Σαπφώ· αὐτὰ κόραις στεφαναπλόκεο\\
Sappho : «~toi qui te faisais tresser des couronnes par des jeunes filles~».}
\end{multicols}
}%
\end{adjustbox}%
\vfill
\end{adjustwidth}
\newpage
% --- PAGE DE DROITE : Traduction française + Commentaire (2/4) ---
\begin{adjustwidth}{-2cm}{-2cm}
\noindent\textbf{\footnotesize Traduction française --- Folio 10 (\texttt{f. 167r}, vv. 380–422 [1/2]) \hfill ΕΠΕΙΣΟΔΙΟΝ Α'}\\[-2pt]
\noindent\rule{16.5cm}{0.5pt}\\[2pt]
{\fontsize{8.4pt}{9.7pt}\selectfont
\ligneEst{\textbf{Mika}}{\francais{Moi.}}{380}
\ligneEst{\textbf{Coryphée}}{\francais{\espace{1cm}Alors commence par ceindre cette couronne, avant que parler. Paix !}}{}
\ligneEst{}{\francais{Silence! Fais attention ! Car voici qu'elle tousse}}{}
\ligneEst{}{\francais{comme font les orateurs. Elle paraît en avoir long à dire.}}{}
\ligneEst{\textbf{Mika}}{\francais{Aucun motif d'ambition, non par les deux déesses,}}{}
\ligneEst{}{\francais{ne m'a fait lever pour prendre la parole, ô femmes. Mais}}{}
\ligneEst{}{\francais{je suis excédée, depuis longtemps, hélas !}}{385}
\ligneEst{}{\francais{de nous voir traînées dans la boue par}}{}
\ligneEst{}{\francais{Euripide, le fils de la vendeuse d'herbes,}}{}
\ligneEst{}{\francais{et constamment en butte à des injures de tout genre.}}{}
\ligneEst{}{\francais{Car quelle est l'injure dont il ne nous salisse ?}}{}
\ligneEst{}{\francais{Où s'est-il fait faute de nous calomnier, pour peu simplement}}{390}
\ligneEst{}{\francais{qu'il y eût des spectateurs, des tragédiens et des choeurs ?}}{}
\ligneEst{}{\francais{nous traitant de galantes, passionnées d'hommes,}}{}
\ligneEst{}{\francais{sujettes au vin, traitresses, bavardes,}}{}
\ligneEst{}{\francais{rien qui vaille, grand fléau des hommes.}}{}
\ligneEst{}{\francais{Si bien qu'à peine rentrés du théâtre,}}{395}
\ligneEst{}{\francais{ils nous soupçonnent, et guettent aussitôt}}{}
\ligneEst{}{\francais{s'il n'est pas quelque amant caché.}}{}
\ligneEst{}{\francais{Et plus rien ne nous est permis de ce que nous faisions autrefois,}}{}
\ligneEst{}{\francais{tant il a enseigné de vilaines choses}}{}
\ligneEst{}{\francais{à nos époux. Une femme tresse-t-elle}}{400}
}
\vspace{2pt}
\noindent\rule{16.5cm}{0.6pt}\\[2pt]
\noindent\textbf{\small Commentaire philologique, métrique et scénique (vv. 380–422 --- partie 2/4)}\\[2pt]
\begin{adjustbox}{max totalheight=13.0cm, max width=16.5cm, keepaspectratio}%
\parbox[t]{16.5cm}{
\setlength{\columnsep}{9pt}%
\begin{multicols}{2}
\fontsize{7.6pt}{8.8pt}\selectfont
\setlength{\parindent}{0pt}%
\setlength{\parskip}{1.2pt}%
\noindent Athénée (XV, 670d) confirme la coutume d'offrir des couronnes à son être aimé :
\citComm{στεφανοῦσιν δὲ τὰς τῶν ἐρωμένων θύρας ἤτοι τιμῆς χάριν καθαπερεί τινος θεοῦ τὰ πρόθυρα [στεφανοῦσιν αὐτῶν], ἢ οὐ τοῖς ἐρωμένοις ἀλλὰ τῷ Ἔρωτι ποιούμενοι τὴν τῶν στεφάνων ἀνάθεσιν τοῦ μὲν Ἔρωτος τὸν ἐρώμενον ἄγαλμα, τούτου δὲ ναὸν ὄντα τὴν οἴκησιν στεφανοῦσι [τὰ τῶν ἐρωμένων πρόθυρα]: διὰ ταῦτα δὲ καὶ θύουσιν ἔνιοι ἐπὶ ταῖς τῶν ἐρωμένων θύραις.\\
Quant aux couronnes et aux guirlandes dont les amants ornent les portes des personnes qu'ils aiment, ils les y mettent ou pour leur rendre des honneurs, comme couronnant le porche d'une divinité ; ou, c'est à l'amour qu'ils les offrent en les posant, et non aux personnes qui sont aimées ; car ces personnes sont pour les amants les images mêmes du dieu d'amour; et la maison où elles demeurent en étant le temple, c'en est aussi le porche qu'ils ornent de couronnes, en les mettant aux portes de ces personnes. Voilà aussi pourquoi on immole quelquefois des victimes aux portes des personnes qu'on aime.}
\noindent\textbf{401-4} \textbf{κἂν ἐκβάλῃ ... τῷ Κορινθίῳ ξένῳ.} \italique{L'étranger de Corinthe}, c'est Bellérophon, objet d'une coupable passion de la part de la princesse Sthénébée dans la pièce éponyme d'Euripide. Cette dernière, voyant ses avances repoussées par le chaste héros, l'accuse devant le roi son mari d'avoir essayé d'abuser d'elle. Le roi ordonne à Bellérophon de partir pour la Carie. Il y reçoit pour mission d'affronter la Chimère. C'est pendant le temps que Bellérophon est en Carie, et que Sthénébée le croit mort, qu'elle verse des sortes de libations à la mémoire du héros, comme l'on fait au jeu du cottabe avec les personnes aimées. Athénée (X, 427e) nous a conservé ce fragment\marqueur{1} :
\citComm{πεσὸν δέ νιν λέληθεν οὐδὲν ἐκ χερός,\\
ἀλλ᾽ εὐθὺς αὐδᾷ· `τῷ Κορινθίῳ ξένῳ'\\
<La Nourrice>\\
\francais{Rien ne tombe, en lui échappant des mains,\\
qu'elle ne s'écrie aussitôt : « Pour l'hôte de Corinthe ! »}}
\noindent Les maris athéniens prenaient leur femme pour autant de Sthénébées.\\ Ce passage est à rapprocher d'un fragment d'Aristophane, rapporté par Diogène Laërte\marqueur{2} :
\citComm{Τὰ δὲ πεσόντα μὴ ἀναιρεῖσθαι, ὑπὲρ τοῦ ἐθίζεσθαι μὴ ἀκολάστως ἐσθίειν ἢ ὅτι ἐπὶ τελευτῇ τινος· καὶ Ἀριστοφάνης δὲ τῶν ἡρώων φησὶν εἶναι τὰ πίπτοντα, λέγων ἐν τοῖς Ἥρωσι,\\
Μηδὲ γεύεσθ' ἅττ' ἂν ἐντὸς τῆς τραπέζης καταπέσῃ.\\
$[$ Pythagore $]$ interdit de ramasser ce qui tombe de la table pour habituer à manger avec tempérance, ou bien encore parce que cela est destiné aux morts. Ce qui tombe de la table est pour les héros, suivant Aristophane ; car il dit dans les \italique{Héros} :
\francais{« Ne goûtez point à ce qui tombe de la table. »}}
\end{multicols}
}%
\end{adjustbox}%
\vfill
\end{adjustwidth}
\newpage

% --- FOLIO 10 (f. 167r) : SPREAD 2.2 (Texte grec établi + Traduction française + Commentaire 3/4 & 4/4) ---
\clearpage\ifodd\value{page}\hbox{}\newpage\fi
\enTeteG{ΘΕΣΜΟΦΟΡΙΑΖΟΥΣΑΙ --- Texte grec établi (f. 167r, vv. 380–422 [2/2])}
\enTeteD{LES THESMOPHORIES --- Traduction française (f. 167r, vv. 380–422 [2/2])}
\begin{adjustwidth}{-2cm}{-2cm}
\noindent\textbf{\footnotesize Texte grec établi --- Folio 10 (\texttt{f. 167r}, vv. 380–422 [2/2]) \hfill ΕΠΕΙΣΟΔΙΟΝ Α'}\\[-2pt]
\noindent\rule{16.5cm}{0.5pt}\\[2pt]
{\fontsize{8.4pt}{9.7pt}\selectfont
\ligneEst{}{γυνὴ στέφανον, ἐρᾶν δοκεῖ· κἂν ἐκβάλῃ}{}
\ligneEst{}{σκεῦός τι κατὰ τὴν οἰκίαν πλανωμένη,}{}
\ligneEst{}{ἁνὴρ ἐρωτᾷ· `τῷ κατέαγεν ἡ χύτρα;}{}
\ligneEst{}{οὐκ ἔσθ᾽ ὅπως οὐ τῷ Κορινθίῳ ξένῳ.'}{}
\ligneEst{}{κάμνει κόρη τις, εὐθὺς ἁδελφὸς λέγει·}{405}
\ligneEst{}{`τὸ χρῶμα τοῦτό μ᾽ οὐκ ἀρέσκει τῆς κόρης.'}{}
\ligneEst{}{εἶἑν, γυνή τις ὑποβαλέσθαι βούλεται}{}
\ligneEst{}{ἀποροῦσα παίδων, οὐδὲ τοῦτ᾽ ἔστιν λαθεῖν.}{}
\ligneEst{}{ἅνδρες γὰρ ἤδη παρακάθηνται πλησίον·}{}
\ligneEst{}{πρὸς τοὺς γέροντάς θ᾽ οἳ πρὸ τοῦ τὰς μείρακας}{410}
\ligneEst{}{ἤγοντο, διαβέβληκεν, ὥστ᾽ οὐδεὶς γέρων}{}
\ligneEst{}{γαμεῖν ἐθέλει γυναῖκα διὰ τοὔπος τοδὶ}{}
\ligneEst{}{`δέσποινα γὰρ γέροντι νυμφίῳ γυνή'.}{}
\ligneEst{}{εἶτα διὰ τοῦτον ταῖς γυναικωνίτισιν}{}
\ligneEst{}{σφραγῖδας ἐπιβάλλουσιν ἤδη καὶ μοχλοὺς}{415}
\ligneEst{}{τηροῦντες ἡμᾶς, καὶ προσέτι Μολοττικοὺς}{}
\ligneEst{}{τρέφουσι μορμολυκεῖα τοῖς μοιχοῖς κύνας.}{}
\ligneEst{}{καὶ ταῦτα μὲν ξυγγνώσθ᾽. ἃ δ᾽ ἦν ἡμῖν πρὸ τοῦ}{}
\ligneEst{}{αὐταῖς ταμιεῦσαι καὶ προαιρούσαις λαβεῖν}{}
\ligneEst{}{ἄλφιτον ἔλαιον οἶνον, οὐδὲ ταῦτ᾽ ἔτι}{420}
\ligneEst{}{ἔξεστιν. οἱ γὰρ ἄνδρες ἤδη κλῄδια}{}
\ligneEst{}{αὐτοὶ φοροῦσι κρυπτὰ κακοηθέστατα}{422}
}
\vspace{2pt}
\noindent\rule{16.5cm}{0.6pt}\\[2pt]
\noindent\textbf{\small Commentaire philologique, métrique et scénique (vv. 380–422 --- partie 3/4)}\\[2pt]
\begin{adjustbox}{max totalheight=13.0cm, max width=16.5cm, keepaspectratio}%
\parbox[t]{16.5cm}{
\setlength{\columnsep}{9pt}%
\begin{multicols}{2}
\fontsize{7.6pt}{8.8pt}\selectfont
\setlength{\parindent}{0pt}%
\setlength{\parskip}{1.2pt}%
\noindent\textbf{[Notes ad 400-1]} \textbf{(1)} Citation corrompue ; le manuscrit de Ravenne donne : αυταόρα αἱστεφανηπλόκουν. La lecture proposée est de mon cru. Théodore Reinach, dubitatif, reprenait la conjecture de Hermann : αὐτὰρ ὀραῖαι στεφανηπλόκευν. Ce fragment de Sappho est le \nieme{70} de l'édition Reinach-Puech aux Belles Lettres (1937) et le \nieme{125} de l'édition d'Éva-Maria Voigt, \italique{Sappho et Alcaeus} (1971).
\noindent \textbf{(1)} fragment 3, \italique{Sthénébée}, de l'édition Jouan-van Looy, Belles Lettres, t. VIII, \nieme{3} partie $=$ fr. 664 de l'édition Kannicht $=$ fr. 883 de l'édition Mette. Traduction Jouan \& van Looy.
\noindent \textbf{(2)} \italique{Vies}, VIII, 1 (Pythagore), 34. Correspond au fragment 320(305) de l'édition Kassel-Austin (\italique{Poetae Comici Graeci}) des fragments d'Aristophane.
\noindent\textbf{405-6} \textbf{εὐθὺς ἁδελφὸς λέγει· / `τὸ χρῶμα τοῦτό μ᾽ οὐκ ἀρέσκει τῆς κόρης.'} Il la soupçonne d'être enceinte. La pièce d'Euripide à laquelle se réfère Aristophane n'est pas identifiée ; pourtant il semble clair qu'il s'agit d'une citation (ne serait-ce que par l'absence de résolution au vers 406).
\noindent\textbf{406-7} \textbf{γυνή τις ὑποβαλέσθαι βούλεται / ἀποροῦσα παίδων, οὐδὲ τοῦτ᾽ ἔστιν λαθεῖν.} Le thème de la \italique{supposition d'enfants}\marqueur{1} est récurrent dans les \italique{Thesmophories} (cf. vv. 340, 502-516, 564-565). À en croire Satyros dans sa \italique{Vita Euripidis}\marqueur{2}, c'est à Euripide qu'on doit son omniprésence dans la Comédie Nouvelle... et le biographe d'Euripide de donner du grain à moudre à ses contempteurs :
\citComm{πρὸς γ[υ]ναῖκα καὶ πατρὶ πρὸς υ[ἱὸ]ν καὶ θερά[πον]τι πρὸς δ[εσ]πότην, ἢ τ[ὰ] κατὰ τὰς π[ερι]πετείας, β[ια]σμοὺς παρθένων, \textbf{ὑποβολὰς παιδίων}, ἀναγνωρισμοὺς διά τε δακτυλίων καὶ διὰ δεραίων· ταῦτα γάρ ἐστι δήπου τὰ συνέχοντα τα τὴν νεωτέραν κωμῳδίαν, ἃ πρὸς ἄκρον ἤγα[γ]εν Εὐρπίδης, Ὁμήρου [ὄ]ντος ἀρχῆς\\
\francais{$[$le mari$]$ contre sa femme, le père contre son fils, le serviteur contre son maître ; ou lors des péripéties : violences de jeunes filles, \textbf{substitutions d'enfants}, reconnaissances par des anneaux ou des colliers ; hé bien c'est tous ces éléments qui composent la trame de la Nouvelle Comédie, éléments qu'Euripide avait conduit au sommet, Homère en étant le point de départ}}
\noindent La tradition ne nous a rapporté qu'une seule supposition d'enfants chez Euripide : dans les \italique{Phéniciennes} (v. 28-31) :
\citComm{Πολύβου δέ νιν λαβόντες ἱπποβουκόλοι\\
φέρουσ᾽ ἐς οἴκους ἔς τε δεσποίνης χέρας\\
ἔθηκαν. ἣ δὲ τὸν ἐμὸν ὠδίνων πόνον\\
μαστοῖς ὑφεῖτο καὶ πόσιν πείθει τεκεῖν,}
\end{multicols}
}%
\end{adjustbox}%
\vfill
\end{adjustwidth}
\newpage
% --- PAGE DE DROITE : Traduction française + Commentaire (4/4) ---
\begin{adjustwidth}{-2cm}{-2cm}
\noindent\textbf{\footnotesize Traduction française --- Folio 10 (\texttt{f. 167r}, vv. 380–422 [2/2]) \hfill ΕΠΕΙΣΟΔΙΟΝ Α'}\\[-2pt]
\noindent\rule{16.5cm}{0.5pt}\\[2pt]
{\fontsize{8.4pt}{9.7pt}\selectfont
\ligneEst{}{\francais{une couronne, on la croit amoureuse ; laisse-t-elle tomber}}{}
\ligneEst{}{\francais{un ustensile en allant et venant par la maison,}}{}
\ligneEst{}{\francais{le mari demande : « Pour qui brise-t-on la marmite ?}}{}
\ligneEst{}{\francais{Nul doute que ce ne soit pour l'étranger de Corinthe ».}}{}
\ligneEst{}{\francais{Une fille est-elle souffrante, aussitôt son frère de dire :}}{405}
\ligneEst{}{\francais{« Ce teint ne me plaît guère chez la fillette ».}}{}
\ligneEst{}{\francais{Mais passons. Une femme privée d'enfants veut}}{}
\ligneEst{}{\francais{en supposer un ; même cela ne se peut celer,}}{}
\ligneEst{}{\francais{car les maris se tiennent tout proche.}}{}
\ligneEst{}{\francais{Et auprès des vieillards, qui auparavant épousaient les jeunesses,}}{410}
\ligneEst{}{\francais{il nous a calomniées, tellement que pas un homme d'âge}}{}
\ligneEst{}{\francais{ne consent à prendre femme, en raison de ce vers :}}{}
\ligneEst{}{\francais{« Pour un époux âgé, la femme est un tyran ».}}{}
\ligneEst{}{\francais{En outre, à cause de lui, sur les appartements de femmes,}}{}
\ligneEst{}{\francais{ils appliquent des cachets et des barres}}{415}
\ligneEst{}{\francais{pour nous garder, et par là-dessus ils entretiennent des chiens molosses,}}{}
\ligneEst{}{\francais{épouvantails pour les galants.}}{}
\ligneEst{}{\francais{Encore lui passerait-on cela : Mais ce qu'on nous permettait autrefois,}}{}
\ligneEst{}{\francais{de diriger nous-mêmes le ménage et de prendre dans le cellier}}{}
\ligneEst{}{\francais{de la farine, de l'huile, du vin, cela même ne nous est plus}}{420}
\ligneEst{}{\francais{possible : Car les maris portent sur eux de petites clefs}}{}
\ligneEst{}{\francais{secrètes, les plus malicieuses du monde,}}{422}
}
\vspace{2pt}
\noindent\rule{16.5cm}{0.6pt}\\[2pt]
\noindent\textbf{\small Commentaire philologique, métrique et scénique (vv. 380–422 --- partie 4/4)}\\[2pt]
\begin{adjustbox}{max totalheight=13.0cm, max width=16.5cm, keepaspectratio}%
\parbox[t]{16.5cm}{
\setlength{\columnsep}{9pt}%
\begin{multicols}{2}
\fontsize{7.6pt}{8.8pt}\selectfont
\setlength{\parindent}{0pt}%
\setlength{\parskip}{1.2pt}%
\noindent ainsi qu'une allusion dans l'\italique{Alceste} (v. 639) : Admète, désapprouvant l'attitude de son père qui a refusé de donner sa vie pour lui, suppose qu'il a été supposé.
\noindent\textbf{413} \textbf{δέσποινα γὰρ γέροντι νυμφίῳ γυνή.} Stobée (IV, 22a109) nous a conservé la fameuse citation d'Euripide ; elle est tirée du \italique{Phoenix}\marqueur{1} :
\citComm{μοχθηρόν ἐστιν ἀνδρὶ πρεσβύτῃ τέκνα\\
δίδωσιν†, ὅστις οὐκέθ᾽ ὡραῖος γαμεῖ·\\
δέςποινα γὰρ γέροντι νυμφίῳ γυνή\\
\francais{C'est un tracas pour un vieillard <qui est doté d'enfants> de se marier alors qu'il n'en a plus l'âge ; avec un vieil époux, une femme est la maîtresse.}}
\noindent \italique{Quand on est jeune, on aime en fou ; Quand on est vieux, qui aime est fou.}
\noindent\textbf{[Notes ad 406-7]} \textbf{(1)} Quand les femmes font passer pour leur des enfants – mâles le plus souvent ! – auprès de leur mari.
\noindent \textbf{(2)} \italique{The Oxyrhynchus Papyri}, IX.1776, fr. 39, Col. vii, 1-24
\noindent \textbf{(1)} fr. 3, \italique{Phoenix} de l'édition Jouan-van Looy aux Belles-lettres, tome VIII, \nieme{3} partie. Je reprends leur traduction. Ce fragment correspond au n°804 de l'édition Kannicht et 1139 de celle de Mette.
\noindent\textbf{416-7} \textbf{Μολοττικοὺς / τρέφουσι μορμολυκεῖα τοῖς μοιχοῖς κύνας.} C'est le conseil et le vœu d'Hippolyte dans la pièce d'Euripide (vv. 645 et suiv. ; note de Willems) :
\citComm{χρῆν δ᾽ ἐς γυναῖκα πρόσπολον μὲν οὐ περᾶν,\\
ἄφθογγα δ᾽ αὐταῖς συγκατοικίζειν δάκη\\
θηρῶν, ἵν᾽ εἶχον μήτε προσφωνεῖν τινα\\
μήτ᾽ ἐξ ἐκείνων φθέγμα δέξασθαι πάλιν.\\
\francais{Il faudrait que les femmes n'eussent point auprès d'elles de servantes,\\
mais qu'elles fussent servies par de muets animaux,\\
pour qu'elles n'eussent personne à qui parler,\\
ni qui pût à son tour leur adresser la parole.}}
\noindent\textbf{418-20} \textbf{ἃ δ᾽ ἦν ἡμῖν πρὸ τοῦ / αὐταῖς ταμιεύεσθαι προαιρούσαις λαβεῖν / ἄλφιτον ἔλαιον οἶνον.} Je renvoie à la courte mais incisive note (la \nieme{12}) de Willems reproduite à la fin : il y fait un sort à la « déplorable conjecture de Scaliger » (qui avait substitué λαθεῖν à λαβεῖν au vers 419, donnant au passage un sens qui fait contre-sens).
\end{multicols}
}%
\end{adjustbox}%
\vfill
\end{adjustwidth}
\newpage
\clearpage\ifodd\value{page}\hbox{}\newpage\fi
\enTeteG{ΕΠΕΙΣΟΔΙΟΝ Α' --- Codex Ravennas 429, f. 167v (Classense 169v) --- vv. 423–465}
\begin{adjustwidth}{-1.5cm}{-3.0cm}
\begin{adjustbox}{max totalheight=21.0cm, max width=16.5cm, keepaspectratio}
\begin{minipage}{16.5cm}
{\fontsize{7.2pt}{8.5pt}\selectfont
\noindent\textbf{Folio 11 --- Codex Ravennas 429, \texttt{f. 167v} (\texttt{Classense 169v}) : ΕΠΕΙΣΟΔΙΟΝ Α', vv. 423–465}\hfill \textit{Témoins :} $\text{R}^1, \Rpc, \Rdeux$ (\Mon)\par\vspace{2pt}
\hrule height 0.4pt\vspace{2.5pt}
}
{\fontsize{8.1pt}{9.5pt}\selectfont
\ligneRav{423}{\auDessus{γυ}{ν}}{Λακωνίκ᾽ ἄττα, τρεῖς ἔχοντα γομφίους·}{}
\ligneRav{}{}{πρὸ τοῦ μὲν οὖν ἦν ἀλλ᾽ ὑποῖξαι τὴν θύραν}{}
\ligneRav{425}{}{ποιησαμέναισι δακτύλιον τριωβόλου,}{}
\ligneRav{426}{}{νῦν δ᾽ οὗτος αὐτοὺς \difference{ᾡκότριψ} Εὐριπίδης}{\textbf{426} ὡκότριψ R : \lecon{ᾡκότριψ} Bentley : οἱκότριψ Kühner-Blass}
\ligneRav{427}{}{\difference{ἐδίδαξε} θριπήδεστ᾽ ἔχειν σφραγίδια}{\textbf{427} ἐδίδαξεν R : \lecon{ἐδίδαξε} Suid. λ 64}
\ligneRav{}{}{ἐξαψαμένους. νῦν οὖν ἐμοὶ τούτῳ δοκεῖ}{}
\ligneRav{429}{}{ὄλεθρόν τιν᾽ ἡμᾶς κυρκανᾶν \difference{ἁμωσγέπως},}{\textbf{429} ἀμωσγέπως R, Suid. : \lecon{ἁμωσγέπως} Blaydes}
\ligneRav{430}{}{ἢ φαρμάκοισιν ἢ μιᾷ γέ τῳ τέχνῃ,}{}
\ligneRav{}{}{ὅπως ἀπολεῖται. ταῦτ᾽ ἐγὼ φανερῶς λέγω,}{}
\ligneRav{}{}{τὰ δ᾽ ἄλλα μετὰ τῆς γραμματέως συγγράψομαι·}{}
\ligneRav{433}{—}{\hspace{1.1em}ἐγὼ γὰρ αὐτὴ πρῶτον, ἵνα μάλιστ᾽ \difference{οὔπω} (οὐπώ\inRas{τε}{ποτε})}{\textbf{433} οὔπω τε \Rac : οὐπώποτε \Rpc : οὔπω Dindorf}
\ligneRav{}{}{\hspace{1.1em}πολυπλοκωτέρας γυναικὸς}{}
\ligneRav{}{}{\hspace{1.1em}οὐδὲ δεινότερον λεγούσης·}{}
\ligneRav{437}{}{\hspace{1.1em}πάντα γὰρ λέγει δίκαια,}{\textbf{437} \lecon{πάσας} εἰδέας \lecon{ἐξήτασεν} R (sed \lecon{ἰδέας} Suid. ε 22) : πᾶσαν ἰδέαν ἐξετάζει Hermann}
\ligneRav{437}{}{\hspace{1.1em}†\difference{πάσας} δ᾽ \difference{ἰδέας} ἐξήτασεν†,}{\textbf{437} \lecon{πάσας} εἰδέας \lecon{ἐξήτασεν} R (sed \lecon{ἰδέας} Suid. ε 22) : πᾶσαν ἰδέαν ἐξετάζει Hermann}
\ligneRav{438}{}{\hspace{1.1em}πάντα δ᾽ ἐβάστασεν φρενὶ πυκνῶς τε}{\textbf{438} ἐβάστασεν R : \lecon{-σε} Suid.}
\ligneRav{439}{}{\hspace{1.1em}ποικίλους λόγους \difference{ἀνηῦρεν}}{\textbf{439} ἀνεῦρεν R, Suid. : \lecon{ἀνηῦρεν} Meineke}
\ligneRav{439}{}{\hspace{1.1em}εὖ διεζητημένους·}{\textbf{439} ἀνεῦρεν R, Suid. : \lecon{ἀνηῦρεν} Meineke}
\ligneRav{440}{}{\hspace{1.1em}ὥστ᾽ ἂν εἰ λέγοι παρ᾽ αὐτῆς}{\textbf{440} \lecon{αὐτῆς} R : αὐτὴν Zanetti}
\ligneRav{441}{}{\hspace{1.1em}\difference{Ξενοκλέης} ὁ Καρκίνου, δο-}{\textbf{441} Ξενοκλῆς R : \lecon{Ξενοκλέης} Bentley}
\ligneRav{441}{}{\hspace{1.1em}κεῖν ἂν αὐτόν, ὡς ἐγᾦμαι,}{\textbf{441} Ξενοκλῆς R : \lecon{Ξενοκλέης} Bentley}
\ligneRav{}{}{\hspace{1.1em}πᾶσιν ὑμῖν}{}
\ligneRav{}{}{\hspace{1.1em}ἄντικρυς μηδὲν λέγειν·}{}
\ligneRav{443}{γυνὴ ἑτέρα}{ὀλίγων \difference{ἕνεκα καὐτὴ} παρῆλθον ῥημάτων·}{\textbf{443} γυνὴ ἑτέρα \Rpc \espApp ἕνεκ᾽ αὐτὴ R : \lecon{ἕνεκα καὐτὴ} Porson}
\ligneRav{}{}{τὰ μὲν γὰρ ἄλλ᾽ αὕτη κατηγόρηκεν εὖ·}{}
\ligneRav{445}{}{ἃ δ᾽ ἐγὼ πέπονθα, ταῦτα λέξαι βούλομαι·}{}
\ligneRav{446}{}{ἐμοὶ γὰρ \difference{ἁνὴρ} ἀπέθανεν μὲν \difference{ἔγκυος}}{\textbf{446} ἀνὴρ R : \lecon{ἁνὴρ} Brunck \espApp ἐν Κύπρῳ R : \lecon{ἔγκυος} Bibi}
\ligneRav{}{}{παιδάρια πέντε καταλιπών, ἁγὼ μόλις}{}
\ligneRav{}{}{στεφανηπλοκοῦσ᾽ ἔβοσκον ἐν ταῖς μυρρίναις·}{}
\ligneRav{}{}{τέως μὲν οὖν ἀλλ᾽ ἡμικάκως ἐβοσκόμην·}{}
\ligneRav{450}{}{νῦν δ᾽ οὗτος ἐν \difference{ταῖσιν} τραγῳδίαις ποιῶν}{\textbf{450} ταῖσι R : \lecon{ταῖσιν} Brunck}
\ligneRav{}{}{τοὺς ἄνδρας ἀναπέπεικεν οὐκ εἶναι θεούς·}{}
\ligneRav{}{}{ὥστ᾽ οὐκέτ᾽ ἐμπολῶμεν οὐδ᾽ εἰς ἥμισυ·}{}
\ligneRav{453}{}{νῦν οὖν \difference{ἁπάσαισιν} παραινῶ καὶ λέγω}{\textbf{453} ἁπάσαισι R : \lecon{ἁπάσαισιν} Gellius \italique{Ν.Α.} 15.20}
\ligneRav{}{}{τοῦτον κολάσαι τὸν ἄνδρα πολλῶν οὕνεκα·}{}
\ligneRav{455}{}{ἄγρια γὰρ ἡμᾶς ὦ γυναῖκες δρᾷ κακά,}{}
\ligneRav{456}{}{ἅτ᾽ ἐν ἀγρίοισι \difference{τοῖς} λαχάνοις αὐτὸς τραφείς·}{\textbf{456} \lecon{τοῖς} Plut., \italique{Mor.}, 853C, Gellius \italique{N.Α.}15.20 : om. R}
\ligneRav{}{}{ἀλλ᾽ εἰς ἀγορὰν ἄπειμι· δεῖ γὰρ ἀνδράσιν}{}
\ligneRav{}{}{πλέξαι στεφάνους συνθηματιαίους εἴκοσιν·}{}
\ligneRav{}{χορ/}{\hspace{1.1em}ἕτερον αὖ τι λῆμα τοῦτο}{}
\ligneRav{460}{}{\hspace{1.1em}κομψότερον ἔτ᾽ ἢ τὸ πρότερον}{}
\ligneRav{460}{}{\hspace{1.1em}ἀναπέφηνεν·}{}
\ligneRav{461}{}{\hspace{1.1em}οἷα κατεστωμύλατο}{\textbf{461} \lecon{κατεστωμύλατο} R : κἀστωμύλατο Dobree}
\ligneRav{462}{}{\hspace{1.1em}οὐκ \difference{ἄκαιρα}, φρένας ἔχουσα}{\textbf{462} ἄκερα R : \lecon{ἄκαιρα} Zanetti}
\ligneRav{463}{}{\hspace{1.1em}καὶ πολύπλοκον νόημ᾽, οὐδ᾽}{\textbf{463} \lecon{πολύπλοκον} R : τι πολύπλοκον Dobree : πολύστροφον Hermann}
\ligneRav{}{}{\hspace{1.1em}ἀσύνετ᾽ ἀλλὰ πιθανὰ πάντα·}{}
\ligneRav{465}{}{\hspace{1.1em}δεῖ δὲ ταύτης τῆς ὕβρεως ἡ-}{}
\ligneRav{465}{}{\hspace{1.1em}μῖν τὸν ἄνδρα}{}
\ligneRav{465}{}{\hspace{1.1em}περιφανῶς δοῦναι δίκην·}{}
}
\vspace{3pt}\noindent\rule{4.5cm}{0.4pt}\par\vspace{1.5pt}
{\fontsize{7.0pt}{8.3pt}\selectfont
\noindent\textbf{Loci paralleli (vv. 423–465) ---} \textbf{423} (γομφίους) \Phryn 61.1 \espTest \textbf{426} (οἰκότριψ $=$ οἰκοτραφής) Moeris ο 25 ; (ᾡκότριψ) Hesychius ο 265-266 ; \AelDion ο 8-9 \espTest \textbf{428-9} (ἐμοὶ -- ἁμωσγέπως) Suidas κ 2771 (τινα pro τιν᾽ ἡμᾶς) \espTest \textbf{429} (κυρκανᾶν) Photius κ 1254 κυρκανᾶν $=$ Hesychius κ 4691 κυρκανᾷ ; \Eustathe 1213.50 \espTest (ἁμωσγέπως) Hesychius α 4182 ἁμωσγέπως ; \Syn α 449 $=$ Suidas α 1645 \espTest \textbf{432} (μετὰ τῆς γραμματέως συγγράψομαι) \Eustathe 1484.53 \espTest (τῆς γραμματέως) Pollux IV 19 \espTest \textbf{437-9} (πάσας -- ἀνηῦρεν) Suidas ε 22 πάσας δ' ἰδέας ἐξήτασε πάντα δ' ἐβάστασε φρενί, πυκνῶς τε ποικίλως λόγους ἀνεῦρεν \espTest \textbf{438} (πυκνῶς) Hesychius π 4335 πυκνῶς \espTest \textbf{442} (ἄντικρυς) Photius α 2210 ἄντικρυς ; \AelDion α 107 \espTest \textbf{446} Hesychius ε 239 ἔγκυος· ἐγγαστρωμένη γυνή \espTest \textbf{447-9} (μόλις -- τέως μὲν οὖν ἀλλ᾽) scholium\expo{T} Plato \italique{Hipparch.} 229d (μυρίναις pro μυρρίναις et μέν οὐκ ἄν pro μὲν οὖν ἀλλ᾽) \espTest \textbf{448} (μυρρίνη $=$ μυρσίνη) Moeris μ 23 \espTest \textbf{449} Suidas τ 321 \espTest (ἡμικάκως) Pollux VI 161 \espTest \textbf{453-6} Aulus Gellius XV 20.7 \espTest \textbf{455-6} Plut. \italique{Moral.}, \italique{Compar. Arist. et Menandr.} (853c) \espTest \textbf{457-8} (δεῖ γὰρ -- εἴκοσιν) Suidas σ 1591 (sed ἀνδράσι et εἴκοσι pro ἀνδράσιν et εἴκοσιν) \espTest \textbf{458} Athenaeus XV 680c \espTest (στεφάνους συνθηματιαίους) Pollux VII 200 \espTest (συνθηματιαίους) Hesychius σ 2624 συνθηματιαῖοι \espTest \textbf{459-60} (ἕτερον -- ἀναπέφηνεν) Suidas λ 441 (ἀναπέφηνε pro ἀναπέφηνεν)\par
}
\vspace{2.5pt}\noindent\rule{4.5cm}{0.4pt}\par\vspace{1.5pt}
{\fontsize{6.9pt}{8.2pt}\selectfont
\noindent\textbf{Suite (en commentaire) de l'apparat critique \& notes paléographiques (\texttt{f. 167v}) ---} \textbf{Correction au v. 433 :} \Rac{} avait écrit οὔπω τε, corrigé par \Rpc{} en οὐπώποτε (Dindorf restitue οὔπω). --- \textbf{Note d'interlocuteur au v. 443 :} \Rpc{} inscrit γυνὴ ἑτέρα. --- \textbf{Conjecture remarquable de \textbf{Bibi} au v. 446 :} ἀνὴρ μὲν ἐν Κύπρῳ R $\to$ \lecon{ἁνὴρ μὲν ἔγκυος} (ἔγκυος se rapportant à la femme ou construit avec l'idée que son mari l'a laissée enceinte de son 5\ieme{} enfant avant de mourir, au lieu de la localisation géographique hors de propos ἐν Κύπρῳ). --- \textbf{v. 437 :} \lecon{πάσας} εἰδέας \lecon{ἐξήτασεν} R (mais \lecon{ἰδέας} Souda $\varepsilon\ 22$, sans iotacisme ; πᾶσαν ἰδέαν ἐξετάζει Hermann). --- \textbf{v. 443 :} \lecon{γυνὴ ἑτέρα} \Rpc{} \espApp ἕνεκ' αὐτὴ R (haplographie de $\kappa\alpha-$) : \lecon{ἕνεκα καὐτὴ} Porson. --- \textbf{v. 446 :} ἀνὴρ R : \lecon{ἁνὴρ} Brunck \espApp ἐν Κύπρῳ R : \lecon{ἔγκυος} \lecon{Bibi}. --- \textbf{v. 456 :} \lecon{τοῖς} Plutarque, \italique{Moralia} 853C, Aulu-Gelle 15.20 : om. R (devant λαχάνοις).\par
}
\end{minipage}
\end{adjustbox}
\end{adjustwidth}
\newpage
\enTeteD{Codex Ravennas 429 (Biblioteca Classense) --- Fac-similé du f. 167v (vv. 423–465)}
\begin{adjustwidth}{-1.5cm}{-3.0cm}
\begin{center}
\includegraphics[width=16.5cm,height=20.5cm,keepaspectratio]{/workspaces/Thesmophories/dossier_ravennas_thesmophories/folios_individuels_ff162v-177r/11_f167v_classense169v_vv0423-0465.jpg}\par
\vspace{2pt}
{\scriptsize\textbf{Codex Ravennas 429, f. 167v (Classense 169v)} --- ΕΠΕΙΣΟΔΙΟΝ Α', \textbf{vv. 423–465} ($\text{R}^1, \Rpc, \Rdeux$ (\Mon)).}
\end{center}
\end{adjustwidth}
\newpage
\clearpage\ifodd\value{page}\hbox{}\newpage\fi% --- FOLIO 11 (f. 167v) : SPREAD 2.1 (Texte grec établi + Traduction française + Commentaire 1/8 & 2/8) ---
\clearpage\ifodd\value{page}\hbox{}\newpage\fi
\enTeteG{ΘΕΣΜΟΦΟΡΙΑΖΟΥΣΑΙ --- Texte grec établi (f. 167v, vv. 423–465 [1/4])}
\enTeteD{LES THESMOPHORIES --- Traduction française (f. 167v, vv. 423–465 [1/4])}
\begin{adjustwidth}{-2cm}{-2cm}
\noindent\textbf{\footnotesize Texte grec établi --- Folio 11 (\texttt{f. 167v}, vv. 423–465 [1/4]) \hfill ΕΠΕΙΣΟΔΙΟΝ Α'}\\[-2pt]
\noindent\rule{16.5cm}{0.5pt}\\[2pt]
{\fontsize{8.8pt}{10.2pt}\selectfont
\ligneEst{}{Λακωνίκ᾽ ἄττα, τρεῖς ἔχοντα γομφίους.}{423}
\ligneEst{}{πρὸ τοῦ μὲν οὖν ἦν ἀλλ᾽ ὑποῖξαι τὴν θύραν}{}
\ligneEst{}{ποιησαμέναισι δακτύλιον τριωβόλου,}{425}
\ligneEst{}{νῦν δ᾽ οὗτος αὐτοὺς ᾡκότριψ Εὐριπίδης}{}
\ligneEst{}{ἐδίδαξε θριπήδεστ᾽ ἔχειν σφραγίδια}{}
\ligneEst{}{ἐξαψαμένους. νῦν οὖν ἐμοὶ τούτῳ δοκεῖ}{}
\ligneEst{}{ὄλεθρόν τιν᾽ ἡμᾶς κυρκανᾶν ἁμωσγέπως,}{}
\ligneEst{}{ἢ φαρμάκοισιν ἢ μιᾷ γέ τῳ τέχνῃ,}{430}
\ligneEst{}{ὅπως ἀπολεῖται. ταῦτ᾽ ἐγὼ φανερῶς λέγω,}{}
\ligneEst{}{τὰ δ᾽ ἄλλα μετὰ τῆς γραμματέως συγγράψομαι.}{}
\ligneEst{\textbf{Χορός} {\tiny\textit{στροφή / par.}}}{οὔπω ταύτης ἤκουσα}{}
\ligneEst{{\tiny\textit{$2$ troch.}}}{πολυπλοκωτέρας γυναικὸς}{}
}
\vspace{2pt}
\noindent\rule{16.5cm}{0.6pt}\\[2pt]
\noindent\textbf{\small Commentaire philologique, métrique et scénique (vv. 423–465 --- partie 1/8)}\\[2pt]
\begin{adjustbox}{max totalheight=15.2cm, max width=16.5cm, keepaspectratio}%
\parbox[t]{16.5cm}{
\setlength{\columnsep}{9pt}%
\begin{multicols}{2}
\fontsize{7.0pt}{8.0pt}\selectfont
\setlength{\parindent}{0pt}%
\setlength{\parskip}{1.2pt}%
\noindent\textbf{424} \textbf{ὑποῖξαι τὴν θύραν.} Il s'agit de la porte de l'appartement (cf. v. 414) et non de celle du cellier.
\noindent\textbf{426} \textbf{ᾡκότριψ Εὐριπίδης.} Ce terme, \italique{tâte-poule}, dont un Grec aurait saisi le sens métaphorique (cf. les verbes βλιμάζειν, σκιμαλίζειν), rend bien le mot ᾡκότριψ. Cf. ἀστύτριψ, \italique{pilier de cité}, Θησειότριψ, \italique{pilier du Théseion} (fragm. du Polyides), πορνότριψ, etc. (Note de Willems.)
\noindent\textbf{429} \textbf{ὄλεθρόν τιν᾽ ἡμᾶς κυρκανᾶν ἁμωσγέπως.} Κυρκανᾶν est un synonyme très rare de κυκᾶν et se dit proprement d'un cuisinier préparant une mixture\marqueur{1}. Il est pris ici au sens de \italique{comploter la perte} de quelqu'un. Il est clair que l'expression ὄλεθρόν κυρκανᾶν est faite sur le modèle de κυρκανᾶν φάρμακα (cf. \italique{Ploutos} 302, 309 : τὰ φάρμακ᾽ ἀνακυκῶσαν) [...] Pour le ton, le verbe κυρκανᾶν paraît correspondre à nos métaphores familières \italique{cuisiner}, \italique{mitonner une affaire}. (Note de Taillardat, §421.)
\noindent\textbf{432} \textbf{τῆς γραμματέως.} Aristophane faisait-il sourire sa mâle assemblée lorsqu'il parlait de \italique{la} greffièr\italique{e} ? c'est ce que suggère Pollux (IV, 19) :
\citComm{Ἀριστοφάνης δὲ παίζει ἐν Θεσμοφοριαζούσαις, λέγων ἡ γραμματεύς.}
\noindent\textbf{[Notes ad 429]} \textbf{(1)} Cf. Épinicus 1,8 : συγκυρκανᾶν (préparation d'une boisson faite de farine, de miel et de vin mêlés).
\noindent\textbf{433-42} \textbf{} Ce chant du chœur trouvera son écho aux vers 520-530 : semblable réplique au point de vue métrique\marqueur{1} qui permettait aux deux passages d'être chantés sur la même mélodie. Là comme ici, le chœur s'exprime à la première personne.
\noindent\textbf{[Notes ad 433-42]} \textbf{(1)} À une exception près : au vers 437b où le texte est probablement corrompu.
\noindent\textbf{441} \textbf{Ξενοκλῆς ὁ Καρκίνου.} Xénoclès, le plus célèbre des trois carcinites était un poète tragique non dépourvu de talent puisqu'en 415 il l'emporta sur Euripide\marqueur{1}. Mais on lui reprochait de chercher à captiver son auditoire par des moyens peu relevés. Platon le comique l'appelle δωδεκαμήχανος [l'homme] \italique{aux douze postures}, surnom qu'on donnait à la courtisane Cyrène\marqueur{2}. Aristophane le traite de μηχανοδίφης, inventeur de trucs. Il résulte d'un passage d'Athénée (IV, 134d), mis en rapport avec le nôtre, que Xénoclès cumulait avec la profession d'acteur et de poète tragique, celle d'orateur. Ce fut aussi le cas d'Eschine. Cela n'a rien d'étonnant, le métier d'acteur étant une excellente préparation à celui d'orateur. Il n'y a qu'à se souvenir de Démosthène et de ses cailloux. (Note de Willems ; voir aussi la note (t. I pp. 553-558) qu'il a consacrée au banquet de Philocléon conté par Xanthias (\italique{Guêpes}, vv. 1299-1323) où il est question des carcinites.)
\noindent\textbf{[Notes ad 441]} \textbf{(1)} Aristophane se montre toutefois sévère avec Xénoclès : il est méchant comme ses vers (v. 169) et son art comme sa personne ne méritent pas de survivre (\italique{Grenouilles}, v. 86).
\noindent \textbf{(2)} Voir la note au vers 235.
\noindent\textbf{442} \textbf{πᾶσιν ὑμῖν.} Les spectateurs.
\noindent\textbf{446} \textbf{ἐμοὶ γὰρ ἁνὴρ ἀπέθανεν μὲν ἐν Κύπρῳ.} La précision du lieu de décès du mari (ἁνὴρ ἀπέθανεν ἐν Κύπρῳ) n'a pas manqué d'intriguer les modernes scholiastes.
\noindent Leur inclination naturelle étant portée au scabreux, ils ont pour une bonne partie voulu y voir un tour signifiant que la matrone athénienne, mère de cinq enfants, avait tué son mari au déduit avec ses appétits démesurés ; «~mourir de sa belle mort~», comme l'heureux \italique{Félix Faurenicator}.
\noindent L'hypothèse est plaisante mais gratuite et donc à rejeter ! tous les lexiques consultés ne soufflent mot d'une telle expression qu'on aurait nécessairement relevée.
\noindent D'autres se sont essayés à prendre ce vers littéralement. Là encore, ils se sont heurtés à un mur : les dernière expéditions sur l'île de Chypre feraient de notre procurette une septuagénaire au moins !.. et comme il ne semble pas y avoir de jeux de mots, il ne reste qu'une seule voie : une erreur de transcription du texte.
\end{multicols}
}%
\end{adjustbox}%
\vfill
\end{adjustwidth}
\newpage
% --- PAGE DE DROITE : Traduction française + Commentaire (2/8) ---
\begin{adjustwidth}{-2cm}{-2cm}
\noindent\textbf{\footnotesize Traduction française --- Folio 11 (\texttt{f. 167v}, vv. 423–465 [1/4]) \hfill ΕΠΕΙΣΟΔΙΟΝ Α'}\\[-2pt]
\noindent\rule{16.5cm}{0.5pt}\\[2pt]
{\fontsize{8.8pt}{10.2pt}\selectfont
\ligneEst{}{\francais{je ne sais quelles laconiennes à trois dents.}}{423}
\ligneEst{}{\francais{Autrefois du moins il tenait à nous d'ouvrir furtivement la porte,}}{}
\ligneEst{}{\francais{en nous faisant graver un anneau de trois oboles.}}{425}
\ligneEst{}{\francais{Aujourd'hui ce tâte-poule d'Euripide}}{}
\ligneEst{}{\francais{leur a enseigné l'emploi de cachets vermoulus}}{}
\ligneEst{}{\francais{qu'ils portent suspendus sur eux. Maintenant donc j'opine}}{}
\ligneEst{}{\francais{que nous combinions de manière ou d'autre un procédé pour le détruire,}}{}
\ligneEst{}{\francais{soit par des drogues, soit par un moyen quelconque,}}{430}
\ligneEst{}{\francais{afin qu'il meure. Ceci, je le dis à pleine bouche.}}{}
\ligneEst{}{\francais{Le reste, je le coucherai par écrit, avec l'aide de la greffière.}}{}
\ligneEst{\textbf{Chœur}}{\francais{Jamais je n'entendis}}{}
\ligneEst{}{\francais{femme plus féconde en retours}}{}
}
\vspace{2pt}
\noindent\rule{16.5cm}{0.6pt}\\[2pt]
\noindent\textbf{\small Commentaire philologique, métrique et scénique (vv. 423–465 --- partie 2/8)}\\[2pt]
\begin{adjustbox}{max totalheight=15.2cm, max width=16.5cm, keepaspectratio}%
\parbox[t]{16.5cm}{
\setlength{\columnsep}{9pt}%
\begin{multicols}{2}
\fontsize{7.0pt}{8.0pt}\selectfont
\setlength{\parindent}{0pt}%
\setlength{\parskip}{1.2pt}%
\noindent (C'est d'autant plus vraisemblable que l'expression « ἀπέθανεν ἐν Κύπρῳ », si elle s'était véritablement trouvée dans le texte original, n'aurait pas manqué d'intriguer les commentateurs anciens qui, même sans y rien connaître, auraient ouvert les vannes de leur verve intarissable... au contraire : ce vers n'a étrangement pas laissé la moindre scholie ancienne derrière lui.)\\ L'intervention de la tresseuse de couronnes n'a donc rien de burlesque : elle essaie au contraire d'introduire du pathétique. Or, qu'y a-t-il de plus banal pour apitoyer son auditoire à peu de frais que de sortir le lieu-commun de la femme qui se retrouve veuve alors qu'elle est enceinte (ἔγκυος) ?
\citComm{ἐμοὶ γὰρ ἁνὴρ ἀπέθανεν μὲν \textbf{ἔγκυος}\\
Mon mari mourut alors que j'étais enceinte \\
Me laissant cinq enfants à charge etc. etc.}
\noindent\textbf{450-1} \textbf{ἐν ταῖσιν τραγῳδίαις ποιῶν / τοὺς ἄνδρας ἀναπέπεικεν οὐκ εἶναι θεούς·} Petit florilège des \italique{\grec{γνῶμαι}} théoclastes d'Euripide :\\ \italique{Hécube}, v. 489-492 :
\citComm{Ταλθύβιος — ὦ Ζεῦ, τί λέξω; πότερά σ᾽ ἀνθρώπους ὁρᾶν;\\
ἢ δόξαν ἄλλως τήνδε κεκτῆσθαι μάτην,\\
ψευδῆ, δοκοῦντας δαιμόνων εἶναι γένος\\
τύχην δὲ πάντα τἀν βροτοῖς ἐπισκοπεῖν;}
\noindent \italique{Béllérophon}\marqueur{1} :
\citComm{Βελλεροφόντης — φῆσίν τις εἶναι δῆτ᾽ ἐν οὐρανῷ θεούς;\\
οὐκ εἰσίν, οὐκ εἴσ᾽,}
\noindent Prologue de \italique{Mélanippe la sage}\marqueur{2} (douteux) :
\citComm{Ζεὺς ὅςτις ὁ Ζεύς, οὐ γὰρ οἶδα πλὴν λόγῳ\\
La véritable version euripidéenne étant peut-être\marqueur{3} :\\
Μελάνιππη — Ζεύς, ὡς λέλεκται τῆς ἀληθείας ὕπο,\\
Ἕλλην᾽ ἔτιχθ᾽, ὃς ἐξέφυσεν Αἴολον·}
\noindent\textbf{[Notes ad 450-1]} \textbf{(1)} Deux premières lignes du \nieme{10} fragment de l'édition Jouan-van Looy aux Belles Lettres ; fr. 286 de l'édition Kannicht ; fr. 385 Mette.
\noindent \textbf{(2)} Fragment douteux car on recense différentes versions de ce début de prologue. Citation apocryphe ? retouche d'Euripide d'un vers ayant fait scandale ? quoi qu'il en soit, Jouan \& van Looy le mettent en tête de leur \italique{Mélanippe} sans le comptabiliser ; fr. 480 de l'édition Kannicht.
\noindent \textbf{(3)} Premier fragment de l'édition Jouan-van Looy ; 481 chez Kannicht et 665a-c chez Mette.
\noindent\textbf{455-6} \textbf{Dossier sur la mère d'Euripide, la vendeuse d'herbes (λαχανοπωλήτρια).} (\italique{Acharn.} vv. 478-479 ; \italique{Cavaliers} vv. 19-20 ; \italique{Thesm.} vv. 455-456, 910 ; \italique{Grenouilles} vv. 840, 946-947.)
\noindent Aristophane plaisante Euripide à maintes reprises sur la profession de sa mère, vendeuse d'herbes\marqueur{1}. Il fit ses choux gras d'une herbe en particulier : ἡ σκάνδιξ, \italique{le scandix peigne de Vénus}. Les commentateurs, anciens comme modernes, se sont interrogés depuis lors pour savoir s'il fallait prendre littéralement cette information, ou bien de façon imagée.
\noindent Sans parvenir à aucune conclusion convaincante ; la raison en est que, pour comprendre les attaques d'Aristophane, il fallait séparer le problème en deux :\\ --- quand il évoque le σκάνδιξ, Aristophane moque le style d'Euripide, tout en finasseries (nous allons voir que le verbe διασκανδικίζειν a le même double sens que \italique{épicer} en français ; le sens propre : on épice une galette ; le sens imagé : on épice un discours) ;\\ --- quand il évoque la mère d'Euripide, la verdurière, Aristophane fait entendre que ce dernier est un \italique{fils de la terre}, \latin{terrae filius} ; dont les origines sont obscures et l'arbre généalogique impossible à établir ; bref : Euripide est un bâtard !
\end{multicols}
}%
\end{adjustbox}%
\vfill
\end{adjustwidth}
\newpage

% --- FOLIO 11 (f. 167v) : SPREAD 2.2 (Texte grec établi + Traduction française + Commentaire 3/8 & 4/8) ---
\clearpage\ifodd\value{page}\hbox{}\newpage\fi
\enTeteG{ΘΕΣΜΟΦΟΡΙΑΖΟΥΣΑΙ --- Texte grec établi (f. 167v, vv. 423–465 [2/4])}
\enTeteD{LES THESMOPHORIES --- Traduction française (f. 167v, vv. 423–465 [2/4])}
\begin{adjustwidth}{-2cm}{-2cm}
\noindent\textbf{\footnotesize Texte grec établi --- Folio 11 (\texttt{f. 167v}, vv. 423–465 [2/4]) \hfill ΕΠΕΙΣΟΔΙΟΝ Α'}\\[-2pt]
\noindent\rule{16.5cm}{0.5pt}\\[2pt]
{\fontsize{8.8pt}{10.2pt}\selectfont
\ligneEst{{\tiny\textit{$2$ troch.}}}{οὐδὲ δεινότερον λεγούσης.}{}
\ligneEst{{\tiny\textit{$2$ troch.}}}{πάντα γὰρ λέγει δίκαια,}{}
\ligneEst{{\tiny\textit{$2$ iamb.}}}{†πάσας δ᾽ ἰδέας ἐξήτασεν†,}{}
\ligneEst{{\tiny\textit{$2$ troch.}}}{πάντα δ᾽ ἐβάστασεν φρενὶ πυκνῶς τε}{}
\ligneEst{{\tiny\textit{$2$ troch.}}}{ποικίλους λόγους ἀνηῦρεν}{}
\ligneEst{{\tiny\textit{$2$ troch\cat}}}{εὖ διεζητημένους·}{}
\ligneEst{{\tiny\textit{$2$ troch.}}}{ὥστ᾽ ἂν εἰ λέγοι παρ᾽ αὐτῆς}{440}
\ligneEst{{\tiny\textit{$2$ troch.}}}{Ξενοκλέης ὁ Καρκίνου, δο-}{}
\ligneEst{{\tiny\textit{$2$ troch.}}}{κεῖν ἂν αὐτόν, ὡς ἐγᾦμαι,}{}
\ligneEst{{\tiny\textit{troch.}}}{πᾶσιν ὑμῖν}{}
\ligneEst{{\tiny\textit{$2$ troch\cat}}}{ἄντικρυς μηδὲν λέγειν.}{}
\ligneEst{\textbf{Στεφανόπωλις} {\tiny\textit{trim. iamb.}}}{ὀλίγων ἕνεκα καὐτὴ παρῆλθον ῥημάτων.}{}
}
\vspace{2pt}
\noindent\rule{16.5cm}{0.6pt}\\[2pt]
\noindent\textbf{\small Commentaire philologique, métrique et scénique (vv. 423–465 --- partie 3/8)}\\[2pt]
\begin{adjustbox}{max totalheight=15.2cm, max width=16.5cm, keepaspectratio}%
\parbox[t]{16.5cm}{
\setlength{\columnsep}{9pt}%
\begin{multicols}{2}
\fontsize{7.0pt}{8.0pt}\selectfont
\setlength{\parindent}{0pt}%
\setlength{\parskip}{1.2pt}%
\noindent Cette note risque fort de trop tirer en longueur ; le dossier sur la mère d'Euripide et son σκάνδιξ est en effet assez épais. En voici le plan :\\ --- je montre d'abord que le vers 910 des \italique{Thesmophories} ne fait pas partie de ce dossier et que les `ἰφύων', \italique{lavandes}\marqueur{2}, n'ont jamais apparu dans le texte d'Aristophane ;\\ --- il y a par contre une autre herbe que le scandix à rétablir au vers 20 des \italique{Cavaliers} : l'ἀπόκυνον (Hésychius α 6416) ;\\ --- le σκάνδιξ et le verbe διασκανδικίζειν renvoient au style et à la façon de parler ;\\ --- descendre d'une verdurière, c'est être un \italique{fils de la terre}, \latin{terrae filius} : un bâtard ;\\ --- c'est la bâtardise d'Euripide qui explique sa rancœur vis-à-vis des femmes ; cela ressort de \italique{Thesm.} vv. 455-456 grâce à l'identification d'un proverbe (αὐτὸς τραφείς...) analogue à notre \italique{tel père, tel fils} ;\\ --- la note se conclue par une tentative pseudétymologique sur le σκάνδιξ pour expliquer les images qu'il a générées.
\noindent Pour le vers 910 \italique{Thesmophories}, je renvoie le lecteur au commentaire \italique{ad locum}.
\noindent Il n'y a aucune allusion aux herbes d'Euripide, qui ont envahi bien mal à propos jusqu'aux sources les plus fiables de ce vers.
\noindent Par contre, il faut en rétablir une aux vers 19-20 des \italique{Cavaliers}. La pièce débute par la traditionnelle scène des deux esclaves -- les généraux Nicias et Démosthène en l'occurrence, généraux qui avaient préparé la victoire de Sphactérie (victoire éclatante par la capture de la garnison spartiate).
\noindent Mais ces deux généraux s'étaient fait voler la vedette par le braillard Cléon débarqué avec un contingent sur place au dernier moment et qui – coup du sort – l'avait emporté\marqueur{3}. Ils ressassent leur amertume au début de la pièce en se renvoyant la balle pour trouver une idée de salut : que faire ? — Dis-le ! — Non, toi dis-le\marqueur{4}... Nicias s'interroge : comment tourner cela précieusement à la manière d'Euripide (κομψευριπικῶς). Démosthène rétorque alors (vers 19-20) :
\citComm{Δημοσθένης --- μὴ 'μοί γε, μὴ 'μοί, μὴ διασκανδικίσῃς·\\
ἀλλ᾽ εὑρέ τιν᾽ †ἀπόκινον† ἀπὸ τοῦ δεσπότου.\\
Démosthène --- Non pas, non pas, ne saupoudre pas de scandix !\\
mais va trouver quelque .......  qui le décroche du maître.}
\noindent Le « maître », c'est Δῆμος Πυκνίτης, le Populo-citoyen-roi-de-l'Assemblée, mené par le bout du nez par son « chien » : l'infâme Cléon.
\noindent (Tout au long des \italique{Cavaliers}, Cléon est renvoyé à son canin statut de \italique{Chien de Dêmos}, animal flatteur et cajoleur – sûrement par reprise ironique de quelque propos du véritable Cléon, qui se prenait pour le chien de garde de la cité. Voir aussi les \italique{Guêpes}, vv. 852-853, avec le procès aux deux chiens.)
\noindent Que le lecteur m'accorde pour l'instant que le verbe διασκανδικίζω signifie « assaisonner à l'aide de σκάνδιξ » et, dans un sens imagé, « agrémenter », « enjoliver [un discours de circonlocutions] » ; le σκάνδιξ jouant un rôle analogue à notre \italique{piment}, et les multiples sens du verbe « pimenter » (Nous l'allons montrer tout-à-l'heure.)
\noindent L'ἀπόκινος des manuscrits, défini par le scholiaste comme une danse licencieuse et qui fait entendre par son étymologie le verbe ἀπο-κινέω, \italique{se retirer}, n'a rien à faire ici\marqueur{5}.
\noindent Comme toujours devant les difficultés, nous savons que nous pouvons compter sur notre bon vieil et fidèle ami : Hésychius. Voici ce que grand érudit a répondu à nos questions (α 6416) :
\citComm{ἀπόκυνον: μάζα μεμαγμένη φαρμάκῳ, πρὸς ἀναίρεσιν κυνῶν. ἢ εἶδος βοτάνης\\
apokynon : galette préparée avec du poison [$=$ l'herbe dont il est question juste après], pour la destruction des chiens. \italique{Item} : sorte d'herbe.}
\end{multicols}
}%
\end{adjustbox}%
\vfill
\end{adjustwidth}
\newpage
% --- PAGE DE DROITE : Traduction française + Commentaire (4/8) ---
\begin{adjustwidth}{-2cm}{-2cm}
\noindent\textbf{\footnotesize Traduction française --- Folio 11 (\texttt{f. 167v}, vv. 423–465 [2/4]) \hfill ΕΠΕΙΣΟΔΙΟΝ Α'}\\[-2pt]
\noindent\rule{16.5cm}{0.5pt}\\[2pt]
{\fontsize{8.8pt}{10.2pt}\selectfont
\ligneEst{}{\francais{ni qui sache mieux parler.}}{}
\ligneEst{}{\francais{Tout ce qu'elle dit est juste.}}{}
\ligneEst{}{\francais{Elle a tout envisagé sous tous égards,}}{}
\ligneEst{}{\francais{tout soupesé dans son esprit : et avec discrétion,}}{}
\ligneEst{}{\francais{elle a trouvé un heureux}}{}
\ligneEst{}{\francais{choix d'adroites raisons.}}{}
\ligneEst{}{\francais{À tel point que si Xénoclès,}}{440}
\ligneEst{}{\francais{fils de Carkinos, parlait tout de suite après elle,}}{}
\ligneEst{}{\francais{il vous paraîtrait je pense,}}{}
\ligneEst{}{\francais{à tous}}{}
\ligneEst{}{\francais{parler sans rien dire.}}{}
\ligneEst{\textbf{Marchande de couronnes}}{\francais{Ce n'est que pour ajouter quelques mots que je m'avance à mon tour.}}{}
}
\vspace{2pt}
\noindent\rule{16.5cm}{0.6pt}\\[2pt]
\noindent\textbf{\small Commentaire philologique, métrique et scénique (vv. 423–465 --- partie 4/8)}\\[2pt]
\begin{adjustbox}{max totalheight=15.2cm, max width=16.5cm, keepaspectratio}%
\parbox[t]{16.5cm}{
\setlength{\columnsep}{9pt}%
\begin{multicols}{2}
\fontsize{7.6pt}{8.8pt}\selectfont
\setlength{\parindent}{0pt}%
\setlength{\parskip}{1.2pt}%
\noindent Hésychius nous explique que l'ἀπόκυνον, comme le σκάνδιξ, est une \italique{herbe sauvage}... à cette différence que celui-ci sert de condiment aux pauvres\marqueur{6} quand celui-là nous débarrasse des nuisibles : chiens selon Hésychius, Cléon pour Aristophane.
\noindent Je propose donc la lecture suivante pour les vers 19-20 des \italique{Cavaliers} :
\citComm{Δημοσθένης --- μὴ 'μοί γε, μὴ 'μοί, μὴ διασκανδικίσῃς·\\
ἀλλ᾽ εὑρέ τιν᾽ ἀπόκυνον ἀπὸ τοῦ δεσπότου.\\
(ἀπόκινον codd. omn. : \lecon{ἀπόκυνον} Bibi ap. Hesych.)\\
Démosthène --- Arrête-moi ça : ne condiments pas !\\
mais trouve-moi plutôt quelque mort-aux-chiens qui en débarrasse son maître.}
\noindent Ce passage des \italique{Cavaliers} est intéressant pour la bonne interprétation du verbe διασκανδικίζειν : Démosthène l'utilise d'abord comme synonyme de parler précieusement à la manière d'Euripide (κομψευριπικῶς) ; on dirait en français « épicer un discours ».  Mais au vers 20, il le sous-entend avec son sens premier : « Va me trouver de l'ἀπόκυνον [pour m'en épicer une galette] qui nous débarrassera le maître [de ce chien de Cléon]. »
\noindent Le verbe διασκανδικίζω n'est pas un hapax d'Aristophane ; il se trouve aussi dans un fragment de Téléclide, un autre poète comique contemporain d'Euripide. Ce fragment\marqueur{7} nous a été conservé par Athénée (\italique{Deipnosophistes}, VI, 56d.) :
\citComm{Ξυγγενέσθαι διὰ χρόνου †λιπαρείτω† με\\
δρυπεπέσι μάζαις καὶ διασκανδικίσαι.}
\noindent Tout ce que l'on peut inférer du contexte chez Athénée, c'est que l'adjectif δρυπεπής, \italique{bien mûr}, lorsqu'il est substantivé comme ici, désigne une \italique{olive}.
\noindent Le sens que nous avons donné à διασκανδικίζω, « assaisonner à l'aide de σκάνδιξ », rendrait bien mieux compte aussi de ce second passage que celui indiqué dans les dictionnaires : « se nourrir de σκάνδιξ »\marqueur{8}.
\noindent Je traduis donc : « Il m'importune (λιπαρεῖ ?) depuis longtemps pour ajouter des olives aux galettes ainsi que les agrémenter de scandix ».
\noindent Nous n'avons pas encore parlé d'une citation non identifiée à l'entrée Σκάνδιξ de la \italique{Souda} (σ 536) qui montre clairement que le scandix, en un sens imagé, désigne un \italique{euphémisme} :
\citComm{καὶ μή μοι τὰ τοῦ σκανδικοπώλου παρατραγῴδει, σεμνὸν τὸ τῆς Πειρήνης ὕδωρ ἀποκαλοῦντος,\\
« et ne m'use pas de l'emphase du vendeur de scandix qui appelle `sacrée' la fontaine de Pirène ».}
\noindent La cible est clairement Euripide : la citation peut faire allusion au vers 69 de sa \italique{Médée} ou aux vers 205-206 des \italique{Troyennes}. Il est toujours délicat de proposer des hypothèses sur des fragments aussi brefs... proposons tout-de-même notre hypothèse qui vaut ce qu'elle vaut – c'est-à-dire pas grand-chose : cette fontaine située sur l'Acrocorinthe délivrait effectivement une eau pure et abondante\marqueur{9} ; ce n'est donc pas sur la qualité de l'eau qu'Euripide enrobe de scandix la vérité... peut-être sur ce qu'on y pratiquait ? La fontaine de Pirène se trouvait juste en contrebas du temple d'Aphrodite dont certaines de ses préposées se livraient à une prostitution « sacrée »\marqueur{10}. Imagine-t-on un tragique contemporain, si peu tragique qu'il nous parlerait du « \italique{sacré} bois de... Boulogne » (!?) ce serait un peu fort de scandix !.. Cet exemple nous dessine très clairement le sens imagé du scandix : Euripide \italique{dore la pilule} en saupoudrant de scandix des choses basses ; il nous fait passer des vessies pour des lanternes.
\end{multicols}
}%
\end{adjustbox}%
\vfill
\end{adjustwidth}
\newpage

% --- FOLIO 11 (f. 167v) : SPREAD 2.3 (Texte grec établi + Traduction française + Commentaire 5/8 & 6/8) ---
\clearpage\ifodd\value{page}\hbox{}\newpage\fi
\enTeteG{ΘΕΣΜΟΦΟΡΙΑΖΟΥΣΑΙ --- Texte grec établi (f. 167v, vv. 423–465 [3/4])}
\enTeteD{LES THESMOPHORIES --- Traduction française (f. 167v, vv. 423–465 [3/4])}
\begin{adjustwidth}{-2cm}{-2cm}
\noindent\textbf{\footnotesize Texte grec établi --- Folio 11 (\texttt{f. 167v}, vv. 423–465 [3/4]) \hfill ΕΠΕΙΣΟΔΙΟΝ Α'}\\[-2pt]
\noindent\rule{16.5cm}{0.5pt}\\[2pt]
{\fontsize{8.8pt}{10.2pt}\selectfont
\ligneEst{}{τὰ μὲν γὰρ ἄλλ᾽ αὕτη κατηγόρηκεν εὖ·}{}
\ligneEst{}{ἃ δ᾽ ἐγὼ πέπονθα, ταῦτα λέξαι βούλομαι.}{445}
\ligneEst{}{ἐμοὶ γὰρ ἁνὴρ ἀπέθανεν μὲν ἔγκυος}{}
\ligneEst{}{παιδάρια πέντε καταλιπών, ἁγὼ μόλις}{}
\ligneEst{}{στεφανηπλοκοῦσ᾽ ἔβοσκον ἐν ταῖς μυρρίναις.}{}
\ligneEst{}{τέως μὲν οὖν ἀλλ᾽ ἡμικάκως ἐβοσκόμην·}{}
\ligneEst{}{νῦν δ᾽ οὗτος ἐν ταῖσιν τραγῳδίαις ποιῶν}{450}
\ligneEst{}{τοὺς ἄνδρας ἀναπέπεικεν οὐκ εἶναι θεούς·}{}
\ligneEst{}{ὥστ᾽ οὐκέτ᾽ ἐμπολῶμεν οὐδ᾽ εἰς ἥμισυ.}{}
\ligneEst{}{νῦν οὖν ἁπάσαισιν παραινῶ καὶ λέγω}{}
\ligneEst{}{τοῦτον κολάσαι τὸν ἄνδρα πολλῶν οὕνεκα·}{}
\ligneEst{}{ἄγρια γὰρ ἡμᾶς ὦ γυναῖκες δρᾷ κακά,}{455}
}
\vspace{2pt}
\noindent\rule{16.5cm}{0.6pt}\\[2pt]
\noindent\textbf{\small Commentaire philologique, métrique et scénique (vv. 423–465 --- partie 5/8)}\\[2pt]
\begin{adjustbox}{max totalheight=15.2cm, max width=16.5cm, keepaspectratio}%
\parbox[t]{16.5cm}{
\setlength{\columnsep}{9pt}%
\begin{multicols}{2}
\fontsize{7.0pt}{8.0pt}\selectfont
\setlength{\parindent}{0pt}%
\setlength{\parskip}{1.2pt}%
\noindent Voici pour le style d'Euripide. Passons à sa généalogie. L'on connaît le \italique{topos} comique : chaque fois qu'un poète voulait discréditer l'un de ses contemporains, il l'accusait de n'être pas de pure race (εὐγενής) : soit qu'il fut Athénien d'importation récente\marqueur{11}, soit qu'on puisse émettre des doutes sur sa lignée officielle. Euripide n'y échappe pas. Dans les \italique{Grenouille} (vv. 946-947), alors qu'il expose sa théorie dramatique, il est interrompu :
\citComm{Εὐριπίδης --- εἶτ’ οὐκ ἐλήρουν ὅ τι τύχοιμ’ οὐδ’ ἐμπεσὼν ἔφυρον,\\
ἀλλ’ οὑξιὼν πρώτιστα μέν μοι τὸ γένος εἶπ’ ἂν εὐθὺς\\
τοῦ δράματος.\\
Διόνυσος --- \decalage{τοῦ }κρεῖττον γὰρ ἦν σοι νὴ Δί’ ἢ τὸ σαυτοῦ.\\
Euripide --- Je ne débitais pas toutes les niaiseries qui me venaient, et je ne brouillais pas tout, sans crier gare ;
mais le premier qui entrait en scène exposait de prime-saut la filiation du drame.\\
Dionysos --- Cela valait mieux pour toi, par Zeus, que d'exposer la tienne.}
\noindent L'attaque n'est pas neuve : quelques vers plus tôt, Euripide avait provoqué Eschyle ; la réplique ne s'était pas fait attendre (vv. 836-843) :
\citComm{Εὐριπίδης --- ἐγᾦδα τοῦτον καὶ διέσκεμμαι πάλαι,\\
ἄνθρωπον ἀγριοποιὸν αὐθαδόστομον,\\
ἔχοντ᾽ ἀχάλινον ἀκρατὲς ἀπύλωτον στόμα,\\
ἀπεριλάλητον κομποφακελορρήμονα.\\
Αἰσχύλος --- `ἄληθες ὦ παῖ τῆς ἀρουραίας θεοῦ;'\\
σὺ δή με ταῦτ᾽ ὦ στωμυλιοσυλλεκτάδη\\
καὶ πτωχοποιὲ καὶ ῥακιοσυρραπτάδη;\\
ἀλλ᾽ οὔ τι χαίρων αὔτ᾽ ἐρεῖς.\\
Euripide --- Je le connais et l'ai de longtemps percé à jour,\\
cet homme aux héros sauvages, au verbe présomptueux,\\
à la langue sans frein, sans retenue, sans barrières,\\
à l'esprit rétif, ce fagoteur de mots ronflants.\\
Eschyle --- « Ah, par exemple ! fils de la déesse agreste », \\
c'est ainsi que tu parles de moi, raccoleur de fadaises,\\
faiseur de mendiants, ravaudeur de haillons.\\
Mais tu n'auras pas lieu de t'en applaudir.}
\noindent (Le scholiaste nous indique que la réaction cinglante d'Eschyle au vers 840 est une parodie d'un vers d'Euripide avec ἀρουραίας, \italique{agreste}, en place de θαλασσίας, \italique{marine}.)
\noindent Or on désigne comme fils de la (déesse) terre dans la langue proverbiale grecque (et latine) les personnes dont la généalogie est obscure.
\noindent Dans la pièce d'Euripide qui porte son nom, Ion interroge Xouthos sur l'identité de sa mère ; comme Xouthos ignore la réponse, Ion s'exclame alors (v. 542) :
\citComm{Ἴων --- γῆς ἄρ᾽ ἐκπέφυκα μητρός;\\
Ion --- J'ai donc eu la terre pour mère ?}
\noindent Érasme a recensé l'adage (n° 786) latin \latin{terrae filius} employé avec le même sens par Tertullien, Cicéron, \&c.
\noindent C'est bien le propre de l'αὐτό-χθων que de ne pas savoir à quand remonte l'installation de ses ancêtres dans le pays : sa filiation se perd dans la nuit des temps.
\noindent L'on comprend maintenant où veut en venir Aristophane avec ses légumes et autres herbes, étroitement associés à la terre : si nul Français ne s'offusquera qu'on dise de lui qu'il est né dans un chou ou une rose, en grec cela revient à émettre des doutes sur la naissance.
\noindent Ainsi s'explique notre passage des \italique{Thesmophories} (vv. 455-456).
\end{multicols}
}%
\end{adjustbox}%
\vfill
\end{adjustwidth}
\newpage
% --- PAGE DE DROITE : Traduction française + Commentaire (6/8) ---
\begin{adjustwidth}{-2cm}{-2cm}
\noindent\textbf{\footnotesize Traduction française --- Folio 11 (\texttt{f. 167v}, vv. 423–465 [3/4]) \hfill ΕΠΕΙΣΟΔΙΟΝ Α'}\\[-2pt]
\noindent\rule{16.5cm}{0.5pt}\\[2pt]
{\fontsize{8.8pt}{10.2pt}\selectfont
\ligneEst{}{\francais{Les autres chefs d'accusation ont été traités à souhait.}}{}
\ligneEst{}{\francais{Mais ce que j'ai souffert personnellement, je tiens à le dire.}}{445}
\ligneEst{}{\francais{Mon mari mourut alors que j'étais enceinte,}}{}
\ligneEst{}{\francais{me laissant cinq petits enfants, que je nourrissais à grand'peine}}{}
\ligneEst{}{\francais{en tressant des couronnes sur le marché aux myrtes.}}{}
\ligneEst{}{\francais{En attendant, je gagnais ma vie vaille que vaille.}}{}
\ligneEst{}{\francais{Mais aujourd'hui cet homme, avec ses tragédies,}}{450}
\ligneEst{}{\francais{a persuadé aux gens qu'il n'y a pas de dieux ;}}{}
\ligneEst{}{\francais{de sorte que notre commerce est réduit environ de moitié.}}{}
\ligneEst{}{\francais{C'est pourquoi je vous recommande à toutes et vous requiers}}{}
\ligneEst{}{\francais{de le punir, cet homme, pour quantité de raisons.}}{}
\ligneEst{}{\francais{Car il se porte contre nous, ô femmes, à de sauvages violences,}}{455}
}
\vspace{2pt}
\noindent\rule{16.5cm}{0.6pt}\\[2pt]
\noindent\textbf{\small Commentaire philologique, métrique et scénique (vv. 423–465 --- partie 6/8)}\\[2pt]
\begin{adjustbox}{max totalheight=15.2cm, max width=16.5cm, keepaspectratio}%
\parbox[t]{16.5cm}{
\setlength{\columnsep}{9pt}%
\begin{multicols}{2}
\fontsize{7.0pt}{8.0pt}\selectfont
\setlength{\parindent}{0pt}%
\setlength{\parskip}{1.2pt}%
\noindent Le procès d'Euripide a été ouvert et la parole est à l'accusation. Après un premier réquisitoire-fleuve, une seconde oratrice s'avance et prend la parole (vers 443-458).
\noindent Le point culminant de l'attaque se trouve dans le distique des vers 455-456 :
\citComm{ἄγρια γὰρ ἡμᾶς ὦ γυναῖκες δρᾷ κακά,\\
ἅτ᾽ ἐν ἀγρίοισι τοῖς λαχάνοις αὐτὸς τραφείς.}
\noindent Aristophane joue sur la répétition d'ἄγριος, qualifiant les mauvais comportements qu'Euripide impute à l'engeance féminine dans ses pièces\marqueur{12} et, au vers suivant, les λάχανα, les \italique{légumes} qui, lorsqu'ils sont ἄγρια, « sauvages », désignent des \italique{herbes sauvages}, de celles qui poussent sans qu'on les cultive (par opposition à ἥμερα).
\noindent Si Euripide nourrit une telle haine vis-à-vis des femmes et les présente comme dépravées, c'est que sa propre mère était telle... Sur le fond, ce genre de psychanalyse de comptoir n'est pas d'une grande finesse, mais le coup sous la ceinture porte auprès des femmes et suffit pour déchaîner l'enthousiasme du chœur quelques vers plus loin\marqueur{13}.
\noindent Ces éloges concernent davantage la forme que le fond de la diatribe\marqueur{14}.
\noindent Le tour `αὐτὸς τραφείς' est une locution proverbiale qu'on retrouve ailleurs : le parémiographe Michel Apostolius\marqueur{15} a relevé dans l'\italique{Anthologie} de Stobée (IV, 25, 6), au chapitre consacré au respect filial, cette citation du poète Dikaiogène\marqueur{16} :
\citComm{Αὐτὸς τραφεὶς δὲ τῶν φυτευσάντων ὕπο\\
καλῶς τὸν αὐτὸν ἔρανον\marqueur{17} αὐτοῖσιν νέμεις\\
Toi-même élevé [litt. : nourri] par tes parents [litt. : ceux qui t'ont semé],\\
tu leur partages comme de juste ton propre repas.}
\noindent La locution proverbiale complète était-elle `αὐτὸς τραφεὶς ὑπὸ τῶν φυτευσάντων' ? On comprendrait où veut en venir l'oratrice : si les parents d'Euripide, ses \italique{semeurs}, l'ont semé au milieu des légumes, il n'est pas douteux que cet homme de la terre y ait récolté tout le fiel qu'il déverse sur les femmes dans ses tragédies.
\noindent Dion Cassius (\italique{Histoire romaine}, XLVI, 4) met dans la bouche d'un orateur la même figure de rhétorique, semblable par la forme (ἐν $+$ datif $+$ αὐτὸς τραφείς) et par le ton. C'est une réponse à Cicéron et ses violentes \italique{Philippiques} contre Marc Antoine. Cicéron... qui est-ce, au juste ?
\citComm{ὁ πατὴρ αὐτῷ ὁ κναφεύς [ ... ] ἄνθρωπος ἀγαπητῶς ἔκ τε τούτων καὶ ἐκ τῶν πλυνῶν διατρεφόμενος [ ... ] \textbf{Ἐν οἷς αὐτὸς τραφεὶς} οὐκ ἀπεικότως τοὺς κρείσσονας αὐτοῦ καὶ πατεῖ καὶ πλύνει.\\
\francais{Son père était foulon, lui qui avait grand-peine à vivre de ce métier et de son lavoir. À la suite d'une telle éducation, non sans invraisemblance, ce personnage foule aux pieds ceux qui valent mieux que lui et leur lave la tête.}}
\noindent Ces clefs d'interprétation permettent d'expliquer le vers 478 des \italique{Acharniens}, alors que Dikaiopolis emprunte à Euripide tout un fatras de matériel puisé dans ses propres drames, voilà qu'il déclenche le courroux d'Euripide en lui demandant :
\citComm{Δικαιόπολις --- σκάνδικά μοι δὸς μητρόθεν δεδεγμένος.\\
Εὐριπίδης --- ἁνὴρ ὑβρίζει· κλῇε πηκτὰ δωμάτων.\\
Dikaiopolis --- Donne-moi de ce scandix que tu tiens de ta mère.\\
Euripide --- Il m'insulte. « Ferme sur lui les ais solides de la porte. »}
\noindent Aristophane fait coup double : l'on comprend bien pourquoi Dikaiopolis veut emprunter à Euripide son scandix : il a besoin de la faconde du tragédien pour la joute verbale qui va l'opposer au chœur des vieillards d'Acharnes\marqueur{18}.
\noindent On pourrait dire que Dikaiopolis compte les « enfariner » (διασκανδικίζειν) avec l'art oratoire d'Euripide... mais en précisant que ce scandix lui vient de sa mère, Dikaiopolis fait allusion au sobriquet de fils de verdurière qui colle au poète.
\end{multicols}
}%
\end{adjustbox}%
\vfill
\end{adjustwidth}
\newpage

% --- FOLIO 11 (f. 167v) : SPREAD 2.4 (Texte grec établi + Traduction française + Commentaire 7/8 & 8/8) ---
\clearpage\ifodd\value{page}\hbox{}\newpage\fi
\enTeteG{ΘΕΣΜΟΦΟΡΙΑΖΟΥΣΑΙ --- Texte grec établi (f. 167v, vv. 423–465 [4/4])}
\enTeteD{LES THESMOPHORIES --- Traduction française (f. 167v, vv. 423–465 [4/4])}
\begin{adjustwidth}{-2cm}{-2cm}
\noindent\textbf{\footnotesize Texte grec établi --- Folio 11 (\texttt{f. 167v}, vv. 423–465 [4/4]) \hfill ΕΠΕΙΣΟΔΙΟΝ Α'}\\[-2pt]
\noindent\rule{16.5cm}{0.5pt}\\[2pt]
{\fontsize{8.8pt}{10.2pt}\selectfont
\ligneEst{}{ἅτ᾽ ἐν ἀγρίοισι τοῖς λαχάνοις αὐτὸς τραφείς.}{}
\ligneEst{}{ἀλλ᾽ εἰς ἀγορὰν ἄπειμι· δεῖ γὰρ ἀνδράσιν}{}
\ligneEst{}{πλέξαι στεφάνους συνθηματιαίους εἴκοσιν.}{}
\ligneEst{\textbf{Χορός} {\tiny\textit{$2$ toch.}}}{ἕτερον αὖ τι λῆμα τοῦτο}{}
\ligneEst{{\tiny\textit{$2$ troch.}}}{κομψότερον ἔτ᾽ ἢ τὸ πρότερον}{460}
\ligneEst{{\tiny\textit{troch.}}}{ἀναπέφηνεν.}{}
\ligneEst{{\tiny\textit{choria. $+$ ia.}}}{οἷα κατεστωμύλατο}{}
\ligneEst{{\tiny\textit{$2$ troch.}}}{οὐκ ἄκαιρα, φρένας ἔχουσα}{}
\ligneEst{{\tiny\textit{$2$ troch.}}}{καὶ πολύπλοκον νόημ᾽, οὐδ᾽}{}
\ligneEst{{\tiny\textit{$2$ troch.}}}{ἀσύνετ᾽ ἀλλὰ πιθανὰ πάντα.}{}
\ligneEst{{\tiny\textit{$2$ troch.}}}{δεῖ δὲ ταύτης τῆς ὕβρεως ἡ-}{465}
\ligneEst{{\tiny\textit{troch.}}}{μῖν τὸν ἄνδρα}{}
\ligneEst{{\tiny\textit{$2$ troch\cat}}}{περιφανῶς δοῦναι δίκην.}{}
}
\vspace{2pt}
\noindent\rule{16.5cm}{0.6pt}\\[2pt]
\noindent\textbf{\small Commentaire philologique, métrique et scénique (vv. 423–465 --- partie 7/8)}\\[2pt]
\begin{adjustbox}{max totalheight=15.2cm, max width=16.5cm, keepaspectratio}%
\parbox[t]{16.5cm}{
\setlength{\columnsep}{9pt}%
\begin{multicols}{2}
\fontsize{7.0pt}{8.0pt}\selectfont
\setlength{\parindent}{0pt}%
\setlength{\parskip}{1.2pt}%
\noindent Euripide le comprend bien qui chasse l'insolent : ἁνὴρ ὑβρίζει· κλῇε πηκτὰ δωμάτων.\\
\noindent Cette trop longue note touchant enfin à sa fin, nous nous permettrons quelque aventureuse conjecture sur l'origine de cette utilisation imagée du σκάνδιξ : il nous semble que si le grec a choisi cette plante précisément et non une autre, c'est non pas pour ses vertus (réelles ou supposées), mais bien plutôt pour des raisons de sonorité. Cette intuition m'est venue de la consultation de l'\latin{Etymologicum Genuium}\marqueur{19} à l'entrée σκάνδαλον, la \italique{chausse-trappe}, pris ici en un sens imagé : σκάνδαλον: τὸ σκάζειν περὶ τὴν ἀλήθειαν.
\noindent Image saisissante et évocatrice s'il en est ! « boîter autour de la vérité », « tourner autour du pot-aux-roses ».
\noindent Elle s'appuie sur un rapprochement pseudo-étymologique\marqueur{20} entre le mot σκάνδ-αλον et le verbe σκάζειν, boîter... le σκάνδ-ιξ ne pouvait-il pas suggérer la même image : celle de \italique{tordre}, \italique{rendre boiteuse} ou de \italique{camoufler} la vérité en la saupoudrant d'une herbe de vil prix ?
\noindent Voilà qui donne quelques nouveaux éléments d'appréciation sur le problème d'Euripide, son σκάνδιξ et sa maman la λαχανοπωλήτρια.
\noindent\textbf{458} \textbf{στεφάνους συνθηματιαίους.} Athénée (XV, 680c), dans le long chapitre consacré aux couronnes de ses \italique{Deipnosophistes}, parle de ces couronnes \italique{synthématiées} :
\citComm{συνθηματιαῖοι στέφανοι, ἠργολαβημένοι καὶ ἐκδόσιμοι. Ἀριστοφάνης ἐν Θεσμοφοριαζούσαις.}
\noindent Glose identique chez Hésychius (σ 2624) :
\citComm{συνθηματιαῖοι: οἱ ἠργολαβημένοι καὶ ἐκδεδομένοι, ἐκδόσιμοι.}
\noindent\textbf{[Notes ad 455-6]} \textbf{(1)} \italique{Thesm.} vv. 387, 456 ; \italique{Acharn.} v. 478 ; \italique{Grenouilles} v. 840.
\noindent \textbf{(2)} Dans un article passionnant sur quelques plantes apparaissant chez Aristophane, mêlant avec un rare bonheur de savantes études de textes et des connaissances botaniques très pointues, la regrettée Suzanne Amigues identifiait en fait l'ἴφυον avec la nielle des blés (\latin{Agrostemma githago L.}). (« Quelques légumes de disette chez Aristophane et Plutarque » paru en 1988 dans le \italique{Journal des Savants}, pp. 157-171.)
\noindent \textbf{(3)} Pour plus de détails, le lecteur pourra utilement se reporter au quart livre de Thucydide, chapitres 28 et suivants ainsi qu'à la longue tirade du Démosthène des \italique{Cavaliers}, à partir du vers 40.
\noindent \textbf{(4)} Aristophane fait-il allusion à un complot qui aurait été tramé contre Cléon ? contre ce va-t-en-guerre de Cléon qui empêchait toute négociation de paix au moment-même où Athènes était en position de force pour se sortir favorablement de cette guerre fratricide qui n'avait que trop duré.
\noindent \textbf{(5)} Il faut donc rejeter la jolie trouvaille du grand Alphonse Willems : le « branle de sortie »...
\noindent \textbf{(6)} Le scandix semble être une épice peu appréciée : Andocide parle ainsi des privations des premières années de la guerre du Péloponnèse : μηδὲ ἄγρια λάχανα καὶ σκάνδικας ἔτι φάγοιμεν, « puissions-nous ne plus manger des herbes sauvages ni de scandix ».
\noindent \textbf{(7)} Fragment 40(38) dans l'édition de Kassel-Austin, \italique{Poetae Comici Graeci}, vol. VII.
\noindent \textbf{(8)} « manger du cerfeuil » Bailly ; « \espagnol{dar de comer perifollo} » DGE ; « \anglais{feed on chervil} » LSJ ; « Kerbel essen » Pape.
\noindent \textbf{(9)} Cf. ce qu'en disent Pausanias, Strabon et Suidas à l'entrée Πειρήνη ainsi que cette citation d'Antipater prise dans l'\italique{Anthologie grecque} : Λαί̈δ' ἔχω πολιῆτιν ἐϋζώνοιο Κορίνθου, Πειρήνης λευκῶν φαιδροτέρην λιβάδων, « Je tiens Laïs, citoyenne de Corinthe, bien ceinturée, plus brillante que les gouttes blanches de Pirène ».
\noindent \textbf{(10)} La réputation des Corinthiennes était telle que les Grecs en ont inventé le verbe κορινθιάζειν.
\noindent \textbf{(11)} Voir par exemple \italique{Grenouilles} vv. 717-733 qui est sans ambiguïté. L'on s'étonne de trouver dans les modernes traductions des \italique{Thesmophories} (v. 330) εὐγενεῖς γυναῖκες rendu par `nobles femmes', comme si la possibilité de célébrer les Thesmophores était conditionnée à une forme quelconque de noblesse. Il serait bon avant d'aborder Aristophane de se défaire de nos préjugés modernes, faute de quoi l'on s'expose à de tels contre-sens. Lire là-dessus l'indispensable article d'Alphonse Willems « Aristophane et la démocratie athénienne » reproduit en annexe et qu'on peut trouver dans son \italique{Aristophane} posthume (1919), t. III, pp. 445-473.
\end{multicols}
}%
\end{adjustbox}%
\vfill
\end{adjustwidth}
\newpage
% --- PAGE DE DROITE : Traduction française + Commentaire (8/8) ---
\begin{adjustwidth}{-2cm}{-2cm}
\noindent\textbf{\footnotesize Traduction française --- Folio 11 (\texttt{f. 167v}, vv. 423–465 [4/4]) \hfill ΕΠΕΙΣΟΔΙΟΝ Α'}\\[-2pt]
\noindent\rule{16.5cm}{0.5pt}\\[2pt]
{\fontsize{8.8pt}{10.2pt}\selectfont
\ligneEst{}{\francais{élevé qu'il fut parmi les herbes sauvages.}}{}
\ligneEst{}{\francais{Mais je m'en vas à l'agora.}}{}
\ligneEst{}{\francais{Car j'ai vingt couronnes à tresser à prix convenu.}}{}
\ligneEst{\textbf{Chœur}}{\francais{Cette autre a montré un tour d'esprit}}{}
\ligneEst{}{\francais{plus piquant encore}}{460}
\ligneEst{}{\francais{que la première.}}{}
\ligneEst{}{\francais{Comme elle a débité,}}{}
\ligneEst{}{\francais{non du hors-d'oœvre — elle a du sens}}{}
\ligneEst{}{\francais{et un esprit délié — ni de l'abstrus,}}{}
\ligneEst{}{\francais{mais toutes choses emportant conviction !}}{}
\ligneEst{}{\francais{Oui, il faut que de cet outrage}}{465}
\ligneEst{}{\francais{l'auteur}}{}
\ligneEst{}{\francais{subisse un châtiment exemplaire.}}{}
}
\vspace{2pt}
\noindent\rule{16.5cm}{0.6pt}\\[2pt]
\noindent\textbf{\small Commentaire philologique, métrique et scénique (vv. 423–465 --- partie 8/8)}\\[2pt]
\begin{adjustbox}{max totalheight=15.2cm, max width=16.5cm, keepaspectratio}%
\parbox[t]{16.5cm}{
\setlength{\columnsep}{9pt}%
\begin{multicols}{2}
\fontsize{7.0pt}{8.0pt}\selectfont
\setlength{\parindent}{0pt}%
\setlength{\parskip}{1.2pt}%
\noindent \textbf{(12)} Aristophane, au vers 349 des \italique{Nuées}, qualifie ainsi d' ἄγριος le chevelu Hiéronyme, un poète dithyrambique, fils de Xénophane qui passait pour pédéraste ; et le scholiaste de noter qu'on appelait ἀγρίους (et débauchés) les pédérastes : ἀγρίους (καὶ κολλοποδιώκτας) ἐκάλουν τοὺς παιδεραστάς. De même, à l'entrée Ἀγρίους de la \italique{Souda}, on peut lire : Αἰσχίνης ἐν τῷ κατὰ Τιμάρχου τοὺς σφόδρα ἐπτοημένους περὶ τὰ παιδικὰ οὕτως ὠνόμασε. καὶ Μένανδρος δὲ ἄγριον εἶπε κυβευτὴν τὸν λίαν περὶ τὸ κυβεύειν ἐσπουδακότα, « Eschine dans son \italique{Contre Timarque} nomme ainsi les hommes en proie à une excessive attirance pour les enfants. Et Ménandre dit qu'un joueur est ἄγριον, celui qui est trop porté au jeu ». L'on voit que le terme ἄγριος, au sens moral, désigne un \italique{intempérant}.
\noindent \textbf{(13)} κομψότερον ἔτ᾽ ἢ τὸ πρότερον, « plus intelligent encore que le premier [discours] » [...] οἷα κατεστωμύλατο, « comme elle a bien parlé » [...] φρένας ἔχουσα, « elle a de l'esprit » \&c.
\noindent \textbf{(14)} L'oratrice précédente, au vers 387, avait déjà traité Euripide de fils de la λαχανοπωλήτρια. Plutarque citait ces deux vers comme exemple d'un jeu de mots qu'il goûtait fort peu (voir ses \italique{Œuvres morales}, 853C, \italique{Abrégé de la comparaison d'Aristophane et de Ménandre}). Je ne suis pas sûr que cet inénarrable raseur, le Cicéron des lettres grecques, ait bien saisi intégralement le détournement d'Aristophane.
\noindent \textbf{(15)} \italique{Corpus paroemiographorum graecorum}, II, ed. Leutsch, Apostolii centuria V 30a, p. 341.
\noindent \textbf{(16)} L'on ne possède pas grand-chose sur ce poète grec, auteur de tragédies et dithyrambes, qui vécut probablement au \nieme{IV} siècle av. J.-C. Tout ce que l'on sait de lui est regroupé dans les \italique{Tragicorum Graecorum Fragmenta}, vol. I, ed. Snell \& Kannicht. Il y est classé en tant que « poète mineur » n° 52.
\noindent \textbf{(17)} L'on trouve chez Démosthène dans son \italique{Contre Neæra}, 8, l'expression ironique τὸν αὐτὸν ἔρανον ἀποδοῦναι, \italique{payer quelqu'un de la même monnaie}.
\noindent \textbf{(18)} C'est la raison de sa visite chez Euripide (vers 416-417) : δεῖ γάρ με λέξαι τῷ χορῷ ῥῆσιν μακράν· / αὕτη δὲ θάνατον, ἢν κακῶς λέξω, φέρει.
\noindent « Car j'ai à débiter au chœur une longue tirade. / Il m'en coûte la vie, si je m'en tire mal. »
\noindent \textbf{(19)} Lexique du \nieme{IX} siècle souvent associé à d'autres, postérieurs mais plus complets : l'\latin{Etymologicum magnum} et l'\latin{Etymologicum Symeonis}.
\noindent \textbf{(20)} Je me vois forcé de préciser \italique{pseudo}-étymologique, car mes velléités de rapprocher étymologiquement σκάνδαλον de σκάζειν furent foudroyées par Alexandre Perret, dont j'eus la chance d'assister à des cours de linguistique, cours qui ressemblaient davantage à d'enivrants voyages à travers les contrées et les parlers d'Europe (et ce jusqu'aux confins de l'Asie). Voici ce que m'a très aimablement répondu M. Perret : « Après recherches fructueuses, je pense avoir trouvé une réponse dans le \italique{Lexikon der Indogermanischen Verben} (LIV) mais il ne semble pas que les deux racines soient liées. Σκανδάλον semble en effet rattaché à un thème *skend- qui signifierait « bondir » (on a skándati en sanskrit pour « sauter »), mais le suffixe -αλον me dérange dans la mesure où ça semble pré-grec. Aucune étymologie n'est assez bien solide pour ce mot qui peut être un emprunt. Quelle motivation pour passer de « bondir » à « faire tomber » ? À la limite, un autre thème *sk'end- (avec un k- palatal donc invisible en grec mais pas en sanskrit) avec le sens parfois de « faire tomber » (skr. Chándayati: « tomber »), est en soit beaucoup plus plausible !! Mais pour σκάζω, on rapproche ça d'une racine *skeng- avec effectivement le sens de « boiteux, tors » (All. \italique{Hinken}). Ce qui m'étonne est qu'Hesychius recense à la fois σκάνδαλον: \italique{piège à souris} et σκάνδαλος: \italique{obstacle}, mais est-ce que les pièges à souris étaient des trappes où tomber, ou bien des claquoirs comme à notre époque ? Ça mérite investigation. En soit passer de *skend- à σκάζω avec un suffixe -yo de présent n'est pas choquant. »
\noindent\textbf{463} \textbf{πολύπλοκον.} Les diverses corrections qui ont été proposées pour faire rentrer de force ce vers dans le mètre trochaïque ne me paraissent pas indispensables : on peut simplement considérer qu'il n'y a pas de \italique{correptio attica} et scander :
\citComm{καὶ πολύπ|λοκον νόημ᾽, οὐδ᾽\\
$ - \cup - \cup - \cup - - $}
\end{multicols}
}%
\end{adjustbox}%
\vfill
\end{adjustwidth}
\newpage
\clearpage\ifodd\value{page}\hbox{}\newpage\fi
\enTeteG{ΕΠΕΙΣΟΔΙΟΝ Α' --- Codex Ravennas 429, f. 168r (Classense 170r) --- vv. 466–508}
\begin{adjustwidth}{-1.5cm}{-3.0cm}
\begin{adjustbox}{max totalheight=21.0cm, max width=16.5cm, keepaspectratio}
\begin{minipage}{16.5cm}
{\fontsize{7.2pt}{8.5pt}\selectfont
\noindent\textbf{Folio 12 --- Codex Ravennas 429, \texttt{f. 168r} (\texttt{Classense 170r}) : ΕΠΕΙΣΟΔΙΟΝ Α', vv. 466–508}\hfill \textit{Témoins :} $\text{R}^1, \Rpc, \Rdeux$ (\Mon)\par\vspace{2pt}
\hrule height 0.4pt\vspace{2.5pt}
}
{\fontsize{8.1pt}{9.5pt}\selectfont
\ligneRav{466}{μ\auDessus{ν}{Η}}{τὸ μὲν ὦ γυναῖκες ὀξυθυμεῖσθαι σφόδρα}{}
\ligneRav{467}{}{Εὐριπίδῃ, τοιαῦτ᾽ \difference{ἀκουούσας} κακά,}{\textbf{467} ἀκουούσαις R : \lecon{ἀκουούσας} Zanetti}
\ligneRav{}{}{οὐ θαυμάσιόν ἐστ᾽, οὐδ᾽ ἐπιζεῖν τὴν χολήν·}{}
\ligneRav{469}{}{καὐτὴ γὰρ ἔγωγ᾽, οὕτως \difference{ὀναίμην} τῶν τέκνων,}{\textbf{469} ὀνοίμην R : \lecon{ὀναίμην} Suid. ο 330}
\ligneRav{470}{}{μισῶ τὸν ἄνδρ᾽ ἐκεῖνον, εἰ μὴ μαίνομαι·}{}
\ligneRav{471}{}{ὅμως δ᾽ ἐν \difference{ἀλλήλαισι} χρὴ δοῦναι λόγον·}{\textbf{471} ἀλλήλοισι R : \lecon{ἀλλήλαισι} Brunck}
\ligneRav{472}{}{αὐταὶ γάρ ἐσμεν, κοὐδεμί᾽ ἔκφορος λόγου·}{\textbf{472} \lecon{ἔκφορος} R : ἐκφορὰ Valckenaer ad Eur. \italique{Hipp.} 294}
\ligneRav{473}{}{τί ταῦτ᾽ ἔχουσαι \difference{᾽κεῖνον} αἰτιώμεθα}{\textbf{473} κεῖνον R, Suid τ 683 : \lecon{᾽κεῖνον} Meineke}
\ligneRav{474}{}{βαρέως τε φέρομεν, \difference{εἰ δύ᾽} ἡμῶν ἢ τρία}{\textbf{474} ἢ δύ᾽ R : \lecon{εἰ δύ᾽} Grynaeus \espApp ἢ δύ᾽ R : \lecon{εἰ δύ᾽} Grynaeus (\italique{si} Divus)}
\ligneRav{475}{}{κακὰ ξυνειδὼς εἶπε δρώσας μυρία;}{}
\ligneRav{476}{}{ἐγὼ γὰρ αὐτὴ πρῶτον, ἵνα \difference{μἄλλην} λέγω,}{\textbf{476} μὴ ἄλλην R : \lecon{μἄλλην} Dindorf}
\ligneRav{477}{}{ξύνοιδ᾽ ἐμαυτῇ \difference{πολλά γ᾽· ἀλλ᾽} ἐκεῖνο δ᾽ οὖν}{\textbf{477} πολλ᾽ R : δεινὰ πολλ᾽ Bentley : \lecon{πολλά γ᾽· ἀλλ᾽} Bibi}
\ligneRav{478}{}{δεινότατον, ὅτε νύμφη μὲν ἦ τρεῖς ἡμέρας,}{\textbf{478} ἦν R : \lecon{ἦ} Dindorf \espApp ἦν R : \lecon{ἦ} Dindorf (1825)}
\ligneRav{479}{}{ὁ δ᾽ ἀνὴρ παρ᾽ ἐμοὶ \difference{καθηῦδεν}· ἦν δέ μοι φίλος,}{\textbf{479} καθεύδειν R : \lecon{καθηῦδεν} Scaliger}
\ligneRav{480}{}{ὅσπερ με διεκόρησεν \difference{οὖσαν ἑπτ}έτιν·}{\textbf{480} διεκόρευσεν R : \lecon{-ρησεν} Pollux cf. \italique{Pax} 59 \espApp ἑπτ. οὖσ. R : \lecon{οὖσαν ἑπτ.} Pollux}
\ligneRav{481}{}{οὖτος πόθῳ \difference{μου ᾽κνυεν} ἐλθὼν τὴν θύραν·}{\textbf{481} μ᾽ οὔκνυεν R : \lecon{μου ᾽κνυεν} Suid. ε 520}
\ligneRav{482}{}{\difference{κᾆτ᾽} εὐθὺς ἔγνων· εἶτα καταβαίνω λάθρᾳ·}{\textbf{482} κᾶτ᾽ R : \lecon{κᾆτ᾽} Küster}
\ligneRav{}{}{ὁ δ᾽ ἀνὴρ ἐρωτᾷ `ποῖ σὺ καταβαίνεις;' `ὅποι;}{}
\ligneRav{}{}{στρόφος μ᾽ ἔχει τὴν γαστέρ᾽ ὦνερ κὠδύνη·}{}
\ligneRav{485}{}{ἐς τὸν κοπρῶν᾽ οὖν ἔρχομαι.' `βάδιζέ νυν.'}{}
\ligneRav{486}{}{κᾆθ᾽ ὁ μὲν ἔτριβεν κεδρίδας \difference{ἄννηθον} σφάκον·}{\textbf{486} ἄνηθον R, Suid. : \lecon{ἄννηθον} Dindorf (1830)}
\ligneRav{}{}{ἐγὼ δὲ καταχέασα τοῦ στροφέως ὕδωρ}{}
\ligneRav{488}{}{ἐξῆλθον ὡς τὸν μοιχόν· εἶτ᾽ \difference{ἠρειδόμην}}{\textbf{488} ἐρειδόμην R : \lecon{ἠρειδόμην} Küster}
\ligneRav{489}{}{κ\inRas{υν}{ύβ}δ᾽ ἐχομένην τῆς δάφνης·}{\textbf{489} κυνδ᾽ \Rac : \lecon{κύβδ᾽} \Rpc, \Mon}
\ligneRav{490}{}{ταῦτ᾽ οὐδεπώποτ᾽ εἶφ᾽, ὁρᾶτ᾽, Εὐριπίδης·}{}
\ligneRav{}{}{οὐδ᾽ ὡς ὑπὸ τῶν δούλων τε κὠρεωκόμων}{}
\ligneRav{}{}{σποδούμεθ᾽, ἢν μὴ 'χωμεν ἕτερον, οὐ λέγει·}{}
\ligneRav{493}{}{ὑπὸ τῶν δούλων τε κὠρεωκόμων \difference{ληκώμεθα} (\inRas{ληκ}{κιν}ώμεθα)}{\textbf{493} \lecon{ληκώμεθα} \Rac (?), Suid. : κινώ- \Rpc}
\ligneRav{494}{}{τὴν νύχθ᾽, ἕωθεν \difference{σκόροδα δια}μασώμεθα,}{\textbf{494} σκορόδια R : \lecon{σκόροδα δια} Bentley \espApp post νύχθ' add. ὅλην Suid. (reminiscentia Pherecratis ? cf. Photius λ 254)}
\ligneRav{495}{}{ἵν᾽ \difference{ὀσφρόμενος} \difference{ἁνὴρ} \difference{ἀπὸ} τείχους εἰσιὼν}{\textbf{495} ἀπὸ τοῦ R : \lecon{ἀπὸ} Biset \espApp ὀσφραινόμενος R : \lecon{ὀσφρόμενος} Bentley, Küster \espApp ἀνὴρ R : \lecon{ἁνὴρ} Dawes}
\ligneRav{}{}{μηδὲν κακὸν δρᾶν ὑποτοπῆται. ταῦθ᾽, ὁρᾷς,}{}
\ligneRav{}{}{οὐπώποτ᾽ εἶπεν. εἰ δὲ Φαίδραν λοιδορεῖ,}{}
\ligneRav{}{}{ἡμῖν τί τοῦτ᾽ ἔστ᾽; οὐδ᾽ ἐκεῖν᾽ εἴρηκέ πω,}{}
\ligneRav{500}{}{ὡς ἡ γυνὴ δεικνῦσα τἀνδρὶ τοὔγκυκλον}{\textbf{500} ὑπ᾽ αὐγὰς οἷόν ἐστιν ἐγκεκαλυμμένον R : ὑπαυγάσ᾽ κτλ. Hermann : \lecon{del} van Herwerden}
\ligneRav{(500)}{}{$[\![\text{\sout{ὑπ᾽ αὐγὰς οἷόν ἐστιν ἐγκεκαλυμμένον}}]\!]$}{\textbf{v. 500} ὑπ᾽ αὐγὰς οἷόν ἐστιν ἐγκεκαλυμμένον R : \lecon{del.} van Herwerden, Bibi (ex gl. ᾦόν)}
\ligneRav{501}{}{τὸν \difference{μοιχὸν} ἐξέπεμψεν, οὐκ εἴρηκέ πω·}{\textbf{501} μυχὸν R : \lecon{μοιχὸν} Gryaneus}
\ligneRav{502}{}{ἑτέραν δ᾽ ἐγᾦδ᾽ \difference{ἣ 'φασκεν} ὠδίνειν γυνὴ}{\textbf{502} ήφασκεν R : \lecon{ἣ 'φασκεν} Zanetti}
\ligneRav{}{}{δέχ᾽ ἡμέρας, ἔως ἐπρίατο παιδίον·}{}
\ligneRav{504}{}{ὁ δ᾽ ἀνὴρ περιήρχετ᾽ \difference{ὠκυτόκι᾽} ὠνούμενος·}{\textbf{504} \lecon{περιήρχετ᾽} R : περιέτρεχ᾽ Bergk \espApp ὠκυτόκειa R : \lecon{ὠκυτόκι᾽} Pollux II.7}
\ligneRav{505}{}{τὸ δ᾽ εἰσέφερε γραῦς ἐν χύτρᾳ τὸ παιδίον,}{\textbf{505} \lecon{τὸ δ᾽} R : τότ᾽ Biset : τῇδ᾽ Daubuz}
\ligneRav{506}{}{τὸ δ᾽ εἰσέφερε γραῦς ἐν χύτρᾳ τὸ παιδίον, ἵνα μὴ \expo{βοῴη} κηρίῳ βεβυσμένον·}{\textbf{506} βοῴη suppl. \Rpc}
\ligneRav{}{}{εἶθ᾽ ὡς ἔνευσεν ἡ φέρουσ᾽, εὐθὺς βοᾷ,}{}
\ligneRav{}{}{`ἄπελθ᾽ ἄπελθ᾽, ἤδη γὰρ ὦνέρ μοι δοκῶ}{}
}
\vspace{3pt}\noindent\rule{4.5cm}{0.4pt}\par\vspace{1.5pt}
{\fontsize{7.0pt}{8.3pt}\selectfont
\noindent\textbf{Loci paralleli (vv. 466–508) ---} \textbf{466} (ὀξυθυμεῖσθαι) Photius ο 380 ὀξυθυμεῖσθαι \espTest \textbf{468} (θαυμάσιον $=$ θαυμαστόν) Moeris θ 16 $\simeq$ Hesychius θ 146 θαυμάσιος \espTest \textbf{469} Suidas ο 330 (sed ἐγώ pro ἔγωγ᾽) ; cf. (ὀναίμην) Hesychius ο 852 ὀναίμην ; \Syn ο 162 ὀναίμην $=$ Photius ο 332 \espTest \textbf{470} Suidas μ 328 (τἄνδρ᾽ pro τὸν ἄνδρ᾽) \espTest \textbf{473-4} (τί ταῦτ᾽ -- φέρομεν) Suidas τ 683 \espTest \textbf{480} Pollux III 42 (οὗτος pro ὅσπερ et διεκόρησεν cod. A pro διεκόρευσεν) \espTest (διεκόρησεν) Hesychius δ 1558 διεκόρησε \espTest \textbf{480} Pollux III 42 (οὗτος pro ὅσπερ et διεκόρησεν pro διεκόρευσεν) \espTest \textbf{481} Suidas ε 520 ; Zonaras 669 ἔκνυεν; \Eustathe 1746.6 (μοὔκνυεν ἕλκων pro μου ᾽κνυεν ἐλθὼν) \espTest \textbf{484} \Eustathe 1746.45 \espTest (στρόφος) \Eustathe 503.20 \espTest \textbf{486} Suidas σ 1712 \espTest (κεδρίδας) Hesychius κ 1980 κεδρίδες ; Photius κ 1980 κ 517 \espTest \textbf{489} (Ἀγυιᾶ) Stephanus Byzantinus 22.15 ; cf. Moeris α 13 \espTest \textbf{493-4} Suidas λ 434 Ληκώμεθα: κινώμεθα. Ἀριστοφάνης: μάλισθ' ὅταν ὑπό του ληκώμεθα τὴν νύχθ' ὅλην. Post νύχθ' add ὅλην Suid. (reminiscentia Pherecratis ? cf. Photius λ 254) \espTest \textbf{504} (ὠκυτόκι᾽) Pollux II 7 ; \Phryn 129.3 \espTest \textbf{506} (κηρίῳ βεβυσμένον) Hesychius κ 2549 κηρίῳ βύσασα ; Photius κ 675 κηρίῳ βῦσαι\par
}
\vspace{2.5pt}\noindent\rule{4.5cm}{0.4pt}\par\vspace{1.5pt}
{\fontsize{6.9pt}{8.2pt}\selectfont
\noindent\textbf{Suite (en commentaire) de l'apparat critique \& notes paléographiques (\texttt{f. 168r}) ---} \textbf{v. 477 :} Le vers de R (οὐδ' εἴρηκεν οὐδὲν πρὸς ἃ δρῶμεν· πολλ') est amputé d'une syllabe finale ; \textbf{\textbf{Bibi} restitue πολλά γ'· ἀλλ'} (chute de -ά γ'· ἀλλ' par saut du même au même entre les deux ΛΛ). --- \textbf{v. 489 :} \Rac{} avait écrit κυνδ' (confusion minuscule $\beta/\nu$), corrigé en \lecon{κύβδ'} par \Rpc{} (suivi par \Mon{}). --- \textbf{v. 493 :} Comme le précise Louis Hauseux (\texttt{vv 1-641.tex}, l. 3715), le copiste avait d'abord écrit la forme crue \lecon{ληκώμεθα} (attestée par la Souda $\sigma\ 1126$ et adoptée par L. Hauseux), puis a \textbf{gratté les trois premières lettres ληκ-} pour les remplacer par κιν- (\lecon{\inRas{ληκ}{κιν}ώμεθα} \Rpc{}) sous l'influence de la glose κινώμεθα. --- \textbf{v. 500 (genèse paléographique d'un vers interpolé) :} Le v. 500 (ὑπ' αὐγὰς οἷόν ἐστιν ἐγκεκαλυμμένον), athétisé par van Herwerden et L. Hauseux (\texttt{vv 1-641.tex}, ll. 3781–3806), est né d'une glose ou citation marginale [τὸν μοιχόν·] ὑπαυγάσας. ᾦόν ἐστιν ἐγκεκαλυμμένον, dans laquelle ΩΙΟΝ (avec iota muet adscrit = $\text{ᾦόν}$, « un œuf enveloppé ») a été lu par iotacisme et bêtacisme $\omega/\text{ο}$ \textbf{ΟΙΟΝ (οἷόν)}, entraînant la réécriture métrique ὑπ' αὐγὰς οἷόν ἐστιν.... --- \textbf{v. 506 :} L'optatif \lecon{βοῴη}, omis par \Rac{}, est suppléé par \Rpc{}. --- \textbf{v. 477 :} πολλ' R (δεινὰ πολλ' Bentley) : \lecon{πολλά γ'· ἀλλ'} \lecon{Bibi}. --- \textbf{v. 481 :} μ' οὔκνυεν R (médivision et diphtongaison) : \lecon{μου 'κνυεν} Souda $\varepsilon\ 520$ (cf. Hésychius $\kappa\ 3133$). --- \textbf{v. 501 :} μυχὸν R (iotacisme) : \lecon{μοιχὸν} Grynaeus.\par
}
\end{minipage}
\end{adjustbox}
\end{adjustwidth}
\newpage
\enTeteD{Codex Ravennas 429 (Biblioteca Classense) --- Fac-similé du f. 168r (vv. 466–508)}
\begin{adjustwidth}{-1.5cm}{-3.0cm}
\begin{center}
\includegraphics[width=16.5cm,height=20.5cm,keepaspectratio]{/workspaces/Thesmophories/dossier_ravennas_thesmophories/folios_individuels_ff162v-177r/12_f168r_classense170r_vv0466-0508.jpg}\par
\vspace{2pt}
{\scriptsize\textbf{Codex Ravennas 429, f. 168r (Classense 170r)} --- ΕΠΕΙΣΟΔΙΟΝ Α', \textbf{vv. 466–508} ($\text{R}^1, \Rpc, \Rdeux$ (\Mon)).}
\end{center}
\end{adjustwidth}
\newpage
\clearpage\ifodd\value{page}\hbox{}\newpage\fi% --- FOLIO 12 (f. 168r) : SPREAD 2.1 (Texte grec établi + Traduction française + Commentaire 1/6 & 2/6) ---
\clearpage\ifodd\value{page}\hbox{}\newpage\fi
\enTeteG{ΘΕΣΜΟΦΟΡΙΑΖΟΥΣΑΙ --- Texte grec établi (f. 168r, vv. 466–508 [1/3])}
\enTeteD{LES THESMOPHORIES --- Traduction française (f. 168r, vv. 466–508 [1/3])}
\begin{adjustwidth}{-2cm}{-2cm}
\noindent\textbf{\footnotesize Texte grec établi --- Folio 12 (\texttt{f. 168r}, vv. 466–508 [1/3]) \hfill ΕΠΕΙΣΟΔΙΟΝ Α'}\\[-2pt]
\noindent\rule{16.5cm}{0.5pt}\\[2pt]
{\fontsize{8.8pt}{10.2pt}\selectfont
\ligneEst{\textbf{Μνησίλοχος}}{τὸ μὲν ὦ γυναῖκες ὀξυθυμεῖσθαι σφόδρα}{466}
\ligneEst{}{Εὐριπίδῃ, τοιαῦτ᾽ ἀκουούσας κακά,}{}
\ligneEst{}{οὐ θαυμάσιόν ἐστ᾽, οὐδ᾽ ἐπιζεῖν τὴν χολήν.}{}
\ligneEst{}{καὐτὴ γὰρ ἔγωγ᾽, οὕτως ὀναίμην τῶν τέκνων,}{}
\ligneEst{}{μισῶ τὸν ἄνδρ᾽ ἐκεῖνον, εἰ μὴ μαίνομαι.}{470}
\ligneEst{}{ὅμως δ᾽ ἐν ἀλλήλαισι χρὴ δοῦναι λόγον·}{}
\ligneEst{}{αὐταὶ γάρ ἐσμεν, κοὐδεμί᾽ ἔκφορος λόγου.}{}
\ligneEst{}{τί ταῦτ᾽ ἔχουσαι ᾽κεῖνον αἰτιώμεθα}{}
\ligneEst{}{βαρέως τε φέρομεν, εἰ δύ᾽ ἡμῶν ἢ τρία}{}
\ligneEst{}{κακὰ ξυνειδὼς εἶπε δρώσας μυρία;}{475}
\ligneEst{}{ἐγὼ γὰρ αὐτὴ πρῶτον, ἵνα μἄλλην λέγω,}{}
\ligneEst{}{ξύνοιδ᾽ ἐμαυτῇ πολλά γ᾽· ἀλλ᾽ ἐκεῖνο δ᾽ οὖν}{}
\ligneEst{}{δεινότατον, ὅτε νύμφη μὲν ἦ τρεῖς ἡμέρας,}{}
\ligneEst{}{ὁ δ᾽ ἀνὴρ παρ᾽ ἐμοὶ καθηῦδεν· ἦν δέ μοι φίλος,}{}
}
\vspace{2pt}
\noindent\rule{16.5cm}{0.6pt}\\[2pt]
\noindent\textbf{\small Commentaire philologique, métrique et scénique (vv. 466–508 --- partie 1/6)}\\[2pt]
\begin{adjustbox}{max totalheight=15.2cm, max width=16.5cm, keepaspectratio}%
\parbox[t]{16.5cm}{
\setlength{\columnsep}{9pt}%
\begin{multicols}{2}
\fontsize{8.4pt}{9.7pt}\selectfont
\setlength{\parindent}{0pt}%
\setlength{\parskip}{1.2pt}%
\noindent\textbf{466 sqq.} \textbf{} Scholiaste : \grec{ἐπιτηέδες ὀ κηδεστὴς ὑπὲρ τὰς ἄλλας γυναῖκας γυναικίζεται, ἵνα διὰ τούτου λαθῃ}... prudence qui va bientôt le quitter.
\noindent\textbf{473} \textbf{τί ταῦτ᾽ ἔχουσαι κεῖνον αἰτιώμεθα.} Ce vers rappelle celui des \italique{Acharniens} (vers 514) :
\citComm{τί ταῦτα τοὺς Λάκωνας αἰτιώμεθα;}
\noindent Ils pourraient donc bien être tous deux une reprise parodique du \italique{Télèphe} d'Euripide, pièce maintes fois détournée ici et là.
\noindent\textbf{477} \textbf{ξύνοιδ᾽ ἐμαυτῇ πολλ᾽ ἐκεῖνο δ᾽ οὖν.} Il manque un pied au vers 477 ; Bentley y suppléa en ajoutant <δεινὰ> πολλ᾽ pour introduire le superlatif δεινότατον du vers suivant. C'est une hypothèse (Dawes préféra πολλὰ δείν᾽). J'en propose une autre qui a l'avantage d'expliquer paléographiquement la lacune par un \italique{saut du même au }quasi-\italique{même} (il n'y a qu'une patte de mouche de différence entre un π et un γ) : \textbf{ΠΟΛΛ'} à la place de ΠΟΛΛΑ\textbf{Γ'ΑΛΛ'}.
\noindent\textbf{480} \textbf{με διεκόρευσεν ἑπτέτιν οὖσαν.} Mnésiloque n'aura pas mis longtemps à perdre toute retenue : il s'invente bien imprudemment une histoire repoussante qui n'a d'égale que celle de la Quartilla de Pétrone (\italique{Satyrico}, XXV) ou de la marquise de Brinvilliers. (Willems.)
\noindent\textbf{484} \textbf{στρόφος μ᾽ ἔχει τὴν γαστέρ᾽ ὦνερ κὠδύνη·} \grec{στρόφος} est le terme «~propre~» pour désigner la \italique{colique} (cf. Hippocrate, \italique{Des maladies aiguës}, XXXVII ; \italique{Aphorismes}, V, 41).
\noindent\textbf{486} \textbf{κᾆθ᾽ ὁ μὲν ἔτριβεν κεδρίδας ἄννηθον σφάκον·} Le niais mari se jette sur les graines de genévrier, d'aneth et de sauge car elles sont réputées préserver du mal que sa femme prétend avoir. Voir ce que dit Pline l'Ancien au sujet de ces plantes.
\noindent (Pline l'Ancien, \italique{Histoire naturelle}. D'après l'édition et la traduction d'Émile Littré, \italique{Histoire naturelle de Pline}, tome second (1877). XXIV, 36 pour le genévrier (ἡ κεδρίς) ; XX, 74 pour l'aneth (τό ἄννηθον) ; XXII, 71 pour la sauge (ὁ σφάκος).)
\end{multicols}
}%
\end{adjustbox}%
\vfill
\end{adjustwidth}
\newpage
% --- PAGE DE DROITE : Traduction française + Commentaire (2/6) ---
\begin{adjustwidth}{-2cm}{-2cm}
\noindent\textbf{\footnotesize Traduction française --- Folio 12 (\texttt{f. 168r}, vv. 466–508 [1/3]) \hfill ΕΠΕΙΣΟΔΙΟΝ Α'}\\[-2pt]
\noindent\rule{16.5cm}{0.5pt}\\[2pt]
{\fontsize{8.8pt}{10.2pt}\selectfont
\ligneEst{\textbf{Mnésiloque}}{\francais{Que vous soyez, ô femmes, vivement irritées}}{466}
\ligneEst{}{\francais{contre Euripide, qui si mal parla de vous,}}{}
\ligneEst{}{\francais{n'est pas pour étonner, ni que votre bile déborde.}}{}
\ligneEst{}{\francais{Car moi-même, — ainsi puissé-je être heureuse dans mes enfants, —}}{}
\ligneEst{}{\francais{je hais cet homme, ou je serais folle.}}{470}
\ligneEst{}{\francais{Néanmoins il convient que nous nous expliquions entre nous,}}{}
\ligneEst{}{\francais{car nous sommes seules, et rien ne sera rapporté de ce qui se dit ici.}}{}
\ligneEst{}{\francais{Pourquoi allons-nous l'accusant}}{}
\ligneEst{}{\francais{et sommes-nous indignées de ce que, sachant deux ou trois}}{}
\ligneEst{}{\francais{de nos méfaits, il les a divulgués, quand nous en commettons d'innombrables ?}}{475}
\ligneEst{}{\francais{Moi-même toute la première, pour ne rien dire des autres,}}{}
\ligneEst{}{\francais{j'en ai sur la conscience un grand nombre, mais voici}}{}
\ligneEst{}{\francais{le plus fort. J'étais mariée de trois jours,}}{}
\ligneEst{}{\francais{et mon mari dormait à mes côtés. Or j'avais un amant}}{}
}
\vspace{2pt}
\noindent\rule{16.5cm}{0.6pt}\\[2pt]
\noindent\textbf{\small Commentaire philologique, métrique et scénique (vv. 466–508 --- partie 2/6)}\\[2pt]
\begin{adjustbox}{max totalheight=15.2cm, max width=16.5cm, keepaspectratio}%
\parbox[t]{16.5cm}{
\setlength{\columnsep}{9pt}%
\begin{multicols}{2}
\fontsize{8.4pt}{9.7pt}\selectfont
\setlength{\parindent}{0pt}%
\setlength{\parskip}{1.2pt}%
\noindent\textbf{487} \textbf{ἐγὼ δὲ καταχέασα τοῦ στροφέως ὕδωρ.} La perfide est rouée (στρόφις) : elle asperge d'eau les gonds de la porte pour ne pas qu'ils fassent de bruit. Le scholiaste l'avait compris ainsi : ἵνα ὀλισθηρὸς γενόμενος μὴ ψοφῇ. On a le même tour chez Plaute (\italique{Curculio}, I, 3, 1) :
\citComm{\latin{Anus — Placide egredere, et sonitum prohibe forium, et \textbf{crepitum cardinum},\\
ne quod hic agimus, herus percipiat fieri, mea Planesium.\\
mane, \textbf{subfundam aquolam}.}
}
\noindent J'attire l'attention du lecteur sur le jeu de mots entre καταχέω / καταχέζω d'une part et στροφεύς / στρόφος d'autre part. Je n'en dis pas davantage car en voilà-z-assez !
\noindent\textbf{489} \textbf{παρὰ τὸν Ἀγυιᾶ κύβδ᾽ ἐχομένη τῆς δάφνης.} L'Agyieus, surnom tiré d'ἀγυιά, la \italique{rue}, donné à Apollon, considéré comme protecteur des rues, devint par extension le nom des images et autels élevés en l'honneur du dieu, à Athènes, devant les portes des maisons (ἐν προθύροις). (Note d'Émile Saglio\marqueur{1}.) Mnésiloque parle bien d'un autel, ceint selon l'usage de laurier ; Hésychius (α 856) :
\citComm{ἀγυιεύς· ὁ πρὸ τῶν θυρῶν ἑστὼς βωμὸς ἐν σχήματι κίονος\\
\francais{Agyieus : l'autel placé devant les portes, en forme de pilier.}}
\noindent Sachant ce qu'il y fait, Mnésiloque ajoute donc l'insulte au dieu à l'outrage de l'adultère. Au vu de la forme de l'ἀγυιεύς, on peut imaginer comment il prétend s'être tenu `παρὰ τὸν Ἀγυιᾶ κύβδα', mais pourquoi `ἐχομένη τῆς δάφνης', \italique{s'agrippant au laurier} ?! Quoique porté aux \italique{Realia}
\noindent , je me suis dispensé de reconstituer la scène ; cela ne m'empêchera pas d'affirmer : voilà une posture fort incommode !
\noindent Pourtant, cette image érotique d'un \italique{Penseur} de Rodin (κύβδα) se faisant levrer alors qu'il se tient à un arbre se retrouve ailleurs\marqueur{2} ; il y a un lièvre à lever.
\end{multicols}
}%
\end{adjustbox}%
\vfill
\end{adjustwidth}
\newpage

% --- FOLIO 12 (f. 168r) : SPREAD 2.2 (Texte grec établi + Traduction française + Commentaire 3/6 & 4/6) ---
\clearpage\ifodd\value{page}\hbox{}\newpage\fi
\enTeteG{ΘΕΣΜΟΦΟΡΙΑΖΟΥΣΑΙ --- Texte grec établi (f. 168r, vv. 466–508 [2/3])}
\enTeteD{LES THESMOPHORIES --- Traduction française (f. 168r, vv. 466–508 [2/3])}
\begin{adjustwidth}{-2cm}{-2cm}
\noindent\textbf{\footnotesize Texte grec établi --- Folio 12 (\texttt{f. 168r}, vv. 466–508 [2/3]) \hfill ΕΠΕΙΣΟΔΙΟΝ Α'}\\[-2pt]
\noindent\rule{16.5cm}{0.5pt}\\[2pt]
{\fontsize{8.8pt}{10.2pt}\selectfont
\ligneEst{}{ὅσπερ με διεκόρησεν οὖσαν ἑπτέτιν.}{480}
\ligneEst{}{οὖτος πόθῳ μου ᾽κνυεν ἐλθὼν τὴν θύραν·}{}
\ligneEst{}{κᾆτ᾽ εὐθὺς ἔγνων· εἶτα καταβαίνω λάθρᾳ.}{}
\ligneEst{}{ὁ δ᾽ ἀνὴρ ἐρωτᾷ `ποῖ σὺ καταβαίνεις;' `ὅποι;}{}
\ligneEst{}{στρόφος μ᾽ ἔχει τὴν γαστέρ᾽ ὦνερ κὠδύνη·}{}
\ligneEst{}{ἐς τὸν κοπρῶν᾽ οὖν ἔρχομαι.' `βάδιζέ νυν.'}{485}
\ligneEst{}{κᾆθ᾽ ὁ μὲν ἔτριβεν κεδρίδας ἄννηθον σφάκον·}{}
\ligneEst{}{ἐγὼ δὲ καταχέασα τοῦ στροφέως ὕδωρ}{}
\ligneEst{}{ἐξῆλθον ὡς τὸν μοιχόν· εἶτ᾽ ἠρειδόμην}{}
\ligneEst{}{παρὰ τὸν Ἀγυιᾶ κύβδ᾽ ἐχομένη τῆς δάφνης.}{}
\ligneEst{}{ταῦτ᾽ οὐδεπώποτ᾽ εἶφ᾽, ὁρᾶτ᾽, Εὐριπίδης·}{490}
\ligneEst{}{οὐδ᾽ ὡς ὑπὸ τῶν δούλων τε κὠρεωκόμων}{}
\ligneEst{}{σποδούμεθ᾽, ἢν μὴ 'χωμεν ἕτερον, οὐ λέγει·}{}
\ligneEst{}{οὐδ᾽ ὡς ὅταν μάλισθ᾽ ὑπό του ληκώμεθα}{}
}
\vspace{2pt}
\noindent\rule{16.5cm}{0.6pt}\\[2pt]
\noindent\textbf{\small Commentaire philologique, métrique et scénique (vv. 466–508 --- partie 3/6)}\\[2pt]
\begin{adjustbox}{max totalheight=15.2cm, max width=16.5cm, keepaspectratio}%
\parbox[t]{16.5cm}{
\setlength{\columnsep}{9pt}%
\begin{multicols}{2}
\fontsize{8.4pt}{9.7pt}\selectfont
\setlength{\parindent}{0pt}%
\setlength{\parskip}{1.2pt}%
\noindent La réponse se trouve peut-être chez les parémiographes : ils ont recensé une des \italique{formules d'Éphèse} (Ἐφέσια γράμματα), incantations magiques pour conjurer le mauvais sort\marqueur{3}. :
\citComm{Ἐφέσια γράμματα: ἐπῳδαί τινες ἦσαν, ἅσπερ οἱ φωνοῦντες ἐνίκων ἐν παντί. ταὐτὸ τῇ [lire : ταὐτότητι ?]·\
Δαφνίνην φορῶ βακτηρίαν.\\
Lettres éphésiennes : paroles magiques que ceux qui les prononçaient triomphaient en tout. En répétant :
« Je porte la branche de laurier ».}
\noindent Ainsi, les Grecs \italique{touchaient du bois} (de laurier) pour éviter les malheurs ; cela pourrait expliquer l'attitude de notre libidineuse : s'accrochant aux branches et à de superstitieuses croyances, elle pense ainsi échapper au mauvais œil.
\noindent\textbf{[Notes ad 489]} \textbf{(1)} Début de l'article \latin{Agyieus} de l'indispensable \italique{Dictionnaire des antiquités grecques et romaines} de Daremberg \& Saglio ; l'on peut y voir des représentations d'Agyeus.
\noindent \textbf{(2)} Lire Théocrite, \italique{Idylles}, V, 115-119 et voir, par exemple, – ce qui n'est d'ailleurs pas indispensable ! – \italique{Antike Denkmäler}, Band II, 1908, Tafel 41, Figure 3.
\noindent \textbf{(3)} \italique{Corpus paroemiographorum graecorum}, I, ed. E. L. a Leutsch, Diogeniani centuria IV 78, p. 244. On pourra consulter la \italique{Souda} ou Photius à l'entrée Ἀλεξιφάρμακα : ces charmes pouvaient concerner des époux. Aristophane met dans la bouche de Chrémyle qui prétend avoir quelque succès (\italique{Ploutos}, v. 213) : ὁ Φοῖβος αὐτὸς Πυθικὴν σείσας δάφνην.
\noindent\textbf{493} \textbf{μάλισθ᾽ ὑπό του ληκώμεθα.} Le scribe ici encore a hésité entre \grec{ληκώμεθα}, qu'il écrivit probablement en premier, avant de gratter les trois premières lettres et de réécrire par dessus \grec{κινώμεθα}. N'étant pas un aficionado de la \italique{lectio difficilior}, il aura probablement préféré le verbe \grec{κινεῖν}, bien connu de lui, au fort rare \italique{ληκᾶν} (ou \grec{ληκᾶσθαι} dans la bouche d'une femme). Suidas (λ 434) ne s'y est pas trompé qui glose ληκώμεθα par κινώμεθα en citant ce passage d'Aristophane.
\end{multicols}
}%
\end{adjustbox}%
\vfill
\end{adjustwidth}
\newpage
% --- PAGE DE DROITE : Traduction française + Commentaire (4/6) ---
\begin{adjustwidth}{-2cm}{-2cm}
\noindent\textbf{\footnotesize Traduction française --- Folio 12 (\texttt{f. 168r}, vv. 466–508 [2/3]) \hfill ΕΠΕΙΣΟΔΙΟΝ Α'}\\[-2pt]
\noindent\rule{16.5cm}{0.5pt}\\[2pt]
{\fontsize{8.8pt}{10.2pt}\selectfont
\ligneEst{}{\francais{qui m'avait déflorée à sept ans.}}{480}
\ligneEst{}{\francais{Pressé de désirs, il s'en vient gratter à la porte.}}{}
\ligneEst{}{\francais{Je comprends sur-le-champ et descends à petit bruit.}}{}
\ligneEst{}{\francais{Mon mari m'interroge : « Où descends-tu? — Où ?}}{}
\ligneEst{}{\francais{Je suis prise de tranchées, mon homme, et j'ai mal.}}{}
\ligneEst{}{\francais{Aussi vais-je aux latrines. — Va donc ».}}{485}
\ligneEst{}{\francais{Et là-dessus il se met à broyer des grains de genévrier, de l'aneth, de la sauge.}}{}
\ligneEst{}{\francais{Pour moi, ayant versé de l'eau sur les gonds,}}{}
\ligneEst{}{\francais{j'allai trouver mon galant et subir son assaut}}{}
\ligneEst{}{\francais{à côté de l'Aguieus, courbant l'échine et me tenant au laurier.}}{}
\ligneEst{}{\francais{Cela au grand jamais, voyez-vous, Euripide n'en a parlé.}}{490}
\ligneEst{}{\francais{Non plus qu'il ne dit qu'à défaut d'autres,}}{}
\ligneEst{}{\francais{nous nous faisons hocher par nos esclaves et par des muletiers ;}}{}
\ligneEst{}{\francais{ni qu'après avoir paillardé à cœur joie}}{}
}
\vspace{2pt}
\noindent\rule{16.5cm}{0.6pt}\\[2pt]
\noindent\textbf{\small Commentaire philologique, métrique et scénique (vv. 466–508 --- partie 4/6)}\\[2pt]
\begin{adjustbox}{max totalheight=15.2cm, max width=16.5cm, keepaspectratio}%
\parbox[t]{16.5cm}{
\setlength{\columnsep}{9pt}%
\begin{multicols}{2}
\fontsize{8.4pt}{9.7pt}\selectfont
\setlength{\parindent}{0pt}%
\setlength{\parskip}{1.2pt}%
\noindent\textbf{494} \textbf{ἕωθεν σκόροδα διαμασώμεθα.} (Pour ne pas que le mari cocu fleure le sacrifice à Vénus de sa femme.)  À quand une thesmophorienne ecclésie pour châtier Aristophane qui enseigne aux Athéniens toutes les fourberies de leur femme ?.. à moins que le tour ne fût déjà connu (Athénée, XIII, 572c.) :
\citComm{Διὰ ταῦθ' ὁ πόρνος οὗτος οὐδὲ τῶν πράσων\\
ἑκάστοτ' ἐπεδείπνει μεθ' ῆμῶν· τοῦτο δ' ἦν,\\
ἵνα μή τι λυπήσειε τὸν ἐραστὴν φιλῶν.}
\noindent et encore avec de l'oignon (Xénophon, \italique{Symposium}, IV, 8.) :
\citComm{Καὶ ὁ Χαρμίδης εἶπεν· `Ὦ ἄνδρες, ὁ Νικήρατος κρομμύων ὄζων ἐπιθυμεῖ οἴκαδε ἐλθεῖν, ἵν᾽ ἡ γυνὴ αὐτοῦ πιστεύῃ μηδὲ διανοηθῆναι μηδένα ἂν φιλῆσαι αὐτόν'\marqueur{2}.}
\noindent\textbf{495} \textbf{ἀπὸ τοῦ τείχους εἰσιὼν.} Thucydide (VIII, 69) :
\citComm{ἦσαν [δ᾽] Ἀθηναῖοι πάντες αἰεὶ οἱ μὲν ἐπὶ τείχει, οἱ δ᾽ ἐν τάξει, τῶν ἐν Δεκελείᾳ πολεμίων ἕνεκα ἐφ᾽ ὅπλοις.\\
\francais{Tous les Athéniens restaient continuellement, les uns sur les murailles, les autres à leurs postes, à cause des ennemis qui occupaient Décélie.}
}
\noindent\textbf{499} \textbf{τοὔγκυκλον.} L'ἔγκυκλον était une sorte de jupe courte, portée par les femmes à l'intérieur. Cf. la note au vers 643.
\noindent\textbf{500} \textbf{ὑπ᾽ αὐγὰς οἷόν ἐστιν ἐγκεκαλυμμένον.} Le manuscrit de Ravenne comporte :
\citComm{ὡς ἡ γυνὴ δεικνῦσα τἀνδρὶ τοὔγκυκλον\\
ὑπ᾽ αὐγὰς οἷόν ἐστιν ἐγκεκαλυμμένον\\
τὸν μυχὸν ἐξέπεμψεν, οὐκ εἴρηκέ πω.}
\end{multicols}
}%
\end{adjustbox}%
\vfill
\end{adjustwidth}
\newpage

% --- FOLIO 12 (f. 168r) : SPREAD 2.3 (Texte grec établi + Traduction française + Commentaire 5/6 & 6/6) ---
\clearpage\ifodd\value{page}\hbox{}\newpage\fi
\enTeteG{ΘΕΣΜΟΦΟΡΙΑΖΟΥΣΑΙ --- Texte grec établi (f. 168r, vv. 466–508 [3/3])}
\enTeteD{LES THESMOPHORIES --- Traduction française (f. 168r, vv. 466–508 [3/3])}
\begin{adjustwidth}{-2cm}{-2cm}
\noindent\textbf{\footnotesize Texte grec établi --- Folio 12 (\texttt{f. 168r}, vv. 466–508 [3/3]) \hfill ΕΠΕΙΣΟΔΙΟΝ Α'}\\[-2pt]
\noindent\rule{16.5cm}{0.5pt}\\[2pt]
{\fontsize{8.8pt}{10.2pt}\selectfont
\ligneEst{}{τὴν νύχθ᾽, ἕωθεν σκόροδα διαμασώμεθα,}{}
\ligneEst{}{ἵν᾽ ὀσφρόμενος ἁνὴρ ἀπὸ τείχους εἰσιὼν}{495}
\ligneEst{}{μηδὲν κακὸν δρᾶν ὑποτοπῆται. ταῦθ᾽, ὁρᾷς,}{}
\ligneEst{}{οὐπώποτ᾽ εἶπεν. εἰ δὲ Φαίδραν λοιδορεῖ,}{}
\ligneEst{}{ἡμῖν τί τοῦτ᾽ ἔστ᾽; οὐδ᾽ ἐκεῖν᾽ εἴρηκέ πω,}{}
\ligneEst{}{ὡς ἡ γυνὴ δεικνῦσα τἀνδρὶ τοὔγκυκλον}{500}
\ligneEst{}{τὸν μοιχὸν ἐξέπεμψεν, οὐκ εἴρηκέ πω.}{}
\ligneEst{}{ἑτέραν δ᾽ ἐγᾦδ᾽ ἣ 'φασκεν ὠδίνειν γυνὴ}{}
\ligneEst{}{δέχ᾽ ἡμέρας, ἔως ἐπρίατο παιδίον·}{}
\ligneEst{}{ὁ δ᾽ ἀνὴρ περιήρχετ᾽ ὠκυτόκι᾽ ὠνούμενος·}{}
\ligneEst{}{τὸ δ᾽ εἰσέφερε γραῦς ἐν χύτρᾳ τὸ παιδίον,}{505}
\ligneEst{}{ἵνα μὴ βοῴη, κηρίῳ βεβυσμένον·}{}
\ligneEst{}{εἶθ᾽ ὡς ἔνευσεν ἡ φέρουσ᾽, εὐθὺς βοᾷ,}{}
\ligneEst{}{`ἄπελθ᾽ ἄπελθ᾽, ἤδη γὰρ ὦνέρ μοι δοκῶ}{508}
}
\vspace{2pt}
\noindent\rule{16.5cm}{0.6pt}\\[2pt]
\noindent\textbf{\small Commentaire philologique, métrique et scénique (vv. 466–508 --- partie 5/6)}\\[2pt]
\begin{adjustbox}{max totalheight=15.2cm, max width=16.5cm, keepaspectratio}%
\parbox[t]{16.5cm}{
\setlength{\columnsep}{9pt}%
\begin{multicols}{2}
\fontsize{8.4pt}{9.7pt}\selectfont
\setlength{\parindent}{0pt}%
\setlength{\parskip}{1.2pt}%
\noindent Gardons-nous, dans un premier temps, de biffer le \nieme{500} vers, ὑπ᾽ αὐγὰς οἷόν ἐστιν ἐγκεκαλυμμένον, à l'instar de Willems\marqueur{1}. Si effectivement interpolation il y a -- ce que nous allons essayer de démontrer --, celle-ci ne peut être que très ancienne ; en effet, une scholie érudite notait déjà l'abrègement violent (συστολή βίαιος) que devait subir ὑπ᾽ αὐγὰς pour rentrer dans le mètre iambique.
\noindent Willems parlait d'une «~pure cheville~» à propos de `οἷόν ἐστιν'. Le cas me paraît plus compliqué.
\noindent Intéressons-nous y donc !\\ Après avoir retourné l'affaire dans tous les sens, je ne m'explique cet οἷόν ἐστιν que par la présence originelle de : ᾦον ἐστιν\marqueur{2}. Athénée (II, 57e) nous précise que ᾦον pouvait être utilisé pour ὑπερῷον, l'\italique{étage supérieur d'une maison}. Le scholiaste T de l'\italique{Iliade} (XVI, 184) emploie d'ailleurs le mot en ce sens -- était-ce le même que le scribe qui écrivit notre scholie devenue le vers 500\marqueur{3} ? Quoi qu'il en soit, l'adjectif ὑπερῷος signifie \italique{situé à l'étage supérieur}.
\noindent Si l'hypothèse est exacte, le sens se laisse alors clairement dégager : on comprend que l'amant (τὸν μοιχόν) regarde vers le bas (αὐγάζειν, au sens d'\italique{observer}, est préfixé de ὑπ- par décalque sur ὑφ-ορᾶν) la femme embobiner son cocu de mari, bien camouflé (ἐγκεκαλυμμένον). L'amant est donc à l'étage supérieur (ᾦον ἐστιν).
\noindent Une fois cette correction minime apportée, le \nieme{500} vers se peut lire ainsi :
\citComm{[τὸν μοιχόν·] ὑπαυγάσας. ᾦον ἐστιν ἐγκεκαλυμμένον.}
\noindent Ainsi compris, il devient tout-à-fait clair que ce « vers » n'est effectivement qu'une glose du vers 501. Le malheur voulut qu'elle fît peu ou prou un trimètre iambique et soit alors très tôt intégrée au texte même d'Aristophane ; l'inintelligibilité du nouveau texte aidant à sa corruption, ΥΠΑΥΓΑΣΑΣ devint ΥΠΑΥΓΑΣ par haplographie et ΩΙΟΝ ΟΙΟΝ.
\end{multicols}
}%
\end{adjustbox}%
\vfill
\end{adjustwidth}
\newpage
% --- PAGE DE DROITE : Traduction française + Commentaire (6/6) ---
\begin{adjustwidth}{-2cm}{-2cm}
\noindent\textbf{\footnotesize Traduction française --- Folio 12 (\texttt{f. 168r}, vv. 466–508 [3/3]) \hfill ΕΠΕΙΣΟΔΙΟΝ Α'}\\[-2pt]
\noindent\rule{16.5cm}{0.5pt}\\[2pt]
{\fontsize{8.8pt}{10.2pt}\selectfont
\ligneEst{}{\francais{la nuit, nous mâchons de l'ail le matin,}}{}
\ligneEst{}{\francais{afin que saisi par l'odeur, notre mari, à son retour du rempart,}}{495}
\ligneEst{}{\francais{n'ait aucun soupçon de notre conduite. Cela, voyez-vous,}}{}
\ligneEst{}{\francais{il n'en a jamais parlé. Que s'il invective Phèdre,}}{}
\ligneEst{}{\francais{que nous importe ? Il n'a pas raconté non plus le cas}}{}
\ligneEst{}{\francais{de cette femme qui, faisant voir son encycle à son mari,}}{500}
\ligneEst{}{\francais{fit partir son galant ; non, il ne l'a pas raconté.}}{}
\ligneEst{}{\francais{Et j'en sais une autre qui prétendit être en travail}}{}
\ligneEst{}{\francais{dix jours durant, jusqu'à ce qu'elle eût acheté un enfant.}}{}
\ligneEst{}{\francais{Le mari courait à la ronde faire emplette de drogues pour hâter la délivrance.}}{}
\ligneEst{}{\francais{L'enfant fut apporté par une vieille dans une marmite,}}{505}
\ligneEst{}{\francais{la bouche bourrée d'un rayon de miel pour qu'il ne criât point.}}{}
\ligneEst{}{\francais{Puis, sur un signe de la porteuse, la femme aussitôt s'écria :}}{}
\ligneEst{}{\francais{« Va-t'en, va-t'en, mon homme, je crois que je vais}}{508}
}
\vspace{2pt}
\noindent\rule{16.5cm}{0.6pt}\\[2pt]
\noindent\textbf{\small Commentaire philologique, métrique et scénique (vv. 466–508 --- partie 6/6)}\\[2pt]
\begin{adjustbox}{max totalheight=15.2cm, max width=16.5cm, keepaspectratio}%
\parbox[t]{16.5cm}{
\setlength{\columnsep}{9pt}%
\begin{multicols}{2}
\fontsize{8.4pt}{9.7pt}\selectfont
\setlength{\parindent}{0pt}%
\setlength{\parskip}{1.2pt}%
\noindent Lisons donc avec Willems :
\citComm{ὡς ἡ γυνὴ δεικνῦσα τἀνδρὶ τοὔγκυκλον\\
τὸν μοιχὸν ἐξέπεμψεν, οὐκ εἴρηκέ πω.\\
\francais{Comment la femme qui exhibant l'encycle à son mari\\
fit échapper son amant, il n'en a pas soufflé mot.}}
\noindent Le passage, ici comme ailleurs, n'est qu'allusif ; aussi fait-il probablement référence à une histoire bien connue des spectateurs. Victor Coulon\marqueur{4} l'a rapprochée de contes au motif semblable, couchés par écrit au \nieme{XI} siècle, et de source probablement arabe (parfois, la femme est aidée de sa mère dans son entreprise de mystification du cocu)~; les conteurs orientaux ayant eux-mêmes été familiarisés aux contes grecs. Ainsi, il nous apparaît plus que plausible qu'Aristophane rapporte ici non pas un fait-divers de la chronique scandaleuse du temps mais une historiette populaire ; ces «~révélations~» sur les Athéniennes ne seraient pas si compromettantes.\\ Voir le commentaire au vers 643 : l'encycle (ἔγκυκλον) est une courte jupe. Cette donnée finit d'anéantir les hypothèses loufoques  où la femme cacherait l'amant aux yeux du mari en tendant l'encycle devant l'amant.
\noindent\textbf{[Notes ad 500]} \textbf{(1)} Willems argue de sa faute de mètre et de sa construction des plus mal embouchées. Le lecteur pourra utilement se reporter à \nieme{13} note ainsi que son introduction aux \italique{Notes sur la} Paix (t. II, pp. 93-99), toutes deux reproduites en annexe. Le premier à notre connaissance à avoir proposé la suppression du vers 500 est van Herwerden (\italique{Mnemosyne N. S.}, 24, 1896, p.293).
\noindent \textbf{(2)} Rappelons que les iotas souscrits n'apparaissent qu'au cours du \nieme{XIII} siècle ; le manuscrit de Ravenne, du \nieme{IX} ou \nieme{X} siècle est adscrit tout comme les manuscrits qu'il transcrivait. La faute -- si faute il y a -- ne consiste par conséquent qu'en la confusion entre un Ω et un Ο.
\noindent \textbf{(3)} Ce même scholiaste T connaissait bien notre comédie : il en cite des passages comme les vers 536-537 (voir les \latin{testimonia} \italique{ad locum}.).
\noindent \textbf{(4)} «~Interprétation d'un passage d'Aristophane imité par les conteurs du Moyen-Âge~», in \italique{Annuaire de l'institut de philologie et d'histoire orientales}, tome II(1934) : \italique{Mélanges Bidez}, premier volume, p. 121-34.
\end{multicols}
}%
\end{adjustbox}%
\vfill
\end{adjustwidth}
\newpage
\clearpage\ifodd\value{page}\hbox{}\newpage\fi
\enTeteG{ΕΠΕΙΣΟΔΙΟΝ Α' --- Codex Ravennas 429, f. 168v (Classense 170v) --- vv. 509–550}
\begin{adjustwidth}{-1.5cm}{-3.0cm}
\begin{adjustbox}{max totalheight=21.0cm, max width=16.5cm, keepaspectratio}
\begin{minipage}{16.5cm}
{\fontsize{7.2pt}{8.5pt}\selectfont
\noindent\textbf{Folio 13 --- Codex Ravennas 429, \texttt{f. 168v} (\texttt{Classense 170v}) : ΕΠΕΙΣΟΔΙΟΝ Α', vv. 509–550}\hfill \textit{Témoins :} $\text{R}^1, \Rpc, \Rdeux$ (\Mon)\par\vspace{2pt}
\hrule height 0.4pt\vspace{2.5pt}
}
{\fontsize{8.1pt}{9.5pt}\selectfont
\ligneRav{509}{μ\auDessus{ν}{Η}}{τέξειν.' τὸ γὰρ ἦτρον τῆς χύτρας ἐλάκτισεν·}{}
\ligneRav{510}{}{χὠ μὲν γεγηθὼς \difference{ἀπέτρεχ᾽}, ἡ δ᾽ ἐξέσπασεν}{\textbf{510} ἔτρεχεν R : \lecon{ἀπέτρεχ᾽} van Herwerden}
\ligneRav{511}{}{ἐκ τοῦ στόματος τοῦ παιδίου, τὸ δ᾽ ἀνέκραγεν·}{\textbf{511} \lecon{τοῦ παιδίου} R : τὸ κηρίον Hirschig}
\ligneRav{512}{}{τὸ δ᾽ ἐξένευσε κἀνέκραγεν, \difference{ἥ ᾽φερεν} (\auDessus{εἰ}{η}\inRas{σ}{}έφερε)}{\textbf{512} \auDessus{εἰ}{η}\inRas{σ}{}έφερε \Rpc : \lecon{ἥ ᾽φερεν} Bentley}
\ligneRav{}{}{θεῖ μειδιῶσα πρὸς τὸν ἄνδρα καὶ λέγει,}{}
\ligneRav{514}{}{λέων λέων σοι γέγονεν, \difference{αὐτέκμαγμα} (αὖτ᾽ ἔ\inRas{γ}{κ}μαγμα) σόν,}{\textbf{514} αὖτ᾽ ἔκμαγμα (ex έγμαγμα) \Rpc, Pollux IX.131 : ἔκμαγμα Suid. λ 271, π 2114 : \lecon{αὐτέκμαγμα} Reiske}
\ligneRav{515}{}{τά τ᾽ ἄλλ᾽ ἁπαξάπαντα καὶ τὸ πόσθιον}{}
\ligneRav{516}{}{τῷ σῷ προσόμοιον, στρεβλὸν ὥσπερ κύτταρον.'}{\textbf{516} \lecon{κύτταρον} R, Suid. κ 2787 : κύτταρος Fritzsche}
\ligneRav{}{}{ταῦτ᾽ οὐ ποιοῦμεν τὰ κακά; νὴ τὴν Ἄρτεμιν}{}
\ligneRav{}{}{ἡμεῖς γε. κᾆτ᾽ Εὐριπίδῃ θυμούμεθα,}{}
\ligneRav{}{}{οὐδὲν παθοῦσαι μεῖζον ἢ δεδράκαμεν;}{}
\ligneRav{520}{—}{\hspace{1.1em}τουτὶ μέντοι θαυμαστόν,}{}
\ligneRav{521}{}{\hspace{1.1em}ὁπόθεν \difference{ηὑρέθη} τὸ χρῆμα,}{\textbf{521} εὑρέθη R : \lecon{ηὑρέθη} Meineke}
\ligneRav{}{}{\hspace{1.1em}χἤτις ἐξέθρεψε χώρα}{}
\ligneRav{}{}{\hspace{1.1em}τήνδε τὴν θρασεῖαν οὕτω·}{}
\ligneRav{}{}{\hspace{1.1em}τάδε γὰρ εἰπεῖν τὴν πανοῦργον}{}
\ligneRav{525}{}{\hspace{1.1em}κατὰ τὸ φανερὸν ὧδ᾽ ἀναιδῶς}{}
\ligneRav{526}{}{\hspace{1.1em}οὐκ ἂν ᾠόμην ἐν \difference{ἡμῖν}}{\textbf{526} ὑμῖν R : \lecon{ἡμῖν} Bonini}
\ligneRav{526}{}{\hspace{1.1em}οὐδὲ τολμῆσαί ποτ᾽ ἄν·}{\textbf{526} ὑμῖν R : \lecon{ἡμῖν} Bonini}
\ligneRav{527}{}{\hspace{1.1em}\difference{ἀλλὰ πᾶν} γένοιτ᾽ ἂν ἤδη·}{\textbf{527} ἀλλ᾽ ἅπαν R : \lecon{ἀλλὰ πᾶν} Cobet}
\ligneRav{}{}{\hspace{1.1em}τὴν παροιμίαν δ᾽ ἐπαινῶ}{}
\ligneRav{}{}{\hspace{1.1em}τὴν παλαιάν· ὑπὸ λίθῳ γὰρ}{}
\ligneRav{530}{}{\hspace{1.1em}παντί που χρὴ}{}
\ligneRav{530}{}{\hspace{1.1em}μὴ δάκῃ ῥητωρ ἀθρεῖν·}{}
\ligneRav{531}{χορ/}{ἀλλ᾽ οὐ γάρ ἐστι τῶν ἀναισχύτων φύσει γυναικῶν}{\textbf{531} (Κορ.) add. von Velsen}
\ligneRav{532}{}{οὐδὲν κάκιον εἰς ἅπαντα πλὴν ἄρ᾽ εἰ γυναῖκες·}{\textbf{532} \lecon{ἄρ᾽ εἰ} R : ἄλλαι Suid. α 1349}
\ligneRav{533}{γυνή τις}{οὔτοι μὰ τὴν \difference{Ἄγλαυρον} ὦ γυναῖκες εὖ φρονεῖτε,}{\textbf{533} (γυνὴ τις) \Rpc : \lecon{(Μίκα)} Rogers \espApp Ἄγραυλον R : \lecon{Ἄγλαυρον} Daubuz}
\ligneRav{}{}{ἀλλ᾽ ἢ πεφάρμαχθ᾽ ἢ κακόν τι μέγα πεπόνθατ᾽ ἄλλο,}{}
\ligneRav{535}{}{ταύτην ἐῶσαι τὴν φθόρον τοιαῦτα περιυβρίζειν}{}
\ligneRav{536}{}{ἡμᾶς ἁπάσας. εἰ μὲν οὖν τις ἔστιν· εἰ δὲ μή, ἡμεῖς}{\textbf{536} \lecon{μή ἡμεῖς} R : μἠμεῖς Bentley}
\ligneRav{537}{}{\difference{αὐταί} γε καὶ τὰ δουλάρια τέφραν ποθὲν λαβοῦσαι}{\textbf{537} αὐτοί R : \lecon{αὐταί} Bonini \espApp \lecon{γε} R : τε Σ T ad Hom.\italique{Il}.I.135}
\ligneRav{}{}{ταύτης ἀποψιλώσομεν τὸν χοῖρον, ἵνα διδαχθῇ}{}
\ligneRav{}{}{γυνὴ γυναῖκας οὖσα μὴ κακῶς λέγειν τὸ λοιπόν·}{}
\ligneRav{540}{μ\auDessus{ν}{Η}}{μὴ δῆτα τόν γε χοῖρον ὦ γυναῖκες. εἰ γὰρ οὔσης}{}
\ligneRav{}{}{παρρησίας κἀξὸν λέγειν ὅσαι πάρεσμεν ἀσταί,}{}
\ligneRav{542}{}{εἶτ᾽ εἶπον \difference{ἅ ᾽γίγνωσκον} ὑπὲρ Εὐριπίδου δίκαια,}{\textbf{542} ἅ γιγνωσκον R : \lecon{ἅ ᾽γίγνωσκον} Bentley}
\ligneRav{}{}{διὰ τοῦτο τιλλομένην με δεῖ δοῦναι δίκην ὑφ᾽ ὑμῶν;}{}
\ligneRav{}{\auDessus{γυ}{ν}}{οὐ γάρ σε δεῖ δοῦναι δίκην; ἥτις μόνη τέτληκας}{}
\ligneRav{545}{}{ὑπὲρ ἀνδρὸς ἀντειπεῖν, \difference{ὃς} ἡμᾶς πολλὰ κακὰ δέδρακεν}{\textbf{545} ὡς R : \lecon{ὃς} Grynaeus}
\ligneRav{546}{}{\difference{ἐπίτηδες} εὑρίσκων λόγους, ὅπου γυνὴ πονηρὰ}{\textbf{546} ἐξεπίτηδες R : \lecon{ἐπίτηδες} Scaliger}
\ligneRav{}{}{ἐγένετο, Μελανίππας ποιῶν Φαίδρας τε· Πηνελόπην δὲ}{}
\ligneRav{548}{}{οὐπώποτ᾽ \difference{ἐποίησ᾽}, ὅτι γυνὴ σώφρων ἔδοξεν εἶναι·}{\textbf{548} ἐποίησεν R : \lecon{ἐποίησ᾽} Scaliger}
\ligneRav{549}{μ\auDessus{ν}{Η}}{ἀλλ᾽ οὖν οὕνεκα τούτου γ᾽ \inRas{ἂρ}{ἂν} ὤφελες δοῦναι δίκην·}{\textbf{549} ἂν ex ἂρ corr. \Rpc}
\ligneRav{550}{}{τῶν νῦν γυναικῶν Πηνελόπην, Φαίδρας \difference{δ᾽ ἁπαξαπάσας}·}{\textbf{550} ἁπαξαπάσας R, M\expo{pc} : \lecon{δ᾽ ἁπαξαπάσας} M\expo{ac}}
}
\vspace{3pt}\noindent\rule{4.5cm}{0.4pt}\par\vspace{1.5pt}
{\fontsize{7.0pt}{8.3pt}\selectfont
\noindent\textbf{Loci paralleli (vv. 509–550) ---} \textbf{509} (τὸ γὰρ ἦτρον -- ἐλάκτισεν) Suidas η 637 ; Zonaras 1012 ἦτρον \espTest (ἦτρον) Moeris η 14 \espTest \textbf{514-6} Suid. λ 271 $\simeq$ Suid. π 2114 \espTest \textbf{514} Pollux IX 131 (λέων\expo{2} omisso et αὐτὸ ἔκμαγμα pro αὐτέκμαγμα) ; \Eustathe 1857.10, 462.40 ; cf. \AelDion ε 26 \espTest \textbf{515-6} (πόσθιον -- κύτταρον) Suidas κ 2787 \espTest \textbf{516} (στρεβλὸν) Hesychius σ 1977 στρεβλὸν \espTest {(κύτταρον)} Hesychius κ 4747 ; cf. Photius κ 1270 \espTest \textbf{527-30} Suidas α 1073 \espTest \textbf{528-30} \Stobee III 13.29 \espTest \textbf{531-2} Suidas α 1349 (sed ἄλλαι pro ἄρ᾽ εἰ) \espTest \textbf{536-7} (εἰ μὲν -- λαβοῦσαι) scholium\expo{bT} Homerus \italique{Il.} I 135-7 (ἔστι τις pro τις ἔστιν et αὐταί omisso)\par
}
\vspace{2.5pt}\noindent\rule{4.5cm}{0.4pt}\par\vspace{1.5pt}
{\fontsize{6.9pt}{8.2pt}\selectfont
\noindent\textbf{Suite (en commentaire) de l'apparat critique \& notes paléographiques (\texttt{f. 168v}) ---} \textbf{Correction complexe \italique{in rasura} et supralinéaire au v. 512 :} \Rac{} avait répété εἰσέφερε d'après le v. 509 ; \Rpc{} a gratté le $\sigma$ et tracé un $\eta$ au-dessus de εἰ (\lecon{\auDessus{εἰ}{η}\inRas{σ}{}έφερε}), visant ἣ 'φερεν (Bentley). --- \textbf{Correction au v. 514 :} αὖτ' ἔκμαγμα corrigé par \Rpc{} à partir de έγμαγμα. --- \textbf{Correction au v. 549 :} ἂν corrigé par \Rpc{} à partir de ἂρ. --- \textbf{v. 533 :} (γυνὴ τις) \Rpc{} : \lecon{(Μίκα)} Rogers \espApp Ἄγραυλον R (métathèse de liquides $\rho/\lambda$) : \lecon{Ἄγλαυρον} Daubuz. --- \textbf{v. 537 :} αὐτοί R : \lecon{αὐταί} Bonini \espApp \lecon{γε} R (τε $\Sigma$ T ad Hom. \italique{Il.} I.135).\par
}
\end{minipage}
\end{adjustbox}
\end{adjustwidth}
\newpage
\enTeteD{Codex Ravennas 429 (Biblioteca Classense) --- Fac-similé du f. 168v (vv. 509–550)}
\begin{adjustwidth}{-1.5cm}{-3.0cm}
\begin{center}
\includegraphics[width=16.5cm,height=20.5cm,keepaspectratio]{/workspaces/Thesmophories/dossier_ravennas_thesmophories/folios_individuels_ff162v-177r/13_f168v_classense170v_vv0509-0550.jpg}\par
\vspace{2pt}
{\scriptsize\textbf{Codex Ravennas 429, f. 168v (Classense 170v)} --- ΕΠΕΙΣΟΔΙΟΝ Α', \textbf{vv. 509–550} ($\text{R}^1, \Rpc, \Rdeux$ (\Mon)).}
\end{center}
\end{adjustwidth}
\newpage
\clearpage\ifodd\value{page}\hbox{}\newpage\fi% --- FOLIO 13 (f. 168v) : SPREAD 2.1 (Texte grec établi + Traduction française + Commentaire 1/4 & 2/4) ---
\clearpage\ifodd\value{page}\hbox{}\newpage\fi
\enTeteG{ΘΕΣΜΟΦΟΡΙΑΖΟΥΣΑΙ --- Texte grec établi (f. 168v, vv. 509–550 [1/2])}
\enTeteD{LES THESMOPHORIES --- Traduction française (f. 168v, vv. 509–550 [1/2])}
\begin{adjustwidth}{-2cm}{-2cm}
\noindent\textbf{\footnotesize Texte grec établi --- Folio 13 (\texttt{f. 168v}, vv. 509–550 [1/2]) \hfill ΕΠΕΙΣΟΔΙΟΝ Α'}\\[-2pt]
\noindent\rule{16.5cm}{0.5pt}\\[2pt]
{\fontsize{8.4pt}{9.7pt}\selectfont
\ligneEst{}{τέξειν.' τὸ γὰρ ἦτρον τῆς χύτρας ἐλάκτισεν·}{509}
\ligneEst{}{χὠ μὲν γεγηθὼς ἀπέτρεχ᾽, ἡ δ᾽ ἐξέσπασεν}{510}
\ligneEst{}{ἐκ τοῦ στόματος τοῦ παιδίου, τὸ δ᾽ ἀνέκραγεν.}{}
\ligneEst{}{εἶθ᾽ ἡ μιαρὰ γραῦς, ἣ ᾽φερεν τὸ παιδίον,}{}
\ligneEst{}{θεῖ μειδιῶσα πρὸς τὸν ἄνδρα καὶ λέγει,}{}
\ligneEst{}{`λέων λέων σοι γέγονεν, αὐτέκμαγμα σόν,}{}
\ligneEst{}{τά τ᾽ ἄλλ᾽ ἁπαξάπαντα καὶ τὸ πόσθιον}{515}
\ligneEst{}{τῷ σῷ προσόμοιον, στρεβλὸν ὥσπερ κύτταρον.'}{}
\ligneEst{}{ταῦτ᾽ οὐ ποιοῦμεν τὰ κακά; νὴ τὴν Ἄρτεμιν}{}
\ligneEst{}{ἡμεῖς γε. κᾆτ᾽ Εὐριπίδῃ θυμούμεθα,}{}
\ligneEst{}{οὐδὲν παθοῦσαι μεῖζον ἢ δεδράκαμεν;}{}
\ligneEst{\textbf{Χορός} {\tiny\textit{ἀντιστρ. / par.}}}{τουτὶ μέντοι θαυμαστόν,}{520}
\ligneEst{{\tiny\textit{$2$ troch.}}}{ὁπόθεν ηὑρέθη τὸ χρῆμα,}{}
\ligneEst{{\tiny\textit{$2$ troch.}}}{χἤτις ἐξέθρεψε χώρα}{}
\ligneEst{{\tiny\textit{$2$ troch.}}}{τήνδε τὴν θρασεῖαν οὕτω.}{}
\ligneEst{{\tiny\textit{$2$ troch.}}}{τάδε γὰρ εἰπεῖν τὴν πανοῦργον}{}
\ligneEst{{\tiny\textit{$2$ troch.}}}{κατὰ τὸ φανερὸν ὧδ᾽ ἀναιδῶς}{525}
\ligneEst{{\tiny\textit{$2$ troch.}}}{οὐκ ἂν ᾠόμην ἐν ἡμῖν}{}
\ligneEst{{\tiny\textit{$2$ troch\cat}}}{οὐδὲ τολμῆσαί ποτ᾽ ἄν.}{}
\ligneEst{{\tiny\textit{$2$ troch.}}}{ἀλλὰ πᾶν γένοιτ᾽ ἂν ἤδη·}{}
\ligneEst{{\tiny\textit{$2$ troch.}}}{τὴν παροιμίαν δ᾽ ἐπαινῶ}{}
\ligneEst{{\tiny\textit{$2$ troch.}}}{τὴν παλαιάν· ὑπὸ λίθῳ γὰρ}{}
}
\vspace{2pt}
\noindent\rule{16.5cm}{0.6pt}\\[2pt]
\noindent\textbf{\small Commentaire philologique, métrique et scénique (vv. 509–550 --- partie 1/4)}\\[2pt]
\begin{adjustbox}{max totalheight=13.0cm, max width=16.5cm, keepaspectratio}%
\parbox[t]{16.5cm}{
\setlength{\columnsep}{9pt}%
\begin{multicols}{2}
\fontsize{7.6pt}{8.8pt}\selectfont
\setlength{\parindent}{0pt}%
\setlength{\parskip}{1.2pt}%
\noindent\textbf{509} \textbf{τέξειν.} Cf. sur le mot τέξειν la note 14 de Willems à lire dans les annexes.
\noindent\textbf{509} \textbf{τὸ γὰρ ἦτρον τῆς χύτρας ἐλάκτισεν·} Il y a une plaisanterie \italique{contre l'attente} (παρὰ προσδοκίαν) : τῆς χύτρας est mis pour τῆς μητρός.
\noindent\textbf{515} \textbf{ἁπαξάπαντα.} La Fontaine, \italique{La chatte métamorphosée en femme} (II, 18) :
\citComm{\francais{La croit femme en tout et partout.}}
\noindent\textbf{516} \textbf{κύτταρον,} synonyme de κῶνος, \italique{pomme de pin}, que les Grecs appelaient aussi στρόβιλος (στρόμβος dans Nicandre. Hehn., p.292). (Note de Willems.)
\noindent\textbf{518-9} \textbf{θυμούμεθα, / οὐδὲν παθοῦσαι μεῖζον ἢ δεδράκαμεν;} On sait grâce au scholiaste que ces vers sont tirés du \italique{Télèphe} :
\citComm{\grec{παρὰ τὰ ἐκ Τηλέφου Εὐριπίδου· `εἰ δὴ θυμούμεθα παθόντες οὐδὲν μᾶλλον ἢ δεδρακότες'\marqueur{1}.}}
\noindent (fragment 20 du \italique{Télèphe} de l'édition Jouan-van Looy aux Belles Lettres, tome VIII, \nieme{3} partie. Correspond aux fragments 711 de l'édition Kannicht et 965 de celle de Mette.)
\noindent\textbf{520-30} \textbf{} Antistrophe du chant du chœur des vers 433-442.
\noindent\textbf{527} \textbf{ἀλλ᾽ ἅπαν γένοιτ᾽ ἂν ἤδη·} Le lieu-commun du \latin{omnia mutantur}\marqueur{1} avait bien-sûr cours aussi chez les Grecs.
\noindent Brève recension : \\ Hérodote (V, 9) :
\citComm{γένοιτο δ̓ ἂν πᾶν ἐν τῷ μακρῷ χρόνῳ.}
\noindent Sophocle (\italique{Philoctète}, vers 305-306) :
\citComm{πολλὰ γὰρ τάδε / ἐν τῷ μακρῷ γένοιτ᾽ ἂν ἀνθρώπων χρόνῳ.}
\noindent Eupolis\marqueur{2} :
\citComm{ἦ πολλά γ' ἐν μακρῷ χρόνῳ γίγνεται μεταλλαγῇ / (τῶν) πραγμάτων· μένει δὲ χρῆμ' οὐδὲν ἐν ταὐτῷ ῥυθμῷ.}
\noindent Ménandre (fragment 685(466)\marqueur{3}) :
\citComm{οὐκ ἔστ' ἄπιστον οὐδὲν ἐν θνητῷ βίῳ / οὐδ' ἂν γένοιτο.}
\noindent Et dans la pièce d'Aristophane d'antan, \italique{Lysistrata} (vers 258) :
\citComm{Χορὸς γερόντων --- ἦ πόλλ᾽ ἄελπτ᾽ ἔνεστιν ἐν τῷ μακρῷ βίῳ φεῦ.}
\end{multicols}
}%
\end{adjustbox}%
\vfill
\end{adjustwidth}
\newpage
% --- PAGE DE DROITE : Traduction française + Commentaire (2/4) ---
\begin{adjustwidth}{-2cm}{-2cm}
\noindent\textbf{\footnotesize Traduction française --- Folio 13 (\texttt{f. 168v}, vv. 509–550 [1/2]) \hfill ΕΠΕΙΣΟΔΙΟΝ Α'}\\[-2pt]
\noindent\rule{16.5cm}{0.5pt}\\[2pt]
{\fontsize{8.4pt}{9.7pt}\selectfont
\ligneEst{}{\francais{accoucher ». Car il avait donné des pieds contre la panse de la marmite.}}{509}
\ligneEst{}{\francais{Tandis qu'il s'éloigne tout joyeux, elle débarrasse}}{510}
\ligneEst{}{\francais{la bouche de l'enfant qui se met à crier.}}{}
\ligneEst{}{\francais{La vieille canaille qui l'a apporté}}{}
\ligneEst{}{\francais{court en souriant à l'époux, et lui dit :}}{}
\ligneEst{}{\francais{« un lion, un lion t'est né, ton portrait moulé ;}}{}
\ligneEst{}{\francais{outre tout le reste absolument, son membre}}{515}
\ligneEst{}{\francais{est pareil au tien et tortu comme une pigne ».}}{}
\ligneEst{}{\francais{Ne sont-ce pas là de nos tours? Si fait, par Artémis,}}{}
\ligneEst{}{\francais{c'en sont. Et nous nous fâchons contre Euripide,}}{}
\ligneEst{}{\francais{« sans avoir enduré par-delà nos mérites ».;}}{}
\ligneEst{\textbf{Chœur}}{\francais{Une chose me confond :}}{520}
\ligneEst{}{\francais{où a-t-on trouvé cette créature ?}}{}
\ligneEst{}{\francais{Quel pays a formé}}{}
\ligneEst{}{\francais{une telle effrontée ?}}{}
\ligneEst{}{\francais{La coquine ! tenir ce langage}}{}
\ligneEst{}{\francais{au grand jour, avec cette impudence ?}}{525}
\ligneEst{}{\francais{je n'aurais pas cru que parmi nous}}{}
\ligneEst{}{\francais{elle l'eût jamais osé.}}{}
\ligneEst{}{\francais{Mais maintenant tout arrive,}}{}
\ligneEst{}{\francais{et j'approuve l'antique proverbe,}}{}
\ligneEst{}{\francais{quasiment sous chaque pierre,}}{}
}
\vspace{2pt}
\noindent\rule{16.5cm}{0.6pt}\\[2pt]
\noindent\textbf{\small Commentaire philologique, métrique et scénique (vv. 509–550 --- partie 2/4)}\\[2pt]
\begin{adjustbox}{max totalheight=13.0cm, max width=16.5cm, keepaspectratio}%
\parbox[t]{16.5cm}{
\setlength{\columnsep}{9pt}%
\begin{multicols}{2}
\fontsize{7.6pt}{8.8pt}\selectfont
\setlength{\parindent}{0pt}%
\setlength{\parskip}{1.2pt}%
\noindent Remarquons que dans tous ces extraits, le grec insiste sur la longueur de temps que prennent les changements pour s'opérer (ἐν τῷ μακρῷ χρόνῳ / ἐν θνητῷ βίῳ / ἐν τῷ μακρῷ βίῳ)... or c'est tout le contraire ici : ἤδη, \italique{aujourd'hui}, \italique{actuellement}. Ce contraste est certainement une allusion d'Aristophane aux temps mouvementés politiquement de la représentation des \italique{Thesmophories} -- où tout semble s'accélérer.
\noindent\textbf{[Notes ad 527]} \textbf{(1)} Cf. Tosi, \italique{Dictionnaire des sentences latines et grecques}, §26. La citation est d'Ovide (\italique{Métamorphoses}, XV, 165). Érasme (\italique{Adages}, 663) recense la formule équivalente de Térence : \latin{Omnium rerum vicissitudo est}.
\noindent \textbf{(2)} Fragment 391(356) des \italique{Poetae Comici Graeci}, ed. Kassel \& Austin, V (1986).
\noindent \textbf{(3)} \italique{Poetae Comici Graeci}, VI 2 (1998). (Le fragment est extrait de Stobée I, 8, 31.)
\noindent\textbf{528-30} \textbf{τὴν παροιμίαν δ᾽ ἐπαινῶ / τὴν παλαιάν· ὑπὸ λίθῳ γὰρ / παντί που χρὴ / μὴ δάκῃ ῥητωρ ἀθρεῖν.} Le proverbe – avec un scorpion à la place de l'orateur – était effectivement populaire puisqu'il nous est diversement parvenu dans les recueils de Proverbes. Ainsi des parémiographes Zénobie\marqueur{1} ou Diogénianus\marqueur{2} et d'Hésychius\marqueur{3} qui rapportent la version la plus courante :
\citComm{ὑπὸ παντὶ λίθῳ σκόρπιος εὕδει: ἐπὶ τῶν κακοήθων λέγεται.}
\noindent Macaire\marqueur{4} en précise le sens : \grec{ἐπὶ πολυπληθίᾳ τῶν πανταχοῦ ὄντων πονηρῶν ἀνθρώπων}, « rapport au grand nombre d'hommes qui sont méchants sous tout rapport ».
\noindent Le proverbe est attesté aussi par Praxilla\marqueur{5} et Sophocle\marqueur{6}.\\ Ici, la \italique{pierre} (λίθος) dont parle la coryphée est le piédestal (βῆμα) de la Pnyx sur lequel les orateurs -- Mnésiloque, en l'occurrence -- parlaient. (Aristophane use du mot dans le même sens aux vers 683 des \italique{Acharniens} ; 680 de la \italique{Paix} et 87 de l'\italique{Assemblée des femmes}.) Hésychius (λ 999) :
\citComm{λίθος· βῶλος. βωμός. καὶ βάσις. τὸ ἐν τῇ Ἀθηναίων ἐκκλησίᾳ βῆμα.}
\noindent Enfin, la substitution au scorpion de l'orateur n'est pas gratuite : on parlait couramment de la \italique{morsure de la langue}\marqueur{7}.
\noindent Ainsi de la \italique{morsure du sycophante}, δῆγμα συκοφάντου, qui paraît proverbiale\marqueur{8}. Diogène le cynique eut cette réplique à l'ironie \italique{mordante} (rapportée par son homonyme laërtien) :
\citComm{ἐρωτηθεὶς τί τῶν θηρίων κάκιστα δάκνει ἔφη: τῶν μὲν ἀγρίων συκοφάντης, τῶν δὲ ἡμέρων κόλαξ.\\
Interrogé sur la bête dont la morsure est la plus mauvaise, [il répondit] que des bêtes sauvages c'était le sycophante, et des animaux domestiques, c'était le flatteur.}
\end{multicols}
}%
\end{adjustbox}%
\vfill
\end{adjustwidth}
\newpage

% --- FOLIO 13 (f. 168v) : SPREAD 2.2 (Texte grec établi + Traduction française + Commentaire 3/4 & 4/4) ---
\clearpage\ifodd\value{page}\hbox{}\newpage\fi
\enTeteG{ΘΕΣΜΟΦΟΡΙΑΖΟΥΣΑΙ --- Texte grec établi (f. 168v, vv. 509–550 [2/2])}
\enTeteD{LES THESMOPHORIES --- Traduction française (f. 168v, vv. 509–550 [2/2])}
\begin{adjustwidth}{-2cm}{-2cm}
\noindent\textbf{\footnotesize Texte grec établi --- Folio 13 (\texttt{f. 168v}, vv. 509–550 [2/2]) \hfill ΕΠΕΙΣΟΔΙΟΝ Α'}\\[-2pt]
\noindent\rule{16.5cm}{0.5pt}\\[2pt]
{\fontsize{8.4pt}{9.7pt}\selectfont
\ligneEst{{\tiny\textit{troch.}}}{παντί που χρὴ}{530}
\ligneEst{{\tiny\textit{$2$ troch\cat}}}{μὴ δάκῃ ῥητωρ ἀθρεῖν.}{}
\ligneEst{\textbf{Κορυφαῖος} {\tiny\textit{$4$ia\cat}}}{ἀλλ᾽ οὐ γάρ ἐστι τῶν ἀναισχύτων φύσει γυναικῶν}{}
\ligneEst{}{οὐδὲν κάκιον εἰς ἅπαντα πλὴν ἄρ᾽ εἰ γυναῖκες.}{}
\ligneEst{\textbf{Μίκα}}{οὔτοι μὰ τὴν Ἄγλαυρον ὦ γυναῖκες εὖ φρονεῖτε,}{}
\ligneEst{}{ἀλλ᾽ ἢ πεφάρμαχθ᾽ ἢ κακόν τι μέγα πεπόνθατ᾽ ἄλλο,}{}
\ligneEst{}{ταύτην ἐῶσαι τὴν φθόρον τοιαῦτα περιυβρίζειν}{535}
\ligneEst{}{ἡμᾶς ἁπάσας. εἰ μὲν οὖν τις ἔστιν· εἰ δὲ μή, ἡμεῖς}{}
\ligneEst{}{αὐταί γε καὶ τὰ δουλάρια τέφραν ποθὲν λαβοῦσαι}{}
\ligneEst{}{ταύτης ἀποψιλώσομεν τὸν χοῖρον, ἵνα διδαχθῇ}{}
\ligneEst{}{γυνὴ γυναῖκας οὖσα μὴ κακῶς λέγειν τὸ λοιπόν.}{}
\ligneEst{\textbf{Μνησίλοχος}}{μὴ δῆτα τόν γε χοῖρον ὦ γυναῖκες. εἰ γὰρ οὔσης}{540}
\ligneEst{}{παρρησίας κἀξὸν λέγειν ὅσαι πάρεσμεν ἀσταί,}{}
\ligneEst{}{εἶτ᾽ εἶπον ἅ ᾽γίγνωσκον ὑπὲρ Εὐριπίδου δίκαια,}{}
\ligneEst{}{διὰ τοῦτο τιλλομένην με δεῖ δοῦναι δίκην ὑφ᾽ ὑμῶν;}{}
\ligneEst{\textbf{Μίκα}}{οὐ γάρ σε δεῖ δοῦναι δίκην; ἥτις μόνη τέτληκας}{}
\ligneEst{}{ὑπὲρ ἀνδρὸς ἀντειπεῖν, ὃς ἡμᾶς πολλὰ κακὰ δέδρακεν}{545}
\ligneEst{}{ἐπίτηδες εὑρίσκων λόγους, ὅπου γυνὴ πονηρὰ}{}
\ligneEst{}{ἐγένετο, Μελανίππας ποιῶν Φαίδρας τε· Πηνελόπην δὲ}{}
\ligneEst{}{οὐπώποτ᾽ ἐποίησ᾽, ὅτι γυνὴ σώφρων ἔδοξεν εἶναι.}{}
\ligneEst{\textbf{Μνησίλοχος}}{ἐγὼ γὰρ οἶδα ταἴτιον. Μίαν γὰρ οὐκ ἂν εἴποις}{}
\ligneEst{}{τῶν νῦν γυναικῶν Πηνελόπην, Φαίδρας δ᾽ ἁπαξαπάσας.}{550}
}
\vspace{2pt}
\noindent\rule{16.5cm}{0.6pt}\\[2pt]
\noindent\textbf{\small Commentaire philologique, métrique et scénique (vv. 509–550 --- partie 3/4)}\\[2pt]
\begin{adjustbox}{max totalheight=13.0cm, max width=16.5cm, keepaspectratio}%
\parbox[t]{16.5cm}{
\setlength{\columnsep}{9pt}%
\begin{multicols}{2}
\fontsize{7.6pt}{8.8pt}\selectfont
\setlength{\parindent}{0pt}%
\setlength{\parskip}{1.2pt}%
\noindent\textbf{[Notes ad 528-30]} \textbf{(1)} \latin{Zenobii centuria} VI, 20 (tomus I). Je cite les proverbes d'après le recensement d'E. Leutsch et F. G. Schneidewin, \italique{Corpus paroemiographorum Graecorum}, en deux tomes. Signalons qu'une nouvelle édition critique des recensions de Zénobius et Diogénianus avec traduction en italien a paru en 2019 sous le titre \italique{I proverbi greci}, dir. Emanuelle Lelli.
\noindent \textbf{(2)} \latin{Diogeniani centuria} VIII, 59 (tomus I).
\noindent \textbf{(3)} Hésychius (υ 717). Photius et Suidas rapportent aussi ce proverbe : il est partout !
\noindent \textbf{(4)} \latin{Macarii centuria} VIII, 69 (tomus II).
\noindent \textbf{(5)} Le scholiaste attribue même le proverbe à Praxilla. Athénée (XV, 695d) donne un exemple de scolie (σκόλιον), sorte de chanson à boire, écrite sur ce proverbe ; comme Praxilla excella dans ce genre, il n'est pas impossible que la scolie fût d'elle, et que le scholiaste l'ait en tête lorsqu'il attribue un peu hâtivement la maternité du proverbe à Praxilla. Cette scolie correspond au fragment 903 des \italique{Poetae Melici Graeci}.
\noindent \textbf{(6)} Fragment 37 de l'édition Radt.
\noindent \textbf{(7)} Voir les sens figurés du verbe δάκνειν.
\noindent \textbf{(8)} Voir le \italique{Ploutos}, vers 885, et le commentaire du scholiaste \italique{ad. loc.}
\noindent\textbf{531-73} \textbf{} Tout ce passage est composé en tétramètres iambiques catalectiques ; mètre du conflit et des débats animés. Cf. aussi les \italique{Cavaliers} v. 335-366 ; \italique{Lysistrata} v. 350-381, 467-475.
\noindent\textbf{531-2} \textbf{ἀλλ᾽ οὐ γάρ ἐστι τῶν ἀναισχύτων φύσει γυναικῶν / οὐδὲν κάκιον εἰς ἅπαντα πλὴν ἄρ᾽ εἰ γυναῖκες.} Oct. Feuillet, \italique{Dalila}, partie IV :
\citComm{\francais{Rien de plus semblable sous le soleil qu'une femme – et une femme.}}
\noindent C'est la traduction presque textuelle de Ménandre (\italique{Sentences monostiques}, Γνῶμαι μονόστιχοι 109) :
\citComm{Γυνὴ γυναικὸς πώποτ᾽ οὐδὲν διαφέρει.}
\noindent (Note de Willems ; pour voir ce que Willems pense de la « manie des conjectures » visant à retoucher le texte d'Aristophane, on lira avec un plaisir non dissimulé la note entière n° 15, reproduite à la fin.)\\ Il serait piquant qu'Aristophane mette dans la bouche de la coryphée une citation d'Euripide lui-même pour dénigrer le comportement de certaines femmes ; le scholiaste n'en a malheureusement pas soufflé mot... on dispose néanmoins de citations d'Euripide très approchantes (fr. 494.27 Kannicht $=$ fr. 14.27 \italique{Mélanippe} Jouan-van Looy) :
\citComm{τῆς μὲν κακῆς κάκι[ο]ν οὐδὲν γίγνεται / γυναικός.}
\end{multicols}
}%
\end{adjustbox}%
\vfill
\end{adjustwidth}
\newpage
% --- PAGE DE DROITE : Traduction française + Commentaire (4/4) ---
\begin{adjustwidth}{-2cm}{-2cm}
\noindent\textbf{\footnotesize Traduction française --- Folio 13 (\texttt{f. 168v}, vv. 509–550 [2/2]) \hfill ΕΠΕΙΣΟΔΙΟΝ Α'}\\[-2pt]
\noindent\rule{16.5cm}{0.5pt}\\[2pt]
{\fontsize{8.4pt}{9.7pt}\selectfont
\ligneEst{}{\francais{il faut regarder}}{530}
\ligneEst{}{\francais{de crainte d'être mordu... par un orateur..}}{}
\ligneEst{\textbf{Coryphée}}{\francais{Mais les femmes au naturel impudent,}}{}
\ligneEst{}{\francais{il n'est rien de pire à tous égards, si ce n'est des femmes.}}{}
\ligneEst{\textbf{Mika}}{\francais{Non, par Aglaure, ô femmes, vous n'êtes pas dans votre bon sens,}}{}
\ligneEst{}{\francais{mais vous êtes maléficiées, ou vous avez pris quelque autre mal,}}{}
\ligneEst{}{\francais{pour que vous laissiez cette peste nous braver toutes à outrance.}}{535}
\ligneEst{}{\francais{(Elle se tourne vers le bureau de l'assemblée.) Ainsi qu'on se décide...}}{}
\ligneEst{}{\francais{sinon nous-mêmes et nos esclaves irons chercher quelque part de la cendre,}}{}
\ligneEst{}{\francais{et lui épilerons le bas-ventre, afin que, femme, elle apprenne}}{}
\ligneEst{}{\francais{à ne plus médire des femmes à l'avenir.}}{}
\ligneEst{\textbf{Mnésiloque}}{\francais{Le bas-ventre ? Non pas, ô femmes. Car la franchise}}{540}
\ligneEst{}{\francais{étant de droit et nous toutes citoyennes ci-présentes ayant congé de parler,}}{}
\ligneEst{}{\francais{si j'ai déduit les raisons que je savais en faveur d'Euripide,}}{}
\ligneEst{}{\francais{est-ce un motif pour m'épiler, afin de m'en faire porter la peine ?}}{}
\ligneEst{\textbf{Mika}}{\francais{Quoi! tu n'en porterais pas la peine ? toi qui seule osas}}{}
\ligneEst{}{\francais{répliquer en faveur d'un homme qui nous a tant malmenées,}}{545}
\ligneEst{}{\francais{trouvant à plaisir des sujets où une femme se montra mauvaise,}}{}
\ligneEst{}{\francais{produisant des Mélanippes, des Phèdres ; mais jamais il ne produisit une Pénélope,}}{}
\ligneEst{}{\francais{vu qu'elle passait pour sage.}}{}
\ligneEst{\textbf{Mnésiloque}}{\francais{J'en sais bien la raison. Tu n'en citerais pas une seule,}}{}
\ligneEst{}{\francais{parmi les femmes d'aujourd'hui, de Pénélope. Toutes Phèdres, sans exception.}}{550}
}
\vspace{2pt}
\noindent\rule{16.5cm}{0.6pt}\\[2pt]
\noindent\textbf{\small Commentaire philologique, métrique et scénique (vv. 509–550 --- partie 4/4)}\\[2pt]
\begin{adjustbox}{max totalheight=13.0cm, max width=16.5cm, keepaspectratio}%
\parbox[t]{16.5cm}{
\setlength{\columnsep}{9pt}%
\begin{multicols}{2}
\fontsize{7.6pt}{8.8pt}\selectfont
\setlength{\parindent}{0pt}%
\setlength{\parskip}{1.2pt}%
\noindent\textbf{533} \textbf{μὰ τὴν Ἄγραυλον.} Scholie d'un érudit :
\citComm{κατὰ τῆς ᾿Αγραύλου ὦμνυον, κατὰ δὲ τῆς Πανδρόσου σπανιώτερον· κατὰ δὲ τῆς Ἕρσης οὐχ εὑρήχαμεν.\\
\francais{L'on jurait sur Agraule, plus rarement sur Pandrosos ; nous n'avons pas trouvé d'exemple sur Ersès.}}
\noindent Aglaure (correction d'Agraule confirmée par l'épigraphie), Pandrose et Hersé étaient les trois filles de Kécrops, mythique fondateur et premier roi d'Athènes.
\noindent\textbf{549-50} \textbf{Μίαν γὰρ οὐκ ἂν εἴποις / τῶν νῦν γυναικῶν Πηνελόπην, Φαίδρας ἁπαξαπάσας.} Le comique Euboulos paraît s'être inspiré de ce passage :
\citComm{ὦ Ζεῦ πολυτίμητ', εἶτ' ἐγὼ κακῶς ποτε\\
ἐρῶ γυναῖκας ; νὴ Δί' ἀπολοίμην ἄρα\\
πάντων ἄριστον κτημάτων. Εἰ δ' ἐγένετο\\
κακὴ γυνὴ Μήδεια, Πηνελόπη δέ <γε>\\
μέγα πρᾶγμ'. ᾿Ερεῖ τις ὡς Κλυταιμνήστρα κακή·\\
Ἄλκηστιν ἀντέθηκα χρηστήν. ᾿Αλλ' ἴσως\\
Φαίδραν ἐρεῖ κακῶς τις · ἀλλὰ νὴ Δία\\
χρηστή τις ἦν μέντοι - τίς ; οἴμοι δείλαιος,\\
ταχέως γέ μ' αἱ χρησταὶ γυναῖκες ἐπέλιπον,\\
τῶν δ' αὖ πονηρῶν ἔτι λέγειν πολλὰς ἔχω.\\
\francais{Zeus très honoré, moi, dire du mal des femmes ? Ah! plutôt mourir ! Des femmes, le plus précieux de tous les biens ! Car si Médée s'est montrée méchante femme, Pénélope, elle, fut un vrai trésor. On dira que Clytemnestre fut mauvaise ? Je mets en regard la vertueuse Alceste. Mais peut-être médira-t-on de Phèdre ? Eh, par Zeus, vertueuse sans contredit fut.... qui donc ? Ah! malheureux ! les femmes vertueuses m'ont de suite fait faute, quand des mauvaises j'en ai encore un tas à citer.}}
\noindent (Meineke, \italique{Fragm. Com. Gr}., t. III, p. 260. [$=$ fragment 115 (116. 117) de l'édition Kassel-Austin, \italique{Poet. Com. Gr.}]) (Note de Willems.)
\noindent\textbf{550} \textbf{Πηνελόπην, Φαίδρας.} Si Phèdre est souvent citée par Aristophane comme (anti-)modèle d'impudicité (voir le vers 153 ; les \italique{Grenouilles} v. 1043 ; fr. 469(453) K.-A.(Kock)), Mélanippe apparaît seulement ici dans les œuvres conservées d'Aristophane.
\end{multicols}
}%
\end{adjustbox}%
\vfill
\end{adjustwidth}
\newpage
\clearpage\ifodd\value{page}\hbox{}\newpage\fi
\enTeteG{ΕΠΕΙΣΟΔΙΟΝ Α' --- Codex Ravennas 429, f. 169r (Classense 171r) --- vv. 551–589}
\begin{adjustwidth}{-1.5cm}{-3.0cm}
\begin{adjustbox}{max totalheight=21.0cm, max width=16.5cm, keepaspectratio}
\begin{minipage}{16.5cm}
{\fontsize{7.2pt}{8.5pt}\selectfont
\noindent\textbf{Folio 14 --- Codex Ravennas 429, \texttt{f. 169r} (\texttt{Classense 171r}) : ΕΠΕΙΣΟΔΙΟΝ Α', vv. 551–589}\hfill \textit{Témoins :} $\text{R}^1, \Rpc, \Rdeux$ (\Mon)\par\vspace{2pt}
\hrule height 0.4pt\vspace{2.5pt}
}
{\fontsize{8.1pt}{9.5pt}\selectfont
\ligneRav{551}{\auDessus{γυ}{ν}}{ἀκούετ᾽ ὦ γυναῖκες οἷ᾽ εἴρηκεν ἡ πανοῦργος}{}
\ligneRav{}{—}{ἡμᾶς ἁπάσας αὖθις αὖ. : καὶ νὴ Δί᾽ οὐδέπω γε}{}
\ligneRav{553}{}{εἴρηχ᾽ ὅσα ξύνοιδ᾽· ἐπεὶ βούλεσθε \difference{πλείον᾽} εἴπω;}{\textbf{553} πλεῖον R : \lecon{πλείον᾽} Bentley}
\ligneRav{554}{\auDessus{γυ}{ν}}{ἀλλ᾽ οὐκ ἂν ἔτ᾽ ἔχοις· ὅσα γὰρ \difference{ᾔδησθ᾽} ἐξέχεας ἅπαντα·}{\textbf{554} ᾔδεις R : \lecon{ᾔδησθ᾽} Hirschig (Ar. \italique{Vesp}.146)}
\ligneRav{555}{μ\auDessus{ν}{Η}}{μὰ Δί᾽ οὐδὲ τήν γε μυριοστὴν μοῖραν ὧν ποιοῦμεν·}{\textbf{555} \lecon{add. γε} Faber}
\ligneRav{556}{}{\difference{ἐπεὶ τάδ᾽} \difference{οὐκ} εἴρηχ᾽, ὁρᾷς, ὡς \difference{στλεγγίδας} λαβοῦσαι}{\textbf{556} ἔπειτά γ᾽ R : \lecon{ἐπεὶ τάδ᾽} Dindorf \espApp οὐχ R : \lecon{οὐκ} \Mon \espApp τλεγγίδας R : \lecon{στλεγγίδας} \Mon}
\ligneRav{557}{—}{ἔπειτα σιφωνίζομεν \difference{τὸν κύσθον}. : ἐπιτριβείης·}{\textbf{557} τὸν σῖτον R, Σ, Suid. σ 512 : \lecon{τὸν κύσθον} Bibi cf. v.1114}
\ligneRav{558}{μ\auDessus{ν}{Η}}{\difference{ὥς τ᾽ αὖ τὰ} κρέ᾽ ἐξ Ἀπατουρίων ταῖς μαστροποῖς διδοῦσαι}{\textbf{558} ὥς τ᾽ αὐτὰ R : ὥς ταῦτα \Mon : \lecon{ὥς τ᾽ αὖ τὰ} Zanetti}
\ligneRav{558}{—}{ἔπειτα τὴν γαλῆν φαμεν– : τάλαιν᾽ ἐγώ· φλυαρεῖς·}{\textbf{558} ὥς τ᾽ αὐτὰ R : ὥς ταῦτα \Mon : \lecon{ὥς τ᾽ αὖ τὰ} Zanetti}
\ligneRav{560}{μ\auDessus{ν}{Η}}{οὐδ᾽ ὡς τὸν ἄνδρα τῷ πελέκει γυνὴ κατεσπόδησεν,}{\textbf{560} \lecon{τὸν ἄνδρα τῷ πελέκει γυνὴ} R : ἑτέρα τὸν ἄνδρα τῷ πελέκει Suid. κ 996}
\ligneRav{}{}{οὐκ εἶπον· οὐδ᾽ ὡς φαρμάκοις ἑτέρα τὸν ἄνδρ᾽ ἔμηνεν,}{}
\ligneRav{562}{—}{οὐδ᾽ ὡς ὑπὸ τῇ πυέλῳ κατώρυξέν ποτε— : ἐξόλοιο·}{\textbf{562} \lecon{ποτε} R : ποτ᾽ ed. pr.}
\ligneRav{563}{μ\auDessus{ν}{Η}}{Ἀχαρνικὴ τὸν πατέρα. : ταυτὶ δῆτ᾽ ἀνέκτ᾽ ἀκούειν;}{\textbf{563} \lecon{ἀχαρνικὴ} R : Ἁχαρνικὴ Dobree}
\ligneRav{}{μ\auDessus{ν}{Η}}{οὐδ᾽ ὡς σὺ τῆς δούλης τεκούσης ἄρρεν εἶτα σαυτῇ}{}
\ligneRav{565}{}{τοῦθ᾽ ὑπεβάλου, τὸ σὸν δὲ θυγάτριον παρῆκας αὐτῇ·}{}
\ligneRav{}{\auDessus{γυ}{ν}}{οὔ τοι μὰ τὼ θεὼ σὺ καταπροίξει λέγουσα ταυτί,}{}
\ligneRav{567}{—}{ἀλλ᾽ ἐκποκιῶ σου τὰς ποκάδας. : οὐ \difference{δὴ} μὰ Δία σύ γ᾽ ἅψει·}{\textbf{567} \lecon{ποκάδας} R, Suid. : πλοκ- Σ \espApp δὲ R : \lecon{δὴ} Bothe}
\ligneRav{}{\auDessus{γυ}{ν}}{καὶ μὴν ἰδού. : καὶ μὴν ἰδού. : λαβὲ θοἰμάτιον Φιλίστη·}{}
\ligneRav{569}{μ\auDessus{ν}{Η}}{\difference{πρόσθιγε} μόνον, κἀγώ σε νὴ τὴν Ἄρτεμιν— : τί δράσεις;}{\textbf{569} πρόσθες R : πρόσιθι Cobet : \lecon{πρόσθιγε} Willems}
\ligneRav{570}{μ\auDessus{ν}{Η}}{τὸν σησαμοῦνθ᾽ ὃν κατέφαγες, τοῦτον χεσεῖν ποιήσω·}{}
\ligneRav{571}{χορ/}{\difference{παύσασθε} λοιδορούμεναι· καὶ γὰρ γυνή τις ἡμῖν}{\textbf{571} (Χορ.) R : \lecon{(Κορ.)} von Velsen \espApp παύσασθαι R : \lecon{παύσασθε} Μυ2}
\ligneRav{}{}{ἐσπουδακυῖα προστρέχει. πρὶν οὖν ὁμοῦ γενέσθαι,}{}
\ligneRav{}{}{σιγᾶθ᾽, ἵν᾽ αὐτῆς κοσμίως πυθώμεθ᾽ ἅττα λέξει·}{}
\ligneRav{}{κλεισθ/}{φίλαι γυναῖκες ξυγγενεῖς τοὐμοῦ τρόπου·}{}
\ligneRav{575}{}{ὅτι μὲν φίλος εἴμ᾽ ὑμῖν, ἐπίδηλος ταῖς γνάθοις·}{}
\ligneRav{576}{}{γυναικομανῶ γὰρ προξενῶ θ᾽ ὑμῶν ἀεί·}{\textbf{576} \lecon{ὑμῶν} R : ὑμᾶς Suid. π 2541}
\ligneRav{}{}{καὶ νῦν ἀκούσας πρᾶγμα περὶ ὑμῶν μέγα}{}
\ligneRav{}{}{ὀλίγῳ τι πρότερον κατ᾽ ἀγορὰν λαλούμενον,}{}
\ligneRav{}{}{ἥκω φράσων τοῦτ᾽ ἀγγελῶν θ᾽ ὑμῖν, ἵνα}{}
\ligneRav{580}{}{\difference{σκοπῆτε καὶ τηρῆτε} \difference{μὴ καὶ} προσπέσῃ}{\textbf{580} σκοπεἴτε καὶ τηρεἴτε R : \lecon{σκοπῆτε καὶ τηρῆτε} Biset \espApp καὶ μὴ R : \lecon{μὴ καὶ} Porson (\italique{Aristophanica} p. 216) \espApp ἡμῖν R : \lecon{ὑμῖν} \Mon \espApp ἀφράκτοις R : \lecon{ἀφάρκτοις} Dindorf (1830)}
\ligneRav{581}{}{\difference{ὑμῖν} \difference{ἀφάρκτοις} πρᾶγμα δεινὸν καὶ μέγα·}{\textbf{581} ἡμῖν R : \lecon{ὑμῖν} \Mon{} \espApp ἀφράκτοις R : \lecon{ἀφάρκτοις} Dindorf (1830, métathèse attique ancienne $-\rho\kappa-$)}
\ligneRav{}{χορ/}{τί δ᾽ ἔστιν ὦ παῖ; παῖδα γάρ σ᾽ εἰκὸς καλεῖν,}{}
\ligneRav{}{}{ἕως ἂν οὕτως τὰς γνάθους ψιλὰς ἔχῃς·}{}
\ligneRav{584}{κλεισθ/}{Εὐριπίδην \difference{φάσ᾽} ἄνδρα κηδεστήν τινα}{\textbf{584} ἔφας᾽ R : \lecon{φάσ᾽} Zanetti}
\ligneRav{585}{}{\difference{αὑτοῦ} γέροντα δεῦρ᾽ ἀναπέμψαι τήμερον·}{\textbf{585} αὐτου R : \lecon{αὑτοῦ} Biset}
\ligneRav{}{χορ/}{πρὸς ποῖον ἔργον ἢ τίνος γνώμης χάριν;}{}
\ligneRav{}{κλεισθ/}{ἵν᾽ ἅττα βουλεύοισθε καὶ μέλλοιτε δρᾶν,}{}
\ligneRav{}{}{ἐκεῖνος εἴη τῶν λόγων κατάσκοπος·}{}
\ligneRav{}{χορ/}{καὶ πῶς λέληθεν ἐν γυναιξὶν ὢν ἀνήρ;}{}
}
\vspace{3pt}\noindent\rule{4.5cm}{0.4pt}\par\vspace{1.5pt}
{\fontsize{7.0pt}{8.3pt}\selectfont
\noindent\textbf{Loci paralleli (vv. 551–589) ---} \textbf{554} (ᾔδησθ᾽) \EtymGen (A) ᾔδεισθα \espTest \textbf{557} (σιφωνίζομεν τὸν σῖτον) Suidas σ 512 \espTest (σιφωνίζομεν) Pollux VI 19 \espTest \textbf{560} Suidas κ 996 (ἑτέρα pro γυνὴ et κατεσπόδησε vel κατεσπόδωσεν) \espTest \textbf{566} Suidas κ 730 (sed καταπροίξῃ pro καταπροίξει) \espTest \textbf{567} (ἐκποκιῶ σου τὰς ποκάδας) Suidas ε 588 (ἐκπιῶ cod. A) ; Zonaras 672 ἐκποκιῶ \espTest (ἐκποκιῶ) Heschius ε 1637 ἐκποκιῶ ; \EtymMagnum 323.25 \espTest (ποκάδας) Hesychius π 2750 ποκάδες $=$ Photius π 1010 π\{λ\}οκάδες \espTest \textbf{570} (χεσεῖν) \EpimAlphab ε 201 \espTest \textbf{571-2} (πρὶν οὖν -- πυθώμεθ᾽) Suidas ο 292 (κοσμίως omisso) \espTest \textbf{574} (ξυγγενεῖς τοὐμοῦ τρόπου) Suidas ξ 83 \espTest \textbf{576} Suidas π 2541 προξενῶ θ' ὑμᾶς ἀεί\par
}
\vspace{2.5pt}\noindent\rule{4.5cm}{0.4pt}\par\vspace{1.5pt}
{\fontsize{6.9pt}{8.2pt}\selectfont
\noindent\textbf{Suite (en commentaire) de l'apparat critique \& notes paléographiques (\texttt{f. 169r}) ---} \textbf{vv. 556–557 (locus classicus de l'édition Louis Hauseux) :} --- \textbf{v. 557 :} τὸν σῖτον R, $\Sigma$, Souda $\sigma\ 512$ : \lecon{τὸν κύσθον} \lecon{Bibi} (cf. v. 1114). --- \textbf{v. 567 :} \lecon{ποκάδας} R, Souda (πλοκάδας $\Sigma$) \espApp δὲ R : \lecon{δὴ} Bothe. --- \textbf{v. 569 :} πρόσθες R (πρόσιθι Cobet) : \lecon{πρόσθιγε} Alphonse Willems (adopté par L. Hauseux). --- \textbf{v. 571 :} (Χορ.) R : \lecon{(Κορ.)} von Velsen \espApp παύσασθαι R (iotacisme $-\alpha\iota/-\varepsilon$) : \lecon{παύσασθε} \Mon$^{2}$. --- \textbf{v. 581 :} ἡμῖν R : \lecon{ὑμῖν} \Mon{} \espApp ἀφράκτοις R : \lecon{ἀφάρκτοις} Dindorf (1830, métathèse attique ancienne $-\rho\kappa-$).\par
}
\end{minipage}
\end{adjustbox}
\end{adjustwidth}
\newpage
\enTeteD{Codex Ravennas 429 (Biblioteca Classense) --- Fac-similé du f. 169r (vv. 551–589)}
\begin{adjustwidth}{-1.5cm}{-3.0cm}
\begin{center}
\includegraphics[width=16.5cm,height=20.5cm,keepaspectratio]{/workspaces/Thesmophories/dossier_ravennas_thesmophories/folios_individuels_ff162v-177r/14_f169r_classense171r_vv0551-0589.jpg}\par
\vspace{2pt}
{\scriptsize\textbf{Codex Ravennas 429, f. 169r (Classense 171r)} --- ΕΠΕΙΣΟΔΙΟΝ Α', \textbf{vv. 551–589} ($\text{R}^1, \Rpc, \Rdeux$ (\Mon)).}
\end{center}
\end{adjustwidth}
\newpage
\clearpage\ifodd\value{page}\hbox{}\newpage\fi% --- FOLIO 14 (f. 169r) : SPREAD 2.1 (Texte grec établi + Traduction française + Commentaire 1/8 & 2/8) ---
\clearpage\ifodd\value{page}\hbox{}\newpage\fi
\enTeteG{ΘΕΣΜΟΦΟΡΙΑΖΟΥΣΑΙ --- Texte grec établi (f. 169r, vv. 551–589 [1/4])}
\enTeteD{LES THESMOPHORIES --- Traduction française (f. 169r, vv. 551–589 [1/4])}
\begin{adjustwidth}{-2cm}{-2cm}
\noindent\textbf{\footnotesize Texte grec établi --- Folio 14 (\texttt{f. 169r}, vv. 551–589 [1/4]) \hfill ΕΠΕΙΣΟΔΙΟΝ Α'}\\[-2pt]
\noindent\rule{16.5cm}{0.5pt}\\[2pt]
{\fontsize{8.8pt}{10.2pt}\selectfont
\ligneEst{\textbf{Μίκα}}{ἀκούετ᾽ ὦ γυναῖκες οἷ᾽ εἴρηκεν ἡ πανοῦργος}{551}
\ligneEst{}{ἡμᾶς ἁπάσας αὖθις αὖ.}{}
\ligneEst{\textbf{Μνησίλοχος}}{\decalage{ἡμᾶς ἁπάσας αὖθις αὖ.}καὶ νὴ Δί᾽ οὐδέπω γε}{}
\ligneEst{}{εἴρηχ᾽ ὅσα ξύνοιδ᾽· ἐπεὶ βούλεσθε πλείον᾽ εἴπω;}{}
\ligneEst{\textbf{Μίκα}}{ἀλλ᾽ οὐκ ἂν ἔτ᾽ ἔχοις· ὅσα γὰρ ᾔδησθ᾽ ἐξέχεας ἅπαντα.}{}
\ligneEst{\textbf{Μνησίλοχος}}{μὰ Δί᾽ οὐδὲ τήν γε μυριοστὴν μοῖραν ὧν ποιοῦμεν.}{555}
\ligneEst{}{ἐπεὶ τάδ᾽ οὐκ εἴρηχ᾽, ὁρᾷς, ὡς στλεγγίδας λαβοῦσαι}{}
\ligneEst{}{ἔπειτα σιφωνίζομεν τὸν κύσθον.}{}
\ligneEst{\textbf{Μίκα}}{\decalage{ἔπειτα σιφωνίζομεν τὸν κύσθον.}ἐπιτριβείης.}{}
\ligneEst{\textbf{Μνησίλοχος}}{ὥς τ᾽ αὖ τὰ κρέ᾽ ἐξ Ἀπατουρίων ταῖς μαστροποῖς διδοῦσαι}{}
\ligneEst{}{ἔπειτα τὴν γαλῆν φαμεν–}{}
\ligneEst{\textbf{Μίκα}}{\decalage{ἔπειτα τὴν γαλῆν φαμεν–}τάλαιν᾽ ἐγώ· φλυαρεῖς.}{}
}
\vspace{2pt}
\noindent\rule{16.5cm}{0.6pt}\\[2pt]
\noindent\textbf{\small Commentaire philologique, métrique et scénique (vv. 551–589 --- partie 1/8)}\\[2pt]
\begin{adjustbox}{max totalheight=15.2cm, max width=16.5cm, keepaspectratio}%
\parbox[t]{16.5cm}{
\setlength{\columnsep}{9pt}%
\begin{multicols}{2}
\fontsize{7.6pt}{8.8pt}\selectfont
\setlength{\parindent}{0pt}%
\setlength{\parskip}{1.2pt}%
\noindent\textbf{556-7} \textbf{ὁρᾷς ὡς τλεγγίδας λαβοῦσαι / ἔπειτα σιφωνίζομεν τὸν σῖτον.} Voilà un trait sur les mœurs féminines qui n'a pas été compris des différents exégètes ; c'est donc à nous que revient la tâche – ô combien délicate ! – d'expliquer ce qui répugne aux oreilles. Muses, venez m'aider... mais vous êtes pucelles ; Muses, ne bougez donc. Ah ! que l'on se sent seul ; je ne peux même pas lancer en avant l'honorable professeur Jeffrey Henderson, le plus fin limier des obscénités d'Aristophane et dont le groin fureteur parvient même à les débusquer là où elles ne sont pas (celle-ci, du moins, lui aura échappé.) Mais la recherche est un plaisir solitaire. Attelons-nous y donc ! Pris littéralement, le texte a de quoi surprendre : « tu vois comme nous saisissons les strigiles, et qu'ensuite nous siphonnons le pain. (sic) » Si vous êtes pris de l'envie de corriger, avec Küster, τὸν σῖτον en τὸν οἶνον (l'on siphonnera plutôt du vin que du pain), renseignez-vous sur ce qu'est un strigile (une sorte de spatule ou plaque rectangulaire tout en longueur, légèrement recourbée pour servir de racloir lors de la toilette) ; j'offre l'édition d'Aristophane de Küster à qui parviendra à se saouler en buvant avec ce genre de coupe. (Il faudra d'ailleurs un jour que l'on m'explique l'intérêt d'introduire des émendations si c'est pour laisser un texte tout aussi obscur – et retouché en prime.) Qu'à cela ne tienne, ce n'est pas ce genre de peccadille qui doit arrêter Austin \& Olson : Aristote a dit un jour (\italique{Tropiques},145a24) qu'on pouvait se servir d'un strigile pour puiser de l'eau (hé ! ce sont des choses qui arrivent lorsqu'on est au bain ; A.\&O. auraient pu citer la phrase suivante qui leur était adressée~: Ἀλλ᾽ ὅμως εἴ τις ὁρίσαιτο τὴν στλεγγίδα ὄργανον πρὸς τὸ ἀρύειν, ἡμάρτηκεν, « Mais cependant, celui qui définirait le strigile comme instrument pour puiser se trompe »). Non, l'hypothèse d'ivrognesses s'enivrant à grandes rasades de strigiles est à écarter, sauf à m'expliquer : 1° comment boire au strigile \ \ \ 2° pourquoi aller chercher dans la salle de bain un instrument qui n'est pas censé en bouger \ \ \ 3° et quand bien même, qu'y aurait-il de déshonorant dans cette pratique ? \ \ \ Il faut chercher ailleurs, car les tentations des bains ne sont pas d'ordre biberonnesque. Aristophane vitupère ailleurs\marqueur{1} contre cette propension des Athéniens décadents à se délicater : les bains chauds sont un exemple parmi d'autres de ces mauvaises habitudes qui amollissent les cœurs et durcissent les vices. Pour comprendre exactement où il veut en venir, c'est Euripide lui-même qui nous donnera la clef de la salle de bain : dans son drame satyrique du \italique{Cyclope}, le héros aux-mille-tours-lubriques, se désole sur « son cher siphon » (v. 339, τὸν σίφωνα τὸν φίλον) qui depuis trop longtemps vit dans la solitude – solitude vis-à-vis du beau sexe ; le verbe χηρεύω servant aussi à désigner l'état de veuvage. Il parle évidemment de son πέος~; une fois attesté ce sens pour σίφων, celui du verbe σιφωνίζω s'éclaire d'une lumière nouvelle (le verbe est d'autant plus amusant qu'il joue sur le rapprochement, cher à Aristophane, entre βινεῖν et πίνειν).\\ La présence étonnante de τὸν σῖτον s'explique alors aisément : la leçon originale était τὸν κύσθον $=$ \latin{pudenda muliebria}, très proche phonétiquement de τὸν σῖτον.
\noindent Iotacisme aidant, il s'agit d'une simple \italique{métathèse} du `s' (avec chute du `k').
\end{multicols}
}%
\end{adjustbox}%
\vfill
\end{adjustwidth}
\newpage
% --- PAGE DE DROITE : Traduction française + Commentaire (2/8) ---
\begin{adjustwidth}{-2cm}{-2cm}
\noindent\textbf{\footnotesize Traduction française --- Folio 14 (\texttt{f. 169r}, vv. 551–589 [1/4]) \hfill ΕΠΕΙΣΟΔΙΟΝ Α'}\\[-2pt]
\noindent\rule{16.5cm}{0.5pt}\\[2pt]
{\fontsize{8.8pt}{10.2pt}\selectfont
\ligneEst{\textbf{Mika}}{\francais{Vous entendez, femmes, comment la coquine}}{551}
\ligneEst{}{\francais{nous arrange toutes, une fois de plus.}}{}
\ligneEst{\textbf{Mnésiloque}}{\francais{\espace{2.8cm}Et, par Zeus, je n'ai pas encore}}{}
\ligneEst{}{\francais{dit tout ce que je sais. En voulez-vous davantage ?}}{}
\ligneEst{\textbf{Mika}}{\francais{Tu ne pourrais. Ce que tu savais, tu l'as épanché.}}{}
\ligneEst{\textbf{Mnésiloque}}{\francais{Non, par Zeus, pas le dix-millième de ce que nous faisons.}}{555}
\ligneEst{}{\francais{Car enfin, je n'ai pas dit – tu vois ? – comment nous prenons des strigiles}}{}
\ligneEst{}{\francais{pour ensuite nous siphonner le con...}}{}
\ligneEst{\textbf{Mika}}{\francais{\espace{4.5cm}La peste soit de toi.}}{}
\ligneEst{\textbf{Mnésiloque}}{\francais{... que les viandes des Apaturies, nous les distribuons à nos procureuses,}}{}
\ligneEst{}{\francais{puis que nous accusons la belette...}}{}
\ligneEst{\textbf{Mika}}{\francais{\espace{3.9cm}Merci de moi ! Tu radotes.}}{}
}
\vspace{2pt}
\noindent\rule{16.5cm}{0.6pt}\\[2pt]
\noindent\textbf{\small Commentaire philologique, métrique et scénique (vv. 551–589 --- partie 2/8)}\\[2pt]
\begin{adjustbox}{max totalheight=15.2cm, max width=16.5cm, keepaspectratio}%
\parbox[t]{16.5cm}{
\setlength{\columnsep}{9pt}%
\begin{multicols}{2}
\fontsize{7.6pt}{8.8pt}\selectfont
\setlength{\parindent}{0pt}%
\setlength{\parskip}{1.2pt}%
\noindent Il est remarquable de noter que la même métathèse se produit au vers 1114 : le rustre archer Scythe, dans l'optique de prouver que Mnésiloque n'est pas une femme, conseille d'aller y voir de plus près :
\citComm{σκέψαι †τὸ σκύτο†.}
\noindent Scaliger \italique{à la grâce du mètre} -- et surtout à la grâce du sens -- avait rétabli :
\citComm{σκέψαι τὸ κύστο ($=$ τὸν κύσθον dans le parler scythe).}
\noindent Il est à croire qu'un scribe des \italique{Thesmophories} confondait avec un autre le mot grec pour désigner le « plus bel apanage » du corps féminin, qui tentait tendrement tant Brassens \& Villon : il s'y sera trompé par deux fois dans notre pièce (\latin{bis repetita...}).
\noindent J'arrête mes explications ici, de peur qu'un Calliclès de passage ne m'assène :
\citComm{Καλλίκλης-Φιλόλογος --- οὐκ αἰσχύνῃ εἰς τοιαῦτα ἄγων, ὦ Πέπωκες, τοὺς λόγους;\marqueur{2}}
\noindent Profitons-en pour remarquer à quel point les griefs des femmes envers Euripide ne sont pas sérieux dans la bouche d'Aristophane (si Euripide est misogyne, c'est bien là le moindre de ses défauts) ; car si les maris athéniens, s'en revenant d'une pièce d'Euripide, regardaient d'un œil soupçonneux leur femme tenant une marmite, avec quels yeux vont-ils dorénavant observer leur femme se saisir d'un strigile ou de tout instrument de forme analogue ?! Quoi qu'il en soit, comme pour l'histoire de l'amant qui s'échappe (v. 501), Mnésiloque n'invente rien mais reprend des thèmes certainement populaires du Landerneau athénien : le détournement à fin lubrique du strigile se retrouve illustrée, ce me semble, sur une coupe d'Onésimos\marqueur{4}, antérieure à la pièce d'Aristophane. Pour finir – était-il besoin de préciser ?.. – précisons tout-de-même : Mnésiloque mime en lançant son « ὁρᾷς » à Mika : « Ainsi font, font, font, les petites midinettes. Un siphon, fond, fond, on se siphonne toutes à fond.»
\noindent\textbf{[Notes ad 556-7]} \textbf{(1)} C'est un point de la querelle du Juste et de l'Injuste des \italique{Nuées} (v. 1044-1046) :
\citComm{Ἄδικος Λόγος — σε θερμῷ φησι λοῦσθαι πρῶτον οὐκ ἐάσειν.\\
καίτοι τίνα γνώμην ἔχων ψέγεις τὰ θερμὰ λουτρά;\\
Δίκαιος Λόγος — ὁτιὴ κάκιστόν ἐστι καὶ δειλὸν ποιεῖ τὸν ἄνδρα.\\
\francais{Raisonnement Injuste — il ne te permettra pas, dit-il, de te baigner dans l'eau chaude.\\
(Au Juste :) Mais d'après quel principe improuves-tu les bains chauds? \\
Raisonnement Juste — Parce qu'ils sont très pernicieux et qu'ils rendent l'homme lâche.}}
\noindent Platon, d'accord sur ce point avec Aristophane comme sur beaucoup d'autres, nomme les bains chauds γεροντικὰ λουτρά, « des bains de vieillards » (\italique{Lois}, VI, 761c). Érasme en a fait un adage (2061): ἁπαλοὶ θερμολουσίαις, \italique{amollis par les bains chauds}.
\end{multicols}
}%
\end{adjustbox}%
\vfill
\end{adjustwidth}
\newpage

% --- FOLIO 14 (f. 169r) : SPREAD 2.2 (Texte grec établi + Traduction française + Commentaire 3/8 & 4/8) ---
\clearpage\ifodd\value{page}\hbox{}\newpage\fi
\enTeteG{ΘΕΣΜΟΦΟΡΙΑΖΟΥΣΑΙ --- Texte grec établi (f. 169r, vv. 551–589 [2/4])}
\enTeteD{LES THESMOPHORIES --- Traduction française (f. 169r, vv. 551–589 [2/4])}
\begin{adjustwidth}{-2cm}{-2cm}
\noindent\textbf{\footnotesize Texte grec établi --- Folio 14 (\texttt{f. 169r}, vv. 551–589 [2/4]) \hfill ΕΠΕΙΣΟΔΙΟΝ Α'}\\[-2pt]
\noindent\rule{16.5cm}{0.5pt}\\[2pt]
{\fontsize{8.8pt}{10.2pt}\selectfont
\ligneEst{\textbf{Μνησίλοχος}}{οὐδ᾽ ὡς τὸν ἄνδρα τῷ πελέκει γυνὴ κατεσπόδησεν,}{560}
\ligneEst{}{οὐκ εἶπον· οὐδ᾽ ὡς φαρμάκοις ἑτέρα τὸν ἄνδρ᾽ ἔμηνεν,}{}
\ligneEst{}{οὐδ᾽ ὡς ὑπὸ τῇ πυέλῳ κατώρυξέν ποτε—}{}
\ligneEst{\textbf{Μίκα}}{\decalage{οὐδ᾽ ὡς ὑπὸ τῇ πυέλῳ κατώρυξέν ποτε—}ἐξόλοιο.}{}
\ligneEst{\textbf{Μνησίλοχος}}{Ἀχαρνικὴ τὸν πατέρα.}{}
\ligneEst{\textbf{Μίκα}}{\decalage{Ἀχαρνικὴ τὸν πατέρα.}ταυτὶ δῆτ᾽ ἀνέκτ᾽ ἀκούειν;}{}
\ligneEst{\textbf{Μνησίλοχος}}{οὐδ᾽ ὡς σὺ τῆς δούλης τεκούσης ἄρρεν εἶτα σαυτῇ}{}
\ligneEst{}{τοῦθ᾽ ὑπεβάλου, τὸ σὸν δὲ θυγάτριον παρῆκας αὐτῇ.}{565}
\ligneEst{\textbf{Μίκα}}{οὔ τοι μὰ τὼ θεὼ σὺ καταπροίξει λέγουσα ταυτί,}{}
\ligneEst{}{ἀλλ᾽ ἐκποκιῶ σου τὰς ποκάδας.}{}
\ligneEst{\textbf{Μνησίλοχος}}{\decalage{ἀλλ᾽ ἐκποκιῶ σου τὰς ποκάδας.}οὐ δὴ μὰ Δία σύ γ᾽ ἅψει.}{}
\ligneEst{\textbf{Μίκα}}{καὶ μὴν ἰδού.}{}
}
\vspace{2pt}
\noindent\rule{16.5cm}{0.6pt}\\[2pt]
\noindent\textbf{\small Commentaire philologique, métrique et scénique (vv. 551–589 --- partie 3/8)}\\[2pt]
\begin{adjustbox}{max totalheight=15.2cm, max width=16.5cm, keepaspectratio}%
\parbox[t]{16.5cm}{
\setlength{\columnsep}{9pt}%
\begin{multicols}{2}
\fontsize{7.6pt}{8.8pt}\selectfont
\setlength{\parindent}{0pt}%
\setlength{\parskip}{1.2pt}%
\noindent \textbf{(2)} Platon, \italique{Gorgias} 494e.
\noindent \textbf{(3)} Bâle, Antikenmuseum BS 440. ARV 326, 86 \italique{bis}. \italique{Para} 359. Elle est reproduite par John Boardman, fig. 233 in \italique{Les vases athéniens à figures rouges}, \italique{La période archaïque} ; première édition sous le titre \italique{Athenian Red Figure Vases}, \italique{The Archaic Period}, 1975. (Le lecteur pourra la retrouver aux annexes de ce livre.) En parcourant le livre de Boardmann, cette coupe a retenu mon attention : elle me paraît parfaitement illustrer l'anecdote de Mnésiloque. (Reconnaissons tout d'abord que l'interprétation est sujette à caution : j'ai montré cette coupe à différents amis et ai eu droit à autant d'interprétations qu'il y avait d'amis... interprétations qui en disaient d'ailleurs plus long sur leur propre psychologie que sur les mœurs des Athéniens ; profitons-en donc pour suggérer à l'éventuel psychanalyste qui se serait égaré à lire ces lignes d'abandonner fissa ses infantiles tests de Rorschach et de passer aux coupes athéniennes du \nieme{VI} siècle : dix vases athéniens d'Onésime seraient aussi révélateurs des obsessions ou perversions des supposés psychotiques que les taches d'encre du regretté Hermann, arraché trop tôt à l'affection des siens... Sans parler du gain en valeur artistique : incommensurable ! Refermons cette parenthèses et retournons à notre coupe :) qu'y ai-je donc vu ? À droite, un satyre (sa calvitie naissante est caractéristique des représentations de satyres de cette époque ; c'est entre autre sur ce détail de la laideur physique de Socrate, visible sur ses bustes, qu'Alcibiade rapproche selon moi Socrate du satyre Silène), un satyre donc, en train de se strigiler le fondement, pointe (du doigt) le Mont de Vénus d'une femme nue, lui désignant par là l'usage impropre du strigile auquel il lui arrive de céder. La clef qui me manque pour être tout-à-fait sûr de la bonne compréhension de cette scène est l'identification de l'objet dans la main de la femme : je crois y reconnaître un strigile enroulé dans du tissu. (plus recommandé pour atteindre le plaisir que la lame de bronze brute.) D'aucun y a vu le pagne du jeune homme à gauche. J'interprète cette scène comme un résumé des vices féminins de l'imaginaire grec : la femme est incapable de se contenir, elle cède à tous les plaisirs : le plaisir de la vigne (la femme est adossée à une outre) ; le plaisir de la chair (l'adultère est suggéré avec l'amant à gauche qui se rhabille) et -- toujours insatiable ! -- le plaisir des dieux. (nue, la voilà habillée pour l'hiver !..) Une autre hypothèse d'un mien ami quant à cette scène vaut d'être rapportée : il l'avait interprétée quelque peu différemment : tout en restant dans le domaine d'Onan, il proposait de garder τὸν σῖτον et de comprendre : \italique{le pain} [d'argile] servant à récupérer les peaux mortes \& autres hormones raclées avec l'huile à l'aide du strigile (les στλεγγίσματα et autre γλοιός), raclures qui étaient ensuite parfois monnayées (notamment lorsqu'elle provenaient de virils athlètes)... et pourquoi pas par quelque vicelarde pour s'en foutre plein les zones érogènes ? Son explication avait l'avantage indéniable de conserver la leçon du manuscrit... Mais je crains que son imagination ne se soit quelque peu emballée : je n'ai pas trouvé d'exemple d'acception du mot σῖτος pour désigner ce «~pain d'argile~». Et si Pline nous confirme un usage médical des στλεγγίσματα, l'usage détourné, qui lui a été suggéré par certains papiers de « recherche », est un pur fantasme des modernes scholiastes.
\end{multicols}
}%
\end{adjustbox}%
\vfill
\end{adjustwidth}
\newpage
% --- PAGE DE DROITE : Traduction française + Commentaire (4/8) ---
\begin{adjustwidth}{-2cm}{-2cm}
\noindent\textbf{\footnotesize Traduction française --- Folio 14 (\texttt{f. 169r}, vv. 551–589 [2/4]) \hfill ΕΠΕΙΣΟΔΙΟΝ Α'}\\[-2pt]
\noindent\rule{16.5cm}{0.5pt}\\[2pt]
{\fontsize{8.8pt}{10.2pt}\selectfont
\ligneEst{\textbf{Mnésiloque}}{\francais{... ni que telle autre occit son époux d'un coup de hache,}}{560}
\ligneEst{}{\francais{je ne l'ai pas dit ; ni qu'une autre rendit fou son mari par des drogues,}}{}
\ligneEst{}{\francais{ni que naguère sous la baignoire...}}{}
\ligneEst{\textbf{Mika}}{\francais{\espace{6cm}Crève...}}{}
\ligneEst{\textbf{Mnésiloque}}{\francais{... cette Acharnienne enfouit son père.}}{}
\ligneEst{\textbf{Mika}}{\francais{\espace{3.5cm}Supporterons-nous ce langage ?}}{}
\ligneEst{\textbf{Mnésiloque}}{\francais{Ni que toi, quand ta servante mit au jour un garçon,}}{}
\ligneEst{}{\francais{tu substituas l'enfant, et livras ta petite fille.}}{565}
\ligneEst{\textbf{Mika}}{\francais{Non, par les deux déesses, tu ne l'auras pas dit gratuitement.}}{}
\ligneEst{}{\francais{Je t'arracherai la toison.}}{}
\ligneEst{\textbf{Mnésiloque}}{\francais{\espace{3.3cm}Non, par Zeus, tu ne me toucheras pas.}}{}
\ligneEst{\textbf{Mika}}{\francais{Eh bien, tiens !}}{}
}
\vspace{2pt}
\noindent\rule{16.5cm}{0.6pt}\\[2pt]
\noindent\textbf{\small Commentaire philologique, métrique et scénique (vv. 551–589 --- partie 4/8)}\\[2pt]
\begin{adjustbox}{max totalheight=15.2cm, max width=16.5cm, keepaspectratio}%
\parbox[t]{16.5cm}{
\setlength{\columnsep}{9pt}%
\begin{multicols}{2}
\fontsize{7.6pt}{8.8pt}\selectfont
\setlength{\parindent}{0pt}%
\setlength{\parskip}{1.2pt}%
\noindent\textbf{558-9} \textbf{ὥς τ᾽ αὐτὰ κρέ᾽ ἐξ Ἀπατουρίων ταῖς μαστροποῖς διδοῦσαι / ἔπειτα τὴν γαλῆν φαμεν.} Si Mnésiloque avait commencé un discours tout en nuance, il a bien vite perdu de sa prudence et se met à surpasser le maître Euripide dans la diffamation du genre féminin. Le ton monte et Mnésiloque en est maintenant à dénoncer les pratiques les plus honteuses des Athéniennes les plus dissolues.\\ Nous ne nous attarderons pas sur ce que pouvaient bien procurer les procureuses en échange de ces détournements de viandes\marqueur{1}.
\noindent Plus intéressante est l'excuse alléguée pour expliquer cette soudaine disparition carnée. La \italique{team 1\expo{er}} degré des commentateurs d'Aristophane n'a pas manqué de faire remarquer que la belette était un animal chapardeur\marqueur{2}. Certes ! mais ce pauvre animal a surtout \italique{bon dos} : on l'accuse bien facilement ; ainsi aux vers 1150-1152 de la \italique{Paix} :
\citComm{ἦν δὲ καὶ πυός τις ἔνδον καὶ λαγῷα τέτταρα,\\
εἴ τι μὴ 'ξήνεγκεν αὐτῶν ἡ γαλῆ τῆς ἑσπέρας·\\
ἐψόφει γοῦν ἔνδον οὐκ οἶδ᾽ ἅττα κἀκυδοιδόπα\\
Il y avait aussi au logis de l'amouille et quatre morceaux de lièvre ; 
à moins que la belette n'en ait volé hier soir ; 
car il se passait là-dedans je ne sais quel bruit et quel remuement.}
\noindent L'allusion au remue-ménage n'est pas limpide mais il est clair que le sens ne peut être pris littéralement : qui, entendant une belette le dévaliser, ne s'armerait aussitôt d'un bâton pour châtier la voleuse ? (Noter d'ailleurs ici aussi la présence de l'article : ἡ γαλῆ.)  Aristophane joue certainement sur un détournement d'une locution proverbiale : la belette attique étant notre \italique{bouc-émissaire} à qui l'on fait \italique{porter le chapeau}.
\noindent Conformité du grec et du français : la belette joue le même rôle que notre chat, animal qui n'a pas encore été introduit dans l'Athènes du \nieme{V} siècle : il chasse les rongeurs. Or, que lit-on chez La Fontaine ?
\citComm{Le sommeil quitta son logis, /
Il eut pour hôte les soucis, /
Les soupçons, les alarmes vaines. / 
Tout le jour il avait l'œil au guet ; \textbf{et la nuit, /
Si quelque chat faisait du bruit, /
Le chat prenait l'argent.}\marqueur{3}}
\noindent Ce n'est pas moi qui m'abuse : remplacez l'argent par la nourriture et les deux passages disent exactement la même chose.
\noindent Voici ce que l'on trouve dans le \italique{Dictionnaire de la langue verte} d'Alfred Delvau :
\citComm{« C'est le chat ! » : Expression de l'argot du peuple, qui souligne ironiquement un doute, une dénégation. Ainsi, quelqu'un disant : Ce n'est pas moi qui ai fait cela. – Non ! c'est le chat, lui répondra-t-on.}
\noindent Un autre passage de la \italique{Paix} montre encore mieux ce sens imagé : aux vers 793-795 :
\citComm{καὶ γὰρ ἔφασχ᾽ ὁ πατὴρ ὃ παρ᾽ ἐλπίδας /
εἶχε τὸ δρᾶμα γαλῆν τῆς /
ἑσπέρας ἀπάγξαι.\\
Aussi bien leur père prétend-il qu'un drame qu'il pensait tenir contre toute attente fut étranglé le soir par une belette.}
\end{multicols}
}%
\end{adjustbox}%
\vfill
\end{adjustwidth}
\newpage

% --- FOLIO 14 (f. 169r) : SPREAD 2.3 (Texte grec établi + Traduction française + Commentaire 5/8 & 6/8) ---
\clearpage\ifodd\value{page}\hbox{}\newpage\fi
\enTeteG{ΘΕΣΜΟΦΟΡΙΑΖΟΥΣΑΙ --- Texte grec établi (f. 169r, vv. 551–589 [3/4])}
\enTeteD{LES THESMOPHORIES --- Traduction française (f. 169r, vv. 551–589 [3/4])}
\begin{adjustwidth}{-2cm}{-2cm}
\noindent\textbf{\footnotesize Texte grec établi --- Folio 14 (\texttt{f. 169r}, vv. 551–589 [3/4]) \hfill ΕΠΕΙΣΟΔΙΟΝ Α'}\\[-2pt]
\noindent\rule{16.5cm}{0.5pt}\\[2pt]
{\fontsize{8.8pt}{10.2pt}\selectfont
\ligneEst{\textbf{Μνησίλοχος}}{\decalage{καὶ μὴν ἰδού.}καὶ μὴν ἰδού.}{}
\ligneEst{\textbf{Μίκα}}{\decalage{καὶ μὴν ἰδού.καὶ μὴν ἰδού.}λαβὲ θοἰμάτιον Φιλίστη.}{}
\ligneEst{\textbf{Μνησίλοχος}}{πρόσθιγε μόνον, κἀγώ σε νὴ τὴν Ἄρτεμιν—}{}
\ligneEst{\textbf{Μίκα}}{\decalage{πρόσθιγε μόνον, κἀγώ σε νὴ τὴν Ἄρτεμιν—}τί δράσεις;}{}
\ligneEst{\textbf{Μνησίλοχος}}{τὸν σησαμοῦνθ᾽ ὃν κατέφαγες, τοῦτον χεσεῖν ποιήσω.}{570}
\ligneEst{\textbf{Κορυφαῖος}}{παύσασθε λοιδορούμεναι· καὶ γὰρ γυνή τις ἡμῖν}{}
\ligneEst{}{ἐσπουδακυῖα προστρέχει. πρὶν οὖν ὁμοῦ γενέσθαι,}{}
\ligneEst{}{σιγᾶθ᾽, ἵν᾽ αὐτῆς κοσμίως πυθώμεθ᾽ ἅττα λέξει.}{}
\ligneEst{\textbf{Κλεισθένης} {\tiny\textit{trim. iamb.}}}{φίλαι γυναῖκες ξυγγενεῖς τοὐμοῦ τρόπου·}{}
\ligneEst{}{ὅτι μὲν φίλος εἴμ᾽ ὑμῖν, ἐπίδηλος ταῖς γνάθοις·}{575}
\ligneEst{}{γυναικομανῶ γὰρ προξενῶ θ᾽ ὑμῶν ἀεί.}{}
\ligneEst{}{καὶ νῦν ἀκούσας πρᾶγμα περὶ ὑμῶν μέγα}{}
}
\vspace{2pt}
\noindent\rule{16.5cm}{0.6pt}\\[2pt]
\noindent\textbf{\small Commentaire philologique, métrique et scénique (vv. 551–589 --- partie 5/8)}\\[2pt]
\begin{adjustbox}{max totalheight=15.2cm, max width=16.5cm, keepaspectratio}%
\parbox[t]{16.5cm}{
\setlength{\columnsep}{9pt}%
\begin{multicols}{2}
\fontsize{7.6pt}{8.8pt}\selectfont
\setlength{\parindent}{0pt}%
\setlength{\parskip}{1.2pt}%
\noindent Juste avant ce passage, Aristophane faisait un honnête résumé des bienfaits qu'il a apportés à la Comédie, il en vient ensuite à étriller les nullités de ses poétereaux de contemporains. Nous en sommes à Carcinos \& Fils qui font des pieds et des mains pour pouvoir danser – c'est-à-dire obtenir la subvention d'un chœur.
\noindent Mais une fois le chœur accordé – à rebours de ce qui était attendu ! (παρ᾽ ἐλπίδας) [du bon goût ?] – Carcinos-père fait volte-face et prétend que son drame « fut étranglé le soir par une belette » (c'est encore le soir que la belette commet son forfait).
\noindent Le scholiaste nous est ici précieux puisqu'il précise que Carcinos voulait représenter une pièce intitulée les \italique{Souris} (Μῦες\marqueur{4}) ; encore que la chronique du temps ne nous ait pas conservé l'affaire, il paraît aisé de reconstituer peu ou prou ce qui s'est passé. Les poètes devaient présenter à l'archonte éponyme leur projet.
\noindent Pour les Dionysies, l'archonte sélectionnait trois tétralogies qui recevaient chacune le soutien financier d'un riche Athénien : le chorège.
\noindent Ce n'est qu'alors que le poète se mettait à écrire. Et c'est à cette étape que les forces auront probablement manqué à notre Carcinos… ce dernier croyant s'en tirer avec un trait d'esprit : s'il ne peut pas présenter ses \italique{Souris}, c'est que la belette les a étranglées.
\noindent Le mot ne peut être \italique{bon} que s'il existait effectivement une expression pour faire endosser la responsabilité d'un méfait à cette belette proverbiale.
\noindent\textbf{[Notes ad 558-9]} \textbf{(1)} Contentons-nous de renvoyer à la magistrale note d'Alphonse Willems sur la fête des Apaturies, t. II, pp. 118-120 ; le lecteur y découvrira l'importance démocratique de la κρεανομία, l'institution de la distribution de viande. Sur les débordements que de telles orgies pouvaient entraîner, l'on pourra – un exemple parmi d'autres – essayer de comprendre le jeu de mots d'Aristophane sur l'ἀνάρρυσις (au vers 890 de la \italique{Paix}), le second jour des Apaturies.
\noindent \textbf{(2)} Voir par exemple les vers 363-364 des \italique{Guêpes}. Lucien, \italique{Piscator}, 34 : ἁρπακτικώτεροι τῶν γαλῶν.
\noindent \textbf{(3)} Il s'agit – vous l'aurez bien sûr reconnue ; sinon, cessez immédiatement votre lecture et ne revenez que lorsque vous la connaîtrez par cœur ! – de la seconde fable du huitième livre : \italique{Le Savetier et le Financier}.
\noindent \textbf{(4)} Probablement le drame satyrique. Μῦς étant un nom propre, comme chez Hérodote (VIII, 133), dans la \italique{Souda} (ο 680) ou chez Athénée (XI, 782b), les Μῦες désignant les enfants de ce Μῦς. Je m'inscris en faux contre l'interprétation de Paul Mazon (\italique{La} Paix\italique{, texte grec publié avec une introduction, des notes critiques et explicatives} (1904), v. 793) reprise par Taillardat (§792, p.463 n.5) qui pensent que Μῦες était le surnom de ces débauchés de Carcinites. L'on s'expliquerait difficilement la confusion du scholiaste (présent sur tous les bons manuscrits et repris par Suidas κ 396). Objection plus forte : cela n'expliquerait pas le texte de la \italique{Paix} : la belette (γαλῆ) est accusé d'avoir étranglé nuitamment non les choreutes mais la pièce elle-même (τὸ δρᾶμα).
\end{multicols}
}%
\end{adjustbox}%
\vfill
\end{adjustwidth}
\newpage
% --- PAGE DE DROITE : Traduction française + Commentaire (6/8) ---
\begin{adjustwidth}{-2cm}{-2cm}
\noindent\textbf{\footnotesize Traduction française --- Folio 14 (\texttt{f. 169r}, vv. 551–589 [3/4]) \hfill ΕΠΕΙΣΟΔΙΟΝ Α'}\\[-2pt]
\noindent\rule{16.5cm}{0.5pt}\\[2pt]
{\fontsize{8.8pt}{10.2pt}\selectfont
\ligneEst{\textbf{Mnésiloque}}{\francais{\espace{2.1cm}Eh bien, tiens !}}{}
\ligneEst{\textbf{Mika}}{\francais{\espace{4.3cm}Prends mon manteau, Philista.}}{}
\ligneEst{\textbf{Mnésiloque}}{\francais{Atteins-moi seulement, et par Artémis...}}{}
\ligneEst{\textbf{Mika}}{\francais{\espace{6.2cm}Que feras-tu ?}}{}
\ligneEst{\textbf{Mnésiloque}}{\francais{Le gâteau de sésame que tu as dévoré, je te le ferai évacuer.}}{570}
\ligneEst{\textbf{Coryphée}}{\francais{Trêve d'invectives. Car une femme}}{}
\ligneEst{}{\francais{accourt vers nous en toute hâte. Avant qu'elle soit proche,}}{}
\ligneEst{}{\francais{faites silence. Écoutons en bonne contenance ce qu'elle va nous dire.}}{}
\ligneEst{\textbf{Clisthène}}{\francais{Femmes aimées, qui m'êtes unies par la parenté des mœurs,}}{}
\ligneEst{}{\francais{que je vous suis ami, mes joues le prouvent assez.}}{575}
\ligneEst{}{\francais{Oui, j'aime les femmes à la rage, et toujours je m'emploie pour vous.}}{}
\ligneEst{}{\francais{Aussi, ayant appris sur votre sujet une chose capitale,}}{}
}
\vspace{2pt}
\noindent\rule{16.5cm}{0.6pt}\\[2pt]
\noindent\textbf{\small Commentaire philologique, métrique et scénique (vv. 551–589 --- partie 6/8)}\\[2pt]
\begin{adjustbox}{max totalheight=15.2cm, max width=16.5cm, keepaspectratio}%
\parbox[t]{16.5cm}{
\setlength{\columnsep}{9pt}%
\begin{multicols}{2}
\fontsize{7.6pt}{8.8pt}\selectfont
\setlength{\parindent}{0pt}%
\setlength{\parskip}{1.2pt}%
\noindent\textbf{568} \textbf{λαβὲ θοἰμάτιον Φιλίστη.} Mettre bas les manteaux était de mise avant d'en venir aux mains. Ainsi au vers 408 des \italique{Guêpes} ; autre exemple : le fragment d'Hipponax\marqueur{1} contre Boupalos\marqueur{2} : Boupalos, sculpteur de Chios, avait fait d'Hipponax un portrait où il exagérait encore la laideur du modèle. Pour se venger, le poète écrivit contre lui des vers mordants et pleins de menaces. Suidas (β 452) nous a conservé celui-ci, qui paraît avoir passé en proverbe :
\citComm{λάβετε μευ ταἰμάτια [θοιμάτιον codd.], κόψω Βουπάλου τον οφθαλμόν.\\
\francais{Prenez mon manteau, que je fasse sauter un œil à Boupalos.}}
\noindent\textbf{[Notes ad 568]} \textbf{(1)} Fr. 121 de l'édition Degani : \italique{Hipponactis testimonia et fragmenta}, Teubner, 1991 ; fr. 120 chez Olivier Masson, voir aussi son introduction p. 13.
\noindent \textbf{(2)} Il est question de ce Boupalos au vers 361 de \italique{Lysistrata} -- le fragment sus-cité d'Hipponax nous a d'ailleurs été transmis par le scholiaste à ce passage ; je reprends la note de Willems de ce passage, p. 351 du second tome.
\noindent\textbf{569} \textbf{πρόσθες μόνον.} Willems substitue πρόσθιγε (du verbe προσθιγγάνω, \italique{toucher}) à πρόσθες qu'il soupçonne de ne pas être le mot propre. Il s'en explique succinctement dans la \nieme{16} note reproduite à la fin. D'aucuns ont essayé de justifier la leçon du texte en avançant une tournure elliptique : πρόσθες [τὴν χεῖρα] (expression qui se trouve effectivement au vers 1669 des \italique{Phéniciennes} d'Euripide\marqueur{1}) mais tous les exemples proposés comportent χεῖρα, qui semble indispensable au sens. Cobet avait donc proposé de lire πρόσιθι μόνον, \italique{approche seulement} ; ce que l'on dirait effectivement dans nos langues modernes mais le grec use d'une image différente : il dit \italique{jette(-moi) seulement ta main} (cf. les exemples déjà cités) ou bien \italique{touche(-moi) seulement} (cf. les exemples proposés par Willems auxquels on peut ajouter \italique{Lysistrata} v. 365, qui est l'exacte correspondance de son émendation : ἅψαι μόνον).
\noindent\textbf{[Notes ad 569]} \textbf{(1)} L'on cite d'autres exemples équivalents : le vers 933 des \italique{Nuées}: κλαύσει, τὴν χεῖρ᾽ ἢν ἐπιβάλλῃς ; le fragment 309(277) (éd. K.-A. = \italique{Zonarae lexikon} p.1168) : τὴν χεῖρα μὴ ᾿πίβαλλε, μὴ κλάων καθῇ ; \italique{Lysistrata}, vers 435-436 \& 439-440 : τὴν χεῖρά μοι ... προσοίσει ... κλαύσεται \& εἴ ... μόνον τὴν χεῖρ᾽ ἐπιβαλεῖς, ἐπιχεσεῖ πατούμενος.
\noindent\textbf{570} \textbf{τὸν σησαμοῦνθ᾽ ὃν κατέφαγες, τοῦτον χεσεῖν ποιήσω.} «~Le gâteau de sésame que tu as engoulé, je vais te le faire chier.~» Mnésiloque aurait pu ajouter : que tu viens d'engouler \italique{en cachette} ; car ce \grec{σησαμοῦς}\marqueur{1} a été avalé en plein milieu des Thesmophories, le jour du \italique{jeûne}, \grec{νηστεία}.
\noindent Athénée (XIV, 647a) nous rapporte que les Syracusaines préparaient des gâteaux semblables de miel et de sésame, les \italique{mylles} pour leurs Thesmophories~:
\citComm{ΜΥΛΛΟΙ. Ἡρακλείδης ὁ Συρακόσιος ἐν τῷ περὶ Θεσμῶν ἐν Συρακούσαις φησὶ τοῖς Παντελείοις τῶν Θεσμοφορίων ἐκ σησάμου καὶ μέλιτος κατασκευάζεσθαι ἐφήβαια γυναικεῖα, ἃ καλεῖσθαι κατὰ πᾶσαν Σικελίαν μυλλοὺς καὶ περιφέρεσθαι ταῖς θεαῖς.\\
\francais{Mylloi. Héraclide de Syracuse rapporte ce qui suit, dans son \italique{Traité des Rites religieux} : « Le jour des grandes Thesmophories, on fait à Syracuse avec de la farine, du miel et du sésame, la figure de la partie sexuelle d'une femme, que l'on porte en pompe, pour l'offrir aux déesses. On l'appelle \italique{mylle} par toute la Sicile. »}}
\end{multicols}
}%
\end{adjustbox}%
\vfill
\end{adjustwidth}
\newpage

% --- FOLIO 14 (f. 169r) : SPREAD 2.4 (Texte grec établi + Traduction française + Commentaire 7/8 & 8/8) ---
\clearpage\ifodd\value{page}\hbox{}\newpage\fi
\enTeteG{ΘΕΣΜΟΦΟΡΙΑΖΟΥΣΑΙ --- Texte grec établi (f. 169r, vv. 551–589 [4/4])}
\enTeteD{LES THESMOPHORIES --- Traduction française (f. 169r, vv. 551–589 [4/4])}
\begin{adjustwidth}{-2cm}{-2cm}
\noindent\textbf{\footnotesize Texte grec établi --- Folio 14 (\texttt{f. 169r}, vv. 551–589 [4/4]) \hfill ΕΠΕΙΣΟΔΙΟΝ Α'}\\[-2pt]
\noindent\rule{16.5cm}{0.5pt}\\[2pt]
{\fontsize{8.8pt}{10.2pt}\selectfont
\ligneEst{}{ὀλίγῳ τι πρότερον κατ᾽ ἀγορὰν λαλούμενον,}{}
\ligneEst{}{ἥκω φράσων τοῦτ᾽ ἀγγελῶν θ᾽ ὑμῖν, ἵνα}{}
\ligneEst{}{σκοπῆτε καὶ τηρῆτε μὴ καὶ προσπέσῃ}{580}
\ligneEst{}{ὑμῖν ἀφάρκτοις πρᾶγμα δεινὸν καὶ μέγα.}{}
\ligneEst{\textbf{Κορυφαῖος}}{τί δ᾽ ἔστιν ὦ παῖ; παῖδα γάρ σ᾽ εἰκὸς καλεῖν,}{}
\ligneEst{}{ἕως ἂν οὕτως τὰς γνάθους ψιλὰς ἔχῃς.}{}
\ligneEst{\textbf{Κλεισθένης}}{Εὐριπίδην φάσ᾽ ἄνδρα κηδεστήν τινα}{}
\ligneEst{}{αὑτοῦ γέροντα δεῦρ᾽ ἀναπέμψαι τήμερον.}{585}
\ligneEst{\textbf{Κορυφαῖος}}{πρὸς ποῖον ἔργον ἢ τίνος γνώμης χάριν;}{}
\ligneEst{\textbf{Κλεισθένης}}{ἵν᾽ ἅττα βουλεύοισθε καὶ μέλλοιτε δρᾶν,}{}
\ligneEst{}{ἐκεῖνος εἴη τῶν λόγων κατάσκοπος.}{}
\ligneEst{\textbf{Κορυφαῖος}}{καὶ πῶς λέληθεν ἐν γυναιξὶν ὢν ἀνήρ;}{589}
}
\vspace{2pt}
\noindent\rule{16.5cm}{0.6pt}\\[2pt]
\noindent\textbf{\small Commentaire philologique, métrique et scénique (vv. 551–589 --- partie 7/8)}\\[2pt]
\begin{adjustbox}{max totalheight=15.2cm, max width=16.5cm, keepaspectratio}%
\parbox[t]{16.5cm}{
\setlength{\columnsep}{9pt}%
\begin{multicols}{2}
\fontsize{7.6pt}{8.8pt}\selectfont
\setlength{\parindent}{0pt}%
\setlength{\parskip}{1.2pt}%
\noindent C'est de peur que Mica risque de relâcher son gâteau. Chez Aristophane la peur se manifeste presque toujours physiquement par des troubles du bas-ventre\marqueur{2}. On dira en français \italique{avoir les pétoches}.
\noindent \latin{Pulchre, bene, recte.} Mais l'essentiel nous échappe encore : Mnésiloque ne promet pas à Mica un petit coup en cas d'agression de sa part, mais une riposte non graduée : une pluie de horions à lui assainir le côlon ! Car il faut comprendre : «~tu vas en chier [de peur] \italique{malgré} le gâteau de sésame par toi engueulé~».
\noindent Ce genre de gâteaux avait la vertu d'être astringent\marqueur{3}. C'est ce qu'il ressort d'un passage des \italique{Nuées} (vers 506-508) : Strepsiade hésite à rentrer dans l'antre de Socrate, qu'il compare aimablement à celui de Trophonius :
\citComm{\ \ \ \ \ \ \ \ \ \ \ \ \ \ \ \ \ \ \ \ \ \ \ \ \ \ \ \ \ \ εἰς τὼ χεῖρέ νυν\\
δός μοι μελιτοῦτταν πρότερον, ὡς δέδοικ᾽ ἐγὼ\\
εἴσω καταβαίνων ὥσπερ εἰς Τροφωνίου.\\
Mets-moi d'abord en main un gâteau au miel ; car j'ai peur de descendre, comme si je pénétrais dans l'antre de Trophonios.}
\noindent Il faut bien-sûr prononcer d'un trait `εἰς Τροφωνίου' pour faire entendre στρόφος, la \italique{colique}\marqueur{4}.\\
\noindent C'est aussi ce qui explique le bon mot de table qu'Athénée (I, 4d) rapporte :
\citComm{ἀμίδα δότω τις ἢ πλακοῦντά τις δότω\\
Qu'on m'apporte un pot de chambre ou bien un gâteau !}
\noindent Au contraire de tous ces gâteaux de miel, la συρμαία (sorte de rave noire que prisaient les Égyptiens pour leur régime ; voir le commentaire \italique{infra} aux vers 855-857) possède des vertus purgatives.
\noindent Cette opposition est vraisemblablement la raison de la confrontation dans la pièce perdue des \italique{Danaïdes}\marqueur{5} entre deux vendeurs : celui de syrméa et celui d'ἐγκρίδες\marqueur{6} :
\citComm{ἐν δ᾿ Ἀριστοφάνης Δαναΐσι συρμαιοπῶλαι καὶ ἐγκριδοπῶλαι.}
\noindent\textbf{[Notes ad 570]} \textbf{(1)} Suidas (π 3192) nous apprend que c'est une sorte de gâteau fait à partir de miel bouilli et de sésame : Πυραμοῦς: εἶδος πλακοῦντος ἐκ μέλιτος ἑφθοῦ καὶ πυροῦ πεφρυγμένου, ὡς σησαμοῦς ὁ διὰ σησαμῆς. ταῦτα δὲ ἐτίθεσαν ἆθλα τοῖς διαγρυπνηταῖς ἐν τοῖς συμποσίοις, καὶ ὁ διαγρυπνήσας μέχρι τῆς ἔω ταῦτα ἐλάμβανεν ἔπαθλον. ἡμέτερος ὁ πυραμοῦς.
\noindent \textbf{(2)} On en ramasse à la pelle de ces exemples de déjections consécutives à une frayeur ; Taillardat (§309) en cite bon nombre : \italique{Assemblée des femmes}, v. 1060 ; \italique{Acharniens}, v. 350 ; et d'autres, mais la liste n'est toujours pas exhaustive : on pourrait encore la compléter avec les \italique{Guêpes}, v. 627 ; \italique{Oiseaux}, v. 66 ; etc. etc. Signalons la même promesse de relâchement des intestins après correction dans \italique{Lysistrata}, v. 440 : \grec{ἐπιχεσεῖ πατούμενος}.
\end{multicols}
}%
\end{adjustbox}%
\vfill
\end{adjustwidth}
\newpage
% --- PAGE DE DROITE : Traduction française + Commentaire (8/8) ---
\begin{adjustwidth}{-2cm}{-2cm}
\noindent\textbf{\footnotesize Traduction française --- Folio 14 (\texttt{f. 169r}, vv. 551–589 [4/4]) \hfill ΕΠΕΙΣΟΔΙΟΝ Α'}\\[-2pt]
\noindent\rule{16.5cm}{0.5pt}\\[2pt]
{\fontsize{8.8pt}{10.2pt}\selectfont
\ligneEst{}{\francais{dont on jasait tout à l'heure sur l'agora,}}{}
\ligneEst{}{\francais{je suis venu la raconter et vous donner avis}}{}
\ligneEst{}{\francais{de veiller et d'être en garde, de crainte}}{580}
\ligneEst{}{\francais{qu'une grave et fâcheuse affaire ne vous tombe sur les bras au dépourvu.}}{}
\ligneEst{\textbf{Coryphée}}{\francais{Qu'y a-t-il mon enfant ? Car il est naturel que je t'appelle enfant,}}{}
\ligneEst{}{\francais{tant que tu auras ainsi les joues lisses.}}{}
\ligneEst{\textbf{Clisthène}}{\francais{On dit qu'Euripide a envoyé ici même, aujourd'hui,}}{}
\ligneEst{}{\francais{un de ses parents, un vieillard.}}{585}
\ligneEst{\textbf{Coryphée}}{\francais{Pourquoi faire, dans quel dessein ?}}{}
\ligneEst{\textbf{Clisthène}}{\francais{À cette fin que, quoi que vous agitiez ou décidiez de faire,}}{}
\ligneEst{}{\francais{il puisse épier vos discours.}}{}
\ligneEst{\textbf{Coryphée}}{\francais{Et comment s'est-il dissimulé au milieu de femmes ? Un homme !}}{589}
}
\vspace{2pt}
\noindent\rule{16.5cm}{0.6pt}\\[2pt]
\noindent\textbf{\small Commentaire philologique, métrique et scénique (vv. 551–589 --- partie 8/8)}\\[2pt]
\begin{adjustbox}{max totalheight=15.2cm, max width=16.5cm, keepaspectratio}%
\parbox[t]{16.5cm}{
\setlength{\columnsep}{9pt}%
\begin{multicols}{2}
\fontsize{7.6pt}{8.8pt}\selectfont
\setlength{\parindent}{0pt}%
\setlength{\parskip}{1.2pt}%
\noindent \textbf{(3)} Hypothèse, certes corroborée par les textes d'Aristophane que nous allons voir mais point testée sur mon propre appareil digestif. Est-ce digne d'un disciple de Willems, ce grand amateur de \italique{Realia} ?
\noindent \textbf{(4)} En passant, Euripide use du même jeu de mots scatologique : au vers 300 de son \italique{Ion} : σὺν ἀνδρί $[$Créuse dit à Ion qu'elle est venue avec son mari, mais... ] σηκοὺς δ᾽ ἐκστρέφει Τροφωνίου... que le mari, Xuthus, était d'ailleurs véritablement allé consulter. Le tour était proverbial et Érasme l'a recensé (\italique{Adages} n° 677)
\noindent \textbf{(5)} Pollux (VII 199 $=$ fr. 276(265) Kassel-Austin). Cf. aussi le passage -- corrompu -- d'Athénée (XIV, 645e) où il est question de l'ἐγκριδοπώλης dans les \italique{Danaïdes}.
\noindent \textbf{(6)} Athénée (XIV, 645e) : ΕΓΚΡΙΔΕΣ. πεμμάτιον ἑψόμενον ἐν ἐλαίῳ καὶ μετὰ τοῦτο μελιτούμενον.
\noindent\textbf{575-6} \textbf{} Sur la conservation de ces deux vers que d'aucuns auraient voulu biffer, voir la note de Willems n° 17.
\noindent\textbf{578} \textbf{ὀλίγῳ τι πρότερον κατ᾽ ἀγορὰν λαλούμενον.} Clisthène rapporte les cancans de l'Agora, qui font grand cas de ce beau-père d'Euripide introduit chez les femmes pou plaider sa cause. Ce n'est certainement pas Euripide qui est à l'origine de la fuite et ce pour des raisons évidentes : \ \ \ 1° Si Aristophane, l'avait imaginé ainsi, il nous l'aurait fait comprendre directement ou par quelque allusion \ \ \ 2° Euripide n'a aucun intérêt à voir son complice être démasqué, qui ne pourra donc plus plaider sa cause \ \ \ 3° Euripide a beaucoup de défauts, mais pas celui d'être ingrat au point d'envoyer à une mort certaine son dévoué parent. \ \ \ Savoir quel domestique d'Agathon, quel voisin, ayant aperçu Mnésiloque attifé en femme, sera allé colporter la nouvelle sur l'Agora... voilà bien des détails sans importance sur lesquels Aristophane n'a pas daigné s'arrêter ; nous aurions fait de même si les professeurs parmi les plus admis ne s'étaient lancés dans ce genre d'extravagantes et vaines conjectures.
\noindent\textbf{579-87} \textbf{ἥκω φράσων...} Passées les retrouvailles, Clisthène prend un ton paratragique -- l'heure est grave ! Noter à ce propos les redoublements\marqueur{1} : φράσων/ἀγγελῶν, σκοπῆτε/τηρῆτε, δεινὸν/μέγα, βουλεύοισθε/μέλλοιτε.
\noindent\textbf{[Notes ad 579-87]} \textbf{(1)} Redoublements propres au style eschyléen et raillés par l'Euripide des \italique{Grenouilles} (vv. 1153 sqq.).
\noindent\textbf{582-3} \textbf{τί δ᾽ ἔστιν ὦ παῖ; παῖδα γάρ σ᾽ εἰκὸς καλεῖν, /	ἕως ἂν οὕτως τὰς γνάθους ψιλὰς ἔχῃς.} Pique de la coryphée analogue à celle du chœur des \italique{Guêpes}, aux vers 1297-1298 :
\citComm{Χορός – τί δ᾽ ἔστιν ὦ παῖ; παῖδα γάρ, κἂν ᾖ γέρων,\\
καλεῖν δίκαιον ὅστις ἂν πληγὰς λάβῃ.}
\noindent\textbf{585} \textbf{ἀναπέμψαι.} Le suffixe `ἀνα' s'explique par le fait que le Thesmophorion se trouvait probablement sur la colline de la Pnyx. (Cf. le vers 658.)
\end{multicols}
}%
\end{adjustbox}%
\vfill
\end{adjustwidth}
\newpage
\clearpage\ifodd\value{page}\hbox{}\newpage\fi
\enTeteG{ΕΠΕΙΣΟΔΙΟΝ Β' --- Codex Ravennas 429, f. 169v (Classense 171v) --- vv. 590–634}
\begin{adjustwidth}{-1.5cm}{-3.0cm}
\begin{adjustbox}{max totalheight=21.0cm, max width=16.5cm, keepaspectratio}
\begin{minipage}{16.5cm}
{\fontsize{7.2pt}{8.5pt}\selectfont
\noindent\textbf{Folio 15 --- Codex Ravennas 429, \texttt{f. 169v} (\texttt{Classense 171v}) : ΕΠΕΙΣΟΔΙΟΝ Β', vv. 590–634}\hfill \textit{Témoins :} $\text{R}^1, \Rpc, \Rdeux$ (\Mon)\par\vspace{2pt}
\hrule height 0.4pt\vspace{2.5pt}
}
{\fontsize{8.1pt}{9.5pt}\selectfont
\ligneRav{590}{κλεισθ/}{ἄφευσεν αὐτὸν \difference{κἀπέτιλ᾽} Εὐριπίδης}{\textbf{590} \lecon{ἄφευσεν} R : ἀφηῦσεν Meineke \espApp καπέτιλλ᾽ R : \lecon{κἀπέτιλ᾽} Bekker}
\ligneRav{}{}{καὶ τἄλλ᾽ ἅπανθ᾽ ὥσπερ γυναῖκ᾽ ἐσκεύασεν·}{}
\ligneRav{}{μ\auDessus{ν}{Η}}{πείθεσθε τούτῳ ταῦτα; τίς δ᾽ οὕτως ἀνὴρ}{}
\ligneRav{}{}{ἠλίθιος ὅστις τιλλόμενος ἠνείχετο;}{}
\ligneRav{594}{}{οὐκ \difference{οἴομαι 'γωγ᾽} ὦ πολυτιμήτω θεώ·}{\textbf{594} οἴομ᾽ ἔγωγ᾽ R : \lecon{οἴομαι 'γωγ᾽} anonymus Parisinus}
\ligneRav{595}{κλεισθ/}{ληρεῖς· ἐγὼ γὰρ οὐκ ἂν ἦλθον ἀγγελῶν,}{}
\ligneRav{596}{}{εἰ μὴ \difference{᾽πεπύσμην} \difference{ταῦτα} τῶν σάφ᾽ εἰδότων·}{\textbf{596} πεπύσμην R, \PsiUn : \lecon{᾽πεπύσμην} Faber \espApp ταυτὶ R : \lecon{ταῦτα} Bentley}
\ligneRav{}{χορ/}{τὸ πρᾶγμα τουτὶ δεινὸν εἰσαγγέλλεται·}{}
\ligneRav{598}{}{ἀλλ᾽ ὦ γυναῖκες οὐκ \difference{ἐλινύειν} ἐχρῆν,}{\textbf{598} ἐλιννύειν R : \lecon{ἐλινύειν} Bourdin}
\ligneRav{}{}{ἀλλὰ σκοπεῖν τὸν ἄνδρα καὶ ζητεῖν ὅπου}{}
\ligneRav{600}{}{λέληθεν \difference{ἡμᾶς} κρυπτὸς ἐγκαθήμενος·}{\textbf{600} ὑμᾶς R : \lecon{ἡμᾶς} \Mon}
\ligneRav{601}{}{καὶ σὺ \difference{ξυνέξευρ᾽} αὐτόν, ὡς ἂν τὴν χάριν}{\textbf{601} συνέξευρ᾽ R, Suid. π 2540 codd. AG : \lecon{ξυνέξευρ᾽} Suid. codd. FVM}
\ligneRav{602}{}{ταύτην τε κἀκείνην ἔχῃς ὦ πρόξενε·}{\textbf{602} \lecon{ταύτην} R : αὐτὴν Suid.}
\ligneRav{}{κλεισθ/}{φέρ᾽ ἴδω· τίς ἡ πρώτη σύ; : ποῖ τις τρέψεται;}{}
\ligneRav{604}{κλεισθ/}{ζητητέαι γάρ ἐστε. : κακοδαίμων ἐγώ·}{\textbf{604} \lecon{ζητητέαι γάρ ἐστε} R : ζητητέον τἄρ᾽ ἐστί Maas \espApp 604 ante 603 transp. Maas, post 607 collocare malui}
\ligneRav{605}{\auDessus{γυ}{ν}}{ἔμ᾽ \difference{ἥτις} \difference{εἴμ᾽ ἤρου}; Κλεωνύμου γυνή·}{\textbf{605} εἰτις ἤρου; R : \lecon{ἥτις} (Zanetti) \lecon{εἴμ᾽ ἤρου;} Bentley}
\ligneRav{606}{κλεισθ/}{γιγνώσκεθ᾽ ὑμεῖς ἥτις ἔσθ᾽ \difference{ἡδὶ} γυνή;}{\textbf{606} ἥδε R : ἥδ᾽ ἡ \Mon : \lecon{ἡδὶ} Grynaeus}
\ligneRav{}{χορ/}{γιγνώσκομεν δῆτ᾽. ἀλλὰ τὰς ἄλλας ἄθρει·}{}
\ligneRav{}{κλεισθ/}{ἡδὶ δὲ δὴ τίς ἐστιν ἡ τὸ παιδίον}{}
\ligneRav{609}{—}{ἔχουσα; : \difference{τίτθη νὴ Δί᾽} ἐμή. : διοίχομαι·}{\textbf{609} τίτθην η δί᾽ R : \lecon{τίτθη νὴ Δί᾽} Ellebodius}
\ligneRav{610}{κλεισθ/}{αὕτη σὺ ποῖ στρέφει; μέν᾽ αὐτοῦ. τί τὸ κακόν;}{\textbf{610} τί τὸ κακόν; continuat Clistheni R : Parenti tribuit (cum versu sequenti) \Rpc}
\ligneRav{611}{μ\auDessus{ν}{Η}}{ἔασον οὐρῆσαί μ᾽· ἀναίσχυντός \difference{τις εἶ}·}{\textbf{611} εἶ R : \lecon{τις εἶ} Küster}
\ligneRav{612}{κλεισθ/}{σὺ δ᾽ οὖν ποίει τοῦτ᾽. \difference{ἀναμενῶ} γὰρ ἐνθάδε·}{\textbf{612} ἀναμένω R : \lecon{ἀναμενῶ} Grynaeus}
\ligneRav{}{χορ/}{ἀνάμενε δῆτα καὶ σκόπει γ᾽ αὐτὴν σφόδρα·}{}
\ligneRav{}{}{μόνην γὰρ αὐτὴν ὦνερ οὐ γιγνώσκομεν·}{}
\ligneRav{615}{κλεισθ/}{\difference{πολύν} γε χρόνον οὐρεῖς σύ. : νὴ Δί᾽ ὦ μέλε·}{\textbf{615} πολύ R : \lecon{πολύν} Gelenius}
\ligneRav{}{}{στραγγουριῶ γάρ· ἐχθὲς ἔφαγον κάρδαμα·}{}
\ligneRav{}{κλεισθ/}{τί καρδαμίζεις; οὐ βαδιεῖ δεῦρ᾽ ὡς ἐμέ;}{}
\ligneRav{}{μ\auDessus{ν}{Η}}{τί δῆτά μ᾽ ἕλκεις ἀσθενοῦσαν; : εἰπέ μοι,}{}
\ligneRav{619}{—}{τίς \difference{ἔστ᾽} ἀνήρ σοι; : τὸν ἐμὸν ἄνδρα πυνθάνει;}{\textbf{619} ἔσθ᾽ R : \lecon{ἔστ᾽} \Mon}
\ligneRav{620}{}{τὸν δεῖνα γιγνώσκεις, τὸν ἐκ Κοθωκιδῶν·}{\textbf{620} add. signum interrogationis Canini}
\ligneRav{621}{κλεισθ/}{τὸν δεῖνα; ποῖον; : ἔσθ᾽ ὁ δεῖν᾽, ὃς καί ποτε}{\textbf{621} (Κλεισθ.) continuat R : \lecon{(Μνη.)} \Mon}
\ligneRav{}{—}{τὸν δεῖνα τὸν τοῦ δεῖνα— : ληρεῖν μοι δοκεῖς·}{}
\ligneRav{}{—}{ἀνῆλθες ἤδη δεῦρο πρότερον; : νὴ Δία}{}
\ligneRav{624}{—}{\difference{ὁσέτη} γε. : καὶ τίς \difference{σοὐστὶ} συσκηνήτρια;}{\textbf{624} ὅς᾽ ἔτη R : \lecon{ὁσέτη} G. H. Schaefer ad Gregor. Cor. p.65 \espApp σοὐστὶν R : \lecon{σοὐστὶ} Dindorf (σοῦ ᾽στὶ iam Biset)}
\ligneRav{625}{μ\auDessus{ν}{Η}}{ἡ δεῖν᾽ \difference{ἔμοιγ᾽}. : οἴμοι τάλας. οὐδὲν λέγεις·}{\textbf{625} ἔμοιγε R : \lecon{ἔμοιγ᾽} edd. pr.}
\ligneRav{}{\auDessus{γυ}{ν}}{ἄπελθ᾽. ἐγὼ γὰρ βασανιῶ ταύτην καλῶς}{}
\ligneRav{}{}{ἐκ τῶν ἱερῶν τῶν πέρυσι· σὺ δ᾽ ἀπόστηθί μοι,}{}
\ligneRav{}{}{ἵνα μὴ 'πακούσῃς ὢν ἀνήρ. σὺ δ᾽ εἰπέ μοι}{}
\ligneRav{}{}{ὅ τι πρῶτον ἡμῖν τῶν ἱερῶν ἐδείκνυτο·}{}
\ligneRav{630}{μ\auDessus{ν}{Η}}{τί δ᾽ \expo{μέντοι} πρῶτον ἦν; : ἐπίνομεν·}{\textbf{630} \lecon{μέντοι} \Rpc in mg., Suid. π 2563 : om. \Rac, \Mon}
\ligneRav{631}{\auDessus{γυ}{ν}}{τί δὲ \difference{μετὰ} τοῦτο δεύτερον; : προὐπίνομεν·}{\textbf{631} \lecon{δὲ} R, Suid. π 2563 : δαὶ Bentley \espApp με R : \lecon{μετὰ} Suid.}
\ligneRav{632}{\auDessus{γυ}{ν}}{ταυτὶ μὲν ἤκουσάς τινος· \difference{τί δ᾽ ἦν τρίτον};}{\textbf{632} τί δὲ τρίτον R : τί δ᾽ αὖ Bentley : τί δαὶ τρίτον Elmsley : \lecon{τί δ᾽ ἦν τρίτον} Teubner}
\ligneRav{}{μ\auDessus{ν}{Η}}{σκάφιον Ξένυλλ᾽ ᾔτησεν· οὐ γὰρ ἦν ἀμίς·}{}
\ligneRav{634}{\auDessus{γυ}{ν}}{ξυνέπινε \difference{δεῦρ᾽ ὦ} (δευρ ὦ \Rdeux) Ξενύλλ᾽· οὐκ ἦν ἀμίς·}{\textbf{634} δεῦρο \Rac, \Mon : \lecon{δευρ ὦ} \Rdeux \espApp Κλεισόσθενες R : \lecon{Κλείσθενες} Bonini}
}
\vspace{3pt}\noindent\rule{4.5cm}{0.4pt}\par\vspace{1.5pt}
{\fontsize{7.0pt}{8.3pt}\selectfont
\noindent\textbf{Loci paralleli (vv. 590–634) ---} \textbf{598} (ἐλινύειν) Hesychius ε 2095 ; cf. Moeris ε 61 (ἐλινύων $=$ ἀναπαυόμενος) $\sim$ Hesychius ε 2097 \espTest \textbf{601-2} Suidas π 2540 \espTest \textbf{616} Suidas σ 1163 (χθὲς γὰρ pro γάρ ἐχθὲς) \espTest \textbf{620} (Κοθωκιδῶν) Photius κ 859 Κοθωκίδαι (ap Scholium) $\sim$ Hesychius κ 3218 \espTest \textbf{630-1} Suidas π 2563 (sed προπίνομεν pro προὐ- cod. V) \espTest \textbf{633} Pollux X 45 ἐν Πολυίδῳ (sic) Ἀριετοφάνης· σκάφιον Ξένυλλ᾽ ἤτησεν· οὐ γὰρ ἦν ἁμίς \espTest Photius σ 286 σκάφιον\par
}
\vspace{2.5pt}\noindent\rule{4.5cm}{0.4pt}\par\vspace{1.5pt}
{\fontsize{6.9pt}{8.2pt}\selectfont
\noindent\textbf{Suite (en commentaire) de l'apparat critique \& notes paléographiques (\texttt{f. 169v}) ---} \textbf{Témoin papyrologique :} \textbf{\PsiUn{} (\italique{PSI} XI 1194)} couvre les \textbf{vv. 594–596}. --- \textbf{Supplément marginal au v. 630 :} La particule \lecon{μέντοι} (confirmée par la Souda $\pi\ 2563$) avait été omise dans la colonne par \Rac{} (et reste \textbf{omise dans \Mon{}}) ; elle a été rétablie dans la marge par \Rpc{}. --- \textbf{Dittographie syllabique au v. 634 :} En bas de page, $\text{R}^1$ écrit Κλεισόσθενες (dittographie de $-\sigma\text{ο}-$ dans Κλείσθενες, corrigée par Bonini), tandis que \Rdeux{} corrige δεῦρο (\Rac{}, \Mon{}) en \lecon{δευρ ὦ}. --- \textbf{v. 609 :} τίτθην η δί' R (médivision de $\nu$ final) : \lecon{τίτθη νὴ Δί'} Ellebodius. --- \textbf{v. 630 :} \lecon{μέντοι} \Rpc{} in mg., Souda $\pi\ 2563$ : om. $\text{R}^{\text{ac}}, \text{M}$.\par
}
\end{minipage}
\end{adjustbox}
\end{adjustwidth}
\newpage
\enTeteD{Codex Ravennas 429 (Biblioteca Classense) --- Fac-similé du f. 169v (vv. 590–634)}
\begin{adjustwidth}{-1.5cm}{-3.0cm}
\begin{center}
\includegraphics[width=16.5cm,height=20.5cm,keepaspectratio]{/workspaces/Thesmophories/dossier_ravennas_thesmophories/folios_individuels_ff162v-177r/15_f169v_classense171v_vv0590-0634.jpg}\par
\vspace{2pt}
{\scriptsize\textbf{Codex Ravennas 429, f. 169v (Classense 171v)} --- ΕΠΕΙΣΟΔΙΟΝ Β', \textbf{vv. 590–634} ($\text{R}^1, \Rpc, \Rdeux$ (\Mon)).}
\end{center}
\end{adjustwidth}
\newpage
\clearpage\ifodd\value{page}\hbox{}\newpage\fi% --- FOLIO 15 (f. 169v) : SPREAD 2.1 (Texte grec établi + Traduction française + Commentaire 1/6 & 2/6) ---
\clearpage\ifodd\value{page}\hbox{}\newpage\fi
\enTeteG{ΘΕΣΜΟΦΟΡΙΑΖΟΥΣΑΙ --- Texte grec établi (f. 169v, vv. 590–634 [1/3])}
\enTeteD{LES THESMOPHORIES --- Traduction française (f. 169v, vv. 590–634 [1/3])}
\begin{adjustwidth}{-2cm}{-2cm}
\noindent\textbf{\footnotesize Texte grec établi --- Folio 15 (\texttt{f. 169v}, vv. 590–634 [1/3]) \hfill ΕΠΕΙΣΟΔΙΟΝ Β'}\\[-2pt]
\noindent\rule{16.5cm}{0.5pt}\\[2pt]
{\fontsize{8.4pt}{9.7pt}\selectfont
\ligneEst{\textbf{Κλεισθένης}}{ἄφευσεν αὐτὸν κἀπέτιλ᾽ Εὐριπίδης}{590}
\ligneEst{}{καὶ τἄλλ᾽ ἅπανθ᾽ ὥσπερ γυναῖκ᾽ ἐσκεύασεν.}{}
\ligneEst{\textbf{Μνησίλοχος}}{πείθεσθε τούτῳ ταῦτα; τίς δ᾽ οὕτως ἀνὴρ}{}
\ligneEst{}{ἠλίθιος ὅστις τιλλόμενος ἠνείχετο;}{}
\ligneEst{}{οὐκ οἴομαι 'γωγ᾽ ὦ πολυτιμήτω θεώ.}{}
\ligneEst{\textbf{Κλεισθένης}}{ληρεῖς· ἐγὼ γὰρ οὐκ ἂν ἦλθον ἀγγελῶν,}{595}
\ligneEst{}{εἰ μὴ ᾽πεπύσμην ταῦτα τῶν σάφ᾽ εἰδότων.}{}
\ligneEst{\textbf{Κορυφαῖος}}{τὸ πρᾶγμα τουτὶ δεινὸν εἰσαγγέλλεται.}{}
\ligneEst{}{ἀλλ᾽ ὦ γυναῖκες οὐκ ἐλινύειν ἐχρῆν,}{}
\ligneEst{}{ἀλλὰ σκοπεῖν τὸν ἄνδρα καὶ ζητεῖν ὅπου}{}
\ligneEst{}{λέληθεν ἡμᾶς κρυπτὸς ἐγκαθήμενος.}{600}
\ligneEst{}{καὶ σὺ ξυνέξευρ᾽ αὐτόν, ὡς ἂν τὴν χάριν}{}
\ligneEst{}{ταύτην τε κἀκείνην ἔχῃς ὦ πρόξενε.}{}
\ligneEst{\textbf{Κλεισθένης}}{φέρ᾽ ἴδω· τίς ἡ πρώτη σύ;}{}
\ligneEst{\textbf{Μνησίλοχος}}{\decalage{φέρ᾽ ἴδω· τίς ἡ πρώτη σύ;}ποῖ τις τρέψεται;}{}
\ligneEst{\textbf{Κλεισθένης}}{ζητητέαι γάρ ἐστε.}{}
\ligneEst{\textbf{Μνησίλοχος}}{\decalage{ζητητέαι γάρ ἐστε.}κακοδαίμων ἐγώ.}{}
\ligneEst{\textbf{Μίκα}}{ἔμ᾽ ἥτις εἴμ᾽ ἤρου; Κλεωνύμου γυνή.}{605}
\ligneEst{\textbf{Κλεισθένης}}{γιγνώσκεθ᾽ ὑμεῖς ἥτις ἔσθ᾽ ἡδὶ γυνή;}{}
}
\vspace{2pt}
\noindent\rule{16.5cm}{0.6pt}\\[2pt]
\noindent\textbf{\small Commentaire philologique, métrique et scénique (vv. 590–634 --- partie 1/6)}\\[2pt]
\begin{adjustbox}{max totalheight=13.0cm, max width=16.5cm, keepaspectratio}%
\parbox[t]{16.5cm}{
\setlength{\columnsep}{9pt}%
\begin{multicols}{2}
\fontsize{8.4pt}{9.7pt}\selectfont
\setlength{\parindent}{0pt}%
\setlength{\parskip}{1.2pt}%
\noindent\textbf{593} \textbf{ἠλίθιος ὅστις τιλλόμενος ἠνείχετο;} Willems, chevalier servant du \italique{Ravennas}, défend encore une fois la leçon du manuscrit dans sa \nieme{18} note, à retrouver à la fin.
\noindent\textbf{602} \textbf{ὦ πρόξενε.} Les proxènes ou hôtes publics, agents officiels dont les fonctions consistaient principalement à recevoir et à héberger les délégués de l'État dont ils exerçaient la proxénie. (Note de Willems.)
\noindent\textbf{605} \textbf{Κλεωνύμου γυνή.} L'évocation de Cléonyme est l'occasion de dresse un bref résumé des railleries à l'encontre de ce personnage parmi les plus plaisantés de l'Ancienne comédie. Cet Athénien nous est autrement connu par trois décrets qu'il proposa à la \grec{βουλή}\marqueur{1} l'an 426/5 ainsi qu'une proposition de fixer à 1000 drachmes la récompense contre des informations touchant à l'affaire de la parodie des Mystères d'Éleusis\marqueur{2} (contemporaine de celle des Hermocopides, l'an 415). Cléonyme est victime d'un tir de barrage d'Aristophane portant à la fois :\\ — au moral : sa lâcheté (il lui était arrivé un jour de jeter son bouclier pour mieux fuir. C'est surtout là-dessus que le poète ne tarit point : \italique{Cavaliers} v. 1372, \italique{Nuées} v. 353, \italique{Guêpes} v. 19, 592, \italique{Paix} v. 446, 672, 1186 (non nommé), 1275, 1295, \italique{Oiseaux} v. 290, 1475, 1481) ; sa duplicité (\italique{Acharniens} v. 88, il est traité de parjure v. 399 des \italique{Nuées}, de flatteur au v. 592 des \italique{Guêpes} (il est appelé Κολακώνυμος), de sycophante au v. 1479 des \italique{Oiseaux}) ; ses manières brusques de grand dadais (\italique{Acharniens} v. 844) ; ses mœurs peu ragoûtantes (accusé d'être un \grec{καταπύγων} au v. 675 des \italique{Nuées}) ;\\ — au physique : une laideur repoussante (v. 822 des \italique{Guêpes}, ) ; son androgynie (Socrate, dans une intention de parler exactement, emploie le féminin \grec{Κλεωνύμη} au v. 680 des \italique{Nuées}) ; une taille excessive (\italique{Acharniens} v. 88, \italique{Oiseaux}, v. 1476) ; sa gloutonnerie (\italique{Cavaliers} v. 1290, \italique{Oiseaux} v. 289).\\ Jouant sur les contrastes, Aristophane nous présente une femme de Cléonyme \italique{petite} de nom (c'est le sens de \grec{Μίκα}), virile et hardie dans les débats.\\
\end{multicols}
}%
\end{adjustbox}%
\vfill
\end{adjustwidth}
\newpage
% --- PAGE DE DROITE : Traduction française + Commentaire (2/6) ---
\begin{adjustwidth}{-2cm}{-2cm}
\noindent\textbf{\footnotesize Traduction française --- Folio 15 (\texttt{f. 169v}, vv. 590–634 [1/3]) \hfill ΕΠΕΙΣΟΔΙΟΝ Β'}\\[-2pt]
\noindent\rule{16.5cm}{0.5pt}\\[2pt]
{\fontsize{8.4pt}{9.7pt}\selectfont
\ligneEst{\textbf{Clisthène}}{\francais{Euripide l'a flambé et épilé ;}}{590}
\ligneEst{}{\francais{et pour tout le reste l'a équipé en femme.}}{}
\ligneEst{\textbf{Mnésiloque}}{\francais{Vous en fiez-vous à lui? Quel homme est assez}}{}
\ligneEst{}{\francais{niais pour s'être laissé épiler.}}{}
\ligneEst{}{\francais{Je n'en crois rien, pour ma part, ô déesses très honorées.}}{}
\ligneEst{\textbf{Clisthène}}{\francais{Tu veux rire. Je ne serais pas venu l'annoncer}}{595}
\ligneEst{}{\francais{si je ne le tenais de ceux qui le savent bien.}}{}
\ligneEst{\textbf{Coryphée}}{\francais{C'est une grave affaire qu'on nous annonce là.}}{}
\ligneEst{}{\francais{Or ça, femmes, il s'agit, non de se morfondre,}}{}
\ligneEst{}{\francais{mais de guetter l'homme, de chercher où}}{}
\ligneEst{}{\francais{il se dérobe à nous et se tient caché.}}{600}
\ligneEst{}{\francais{Toi, aide-nous à le trouver, que nous t'ayons}}{}
\ligneEst{}{\francais{encore cette autre obligation, ô notre patron.}}{}
\ligneEst{\textbf{Clisthène}}{\francais{Voyons, toi d'abord, qui es-tu ?}}{}
\ligneEst{\textbf{Mnésiloque}}{\francais{\espace{4cm}(À part :) Où se tourner ?}}{}
\ligneEst{\textbf{Clisthène}}{\francais{... Car on va vous passer en revue.}}{}
\ligneEst{\textbf{Mnésiloque}}{\francais{\espace{2.5cm}(À part :) Malheureux moi !}}{}
\ligneEst{\textbf{Mika}}{\francais{Tu me demandais qui je suis ? la femme de Cléonyme.}}{605}
\ligneEst{\textbf{Clisthène}}{\francais{Vous savez, vous autres, qui est cette femme ?}}{}
}
\vspace{2pt}
\noindent\rule{16.5cm}{0.6pt}\\[2pt]
\noindent\textbf{\small Commentaire philologique, métrique et scénique (vv. 590–634 --- partie 2/6)}\\[2pt]
\begin{adjustbox}{max totalheight=13.0cm, max width=16.5cm, keepaspectratio}%
\parbox[t]{16.5cm}{
\setlength{\columnsep}{9pt}%
\begin{multicols}{2}
\fontsize{8.4pt}{9.7pt}\selectfont
\setlength{\parindent}{0pt}%
\setlength{\parskip}{1.2pt}%
\noindent Revenons sur le grief le plus reproché à Cléonyme : avoir jeté son bouclier pour mieux prendre ses jambes à son cou, sa \grec{ῥίψασπις}. Il n'y a rien à y redire : le crime était théoriquement puni de mort\marqueur{3} et l'on comprend aisément que le poète ait exposé à l'opprobre publique le couard qui avait encore l'outrecuidance de venir palabrer après son forfait. Exégèse par trop simple pour nos éminents professeurs – dont le délassement consiste à extravaguer sur les passages les plus obscurs et remettre systématiquement en doute les passages les plus clairs\marqueur{4}. On ne la fait pas à ces \italique{viri docti}, comme disait Fritzsche : le dernier en date\marqueur{5} argue de ce que Cléonyme serait mort – condamné à mort – bien avant les dernière traces archéologiques de sa vie démocratique s'il avait effectivement commis ce qu'Aristophane lui reproche... avec le même raisonnement, les historiens des prochains millénaires déduiront que l'\italique{Élysée}, \italique{Matignon} et le \italique{palais Bourbon} ne pouvaient désigner que des prisons au vu des canailles hors-la-loi qui les fréquentaient ; prisons flanquées, semblera-t-il, d'une maison de santé pour ceux dont la panse menacait de crever à tout instant : le \italique{palais du Luxembourg}\marqueur{6}. Mais il n'y a pas besoin d'extrapoler : le même raisonnement intenable entraînait déjà de refuser le vers 1166 des \italique{Acharniens}\marqueur{7} : comment un bandit (un \italique{Oreste}, nom générique ; un \italique{Cartouche} ou \italique{Mandrin} chez nous) aurait-il pu sévir la nuit à Athènes alors que le banditisme était aussi puni de mort ?..
\noindent Certes, « \anglais{in interpreting the personal jokes in Aristophanes one must always be wary of the too serious or too literal explanation of the text} »\marqueur{8}, mais l'\italique{on ne détruit que ce que l'on remplace}, et pour remplacer une explication qui coule de source comme celle de nos passages, il vaut mieux présenter de solides arguments. Or M. Storey ne me fera pas accroire sans preuve\marqueur{9} que le motif pour lequel Aristophane poursuivit inlassablement de sa verve pendant des années Cléonyme fût, non pas une désertion devant l'ennemi doublée d'un abandon d'armes, mais une banale brigue pour échapper au service militaire\marqueur{10}.
\noindent\textbf{[Notes ad 605]} \textbf{(1)} \italique{Inscriptiones Graecae}, I\expo{3}, 61.34, 68.5, 69.3-4.
\end{multicols}
}%
\end{adjustbox}%
\vfill
\end{adjustwidth}
\newpage

% --- FOLIO 15 (f. 169v) : SPREAD 2.2 (Texte grec établi + Traduction française + Commentaire 3/6 & 4/6) ---
\clearpage\ifodd\value{page}\hbox{}\newpage\fi
\enTeteG{ΘΕΣΜΟΦΟΡΙΑΖΟΥΣΑΙ --- Texte grec établi (f. 169v, vv. 590–634 [2/3])}
\enTeteD{LES THESMOPHORIES --- Traduction française (f. 169v, vv. 590–634 [2/3])}
\begin{adjustwidth}{-2cm}{-2cm}
\noindent\textbf{\footnotesize Texte grec établi --- Folio 15 (\texttt{f. 169v}, vv. 590–634 [2/3]) \hfill ΕΠΕΙΣΟΔΙΟΝ Β'}\\[-2pt]
\noindent\rule{16.5cm}{0.5pt}\\[2pt]
{\fontsize{8.4pt}{9.7pt}\selectfont
\ligneEst{\textbf{Κορυφαῖος}}{γιγνώσκομεν δῆτ᾽. ἀλλὰ τὰς ἄλλας ἄθρει.}{}
\ligneEst{\textbf{Κλεισθένης}}{ἡδὶ δὲ δὴ τίς ἐστιν ἡ τὸ παιδίον}{}
\ligneEst{}{ἔχουσα;}{}
\ligneEst{\textbf{Μίκα}}{\decalage{ἔχουσα;}τίτθη νὴ Δί᾽ ἐμή.}{}
\ligneEst{\textbf{Μνησίλοχος}}{\decalage{ἔχουσα;τίτθη νὴ Δί᾽ ἐμή.}διοίχομαι.}{}
\ligneEst{\textbf{Κλεισθένης}}{αὕτη σὺ ποῖ στρέφει; μέν᾽ αὐτοῦ. τί τὸ κακόν;}{610}
\ligneEst{\textbf{Μνησίλοχος}}{ἔασον οὐρῆσαί μ᾽· ἀναίσχυντός τις εἶ.}{}
\ligneEst{\textbf{Κλεισθένης}}{σὺ δ᾽ οὖν ποίει τοῦτ᾽. ἀναμενῶ γὰρ ἐνθάδε.}{}
\ligneEst{\textbf{Κορυφαῖος}}{ἀνάμενε δῆτα καὶ σκόπει γ᾽ αὐτὴν σφόδρα·}{}
\ligneEst{}{μόνην γὰρ αὐτὴν ὦνερ οὐ γιγνώσκομεν.}{}
\ligneEst{\textbf{Κλεισθένης}}{πολύν γε χρόνον οὐρεῖς σύ.}{615}
\ligneEst{\textbf{Μνησίλοχος}}{\decalage{πολύν γε χρόνον οὐρεῖς σύ.}νὴ Δί᾽ ὦ μέλε·}{}
\ligneEst{}{στραγγουριῶ γάρ· ἐχθὲς ἔφαγον κάρδαμα.}{}
\ligneEst{\textbf{Κλεισθένης}}{τί καρδαμίζεις; οὐ βαδιεῖ δεῦρ᾽ ὡς ἐμέ;}{}
\ligneEst{\textbf{Μνησίλοχος}}{τί δῆτά μ᾽ ἕλκεις ἀσθενοῦσαν;}{}
\ligneEst{\textbf{Κλεισθένης}}{\decalage{τί δῆτά μ᾽ ἕλκεις ἀσθενοῦσαν;}εἰπέ μοι,}{}
\ligneEst{}{τίς ἔστ᾽ ἀνήρ σοι;}{}
\ligneEst{\textbf{Μνησίλοχος}}{\decalage{τίς ἔστ᾽ ἀνήρ σοι;}τὸν ἐμὸν ἄνδρα πυνθάνει;}{}
\ligneEst{}{τὸν δεῖνα γιγνώσκεις, τὸν ἐκ Κοθωκιδῶν.}{620}
}
\vspace{2pt}
\noindent\rule{16.5cm}{0.6pt}\\[2pt]
\noindent\textbf{\small Commentaire philologique, métrique et scénique (vv. 590–634 --- partie 3/6)}\\[2pt]
\begin{adjustbox}{max totalheight=13.0cm, max width=16.5cm, keepaspectratio}%
\parbox[t]{16.5cm}{
\setlength{\columnsep}{9pt}%
\begin{multicols}{2}
\fontsize{8.4pt}{9.7pt}\selectfont
\setlength{\parindent}{0pt}%
\setlength{\parskip}{1.2pt}%
\noindent \textbf{(2)} Andocide, I-\italique{\grec{Περὶ τῶν μυστηρίων}}, 27. La démagogie de Cléonyme se voit elle-même débordée par celle de Pisandre qui renchérit à 10000 drachmes !..
\noindent \textbf{(3)} Andocide, I-\italique{\grec{Περὶ τῶν μυστηρίων}}, 74.
\noindent \textbf{(4)} Willems, tome I, p. 108, s'en plaignait déjà en son temps.
\noindent \textbf{(5)} I. C. Storey, «~The `Blameless Shield' of Kleonymus~», in \italique{Rheinisches Museum für Philologie}, n°132 (1989), p. 247-261.
\noindent \textbf{(6)} ``Palais'' est à prendre dans son acception judiciaire ; suite logique des \italique{palais} de justice.
\noindent \textbf{(7)} voir Willems, t. I p. 106 sq.
\noindent \textbf{(8)} «~The `Blameless Shield' of Kleonymus~», p. 248.
\noindent \textbf{(9)} Car je tiens pour nulle et non avenue sa citation du passage des \italique{Cavaliers} (v. 1369-1372) qui \italique{prouve} simplement que M. Storey n'a pas saisi l'image et le jeu de mots d'Aristophane sur \grec{κατὰ σπουδὰς} (mis à dessein pour faire entendre κατ᾿ ἀσπίδας en lieu et place de l'expression figée `κατὰ σπουδήν' qu'on trouve chez Thucydide, Xénophon, Élien, Plutarque, etc.).
\noindent Au reste, si M. Storey convient  que les vers 827-829 des \italique{Thesmophories} visent Cléonyme, je lui saurais gré d'expliquer le tour `ἀπὸ τῶν ὤμων' à l'aune de sa brillante exégèse.
\noindent \textbf{(10)} Si banale qu'on en trouve à la pelle parmi nos  modernes brigands. Citons seulement l'exemple le plus caricatural : M. Ciotti-le-bien-nommé, alors qu'il n'était encore que vil attaché-parlementaire fit jouer ses relations pour échapper au service militaire... maintenant que son derrière est confortablement installé au palais bourbeux, bien loin des fangeuses casernes, Chiotti n'a de cesse de haranguer le veau français pour un rétablissement dudit service militaire.
\noindent\textbf{609} \textbf{τίτθη.} Remarquons en passant qu'une nourrice à Athènes n'était pas nécessairement une esclave. Cf. Démosth., LVII, 35. (Note de Willems.)
\end{multicols}
}%
\end{adjustbox}%
\vfill
\end{adjustwidth}
\newpage
% --- PAGE DE DROITE : Traduction française + Commentaire (4/6) ---
\begin{adjustwidth}{-2cm}{-2cm}
\noindent\textbf{\footnotesize Traduction française --- Folio 15 (\texttt{f. 169v}, vv. 590–634 [2/3]) \hfill ΕΠΕΙΣΟΔΙΟΝ Β'}\\[-2pt]
\noindent\rule{16.5cm}{0.5pt}\\[2pt]
{\fontsize{8.4pt}{9.7pt}\selectfont
\ligneEst{\textbf{Coryphée}}{\francais{Nous le savons. Examine les autres.}}{}
\ligneEst{\textbf{Clisthène}}{\francais{Et celle-là, qui est-elle ? Celle qui porte}}{}
\ligneEst{}{\francais{l'enfant ?}}{}
\ligneEst{\textbf{Mika}}{\francais{\espace{1.3cm}C'est ma nourrice, par Zeus.}}{}
\ligneEst{\textbf{Mnésiloque}}{\francais{\espace{3.8cm}Je suis perdu. (Il va pour sortir.)}}{}
\ligneEst{\textbf{Clisthène}}{\francais{Holà; toi, où vas-tu ? Reste ici. (Mnésiloque s'accroupit.) Qu'est-ce qui t'arrive ?}}{610}
\ligneEst{\textbf{Mnésiloque}}{\francais{Laisse-moi lâcher l'eau. Tu es un impudent.}}{}
\ligneEst{\textbf{Clisthène}}{\francais{Eh bien, fais. J'attendrai ici.}}{}
\ligneEst{\textbf{Coryphée}}{\francais{Oui, attends, et surveille-la bien.}}{}
\ligneEst{}{\francais{Elle est la seule que nous ne connaissions pas.}}{}
\ligneEst{\textbf{Clisthène}}{\francais{Tu es bien longue à faire de l'eau.}}{615}
\ligneEst{\textbf{Mnésiloque}}{\francais{\espace{4.2cm}Par Zeus, l'ami,}}{}
\ligneEst{}{\francais{c'est que j'ai une rétention : hier j'ai mangé du cardame.}}{}
\ligneEst{\textbf{Clisthène}}{\francais{Que parles-tu de cardame ? Veux-tu venir ici, près de moi ? (Il la saisit)}}{}
\ligneEst{\textbf{Mnésiloque}}{\francais{Pourquoi me tirer ainsi ? une valétudinaire !}}{}
\ligneEst{\textbf{Clisthène}}{\francais{\espace{4.5cm}Réponds-moi.}}{}
\ligneEst{}{\francais{Quel est ton mari ?}}{}
\ligneEst{\textbf{Mnésiloque}}{\francais{\espace{2.7cm}Mon mari, dis-tu ?}}{}
\ligneEst{}{\francais{Connais-tu... chose, du dème de Cothocides.}}{620}
}
\vspace{2pt}
\noindent\rule{16.5cm}{0.6pt}\\[2pt]
\noindent\textbf{\small Commentaire philologique, métrique et scénique (vv. 590–634 --- partie 4/6)}\\[2pt]
\begin{adjustbox}{max totalheight=13.0cm, max width=16.5cm, keepaspectratio}%
\parbox[t]{16.5cm}{
\setlength{\columnsep}{9pt}%
\begin{multicols}{2}
\fontsize{8.4pt}{9.7pt}\selectfont
\setlength{\parindent}{0pt}%
\setlength{\parskip}{1.2pt}%
\noindent\textbf{616} \textbf{κάρδαμον} Les botanistes s'accordent, je crois, à reconnaître dans le κάρδαμον, non le \italique{cresson}, comme on l'a cru jusqu'en ces derniers temps, mais une autre crucifère, connue des savants sous le nom de \latin{erucaria aleppica}. Suivant Fraas, \italique{Synopsis florae classicae}, les Grecs lui donnent encore aujourd'hui le même nom. Hippocrate dit positivement que le \grec{κάρδαμον} provoque la dysurie (\italique{Du Régime}, II, 54). On cite à ce sujet une singulière anecdote, recueillie par Xénophon. Chez les Perses, où il passait pour indécent de cracher, se moucher ou s'écarter pour faire de l'eau, on ne donnait aux enfants à leur repas du midi que du pain, avec du cardame pour tout assaisonnement. (\italique{Cyrop.}, I, 2, 8, 11 et 16) Mettez cela d'accord avec les traités de thérapeutique où le cresson est donné pour un sialagogue. (Note de Willems. Les dictionnaires  \LSJ , Bailly et Pape désignent le \grec{κάρδαμον} respectivement comme « \anglais{nose-smart}, \latin{Lepidium sativum} »,\italique{cresson alénois} ; « cresson » et « \allemand{eine Art Kresse mit bitterem Kraut} ».
\noindent Pour aller dans le sens de Willems, ajoutons que le cresson, outre ses vertus sialagogiques, est loin d'être dysurique : il a des vertus au contraire diurétiques. Cf. le \italique{Dictionnaire des plantes médicinales et vénéneuses de France} de Paul-Victor Fournier, p. 731. Mais la raison de sa consommation chez les Perses est peut-être à chercher ailleurs : cette plante était réputée donner du courage.)
\noindent\textbf{619} \textbf{τὸν ἐκ Κοθωκιδῶν;} « Du $[$dème$]$ des Cothôcides ». Ce dème faisait partie de la tribu des \grec{Οἰνηίς} au nord-ouest de la ville\marqueur{1}. Y a-t-il une raison particulière pour que Mnésiloque évoque ce dème ? certainement... mais cette raison nous reste obscure et aucune explication satisfaisante n'a jusqu'ici été avancée.
\noindent\textbf{[Notes ad 619]} \textbf{(1)} L'on pourra consulter le travail de John S. Traill, \italique{The Political Organization of Attica}, \italique{Hesperia} supplement XIV (1975). Voir notamment les pages 49, 62 et 68 ainsi que les cartes d'Athènes à la fin.
\noindent\textbf{624} \textbf{συσκηνήτρια.} Il faut croire d'après ce passage et le vers 658 que lors des Thesmophories, des tentes pour deux femmes étaient installées à proximité du Pnyx.
\end{multicols}
}%
\end{adjustbox}%
\vfill
\end{adjustwidth}
\newpage

% --- FOLIO 15 (f. 169v) : SPREAD 2.3 (Texte grec établi + Traduction française + Commentaire 5/6 & 6/6) ---
\clearpage\ifodd\value{page}\hbox{}\newpage\fi
\enTeteG{ΘΕΣΜΟΦΟΡΙΑΖΟΥΣΑΙ --- Texte grec établi (f. 169v, vv. 590–634 [3/3])}
\enTeteD{LES THESMOPHORIES --- Traduction française (f. 169v, vv. 590–634 [3/3])}
\begin{adjustwidth}{-2cm}{-2cm}
\noindent\textbf{\footnotesize Texte grec établi --- Folio 15 (\texttt{f. 169v}, vv. 590–634 [3/3]) \hfill ΕΠΕΙΣΟΔΙΟΝ Β'}\\[-2pt]
\noindent\rule{16.5cm}{0.5pt}\\[2pt]
{\fontsize{8.4pt}{9.7pt}\selectfont
\ligneEst{\textbf{Κλεισθένης}}{τὸν δεῖνα; ποῖον;}{}
\ligneEst{\textbf{Μνησίλοχος}}{\decalage{τὸν δεῖνα; ποῖον;}ἔσθ᾽ ὁ δεῖν᾽, ὃς καί ποτε}{}
\ligneEst{}{τὸν δεῖνα τὸν τοῦ δεῖνα—}{}
\ligneEst{\textbf{Κλεισθένης}}{\decalage{τὸν δεῖνα τὸν τοῦ δεῖνα—}ληρεῖν μοι δοκεῖς.}{}
\ligneEst{}{ἀνῆλθες ἤδη δεῦρο πρότερον;}{}
\ligneEst{\textbf{Μνησίλοχος}}{\decalage{ἀνῆλθες ἤδη δεῦρο πρότερον;}νὴ Δία}{}
\ligneEst{}{ὁσέτη γε.}{}
\ligneEst{\textbf{Κλεισθένης}}{\decalage{ὁσέτη γε.}καὶ τίς σοὐστὶ συσκηνήτρια;}{}
\ligneEst{\textbf{Μνησίλοχος}}{ἡ δεῖν᾽ ἔμοιγ᾽.}{625}
\ligneEst{\textbf{Κλεισθένης}}{\decalage{ἡ δεῖν᾽ ἔμοιγ᾽.}οἴμοι τάλας. οὐδὲν λέγεις.}{}
\ligneEst{\textbf{Μίκα}}{ἄπελθ᾽. ἐγὼ γὰρ βασανιῶ ταύτην καλῶς}{}
\ligneEst{}{ἐκ τῶν ἱερῶν τῶν πέρυσι· σὺ δ᾽ ἀπόστηθί μοι,}{}
\ligneEst{}{ἵνα μὴ 'πακούσῃς ὢν ἀνήρ. σὺ δ᾽ εἰπέ μοι}{}
\ligneEst{}{ὅ τι πρῶτον ἡμῖν τῶν ἱερῶν ἐδείκνυτο.}{}
\ligneEst{\textbf{Μνησίλοχος}}{φέρ᾽ ἴδω, τί μέντοι πρῶτον ἦν; ἐπίνομεν.}{630}
\ligneEst{\textbf{Μίκα}}{τί δὲ μετὰ τοῦτο δεύτερον;}{}
\ligneEst{\textbf{Μνησίλοχος}}{\decalage{τί δὲ μετὰ τοῦτο δεύτερον;}προὐπίνομεν.}{}
\ligneEst{\textbf{Μίκα}}{ταυτὶ μὲν ἤκουσάς τινος· τί δ᾽ ἦν τρίτον;}{}
\ligneEst{\textbf{Μνησίλοχος}}{σκάφιον Ξένυλλ᾽ ᾔτησεν· οὐ γὰρ ἦν ἀμίς.}{}
\ligneEst{\textbf{Μίκα}}{οὐδὲν λέγεις. δεῦρ᾽ ἐλθὲ δεῦρ᾽ ὦ Κλείσθενες·}{634}
}
\vspace{2pt}
\noindent\rule{16.5cm}{0.6pt}\\[2pt]
\noindent\textbf{\small Commentaire philologique, métrique et scénique (vv. 590–634 --- partie 5/6)}\\[2pt]
\begin{adjustbox}{max totalheight=13.0cm, max width=16.5cm, keepaspectratio}%
\parbox[t]{16.5cm}{
\setlength{\columnsep}{9pt}%
\begin{multicols}{2}
\fontsize{8.4pt}{9.7pt}\selectfont
\setlength{\parindent}{0pt}%
\setlength{\parskip}{1.2pt}%
\noindent\textbf{633} \textbf{σκάφιον Ξένυλλ᾽ ᾔτησεν· οὐ γὰρ ἦν ἀμίς.} Acculé, c'est à ce vers que Mnésiloque fait son hara-kiri : il lance une dernière boutade en guise de baroud d'honneur ! L'explication communément admise laisse songeur quant à l'humour des commentateurs d'Aristophane : Mnésiloque emploierait selon eux un terme impropre (l'\italique{amis} serait une sorte de pot de chambre à usage exclusivement masculin en raison d'un col étroit... mais sa forme varie et le terme pourrait très bien aussi désigner ce genre d'ustensile pour les femmes\marqueur{1}). Quoi qu'il en soit, c'est un peu \italique{gros} : tout rustre qu'il soit, Mnésiloque ne peut pas commettre une telle bévue ; je dirais même plus~: dans ce domaine, il est certainement tout aussi savant que l'\italique{Athénienne de la rue}.
\noindent Voici une explication bien meilleure : le \grec{σκάφιον} désignait à l'origine une \italique{coupe} (en forme de \italique{barque} selon l'étymologie). Il n'acquiert le sens quasi-exclusif d'\italique{urinoir} qu'en passant au latin \latin{scaphium}. À l'époque d'Aristophane, les deux sens devaient encore co-exister. La première partie du vers « σκάφιον Ξένυλλ᾽ ᾔτησεν » était donc interprétée par Mica favorablement~: Xénylla a demandé une coupe $[$pour faire une libation$]$... la chute est rude et Mnésiloque dans la seconde partie accuse les femmes d'un usage détourné des coupes consacrées qui passe la mesure : cet individu n'a rien des manières d'une femme : après avoir uriné publiquement, voilà qu'il prononce maintenant des \italique{insanités} (οὐδὲν ὑγιὲς v. 636).\\
\noindent Voici pour le sens propre -- si l'on peut dire. Mais, comme de coutume, Aristophane joue à différents niveaux : la réplique est aussi un clin d'œil au proverbe grec bien connu\marqueur{2} :
\citComm{Ἄμας ἀπῄτουν, οἱ δ' ἀπηρνοῦντο σκάφας\\
Je demandais des bêches, ils me refusent des seaux.}
\end{multicols}
}%
\end{adjustbox}%
\vfill
\end{adjustwidth}
\newpage
% --- PAGE DE DROITE : Traduction française + Commentaire (6/6) ---
\begin{adjustwidth}{-2cm}{-2cm}
\noindent\textbf{\footnotesize Traduction française --- Folio 15 (\texttt{f. 169v}, vv. 590–634 [3/3]) \hfill ΕΠΕΙΣΟΔΙΟΝ Β'}\\[-2pt]
\noindent\rule{16.5cm}{0.5pt}\\[2pt]
{\fontsize{8.4pt}{9.7pt}\selectfont
\ligneEst{\textbf{Clisthène}}{\francais{Chose ? Qui ça ?}}{}
\ligneEst{\textbf{Mnésiloque}}{\francais{\espace{2cm}C'est chose, qui naguère...}}{}
\ligneEst{}{\francais{avec chose, fils de chose...}}{}
\ligneEst{\textbf{Clisthène}}{\francais{\espace{3.7cm}Tu m'as l'air de baliverner.}}{}
\ligneEst{}{\francais{Es-tu déjà venue ici ?}}{}
\ligneEst{\textbf{Mnésiloque}}{\francais{\espace{4.5cm}Oui, par Zeus,}}{}
\ligneEst{}{\francais{tous les ans.}}{}
\ligneEst{\textbf{Clisthène}}{\francais{\espace{1.5cm}Et quelle est ta compagne de tente ?}}{}
\ligneEst{\textbf{Mnésiloque}}{\francais{La mienne, c'est chose...}}{625}
\ligneEst{\textbf{Clisthène}}{\francais{\espace{2.1cm}Misère de moi. Tu ne dis rien.}}{}
\ligneEst{\textbf{Mika}}{\francais{Retire-toi. C'est moi qui vais l'éplucher comme il faut}}{}
\ligneEst{}{\francais{sur les cérémonies de l'an dernier. — Tiens-toi à distance !}}{}
\ligneEst{}{\francais{tu es un homme, tu ne dois rien entendre. — Toi, dis-moi}}{}
\ligneEst{}{\francais{lequel des rites sacrés nous fut révélé le premier.}}{}
\ligneEst{\textbf{Mnésiloque}}{\francais{Voyons. Que fût-ce encore que le premier? ... Nous bûmes.}}{630}
\ligneEst{\textbf{Mika}}{\francais{Et après ? le second ?}}{}
\ligneEst{\textbf{Mnésiloque}}{\francais{\espace{4cm}Nous bûmes à nos santés.}}{}
\ligneEst{\textbf{Mika}}{\francais{Tu l'as appris de quelqu'une. Et le troisième ?}}{}
\ligneEst{\textbf{Mnésiloque}}{\francais{Xenylla demanda une nef, car il n'y avait pas de pot de chambre.}}{}
\ligneEst{\textbf{Mika}}{\francais{Cela ne veut rien dire. Ici, viens ici, ô Clisthène.}}{634}
}
\vspace{2pt}
\noindent\rule{16.5cm}{0.6pt}\\[2pt]
\noindent\textbf{\small Commentaire philologique, métrique et scénique (vv. 590–634 --- partie 6/6)}\\[2pt]
\begin{adjustbox}{max totalheight=13.0cm, max width=16.5cm, keepaspectratio}%
\parbox[t]{16.5cm}{
\setlength{\columnsep}{9pt}%
\begin{multicols}{2}
\fontsize{8.4pt}{9.7pt}\selectfont
\setlength{\parindent}{0pt}%
\setlength{\parskip}{1.2pt}%
\noindent (Jeux de consonances entre ἄμη/ἀμίς et σκάφη/σκάφιον.) Apprécions à sa juste valeur l'à-propos de l'allusion au proverbe : il repose sur un quiproquo dû à la polysémie d'ἄμη -- comme est polysémique le σκάφιον dont parle Mnésiloque. De plus, le proverbe évoque ceux qui fournissent des réponses sans rapport avec la question posée -- ce qui est aussi le cas de notre pauvre Mnésiloque, bien embarrassé.
\noindent ἄμη et σκάφη pouvaient en effet être tous deux synonymes de \italique{seau}\marqueur{3}... mais c'était des \italique{houes} (autre sens de ἄμη) que demandait le paysan du proverbe...
\noindent\textbf{[Notes ad 633]} \textbf{(1)} Pour toutes les remarques de cette note, je renvoie aux articles \latin{amis} et \latin{scaphium} du \italique{Dictionnaire des antiquités grecques et romaines} de Daremberg \& Saglio.
\noindent \textbf{(2)} Il a été relevé par Suidas en α 1504 et α 1574. Érasme y consacre une entrée dans ses \italique{Adages} (au n°1149) et l'évoque aussi au n°1880. Maintes fois recensé par les parémiographes (par exemple Zénobius I, 83), il était connu de Plutarque (\italique{Du bavardage}, OM, 6, 35, 512e) et le scholiaste de la \italique{Paix} d'Aristophane le cite encore à propos de la chanson de Datis le Lydien aux vers 289-291 et de l'équivoque en grec que suggère la forme barbaristique χαίρομαι -- au passage : Aristophane joue sur le stéréotype du débauché « Lydien ($=$ Datis) à midi », Λυδὸς ἐν μεσημβρίᾳ (rapporté par Athénée, XII, 515d-517a et recensé par Érasme dans ses \italique{Adages}, n° 1594). Enfin, le scholiaste de Théocrite (\italique{Idylles}, IV, 10) cite aussi de ce proverbe. Pour finir, Aristophane y fait de nouveau allusion au fr. 927 Kassel-Austin : ἄγροικός εἰμι· τὴν σκάφην σκάφην λέγω.
\noindent \textbf{(3)} Comme dans le proverbe pour évoquer abondance de biens : ἄμαις καὶ σκάφαις ἀρύσασθαι. Abondance de biens nuit, soit dit en passant, à la santé du corps comme celle de l'esprit.
\noindent Bloy n'a, à ma connaissance, étrangement pas fait l'exégèse de ce lieu commun -- Flaubert aurait parlé d'une idée reçue.
\noindent Les latins et leur solide bon sens savaient que \latin{Copia parit fastidium} et \latin{Neque laus in copia, neque culpa in penuria constitit}.)
\end{multicols}
}%
\end{adjustbox}%
\vfill
\end{adjustwidth}
\newpage
\clearpage\ifodd\value{page}\hbox{}\newpage\fi
\enTeteG{ΕΠΕΙΣΟΔΙΟΝ Β' --- Codex Ravennas 429, f. 170r (Classense 172r) --- vv. 635–678}
\begin{adjustwidth}{-1.5cm}{-3.0cm}
\begin{adjustbox}{max totalheight=21.0cm, max width=16.5cm, keepaspectratio}
\begin{minipage}{16.5cm}
{\fontsize{7.2pt}{8.5pt}\selectfont
\noindent\textbf{Folio 16 --- Codex Ravennas 429, \texttt{f. 170r} (\texttt{Classense 172r}) : ΕΠΕΙΣΟΔΙΟΝ Β', vv. 635–678}\hfill \textit{Témoins :} $\text{R}^1, \Rpc, \Rdeux$ (\Mon)\par\vspace{2pt}
\hrule height 0.4pt\vspace{2.5pt}
}
{\fontsize{8.1pt}{9.5pt}\selectfont
\ligneRav{635}{\auDessus{γυ}{ν}}{ὅδ᾽ ἐστὶν \difference{ἁνὴρ} ὃν λέγεις. : τί οὖν ποιῶ;}{\textbf{635} ἀνὴρ R : \lecon{ἁνὴρ} Bentley}
\ligneRav{}{\auDessus{γυ}{ν}}{ἀπόδυσον αὐτόν· οὐδὲν ὑγιὲς γὰρ λέγει·}{}
\ligneRav{637}{μ\auDessus{ν}{Η}}{κἄπειτ᾽ ἀποδύσετ᾽ ἐννέα παίδων μητέρα;}{\textbf{637-40} (Μνη.) R : \lecon{(Μνη.) (637-8) et (Μίκα) (639-40)} Bibi}
\ligneRav{638}{}{\difference{χαλᾷ} ταχέως τὸ στρόφιον ὦναίσχυντε σύ·}{\textbf{638} χάλᾳ R : χάλα Grynaeus : \lecon{χαλᾷ} Bibi}
\ligneRav{639}{\auDessus{γυ}{ν}}{ὡς καὶ στιβαρά τις φαίνεται καὶ καρτερά·}{\textbf{639-40} (Μνη.) 639 et (Γυνή) 640 R : \lecon{(Μίκα)} Bibi}
\ligneRav{640}{}{καὶ νὴ Δία \difference{τιτθούς} γ᾽ ὥσπερ ἡμεῖς οὐκ ἔχει·}{\textbf{640} τίτθος R : \lecon{τιτθούς} ed. Iuntina}
\ligneRav{641}{μ\auDessus{ν}{Η}}{\difference{στερίφη} γάρ εἰμι κοὐκ ἐκύησα πώποτε·}{\textbf{641} στεριφῇ R : \lecon{στερίφη} (sed στεριφὴ) \Mon}
\ligneRav{642}{κλεισθ/}{νῦν· τότε \difference{δὲ} μήτηρ ἦσθα παίδων ἐννέα·}{\textbf{642-3} (Μνησ.) (641-3) R : \lecon{(Κλεισθ.) (642-3)} Bibi \espApp δὴ R : \lecon{δὲ} \Mon}
\ligneRav{}{}{ἀνίστασ᾽ ὀρθός. ποῖ τὸ πέος ὠθεῖς κάτω;}{}
\ligneRav{644}{\auDessus{γυ}{ν}}{\difference{τοδὶ διέκυψε} καὶ μάλ᾽ εὔχρων ὦ τάλαν·}{\textbf{644} τὸ δὶ δὴ ἔκυψε R : \lecon{τοδὶ διέκυψε} Dobree, Dindorf}
\ligneRav{645}{κλεισθ/}{καὶ ποῦ 'στιν; : αὖθις εἰς τὸ πρόσθεν οἴχεται·}{}
\ligneRav{646}{κλεισθ/}{οὐκ ἐνγεταυθί. : ἀλλὰ δεῦρ᾽ ἥκει πάλιν·}{\textbf{646} \lecon{ἐνγεταυθί} R : ἐγγεταυθί Thiersch \espApp \lecon{ἀλλὰ} R : μἀλλὰ Bentley}
\ligneRav{647}{κλεισθ/}{ἰσθμόν τιν᾽ ἔχεις \difference{ὤνθρωπ᾽}· ἄνω τε καὶ κάτω}{\textbf{647} ῶνθρωπε R : ἄνθρωπε \Mon : ἄνθρωπ᾽ Suid. ι 640 : \lecon{ὤνθρωπ᾽} Ellebodius}
\ligneRav{}{}{τὸ πέος διέλκεις πυκνότερον Κορινθίων·}{}
\ligneRav{}{\auDessus{γυ}{ν}}{ὤ μιαρὸς οὗτος· ταῦτ᾽ ἄρ᾽ ὑπὲρ Εὐριπίδου}{}
\ligneRav{650}{—}{ἡμῖν ἐλοιδορεῖτο. : κακοδαίμων ἐγώ,}{}
\ligneRav{651}{}{εἰς οἷ᾽ ἐμαυτὸν \difference{εἰσεκύλισα} πράγματα·}{\textbf{651} εἰσεκύλησα R : \lecon{εἰσεκύλισα} Scaliger}
\ligneRav{}{\auDessus{γυ}{ν}}{ἄγε δὴ τί δρῶμεν; : τουτονὶ φυλάττετε}{}
\ligneRav{653}{}{καλῶς, ὅπως μὴ διαφυγὼν \difference{οἰχήσεται}·}{\textbf{653} οἴχεται R : \lecon{οἰχήσεται} Faber p.788}
\ligneRav{654}{}{ἐγὼ δὲ ταῦτα τοῖς \difference{πρυτάνεσιν} ἀγγελῶ·}{\textbf{654} πρυτάνεσσιν R : \lecon{-εσιν} \Mon \espApp πρυτάνεσσιν R : \lecon{πρυτάνεσιν} \Mon}
\ligneRav{655}{χορ/}{ἡμᾶς τοίνυν μετὰ τοῦτ᾽ ἤδη τὰς λαμπάδας ἁψαμένας χρὴ}{}
\ligneRav{}{}{ξυζωσαμένας εὖ κἀνδρείως τῶν θ᾽ ἱματίων ἀποδύσας}{}
\ligneRav{657}{}{ζητεῖν, εἴ που κἄλλος τις ἀνὴρ \difference{ἐπελήλυθε}, καὶ περιθρέξαι}{\textbf{657} εἰσελήλυθεν R : \lecon{ἐπελήλυθε} Handley \textbf{658} ἀθρῆσαι R : \lecon{διαθρῆσαι} Daubuz}
\ligneRav{658}{}{τὴν Πύκνα πᾶσαν καὶ τὰς σκηνὰς καὶ τὰς διόδους \difference{διαθρῆσαι}·}{\textbf{658} ἀθρῆσαι R : \lecon{διαθρῆσαι} Daubuz}
\ligneRav{659}{ἡμιχ/}{εἶα δὴ πρώτιστα μὲν χρὴ κοῦφον ἐξορμᾶν πόδα}{\textbf{659} (Ἡμιχ. γυν.) \Rpc}
\ligneRav{660}{}{καὶ διασκοπεῖν σιωπῇ πανταχῇ· μόνον δὲ χρὴ}{\textbf{660} πανταχῇ post χρὴ R : \lecon{ante μόνον} Küster}
\ligneRav{}{}{μὴ βραδύνειν, ὡς ὁ καιρός ἐστι μὴ μέλλειν ἔτι,}{}
\ligneRav{662}{}{ἀλλὰ τὴν πρώτην τρέχειν \difference{χρῆν} ὡς τάχιστ᾽ ἤδη κύκλῳ·}{\textbf{662} χρὴ R : \lecon{χρῆν} Bentley : χρῆ μ᾽ Austin}
\ligneRav{}{χορ/}{\hspace{1.1em}εἶά νυν ἴχνευε καὶ μά-}{}
\ligneRav{}{}{\hspace{1.1em} τευε ταχὺ πάντ᾽, εἴ τις ἐν τόποις ἑδραῖος}{}
\ligneRav{}{}{\hspace{1.1em}ἄλλος αὖ λέληθεν ὤν·}{}
\ligneRav{665}{}{\hspace{1.1em}πανταχῇ \difference{δὲ ῥῖψον} ὄμμα,}{\textbf{665} διάρριψον R : \lecon{δὲ ῥῖψον} Hermann}
\ligneRav{666}{}{\hspace{1.1em}καὶ τὰ τῇδε καὶ τὰ δεῦρο}{\textbf{666} post καὶ τὰ τῇδε add. καὶ τὰ κεῖσε Kaibel \espApp Post καὶ τὰ τῇδε add. καὶ τὰ κεῖσε Kaibel}
\ligneRav{666}{}{\hspace{1.1em}πάντ᾽ ἀνασκόπει καλῶς·}{\textbf{666} post καὶ τὰ τῇδε add. καὶ τὰ κεῖσε Kaibel \espApp Post καὶ τὰ τῇδε add. καὶ τὰ κεῖσε Kaibel}
\ligneRav{667}{}{\hspace{1.1em}ἢν γάρ ληφθῇ δράσας ἀνόσια,}{\textbf{667} μὴ λάθῃ R : \lecon{ληφθῇ} Reisig : με λάθῃ Bergk}
\ligneRav{}{}{\hspace{1.1em}δώσει τε δίκην καὶ πρὸς τούτῳ}{}
\ligneRav{669}{}{\hspace{1.1em}ἦν δὲ μὴ λάθῃ, δώσει δίκην, παράδειγμά τ᾽ \difference{ἀνδράσιν} (ἅπασιν \inRas{ἅπασιν}{}) ἔσται}{\textbf{669} ἅπασιν \inRas{ἅπασιν}{} ἔσται \Rpc : ἅπασιν ἔσται \Mon : \lecon{ἀνδράσιν ἔσται} Beer ap. Hermann (1845) : ἔσται ἅπασιν Bergk : alii alia}
\ligneRav{670}{}{\hspace{1.1em}παράδειγμ᾽ ὕβρεως}{}
\ligneRav{}{}{\hspace{1.1em}ἀδίκων τ᾽ ἔργων ἀθέων τε τρόπων·}{}
\ligneRav{}{}{\hspace{1.1em}φήσει δ᾽ εἶναί τε θεοὺς φανερῶς,}{}
\ligneRav{}{}{\hspace{1.1em}δείξει τ᾽ ἤδη}{}
\ligneRav{}{}{\hspace{1.1em}πᾶσιν ἀνθρώποις σεβίζειν δαίμονας}{}
\ligneRav{675}{}{\hspace{1.1em}†δικαίως τ᾽ ἐφέποντας † ὅσια καὶ νόμιμα}{\textbf{675} \lecon{ἐφέποντας} R : ἐφέπειν Hermann}
\ligneRav{}{}{\hspace{1.1em}μηδομένους ποιεῖν ὅ τι καλῶς ἔχει·}{}
\ligneRav{}{}{\hspace{1.1em}κἂν μὴ ποιῶσι ταῦτα τοιάδ᾽ ἔσται·}{}
}
\vspace{3pt}\noindent\rule{4.5cm}{0.4pt}\par\vspace{1.5pt}
{\fontsize{7.0pt}{8.3pt}\selectfont
\noindent\textbf{Loci paralleli (vv. 635–678) ---} \textbf{641} (στερίφη) Hesychius σ 1769 στερίφη· στεῖρα, μὴ τεκοῦσα, μηδὲ τίκτουσα $=$ Photius σ 532 ; cf. Moeris σ 46 (στερίφη $=$ στερρά) \espTest \textbf{647-8} Suidas ι 640 \espTest \textbf{659} (εἶα δὴ) Hesychius ε 8 ἔα δή· ἄγε δή $=$ Photius ε 1 $=$ \Syn ε 1 $=$ Suidas ε 3\par
}
\vspace{2.5pt}\noindent\rule{4.5cm}{0.4pt}\par\vspace{1.5pt}
{\fontsize{6.9pt}{8.2pt}\selectfont
\noindent\textbf{Suite (en commentaire) de l'apparat critique \& notes paléographiques (\texttt{f. 170r}) ---} \textbf{Sigle tachygraphique marginal μν\auDessus{η}{σ} au v. 638 et redistribution des répliques par \textbf{Bibi} (vv. 637–643) :} Comme le précise Louis Hauseux (\texttt{vv 1-641.tex}, l. 4618), R comporte en marge du v. 638 le sigle tachygraphique \lecon{μν\auDessus{η}{σ}} et aucune indication de changement d'interlocuteur jusqu'au v. 640 (Γυνή). Aux vv. 637–640, Louis Hauseux donne 637–638 à (Μνη.) (en lisant l'indicatif 3\ieme{} pers. \lecon{χαλᾷ} au v. 638) et 639–640 à \lecon{(Μίκα)} ; aux vv. 642–643, il restitue la réplique à \lecon{(Κλεισθ.)}. --- \textbf{Confusion onciale ΙΣ $\leftrightarrow$ Π au v. 657 :} εἰσελήλυθεν R corrigé par \textbf{Handley} en \lecon{ἐπελήλυθε}. --- \textbf{Grattage d'un mot entier répété par dittographie au v. 669 :} \Rac{} avait écrit deux fois de suite ἅπασιν ; \Rpc{} a gratté le second (\lecon{ἅπασιν \inRas{ἅπασιν}{} ἔσται}), laissant sur le parchemin un espace blanc très net entre ἅπασιν et ἔσται (Beer ap. Hermann restitue ἀνδράσιν ἔσται). --- \textbf{vv. 637–640 :} μν\auDessus{η}{σ} (ad v. 638, nulla mut. pers. usque ad 640 Γυνή) R : \lecon{(Μνη.) (637–638) et (Μίκα) (639–640)} \lecon{Bibi}. --- \textbf{v. 638 :} χάλᾳ R (χάλα Grynaeus) : \lecon{χαλᾷ} \lecon{Bibi}. --- \textbf{vv. 642–643 :} (Μνησ.) (641–643) R : \lecon{(Κλεισθ.) (642–643)} \lecon{Bibi} \espApp \textbf{v. 642 :} δὴ R : \lecon{δὲ} \Mon{}. --- \textbf{v. 644 :} τὸ δὶ δὴ ἔκυψε R (médivision et iotacisme) : \lecon{τοδὶ διέκυψε} Dobree, Dindorf. --- \textbf{v. 651 :} εἰσεκύλησα R (iotacisme) : \lecon{εἰσεκύλισα} Scaliger. --- \textbf{v. 657 :} εἰσελήλυθεν R : \lecon{ἐπελήλυθε} Handley (confusion onciale ΙΣ/Π).\par
}
\end{minipage}
\end{adjustbox}
\end{adjustwidth}
\newpage
\enTeteD{Codex Ravennas 429 (Biblioteca Classense) --- Fac-similé du f. 170r (vv. 635–678)}
\begin{adjustwidth}{-1.5cm}{-3.0cm}
\begin{center}
\includegraphics[width=16.5cm,height=20.5cm,keepaspectratio]{/workspaces/Thesmophories/dossier_ravennas_thesmophories/folios_individuels_ff162v-177r/16_f170r_classense172r_vv0635-0678.jpg}\par
\vspace{2pt}
{\scriptsize\textbf{Codex Ravennas 429, f. 170r (Classense 172r)} --- ΕΠΕΙΣΟΔΙΟΝ Β', \textbf{vv. 635–678} ($\text{R}^1, \Rpc, \Rdeux$ (\Mon)).}
\end{center}
\end{adjustwidth}
\newpage
\clearpage\ifodd\value{page}\hbox{}\newpage\fi% --- FOLIO 16 (f. 170r) : SPREAD 2.1 (Texte grec établi + Traduction française + Commentaire 1/8 & 2/8) ---
\clearpage\ifodd\value{page}\hbox{}\newpage\fi
\enTeteG{ΘΕΣΜΟΦΟΡΙΑΖΟΥΣΑΙ --- Texte grec établi (f. 170r, vv. 635–678 [1/4])}
\enTeteD{LES THESMOPHORIES --- Traduction française (f. 170r, vv. 635–678 [1/4])}
\begin{adjustwidth}{-2cm}{-2cm}
\noindent\textbf{\footnotesize Texte grec établi --- Folio 16 (\texttt{f. 170r}, vv. 635–678 [1/4]) \hfill ΕΠΕΙΣΟΔΙΟΝ Β'}\\[-2pt]
\noindent\rule{16.5cm}{0.5pt}\\[2pt]
{\fontsize{8.8pt}{10.2pt}\selectfont
\ligneEst{}{ὅδ᾽ ἐστὶν ἁνὴρ ὃν λέγεις.}{635}
\ligneEst{\textbf{Κλεισθένης}}{\decalage{ὅδ᾽ ἐστὶν ἁνὴρ ὃν λέγεις.}τί οὖν ποιῶ;}{}
\ligneEst{\textbf{Μίκα}}{ἀπόδυσον αὐτόν· οὐδὲν ὑγιὲς γὰρ λέγει.}{}
\ligneEst{\textbf{Μνησίλοχος}}{κἄπειτ᾽ ἀποδύσετ᾽ ἐννέα παίδων μητέρα;}{}
\ligneEst{}{χαλᾷ ταχέως τὸ στρόφιον ὦναίσχυντε σύ.}{}
\ligneEst{\textbf{Μίκα}}{ὡς καὶ στιβαρά τις φαίνεται καὶ καρτερά·}{}
\ligneEst{}{καὶ νὴ Δία τιτθούς γ᾽ ὥσπερ ἡμεῖς οὐκ ἔχει.}{640}
\ligneEst{\textbf{Μνησίλοχος}}{στερίφη γάρ εἰμι κοὐκ ἐκύησα πώποτε.}{}
\ligneEst{\textbf{Κλεισθένης}}{νῦν· τότε δὲ μήτηρ ἦσθα παίδων ἐννέα.}{}
\ligneEst{}{ἀνίστασ᾽ ὀρθός. ποῖ τὸ πέος ὠθεῖς κάτω;}{}
\ligneEst{\textbf{Μίκα}}{τοδὶ διέκυψε καὶ μάλ᾽ εὔχρων ὦ τάλαν.}{}
\ligneEst{\textbf{Κλεισθένης}}{καὶ ποῦ 'στιν;}{645}
}
\vspace{2pt}
\noindent\rule{16.5cm}{0.6pt}\\[2pt]
\noindent\textbf{\small Commentaire philologique, métrique et scénique (vv. 635–678 --- partie 1/8)}\\[2pt]
\begin{adjustbox}{max totalheight=15.2cm, max width=16.5cm, keepaspectratio}%
\parbox[t]{16.5cm}{
\setlength{\columnsep}{9pt}%
\begin{multicols}{2}
\fontsize{7.6pt}{8.8pt}\selectfont
\setlength{\parindent}{0pt}%
\setlength{\parskip}{1.2pt}%
\noindent\textbf{637-40} \textbf{κἄπειτ᾽ ἀποδύσετ᾽ ἐννέα παίδων μητέρα; ... καὶ νὴ Δία τίτθος γ᾽ ὥσπερ ἡμεῖς οὐκ ἔχει.} Le manuscrit de Ravenne porte en marge du vers 638 μν\auDessus{η}{σ}, indiquant ainsi que Mnésiloque prend la parole ; le manuscrit ne présente alors plus aucun signe de changement de locuteur jusqu'au vers 641\marqueur{1}. Or Mnésiloque  ne peut aucunement prononcer le vers 639, commentaire sur sa poitrine étonnamment plate (ὡς καὶ στιβαρά τις φαίνεται καὶ καρτερά).
\noindent Il y a donc un \italique{paragraphos} à rétablir.
\noindent Je dis \italique{un} car les éditeurs suivent quasi-unanimement Brunck en `rétablissent' deux : au début du vers 638 (Clistène) et au début du vers 639 (la femme -- Mika -- reprend la parole).
\noindent Une autre solution me semble plus légère : laisser à Mnésiloque les deux vers 637-638 et à Mika les vers 639-640.
\noindent Cette retouche comporte surtout l'avantage d'être bien meilleure quant à l'action et au jeu de scène.\\ Presque tous les éditeurs depuis Brunck, disais-je, retirent le vers 638 à Mnésiloque pour l'accorder à Clisthène. Double erreur ! \ \ \ \italique{Primo}, parce qu'ils sont alors contraints de surajouter une seconde émendation : transformer le χάλᾳ du manuscrit en un χάλα (impératif lancé par Clisthène) alors que le iota adscrit est nettement écrit. \ \ \ \italique{Secundo}, parce que Mika, au vers 636, avait déjà ordonné à Clisthène de déshabiller Mnésiloque (ἀπόδυσον αὐτόν) ; il n'y a aucune de raison de douter qu'il ne soit en train de lui obéir. Un peu d'action, que diable ! une pièce d'Aristophane ne doit pas avoir la même pesanteur que l'administration française. (Selon nos éditeurs, Mika ordonnerait à Clisthène de déshabiller Mnésiloque qui transmettrait alors l'ordre à l'intéressé... qui finirait par ne rien faire ?..)\\ Le vers 638, χαλᾷ ταχέως τὸ στρόφιον ὦναίσχυντε σύ, était pourtant excellent dans la bouche de Mnésiloque : Clisthène est en train de le déshabiller, il feint l'indignation (χαλᾷ est un subjonctif dit \italique{délibératif}) devant l'empressement (ταχέως) que cet homme -- car Clisthène reste un homme -- met à enlever le \italique{soutien gorge} d'une femme. « Il défait mon strophion avec promptitude. Impudent ! » Mnésiloque continue de jouer une femme -- vertueuse même, mais un peu tard. Qu'y a-t-il de plus limpide ? l'apostrophe ὦναίσχυντε σύ conclue logiquement la tirade de mère effarouchée qu'affecte le peu crédible Mnésiloque (et qui fait d'ailleurs écho au précédent ἀναίσχυντός τις εἶ (v. 611) lancé par le même Mnésiloque au même Clisthène\marqueur{2}).
\noindent\textbf{[Notes ad 637-40]} \textbf{(1)} Il est d'usage dans notre manuscrit que le \italique{signe} soit un παράγραφος semblable à un tiret cadratin ` --- ' lorsque le changement de locuteur se fait en début de vers. Si le changement intervient en milieu de vers, il y a plus souvent un \latin{punctum duplex} ` : ' ou bien une abréviation indiquant le nom du personnage qui prend la parole. Par abus de langage, nous utiliserons le terme de \italique{paragraphos} pour désigner indifféremment l'un de ces signes.
\noindent \textbf{(2)} Quelques éditeurs avaient même prétexté de ce double ἀναίσχυντος pour attribuer l'ἀναίσχυντός τις εἶ du vers 611 aussi à Clisthène, histoire que leurs multiples tripatouillages présentent une apparence d'homogénéité...
\end{multicols}
}%
\end{adjustbox}%
\vfill
\end{adjustwidth}
\newpage
% --- PAGE DE DROITE : Traduction française + Commentaire (2/8) ---
\begin{adjustwidth}{-2cm}{-2cm}
\noindent\textbf{\footnotesize Traduction française --- Folio 16 (\texttt{f. 170r}, vv. 635–678 [1/4]) \hfill ΕΠΕΙΣΟΔΙΟΝ Β'}\\[-2pt]
\noindent\rule{16.5cm}{0.5pt}\\[2pt]
{\fontsize{8.8pt}{10.2pt}\selectfont
\ligneEst{}{\francais{Voilà l'homme dont tu parles.}}{635}
\ligneEst{\textbf{Clisthène}}{\francais{\espace{3.7cm}Que dois-je faire ?}}{}
\ligneEst{\textbf{Mika}}{\francais{Déshabille-le. Ce qu'il dit n'a pas de sens.}}{}
\ligneEst{\textbf{Mnésiloque}}{\francais{Eh quoi! déshabiller une mère de neuf enfants !}}{}
\ligneEst{}{\francais{Il défait mon strophion avec promptitude. Impudent !}}{}
\ligneEst{\textbf{Mika}}{\francais{Qu'elle a l'air robuste et vigoureuse !}}{}
\ligneEst{}{\francais{Et, par Zeus, elle n'a pas de tétons comme nous.}}{640}
\ligneEst{\textbf{Mnésiloque}}{\francais{C'est que je suis stérile. Je n'ai jamais conçu.}}{}
\ligneEst{\textbf{Clisthène}}{\francais{À présent. Mais tout à l'heure tu étais mère de neuf enfants.}}{}
\ligneEst{}{\francais{Tiens-toi debout. Où pousses-tu ton cas là-dessous ?}}{}
\ligneEst{\textbf{Mika}}{\francais{(Derrière Mnésiloque). Le voilà qui émerge, et beau de couleur, le coquin.}}{}
\ligneEst{\textbf{Clisthène}}{\francais{(Échangeant de place avec Mika). Où est-il ?}}{645}
}
\vspace{2pt}
\noindent\rule{16.5cm}{0.6pt}\\[2pt]
\noindent\textbf{\small Commentaire philologique, métrique et scénique (vv. 635–678 --- partie 2/8)}\\[2pt]
\begin{adjustbox}{max totalheight=15.2cm, max width=16.5cm, keepaspectratio}%
\parbox[t]{16.5cm}{
\setlength{\columnsep}{9pt}%
\begin{multicols}{2}
\fontsize{7.6pt}{8.8pt}\selectfont
\setlength{\parindent}{0pt}%
\setlength{\parskip}{1.2pt}%
\noindent\textbf{641} \textbf{στερίφη γάρ εἰμι.} Jeu de mots d'Aristophane : comment Mnésiloque explique-t-il la platitude de sa poitrine ? tout simplement parce qu'il a les chairs \italique{fermes} (sens premier de στέριφος, le Scythe l'emploie avec cette acception au vers 1185). Mais στέριφος peut aussi avoir le sens dérivé de \italique{stérile} ; sens que le spectateur ne privilégiera qu'en entendant la seconde partie de la réplique de Mnésiloque (κοὐκ ἐκύησα). C'est la raison pour laquelle Aristophane recourt à un mot rare, στέριφος. (Le scholiaste de noter : ἀντὶ τοῦ στείρα, commentaire qui se retrouve dans de nombreux lexiques, cf. l'apparat des témoins\marqueur{1}.)
\noindent\textbf{[Notes ad 641]} \textbf{(1)} Et pour l'apparat le plus fourni, se reporter à celui de l'édition Theodoridis de Photius.
\noindent\textbf{642-3} \textbf{νῦν· τότε δὲ μήτηρ ἦσθα παίδων ἐννέα. / ἀνίστασ᾽ ὀρθός. ποῖ τὸ πέος ὠθεῖς κάτω;} Il n'y a aucun \italique{paragraphos} aux vers 642 et 643.
\noindent Pourtant, ce n'est assurément pas Mnésiloque (qui a la parole au vers précédent) qui peut prononcer ces vers. On ajoute souvent deux changement d'interlocuteurs : Mika pour le vers 642 et Clisthène pour le vers 643. Par simplicité, je préfère n'en ajouter qu'un et attribue donc le distique entier à Clisthène.
\noindent\textbf{643} \textbf{ἀνίστασ᾽ ὀρθός.} «~Redresse-toi !~» L'injonction adressée à Mnésiloque est bien-sûr à double entente, et Clisthène se fera doublement obéir. Mais interrogeons-nous : pourquoi Mnésiloque s'était-il assis ou recroquevillé ? réponse simple : pour cacher sa virilité.
\noindent Ils se redressent. Mnésiloque tente une dernière fois de cacher ce qui doit être caché et Clisthène : ποῖ τὸ πέος ὠθεῖς κάτω; «~Où repousses-tu ce \italique{péos} ?!~» Le \italique{péos} passe alors au verso de Mnésiloque, d'où peut le contempler Mika, admiratrice. Avant que ne s'ensuive la scène burlesque où Clisthène paraît fort intéressé aussi.\\ Mnésiloque s'était habillé dans cet ordre : κροκωτός, στρόφιον et ἔγκυκλον. Le στρόφιον et l'ἔγκυκλον viennent d'être retirés, mais pas le κροκωτός (au vers 941, Mnésiloque le porte toujours) ; comment se fait-il que son \italique{olisbos} soit apparent malgré la présence de cette robe safranée ? parce que la crocote est composée d'un tissu très fin qui la rend transparente\marqueur{1}
\noindent\textbf{643} \textbf{Note sur τὸ ἔγκυκλον $=$ \italique{une jupe}.} Clisthène vient de retirer le στρόφιον de Mnésiloque, ce qui entraîna quelques remarques sur sa poitrine. Or, lors de la scène d'habillement (vers 253-261), Mnésiloque avait successivement enfilé le κροκωτός, le στρόφιον \italique{puis, à la fin,} l'ἔγκυκλον. Il résulte donc que l'ἔγκυκλον \italique{est un vêtement qui ne peut pas recouvrir la poitrine}. (Sinon Clisthène aurait dû commencer par l'encycle avant le στρόφιον.)  Ce ne peut donc être en aucun cas un \italien{mantello ampio, pesante}\marqueur{1}.
\noindent Allongeons quelque peu cette note afin que soit irréfragable l'identification de l'ἔγκυκλον avec une sorte de jupe courte.\\
\end{multicols}
}%
\end{adjustbox}%
\vfill
\end{adjustwidth}
\newpage

% --- FOLIO 16 (f. 170r) : SPREAD 2.2 (Texte grec établi + Traduction française + Commentaire 3/8 & 4/8) ---
\clearpage\ifodd\value{page}\hbox{}\newpage\fi
\enTeteG{ΘΕΣΜΟΦΟΡΙΑΖΟΥΣΑΙ --- Texte grec établi (f. 170r, vv. 635–678 [2/4])}
\enTeteD{LES THESMOPHORIES --- Traduction française (f. 170r, vv. 635–678 [2/4])}
\begin{adjustwidth}{-2cm}{-2cm}
\noindent\textbf{\footnotesize Texte grec établi --- Folio 16 (\texttt{f. 170r}, vv. 635–678 [2/4]) \hfill ΕΠΕΙΣΟΔΙΟΝ Β'}\\[-2pt]
\noindent\rule{16.5cm}{0.5pt}\\[2pt]
{\fontsize{8.8pt}{10.2pt}\selectfont
\ligneEst{\textbf{Μίκα}}{\decalage{καὶ ποῦ 'στιν;}αὖθις εἰς τὸ πρόσθεν οἴχεται.}{}
\ligneEst{\textbf{Κλεισθένης}}{οὐκ ἐνγεταυθί.}{}
\ligneEst{\textbf{Μίκα}}{\decalage{οὐκ ἐνγεταυθί.}ἀλλὰ δεῦρ᾽ ἥκει πάλιν.}{}
\ligneEst{\textbf{Κλεισθένης}}{ἰσθμόν τιν᾽ ἔχεις ὤνθρωπ᾽· ἄνω τε καὶ κάτω}{}
\ligneEst{}{τὸ πέος διέλκεις πυκνότερον Κορινθίων.}{}
\ligneEst{\textbf{Μίκα}}{ὤ μιαρὸς οὗτος· ταῦτ᾽ ἄρ᾽ ὑπὲρ Εὐριπίδου}{}
\ligneEst{}{ἡμῖν ἐλοιδορεῖτο.}{650}
\ligneEst{\textbf{Μνησίλοχος}}{\decalage{ἡμῖν ἐλοιδορεῖτο.}κακοδαίμων ἐγώ,}{}
\ligneEst{}{εἰς οἷ᾽ ἐμαυτὸν εἰσεκύλισα πράγματα.}{}
\ligneEst{\textbf{Μίκα}}{ἄγε δὴ τί δρῶμεν;}{}
\ligneEst{\textbf{Κλεισθένης}}{\decalage{ἄγε δὴ τί δρῶμεν;}τουτονὶ φυλάττετε}{}
\ligneEst{}{καλῶς, ὅπως μὴ διαφυγὼν οἰχήσεται·}{}
\ligneEst{}{ἐγὼ δὲ ταῦτα τοῖς πρυτάνεσιν ἀγγελῶ.}{}
}
\vspace{2pt}
\noindent\rule{16.5cm}{0.6pt}\\[2pt]
\noindent\textbf{\small Commentaire philologique, métrique et scénique (vv. 635–678 --- partie 3/8)}\\[2pt]
\begin{adjustbox}{max totalheight=15.2cm, max width=16.5cm, keepaspectratio}%
\parbox[t]{16.5cm}{
\setlength{\columnsep}{9pt}%
\begin{multicols}{2}
\fontsize{7.6pt}{8.8pt}\selectfont
\setlength{\parindent}{0pt}%
\setlength{\parskip}{1.2pt}%
\noindent Que pourrait-il être d'autre ce vêtement qui sitôt retiré laisse tout émoustillé le mari de Praxagora ; voyez plutôt (\italique{Assemblée des femmes}, vers 535-538) :
\citComm{Βλέπυρος -- εἶτ' οὐ τὸ σαυτῆς ἱμάτιον ἐχρῆν σ' ἔχειν;\\
ἀλλ' ἔμ' ἀποδύσασ' ἐπιβαλοῦσα \textbf{τοὔγκυκλον}\\
ᾤχου καταλιποῦσ' ὡσπερεὶ προκείμενον,\\
μόνον οὐ στεφανώσασ' οὐδ' ἐπιθεῖσα λήκυθον.\\
Alors que n'as-tu pris ton propre manteau [pour sortir] ? Mais tu me dévêts et jettes sur moi \italique{ton encycle}, puis tu t'en vas, me laissant là comme un cadavre exposé, hors que tu as omis de me couronner et de placer près de moi un lécythe.}
\noindent Le champ lexical est celui d'une exposition funèbre\marqueur{2} comme la traduction de Willems le laisse entendre, mais le λήκυθον du vers 538 doit être pris \italique{sensu obsceno} selon la matrice métaphorique : λήκυθος $=$ ἀλάβαστος\marqueur{3} $=$ \latin{phallus}.
\noindent Blépyros reproche à sa femme de ne pas s'être acquitté du \italique{nome de mâtines} (ὄρθριος νόμος\marqueur{4}).
\noindent Ce vêtement permet à Aristophane un jeu de mot avec l'\italique{enceinte} de la ville fortifiée de Pylos\marqueur{5} -- oui ! --, la Πύλος, cette ville qui déclenche intarissablement la verve éroticomique de notre poète, jonglant avec l'autre matrice métaphorique (féminine cette fois) : πύλος $=$ πύλη $=$ θύρα $=$ \latin{pudenda muliebria} (raison pour laquelle l'Athénien ne veut pas la céder).
\noindent Il y avait certainement à l'époque de la représentation de \italique{Lysistrata} -- sans doute d'une année antérieure à celle des \italique{Thesmophories} -- des tractations pour que la paix fût conclue en échange de la remise de Pylos ou, par synecdoque, de son \italique{encycle}.
\noindent Ce n'est pas autrement qu'il faut comprendre le passage au début de la même pièce (vers 111-114) : Lysistrata demande :
\citComm{ἐθέλοιτ᾽ ἂν οὖν, εἰ μηχανὴν εὕροιμ᾽ ἐγώ,\\
μετ᾽ ἐμοῦ καταλῦσαι τὸν πόλεμον;\\
Consentirez-vous donc, si j'en trouve le moyen, à mettre fin avec moi à la guerre ?}
\noindent et Aristophane de répondre, par la voix d'une femme :
\citComm{\decalage{μετ᾽ ἐμοῦ καταλῦσαι τὸν πόλεμον;}νὴ τὼ θεώ\\
ἔγωγ᾽ ἂν οὖν κἂν εἴ με χρείη \textbf{τοὔγκυκλον}\\
τουτὶ καταθεῖσαν ἐκπιεῖν αὐθημερόν.\\
Oui, par les deux déesses. J'en suis, quand il me faudrait déposer l'\italique{encycle} que voici en gage pour boire dès aujourd'hui [à la paix : les traités, σπονδαί, sont en grec, littéralement, des \italique{libations}].}
\noindent Il faut donc bien-sûr aussi attribuer à cette même femme -- et comme ne le font pas les éditeurs d'Aristophane -- le vers 104 :
\citComm{ὁ δ᾽ ἐμός γε τελέους ἑπτὰ μῆνας ἐν Πύλῳ.}
\end{multicols}
}%
\end{adjustbox}%
\vfill
\end{adjustwidth}
\newpage
% --- PAGE DE DROITE : Traduction française + Commentaire (4/8) ---
\begin{adjustwidth}{-2cm}{-2cm}
\noindent\textbf{\footnotesize Traduction française --- Folio 16 (\texttt{f. 170r}, vv. 635–678 [2/4]) \hfill ΕΠΕΙΣΟΔΙΟΝ Β'}\\[-2pt]
\noindent\rule{16.5cm}{0.5pt}\\[2pt]
{\fontsize{8.8pt}{10.2pt}\selectfont
\ligneEst{\textbf{Mika}}{\francais{\espace{2.2cm}Il est reparti en avant.}}{}
\ligneEst{\textbf{Clisthène}}{\francais{(Revenant devant.) Non, pas de ce côté.}}{}
\ligneEst{\textbf{Mika}}{\francais{\espace{2.2cm}Eh! le voilà revenu par ici.}}{}
\ligneEst{\textbf{Clisthène}}{\francais{Une façon d'isthme que tu as là, homme : tu fais passer ton cas de-ci de-là,}}{}
\ligneEst{}{\francais{plus dru que les Corinthiens.}}{}
\ligneEst{\textbf{Mika}}{\francais{Ah! la canaille! Voilà donc pourquoi au profit d'Euripide}}{}
\ligneEst{}{\francais{il nous invectivait.}}{650}
\ligneEst{\textbf{Mnésiloque}}{\francais{\espace{2.5cm}Malheureux que je suis !}}{}
\ligneEst{}{\francais{dans quelle méchante affaire je me suis empêtré !}}{}
\ligneEst{\textbf{Mika}}{\francais{Voyons, que faisons-nous ?}}{}
\ligneEst{\textbf{Clisthène}}{\francais{\espace{2.6cm}Veillez bien sur lui,}}{}
\ligneEst{}{\francais{qu'il ne se dérobe point par la fuite.}}{}
\ligneEst{}{\francais{Moi, je vais rapporter la chose aux prytanes.}}{}
}
\vspace{2pt}
\noindent\rule{16.5cm}{0.6pt}\\[2pt]
\noindent\textbf{\small Commentaire philologique, métrique et scénique (vv. 635–678 --- partie 4/8)}\\[2pt]
\begin{adjustbox}{max totalheight=15.2cm, max width=16.5cm, keepaspectratio}%
\parbox[t]{16.5cm}{
\setlength{\columnsep}{9pt}%
\begin{multicols}{2}
\fontsize{7.6pt}{8.8pt}\selectfont
\setlength{\parindent}{0pt}%
\setlength{\parskip}{1.2pt}%
\noindent (Elle se dit prête à faire tomber l'\italique{encycle} dès le retour de son mari en garnison depuis sept longs mois à Pylos.) Aristophane, par deux fois dans la même pièce, exhorte ses contemporains à renoncer à Pylos, c'est dire l'importance qu'il accordait à l'éventualité de la paix ! On connaît la suite...
\noindent Ce vêtement, donc, qui enclot le sadinet, puisqu'il faut parler net... Alors ? que peut-il bien être ce vêtement ? S'il vous plaît, ne me répondez pas de chœur avec tous les fonctionnaires de la recherche : «~un lourd et pesant manteau~».\\ Mais délaissons un peu nos primes amours, les comédies d'Aristophane, pour aborder un terrain plus aride et cependant qui n'est certainement pas à négliger, je veux parler de la lexicologie.
\noindent Ce sens de \italique{jupe} pour l'ἔγκυκλον est confirmé, ce me semble, par le commentaire le plus précis : celui du scholiaste de \italique{Lysistrata} (vers 113), qui glose ἔγκυκλον ainsi :
\citComm{τὸ περιβόλαιον τὸ ζωνάριον.}
\noindent J'interprète ζωνάριον comme un adjectif se rapportant à περιβόλαιον\marqueur{6} : \italique{le vêtement / la couverture}, celle qui enceint, \italique{au niveau de la taille} ; une jupe donc. Citons aussi ce passage du copieux commentaire à l'\italique{Iliade} (976.13) d'Eustathe de Thessalonique (il commente \italique{Iliade} XIV, v. 178) :
\citComm{ἰστέον δὲ ὅτι, καθάπερ ὁ ἑανὸς οὕτω καὶ τὸ ἔκκυκλον ἢ ἔγκυκλον, οὗ καὶ ὁ κωμικὸς μέμνηται, γυναικεῖον ἦν ποτε φόρημα, φησὶ γοῦν Παυσανίας ὅτι ἔγκυκλον περιπόρφυρον\marqueur{7} ἱμάτιον καὶ χιτὼν γυναικεῖος, ὃν ἔνδοθεν ἐνδύονται γυναῖκες, εἶτα τὸ ἔνδυμα. λέγεται δὲ φησίν [sic, je lis : σφῇσιν] ὑπ' ἐνίων ἐπιγονατίς, διότι μέχρι γονάτων διήκει.\\
Il faut savoir que, de même que l'ἑανὸς [= riche robe de femme dans le passage d'Homère qu'Eustathe commente], il y a aussi l'ἔκκυκλον (ou ἔγκυκλον), dont s'est aussi souvenu le Comique [= Aristophane]. C'était un vêtement (φόρημα) féminin porté jadis. Ce qui est sûr, c'est que Pausanias indique que l'ἔγκυκλον était un vêtement (ἱμάτιον) aux franges pourpres (περιπόρφυρον) et une tunique (χιτὼν) féminine, que les femmes revêtaient à l'intérieur, et ensuite un habit (ἔνδυμα). Il est qualifié par elles-mêmes, d'après certains, d'ἐπιγονατίς, parce qu'il atteignait les genoux.}
\noindent Ce vêtement érotique d'intérieur peut-il être autre chose qu'une \italique{jupe sexy} ?
\noindent\textbf{[Notes ad 643]} \textbf{(1)} Cf. le cratère à figures rouges (ca. -380) représentant une scène de comédie : un vieillard entouré de deux gardes et vêtu de la crocote... qui laisse tout transparaître ; il est reproduit en annexe. Cratère «~St. Agata Antigone~», St. Agata dei Goti I, collection Rainone-Mustilli. Inventorié par Arthur D. Trendall , \italique{Phlyax Vases}, «~BICS~» suppl. LXIX (1967).
\noindent \textbf{(1)} Comme le conclue Franca Perusino dans son article «~\italien{L'ἔγκυκλον, un mantello femminile nelle commedie di Aristofane}~», \italique{Quaderni Urbinati di Cultura Classica}, 72.3 (2002), pp. 129-133.
\end{multicols}
}%
\end{adjustbox}%
\vfill
\end{adjustwidth}
\newpage

% --- FOLIO 16 (f. 170r) : SPREAD 2.3 (Texte grec établi + Traduction française + Commentaire 5/8 & 6/8) ---
\clearpage\ifodd\value{page}\hbox{}\newpage\fi
\enTeteG{ΘΕΣΜΟΦΟΡΙΑΖΟΥΣΑΙ --- Texte grec établi (f. 170r, vv. 635–678 [3/4])}
\enTeteD{LES THESMOPHORIES --- Traduction française (f. 170r, vv. 635–678 [3/4])}
\begin{adjustwidth}{-2cm}{-2cm}
\noindent\textbf{\footnotesize Texte grec établi --- Folio 16 (\texttt{f. 170r}, vv. 635–678 [3/4]) \hfill ΕΠΕΙΣΟΔΙΟΝ Β'}\\[-2pt]
\noindent\rule{16.5cm}{0.5pt}\\[2pt]
{\fontsize{8.8pt}{10.2pt}\selectfont
\ligneEst{\textbf{Κορυφαῖος} {\tiny\textit{\decalageN{4}$4$ an\cat}}}{ἡμᾶς τοίνυν μετὰ τοῦτ᾽ ἤδη τὰς λαμπάδας ἁψαμένας χρὴ}{655}
\ligneEst{}{ξυζωσαμένας εὖ κἀνδρείως τῶν θ᾽ ἱματίων ἀποδύσας}{}
\ligneEst{}{\Automatique{ζητεῖν, εἴ που κἄλλος τις ἀνὴρ ἐπελήλυθε, καὶ περιθρέξαι}}{}
\ligneEst{}{\Automatique{τὴν Πύκνα πᾶσαν καὶ τὰς σκηνὰς καὶ τὰς διόδους διαθρῆσαι.}}{}
\ligneEst{\textbf{Κορυφ. β} {\tiny\textit{$4$ troch.\cat}}}{εἶα δὴ πρώτιστα μὲν χρὴ κοῦφον ἐξορμᾶν πόδα}{}
\ligneEst{}{καὶ διασκοπεῖν σιωπῇ πανταχῇ· μόνον δὲ χρὴ}{660}
\ligneEst{}{μὴ βραδύνειν, ὡς ὁ καιρός ἐστι μὴ μέλλειν ἔτι,}{}
\ligneEst{}{\Automatique{ἀλλὰ τὴν πρώτην τρέχειν χρῆν ὡς τάχιστ᾽ ἤδη κύκλῳ.}}{}
\ligneEst{\textbf{Χορός} {\tiny\textit{$5$ troch.}}}{εἶά νυν ἴχνευε καὶ μά-}{}
\ligneEst{}{\alinea τευε ταχὺ πάντ᾽, εἴ τις ἐν τόποις ἑδραῖος}{}
\ligneEst{{\tiny\textit{lecyth.}}}{ἄλλος αὖ λέληθεν ὤν.}{}
\ligneEst{{\tiny\textit{$2$ troch.}}}{πανταχῇ δὲ ῥῖψον ὄμμα,}{665}
}
\vspace{2pt}
\noindent\rule{16.5cm}{0.6pt}\\[2pt]
\noindent\textbf{\small Commentaire philologique, métrique et scénique (vv. 635–678 --- partie 5/8)}\\[2pt]
\begin{adjustbox}{max totalheight=15.2cm, max width=16.5cm, keepaspectratio}%
\parbox[t]{16.5cm}{
\setlength{\columnsep}{9pt}%
\begin{multicols}{2}
\fontsize{7.6pt}{8.8pt}\selectfont
\setlength{\parindent}{0pt}%
\setlength{\parskip}{1.2pt}%
\noindent Ses conclusions erronées ont malheureusement été reproduites dans la plupart des travaux portant sur le costume de l'Ancienne comédie... espérons que cette note rectificatrice mettra enfin un terme à la contagion !
\noindent \textbf{(2)} Cf. par exemple \italique{Lysistrata}, vv. 599 sqq.
\noindent \textbf{(3)} Cf. la saillie du chaud-lapin Kinésias soumis à la torture par Myrrhine (\italique{Lysistrata}, v. 947) :
\citComm{\textbf{Μυρρίνη} -- λαβὲ τόνδε τὸν ἀλάβαστον. --- \textbf{Κινησίας} -- ἀλλ᾽ ἕτερον ἔχω.}
\noindent \textbf{(4)} \italique{Assemblée des femmes}, vers 741. Par jeu de mots avec l'ὄρθιος νόμος, le \italique{mode orthien}, qu'Hésychius (ο 1188) n'a pas manqué de nous définir :
\citComm{ὄρθιον νόμον· οὕτως ὠνομάσθαι φασὶν αὐτὸν ἀπὸ τοῦ ῥυθμοῦ, καθάπερ καὶ τὸν τροχαῖον.}
\noindent (Cf. aussi sur le mode orthien \italique{Cavaliers} 1280 ; \italique{Acharniens} 16.)\\ Les allusions d'Aristophane à cette tradition matinale à laquelle sacrifient ses personnages sont légions ; ne citons que \italique{Lysistrata} v. 966 puisque le passage est encore sous nos yeux.
\noindent \textbf{(5)} \italique{Lysistrata}, vers 1162-1163. Pour conclure la paix, l'ambassadeur laconien demande à récupérer l'\italique{encyle} :
\citComm{Λάκων -- ἁμές γε λῶμες, αἴ τις ἁμὶν τὤγκυκλον\\
λῇ τοῦτ᾽ ἀποδόμεν. --- Λυσιστράτη -- ποῖον ὦ τᾶν; --- Λάκων -- τὰν Πύλον,\\
ἇσπερ πάλαι δεόμεθα καὶ βλιμάττομες.}
\noindent \textbf{(6)} Chantraine, page 74 de sa \italique{Formation des noms en grec ancien}, précise : «~la finale -άριον qui a fourni essentiellement des diminutifs s'observe parfois dans des adjectifs, cf. σωμα παιδαριον dans des inscription de Delphes~».
\noindent \textbf{(7)} On lit en épigraphie : ἔνκυκλον περιποίκιλο, ἔγκυκλον ποικίλον. À l'inverse, on croise aussi un ἔνκυκλον λευκ[όν].
\noindent\textbf{644} \textbf{τάλαν.} Ce mot, le vocatif de τάλας, tant masculin que féminin, n'est mis par Aristophane que dans la bouche de femmes\marqueur{1}.
\noindent\textbf{[Notes ad 644]} \textbf{(1)} \italique{Lysistrata}, vers 102, 910 et 914 ; \italique{Grenouilles}, vers 559 ; \italique{Assemblée des femmes}, vers 124, 658 et 1005 ; \italique{Ploutos}, vers 706 et 1055.
\noindent\textbf{647-8} \textbf{ἰσθμόν τιν᾽ ἔχεις ὦνθρωπ᾽· ἄνω τε καὶ κάτω / τὸ πέος διέλκεις πυκνότερον Κορινθίων.} Clisthène fait référence au \grec{δίολκος} : un chemin dallé de 8 kilomètres qui permettait de faire traverser aux bateaux l'isthme de Corinthe grâce à une sorte de rail rudimentaire et d'éviter ainsi le long et périlleux contournement du Péloponnèse\marqueur{1}. Pendant la scène burlesque de déshabillage qui vient d'avoir lieu, Mnésiloque, pour cacher sa virilité – et certainement aussi émoustillé d'être effeuillé de la sorte –, a actionné à plusieurs reprises le mécanisme permettant de lever et baisser son \italique{phallus}. Admiratif devant de telles performances, Clisthène le flatte d'être encore mieux outillé que les machines qui permettaient le transbordement des bateaux.\\ Voilà pour l'image `τὸ πέος διέλκεις πυκνότερον Κορινθίων' ! il reste à expliquer la première partie de la réplique de Clisthène : ἰσθμόν τιν᾽ ἔχεις, «~tu possèdes un [drôle de] genre d'\italique{isthme}~».
\end{multicols}
}%
\end{adjustbox}%
\vfill
\end{adjustwidth}
\newpage
% --- PAGE DE DROITE : Traduction française + Commentaire (6/8) ---
\begin{adjustwidth}{-2cm}{-2cm}
\noindent\textbf{\footnotesize Traduction française --- Folio 16 (\texttt{f. 170r}, vv. 635–678 [3/4]) \hfill ΕΠΕΙΣΟΔΙΟΝ Β'}\\[-2pt]
\noindent\rule{16.5cm}{0.5pt}\\[2pt]
{\fontsize{8.8pt}{10.2pt}\selectfont
\ligneEst{\textbf{Chœur}}{\francais{À nous maintenant, après avoir allumé nos torches,}}{655}
\ligneEst{}{\francais{serré bien et résolument nos tuniques et mis bas nos manteaux,}}{}
\ligneEst{}{\francais{de chercher si d'aventure il n'est pas entré quelqu'autre homme. Tournons et virons}}{}
\ligneEst{}{\francais{tout à l'entour du Pnyx, inspectons avec soin les tentes et les passages.}}{}
\ligneEst{}{\francais{Or sus! tout d'abord élançons-nous au dehors d'un pied léger}}{}
\ligneEst{}{\francais{et épions partout en silence. Surtout ne tardons pas : le cas ne souffre plus de délai.}}{660}
\ligneEst{}{\francais{Il faut que notre coryphée coure au plus tôt faire la ronde.}}{}
\ligneEst{}{\francais{(Ici commence un chœur dansé.)}}{}
\ligneEst{}{\francais{Or sus ! suis chaque piste, quête vivement,}}{}
\ligneEst{}{\francais{au cas où quelqu'autre en ces lieux}}{}
\ligneEst{}{\francais{se serait tapi à notre insu.}}{}
\ligneEst{}{\francais{Jette les yeux partout.}}{665}
}
\vspace{2pt}
\noindent\rule{16.5cm}{0.6pt}\\[2pt]
\noindent\textbf{\small Commentaire philologique, métrique et scénique (vv. 635–678 --- partie 6/8)}\\[2pt]
\begin{adjustbox}{max totalheight=15.2cm, max width=16.5cm, keepaspectratio}%
\parbox[t]{16.5cm}{
\setlength{\columnsep}{9pt}%
\begin{multicols}{2}
\fontsize{7.6pt}{8.8pt}\selectfont
\setlength{\parindent}{0pt}%
\setlength{\parskip}{1.2pt}%
\noindent ἰσθμός désigne proprement un \italique{passage étroit} et, par analogie, le \italique{cou}. Le spectateur athénien saisissait instamment l'image : par \italique{isthme} il fallait entendre \italique{pénis} grâce à la matrice métaphorique suivante : ἰσθμός = αὐχήν\marqueur{2} = τράχηλος\marqueur{3} = πέος.
\noindent\textbf{[Notes ad 647-8]} \textbf{(1)} Voir à ce sujet l'article de M. J. T. Lewis, « Railways in the Greek and Roman world », in \italique{Early Railways : a Selection of Papers from the First International Early Railways Conference}, 2001, p. 8-19. Thucydide (III, 15 ; VIII, 7) évoque le δίολκος dans sa \italique{Guerre du Péloponnèse}.
\noindent \textbf{(2)} αὐχήν a le sens de πέος par exemple au vers 681 de \italique{Lysistrata} : le chœur des vieillards a maille à partir avec l'engeance féminine ; après avoir rappelé par une métaphore hippique combien la femme est portée à la luxure, il en vient à la conclusion que le meilleur moyen pour dominer ces fougueuses juments est de les chevaucher -- la fameuse \italique{phallocratie} ? Il formule cette idée à l'aide d'une image carcérale : il faut assujettir leur \italique{bois percé} (= \italique{carcan}) au \italique{cou} (αὐχήν) :
\citComm{ἀλλὰ τούτων χρῆν ἁπασῶν ἐς τετρημένον ξύλον\\
ἐγκαθαρμόσαι λαβόντας \textbf{τουτονὶ τὸν αὐχένα}.\\
Eh bien, il importe de s'assurer d'elles toutes, et d'assujettir dans un carcan \italique{ce cou-là}.}
\noindent On peut avoir un autre exemple d'αὐχήν avec ce sens-là au vers 184 du \italique{Cyclope} d'Euripide.
\noindent \textbf{(3)} τράχηλος est phallossémique au vers 20 des \italique{Grenouilles}. Aristophane commence la pièce avec un rapide survol des blagues lourdes et vulgaires de ses concurrents aux concours comiques ; ce n'est pas sa muse qui s'abaisserait à de si grossières manœuvres ! non, \italique{on n'en soufflera mot}, οὐκ ἐρεῖ ! οὐκ ἐρεῖ !.. et puis au vers 19 :
\citComm{ὢ τρισκακοδαίμων ἄρ᾽ \textbf{ὁ τράχηλος οὑτοσί},\\
ὅτι θλίβεται μέν, τὸ δὲ γέλοιον οὐκ ἐρεῖ.\\
Ô trois fois malencontreuse \italique{nuque}, qui est moulue, et n'a pas le mot pour rire !}
\noindent Il faut bien-sûr lire -- je dis «~bien-sûr~», mais personne ne l'avait proposé jusqu'ici... --, il faut donc lire :
\citComm{ἄρ᾽ ὁ τράχηλος οὑτοσί,\\
ὅτι θλίβεται μέν, τὸ δὲ γέλοιον \textbf{οὐ κε ῥεῖ;}\\
Est-ce que le τράχηλος que voici, puisqu'il est pressé, pour la farce ne devrait pas couler ?}
\end{multicols}
}%
\end{adjustbox}%
\vfill
\end{adjustwidth}
\newpage

% --- FOLIO 16 (f. 170r) : SPREAD 2.4 (Texte grec établi + Traduction française + Commentaire 7/8 & 8/8) ---
\clearpage\ifodd\value{page}\hbox{}\newpage\fi
\enTeteG{ΘΕΣΜΟΦΟΡΙΑΖΟΥΣΑΙ --- Texte grec établi (f. 170r, vv. 635–678 [4/4])}
\enTeteD{LES THESMOPHORIES --- Traduction française (f. 170r, vv. 635–678 [4/4])}
\begin{adjustwidth}{-2cm}{-2cm}
\noindent\textbf{\footnotesize Texte grec établi --- Folio 16 (\texttt{f. 170r}, vv. 635–678 [4/4]) \hfill ΕΠΕΙΣΟΔΙΟΝ Β'}\\[-2pt]
\noindent\rule{16.5cm}{0.5pt}\\[2pt]
{\fontsize{8.8pt}{10.2pt}\selectfont
\ligneEst{{\tiny\textit{$2$ troch.}}}{καὶ τὰ τῇδε καὶ τὰ δεῦρο}{}
\ligneEst{{\tiny\textit{lecyth.}}}{πάντ᾽ ἀνασκόπει καλῶς.}{}
\ligneEst{{\tiny\textit{$2$ anap.}}}{ἢν γάρ ληφθῇ δράσας ἀνόσια,}{}
\ligneEst{{\tiny\textit{$2$ anap.}}}{δώσει τε δίκην καὶ πρὸς τούτῳ}{}
\ligneEst{{\tiny\textit{$2$ anap\cat}}}{τοῖς ἄλλοις ἀνδράσιν ἔσται}{}
\ligneEst{{\tiny\textit{anap.}}}{παράδειγμ᾽ ὕβρεως}{670}
\ligneEst{{\tiny\textit{$2$ anap.}}}{ἀδίκων τ᾽ ἔργων ἀθέων τε τρόπων·}{}
\ligneEst{{\tiny\textit{$2$ anap.}}}{φήσει δ᾽ εἶναί τε θεοὺς φανερῶς,}{}
\ligneEst{{\tiny\textit{anap.}}}{δείξει τ᾽ ἤδη}{}
\ligneEst{{\tiny\textit{$3$ troch\cat}}}{πᾶσιν ἀνθρώποις σεβίζειν δαίμονας}{}
\ligneEst{{\tiny\textit{$2$ dochm. ?}}}{†δικαίως τ᾽ ἐφέποντας † ὅσια καὶ νόμιμα}{675}
\ligneEst{{\tiny\textit{$2$ dochm.}}}{μηδομένους ποιεῖν ὅ τι καλῶς ἔχει.}{}
\ligneEst{{\tiny\textit{$3$ iamb\cat}}}{κἂν μὴ ποιῶσι ταῦτα τοιάδ᾽ ἔσται·}{678}
}
\vspace{2pt}
\noindent\rule{16.5cm}{0.6pt}\\[2pt]
\noindent\textbf{\small Commentaire philologique, métrique et scénique (vv. 635–678 --- partie 7/8)}\\[2pt]
\begin{adjustbox}{max totalheight=15.2cm, max width=16.5cm, keepaspectratio}%
\parbox[t]{16.5cm}{
\setlength{\columnsep}{9pt}%
\begin{multicols}{2}
\fontsize{7.6pt}{8.8pt}\selectfont
\setlength{\parindent}{0pt}%
\setlength{\parskip}{1.2pt}%
\noindent\textbf{651} \textbf{εἰσεκύλισα.} Homère emploie le verbe κυλίνδειν en parlant de malheurs que le ciel \italique{fait rouler} sur quelqu'un : πῆμα θεὸς Δαναοῖσι κυλίνδει (\italique{Il}. 17, 688)\marqueur{1}. La même image se trouve chez Aristophane : νὴ τὸν Διόνυσον ἄπορά γ᾽ ἡμῖν πράγματα / δαίμων τις ἐσκεκύκληκεν ἐς τὴν οἰκίαν (\italique{Guêp.} 1474 s.), « ah! les inextricables tracas qu'un dieu a fait dévaler sur notre maison. » Le même verbe est employé d'une manière un peu différente, \italique{Thesm}. 767, mais toujours à propos de malheurs : ὁ μὲν γὰρ αἴτιος / κἄμ᾽ ἐσκυλίσας εἰς τοιαυτὶ πράγματα / οὐ φαίνεταί πω, « le coupable, celui qui m'a fait rouler dans de tels malheurs, n'apparaît toujours pas. » Cf. \italique{Thesm}. 650 s. : κακοδαίμων ἐγώ, / εἰς οἷ᾽ ἐμαυτὸν εἰσεκύλισα πράγματα. Il est clair que, dans ce cas, ἐμβάλλειν τινὰ εἰς συμφοράς n'est qu'un synonyme expressif\marqueur{2} de ἐμβάλλειν (p. ex. ἐμβάλλειν τινὰ εἰς συμφοράς, Antiphon 3, 4, 10 ; εἰς ἀτυχίας, Eschin. 3, 79 ; εἰς κακόν, Ar. \italique{Nuées} 1460). (Note de Taillardat, §558.)
\noindent\textbf{[Notes ad 651]} \textbf{(1)} Le moyen κυλίνδεσθαι se dit de malheur lui-même «~roulant~» sur quelqu'un : νῶιν δὴ τόδε πῆμα κυλίνδεσθαι, \italique{Il.}, 11, 347 ; cf. \italique{Od}. 2, 163 ; 8, 81.
\noindent \textbf{(2)} Connu d'autres auteurs ; Phérécr. 146 [K = 156 K.-A.] : εἰκῆ μ᾿ ἐπῆρας ὄντα τηλικουτονὶ / πολλοῖς ἐμαυτὸν ἐγκυλῖσαι πράγμασιν ; Cf. Xén. \italique{Mém.} I, 2, 22 : τοὺς εἰς ἔρωτας ἐγκυλισθέντας ; Ἀπολλόδωρος ὁ Καρύστιος 5, 7 [ $=$ Athénée 280e] : εἰκῆ τέ πως / ἡμᾶς κυλίνδους᾽ ὅντιν᾽ ἂν τύχῃ τρόπον (sc. τύχη ἄγροικος).
\noindent\textbf{655-86} \textbf{Le chœur compris entre les vv. 655-686.} La \nieme{19} note d'Alphonse Willems est d'une telle richesse d'explication que je ne peux que la reproduire intégralement ci-dessous. Pour ce qui touche à la colométrie, je reprends celle de Lætitia\marqueur{1}. (Le passage des vers 667-686 a de fortes correspondances métriques avec les vers 707-725).\\ Ce petit chœur n'a pas été examiné avec assez d'attention. On n'a pas remarqué qu'il s'accompagne d'un jeu de scène unique dans l'œuvre d'Aristophane. Trois morceaux le composent : \\ \ \ \ 1° Les vv. 655 à 658 appartiennent au coryphée. Mnésiloque n'est peut-être pas le seul homme qui ait pénétré parmi les femmes. Il importe de s'en assurer, et, après avoir déposé les manteaux et allumé les torches, de « tourner et virer tout à l'entour du Pnyx, et d'inspecter avec soin les tentes et les passages ».\\ \ \ \ 2° Les quatre suivants (659-662) sont des trochées, et le changement de mètre indique un changement d'interlocuteur. C'est, en effet, la réponse du \italique{parastate} ou chef du demi-chœur [le \italique{Ravennas} porte sur sa marge, η\auDessus{μ}{χ}ι γ\auDessus{υ}{ν}, écrite d'une seconde main, mais dont l'indication se retrouve sur le manuscrit de Münich]. En disant qu'il y avait lieu d'inspecter le Pnyx, les tentes et les passages, la coryphée indiquait clairement qu'il fallait quitter l'orchestre, c'est-à-dire l'enceinte sacrée adjacente au temple\marqueur{2}. La parastate annonce donc que le chœur va « s'élancer dehors (ἐξορμᾶν) d'un pied léger\marqueur{3} ». Puis elle engage celle qu'elle appelle τὴν πρώτην à commencer tout de suite la reconnaissance. Comment faut-il interpréter τὴν πρώτην ? D'y voir, comme on le fait généralement, une locution adverbiale (\latin{quam primum} ou \latin{prima quaque via}) me semble inadmissible. Avec M. Mazon, je crois que ce terme désigne une personne déterminée\marqueur{4} ; suivant moi, ce serait τὴν πρώτην ἱέρειαν, appelée aussi τὴν προστάτιν, savoir la première, il y en avait beaucoup\marqueur{5}, des prêtresses temporaires, ἀρήτειραι, de Déméter, celle qui, tout à l'heure, avait prononcé les prières d'usage (vv. 295 et suiv.), en un mot, la coryphée.
\end{multicols}
}%
\end{adjustbox}%
\vfill
\end{adjustwidth}
\newpage
% --- PAGE DE DROITE : Traduction française + Commentaire (8/8) ---
\begin{adjustwidth}{-2cm}{-2cm}
\noindent\textbf{\footnotesize Traduction française --- Folio 16 (\texttt{f. 170r}, vv. 635–678 [4/4]) \hfill ΕΠΕΙΣΟΔΙΟΝ Β'}\\[-2pt]
\noindent\rule{16.5cm}{0.5pt}\\[2pt]
{\fontsize{8.8pt}{10.2pt}\selectfont
\ligneEst{}{\francais{Que rien ici, rien là-bas,}}{}
\ligneEst{}{\francais{n'échappe à ton \oe il vigilant.}}{}
\ligneEst{}{\francais{Car s'il est pris après ce sacrilège,}}{}
\ligneEst{}{\francais{il paiera son crime, et de surcroît}}{}
\ligneEst{}{\francais{prémunira tous les autres,}}{}
\ligneEst{}{\francais{par son exemple, contre la violence,}}{670}
\ligneEst{}{\francais{les actes injurieux}}{}
\ligneEst{}{\francais{et les façons impies.}}{}
\ligneEst{}{\francais{Il confessera hautement qu'il est des dieux,}}{}
\ligneEst{}{\francais{il apprendra}}{}
\ligneEst{}{\francais{à tous les hommes à les honorer,}}{675}
\ligneEst{}{\francais{à s'efforcer, par une conduite pieuse et un constant}}{}
\ligneEst{}{\francais{souci des lois, de ne rien faire que de bien.}}{678}
}
\vspace{2pt}
\noindent\rule{16.5cm}{0.6pt}\\[2pt]
\noindent\textbf{\small Commentaire philologique, métrique et scénique (vv. 635–678 --- partie 8/8)}\\[2pt]
\begin{adjustbox}{max totalheight=15.2cm, max width=16.5cm, keepaspectratio}%
\parbox[t]{16.5cm}{
\setlength{\columnsep}{9pt}%
\begin{multicols}{2}
\fontsize{7.6pt}{8.8pt}\selectfont
\setlength{\parindent}{0pt}%
\setlength{\parskip}{1.2pt}%
\noindent Callimaque, parlant d'une de ces prêtresses, ἱερέη Δήμητρος, se sert d'une désignation analogue : πολλῶν προστασίη νέων γυναικῶν\marqueur{6}. \\ \ \ \ 3° Le morceau qui suit (vv. 663-684) est un \italique{hyporchême}. Une partie du chœur reste en place et chante, tandis que l'autre, en quête de quelque complice de Mnésiloque\marqueur{7}, évolue silencieusement, σιωπῇ (v. 660), va et vient, entre et sort, cherchant partout, dans l'orchestre et dehors, καὶ τὰ τῇδε καὶ τὰ δεῦρο (v. 666), jusqu'à ce qu'il déclare n'avoir rien trouvé (687-688).\\
\noindent\textbf{656} \textbf{τῶν θ᾽ ἱματίων ἀποδύσας.} Willems, fidèle à son poste de \italique{gardien de la tradition}, justifie contre les retoucheurs de texte l'expression ἀποδῦναι $+$ gén. (C'est sa \nieme{20} note aux \italique{Thesmophories}.)
\noindent\textbf{657} \textbf{εἰσελήλυθεν.} La leçon du manuscrit de Ravenne est \italique{contra metrum}. Tannegui Le Fèvre (Faber dans l'apparat), le père de Mme Dacier, avait proposé la simple correction ἐσελήλυθε. Austin\marqueur{1} la réfuta : ἐς ne précède jamais une voyelle chez Aristophane. Restent les conjectures de Fritzsche et de Handley. Fritzsche proposa de lire ἀνελήλυθε, l'expliquant par une haplographie ἀνὴρ (ἀν)ἐλήλυθε suivie d'une complétion postérieure εἰσ-.
\noindent Handley proposa ἐπελήλυθε (la confusion entre ΙC et Π en onciales étant une erreur courante) -- hypothèse beaucoup plus simple et, partant, plus satisfaisante.
\noindent\textbf{[Notes ad 655-86]} \textbf{(1)} \italique{The songs of Aristophanes}, pp. 418-422.
\noindent \textbf{(2)} On peut conclure de là que le Thesmophorion, qu'on savait situé sur une hauteur, l'était sur la colline du Pnyx et voisin du lieu d'assemblée.
\noindent \textbf{(3)} Le simple, ὀρμᾶν, qui se lit plus loin, v. 953, et \italique{Gren}., 478, signifierait « prendre sa course ».
\noindent \textbf{(4)} \italique{Essai sur la composition des comédies d'Aristophane}, Paris (1904), p. 131, n. 4.
\noindent \textbf{(5)} Il y en avait deux par dème (Isée, \italique{Succession de Kiron}, 19-20), donc plus de deux cents.
\noindent \textbf{(6)} \italique{Anthol. Palat}., VII, 728.
\noindent \textbf{(7)} εἴ τις ἐν τόποις ἑδραῖος. Qu'avons-nous affaire d'ajouter, avec MM. Blaydes et van Leeuwen, τούτοισι τοῖς devant τόποις ? La même locution adverbiale se lit dans un fragment de Nausicratès : [τρίγκλαι] ἃς κλύδων Αἰξωνικὸς ἐν τόποις παιδεύεται (\italique{Fragm. Com. Gr}., t. IV, p. 575 [= Athénée 325e = Nausicratès, fr. 1.6-11 Kassel-Austin ; le texte n'est pas établi avec certitude : on trouve aussi : ἐντόπους à la place de ἐν τόποις. Même si l'exemple n'est pas probant, Willems n'en a pas moins raison ; Lætitia remarque que l'adjonction de Blaydes accroît la disparité entre les deux stances des vers 663-666 -- là où l'on chercherait à rétablir une symétrie.]) et la locution inverse se lit dans Sophocle : ἥτις μ᾽ ἀποικίσειεν ἐκ τόπων (\italique{Trach}., 955), où jamais personne n'a trouvé à redire (ἐκ τούτων τῶν τόπων.). Je ne vois pas que ces tournures soient plus bizarres que ἐν ποσί et ἐκ ποδῶν, ou bien que ἐν χώρᾳ et ἐκ χώρας, πρὸ ou παρὰ χειρῶν, etc.
\noindent \textbf{(1)} Voir ses notes : « Textual problems in Ar. \italique{Thesm}. » (1987) ; je reprends la trame de celle qui nous intéresse présentement.
\end{multicols}
}%
\end{adjustbox}%
\vfill
\end{adjustwidth}
\newpage
\clearpage\ifodd\value{page}\hbox{}\newpage\fi
\enTeteG{ΕΠΕΙΣΟΔΙΟΝ Γ' --- Codex Ravennas 429, f. 170v (Classense 172v) --- vv. 679–720}
\begin{adjustwidth}{-1.5cm}{-3.0cm}
\begin{adjustbox}{max totalheight=21.0cm, max width=16.5cm, keepaspectratio}
\begin{minipage}{16.5cm}
{\fontsize{7.2pt}{8.5pt}\selectfont
\noindent\textbf{Folio 17 --- Codex Ravennas 429, \texttt{f. 170v} (\texttt{Classense 172v}) : ΕΠΕΙΣΟΔΙΟΝ Γ', vv. 679–720}\hfill \textit{Témoins :} $\text{R}^1, \Rpc, \Rdeux$ (\Mon)\par\vspace{2pt}
\hrule height 0.4pt\vspace{2.5pt}
}
{\fontsize{8.1pt}{9.5pt}\selectfont
\ligneRav{679}{χορ/}{\hspace{1.1em}αὐτῶν ὅταν ληφθῇ τις †ὅσια δρῶν†,}{\textbf{679} \lecon{αὐτῶν} R : τούτων Richards \espApp \lecon{ὅσια} R : ἀνόσια Zanetti : ὅσια μὴ Burges : ἀνόσιόν τι Hermann}
\ligneRav{680}{}{\hspace{1.1em}μανίαις φλέγων λύσσῃ παράκο-}{}
\ligneRav{}{}{\hspace{1.1em} πος, εἴ τι δρῴη πᾶσιν ἐμφανὴς ὁρᾶν}{}
\ligneRav{683}{}{\hspace{1.1em}\difference{ἔσται} γυναιξὶ καὶ βροτοῖς,}{\textbf{683} ἐστὶν R : \lecon{ἔσται} Bothe \espApp γυναιξὶν : \lecon{-ξὶ} Brunck}
\ligneRav{}{}{\hspace{1.1em}ὅτι τά παράνομα τά τ᾽ ἀνόσια παρὼν}{}
\ligneRav{686}{}{\hspace{1.1em}θεὸς \difference{παραχρῆμ᾽ ἀποτίνεται}·}{\textbf{686} ἀποτίνεται· | παραχρῆμά τε τίνεται R : \lecon{παραχρῆμ᾽ ἀποτίνεται} Hermann}
\ligneRav{687}{χορ/}{ἀλλ᾽ ἔοιχ᾽ ἡμῖν ἅπαντά πως διεσκέφθαι καλῶς·}{\textbf{687} (Κορ.) add. von Velsen \espApp Κορ. add. Velsen \espApp \lecon{πως} R : που Blaydes}
\ligneRav{}{}{οὐχ ὁρῶμεν γοῦν ἔτ᾽ ἄλλον οὐδέν᾽ ἐγκαθήμενον}{}
\ligneRav{689}{\auDessus{γυ}{ν}}{\difference{ἆ ἆ}}{\textbf{689} ἀά supra lin. R : \lecon{ἆ ἆ} Thiersch \espApp ποῖ R : \lecon{ποῖ ποῖ} Biset}
\ligneRav{689}{}{\difference{ποῖ ποῖ} σὺ φεύγεις; οὗτος οὗτος οὐ μενεῖς;}{\textbf{689} ἀά supra lin. R : \lecon{ἆ ἆ} Thiersch \espApp ποῖ R : \lecon{ποῖ ποῖ} Biset}
\ligneRav{690}{}{τάλαιν᾽ ἐγὼ τάλαινα, καὶ τὸ παιδίον}{}
\ligneRav{}{}{ἐξαρπάσας μοι φροῦδος ἀπὸ τοῦ τιτθίου·}{}
\ligneRav{}{μ\auDessus{ν}{Η}}{κέκραχθι· τοῦτο δ᾽ οὐδέποτε σὺ ψωμιεῖς,}{}
\ligneRav{693}{}{τάδε δ᾽ \difference{ἀφῆτ᾽} (ἀφε\auDessus{ῖ}{ῆ}τ᾽ \Rdeux) ἤδη τὸ παιδίον;}{\textbf{693} ἀφε\auDessus{ῖ}{ῆ}τ᾽ \Rdeux : ἀφεῖτ᾽ \Mon}
\ligneRav{}{}{πληγὲν μαχαίρᾳ τῇδε φοινίας φλέβας}{}
\ligneRav{695}{—}{καθαιματώσει βωμόν. : ὢ τάλαιν᾽ ἐγώ·}{}
\ligneRav{}{}{γυναῖκες, οὐκ ἀρήξετ᾽; οὐ πολλὴν βοὴν}{}
\ligneRav{697}{}{στήσεσθε καὶ τροπαῖον, ἀλλὰ τοῦ μόνου}{\textbf{697} καὶ om. R : καὶ suppl. Σ ad \italique{Pl.} v.453 et Greg. Cor. \italique{de dialectis} p.22}
\ligneRav{}{—}{τέκνου με περιόψεσθ᾽ ἀποστερουμένην; : ἔα ἔα·}{}
\ligneRav{700}{}{\hspace{1.1em}ὦ πότνιαι Μοῖραι τί \difference{τόδε} δέρκομαι}{\textbf{700} δὲ R : \lecon{τόδε} Dobree : δὴ Reisig}
\ligneRav{}{}{\hspace{1.1em}νεοχμὸν αὖ τέρας;}{}
\ligneRav{702}{}{ὡς \difference{ἅπαντ᾽ ἄρ᾽ ἐστὶ} τόλμης †\difference{πλέα}† κἀναισχυντίας·}{\textbf{702} ἅπαν γὰρ ἐστὶν R : ἅπαντα Suid., cod. Laur. : \lecon{ἅπαντ᾽ ἄρ᾽ ἐστὶ} Bentley \espApp ἔργα R : \lecon{πλέα} Suid. : δοῦλα cod. Laur. : μεστὰ Porson}
\ligneRav{}{}{οἷον αὖ δέδρακεν ἔργον, οἷον αὖ φίλαι τόδε·}{}
\ligneRav{704}{μ\auDessus{ν}{Η}}{\difference{οἷον}; ὑμῶν \difference{ἐξαρήξω} τὴν ἄγαν αὐθαδίαν·}{\textbf{704} οἷον R : \lecon{οἷον;} Ellebodius \espApp ἐξάρξω R : ἐξαράξει Ellebodius : ἐξαράξω Bentley : \lecon{ἐξαρήξω} Bibi}
\ligneRav{705}{χορ/}{ταῦτα δῆτ᾽ οὐ δεινὰ πράγματ᾽ ἐστὶ καὶ περαιτέρω;}{}
\ligneRav{706}{\auDessus{γυ}{ν}}{δεινὰ δῆθ᾽, \difference{ὁτιή} γ᾽ ἔχει μου 'ξαρπάσας τὸ παιδίον·}{\textbf{706} ὅτι R : \lecon{ὁτιή} Faber : ὅστις Porson}
\ligneRav{}{χορ/}{\hspace{1.1em}τί ἂν οὖν εἴποι πρὸς ταῦτά τις, ὅτε}{}
\ligneRav{}{}{\hspace{1.1em}τοιαῦτα ποιῶν ὅδ᾽ ἀναισχυντεῖ;}{}
\ligneRav{}{μ\auDessus{ν}{Η}}{\hspace{1.1em}κοὔπω μέντοι γε πέπαυμαι·}{}
\ligneRav{710}{χορ/}{\hspace{1.1em}ἀλλ᾽ οὖν \difference{ἥκειν γ᾽} ὅθεν ἥκεις·}{\textbf{710} ἥκεις τ᾿ R : ἥκεις γ᾽ Suid. η 170 : ἥκων γ᾽ Willems : \lecon{ἥκειν γ᾽} Bibi ap. Phot. η 88}
\ligneRav{711}{}{\hspace{1.1em}φαύλως \difference{δ᾿} ἀποδρὰς \difference{οὐκ ἔξει}·}{\textbf{711} τ᾿ R, Suid : \lecon{δ᾿} Bentley \espApp οὐ λέξεις R, Suid. : \lecon{οὐκ ἔξει.} Bibi ap. Σ}
\ligneRav{}{}{\hspace{1.1em}οἷον δράσας διέδυς ἔργον·}{}
\ligneRav{}{}{\hspace{1.1em}λήψει δὲ κακόν·}{}
\ligneRav{}{μ\auDessus{ν}{Η}}{τοῦτο μέντοι μὴ γένοιτο μηδαμῶς, ἀπεύχομαι·}{}
\ligneRav{715}{χορ/}{\hspace{1.1em}τίς οὖν σοι, τίς ἂν σύμμαχος ἐκ θεῶν}{\textbf{715} \lecon{οὖν σοι} R : ἂν σοι Hamaker}
\ligneRav{}{}{\hspace{1.1em}ἀθανάτων ἔλθοι ξὺν ἀδίκοις ἔργοις;}{}
\ligneRav{717}{μ\auDessus{ν}{Η}}{μάτην λαλεῖτε· τὴνδ᾽ ἐγὼ οὐκ ἀφήσω·}{\textbf{717} τὴν δ᾽ R : \lecon{τήνδ᾽} Biset}
\ligneRav{}{χορ/}{\hspace{1.1em}ἀλλ᾽ οὐ μὰ τὼ θεὼ τάχ᾽ οὐ}{}
\ligneRav{719}{}{\hspace{1.1em}χαίρων ἴσως \difference{ἐνυβριεῖς}}{\textbf{719} ἐνυβρίσεις R : \lecon{ἐνυβριεῖς} anonymus Romanus}
\ligneRav{720}{}{\hspace{1.1em}λόγους τε λέξεις ἀνοσίους}{}
}
\vspace{3pt}\noindent\rule{4.5cm}{0.4pt}\par\vspace{1.5pt}
{\fontsize{7.0pt}{8.3pt}\selectfont
\noindent\textbf{Loci paralleli (vv. 679–720) ---} \textbf{696-7} (γυναῖκες -- τροπαῖον) \GregCor p.20-22 cf. schol.\expo{Ald} \italique{Plutus} 453c \espTest \textbf{702} (ἅπαντ' -- κἀναισχυντίας) Suidas α 2902 (sed ἅπαντα pro ἅπαν γὰρ ἐστὶν et πλέα pro ἔργα) ; \italique{Corpus paroemiographorum Graecorum}, Appendices, I, 35 (sed ἅπαντα pro ἅπαν γὰρ ἐστὶν et ἔργα omisso) ; \italique{Laurentiamus} LVIII 24 p. 40 (ed. Cohn) (sed ἅπαντα pro ἅπαν γὰρ ἐστὶν et δοῦλα pro ἔργα) \espTest \textbf{710-3} Suidas η 170 (sed λήψῃ pro λήψει) \espTest \textbf{710} Photius η 88 ἥκειν· ἀντὶ τοῦ ἥκων\par
}
\vspace{2.5pt}\noindent\rule{4.5cm}{0.4pt}\par\vspace{1.5pt}
{\fontsize{6.9pt}{8.2pt}\selectfont
\noindent\textbf{Suite (en commentaire) de l'apparat critique \& notes paléographiques (\texttt{f. 170v}) ---} \textbf{Page particulièrement riche en corruptions d'onciale et de minuscule élucidées par Louis Hauseux (\textbf{Bibi}) :} --- \textbf{v. 679 :} \lecon{αὐτῶν} R (τούτων Richards) \espApp \lecon{ὅσια} R (maintenu par L. Hauseux avec interrogation ; ἀνόσια Zanetti ; ὅσια μὴ Burges ; ἀνόσιόν τι Hermann). --- \textbf{v. 697 :} καὶ om. R : \lecon{καὶ} suppl. $\Sigma$ ad \italique{Plut.} 453 et Grégoire de Corinthe, \italique{De dialectis}, p. 22. --- \textbf{v. 704 :} οἷον R : \lecon{οἷον;} Ellebodius \espApp ἐξάρξω R (ἐξαράξει Ellebodius ; ἐξαράξω Bentley) : \lecon{ἐξαρήξω} \lecon{Bibi} ap. Hésychius $\varepsilon\ 3599$. --- \textbf{v. 710 :} ἥκεις τ' R (ἥκεις γ' Souda $\eta\ 170$ ; ἥκων γ' Willems) : \lecon{ἥκειν γ'} \lecon{Bibi} ap. Photius $\eta\ 88$ (permutation minuscule $\nu \to \varsigma$). --- \textbf{v. 711 :} τ' R, Souda : \lecon{δ'} Bentley \espApp οὐ λέξεις R, Souda : \lecon{οὐκ ἔξει.} \lecon{Bibi} ap. $\Sigma^{\text{R}}$ (ὅθεν οὐ φεύξῃ, confusion onciale Κ/Λ).\par
}
\end{minipage}
\end{adjustbox}
\end{adjustwidth}
\newpage
\enTeteD{Codex Ravennas 429 (Biblioteca Classense) --- Fac-similé du f. 170v (vv. 679–720)}
\begin{adjustwidth}{-1.5cm}{-3.0cm}
\begin{center}
\includegraphics[width=16.5cm,height=20.5cm,keepaspectratio]{/workspaces/Thesmophories/dossier_ravennas_thesmophories/folios_individuels_ff162v-177r/17_f170v_classense172v_vv0679-0720.jpg}\par
\vspace{2pt}
{\scriptsize\textbf{Codex Ravennas 429, f. 170v (Classense 172v)} --- ΕΠΕΙΣΟΔΙΟΝ Γ', \textbf{vv. 679–720} ($\text{R}^1, \Rpc, \Rdeux$ (\Mon)).}
\end{center}
\end{adjustwidth}
\newpage
\clearpage\ifodd\value{page}\hbox{}\newpage\fi% --- FOLIO 17 (f. 170v) : SPREAD 2.1 (Texte grec établi + Traduction française + Commentaire 1/8 & 2/8) ---
\clearpage\ifodd\value{page}\hbox{}\newpage\fi
\enTeteG{ΘΕΣΜΟΦΟΡΙΑΖΟΥΣΑΙ --- Texte grec établi (f. 170v, vv. 679–720 [1/4])}
\enTeteD{LES THESMOPHORIES --- Traduction française (f. 170v, vv. 679–720 [1/4])}
\begin{adjustwidth}{-2cm}{-2cm}
\noindent\textbf{\footnotesize Texte grec établi --- Folio 17 (\texttt{f. 170v}, vv. 679–720 [1/4]) \hfill ΕΠΕΙΣΟΔΙΟΝ Γ'}\\[-2pt]
\noindent\rule{16.5cm}{0.5pt}\\[2pt]
{\fontsize{8.8pt}{10.2pt}\selectfont
\ligneEst{{\tiny\textit{$2$ iamb...}}}{αὐτῶν ὅταν ληφθῇ τις †ὅσια δρῶν†,}{679}
\ligneEst{{\tiny\textit{$5$ iamb.}}}{μανίαις φλέγων λύσσῃ παράκο-}{680}
\ligneEst{}{\alinea πος, εἴ τι δρῴη πᾶσιν ἐμφανὴς ὁρᾶν}{}
\ligneEst{{\tiny\textit{$2$ iamb.}}}{ἔσται γυναιξὶ καὶ βροτοῖς,}{}
\ligneEst{{\tiny\textit{$2$ iamb.}}}{ὅτι τά παράνομα τά τ᾽ ἀνόσια παρὼν}{}
\ligneEst{{\tiny\textit{dochm...}}}{θεὸς παραχρῆμ᾽ ἀποτίνεται.}{}
\ligneEst{\textbf{Κορυφαῖος} {\tiny\textit{tetr. tr\cat}}}{ἀλλ᾽ ἔοιχ᾽ ἡμῖν ἅπαντά πως διεσκέφθαι καλῶς.}{}
\ligneEst{}{οὐχ ὁρῶμεν γοῦν ἔτ᾽ ἄλλον οὐδέν᾽ ἐγκαθήμενον}{}
\ligneEst{\textbf{Μίκα} {\tiny\textit{extr. metr.}}}{ἆ ἆ}{}
\ligneEst{{\tiny\textit{trim. iamb.}}}{ποῖ ποῖ σὺ φεύγεις; οὗτος οὗτος οὐ μενεῖς;}{}
}
\vspace{2pt}
\noindent\rule{16.5cm}{0.6pt}\\[2pt]
\noindent\textbf{\small Commentaire philologique, métrique et scénique (vv. 679–720 --- partie 1/8)}\\[2pt]
\begin{adjustbox}{max totalheight=15.2cm, max width=16.5cm, keepaspectratio}%
\parbox[t]{16.5cm}{
\setlength{\columnsep}{9pt}%
\begin{multicols}{2}
\fontsize{8.4pt}{9.7pt}\selectfont
\setlength{\parindent}{0pt}%
\setlength{\parskip}{1.2pt}%
\noindent\textbf{680} \textbf{μανίαις φλέγων λύσσῃ παράκοπος.} Παράκοπος, \italique{délirant}, est une métaphore usée\marqueur{1} et d'origine incertaine. Elle paraît sans rapport avec la frappe de la monnaie\marqueur{2} ; mais elle en a peut-être un avec la notion d'« ébranler le cerveau »\marqueur{3}, cf. Euripide \italique{Hippolyte}, v. 238, etc... παρακόπτειν φρένας. (Note de Taillardat §476.)
\noindent\textbf{[Notes ad 680]} \textbf{(1)} Eschyle \italique{Prom.} 581 ; Eur. \italique{Bacch.} 33, 1000, etc.
\noindent \textbf{(2)} Malgré παρακόπτειν (\italique{Ach.} 517), « battre de la fausse monnaie ».
\noindent \textbf{(3)} Cf. §474 du même livre : \italique{Les images d'Aristophane} de Jean Taillardat.
\noindent\textbf{683} \textbf{γυναιξὶ καὶ βροτοῖς.} L'on s'attendait à la formule consacrée : θεοῖσι καὶ βροτοῖς\marqueur{1}.
\noindent\textbf{[Notes ad 683]} \textbf{(1)} Homère, \italique{Iliade} XXIV, 553 ; Pindare, \italique{Pythiques} 3.30 ; Sophocle, \italique{Électre} v. 199.
\noindent\textbf{689-761} \textbf{Parodie du \italique{Télèphe} d'Euripide.} Comme dans ses \italique{Acharniens} représentés une quinzaine d'années plus tôt -- il y a de nombreuses analogies entre ces deux pièces --, Aristophane détourne le \italique{Télèphe} d'Euripide. On ne sait pas grand'chose de cette dernière pièce, la majorité des fragments conservés étant eux-mêmes tirés d'Aristophane et de ses anciens commentateurs\marqueur{1}.
\noindent Télèphe était un héros populaire puisque les trois grands tragiques lui ont consacré pas moins de sept pièces ; il fut blessé par un coup de lance d'Achille.
\noindent Comme sa blessure ne guérissait pas, il finit par consulter l'oracle d'Apollon, réponse de ce dernier : ὁ τρώσας ἰάσεται, « Qui t'a blessé te guérira ». Il se rend donc, déguisé en mendiant, dans le camp de l'armée achéenne préparant son expédition sur Troie. Allait-il d'abord trouver Clytemnestre qui lui confiait son petit Oreste pour appuyer son plaidoyer ? toujours est-il que Télèphe se retrouva la proie exposée à la colère des chefs grecs, il se réfugiait alors en suppliant sur l'autel en prenant en otage Oreste. Le sort de l'enfant était plus heureux que dans sa thesmophorienne parodie : grâce à l'entremise d'Agamemnon, Achille se laissait convaincre de soigner Télèphe et tout était bien qui finissait bien. Outre cette scène d'enlèvement du petit Oreste, de nombreux autres passages des \italique{Thesmophories} sont sans doute aussi inspirés du \italique{Télèphe} (v. 39-70 : où l'altercation entre le héros et un gardien attirait Clytemnestre à l'extérieur et permettait sa rencontre avec Télèphe ; v. 466-470 : lorsqu'au début de son plaidoyer, le héros -- sous une fausse identité -- usait de précautions oratoires ; v. 598-600, 659-667 avec l'annonce de la présence d'un espion dans le camp grec.)
\end{multicols}
}%
\end{adjustbox}%
\vfill
\end{adjustwidth}
\newpage
% --- PAGE DE DROITE : Traduction française + Commentaire (2/8) ---
\begin{adjustwidth}{-2cm}{-2cm}
\noindent\textbf{\footnotesize Traduction française --- Folio 17 (\texttt{f. 170v}, vv. 679–720 [1/4]) \hfill ΕΠΕΙΣΟΔΙΟΝ Γ'}\\[-2pt]
\noindent\rule{16.5cm}{0.5pt}\\[2pt]
{\fontsize{8.8pt}{10.2pt}\selectfont
\ligneEst{}{\francais{S'ils se comportent autrement, voici ce qui se passera :}}{679}
\ligneEst{}{\francais{celui d'eux qui sera pris à faire un sacrilège,}}{680}
\ligneEst{}{\francais{enflammé de fureur, délirant de rage,}}{}
\ligneEst{}{\francais{témoignera aux femmes et à tous les mortels}}{}
\ligneEst{}{\francais{que la divinité châtie les atteintes aux lois}}{}
\ligneEst{}{\francais{et à la piété et ne se fait pas attendre.}}{}
\ligneEst{\textbf{Coryphée}}{\francais{Mais il nous semble avoir bien cherché partout et nous ne voyons pas d'autre homme caché parmi nous.}}{}
\ligneEst{}{\francais{(Mnésiloque enlève l'enfant et se réfugie sur l'autel.)}}{}
\ligneEst{\textbf{Mika}}{\francais{Ah ! Ah !}}{}
\ligneEst{}{\francais{Où ? Où te sauves-tu ? holà ! holà ! Arrête.}}{}
}
\vspace{2pt}
\noindent\rule{16.5cm}{0.6pt}\\[2pt]
\noindent\textbf{\small Commentaire philologique, métrique et scénique (vv. 679–720 --- partie 2/8)}\\[2pt]
\begin{adjustbox}{max totalheight=15.2cm, max width=16.5cm, keepaspectratio}%
\parbox[t]{16.5cm}{
\setlength{\columnsep}{9pt}%
\begin{multicols}{2}
\fontsize{8.4pt}{9.7pt}\selectfont
\setlength{\parindent}{0pt}%
\setlength{\parskip}{1.2pt}%
\noindent\textbf{[Notes ad 689-761]} \textbf{(1)} Pour une présentation détaillée du \italique{Télèphe}, le lecteur pourra consulter l'édition de François Jouan \& Herman van Looy aux Belles Lettres, tome VIII, \nieme{3} partie.
\noindent\textbf{693-5} \textbf{ἐπὶ τῶν μηρίων ... βωμόν.} Vers tragiques, très sûrement empruntés à Euripide. (Note de Willems.)
\noindent\textbf{696-7} \textbf{οὐ πολλὴν βοὴν / στήσεσθε καὶ τροπαῖον.} Dans le sens de μήστωρ φόβοιο d'Homère. « Un vacarme à faire sauver ». (Musset, \italique{Confess}., p. 111 ; note de Willems.)
\noindent\textbf{702} \textbf{ὡς ἅπαν γὰρ ἐστὶν τόλμης ἔργα κἀναισχυντίας.} Le scholiaste met en avant le tour proverbial de la formule en commentant d'un sobre «~παροιμία~». Heureusement, d'autres glosateurs furent davantage prolixes.
\noindent On retrouve la formule chez Suidas (α 2902) :
\citComm{ἅπαντα τόλμης πλέα κἀναισχυντίας: ἐπὶ τῶν ἰταμῶν καὶ ἀναιδῶν.}
\noindent Et dans le \italique{Corpus paroemiographorum Graecorum}\marqueur{1} :
\citComm{ἅπαντα τόλμης καὶ ἀναισχυντίας: ὅτι ἡ ἀναισχυντία νικᾷ τὰ πολλά.}
\noindent Enfin dans le codex \italique{Laurentiamus} LVIII 24\marqueur{2}, l'on tombe sur la formulation suivante :
\citComm{ἅπαντα τόλμης δοῦλα καὶ ἀναισχυντίας.}
\noindent Ces trois leçons sont évidemment tirées du vers 702 des \italique{Thesmophories} et des anciens commentaires afférents.
\noindent Le début du vers dans le \italique{Ravennas}, ὡς ἅπαν γὰρ ... ἔργα κτλ., est fautif. La correction la plus simple à apporter est de lire : ὡς ἅπαντ᾽ ἄρ᾽ κτλ.
\noindent Il n'y aurait alors qu'une simple confusion d'un Τ lu en Γ : ΑΠΑΝΤΑΡ $ \to $ ΑΠΑΝΓΑΡ.
\noindent Cette correction mineure fait concorder le texte d'Aristophane avec les parémiographes.
\noindent Les variantes ne différent plus que sur un seul mot : ἔργα (dans le \italique{Ravennas}) ; πλέα (\italique{Souda}) ; δοῦλα (Laurent.) ; $\emptyset$ (\italique{Corpus paroemiographorum}).
\end{multicols}
}%
\end{adjustbox}%
\vfill
\end{adjustwidth}
\newpage

% --- FOLIO 17 (f. 170v) : SPREAD 2.2 (Texte grec établi + Traduction française + Commentaire 3/8 & 4/8) ---
\clearpage\ifodd\value{page}\hbox{}\newpage\fi
\enTeteG{ΘΕΣΜΟΦΟΡΙΑΖΟΥΣΑΙ --- Texte grec établi (f. 170v, vv. 679–720 [2/4])}
\enTeteD{LES THESMOPHORIES --- Traduction française (f. 170v, vv. 679–720 [2/4])}
\begin{adjustwidth}{-2cm}{-2cm}
\noindent\textbf{\footnotesize Texte grec établi --- Folio 17 (\texttt{f. 170v}, vv. 679–720 [2/4]) \hfill ΕΠΕΙΣΟΔΙΟΝ Γ'}\\[-2pt]
\noindent\rule{16.5cm}{0.5pt}\\[2pt]
{\fontsize{8.8pt}{10.2pt}\selectfont
\ligneEst{}{τάλαιν᾽ ἐγὼ τάλαινα, καὶ τὸ παιδίον}{690}
\ligneEst{}{ἐξαρπάσας μοι φροῦδος ἀπὸ τοῦ τιτθίου.}{}
\ligneEst{\textbf{Μνησίλοχος}}{κέκραχθι· τοῦτο δ᾽ οὐδέποτε σὺ ψωμιεῖς,}{}
\ligneEst{}{ἢν μή μ᾽ ἀφῆτ᾽· ἀλλ᾽ ἐνθάδ᾽ ἐπὶ τῶν μηρίων}{}
\ligneEst{}{πληγὲν μαχαίρᾳ τῇδε φοινίας φλέβας}{}
\ligneEst{}{καθαιματώσει βωμόν.}{695}
\ligneEst{\textbf{Μίκα}}{\decalage{καθαιματώσει βωμόν.}ὢ τάλαιν᾽ ἐγώ.}{}
\ligneEst{}{γυναῖκες, οὐκ ἀρήξετ᾽; οὐ πολλὴν βοὴν}{}
\ligneEst{}{στήσεσθε καὶ τροπαῖον, ἀλλὰ τοῦ μόνου}{}
\ligneEst{}{τέκνου με περιόψεσθ᾽ ἀποστερουμένην;}{}
\ligneEst{\textbf{Χορός} {\tiny\textit{προαναφώνημα}}}{\decalage{aa}ἔα ἔα.}{}
}
\vspace{2pt}
\noindent\rule{16.5cm}{0.6pt}\\[2pt]
\noindent\textbf{\small Commentaire philologique, métrique et scénique (vv. 679–720 --- partie 3/8)}\\[2pt]
\begin{adjustbox}{max totalheight=15.2cm, max width=16.5cm, keepaspectratio}%
\parbox[t]{16.5cm}{
\setlength{\columnsep}{9pt}%
\begin{multicols}{2}
\fontsize{8.4pt}{9.7pt}\selectfont
\setlength{\parindent}{0pt}%
\setlength{\parskip}{1.2pt}%
\noindent Rejetons d'emblée la leçon du \italique{Ravennas}, qui sent son interpolation ; ἔργα n'est qu'une glose d'ἅπαντα qui aura pénétré le texte.
\noindent La même erreur se produit au vers 717 des \italique{Oiseaux} : les manuscrits comportent `πρὸς ἅπαντα τρέπεσθε' ; une glose du scholiaste substitue ἔργα à ἅπαντα et Suidas (α 1638) finit par citer le vers ainsi : `πρὸς ἔργα τράπεσθε'.
\noindent La leçon de la \italique{Souda}, πλέα, est satisfaisante quant au sens (tout \italique{plein} d'effronterie et d'audace) ; elle n'est cependant pas métriquement correcte, raison pour laquelle Porson lui substitua μεστά.
\noindent À défaut d'explication satisfaisante, j'ai préféré gardé la leçon fautive de Suidas, non sans la frapper de l'obèle du désespoir (†).
\noindent Finissons cette note aporétique en remarquant que les termes τόλμη et ἀναισχυντία sont souvent mis en collocation : Platon\marqueur{3} et Démosthène\marqueur{4} associent les deux mots avec un sens analogue. Isocrate nous apporte quelque précision, l'on comprend en le lisant\marqueur{5} que l'audace (τόλμα) pouvait être interprétée comme un signe (σημεῖον) d'impudence (ἀναισχυντία). Les Grecs avaient fait d'\italique{Implacabilité} un dieu : θεὸς ἡ Ἀναίδεια\marqueur{6}.
\noindent\textbf{[Notes ad 702]} \textbf{(1)} \italique{Corpus paroemiographorum Graecorum}, vol. I, Appendices, I, 35, p. 384 ($=$ Vat. K II 36).
\noindent \textbf{(2)} Il s'agit d'un manuscrit conservé à la Biblioteca Medicea Laurenziana (BML) de Florence, dépôt Plutarque, sous la cote 58. 24  et édité par Leopold Cohn en 1887 : \italique{Zu den Paroemoigraphen} (notre proverbe y figure à la page 40).
\noindent \textbf{(3)} Platon, \italique{Apologie de Socrate}, 38d. Voir aussi son \italique{Théétète} (196d) où le tour proverbial est encore plus sensible : ὅμως δέ — πάντα γὰρ τολμητέον — τί εἰ ἐπιχειρήσαιμεν ἀναισχυντεῖν; «~De l'audace ! encore de l'audace ! toujours de l'audace !~»
\noindent \textbf{(4)} Παραγραφὴ πρὸς Ναυσίμακον καὶ Ξενοπείθην, 5 (Νόμος).
\noindent \textbf{(5)} Περὶ ἀντιδόσεως, 190 : ἔτι δὲ τὴν τόλμαν μὴ τὴν ἀναισχυντίας σημεῖον γιγνομένην, ἀλλὰ τὴν μετὰ σωφροσύνης. Voir aussi Démosthène, \italique{In Aristogitonem 1}, 68.8 : ταῦτα τολμᾷ λέγειν δι' ἀναισχυντίαν.
\end{multicols}
}%
\end{adjustbox}%
\vfill
\end{adjustwidth}
\newpage
% --- PAGE DE DROITE : Traduction française + Commentaire (4/8) ---
\begin{adjustwidth}{-2cm}{-2cm}
\noindent\textbf{\footnotesize Traduction française --- Folio 17 (\texttt{f. 170v}, vv. 679–720 [2/4]) \hfill ΕΠΕΙΣΟΔΙΟΝ Γ'}\\[-2pt]
\noindent\rule{16.5cm}{0.5pt}\\[2pt]
{\fontsize{8.8pt}{10.2pt}\selectfont
\ligneEst{}{\francais{Malheur ! Ah ! malheur ! le voilà parti avec mon enfant}}{690}
\ligneEst{}{\francais{arraché de ma mamelle.}}{}
\ligneEst{\textbf{Mnésiloque}}{\francais{Hurle. Jamais plus tu ne l'appateras,}}{}
\ligneEst{}{\francais{si vous ne me laissez pas libre. Mais ici même « sur ces oblations,}}{}
\ligneEst{}{\francais{frappé de ce couteau, du pur sang de ses veines}}{}
\ligneEst{}{\francais{il rougira l'autel ».}}{695}
\ligneEst{\textbf{Mika}}{\francais{\espace{3cm}Malheur à moi !}}{}
\ligneEst{}{\francais{Femmes, au secours ! Poussez une puissante clameur,}}{}
\ligneEst{}{\francais{une clameur à faire lâcher pied. Souffrirez-vous}}{}
\ligneEst{}{\francais{qu'on me ravisse mon unique enfant?}}{}
\ligneEst{\textbf{Chœur}}{\francais{Ah! Ah !}}{}
}
\vspace{2pt}
\noindent\rule{16.5cm}{0.6pt}\\[2pt]
\noindent\textbf{\small Commentaire philologique, métrique et scénique (vv. 679–720 --- partie 4/8)}\\[2pt]
\begin{adjustbox}{max totalheight=15.2cm, max width=16.5cm, keepaspectratio}%
\parbox[t]{16.5cm}{
\setlength{\columnsep}{9pt}%
\begin{multicols}{2}
\fontsize{8.4pt}{9.7pt}\selectfont
\setlength{\parindent}{0pt}%
\setlength{\parskip}{1.2pt}%
\noindent \textbf{(6)} \italique{Implacabilité} et non \italique{Impudence}. Cf. le \italique{Dictionnaire des antiquités grecques et romaines} (Daremberg \& Saglio), t. I, p. 398b (\latin{Areopagus}). Les Athéniens avaient dressé des autels à Ὕβρις et Ἀναίδεια dans l'Aréopage, contre lesquels s'était insurgé ce rabat-joie de Cicéron -- non sans faire un contre-sens comme le souligne l'auteur de l'article. Dans sa querelle avec Agamemnon, l'attitude d'Achille -- qui par bien des aspects est celle d'un adolescent attardé -- nous révèle en fait son côté divin : il est \italique{implacable}, tout autant que le destin qu'il connaît et a accepté. Pour revenir au proverbe `θεὸς ἡ Ἀναίδεια', certainement inspiré par ces autels de l'Aréopage, voir Zénobius (IV, 36). Érasme l'a aussi recensé dans ses \italique{Adages} (n° 2063).
\noindent\textbf{704} \textbf{οἷον ὑμῶν ἐξάρξω τὴν ἄγαν αὐθαδίαν.} Il manque une syllabe à ἐξάρξω, la leçon du manuscrit de Ravenne. Les différents éditeurs y suppléent avec le verbe ἐξαράττειν (ἐξαράξω au futur ou, en s'éloignant davantage de la leçon -- quitte à corriger le texte... --, ἐξαράξει, pour faire de οἷον le sujet). La correction est légère et pourtant elle ne manque pas de laisser dubitatif : si l'image de \italique{briser} (ἐξαράξω) l'\italique{audace} (αὐθαδίαν) de \italique{quelqu'un} (ὑμῶν) ne paraît aucunement étrangère à une imagination française, il n'en allait peut-être pas de même pour les Grecs. L'incontournable Taillardat\marqueur{1} l'avait déjà noté : considérer l'arrogance «~comme un objet matériel qu'on peut briser [...] est une image unique~». En revanche, en grec comme en français, l'image de \italique{chasser}, \italique{pousser} hors (ἐξ-), \italique{repousser} un sentiment ou un trait de caractère est courante. Ainsi, l'on chassera l'audace par l'audace\marqueur{2} :
\citComm{Ἀντιφάνης δέ· οἴνῳ .... τὸν οἶνον ἐξελαύνειν, [...] τριωβόλῳ δὲ πόρνην, αὐθαδίαν αὐθαδίᾳ.}
\noindent L'on se débarrassera de son chagrin\marqueur{3} :
\citComm{ἐκκρούει τὴν λύπην}
\noindent comme les vieillards cagneux secouent leurs ennuis chez Aristophane\marqueur{4} :
\citComm{ἀποσείονται δὲ λύπας.}
\end{multicols}
}%
\end{adjustbox}%
\vfill
\end{adjustwidth}
\newpage

% --- FOLIO 17 (f. 170v) : SPREAD 2.3 (Texte grec établi + Traduction française + Commentaire 5/8 & 6/8) ---
\clearpage\ifodd\value{page}\hbox{}\newpage\fi
\enTeteG{ΘΕΣΜΟΦΟΡΙΑΖΟΥΣΑΙ --- Texte grec établi (f. 170v, vv. 679–720 [3/4])}
\enTeteD{LES THESMOPHORIES --- Traduction française (f. 170v, vv. 679–720 [3/4])}
\begin{adjustwidth}{-2cm}{-2cm}
\noindent\textbf{\footnotesize Texte grec établi --- Folio 17 (\texttt{f. 170v}, vv. 679–720 [3/4]) \hfill ΕΠΕΙΣΟΔΙΟΝ Γ'}\\[-2pt]
\noindent\rule{16.5cm}{0.5pt}\\[2pt]
{\fontsize{8.8pt}{10.2pt}\selectfont
\ligneEst{{\tiny\textit{2 dochm.}}}{ὦ πότνιαι Μοῖραι τί τόδε δέρκομαι}{700}
\ligneEst{{\tiny\textit{dochm.}}}{νεοχμὸν αὖ τέρας;}{}
\ligneEst{{\tiny\textit{$4$ troch\cat}}}{\Automatique{ὡς ἅπαντ᾽ ἄρ᾽ ἐστὶ τόλμης †πλέα† κἀναισχυντίας.}}{}
\ligneEst{}{οἷον αὖ δέδρακεν ἔργον, οἷον αὖ φίλαι τόδε.}{}
\ligneEst{\textbf{Μνησίλοχος}}{οἷον; ὑμῶν ἐξαρήξω τὴν ἄγαν αὐθαδίαν.}{}
\ligneEst{\textbf{Χορός}}{\Automatique{ταῦτα δῆτ᾽ οὐ δεινὰ πράγματ᾽ ἐστὶ καὶ περαιτέρω;}}{705}
\ligneEst{\textbf{Μίκα}}{δεινὰ δῆθ᾽, ὁτιή γ᾽ ἔχει μου 'ξαρπάσας τὸ παιδίον.}{}
\ligneEst{\textbf{Χορός} {\tiny\textit{dim. anap.}}}{τί ἂν οὖν εἴποι πρὸς ταῦτά τις, ὅτε}{}
\ligneEst{}{τοιαῦτα ποιῶν ὅδ᾽ ἀναισχυντεῖ;}{}
\ligneEst{\textbf{Μνησίλοχος} {\tiny\textit{paroem.}}}{κοὔπω μέντοι γε πέπαυμαι.}{}
}
\vspace{2pt}
\noindent\rule{16.5cm}{0.6pt}\\[2pt]
\noindent\textbf{\small Commentaire philologique, métrique et scénique (vv. 679–720 --- partie 5/8)}\\[2pt]
\begin{adjustbox}{max totalheight=15.2cm, max width=16.5cm, keepaspectratio}%
\parbox[t]{16.5cm}{
\setlength{\columnsep}{9pt}%
\begin{multicols}{2}
\fontsize{7.6pt}{8.8pt}\selectfont
\setlength{\parindent}{0pt}%
\setlength{\parskip}{1.2pt}%
\noindent Euripide met dans la bouche du fourbe Polydectès qui veut s'accaparer Danaé que\marqueur{5} :
\citComm{πατέρα τε παισὶν ἡδέως συνεκφέρειν\\
φίλους ἔρωτας \textbf{ἐκβαλόντ' αὐθαδίαν}.\\
Un père doit supporter avec indulgence les affaires amoureuses de son fils et \italique{repousser l'entêtement}.}
\noindent L'occasion fait le larron : profitons-en pour placer cet aphorisme rapporté par Athénée\marqueur{6} illustrant le sens de la mesure grec, même quand il s'agit de boire...
\citComm{Πάσας δ᾽ ἐκ κραδίας ἀνίας ἀνδρῶν ἀλαπάζει\\
πινόμενος κατὰ μέτρον· ὑπὲρ μέτρον δὲ χερείων.\\
\francais{Le vin est souverain pour chasser les soucis du cœur de qui n'en boit pas trop, mais l'excès nuit.}}
\noindent (Un franc-buveur espiègle répondrait que la mesure du boire est de boire sans mesure\marqueur{7} !..) Le verbe ἀρήγω (futur : ἀρήξω) signifie \italique{secourir} et, par suite, \italique{écarter}, \italique{repousser}. Il conviendrait tant pour le sens que paléographiquement s'il était préfixé de ἐξ-.
\noindent Comme pour nous conforter, Hésychius (ε 3599) nous atteste justement l'existence du composé :
\citComm{ἐξαρῆξαι· ἐκφορῆσαι.}
\noindent Lisons donc :
\citComm{ὑμῶν ἐξαρήξω τὴν ἄγαν αὐθαδίαν\\
Je vous extirperai votre trop-plein d'audace.}
\noindent P.-S. : l'idée de recourir aux matrices métaphoriques pour le bon établissement d'un texte revient à l'immense Jean Taillardat. Je ne peux que renvoyer à son coruscant article « Images et matrices métaphoriques »\marqueur{8}. Il le termine par ce vœu :
\citComm{En guise de conclusion, je formulerai un souhait, celui qu'on puisse, un jour, consulter un dictionnaire des matrices métaphoriques pour les langues indo-européennes vivantes (langues populaires et argots compris) et mortes. On l'imaginerait volontiers comme une sorte de dictionnaire analogique, mais réduit aux images, les notions à exprimer (par exemple « colère », « inconsistance », « jeunesse », etc.) étant classées en ordre alphabétique ; quel merveilleux outil serait ainsi offert à qui s'intéresse à la recherche étymologique et combien de tâtonnements, d'erreurs même, seraient alors évités ! Mais quelle équipe aura assez d'abnégation et de courage pour se lancer jamais dans cette énorme entreprise ?}
\end{multicols}
}%
\end{adjustbox}%
\vfill
\end{adjustwidth}
\newpage
% --- PAGE DE DROITE : Traduction française + Commentaire (6/8) ---
\begin{adjustwidth}{-2cm}{-2cm}
\noindent\textbf{\footnotesize Traduction française --- Folio 17 (\texttt{f. 170v}, vv. 679–720 [3/4]) \hfill ΕΠΕΙΣΟΔΙΟΝ Γ'}\\[-2pt]
\noindent\rule{16.5cm}{0.5pt}\\[2pt]
{\fontsize{8.8pt}{10.2pt}\selectfont
\ligneEst{}{\francais{Augustes M\oe res, quel prodige}}{700}
\ligneEst{}{\francais{nouveau vient frapper mes yeux !}}{}
\ligneEst{}{\francais{Tout n'est donc qu'audace doublée d'effronterie !}}{}
\ligneEst{}{\francais{Quel forfait encore, amies, quel forfait il a commis !}}{}
\ligneEst{\textbf{Mnésiloque}}{\francais{Quel forfait ? je vous extirperai votre outrecuidance !}}{}
\ligneEst{\textbf{Chœur}}{\francais{N'est-ce pas une abomination, et pis encore ?}}{705}
\ligneEst{\textbf{Mika}}{\francais{Une abomination, oui, qu'on m'ait arraché mon enfant.}}{}
\ligneEst{\textbf{Chœur}}{\francais{Que dire après cela, quand}}{}
\ligneEst{}{\francais{l'auteur d'une telle action a perdu toute honte ?}}{}
\ligneEst{\textbf{Mnésiloque}}{\francais{Et je n'ai pas fini.}}{}
}
\vspace{2pt}
\noindent\rule{16.5cm}{0.6pt}\\[2pt]
\noindent\textbf{\small Commentaire philologique, métrique et scénique (vv. 679–720 --- partie 6/8)}\\[2pt]
\begin{adjustbox}{max totalheight=15.2cm, max width=16.5cm, keepaspectratio}%
\parbox[t]{16.5cm}{
\setlength{\columnsep}{9pt}%
\begin{multicols}{2}
\fontsize{8.4pt}{9.7pt}\selectfont
\setlength{\parindent}{0pt}%
\setlength{\parskip}{1.2pt}%
\noindent Puisse ce grand et beau dictionnaire un jour voir le jour !
\noindent\textbf{[Notes ad 704]} \textbf{(1)} §628 de ses \italique{Images d'Aristophane}, ouvrage dont on ne soulignera jamais assez l'importance pour toute personne désireuse d'avoir un avant-goût du génie de la langue aristophanesque.
\noindent \textbf{(2)} Athénée, \italique{Deipnosophistes}, II, 44a.
\noindent \textbf{(3)} Aristote, \italique{Éthique à Nicomaque, 1154a27}.
\noindent \textbf{(4)} \italique{Grenouilles}, v. 346.
\noindent \textbf{(5)} Fragment du \italique{Dictys}. Stobée, IV, 26, 16 $=$ fr. 7 \italique{Dictys} Jouan \& van Looy $=$ fr. 339 Kannicht $=$ fr. 433 Mette.
\noindent \textbf{(6)} Athénée, II, 37b.
\noindent \textbf{(7)} Aristophane, \italique{Guêpes}, vers 79-80 :
\citComm{εἶναι φιλοπότην αὐτόν. [ ... ]\\
αὕτη γε χρηστῶν ἐστιν ἀνδρῶν ἡ νόσος.}.
\noindent \textbf{(8)} \italique{Bulletin de l'Association Guillaume Budé} (décembre 1977), Lettre d'humanité n° 36, pp. 344-354.
\noindent\textbf{707-25} \textbf{} Ce passage, sans correspondre parfaitement du point de vue métrique à celui des vers 667-686, présente néanmoins de troublantes ressemblances avec ce dernier. Se reporter à Lætitia\marqueur{1} pour les parallélismes qu'on peut faire avec le texte tel qu'il nous a été conservé.
\noindent\textbf{[Notes ad 707-25]} \textbf{(1)} \italique{The songs of Aristophanes}, pp. 421 sqq. Je reprends pour une bonne part sa colométrie.
\noindent\textbf{710-2} \textbf{ἀλλ᾽ οὖν ἥκεις τ᾿ ὅθεν ἥκεις.} Je reproduis ci-dessous la \nieme{21} note d'Alphonse Willems. La correction qu'il y propose fut sauvée du néant de l'oubli et adoptée par la plupart des éditeurs modernes grâce à la diligente entremise de Coulon :
\noindent «~Ce passage a été refait de cinq ou six façons, plus bizarres les unes que les autres.
\end{multicols}
}%
\end{adjustbox}%
\vfill
\end{adjustwidth}
\newpage

% --- FOLIO 17 (f. 170v) : SPREAD 2.4 (Texte grec établi + Traduction française + Commentaire 7/8 & 8/8) ---
\clearpage\ifodd\value{page}\hbox{}\newpage\fi
\enTeteG{ΘΕΣΜΟΦΟΡΙΑΖΟΥΣΑΙ --- Texte grec établi (f. 170v, vv. 679–720 [4/4])}
\enTeteD{LES THESMOPHORIES --- Traduction française (f. 170v, vv. 679–720 [4/4])}
\begin{adjustwidth}{-2cm}{-2cm}
\noindent\textbf{\footnotesize Texte grec établi --- Folio 17 (\texttt{f. 170v}, vv. 679–720 [4/4]) \hfill ΕΠΕΙΣΟΔΙΟΝ Γ'}\\[-2pt]
\noindent\rule{16.5cm}{0.5pt}\\[2pt]
{\fontsize{8.8pt}{10.2pt}\selectfont
\ligneEst{\textbf{Χορός} {\tiny\textit{paroem.}}}{ἀλλ᾽ οὖν ἥκειν γ᾽ ὅθεν ἥκεις.}{710}
\ligneEst{{\tiny\textit{paroem.}}}{φαύλως δ᾿ ἀποδρὰς οὐκ ἔξει.}{}
\ligneEst{{\tiny\textit{2 anap.}}}{οἷον δράσας διέδυς ἔργον·}{}
\ligneEst{{\tiny\textit{anap.}}}{λήψει δὲ κακόν.}{}
\ligneEst{\textbf{Μνησίλοχος} {\tiny\textit{$4$ troch\cat}}}{τοῦτο μέντοι μὴ γένοιτο μηδαμῶς, ἀπεύχομαι.}{}
\ligneEst{\textbf{Χορός} {\tiny\textit{2 dochm.}}}{τίς οὖν σοι, τίς ἂν σύμμαχος ἐκ θεῶν}{715}
\ligneEst{}{ἀθανάτων ἔλθοι ξὺν ἀδίκοις ἔργοις;}{}
\ligneEst{\textbf{Μνησίλοχος} {\tiny\textit{trim. ia\cat}}}{μάτην λαλεῖτε· τὴνδ᾽ ἐγὼ οὐκ ἀφήσω.}{}
\ligneEst{\textbf{Χορός} {\tiny\textit{2 iamb.}}}{ἀλλ᾽ οὐ μὰ τὼ θεὼ τάχ᾽ οὐ}{}
\ligneEst{}{χαίρων ἴσως ἐνυβριεῖς}{}
\ligneEst{}{λόγους τε λέξεις ἀνοσίους}{720}
}
\vspace{2pt}
\noindent\rule{16.5cm}{0.6pt}\\[2pt]
\noindent\textbf{\small Commentaire philologique, métrique et scénique (vv. 679–720 --- partie 7/8)}\\[2pt]
\begin{adjustbox}{max totalheight=15.2cm, max width=16.5cm, keepaspectratio}%
\parbox[t]{16.5cm}{
\setlength{\columnsep}{9pt}%
\begin{multicols}{2}
\fontsize{8.4pt}{9.7pt}\selectfont
\setlength{\parindent}{0pt}%
\setlength{\parskip}{1.2pt}%
\noindent La correction me paraît élémentaire. Les manuscrits concordent, sauf sur le premier ἥκεις (ἥκεις γ' dans Suidas, ἥκεις τ' dans le \italique{Ravennas}, ἥκισθ' dans le manuscrit de Munich [qui n'est qu'une copie du manuscrit de Ravenne, comportant quelques émendations parfois bien senties]). Lisez ἥκων γ', le sens devient parfait :
\citComm{ἀλλ᾽ οὖν ἥκων γ᾽ ὅθεν ἥκεις}
\noindent Il est à peine besoin d'ajouter que ἥκων est pris dans le sens de ἀνήκων, \italique{redux}, ce qui est des plus fréquents\marqueur{1}. Et quant à ἀλλ᾽ οὖν..γε, chacun sait qu'il signifie \italique{toujours}, \italique{tout au moins}, \italique{en tous cas}\marqueur{2}.\\
\noindent C'est le mot d'Épicharme, fr. 126\marqueur{3} : κἀπῆνθεν ὅθεν ἦνθεν πάλιν ; ou de Ménandre (t. IV, p. 211 [$=$ fr. 373 K.-A.]) : ἀπῆλθεν ὅθεν ἦλθεν ταχύ. Cf. aussi Soph., \italique{Œd. Col.}, 1227 : βῆναι κεῖσ' ὁπόθεν περ ἥκει ; Eur., \italique{Her. F.}, 1247 : θανών, ὅθενπερ ἦλθον, εἶμι γῆς ὕπο, etc.~»
\noindent Le chœur enjoint Mnésiloque de cesser cette escalade de l'impudence : qu'il quitte au moins l'autel et les rejoigne\marqueur{4}.\\ À la lumière de la démonstration de Willems, je me permets d'ajouter la glose suivante du \latin{Lexicon} de Photius (η 88) :
\citComm{ἥκειν· ἀντὶ τοῦ ἥκων.}
\noindent Supposer que la glose de Photius se rapporte à notre vers offre l'avantage de rester plus proche du texte des manuscrits : il y aurait eu simple permutation d'un ν par un ς.
\noindent L'on peut donc lire après Willems :
\citComm{ἀλλ᾽ οὖν ἥκειν γ᾽ ὅθεν ἥκεις.\\
\francais{Au moins, de revenir d'où tu es parti.}}
\noindent\textbf{[Notes ad 710-2]} \textbf{(1)} \italique{Paix}, 265, 275; \italique{Lysistr}., 595, 731, 924 ; \italique{Gren}., 503 ; \italique{Ploutos}, 521, 1179, etc.
\noindent \textbf{(2)} \italique{Nuées}, 985, 1002 ; \italique{Guêpes}, 1129, 1190 ; Esch., \italique{Prom}., 1058 ; Soph., \italique{El}., 233 ; \italique{Ant}., 84, etc. M'éloignant de Willems quant au sens à donner à ce vers, je ne retranscris pas ici la traduction qu'il en propose. Le lecteur la trouvera en annexe où toutes les notes d'Alphonse Willems sur les \italique{Thesmophories} sont intégralement reproduites.
\end{multicols}
}%
\end{adjustbox}%
\vfill
\end{adjustwidth}
\newpage
% --- PAGE DE DROITE : Traduction française + Commentaire (8/8) ---
\begin{adjustwidth}{-2cm}{-2cm}
\noindent\textbf{\footnotesize Traduction française --- Folio 17 (\texttt{f. 170v}, vv. 679–720 [4/4]) \hfill ΕΠΕΙΣΟΔΙΟΝ Γ'}\\[-2pt]
\noindent\rule{16.5cm}{0.5pt}\\[2pt]
{\fontsize{8.8pt}{10.2pt}\selectfont
\ligneEst{\textbf{Chœur}}{\francais{Au moins, reviens d'où tu es parti ; ce n'est pas}}{710}
\ligneEst{}{\francais{avec une fuite facile que tu t'en sortiras.}}{}
\ligneEst{}{\francais{Avec tout ce que tu as fait en t'immisçant ici,}}{}
\ligneEst{}{\francais{malheur t'en prendra !}}{}
\ligneEst{\textbf{Mnésiloque}}{\francais{Puisse-t-il n'en être rien, j'en conjure le ciel.}}{}
\ligneEst{\textbf{Chœur}}{\francais{Et qui donc, qui des dieux immortels s'associera}}{715}
\ligneEst{}{\francais{à des actes injurieux ?}}{}
\ligneEst{\textbf{Mnésiloque}}{\francais{Inutiles bavardages ! Je ne la lâcherai pas.}}{}
\ligneEst{\textbf{Chœur}}{\francais{Ah! par les deux déesses, peut-être}}{}
\ligneEst{}{\francais{n'auras-tu pas lieu de rire de tes outrages}}{}
\ligneEst{}{\francais{et des discours insensés que tu profères.}}{720}
}
\vspace{2pt}
\noindent\rule{16.5cm}{0.6pt}\\[2pt]
\noindent\textbf{\small Commentaire philologique, métrique et scénique (vv. 679–720 --- partie 8/8)}\\[2pt]
\begin{adjustbox}{max totalheight=15.2cm, max width=16.5cm, keepaspectratio}%
\parbox[t]{16.5cm}{
\setlength{\columnsep}{9pt}%
\begin{multicols}{2}
\fontsize{8.4pt}{9.7pt}\selectfont
\setlength{\parindent}{0pt}%
\setlength{\parskip}{1.2pt}%
\noindent \textbf{(3)} $[=$ Plutarque, \italique{Moralia}, Παραμυθητικὸς πρὸς Ἀπολλώνιον, 110a $=$ fr. 213(245) Epich. Kassel-Austin, \italique{Poet. Com. Gr.}, I, \italique{Comoedia Dorica}.$]$
\noindent \textbf{(4)} C'est l'explication de la scholie aussi conservée par Suidas (η 170) : ἥκεις φησίν εἰς τὰς χεῖρας ἡμῶν, ὅθεν οὐ φεύξῃ.
\noindent\textbf{711-2} \textbf{φαύλως τ᾽ ἀποδρὰς οὐ λέξεις / οἷον δράσας διέδυς ἔργον.} Où se trouve le complément réclamé par λέγειν ? Je ne suis pas sûr qu'il soit de bonne grécité de construire la phrase ainsi  : λέξεις οἷον δράσας, «~tu raconteras tout ce que tu as fait~». (Sauf erreur de ma part, le grec aurait recouru à un verbe conjugué et non au participe de δράω.)
\noindent Une légère retouche du texte permettrait une bien meilleure lecture, tant pour la syntaxe grecque qu'au niveau du sens.
\noindent Lisez :
\citComm{φαύλως γ᾽ ἀποδρὰς οὐκ ἔξει.\\
\francais{Ce n'est pas avec une fuite facile que tu t'en iras !}}
\noindent (Confusion entre un Κ et un Λ.) Le vers 712, οἷον δράσας διέδυς ἔργον, est en fait apposé au vers 713 et en explique la sentence : λήψει δὲ κακόν, « pour tout ce que tu as fait, tu recevras châtiment ».
\noindent Pour justifier cette correction, je m'appuierai sur un débris de scholie qui se trouve à l'intérieur même du texte du \italique{Ravennas}, et que j'interprète ainsi :
\citComm{ἥκεις, φησίν, εἰς τὰς χεῖρας ἡμῶν, ὅθεν οὐ φεύξῃ.\\
Tu [re]viens, dit-elle, dans nos bras, et d'où tu ne fuiras pas.}
\noindent Comprise ainsi, cette scholie est clairement une paraphrase du vers 710 pour la première partie (ἥκεις εἰς τὰς χεῖρας ἡμῶν $=$ ἥκειν γ᾽ ὅθεν ἥκεις). C'est la seconde partie de la scholie qui nous intéresse : ὅθεν οὐ φεύξῃ, «~d'où tu ne fuiras pas~». Peut-on espérer meilleure glose que `οὐ φεύξῃ' pour le `οὐκ ἔξει' que nous rétablissons ?
\end{multicols}
}%
\end{adjustbox}%
\vfill
\end{adjustwidth}
\newpage
\clearpage\ifodd\value{page}\hbox{}\newpage\fi
\enTeteG{ΕΠΕΙΣΟΔΙΟΝ Γ' --- Codex Ravennas 429, f. 171r (Classense 173r) --- vv. 721–764}
\begin{adjustwidth}{-1.5cm}{-3.0cm}
\begin{adjustbox}{max totalheight=21.0cm, max width=16.5cm, keepaspectratio}
\begin{minipage}{16.5cm}
{\fontsize{7.2pt}{8.5pt}\selectfont
\noindent\textbf{Folio 18 --- Codex Ravennas 429, \texttt{f. 171r} (\texttt{Classense 173r}) : ΕΠΕΙΣΟΔΙΟΝ Γ', vv. 721–764}\hfill \textit{Témoins :} $\text{R}^1, \Rpc, \Rdeux$, \PsiDeux\par\vspace{2pt}
\hrule height 0.4pt\vspace{2.5pt}
}
{\fontsize{8.1pt}{9.5pt}\selectfont
\ligneRav{721}{χορ/}{\hspace{1.1em}ξὺν ἀθέοις ἔργοις}{\textbf{721} ἀθέοις ἔργοις γὰρ R : ἀθ. γὰρ ἔργ. Burges : ἐπ' ἀθ. ἔρ. γὰρ Enger : \lecon{ξὺν ἀθ. ἔργ. γὰρ} Bibi ap. Taillardat cf. v.716 \espApp γὰρ R : \lecon{καὶ γὰρ} Hermann}
\ligneRav{}{}{καὶ γὰρ ἀνταμειψόμεσθά σ᾽ ὥσπερ εἰκὸς ἀντὶ τῶνδε·}{}
\ligneRav{724}{}{\hspace{1.1em}\difference{τάχα δὲ} μεταβαλοῦσ᾽ ἐπὶ κακὸν ἑτερότρο-}{\textbf{724} τάχα δέ σε R : \lecon{τάχα δὲ} Fritzsche}
\ligneRav{725}{}{\hspace{1.1em} πον ἐπέχει \difference{τύχη}·}{\textbf{725} τις τύχη R : \lecon{τύχη} Bergk}
\ligneRav{726}{χορ/}{ἀλλὰ τάσδε μὲν λαβεῖν χρῆν σ᾽ ἐκφέρειν τε τῶν ξύλων,}{\textbf{726} \lecon{\expo{χ}ρῆν σ᾽} \Rdeux : χρὴ μ᾽ \Mon}
\ligneRav{}{}{καὶ καταίθειν τὸν πανοῦργον πυρπολεῖν θ᾽ ὅσον τάχος·}{}
\ligneRav{}{\auDessus{γυ}{ν}}{ἴωμεν ἐπὶ τὰς κληματίδας ὦ Μανία·}{}
\ligneRav{729}{}{κἀγώ σ᾽ ἀποδείξω θυμάλωπα τήμερον·}{\textbf{729} \lecon{κἀγώ} R : ἦ ᾽γώ Sommerstein}
\ligneRav{730}{μ\auDessus{ν}{Η}}{ὕφαπτε καὶ κάταιθε· \difference{σὺ δὲ τὸ} Κρητικὸν}{\textbf{730} σὺ τόδε τὸ R : \lecon{σὺ δὲ τὸ} Suid. κ 2408}
\ligneRav{}{}{ἀπόδυθι ταχέως· τοῦ θανάτου δ᾽ ὦ παιδίον}{}
\ligneRav{}{}{μόνην γυναικῶν αἰτιῶ τὴν μητέρα·}{}
\ligneRav{733}{}{τουτὶ τί ἔστιν; ἀσκὸς \difference{ἐγένεθ᾽} ἡ κόρη}{\textbf{733} ἐγένηθ᾽ R : \lecon{ἐγένεθ᾽} Gryaneaus}
\ligneRav{}{}{οἴνου πλέως καὶ ταῦτα Περσικὰς ἔχων·}{}
\ligneRav{735}{}{ὦ θερμόταται γυναῖκες, ὦ ποτίσταται}{\textbf{735} \lecon{θερμόταται} R, Suid. θ 251}
\ligneRav{}{}{κἀκ παντὸς ὑμεῖς μηχανώμεναι πιεῖν,}{}
\ligneRav{}{}{ὦ μέγα καπήλοις ἀγαθὸν ἡμῖν δ᾽ αὖ κακόν,}{}
\ligneRav{}{}{κακὸν δὲ καὶ τοῖς σκευαρίοις καὶ τῇ κρόκῃ·}{}
\ligneRav{}{\auDessus{γυ}{ν}}{παράβαλλε πολλὰς κληματίδας ὦ Μανία·}{}
\ligneRav{740}{μ\auDessus{ν}{Η}}{παράβαλλε δῆτα· σὺ δ᾽ \difference{ἀπόκριναί} μοι τοδί,}{\textbf{740} ἀπόκρινέ R : \lecon{ἀπόκριναί} Μυ2}
\ligneRav{741}{—}{τουτὶ τεκεῖν φῄς; : \difference{καὶ δέκα μῆνας} αὔτ᾽ ἐγὼ}{\textbf{741} δέκα μῆνας R : \lecon{καὶ δέκα μῆνας} Biset : δέκα μῆνας γε Scaliger}
\ligneRav{}{—}{ἤνεγκον. : ἤνεγκας σύ; : νὴ τὴν Ἄρτεμιν·}{}
\ligneRav{743}{μ\auDessus{ν}{Η}}{τρικότυλον; ἢ πῶς. : εἰπέ μοι· τί μ᾽ \difference{ἠργάσω};}{\textbf{743} \_τρικο\douteux{τ}\douteux{υ}\douteux{λ}\douteux{ο}ν·η\douteux{π}\douteux{ω}\douteux{c} \  ειπ\douteux{ε}\douteux{μ}οι\douteux{τ}[ \PsiDeux : τρικότυλον ἢ πῶς; εἰπέ μοι --- τί μ᾽ εἰργάσω; R : \lecon{ἠργάσω} Hall-Geldart}
\ligneRav{}{}{ἀπέδυσας ὦναίσχυντέ μου τὸ παιδίον}{}
\ligneRav{745}{—}{\difference{\difference{τυννοῦτο}ν} ὄν. : τυννοῦτο; μικρὸν νὴ Δία·}{\textbf{745} \lecon{τυννοῦτον} et \lecon{τυννοῦτο} \PsiDeux, Brunck : τυνοῦτον et τυνοῦτο R \espApp Γϋ. ante μικρὸν R : \lecon{nulla nota mutat. pers. } \PsiDeux}
\ligneRav{746}{}{πόσ᾽ ἔτη δὲ \difference{γέγονε}; τρεῖς χοᾶς ἢ τέτταρας;}{\textbf{746} \lecon{γέγονε} \PsiDeux, Brunck : γέγονεν R}
\ligneRav{747}{\auDessus{γυ}{ν}}{σχεδὸν τοσοῦτον \difference{χὤσον} ἐκ Διονυσίων·}{\textbf{747} καὶ ὅσον \PsiDeux, R : \lecon{χὤσον} Bentley \espApp καὶ ὅσον \PsiDeux, R : ὅσον Suid. δ 1168 : \lecon{χὤσον} Bentley}
\ligneRav{}{—}{ἀλλ᾽ ἀπόδος αὐτό. : μὰ τὸν Ἀπόλλω τουτονί·}{}
\ligneRav{749}{\auDessus{γυ}{ν}}{ἐμπρήσομεν τοίνυν σε. : πάνυ γ᾽ ἐμπίμπρατε·}{\textbf{749} ἐμπιμπρᾶται R : \lecon{ἐμπιμπράτε} (sic) \Mon}
\ligneRav{750}{}{αὕτη δ᾽ ἀποσφαγήσεται μάλ᾽ αὐτίκα·}{}
\ligneRav{}{\auDessus{γυ}{ν}}{μὴ δῆθ᾽, ἱκετεύω σ᾽· ἀλλ᾽ ἔμ᾽ ὅ τι χρῄζεις ποίει}{}
\ligneRav{}{—}{ὑπέρ γε τούτου. : φιλότεκνός τις εἶ φύσει·}{}
\ligneRav{}{}{ἀλλ᾽ οὐδὲν ἧττον ἥδ᾽ ἀποσφαγήσεται}{}
\ligneRav{754}{\auDessus{γυ}{ν}}{δὸς τὸ σφαγεῖον (σφαγ\inRas{ί}{εῖ}ον \Rpc), ὦ γύναι,}{\textbf{754} \lecon{δός τὸ} \PsiDeux, Lobeck ad Soph. \italique{Ai.} v.1066 : δός μοι τὸ R \espApp σφαγίον \Rac : \lecon{σφαγεῖον} \Rpc, \Mon}
\ligneRav{755}{}{ἵν᾽ οὖν τό γ᾽ αἷμα τοῦ τέκνου τοὐμοῦ λάβω·}{}
\ligneRav{}{μ\auDessus{ν}{Η}}{ὕπεχ᾽ αὐτό, χαριοῦμαι γὰρ ἕν γε τοῦτό σοι·}{}
\ligneRav{}{\auDessus{γυ}{ν}}{κακῶς ἀπόλοι᾽, ὡς φθονερὸς εἶ καὶ δυσμενής·}{}
\ligneRav{}{μ\auDessus{ν}{Η}}{τουτὶ τὸ δέρμα τῆς ἱερείας γίγνεται·}{}
\ligneRav{758}{\auDessus{γυ}{ν}}{τί τῆς \difference{ἱερείας} γίγνεται; : τουτί. λαβέ·}{\textbf{758-9} \lecon{ἱερείας} \PsiDeux, R : ἱερέας Austin \& Olson}
\ligneRav{760}{ἄλλη γυνή}{ταλαντάτη Μίκα τίς ἐξεκόρησέ σε;}{\textbf{760} ἐξεκόρησεν R : \lecon{ἐξεκόρησε} \Mon}
\ligneRav{761}{}{τίς τὴν ἀγαπητὴν παῖδά \difference{σοὐξηράσατο};}{\textbf{761} [\douteux{ }\douteux{ }\douteux{ }\douteux{ }\douteux{ }\douteux{ }\douteux{ }\douteux{ }]απη\douteux{τ}η\douteux{ν}\douteux{π}\douteux{ }ι\douteux{ }α\douteux{ }\douteux{ }υξ[ \PsiDeux \espApp σοὐξῃρήσατο R : \lecon{σοὐξηράσατο} Fritzsche}
\ligneRav{}{\auDessus{γυ}{ν}}{ὁ πανοῦργος οὗτος. ἀλλ᾽ ἐπειδήπερ πάρει,}{}
\ligneRav{}{}{φύλαξον αὐτόν, ἵνα λαβοῦσα Κλεισθένη}{}
\ligneRav{}{}{τοῖσιν πρυτάνεσιν ἃ πεποίηχ᾽ οὗτος φράσω·}{}
}
\vspace{3pt}\noindent\rule{4.5cm}{0.4pt}\par\vspace{1.5pt}
{\fontsize{7.0pt}{8.3pt}\selectfont
\noindent\textbf{Loci paralleli (vv. 721–764) ---} \textbf{729} (θυμάλωπα) Hesychius θ 862 θυμάλωψ ; cf. Photius θ 256 Θυμάλωψ, Hesychius θ 861 \espTest \textbf{730-1} (σὺ δὲ -- ταχέως) Suidas κ 2408 \espTest \textbf{731} (Κρητικὸν) Hesychius κ 4087 Κρητικόν· ἱματίδιον λεπτὸν καὶ βραχύ \espTest \textbf{734} (Περσικὰς) \Eustathe 1522.12 \espTest \textbf{735-7} Suidas θ 251 (ὦ et αὖ omissis, καλόν et δὲ pro ἀγαθόν et δ᾽) \espTest \textbf{735} \Eustathe 1441.23 ; \EtymSym α 241 (ed. Lasserre-Livadaras) \espTest (ποτίσταται) \EtymMagnum 31.13 \espTest \textbf{742} (ἤνεγκον) \Antiatt 98.11 ; \Phryn 73.1 \espTest \textbf{745} (τυννοῦτο) Photius τ 552 τυννοῦτο $\sim$ Hesychius τ 1645 \espTest \textbf{747} Suidas δ 1168 (καὶ omisso) \espTest \textbf{749} (ἐμπίμπρατε) Hesychius ε 2458 ἐμπίπραται $\sim$ Suidas ε 1018 $\sim$ Zonaras 709 \espTest \textbf{754} Photius σ 850 σφαγεῖον· [...] περισπᾶται ἡ δευτέρα συλλαβή\par
}
\vspace{2.5pt}\noindent\rule{4.5cm}{0.4pt}\par\vspace{1.5pt}
{\fontsize{6.9pt}{8.2pt}\selectfont
\noindent\textbf{Suite (en commentaire) de l'apparat critique \& notes paléographiques (\texttt{f. 171r}) ---} \textbf{Contexte :} Suite et dénouement de la parodie du \italique{Télèphe} : les femmes apportent du bois pour brûler Mnésiloque ; en démaillotant l'otage, Mnésiloque découvre que le « bébé » est une outre pleine de vin chaussée de petites bottines perses (\italique{persikai}), qu'il « égorge » et boit goulûment au-dessus du bassin sacrificiel (\italique{sphageion}, vv. 721–764). --- \textbf{Témoin papyrologique :} \textbf{\PsiDeux{} (\italique{P.Oxy.} LVI 3839, fr. 1)} couvre les \textbf{vv. 742–766} et apporte quatre confirmations décisives face à R (vv. 743, 745, 746, 754). --- \textbf{Confrontation directe entre R et le papyrus \PsiDeux{} (vv. 743–754) :} --- \textbf{v. 721 :} ἀθέοις ἔργοις γὰρ R (ἀθ. γὰρ ἔργ. Burges ; ἐπ' ἀθ. ἔρ. γὰρ Enger) : \lecon{ξὺν ἀθέοις ἔργοις γὰρ} \lecon{Bibi} ap. Taillardat (cf. v. 716) \espApp γὰρ R : \lecon{καὶ γὰρ} Hermann.\par
}
\end{minipage}
\end{adjustbox}
\end{adjustwidth}
\newpage
\enTeteD{Codex Ravennas 429 (Biblioteca Classense) --- Fac-similé du f. 171r (vv. 721–764)}
\begin{adjustwidth}{-1.5cm}{-3.0cm}
\begin{center}
\includegraphics[width=16.5cm,height=20.5cm,keepaspectratio]{/workspaces/Thesmophories/dossier_ravennas_thesmophories/folios_individuels_ff162v-177r/18_f171r_classense173r_vv0721-0764.jpg}\par
\vspace{2pt}
{\scriptsize\textbf{Codex Ravennas 429, f. 171r (Classense 173r)} --- ΕΠΕΙΣΟΔΙΟΝ Γ', \textbf{vv. 721–764} ($\text{R}^1, \Rpc, \Rdeux$, \PsiDeux).}
\end{center}
\end{adjustwidth}
\newpage
\clearpage\ifodd\value{page}\hbox{}\newpage\fi% --- FOLIO 18 (f. 171r) : SPREAD 2.1 (Texte grec établi + Traduction française + Commentaire 1/8 & 2/8) ---
\clearpage\ifodd\value{page}\hbox{}\newpage\fi
\enTeteG{ΘΕΣΜΟΦΟΡΙΑΖΟΥΣΑΙ --- Texte grec établi (f. 171r, vv. 721–764 [1/4])}
\enTeteD{LES THESMOPHORIES --- Traduction française (f. 171r, vv. 721–764 [1/4])}
\begin{adjustwidth}{-2cm}{-2cm}
\noindent\textbf{\footnotesize Texte grec établi --- Folio 18 (\texttt{f. 171r}, vv. 721–764 [1/4]) \hfill ΕΠΕΙΣΟΔΙΟΝ Γ'}\\[-2pt]
\noindent\rule{16.5cm}{0.5pt}\\[2pt]
{\fontsize{8.8pt}{10.2pt}\selectfont
\ligneEst{{\tiny\textit{dochm.}}}{ξὺν ἀθέοις ἔργοις}{721}
\ligneEst{{\tiny\textit{tetram. troch.}}}{καὶ γὰρ ἀνταμειψόμεσθά σ᾽ ὥσπερ εἰκὸς ἀντὶ τῶνδε.}{}
\ligneEst{{\tiny\textit{2 dochm.}}}{τάχα δὲ μεταβαλοῦσ᾽ ἐπὶ κακὸν ἑτερότρο-}{}
\ligneEst{{\tiny\textit{dochm.}}}{\alinea πον ἐπέχει τύχη.}{725}
\ligneEst{\textbf{Κορυφαῖος} {\tiny\textit{$4$ troch\cat}}}{ἀλλὰ τάσδε μὲν λαβεῖν χρῆν σ᾽ ἐκφέρειν τε τῶν ξύλων,}{}
\ligneEst{}{καὶ καταίθειν τὸν πανοῦργον πυρπολεῖν θ᾽ ὅσον τάχος.}{}
\ligneEst{\textbf{Μίκα} {\tiny\textit{trim. iamb.}}}{ἴωμεν ἐπὶ τὰς κληματίδας ὦ Μανία.}{}
\ligneEst{}{κἀγώ σ᾽ ἀποδείξω θυμάλωπα τήμερον.}{}
\ligneEst{\textbf{Μνησίλοχος}}{ὕφαπτε καὶ κάταιθε· σὺ δὲ τὸ Κρητικὸν}{730}
\ligneEst{}{ἀπόδυθι ταχέως· τοῦ θανάτου δ᾽ ὦ παιδίον}{}
\ligneEst{}{μόνην γυναικῶν αἰτιῶ τὴν μητέρα.}{}
\ligneEst{}{τουτὶ τί ἔστιν; ἀσκὸς ἐγένεθ᾽ ἡ κόρη}{}
\ligneEst{}{οἴνου πλέως καὶ ταῦτα Περσικὰς ἔχων.}{}
}
\vspace{2pt}
\noindent\rule{16.5cm}{0.6pt}\\[2pt]
\noindent\textbf{\small Commentaire philologique, métrique et scénique (vv. 721–764 --- partie 1/8)}\\[2pt]
\begin{adjustbox}{max totalheight=15.2cm, max width=16.5cm, keepaspectratio}%
\parbox[t]{16.5cm}{
\setlength{\columnsep}{9pt}%
\begin{multicols}{2}
\fontsize{7.0pt}{8.0pt}\selectfont
\setlength{\parindent}{0pt}%
\setlength{\parskip}{1.2pt}%
\noindent\textbf{721-5} \textbf{ἀθέοις ἔργοις γὰρ ἀνταμει- / ψόμεσθά σ᾽ ὥσπερ εἰκὸς ἀντὶ τῶνδε / τάχα δὲ \{σε\} μεταβαλοῦσ᾽ ἐπὶ κακὸν ἑτερότροπον / ἐπέχει \{τις\} τύχη.} Le grand Jean Taillardat avait donné un exposé sur son interprétation des vers 721-725 du manuscrit de Ravenne. Il nous offre une explication limpide de ce passage fort débattu ; je transcris pieusement ci-dessous le résumé de son intervention, tout en reprenant sa traduction\marqueur{1} : \\ « Car toi, avec tes impiétés, nous te châtierons de ta conduite, comme de juste. Et puis la Fortune, changeant brusquement d'humeur, aura vite fait de poursuivre un méchant homme qui a pris le chemin du crime ».\\ Cette traduction appelle les justifications suivantes : \\ — Ἀθέοις ἔργοις : selon Jean Taillardat, il est impossible d'interpréter ce datif comme le complément de moyen ἀνταμειψόμεσθα et de traduire « \latin{\italique{Impiis factis} enim remunerabimus te sicut par est pro eis} » (Andrea Divo) ou « \latin{\italique{Nefariis} enim \italique{factis} vicissim, ut par est, haec tua ulciscemur} » (Brunck) ou « Nous y répondrons, comme de raison, \italique{par des actes sacrilèges} » (Willems)\marqueur{2}. En effet, on ne pouvait pas considérer comme sacrilège « \italique{la coutume} (qui) était d'apporter du feu près des suppliants réfugiés sur l'autel » afin de les contraindre à le lâcher (Σ Eur. \italique{Androm}. 256 : ἔθος ἦν πῦρ προσφέρειν τοῖς εἰς βωμὸν καταφεύγουσιν ; cf. Eur. \italique{Herc. fur}. 243, 285). Ἀθέοις ἔργοις est donc : \ \ \ 1) \ soit un datif-instrumental d'accompagnement lié à σ(ε), cf. Kühner-Gerth, I, p. 435; ἀθέοις ἔργοις a, dans ce cas, la même fonction syntaxique que ξὺν ἀδίκοις ἔργοις\marqueur{3} (v. 716) lié à σοι (v. 715) ; \ \ \ 2) \ soit un datif-instrumental de cause lié à ἀνταμειψόμεσθα. Jean Taillardat penche pour la première solution.\\ — Ἀντὶ τῶνδε (v. 723) $=$ ἀντὶ τῶνδε τῶν ἀθέων ἔργων « nous te payerons en retour de tes impiétés».\\ — Μεταβαλοῦσ(α), v. 724, est employé absolument, ce qui n'est pas rare, cf. L.S.J. \italique{s. v.} μεταβάλλω, A. III. 2.\\ On reconnaîtra d'autre part la locution ἐπέχειν ἐπί τινα (Hérodote 9.59) à peu près synonyme de διώκειν τινά (Id. \italique{ibid}.), car τάχα ... ἐπὶ κακὸν ἑτερότροπον ἐπέχει (présent de prophétie) τύχη équivaut à τάχα ... ἐπὶ κακὸν ἄνδρ᾽ ὃς τὴν ἑτέραν ὁδὸν ἐτράπετ' ἐπέχει τύχη.\\ Pour κακός $=$ κακός τις ἀνήρ, cf. Théognis 189, 190, 577, 972, 1112 (de même κακὴ $=$ κακή τις γυνή, id. 185). Le composé ἑτερότροπος « qui a pris la mauvaise route » (avec ἕτερος euphémique pour κακός, cf. L.S.J. III. 2) est une allusion au mythe bien connu des deux chemins, celui du bien et celui du mal (Hésiode \italique{Op}. 287-292 ; Prodicos in Xénophon \italique{Mém}. 2. 1. 21-27). \\ Quant à σε de τάχα δέ σε μεταβαλοῦσα, κτλ. (cod. R), Jean Taillardat le considère comme une glose qui expliquait κακόν ἑτερότροπον et qui s'est glissée dans le texte même. \\ Quoi qu'on en ait dit, le sens de ce passage est très clair : le chœur fait entendre au Parent d'Euripide qu'à défaut d'une punition donnée par les femmes, il recevra un châtiment infligé par le Ciel.
\noindent\textbf{[Notes ad 721-5]} \textbf{(1)} À retrouver dans la \italique{Revue des Études Grecques}, tome 101, fascicule 482-484 (juillet-décembre 1988). Actes de l'association, p. XVIII-XIX.
\noindent \textbf{(2)} Ou « \anglais{Since your crimes are godless, it is fitting that our punishment of you be \italique{godless} as well} » (Austin \& Olson.)
\noindent \textbf{(3)} Fort de cette remarque  , je [Bibi] n'ai pas hésité à introduire ce ξύν dans le texte  au vers 721, qui avait un problème de mètre : [ξὺν] ἀθέοις ἔργοις. C'est aussi \italique{gratia metri} qu'Enger proposait en son temps l'ajout d' [ἐπ'] ἀθέοις ἔργοις, conjecture acceptée par Lætitia (\italique{The Songs of Aristophanes}, p. 427). N.D.É.
\noindent\textbf{724-5} \textbf{τάχα δὲ \{σε\} μεταβαλοῦσ᾽ ἐπὶ κακὸν ἑτερότροπον ἐπέχει \{τις\} τύχη.} Il suffit que le chœur dise `τάχα δὲ' ... et le rythme même du vers s'accélère : les trois dochmiaques des vers 724-725 contiennent force résolutions.
\noindent\textbf{726-7} \textbf{ἀλλὰ τάσδε μὲν λαβεῖν χρῆν σ᾽ ἐκφέρειν τε τῶν ξύλων, / καὶ καταίθειν τὸν πανοῦργον πυρπολεῖν θ᾽ ὅσον τάχος.} C'était, en effet, l'habitude de faire périr par le feu ceux qui avaient cherché refuge sur un autel. Cf. Euripide, \italique{Androm}., v. 256 ; \italique{Herc. Fur}., v. 243 et 285 et les interprètes. (Note de Willems.)\\ Cette \italique{exhortation} -- κατακελευσμός selon la dénomination consacrée de Zieliński\marqueur{1} -- de la Coryphée pour appeler à la reprise de l'action est selon l'usage un distique de tétramètres trochaïques. (Et, comme au vers 531-532, débute par `ἀλλά'.)
\end{multicols}
}%
\end{adjustbox}%
\vfill
\end{adjustwidth}
\newpage
% --- PAGE DE DROITE : Traduction française + Commentaire (2/8) ---
\begin{adjustwidth}{-2cm}{-2cm}
\noindent\textbf{\footnotesize Traduction française --- Folio 18 (\texttt{f. 171r}, vv. 721–764 [1/4]) \hfill ΕΠΕΙΣΟΔΙΟΝ Γ'}\\[-2pt]
\noindent\rule{16.5cm}{0.5pt}\\[2pt]
{\fontsize{8.8pt}{10.2pt}\selectfont
\ligneEst{}{\francais{Car toi, avec tes impiétés, nous te châtierons}}{721}
\ligneEst{}{\francais{de ta conduite, comme de juste.}}{}
\ligneEst{}{\francais{Et puis la Fortune, changeant brusquement d'humeur aura vite fait de poursuivre un méchant homme}}{}
\ligneEst{}{\francais{qui a pris le chemin du crime.}}{725}
\ligneEst{\textbf{Coryphée}}{\francais{Mais prends ces femmes avec toi, et va chercher du bois.}}{}
\ligneEst{}{\francais{Il faut le brûler vif, ce fourbe, l'anéantir au plus tôt par le feu.}}{}
\ligneEst{\textbf{Mika}}{\francais{Allons chercher les fagots, Mania.}}{}
\ligneEst{}{\francais{Je ferai de toi un tison aujourd'hui..}}{}
\ligneEst{\textbf{Mnésiloque}}{\francais{Allumez, brûlez. Toi, ma petite, ôte vite ta crétique}}{730}
\ligneEst{}{\francais{et n'accuse de ta mort que ta mère.}}{}
\ligneEst{}{\francais{Qu'est ceci ? La fillette n'est plus qu'une outre}}{}
\ligneEst{}{\francais{pleine de vin, ... et avec cela chaussée de persiques.}}{}
\ligneEst{}{\francais{Ô femmes archi-enfiévrées et archi-biberonnes,}}{}
}
\vspace{2pt}
\noindent\rule{16.5cm}{0.6pt}\\[2pt]
\noindent\textbf{\small Commentaire philologique, métrique et scénique (vv. 721–764 --- partie 2/8)}\\[2pt]
\begin{adjustbox}{max totalheight=15.2cm, max width=16.5cm, keepaspectratio}%
\parbox[t]{16.5cm}{
\setlength{\columnsep}{9pt}%
\begin{multicols}{2}
\fontsize{7.0pt}{8.0pt}\selectfont
\setlength{\parindent}{0pt}%
\setlength{\parskip}{1.2pt}%
\noindent\textbf{726} \textbf{Τάσδε} renvoie à des compagnes de Mica et non pas à des torches, selon l'hypothèse invraisemblable de Fritzsche -- qui a malheureusement fait florès. Logique pratique élémentaire : avant d'aller recueillir du bois, il faut d'abord se \italique{débarrasser} de ce qui encombre -- en l'occurrence : les torches -- et non s'en embarrasser. Ce n'est qu'une fois le bois entassé que Mica pourra effectivement se saisir d'une torche et bouter le feu au bûcher.
\noindent\textbf{[Notes ad 726-7]} \textbf{(1)} Voir \italique{The Verse of Greek Comedy} (1912), p. 316, de John W. White. Le lecteur pourra utilement consulter ce livre de référence, entre autre s'il s'intéresse à la structure des comédies d'Aristophane.
\noindent\textbf{728} \textbf{ἴωμεν ἐπὶ τὰς κληματίδας.} Blaydes, jugeant l'article de trop, propose : ἴωμεν οὖν ἐπὶ κληματίδας. Nous répondrons par le texte suivant de Plutarque, où il s'agit, comme ici, d'un sanctuaire où des femmes se trouvent réunies pour la célébration de leurs mystères :
\citComm{$[$ ἡ Δαμοκρίτα $]$ ξύλα ταῖς θύραις πολλὰ προσνήσασα ταῦτα δ᾽ ἦν \textbf{εἰς τὴν τῆς ἑορτῆς θυσίαν ὑπ᾽ ἐκείνων παρεσκευασμένα}, πῦρ ἐνῆκε.\\
Damocrita ayant entassé contre les portes une quantité de bois, \italique{que les femmes avaient préparé pour le sacrifice de la fête}, elle y mit le feu.}
\noindent (Plutarque, \italique{Narrat. amour}., 5, p. 775e ; Note de Willems.)
\noindent\textbf{730} \textbf{τὸ Κρητικὸν.} Hésychius (κ 4087) :
\citComm{Κρητικόν· ἱματίδιον λεπτὸν καὶ βραχύ.\\
\francais{Crétique : petit habit fin et court.}}
\noindent\textbf{733} \textbf{ἀσκὸς ἐγένεθ᾽ ἡ κόρη.} Il était probablement interdit d'apporter du vin lors des Thesmophories, fêtes consacrées à Déméter et sa fille. (Cf. par exemple l'adage latin \latin{Cereri sacrificant}\marqueur{1}, \italique{sacrifier à Cérès}, c'est-à-dire faire abstinence de vin.) Cette interdiction théorique devait souffrir quelques écarts pratiques, si l'on en croit le compte-rendu de Mnésiloque (vers 626-633).
\noindent\textbf{734} \textbf{Περσικὰς ἔχων.} On comprend d'après d'autres passages d'Aristophane que les \italique{persiques} désignent des chaussures légères et recouvrant le pied\marqueur{1}, peut-être des \italique{chaussons}\marqueur{2} ; en tout cas, ces \italique{persiques} sont portées par les femmes\marqueur{3}, ou comme ici les enfants.\\ Les «~persiques~» dont parle Mnésiloque sont visibles sur le fameux cratère de Würzbourg\marqueur{4}. Il faut savoir que les outres à l'époque étaient faites de la peau d'un animal évidé\marqueur{5}, bouc ou chèvre la plupart du temps. Ces «~persiques~» sont en réalité les pattes arrières de l'animal outré, fermées par un lacet.
\noindent\textbf{735} \textbf{ὦ θερμόταται γυναῖκες.} θερμός évoque l'\italique{impétuosité} féminine : \anglais{hot-headed, hasty}, l'\italique{audace} des femmes\marqueur{1}.
\noindent Le mot peut aussi désigner la \italique{fièvre} ; il n'y a cependant pas d'équivalent grec à notre \italique{fièvre de Bercy} : à ma connaissance, le grec n'a jamais rapproché la fièvre de l'ivrognerie.
\noindent Quand donc le scholiaste à notre vers écrit `ὦ μέθυσοι', il glose `ὦ ποτίσταται' et non `ὦ θερμόταται'.
\noindent Suidas (θ 251) s'y est trompé qui recopie la glose en la faisant précéder du lemme θερμόταται :
\citComm{Θερμόταται: μέθυσοι. ὦ θερμόταται γυναῖκες, μέγα καπήλοις καλόν, ἡμῖν δὲ κακόν [...]}
\noindent Le cas n'est pas isolé : Suidas avait déjà commis un précédent du même genre : au vers 178 (\italique{confer} le commentaire \italique{ad locum}), lisant κοινῇ à place de καινῇ, il n'hésita pas à utiliser la glose afférente, μεγάλῃ καὶ θαυμαστῇ, comme nouvelles acceptions de κοινῇ (en y ajoutant le véritable sens de κοινῇ : ἀνθρωπίνῃ).
\noindent Ce sont là des méthodes qui montrent que la \italique{Souda} est un lexique à manipuler avec précaution.
\noindent\textbf{[Notes ad 733]} \textbf{(1)} Recensé par Érasme, \italique{Adages} n° 3071. Voir aussi les adage n° 1895, 1896.
\noindent \textbf{(1)} La puce des socratiques expériences des \italique{Nuées}, v. 151, se retrouve chaussée de `persiques' après que Socrate lui ait plongé les pattes dans la cire – voir aussi Hésychius (π 2005) : \grec{περσικά· εὐτελῆ ὑποδήματα}.
\noindent \textbf{(2)} Les Athéniennes pouvaient les garder au lit, en prenant garde toutefois de maintenir les jambes en l'air, cf. \italique{Lysistrata} v. 229.
\noindent \textbf{(3)} \italique{Assemblée des femmes} v. 319 et Hésychius (μ 929) : \grec{μεσοπερσικαί· ὑποδήματα γυναικεῖα}.
\noindent \textbf{(4)} Reproduit en annexe. Voir la note aux vers 754-757.
\noindent \textbf{(5)} Pour plus d'informations, aller consulter l'indispensable Daremberg \& Saglio, \italique{Dictionnaire des antiquités grecques et romaines} à l'entrée \latin{Uter}.
\end{multicols}
}%
\end{adjustbox}%
\vfill
\end{adjustwidth}
\newpage

% --- FOLIO 18 (f. 171r) : SPREAD 2.2 (Texte grec établi + Traduction française + Commentaire 3/8 & 4/8) ---
\clearpage\ifodd\value{page}\hbox{}\newpage\fi
\enTeteG{ΘΕΣΜΟΦΟΡΙΑΖΟΥΣΑΙ --- Texte grec établi (f. 171r, vv. 721–764 [2/4])}
\enTeteD{LES THESMOPHORIES --- Traduction française (f. 171r, vv. 721–764 [2/4])}
\begin{adjustwidth}{-2cm}{-2cm}
\noindent\textbf{\footnotesize Texte grec établi --- Folio 18 (\texttt{f. 171r}, vv. 721–764 [2/4]) \hfill ΕΠΕΙΣΟΔΙΟΝ Γ'}\\[-2pt]
\noindent\rule{16.5cm}{0.5pt}\\[2pt]
{\fontsize{8.8pt}{10.2pt}\selectfont
\ligneEst{}{ὦ θερμόταται γυναῖκες, ὦ ποτίσταται}{735}
\ligneEst{}{κἀκ παντὸς ὑμεῖς μηχανώμεναι πιεῖν,}{}
\ligneEst{}{ὦ μέγα καπήλοις ἀγαθὸν ἡμῖν δ᾽ αὖ κακόν,}{}
\ligneEst{}{κακὸν δὲ καὶ τοῖς σκευαρίοις καὶ τῇ κρόκῃ.}{}
\ligneEst{\textbf{Μίκα}}{παράβαλλε πολλὰς κληματίδας ὦ Μανία.}{}
\ligneEst{\textbf{Μνησίλοχος}}{παράβαλλε δῆτα· σὺ δ᾽ ἀπόκριναί μοι τοδί,}{}
\ligneEst{}{τουτὶ τεκεῖν φῄς;}{740}
\ligneEst{\textbf{Μίκα}}{\decalage{τουτὶ τεκεῖν φῄς;}καὶ δέκα μῆνας αὔτ᾽ ἐγὼ}{}
\ligneEst{}{ἤνεγκον.}{}
\ligneEst{\textbf{Μνησίλοχος}}{\decalage{ἤνεγκον.}ἤνεγκας σύ;}{}
\ligneEst{\textbf{Μίκα}}{\decalage{ἤνεγκον.ἤνεγκας σύ;}νὴ τὴν Ἄρτεμιν.}{}
\ligneEst{\textbf{Μνησίλοχος}}{τρικότυλον; ἢ πῶς.}{}
\ligneEst{\textbf{Μίκα}}{\decalage{τρικότυλον; ἢ πῶς.}εἰπέ μοι· τί μ᾽ ἠργάσω;}{}
}
\vspace{2pt}
\noindent\rule{16.5cm}{0.6pt}\\[2pt]
\noindent\textbf{\small Commentaire philologique, métrique et scénique (vv. 721–764 --- partie 3/8)}\\[2pt]
\begin{adjustbox}{max totalheight=15.2cm, max width=16.5cm, keepaspectratio}%
\parbox[t]{16.5cm}{
\setlength{\columnsep}{9pt}%
\begin{multicols}{2}
\fontsize{7.0pt}{8.0pt}\selectfont
\setlength{\parindent}{0pt}%
\setlength{\parskip}{1.2pt}%
\noindent \textbf{(1)} Cf. le \LSJ , Taillardat §§313, 351.
\noindent\textbf{738} \textbf{κακὸν δὲ καὶ τοῖς σκευαρίοις καὶ τῇ κρόκῃ.} Mnésiloque reproche à ces soûlaudes d'Athéniennes de contribuer à la ruine des ménages par leur inconduite biberonnesque.
\noindent Montrons-le !
\noindent \grec{τό σκευάριον}, diminutif de σκεῦος, désigne littéralement un \italique{petit récipient} (\italique{menue vaisselle}) ou \italique{petit meuble}.
\noindent Employé au pluriel, il peut désigner plus particulièrement le meuble où l'on entasse son argent. \italique{Confer} le passage du \italique{Ploutos} ci-dessous (vers 806-811) qui nous renseigne sur l'usage du mobilier grec :
\citComm{Ἡ μὲν σιπύη μεστή ᾽στι λευκῶν ἀλφίτων,\\
οἱ δ᾽ ἀμφορῆς οἴνου μέλανος ἀνθοσμίου.\\
Ἅπαντα δ᾽ ἡμῖν ἀργυρίου καὶ χρυσίου\\
\textbf{τὰ σκευάρια} πλήρη ᾽στίν, ὥστε θαυμάσαι.\\
Τὸ φρέαρ δ᾽ ἐλαίου μεστόν· αἱ δὲ λήκυθοι\\
μύρου γέμουσι, τὸ δ᾽ ὑπερῷον ἰσχάδων.\\
\francais{Notre huche au pain regorge d'une blanche farine, /
nos amphores d'un vin noir qui embaume. /
Tous nos \italique{meubles}} (τὰ σκευάρια) \francais{sont pleins /
d'argent et d'or, que c'est merveille. /
La citerne regorge d'huile, les lécythes /
sont emplis d'essences, le grenier de figues sèche}}
\noindent (C'est bien de \italique{pièces} d'or et d'argent dont il est question puisque les plats et autres assiettes d'airain ou d'argent sont évoqués aux vers suivants.) Il est tout naturel que ces meubles (coffrets ?) servant de porte-monnaie aient pris par métonymie le sens de \italique{biens}, \italique{économies}.
\noindent C'est d'ailleurs si naturel que c'est effectivement le cas : aux vers 838-839, le Juste, contant ses misères passées, emploie τὰ σκευάρια en ce sens métonymique :
\citComm{Δίκαιος — κομιδῇ μὲν οὖν· / αὐχμὸς γὰρ ὢν \textbf{τῶν σκευαρίων} μ᾽ ἀπώλεσεν.\\
\francais{Le Juste — Précisément. Une sécheresse s'étant mise dans \italique{mon mobilier} ($=$ mes économies) fut ce qui me perdit.}}
\noindent Cette acception de τὰ σκευάρια est encore plus manifeste dans un passage au début des \italique{Grenouilles}. Dionysos y cherche un moyen d'expédier en Enfer ses \italique{bagages} ; son esclave Xanthias regimbant, il hèle un mort en route pour l'Hadès (vers 170-172) :
\citComm{Καὶ γάρ τιν᾽ ἐκφέρουσι τουτονὶ νεκρόν.\\
Οὗτος, σὲ λέγω μέντοι, σὲ τὸν τεθνηκότα.\\
Ἄνθρωπε, βούλει \textbf{σκευάρι᾽} εἰς Ἅιδου φέρειν;\\
\francais{Car voilà justement un mort qu'on emporte.\\
Holà ! c'est à toi que je parle, toi le mort.\\
Homme, veux-tu porter à l'Hadès quelques \italique{petits colis} ?}}
\noindent Il y a un évident détournement dans la plus pure veine aristophanesque de l'universel adage : on n'emporte (φέρειν) pas son \italique{argent} (σκευάρι᾽) dans l'autre monde (εἰς Ἅιδου) ; les σκευάρια de l'adage sont à prendre au sens figuré, \italique{argent} mais se retrouvent tout matériellement incarnés en \italique{bagages} dans la version illustrée d'Aristophane\marqueur{1}. Voilà pour τοῖς σκευαρίοις, passons à τῇ κρόκῃ.\\ \grec{ἡ κρόκη} est au propre \italique{le fil de trame}.
\noindent Il possède, comme nous l'allons voir, le même sens imagé de \italique{richesses} que τὰ σκευάρια.
\noindent Le grec, comme le français, a prodigalement tiré maintes images et maints proverbes de l'activité filandière et tisserande. Observons l'un deux : σπαθᾶν τὴν κρόκην.
\noindent Au vers 55 des \italique{Nuées}, Strepsiade énumère les oppositions entre lui, le rustre économe, et sa femme, la raffinée dilapidatrice. Il pousse la comparaison jusque dans l'alcôve où sa femme souffre de besoins tyranniques ; il lui reproche de \italique{tirer trop sur la corde} -- le grec dit \italique{frapper sur le fil} :
\citComm{ὦ γύναι, λίαν σπαθᾷς.}
\noindent σπαθᾶν [τὴν κρόκην] désigne l'action de battre chaque fil de trame (κρόκη) au moyen de la σπάθη, ou \italique{battant}, pour rendre le tissu compact\marqueur{2} ; ce faisant, l'on utilise davantage de laine, d'où le fait que l'expression signifie \italique{dissiper l'argent}, \italique{gaspiller}.
\noindent Or l'on trouve les expressions σπαθᾶν τὰ πατρῷα\marqueur{3}, τὰ χρήματα\marqueur{4}, τάλαντα\marqueur{5}.
\end{multicols}
}%
\end{adjustbox}%
\vfill
\end{adjustwidth}
\newpage
% --- PAGE DE DROITE : Traduction française + Commentaire (4/8) ---
\begin{adjustwidth}{-2cm}{-2cm}
\noindent\textbf{\footnotesize Traduction française --- Folio 18 (\texttt{f. 171r}, vv. 721–764 [2/4]) \hfill ΕΠΕΙΣΟΔΙΟΝ Γ'}\\[-2pt]
\noindent\rule{16.5cm}{0.5pt}\\[2pt]
{\fontsize{8.8pt}{10.2pt}\selectfont
\ligneEst{}{\francais{ingénieuses à trouver partout matière à boire !}}{735}
\ligneEst{}{\francais{Ô providence des débitants, et notre fléau à nous !}}{}
\ligneEst{}{\francais{Fléau aussi des ustensiles de ménage, et fléau de la trame !}}{}
\ligneEst{}{\francais{(Les femmes reviennent avec des fagots.)}}{}
\ligneEst{\textbf{Mika}}{\francais{Entasse force fagots autour de lui, Mania.}}{}
\ligneEst{\textbf{Mnésiloque}}{\francais{Entasse, c'est bon. Mais toi, réponds à ceci.}}{}
\ligneEst{}{\francais{Tu te dis mère de cet enfant ?}}{740}
\ligneEst{\textbf{Mika}}{\francais{\espace{2.8cm}Dix mois durant}}{}
\ligneEst{}{\francais{je l'ai porté.}}{}
\ligneEst{\textbf{Mnésiloque}}{\francais{\espace{1.5cm}Tu l'as porté, toi ?}}{}
\ligneEst{\textbf{Mika}}{\francais{\espace{3.5cm}Oui, par Artémis.}}{}
\ligneEst{\textbf{Mnésiloque}}{\francais{(Il porte l'outre à ses lèvres.) Du trois cotyles ? Ou de quel prix ?}}{}
\ligneEst{\textbf{Mika}}{\francais{\espace{4.2cm}Dis-moi : quel tour vas-tu me faire ?}}{}
}
\vspace{2pt}
\noindent\rule{16.5cm}{0.6pt}\\[2pt]
\noindent\textbf{\small Commentaire philologique, métrique et scénique (vv. 721–764 --- partie 4/8)}\\[2pt]
\begin{adjustbox}{max totalheight=15.2cm, max width=16.5cm, keepaspectratio}%
\parbox[t]{16.5cm}{
\setlength{\columnsep}{9pt}%
\begin{multicols}{2}
\fontsize{7.0pt}{8.0pt}\selectfont
\setlength{\parindent}{0pt}%
\setlength{\parskip}{1.2pt}%
\noindent Le sème commun à ces différents mots est celui de \italique{richesse} (matérielle ou pécuniaire) ; puisqu'ils ont été substitués à τὴν κρόκην, c'est donc la preuve -- taillardesque\marqueur{6} -- qu'ils partagent avec κρόκη ce sème.
\noindent Par exemple, en français l'on rencontre les expressions \italique{ourdir un complot}, \italique{une intrigue}, \italique{une conspiration}, \italique{une vengeance}. Pris dans leur contexte, le sème commun aux différents compléments d'ourdir est : \italique{machination secrète}. Or l'expression tisserandière propre est «~ourdir une trame~», qui consiste à entrecroiser les fils de trame avec les fils ourdis, c'est-à-dire disposés sur l'ourdissoir. Les termes de complot, intrigue, \&c. ne sont substituables à la trame que parce qu'ils ont en commun avec elle le sème de \italique{machination secrète} (sens figuré de trame).\\ Finissons avec un amusant proverbe sur ce fléau qu'est l'alcool, surtout pour le sexe filandier\marqueur{7} :
\citComm{ἡ γραῦς τὰς ἀγαθίδας ἐκπίνουσα\\
La vieille boit\marqueur{8} ses pelotes.}
\noindent ἡ ἀγαθίς est une \italique{pelote de fil}. On constate que ce mot a aussi pris le sens de \italique{bien}, \italique{richesse} (on parlerait d'un \italique{bas de laine} en français) ; par permutation métaphorique avec κρόκη ? peut-être aussi par paronomase : ἀγαθῶν ἀγαθίδες est une expression désignant une grande quantité de biens.
\noindent Notre tradition populaire n'est d'ailleurs pas en reste : voici ce que l'on chantait à l'époque où l'on chantait encore\marqueur{9} :
\citComm{\francais{Parlons un peu des filandières.\\
Celles qui filaient l'étoupe\\
Ont mangé toute la soupe.\\
Celles qui filaient le lin\\
Ont lampé tout le vin.}}
\noindent Et encore, à propos de filles qui se sont enivrées :
\citComm{\francais{Elles ne l'ont pas fait en filant du lin,\\
Mais elles l'ont bien fait en buvant du vin.}}
\noindent\textbf{[Notes ad 738]} \textbf{(1)} Même procédé dans les \italique{Nuées} : quand Strepsiade désabusé songe au \italique{Tourbillon} (Δῖνος) de Socrate, il ne voit plus en cette nouvelle divinité qu'une vile \italique{poterie} (δῖνος). Si l'on veut savoir ce que le véritable Socrate entendait par ce \italique{Tourbillon}, il suffit d'interroger un autre grand esprit socratique : Pierre Boutang : «~La lignée [de race et d'esprit], c'est merveilleux !.. mais un peu trop linéaire pour mon goût. Donc il me faut quelque chose de plus circulaire : des lignes qui se répètent, qui reviennent sur elles-mêmes, qui se retournent ; des tourbillons ! [ ... ] Un grand tourbillon créateur !~». Quoi qu'on en ait dit, le Socrate croqué par Aristophane est loin d'être si caricatural et malgré ses défauts qu'Aristophane grossit à la loupe, il est encore à mille coudées au-dessus Cléon, par exemple.
\noindent \textbf{(2)} Avec le commentaire du scholiaste repris dans la \italique{Souda} (ε 3176) : καὶ γὰρ ἐπὶ τῶν ὑφαινομένων ἱματίων λέγομεν σπαθᾶν τὸ ἄγαν κρούειν τὴν κρόκην. Taillardat (§440) développe l'exégèse de l'image. Voir aussi le lexique de Photius (ε 2014) : Ἐσπαθημένος· πεπυκνωμένος· ἀπὸ τῶν κατακρουουσῶν σπαθῶν τὴν κρόκην.
\noindent \textbf{(3)} Diphil. 43, 27.
\noindent \textbf{(4)} Plut., \italique{Per.} 14. Il s'agit ici des \italique{deniers publics} : τῶν δὲ περὶ τὸν Θουκυδίδην ῥητόρων καταβοώντων τοῦ Περικλέους ὡς σπαθῶντος τὰ χρήματα καὶ τὰς προσόδους ἀπολλύντος.
\noindent \textbf{(5)} Lucien, \italique{Cat.} 20.
\noindent \textbf{(6)} Lire et relire l'excellentissime article du regretté Jean, « Images et matrices métaphoriques » (1977), in \italique{Bulletin de l'Association Guillaume Budé} 4, pp. 344-354.
\noindent \textbf{(7)} \italique{Corpus paroemiographorum Graecorum}, Appendices, III, 3.
\noindent \textbf{(8)} Sur l'image de «~manger~» son bien, voir Taillardat §439.
\noindent \textbf{(9)} Doit-on dire de cette époque qu'«~elle est révolue, qu'elle a fait son temps ; qu'on n'la reverra plus... et qu'c'est bien attristant~» ? Toujours est-il que j'emprunte ces chansons à une tradition qui a depuis longtemps défunté, mais qui fut néanmoins -- et heureusement ! -- bien embaumée par notre grand Patrice Coirault, aux cotes 110-05 et 110-06 de son \italique{Répertoire des chansons françaises de tradition orale}.
\noindent\textbf{741} \textbf{δέκα μῆνας.} Les mois lunaires grecs étant de 29 ou 30 jours, la femme portait bien environ dix mois l'enfant. Aulu-Gelle dans ses \italique{Nuits attiques} traite pendant tout un chapitre (III, 16) de la question du temps de gestation. N'en déplaise au scholiaste, la durée de 10 mois est bien attestée, notamment par Ménandre ; Aulu-Gelle cite aussi Varron expliquant que les noms des Parques \italique{Nona} et \italique{Decima} viennent des nombres de mois fixés pour l'accouchement (\latin{item Nona et Decima a partus tempestivi tempore}).
\end{multicols}
}%
\end{adjustbox}%
\vfill
\end{adjustwidth}
\newpage

% --- FOLIO 18 (f. 171r) : SPREAD 2.3 (Texte grec établi + Traduction française + Commentaire 5/8 & 6/8) ---
\clearpage\ifodd\value{page}\hbox{}\newpage\fi
\enTeteG{ΘΕΣΜΟΦΟΡΙΑΖΟΥΣΑΙ --- Texte grec établi (f. 171r, vv. 721–764 [3/4])}
\enTeteD{LES THESMOPHORIES --- Traduction française (f. 171r, vv. 721–764 [3/4])}
\begin{adjustwidth}{-2cm}{-2cm}
\noindent\textbf{\footnotesize Texte grec établi --- Folio 18 (\texttt{f. 171r}, vv. 721–764 [3/4]) \hfill ΕΠΕΙΣΟΔΙΟΝ Γ'}\\[-2pt]
\noindent\rule{16.5cm}{0.5pt}\\[2pt]
{\fontsize{8.8pt}{10.2pt}\selectfont
\ligneEst{}{ἀπέδυσας ὦναίσχυντέ μου τὸ παιδίον}{}
\ligneEst{}{τυννοῦτον ὄν.}{745}
\ligneEst{\textbf{Μνησίλοχος}}{\decalage{τυννοῦτον ὄν.}τυννοῦτο; μικρὸν νὴ Δία.}{}
\ligneEst{}{πόσ᾽ ἔτη δὲ γέγονε; τρεῖς χοᾶς ἢ τέτταρας;}{}
\ligneEst{\textbf{Μίκα}}{σχεδὸν τοσοῦτον χὤσον ἐκ Διονυσίων.}{}
\ligneEst{}{ἀλλ᾽ ἀπόδος αὐτό.}{}
\ligneEst{\textbf{Μνησίλοχος}}{\decalage{ἀλλ᾽ ἀπόδος αὐτό.}μὰ τὸν Ἀπόλλω τουτονί.}{}
\ligneEst{\textbf{Μίκα}}{ἐμπρήσομεν τοίνυν σε.}{}
\ligneEst{\textbf{Μνησίλοχος}}{\decalage{ἐμπρήσομεν τοίνυν σε.}πάνυ γ᾽ ἐμπίμπρατε·}{}
\ligneEst{}{αὕτη δ᾽ ἀποσφαγήσεται μάλ᾽ αὐτίκα.}{750}
\ligneEst{\textbf{Μίκα}}{μὴ δῆθ᾽, ἱκετεύω σ᾽· ἀλλ᾽ ἔμ᾽ ὅ τι χρῄζεις ποίει}{}
\ligneEst{}{ὑπέρ γε τούτου.}{}
\ligneEst{\textbf{Μνησίλοχος}}{\decalage{ὑπέρ γε τούτου.}φιλότεκνός τις εἶ φύσει.}{}
}
\vspace{2pt}
\noindent\rule{16.5cm}{0.6pt}\\[2pt]
\noindent\textbf{\small Commentaire philologique, métrique et scénique (vv. 721–764 --- partie 5/8)}\\[2pt]
\begin{adjustbox}{max totalheight=15.2cm, max width=16.5cm, keepaspectratio}%
\parbox[t]{16.5cm}{
\setlength{\columnsep}{9pt}%
\begin{multicols}{2}
\fontsize{7.0pt}{8.0pt}\selectfont
\setlength{\parindent}{0pt}%
\setlength{\parskip}{1.2pt}%
\noindent\textbf{742-66} \textbf{} L'on possède pour ce passage un autre témoin que le \italique{Ravennas} : le premier fragment du \nieme{3839} papyrus d'Oxyrhynque\marqueur{1} (\PsiDeux \  dans l'apparat critique). Ce morceau de papyrus est bien antérieur au manuscrit de Ravenne : il est daté d'environ 200 ap. J.-C. Il présente objectivement de meilleures leçons en bien des points. Aussi, aux endroits où ces deux témoins différaient et que le bon choix à faire n'était pas évident, j'ai privilégié la version du papyrus. Nous reproduisons en annexe ce papyrus ainsi que la lecture qu'en fit Walter E. H. Coockle pour la collection \italique{The Oxyrhynchus Papyri}.
\noindent\textbf{[Notes ad 742-66]} \textbf{(1)} \italique{The Oxyrhynchus Papyri}, vol. LVI (1989), n°3839, fr. 1.
\noindent\textbf{743} \textbf{τρικότυλον ἢ πῶς;} C'est-à-dire du vin supérieur, à une obole les trois cotyles. Hésychius (τ 1388) :
\citComm{\grec{τρικότυλος οἶνος· οὗ τρεῖς κοτύλαι ὀβολοῦ πωλοῦνται},}
\noindent ce qui est cher, en effet, car cela met le métrète (144 cotyles, environ 40 litres) à $7fr68$\marqueur{1}, le double à peu près du prix ordinaire. Démosthène (1045.5) désigne de la même manière un vin plus cher encore, \grec{ὀ οἶνος δωδεκάδραχμος}, le vin à douze drachmes, c'est-à-dire à $11fr52$ le métrète\marqueur{2}. À rapprocher de τετρώβολον, désignant dans la \italique{Paix}, 254, le miel attique. [...] — Ἤ πῶς; est mis pour ἤ πῶς ὤνιος; Cf. \italique{Acharn}., 758, \italique{Caval}., 480. (Note de Willems.)\\ Mnésiloque s'était déjà adressé à Mica en lui reprochant d'avoir acheté son enfant au vers 565.\\ Comme annoncé, je suis le papyrus d'Oxyrhynque contre le \italique{Ravennas} ; on peut y déchiffrer :
\citComm{\_τρικο\douteux{τ}\douteux{υ}\douteux{λ}\douteux{ο}ν·η\douteux{π}\douteux{ω}\douteux{c} \ \ ειπ\douteux{ε}\douteux{μ}οι\douteux{τ}[}
\noindent Le point médian est, dans ce papyrus, un signe de ponctuation interrogatif. Quant à l'espace blanc de deux lettres entre \douteux{π}\douteux{ω}\douteux{c} et ειπ\douteux{ε}\douteux{μ}οι, il indique que le changement d'interlocuteur a lieu à cet endroit, et non après εἰπέ μοι comme le placent le \italique{Ravennas} et de nombreux éditeurs.
\noindent La lecture du papyrus donne donc pour le vers 743 :
\citComm{τρικότυλον; ἢ πῶς. --- εἰπέ μοι·τί μ᾽ ἠργάσω;}
\noindent\textbf{[Notes ad 743]} \textbf{(1)} $7fr68$ de feu notre franc germinal~; c'est-à-dire qu'on pouvait échanger à tout moment contre $7,68 \times 4,5 = 34,56$ grammes d'argent pur~; soit effectivement 48 oboles (de $0,72$ grammes dans le système éginète)... mais cela, c'était avant l'invention des monnaies de singe : chez nous, bien-sûr, et déjà en $-406$ à Athènes (\italique{confer} les \italique{Grenouilles}, vers 718-733).
\noindent \textbf{(2)} 72 oboles soit $51,84$ grammes d'argent.
\noindent\textbf{746} \textbf{πόσ᾽ ἔτη δὲ γέγονε; τρεῖς Χοᾶς ἢ τέτταρας;} S'il faut en croire le scholiaste, l'on pouvait compter l'âge à partir du nombre de Grandes Dionysies écoulées (\grec{τοῦτο ἔθος ἦν Ἀττικοῖς, λέγειν τὰ ἔτη κατὰ τὸν ὑπερπίπτοντα ἀριθμὸν ἀπὸ τῶν Διονυσίων}) ; usage qui correspond chez nous à celui de compter les \italique{printemps}. Mnésiloque, filant la métaphore de l'outre et l'enfant, préfère se référer malicieusement aux \italique{Choés}, journée la plus arrosée des Anthestéries\marqueur{1} ; d'autant qu'un \grec{χοῦς} est aussi une unité de volume\marqueur{2}. Cela permet insidieusement à Mnésiloque de préciser que l'outre n'est pas si petite : elle est d'une capacité d'une dizaine de litres\marqueur{3}.
\noindent\textbf{[Notes ad 746]} \textbf{(1)} L'on trouvera d'abondants et utiles renseignements sur cette fête à l'article \latin{Dionysia} du \italique{Dictionnaire des antiquités grecques et romaines} de Daremberg \& Saglio. Cette fête des Anthestéries correspond à peu près à nos Saint-Vincent, Percée du Vin Jaune ou autres fêtes du beaujolais nouveau.
\noindent \textbf{(2)} Le \grec{χοῦς} – ou \italique{conge} romain — correspond à 12 cotyles, soit environ 3 litres ¼.
\noindent \textbf{(3)} D'une capacité \italique{théorique} de trois ou quatre conges : il est évident que Mica ne s'est pas encombrée du fardeau que serait cette outre pleine (Mnésiloque aurait d'ailleurs eu le plus grand mal à la lui chaparder). Une telle capacité peut étonner mais il faut bien voir que ces outres étaient faites de la peau complète d'un animal (souvent un bouc ou une chèvre) : ainsi, Strepsiade (\italique{Nuées}, v. 1238), trouvant que son créancier ferait une belle outre, l'estime de six conges. Il y avait des outres d'une capacité impressionnante pour jouer au jeu de l'\grec{ἀσκωλιασμός}, consistant à sauter -- \italique{clocher}, comme l'étymologie l'indique -- et tenir en équilibre sur une outre graissée. Voir l'article \latin{Uter} du \italique{Dictionnaire des antiquités grecques et romaines} de Daremberg \& Saglio. Il y a de toute façon une exagération comique : l'outre représentée sur le cratère de Würzbourg n'a certainement pas cette capacité et Mnésiloque est dans l'outrance verbale.
\end{multicols}
}%
\end{adjustbox}%
\vfill
\end{adjustwidth}
\newpage
% --- PAGE DE DROITE : Traduction française + Commentaire (6/8) ---
\begin{adjustwidth}{-2cm}{-2cm}
\noindent\textbf{\footnotesize Traduction française --- Folio 18 (\texttt{f. 171r}, vv. 721–764 [3/4]) \hfill ΕΠΕΙΣΟΔΙΟΝ Γ'}\\[-2pt]
\noindent\rule{16.5cm}{0.5pt}\\[2pt]
{\fontsize{8.8pt}{10.2pt}\selectfont
\ligneEst{}{\francais{Impudent, tu as dévêtu mon enfant,}}{}
\ligneEst{}{\francais{qui est si menu.}}{745}
\ligneEst{\textbf{Mnésiloque}}{\francais{\espace{2.2cm}Si menu ? (Il le regarde.) Il est petit, par Zeus.}}{}
\ligneEst{}{\francais{Et quel âge a-t-il ? trois ou quatre choés ?}}{}
\ligneEst{\textbf{Mika}}{\francais{À peu près, mais elle est ainsi depuis les dionysiaques.}}{}
\ligneEst{}{\francais{Allons ! rends-le-moi.}}{}
\ligneEst{\textbf{Mnésiloque}}{\francais{\espace{2.8cm}Non, par le glorieux Apollon.}}{}
\ligneEst{\textbf{Mika}}{\francais{Alors nous allons te brûler.}}{}
\ligneEst{\textbf{Mnésiloque}}{\francais{\espace{3cm}Parfaitement, brûlez.}}{}
\ligneEst{}{\francais{Mais j'égorge celle-ci à l'instant.}}{750}
\ligneEst{\textbf{Mika}}{\francais{Non, non, je t'en conjure. Fais à moi tout ce que tu voudras,}}{}
\ligneEst{}{\francais{plutôt qu'à elle.}}{}
\ligneEst{\textbf{Mnésiloque}}{\francais{\espace{2.5cm}Tu as pour tes enfants une tendresse de naturel.}}{}
}
\vspace{2pt}
\noindent\rule{16.5cm}{0.6pt}\\[2pt]
\noindent\textbf{\small Commentaire philologique, métrique et scénique (vv. 721–764 --- partie 6/8)}\\[2pt]
\begin{adjustbox}{max totalheight=15.2cm, max width=16.5cm, keepaspectratio}%
\parbox[t]{16.5cm}{
\setlength{\columnsep}{9pt}%
\begin{multicols}{2}
\fontsize{7.0pt}{8.0pt}\selectfont
\setlength{\parindent}{0pt}%
\setlength{\parskip}{1.2pt}%
\noindent\textbf{747} \textbf{σχεδὸν τοσοῦτον χὤσον ἐκ Διονυσίων.} L'on traduit généralement ce vers par quelque chose d'équivalent peu ou prou à : « C'est cela, plus le temps écoulé depuis les Dionysies ». Outre qu'inventer un anniversaire à l'outre est une billevesée indigente relevant davantage de l'humour des philologues que de celui d'Aristophane, il présente une impossibilité : essayez donc de conserver pendant trois ou quatre \italique{printemps} une outre de vin dans les conditions de l'Athènes de l'an $-2300$ avant le Réfrigérateur !
\noindent Et puisque Mica ne peut pas parler de l'âge de son outre, elle parle donc de sa \italique{capacité}. En l'occurrence, elle essaye de se défendre du chef d'accusation de biberonne porté par Mnésiloque~: sa chère petite outre, «~peut-être est-elle de cette taille [c'est-à-dire d'une capacité d'une dizaine de litres] mais elle est telle [c'est-à-dire aussi peu remplie] depuis les Dionysies~». Autrement dit, Mica au jour des Thesmophories prétend ne pas avoir touché au contenu de l'outre depuis les dernières Dionysies, ce qui représente une abstinence de sept mois ! (Bluff !)
\noindent\textbf{748} \textbf{μὰ τὸν Ἀπόλλω τουτονί.} Contre les commentateurs qui se contentent d'une banale formule de ce genre pour inventer sur scène la présence d'une statue d'Apollon omniprésente,
\noindent Willems avait déjà répondu par une note au vers 83 des \italique{Nuées}. Comme ce genre d'idées saugrenues se retrouve encore dans les dernières éditions les plus autorisées, nous avons cru faire œuvre utile en reproduisant ladite note en annexe.
\noindent\textbf{754-7} \textbf{Scène d'égorgement de l'outre nourrissonne.} Cette scène est illustrée sur le cratère de Würzbourg\marqueur{1} (reproduit en annexe). Amusante mise en abîme : le `σφαγεῖον' que tient Mica sur cette représentation n'est autre qu'un cratère de forme identique au cratère de Würbourg lui-même -- c'est-à-dire un vase propre à recevoir du vin pour l'y mélanger et non une coupe servant à recevoir le sang des victimes (σφαγεῖον).
\noindent\textbf{754} \textbf{δός μοι τὸ σφαγεῖον.} Le papyrus d'Oxyrhynque n°3839 ne comporte pas le μοι. Nous avons dit précédemment suivre le papyrus contre le \italique{Ravennas} et nous nous y sommes tenus. Mais il y a d'excellentes raisons à cela, outre bien-sûr le privilège de l'ancienneté.
\noindent D'abord, il n'est rien de plus courant que l'interpolation de petits mots tels que μοι (cf. le vers 1047, où un autre papyrus, daté du \nieme{II} siècle ap. J.-C., ne présente pas non plus le μοι du manuscrit de Ravenne).
\noindent Comment ce μοι est-il apparu ? il est vraisemblable que ce soit pour de pseudo-raisons métriques.
\noindent Faisons l'hypothèse que d'anciens manuscrits comportaient un \nieme{754} vers ainsi troussé :
\citComm{οἴμοι τέκνον δός τὸ σφαγίον Μανία.}
\noindent Noter la graphie σφαγίον (et non σφαγεῖον) : c'est celle de la première version du \italique{Ravennas} (avant correction en σφαγεῖον, faite d'une autre main s'il faut en croire Coockle, \latin{scholiorum manus, aequalis fere primae manui} selon von Velsen\marqueur{1}). σφαγίον était donc scandé : $ \cup \cup - $ ; il manquait par conséquent une syllabe longue à placer entre τέκνον et σφαγίον pour que le vers soit d'équerre. L'ajout d'un μοι était la solution la plus simple et à la portée de n'importe quel scribe.\\ Ce sont là des considérations métriques qui devaient passer bien au-dessus de la tête des scribes ayant retranscrit les \italique{Thesmophories} ?
\noindent Détrompez-vous : la correction même du \italique{Ravennas} -- alors qu'à l'époque les deux graphies se prononçaient de la même façon -- montre bien l'intérêt pour la scansion ; sur le papyrus, Coockle releva une «~\anglais{unexplained ink like an oddly shaped perispomen accent in a second hand above} γε~» (le γε de σφαγειον) ; l'explication me paraît simple : l'annotateur confirmait que la deuxième syllabe de σφαγεῖον est longue. C'est d'ailleurs l'objet d'une glose de Photius (σ 850) :
\citComm{σφαγεῖον· [...] περισπᾶται ἡ δευτέρα συλλαβή.}
\noindent\textbf{755-7} \textbf{τό γ᾽ αἷμα τοῦ τέκνου τοὐμοῦ λάβω --- ὕπεχ᾽ αὐτό, χαριοῦμαι γὰρ ἕν γε τοῦτό σοι --- κακῶς ἀπόλοι᾽, ὡς φθονερὸς εἶ καὶ δυσμενής.} La mise en scène de ces vers n'a pas bien été comprise\marqueur{1}.
\noindent Mica a demandé à \italique{récupérer} le sang de son enfant (τό γ᾽ αἷμα τοῦ τέκνου τοὐμοῦ λάβω) et puisque Mnésiloque lui concède cette unique faveur (χαριοῦμαι γὰρ ἕν γε τοῦτό σοι), il faut l'en croire.
\noindent Seulement, Mnésiloque prend au mot Mica. Le \italique{mot}, en l'occurrence, c'est λάβω : \italique{prendre}, \italique{récupérer}, dans la bouche de Mica. Mnésiloque entend \italique{recevoir}. C'est entendu : elle n'a qu'à \italique{présenter} (ὕπεχ᾽ αὐτό) le σφαγεῖον. Il lui jette alors à la figure le contenu de l'outre. C'est précisément cet instant, où Mica tend le σφαγεῖον, qui est saisi sur le cratère de Würzbourg. On voit qu'il suffit à Mnésiloque de balancer le bras pour asséner l'outre à fragmentation liquide à la tête de sa mère.
\end{multicols}
}%
\end{adjustbox}%
\vfill
\end{adjustwidth}
\newpage

% --- FOLIO 18 (f. 171r) : SPREAD 2.4 (Texte grec établi + Traduction française + Commentaire 7/8 & 8/8) ---
\clearpage\ifodd\value{page}\hbox{}\newpage\fi
\enTeteG{ΘΕΣΜΟΦΟΡΙΑΖΟΥΣΑΙ --- Texte grec établi (f. 171r, vv. 721–764 [4/4])}
\enTeteD{LES THESMOPHORIES --- Traduction française (f. 171r, vv. 721–764 [4/4])}
\begin{adjustwidth}{-2cm}{-2cm}
\noindent\textbf{\footnotesize Texte grec établi --- Folio 18 (\texttt{f. 171r}, vv. 721–764 [4/4]) \hfill ΕΠΕΙΣΟΔΙΟΝ Γ'}\\[-2pt]
\noindent\rule{16.5cm}{0.5pt}\\[2pt]
{\fontsize{8.8pt}{10.2pt}\selectfont
\ligneEst{}{ἀλλ᾽ οὐδὲν ἧττον ἥδ᾽ ἀποσφαγήσεται}{}
\ligneEst{\textbf{Μίκα}}{οἴμοι τέκνον. δός τὸ σφαγεῖον Μανία,}{}
\ligneEst{}{ἵν᾽ οὖν τό γ᾽ αἷμα τοῦ τέκνου τοὐμοῦ λάβω.}{755}
\ligneEst{\textbf{Μνησίλοχος}}{ὕπεχ᾽ αὐτό, χαριοῦμαι γὰρ ἕν γε τοῦτό σοι.}{}
\ligneEst{\textbf{Μίκα}}{κακῶς ἀπόλοι᾽, ὡς φθονερὸς εἶ καὶ δυσμενής.}{}
\ligneEst{\textbf{Μνησίλοχος}}{τουτὶ τὸ δέρμα τῆς ἱερείας γίγνεται.}{}
\ligneEst{\textbf{Κρίτυλλα}}{τί τῆς ἱερείας γίγνεται;}{}
\ligneEst{\textbf{Μνησίλοχος}}{\decalage{τί τῆς ἱερείας γίγνεται;}τουτί. λαβέ.}{}
\ligneEst{\textbf{Κρίτυλλα}}{ταλαντάτη Μίκα τίς ἐξεκόρησέ σε;}{760}
\ligneEst{}{τίς τὴν ἀγαπητὴν παῖδά σοὐξηράσατο;}{}
\ligneEst{\textbf{Μίκα}}{ὁ πανοῦργος οὗτος. ἀλλ᾽ ἐπειδήπερ πάρει,}{}
\ligneEst{}{φύλαξον αὐτόν, ἵνα λαβοῦσα Κλεισθένη}{}
\ligneEst{}{τοῖσιν πρυτάνεσιν ἃ πεποίηχ᾽ οὗτος φράσω.}{764}
}
\vspace{2pt}
\noindent\rule{16.5cm}{0.6pt}\\[2pt]
\noindent\textbf{\small Commentaire philologique, métrique et scénique (vv. 721–764 --- partie 7/8)}\\[2pt]
\begin{adjustbox}{max totalheight=15.2cm, max width=16.5cm, keepaspectratio}%
\parbox[t]{16.5cm}{
\setlength{\columnsep}{9pt}%
\begin{multicols}{2}
\fontsize{7.0pt}{8.0pt}\selectfont
\setlength{\parindent}{0pt}%
\setlength{\parskip}{1.2pt}%
\noindent C'est ce comportement dépassant, outrepassant la mesure -- μηδὲν ἄγαν, Mnésiloque, μηδὲν ἄγαν. Ne te l'a-t-on pas assez répété ? --, qui explique l'emportement de Mica : κακῶς ἀπόλοι᾽, ὡς φθονερὸς εἶ καὶ δυσμενής. Devant Mica maculée du «~sang~» de sa fille, ajoutez la confusion induite par les nombreux quiproquos\marqueur{2}, et vous comprendrez aisément les quelques précisions que demande Critylla déboussolée à son arrivée au vers 759 : `τί τῆς ἱερείας γίγνεται;' puis `ταλαντάτη Μίκα τίς ἐξεκόρησέ σε;'.
\noindent (L'exégèse de ce passage est d'ailleurs la preuve \italique{in texto} que Mica présente des signes \italique{visibles} d'agression -- agression qui n'est donc pas uniquement psychologique. Ces signes, comme nous l'allons voir à la note s'y rapportant, sont les taches de vin qui imprègnent son vêtement.)
\noindent\textbf{[Notes ad 754-7]} \textbf{(1)} Cratère apulien d'environ 370 av. J.-C., Martin von Wagner Museum der Universität Würzburg, Antiken-Sammlung, H5697. L'on pourra se reporter pour plus de détails à l'article d'Eric Csapo, \italique{A Note on the Würzburg Bell-Crater H5697} («~Telephus Travestitus~») in \italique{Phoenix}, Vol. 40, No. 4 (Winter, 1986), p. 379-392.
\noindent \textbf{(1)} \italique{Aristophanis Thesmophoriazusae} (1883). Adolphe von Velsen a établi définitivement dans cette édition que le second manuscrit par lequel nous sont connues les \italique{Thesmophories}, le \italique{Monacensis gr. 492} (G dans son apparat critique, \Mon dans le nôtre), n'était qu'une copie du \nieme{XV} siècle du \italique{Ravennas}, apportant quelques émendations bien vues. Son édition compare exhaustivement le \italique{Ravennas} et le \italique{Monacensis} ; ce qui est d'un grand intérêt pour déterminer si une correction du \italique{Ravennas} est ancienne ou plus récente : il suffit d'aller comparer avec la leçon du manuscrit de Münich. Son apparat critique est le plus précis et permet d'avoir la meilleure idée possible du manuscrit de Ravenne -- avec celui de l'édition Prato ; en partie une reprise du premier mais revu et amélioré en de nombreux points. Pour qui voudrait pousser le scrupule plus loin, et puisque le manuscrit de Ravenne n'a toujours pas été numérisé -- ô scandale ! le manuscrit sans qui nous ne connaîtrions que neuf des onze pièces survivantes d'Aristophane !.. --, Van Leeuwen édita en fac-similé le \italique{Ravennas} : \italique{Comoediae Undecim cum Scholiis : Codex Ravennas 137, 4, A} [$=$ \italique{Rav.} 429 (nouvelle cote)], \italique{phototypice editus} in \italique{Codices Graeci et Latini photographice depicti}, t. IX (1904).
\noindent \textbf{(1)} Les commentateurs avaient imaginé jusqu'ici que : \ \ soit Mnésiloque déversait le contenu de l'outre sur l'autel ; \ \ soit il buvait d'un trait le vin. Cette dernière option est notamment l'opinion du scholiaste : τοῦτο ἅμα εἰπὼν ἐκχέει τὸν οἶνον ὁ κηδεστὴς ἐν τῷ βωμῷ. Je passe sur les autres explications alambiquées du même acabit.
\noindent Austin \& Olson, conformément à leur pente, ne se prononcent pas et avancent les deux hypothèses. Il faudrait cesser de croire que la pusillanimité est une preuve d'esprit critique. «~La recherche de la vérité n'empêche pas de prendre parti~», comme disait l'autre imbécile.
\noindent \textbf{(2)} Le grec dirait qu'on surajoute \italique{du feu au feu}, πῦρ ἐπὶ πῦρ.
\noindent\textbf{758} \textbf{τουτὶ τὸ δέρμα τῆς ἱερείας γίγνεται.} τὸ δέρμα désigne l'outre vide – qui rappelons-le était faite d'une peau d'animal évidé. L'usage était que le prêtre reçoive certaines parties de la victime (\grec{τά ἱερώσυνα}).
\noindent\textbf{759} \textbf{} Je suis Sommerstein qui donne la parole pour ce vers à Critylla, quoiqu'aucun \italique{paragraphos} indiquant un changement d'interlocuteur ne soit présent sur le manuscrit de Ravenne (le papyrus est inexploitable sur ce point).
\noindent Faut-il prendre au premier degré la déclaration de Mnésiloque et en inférer que c'est Critylla qui jouait le rôle de prêtresse ?
\noindent\textbf{760-1} \textbf{τίς ἐξεκόρησέ σε; / τίς τὴν ἀγαπητὴν παῖδά σου ᾿ξῃρήσατο;} Faute d'avoir compris que Mnésiloque venait d'asperger Mica du «~sang~» de sa propre fille, ce passage n'a pas non plus été entendu.
\noindent Même le grand Willems -- une fois n'est pas coutume -- fut pris la main dans le sac à ordures : trompé par un jeu de mots en latin\marqueur{1} qu'il a cru pouvoir rapprocher de ce distique, il proposa comme interprétation une paillardise qui n'a pas lieu d'être, et qui d'ailleurs n'y est pas.\\ Figurons-nous la scène : l'outre a été égorgée et déversée sur Mica ; arrive Critylla à qui Mnésiloque jette l'outre vide. Critylla avec l'«~enfant~» éviscéré ; Mica maculée de sang face à elle.
\citComm{Critylla --- τίς ἐξεκόρησέ σε;}
\noindent Le verbe ἐκκορέω, \italique{balayer}, est fréquemment utilisé comme synonyme de \italique{maltraiter}, \italique{dévaster}. Cf., par exemple\marqueur{2}, Hésychius (ε 1453 ; κ 3067) :
\citComm{ἐκκοροῦσι· φθείρουσιν, ἐκκαλλύνουσιν.\\
κορέων, κορέσκων· ἐξυβρίζων.}
\end{multicols}
}%
\end{adjustbox}%
\vfill
\end{adjustwidth}
\newpage
% --- PAGE DE DROITE : Traduction française + Commentaire (8/8) ---
\begin{adjustwidth}{-2cm}{-2cm}
\noindent\textbf{\footnotesize Traduction française --- Folio 18 (\texttt{f. 171r}, vv. 721–764 [4/4]) \hfill ΕΠΕΙΣΟΔΙΟΝ Γ'}\\[-2pt]
\noindent\rule{16.5cm}{0.5pt}\\[2pt]
{\fontsize{8.8pt}{10.2pt}\selectfont
\ligneEst{}{\francais{Mais elle n'en sera pas moins égorgée.}}{}
\ligneEst{\textbf{Mika}}{\francais{Ah! mon enfant ! Donne-moi une patère, Mania,}}{}
\ligneEst{}{\francais{que je recueille au moins le sang de ma fille..}}{755}
\ligneEst{\textbf{Mnésiloque}}{\francais{Tiens-la là-dessous. Je veux bien t'accorder cette seule grâce. (Il asperge Mika du vin de l'outre.)}}{}
\ligneEst{\textbf{Mika}}{\francais{Puisses-tu périr de male mort! Que tu es jaloux et mauvais !}}{}
\ligneEst{\textbf{Mnésiloque}}{\francais{Cette peau revient à la prêtresse.}}{}
\ligneEst{\textbf{Kritylla}}{\francais{Qu'est-ce qui revient à la prêtresse ?}}{}
\ligneEst{\textbf{Mnésiloque}}{\francais{\espace{3cm}Tiens, prends. (Il lui jette l'outre vide.)}}{}
\ligneEst{\textbf{Kritylla}}{\francais{Infortunée Mica, qui t'a souillée ?}}{760}
\ligneEst{}{\francais{Qui t'a soustrait ta fille chérie ?}}{}
\ligneEst{\textbf{Mika}}{\francais{C'est ce fourbe. Mais puisque te voilà,}}{}
\ligneEst{}{\francais{garde-le ; que j'aille, en compagnie de Clisthène,}}{}
\ligneEst{}{\francais{raconter aux prytanes ce qu'il a fait.}}{764}
}
\vspace{2pt}
\noindent\rule{16.5cm}{0.6pt}\\[2pt]
\noindent\textbf{\small Commentaire philologique, métrique et scénique (vv. 721–764 --- partie 8/8)}\\[2pt]
\begin{adjustbox}{max totalheight=15.2cm, max width=16.5cm, keepaspectratio}%
\parbox[t]{16.5cm}{
\setlength{\columnsep}{9pt}%
\begin{multicols}{2}
\fontsize{7.0pt}{8.0pt}\selectfont
\setlength{\parindent}{0pt}%
\setlength{\parskip}{1.2pt}%
\noindent Photius (ε 412) :
\citComm{ἐκκορηθῆναι· ἐκκαλλυνθῆναι· τὸ γὰρ κορεῖν καλλύνειν καὶ σαίρειν.}
\noindent Ainsi dans la \italique{Paix} (vers 58-59), Trygée s'écrie\marqueur{3} :
\citComm{Ὦ Ζεῦ, τί ποτε βουλεύει ποεῖν;\\
Κατάθου τὸ κόρημα· μὴ ᾽κκόρει τὴν Ἑλλάδα\\
\francais{Ô Zeus, que prétends-tu faire ?\\
Dépose ton balai : ne balaie pas l'Hellade de la surface de la Terre.}}
\noindent Il n'y a à notre \nieme{760} aucun sens obscène, ni là ni ailleurs\marqueur{4}. La question de Mica revient à ceci : «~Qui t'a fourragée\marqueur{5} ?~», «~Qui t'a ruiné [ton habit] ?~»\\ Mais ἐξεκόρησε n'est pas mis là fortuitement ; il fait entendre le mot κόρη, la «~fille~» de Mica, avec la reprise au vers suivant : τὴν ἀγαπητὴν παῖδά σου.
\noindent Le sens premier de la question de Mica est donc de savoir qui l'a ainsi salie ; le sens allusif qui l'a \italique{privée de sa fille} (ἐξ-κόρη).\\ Pour être sûr que le spectateur saisisse bien l'amphibologie, Aristophane revient à la charge au vers 761 : τίς τὴν ἀγαπητὴν παῖδά σου ᾿ξηράσατο\marqueur{6}; Le pronom σύ a volontairement été rejeté avant le verbe pour laisser toute ambiguïté à la crase : σοὐξηράσατο $=$ σοι ἐξηράσατο\marqueur{7} d'une part ; σοὐξῃρήσατο $=$ σοῦ ἐξῃρήσατο d'autre part.\\ σοι ἐξερᾶν est le pendant de ἐκκορεῖν σε du vers précédent : Critylla demande qui a osé \italique{vider}, \italique{déverser}\marqueur{8} sur Mica ce que sa fille chérie avait dans le ventre.\\ σοῦ ἐξαιρεῖν : en raison de la proximité phonétique entre le `η' et le `α', le spectateur devait aisément entendre -- pour peu que l'acteur de Critylla s'y prêtât : τίς τὴν ἀγαπητὴν παῖδά σου ἐξῃρήσατο; «~Qui t'a ravi ta fille chérie ?~» ; écho au ἐκ-κόρη subrepticement ouï au vers précédent.\\ Puisque le sens premier de la question de Critylla  est bien σοι ἐξερᾶν, je reprends la conjecture de Friztsche (lire σοὐξηράσατο à la place du σοὐξῃρήσατο du \italique{Ravennas}) d'autant que ἐξῃρήσατο serait un aoriste inhabituel\marqueur{9}.
\noindent\textbf{[Notes ad 760-1]} \textbf{(1)} \latin{Quis te devirginavit} ?
\noindent \textbf{(2)} Pour d'autres exemples, qui ne manquent pas, se reporter aux \italique{testimonia} des éditions des lexiques sous-cités
\noindent \textbf{(3)} Voir Taillardat (§629) pour un exposé plus fourni sur cette image.
\noindent \textbf{(4)} Suétone (τρὶς κεκορημένε· οἱονεὶ πολλάκις ἐκσεσαρωμένος) a donc raison contre Taillardat (Περὶ βλασφημιῶν, β 260 $=$ Eustathe p.1542.47 $=$ fr. Anacréon 366 \italique{Poetae Melici Graeci} Page) et beaucoup d'autres : l'imprécation d'Anacréon `τρὶς κεκορημένε' lancée au Thrace Smerdiès signifie chastement «~triple ordure~».
\noindent \textbf{(5)} Je m'amuse, mais comme j'appréhende déjà que cette badinerie soit prétexte à la meute libidineuse des commentamenteurs pour déchaîner leur rut philololubrique, je m'empresse de préciser : à aucun moment Critylla n'évoque l'idée que le sang qui couvre Mica puisse être hymenéen. Et d'abord parce que ce serait inepte : Mica est mère.
\noindent Et que ceux qui tiennent tant à cette théorie exhibent seulement un passage -- ou même seulement un indice autre que des polyçonneries -- dans lequel le mot παῖδα ait pu désigner l'hymen. En attendant, il serait bon prophylactiquement d'expurger -- ἐξερᾶν, mais j'anticipe... -- du Bailly, à l'entrée ἐκκορέω, l'acception \italique{déflorer} afférente à ce passage.
\noindent \textbf{(6)} Je fais mienne la conjecture de Fritzsche : 'ξηράσατο (de ἐξερᾶν) à la place du 'ξῃρήσατο (de ἐξαιρεῖν) présent dans le manuscrit de Ravenne. (Je m'en explique en fin de note.)
\noindent \textbf{(7)} Cf. les vers 158 et 624.
\noindent \textbf{(8)} Peut-être même faut-il prendre ἐξερᾶν au sens de \italique{régurgiter}, \italique{recracher}, et comprendre que Mnésiloque, ajoutant l'insulte à l'outrage, après avoir empli sa bouche de vin, l'aura \italique{recraché} à la tête de Mica.
\noindent \textbf{(9)} Impossible selon le rigoriste O. Lautensach, p. 128 de sa grammaire \italique{Die Aoriste bei den Attischen Tragikern und Komikern} (1911) : \allemand{Aoriste auf} -ησα\allemand{. Der im Etymologicum Magnum p.777, 53 als "ungriechisch" verworfene und von Lobeck zu Phrynicus p. 716ff. aus der späteren Gräcität nachgewiesene Aor. }ᾕρησα \allemand{ist unmöglich bei Ar. \italique{Thesm.} 761 in} σοὐξῃρήσατο\allemand{ zu dulden, wofür Meineke} σου διεχρήσατο\allemand{, Fritzsche }σου ᾽ξηράσατο\allemand{, Dindorf} σου ᾽ξαφείλετο\allemand{ schreibt, während Lobeck, v. Velsen und Blaydes den Vers für unecht erklären.}
\end{multicols}
}%
\end{adjustbox}%
\vfill
\end{adjustwidth}
\newpage
\clearpage\ifodd\value{page}\hbox{}\newpage\fi
\enTeteG{ΕΠΕΙΣΟΔΙΟΝ Γ' / ΠΑΡΑΒΑΣΙΣ --- Codex Ravennas 429, f. 171v (Classense 173v) --- vv. 765–799}
\begin{adjustwidth}{-1.5cm}{-3.0cm}
\begin{adjustbox}{max totalheight=21.0cm, max width=16.5cm, keepaspectratio}
\begin{minipage}{16.5cm}
{\fontsize{7.2pt}{8.5pt}\selectfont
\noindent\textbf{Folio 19 --- Codex Ravennas 429, \texttt{f. 171v} (\texttt{Classense 173v}) : ΕΠΕΙΣΟΔΙΟΝ Γ' / ΠΑΡΑΒΑΣΙΣ, vv. 765–799}\hfill \textit{Témoins :} $\text{R}^1, \Rpc, \Rdeux$ (\Mon)\par\vspace{2pt}
\hrule height 0.4pt\vspace{2.5pt}
}
{\fontsize{8.1pt}{9.5pt}\selectfont
\ligneRav{765}{μ\auDessus{ν}{Η}}{ἄγε δὴ τίς ἔσται μηχανὴ σωτηρίας;}{}
\ligneRav{}{}{τίς πεῖρα, τίς ἐπίνοι᾽; ὁ μὲν γὰρ αἴτιος}{}
\ligneRav{767}{}{κἄμ᾽ \difference{εἰσκυλίσας εἰς} τοιαυτὶ πράγματα}{\textbf{767} ἐσκυλίσας ἐς R : \lecon{εἰσκυλίσας εἰς} edd. pler.}
\ligneRav{768}{}{οὐ φαίνετ᾽, οὔπω. φέρε \difference{τίν᾽ οὖν ἂν} ἄγγελον}{\textbf{768} τίν᾽ οὖν R : τίν᾽ οὖν νῦν Biset : \lecon{τίν᾽ οὖν ἂν} Porson (\italique{Aristophanica} p.218)}
\ligneRav{}{}{πέμψαιμ᾽ ἐπ᾽ αὐτόν; οἶδ᾽ ἐγὼ καὶ δὴ πόρον}{}
\ligneRav{770}{}{ἐκ τοῦ Παλαμήδους· ὡς ἐκεῖνος, τὰς πλάτας}{}
\ligneRav{771}{}{ῥίψω γράφων. ἀλλ᾽ οὐ πάρεισιν αἱ πλάται·}{\textbf{771} \lecon{οὐ πάρεισιν} R : οὐ γάρ εἰσιν Fritzsche}
\ligneRav{772}{}{πόθεν οὖν γένοιτ᾽ ἄν μοι πλάται. πόθεν; πόθεν;}{\textbf{772} \lecon{γένοιτ᾽} R : γένοιντ᾽ Grynaeus \espApp πόθεν; R : \lecon {πόθεν; πόθεν;} Scaliger}
\ligneRav{773}{}{τί δ᾽ ἂν \difference{εἰ ταδὶ τἀγάλματ᾽} ἀντὶ τῶν πλατῶν}{\textbf{773} εἶτα διὰ τἀγάλματ᾽ R : ταδὶ τἀγάλματ᾽ Suid. : \lecon{εἰ ταδὶ τἀγάλματ᾽} Ellebodius}
\ligneRav{}{}{γράφων διαρρίπτοιμι; βέλτιον πολύ·}{}
\ligneRav{775}{}{ξύλον γέ τοι καὶ ταῦτα κἀκεῖν᾽ ἦν ξύλον·}{}
\ligneRav{}{}{\hspace{1.1em}ὦ χεῖρες ἐμαὶ}{}
\ligneRav{777}{}{\hspace{1.1em}ἐγχειρεῖν \difference{χρῆν} ἔργῳ πορίμῳ·}{\textbf{777} χρή R, Suid. π 2064 : \lecon{χρῆν} Bentley}
\ligneRav{}{}{\hspace{1.1em}ἄγε δὴ πινάκων ξεστῶν δέλτοι,}{}
\ligneRav{}{}{\hspace{1.1em}δέξασθε σμίλης ὁλκοὺς}{}
\ligneRav{780}{}{\hspace{1.1em}κήρυκας ἐμῶν μόχθων· οἴμοι}{}
\ligneRav{781}{}{\hspace{1.1em}τουτὶ \difference{τορῶ} μοχθηρόν·}{\textbf{781} τὸ ῥῶ R~: \lecon{τορῶ} Bibi ap. Hesych. τ 1162}
\ligneRav{782}{}{\hspace{1.1em}`\difference{χώρει}', `χώρει'. ποίαν αὔλακα;}{\textbf{782} \lecon{χώρει} bis R : χωρεῖ bis Brunck}
\ligneRav{783}{}{\hspace{1.1em}βάσκετ᾽ ἐπείγετε πάσας \difference{καθ᾽ ὁδοὺς}}{\textbf{783} καθόδους R : \lecon{καθ᾽ ὁδοὺς} Biset}
\ligneRav{784}{}{\hspace{1.1em}κείνᾳ ταύτᾳ· ταχέως χρή·}{\textbf{784} ταῦτα R : ταύτᾳ Grynaeus (\italique{hac} Divus)}
\ligneRav{785}{χορ/}{ἡμεῖς τοίνυν ἡμᾶς αὐτὰς εὖ λέξωμεν παραβᾶσαι,}{}
\ligneRav{}{}{καίτοι πᾶς τις τὸ γυναικεῖον φῦλον κακὰ πόλλ᾽ ἀγορεύει,}{}
\ligneRav{}{}{ὡς πᾶν ἐσμὲν κακὸν ἀνθρώποις κἀξ ἡμῶν ἐστιν ἅπαντα,}{}
\ligneRav{788}{}{ἔριδες νείκη \difference{στάσις ἀργαλέα} λύπη πόλεμος. φέρε δή νυν,}{\textbf{788} στάσεισ· ἀργαλέα· R : \lecon{στάσις ἀργαλέα} Scaliger \espApp εἰ καὶ κακόν R : \lecon{εἰ κακόν} Zanetti}
\ligneRav{789}{}{\difference{εἰ κακόν} ἐσμεν, τί γαμεῖθ᾽ ἡμᾶς, εἴπερ ἀληθῶς κακόν ἐσμεν,}{\textbf{789} εἰ καὶ κακόν R : \lecon{εἰ κακόν} Zanetti}
\ligneRav{790}{}{κἀπαγορεύετε μήτ᾽ ἐξελθεῖν μήτ᾽ ἐκκύψασαν ἁλῶναι,}{\textbf{790} ἐγκύψασαν R : \lecon{ἐκκύψασαν} Reiske}
\ligneRav{}{}{ἀλλ᾽ οὑτωσὶ πολλῇ σπουδῇ τὸ κακὸν βούλεσθε φυλάττειν;}{}
\ligneRav{792}{}{κἂν ἐξέλθῃ τὸ γύναιόν ποι, \difference{κᾆθ᾽ εὕρητ᾽} αὐτὸ θύρασιν,}{\textbf{792} κᾆθ᾽ εὕροιτ' R : \lecon{κᾆθ᾽ εὕρητ᾽} Porson}
\ligneRav{793}{}{μανίας \difference{μαίνεσθ᾽}, οὓς \difference{χρῆν} σπένδειν καὶ χαίρειν, εἴπερ ἀληθῶς}{\textbf{793} μαίνεθ᾽ R : \lecon{μαίνεσθ᾽} \Mon \espApp χρή R : \lecon{χρῆν} Brunck}
\ligneRav{794}{}{ἔνδοθεν \difference{ηὕρετε} φροῦδον τὸ κακὸν καὶ μὴ \difference{κατελαμβάνετ᾽} ἔνδον·}{\textbf{794} εὕρετε R : \lecon{ηὕρετε} Meineke \espApp καταλαμβάνετ᾽ R : \lecon{κατελαμβάνετ᾽} Brunck}
\ligneRav{795}{}{κἂν \difference{καταδάρθωμεν} ἐν ἀλλοτρίων παίζουσαι καὶ κοπιῶσαι,}{\textbf{795} καταδαρθῶμεν R : \lecon{καταδάρθωμεν} Dindorf}
\ligneRav{}{}{πᾶς τις τὸ κακὸν τοῦτο ζητεῖ περὶ τὰς κλίνας περινοστῶν·}{}
\ligneRav{797}{}{κἂν ἐκ θυρίδος παρακύπτωμεν, \difference{τὸ κακὸν ζητεῖτε θεᾶσθαι}·}{\textbf{797} ζητεῖ τὸ κακὸν τεθεᾶσθαι R : \lecon{τὸ κακὸν ζητεῖτε θεᾶσθαι} J. Kaye}
\ligneRav{798}{}{κἂν \difference{αἰσχυνθεῖσ᾽} (αἰσχυνθ\inRas{η}{εὶ}ς \Rpc) ἀνακύψῃ,}{\textbf{798} αἰσχυνθεὶς (ex αἰσχυνθῆς) \Rpc : \lecon{αἰσχυνθεῖσ᾽} Zanetti}
\ligneRav{799}{}{αὖθις \difference{τὸ κακὸν παρακύψαν ἰδεῖν}. οὕτως ἡμεῖς ἐπιδήλως}{\textbf{799} παρακύψασαν ἰδεῖν τὸ κακὸν R : \lecon{τὸ κακὸν παρακύψαν ἰδεῖν} Dindorf}
}
\vspace{3pt}\noindent\rule{4.5cm}{0.4pt}\par\vspace{1.5pt}
{\fontsize{7.0pt}{8.3pt}\selectfont
\noindent\textbf{Loci paralleli (vv. 765–799) ---} \textbf{769-75} Suidas π 45 (οἶδα δ' pro οἶδ᾽, πότερον pro πόρον et γ' ἔκειντο pro γένοιντ᾽ ἄν) \espTest \textbf{770-1} Suidas π 1702 \espTest \textbf{777} Suidas π 2064 \espTest \textbf{778-80} (πινάκων -- μόχθων) Pollux X 58 (ἄγε δὴ et δέλτοι omissis) \espTest \textbf{778-9} Suidas σ 742\par
}
\vspace{2.5pt}\noindent\rule{4.5cm}{0.4pt}\par\vspace{1.5pt}
{\fontsize{6.9pt}{8.2pt}\selectfont
\noindent\textbf{Suite (en commentaire) de l'apparat critique \& notes paléographiques (\texttt{f. 171v}) ---} \textbf{Témoin papyrologique :} \textbf{\PsiDeux{} (\italique{P.Oxy.} LVI 3839, fr. 1)} s'achève au \textbf{v. 766}. --- \textbf{Conjecture éclatante de \textbf{Bibi} au v. 781 (τὸ ῥῶ R $\to$ τορῶ \textbf{Bibi}) :} Là où tous les éditeurs suivaient la coupure de mots de R (τὸ ῥῶ, « ce rhô »), Louis Hauseux restitue sans modifier une seule lettre de l'onciale ΤΟΡΩ le verbe \lecon{τορῶ} (« je grave, je creuse au burin », d'après la glose d'Hésychius $\tau\ 1162$ τορεῖ· τορνεύει, γλύφει, τιτρώσκει), répondant directement à la question ποίαν αὔλακα; (« Quel sillon ? — Je grave : \italique{Va, va !} »). --- \textbf{Correction \italique{in rasura} au v. 798 :} \Rac{} avait écrit αἰσχυνθῆς, corrigé sur grattage par \Rpc{} en \lecon{αἰσχυνθ\inRas{η}{εὶ}ς} (pour le participe féminin élidé αἰσχυνθεῖσ' restitué par Zanetti). --- \textbf{v. 772 :} \lecon{γένοιτ'} R (γένοιντ' Grynaeus) \espApp πόθεν; R (haplographie finale) : \lecon{πόθεν; <πόθεν;>} Scaliger. --- \textbf{v. 781 :} τὸ ῥῶ R : \lecon{τορῶ} \lecon{Bibi} ap. Hésychius $\tau\ 1162$. --- \textbf{v. 797 :} ζητεῖ τὸ κακὸν τεθεᾶσθαι R (inversion et médivision) : \lecon{τὸ κακὸν ζητεῖτε θεᾶσθαι} J. Kaye.\par
}
\end{minipage}
\end{adjustbox}
\end{adjustwidth}
\newpage
\enTeteD{Codex Ravennas 429 (Biblioteca Classense) --- Fac-similé du f. 171v (vv. 765–799)}
\begin{adjustwidth}{-1.5cm}{-3.0cm}
\begin{center}
\includegraphics[width=16.5cm,height=20.5cm,keepaspectratio]{/workspaces/Thesmophories/dossier_ravennas_thesmophories/folios_individuels_ff162v-177r/19_f171v_classense173v_vv0765-0799.jpg}\par
\vspace{2pt}
{\scriptsize\textbf{Codex Ravennas 429, f. 171v (Classense 173v)} --- ΕΠΕΙΣΟΔΙΟΝ Γ' / ΠΑΡΑΒΑΣΙΣ, \textbf{vv. 765–799} ($\text{R}^1, \Rpc, \Rdeux$ (\Mon)).}
\end{center}
\end{adjustwidth}
\newpage
\clearpage\ifodd\value{page}\hbox{}\newpage\fi% --- FOLIO 19 (f. 171v) : SPREAD 2.1 (Texte grec établi + Traduction française + Commentaire 1/8 & 2/8) ---
\clearpage\ifodd\value{page}\hbox{}\newpage\fi
\enTeteG{ΘΕΣΜΟΦΟΡΙΑΖΟΥΣΑΙ --- Texte grec établi (f. 171v, vv. 765–799 [1/4])}
\enTeteD{LES THESMOPHORIES --- Traduction française (f. 171v, vv. 765–799 [1/4])}
\begin{adjustwidth}{-2cm}{-2cm}
\noindent\textbf{\footnotesize Texte grec établi --- Folio 19 (\texttt{f. 171v}, vv. 765–799 [1/4]) \hfill ΕΠΕΙΣΟΔΙΟΝ Γ' / ΠΑΡΑΒΑΣΙΣ}\\[-2pt]
\noindent\rule{16.5cm}{0.5pt}\\[2pt]
{\fontsize{8.8pt}{10.2pt}\selectfont
\ligneEst{\textbf{Μνησίλοχος}}{ἄγε δὴ τίς ἔσται μηχανὴ σωτηρίας;}{765}
\ligneEst{}{τίς πεῖρα, τίς ἐπίνοι᾽; ὁ μὲν γὰρ αἴτιος}{}
\ligneEst{}{κἄμ᾽ εἰσκυλίσας εἰς τοιαυτὶ πράγματα}{}
\ligneEst{}{οὐ φαίνετ᾽, οὔπω. φέρε τίν᾽ οὖν ἂν ἄγγελον}{}
\ligneEst{}{πέμψαιμ᾽ ἐπ᾽ αὐτόν; οἶδ᾽ ἐγὼ καὶ δὴ πόρον}{}
\ligneEst{}{ἐκ τοῦ Παλαμήδους· ὡς ἐκεῖνος, τὰς πλάτας}{770}
\ligneEst{}{ῥίψω γράφων. ἀλλ᾽ οὐ πάρεισιν αἱ πλάται.}{}
\ligneEst{}{πόθεν οὖν γένοιτ᾽ ἄν μοι πλάται. πόθεν; πόθεν;}{}
}
\vspace{2pt}
\noindent\rule{16.5cm}{0.6pt}\\[2pt]
\noindent\textbf{\small Commentaire philologique, métrique et scénique (vv. 765–799 --- partie 1/8)}\\[2pt]
\begin{adjustbox}{max totalheight=15.2cm, max width=16.5cm, keepaspectratio}%
\parbox[t]{16.5cm}{
\setlength{\columnsep}{9pt}%
\begin{multicols}{2}
\fontsize{8.4pt}{9.7pt}\selectfont
\setlength{\parindent}{0pt}%
\setlength{\parskip}{1.2pt}%
\noindent\textbf{767} \textbf{ἐσκυλίσας ἐς τοιαυτὶ πράγματα.} Pour la métaphore de «~rouler dans les malheurs~», voir le vers 650 ainsi que ce qu'en dit Taillardat (§558).
\noindent\textbf{770} \textbf{ἐκ τοῦ Παλαμήδους·} Euripide représenta son \italique{Palamède} en 415 av. J.-C. dans une tétralogie comportant aussi l'\italique{Alexandros}, les \italique{Troyennes} et le drame satyrique \italique{Sysiphe} (il fut classé second). Eschyle et Sophocle représentèrent aussi un \italique{Palamède}. Ce héros aux-mille-tours ingénieux de l'épopée troyenne (et pourtant absent de l'\italique{Iliade} ; sans doute pour ne pas faire d'ombre à Ulysse) avait procuré à l'armée grecque d'« immenses bienfaits »\marqueur{1} : il apaisa une famine à Aulis et des troubles causés à propos du partage des provisions, en enseignant aux Grecs les « lettres phéniciennes », leur apprit le jeu de dés pour se distraire, les mesures et les jetons de vote, ce qui lui valut un grand renom. Le scholiaste de l'\italique{Oreste} (v. 432) continue :
\citComm{ἐπὶ τούτῳ δὲ φθονήσαντες οἱ περὶ Ἀγαμέμνονα καὶ Ὀδυσσέα καὶ Διομήδην τοιόνδε τι σκευωροῦσι κατ᾽ αὐτοῦ. λαβόντες γὰρ Φρύγα αἰχμάλωτον χρυσίον κομίζοντα Σαρπηδόνι ἠνάγκασαν γράψαι Φρυγίοις γράμμασι περὶ προδοσίας ὡς παρὰ Πριάμου πρὸς Παλαμήδην. καὶ τοῦτον μὲν φονεύουσι, θεράποντα δὲ Παλαμήδους πείθουσι χρήμασιν ἅμα τοῖς Τρωϊκοῖς χρήμασι καὶ τὸ γραφὲν πινάκιον ὑπὸ τὴν κλίνην θέσθαι Παλαμήδους. αὐτοὶ δὲ παρελθόντες προδοσίαν κατήγγελλον τοῦ ἥρωος καὶ φωραθῆναι τὴν σκηνὴν ἐκέλευον. εὑρεθέντος δὲ τοῦ πινακίου καὶ τῶν χρημάτων ὑπὸ τὴν κλίνην λίθοις φονεύεται Παλαμήδης\\
\francais{Devenus jaloux de ce fait, Agamemnon, Ulysse et Diomède imaginèrent ce genre de procédé : ayant fait prisonnier un Phrygien qui apportait de l'or à Sarpédon, ils le forcèrent à écrire en caractères phrygiens une lettre parlant de trahison, comme si elle avait été adressée par Priam à Palamède, puis ils le tuèrent. Ils corrompirent alors un serviteur de Palamède, pour qu'il place en même temps que l'or phrygien la tablette écrite sous la tente de son maître. Eux-mêmes, ayant comparu [dans l'assemblée des chefs], accusèrent le héros de trahison et demandèrent qu'on fouille la tente. Après la découverte de la tablette et de l'or sous la tente, Palamède fut lapidé.}}
\end{multicols}
}%
\end{adjustbox}%
\vfill
\end{adjustwidth}
\newpage
% --- PAGE DE DROITE : Traduction française + Commentaire (2/8) ---
\begin{adjustwidth}{-2cm}{-2cm}
\noindent\textbf{\footnotesize Traduction française --- Folio 19 (\texttt{f. 171v}, vv. 765–799 [1/4]) \hfill ΕΠΕΙΣΟΔΙΟΝ Γ' / ΠΑΡΑΒΑΣΙΣ}\\[-2pt]
\noindent\rule{16.5cm}{0.5pt}\\[2pt]
{\fontsize{8.8pt}{10.2pt}\selectfont
\ligneEst{\textbf{Mnésiloque}}{\francais{Voyons, quel sera l'instrument de mon salut ?}}{765}
\ligneEst{}{\francais{Quelle tentative? Quelle combinaison ? Car l'auteur vrai,}}{}
\ligneEst{}{\francais{celui qui m'impliqua dans cette méchante affaire,}}{}
\ligneEst{}{\francais{n'apparaît pas encore. Allons, quel messager}}{}
\ligneEst{}{\francais{lui dépêcherai-je ? (Il réfléchit)... Il me vient à l'esprit un expédient}}{}
\ligneEst{}{\francais{tiré du \italique{Palamède}. Comme cet autre, je lancerai mes avirons}}{770}
\ligneEst{}{\francais{chargés d'écritures... Mais je ne les ai pas sous la main, ces avirons...}}{}
\ligneEst{}{\francais{Comment donc faire pour me procurer des avirons ? Comment ?...}}{}
}
\vspace{2pt}
\noindent\rule{16.5cm}{0.6pt}\\[2pt]
\noindent\textbf{\small Commentaire philologique, métrique et scénique (vv. 765–799 --- partie 2/8)}\\[2pt]
\begin{adjustbox}{max totalheight=15.2cm, max width=16.5cm, keepaspectratio}%
\parbox[t]{16.5cm}{
\setlength{\columnsep}{9pt}%
\begin{multicols}{2}
\fontsize{8.4pt}{9.7pt}\selectfont
\setlength{\parindent}{0pt}%
\setlength{\parskip}{1.2pt}%
\noindent Anxieux de faire connaître la mort injuste de son fils à Nauplios, Œax, le frère de Palamède, imaginait le stratagème ici repris par Mnésiloque\marqueur{2} :
\citComm{ὁ γὰρ Ἐὐριπίδης ἐν τῷ Παλαμήδει ἐποίησε τὸν Οἴακα τὸν ἀδελφὸν Παλαμήδους ἐπιγράψαι εἰς τὰς ναῦς τὸν θάνατον αὐτοῦ, ἵνα φερόμεναι αὗται ἔλθωσιν εἰς τὸν Ναύπλιον τὸν πατέρα αὐτοῦ καὶ ἀπαγγείλωσι τὸν θάνατον αὐτοῦ.\\
\francais{Euripide, dans le \italique{Palamède}, a représenté Oiax, frère de Palamède, gravant sur des navires [le récit de] sa mort, afin qu'ils soient portés jusque chez Nauplios et lui apprennent sa mort.}}
\noindent\textbf{770} \textbf{ὡς ἐκεῖνος,} Œax, frère de Palamède.
\noindent\textbf{772} \textbf{πόθεν οὖν γένοιντ᾽ ἄν μοι πλάται. πόθεν;} Le vers est incomplet. Scaliger proposa la solution la plus instinctive : redoubler le πόθεν en fin de vers : πόθεν ; <πόθεν;> Cette solution est vraisemblablement la bonne. \ \ \ 1° \ Une simple haplographie l'explique. \ \ \ 2° \ Elle est une nouvelle occurrence des railleries artiphanesques quant à la manie du redoublement chez les tragiques. Le double πόθεν se trouve d'ailleurs dans l'\italique{Ion} (vers 1453 sq.) :
\citComm{ἰὼ γύναι, πόθεν πόθεν ἔλαβες ἐμὸν\\
βρέφος ἐς ἀγκάλας;}
\noindent Même si ce qui faisait vraiment battre le cœur d'Euripide, c'était le redoublement du ποῦ (\italique{Électre}, v. 487 sq.) :
\citComm{ποῦ ποῦ νεᾶνις πότνι᾽ ἐμὴ δέσποινά τε,\\
Ἀγαμέμνονος παῖς, ὅν ποτ᾽ ἐξέθρεψ᾽ ἐγώ;}
\noindent Dans les \italique{Oiseaux} (vers 921), c'est le redoublement de πάλαι qui est moqué : l'ennuyeux poète qui débite
\citComm{πάλαι, πάλαι δὴ τήνδ᾽ ἐγὼ κλῄζω πόλιν}
\noindent ne pénétrera pas plus dans Νεφελοκοκκυγίᾳ qu'il n'aura de chance de demeurer dans la Cité de Platon.
\end{multicols}
}%
\end{adjustbox}%
\vfill
\end{adjustwidth}
\newpage

% --- FOLIO 19 (f. 171v) : SPREAD 2.2 (Texte grec établi + Traduction française + Commentaire 3/8 & 4/8) ---
\clearpage\ifodd\value{page}\hbox{}\newpage\fi
\enTeteG{ΘΕΣΜΟΦΟΡΙΑΖΟΥΣΑΙ --- Texte grec établi (f. 171v, vv. 765–799 [2/4])}
\enTeteD{LES THESMOPHORIES --- Traduction française (f. 171v, vv. 765–799 [2/4])}
\begin{adjustwidth}{-2cm}{-2cm}
\noindent\textbf{\footnotesize Texte grec établi --- Folio 19 (\texttt{f. 171v}, vv. 765–799 [2/4]) \hfill ΕΠΕΙΣΟΔΙΟΝ Γ' / ΠΑΡΑΒΑΣΙΣ}\\[-2pt]
\noindent\rule{16.5cm}{0.5pt}\\[2pt]
{\fontsize{8.8pt}{10.2pt}\selectfont
\ligneEst{}{τί δ᾽ ἂν εἰ ταδὶ τἀγάλματ᾽ ἀντὶ τῶν πλατῶν}{}
\ligneEst{}{γράφων διαρρίπτοιμι; βέλτιον πολύ.}{}
\ligneEst{}{ξύλον γέ τοι καὶ ταῦτα κἀκεῖν᾽ ἦν ξύλον.}{775}
\ligneEst{{\tiny\textit{anap.}}}{ὦ χεῖρες ἐμαὶ}{}
\ligneEst{{\tiny\textit{dim. anap.}}}{ἐγχειρεῖν χρῆν ἔργῳ πορίμῳ.}{}
\ligneEst{}{ἄγε δὴ πινάκων ξεστῶν δέλτοι,}{}
\ligneEst{{\tiny\textit{(catalect.)}}}{δέξασθε σμίλης ὁλκοὺς}{}
\ligneEst{}{κήρυκας ἐμῶν μόχθων· οἴμοι}{780}
\ligneEst{{\tiny\textit{(catalect.)}}}{τουτὶ τορῶ μοχθηρόν·}{}
}
\vspace{2pt}
\noindent\rule{16.5cm}{0.6pt}\\[2pt]
\noindent\textbf{\small Commentaire philologique, métrique et scénique (vv. 765–799 --- partie 3/8)}\\[2pt]
\begin{adjustbox}{max totalheight=15.2cm, max width=16.5cm, keepaspectratio}%
\parbox[t]{16.5cm}{
\setlength{\columnsep}{9pt}%
\begin{multicols}{2}
\fontsize{8.4pt}{9.7pt}\selectfont
\setlength{\parindent}{0pt}%
\setlength{\parskip}{1.2pt}%
\noindent\textbf{773-4} \textbf{τί δ᾽ ἂν εἶτα διὰ τἀγάλματ᾽ ἀντὶ τῶν πλατῶν / γράφων διαρρίπτοιμι;} Willems (\nieme{22} note, reproduite en annexe) a expliqué en quel sens il fallait prendre \grec{τὰ ἀγάλματα} : \grec{πᾶν ἐφ᾽ ᾧ τις ἀγάλλεται} (scholie à notre vers), \italique{tout ce qu'on offre et consacre aux dieux}. En l'occurrence, ce sont les \grec{πόπανα} que Mnésiloque a apportés (au vers 285) et non pas des \italique{ex-voto} comme on en voit encore apparaître dans les modernes traductions\marqueur{1}.\\ Y a-t-il ici une allusion au proverbe grec ἄλλο γένος κώπης, «~un autre genre de rame~»\marqueur{2} ?
\noindent\textbf{[Notes ad 770]} \textbf{(1)} On retrouve dans une scholie à son \italique{Oreste}, au vers 432, les grandes lignes de la version euripidéenne du drame. J'en reprends la traduction de Jouan \& van Looy, à retrouver dans leur édition des fragments d'Euripide aux Belles Lettres, tome VIII, \nieme{2} partie.
\noindent \textbf{(2)} Scholie au vers 771 des \italique{Thesmophories}.
\noindent \textbf{(1)} Après avoir encombré la scène d'une statue d'Apollon permanente d'après le vers 748, les commentateurs rajoutent ici des tablettes votives pour le décor.
\noindent \textbf{(2)} Rapporté par Plutarque, \italique{Proverbes alexandrins}, 2, 52. Ou aussi dans \italique{Zu den Paroemiographen} de Cohn, p.42. Érasme l'a incorporé à ses \italique{Adages} (n°1621).
\noindent\textbf{781} \textbf{τουτὶ τὸ ῥῶ μοχθηρόν·} « Voilà un `Rhô' défectueux ! »... telle est la leçon de notre unique manuscrit et telles aussi les traductions embarrassées. Car la question fuse : pourquoi Mnésiloque nous parle-t-il d'un `ρ' ? Le scholiaste toujours prompt à lancer une ânerie propose d'y voir la troisième lettre du destinataire, Εὐρῑπίδη... ce qui ne répond pas à la question : pourquoi la troisième lettre ?
\noindent Passons sur les autres explications du même acabit (si Aristophane avait un peu plus de glosateurs, tous les mots du dictionnaire comportant un rhô y seraient effectivement passés : Εὐριπίδη, χώρει, χαίρειν, \&c.) Lisez simplement :
\citComm{τουτὶ τορῶ μοχθηρόν.}
\end{multicols}
}%
\end{adjustbox}%
\vfill
\end{adjustwidth}
\newpage
% --- PAGE DE DROITE : Traduction française + Commentaire (4/8) ---
\begin{adjustwidth}{-2cm}{-2cm}
\noindent\textbf{\footnotesize Traduction française --- Folio 19 (\texttt{f. 171v}, vv. 765–799 [2/4]) \hfill ΕΠΕΙΣΟΔΙΟΝ Γ' / ΠΑΡΑΒΑΣΙΣ}\\[-2pt]
\noindent\rule{16.5cm}{0.5pt}\\[2pt]
{\fontsize{8.8pt}{10.2pt}\selectfont
\ligneEst{}{\francais{Mais quoi, si à défaut d'avirons, j'écrivais sur ces offrandes}}{}
\ligneEst{}{\francais{et les jetais de tous côtés ? Cela vaut bien mieux.}}{}
\ligneEst{}{\francais{Elles sont de bois après tout; ceux-là aussi étaient de bois.}}{775}
\ligneEst{}{\francais{À l'œuvre, ô mes mains,}}{}
\ligneEst{}{\francais{il s'agit d'un travail expédient.}}{}
\ligneEst{}{\francais{Allons, feuillets de mes tablettes polies,}}{}
\ligneEst{}{\francais{recevez les entailles du ciseau,}}{}
\ligneEst{}{\francais{messagères de mes misères. Aïe,}}{780}
\ligneEst{}{\francais{c'est pénible à percer !..}}{}
}
\vspace{2pt}
\noindent\rule{16.5cm}{0.6pt}\\[2pt]
\noindent\textbf{\small Commentaire philologique, métrique et scénique (vv. 765–799 --- partie 4/8)}\\[2pt]
\begin{adjustbox}{max totalheight=15.2cm, max width=16.5cm, keepaspectratio}%
\parbox[t]{16.5cm}{
\setlength{\columnsep}{9pt}%
\begin{multicols}{2}
\fontsize{8.4pt}{9.7pt}\selectfont
\setlength{\parindent}{0pt}%
\setlength{\parskip}{1.2pt}%
\noindent L'aoriste τορεῖν est homérique et fait partie d'un groupe dialectal, également homérique, d'aoristes radicaux thématiques où s'est dégagé un vocalisme en `o' : ἔμολον de βλώσκω, ἔθορον de θρώσκω et ἔτορον de τείρω. Ces homérismes se sont introduits chez les Tragiques, et peuvent aussi se rencontrer dans certains passages d'Aristophane. On l'avait déjà noté pour ἔμολον : on peut maintenant y ajouter ἔτορον\marqueur{1}. Hésychius (τ 1162) -- qui nous fut d'un grand secours comme à son habitude – l'avait déjà relevé :
\citComm{τορεῖν· τορῆσαι, τρῆσαι, τεμεῖν ;}
\noindent le verbe signifie \italique{percer}, \italique{perforer}, \italique{couper}.
\noindent Reconstituons la scène : Mnésiloque cherche un expédient pour envoyer son message. Heureusement, il connaît son euripide par cœur : il décide de procéder à la façon d'Œax dans le \italique{Palamède}. Manque de chance : il manque les rames (ἀλ ̓ οὐ πάρεισιν αἱ πλάται) pour graver dessus son message !.. qu'à cela ne tienne : les ustensiles de théâtres pour figurer les πόπανα feront bien l'affaire (ils sont de bois, après tout !) ; voilà donc nos morceaux de bois promus en tablettes (δέλτοι). Mnésiloque peut maintenant se mettre à graver la “cire” avec son stylet – un stylet ? quel stylet ? Nouveau détournement : il appelle `stylet' le coutelas avec lequel il vient d'égorger l'outre de vin. Seulement, à force de prendre des vessies pour des lanternes, Mnésiloque se fait rattraper par la dure sylvestre réalité : graver un message au couteau sur du bois se révèle bien plus difficile qu'avec un stylet sur une tablette de cire ; cela explique son cri de découragement : τουτὶ τορῶ μοχθηρόν, « ceci que je perce [il ne précise pas quoi, et pour cause : c'est à la fois un πόπανον, une πλάτη et un δέλτος]... est de mauvaise qualité ! »
\noindent\textbf{782} \textbf{χώρει, χώρει.} Mnésiloque est en train de graver son message : χώρει, «~accours !~», enjoint-il à Euripide ; χώ...ρει χώ...ρει, «~Ac...cours ... Ac...cours~», se dicte-t-il à lui-même. Il l'écrit sur deux rames différentes, comme le pluriel αἱ πλάται le laissait présager et comme c'était aussi le cas dans la tragédie originale ; le scholiaste de notre \italique{Ravennas} précise bien :
\citComm{ὁ γὰρ Οἴαξ ἐγχαράττει \textbf{πολλαῖς πλάταις} τὰ περὶ τὸν Παλαμήδην καὶ ἀφίησιν εἰς θάλασσαν, ὥστε μιᾷ γέ τινι τὸν Ναύπλιον περιπεσεῖν.\\
Œax grave effectivement de \italique{nombreuses rames} avec les circonstances touchant à $[$la mort de$]$ Palamède et les lance vers la mer de sorte que Nauplios tombe sur l'une d'elles.}
\end{multicols}
}%
\end{adjustbox}%
\vfill
\end{adjustwidth}
\newpage

% --- FOLIO 19 (f. 171v) : SPREAD 2.3 (Texte grec établi + Traduction française + Commentaire 5/8 & 6/8) ---
\clearpage\ifodd\value{page}\hbox{}\newpage\fi
\enTeteG{ΘΕΣΜΟΦΟΡΙΑΖΟΥΣΑΙ --- Texte grec établi (f. 171v, vv. 765–799 [3/4])}
\enTeteD{LES THESMOPHORIES --- Traduction française (f. 171v, vv. 765–799 [3/4])}
\begin{adjustwidth}{-2cm}{-2cm}
\noindent\textbf{\footnotesize Texte grec établi --- Folio 19 (\texttt{f. 171v}, vv. 765–799 [3/4]) \hfill ΕΠΕΙΣΟΔΙΟΝ Γ' / ΠΑΡΑΒΑΣΙΣ}\\[-2pt]
\noindent\rule{16.5cm}{0.5pt}\\[2pt]
{\fontsize{8.8pt}{10.2pt}\selectfont
\ligneEst{}{`χώρει', `χώρει'. ποίαν αὔλακα;}{}
\ligneEst{}{βάσκετ᾽ ἐπείγετε πάσας καθ᾽ ὁδοὺς}{}
\ligneEst{{\tiny\textit{(catalect.)}}}{κείνᾳ ταύτᾳ· ταχέως χρή.}{}
\ligneEst{\textbf{Χορός} {\tiny\textit{\decalageN{a}$4$ an\cat}}}{ἡμεῖς τοίνυν ἡμᾶς αὐτὰς εὖ λέξωμεν παραβᾶσαι,}{785}
\ligneEst{}{καίτοι πᾶς τις τὸ γυναικεῖον φῦλον κακὰ πόλλ᾽ ἀγορεύει,}{}
\ligneEst{}{ὡς πᾶν ἐσμὲν κακὸν ἀνθρώποις κἀξ ἡμῶν ἐστιν ἅπαντα,}{}
\ligneEst{}{\Automatique{ἔριδες νείκη στάσις ἀργαλέα λύπη πόλεμος. φέρε δή νυν,}}{}
\ligneEst{}{\Automatique{εἰ κακόν ἐσμεν, τί γαμεῖθ᾽ ἡμᾶς, εἴπερ ἀληθῶς κακόν ἐσμεν,}}{}
\ligneEst{}{\Automatique{κἀπαγορεύετε μήτ᾽ ἐξελθεῖν μήτ᾽ ἐκκύψασαν ἁλῶναι,}}{790}
}
\vspace{2pt}
\noindent\rule{16.5cm}{0.6pt}\\[2pt]
\noindent\textbf{\small Commentaire philologique, métrique et scénique (vv. 765–799 --- partie 5/8)}\\[2pt]
\begin{adjustbox}{max totalheight=15.2cm, max width=16.5cm, keepaspectratio}%
\parbox[t]{16.5cm}{
\setlength{\columnsep}{9pt}%
\begin{multicols}{2}
\fontsize{8.4pt}{9.7pt}\selectfont
\setlength{\parindent}{0pt}%
\setlength{\parskip}{1.2pt}%
\noindent À propos des ἀγάλματα, un scholiaste avait aussi écrit :
\citComm{πίνακες γὰρ ἦσαν ἐν τῷ ἱερῷ, καὶ ἀπὸ τούτων λαβὼν γράφει, καὶ ῥίπτει αὐτὰς λέγων· `ἀπέλθετε, σημάνατε Ἐὐριπίδῄ.\\
C'était des tablettes dans le temple. Les prenant, il écrivit puis les lança, en disant : «~Allez ! Informez Euripide !~»}
\noindent\textbf{782-4} \textbf{ποίαν αὔλακα / βάσκετ᾽ ἐπείγετε πάσας καθόδους / κείνᾳ ταύτᾳ· ταχέως χρή.} Les différents traducteurs voient dans la formule `ποίαν αὔλακα' une exclamation d'auto-contentement de Mnésiloque : «~Quelle belle lettre ! Quel sillon !~» Αὖλαξ est en effet attesté dans l'\italique{Anthologie palatine} (6, 68) au sens de sillon creusé par le stylet, c'est-à-dire le \italique{trait d'écriture}. Soit. Mais cette interprétation confère à ποίαν un sens qu'il n'a pas. Cet adjectif interrogatif sert à demander des précisions : \italique{de quelle nature ?} \italique{de quel genre ?} (Aristophane l'emploie fréquemment dans ses dialogues pour débrouiller -- ou souligner ? -- un quiproquo.) Il est utilisé en reprise d'un terme, reprise directe : Ἀγάθων ὁ κλεινὸς --- Ποῖος οὗτος Ἁγάθων; (vers 30) ; Τὸν δεῖνα γιγνώσκεις --- Τὸν δεῖνα; ποῖον; (vers 621) ;  Πρωτέως τάδ᾽ ἐστὶ μέλαθρα --- Ποίου Πρωτέως; (vers 874) ou implicite : Πρὸς ποῖον ἔργον; (vers 586) ; Χρείᾳ δὲ ποίᾳ τόνδ᾽ ἐπεισφέρεις λόγον; (vers 1163). Mnésiloque mentionne pour la première fois l'αὖλαξ ; mais reprise il y a bien\marqueur{1} : le mot devait sans aucun doute apparaître dans la tragédie d'Euripide.
\noindent Où en est Mnésiloque de l'action ? Il vient de graver en deux exemplaires l'appel au secours, χώρει, qu'il compte faire parvenir à Euripide... comment ?
\noindent Assurément, cette αὖλαξ désigne l'expédient qu'Œnax avait trouvé pour emporter ses messages.
\end{multicols}
}%
\end{adjustbox}%
\vfill
\end{adjustwidth}
\newpage
% --- PAGE DE DROITE : Traduction française + Commentaire (6/8) ---
\begin{adjustwidth}{-2cm}{-2cm}
\noindent\textbf{\footnotesize Traduction française --- Folio 19 (\texttt{f. 171v}, vv. 765–799 [3/4]) \hfill ΕΠΕΙΣΟΔΙΟΝ Γ' / ΠΑΡΑΒΑΣΙΣ}\\[-2pt]
\noindent\rule{16.5cm}{0.5pt}\\[2pt]
{\fontsize{8.8pt}{10.2pt}\selectfont
\ligneEst{}{\francais{«~Ac...cours~», «~Ac...cours !~» et quel sera notre cours d'eau ? (Il se tourne et lance les tablettes des deux côtés.)}}{}
\ligneEst{}{\francais{Partez, volez par toutes les voies}}{}
\ligneEst{}{\francais{avec cette nouvelle façon de faire : il faut aller vite.}}{}
\ligneEst{\textbf{Chœur}}{\francais{Maintenant nous allons faire nous-mêmes notre éloge dans la parabase.}}{785}
\ligneEst{}{\francais{Chacun pourtant médit à l'envi de l'espèce féminine :}}{}
\ligneEst{}{\francais{nous sommes un pur fléau pour les hommes ; de nous viennent tous les maux,}}{}
\ligneEst{}{\francais{disputes, querelles, affreuses séditions, chagrins, guerres. Çà, voyons,}}{}
\ligneEst{}{\francais{si nous sommes un fléau, pourquoi nous épousez-vous, ce fléau que nous sommes ?}}{}
\ligneEst{}{\francais{Pourquoi nous interdire de sortir, de nous laisser prendre à pousser la tête dehors ?}}{790}
}
\vspace{2pt}
\noindent\rule{16.5cm}{0.6pt}\\[2pt]
\noindent\textbf{\small Commentaire philologique, métrique et scénique (vv. 765–799 --- partie 6/8)}\\[2pt]
\begin{adjustbox}{max totalheight=15.2cm, max width=16.5cm, keepaspectratio}%
\parbox[t]{16.5cm}{
\setlength{\columnsep}{9pt}%
\begin{multicols}{2}
\fontsize{8.4pt}{9.7pt}\selectfont
\setlength{\parindent}{0pt}%
\setlength{\parskip}{1.2pt}%
\noindent On peut raisonnablement penser qu'il les déposait dans un \italique{cours d'eau} (qui envoyait ensuite les πλάται où l'une d'entre elles finissait par arriver à bon port, comme l'indique le scholiaste. Ça ne paraît pas très vraisemblable, mais la vraisemblance et une pièce d'Euripide font deux.) C'est bien en ce sens qu'il faut prendre αὖλαξ, le \italique{sillon}, et, par extension, la \italique{rigole} ou le \italique{ruisseau}\marqueur{2}.\\ Mnésiloque ne s'embarrasse pas vraiment du lièvre qu'il vient de soulever : la question n'est là que pour souligner les invraisemblances de la mise en scène, comme Aristophane aime bien les relever. Il décide d'aller au plus rapide : il jette avec force les deux πλάται. (Une de chaque côté, vers chaque εἴσοδος, comme le laisse soupçonner πάσας καθ᾽ ὁδοὺς.) Cette \italique{nouveauté} (κείνᾳ ταύτᾳ) par rapport au modèle original d'Euripide a l'avantage d'accélérer l'action ; c'est ainsi qu'il faut comprendre le vers 784 : κείνᾳ ταύτᾳ· ταχέως χρή.
\noindent\textbf{[Notes ad 781]} \textbf{(1)} Je dois cette remarque à Jean-Victor Vernhes qui, entre autres qualités, est un très fin connaisseur de la langue homérique. Que le lecteur n'hésite pas à consulter le Dictionnaire étymologique de la langue grecque de Pierre Chantraine à l'entrée τορεῖν pour plus de détails.
\noindent \textbf{(1)} Et si l'on en doutait, l'on pourrait toujours corriger le texte -- ce dont on se gardera bien ! Lire ποῖ ποτ' ἂν αὔλακα revient au même quant à l'action et dispense de supposer que l'αὖλαξ était le terme employé par Euripide. (ποτε \italique{gratia metri} mais dont la disparition relèverait d'une haplographie bien explicable. Au reste, la collocation `ποῖ ποτε' est très fréquente.)
\noindent \textbf{(2)} Voir le \LSJ \ et le \espagnol{Diccionario Griego-Español} IV à l'entrée αὖλαξ.
\noindent\textbf{785-845} \textbf{Parabase.} La parabase des \italique{Thesmophories}, à l'instar de \italique{Lysistrata} et des \italique{Grenouilles}, ne présente ni ᾠδή/ἀντῳδή ni ἀντεπίρρημα\marqueur{1}.
\noindent\textbf{788} \textbf{ἔριδες νείκη στάσις ἀργαλέα} Le vocabulaire est Homérique( \italique{Iliade}, XVII, v. 384-5) :
\citComm{τοῖς δὲ πανημερίοις ἔριδος μέγα νεῖκος ὀρώρει / ἀργαλέης·\\
\francais{Durant tout le jour se prolongea cette lutte sanglante et terrible.}}
\end{multicols}
}%
\end{adjustbox}%
\vfill
\end{adjustwidth}
\newpage

% --- FOLIO 19 (f. 171v) : SPREAD 2.4 (Texte grec établi + Traduction française + Commentaire 7/8 & 8/8) ---
\clearpage\ifodd\value{page}\hbox{}\newpage\fi
\enTeteG{ΘΕΣΜΟΦΟΡΙΑΖΟΥΣΑΙ --- Texte grec établi (f. 171v, vv. 765–799 [4/4])}
\enTeteD{LES THESMOPHORIES --- Traduction française (f. 171v, vv. 765–799 [4/4])}
\begin{adjustwidth}{-2cm}{-2cm}
\noindent\textbf{\footnotesize Texte grec établi --- Folio 19 (\texttt{f. 171v}, vv. 765–799 [4/4]) \hfill ΕΠΕΙΣΟΔΙΟΝ Γ' / ΠΑΡΑΒΑΣΙΣ}\\[-2pt]
\noindent\rule{16.5cm}{0.5pt}\\[2pt]
{\fontsize{8.8pt}{10.2pt}\selectfont
\ligneEst{}{\Automatique{ἀλλ᾽ οὑτωσὶ πολλῇ σπουδῇ τὸ κακὸν βούλεσθε φυλάττειν;}}{}
\ligneEst{}{\Automatique{κἂν ἐξέλθῃ τὸ γύναιόν ποι, κᾆθ᾽ εὕρητ᾽ αὐτὸ θύρασιν,}}{}
\ligneEst{}{\Automatique{μανίας μαίνεσθ᾽, οὓς χρῆν σπένδειν καὶ χαίρειν, εἴπερ ἀληθῶς}}{}
\ligneEst{}{\Automatique{ἔνδοθεν ηὕρετε φροῦδον τὸ κακὸν καὶ μὴ κατελαμβάνετ᾽ ἔνδον.}}{}
\ligneEst{}{\Automatique{κἂν καταδάρθωμεν ἐν ἀλλοτρίων παίζουσαι καὶ κοπιῶσαι,}}{795}
\ligneEst{}{\Automatique{πᾶς τις τὸ κακὸν τοῦτο ζητεῖ περὶ τὰς κλίνας περινοστῶν.}}{}
\ligneEst{}{\Automatique{κἂν ἐκ θυρίδος παρακύπτωμεν, τὸ κακὸν ζητεῖτε θεᾶσθαι·}}{}
\ligneEst{}{\Automatique{κἂν αἰσχυνθεῖσ᾽ ἀναχωρήσῃ, πολὺ μᾶλλον πᾶς ἐπιθυμεῖ}}{}
\ligneEst{}{\Automatique{αὖθις τὸ κακὸν παρακύψαν ἰδεῖν. οὕτως ἡμεῖς ἐπιδήλως}}{799}
}
\vspace{2pt}
\noindent\rule{16.5cm}{0.6pt}\\[2pt]
\noindent\textbf{\small Commentaire philologique, métrique et scénique (vv. 765–799 --- partie 7/8)}\\[2pt]
\begin{adjustbox}{max totalheight=15.2cm, max width=16.5cm, keepaspectratio}%
\parbox[t]{16.5cm}{
\setlength{\columnsep}{9pt}%
\begin{multicols}{2}
\fontsize{8.4pt}{9.7pt}\selectfont
\setlength{\parindent}{0pt}%
\setlength{\parskip}{1.2pt}%
\noindent\textbf{789} \textbf{εἰ καὶ κακόν ἐσμεν, τί γαμεῖθ᾽ ἡμᾶς, εἴπερ ἀληθῶς κακόν ἐσμεν} La perfidie d'Aristophane est à son comble quand il donne la parole aux femmes dans la parabase où elles assurent leur propre défense. « Si nous sommes un fléau, disent-elles aux hommes, pourquoi nous épousez-vous ? » Mais le spectateur, qui connaît, depuis l'école, son Hésiode par cœur et qui a déjà vu bien des comédies, répond intérieurement, comme mu par un ressort : « Parce qu'on n'échappe pas si facilement au piège tendu par Zeus, parce que la femme est un fléau nécessaire pour la perpétuation de l'espèce\marqueur{1} ». Le poète comique ne fait ici, sous des apparences trompeuses, qu'entretenir le vieux réflexe de misogynie. Autre argument employé par les femmes pour montrer leur valeur : celui de la jalousie qu'elles excitent\marqueur{2}. Tout commentaire est ici inutile... (Note d'Alain Blanchard\marqueur{3}.)
\noindent\textbf{[Notes ad 785-845]} \textbf{(1)} Mon but n'est pas de proposer une description exhaustive de la structure des comédies d'Aristophane, ni même des \italique{Thesmophories} en particulier. Ce travail a déjà été très bien fait par White, \italique{The Verse of Greek Comedy} (1912). On lira aussi avec un grand intérêt les travaux de Jean Irigoin regroupés dans son anthologie posthume \italique{Le poète grec au travail} (2009) et notamment son article «~Prologue et \italique{parodos} dans la comédie d'Aristophane (\italique{Acharniens}, \italique{Paix}, \italique{Thesmophories})~» :  «~[Aristophane], ce poète comique à l'imagination débridée, s'astreignait, dans certaines parties de ses pièces, à respecter des règles concernant le choix des vers, leur disposition, et même leur nombre~». Le lecteur intéressé se reportera à ces travaux -- et à nombre d'autres -- auxquels je ne saurais me substituer.
\noindent \textbf{(1)} Hésiode, les \italique{Travaux et les jours}, v. 83-84 ; Susarion, fr. 1, 3-4 Kassel-Austin. Voir aussi ce que dit Renzo Tosi (§284 et §285 de son monumental \italique{Dictionnaire des sentences latines et grecques}, 2010) : Κακὸν ἀναγκαῖον γυνή ou \latin{Femina dulce malum} : la femme est un καλὸν κακόν (\italique{Théogonie}, v. 587). Quant au vers précédent, il y une allusion à peine voilée aux pandoriennes calamités déclenchées par la mythique -- et belle, surtout, pour nos malheurs ! -- Hélène. (Voir à ce propos Tosi, §350 : \latin{Nam fuit ante Helenam cunnus taeterrima belli causa}.) Il est donc tout naturel que Mnésiloqie réponde à cette parabase par une reprise de l'\italique{Hélène}.
\end{multicols}
}%
\end{adjustbox}%
\vfill
\end{adjustwidth}
\newpage
% --- PAGE DE DROITE : Traduction française + Commentaire (8/8) ---
\begin{adjustwidth}{-2cm}{-2cm}
\noindent\textbf{\footnotesize Traduction française --- Folio 19 (\texttt{f. 171v}, vv. 765–799 [4/4]) \hfill ΕΠΕΙΣΟΔΙΟΝ Γ' / ΠΑΡΑΒΑΣΙΣ}\\[-2pt]
\noindent\rule{16.5cm}{0.5pt}\\[2pt]
{\fontsize{8.8pt}{10.2pt}\selectfont
\ligneEst{}{\francais{Pourquoi vous piquer de le garder avec un soin si jaloux, ce fléau ?}}{}
\ligneEst{}{\francais{Arrive-t-il à votre petite femme d'être en course, et que vous la trouviez absente,}}{}
\ligneEst{}{\francais{vous entrez dans des fureurs, quand il vous faudrait faire libation et vous réjouir}}{}
\ligneEst{}{\francais{de ce qu'il a disparu de chez vous, le fléau, et que vous ne le trouviez pas au logis.}}{}
\ligneEst{}{\francais{Nous endormons-nous chez les autres, fatiguées de danser,}}{795}
\ligneEst{}{\francais{chacun se met après le fléau et rode autour des lits.}}{}
\ligneEst{}{\francais{Nous penchons-nous par une fenêtre, vous cherchez à le dévisager, le fléau,}}{}
\ligneEst{}{\francais{et si par pudeur elle se retire, chacun n'en est que plus impatient}}{}
\ligneEst{}{\francais{de la voir se pencher de nouveau, ce fléau. Tant il est vrai que nous valons bien mieux}}{799}
}
\vspace{2pt}
\noindent\rule{16.5cm}{0.6pt}\\[2pt]
\noindent\textbf{\small Commentaire philologique, métrique et scénique (vv. 765–799 --- partie 8/8)}\\[2pt]
\begin{adjustbox}{max totalheight=15.2cm, max width=16.5cm, keepaspectratio}%
\parbox[t]{16.5cm}{
\setlength{\columnsep}{9pt}%
\begin{multicols}{2}
\fontsize{8.4pt}{9.7pt}\selectfont
\setlength{\parindent}{0pt}%
\setlength{\parskip}{1.2pt}%
\noindent \textbf{(2)} Aux vers 792-794.
\noindent \textbf{(3)} Tirée de son remarquable article «~L'art de la mauvaise foi dans les \italique{Thesmophories} d'Aristophane~», in \italique{Dieux, héros et médecins grecs. Hommage à Fernand Robert.} (2001), p. 97-106. Tout amateur d'Aristophane ne peut pas ne pas aller prestement lire ces pages si riches (j'émets seulement quelques réserves sur le tout dernier paragraphe).
\noindent Blanchard y montre que « la mauvaise foi d'Aristophane s'est exercée, en ce qui concerne Euripide, sur un thème moins traditionnel que celui de la misogynie et qu'elle permet de mieux comprendre le sujet même de sa comédie des \italique{Thesmophories}. » De même que la défaite athénienne à Délion avait profité à la popularité de Socrate, qui s'était conduit héroïquement dans la débâcle comme on le sait par Platon, et que c'est ce bonheur d'une réputation rehaussée dans le malheur athénien qui avait provoqué Aristophane à l'écriture vengeresse des \italique{Nuées} ; de même, les \italique{Thesmophories} furent écrites peu après la désastreuse expédition de Sicile : alors qu'Athènes vivait ses heures les plus sombres, Euripide en profita pour redorer son blason : on fit appel à lui pour écrire un chant funèbre en l'honneur des Athéniens morts à Syracuse (Plutarque nous en a conservé deux vers).
\noindent Anecdote plus significative encore : certains Athéniens, réduits en esclavage, ne durent leur salut qu'à Euripide, en récitant de ses prologues, si prisés des Syracusains. Que la poésie d'Euripide pût sauver qui que ce soit ?.. alors qu'elle est justement l'une des causes de la dégénérescence athénienne !? voilà une nouvelle fumisterie qu'Aristophane mettra à mal tout au long des \italique{Thesmophories} : les différents détournements des tragédies d'Euripide se heurteront toutes au mur de l'incompréhension des geôliers de Mnésiloque. Une fois bien cerné le contexte historique, l'on comprend mieux Aristophane : ce « jeu de massacre » n'est motivé que par la passion -- brûlante chez Aristophane -- de la justice : il cogne, cogne, et cogne encore contre tous ces charognards des misères athéniennes. \anglais{Now is a time to storm !} Pour conclure : si Euripide est misogyne, c'est là son moindre défaut.
\noindent\textbf{797} \textbf{τὸ κακὸν ζητεῖτε θεᾶσθαι} Suivant la belle correction de Kaye : τὸ κακὸν ζητεῖτε θεᾶσθαι, v. 797, au lieu de : ζητεῖ τὸ κακὸν τε θεᾶσθαι. (Note de Willems.)
\end{multicols}
}%
\end{adjustbox}%
\vfill
\end{adjustwidth}
\newpage
\clearpage\ifodd\value{page}\hbox{}\newpage\fi
\enTeteG{ΠΑΡΑΒΑΣΙΣ --- Codex Ravennas 429, f. 172r (Classense 174r) --- vv. 800–837}
\begin{adjustwidth}{-1.5cm}{-3.0cm}
\begin{adjustbox}{max totalheight=21.0cm, max width=16.5cm, keepaspectratio}
\begin{minipage}{16.5cm}
{\fontsize{7.2pt}{8.5pt}\selectfont
\noindent\textbf{Folio 20 --- Codex Ravennas 429, \texttt{f. 172r} (\texttt{Classense 174r}) : ΠΑΡΑΒΑΣΙΣ, vv. 800–837}\hfill \textit{Témoins :} $\text{R}^1, \Rpc, \Rdeux$, \PsiUn\par\vspace{2pt}
\hrule height 0.4pt\vspace{2.5pt}
}
{\fontsize{8.1pt}{9.5pt}\selectfont
\ligneRav{800}{χορ/}{ὑμῶν ἐσμεν πολὺ βελτίους, βάσανός τε πάρεστιν ἰδέσθαι·}{}
\ligneRav{}{}{βάσανον δῶμεν πότεροι χείρους. ἡμεῖς μὲν γάρ φαμεν ὑμᾶς,}{}
\ligneRav{802}{}{\difference{ὑμεῖς δ᾽ ἡμᾶς}. \difference{σκεψώμεθα} δὴ κἀντιτιθῶμεν πρὸς ἕκαστον,}{\textbf{802} ὑμεῖς δ᾽ ὑμᾶς R : \lecon{ὑμεῖς δ᾽ ἡμᾶς} \Mon \espApp σκεψώμεσθα R : \lecon{σκεψώμεθα} \Mon}
\ligneRav{803}{}{παραβάλλουσαι τῆς τε γυναικὸς καὶ τἀνδρὸς τοὔνομ᾽ \difference{ἑκάστου}·}{\textbf{803} ἕκαστον R : \lecon{ἑκάστου} Grynaeus}
\ligneRav{804}{}{Ναυσιμάχης \difference{μέν γ᾽} ἥττων \difference{ἐστὶν} Χαρμῖνος· \difference{δῆλα δὲ} τἄργα·}{\textbf{804} μέν R : \lecon{μέν γ᾽} Dobree \espApp ἐστὶ R : \lecon{ἐστὶν} Küster \espApp δηλαδῆ R : \lecon{δῆλα δὲ} Grynaeus}
\ligneRav{805}{}{καὶ μὲν δὴ καὶ Κλεοφῶν \difference{χείρων} πάντως δήπου Σαλαβακχοῦς·}{\textbf{805} χεῖρον R : \lecon{χείρων} Zanetti}
\ligneRav{}{}{πρὸς Ἀριστομάχην δὲ χρόνου πολλοῦ, πρὸς ἐκείνην τὴν Μαραθῶνι,}{}
\ligneRav{807}{}{καὶ \difference{Στρατονίκην} ὑμῶν οὐδεὶς οὐδ᾽ ἐγχειρεῖ πολεμίζειν·}{\textbf{807} \lecon{Στρατονίκην} \PsiUn : στρατωνίκην R}
\ligneRav{808}{}{ἀλλ᾽ Εὐβούλης τῶν πέρυσίν \difference{τις βουλευτής} ἐστιν ἀμείνων}{\textbf{808} τίς βουλευτής R : \lecon{τις βουλευτής} Küster}
\ligneRav{809}{}{παραδοὺς ἑτέρῳ τὴν βουλείαν; οὐδ᾽ αὐτὸς τοῦτό γε \difference{φήσει}·}{\textbf{809} \lecon{αὐτὸς} R : Ἄνυτος Maas \espApp φήσεις R : \lecon{φήσει} Küster}
\ligneRav{810}{}{οὕτως ἡμεῖς πολὺ βελτίους τῶν ἀνδρῶν \difference{εὐχόμεθ᾽} εἶναι·}{\textbf{810} εὐχόμεσθ᾽ R : \lecon{εὐχόμεθ᾽} \Mon}
\ligneRav{811}{}{οὐδ᾽ ἂν κλέψασα γυνὴ ζεύγει κατὰ πεντήκοντα τάλαντα}{\textbf{811} ζεύγει post γυνὴ R : post οὐδ᾽ ἂν van Herwerden}
\ligneRav{}{}{εἰς πόλιν ἔλθοι τῶν δημοσίων· ἀλλ᾽ ἢν τὰ μέγισθ᾽ ὑφέληται}{}
\ligneRav{813}{}{φορμὸν πυρῶν τἀνδρὸς κλέψασ᾽, αὐθημερὸν αὔτ᾽ ἀπέδωκεν}{\textbf{813} \lecon{αὔτ᾽ ἀπέδωκεν} R, M\expo{pc} : ἀνταπέδωκεν M\expo{ac}, Bentley}
\ligneRav{}{—}{\hspace{1.1em}ἀλλ᾽ ἡμεῖς ἂν πολλοὺς τούτων}{}
\ligneRav{815}{}{\hspace{1.1em}\difference{ἀποδείξαιμεν} ταῦτα \difference{ποιοῦντας}·}{\textbf{815} ἀποδείξομεν R : \lecon{ἀποδείξαιμεν} Bentley \espApp ποθοῦντας R : \lecon{ποιοῦντας} Zanetti (\italique{facientes} Divus)}
\ligneRav{}{}{\hspace{1.1em}καὶ πρὸς τούτοις γάστριδας ἡμῶν}{}
\ligneRav{}{}{\hspace{1.1em}ὄντας μᾶλλον καὶ λωποδύτας}{}
\ligneRav{}{}{\hspace{1.1em}καὶ βωμολόχους κἀνδραποδιστάς·}{}
\ligneRav{819}{}{\hspace{1.1em}καὶ μὲν δήπου καὶ τὰ πατρῷά γε}{\textbf{819} δήπου R : δή που Denniston}
\ligneRav{820}{}{\hspace{1.1em}χείρους ἡμῶν εἰσιν σῴζειν·}{}
\ligneRav{}{}{\hspace{1.1em}ἡμῖν μὲν γὰρ σῶν ἔτι καὶ νῦν}{}
\ligneRav{822}{}{\hspace{1.1em}\difference{ταὐτὸ κανοῦν καὶ τὸ σκιάδειον}·}{\textbf{822-3} τἀντίον ὁ κανὼν οἱ καλαθίσκοι / τὸ σκιάδειον· R : \lecon{ταὐτὸ κανοῦν καὶ τὸ σκιάδειον·} Bibi}
\ligneRav{824}{}{\hspace{1.1em}τοῖς δ᾽ ἡμετέροις \difference{ἀνδράσι} τούτοις}{\textbf{824} ἀνδράσιν R : \lecon{ἀνδράσι} \Mon}
\ligneRav{825}{}{\hspace{1.1em}\difference{ἀπόλωλεν} μὲν πολλοῖς ὁ κανὼν}{\textbf{825} ἀπόλωλε R : \lecon{ἀπόλωλεν} Daubuz, Bentley}
\ligneRav{}{}{\hspace{1.1em}ἐκ τῶν οἴκων αὐτῇ λόγχῃ,}{}
\ligneRav{}{}{\hspace{1.1em}πολλοῖς δ᾽ ἑτέροις ἀπὸ τῶν ὤμων}{}
\ligneRav{}{}{\hspace{1.1em}\{ ἐν ταῖς στρατιαῖς \}}{}
\ligneRav{}{}{\hspace{1.1em}ἔρριπται τὸ σκιάδειον·}{}
\ligneRav{830}{χορ/}{πόλλ᾽ ἂν αἱ γυναῖκες ἡμεῖς ἐν δίκῃ μεμψαίμεθ᾽ ἂν}{}
\ligneRav{}{}{τοῖσιν ἀνδράσιν δικαίως, ἓν δ᾽ ὑπερφυέστατον·}{}
\ligneRav{832}{}{\difference{χρῆν} γάρ, ἡμῶν εἰ τέκοι τις ἄνδρα χρηστὸν τῇ πόλει,}{\textbf{832} χρὴ R : \lecon{χρῆν} Brunck}
\ligneRav{}{}{ταξίαρχον ἢ στρατηγόν, λαμβάνειν τιμήν τινα,}{}
\ligneRav{834}{}{προεδρίαν τ᾽ αὐτῇ δίδοσθαι \difference{Στηνίοισι} καὶ Σκίροις}{\textbf{834} Τηνίοισι R : \lecon{Στηνίοισι} Σ, Suid. σ 1089}
\ligneRav{835}{}{ἔν τε ταῖς ἄλλαις ἑορταῖς αἷσιν ἡμεῖς ἤγομεν·}{}
\ligneRav{836}{}{\difference{εἰ δὲ} δειλὸν καὶ πονηρὸν ἄνδρα τις τέκοι γυνή,}{\textbf{836} ἤν δὲ R : \lecon{εἰ δὲ} Brunck}
\ligneRav{}{}{ἢ τριήραρχον πονηρὸν ἢ κυβερνήτην κακόν,}{}
}
\vspace{3pt}\noindent\rule{4.5cm}{0.4pt}\par\vspace{1.5pt}
{\fontsize{7.0pt}{8.3pt}\selectfont
\noindent\textbf{Loci paralleli (vv. 800–837) ---} \textbf{800, 810} (βελτίους $=$ βελτίονες) Moeris β 8 \espTest \textbf{805} (Σαλαβακχοῦς) Hésychius (σ 88) Σαλαβακχώ· πόρνης ὄνομα. ἀπὸ δὲ ταύτης καὶ τὰς καταφερεῖς εἰς τὰ ἀφροδίσια οὕτως ἔλεγον Ἀττικοί \espTest \textbf{808-9} (τις βουλευτής -- βουλείαν) Suidas β 431 (ἐστι pro ἐστιν et ἀμείνων omisso) \espTest \textbf{809} (παραδοὺς ἑτέρῳ τὴν βουλείαν) \Antiatt 84.30 \espTest (βουλείαν) \Harp B 19 ; β Photius β 233 βουλία ; Hesychius β 921 βουλεῖα \espTest \textbf{816} (γάστριδας) Photius γ 40 γάστρις $=$ Suidas γ 72 ; cf. \AelDion γ 3 \espTest \textbf{822} %(τἀντίον ὁ κανὼν) Pollux VII 35 κανὼν ἱστοῦ τὸ καλούμενον ἀντίον \espTest \textbf{834} (Στηνίοισι) Suidas σ 1089 Στήνια καὶ Σκίρα ($=$ Scholium) ; Hesychius σ 1825 στήνια\par
}
\vspace{2.5pt}\noindent\rule{4.5cm}{0.4pt}\par\vspace{1.5pt}
{\fontsize{6.9pt}{8.2pt}\selectfont
\noindent\textbf{Suite (en commentaire) de l'apparat critique \& notes paléographiques (\texttt{f. 172r}) ---} \textbf{Témoin papyrologique :} \textbf{\PsiUn{} (\italique{PSI} XI 1194)} couvre les \textbf{vv. 804–809} et confirme Στρατονίκην au v. 807. --- \textbf{v. 815 (ποθοῦντας R pour πο(ι)οῦντας) :} L. Hauseux consacre un important excursus paléographique à cette faute d'onciale (ΠΟΟΥΝΤΑΣ $\to$ ΠΟΘΟΥΝΤΑΣ, confusion des lettres rondes Ο/Θ sur la forme attique sans iota ποοῦντας). --- \textbf{vv. 822–823 (restitution de \textbf{Bibi}) :} τἀντίον ὁ κανὼν οἱ καλαθίσκοι / τὸ σκιάδειον· R $\to$ \lecon{ταὐτὸ κανοῦν καὶ τὸ σκιάδειον·} \lecon{Bibi} (οἱ καλαθίσκοι est une glose marginale intruse expliquant κανοῦν, corrompu en κανὼν). --- \textbf{v. 834 :} Τηνίοισι R (chute du $\Sigma$ initial par haplographie après δίδοσθαι) $\to$ \lecon{Στηνίοισι} $\Sigma^{\text{R}}$, Souda $\sigma\ 1089$ (fête des \italique{Stènia}). --- \textbf{vv. 822–823 :} τἀντίον ὁ κανὼν οἱ καλαθίσκοι / τὸ σκιάδειον· R : \lecon{ταὐτὸ κανοῦν καὶ τὸ σκιάδειον·} \lecon{Bibi}. --- \textbf{v. 834 :} Τηνίοισι R : \lecon{Στηνίοισι} $\Sigma$, Souda $\sigma\ 1089$.\par
}
\end{minipage}
\end{adjustbox}
\end{adjustwidth}
\newpage
\enTeteD{Codex Ravennas 429 (Biblioteca Classense) --- Fac-similé du f. 172r (vv. 800–837)}
\begin{adjustwidth}{-1.5cm}{-3.0cm}
\begin{center}
\includegraphics[width=16.5cm,height=20.5cm,keepaspectratio]{/workspaces/Thesmophories/dossier_ravennas_thesmophories/folios_individuels_ff162v-177r/20_f172r_classense174r_vv0800-0837.jpg}\par
\vspace{2pt}
{\scriptsize\textbf{Codex Ravennas 429, f. 172r (Classense 174r)} --- ΠΑΡΑΒΑΣΙΣ, \textbf{vv. 800–837} ($\text{R}^1, \Rpc, \Rdeux$, \PsiUn).}
\end{center}
\end{adjustwidth}
\newpage
\clearpage\ifodd\value{page}\hbox{}\newpage\fi% --- FOLIO 20 (f. 172r) : SPREAD 2.1 (Texte grec établi + Traduction française + Commentaire 1/8 & 2/8) ---
\clearpage\ifodd\value{page}\hbox{}\newpage\fi
\enTeteG{ΘΕΣΜΟΦΟΡΙΑΖΟΥΣΑΙ --- Texte grec établi (f. 172r, vv. 800–837 [1/4])}
\enTeteD{LES THESMOPHORIES --- Traduction française (f. 172r, vv. 800–837 [1/4])}
\begin{adjustwidth}{-2cm}{-2cm}
\noindent\textbf{\footnotesize Texte grec établi --- Folio 20 (\texttt{f. 172r}, vv. 800–837 [1/4]) \hfill ΠΑΡΑΒΑΣΙΣ}\\[-2pt]
\noindent\rule{16.5cm}{0.5pt}\\[2pt]
{\fontsize{8.8pt}{10.2pt}\selectfont
\ligneEst{}{\Automatique{ὑμῶν ἐσμεν πολὺ βελτίους, βάσανός τε πάρεστιν ἰδέσθαι.}}{800}
\ligneEst{}{\Automatique{βάσανον δῶμεν πότεροι χείρους. ἡμεῖς μὲν γάρ φαμεν ὑμᾶς,}}{}
\ligneEst{}{\Automatique{ὑμεῖς δ᾽ ἡμᾶς. σκεψώμεθα δὴ κἀντιτιθῶμεν πρὸς ἕκαστον,}}{}
\ligneEst{}{\Automatique{παραβάλλουσαι τῆς τε γυναικὸς καὶ τἀνδρὸς τοὔνομ᾽ ἑκάστου.}}{}
\ligneEst{}{\Automatique{Ναυσιμάχης μέν γ᾽ ἥττων ἐστὶν Χαρμῖνος· δῆλα δὲ τἄργα.}}{}
\ligneEst{}{\Automatique{καὶ μὲν δὴ καὶ Κλεοφῶν χείρων πάντως δήπου Σαλαβακχοῦς.}}{805}
\ligneEst{}{\Automatique{πρὸς Ἀριστομάχην δὲ χρόνου πολλοῦ, πρὸς ἐκείνην τὴν Μαραθῶνι,}}{}
\ligneEst{}{\Automatique{καὶ Στρατονίκην ὑμῶν οὐδεὶς οὐδ᾽ ἐγχειρεῖ πολεμίζειν.}}{}
\ligneEst{}{\Automatique{ἀλλ᾽ Εὐβούλης τῶν πέρυσίν τις βουλευτής ἐστιν ἀμείνων}}{}
}
\vspace{2pt}
\noindent\rule{16.5cm}{0.6pt}\\[2pt]
\noindent\textbf{\small Commentaire philologique, métrique et scénique (vv. 800–837 --- partie 1/8)}\\[2pt]
\begin{adjustbox}{max totalheight=15.2cm, max width=16.5cm, keepaspectratio}%
\parbox[t]{16.5cm}{
\setlength{\columnsep}{9pt}%
\begin{multicols}{2}
\fontsize{7.6pt}{8.8pt}\selectfont
\setlength{\parindent}{0pt}%
\setlength{\parskip}{1.2pt}%
\noindent\textbf{801-2} \textbf{ἡμεῖς μὲν γάρ φαμεν ὑμᾶς, ὑμεῖς δ᾽ ἡμᾶς} Molière, \italique{Impromptu de Versailles}.
\citComm{\francais{Il gage que c'est moi, et moi je gage que c'est lui.}}
\noindent\textbf{804} \textbf{Ναυσιμάχης μέν ἥττων ἐστὶν Χαρμῖνος·} Charminos, stratège athénien, qui venait de se faire battre sur mer près de Symé\marqueur{1}, est rapproché de Nausimaché (\italique{bataille navale}). En quoi Charminos est-il inférieur (ἥττων) à Nausimaché ? Certes par le nom -- comme Aristophane n'avait pas manqué de prévenir : après les déboires qu'il fit subir à la marine athénienne, il serait bien indigne d'en porter un de ce genre !.. mais en rester à cet unique rapprochement serait bien léger.
\noindent La précision de la scholie -- Ναυσιμάχη πόρνη καὶ Σαλαβακχώ -- n'est pour une fois pas une fantaisie de Byzantin ; Nausimaché ne surpasse pas Charminos dans le \italique{combat naval} que par le nom : de par sa profession, elle devait aussi y exceller \italique{en pratique} (τἄργα) -- métaphoriquement, s'entend\marqueur{2} --, ce qui lui valut à la fois son renom et son surnom.
\noindent\textbf{805} \textbf{Κλεοφῶν χείρων πάντως δήπου Σαλαβακχοῦς.} Salabachô est un nom générique de prostituée\marqueur{1}. Aristophane le mentionne aussi au vers 765 des \italique{Cavaliers} -- c'est le Paphlagonien qui parle :
\citComm{Κλέων --- τῇ μὲν δεσποίνῃ Ἀθηναίᾳ τῇ τῆς πόλεως μεδεούσῃ\\
εὔχομαι, εἰ μὲν περὶ τὸν δῆμον τὸν Ἀθηναίων γεγένημαι\\
βέλτιστος ἀνὴρ μετὰ Λυσικλέα καὶ Κύνναν καὶ Σαλαβακχώ,\\
ὥσπερ νυνὶ μηδὲν δράσας δειπνεῖν ἐν πρυτανείῳ·\\
Cléon --- \francais{À notre maîtresse Athéna, qui régit cette cité,\\
je fais cette prière ; si pour le peuple des Athéniens je me suis montré\\
le plus méritant des hommes après Lysiclès, Cynna et Salabacchô,\\
puissé-je comme maintenant, sans avoir rien fait, dîner dans le Prytanée.}}
\noindent L'exégèse de ce passage nous permettra de comprendre du même coup celui qui nous intéresse.
\noindent Willems avait justement remarqué que les trois noms étaient mis ici par jeu de consonance avec les trois représentants de l'Athènes impérialiste : Périclès, Cimon et Thémistocle\marqueur{2}.
\noindent Pour le sens, voici ce qu'il en est : la modestie -- toujours fausse -- de Cléon, qui se place après ces trois-là, n'est que pure affectation : il fait comprendre qu'il cumule en fait les mérites des trois.
\noindent Λυσι-κλῆς, il l'est depuis sa victoire de Sphactérie : il acquit la \italique{célébrité} en \italique{déliant} la situation empêtrée du siège de Sphactérie.
\noindent Κύννα, il prétend l'être, lui qui ne cesse de proclamer être le \italique{chien} de garde de la cité\marqueur{3}.
\noindent Σαλα-βακχώ se laisse aisément comprendre ; on rapprochera la première partie du nom d'une glose de Photius (σ 43) :
\citComm{σάλα· φροντίς. οὕτως Αἰσχύλος.}
\noindent On reprochait au braillard Cléon d'être un excité ; il a beau jeu de prétendre l'être en vertu de sa \italique{préoccupation}, du \italique{soin} qu'il porte à la cité. La seconde partie du nom, βακχώ, étant une allusion à l'agitation de tous les fidèles de Bacchus, et notamment les bacchantes. Suétone, dans son Περὶ βλασφημιῶν, \italique{Des termes injurieux}, cite βακχός au chapitre εἰς ἀλαζόνας, «~envers les bavards, les prétentieux~»\marqueur{4}.
\end{multicols}
}%
\end{adjustbox}%
\vfill
\end{adjustwidth}
\newpage
% --- PAGE DE DROITE : Traduction française + Commentaire (2/8) ---
\begin{adjustwidth}{-2cm}{-2cm}
\noindent\textbf{\footnotesize Traduction française --- Folio 20 (\texttt{f. 172r}, vv. 800–837 [1/4]) \hfill ΠΑΡΑΒΑΣΙΣ}\\[-2pt]
\noindent\rule{16.5cm}{0.5pt}\\[2pt]
{\fontsize{8.8pt}{10.2pt}\selectfont
\ligneEst{}{\francais{que vous. Si l'on veut une preuve, il n'y a qu'à regarder.}}{800}
\ligneEst{}{\francais{Venons-en à la preuve et sachons qui sont les pires. Car nous disons que c'est vous,}}{}
\ligneEst{}{\francais{et vous dites que c'est nous. Vérifions donc et mettons-les en présence isolément}}{}
\ligneEst{}{\francais{en rapprochant chaque fois le nom de la femme de celui de l'homme.}}{}
\ligneEst{}{\francais{Charminos est inférieur à Nausimaché : les faits parlent d'eux-mêmes,}}{}
\ligneEst{}{\francais{et à coup sûr aussi Cléophon est pire en somme que Salabaccho.}}{805}
\ligneEst{}{\francais{Avec Aristomaché, l'héroïne de Marathon,}}{}
\ligneEst{}{\francais{et avec Stratonice, il y a beau jour que nul d'entre vous ne tenterait de se mesurer.}}{}
\ligneEst{}{\francais{Et Euboulé, parmi les bouleutes de l'an dernier, en est-il un qui,}}{}
}
\vspace{2pt}
\noindent\rule{16.5cm}{0.6pt}\\[2pt]
\noindent\textbf{\small Commentaire philologique, métrique et scénique (vv. 800–837 --- partie 2/8)}\\[2pt]
\begin{adjustbox}{max totalheight=15.2cm, max width=16.5cm, keepaspectratio}%
\parbox[t]{16.5cm}{
\setlength{\columnsep}{9pt}%
\begin{multicols}{2}
\fontsize{7.0pt}{8.0pt}\selectfont
\setlength{\parindent}{0pt}%
\setlength{\parskip}{1.2pt}%
\noindent \textbf{Κλεοφῶν} était un démagogue dont la trop longue `carrière' politique fut brutalement interrompue en 404 av. J.-C. par une condamnation à mort. Dans les \italique{Grenouilles}, Aristophane charge ce va-t-en-guerre\marqueur{5} dont la langue par trop pendue et la lippe par trop pendante n'émettaient qu'un babil barbarisant criant ses barbaresques origines\marqueur{6}. (Pour un dossier complet sur Cléophon, le lecteur pourra consulter le commentaire de Fritzsche \italique{ad loc.})
\noindent\textbf{806-8} \textbf{Ἀριστομάχη, Στρατονίκη et Εὐβούλη} sont des noms attestés au moment de la représentation des \italique{Thesmophories}\marqueur{1}. Ces noms sont parlants, bien-sûr, mais l'allusion ; la raison pour laquelle Aristophane les cite, m'échappe. Les deux premières étaient-elles de puissantes princesses contre lesquelles les Athéniens n'auraient pas osé se frotter\marqueur{2} ?
\noindent\textbf{[Notes ad 804]} \textbf{(1)} Cf. Thucydide, VIII, 41.
\noindent \textbf{(2)} La langue grecque, et celle d'Aristophane en particulier, fourmille de métaphores nautiques pour rendre de façon imagée les corps-à-corps charnels ; ces images remplissent les sections 179 à 184 des \italique{Images d'Aristophane} de Taillardat et les pages 161-166 de la \italique{Maculate Muse} d'Henderson. Pour faire écho à Nausimaché, ne citons que le verbe ναυμαχεῖν usité par Aristophane aux vers 430 des \italique{Grenouilles} et 675 de \italique{Lysistrata}. Quant à τἄργα, on a déjà souligné l'acception scabreuse que peut avoir ce mot (cf. le vers 204 : ἔργα νυκτερείσια).
\noindent \textbf{(1)} Hésychius (σ 88) : Σαλαβακχώ· πόρνης ὄνομα. ἀπὸ δὲ ταύτης καὶ τὰς καταφερεῖς εἰς τὰ ἀφροδίσια οὕτως ἔλεγον Ἀττικοί.
\noindent \textbf{(2)} Platon prit les trois mêmes comme exemples d'hommes d'état honorés alors qu'ils sont en partie responsables de la crise athénienne (\italique{Gorgias}, 519a) :
\citComm{ἄνευ γὰρ σωφροσύνης καὶ δικαιοσύνης λιμένων καὶ νεωρίων καὶ τειχῶν καὶ φόρων καὶ τοιούτων φλυαριῶν ἐμπεπλήκασι τὴν πόλιν.\\
sans se préoccuper de la sagesse ni de la justice, ils ont gorgé la ville de ports, d'arsenaux, de murs, de tributs \& autres niaiseries.}
\noindent \textbf{(3)} Même jeu de mots avec le nom de Κύννα au vers 1032 des \italique{Guêpes} : ἀπ᾽ ὀφθαλμῶν Κύννης rappelle l'expression homérique des \italique{yeux d'un chien} (\italique{Iliade}, I, 225).
\noindent \textbf{(4)} Περὶ βλασφημιῶν, Περὶ παιδιῶν, édité en 1967 (aux Belles Lettres) par Jean Taillardat -- taillardesque édition, c'est tout dire ! βακχός est la \nieme{127} βλασφημία.
\noindent \textbf{(5)} Vers 1504 et 1532-3.
\noindent \textbf{(6)} Vers 678-682. L'on se reportera pour l'exégèse au toujours brillant Alphonse Willems, t. III, p. 115-117.
\noindent \textbf{(1)} Voir \italique{A Lexicon of Greek Personal Names} de P. M. Fraser et E. Matthews.
\noindent \textbf{(2)} Στρατονίκη est porté par plusieurs membres de la famille royale macédonienne, dont l'une fut donnée en mariage à l'un des rois de Thrace Σεύθης, contemporain des \italique{Thesmophories} (voir Thucydide, II, 101 et IV, 101). Cicéron (\italique{Tusculanes}, V, 59) nous apprend que Denys lorsqu'il devint Tyran de Syracuse en 406 comptait une certaine \latin{Aristomache} parmi ses épouses.
\noindent\textbf{808-9} \textbf{ἀλλ᾽ Εὐβούλης τῶν πέρυσίν τις βουλευτής ἐστιν ἀμείνων / παραδοὺς ἑτέρῳ τὴν βουλείαν; οὐδ᾽ αὐτὸς τοῦτό γε φήσεις.} Ces deux vers ne sont qu'allusifs et ont été diversement interprétés. Ils sont d'autant plus importants que c'est peut-être le seul passage de la comédie qui puisse nous aider véritablement à dater les \italique{Thesmophories}. (Je renvoie à l'Introduction sur la question de la date de notre pièce.) À quel événement font-ils allusion ? Deux écoles s'affrontent :\\ \ \ \ 1°) \ celle de Wilamowitz-Willems\marqueur{1} qui situe la représentation des \italique{Thesmophories} aux Dionysies de 411 : le Conseil des Cinq-Cents bouleutes fut rendu tout-à-fait caduc par l'instauration d'un conseil de dix vieillards, les πρόβουλοι\marqueur{2}, magistrats à compétence presque illimitée qui prirent en main la direction des affaires. Willems continue : «~C'était délivrer au Conseil des Cinq-Cents un certificat d'incapacité  $[\ldots]$. Aussi, quand il sortit de charge en skirophorion (juin) 412, n'existait-il plus que de nom.
\end{multicols}
}%
\end{adjustbox}%
\vfill
\end{adjustwidth}
\newpage

% --- FOLIO 20 (f. 172r) : SPREAD 2.2 (Texte grec établi + Traduction française + Commentaire 3/8 & 4/8) ---
\clearpage\ifodd\value{page}\hbox{}\newpage\fi
\enTeteG{ΘΕΣΜΟΦΟΡΙΑΖΟΥΣΑΙ --- Texte grec établi (f. 172r, vv. 800–837 [2/4])}
\enTeteD{LES THESMOPHORIES --- Traduction française (f. 172r, vv. 800–837 [2/4])}
\begin{adjustwidth}{-2cm}{-2cm}
\noindent\textbf{\footnotesize Texte grec établi --- Folio 20 (\texttt{f. 172r}, vv. 800–837 [2/4]) \hfill ΠΑΡΑΒΑΣΙΣ}\\[-2pt]
\noindent\rule{16.5cm}{0.5pt}\\[2pt]
{\fontsize{8.8pt}{10.2pt}\selectfont
\ligneEst{}{\Automatique{παραδοὺς ἑτέρῳ τὴν βουλείαν; οὐδ᾽ αὐτὸς τοῦτό γε φήσει.}}{}
\ligneEst{}{\Automatique{οὕτως ἡμεῖς πολὺ βελτίους τῶν ἀνδρῶν εὐχόμεθ᾽ εἶναι.}}{810}
\ligneEst{}{\Automatique{οὐδ᾽ ἂν κλέψασα γυνὴ ζεύγει κατὰ πεντήκοντα τάλαντα}}{}
\ligneEst{}{\Automatique{εἰς πόλιν ἔλθοι τῶν δημοσίων· ἀλλ᾽ ἢν τὰ μέγισθ᾽ ὑφέληται}}{}
\ligneEst{}{\Automatique{φορμὸν πυρῶν τἀνδρὸς κλέψασ᾽, αὐθημερὸν αὔτ᾽ ἀπέδωκεν}}{}
\ligneEst{\textbf{} {\tiny\textit{πνῖγος / \decalage{πνῖγος \ }dim. anap.}}}{ἀλλ᾽ ἡμεῖς ἂν πολλοὺς τούτων}{}
\ligneEst{}{ἀποδείξαιμεν ταῦτα ποιοῦντας.}{815}
\ligneEst{}{καὶ πρὸς τούτοις γάστριδας ἡμῶν}{}
\ligneEst{}{ὄντας μᾶλλον καὶ λωποδύτας}{}
}
\vspace{2pt}
\noindent\rule{16.5cm}{0.6pt}\\[2pt]
\noindent\textbf{\small Commentaire philologique, métrique et scénique (vv. 800–837 --- partie 3/8)}\\[2pt]
\begin{adjustbox}{max totalheight=15.2cm, max width=16.5cm, keepaspectratio}%
\parbox[t]{16.5cm}{
\setlength{\columnsep}{9pt}%
\begin{multicols}{2}
\fontsize{7.6pt}{8.8pt}\selectfont
\setlength{\parindent}{0pt}%
\setlength{\parskip}{1.2pt}%
\noindent [...] C'est à cette transmission annuelle de pouvoirs (παράδοσις, c'est le terme consacré), par laquelle chaque bouleute sortant cédait sa place à son successeur, que le chœur fait allusion~» ;\\ \ \ \ 2°) \ celle de Rogers -- et à laquelle je me rattache -- qui considère que la pièce est d'un an plus tardive : en mai 411, les Quatre Cents conjurés favorables à l'oligarchie se rendirent au βουλευτήριον et versèrent leur salaire aux bouleutes alors en séance. Les bouleutes se virent intimer l'ordre de se retirer, ce qu'ils firent sans même que les conjurés eussent à recourir à la violence\marqueur{3}. Ainsi périt une première fois la démocratie athénienne. C'est ce manque de courage (ἀμείνων) qu'Aristophane brocarde : pas un n'aura défendu la charge qu'ils avaient reçue. Que ces piteux représentants de la démocratie fussent ceux de l'\italique{an dernier} (τῶν πέρυσίν) place donc notre pièce aux Dionysies de l'an 410.
\noindent\textbf{811-2} \textbf{οὐδ᾽ ἂν κλέψασα γυνὴ ζεύγει κατὰ πεντήκοντα τάλαντα / εἰς πόλιν ἔλθοι τῶν δημοσίων·} Aristophane use de la même formule κατὰ πεντήκοντα τάλαντα, «~cinquante talents d'un coup~», dans les \italique{Guêpes} : Bdélycléon tance son paternel berné par les démagogues qui, non contents de dissiper l'argent de la cité, détournent le tribut des autres cités (vers 666-671) :
\citComm{Φιλοκλέων --- καὶ ποῖ τρέπεται δὴ ᾽πειτα τὰ χρήματα τἄλλα;\\
Βδελυκλέων --- εἰς τούτους τοὺς “οὐχὶ προδώσω τὸν Ἀθηναίων κολοσυρτόν,\\
ἀλλὰ μαχοῦμαι περὶ τοῦ πλήθους ἀεί”. σὺ γάρ, ὦ πάτερ, αὐτοὺς\\
ἄρχειν αἱρεῖ σαυτοῦ τούτοις τοῖς ῥηματίοις περιπεφθείς.\\
κᾆθ᾽ οὗτοι μὲν δωροδοκοῦσιν \textbf{κατὰ πεντήκοντα τάλαντα}\\
ἀπὸ τῶν πόλεων ἐπαπειλοῦντες τοιαυτὶ κἀναφοβοῦντες·\\
“δώσετε τὸν φόρον, ἢ βροντήσας τὴν πόλιν ὑμῶν ἀνατρέψω.”\\
Philocléon --- \francais{Et où passe alors le reste de l'argent ?}\\
Bdélycléon --- \francais{À ceux qui nous disent: « Non, je ne trahirai point la masse athénienne, toujours je combattrai pour le populaire ! » Car c'est toi, mon père, qui te les donnes pour maîtres, enjôlé par leurs boniments. Eux cependant extorquent aux cités des cinquante talents à la fois, en les menaçant, en les terrifiant de cette sorte : « Vous paierez le tribut, ou je tonne contre votre ville et la renverse. »}}
\noindent Hérodote (VI, 132-136) rapporte une histoire semblable : celle de Miltiade. Le vainqueur de Marathon, promit aux Athéniens quantité d'or s'ils lui remettaient soixante-dix vaisseaux avec troupe et argent. Les Athéniens lui fournirent le tout et Miltiade de cingler vers Paros pour en faire le siège. Il envoya un héraut réclamer 100 talents avec le même genre de menaces que sus-citées. L'expédition fut un échec et l'on intenta un procès à Miltiade pour avoir trompé la cité (τῆς Ἀθηναίων ἀπάτης εἵνεκεν). Or Miltiade se mourait d'une mauvaise blessure.
\noindent Compte tenu des services rendus à la nation, et que c'était certainement peine perdue de transporter ce semi-cadavre jusqu'au \italique{barathre}\marqueur{1}, l'agonisant ne fut condamné qu'à restituer \italique{cinquante talents}.
\noindent L'histoire de Miltiade et les deux passages d'Aristophane coïncident jusque dans le chiffre du détournement !..
\end{multicols}
}%
\end{adjustbox}%
\vfill
\end{adjustwidth}
\newpage
% --- PAGE DE DROITE : Traduction française + Commentaire (4/8) ---
\begin{adjustwidth}{-2cm}{-2cm}
\noindent\textbf{\footnotesize Traduction française --- Folio 20 (\texttt{f. 172r}, vv. 800–837 [2/4]) \hfill ΠΑΡΑΒΑΣΙΣ}\\[-2pt]
\noindent\rule{16.5cm}{0.5pt}\\[2pt]
{\fontsize{8.8pt}{10.2pt}\selectfont
\ligneEst{}{\francais{après avoir résigné à un autre son office, l'emporte sur elle ? Lui-même s'en défendra.}}{}
\ligneEst{}{\francais{Ainsi nous nous prévalons d'une grande supériorité sur les hommes.}}{810}
\ligneEst{}{\francais{Joignez qu'une femme, après avoir diverti par cinquante talents les deniers publics,}}{}
\ligneEst{}{\francais{ne se rendra pas à l'Acropole sur un char ; mais eût-elle commis le plus grave des larcins,}}{}
\ligneEst{}{\francais{un boisseau de froment volé à son mari, elle le restituera le jour même.}}{}
\ligneEst{}{\francais{Mais nous pourrions en désigner plusieurs parmi ceux-là}}{}
\ligneEst{}{\francais{qui en font autant,}}{815}
\ligneEst{}{\francais{et qui, par là-dessus, sont gloutons}}{}
\ligneEst{}{\francais{plus que nous, tire-laine,}}{}
}
\vspace{2pt}
\noindent\rule{16.5cm}{0.6pt}\\[2pt]
\noindent\textbf{\small Commentaire philologique, métrique et scénique (vv. 800–837 --- partie 4/8)}\\[2pt]
\begin{adjustbox}{max totalheight=15.2cm, max width=16.5cm, keepaspectratio}%
\parbox[t]{16.5cm}{
\setlength{\columnsep}{9pt}%
\begin{multicols}{2}
\fontsize{7.6pt}{8.8pt}\selectfont
\setlength{\parindent}{0pt}%
\setlength{\parskip}{1.2pt}%
\noindent Miltiade, le héros de Marathon -- τὸν Μαραθῶνι, c'est son épithète, quasi-homérique, qui lui est inlassablement accolée --, profita du crédit qu'il avait auprès de Dêmos, le cajolant des souvenirs de Marathon\marqueur{2}, pour mener une expédition dont le but véritable était d'enflouer ses propres caisses.\\ En regard des déboires de Miltiade, la construction de l'argumentation du chœur féminin se comprend parfaitement : «~Devant Aristomaque, notre héroïne de Marathon (τὴν Μαραθῶνι), aucun d'entre vous n'entreprendrait de se battre~», et, parant d'avance à l'objection que les hommes ont aussi leur marathonien héros en la personne de Miltiade, l'on précise aussitôt : «~Une femme ayant volé cinquante talents d'un (bon) coup n'aurait pas l'outrecuidance de venir parader sur un char lors de son retour à la cité~».\\ \italique{Confer} sur ces deux vers 811-812 la note 24 d'Alphonse Willems (reproduite en annexe) : à la fête des Panathénées, la file brillante des chars était ce qui attirait le plus les regards dans la procession (Schoemann, \italique{Antiquités grecques}, t. II, p. 549).
\noindent\textbf{[Notes ad 808-9]} \textbf{(1)} Voir la \nieme{23} note de Willems, reproduite entièrement à la fin.
\noindent \textbf{(2)} Voir Thucydide VIII, 1 et la \italique{Constitution d'Athènes}, XXIX ; Aristophane met en scène un Probule dans \italique{Lysistrata} (v. 387-613), jouée quelques mois auparavant -- là est la question : combien ? deux mois comme le pense Willems ?.. ou quatorze comme le pense Rogers ?
\noindent \textbf{(3)} Événement conté par Thucydide aux chapitre 69-70 de son livre VIII.
\noindent \textbf{(1)} Platon (\italique{Gorgias}, 516e) nous apprend que Miltiade devait être initialement précipité vers le \italique{barathre} (εἰς τὸ βάραθρον ἐμβαλεῖν), illustrant l'adage latin qu'il n'y a pas loin du Capitole à la roche tarpéienne. (Pour en apprendre davantage sur ce gouffre, que le lecteur se reporte au \italique{Dictionnaire des antiquités grecques et romaines}, à l'entrée \latin{Barathron}.) Si Miltiade s'est vu gracier -- la Parque se consola néanmoins très rapidement : la grâce ne fut qu'un sursis de quelques jours --, la raison principale est peut-être qu'il bénéficiait de nombreuses relations, critère qui distingue souvent les condamnés à mort exécutés de ceux qui voient leur peine commuée, selon le mot de Gide.
\noindent \textbf{(2)} «~Flattez-les, payez-les d'agréables mensonges : /	Quelque indignation dont leur cœur soit rempli, / Ils goberont l'appât ; vous serez leur ami.~»
\noindent\textbf{814} \textbf{τούτων.} Les spectateurs.
\noindent\textbf{815} \textbf{ποθοῦντας.} L'erreur manifeste du scribe du \italique{Ravennas}, copiant ποθοῦντας à la place de πο(ι)οῦντας, m'est l'occasion -- ὁ καιρός ! -- d'une courte digression portant sur les différentes graphies ποι- ou πο- lorsque la lettre qui suit est un `ο' ou un `ω'\marqueur{1}.
\noindent Je m'en tiens aux leçons de l'épigraphie attique ainsi qu'à l'usage scolaire : savoir, écrire systématiquement ποι-. Toutefois, le \italique{Ravennas} adopte majoritairement la graphie πο- lorsque la syllabe est brève. Cette convention, reprise par la plupart des éditions actuelles des \italique{Thesmophories}, est celle de nombreux manuscrits remontant à la période hellénistique. La confusion du scribe -- qui aurait lu rapidement une suite de `Ο' en y ajoutant un `Θ' -- tend à faire penser que cette règle était aussi de mise pour le manuscrit qu'il recopiait.
\noindent\textbf{819-23} \textbf{} Le manuscrit présente le texte suivant :
\citComm{καὶ μὲν δήπου καὶ τὰ πατρῷά γε\\
χείρους ἡμῶν εἰσιν σῴζειν·\\
ἡμῖν μὲν γὰρ σῶν ἔτι καὶ νῦν\\
τἀντίον ὁ κανὼν οἱ καλαθίσκοι\\
τὸ σκιάδειον·}
\end{multicols}
}%
\end{adjustbox}%
\vfill
\end{adjustwidth}
\newpage

% --- FOLIO 20 (f. 172r) : SPREAD 2.3 (Texte grec établi + Traduction française + Commentaire 5/8 & 6/8) ---
\clearpage\ifodd\value{page}\hbox{}\newpage\fi
\enTeteG{ΘΕΣΜΟΦΟΡΙΑΖΟΥΣΑΙ --- Texte grec établi (f. 172r, vv. 800–837 [3/4])}
\enTeteD{LES THESMOPHORIES --- Traduction française (f. 172r, vv. 800–837 [3/4])}
\begin{adjustwidth}{-2cm}{-2cm}
\noindent\textbf{\footnotesize Texte grec établi --- Folio 20 (\texttt{f. 172r}, vv. 800–837 [3/4]) \hfill ΠΑΡΑΒΑΣΙΣ}\\[-2pt]
\noindent\rule{16.5cm}{0.5pt}\\[2pt]
{\fontsize{8.8pt}{10.2pt}\selectfont
\ligneEst{}{καὶ βωμολόχους κἀνδραποδιστάς.}{}
\ligneEst{}{καὶ μὲν δήπου καὶ τὰ πατρῷά γε}{}
\ligneEst{}{χείρους ἡμῶν εἰσιν σῴζειν·}{820}
\ligneEst{}{ἡμῖν μὲν γὰρ σῶν ἔτι καὶ νῦν}{}
\ligneEst{}{ταὐτὸ κανοῦν καὶ τὸ σκιάδειον·}{}
\ligneEst{}{τοῖς δ᾽ ἡμετέροις ἀνδράσι τούτοις}{}
\ligneEst{}{ἀπόλωλεν μὲν πολλοῖς ὁ κανὼν}{825}
\ligneEst{}{ἐκ τῶν οἴκων αὐτῇ λόγχῃ,}{}
\ligneEst{}{πολλοῖς δ᾽ ἑτέροις ἀπὸ τῶν ὤμων}{}
}
\vspace{2pt}
\noindent\rule{16.5cm}{0.6pt}\\[2pt]
\noindent\textbf{\small Commentaire philologique, métrique et scénique (vv. 800–837 --- partie 5/8)}\\[2pt]
\begin{adjustbox}{max totalheight=15.2cm, max width=16.5cm, keepaspectratio}%
\parbox[t]{16.5cm}{
\setlength{\columnsep}{9pt}%
\begin{multicols}{2}
\fontsize{7.6pt}{8.8pt}\selectfont
\setlength{\parindent}{0pt}%
\setlength{\parskip}{1.2pt}%
\noindent Les trois premiers vers ne posent pas de difficulté. Et encore : à condition de garder à l'esprit les divers sens de τὰ πατρῷα, \italique{l'héritage} : désigne les \italique{biens transmis par les ancêtres} ; \italique{héritage} tout matériel que seul une femme sera à même de conserver\marqueur{1}, mais héritage aussi immatériel, \italique{les traditions}\marqueur{2}, que les femmes perpétuent tout-de-même moins mal que les hommes (χείρους ἡμῶν εἰσιν σῴζειν).
\noindent Suit une liste de ces biens que les femmes sauvent mais que leur mari, eux, perdent. C'est là que le bât blesse.
\noindent D'abord, cette liste est étonnamment hétéroclite ; si les trois premiers éléments qu'elle comporte sont en rapport avec le filage et le tissage : τὸ ἀντίον, \italique{l'ensouple} (cyclindre du métier à tisser) ; ὁ κανὼν, \italique{barre}, \italique{règle} (avec le sens spécialisé de \italique{tige de quenouille}), οἱ καλαθίσκοι, \italique{petits paniers} (où les femmes déposaient la laine filée), l'on se demande par quel ressort \italique{l'ombrelle}, τὸ σκιάδειον, vient clôturer cette liste.
\noindent Plus important : sur ces quatre éléments, seuls deux, κανών et σκιάδειον, subsistent dans la comparaison avec ce qu'en font les hommes. Pourquoi Aristophane nous parle-t-il alors de l'ἀντίον et des καλαθίσκοι, inutiles -- donc nuisibles -- si ce n'est pas pour filer la métaphore quelques vers plus loin ? (Je ne parle que de ce qui achoppe au niveau du sens. Quant à la forme, l'on avait déjà relevé l'inhabituelle succession de quatre brèves au vers 822.)\\
\noindent Aristophane veut nous peindre la sotte dissipation des hommes ; Interrogeons les parémiographes\marqueur{3} pour savoir comment on évoque une prodigalité de cigale ou une gestion de fourmi dans sa langue imagée et proverbiale :
\citComm{Αὐτῷ κανῷ: αὕτη τάττεται κατὰ τῶν ἄρδην ὁτιοῦν λυμαινομένων. Φησὶ γὰρ ὁ Δίδυμος ὅτι τὰ δεῖπνα ἐπὶ κανῶν ἐκομίζετο. Καθάπερ ἐν Γρυψὶ λέγει καὶ Πλάτων·\\
Ἀνήρπακας ἁπαξάπαντ᾿ \textbf{αὐτῷ κανῷ}.\\
Et la corbeille elle-même : ce [proverbe] se range contre ceux prompts à se ruiner de quelque manière. Didyme dit que les repas étaient emportés dans des paniers. Ainsi Platon dans ses \italique{Griffons} :\\
Tu as tout emporté, tout, \italique{et le panier avec} !..}
\noindent Le grec ajoute «~et le panier~» pour désigner tous les biens jusqu'au dernier\marqueur{4}, comme en français on ajoute «~et son père~» pour désigner «~tout le monde~». Dans les \italique{Grenouilles} (vers 560), ce n'est pas autrement qu'il faut entendre le désarroi de l'hôtesse face à Hercule, cet ogre qui a liquidé toutes ses provisions :
\citComm{Μὰ Δί᾽ οὐδὲ τὸν τυρόν γε τὸν χλωρόν, τάλαν,\\
ὃν οὗτος \textbf{αὐτοῖς τοῖς ταλάροις} κατήσθιεν.\\
Par Zeus, et ce fromage frais, pauvre amie, qu'il dévorait \italique{avec les éclisses elles-mêmes}.}
\noindent C'est peut-être encore ce même proverbe qu'on retrouve au vers 170 des \italique{Guêpes} dans la bouche de Philocléon :
\citComm{μὰ τὸν Δί᾽ οὐ δῆτ᾽, ἀλλ᾽ ἀποδόσθαι βούλομαι\\
τὸν ὄνον ἄγων \textbf{αὐτοῖσι τοῖς κανθηλίοις}·\\
νουμηνία γάρ ἐστιν.\\
Par Zeus, certes non ! Je veux aller vendre l'âne \italique{avec ses paniers} ; car c'est nouvelle lune.}
\end{multicols}
}%
\end{adjustbox}%
\vfill
\end{adjustwidth}
\newpage
% --- PAGE DE DROITE : Traduction française + Commentaire (6/8) ---
\begin{adjustwidth}{-2cm}{-2cm}
\noindent\textbf{\footnotesize Traduction française --- Folio 20 (\texttt{f. 172r}, vv. 800–837 [3/4]) \hfill ΠΑΡΑΒΑΣΙΣ}\\[-2pt]
\noindent\rule{16.5cm}{0.5pt}\\[2pt]
{\fontsize{8.8pt}{10.2pt}\selectfont
\ligneEst{}{\francais{charlatans et marchands d'hommes.}}{}
\ligneEst{}{\francais{Et assurément aussi ils savent moins que nous}}{}
\ligneEst{}{\francais{conserver les traditions}}{820}
\ligneEst{}{\francais{qui sont par notre fait sauves encore aujourd'hui}}{}
\ligneEst{}{\francais{– panier et ombrelle compris.}}{}
\ligneEst{}{\francais{Au contraire, du fait de nos nombreux maris,}}{}
\ligneEst{}{\francais{la lance a ruiné}}{825}
\ligneEst{}{\francais{notre capital – lance comprise.}}{}
\ligneEst{}{\francais{Et maint autre}}{}
}
\vspace{2pt}
\noindent\rule{16.5cm}{0.6pt}\\[2pt]
\noindent\textbf{\small Commentaire philologique, métrique et scénique (vv. 800–837 --- partie 6/8)}\\[2pt]
\begin{adjustbox}{max totalheight=15.2cm, max width=16.5cm, keepaspectratio}%
\parbox[t]{16.5cm}{
\setlength{\columnsep}{9pt}%
\begin{multicols}{2}
\fontsize{7.6pt}{8.8pt}\selectfont
\setlength{\parindent}{0pt}%
\setlength{\parskip}{1.2pt}%
\noindent L'«~αὐτῇ λόγχῃ~» du vers 826 -- noter la présence du pronom --, mis comme permutant métaphorique d'«~αὐτῷ κανόνι~», rappelle le détournement du proverbe avec jeu de mots sur κανών / κανοῦν.\\ Notre conclusion est que le texte original des \italique{Thesmophories} -- l'\allemand{Urtext} -- comportait certainement au vers 822 le terme κανοῦν, \italique{corbeille}, et non κανών.
\noindent La singulière présence des καλαθίσκοι, \italique{les corbeilles à laine}, s'explique alors par l'insertion fort importune dans le texte d'une glose de κανοῦν.
\noindent \\ L'image ne porte pas sur l'activité filandière ; la présence simultanée de corbeilles et d'une ombrelle ne laisse pas de doute possible : Aristophane nous décrit une procession religieuse\marqueur{5}.
\noindent Au sein de ces cortèges traditionnels, on voyait défiler les \italique{cané-}phores, qui avaient la charge des corbeilles à offrandes, et qu'elles portaient sur la tête\marqueur{6}, quand d'autres, les \italique{skiro-}phores, les accompagnaient\marqueur{7} en tenant de deux mains fermes et assurées, la lourde \italique{ombrelle} consacrée, τὸ σκίρον ou τὸ σκιάδειον.
\noindent Et pour Aristophane, sa foi, ces parades d'Athéniennes défilant cahin-caha n'étaient pas tout-à-fait indignes des traditions reçues des ancêtres (τὰ πατρῷα).
\noindent À la place des vers 822-823, je propose donc de lire :
\citComm{καὶ μὲν δήπου καὶ τὰ πατρῷά γε\\
χείρους ἡμῶν εἰσιν σῴζειν·\\
ἡμῖν μὲν γὰρ σῶν ἔτι καὶ νῦν\\
ταὐτὸ κανοῦν\marqueur{8} [καὶ] τὸ σκιάδειον·\\
Et assurément aussi, ils sont moins aptes que nous\\
à garder les traditions\\
qui sont, par notre fait, sauves encore aujourd'hui\\
-- même le panier [et] l'ombrelle.}
\noindent\textbf{[Notes ad 815]} \textbf{(1)} Débat qui a fait déjà couler beaucoup d'encre et que je ne compte ni trancher, ni même exposer en détail. Le lecteur intéressé pourra se reporter aux pages 226-230 que lui consacre Willem J. W. Koster dans son livre \italique{Autour d'un manuscrit d'Aristophane écrit par Demetrius Triclinius} (1957).
\noindent \textbf{(1)} L'homme qui avait aliéné son patrimoine, ὁ τὰ πατρῷα κατεδηδοκώς, était frappé d'ἀτιμία. Eschine, \italique{In Timarchum}, 30 ; Diogène Laërce, \italique{Solon}, I, 55. Cette loi qui n'était certainement plus observée au temps d'Eschine subsistait pour la forme, comme un vestige de l'antique règle ; il y eut toujours une δίκη κατεδηδοκέναι τὰ πατρῷα (Bekker, \italique{Anecdota}, p. 199 et 310 ; note de Fustel de Coulanges).
\noindent \textbf{(2)} Certains commentateurs n'ont pas manqué de noter que les Athéniennes n'auraient pas eu légalement de droit de propriété. Cette remarque aurait pu -- aurait dû -- leur mettre la puce à l'oreille : c'est bien en tant que championnes des \italique{traditions} que se posent les femmes du chœur dans ce passage.
\noindent \textbf{(3)} Leutsch \& Schneidewim, \italique{Corpus Paroemiographorum Graecorum}, vol. I, Zenobius, II, 31. Recensé par Érasme (\italique{Adages}, n° 1088). Kassel \& Austin l'ont intégré comme \nieme{15} fragment de Platon le Comique. Suidas (α 4465 ; α 4466 ; α 4540) a relevé des expressions analogues : αὐταῖς ἀμάλαις ; αὐταῖσι διαβολαῖς ; αὐτοῖς ὅπλοις ; αὐτοῖς ταλάροις.
\end{multicols}
}%
\end{adjustbox}%
\vfill
\end{adjustwidth}
\newpage

% --- FOLIO 20 (f. 172r) : SPREAD 2.4 (Texte grec établi + Traduction française + Commentaire 7/8 & 8/8) ---
\clearpage\ifodd\value{page}\hbox{}\newpage\fi
\enTeteG{ΘΕΣΜΟΦΟΡΙΑΖΟΥΣΑΙ --- Texte grec établi (f. 172r, vv. 800–837 [4/4])}
\enTeteD{LES THESMOPHORIES --- Traduction française (f. 172r, vv. 800–837 [4/4])}
\begin{adjustwidth}{-2cm}{-2cm}
\noindent\textbf{\footnotesize Texte grec établi --- Folio 20 (\texttt{f. 172r}, vv. 800–837 [4/4]) \hfill ΠΑΡΑΒΑΣΙΣ}\\[-2pt]
\noindent\rule{16.5cm}{0.5pt}\\[2pt]
{\fontsize{8.8pt}{10.2pt}\selectfont
\ligneEst{}{\{ ἐν ταῖς στρατιαῖς \}}{}
\ligneEst{{\tiny\textit{(catalect.)}}}{ἔρριπται τὸ σκιάδειον.}{}
\ligneEst{\textbf{Κορυφαῖος} {\tiny\textit{$4$ troch\cat}}}{πόλλ᾽ ἂν αἱ γυναῖκες ἡμεῖς ἐν δίκῃ μεμψαίμεθ᾽ ἂν}{830}
\ligneEst{}{τοῖσιν ἀνδράσιν δικαίως, ἓν δ᾽ ὑπερφυέστατον.}{}
\ligneEst{}{\Automatique{χρῆν γάρ, ἡμῶν εἰ τέκοι τις ἄνδρα χρηστὸν τῇ πόλει,}}{}
\ligneEst{}{ταξίαρχον ἢ στρατηγόν, λαμβάνειν τιμήν τινα,}{}
\ligneEst{}{\Automatique{προεδρίαν τ᾽ αὐτῇ δίδοσθαι Στηνίοισι καὶ Σκίροις}}{}
\ligneEst{}{ἔν τε ταῖς ἄλλαις ἑορταῖς αἷσιν ἡμεῖς ἤγομεν·}{835}
\ligneEst{}{εἰ δὲ δειλὸν καὶ πονηρὸν ἄνδρα τις τέκοι γυνή,}{}
\ligneEst{}{ἢ τριήραρχον πονηρὸν ἢ κυβερνήτην κακόν,}{837}
}
\vspace{2pt}
\noindent\rule{16.5cm}{0.6pt}\\[2pt]
\noindent\textbf{\small Commentaire philologique, métrique et scénique (vv. 800–837 --- partie 7/8)}\\[2pt]
\begin{adjustbox}{max totalheight=15.2cm, max width=16.5cm, keepaspectratio}%
\parbox[t]{16.5cm}{
\setlength{\columnsep}{9pt}%
\begin{multicols}{2}
\fontsize{7.6pt}{8.8pt}\selectfont
\setlength{\parindent}{0pt}%
\setlength{\parskip}{1.2pt}%
\noindent \textbf{(4)} Biens dévorés, en l'occurrence. Mais l'image du panier comme signe d'abondance se retrouve dans un autre proverbe, transmis par Hésychius (ε 7077) : εὔπλουτον κανοῦν.
\noindent \textbf{(5)} La fête des Σκίρα est évoquée quelques vers plus loin. Cf. les articles \latin{Umbella} et \latin{Skirophoria} du \italique{Dictionnaire des antiquités grecques et romaines} de Daremberg \& Saglio.
\noindent \textbf{(6)} Ces canéphores ont certainement servi de modèle aux fameuses Caryatides, portant, elles, un chapiteau sur la tête.
\noindent \textbf{(7)} Voir, par exemple, le vers 1550 des \italique{Oiseaux}. Les ombrelles tenaient effectivement un rôle important dans la liturgie.
\noindent \textbf{(8)} Au vers 164, il y a confusion de finale -ων / -ουν du même genre entre les leçons du \italique{Ravennas} (διεκίνων) et de la \italique{Souda} ε 989 (διεκίνουν).
\noindent\textbf{826} \textbf{ἐκ τῶν οἴκων.} Précisons le sens «~ἐκ τῶν οἴκων~». Ce ne sera pas du luxe. Car il ne peut être question de rompre des lances \italique{à l'extérieur des maisons} plutôt qu'à \italique{à l'intérieur} ?
\noindent Devrons-nous aussi commenter les commentaires qui voient en cette panoplie militaire le fameux héritage paternel, τὰ πατρῷα, que les fils prodigues indignes dissiperaient -- ô scandale !
\noindent Au garde-à-vous, je félicite ces brillants universitaires pour leurs éclaircissements : sans eux, je crois que les douceurs martiales de la poésie d'Aristophane m'auraient échappé. Notre poète roulant des variations sur le thème : «~Tu seras un soldat, mon fils !~», c'est à vous tirer des larmes (de guerre, bien-sûr !)\\
\noindent Οἶκος, ce mot rare et qui renvoie à des réalités n'existant pas ailleurs qu'en Grèce classique a tout-de-même donné le terme οἰκο-νομία, qui me semble être la bonne gestion d'une \italique{maison}, perçue comme un \italique{fonds} ou un \italique{capital}. «~ἐκ τῶν οἴκων~» désigne tout ce que l'on \italique{tire de} ce fonds, ou de ce capital. La bonne équation est évidemment celle-ci : τὰ πατρῷα $ = $ ἐκ τῶν οἴκων.
\noindent La même tournure avec le même sens se trouve chez Xénophon dans son ... \italique{Économique} (XX, 21) :
\citComm{τὰ οὖν συντρίβοντα τοὺς οἴκους πολῦ μᾶλλον ταῦτά ἐστιν ἢ αἱ λίαν ἀνεπιστημοσύναι. τὸ γὰρ τὰς μὲν δαπάνας χωρεῖν ἐντελεῖς \textbf{ἐκ τῶν οἴκων}, τὰ δὲ ἔργα μὴ τελεῖσθαι λυσιτελούντως πρὸς τὴν δαπάνην\\
Voilà qui ruine une maison beaucoup plus sûrement qu'une excessive ignorance : d'un côté, dépenser complètement \italique{ce qu'on tire des maisons}, de l'autre des travaux qui ne sont pas profitables à la dépense}
\noindent Traduisez donc :
\citComm{τοῖς δ᾽ ἡμετέροις ἀνδράσι τούτοις\\
ἀπόλωλεν μὲν πολλοῖς ὁ κανὼν\\
ἐκ τῶν οἴκων αὐτῇ λόγχῃ\marqueur{1}.\\
Du fait de nos ci-devant maris (elles les désignent),\\
la lance a anéanti pour beaucoup\\
ce qu'on tirait de nos maisons -- lance comprise\marqueur{1}.}
\end{multicols}
}%
\end{adjustbox}%
\vfill
\end{adjustwidth}
\newpage
% --- PAGE DE DROITE : Traduction française + Commentaire (8/8) ---
\begin{adjustwidth}{-2cm}{-2cm}
\noindent\textbf{\footnotesize Traduction française --- Folio 20 (\texttt{f. 172r}, vv. 800–837 [4/4]) \hfill ΠΑΡΑΒΑΣΙΣ}\\[-2pt]
\noindent\rule{16.5cm}{0.5pt}\\[2pt]
{\fontsize{8.8pt}{10.2pt}\selectfont
\ligneEst{}{\francais{$\{$dans ses campagnes$\}$}}{}
\ligneEst{}{\francais{a jeté à bas de ses épaules... l'ombrelle.}}{}
\ligneEst{\textbf{Coryphée}}{\francais{Nous autres femmes aurions droit en toute justice}}{830}
\ligneEst{}{\francais{d'imputer aux hommes bien des torts, un surtout passant toute mesure.}}{}
\ligneEst{}{\francais{Il serait de votre devoir, quand l'une de nous a donné à la patrie un homme utile,}}{}
\ligneEst{}{\francais{taxiarque ou stratège, d'avoir pour elle quelque considération,}}{}
\ligneEst{}{\francais{de lui accorder la proédrie dans les Sténies, les Scirophories}}{}
\ligneEst{}{\francais{et les autres fêtes que nous célébrons.}}{835}
\ligneEst{}{\francais{Celle, au contraire, qui aurait donné le jour à un pleutre, à un propre à rien,}}{}
\ligneEst{}{\francais{méchant triérarque ou lâche commandant de navire,}}{837}
}
\vspace{2pt}
\noindent\rule{16.5cm}{0.6pt}\\[2pt]
\noindent\textbf{\small Commentaire philologique, métrique et scénique (vv. 800–837 --- partie 8/8)}\\[2pt]
\begin{adjustbox}{max totalheight=15.2cm, max width=16.5cm, keepaspectratio}%
\parbox[t]{16.5cm}{
\setlength{\columnsep}{9pt}%
\begin{multicols}{2}
\fontsize{7.6pt}{8.8pt}\selectfont
\setlength{\parindent}{0pt}%
\setlength{\parskip}{1.2pt}%
\noindent\textbf{[Notes ad 826]} \textbf{(1)} Écho au proverbe : αὐτῷ κανῷ... sauf que le \italique{canon} est remplacé par la \italique{lance}. Autre sexe, autre proverbe :
\citComm{ἦν γάρ τις αἶνος, ὡς γυναιξὶ μὲν τέχναι\\
μέλουσι, \textbf{λόγχῃ δ ̓ ἄνδρες} εὐστοχώτεροι.\\
εἰ γὰρ δόλοισιν ἦν τὸ νικητήριον,\\
ἡμεῖς ἂν ἀνδρῶν εἴχομεν τυραννίδα.\\
Il existe en effet un proverbe qui dit qu'à la femme les artifices reviennent, tandis que \italique{les hommes sont plus habiles à atteindre le but par la lance}. Si le prix de la victoire allait à la ruse, c'est nous qui régnerions sur les hommes.}
\noindent (Euripide, fr. 6 \italique{Danaé} Jouan-van Looy $=$ fr. 321 Kannicht $=$ fr. 413 Mette.)
\noindent\textbf{827-9} \textbf{πολλοῖς δ᾽ ἑτέροις ἀπὸ τῶν ὤμων / ἐν ταῖς στρατιαῖς / ἔρριπται τὸ σκιάδειον.} Un bouclier est une \italique{ombrelle}. Cette métaphore unique est très claire : un bouclier est rond (cf. κύκλος, Taillardat § 642) et bombé comme une ombrelle. Mais on peut se demander si l'argot des soldats athéniens ne connaissait aussi pas cette métaphore unique\marqueur{1} : le bouclier protège des coups comme l'ombrelle du soleil, cf. Σ : τὴν ἀσπίδα λέγει· σκιάζει\marqueur{2} γὰρ τὸν ἄνδρα. Cette hypothèse est malheureusement invérifiable.
\noindent (Note de Taillardat, §643. L'utilisation deux vers plus haut de κανών comme synonyme de \italique{lance}, λόγχη, pourrait relever de ce même argot militaire dont Taillardat fait la très astucieuse hypothèse.)\\ La charge soldatesque des hommes ? elle leur passe par-dessus la tête et les épaules lorsqu'ils en délestent leur bouclier (ἀπὸ τῶν ὤμων ἔρριπται τὸ σκιάδειον), bien encombrant il est vrai lorsqu'il faut prendre la poudre d'escampette en plein milieu des combats...
\noindent Mica, la femme de Cléonyme -- Cléonyme qui semble avoir fait école avec sa célèbre ῥίψασπις (voir le commentaire au vers 605) -- devait probablement à ce moment être pointée, au moins des yeux.
\noindent\textbf{[Notes ad 827-9]} \textbf{(1)} À propos de θύλακος (§268), on [l'«~on~», ou l'«~om~» en question n'est autre que le grand Taillardat !] a déjà relevé une métaphore ironique qui a dû appartenir à l'argot des soldats athéniens. Cette tendance à ironiser sur ses armes a toujours existé : on sait que les soldats de César appelaient, par plaisanterie, \latin{cippi} et \latin{lilia} certaines défenses qu'ils construisaient devant Alésia (\italique{Bell. Gall.} 7, 73, 4 et 8). C'est de la même façon que dans l'argot militaire français, un fusil peut s'appeler une \italique{arbalète}, une \italique{seringue}, une \italique{canne à pêche}, etc. (voir A. Dauzat, \italique{L'argot de la guerre}\expo{2} (1919).)
\noindent \textbf{(2)} Σκιαζειν, \italique{couvrir} : Hésiod. \italique{Théog.} 716 ; Hdt. 6, 117, etc.
\noindent\textbf{834} \textbf{Στηνίοισι καὶ Σκίροις.} Fêtes athéniennes réservées aux femmes. C'est notamment aux Sciries que fut élaboré le projet politique de Praxagora (\italique{Assemblée des femmes}, vers 18).
\noindent\textbf{833-7} \textbf{ταξίαρχον ἢ στρατηγόν ... ἢ τριήραρχον πονηρὸν ἢ κυβερνήτην κακόν} Les exemples d'Aristophane sont tirés de l'actualité la plus brûlante.
\noindent Les mauvais triérarques et commandants de vaisseaux sont ceux responsables de la défaite navale d'Érétrie à l'été 411, défaite qui plongea les Athéniens dans une consternation, nous rapporte Thucydide (VIII, 95 et 96), plus grande encore que celle consécutive au désastre de Sicile. Défaite cinglante qui ne rattrape pas aux yeux d'Aristophane la récente et lointaine victoire de Κυνὸς Σῆμα (VIII, 104-108) dans l'Hellespont. (Cette victoire devait être bien amère à Aristophane, puisque son contemporain Eupolis y mourut. Nous y reviendrons.)
\end{multicols}
}%
\end{adjustbox}%
\vfill
\end{adjustwidth}
\newpage
\clearpage\ifodd\value{page}\hbox{}\newpage\fi
\enTeteG{ΠΑΡΑΒΑΣΙΣ / ΕΠΕΙΣΟΔΙΟΝ Δ' --- Codex Ravennas 429, f. 172v (Classense 174v) --- vv. 838–875}
\begin{adjustwidth}{-1.5cm}{-3.0cm}
\begin{adjustbox}{max totalheight=21.0cm, max width=16.5cm, keepaspectratio}
\begin{minipage}{16.5cm}
{\fontsize{7.2pt}{8.5pt}\selectfont
\noindent\textbf{Folio 21 --- Codex Ravennas 429, \texttt{f. 172v} (\texttt{Classense 174v}) : ΠΑΡΑΒΑΣΙΣ / ΕΠΕΙΣΟΔΙΟΝ Δ', vv. 838–875}\hfill \textit{Témoins :} $\text{R}^1, \Rpc, \Rdeux$, \PsiUn\par\vspace{2pt}
\hrule height 0.4pt\vspace{2.5pt}
}
{\fontsize{8.1pt}{9.5pt}\selectfont
\ligneRav{838}{χορ/}{ὑστέρ\auDessus{ο}{α}ν αὐτὴν καθῆσθαι σκάφιον ἀποκεκαρμένην}{\textbf{838} \lecon{ὑστέρ\auDessus{ο}{α}ν} \Rpc : ὑστάτην Σ}
\ligneRav{839}{}{τῆς τὸν ἀνδρεῖον τεκούσης· τῷ γὰρ εἰκὸς \difference{ωὖπολι}}{\textbf{839} ὦ πόλεις R : ὦ πόλις Gelenius : \lecon{ωὖπολι} Bibi}
\ligneRav{840}{}{τὴν Ὑπερβόλου καθῆσθαι μητέρ᾽ ἠμφιεσμένην}{}
\ligneRav{}{}{λευκὰ καὶ κόμας καθεῖσαν πλησίον τῆς Λαμάχου·}{}
\ligneRav{842}{}{καὶ δανείζειν \difference{χρήμαθ᾽}, ἥν χρῆν, εἰ \difference{δανείσειέν} τινι}{\textbf{842} χρήματα τ' R : \lecon{χρήμαθ᾽} Scaliger \espApp \lecon{ἥν χρῆν} R : ᾗ χρῆν Scaliger \espApp χρῆν ἦν R : \lecon{χρῆν εἰ} Brunck \espApp δανείσειε R : \lecon{δανείσειέν} Biset}
\ligneRav{}{}{καὶ τόκον πράττοιτο, διδόναι μηδέν᾽ ἀνθρώπων τόκον,}{}
\ligneRav{844}{}{ἀφαιρεῖσθαι τ\auDessus{ὰ}{ο}δὶ τὰ χρήματα,}{\textbf{844} ἀφερεῖσθαι R : \lecon{ἀφαιρεῖσθαι} \Mon \espApp χρήματατ᾽ R : \lecon{χρήματ᾽} Grynaeus \espApp τ\auDessus{ὰ}{ο}δὶ \Rpc : \lecon{τοδί} \Mon}
\ligneRav{845}{}{`ἀξία γοῦν εἶ τόκου τεκοῦσα τοιοῦτον τόκον.'}{}
\ligneRav{}{μ\auDessus{ν}{Η}}{ἰλλὸς γεγένημαι προσδοκῶν· ὁ δ᾽ οὐδέπω·}{}
\ligneRav{}{}{τί δῆτ᾽ ἂν εἴη τοὐμποδών; οὐκ ἔσθ᾽ ὅπως}{}
\ligneRav{848}{}{οὐ τὸν \difference{Παλαμήδη} ψυχρὸν ὄντ᾽ αἰσχύνεται·}{\textbf{848} Παλαμήδην R : \lecon{Παλαμήδη} Elmsley (\italique{Quart. Rev.} 7 (1812) p.454)}
\ligneRav{}{}{τῷ δῆτ᾽ ἂν αὐτὸν προσαγαγοίμην δράματι;}{}
\ligneRav{850}{}{ἐγᾦδα· τὴν καινὴν Ἑλένην μιμήσομαι·}{}
\ligneRav{851}{}{\difference{πάντως δ᾽} ὑπάρχει μοι γυναικεία στολή·}{\textbf{851} πάντως R : \lecon{πάντως δ᾽} Bentley, cf. v.1012}
\ligneRav{852}{\auDessus{γυ}{ν}}{τί αὖ σὺ κυρκανᾷς; \difference{τί κοικύλλεις} ἔχων;}{\textbf{852} ἢ τί κοικύλλεις R, Suid. : \lecon{τί κοικύλλεις} Faber}
\ligneRav{853}{}{πικρὰν \difference{ἑλένην} ὄψει τάχ᾽, εἰ μὴ κοσμίως}{\textbf{853} Ἑλένη R : \lecon{ἑλένην} Bibi ap. Hesych. ε 1995}
\ligneRav{}{}{ἕξεις, ἕως ἂν τῶν πρυτάνεών τις φανῇ·}{}
\ligneRav{855}{μ\auDessus{ν}{Η}}{Νείλου μὲν αἵδε καλλιπάρθενοι ῥοαί,}{}
\ligneRav{}{}{ὃς ἀντὶ δίας ψακάδος Αἰγύπτου πέδον}{}
\ligneRav{}{}{λευκῆς νοτίζει μελανοσυρμαῖον λεών·}{}
\ligneRav{}{\auDessus{γυ}{ν}}{πανοῦργος εἶ νὴ τὴν Ἑκάτην τὴν φωσφόρον·}{}
\ligneRav{}{μ\auDessus{ν}{Η}}{ἐμοὶ δὲ γῆ μὲν πατρὶς οὐκ ἀνώνυμος}{}
\ligneRav{860}{—}{ὅστις γε τολμᾷς σῆμα \difference{σοί} (σ\auDessus{ί}{ύ} \Rdeux) κλέπτειν γυνή·}{\textbf{860} σύ (ex σί) \Rdeux : \lecon{σοί} Zanetti (\italique{tibi} Divus) \espApp τ' R : \lecon{γ᾽} Gelenius}
\ligneRav{861}{}{πατὴρ ἐκεῖνός \difference{ἐστι}; Φρυνώνδας μὲν οὖν·}{\textbf{861} ἐστιν R : \lecon{ἐστι} \Mon}
\ligneRav{862}{μ\auDessus{ν}{Η}}{Ἑλένη δ᾽ ἐκλήθην. : αὖθις αὖ \difference{γίγνει} γυνή,}{\textbf{862} γίγνη R : \lecon{γίγνει} Brunck}
\ligneRav{863}{}{πρὶν τῆς ἑτέρας δοῦναι γυναικίσεως δίκην;}{\textbf{863} \lecon{δοῦναι post ἑτέρας} R : post δίκην Suid. γ 496}
\ligneRav{864}{μ\auDessus{ν}{Η}}{ψυχαὶ δὲ πολλαὶ δι᾽ ἔμ᾽ ἐπὶ \difference{Σκαμανδρίοις}}{\textbf{864} Σκαμανδρίαις R : \lecon{Σκαμανδρίοις} Fritzsche ex Eur.\italique{Hel} 52 et 609}
\ligneRav{865}{—}{ῥοαῖσιν ἔθανον. : ὤφελες δὲ καὶ σύ γε·}{}
\ligneRav{}{μ\auDessus{ν}{Η}}{κἀγὼ μὲν ἐνθάδ᾽ εἴμ᾽· ὁ δ᾽ ἄθλιος πόσις}{}
\ligneRav{867}{}{οὑμὸς \difference{Μενέλεως} οὐδέπω προσέρχεται·}{\textbf{867} Μενέλαος R : \lecon{Μενέλεως} Brunck, cf. v.901}
\ligneRav{868}{—}{τί οὖν ἔτι ζῶ; : τῶν κοράκων πονηρίᾳ·}{\textbf{868} τῶν κορ. πον. continuat (Μνη.) R : \lecon{(Γυνή)} Σ, Ellebodius}
\ligneRav{}{μ\auDessus{ν}{Η}}{ἀλλ᾽ ὥσπερ αἰκάλλει τι καρδίαν ἐμήν·}{}
\ligneRav{870}{}{μὴ ψεῦσον ὦ Ζεῦ τῆς ἐπιούσης ἐλπίδος·}{\textbf{870} \lecon{ἐλπίδος} R : ἡμέρας ap. Greg. Cor. p.14-15, Σ ad Eur. \italique{Hec}.225 \espApp post v. 870 < οἶσθ᾽ οὖν ὃ δρᾶσον; $ - \cup - | \times - \cup - $ > \lecon{om.} R : restituit ap. Σ Bibi (cf. Eur. \italique{Hel.}325)}
\ligneRav{871}{\auDessus{εὐ}{ρ}}{τίς τῶνδ᾽ ἐρυμνῶν δωμάτων ἔχει κράτος,}{\textbf{871} ὡς Περσεύς in mg. R : \lecon{ὡς Μενέλαος} Zanetti}
\ligneRav{872}{}{ὅστις \difference{ξένους} δέξαιτο ποντίῳ σάλῳ}{\textbf{872} ξένος R : \lecon{ξένους} Zanetti}
\ligneRav{873}{}{κάμνοντας ἐν χειμῶνι καὶ ναυαγίαις;}{\textbf{873} \lecon{κάμνοντας} R : καμόντας Lenting}
\ligneRav{}{μ\auDessus{ν}{Η}}{Πρωτέως τάδ᾽ ἐστὶ μέλαθρα. : ποίου Πρωτέως;}{}
\ligneRav{875}{}{\difference{ὤ τρισκακοδαίμων}, ψεύδεται νὴ τὼ θεώ,}{\textbf{875} ὦ τρισκακόδαιμον R : \lecon{ὤ τρισκακοδαίμων} Wilamowitz ad \italique{Lys.}449}
}
\vspace{3pt}\noindent\rule{4.5cm}{0.4pt}\par\vspace{1.5pt}
{\fontsize{7.0pt}{8.3pt}\selectfont
\noindent\textbf{Loci paralleli (vv. 838–875) ---} \textbf{838} (σκάφιον ἀποκεκαρμένην) Photius σ 291 σκάφιον ἀποκεκαρμένην ; cf. Hesychius σ 879 σκάφιον ; Photius σ 292 ($\simeq$ Scholium) ; \Eustathe 907.40, 1292.61 ; \AelDion σ 23 \espTest \textbf{846} Suidas ι 323 $\sim$ Zonaras 1103 ἰλλός \espTest (ἰλλός) Pollux II 52 ; \Eustathe 768.9 ; \AelDion ι 6 ; cf. Moeris ι 21 (ἰλλόν $=$ στραβόν) ; Hesychius ι 578 ἰλλός ; Hesychius η 257 (ἠίλλοντο $=$ ἐστράβιζον) \espTest \textbf{852} Suidas κ 2534 \espTest (κοικύλλεις) Hesychius κ 3232 κοικύλλειν $\sim$ Photius κ 860 \espTest \textbf{861} (Φρυνώνδας) Suidas φ 700 ; \Syn φ 211 Φρυνώνδας $\sim$ Hesychius φ 938 \espTest \textbf{863} Suidas γ 496 γυναίκισις et Zonaras 459 πρὶν τῆς ἑτέρας γυναικίσεως δίκην δοῦναι \espTest \textbf{868} Suidas τ 648 \espTest \textbf{869} (αἰκάλλει τι καρδίαν ἐμήν) Suidas αι 161 ; Zonaras 91 (τι omisso) \espTest \textbf{870} \GregCor p.14-15 ($=$ scholium\expo{Mosch.} Euripides, \italique{Hecub.} 225) (sed ἡμέρας pro ἐλπίδος) ; Suidas ψ 62 \espTest (μὴ ψεῦσον) \Antiatt 107.31 \espTest \textbf{870} \GregCor pp. 14-15 ($=$ scholium\expo{Mosch.} Euripides, \italique{Hecub.} 225) (sed ἡμέρας pro ἐλπίδος) ; Suidas ψ 62\par
}
\vspace{2.5pt}\noindent\rule{4.5cm}{0.4pt}\par\vspace{1.5pt}
{\fontsize{6.9pt}{8.2pt}\selectfont
\noindent\textbf{Suite (en commentaire) de l'apparat critique \& notes paléographiques (\texttt{f. 172v}) ---} \textbf{Corrections supralinéaires de $\text{R}^{\text{pc}} / \text{R}^2$ :} Au \textbf{v. 838}, \lecon{ὑστέρ\auDessus{ο}{α}ν} (\Rpc{} trace $\alpha$ au-dessus de $\text{ο}$ ; $\Sigma^{\text{R}}$ lisait ὑστάτην) ; au \textbf{v. 844}, \lecon{τ\auDessus{ὰ}{ο}δὶ} (\Rpc{}, τοδί \Mon{}) ; au \textbf{v. 860}, \lecon{σύ} corrigé par \Rdeux{} à partir de σί (pour σοί Zanetti). --- \textbf{Deux conjectures capitales de \textbf{Bibi} (vv. 839 et 853) :} --- \textbf{v. 838 :} \lecon{ὑστέρ\auDessus{ο}{α}ν} \Rpc{} : ὑστάτην $\Sigma$. --- \textbf{v. 839 :} ὦ πόλεις R (ὦ πόλις Gelenius) : \lecon{ωὖπολι} \lecon{Bibi}. --- \textbf{v. 853 :} Ἑλένη R : \lecon{ἑλένην} \lecon{Bibi} ap. Hésychius $\varepsilon\ 1995$. --- \textbf{v. 868 :} τῶν κορ. πον. continuat (Μνη.) R : \lecon{(Γυνή)} $\Sigma$, Ellebodius.\par
}
\end{minipage}
\end{adjustbox}
\end{adjustwidth}
\newpage
\enTeteD{Codex Ravennas 429 (Biblioteca Classense) --- Fac-similé du f. 172v (vv. 838–875)}
\begin{adjustwidth}{-1.5cm}{-3.0cm}
\begin{center}
\includegraphics[width=16.5cm,height=20.5cm,keepaspectratio]{/workspaces/Thesmophories/dossier_ravennas_thesmophories/folios_individuels_ff162v-177r/21_f172v_classense174v_vv0838-0875.jpg}\par
\vspace{2pt}
{\scriptsize\textbf{Codex Ravennas 429, f. 172v (Classense 174v)} --- ΠΑΡΑΒΑΣΙΣ / ΕΠΕΙΣΟΔΙΟΝ Δ', \textbf{vv. 838–875} ($\text{R}^1, \Rpc, \Rdeux$, \PsiUn).}
\end{center}
\end{adjustwidth}
\newpage
\clearpage\ifodd\value{page}\hbox{}\newpage\fi% --- FOLIO 21 (f. 172v) : SPREAD 2.1 (Texte grec établi + Traduction française + Commentaire 1/8 & 2/8) ---
\clearpage\ifodd\value{page}\hbox{}\newpage\fi
\enTeteG{ΘΕΣΜΟΦΟΡΙΑΖΟΥΣΑΙ --- Texte grec établi (f. 172v, vv. 838–875 [1/4])}
\enTeteD{LES THESMOPHORIES --- Traduction française (f. 172v, vv. 838–875 [1/4])}
\begin{adjustwidth}{-2cm}{-2cm}
\noindent\textbf{\footnotesize Texte grec établi --- Folio 21 (\texttt{f. 172v}, vv. 838–875 [1/4]) \hfill ΠΑΡΑΒΑΣΙΣ / ΕΠΕΙΣΟΔΙΟΝ Δ'}\\[-2pt]
\noindent\rule{16.5cm}{0.5pt}\\[2pt]
{\fontsize{8.8pt}{10.2pt}\selectfont
\ligneEst{}{\Automatique{ὑστέραν αὐτὴν καθῆσθαι σκάφιον ἀποκεκαρμένην}}{838}
\ligneEst{}{τῆς τὸν ἀνδρεῖον τεκούσης· τῷ γὰρ εἰκὸς ωὖπολι}{}
\ligneEst{}{τὴν Ὑπερβόλου καθῆσθαι μητέρ᾽ ἠμφιεσμένην}{840}
\ligneEst{}{\Automatique{λευκὰ καὶ κόμας καθεῖσαν πλησίον τῆς Λαμάχου·}}{}
\ligneEst{}{\Automatique{καὶ δανείζειν χρήμαθ᾽, ἥν χρῆν, εἰ δανείσειέν τινι}}{}
\ligneEst{}{\Automatique{καὶ τόκον πράττοιτο, διδόναι μηδέν᾽ ἀνθρώπων τόκον,}}{}
\ligneEst{}{ἀλλ᾽ ἀφαιρεῖσθαι βίᾳ τὰ χρήματ᾽ εἰπόντας τοδί,}{}
\ligneEst{}{`ἀξία γοῦν εἶ τόκου τεκοῦσα τοιοῦτον τόκον.'}{845}
\ligneEst{\textbf{Μνησίλοχος} {\tiny\textit{trim. iamb.}}}{ἰλλὸς γεγένημαι προσδοκῶν· ὁ δ᾽ οὐδέπω.}{}
\ligneEst{}{τί δῆτ᾽ ἂν εἴη τοὐμποδών; οὐκ ἔσθ᾽ ὅπως}{}
}
\vspace{2pt}
\noindent\rule{16.5cm}{0.6pt}\\[2pt]
\noindent\textbf{\small Commentaire philologique, métrique et scénique (vv. 838–875 --- partie 1/8)}\\[2pt]
\begin{adjustbox}{max totalheight=15.2cm, max width=16.5cm, keepaspectratio}%
\parbox[t]{16.5cm}{
\setlength{\columnsep}{9pt}%
\begin{multicols}{2}
\fontsize{7.0pt}{8.0pt}\selectfont
\setlength{\parindent}{0pt}%
\setlength{\parskip}{1.2pt}%
\noindent\textbf{838} \textbf{σκάφιον ἀποκεκαρμένην.} C'est l'humiliation que le tyran de Cumes, Aristodème, avait fait subir aux filles de son pays : περιτρόχαλα κείρεσθαι, tandis qu'il forçait les garçons à porter des cheveux longs. Plutarque, \italique{De mul. virtut}, 26, p. 361f. (Note de Willems.)\\ La scholie précise à propos de σκάφιον :
\citComm{εἶδος κουρᾶς δουλικῆς.\\
\francais{Genre de coupe rase qui convient à l'esclave.}}
\noindent\textbf{839} \textbf{ὦ πόλεις.} Dans son livre \italique{Aristophanes und Eupolis} (2007), Natalia Kyriakidi a montré qu'Eupolis, le poète comique rival d'Aristophane, est selon toute vraisemblance mort en 411 lors de la victoire navale de Κυνὸς Σῆμα (cf. comm. vv. 833-7). Pour arriver à cette conclusion, elle récuse très justement la plupart des témoignages, qui exhalent l'anecdote recyclée, pour ne conserver que le témoignage de Suidas (ε 3657) -- ou de la \italique{Souda}, si l'on préfère :
\citComm{Εὔπολις, Σωσίπολις, Ἀθηναῖος, κωμικός. ὃς ἑπτακαίδεκα ἔτη γεγονὼς ἤρξατο ἐπιδείκνυσθαι: καὶ ἐδίδαξε μὲν δράματα ιζ, ἐνίκησε δὲ ζ. καὶ ἀπέθανε ναυαγήσας κατὰ τὸν Ἑλλήσποντον ἐν τῷ πρὸς Λακεδαιμονίους πολέμῳ: καὶ ἐκ τούτου ἐκωλύθη στρατεύεσθαι ποιητήν. 
\\
Fils de Sosipolis, Athénien, un poète comique. Il a commencé à monter des pièces de théâtre à l'âge de dix-sept ans. Il a produit 17 pièces et en a gagné 7. Il est mort dans le naufrage de son navire dans l'Hellespont pendant la guerre contre les Lacédémoniens ; à la suite de cela, il fut interdit à un poète de servir en campagne. 
}
\noindent Natalia continue (p. 11) : \allemand{Die zwei großen Schlachten, die am Hellespont während des Peloponnesischen Krieges stattfanden, sind der Sieg der Athener bei} Κυνὸς Σῆμα \allemand{411 und die endgültige Niederlage bei }Αἰγὸς Ποταμοί \allemand{405. Der Tod des Eupolis während einer Seeschlacht könnte auch die Entstehung der Geschichte vom Ertränken des Eupolis durch Alkibiades erklären} [$=$ Eupolis, selon certaines légendes, aurait été noyé en mer par Alcibiade au retour de Sicile], \allemand{weil beide Geschichten als Vorfälle im Meer bei einer Kriegshandlung wenige Jahre vor dem Ende des Peloponnesischen Krieges überliefert sind. Storey vermutet, dass Ailians aitiologische Erklärung des Ortsnamens} Κυνὸς Θρῆνος [$=$ le lieu commun du chien qui meurt de chagrin sur la tombe de son maître] \allemand{auf }Κυνὸς Σῆμα \allemand{zurückzuführen ist. Da kein Werk des Eupolis nach 411 zu datieren ist und für die Zeit danach keine Indizien für eine weitere Tätigkeit des Dichters überliefert sind, erscheint die Information der Suda als das einzige zuverlässige Zeugnis über den Tod des Eupolis. So enden sein Leben und seine dichterische Karriere wahrscheinlich im Sommer 411 während der Seeschlacht bei} Κυνὸς Σῆμα.\\ Le texte du manuscrit de \italique{Ravenne} comporte au vers 839 une adresse à la Cité :
\citComm{τῷ γὰρ εἰκὸς ὦ πόλεις (sic.)}
\noindent La vraisemblable quasi-coïncidence des dates entre la mort d'Eupolis et la représentation des \italique{Thesmophories}\marqueur{1} m'incite à penser qu'il y a là plutôt une évocation d'Eupolis, auteur intarissable sur la mère d'Hyperbolos, et à laquelle il a même consacré une pièce entière : son \italique{Maricas} (voir le commentaire du vers 840).
\noindent S'il faut attribuer l'idée d'accorder les sièges des Athéniennes avec la valeur de leur progéniture à feu Eupolis (et peut-être aussi celle de les tondre le cas décevant échéant), il n'aura pas manqué de s'indigner lui aussi que la mère d'Hyperbolos spécule avec le fruit de la rapine de son fils.
\noindent On ne possède malheureusement que peu de choses d'Eupolis.
\noindent On sait que le \italique{Maricas} dut être représenté en $-421$, que le chœur était formé de deux demi-chœurs : celui des riches et celui des pauvres citoyens , chœurs qui débattaient avec le démagogue Hyperbolos.
\noindent La comédie est une violente charge contre ce dernier.
\noindent Elle se terminait par une danse licencieuse de la mère du démagogue.
\noindent Nombre de commentateurs, à l'instar des scholiastes, ont rapproché les \italique{Cavaliers} du \italique{Maricas}... mais pas les \italique{Thesmophories} -- plus pauvres en scholies anciennes.
\end{multicols}
}%
\end{adjustbox}%
\vfill
\end{adjustwidth}
\newpage
% --- PAGE DE DROITE : Traduction française + Commentaire (2/8) ---
\begin{adjustwidth}{-2cm}{-2cm}
\noindent\textbf{\footnotesize Traduction française --- Folio 21 (\texttt{f. 172v}, vv. 838–875 [1/4]) \hfill ΠΑΡΑΒΑΣΙΣ / ΕΠΕΙΣΟΔΙΟΝ Δ'}\\[-2pt]
\noindent\rule{16.5cm}{0.5pt}\\[2pt]
{\fontsize{8.8pt}{10.2pt}\selectfont
\ligneEst{}{\francais{s'assiérait, les cheveux taillés au bol, derrière}}{838}
\ligneEst{}{\francais{la femme qui enfanta le brave. Comment tolérer, ô Eupolis,}}{}
\ligneEst{}{\francais{que la mère d'Hyperbolos siège, vêtue}}{840}
\ligneEst{}{\francais{de blanc, les cheveux flottants, à côté de la mère de Lamachos,}}{}
\ligneEst{}{\francais{et qu'elle prête de l'argent à usure... le faut-il supporter ? Si elle prête de l'argent}}{}
\ligneEst{}{\francais{et exige un profit, loin de lui payer un intérêt,}}{}
\ligneEst{}{\francais{il faudrait lui ôter de force ses richesses, en lui disant :}}{}
\ligneEst{}{\francais{« Tu vaux bien que ton argent produise, toi qui produisis un tel fruit. »}}{845}
\ligneEst{\textbf{Mnésiloque}}{\francais{Je suis devenu louche à lorgner de ce côté. Et lui ?... pas encore.}}{}
\ligneEst{}{\francais{Qu'est-ce donc qui l'empêche ? Nul doute,}}{}
}
\vspace{2pt}
\noindent\rule{16.5cm}{0.6pt}\\[2pt]
\noindent\textbf{\small Commentaire philologique, métrique et scénique (vv. 838–875 --- partie 2/8)}\\[2pt]
\begin{adjustbox}{max totalheight=15.2cm, max width=16.5cm, keepaspectratio}%
\parbox[t]{16.5cm}{
\setlength{\columnsep}{9pt}%
\begin{multicols}{2}
\fontsize{7.0pt}{8.0pt}\selectfont
\setlength{\parindent}{0pt}%
\setlength{\parskip}{1.2pt}%
\noindent Pourtant, certaines idées d'Eupolis semblent en être reprises. Nous l'allons voir tout-à-l'heure.
\noindent Mettons en garde tout-de-suite : les hypothèses -- comme les explications -- seront fort spéculatives.\\ Un commentaire du \italique{Maricas}, daté du \nieme{II} siècle ap. J.-C., nous est parvenu de façon très morcelaire. On y lit cette citation (fr. 192.96-99 Kassel-Austin) :
\citComm{τοῦτ᾽ ἐκδανείζει καὶ κυκᾷς [τὸ]\douteux{ν} ναυτικὸν \douteux{ἄ}[ργυρον \ \ \ \ \ (\textbf{ἄργυρον} Bibi)\\
ἐπὶ τῷ πέμπτῳ μέρε[ι \ \ \ \ \ \ \ ]τοὺς τόκ[ους οἱ\\
ναυτικοί, ἡμεῖς δ᾽ ἄρ᾽ οἴκ[ \ \ \ \ \ \ \ \ \ \ \ ]εν· ὁ τῶν\\
πλουσίων λέγει χορός.}
\noindent Le chœur des riches (ὁ τῶν πλουσίων χορός) reproche à Maricas ($=$ Hyperbolos) de prêter (ἐκδανείζει) en concoctant (κυκᾷς) des prêts (τὸν ναυτικὸν ἄργυρον) au cinquième (ἐπὶ τῷ πέμπτῳ μέρει\marqueur{2}), c'est-à-dire avec un taux d'intérêt de $20 \%$.
\noindent Il est questions ensuite des τόκους, `enfants' ($=$ intérêts) de ces ναυτικοί (cf. vers 845).
\noindent Continuant le jeu de mot nautique (les prêts sont littéralement de l'argent \italique{marin}), est-ce à ce moment qu'Eupolis, par la voix du chœur, proposait de faire tondre \italique{nautiquement} cet usurier ou sa mère et de lui faire porter une \italique{barque} (σκάφιον vers 838) ?\\ Revenons aux \italique{Thesmophories} : en guise d'oraison funèbre, Aristophane rend hommage à la pertinence des propositions d'Eupolis :
\citComm{τῷ γὰρ εἰκὸς\marqueur{3} ωὖπολι\\
Comment tolérer\marqueur{3}, ô Eupolis, que \&c.}
\noindent Dernier coup de chapeau d'Aristophane à son défunt rival, mort des suites de la folie des Athéniens et de l'incompétence des amiraux.
\noindent\textbf{[Notes ad 839]} \textbf{(1)} Voir l'introduction et la question de la date. Avec Rogers, je penche pour l'hypothèse que les \italique{Thesmophories} furent représentées aux Dionysies de 410. Si nous avons vu juste avec Natalia, Eupolis serait donc mort l'été précédent la première représentation de la pièce.
\noindent \textbf{(2)} Sens bien attesté d'ἐπί. Cf. Dém. 1250, 20 : ἐπ᾽ ὀκτὼ ὀβολοῖς δανείζειν τοῦ μηνὸς τὴν μνᾶν, «~prêter de l'argent moyennant un intérêt de huit oboles la mine par mois~», c'est-à-dire à $16 \%$.
\noindent \textbf{(3)} Le tour elliptique apparaît aussi au vers 703 des \italique{Acharniens} où il est glosé `πῶς δίκαιόν ἐστιν' par une scholie. Le ton est effectivement indigné : «~comment peut-il être juste/raisonnable que...~», c'est-à-dire «~quelle indignité que...~». Voir aussi Sophocle (\italique{Ajax}, vers 1366-1367), non indigné mais ironique :
\citComm{Ἀγαμέμνων -- ἡ πάνθ᾽ ὅμοια πᾶς ἀνὴρ αὑτῷ πονεῖ.\\
Ὀδυσσεύς -- τῷ γάρ με μᾶλλον εἰκὸς ἢ ᾽μαυτῷ πονεῖν;\\
Agamemnon -- En toute chose, chacun ne travaille que pour lui-même.\\
Ulysse -- Comment peut-il être plus raisonnable que de travailler pour moi-même ?}
\noindent\textbf{840-1} \textbf{Ὑπερβόλου ... Λαμάχου,} Hyperbolos, le démagogue athénien -- marchand de lampes de son état\marqueur{1} -- devint démagogue en chef après la mort de Cléon devant Amphipolis en 422 av. J.-C.\marqueur{2} ; sa mère et lui furent lors une cible privilégiée d'Aristophane et des poètes comiques\marqueur{3} (suite sans doute à quelque lapsus à relent freudien du démagogue, lapsus qui fut repris à l'encan par les Comiques). Cet ὀξίνην Ὑπερβόλον\marqueur{4} réunissait tout pour être détesté d'Aristophane : de caractère exécrable, il était un partisan, comme son mentor Cléon, de la guerre à outrance\marqueur{5} ; c'était un intenteur de procès\marqueur{6}, un fourbe\marqueur{7}, qui s'est enrichi par ses coquineries\marqueur{8}. Il eut l'insigne honneur d'être le dernier des Athéniens ostracisés en 417\marqueur{9}. Il fut assassiné à Samos, son lieu d'exil\marqueur{10}, et son corps jeté en mer\marqueur{11}... cette \italique{inhudation} lui valut une dernière raillerie de la part d'Aristophane\marqueur{12}.
\end{multicols}
}%
\end{adjustbox}%
\vfill
\end{adjustwidth}
\newpage

% --- FOLIO 21 (f. 172v) : SPREAD 2.2 (Texte grec établi + Traduction française + Commentaire 3/8 & 4/8) ---
\clearpage\ifodd\value{page}\hbox{}\newpage\fi
\enTeteG{ΘΕΣΜΟΦΟΡΙΑΖΟΥΣΑΙ --- Texte grec établi (f. 172v, vv. 838–875 [2/4])}
\enTeteD{LES THESMOPHORIES --- Traduction française (f. 172v, vv. 838–875 [2/4])}
\begin{adjustwidth}{-2cm}{-2cm}
\noindent\textbf{\footnotesize Texte grec établi --- Folio 21 (\texttt{f. 172v}, vv. 838–875 [2/4]) \hfill ΠΑΡΑΒΑΣΙΣ / ΕΠΕΙΣΟΔΙΟΝ Δ'}\\[-2pt]
\noindent\rule{16.5cm}{0.5pt}\\[2pt]
{\fontsize{8.8pt}{10.2pt}\selectfont
\ligneEst{}{οὐ τὸν Παλαμήδη ψυχρὸν ὄντ᾽ αἰσχύνεται.}{}
\ligneEst{}{τῷ δῆτ᾽ ἂν αὐτὸν προσαγαγοίμην δράματι;}{}
\ligneEst{}{ἐγᾦδα· τὴν καινὴν Ἑλένην μιμήσομαι.}{850}
\ligneEst{}{πάντως δ᾽ ὑπάρχει μοι γυναικεία στολή.}{}
\ligneEst{\textbf{Κρίτυλλα}}{τί αὖ σὺ κυρκανᾷς; τί κοικύλλεις ἔχων;}{}
\ligneEst{}{πικρὰν ἑλένην ὄψει τάχ᾽, εἰ μὴ κοσμίως}{}
\ligneEst{}{ἕξεις, ἕως ἂν τῶν πρυτάνεών τις φανῇ.}{}
\ligneEst{\textbf{Μνησίλοχος}}{Νείλου μὲν αἵδε καλλιπάρθενοι ῥοαί,}{855}
\ligneEst{}{ὃς ἀντὶ δίας ψακάδος Αἰγύπτου πέδον}{}
\ligneEst{}{λευκῆς νοτίζει μελανοσυρμαῖον λεών.}{}
\ligneEst{\textbf{Κρίτυλλα}}{πανοῦργος εἶ νὴ τὴν Ἑκάτην τὴν φωσφόρον.}{}
}
\vspace{2pt}
\noindent\rule{16.5cm}{0.6pt}\\[2pt]
\noindent\textbf{\small Commentaire philologique, métrique et scénique (vv. 838–875 --- partie 3/8)}\\[2pt]
\begin{adjustbox}{max totalheight=15.2cm, max width=16.5cm, keepaspectratio}%
\parbox[t]{16.5cm}{
\setlength{\columnsep}{9pt}%
\begin{multicols}{2}
\fontsize{7.0pt}{8.0pt}\selectfont
\setlength{\parindent}{0pt}%
\setlength{\parskip}{1.2pt}%
\noindent \textbf{Lamachos} était un stratège athénien\marqueur{13} et l'un des représentants athéniens de la Paix de Nicias\marqueur{14}. Avec Nicias et Alcibiabe, il est investi des pleins pouvoirs lors de l'expédition de Sicile\marqueur{15} où il trouve la mort\marqueur{16}. Lamachos présente une originalité unique parmi les personnages historiques d'Aristophane puisque le poète opère un volte-face à son endroit : du va-t-en-guerre au panache ridicule dépeint dans les \italique{Acharniens} et la \italique{Paix}\marqueur{17} (accusé même de s'enrichir grâce à la guerre\marqueur{18}), Aristophane lui rend hommage ici comme dans les \italique{Grenouilles}\marqueur{19}. La mort héroïque de Lamachos\marqueur{20} est sans doute pour beaucoup dans ce changement radical d'opinion.
\noindent\textbf{845} \textbf{ἀξία γοῦν εἶ τόκου τεκοῦσα τοιοῦτον τόκον.} Jeu de mots qui roule sur le double sens de τόκος, \latin{usura} et \latin{proles}. C'est ainsi que l'orateur Eschine disait de Démosthène qui, souffleté par Midias, s'était désisté moyennant 30 mines :
\citComm{ὁ γὰρ ἄνθρωπος οὐ κεφαλήν, ἀλλὰ πρόσοδον κέκτηται.\\
Ce n'est pas une tête (κεφαλήν), mais un capital que possède cet homme.}
\noindent (Cf. \italique{Contre Ctésiphon}, 212 ; note de Willems.)
\noindent\textbf{[Notes ad 840-1]} \textbf{(1)} \italique{Cavaliers}, v. 739 (non nommé) ; \italique{Nuées}, v. 1064 ; \italique{Paix}, v. 690.
\noindent \textbf{(2)} Voir la \italique{Paix}, v. 679-692, Paix qui fout le camp, révoltée que le peuple se soit donné un si indigne patron.
\noindent \textbf{(3)} Cf. les \italique{Nuées}, v. 551-562 : Aristophane daube sur ses plagiaires de contemporains, notamment Eupolis réchauffant leurs \italique{Cavaliers} pour son \italique{Maricas} et y rajoutant un autre plat déjà concocté par Phrynicos en son temps : une mère d'Hyperbolos, ivrognesse et dansant le licencieux cordax. Sur les attaques des comiques contre la mère d'Hyperbolos, présentée comme une débauchée notoire, voir aussi les fragments 7-12 des \italique{Vendeuses de pain} (αἱ Ἀρτοπώλιδες) d'Hermippos (dans l'édition Kassel-Austin) ainsi que les fragments de Platon le Comique dans sa pièce éponyme \italique{Hyperbolos}.
\noindent \textbf{(4)} \italique{Cavaliers}, v. 1303. Thucydide, qui n'estimait guère plus ce μοχθὴρον ἄνθρωπον, ne parle qu'une seule fois d'Hyperbolos : pour nous annoncer sa mort et en brosser un portrait expéditif (VIII, 73.3) :
\citComm{καὶ Ὑπέρβολόν τέ τινα τῶν Ἀθηναίων, μοχθηρὸν ἄνθρωπον, ὠστρακισμένον οὐ διὰ δυνάμεως καὶ ἀξιώματος φόβον, ἀλλὰ διὰ πονηρίαν καὶ αἰσχύνην τῆς πόλεως, ἀποκτείνουσι μετὰ Χαρμίνου\\
S'en prenant à Hyperbolos, un Athénien méprisable qu'on avait ostracisé non par peur de son influence et de son prestige, mais parce que c'était un malhonnête homme qui déshonorait la cité, ils [les conjurés favorables à l'oligarchie] l'assassinèrent avec l'aide de Charminos.}
\noindent \textbf{(5)} \italique{Cavaliers}, v. 1303-1315 ; il cherchait à renforcer la flotte pour que l'expédition de Sicile fût poussée jusqu'à Carthage.
\noindent \textbf{(6)} \italique{Acharniens}, v. 846-847.
\noindent \textbf{(7)} \italique{Guêpes}, v. 1007.
\noindent \textbf{(8)} \italique{Nuées}, v. 1064-1066. Tel fils, telle mère ? L'usure brocardée par Aristophane dont fait preuve la mère d'Hyperbolos n'était-elle pas la continuation des activités louches du fils, depuis que ce dernier se trouvait en exil à Samos ?
\noindent \textbf{(9)} Plut., \italique{Nicias} 11 ; \italique{Alcib.} 13.
\noindent \textbf{(10)} Thucydide, VIII, 73.3.
\noindent \textbf{(11)} Théopompe 115F69, in \italique{Die Fragmente der griechischen Historiker}, Teil 2, \italique{Zeitgeschichte - B}.
\noindent \textbf{(12)} \italique{Grenouilles}, v. 570.
\noindent \textbf{(13)} Thucydide IV, 75. 1 : il commande une dizaine de navires dans le Pont.
\noindent \textbf{(14)} Thucydide V 19.2 et 24.1
\noindent \textbf{(15)} Thucydide (VI, 8) ; il s'opposa à Alcibiade sur la stratégie (V 49) : Lamachos voulait aller droit au but et attaquer Syracuse sitôt le débarquement fait ; il finit par se ranger à l'opinion d'Alcibiade (VI, 50).
\noindent \textbf{(16)} Thucydide (VI, 101) : Lamachos courant à la rescousse de l'aile droite de l'armée athénienne se retrouva isolé avec quelques-uns de ses camarades ; il fut tué par les Syracusains qui récupérèrent son corps avant de le restituer (VI, 103) aux Athéniens. Nicias fut dès lors le seul commandant en Sicile.
\end{multicols}
}%
\end{adjustbox}%
\vfill
\end{adjustwidth}
\newpage
% --- PAGE DE DROITE : Traduction française + Commentaire (4/8) ---
\begin{adjustwidth}{-2cm}{-2cm}
\noindent\textbf{\footnotesize Traduction française --- Folio 21 (\texttt{f. 172v}, vv. 838–875 [2/4]) \hfill ΠΑΡΑΒΑΣΙΣ / ΕΠΕΙΣΟΔΙΟΝ Δ'}\\[-2pt]
\noindent\rule{16.5cm}{0.5pt}\\[2pt]
{\fontsize{8.8pt}{10.2pt}\selectfont
\ligneEst{}{\francais{le \italique{Palamède} est froid et il en a honte.}}{}
\ligneEst{}{\francais{Par quel drame alors l'attirer ...}}{}
\ligneEst{}{\francais{J'y suis : je vais contrefaire sa récente \italique{Hélène}}}{850}
\ligneEst{}{\francais{Aussi bien, je suis accoutré d'un costume de femme.}}{}
\ligneEst{\textbf{Kritylla}}{\francais{Que manigances-tu encore ? Qu'as-tu à écarquiller les yeux ?}}{}
\ligneEst{}{\francais{On t'en fera voir une Hélène, amère celle-là, si tu ne te tiens décemment}}{}
\ligneEst{}{\francais{jusqu'à la venue de l'un des prytanes.}}{}
\ligneEst{\textbf{Mnésiloque}}{\francais{(Rôle d'Helène.) « Du Nil ce sont ici les rives aux belles Nymphes,}}{855}
\ligneEst{}{\francais{Nil qui à la place de l'eau du ciel, humecte le fonde-ment d'Égypte}}{}
\ligneEst{}{\francais{la blanche » et du peuple à la noire syrméa.}}{}
\ligneEst{\textbf{Kritylla}}{\francais{Tu es un effronté, par Hécate la porte-torche. (Elle lui porte un coup de sa torche.)}}{}
}
\vspace{2pt}
\noindent\rule{16.5cm}{0.6pt}\\[2pt]
\noindent\textbf{\small Commentaire philologique, métrique et scénique (vv. 838–875 --- partie 4/8)}\\[2pt]
\begin{adjustbox}{max totalheight=15.2cm, max width=16.5cm, keepaspectratio}%
\parbox[t]{16.5cm}{
\setlength{\columnsep}{9pt}%
\begin{multicols}{2}
\fontsize{7.0pt}{8.0pt}\selectfont
\setlength{\parindent}{0pt}%
\setlength{\parskip}{1.2pt}%
\noindent \textbf{(17)} Lamachos arborait sur son casque trois aigrettes qui lui donnaient des faux airs de Gorgone (\italique{Acharniens}, v. 574 ; v. 964-965). Aristophane rapproche cet accoutrement d'épouvante à d'autre monstre-croquemitaine au nom voisin (Μορμώ, \italique{Acharniens} v. 582 ; \italique{Paix} v. 474). Enfin, cette crinière en bataille -- de circonstance pour un militaire ! -- permettait un jeu analogue en grec entre λόφοι et λόχοι, trop tentant pour qu'Aristophane n'y cèdât pas (\italique{Acharniens}, v. 567, 574, 1074).
\noindent \textbf{(18)} Il est traité de μισθαρχίδης au v. 597 des \italique{Acharniens} ; attaque renouvelée au vers 619.
\noindent \textbf{(19)} \italique{Grenouilles}, v. 1039.
\noindent \textbf{(20)} Voir -- outre le passage de Thucydide (VI, 101) -- Plutarque, \italique{Nicias} 18 (et aussi le \italique{Lachès} (19c) de Platon où il est question du courage de Lamachos.)
\noindent\textbf{850} \textbf{τὴν καινὴν Ἑλένην.} « La nouvelle \italique{Hélène} ». \italique{Nouvelle} parce qu'Euripide l'avait présentée il y a peu de temps, sous l'archontat de Cléocrite (printemps 412). Mais l'adjectif καινή est aussi à prendre au sens d'\italique{innovante}, \italique{imprévue}, \italique{étrange}. Voici ce que dit Victor-Henry Debidour de cette « \italique{Hélène} nouveau style »\marqueur{1} :\\ Dans la tradition légendaire, Hélène de Sparte, femme de Ménélas, était le parangon de la beauté. Mais aussi le plus scandaleux exemple d'infidélité féminine, cause, par son inconduite, des malheurs sans nombre de la guerre de Troie. Le personnage s'était stéréotypé en « femme fatale », d'autant plus haïssable qu'elle était plus prestigieuse ; et \italique{Euripide lui-même ne la traite pas autrement} -- sauf justement dans la tragédie très peu tragique qui porte son nom. Certes, elle y reste, bien-sûr, la belle Hélène, et si charmeuse que le fils du sage Protée l'aime comme un fou, et comme un sot. Mais la voici transformée en épouse sans tache, victime innocente de la longue perfidie des deux déesses, dont l'une l'a prostituée et l'autre séquestrée : elles l'ont engagée dans une rocambolesque histoire où elle mérite toute notre sympathie. Pour se faire désigner par Pâris comme « la plus belle », Cypris avait promis au prince-berger troyen de lui donner Hélène ; Héra, furieuse, contre sa rivale et son juge, fit en sorte que Pâris n'emmenât à Troie qu'un fantôme d'Hélène, une « nuée » toute semblable à la vraie et suffisamment consistante pour leurrer le ravisseur pendant les dix ans qu'il la garda, et le légitime époux après ce qu'il crut être sa reconquête. La vraie Hélène, elle, miraculeusement transportée en Égypte, y était tenue cachée durant toutes ces années, confiée à la vertueuse protection de Protée. Ainsi, c'est pour un « crime sans criminelle » que des milliers de preux se sont entre-tués et que Troie a été anéantie.\\ Telle est l'\italique{étonnante} variante que recueille Euripide sur la plus célèbre entreprise guerrière de la légende grecque. Car il ne l'a pas inventée : Hérodote en fait état bien avant lui\marqueur{2}. De ce point de départ, Euripide a tiré un canevas sur lequel il a brodé une aventure romanesque et humoristique qui souvent ferait croire qu'avant d'être parodiée elle-même par Aristophane, l'\italique{Hélène} donnait dans la parodie.
\noindent\textbf{[Notes ad 850]} \textbf{(1)} Dans l'introduction à sa traduction de l'\italique{Hélène}. Oui ! car Debidour le grand traducteur d'Aristophane, qui a su rendre un peu de la verve de notre poète dans la langue de Rabelais, Debidour ne s'est pas limité à la Comédie Ancienne ; il est aussi l'auteur d'une traduction des pièces conservées des trois grands Tragiques : \italique{Les Tragiques -- Eschyle, Sophocle, Euripide, théâtre complet}, La Pochothèque (1999). Traduction à laquelle s'ajoute l'excellent dossier et les index fort utiles de Paul Demont et Anne Lebeau. Victor-Henry Debidour vérifie ainsi le paradoxe de Socrate dont Platon n'a pas cru utile de nous donner la démonstration à la fin de son \italique{Banquet}.
\noindent \textbf{(2)} Hérodote II, 112-120. Par la suite, du reste, le thème de la réhabilitation d'Hélène fut maintes fois repris en exercice d'école : la paradoxe devient à son tour lieu commun.
\noindent\textbf{852} \textbf{τί κοικύλλεις ἔχων;} Κοικύλλειν est un verbe à redoublement expressif\marqueur{1} qui tire son origine de τὰ κύλα, \italique{les poches sous les yeux}, d'où le sens d'« écarquiller les yeux »\marqueur{2}. Voici pour le sens en lequel l'emploie Critylla.
\noindent Le verbe est particulièrement bien adapté à la situation : notre anti-héros vient de lancer `à la mer' ses `rames' dans l'espoir de rameuter Euripide ; il est là maintenant, le regard hagard, devant cette mer imaginaire, à attendre... pareil à ce personnage comique, surnommé Κοικυλίων, type de l'idiot qui comptait les vagues de la mer\marqueur{3}, le comble de l'inutilité pour un grec : κύματα μετρεῖν était même passé en proverbe (détourné par Aristophane au vers 1257 de la \italique{Paix}).
\end{multicols}
}%
\end{adjustbox}%
\vfill
\end{adjustwidth}
\newpage

% --- FOLIO 21 (f. 172v) : SPREAD 2.3 (Texte grec établi + Traduction française + Commentaire 5/8 & 6/8) ---
\clearpage\ifodd\value{page}\hbox{}\newpage\fi
\enTeteG{ΘΕΣΜΟΦΟΡΙΑΖΟΥΣΑΙ --- Texte grec établi (f. 172v, vv. 838–875 [3/4])}
\enTeteD{LES THESMOPHORIES --- Traduction française (f. 172v, vv. 838–875 [3/4])}
\begin{adjustwidth}{-2cm}{-2cm}
\noindent\textbf{\footnotesize Texte grec établi --- Folio 21 (\texttt{f. 172v}, vv. 838–875 [3/4]) \hfill ΠΑΡΑΒΑΣΙΣ / ΕΠΕΙΣΟΔΙΟΝ Δ'}\\[-2pt]
\noindent\rule{16.5cm}{0.5pt}\\[2pt]
{\fontsize{8.8pt}{10.2pt}\selectfont
\ligneEst{\textbf{Μνησίλοχος}}{ἐμοὶ δὲ γῆ μὲν πατρὶς οὐκ ἀνώνυμος}{}
\ligneEst{}{Σπάρτη, πατὴρ δὲ Τυνδάρεως.}{860}
\ligneEst{\textbf{Κρίτυλλα}}{\decalage{Σπάρτη, πατὴρ δὲ Τυνδάρεως.}σοί γ᾽ ὦλεθρε}{}
\ligneEst{}{πατὴρ ἐκεῖνός ἐστι; Φρυνώνδας μὲν οὖν.}{}
\ligneEst{\textbf{Μνησίλοχος}}{Ἑλένη δ᾽ ἐκλήθην.}{}
\ligneEst{\textbf{Κρίτυλλα}}{\decalage{Ἑλένη δ᾽ ἐκλήθην.}αὖθις αὖ γίγνει γυνή,}{}
\ligneEst{}{πρὶν τῆς ἑτέρας δοῦναι γυναικίσεως δίκην;}{}
\ligneEst{\textbf{Μνησίλοχος}}{ψυχαὶ δὲ πολλαὶ δι᾽ ἔμ᾽ ἐπὶ Σκαμανδρίοις}{}
\ligneEst{}{ῥοαῖσιν ἔθανον.}{865}
\ligneEst{\textbf{Κρίτυλλα}}{\decalage{ῥοαῖσιν ἔθανον.}ὤφελες δὲ καὶ σύ γε.}{}
\ligneEst{\textbf{Μνησίλοχος}}{κἀγὼ μὲν ἐνθάδ᾽ εἴμ᾽· ὁ δ᾽ ἄθλιος πόσις}{}
}
\vspace{2pt}
\noindent\rule{16.5cm}{0.6pt}\\[2pt]
\noindent\textbf{\small Commentaire philologique, métrique et scénique (vv. 838–875 --- partie 5/8)}\\[2pt]
\begin{adjustbox}{max totalheight=15.2cm, max width=16.5cm, keepaspectratio}%
\parbox[t]{16.5cm}{
\setlength{\columnsep}{9pt}%
\begin{multicols}{2}
\fontsize{7.0pt}{8.0pt}\selectfont
\setlength{\parindent}{0pt}%
\setlength{\parskip}{1.2pt}%
\noindent\textbf{853} \textbf{πικρὰν ἑλένη.} Le détournement du \italique{Palamède} ayant échoué (après celui du \italique{Télèphe}), Mnésiloque songe maintenant à recourir à la récente \italique{Hélène} d'Euripide. Seulement, de toutes les \italique{Thesmophories}, aucun des geôliers de Mnésiloque ne rentre jamais dans le jeu de ses parodies d'Euripide.
\noindent Quand Kritylla prononce `πικρὰν ἑλένην ὄψει τάχ᾽', elle ne songe pas un seul instant à la mythique Hélène, mise nouvellement en scène par Euripide avec réécriture du scénario. C'est Hésychius encore une fois – décidément le plus utile des lexicographes pour les \italique{Thesmophories}~! – qui nous le confirme (ε 1995)~:
\citComm{ἑλένη· λαμπάς, δετή\\
\francais{hélène~: flambeau, torche.}}
\noindent Kritylla menace Mnésiloque d'un bon coup de flambeau, flambeau qu'elle avait lors de la procession et l'arrivée au temple. (Coup de flambeau qu'elle assénera au vers 858 en évoquant Hécate \textbf{φωσ}φόρος.)
\noindent\textbf{[Notes ad 852]} \textbf{(1)} Sur le même modèle que μοιμύλλειν, δενδίλλειν, δανδαίνειν, παπταίνειν, \&c.
\noindent \textbf{(2)} Cf. Hésychius κ 3232 : κοικύλλειν· περιβλέπειν· κύλα γὰρ τὰ ἐπάνω τῶν ὀφθαλμῶν [$=$ la partie supérieure des paupières, mais voir aussi la glose d'Hésychius sous κύλα, où le terme désigne la partie \italique{inférieure} des paupières]. οἱ δὲ τὸ περιστρέφεσθαι καὶ διατρίβειν περιβλεπόμενον. ἄλλοι φθονεῖν.
\noindent \textbf{(3)} Ælien, \italique{Varia Historia}, XIII, 15 ($=$ \latin{Adespota comica}, fr. 72 K.-A.) :
\citComm{φασὶ παχύτατον γενέσθαι τὴν διάνοιαν οἱ τῆς κωμῳδίας ποιηταὶ Πολύδωρόν τινα, καὶ ἄλλον Κοικυλίωνα ὄνομα, ὅσπερ τὰ κύματα ἠρίθμει ὑπὸ τῆς ἄγαν ἀνοίας.\\
\francais{Les poètes comiques disent qu'un certain Polydore (aussi [connu sous] un autre nom : Cœcylion) affolit tant qu'il comptait les vagues sous le coup de son trop plein de folie.}}
\noindent Ce passage, déniché par Fritzsche (à l'érudition et la bibliothèque toujours impressionnantes !), est à rapprocher du chapitre du Περὶ βλασφημιῶν de Suétone (édité, pour ce qu'il a pu en reconstituer, en 1967 aux Belles Lettres par le grand Taillardat) consacré à quelques idiots célèbres. Περὶ βλασφημιῶν, β 189 $=$ Eustathe \italique{Comm. Od.} p. 1669.53 : ὥσπερ δὲ τουτούς ἡ ἱστορία λόγου ἠξίωσεν ἐν οἷς καὶ τὸν τὰ κύματα μετροῦντα Πολύ[δ]ωρον.
\noindent\textbf{855-8} \textbf{Début du prologue de l'\italique{Hélène}.} Ainsi commence la véritable pièce d'Euripide :
\citComm{Ἑλένη --- Νείλου μὲν αἵδε καλλιπάρθενοι ῥοαί,\\
ὃς ἀντὶ δίας ψακάδος Αἰγύπτου πέδον\\
λευκῆς \difference{τακείσης χιόνος ὑγραίνει γύας.}\\
\francais{Du Nil voici les rives aux belles Nymphes,\\
Nil qui à la place de la pluie dispensée par Zeus\\
arrose la plaine d'Égypte une fois la blanche neige fondue et irrigue les terres.}}
\noindent Comme tous les semi-habiles – euphémisme pascalien : comprenez `semi-débiles' – Euripide aime à bavasser ; il ne peut s'empêcher dès le troisième vers de la pièce d'exposer son opinion sur la raison des crues du Nil – Ah ça ! se croit-il obligé d'en avoir une quel que soit le sujet ? et croit-il par-dessus le marché qu'elle nous intéresse ? Selon donc le très savant Euripide, ces crues sont dues à la fonte des neiges. Théorie fantaisiste et traitée avec dédain par Hérodote (II, 19-25), condescendant vis-à-vis de cette thèse soutenue par Anaxagore\marqueur{1}, le maître d'Euripide – raison pour laquelle ce dernier se croit tenu de la défendre ?
\noindent En tout cas il récidive, imperturbable, quelques années après l'\italique{Hélène} et les \italique{Thesmophories}, l'an 408/407, dans \italique{Archélaos}\marqueur{2}, tragédie aujourd'hui perdue. (Remarquons que dans le prologue d'\italique{Hélène} comme dans le fragment d'\italique{Archélaos}, les ῥοαί ne doivent pas se traduire par \italique{flots} mais bien par \italique{rives} ou \italique{berges}\marqueur{3} et que l'expression « καλλιπάρθενοι ῥοαί » suggère peut-être un jeu de mots avec les Ῥοαί, les nymphes des grenadiers\marqueur{4}.)
\end{multicols}
}%
\end{adjustbox}%
\vfill
\end{adjustwidth}
\newpage
% --- PAGE DE DROITE : Traduction française + Commentaire (6/8) ---
\begin{adjustwidth}{-2cm}{-2cm}
\noindent\textbf{\footnotesize Traduction française --- Folio 21 (\texttt{f. 172v}, vv. 838–875 [3/4]) \hfill ΠΑΡΑΒΑΣΙΣ / ΕΠΕΙΣΟΔΙΟΝ Δ'}\\[-2pt]
\noindent\rule{16.5cm}{0.5pt}\\[2pt]
{\fontsize{8.8pt}{10.2pt}\selectfont
\ligneEst{\textbf{Mnésiloque}}{\francais{« J'ai pour patrie un lieu non sans renom,}}{}
\ligneEst{}{\francais{c'est Sparte, et mon père est Tyndare. »}}{860}
\ligneEst{\textbf{Kritylla}}{\francais{\espace{4.5cm}Lui ton père, être de malheur ?}}{}
\ligneEst{}{\francais{Phrynondas, veux-tu dire.}}{}
\ligneEst{\textbf{Mnésiloque}}{\francais{« Je fus nommée Hélène. »}}{}
\ligneEst{\textbf{Kritylla}}{\francais{\espace{2.8cm}Voilà qu'encore tu redeviens femme,}}{}
\ligneEst{}{\francais{avant d'avoir expié ton premier déguisement.}}{}
\ligneEst{\textbf{Mnésiloque}}{\francais{« Et à cause de moi, bien des guerriers sont morts}}{}
\ligneEst{}{\francais{aux rives du Scamandre. »}}{865}
\ligneEst{\textbf{Kritylla}}{\francais{\espace{2.4cm}Que n'as-tu fait comme eux ?}}{}
\ligneEst{\textbf{Mnésiloque}}{\francais{« Moi, me voici. Mais lui l'infortuné », mon époux}}{}
}
\vspace{2pt}
\noindent\rule{16.5cm}{0.6pt}\\[2pt]
\noindent\textbf{\small Commentaire philologique, métrique et scénique (vv. 838–875 --- partie 6/8)}\\[2pt]
\begin{adjustbox}{max totalheight=15.2cm, max width=16.5cm, keepaspectratio}%
\parbox[t]{16.5cm}{
\setlength{\columnsep}{9pt}%
\begin{multicols}{2}
\fontsize{7.0pt}{8.0pt}\selectfont
\setlength{\parindent}{0pt}%
\setlength{\parskip}{1.2pt}%
\noindent Attaquons-nous au détournement du prologue d'\italique{Hélène} par Aristophane (\italique{Thesm.} vv. 855-857). Pour traîner dans la boue – ou pis – la muse d'Euripide, notre poète ne recule devant rien. La syrméa dont il est question au vers 857 était une espèce de rave très abondante en Égypte. On en tirait un suc portant le même nom utilisé comme purgatif\marqueur{5}. Suivant Hérodote\marqueur{6}, les Égyptiens se purgeaient menstruellement, trois jours de suite. Ce curieux hygiénisme est assez plaisant pour qu'Aristophane n'ait pas manqué de le railler.
\noindent Aristophane avait certainement mis en scène une confrontation entre le vendeur de syrméa et le vendeur de pâtisseries dans ses \italique{Danaïdes} (comédie aujourd'hui perdue) – les gâteaux de miel étaient en effet réputés astringents\marqueur{7}.
\noindent Dans la \italique{Paix}, Trygée suggère au marchand d'armes menacé par la ruine de couper ses lances pour les proposer à la vente comme échalas. Autre idée : recycler ses casques d'une bien étrange manière (vers 1253-1254) :
\citComm{Τρυγαῖος – Πώλει βαδίζων αὐτὰ τοῖς Αἰγυπτίοις·\\
ἔστιν γὰρ ἐπιτήδεια συρμαίαν μετρεῖν.\\
Trygée – Va les vendre en Égypte.\\
Ils sont bons à mesurer la syrméa.}
\noindent (C'est-à-dire à mesurer les \italique{effets} de la syrméa ; Trygée ne voit donc en ces casques que de vulgaires pots de chambre. En fait ils ne sont même bons à rien ; car Trygée fait entendre le proverbe κύματα μετρεῖν, compter les vagues, comble de l'inutilité\marqueur{8}.)
\noindent Même genre de détournement dans le présent passage des \italique{Thesmophories} : les Égyptiens se servent du Nil pour se \italique{nettoyer} (ici νοτίζειν = βαλανεύειν, σπογγίζειν) des désagréments de leur régime.
\noindent On peut conclure de ce passage qu'en grec le mot πέδον, le sol ; littéralement, ce sur quoi repose le pied, a fourni la même image qu'en français le \italique{fondement} et sert à désigner ἡ πυγή\marqueur{9}.
\noindent Revenons à notre passage des \italique{Thesmophories} et proposons une traduction :
\citComm{Μνησίλοχος – Νείλου μὲν αἵδε καλλιπάρθενοι ῥοαί,\\
ὃς ἀντὶ δίας ψακάδος Αἰγύπτου πέδον\\
λευκῆς νοτίζει μελανοσυρμαῖον λεών.\\
Mnésiloque – (Rôle d'Hélène.) « Du Nil ce sont ici les rives aux belles Nymphes,\\
Nil qui à la place de l'eau du ciel, en humecte le fondement\\
d'Égypte la blanche » et du peuple à la noire syrméa.}
\noindent On comprend la réaction indignée de la gardienne de Mnésiloque (\italique{Thesm}. v. 858) :
\citComm{Κρίτυλλα – πανοῦργος εἶ νὴ τὴν Ἑκάτην τὴν φωσφόρον.\\
Critylla – Tu es un effronté, par Hécate la porte-torche.\\
(Elle lui porte un coup de sa torche.)}
\noindent Cf. le vers 853 et le commentaire \italique{ad loc.} :
\noindent Critylla menaçait Mnésiloque d'un bon coup de flambeau. Et c'est au présent vers 858 qu'elle l'assène en évoquant Hécate \textbf{φωσ}φόρος.
\noindent\textbf{[Notes ad 855-8]} \textbf{(1)} 59A91 D.-K. = fr. 300 Eschyle (ed. Radt). Eschyle avait en son temps aussi évoqué cette fonte des neiges dans les \italique{Suppliantes}, v. 559 sq.
\noindent \textbf{(2)} Fr. 1 \italique{Ἀρχέλαος} Jouan-van Looy = fr. 228 ed. Kannicht.
\noindent \textbf{(3)} Hésychius (ρ 415) : ῥοία· κυλίστρα τῶν ἵππων παρὰ τῷ ποταμῷ καὶ ψάμμῳ. Rien que pour l'\italique{Hélène}, les ῥοαί sont à prendre au sens de rives au moins aux vers 53, 124 (avec jeu de mots sur Εὐρώτας), 250, 492.
\noindent \textbf{(4)} Eustathe, \italique{Commentaire à l'Odyssée}, p. 1572, 35 : ῥοιαὶ ὁμωνύμως τῷ δένδρῳ καὶ αἱ περὶ αὐτὰς μυθικαὶ νύμφαι.
\noindent \textbf{(5)} Souda (σ 1667) : Συρμαία : τὸ καθαρτικὸν δι' ἐμέτων καὶ κοιλίας. καὶ Συρμαίζειν τὸ κενοῦσθαι, ἐκ ζειῶν καὶ ὕδατος. οὕτως Ἡρόδοτος. λάχανον σελίνῳ ἐοικός : ἢ βρῶμα διὰ στέατος καὶ μέλιτος.\\ Et voici ce qu'en dit Hésychius (σ 2782) : σύρμαια· […] ἔστι δὲ βρωμάτιον διὰ στέατος καὶ μέλιτος. λέγεται δὲ καὶ συρμαϊσμός, καὶ ἔστι πρὸς κάθαρσιν. καὶ πόμα δι' ὕδατος καὶ ἁλῶν.
\end{multicols}
}%
\end{adjustbox}%
\vfill
\end{adjustwidth}
\newpage

% --- FOLIO 21 (f. 172v) : SPREAD 2.4 (Texte grec établi + Traduction française + Commentaire 7/8 & 8/8) ---
\clearpage\ifodd\value{page}\hbox{}\newpage\fi
\enTeteG{ΘΕΣΜΟΦΟΡΙΑΖΟΥΣΑΙ --- Texte grec établi (f. 172v, vv. 838–875 [4/4])}
\enTeteD{LES THESMOPHORIES --- Traduction française (f. 172v, vv. 838–875 [4/4])}
\begin{adjustwidth}{-2cm}{-2cm}
\noindent\textbf{\footnotesize Texte grec établi --- Folio 21 (\texttt{f. 172v}, vv. 838–875 [4/4]) \hfill ΠΑΡΑΒΑΣΙΣ / ΕΠΕΙΣΟΔΙΟΝ Δ'}\\[-2pt]
\noindent\rule{16.5cm}{0.5pt}\\[2pt]
{\fontsize{8.8pt}{10.2pt}\selectfont
\ligneEst{}{οὑμὸς Μενέλεως οὐδέπω προσέρχεται.}{}
\ligneEst{}{τί οὖν ἔτι ζῶ;}{}
\ligneEst{\textbf{Κρίτυλλα}}{\decalage{τί οὖν ἔτι ζῶ;}τῶν κοράκων πονηρίᾳ.}{}
\ligneEst{\textbf{Μνησίλοχος}}{ἀλλ᾽ ὥσπερ αἰκάλλει τι καρδίαν ἐμήν.}{}
\ligneEst{}{μὴ ψεῦσον ὦ Ζεῦ τῆς ἐπιούσης ἐλπίδος.}{870}
\ligneEst{\textbf{Εὐριπίδης}}{τίς τῶνδ᾽ ἐρυμνῶν δωμάτων ἔχει κράτος,}{}
\ligneEst{}{ὅστις ξένους δέξαιτο ποντίῳ σάλῳ}{}
\ligneEst{}{κάμνοντας ἐν χειμῶνι καὶ ναυαγίαις;}{}
\ligneEst{\textbf{Μνησίλοχος}}{Πρωτέως τάδ᾽ ἐστὶ μέλαθρα.}{}
\ligneEst{\textbf{Κρίτυλλα}}{\decalage{Πρωτέως τάδ᾽ ἐστὶ μέλαθρα.}ποίου Πρωτέως;}{}
\ligneEst{}{ὤ τρισκακοδαίμων, ψεύδεται νὴ τὼ θεώ,}{875}
}
\vspace{2pt}
\noindent\rule{16.5cm}{0.6pt}\\[2pt]
\noindent\textbf{\small Commentaire philologique, métrique et scénique (vv. 838–875 --- partie 7/8)}\\[2pt]
\begin{adjustbox}{max totalheight=15.2cm, max width=16.5cm, keepaspectratio}%
\parbox[t]{16.5cm}{
\setlength{\columnsep}{9pt}%
\begin{multicols}{2}
\fontsize{7.0pt}{8.0pt}\selectfont
\setlength{\parindent}{0pt}%
\setlength{\parskip}{1.2pt}%
\noindent \textbf{(6)} Hérodote II, 77 ; et aussi II, 88 ; II, 125.
\noindent \textbf{(7)} C'est ce qui ressort d'un passage des \italique{Nuées} (vers 506-508) : Strepsiade hésite à entrer dans l'antre de Socrate, qu'il compare aimablement à celui de Trophonius : εἰς τὼ χεῖρέ νυν / δός μοι μελιτοῦτταν πρότερον, ὡς δέδοικ' ἐγὼ / εἴσω καταβαίνων ὥσπερ εἰς Τροφωνίου. « Mets-moi d'abord en main un gâteau au miel ; car j'ai peur de descendre, comme si je pénétrais dans l'antre de Trophonios. » Il faut bien sûr prononcer d'un trait `εἰς Τροφωνίου' pour faire entendre στρόφος, la colique (la peur se manifestant chez les Grecs par des troubles du bas-ventre ; en français : \italique{avoir les pétoches}). En passant, Euripide suggère le même jeu de mots scatologique : au vers 300 de son \italique{Ion} : Créuse dit à Ion qu'elle est venue avec son mari, mais qu'il a dû s'arrêter en chemin (ἐκστρέφει Τροφωνίου ; Euripide, ainsi qu'il est coutumier, joue à différents niveaux de compréhension)… Le tour était proverbial et Érasme l'a recensé (\italique{Adages} n° 677). Sur ces gâteaux de miel qui préservaient des dérangements intestinaux, voir aussi \italique{Thesm}. 570. On peut maintenant comprendre le bon mot de table qu'Athénée (\italique{Deipnosophistes}, I, 4d) rapporte : ἀμίδα δότω τις ἢ πλακοῦντά τις δότω, « Qu'on m'apporte un pot de chambre ou bien un gâteau ! » Voilà aussi pourquoi Aristophane fait intervenir dans ses \italique{Danaïdes} deux vendeurs opposés : celui de syrméa et celui d'ἐγκρίδες (Pollux, \italique{Onomasticon}, VII 199 = fr. 276(265) Kassel-Austin) : ἐν δ' Ἀριστοφάνης Δαναΐσι συρμαιοπῶλαι καὶ ἐγκριδοπῶλαι. Athénée (XIV, 645e) définit les encrides comme des gâteaux frits à l'huile et emmiellés : ΕΓΚΡΙΔΕΣ. πεμμάτιον ἑψόμενον ἐν ἐλαίῳ καὶ μετὰ τοῦτο μελιτούμενον.
\noindent \textbf{(8)} Voir Érasme, \italique{Adages} (n° 345, \latin{Undas numeras}) ou Suétone, Περὶ βλασφημιῶν (éd. Taillardat), β 188.
\noindent \textbf{(9)} Noter que si πέδον, \italique{la plaine}, renvoie au verso ; πεδίον, \italique{la petite plaine}, elle, fait allusion au recto. Cf. \italique{Lysistrata}, v. 88 ; \italique{Oiseaux}, v. 507 avec le commentaire de Willems. Voir aussi l'emploi de λειμών au vers 171 du \italique{Cyclope} d'Euripide.
\noindent\textbf{858} \textbf{τὴν Ἑκάτην τὴν φωσφόρον.} Ce n'est pas par hasard que Critylla adjoint à Hécate l'épithète de φωσφόρος, \italique{porte-torche} : elle est elle-même équipée d'une telle torche et il est à parier qu'en jurant de la sorte, Critylla met sa menace du vers 853 à exécution en assénant à Mnésiloque le coup de flambeau promis.
\noindent Cf. \italique{Lysistrata} vers 453 ; fr. 592.5 Kassel-Austin [$=$ Athénée XV, 701b] : ἐκφέρετε πεύκας κατ᾽ Ἀγάθωνα φωσφόρους.
\noindent\textbf{859-60} \textbf{} \italique{Hélène}, v. 16-17a :
\citComm{\difference{ἡμῖν} δὲ γῆ μὲν πατρὶς οὐκ ἀνώνυμος\\
Σπάρτη, πατὴρ δὲ Τυνδάρεως.}
\noindent\textbf{861} \textbf{Φρυνώνδας.} Phrynondas et Eurybatos, deux coquins dont le nom était devenu appellatif, comme en français Cartouche et Mandrin. (Note de Willems ; qui a en tête le \italique{Protagoras} (327d) de Platon lorsque ces deux truands de la tradition populaire sont rapprochés.)\\ \italique{Souda}, φ 770 :
\citComm{Φρυνώνδας· τῶν ἐπὶ πονηρίᾳ διαβεβοημένων. ὃς ξένος ὢν κατὰ τὰ Πελοποννησιακὰ διέτριβεν Ἀθήνησιν. Ἀριστοφάνης Ἀμφιάρεῳ· ὦ μιαρὲ καὶ Φρυνώνδα καὶ πονηρὲ σύ. ἐκ τούτου τοὺς πονηροὺς Φρυνώνδας καλοῦσι. κακοήθης ὁ Φρυνώνδας οὗτος καὶ πανοῦργος, μισθοῦ τινας μηχανὰς πραγμάτων ἐπὶ κακοῖς ῥᾳδίως συντιθείς, ὥς φησιν Εὔδημος. καὶ αὖθις· ἀλλ' οἶμαι οὔτε Φρυνώνδας οὔτ' Εὐρύβατος οὐδ' ἄλλος πώποτε τῶν πάλαι πονηρῶν, τοιοῦτος μάγος καὶ γόης.\\
\francais{Phrynondas : de ceux connus pour leur méchanceté. C'était un étranger qui vivait à Athènes au moment de la guerre du Péloponnèse. Aristophane dans son \italique{Amphiaraos} : « ô misérable et Phrynondas et vil que tu es ». On appelle les méchants de son nom : Phrynondas.
Ce Phrynondas [était] mauvais et méchant ; quelqu'un qui, pour de l'argent, élaborait facilement quelque intrigue pour de mauvaises fins, comme le dit Eudemos.
Et ailleurs : « mais je pense que ni Phrynondas ni Eurybatos ni aucun autre des vilains d'autrefois ne se sont jamais comportés comme un filou et un tricheur ».}
}
\noindent\textbf{862} \textbf{} \italique{Hélène}, v. 22 :
\citComm{Ἑλένη δ᾽ ἐκλήθην. ἃ δὲ πεπόνθαμεν κακὰ}
\end{multicols}
}%
\end{adjustbox}%
\vfill
\end{adjustwidth}
\newpage
% --- PAGE DE DROITE : Traduction française + Commentaire (8/8) ---
\begin{adjustwidth}{-2cm}{-2cm}
\noindent\textbf{\footnotesize Traduction française --- Folio 21 (\texttt{f. 172v}, vv. 838–875 [4/4]) \hfill ΠΑΡΑΒΑΣΙΣ / ΕΠΕΙΣΟΔΙΟΝ Δ'}\\[-2pt]
\noindent\rule{16.5cm}{0.5pt}\\[2pt]
{\fontsize{8.8pt}{10.2pt}\selectfont
\ligneEst{}{\francais{Ménélas n'apparaît toujours point.}}{}
\ligneEst{}{\francais{« Ah ! pourquoi vis-je encore? »}}{}
\ligneEst{\textbf{Kritylla}}{\francais{\espace{2.3cm}C'est la faute aux corbeaux.}}{}
\ligneEst{\textbf{Mnésiloque}}{\francais{(Mnésiloque aperçoit Euripide.) Mais quelque chose vient flatter mon cœur.}}{}
\ligneEst{}{\francais{Ô Zeus, ne démens pas ma naissante espérance !}}{870}
\ligneEst{\textbf{Euripide}}{\francais{(Rôle de Ménélas.) « Quel est le chef puissant de ces fortes demeures ? »}}{}
\ligneEst{}{\francais{Voudra-t-il accueillir des hôtes dérangés par la houle marine}}{}
\ligneEst{}{\francais{au cours d'une tempête et d'un naufrage ?}}{}
\ligneEst{\textbf{Mnésiloque}}{\francais{C'est à Protée qu'est ce palais.}}{}
\ligneEst{\textbf{Kritylla}}{\francais{\espace{4.2cm}Qui ça, Protée ?}}{}
\ligneEst{}{\francais{Oh le triple misérable! Il ment par les deux déesses.}}{875}
}
\vspace{2pt}
\noindent\rule{16.5cm}{0.6pt}\\[2pt]
\noindent\textbf{\small Commentaire philologique, métrique et scénique (vv. 838–875 --- partie 8/8)}\\[2pt]
\begin{adjustbox}{max totalheight=15.2cm, max width=16.5cm, keepaspectratio}%
\parbox[t]{16.5cm}{
\setlength{\columnsep}{9pt}%
\begin{multicols}{2}
\fontsize{7.0pt}{8.0pt}\selectfont
\setlength{\parindent}{0pt}%
\setlength{\parskip}{1.2pt}%
\noindent\textbf{864-5} \textbf{} \italique{Hélène}, v. 52-53a :
\citComm{ψυχαὶ δὲ πολλαὶ δι᾽ ἔμ᾽ ἐπὶ Σκαμανδρίοις\\
ῥοαῖσιν ἔθανον· \difference{ἡ δὲ πάντα τλᾶσ᾽ ἐγὼ}}
\noindent\textbf{866-8} \textbf{} \italique{Hélène}, v. 49 :
\citComm{κἀγὼ μὲν ἐνθάδ᾽ εἴμ᾽, ὁ δ᾽ ἄθλιος πόσις}
\noindent \italique{Hélène}, v. 56a :
\citComm{τί δῆτ᾽ ἔτι ζῶ; }
\noindent\textbf{870} \textbf{μὴ ψεῦσον ὦ Ζεῦ τῆς ἐπιούσης ἐλπίδος.} Voici la scholie au vers 870 :
\citComm{σύνηθες τὸ σχῆμα. Μένανδρος·\\
\alinea οἶσθ᾽ ὃ ποίησον.\\
Σοφοκλῆς Πηλεῖ·\\
\alinea μὴ ψεῦσον ὦ Ζεῦ· μή μ᾽ ἕλης ἄνευ δορός.\\
\alinea οἶσθ᾽ οὖν ὃ δρᾶσον.\\
Εὐριπίδης Πολυΐδῳ.\\
\francais{La construction est usuelle. Ménandre : 
« Tu sais quoi faire. » 
Sophocle dans son \italique{Pelée} : 
« Ne me trompe pas, ô Zeus ; tu ne me prendras pas sans lutte. /
Tu sais donc ce qu'il te reste à faire. »
(Euripide dans son \italique{Polyidos}.)}}
\noindent Cette scholie ne se peut comprendre qu'en la mettant en regard d'un passage du Περὶ διαλέκτων de Grégoire de Corinthe (\nieme{XII} siècle, p. 14 ed. Schaefer) ou de la scholie de l'\italique{Hécube} d'Euripide (vers 225 -- nos trois commentaires remontent assurément à la même source ; cf. aussi Suidas ψ 62). À propos de la tournure `οἶσθ᾽ οὖν ὃ δρᾶσον;', voici ce que le scholiaste de l'\italique{Hécube} écrit :
\citComm{δρᾶσον: τὰ καλὰ τῶν ἀντιγράφων δρᾶσον ἔχει. οἱ Ἀττικοὶ χρῶνται τοῖς προστακτικοῖς ἀντὶ ὑποτακτικῶν ἀπαγορευτικῶν, καὶ αὐτοῖς τοῖς προστακτικοῖς ἀντὶ ὁριστικῶν.\\
τοῦ μὲν οὖν πρώτου τὸ Ἀριστοφάνους ἐν Θεσμοφοριαζούσαις (v. 870) παράδειγμα `μὴ ψεῦσον, ὦ Ζεῦ, τῆς ἐπιούσης ἡμέρας (sic).'\\
τοῦ δευτέρου δὲ Μένανδρος `οἶσθ' ὅ τι ποίησον' ἀντὶ 





τοῦ `ὅ τι ποιήσεις'.\\
τὸν αὐτὸν τρόπον καὶ Εὐριπίδης `οἶσθ᾿ οὖν ὃ δρᾶσον' ἀντὶ τοῦ `δράσεις' ἐξήνεγκεν.\\
δρᾶσον: les meilleures copies ont `δρᾶσον'.
Les Attiques recourent aux impératifs à place de subjonctifs de prohibition. Et de nouveau aux (je lis : καὶ αὖ τοῖς) impératifs à la place d'indicatifs.\\
L'exemple d'Aristophane sur le premier point : `μὴ ψεῦσον κτλ.'\\
Ménandre sur le second point : `οἶσθ' ὅ τι ποίησον' à la place de `ὅ τι ποιήσεις'.\\
Euripide a le même tour : `οἶσθ᾿ οὖν ὃ δρᾶσον' à la place de `δράσεις'.}
\noindent Noter que ce tour d'Euripide se trouve justement dans l'\italique{Hélène} (vers 325) :
\citComm{Χορός --- οἶσθ᾽ οὖν ὃ δρᾶσον; μνήματος λιποῦσ᾽ ἕδραν.\\
\francais{Hé bien, sais-tu ce que tu dois faire ? Quitte ton poste près du tombeau.}}
\noindent\textbf{871} \textbf{} \italique{Hélène}, v. 68 :
\citComm{Τεῦκρος --- τίς τῶνδ᾽ ἐρυμνῶν δωμάτων ἔχει κράτος;}
\noindent \italique{Hélène}, v. 449 :
\citComm{Μενελέως --- ναυαγὸς \difference{ἥκω} ξένος, \difference{ἀσύλητον γένος}.}
\noindent\textbf{874-8} \textbf{} \italique{Hélène}, v. 460 :
\citComm{Πρωτεὺς τάδ᾽ οἰκεῖ δώματ᾽, Αἴγυπτος δὲ γῆ.}
\noindent\textbf{874-6} \textbf{ποίου Πρωτέως; ... ἐπεὶ τέθνηκε Πρωτέας ἔτη δέκα.} Elle confond Protée avec Protéas\marqueur{1}, stratège athénien dont il est question dans Thucydide\marqueur{2}. (Note de Willems.)
\noindent\textbf{[Notes ad 874-6]} \textbf{(1)} Identiques au génitif.
\noindent \textbf{(2)} Thucydide I, 45.2 ; II, 23.2.
\end{multicols}
}%
\end{adjustbox}%
\vfill
\end{adjustwidth}
\newpage
\clearpage\ifodd\value{page}\hbox{}\newpage\fi
\enTeteG{ΕΠΕΙΣΟΔΙΟΝ Δ' (Ἑλένη) --- Codex Ravennas 429, f. 173r (Classense 175r) --- vv. 876–919}
\begin{adjustwidth}{-1.5cm}{-3.0cm}
\begin{adjustbox}{max totalheight=21.0cm, max width=16.5cm, keepaspectratio}
\begin{minipage}{16.5cm}
{\fontsize{7.2pt}{8.5pt}\selectfont
\noindent\textbf{Folio 22 --- Codex Ravennas 429, \texttt{f. 173r} (\texttt{Classense 175r}) : ΕΠΕΙΣΟΔΙΟΝ Δ' (Ἑλένη), vv. 876–919}\hfill \textit{Témoins :} $\text{R}^1, \Rpc, \Rdeux$ (\Mon)\par\vspace{2pt}
\hrule height 0.4pt\vspace{2.5pt}
}
{\fontsize{8.1pt}{9.5pt}\selectfont
\ligneRav{876}{\auDessus{γυ}{ν}}{ἐπεὶ τέθνηκε Πρωτέας ἔτη δέκα·}{}
\ligneRav{}{\auDessus{εὐ}{ρ}}{ποίαν δὲ χώραν εἰσεκέλσαμεν σκάφει;}{}
\ligneRav{}{μ\auDessus{ν}{Η}}{Αἴγυπτον. : ὦ δύστηνος οἷ πεπλώκαμεν·}{}
\ligneRav{879}{μ\auDessus{ν}{Η}}{πείθει τι \difference{τῷ κακῷ} κακῶς ἀπολουμένῳ}{\textbf{879} τῷ R : τούτῳ τῷ Grynaeus : \lecon{τῷ κακῷ} Scaliger}
\ligneRav{880}{}{ληροῦντι λῆρον; Θεσμοφόριον τουτογί·}{\textbf{880} \lecon{Θεσμοφόριον} R, cf.278}
\ligneRav{881}{\auDessus{εὐ}{ρ}}{αὐτὸς δὲ Πρωτεὺς ἔνδον ἔστ᾽ \difference{ἢ 'ξώπιος};}{\textbf{881} ἠξώπιος R, id. v.884 : \lecon{ἢ 'ξώπιος;} Biset}
\ligneRav{}{\auDessus{γυ}{ν}}{οὐκ ἔσθ᾽ ὅπως οὐ ναυτιᾷς ἔτ᾽ ὦ ξένε,}{}
\ligneRav{883}{}{\difference{ὅστις γ᾽} ἀκούσας ὅτι τέθνηκε Πρωτέας}{\textbf{883} ὅστις R : \lecon{ὅστις γ᾽} Scaliger}
\ligneRav{}{}{ἔπειτ᾽ ἐρωτᾷς `ἔνδον ἔστ᾽ ἢ 'ξώπιος;'}{}
\ligneRav{885}{\auDessus{εὐ}{ρ}}{αἰαῖ \difference{τέθνηκε}. ποῦ δ᾽ ἐτυμβεύθη τάφῳ;}{\textbf{885} τέθνηκεν R : \lecon{τέθνηκε} Bentley}
\ligneRav{}{μ\auDessus{ν}{Η}}{τόδ᾽ ἐστὶν αὐτοῦ σῆμ᾽, ἐφ᾽ ᾧ καθήμεθα·}{}
\ligneRav{887}{\auDessus{γυ}{ν}}{κακῶς τ᾽ ἄρ᾽ ἐξόλοιο κἀξολεῖ γέ τοι,}{\textbf{887} \lecon{κακῶς τ᾽ ἄρ᾽} R : κακῶς γ᾽ ἄρ᾽ Denniston (\italique{CR} XLIV(1930), p.215)}
\ligneRav{}{}{ὅστις γε τολμᾷς σῆμα τὸν βωμὸν καλεῖν·}{}
\ligneRav{889}{\auDessus{εὐ}{ρ}}{\difference{τί δὴ σὺ} θάσσεις τάσδε τυμβήρεις ἕδρας}{\textbf{889} τί δὲ σὺ R : τί δαὶ σὺ Faber : \lecon{τί δὴ σὺ} Scaliger}
\ligneRav{890}{—}{φάρει καλυπτὸς ὦ ξένη; : βιάζομαι}{}
\ligneRav{891}{}{γάμοισι Πρωτέως παιδὶ \difference{συμμεῖξαι} λέχος·}{\textbf{891} συμμῖξαι R : \lecon{συμμεῖξαι} van Leeuwen}
\ligneRav{892}{\auDessus{γυ}{ν}}{τί ὦ κακόδαιμον ἐξαπατᾷς \difference{αὖ τὸν} ξένον;}{\textbf{892} αὐτὸν (ex αὖτὸν) \Rpc : \lecon{αὖ τὸν} Casaubon}
\ligneRav{}{}{οὗτος πανουργῶν δεῦρ᾽ ἀνῆλθεν ὦ ξένε}{}
\ligneRav{}{}{ὡς τὰς γυναῖκας ἐπὶ κλοπῇ τοῦ χρυσίου·}{}
\ligneRav{895}{μ\auDessus{ν}{Η}}{\difference{κόνυζα} τοὐμὸν σῶμα βάλλουσα ψόγῳ·}{\textbf{895} βαϋζα R : βάϋζε Zanetti : \lecon{κόνυζα} Bibi}
\ligneRav{}{\auDessus{εὐ}{ρ}}{ξένη τίς ἡ γραῦς ἡ κακορροθοῦσά σε;}{}
\ligneRav{}{μ\auDessus{ν}{Η}}{αὕτη Θεονόη Πρωτέως. : μὰ τὼ θεὼ}{}
\ligneRav{}{}{εἰ μὴ Κρίτυλλά γ᾽ Ἀντιθέου Γαργηττόθεν·}{}
\ligneRav{}{—}{σὺ δ᾽ εἶ πανοῦργος. : ὁπόσα τοι βούλει λέγε·}{}
\ligneRav{900}{}{οὐ γὰρ γαμοῦμαι σῷ κασιγνήτῳ ποτὲ}{}
\ligneRav{901}{}{προδοῦσα \difference{Μενέλεων} τὸν ἐμὸν ἐν Τροίᾳ πόσιν·}{\textbf{901} Μενέλαον R : \lecon{Μενέλεων} Scaliger}
\ligneRav{}{\auDessus{εὐ}{ρ}}{γύναι τί εἶπας; στρέψον ἀνταυγεῖς κόρας·}{}
\ligneRav{}{μ\auDessus{ν}{Η}}{αἰσχύνομαί σε τὰς γνάθους ὑβρισμένη·}{}
\ligneRav{}{\auDessus{εὐ}{ρ}}{τουτὶ τί ἔστιν; ἀφασία τίς τοί μ᾽ ἔχει·}{}
\ligneRav{905}{}{ὦ θεοὶ τίν᾽ ὄψιν εἰσορῶ; τίς εἶ γύναι;}{}
\ligneRav{906}{μ\auDessus{ν}{Η}}{σὺ δ᾽ εἶ τίς; \difference{αὑτὸς} γὰρ σὲ κἄμ᾽ ἔχει λόγος·}{\textbf{906} αὐτὸς R : \lecon{αὑτὸς} Boissonade}
\ligneRav{}{\auDessus{εὐ}{ρ}}{Ἑλληνὶς εἶ τις ἢ 'πιχωρία γυνή;}{}
\ligneRav{}{μ\auDessus{ν}{Η}}{Ἑλληνίς. ἀλλὰ καὶ τὸ σὸν θέλω μαθεῖν·}{}
\ligneRav{909}{\auDessus{εὐ}{ρ}}{\difference{Ἑλένῃ σ᾽} ὁμοίαν δὴ \difference{μάλιστ᾽ εἶδον} γύναι·}{\textbf{909} Ἐλένης R : \lecon{Ἑλένῃ σ᾽} Eur. \italique{Hel.}563 \espApp μάλισθ᾽ ἴδον R : \lecon{μάλιστ᾽ εἶδον} \italique{Hel}.563}
\ligneRav{910}{μ\auDessus{ν}{Η}}{ἐγὼ δὲ \difference{Μενελέῳ σ᾽} ὅσα γ᾽ ἐκ τῶν \difference{λαιφίων}·}{\textbf{910} Μενέλαον R : Μενέλεων Suid. : \lecon{Μενελέῳ σ᾽} Thiersch \espApp ἀφύων R : ἰφύων Σ, Suid. : ἀμφίων Grégoire : \lecon{λαιφίων} Bibi \espApp Μενέλαον R : Μενέλεων Suid.}
\ligneRav{}{\auDessus{εὐ}{ρ}}{ἔγνως ἄρ᾽ ὀρθῶς ἄνδρα δυστυχέστατον·}{}
\ligneRav{}{μ\auDessus{ν}{Η}}{ὦ χρόνιος ἐλθὼν σῆς δάμαρτος ἐσχάρας}{}
\ligneRav{914}{}{\hspace{1.1em}λαβέ με λαβέ με πόσι, \difference{περίβαλε} δὲ χέρας·}{\textbf{914} περίβαλλε R : \lecon{περίβαλε} Biset}
\ligneRav{915}{}{\hspace{1.1em}φέρε σὲ κύσω. ἄπαγέ μ᾽ ἄπαγ᾽ ἄπαγέ με}{}
\ligneRav{}{—}{λαβὼν ταχὺ πάνυ. : κλαύσετ᾽ ἄρα νὴ τὼ θεὼ}{}
\ligneRav{}{}{ὅστις σ᾽ ἀπάξει τυπτόμενος τῇ λαμπάδι·}{}
\ligneRav{}{\auDessus{εὐ}{ρ}}{σὺ τὴν ἐμὴν γυναῖκα κωλύεις ἐμέ,}{}
\ligneRav{}{}{τὴν Τυνδάρειον παῖδ᾽, ἐπὶ Σπάρτην ἄγειν;}{}
}
\vspace{3pt}\noindent\rule{4.5cm}{0.4pt}\par\vspace{1.5pt}
{\fontsize{7.0pt}{8.3pt}\selectfont
\noindent\textbf{Loci paralleli (vv. 876–919) ---} \textbf{878} Suidas οι 1 ὦ δύστηνος, οἷ πεπλώκαμεν; \espTest \textbf{882} (ναυττιᾶν $=$ ναυσιᾶν) Moeris ν 9 \espTest \textbf{902} (στρέψον ἀνταυγεῖς κόρας) Suidas α 2617 \espTest (ἀνταυγεῖς κόρας) Photius α 2059 ἀνταυγεῖς κόρας \espTest \textbf{904} \Antiatt 83.9 ; cf. Hesychius α 8572 ἀφασία \espTest \textbf{910} Suidas ι 778 (sed Μενέλεων pro Μενέλαον, ὄσον pro ὅσα γ᾽ et ἰφύων pro ἀφύων) \espTest \textbf{912} (ἐσχάρας) Photius ε 2044 ἐσχάρας $\subset$ Hesychius ε 6447 ἐσχάραι\par
}
\vspace{2.5pt}\noindent\rule{4.5cm}{0.4pt}\par\vspace{1.5pt}
{\fontsize{6.9pt}{8.2pt}\selectfont
\noindent\textbf{Suite (en commentaire) de l'apparat critique \& notes paléographiques (\texttt{f. 173r}) ---} \textbf{Contexte :} Suite et sommet de la \textbf{parodie d'\italique{Hélène}} : stichomythie de reconnaissance entre Euripide-Ménélas et Mnésiloque-Hélène sous le regard incrédule et furieux de la gardienne (\textbf{Kritylla}, vv. 876–919). --- \textbf{Deux émendations paléographiques majeures de \textbf{Bibi} (vv. 895 et 910) :} --- \textbf{v. 895 :} βαϋζα R (βάϋζε Zanetti) : \lecon{κόνυζα} \lecon{Bibi} (confusion minuscule $\kappa \to \beta$). --- \textbf{v. 910 :} Μενέλαον R (Μενέλεων Souda) : \lecon{Μενελέῳ σ'} Thiersch \espApp \textbf{ἀφύων R} (ἰφύων $\Sigma$, Souda ; ἀμφίων Grégoire) : \lecon{λαιφίων} \lecon{Bibi} ap. Hésychius $\lambda\ 171$ (chute de $\Lambda$ devant $\text{Α}$ en onciale).\par
}
\end{minipage}
\end{adjustbox}
\end{adjustwidth}
\newpage
\enTeteD{Codex Ravennas 429 (Biblioteca Classense) --- Fac-similé du f. 173r (vv. 876–919)}
\begin{adjustwidth}{-1.5cm}{-3.0cm}
\begin{center}
\includegraphics[width=16.5cm,height=20.5cm,keepaspectratio]{/workspaces/Thesmophories/dossier_ravennas_thesmophories/folios_individuels_ff162v-177r/22_f173r_classense175r_vv0876-0919.jpg}\par
\vspace{2pt}
{\scriptsize\textbf{Codex Ravennas 429, f. 173r (Classense 175r)} --- ΕΠΕΙΣΟΔΙΟΝ Δ' (Ἑλένη), \textbf{vv. 876–919} ($\text{R}^1, \Rpc, \Rdeux$ (\Mon)).}
\end{center}
\end{adjustwidth}
\newpage
\clearpage\ifodd\value{page}\hbox{}\newpage\fi% --- FOLIO 22 (f. 173r) : SPREAD 2.1 (Texte grec établi + Traduction française + Commentaire 1/8 & 2/8) ---
\clearpage\ifodd\value{page}\hbox{}\newpage\fi
\enTeteG{ΘΕΣΜΟΦΟΡΙΑΖΟΥΣΑΙ --- Texte grec établi (f. 173r, vv. 876–919 [1/4])}
\enTeteD{LES THESMOPHORIES --- Traduction française (f. 173r, vv. 876–919 [1/4])}
\begin{adjustwidth}{-2cm}{-2cm}
\noindent\textbf{\footnotesize Texte grec établi --- Folio 22 (\texttt{f. 173r}, vv. 876–919 [1/4]) \hfill ΕΠΕΙΣΟΔΙΟΝ Δ' (Ἑλένη)}\\[-2pt]
\noindent\rule{16.5cm}{0.5pt}\\[2pt]
{\fontsize{8.8pt}{10.2pt}\selectfont
\ligneEst{}{ἐπεὶ τέθνηκε Πρωτέας ἔτη δέκα.}{876}
\ligneEst{\textbf{Εὐριπίδης}}{ποίαν δὲ χώραν εἰσεκέλσαμεν σκάφει;}{}
\ligneEst{\textbf{Μνησίλοχος}}{Αἴγυπτον.}{}
\ligneEst{\textbf{Εὐριπίδης}}{\decalage{Αἴγυπτον.}ὦ δύστηνος οἷ πεπλώκαμεν.}{}
\ligneEst{\textbf{Μνησίλοχος}}{πείθει τι τῷ κακῷ κακῶς ἀπολουμένῳ}{}
\ligneEst{}{ληροῦντι λῆρον; Θεσμοφόριον τουτογί.}{880}
\ligneEst{\textbf{Εὐριπίδης}}{αὐτὸς δὲ Πρωτεὺς ἔνδον ἔστ᾽ ἢ 'ξώπιος;}{}
\ligneEst{\textbf{Κρίτυλλα}}{οὐκ ἔσθ᾽ ὅπως οὐ ναυτιᾷς ἔτ᾽ ὦ ξένε,}{}
\ligneEst{}{ὅστις γ᾽ ἀκούσας ὅτι τέθνηκε Πρωτέας}{}
\ligneEst{}{ἔπειτ᾽ ἐρωτᾷς `ἔνδον ἔστ᾽ ἢ 'ξώπιος;'}{}
\ligneEst{\textbf{Εὐριπίδης}}{αἰαῖ τέθνηκε. ποῦ δ᾽ ἐτυμβεύθη τάφῳ;}{885}
\ligneEst{\textbf{Μνησίλοχος}}{τόδ᾽ ἐστὶν αὐτοῦ σῆμ᾽, ἐφ᾽ ᾧ καθήμεθα.}{}
}
\vspace{2pt}
\noindent\rule{16.5cm}{0.6pt}\\[2pt]
\noindent\textbf{\small Commentaire philologique, métrique et scénique (vv. 876–919 --- partie 1/8)}\\[2pt]
\begin{adjustbox}{max totalheight=15.2cm, max width=16.5cm, keepaspectratio}%
\parbox[t]{16.5cm}{
\setlength{\columnsep}{9pt}%
\begin{multicols}{2}
\fontsize{7.6pt}{8.8pt}\selectfont
\setlength{\parindent}{0pt}%
\setlength{\parskip}{1.2pt}%
\noindent\textbf{878} \textbf{} \italique{Hélène}, v. 461 :
\citComm{Αἴγυπτος; ὦ δύστηνος, οἷ \difference{πέπλευκ᾽ ἄρα.}}
\noindent\textbf{881} \textbf{} \italique{Hélène}, v. 465 :
\citComm{Μενελέως --- \difference{ποῦ δῆτ᾽ ἂν εἴη; πότερον ἐκτὸς ἢ 'ν δόμοις;}}
\noindent\textbf{886} \textbf{} \italique{Hélène}, v. 466 :
\citComm{Γραῦς --- τόδ᾽ ἐστὶν αὐτοῦ \difference{μνῆμα, παῖς δ᾽ ἄρχει χθονός}.}
\noindent\textbf{894} \textbf{ὡς τὰς γυναῖκας ἐπὶ κλοπῇ τοῦ χρυσίου,} «~Subtiliser les bijoux des femmes~» ; voilà un chef d'accusation nouveau contre Mnésiloque d'autant plus étonnant que rien jusqu'ici ne l'a laissé présager.
\noindent Le fait que Critylla porte un flambeau indique pourtant qu'elle est présente depuis l'arrivée de la procession. Ment-elle pour éviter d'avoir à raconter la trop longue affaire ? il pourrait y avoir une autre explication.
\noindent Il semble que seules les femmes tenant un rôle actif lors de ces cérémonies pouvaient porter des bijoux (prêtresses, canéphores, \&c.) ; c'est ce qui ressort notamment d'un passage de \italique{Lysistrata} :
\citComm{στρωμάτων δὲ ποικίλων καὶ /
χλανιδίων καὶ ξυστίδων καὶ /
\textbf{χρυσίων}, ὅσ᾽ ἐστί μοι, /
οὐ φθόνος ἔνεστί μοι πᾶσι παρέχειν φέρειν /
τοῖς παισίν, \textbf{ὁπόταν} τε θυγάτηρ τινὶ \textbf{κανηφορῇ}. \\
Et de surtouts brodés, petites chlamides, xystides, \italique{bijoux}, tout ce que je possède, il ne me coûte pas d'offrir à chacun de l'emporter, pour ses enfants et pour sa fille, \italique{si elle devenait canéphore}.}
\noindent En fait, le terme χρυσίον, \italique{bijou}, \italique{ornement} semble avoir parfois un sens imagé chez Aristophane et désigner les \italique{appâts} féminins. Ainsi, lorsque Pisthétaire s'exclame à propos de la belle \& court vêtue joueuse de flûte Procné (\italique{Oiseaux}, v. 669) :
\citComm{Πισθέταιρος --- ὅσον δ᾽ ἔχει τὸν χρυσόν, ὥσπερ παρθένος,}
\noindent il est bien évident qu'il fait l'éloge de sa beauté, son galbe, ou que sais-je encore\marqueur{1} et non pas de sa verroterie. Ce blason sera d'ailleurs continué quelques vers plus loin par le chœur (vers 676-684).
\noindent Autre exemple : lorsque Dicéopolis fait venir la jeune canéphore en lui prodiguant ce conseil (vers 257-258) :
\citComm{πρόβαινε, κἀν τὤχλῳ φυλάττεσθαι σφόδρα\\
μή τις λαθών σου περιτράγῃ τὰ χρυσία.\\
\francais{Avance, et surtout prends garde que dans la foule\\
on ne te grignote tes bijoux.}}
\noindent Certes, comme on l'a vu, les canéphores portaient des bijoux ; mais on devine aisément avec Aristophane qu'en voyant cette jeune beauté dont Dicéopolis vient de dire : « Heureux celui qui t'épousera et te procréera des belettes », la foule va ressentir quelque pulsion ne relevant pas uniquement de la cleptomanie.
\end{multicols}
}%
\end{adjustbox}%
\vfill
\end{adjustwidth}
\newpage
% --- PAGE DE DROITE : Traduction française + Commentaire (2/8) ---
\begin{adjustwidth}{-2cm}{-2cm}
\noindent\textbf{\footnotesize Traduction française --- Folio 22 (\texttt{f. 173r}, vv. 876–919 [1/4]) \hfill ΕΠΕΙΣΟΔΙΟΝ Δ' (Ἑλένη)}\\[-2pt]
\noindent\rule{16.5cm}{0.5pt}\\[2pt]
{\fontsize{8.8pt}{10.2pt}\selectfont
\ligneEst{}{\francais{Il est mort, Protéas, depuis dix ans.}}{876}
\ligneEst{\textbf{Euripide}}{\francais{Et quel est ce pays où ma nef a pris terre ?}}{}
\ligneEst{\textbf{Mnésiloque}}{\francais{L'Égypte.}}{}
\ligneEst{\textbf{Euripide}}{\francais{\espace{1.7cm}« Vers quel lieu, malheureux, nous avons navigué ! »}}{}
\ligneEst{\textbf{Kritylla}}{\francais{T'en fies-tu à ce coquin — la peste soit de lui —}}{}
\ligneEst{}{\francais{qui te raconte des bourdes ? Ceci, c'est le Thesmophorion.}}{880}
\ligneEst{\textbf{Euripide}}{\francais{Et lui, Protée, est-il céans, ou s'il est hors ?}}{}
\ligneEst{\textbf{Kritylla}}{\francais{Il faut que tu aies encore mal au c\oe ur, étranger.}}{}
\ligneEst{}{\francais{On te dit qu'il est mort, Protéas,}}{}
\ligneEst{}{\francais{et tu demandes s'il est céans ou dehors.}}{}
\ligneEst{\textbf{Euripide}}{\francais{Hélas! il a vécu. Où donc est le sépulcre où il fut inhumé ?}}{885}
\ligneEst{\textbf{Mnésiloque}}{\francais{C'est ici son tombeau, tu nous y vois assise.}}{}
}
\vspace{2pt}
\noindent\rule{16.5cm}{0.6pt}\\[2pt]
\noindent\textbf{\small Commentaire philologique, métrique et scénique (vv. 876–919 --- partie 2/8)}\\[2pt]
\begin{adjustbox}{max totalheight=15.2cm, max width=16.5cm, keepaspectratio}%
\parbox[t]{16.5cm}{
\setlength{\columnsep}{9pt}%
\begin{multicols}{2}
\fontsize{7.6pt}{8.8pt}\selectfont
\setlength{\parindent}{0pt}%
\setlength{\parskip}{1.2pt}%
\noindent Le verbe περιτρώγειν paraît d'ailleurs bien impropre pour une image de vol à la tire\marqueur{2}.
\noindent Quant à κλοπῇ, Mnésiloque nous a déjà fait remarquer au vers 205 que le verbe κλέπτειν pouvait par euphémisme désigner une activité qui n'a de commun avec le vol que de se passer au plus noir de la nuit.\\ Terminons en indiquant que cette accusation de tentative de viol lors d'une réunion exclusivement féminine se retrouve ailleurs ; ces assemblées féminines à huis-clos et que la rumeur publique voyait comme des séances de débauche sexuelle avaient tout pour attiser la curiosité de quelque luxurieux. C'est à une telle tentative d'un certain Critias qu'un fragment d'Hipponax\marqueur{3} fait allusion :
\citComm{οὔ μοι δικαίως μοιχὸς ἁλῶναι δοκέει\\
Κριτίης ὁ Χῖος ἐν $ \cup - \times - $ δούμῳ\\
Ce n'est pas selon moi par justice\marqueur{4} qu'en tant que débauché fût saisi\\
Critias de Khios dans $[$ … $]$ une assemblée féminine.}
\noindent\textbf{[Notes ad 894]} \textbf{(1)} Hésychius (χ 792) : χρυσίον· τὸ τῶν παιδίων αἰδοῖον. Henderson (\italique{Maculate Muse}, p. 131) y flaire un jeu de mot entre τὸν χρυσόν et τὸν κυσόν.
\noindent \textbf{(2)} Si impropre que l'on en serait tenté de comprendre λαβών à la place de λαθών.
\noindent \textbf{(3)} Fragment 30 de l'édition au fort utile commentaire d'Olivier Masson ($=$ fragment 42 de l'édition Degani). Je reprends la correction proposée par Masson : δούμῳ à la place de δούλῳ, qui s'appuie sur une correction identique chez Hésychius (δ 2256) :
\citComm{δοῦμος· ἡ οἰκία, ἢ τὴν ἐπὶ τὸ αὐτὸ συνέλευσιν τῶν γυναικῶν,\\
doumos : la maison. \italique{Item} : une réunion en commun des femmes.}
\noindent \textbf{(4)} Je lis : οὔ μοι δίκᾳ ὡς μοιχὸς ἁλῶναι δοκέει κτλ.
\noindent\textbf{895} \textbf{βαϋζα τοὐμὸν σῶμα βάλλουσα ψόγῳ.} Taillardat (§569) avait observé que le tour βάλλειν τινὰ ψόγῳ, « accabler quelqu'un de reproches » est d'un ton tragique fait sur le modèle de βάλλειν τινὰ λίθοις\marqueur{1}. (Et réciproquement, les mots sont bien souvent considérés comme des projectiles, cf. Taillardat §502.)  De même, la périphrase τοὐμὸν σῶμα pour ἐγώ est d'un style élevé et particulièrement prisée d'Euripide\marqueur{2}.
\noindent La correction de la leçon du \italique{Ravennas} universellement admise, βαϋζα en βάϋζε, est inintelligible. (Les dictionnaires se sont d'ailleurs bien gardés de citer notre vers 895 à l'entrée βαΰζω.) Faisons l'hypothèse que la fin du mot transmis soit correcte. Le suffixe -υζᾰ a été productif de quelques noms de végétaux en grec\marqueur{3} ; la \italique{conyze}, κόνυζα convient parfaitement : les minuscules `κ' et `β' sont très semblables sur de nombreux manuscrits\marqueur{4} -- dont notre \italique{Ravennas}.
\end{multicols}
}%
\end{adjustbox}%
\vfill
\end{adjustwidth}
\newpage

% --- FOLIO 22 (f. 173r) : SPREAD 2.2 (Texte grec établi + Traduction française + Commentaire 3/8 & 4/8) ---
\clearpage\ifodd\value{page}\hbox{}\newpage\fi
\enTeteG{ΘΕΣΜΟΦΟΡΙΑΖΟΥΣΑΙ --- Texte grec établi (f. 173r, vv. 876–919 [2/4])}
\enTeteD{LES THESMOPHORIES --- Traduction française (f. 173r, vv. 876–919 [2/4])}
\begin{adjustwidth}{-2cm}{-2cm}
\noindent\textbf{\footnotesize Texte grec établi --- Folio 22 (\texttt{f. 173r}, vv. 876–919 [2/4]) \hfill ΕΠΕΙΣΟΔΙΟΝ Δ' (Ἑλένη)}\\[-2pt]
\noindent\rule{16.5cm}{0.5pt}\\[2pt]
{\fontsize{8.8pt}{10.2pt}\selectfont
\ligneEst{\textbf{Κρίτυλλα}}{κακῶς τ᾽ ἄρ᾽ ἐξόλοιο κἀξολεῖ γέ τοι,}{}
\ligneEst{}{ὅστις γε τολμᾷς σῆμα τὸν βωμὸν καλεῖν.}{}
\ligneEst{\textbf{Εὐριπίδης}}{τί δὴ σὺ θάσσεις τάσδε τυμβήρεις ἕδρας}{}
\ligneEst{}{φάρει καλυπτὸς ὦ ξένη;}{890}
\ligneEst{\textbf{Μνησίλοχος}}{\decalage{φάρει καλυπτὸς ὦ ξένη;}βιάζομαι}{}
\ligneEst{}{γάμοισι Πρωτέως παιδὶ συμμεῖξαι λέχος.}{}
\ligneEst{\textbf{Κρίτυλλα}}{τί ὦ κακόδαιμον ἐξαπατᾷς αὖ τὸν ξένον;}{}
\ligneEst{}{οὗτος πανουργῶν δεῦρ᾽ ἀνῆλθεν ὦ ξένε}{}
\ligneEst{}{ὡς τὰς γυναῖκας ἐπὶ κλοπῇ τοῦ χρυσίου.}{}
\ligneEst{\textbf{Μνησίλοχος}}{κόνυζα τοὐμὸν σῶμα βάλλουσα ψόγῳ.}{895}
\ligneEst{\textbf{Εὐριπίδης}}{ξένη τίς ἡ γραῦς ἡ κακορροθοῦσά σε;}{}
\ligneEst{\textbf{Μνησίλοχος}}{αὕτη Θεονόη Πρωτέως.}{}
}
\vspace{2pt}
\noindent\rule{16.5cm}{0.6pt}\\[2pt]
\noindent\textbf{\small Commentaire philologique, métrique et scénique (vv. 876–919 --- partie 3/8)}\\[2pt]
\begin{adjustbox}{max totalheight=15.2cm, max width=16.5cm, keepaspectratio}%
\parbox[t]{16.5cm}{
\setlength{\columnsep}{9pt}%
\begin{multicols}{2}
\fontsize{7.6pt}{8.8pt}\selectfont
\setlength{\parindent}{0pt}%
\setlength{\parskip}{1.2pt}%
\noindent Le sens s'éclaire ainsi : Mnésiloque se compare à une plante -- une fleur dira-t-on en français -- en proie aux jets de reproches de Critylla. C'est là le point : ψόγῳ à la fin du vers est mis πάρα προσδοκίαν pour λίθῳ. Or le proverbe grec « jeter une pierre à une herbe » (Apostolius, XIV, 72) :
\citComm{πόᾳ λίθον ἐκβάλλει· ἐπὶ τῶν τἀναντία τοῖς ἐπιβούλοις ποιούντων·\\
Il jette une pierre à une herbe\marqueur{5}.}
\noindent\textbf{[Notes ad 895]} \textbf{(1)} Βάλλειν κακοῖς (Sophocle, \italique{Ajax} 1244), αἰτίᾳ (\italique{id}. \italique{Trach.} 940), φθόνῳ (Eur. \italique{Élect.} 902). Cf. ἀτιμίῃσιν ἰάλλειν (Homère, \italique{Odyssée} XIII, 142), λόγοις ἰάπτειν (Sophocle, \italique{Ajax} 501).
\noindent \textbf{(2)} Voir le \LSJ \ à l'entrée σῶμα II.
\noindent \textbf{(3)} ἄλυζα $=$ ἄλυπον ; κνύζα	 $=$ κόνυζα $=$ σκόνυζα, \italique{conyze} ou \italique{herbe aux puces} ou \italique{encensière} (\latin{erigeron viscosum L.}) ; ἡ μώλυζα, \italique{la tête d'ail} ; ἡ ὄρυζα, \italique{le riz}.
\noindent \textbf{(4)} Ex. : dans les \italique{Euménides} d'Eschyle, lire νεβρόν au lieu de νεκρόν v. 246.
\noindent \textbf{(5)} L'explication qui suit de Michel Apostolius devrait plutôt nous conduire à traduire ainsi : « Il retire une pierre à l'aide d'une herbe. » Mais la tentative d'Apostolius de trouver une exégèse savante au proverbe est vouée à l'échec -- comme pour tout ce qui touche à la tradition orale. Le rapprochement que fait Apostolius avec un passage du précieux Élien (\italique{De Natura Animalium}, I, 45), pour astucieux qu'il soit, n'est pas convaincant.
\noindent\textbf{896} \textbf{κακορροθοῦσά.} Le vocabulaire est euripidéen -- du véritable Euripide s'entend.
\noindent Le verbe κακορροθεῖν, substitut expressif de κακολογεῖν, est un terme qu'on retrouve exclusivement dans ses tragédies -- sauf reprise aristophanesque (\italique{Acharniens}, v. 576).
\noindent\textbf{898} \textbf{εἰ μὴ Κρίτυλλά γ᾽ Ἀντιθέου Γαργηττόθεν·} Je signale à l'attention d'un éventuel exégète parcourant ces pages que la réplique de Critylla trouve un écho dans un fragment de la comédie perdue d'Aristophane des Δαιταλεῖς, \italique{Convives}, transmise par Pollux (X, 158 $=$ fr. 226(217) Kassel-Austin) :
\citComm{εἰ μὴ δικῶν τε γυργάθους ψηφισμάτων τε θωμοὺς\\
< φέροντες >}
\noindent Le rapport est flagrant : même formule commençant par «~εἰ μὴ~» ; Κρίτ-υλλά, nom formé du suffixe -υλλα\marqueur{1} et du radical κρίτ- qui fait entendre le \italique{jugement} a le même champs sémantique que les δικῶν ; Γαργηττόθεν et γυργάθους sont très proches phonétiquement ainsi que -τιθέου (de Ἀντιθέου) et τε θωμοὺς.\\ Athénée (I, 4d) évoque ce même fragment d'Aristophane dans son premier livre ; ce qui nous permet d'en préciser le contexte. Dans un chapitre assez décousu\marqueur{2} consacré aux propos de table, on comprend qu'il est certainement question de sycophantes -- pique-assiettes et autres croque-lardons --, dont la vie est douce, bonne et grasse :
\citComm{τὸν βίον, φησίν, εὐσταθεῖς· οὐκ ἐγχειρογάστορες, γυργάθους ψηφισμάτων φέροντες, Ἀριστοφάνης φησίν}
\end{multicols}
}%
\end{adjustbox}%
\vfill
\end{adjustwidth}
\newpage
% --- PAGE DE DROITE : Traduction française + Commentaire (4/8) ---
\begin{adjustwidth}{-2cm}{-2cm}
\noindent\textbf{\footnotesize Traduction française --- Folio 22 (\texttt{f. 173r}, vv. 876–919 [2/4]) \hfill ΕΠΕΙΣΟΔΙΟΝ Δ' (Ἑλένη)}\\[-2pt]
\noindent\rule{16.5cm}{0.5pt}\\[2pt]
{\fontsize{8.8pt}{10.2pt}\selectfont
\ligneEst{\textbf{Kritylla}}{\francais{Puisses-tu périr misérablement ! Et tu périras,}}{}
\ligneEst{}{\francais{toi qui oses appeler l'autel une tombe.}}{}
\ligneEst{\textbf{Euripide}}{\francais{D'où vient qu'assise sur ce siège sépulcral,}}{}
\ligneEst{}{\francais{tu dérobes tes traits sous un voile, étrangère ?}}{890}
\ligneEst{\textbf{Mnésiloque}}{\francais{\espace{3.8cm}On prétend sans mon v\oe u}}{}
\ligneEst{}{\francais{que le fils de Protée me reçoive en son lit.}}{}
\ligneEst{\textbf{Kritylla}}{\francais{Pourquoi, misérable, tromper encore cet étranger ?}}{}
\ligneEst{}{\francais{Il est entré en fraude, ô étranger,}}{}
\ligneEst{}{\francais{auprès des femmes afin de toucher subrepticement à leur bijou.}}{}
\ligneEst{\textbf{Mnésiloque}}{\francais{Je suis la conyze qu'on accable d'un blâme.}}{895}
\ligneEst{\textbf{Euripide}}{\francais{Étrangère, dis-moi, qui donc est cette vieille, qui t'incrimine ainsi ?}}{}
\ligneEst{\textbf{Mnésiloque}}{\francais{Elle ? C'est Théonoé, la fille de Protée.}}{}
}
\vspace{2pt}
\noindent\rule{16.5cm}{0.6pt}\\[2pt]
\noindent\textbf{\small Commentaire philologique, métrique et scénique (vv. 876–919 --- partie 4/8)}\\[2pt]
\begin{adjustbox}{max totalheight=15.2cm, max width=16.5cm, keepaspectratio}%
\parbox[t]{16.5cm}{
\setlength{\columnsep}{9pt}%
\begin{multicols}{2}
\fontsize{7.6pt}{8.8pt}\selectfont
\setlength{\parindent}{0pt}%
\setlength{\parskip}{1.2pt}%
\noindent Comme l'oreille grecque percevait comme particulièrement proches les sons `γ' et `κ', faut-il entendre la locution grecque :
\citComm{οὐκ ... εἰ μὴ τὴν δίκην τις θεῶν κατώρθωσεν,\\
À moins qu'un dieu ne rétablisse la justice,}
\noindent conservée par Isocrate\marqueur{3} et auquel Aristophane fait vraisemblablement allusion en \italique{Cavaliers} 1360.
\noindent\textbf{898} \textbf{Γαργηττόθεν·} La scholie :
\citComm{ἀπὸ δήμου τῆς Αἰγηίδος φυλῆς, ἔνθα τὸ Εὐρυσθέως ἀπόκειται σῶμα ἐν τῷ Ἀρατηρίῳ. ἐπίρρημα δὲ τὸ Γαργηττόθεν.\\
D'un dème de la tribu Égéide, c'est là que repose le corps d'Eurysthée dans l'Arêtérion. D'où le nom de `Gargêttothen'.}
\noindent\textbf{[Notes ad 898]} \textbf{(1)} Le suffixe en -υλλα/ιλλα est courant pour les noms de femme : Κύδιλλα, Ἥρυλλα, ou Ξένυλλα qui est apparue plus haut \&c. (voir \italique{La formation des noms en grec ancien} de Chantraine, p. 252.)
\noindent \textbf{(2)} Le premier livre, duquel le passage est tiré, ne nous est connu que par l'\italique{épitomé} des \italique{Deipnosophistes}, ce qui complexifie encore la délimitation et l'attribution des citations... Voir l'article de Benoît Louyest, «~L'épitomé du \italique{Banquet des sophistes} d'Athénée~», in \italique{Rursus} 8 (2012).
\noindent \textbf{(3)} Isocrate, \italique{Orationes}, V (Φίλιππος), 152 ; XV (Περὶ ἀντιδόσεως), 105.
\noindent\textbf{903} \textbf{αἰσχύνομαί σε τὰς γνάθους ὑβρισμένη.} Rappelons que lors de la séance de travestissement, le pauvre et anciennement barbu Mnésiloque s'était fait flamber et raser la barbe. Sur un fameux cratère\marqueur{1} représentant la scène d'égorgement de la `petite' de Mika, on le voit en porter de très visibles stigmates. Et le scholiaste de commenter ce vers :
\citComm{ἐπειδὴ ξυρηθεὶς ἦν,\\
puisqu'il a été rasé.}
\noindent\textbf{[Notes ad 903]} \textbf{(1)} Le « cratère de Würzburg », Martin-von-Wagner-Museum der Universitàt Wurzburg, H 5697. Reproduit en annexe ; cf. la note aux vers 754-757.
\noindent\textbf{904-5} \textbf{} \italique{Hélène}, v. 549 :
\noindent « \difference{ἔκπληξιν ἡμῖν} ἀφασίαν \difference{τε προστίθης} » ; v. 557 : « τίς εἶ; τίν᾽ ὄψιν \difference{σήν}, γύναι, \difference{προσδέρκομαι}; » ; v. 72 : « ὦ θεοί, τίν᾽ \difference{εἶδον} ὄψιν; \difference{ἐχθίστην ὁρῶ} ».
\noindent\textbf{906-12} \textbf{} \italique{Hélène}, v. 558 sq. :
\citComm{Ἑλένη --- σὺ δ᾽ εἶ τίς; αὑτὸς γὰρ σὲ κἄμ᾽ ἔχει λόγος.\\
\difference{Μενελέως --- οὐπώποτ᾽ εἶδον προσφερέστερον δέμας.\\
Ἑλένη --- ὦ θεοί· θεὸς γὰρ καὶ τὸ γιγνώσκειν φίλους.}\\
< Μενελέως --- Ἑλληνὶς εἶ τις ἢ ἐπιχωρία γυνή; > (Ce vers a été restitué à la tragédie d'après Aristophane.)\\
Ἑλένη --- Ἑλληνίς· ἀλλὰ καὶ τὸ σὸν θέλω μαθεῖν.\\
Μενελέως --- Ἑλένῃ σ᾽ ὁμοίαν δὴ μάλιστ᾽ εἶδον, γύναι.\\
Ἑλένη --- ἐγὼ δὲ Μενέλεῴ γε σέ· \difference{οὐδ᾽ ἔχω τί φῶ}.\\
Μενελέως --- ἔγνως γὰρ ὀρθῶς ἄνδρα δυστυχέστατον.\\
Ἑλένη --- ὦ χρόνιος ἐλθὼν σῆς δάμαρτος \difference{ἐς χέρας}.}
\end{multicols}
}%
\end{adjustbox}%
\vfill
\end{adjustwidth}
\newpage

% --- FOLIO 22 (f. 173r) : SPREAD 2.3 (Texte grec établi + Traduction française + Commentaire 5/8 & 6/8) ---
\clearpage\ifodd\value{page}\hbox{}\newpage\fi
\enTeteG{ΘΕΣΜΟΦΟΡΙΑΖΟΥΣΑΙ --- Texte grec établi (f. 173r, vv. 876–919 [3/4])}
\enTeteD{LES THESMOPHORIES --- Traduction française (f. 173r, vv. 876–919 [3/4])}
\begin{adjustwidth}{-2cm}{-2cm}
\noindent\textbf{\footnotesize Texte grec établi --- Folio 22 (\texttt{f. 173r}, vv. 876–919 [3/4]) \hfill ΕΠΕΙΣΟΔΙΟΝ Δ' (Ἑλένη)}\\[-2pt]
\noindent\rule{16.5cm}{0.5pt}\\[2pt]
{\fontsize{8.8pt}{10.2pt}\selectfont
\ligneEst{\textbf{Κρίτυλλα}}{\decalage{αὕτη Θεονόη Πρωτέως.}μὰ τὼ θεὼ}{}
\ligneEst{}{εἰ μὴ Κρίτυλλά γ᾽ Ἀντιθέου Γαργηττόθεν·}{}
\ligneEst{}{σὺ δ᾽ εἶ πανοῦργος.}{}
\ligneEst{\textbf{Μνησίλοχος}}{\decalage{σὺ δ᾽ εἶ πανοῦργος.}ὁπόσα τοι βούλει λέγε.}{}
\ligneEst{}{οὐ γὰρ γαμοῦμαι σῷ κασιγνήτῳ ποτὲ}{900}
\ligneEst{}{προδοῦσα Μενέλεων τὸν ἐμὸν ἐν Τροίᾳ πόσιν.}{}
\ligneEst{\textbf{Εὐριπίδης}}{γύναι τί εἶπας; στρέψον ἀνταυγεῖς κόρας.}{}
\ligneEst{\textbf{Μνησίλοχος}}{αἰσχύνομαί σε τὰς γνάθους ὑβρισμένη.}{}
\ligneEst{\textbf{Εὐριπίδης}}{τουτὶ τί ἔστιν; ἀφασία τίς τοί μ᾽ ἔχει.}{}
\ligneEst{}{ὦ θεοὶ τίν᾽ ὄψιν εἰσορῶ; τίς εἶ γύναι;}{905}
\ligneEst{\textbf{Μνησίλοχος}}{σὺ δ᾽ εἶ τίς; αὑτὸς γὰρ σὲ κἄμ᾽ ἔχει λόγος.}{}
\ligneEst{\textbf{Εὐριπίδης}}{Ἑλληνὶς εἶ τις ἢ 'πιχωρία γυνή;}{}
}
\vspace{2pt}
\noindent\rule{16.5cm}{0.6pt}\\[2pt]
\noindent\textbf{\small Commentaire philologique, métrique et scénique (vv. 876–919 --- partie 5/8)}\\[2pt]
\begin{adjustbox}{max totalheight=15.2cm, max width=16.5cm, keepaspectratio}%
\parbox[t]{16.5cm}{
\setlength{\columnsep}{9pt}%
\begin{multicols}{2}
\fontsize{7.6pt}{8.8pt}\selectfont
\setlength{\parindent}{0pt}%
\setlength{\parskip}{1.2pt}%
\noindent\textbf{906} \textbf{« C'est une question qui m'appartient aussi ».} J'emprunte ce vers à A. de Musset, \italique{À quoi rêvent...}, acte II, scène 3. (Note de Willems.)
\noindent\textbf{910} \textbf{ἐγὼ δὲ Μενέλαον ὅσα γ᾽ ἐκ τῶν ἰφύων.} Scène de reconnaissance. Lorsqu'elle avait reconnu son vieux mari face à elle, la véritable Hélène d'Euripide s'était exclamé (v. 564) :
\citComm{Ἐγὼ δὲ Μενέλεῴ γε σέ· οὐδ' ἔχω τί φῶ.\\
Nul ne fut plus pareil à Ménélas. Que dire ?\\
(traduction d'Henri Grégoire.)}
\noindent Chez Aristophane, à quoi Hélène reconnaît-elle son naufragé de mari ? « Οὐδ' ἔχω τί φῶ » devient dans le manuscrit de Ravenne : ὅσα γ᾽ ἐκ τῶν ἀφύων, « d'après ton menu fretin ».
\noindent On lit dans le lexique de la \italique{Souda} (ι 778) : ὄσον ἐκ τῶν ἰφύων, « à tes lavandes\marqueur{1} ».
\noindent (C'est cette dernière version que commente le scholiaste, recopié en marge du manuscrit de Ravenne.) Qu'il s'agisse de \italique{menu fretin} ou de \italique{lavande}, de telles leçons sont ineptes\marqueur{2} !
\noindent Colin Austin a bien résumé la situation\marqueur{3} :
\noindent « L'alpha d'ἀφύη étant bref, cette leçon est métriquement impossible. Quant à la ``lavande'' de la \italique{Souda}, le scholiaste y voit une allusion aux herbes de la mère d'Euripide, ce qui semble bien tiré par les cheveux\marqueur{4} ! Comme le note H. Grégoire\marqueur{5} « les herbes sont ici plus ridicules encore que les anchois », moins ridicules cependant que les \italique{oignons marins} de Gag***\marqueur{6} (τιφύων, « \italique{your squills}... \anglais{here a metaphor for male genitals} ») : cette saloperie-là\marqueur{7} se réfute d'elle-même. En fait la solution du problème se trouve au v. 935. Là Euripide est accusé d'être un ἱστιορράφος, \italique{un ravaudeur de voiles}. Fritzsche (p. 363)\marqueur{8} avait bien vu qu'Aristophane se moque ici de l'accoutrement de Ménélas dans l'\italique{Hélène} d'Euripide: tout au long de cette pièce il est fait mention des haillons, des lambeaux de voile dont le roi de Sparte est pitoyablement affublé, comme un gueux, après son naufrage en Égypte (v. 416 τὰς ἐμὰς δυσχλαινίας, 422 ναὸς ἐκβόλοις ἃ ἀμπίσχομαι, 544 sq. ἄγριος δέ τις / μορφὴν ὅδ᾽ ἐστίν, 554 στολήν γ᾽ ἄμορφον, 790 ὥσπερ πτωχὸς, 1204 ὡς ἐσθῆτι δυσμόρφῳ πρέπει, 1282 τῆς ἀχλαινίας, 1297 ἐσθῆτά τ᾽ ἐξάλλαξον). En particulier au ν. 423 il s'écrie: « Mes vêtements d'antan, mes habits somptueux, mes anciennes splendeurs, la vague m'a tout pris », πέπλους δὲ τοὺς πρὶν λαμπρά τ᾽ ἀμφιβλήματα / χλιδάς τε πόντος ἥρπασ᾽, et au ν. 1079 « les haillons qui m'affublent d'eux-mêmes parleront, en témoins du naufrage », καὶ μὴν τάδ᾽ ἀμφίβληστρα σώματος ῥάκη / ξυμμαρτυρήσει ναυτικῶν ἐρειπίων. Dans les \italique{Grenouilles} (ν. 842) Eschyle a bien raison d'appeler Euripide \italique{un faiseur de mendiants}, \italique{un ravaudeur de guenille}, καὶ πτωχοποιὲ καὶ ῥακιοσυρραπτάδη. Pour en revenir au v. 910 des « Thesmo », on est bien tenté de croire que c'est à ses \italique{hardes} que le pauvre naufragé est reconnu par sa femme. »\\  Grégoire\marqueur{9} proposa de lire : `ἐκ τῶν ἀμφίων' ; τὸ ἄμφιον étant une abréviation de ἀμφίεσμα, le \italique{vêtement}. Cette proposition expliquerait la présence des ἀφύων, par simple chute d'une lettre suivie d'iotacisme. En revanche, elle passe à l'as la leçon `ἰφύων' conservée par le scholiaste du manuscrit de Ravenne et par Suidas (ι 778), témoins tout aussi dignes de confiance.
\end{multicols}
}%
\end{adjustbox}%
\vfill
\end{adjustwidth}
\newpage
% --- PAGE DE DROITE : Traduction française + Commentaire (6/8) ---
\begin{adjustwidth}{-2cm}{-2cm}
\noindent\textbf{\footnotesize Traduction française --- Folio 22 (\texttt{f. 173r}, vv. 876–919 [3/4]) \hfill ΕΠΕΙΣΟΔΙΟΝ Δ' (Ἑλένη)}\\[-2pt]
\noindent\rule{16.5cm}{0.5pt}\\[2pt]
{\fontsize{8.8pt}{10.2pt}\selectfont
\ligneEst{\textbf{Kritylla}}{\francais{\espace{3.6cm}Non, par les déesses,}}{}
\ligneEst{}{\francais{je ne suis autre que Critylla, fille d'Antithéos, du dème de Gargettos.}}{}
\ligneEst{}{\francais{Toi, tu es un fripon.}}{}
\ligneEst{\textbf{Mnésiloque}}{\francais{\espace{3cm}Dis ce que tu voudras.}}{}
\ligneEst{}{\francais{En épousant ton frère, je ne trahirai point}}{900}
\ligneEst{}{\francais{Ménélas, mon époux, qui reste à Troie.}}{}
\ligneEst{\textbf{Euripide}}{\francais{Qu'as-tu dit ? Femme, tourne vers moi ta brillante prunelle.}}{}
\ligneEst{\textbf{Mnésiloque}}{\francais{J'ai honte, après l'outrage imprimé sur mes joues.}}{}
\ligneEst{\textbf{Euripide}}{\francais{Qu'est-ce là ? Une stupeur me prend.}}{}
\ligneEst{}{\francais{Ô dieux, que vois-je ? Femme, qui donc es-tu?}}{905}
\ligneEst{\textbf{Mnésiloque}}{\francais{Qui donc es-tu toi-même ? « C'est une question qui m'appartient aussi ».}}{}
\ligneEst{\textbf{Euripide}}{\francais{« Es-tu Hellène, ou bien es-tu de ce pays ? »}}{}
}
\vspace{2pt}
\noindent\rule{16.5cm}{0.6pt}\\[2pt]
\noindent\textbf{\small Commentaire philologique, métrique et scénique (vv. 876–919 --- partie 6/8)}\\[2pt]
\begin{adjustbox}{max totalheight=15.2cm, max width=16.5cm, keepaspectratio}%
\parbox[t]{16.5cm}{
\setlength{\columnsep}{9pt}%
\begin{multicols}{2}
\fontsize{8.4pt}{9.7pt}\selectfont
\setlength{\parindent}{0pt}%
\setlength{\parskip}{1.2pt}%
\noindent (Re-)prends ton habit Grégoire ! Mais creusons son idée, qui est la bonne ; il nous suffit de trouver un synonyme de vêtement (avec, de préférence, une connotation péjorative) qui puisse rendre compte des deux leçons\marqueur{10}.
\noindent Ce bon vieil Hésychius ne nous ayant jamais trahi, c'est encore à lui que nous avons posé la question. Cet affable vieillard n'a pas tardé à répondre à notre requête (λ 171) :
\citComm{λαίφια· ῥάκη.\\
Laiphia : hardes.}
\noindent Hésychius nous confirme l'existence d'un diminutif de λαῖφος, \italique{vêtement en lambeaux}, \italique{tissu}, \italique{voile}, ainsi que son sens.\\ Tentons à présent une explication paléographique aux ἀφύων et autres ἰφύων : si l'on fait la simple hypothèse de la disparition de la première lettre du mot λαιφίων\marqueur{11}, le lambda, le scribe se retrouvait alors devant l'inconnu αιφιων, qu'il lisait indifféremment – iotacisme oblige – αιφυων. En butte à ce mot qui n'existe pas, il pouvait alors raisonnablement écrire, selon l'inspiration du moment, ἀφύων ou ἰφύων.
\noindent\textbf{[Notes ad 910]} \textbf{(1)} Suzanne Amigues identifiait en fait l'ἴφυον avec la nielle des blés (\latin{Agrostemma githago L.}). Voir son article « Quelques légumes de disette chez Aristophane et Plutarque » paru en 1988 dans le \italique{Journal des Savants}, pp. 157-171.
\noindent \textbf{(2)} Il se trouve néanmoins de nombreux défenseurs de la leçon « ἐκ τῶν ἰφύων », leçon qui fait beaucoup rire de savants hellénistes (?!) – décidément très bon public !
\noindent \textbf{(3)} «~Observations critiques sur les \italique{Thesmophories} d'Aristophane~», in \italique{Dodone} 19.2 (1990), p. 27.
\noindent \textbf{(4)} C'est pourtant cette leçon « bien tirée par les cheveux » qu'Austin fit imprimer dans son édition des \italique{Thesmophories} ; il faut croire que, ici comme en d'autres occasions, son acolyte Olson fut de mauvais conseil.
\noindent \textbf{(5)} Tome V de l'édition d'Euripide aux Belles Lettres (1950), p. 65 n. 1.
\noindent \textbf{(6)} « Une telle saloperie ?.. ou plutôt un Gag, non ? » C'est ce mot d'ordre que je reprends à Austin, modifiant au passage le nom – \latin{nomen omen} ! – du farceur qui va recevoir sa fessée, conformément à ce que réclame la justice !
\end{multicols}
}%
\end{adjustbox}%
\vfill
\end{adjustwidth}
\newpage

% --- FOLIO 22 (f. 173r) : SPREAD 2.4 (Texte grec établi + Traduction française + Commentaire 7/8 & 8/8) ---
\clearpage\ifodd\value{page}\hbox{}\newpage\fi
\enTeteG{ΘΕΣΜΟΦΟΡΙΑΖΟΥΣΑΙ --- Texte grec établi (f. 173r, vv. 876–919 [4/4])}
\enTeteD{LES THESMOPHORIES --- Traduction française (f. 173r, vv. 876–919 [4/4])}
\begin{adjustwidth}{-2cm}{-2cm}
\noindent\textbf{\footnotesize Texte grec établi --- Folio 22 (\texttt{f. 173r}, vv. 876–919 [4/4]) \hfill ΕΠΕΙΣΟΔΙΟΝ Δ' (Ἑλένη)}\\[-2pt]
\noindent\rule{16.5cm}{0.5pt}\\[2pt]
{\fontsize{8.8pt}{10.2pt}\selectfont
\ligneEst{\textbf{Μνησίλοχος}}{Ἑλληνίς. ἀλλὰ καὶ τὸ σὸν θέλω μαθεῖν.}{}
\ligneEst{\textbf{Εὐριπίδης}}{Ἑλένῃ σ᾽ ὁμοίαν δὴ μάλιστ᾽ εἶδον γύναι.}{}
\ligneEst{\textbf{Μνησίλοχος}}{ἐγὼ δὲ Μενελέῳ σ᾽ ὅσα γ᾽ ἐκ τῶν λαιφίων.}{910}
\ligneEst{\textbf{Εὐριπίδης}}{ἔγνως ἄρ᾽ ὀρθῶς ἄνδρα δυστυχέστατον.}{}
\ligneEst{\textbf{Μνησίλοχος}}{ὦ χρόνιος ἐλθὼν σῆς δάμαρτος ἐσχάρας}{}
\ligneEst{{\tiny\textit{2. dochm.}}}{λαβέ με λαβέ με πόσι, περίβαλε δὲ χέρας.}{}
\ligneEst{}{φέρε σὲ κύσω. ἄπαγέ μ᾽ ἄπαγ᾽ ἄπαγέ με}{915}
\ligneEst{{\tiny\textit{trim. iamb.}}}{λαβὼν ταχὺ πάνυ.}{}
\ligneEst{\textbf{Κρίτυλλα}}{\decalage{λαβὼν ταχὺ πάνυ.}κλαύσετ᾽ ἄρα νὴ τὼ θεὼ}{}
\ligneEst{}{ὅστις σ᾽ ἀπάξει τυπτόμενος τῇ λαμπάδι.}{}
\ligneEst{\textbf{Εὐριπίδης}}{σὺ τὴν ἐμὴν γυναῖκα κωλύεις ἐμέ,}{}
\ligneEst{}{τὴν Τυνδάρειον παῖδ᾽, ἐπὶ Σπάρτην ἄγειν;}{919}
}
\vspace{2pt}
\noindent\rule{16.5cm}{0.6pt}\\[2pt]
\noindent\textbf{\small Commentaire philologique, métrique et scénique (vv. 876–919 --- partie 7/8)}\\[2pt]
\begin{adjustbox}{max totalheight=15.2cm, max width=16.5cm, keepaspectratio}%
\parbox[t]{16.5cm}{
\setlength{\columnsep}{9pt}%
\begin{multicols}{2}
\fontsize{7.6pt}{8.8pt}\selectfont
\setlength{\parindent}{0pt}%
\setlength{\parskip}{1.2pt}%
\noindent \textbf{(7)} Pour réussir à faire perdre son sang-froid au pourtant toujours modéré Colin Austin, il fallait effectivement que ce fut une belle salauderie… Nous aurions volontiers laissé la misérable hypo-foutaise « qui se réfute d'elle-même » de Gag*** dans un oubli bien mérité, malheureusement, il nous a été infligé de la voir exhumée et prise au sérieux dans plusieurs articles récents…  Enterrons une bonne fois pour toute -- espérons-le ! -- cette charogne : la saloperie de Gag*** présente à la fois le quadruple exploit \ \ \ 1°) \ de ne pas rendre compte de la leçon `ἀφύων' du témoin principal de la pièce : le manuscrit de Ravenne ; \ \ \ 2°) \ d'élucubrer sur un soi-disant sens obscène de τίφυον qui n'est attesté nulle part ; \ \ \ 3°) \ d'être intrinsèquement absurde : si Ménélas est réputé avoir bien peu d'ardeur (raison du départ de l'insatisfaite madame ?) ; comment se fait-il qu'il soit en pleine \italique{érection} et que ce soit à cela qu'Hélène le reconnaisse ? \ \ \ 4°) \ \italique{last but not least}, d'être anti-comique à souhait : un cocu impuissant n'est drôle \italique{que} parce qu'il est impuissant : représentez-le en pleine érection et vous gâterez à coup sûr la farce ! \ \ \ \ Reconnaissons au moins à Gag*** le mérite (pédagogique pour l'édification du lecteur) d'avoir réussi à concentrer en un seul passage les quatre vices cardinaux du \italique{misthophore} de la recherche, savant homme qui se paye – littéralement ! – de mots. Qu'y a-t-il d'étonnant à ce qu'un individu salarié pour palabrer palabre ? Rien, à priori, mais on leur saurait tout-de-même gré d'occuper leur jactance à autre chose que de salir Aristophane.
\noindent \textbf{(8)} Voir aussi G. Lange, \italique{Quaest. in Ar. Thesm., Diss. Gott}. (1891), p. 8.
\noindent \textbf{(9)} Au tome V de son édition d'Euripide aux Belles Lettres, déjà mentionnée. Mais son hypothèse est de douze ans plus ancienne : il l'avait présentée en anglais dans l'article « Modern Greek in Aristophanes » paru dans le premier numéro de \italique{The Link, A Review of Mediaeval and Modern Greek} (1938), pp. 16-20. Une anecdote qui montre comme le monde grec est petit : Henri Grégoire n'est autre que celui qui reprit la suite du grand Alphonse Willems à la chaire de philologie grecque de l'Université libre de Bruxelles en 1909, lorsque ce dernier dut prendre sa retraite.
\noindent \textbf{(10)} Il y avait une autre possibilité de recherche : trouver un mot dont les multiples sens auraient pu être glosés à la fois par ἴφυον et ἀφύη ; les gloses explicatives ayant fini par prendre la place de la leçon originale (le cas s'est déjà vu). Creusant dans cette voie, Austin \& Olson avaient suggéré une émendation très ingénieuse : φυκίων. Émendation ingénieuse sur le plan paléographique, beaucoup moins sur le plan comique…
\noindent \textbf{(11)} Erreur courante dans notre manuscrit de Ravenne. Citons pour exemple le vers 834 où le manuscrit comporte « δίδοσθαι \textbf{Τ}ηνίοισι » à la place de « δίδοσθαι \textbf{Στ}ηνίοισι ». Cette erreur est d'autant plus fréquente lorsque il y a deux lettres identiques consécutives, ainsi au vers 555 du manuscrit de Ravenne : « ὡς \textbf{τ}λεγγίδας λαβοῦσαι » en place de « ὡς \textbf{στ}λεγγίδας λαβοῦσαι ».
\noindent Une telle faute s'explique aisément si l'on imagine le manuscrit copié en onciales et où les mots n'étaient pas même espacés. Dans le mot « λαιφίων », les deux premières lettres, le `λ' et le `α' sont de formes semblables en onciales. La `chute' du `λ' me paraît donc une hypothèse loin d'être invraisemblable.
\end{multicols}
}%
\end{adjustbox}%
\vfill
\end{adjustwidth}
\newpage
% --- PAGE DE DROITE : Traduction française + Commentaire (8/8) ---
\begin{adjustwidth}{-2cm}{-2cm}
\noindent\textbf{\footnotesize Traduction française --- Folio 22 (\texttt{f. 173r}, vv. 876–919 [4/4]) \hfill ΕΠΕΙΣΟΔΙΟΝ Δ' (Ἑλένη)}\\[-2pt]
\noindent\rule{16.5cm}{0.5pt}\\[2pt]
{\fontsize{8.8pt}{10.2pt}\selectfont
\ligneEst{\textbf{Mnésiloque}}{\francais{« Hellène. Et à mon tour je veux savoir ta race. »}}{}
\ligneEst{\textbf{Euripide}}{\francais{« Ô femme, je te vois tout pareille à Hélène. »}}{}
\ligneEst{\textbf{Mnésiloque}}{\francais{« Toi-même à Ménélas, si j'en juge par ces hardes. »}}{910}
\ligneEst{\textbf{Euripide}}{\francais{« Oui, tu l'as reconnu, ce grand infortuné. »}}{}
\ligneEst{\textbf{Mnésiloque}}{\francais{« Ô toi, le tard venu dans les bas de ta femme »,}}{}
\ligneEst{}{\francais{prends, prends-moi, mon époux. Presse-moi dans tes bras.}}{}
\ligneEst{}{\francais{Voyons, que je te baise. Emmène-moi, emmène, emmène-moi, et au plus vite.}}{915}
\ligneEst{}{\francais{(Il descend de l'autel, les femmes l'entourent.)}}{}
\ligneEst{\textbf{Kritylla}}{\francais{\espace{2.8cm}Ah! il pleurera, par les deux déesses,}}{}
\ligneEst{}{\francais{celui qui t'emmènera. Je le frapperai de ma torche.}}{}
\ligneEst{\textbf{Euripide}}{\francais{Quoi ? mon épouse à moi, fille de Tyndare,}}{}
\ligneEst{}{\francais{tu voudrais m'empêcher de la conduire à Sparte ?}}{919}
}
\vspace{2pt}
\noindent\rule{16.5cm}{0.6pt}\\[2pt]
\noindent\textbf{\small Commentaire philologique, métrique et scénique (vv. 876–919 --- partie 8/8)}\\[2pt]
\begin{adjustbox}{max totalheight=15.2cm, max width=16.5cm, keepaspectratio}%
\parbox[t]{16.5cm}{
\setlength{\columnsep}{9pt}%
\begin{multicols}{2}
\fontsize{7.6pt}{8.8pt}\selectfont
\setlength{\parindent}{0pt}%
\setlength{\parskip}{1.2pt}%
\noindent\textbf{912} \textbf{ἐλθὼν σῆς δάμαρτος ἐσχάρας / λαβέ.} Le vers correspondant de l'\italique{Hélène} (le vers 566) finit par : ἐς χέρας ; Hélène reproche à Ménélas d'avoir trop tardé à rejoindre ses \italique{mains} (ses \italique{bras}, dira-t-on en français). Occasion par trop belle pour Aristophane de renouer avec la tradition populaire qui ne voit en Hélène qu'une pandorienne dévergondée, à l'origine de tant de maux : il fait donc prononcer à Mnésiloque `ἐσχάρας'. Le mot ἐσχάρα désigne proprement le \italique{foyer} ou le \italique{réchaud}\marqueur{1}. Ici comme ailleurs\marqueur{2}, le terme peut désigner le sexe de la femme. Jean Taillardat\marqueur{3} rapproche cette métaphore d'un passage de la \italique{Paix}\marqueur{4} où le mot ὀπτάνιον, \italique{cuisine} est pris comme synonyme de κύσθος.
\noindent\textbf{[Notes ad 912]} \textbf{(1)} \italique{Acharniens}, v. 888 ; \italique{Guêpes}, v. 938 ; \&c.
\noindent \textbf{(2)} \italique{Cavaliers} v. 1286 : κυκῶν τὰς ἐσχάρας avec le commentaire du scholiaste : τὰ χείλη τῶν γυναικείων αἰδοίων.
\noindent \textbf{(3)} \italique{Les images d'Aristophane}, §117.
\noindent \textbf{(4)} Vers 891.
\noindent\textbf{913-5} \textbf{λαβέ με λαβέ με πόσι, / περίβαλε δὲ χέρας. / φέρε σὲ κύσω. ἄπαγέ μ᾽ ἄπαγ᾽ ἄπαγέ με.} Les quatre dochmiaques sont tous intégralement résolus (et Lætitia\marqueur{1} note la «~\anglais{correption }κύσ\lettreBreve{ω}. ἄπαγε \anglais{harsh and unusual in dochmiacs}~») ; signe d'une accélération -- voire une précipitation en l'occurrence -- du rythme.\\ Aristophane nous présente Hélène pour ce qu'elle est : une nymphomane. (Loin donc de la pseudo-réhabilitation d'Euripide, plus risible que crédible.) Comment Hélène nous fait-elle entendre les besoins tyranniques qui la tiraillent ? par un procédé vieux comme le monde et qu'Aristophane emploie aussi au début des \italique{Cavaliers} (vers 26) : répéter un bout de phrase, en l'accélérant... et faire entendre ainsi un autre motif, bien différent de l'initial.
\noindent La luronne détachait-elle le `πάγε με' de la fin de son exhortation afin de suggérer le verbe πήγνυμι, \italique{ficher} ?\\ La source d'inspiration : \italique{Hélène}, v. 625 sq. :
\citComm{Ἑλένη --- ὦ φίλτατ᾽ ἀνδρῶν Μενέλεως, ὁ μὲν χρόνος\\
παλαιός, ἡ δὲ τέρψις ἀρτίως πάρα.\\
ἔλαβον ἀσμένα πόσιν ἐμόν, φίλαι,\\
περί τ᾽ ἐπέτασα χέρα\\
φίλιον ἐν μακρᾷ φλογὶ φαεσφόρῳ.\\
Μενελέως --- κἀγὼ σέ· πολλοὺς δ᾽ ἐν μέσῳ λόγους ἔχων\\
οὐκ οἶδ᾽ ὁποίου πρῶτον ἄρξωμαι τὰ νῦν.\\
Ἑλένη --- γέγηθα, κρατὶ δ᾽ ὀρθίους ἐθείρας\\
ἀνεπτέρωκα καὶ δάκρυ σταλάσσω,\\
περὶ δὲ γυῖα χέρας ἔβαλον, ἡδονάν,\\
ὦ πόσις, ὡς λάβω.}
\noindent\textbf{[Notes ad 913-5]} \textbf{(1)} \italique{The Songs of Aristophanes}, p. 429.
\end{multicols}
}%
\end{adjustbox}%
\vfill
\end{adjustwidth}
\newpage
\clearpage\ifodd\value{page}\hbox{}\newpage\fi
\enTeteG{ΕΠΕΙΣΟΔΙΟΝ Δ' / ΧΟΡΙΚΟΝ --- Codex Ravennas 429, f. 173v (Classense 175v) --- vv. 920–958}
\begin{adjustwidth}{-1.5cm}{-3.0cm}
\begin{adjustbox}{max totalheight=21.0cm, max width=16.5cm, keepaspectratio}
\begin{minipage}{16.5cm}
{\fontsize{7.2pt}{8.5pt}\selectfont
\noindent\textbf{Folio 23 --- Codex Ravennas 429, \texttt{f. 173v} (\texttt{Classense 175v}) : ΕΠΕΙΣΟΔΙΟΝ Δ' / ΧΟΡΙΚΟΝ, vv. 920–958}\hfill \textit{Témoins :} $\text{R}^1, \Rpc, \Rdeux$, \PsiDeux\par\vspace{2pt}
\hrule height 0.4pt\vspace{2.5pt}
}
{\fontsize{8.1pt}{9.5pt}\selectfont
\ligneRav{920}{\auDessus{γυ}{ν}}{οἴμ᾽ ὡς πανοῦργος καὐτὸς εἶναί μοι δοκεῖς}{}
\ligneRav{}{}{καὶ τοῦδέ τις ξύμβουλος. οὐκ ἐτὸς πάλαι}{}
\ligneRav{}{}{ᾐγυπτιάζετ᾽. ἀλλ᾽ ὅδε μὲν δώσει δίκην·}{}
\ligneRav{}{}{προσέρχεται γὰρ ὁ πρύτανις χὠ τοξότης·}{}
\ligneRav{}{\auDessus{εὐ}{ρ}}{τουτὶ πονηρόν· ἀλλ᾽ ὑπαποκινητέον·}{}
\ligneRav{925}{μ\auDessus{ν}{Η}}{ἐγὼ δ᾽ ὁ κακοδαίμων τί δρῶ; : μέν᾽ ἥσυχος·}{}
\ligneRav{}{}{οὐ γὰρ προδώσω σ᾽ οὐδέποτ᾽, ἤνπερ ἐμπνέω,}{}
\ligneRav{927}{}{ἢν μὴ \difference{προλίπωσ᾽} αἱ μυρίαι με μηχαναί·}{\textbf{927} προλείπωσ᾽ R : \lecon{προλίπωσ᾽} Scaliger}
\ligneRav{928}{μ\auDessus{ν}{Η}}{αὕτη μὲν ἡ μήρινθος οὐδὲν ἔσπασεν·}{\textbf{928} (Γυ.) R : \lecon{(Μνη.)} Scaliger}
\ligneRav{929}{πρυτ/}{ὅδ᾽ \difference{ἔσθ᾽} ὁ πανοῦργος ὃν \difference{ἔλεγ᾽} ἡμῖν Κλεισθένης;}{\textbf{929} ἔστ᾽ R : \lecon{ἔσθ᾽} \Mon \espApp ἔλεγεν R : \lecon{ἔλεγ᾽} \Mon}
\ligneRav{930}{}{οὗτος τί κύπτεις; δῆσον αὐτὸν εἰσάγων}{}
\ligneRav{}{}{ὦ τοξότ᾽ ἐν τῇ σανίδι, κἄπειτ᾽ ἐνθαδὶ}{}
\ligneRav{}{}{στήσας φύλαττε καὶ προσιέναι μηδένα}{}
\ligneRav{}{}{ἔα πρὸς αὐτόν, ἀλλὰ τὴν μάστιγ᾽ ἔχων}{}
\ligneRav{934}{—}{παῖ᾽ ἢν προσίῃ τις. : νὴ Δί᾽ ὡς \difference{νῦν δή γ᾽} ἀνὴρ}{\textbf{934} νῦν δῆτ' R : \lecon{νῦν δή γ᾽} Dobree : νυνδή γ' Cobet}
\ligneRav{935}{}{ὀλίγου μ᾽ ἀφείλετ᾽ αὐτὸν ἱστιορράφος·}{}
\ligneRav{}{μ\auDessus{ν}{Η}}{ὦ πρύτανι πρὸς τῆς δεξιᾶς, ἥνπερ φιλεῖς}{}
\ligneRav{}{}{κοίλην προτείνειν ἀργύριον ἤν τις διδῷ,}{}
\ligneRav{}{}{χάρισαι βραχύ τί μοι καίπερ ἀποθανουμένῳ·}{}
\ligneRav{939}{πρυτ/}{τί σοι \difference{χαρίσωμαι}; : γυμνὸν ἀποδύσαντά με,}{\textbf{939} χαρίσομαι R : \lecon{χαρίσωμαι} Casaubon}
\ligneRav{940}{}{κέλευε πρὸς τῇ σανίδι δεῖν τὸν τοξότην,}{}
\ligneRav{941}{}{ἵνα \difference{μὴ 'ν} κροκωτοῖς καὶ μίτραις γέρων ἀνὴρ}{\textbf{941} μὴν R : \lecon{μὴ 'ν} Küster}
\ligneRav{942}{}{μὴ \inRas{γε}{} τοῖς κροκωτοῖς καὶ τῇ μίτρᾳ}{\textbf{942} \inRas{γε}{} τοῖς \Rdeux}
\ligneRav{943}{πρυτ/}{ἔχοντα ταῦτ᾽ \difference{ἔδοξε} τῇ βουλῇ σε δεῖν,}{\textbf{943} ἔδοξεν R : \lecon{ἔδοξε} Biset}
\ligneRav{944}{}{ἵνα τοῖς παροῦσι δῆλος ᾖς πανοῦργος ὤν·}{\textbf{944} \lecon{παροῦσι} R : παριοῦσι Brunck}
\ligneRav{945}{μ\auDessus{ν}{Η}}{ἰαππαπαιάξ· ὦ κροκώθ᾽ οἷ᾽ εἴργασαι·}{\textbf{945} \lecon{ἰαππαπαιάξ} R, Suid. ε 2807 : ἰατταταιάξ Biset}
\ligneRav{946}{}{κοὐκ \difference{ἔστ᾽} ἔτ᾽ ἐλπὶς οὐδεμία σωτηρίας·}{\textbf{946} ἔστιν R : \lecon{ἔστ᾽} Bentley}
\ligneRav{947}{χορ/}{ἄγε νυν ἡμεῖς \difference{παίσωμεν} ἅπερ νόμος ἐνθάδε \difference{ταῖσι} γυναιξίν,}{\textbf{947} πέσωμεν R : \lecon{παίσωμεν} Zanetti, cf. Σ \espApp ταῖσ R : \lecon{ταῖσι} Daubuz}
\ligneRav{948}{}{παραχώρει ὅταν ὄργια σεμνὰ θεαῖν ἱεραῖς ὥραις ἀνέχωμεν, ἅπερ καὶ}{\textbf{948} \lecon{θεαῖν} R : θεοῖν Cobet}
\ligneRav{}{}{\hspace{1.1em}Παύσων σέβεται καὶ νηστεύει,}{}
\ligneRav{950}{}{\hspace{1.1em}πολλάκις αὐταῖν ἐκ τῶν ὡρῶν}{\textbf{950} \lecon{αὐταῖν} R : αὐτοῖν Cobet}
\ligneRav{}{}{\hspace{1.1em}ἐς τὰς ὥρας ξυνεπευχόμενος}{}
\ligneRav{952}{}{\hspace{1.1em}τοιαῦτα \difference{μέλειν} θάμ᾽ ἑαυτῷ·}{\textbf{952} μέλλειν R : \lecon{μέλειν} Zanetti (\italique{curae esse} Divus) : μολεῖν Casaubon}
\ligneRav{}{χορ/}{\hspace{1.1em}ὅρμα χώρει·}{}
\ligneRav{}{}{\hspace{1.1em}κοῦφα ποσὶν ἄγ᾽ ἐς κύκλον,}{}
\ligneRav{955}{}{\hspace{1.1em}χειρὶ σύναπτε χεῖρα,}{\textbf{955} \lecon{χειρὶ} R : χερὶ Dobree}
\ligneRav{956}{}{\hspace{1.1em}†ῥυθμὸν χορείας ὕπαγε πᾶσα†·}{\textbf{956} \lecon{ὕπαγε πᾶσα} R : πᾶσ' ὕπαγε Austin}
\ligneRav{956}{}{\hspace{1.1em}βαῖνε καρπαλίμοιν ποδοῖν·}{\textbf{956} \lecon{ὕπαγε πᾶσα} R : πᾶσ' ὕπαγε Austin}
\ligneRav{}{}{\hspace{1.1em}ἐπισκοπεῖν δὲ πανταχῇ}{}
\ligneRav{}{}{\hspace{1.1em}κυκλοῦσαν ὄμμα χρὴ χοροῦ κατάστασιν·}{}
}
\vspace{3pt}\noindent\rule{4.5cm}{0.4pt}\par\vspace{1.5pt}
{\fontsize{7.0pt}{8.3pt}\selectfont
\noindent\textbf{Loci paralleli (vv. 920–958) ---} \textbf{922} (ᾐγυπτιάζετ᾽) \SynB α 494 αἰγυπτιάζειν $=$ Photius α 518 $\subset$ Suidas αι 75 \espTest \textbf{926} (ἤνπερ ἐμπνέω) Suidas η 379 (ἤν pro ἤνπερ) \espTest \textbf{928} Suidas α 4475, μ 978 ; \Eustathe 760.44 ; scholium\expo{RV} \italique{Vespae} 175b \espTest \textbf{931} (σανίδι) Photius σ 68 σανίς, σ 65 σανίδα ; Hesychius σ 169 σανίς ; \AelDion σ 7 \espTest \textbf{935} Suidas ι 691 ὀλίγου ταύτην ἱστιορράφος \espTest (ἱστιορράφος) Pollux VII 161 \espTest \textbf{940} (πρὸς τῇ σανίδι δεῖν) Hesychius δ 926 δῆσαι πρὸς σανίδι· Ἀττικῶς· βασανίσαι ἐν τῷ δεσμωτηρίῳ \espTest \textbf{942} Suidas ε 3215 \espTest (ἑστιῶν) Hesychius ε 6411 ἑστιῶν· ἀριστοποιῶν $=$ Photius ε 2031 $=$ \Syn ε 875 $=$ Suidas ε 3227 ; cf. \Phryn 66.7 \espTest \textbf{945} (ἰαππαπαιάξ) Suidas ε 2807 \espTest \textbf{947} (ὅταν -- ἀνέχωμεν) Suidas α 2364 (θεαῖν codd. GFS\expo{a.c.} vel θεῶν codd. ATS\expo{p.c} vel θεεῖν cod. I vel θεαῖνων cod. M pro θεοῖν)\par
}
\vspace{2.5pt}\noindent\rule{4.5cm}{0.4pt}\par\vspace{1.5pt}
{\fontsize{6.9pt}{8.2pt}\selectfont
\noindent\textbf{Suite (en commentaire) de l'apparat critique \& notes paléographiques (\texttt{f. 173v}) ---} \textbf{Témoin papyrologique :} \textbf{\PsiDeux{} (\italique{P.Oxy.} LVI 3839, fr. 2)} couvre les \textbf{vv. 941–953}. --- \textbf{Grattage \italique{in rasura} au v. 942 :} \Rac{} avait inséré la particule γε devant τοῖς ; \textbf{\Rdeux{} l'a grattée (\inRas{γε}{} τοῖς)} pour rétablir le trimètre iambique. --- \textbf{Bêtacisme vocalique au v. 947 :} πέσωμεν R pour \lecon{παίσωμεν} (Zanetti, confirmé par $\Sigma^{\text{R}}$). --- \textbf{v. 927 :} προλείπωσ' R (iotacisme) : \lecon{προλίπωσ'} Scaliger. --- \textbf{v. 947 :} πέσωμεν R : \lecon{παίσωμεν} Zanetti, cf. $\Sigma$ \espApp ταῖσ R : \lecon{ταῖσι} Daubuz.\par
}
\end{minipage}
\end{adjustbox}
\end{adjustwidth}
\newpage
\enTeteD{Codex Ravennas 429 (Biblioteca Classense) --- Fac-similé du f. 173v (vv. 920–958)}
\begin{adjustwidth}{-1.5cm}{-3.0cm}
\begin{center}
\includegraphics[width=16.5cm,height=20.5cm,keepaspectratio]{/workspaces/Thesmophories/dossier_ravennas_thesmophories/folios_individuels_ff162v-177r/23_f173v_classense175v_vv0920-0958.jpg}\par
\vspace{2pt}
{\scriptsize\textbf{Codex Ravennas 429, f. 173v (Classense 175v)} --- ΕΠΕΙΣΟΔΙΟΝ Δ' / ΧΟΡΙΚΟΝ, \textbf{vv. 920–958} ($\text{R}^1, \Rpc, \Rdeux$, \PsiDeux).}
\end{center}
\end{adjustwidth}
\newpage
\clearpage\ifodd\value{page}\hbox{}\newpage\fi% --- FOLIO 23 (f. 173v) : SPREAD 2.1 (Texte grec établi + Traduction française + Commentaire 1/8 & 2/8) ---
\clearpage\ifodd\value{page}\hbox{}\newpage\fi
\enTeteG{ΘΕΣΜΟΦΟΡΙΑΖΟΥΣΑΙ --- Texte grec établi (f. 173v, vv. 920–958 [1/4])}
\enTeteD{LES THESMOPHORIES --- Traduction française (f. 173v, vv. 920–958 [1/4])}
\begin{adjustwidth}{-2cm}{-2cm}
\noindent\textbf{\footnotesize Texte grec établi --- Folio 23 (\texttt{f. 173v}, vv. 920–958 [1/4]) \hfill ΕΠΕΙΣΟΔΙΟΝ Δ' / ΧΟΡΙΚΟΝ}\\[-2pt]
\noindent\rule{16.5cm}{0.5pt}\\[2pt]
{\fontsize{8.8pt}{10.2pt}\selectfont
\ligneEst{\textbf{Κρίτυλλα}}{οἴμ᾽ ὡς πανοῦργος καὐτὸς εἶναί μοι δοκεῖς}{920}
\ligneEst{}{καὶ τοῦδέ τις ξύμβουλος. οὐκ ἐτὸς πάλαι}{}
\ligneEst{}{ᾐγυπτιάζετ᾽. ἀλλ᾽ ὅδε μὲν δώσει δίκην.}{}
\ligneEst{}{προσέρχεται γὰρ ὁ πρύτανις χὠ τοξότης.}{}
\ligneEst{\textbf{Εὐριπίδης}}{τουτὶ πονηρόν· ἀλλ᾽ ὑπαποκινητέον.}{}
\ligneEst{\textbf{Μνησίλοχος}}{ἐγὼ δ᾽ ὁ κακοδαίμων τί δρῶ;}{925}
\ligneEst{\textbf{Εὐριπίδης}}{\decalage{ἐγὼ δ᾽ ὁ κακοδαίμων τί δρῶ;}μέν᾽ ἥσυχος.}{}
\ligneEst{}{οὐ γὰρ προδώσω σ᾽ οὐδέποτ᾽, ἤνπερ ἐμπνέω,}{}
\ligneEst{}{ἢν μὴ προλίπωσ᾽ αἱ μυρίαι με μηχαναί.}{}
\ligneEst{\textbf{Μνησίλοχος}}{αὕτη μὲν ἡ μήρινθος οὐδὲν ἔσπασεν.}{}
}
\vspace{2pt}
\noindent\rule{16.5cm}{0.6pt}\\[2pt]
\noindent\textbf{\small Commentaire philologique, métrique et scénique (vv. 920–958 --- partie 1/8)}\\[2pt]
\begin{adjustbox}{max totalheight=15.2cm, max width=16.5cm, keepaspectratio}%
\parbox[t]{16.5cm}{
\setlength{\columnsep}{9pt}%
\begin{multicols}{2}
\fontsize{8.4pt}{9.7pt}\selectfont
\setlength{\parindent}{0pt}%
\setlength{\parskip}{1.2pt}%
\noindent\textbf{922} \textbf{ᾐγυπτιάζετ᾽.} Αἰγυπτιάζειν, faire l'Égyptien, comme Σιφνιάζειν, Χιάζειν, \italique{faire le Chiote}, Ἑλληνίζειν, \italique{faire l'Hellène}, Λεσβιάζειν, Βακίζειν, \italique{faire le Bakis} (\italique{Paix}, 1072). Les Égyptiens avaient la réputation d'être très rusés. On a conservé un vers d'Eschyle\marqueur{1} :
\citComm{δεινοί πλέκειν τοι μηχανὰς Αἰγύπτιοι\\
Les Égyptiens habiles à ourdir des trames.}
\noindent Cf. Properce (III, 9, 33) : \latin{Noxia Alexandria, dolis aptissima lellus.} (Note de Willems ; cf. aussi Taillardat §409.)
\noindent\textbf{[Notes ad 922]} \textbf{(1)} Tiré d'une scholie à Théocrite (15, 48b). C'est le fragment 373 de l'édition de Radt.
\noindent\textbf{928} \textbf{αὕτη μὲν ἡ μήρινθος οὐδὲν ἔσπασεν.} Le tour est proverbial et l'expression nous est bien connue grâce à une scholie au vers 175 des \italique{Guêpes} où Aristophane joue avec l'halieutique proverbe :
\citComm{παρὰ τὴν παροιμίαν « αὕτη μὲν ἡ μήρινθος οὐδὲν ἔσπασεν ». ἐλέγετο δὲ ἐπὶ τῶν πείρᾳ καθιέντων ἄγκιστρον μετὰ σχοίνου, καὶ μὴ ἐπιτυγχανόντων ἰχθύος τινός.\\
\francais{En regard du proverbe : « voici une ligne qui n'a rien tiré hors de l'eau ». On le dit à propos de ceux qui jettent un appât à la pointe acérée (un hameçon avec un roseau), mais qui ne rencontrent pas un poisson.}}
\noindent Ce proverbe est dit d'une ruse qui échoue\marqueur{1}. En français : un échec sur toute la ligne.
\noindent\textbf{[Notes ad 928]} \textbf{(1)} Pour un exposé plus développé sur ces deux passages avec à la clef l'explication du jargon halieutique, le lecteur pourra aller consulter le §395 des \italique{Images d'Aristophane} de J. Taillardat. Érasme a recensé l'adage (n° 845).
\noindent\textbf{935} \textbf{ἱστιορράφος.} Allusion aux hardes dont était couvert le prétendu Ménélas. (Note de Willems ; ἱστιορράφος est mis παρὰ προσδοκίαν pour μηχανορράφος\marqueur{1}.)
\noindent\textbf{[Notes ad 935]} \textbf{(1)} Cf. Euripide, \italique{Andromaque}, vv. 447, 1116 ; Sophocle, \italique{Œdipe Roi}, v. 387.
\noindent\textbf{936-7} \textbf{ὦ πρύτανι πρὸς τῆς δεξιᾶς, ἥνπερ φιλεῖς / κοίλην προτείνειν ἀργύριον ἤν τις διδῷ.} Précis à l'usage du justiciable confronté à un prytane.
\noindent Les deux répliques que Mnésiloque adresse au prytane nous donnent un bon aperçu de l'estime en laquelle Aristophane tenait ces magistrats -- car nous allons voir que Mnésiloque tente de le soudoyer, comme cela semble l'usage à Athènes.
\end{multicols}
}%
\end{adjustbox}%
\vfill
\end{adjustwidth}
\newpage
% --- PAGE DE DROITE : Traduction française + Commentaire (2/8) ---
\begin{adjustwidth}{-2cm}{-2cm}
\noindent\textbf{\footnotesize Traduction française --- Folio 23 (\texttt{f. 173v}, vv. 920–958 [1/4]) \hfill ΕΠΕΙΣΟΔΙΟΝ Δ' / ΧΟΡΙΚΟΝ}\\[-2pt]
\noindent\rule{16.5cm}{0.5pt}\\[2pt]
{\fontsize{8.8pt}{10.2pt}\selectfont
\ligneEst{\textbf{Kritylla}}{\francais{Ah! toi aussi, que tu m'as l'air d'un fripon}}{920}
\ligneEst{}{\francais{et d'être son conseiller. Rien d'étonnant que depuis un temps}}{}
\ligneEst{}{\francais{vous jouiez les Égyptiens. Mais celui-là en portera la peine.}}{}
\ligneEst{}{\francais{Car voici venir le prytane et l'archer.}}{}
\ligneEst{\textbf{Euripide}}{\francais{Voilà qui ne vaut rien. Il est temps d'esquiver.}}{}
\ligneEst{\textbf{Mnésiloque}}{\francais{Et moi, malheureux, que vais-je faire ?}}{925}
\ligneEst{\textbf{Euripide}}{\francais{\espace{4.2cm}Sois en paix.}}{}
\ligneEst{}{\francais{Jamais je ne te trahirai, tant que j'aurai le souffle,}}{}
\ligneEst{}{\francais{à moins que mes mille stratagèmes ne viennent à me faillir.}}{}
\ligneEst{\textbf{Mnésiloque}}{\francais{Voilà une amorce qui n'a pas pris.}}{}
}
\vspace{2pt}
\noindent\rule{16.5cm}{0.6pt}\\[2pt]
\noindent\textbf{\small Commentaire philologique, métrique et scénique (vv. 920–958 --- partie 2/8)}\\[2pt]
\begin{adjustbox}{max totalheight=15.2cm, max width=16.5cm, keepaspectratio}%
\parbox[t]{16.5cm}{
\setlength{\columnsep}{9pt}%
\begin{multicols}{2}
\fontsize{8.4pt}{9.7pt}\selectfont
\setlength{\parindent}{0pt}%
\setlength{\parskip}{1.2pt}%
\noindent On trouve le même reproche de concussion dans la \italique{Paix} (vers 905-908) :
\citComm{Ἀλλ᾽, ὦ πρυτάνεις, δέχεσθε τὴν Θεωρίαν.\\
Θέασ᾽ ὡς προθύμως ὁ πρύτανις παρεδέξατο.\\
Ἀλλ᾽ οὐκ ἄν, εἴ τι προῖκα προσαγαγεῖν σ᾽ ἔδει,\\
ἀλλ᾽ ηὗρον ἄν σ᾽ ὑπέχοντα τὴν ἐκεχειρίαν.\\
Allons, prytanes, recevez Théoria. Voyez l'accueil empressé que lui fait ce prytane. Quelle différence s'il t'avait fallu introduire une affaire gratis. Je t'aurais vu te retrancher sur les vacances.}
\noindent (ὑπέχοντα τὴν ἐκεχειρίαν, «~alléguant la trêve~», est la version bégayée de ce qu'il fallait vraiment entendre : ὑπέχοντα τὴν χεῖρα, «~tendant la main~».)\\ En effet, pour soumettre une requête à la Boulé l'on devait préalablement obtenir l'aval des prytanes ; le cas échéant, les bouleutes discutaient et rédigeaient une proposition de loi (προβούλευμα) soumise ensuite à l'Ecclésia. L'auteur de la \italique{République des Athéniens} (III, 1) constate qu'il fallait parfois patienter jusqu'à une année avant que les prytanes ne donnent leur accord.
\noindent De même, pour intenter une action en justice, un citoyen devait régler des \italique{frais de justice}, τά πρυτανεῖα\marqueur{1}, dont le nom indique clairement à qui cette somme devait être remise.\\  Par ses fonctions, le prytane \italique{s'engageait} à servir d'intermédiaire : ce genre de serment passait très certainement par une \italique{poignée de main}\marqueur{2}.
\noindent C'est le sens de «~ἥνπερ φιλεῖς / κοίλην προτείνειν~», «~[cette main] que tu as l'habitude de tendre ouverte~».\\ Mais toutes ces transactions étaient autant d'occasions pour les prytanes de détourner à leur profit une partie du bien public.
\noindent C'est entendu : le prytane est concussionnaire mais il serait suicidaire pour Mnésiloque de le lui dire !
\citComm{Flattez-les, payez-les d'agréables mensonges,\\
Quelque indignation dont leur cœur soit rempli,\\
Ils goberont l'appât, vous serez leur ami.}
\end{multicols}
}%
\end{adjustbox}%
\vfill
\end{adjustwidth}
\newpage

% --- FOLIO 23 (f. 173v) : SPREAD 2.2 (Texte grec établi + Traduction française + Commentaire 3/8 & 4/8) ---
\clearpage\ifodd\value{page}\hbox{}\newpage\fi
\enTeteG{ΘΕΣΜΟΦΟΡΙΑΖΟΥΣΑΙ --- Texte grec établi (f. 173v, vv. 920–958 [2/4])}
\enTeteD{LES THESMOPHORIES --- Traduction française (f. 173v, vv. 920–958 [2/4])}
\begin{adjustwidth}{-2cm}{-2cm}
\noindent\textbf{\footnotesize Texte grec établi --- Folio 23 (\texttt{f. 173v}, vv. 920–958 [2/4]) \hfill ΕΠΕΙΣΟΔΙΟΝ Δ' / ΧΟΡΙΚΟΝ}\\[-2pt]
\noindent\rule{16.5cm}{0.5pt}\\[2pt]
{\fontsize{8.8pt}{10.2pt}\selectfont
\ligneEst{\textbf{Πρύτανις}}{ὅδ᾽ ἔσθ᾽ ὁ πανοῦργος ὃν ἔλεγ᾽ ἡμῖν Κλεισθένης;}{}
\ligneEst{}{οὗτος τί κύπτεις; δῆσον αὐτὸν εἰσάγων}{930}
\ligneEst{}{ὦ τοξότ᾽ ἐν τῇ σανίδι, κἄπειτ᾽ ἐνθαδὶ}{}
\ligneEst{}{στήσας φύλαττε καὶ προσιέναι μηδένα}{}
\ligneEst{}{ἔα πρὸς αὐτόν, ἀλλὰ τὴν μάστιγ᾽ ἔχων}{}
\ligneEst{}{παῖ᾽ ἢν προσίῃ τις.}{}
\ligneEst{\textbf{Κρίτυλλα}}{\decalage{παῖ᾽ ἢν προσίῃ τις.}νὴ Δί᾽ ὡς νῦν δή γ᾽ ἀνὴρ}{}
\ligneEst{}{ὀλίγου μ᾽ ἀφείλετ᾽ αὐτὸν ἱστιορράφος.}{935}
\ligneEst{\textbf{Μνησίλοχος}}{ὦ πρύτανι πρὸς τῆς δεξιᾶς, ἥνπερ φιλεῖς}{}
\ligneEst{}{κοίλην προτείνειν ἀργύριον ἤν τις διδῷ,}{}
\ligneEst{}{χάρισαι βραχύ τί μοι καίπερ ἀποθανουμένῳ.}{}
}
\vspace{2pt}
\noindent\rule{16.5cm}{0.6pt}\\[2pt]
\noindent\textbf{\small Commentaire philologique, métrique et scénique (vv. 920–958 --- partie 3/8)}\\[2pt]
\begin{adjustbox}{max totalheight=15.2cm, max width=16.5cm, keepaspectratio}%
\parbox[t]{16.5cm}{
\setlength{\columnsep}{9pt}%
\begin{multicols}{2}
\fontsize{8.4pt}{9.7pt}\selectfont
\setlength{\parindent}{0pt}%
\setlength{\parskip}{1.2pt}%
\noindent Au contraire : Mnésiloque tente de mettre le prytane dans sa poche -- il y parvient d'ailleurs au vu de la bienveillante réponse qu'il reçoit : τί σοι χαρίσωμαι;\\ Pour cela, jeu de scène\marqueur{3} : tout en déclamant «~πρὸς τῆς δεξιᾶς κτλ.~», Mnésiloque, qui n'est pas encore lié, tend sa main vers le prytane et s'offre par ce geste en suppliant\marqueur{4}.
\noindent (Et d'ajouter -- en aparté, évidemment ! mais ce que le prytane a bien compris et qui explique ses bonnes dispositions à l'égard de Mnésiloque -- : «~ἀργύριον ἤν τις διδῷ~», «~[cette main que tu aimes à tendre] lorsque l'on t'offre de l'argent~».)\\ Le prytane accordera bien cette petite faveur à Mnésiloque (χάρισαι βραχύ τί μοι), à lui qui est sur le point de mourir ; et Mnésiloque d'insister sur ce dernier point (καίπερ ἀποθανουμένῳ).
\noindent Aristophane joue avec une locution proverbiale, selon sa pente maintenant bien établie par notre commentaire : les Athéniens, même à l'article de la mort, étaient réputés continuer de quémander la main tendue\marqueur{5} :
\citComm{Ἀττικὸς ὑπέχει τὴν χεῖρα ἀποθνήσκων· ἐπὶ τῶν φιλαργύρων. φιλοκερδεῖς γὰρ οἱ Ἀθηναῖοι.}
\noindent Le \italique{verre du condamné} que demande Mnésiloque (καίπερ ἀποθανουμένῳ) n'est autre que cette faveur mendiée la main tendue vers le prytane.
\noindent\textbf{[Notes ad 936-7]} \textbf{(1)} Voir à ce sujet le Daremberg \& Saglio à l'entrée \latin{Prytaneia}. Ces informations aident à interpréter la glose d'Hésychius (π 1511) :
\citComm{πεπρυτανευκώς· ἐπιμεμηνιευκώς. καὶ προανηλωκώς,}
\noindent rangée parmi les fragments comiques adespotes (408) dans l'édition Kassel-Austin.
\noindent \textbf{(2)} On pourra lire l'article de Francesco Mari : «~Les sens de la poignée de main en Grèce ancienne du VIIIe au Ve siècle avant J.-C.~», in \italique{Ktèma}, 43 (2018), pp. 105-131.
\noindent \textbf{(3)} Il faut rendre aux jeux de scènes toute l'importance qu'ils avaient à l'époque ; empêtrés dans leurs costumes et d'imposants masques de théâtre, il est à parier que les acteurs compensaient cette gêne par des jeux de mains des plus expressifs.
\noindent \textbf{(4)} (La \italique{droite} en question est bien celle de Mnésiloque et non du prytane.) Pour cette signification de la main tendue, cf. \LSJ \ à l'entrée προτείνω II.
\end{multicols}
}%
\end{adjustbox}%
\vfill
\end{adjustwidth}
\newpage
% --- PAGE DE DROITE : Traduction française + Commentaire (4/8) ---
\begin{adjustwidth}{-2cm}{-2cm}
\noindent\textbf{\footnotesize Traduction française --- Folio 23 (\texttt{f. 173v}, vv. 920–958 [2/4]) \hfill ΕΠΕΙΣΟΔΙΟΝ Δ' / ΧΟΡΙΚΟΝ}\\[-2pt]
\noindent\rule{16.5cm}{0.5pt}\\[2pt]
{\fontsize{8.8pt}{10.2pt}\selectfont
\ligneEst{\textbf{Prytane}}{\francais{Est-ce là le fripon dont nous parlait Clisthène ?}}{}
\ligneEst{}{\francais{Holà, pour quoi courbes-tu le front? Archer, rentre avec lui et attache-le}}{930}
\ligneEst{}{\francais{à l'ais. Puis expose-le ici,}}{}
\ligneEst{}{\francais{et garde-le bien. Ne laisse approcher personne,}}{}
\ligneEst{}{\francais{mais frappe de ton fouet}}{}
\ligneEst{}{\francais{quiconque s'avancera.}}{}
\ligneEst{\textbf{Kritylla}}{\francais{\espace{2.9cm}Bien, par Zeus, car tout à l'heure}}{}
\ligneEst{}{\francais{un rapetasseur de voiles a failli me l'enlever.}}{935}
\ligneEst{\textbf{Mnésiloque}}{\francais{Ô prytane, au nom de cette droite dont tu aimes}}{}
\ligneEst{}{\francais{à tendre le creux (À part :) quand on t'offre de l'argent...}}{}
\ligneEst{}{\francais{Accorde-moi une petite faveur, à moi qui suis sur le point de mourir.}}{}
}
\vspace{2pt}
\noindent\rule{16.5cm}{0.6pt}\\[2pt]
\noindent\textbf{\small Commentaire philologique, métrique et scénique (vv. 920–958 --- partie 4/8)}\\[2pt]
\begin{adjustbox}{max totalheight=15.2cm, max width=16.5cm, keepaspectratio}%
\parbox[t]{16.5cm}{
\setlength{\columnsep}{9pt}%
\begin{multicols}{2}
\fontsize{8.4pt}{9.7pt}\selectfont
\setlength{\parindent}{0pt}%
\setlength{\parskip}{1.2pt}%
\noindent \textbf{(5)} Recensé par Diogénianus (III, 12), \italique{Corpus paroemiographorum Graecorum}, I, p. 216. C'est l'adage n° 1646 d'Érasme. À noter qu'Aristophane y fait aussi une possible allusion aux vers 170 et suivants des \italique{Grenouilles}, lorsque Dionysos hèle un mort pour qu'il porte ses bagages, ce dernier lui répond : Δύο δραχμὰς μισθὸν τελεῖς; Voir plus haut le commentaire du vers 738.
\noindent\textbf{939-42} \textbf{γυμνὸν ἀποδύσαντά με / κέλευε πρὸς τῇ σανίδι δεῖν τὸν τοξότην, / ἵνα μὴ 'ν κροκωτοῖς καὶ μίτραις γέρων ἀνὴρ / γέλωτα παρέχω τοῖς κόραξιν ἑστιῶν.} Ces quatre vers sont peut-être parmi les plus mécompris de toute la comédie.
\noindent L'erreur vient de ce que `πρός τῇ σανίδι δεῖν' signifie littéralement « lier \italique{contre} la planche » et non pas « lier à la planche » ; cette dernière expression se construisant en grec avec la préposition ἐν et non πρός\marqueur{1}.
\noindent Que dit en apparence Mnésiloque ? il fait profession honorable de pudeur ; ce pudibond de la onzième heure tient à cacher ce qui doit être caché.
\noindent Voilà pourquoi il demande à être retourné \italique{contre} (πρός) \italique{la planche} du supplice (τῇ σανίδι)\marqueur{2}.
\noindent À qui tient-il à cacher son anatomie ? aux femmes évidemment\marqueur{3} !.. Au vers 942, les `κροκωτοῖς καὶ μίτραις' sont les \italique{crocotes} et \italique{mitres} des Athéniennes réunies, rappelons-le, en assemblée dans le but de perdre Euripide. Dans ce contexte, ἐν + datif a le sens qu'il a toujours\marqueur{4} : \italique{au milieu de}, \italique{parmi}, \italique{devant} : «~ἐν κροκωτοῖς καὶ μίτραις~» $=$ «~au milieu des (en)crocotées et des (en)mitrées~».
\noindent Voici pour le sens clair.
\noindent Mais n'oublions pas que Mnésiloque pratique avec un art consommé le double langage : toute parole qu'il adresse au prytane est à double sens : un clair, où il est dit que le sens véritable est caché\marqueur{5}.
\noindent Que demande réellement Mnésiloque au prytane ?
\noindent L'expression `δῆσαι πρὸς σανίδι' a heureusement été relevée par Hésychius (δ 926) -- et c'est le seul des nombreux lexique consultés à avoir relevé ce tour :
\citComm{δῆσαι πρὸς σανίδι· Ἀττικῶς· βασανίσαι ἐν τῷ δεσμωτηρίῳ.\\
Lier face à la planche : Atticisme. Interroger dans la prison.}
\end{multicols}
}%
\end{adjustbox}%
\vfill
\end{adjustwidth}
\newpage

% --- FOLIO 23 (f. 173v) : SPREAD 2.3 (Texte grec établi + Traduction française + Commentaire 5/8 & 6/8) ---
\clearpage\ifodd\value{page}\hbox{}\newpage\fi
\enTeteG{ΘΕΣΜΟΦΟΡΙΑΖΟΥΣΑΙ --- Texte grec établi (f. 173v, vv. 920–958 [3/4])}
\enTeteD{LES THESMOPHORIES --- Traduction française (f. 173v, vv. 920–958 [3/4])}
\begin{adjustwidth}{-2cm}{-2cm}
\noindent\textbf{\footnotesize Texte grec établi --- Folio 23 (\texttt{f. 173v}, vv. 920–958 [3/4]) \hfill ΕΠΕΙΣΟΔΙΟΝ Δ' / ΧΟΡΙΚΟΝ}\\[-2pt]
\noindent\rule{16.5cm}{0.5pt}\\[2pt]
{\fontsize{8.8pt}{10.2pt}\selectfont
\ligneEst{\textbf{Πρύτανις}}{τί σοι χαρίσωμαι;}{}
\ligneEst{\textbf{Μνησίλοχος}}{\decalage{τί σοι χαρίσωμαι;}γυμνὸν ἀποδύσαντά με,}{}
\ligneEst{}{κέλευε πρὸς τῇ σανίδι δεῖν τὸν τοξότην,}{940}
\ligneEst{}{ἵνα μὴ 'ν κροκωτοῖς καὶ μίτραις γέρων ἀνὴρ}{}
\ligneEst{}{γέλωτα παρέχω τοῖς κόραξιν ἑστιῶν.}{}
\ligneEst{\textbf{Πρύτανις}}{ἔχοντα ταῦτ᾽ ἔδοξε τῇ βουλῇ σε δεῖν,}{}
\ligneEst{}{ἵνα τοῖς παροῦσι δῆλος ᾖς πανοῦργος ὤν.}{}
\ligneEst{\textbf{Μνησίλοχος}}{ἰαππαπαιάξ· ὦ κροκώθ᾽ οἷ᾽ εἴργασαι·}{945}
\ligneEst{}{κοὐκ ἔστ᾽ ἔτ᾽ ἐλπὶς οὐδεμία σωτηρίας.}{}
\ligneEst{\textbf{Κορυφαῖος} {\tiny\textit{\decalageN{4}$4$ an\cat}}}{\Automatique{ἄγε νυν ἡμεῖς παίσωμεν ἅπερ νόμος ἐνθάδε ταῖσι γυναιξίν,}}{}
\ligneEst{}{\Automatique{παραχώρει ὅταν ὄργια σεμνὰ θεαῖν ἱεραῖς ὥραις ἀνέχωμεν, ἅπερ καὶ}}{}
}
\vspace{2pt}
\noindent\rule{16.5cm}{0.6pt}\\[2pt]
\noindent\textbf{\small Commentaire philologique, métrique et scénique (vv. 920–958 --- partie 5/8)}\\[2pt]
\begin{adjustbox}{max totalheight=15.2cm, max width=16.5cm, keepaspectratio}%
\parbox[t]{16.5cm}{
\setlength{\columnsep}{9pt}%
\begin{multicols}{2}
\fontsize{8.4pt}{9.7pt}\selectfont
\setlength{\parindent}{0pt}%
\setlength{\parskip}{1.2pt}%
\noindent Mnésiloque veut un procès équitable (il sent que la justice d'exception risque de lui être par trop expéditive) : il est préférable pour lui de s'en remettre aux masculines institutions. Il demande donc à être conduit en prison pour y être interrogé.
\noindent Réponse négative du prytane : c'est sur place, \italique{devant les personnes présentes} et non pas au δεσμωτήριον, qu'il devra expier son forfait :
\citComm{Πρύτανις --- ἔχοντα ταῦτ᾽ ἔδοξε τῇ βουλῇ σε δεῖν,\\
ἵνα τοῖς παροῦσι δῆλος ᾖς πανοῦργος ὤν.\\}
La σανίς de l'expression rapportée par Hésychius désigne sans doute la \italique{planche} sur laquelle était inscrite l'ordre du jour des tribunaux.
Cf. en \italique{Guêpes} (v. 848) et Photius (σ 68) : 
\citComm{σανίς· καὶ ἐν ᾗ τοὺς κακούργους ἔδουν καὶ ἐν ᾗ τὰς δίκας ἐνεχάραττον· τὸ µὲν ἐν Θεσµοφοριαζούσαις, τὸ δὲ ἐν Σφηξὶ (848) δεδήλωται.}
Aristophane, ainsi qu'il aime à le faire assez souvent, détourne un mot ou une expression imagée qu'il interprète littéralement.

Devant le refus du prytane, Mnésiloque se retourne contre sa crocote -- au cœur des péripéties de la pièce -- qu'il accuse d'être à l'origine de la situation : il la gronde comme on gronderait un enfant\marqueur{6} :
\citComm{Μνησίλοχος --- ἰαππαπαιάξ· ὦ κροκώθ᾽ οἷ᾽ εἴργασαι.}

C'est encore Hésychius (κ 4170) qui nous permet de confirmer qu'il y a probablement là un jeu de mots κρόκαλ' $\sim$ κροκώθ' :
\citComm{κρόκαλον· τὸ πανοῦργον παιδίον.}
\noindent\textbf{[Notes ad 939-42]} \textbf{(1)} Cf. les vers 934 (δῆσον αὐτὸν [...] ἐν τῇ σανίδι) et 1165 (οὑν [$=$ ὁ ἐν] τῇ σανίδι).
\noindent \textbf{(2)} En évitant par la même occasion que les corbeaux ne l'énucléent et ne se repaissent de ses yeux... spectacle qu'on image sans peine peu ragoûtant et fort désagréable pour le supplicié.
\noindent \textbf{(3)} Et non aux corbeaux, comme on peut le lire dans les éditions les plus récentes des \italique{Thesmophories}...
\noindent \textbf{(4)} \italique{Confer} les vers 90-91.
\noindent \textbf{(5)} Quelques vers auparavant, il avait demandé une faveur au prytane, tout en lui faisant bien comprendre que cette faveur mériterait salaire. Sur la concussion des prytanes, voir aussi la \italique{Paix} (vers 905-908).
\end{multicols}
}%
\end{adjustbox}%
\vfill
\end{adjustwidth}
\newpage
% --- PAGE DE DROITE : Traduction française + Commentaire (6/8) ---
\begin{adjustwidth}{-2cm}{-2cm}
\noindent\textbf{\footnotesize Traduction française --- Folio 23 (\texttt{f. 173v}, vv. 920–958 [3/4]) \hfill ΕΠΕΙΣΟΔΙΟΝ Δ' / ΧΟΡΙΚΟΝ}\\[-2pt]
\noindent\rule{16.5cm}{0.5pt}\\[2pt]
{\fontsize{8.8pt}{10.2pt}\selectfont
\ligneEst{\textbf{Prytane}}{\francais{Quelle faveur ?}}{}
\ligneEst{\textbf{Mnésiloque}}{\francais{\espace{2.8cm}ordonne à l'archer}}{}
\ligneEst{}{\francais{de m'attacher nu face à l'ais.}}{940}
\ligneEst{}{\francais{Je crains d'être la risée des crocotes et des mitres}}{}
\ligneEst{}{\francais{avec les corbeaux que je vais régaler.}}{}
\ligneEst{\textbf{Prytane}}{\francais{La boulé a décidé qu'on t'attacherait dans ce costume,}}{}
\ligneEst{}{\francais{pour témoigner aux personnes présentes que tu es un fripon.}}{}
\ligneEst{\textbf{Mnésiloque}}{\francais{Aïe, aïe! aïe, aïe! Ô crocote, que de malheurs tu as causés !}}{945}
\ligneEst{}{\francais{Il n'est plus pour moi le moindre espoir de salut.}}{}
\ligneEst{\textbf{Coryphée}}{\francais{Allons, livrons-nous à nos danses, comme c'est ici la coutume des femmes,}}{}
\ligneEst{}{\francais{quand aux dates saintes nous célébrons les augustes mystères des deux déesses.}}{}
}
\vspace{2pt}
\noindent\rule{16.5cm}{0.6pt}\\[2pt]
\noindent\textbf{\small Commentaire philologique, métrique et scénique (vv. 920–958 --- partie 6/8)}\\[2pt]
\begin{adjustbox}{max totalheight=15.2cm, max width=16.5cm, keepaspectratio}%
\parbox[t]{16.5cm}{
\setlength{\columnsep}{9pt}%
\begin{multicols}{2}
\fontsize{8.4pt}{9.7pt}\selectfont
\setlength{\parindent}{0pt}%
\setlength{\parskip}{1.2pt}%
\noindent \textbf{(6)} Voir par exemple, \italique{Philoctète}, v. 940 : οἷ᾽ ἔργ᾽ ὁ παῖς μ᾽ ἔδρασεν οὑξ Ἀχιλλέως.
\noindent\textbf{943-56.} \textbf{Le papyrus d'Oxyrhynque n° 3839\marqueur{1},} dont il a été déjà question aux vers 742-766 contient contient également un fragment de ce passage. Il est si gâté qu'il ne nous apporte aucune information (si ce n'est une colométrie que Lætitia suivit \anglais{without strong conviction}\marqueur{2}). Le lecteur le trouvera néanmoins reproduit en annexe ; gage à lui d'en tirer quelque chose.
\noindent\textbf{[Notes ad 943-56]} \textbf{(1)} \italique{The Oxyrhynchus Papyri}, vol. LVI (1989), n°3839, fr. 2. Analysé par W. E. H. Coockle avec le concours de Colin Austin.
\noindent \textbf{(2)} L. P. E. Parker, \italique{The Songs of Aristophanes} (1997), p. 430.
\noindent\textbf{945} \textbf{ὦ κροκώθ᾽ οἷ᾽ εἴργασαι·} Y a-t-il un jeu de mots entre κρόκαλον et κροκωτός ?
\noindent Cf. Hésychius (κ 4170 et 4270) :
\citComm{κρόκαλον· τὸ πανοῦργον παιδίον.}
\noindent\textbf{947-1000} \textbf{} Pour une analyse plus fine du chant mêlé de danses de tout ce passage, je renvoie à l'article d'Ester Cerbo\marqueur{1}.
\noindent Quant à sa structure, celle d'Enger, assez communément admise -- et notamment par Lætitia\marqueur{2} --, me semble juste  :
\citComm{--- anapestes (v. 947-952) ;\\
--- proode (v. 953-958) ;\\
--- première paire strophique (v. 959-962 $=$ v. 963-965) ;\\
--- première mésode (v. 966-968) ;\\
--- seconde paire strophique (v. 969-976 $=$ v. 977-984) ;\\
--- seconde mésode (v. 985-989) ;\\
--- tierce paire strophique (v. 990-994 $=$ v. 995-1000).}
\noindent\textbf{[Notes ad 947-1000]} \textbf{(1)} «~Uno spettacolo di danza alla festa delle donne: a proposito di Ar. \italique{Thesm.} 947-1000~» (2021).
\noindent \textbf{(2)} Dans son indispensable \italique{The Songs of Aristophanes}, pp. 428-436.
\end{multicols}
}%
\end{adjustbox}%
\vfill
\end{adjustwidth}
\newpage

% --- FOLIO 23 (f. 173v) : SPREAD 2.4 (Texte grec établi + Traduction française + Commentaire 7/8 & 8/8) ---
\clearpage\ifodd\value{page}\hbox{}\newpage\fi
\enTeteG{ΘΕΣΜΟΦΟΡΙΑΖΟΥΣΑΙ --- Texte grec établi (f. 173v, vv. 920–958 [4/4])}
\enTeteD{LES THESMOPHORIES --- Traduction française (f. 173v, vv. 920–958 [4/4])}
\begin{adjustwidth}{-2cm}{-2cm}
\noindent\textbf{\footnotesize Texte grec établi --- Folio 23 (\texttt{f. 173v}, vv. 920–958 [4/4]) \hfill ΕΠΕΙΣΟΔΙΟΝ Δ' / ΧΟΡΙΚΟΝ}\\[-2pt]
\noindent\rule{16.5cm}{0.5pt}\\[2pt]
{\fontsize{8.8pt}{10.2pt}\selectfont
\ligneEst{{\tiny\textit{dim. anap.}}}{Παύσων σέβεται καὶ νηστεύει,}{}
\ligneEst{}{πολλάκις αὐταῖν ἐκ τῶν ὡρῶν}{950}
\ligneEst{}{ἐς τὰς ὥρας ξυνεπευχόμενος}{}
\ligneEst{{\tiny\textit{paroem.}}}{τοιαῦτα μέλειν θάμ᾽ ἑαυτῷ.}{}
\ligneEst{\textbf{Χορός} {\tiny\textit{anap.}}}{ὅρμα χώρει·}{}
\ligneEst{{\tiny\textit{lecyth.}}}{κοῦφα ποσὶν ἄγ᾽ ἐς κύκλον,}{}
\ligneEst{{\tiny\textit{aristoph.}}}{χειρὶ σύναπτε χεῖρα,}{955}
\ligneEst{}{†ῥυθμὸν χορείας ὕπαγε πᾶσα†·}{}
\ligneEst{{\tiny\textit{glyc.}}}{βαῖνε καρπαλίμοιν ποδοῖν.}{}
\ligneEst{{\tiny\textit{dim. iamb.}}}{ἐπισκοπεῖν δὲ πανταχῇ}{}
\ligneEst{{\tiny\textit{trim. iamb.}}}{κυκλοῦσαν ὄμμα χρὴ χοροῦ κατάστασιν.}{958}
}
\vspace{2pt}
\noindent\rule{16.5cm}{0.6pt}\\[2pt]
\noindent\textbf{\small Commentaire philologique, métrique et scénique (vv. 920–958 --- partie 7/8)}\\[2pt]
\begin{adjustbox}{max totalheight=15.2cm, max width=16.5cm, keepaspectratio}%
\parbox[t]{16.5cm}{
\setlength{\columnsep}{9pt}%
\begin{multicols}{2}
\fontsize{7.6pt}{8.8pt}\selectfont
\setlength{\parindent}{0pt}%
\setlength{\parskip}{1.2pt}%
\noindent\textbf{949-52} \textbf{Παύσων σέβεται καὶ νηστεύει κτλ.} Pauson était connu pour son indigence. Les parémiographes nous ont rapporté à son sujet deux proverbes\marqueur{1} :
\citComm{Παύσωνος πτωχότερος· οὗτος ζωγράφος ἦν καὶ ἐπὶ πενίᾳ διατεθρύλλητο.\\
Παύσων καὶ Ἶρος: ὀνόματα κύρια πενήτων. ὁ Ἀσκληπιὸς Παύσωνα καὶ Ἶρον κἂν ἄλλον τινὰ τῶν ἀπόρων ἰάσαιτο. κτλ.\\
(Iros, le mendiant odieux d'Homère, \italique{Odyssée} XVIII.)}
\noindent L'on trouvera tous les témoignages sur Pauson d'Athènes dans le \italique{Recueil Milliet} d'Adolphe Reinach\marqueur{2}. Pour ce qui est d'Aristophane, la scholie à notre passage indique :
\citComm{Παύσων ὁ ζωγράφος εἰς πενίαν διαβάλλεται, καὶ διὰ τοῦτο εἶπε `νηστεύει'.}
\noindent Pauson l'importun est aussi mentionné au vers 853 des \italique{Acharniens} :
\citComm{οὐδ᾽ αὖθις αὖ σε σκώψεται\\
Παύσων ὁ παμπόνηρος\\
Ni les brocards de Pauson, le gueux achevé, tu ne les essuieras [sur le marché].\\
Σ : οὗτος ὁ Παύσων ζωγράφος πένης σχωπτολόγος.}
\noindent Enfin au vers 602 du \italique{Ploutos} :
\citComm{Χρέμυλος --- Παρὰ τῆς Ἑκάτης ἔξεστιν τοῦτο πυθέσθαι,\\
εἴτε τὸ πλουτεῖν εἴτε τὸ πεινῆν βέλτιον. Φησὶ γὰρ αὕτη\\
τοὺς μὲν ἔχοντας καὶ πλουτοῦντας δεῖπνον κατὰ μῆν᾽ ἀποπέμπειν,\\
τοὺς δὲ πένητας τῶν ἀνθρώπων ἁρπάζειν πρὶν καταθεῖναι.\\
Ἀλλὰ φθείρου καὶ μὴ γρύξῃς\\
ἔτι μηδ᾽ ὁτιοῦν.\\
Οὐ γὰρ πείσεις, οὐδ᾽ ἢν πείσῃς.\\
Πενία --- Ὦ πόλις Ἄργους, κλύεθ᾽ οἷα λέγει.\\
Χρέμυλος --- \textbf{Παύσωνα} κάλει τὸν ξύσσιτον.\\
Chrémyle --- Il n'est que de s'informer chez Hécate, lequel vaut mieux d'être riche ou d'être pauvre. Car elle prétend que les riches et les opulents lui font hommage chaque mois d'un dîner, et que les pauvres s'en emparent avant qu'on l'ait mis là. Mais la peste soit de toi ; ne souffle plus le mot. Car tu ne me persuaderas pas, quand tu me persuaderais.\\
La Pauvreté --- « Cité des Argiens, écoutez ce qu'il dit ! »\\
Chrémyle --- Invoque \textbf{Pauson}, ton commensal.\\
Σ : ὁ Παύσων δὲ ἐπὶ πενίᾳ κωμῳδεῖται ζωγράφος ὦν. --- μετακαλοῦ σύντροφον τὸν Παύσωνα.}
\end{multicols}
}%
\end{adjustbox}%
\vfill
\end{adjustwidth}
\newpage
% --- PAGE DE DROITE : Traduction française + Commentaire (8/8) ---
\begin{adjustwidth}{-2cm}{-2cm}
\noindent\textbf{\footnotesize Traduction française --- Folio 23 (\texttt{f. 173v}, vv. 920–958 [4/4]) \hfill ΕΠΕΙΣΟΔΙΟΝ Δ' / ΧΟΡΙΚΟΝ}\\[-2pt]
\noindent\rule{16.5cm}{0.5pt}\\[2pt]
{\fontsize{8.8pt}{10.2pt}\selectfont
\ligneEst{}{\francais{Pauson aussi les sanctifie. Il jeûne}}{}
\ligneEst{}{\francais{et souvent prie avec nous que d'année}}{950}
\ligneEst{}{\francais{en année il puisse longtemps}}{}
\ligneEst{}{\francais{encore leur témoigner même dévotion.}}{}
\ligneEst{\textbf{Chœur}}{\francais{Pars, élance-toi}}{}
\ligneEst{}{\francais{d'un pied léger, forme le cercle,}}{}
\ligneEst{}{\francais{les mains enlacées aux mains,}}{955}
\ligneEst{}{\francais{marque le rythme de la danse.}}{}
\ligneEst{}{\francais{Avance de tes pieds alertes.}}{}
\ligneEst{}{\francais{Mais il faut que le chœur soit aux aguets}}{}
\ligneEst{}{\francais{et tourne la vue de toutes parts.}}{958}
}
\vspace{2pt}
\noindent\rule{16.5cm}{0.6pt}\\[2pt]
\noindent\textbf{\small Commentaire philologique, métrique et scénique (vv. 920–958 --- partie 8/8)}\\[2pt]
\begin{adjustbox}{max totalheight=15.2cm, max width=16.5cm, keepaspectratio}%
\parbox[t]{16.5cm}{
\setlength{\columnsep}{9pt}%
\begin{multicols}{2}
\fontsize{8.4pt}{9.7pt}\selectfont
\setlength{\parindent}{0pt}%
\setlength{\parskip}{1.2pt}%
\noindent On constate que seuls les commentateurs indiquent que Pauson fût peintre (ζωγράφος) : il n'en n'est jamais question chez Aristophane.
\noindent Il paraît douteux que ce Pauson-là, le mendiant fâcheux à la dévotion fort intéressée\marqueur{4}, fût le peintre connu qu'évoque Aristote\marqueur{3} ; la confusion venant peut-être de ce que la façon de peindre les figures du ζωγράφος ainsi que la seule anecdote nous ayant été rapportée à son sujet témoignent d'un caractère porté à la misanthropie\marqueur{5}.
\noindent\textbf{[Notes ad 949-52]} \textbf{(1)} Respectivement Michel Apostolius (XIII, 2) et Suidas (π 823).
\noindent \textbf{(2)} \italique{Recueil Milliet : Textes grecs et latins relatifs à l'histoire de la peinture ancienne}, tome I, publiés, traduits et commentés par Adolphe Reinach (1921) ; §§ 131-132 ; 173-179.
\noindent \textbf{(3)} Aristote, \italique{Politique}, 1340a36 ; \italique{Poétique}, 1448a6.
\noindent \textbf{(4)} Si sa dévotion prend les formes habituelles du souhait de renouvellement des prières (cf. le vers 287 ; Euripide, \italique{Électre}, vers 805), on se doute bien de ce qui attirait réellement Pauson-crève-la-faim. Voir ce que nous avons déjà dit aux vers 558-9 de la κρεανομία -- ou orgie institutionnalisée -- qui seule explique la réputation de piété des Athéniens.
\noindent \textbf{(5)} Pauson le peintre peignait des figures pires que ce qu'elles étaient en réalité et Aristote le déconseille pour l'éducation des jeunes gens. On lui avait commandé le tableau d'un cheval se roulant dans la poussière ; l'artiste dessine un cheval courant et soulevant un nuage de poussière autour de lui. Il y travaillait quand survient l'amateur. Celui-ci récrimine : ce n'est pas le sujet commandé. Pauson fait retourner par son garçon le tableau sens dessus dessous ; on voit alors un cheval renversé, qui se roule dans la poussière.
\noindent\textbf{955} \textbf{χειρὶ σύναπτε χεῖρα.} La danse que commande le chœur est probablement ce que Maurice Emmanuel\marqueur{1} appelle une \italique{danse des mains jointes}. Il s'agit donc bien d'une danse individuelle et non en couple : chaque danseur joignant la main à son autre main.
\noindent\textbf{[Notes ad 955]} \textbf{(1)} \italique{La danse grecque antique} (1896), §319.
\end{multicols}
}%
\end{adjustbox}%
\vfill
\end{adjustwidth}
\newpage
\clearpage\ifodd\value{page}\hbox{}\newpage\fi
\enTeteG{ΧΟΡΙΚΟΝ --- Codex Ravennas 429, f. 174r (Classense 176r) --- vv. 959–1000}
\begin{adjustwidth}{-1.5cm}{-3.0cm}
\begin{adjustbox}{max totalheight=21.0cm, max width=16.5cm, keepaspectratio}
\begin{minipage}{16.5cm}
{\fontsize{7.2pt}{8.5pt}\selectfont
\noindent\textbf{Folio 24 --- Codex Ravennas 429, \texttt{f. 174r} (\texttt{Classense 176r}) : ΧΟΡΙΚΟΝ, vv. 959–1000}\hfill \textit{Témoins :} $\text{R}^1, \Rpc, \Rdeux$ (\Mon)\par\vspace{2pt}
\hrule height 0.4pt\vspace{2.5pt}
}
{\fontsize{8.1pt}{9.5pt}\selectfont
\ligneRav{959}{χορ/}{\hspace{1.1em}ἅμα δὲ καὶ γένος Ὀλυμπίων θεῶν}{}
\ligneRav{960}{}{\hspace{1.1em}μέλπε καὶ γέραιρε φωνῇ πᾶσα χορομανεῖ τρόπῳ·}{}
\ligneRav{}{}{\hspace{1.1em}εἰ δέ τις προσδοκᾷ κακῶς ἐρεῖν}{}
\ligneRav{965}{}{\hspace{1.1em}ἐν ἱερῷ γυναῖκά μ᾽ οὖσαν ἄνδρας, οὐκ ὀρθῶς φρονεῖ·}{}
\ligneRav{966}{}{\hspace{1.1em}ἀλλὰ χρή}{\textbf{966} \lecon{χρὴ} R : χρῆν Bothe : χρή μ᾽ Meineke}
\ligneRav{967}{}{\hspace{1.1em}ὥσπερ ἔργον αὖ τι καινὸν}{\textbf{967} \lecon{ὥσπερ} R : ὡς πρὸς Austin}
\ligneRav{968}{}{\hspace{1.1em}{\Taille{11.5}πρῶτον εὐκύκλου χορείας \difference{εὐφυᾶ} στῆσαι βάσιν.}}{\textbf{968} εὐφυῆ R : \lecon{εὐφυᾶ} Brunck}
\ligneRav{969}{}{\hspace{1.1em}πρόβαινε \difference{ποσὶ} τὸν \difference{εὐλύραν}}{\textbf{969} ποσὶν R : \lecon{ποσὶ} Bentley \espApp ἐλύραν R : \lecon{εὐλύραν} Zanetti}
\ligneRav{970}{}{\hspace{1.1em}μέλπουσα καὶ τὴν τοξοφόρον}{}
\ligneRav{}{}{\hspace{1.1em}Ἄρτεμιν ἄνασσαν ἁγνήν·}{}
\ligneRav{}{}{\hspace{1.1em}χαῖρ᾽ ὦ ἑκάεργε,}{}
\ligneRav{}{}{\hspace{1.1em}ὄπαζε δὲ νίκην·}{}
\ligneRav{}{}{\hspace{1.1em}Ἥραν τε τὴν τελείαν}{}
\ligneRav{975}{}{\hspace{1.1em}μέλψωμεν ὥσπερ εἰκός,}{}
\ligneRav{}{}{\hspace{1.1em}ἣ πᾶσι τοῖς χοροῖσιν ἐμπαίζει τε καὶ}{}
\ligneRav{}{}{\hspace{1.1em}κλῇδας γάμου φυλάττει·}{}
\ligneRav{}{}{\hspace{1.1em}Ἑρμῆν τε νόμιον ἄντομαι}{}
\ligneRav{}{}{\hspace{1.1em}καὶ Πᾶνα καὶ Νύμφας φίλας}{}
\ligneRav{}{}{\hspace{1.1em}ἐπιγελάσαι προθύμως}{}
\ligneRav{980}{}{\hspace{1.1em}ταῖς \difference{ἡμετέραισι}}{\textbf{980} ἡμετέραις R : \lecon{ἡμετέραισι} Hermann}
\ligneRav{}{}{\hspace{1.1em}χαρέντα χορείαις·}{}
\ligneRav{982}{}{\hspace{1.1em}ἔξαιρε δὴ προθύμως}{\textbf{982} χαίρειν R : \lecon{χάριν} Ellebodius}
\ligneRav{982}{}{\hspace{1.1em}διπλῆν \difference{χάριν} χορείας·}{\textbf{982} χαίρειν R : \lecon{χάριν} Ellebodius}
\ligneRav{}{}{\hspace{1.1em}παίσωμεν ὦ γυναῖκες οἷάπερ νόμος,}{}
\ligneRav{984}{}{\hspace{1.1em}\difference{νηστεύομεν} δὲ πάντως·}{\textbf{984} νηστεύωμεν R : \lecon{νηστεύομεν} Bentley}
\ligneRav{985}{}{\hspace{1.1em}ἀλλ᾽ \difference{εἶα πάλλ᾽} ἀνάστρεφ᾽ εὐρύθμῳ ποδί,}{\textbf{985} ειαπαλλ᾽ R : \lecon{εἶα πάλλ᾽} Blaydes}
\ligneRav{}{}{\hspace{1.1em}τόρευε πᾶσαν ᾠδήν·}{}
\ligneRav{987}{}{\hspace{1.1em}ἡγοῦ δέ γ᾽ ὧδ᾽ αὐτὸς σὺ}{\textbf{987} γ᾽ ὧδ᾽ suspectum}
\ligneRav{}{}{\hspace{1.1em}κισσοφόρε Βακχεῖε}{}
\ligneRav{989}{}{\hspace{1.1em}δέσποτ᾽· ἐγὼ δὲ κώμοις}{\textbf{989} φιλοχόροισιν R : \lecon{φιλοχόροισι} Scaliger}
\ligneRav{989}{}{\hspace{1.1em}σὲ \difference{φιλοχόροισι} μέλψω·}{\textbf{989} φιλοχόροισιν R : \lecon{φιλοχόροισι} Scaliger}
\ligneRav{990}{}{\hspace{1.1em}†εὔιον ὦ Διόνυσε}{\textbf{990} \lecon{Διόνυσε} R : Διὸς σὺ Enger : Διός τε Fritzsche}
\ligneRav{}{}{\hspace{1.1em}Βρόμιε καὶ Σεμέλας παῖ,}{}
\ligneRav{}{}{\hspace{1.1em}χοροῖς τερπόμενος}{}
\ligneRav{993}{}{\hspace{1.1em}κατ᾽ ὄρεα νυμφᾶν ἐρατοῖς ἐν ὕμνοις,}{\textbf{993} \lecon{ἐρατοῖς ἐν ὕμνοις} R : ἐρατοῖσιν ὕμνοις Wilamowitz}
\ligneRav{994}{}{\hspace{1.1em}†εὔιον εὔιον εὐοῖ}{\textbf{994} lacunam statuit Hermann : ἥκεις γὰρ von Velsen : ἡδόμενος Austin cf. \italique{Plut}.288 \espApp \lecon{εὔιον εὔιον} R : ὦ Εὔι Εὔι' Hermann}
\ligneRav{994}{}{\hspace{1.1em}< $ \times \ \widearc{\cup \cup} \ \cup $ > ἀναχορεύων·}{\textbf{994} lacunam statuit Hermann : ἥκεις γὰρ von Velsen : ἡδόμενος Austin cf. \italique{Plut}.288 \espApp \lecon{εὔιον εὔιον} R : ὦ Εὔι Εὔι' Hermann}
\ligneRav{995}{}{\hspace{1.1em}ἀμφὶ δὲ \difference{σοὶ κτυπεῖται}}{\textbf{995} συῒ κτυπεῖται R : \lecon{σοὶ κτυπεῖται} Zanetti}
\ligneRav{996}{}{\hspace{1.1em}\difference{Κιθαιρώνιος} ἠχώ,}{\textbf{996} Κιθαρώνιος R : \lecon{Κιθαιρώνιος} Zanetti}
\ligneRav{}{}{\hspace{1.1em}μελάμφυλλά τ᾽ ὄρη}{}
\ligneRav{998}{}{\hspace{1.1em}δάσκια †καὶ νάπαι πετρώδεις† βρέμονται·}{\textbf{998} \lecon{καὶ νάπαι πετρώδεις} R : πετρώδεις τε νάπαι Enger}
\ligneRav{}{}{\hspace{1.1em}κύκλῳ δὲ περὶ σὲ κισσὸς}{}
\ligneRav{1000}{}{\hspace{1.1em}εὐπέταλος ἕλικι θάλλει·}{}
}
\vspace{3pt}\noindent\rule{4.5cm}{0.4pt}\par\vspace{1.5pt}
{\fontsize{7.0pt}{8.3pt}\selectfont
\noindent\textbf{Loci paralleli (vv. 959–1000) ---} \textbf{968} (εὐκύκλου) Hesychius ε 6929 εὔκυκλος· εὔχωρος $=$ Photius ε 2236 \espTest \textbf{973} (τελείαν) Hesychius τ 400 Τελεία \espTest \textbf{986} Suidas τ 789\par
}
\vspace{2.5pt}\noindent\rule{4.5cm}{0.4pt}\par\vspace{1.5pt}
{\fontsize{6.9pt}{8.2pt}\selectfont
\noindent\textbf{Suite (en commentaire) de l'apparat critique \& notes paléographiques (\texttt{f. 174r}) ---} \textbf{Contexte :} Grande ode chorale de danse en ronde (\italique{choros kyklikos}) célébrant Apollon, Artémis, Héra, Hermès, Pan, les Nymphes et Dionysos couronné de lierre (vv. 959–1000). --- Page entièrement lyrique où la colométrie de $\text{R}^1$ remonte directement à l'édition alexandrine d'Aristophane de Byzance (avec \italique{koronis} et \italique{paragraphoi} strophiques). --- \textbf{Fautes phonétiques pittoresques de $\text{R}^1$ :} Au \textbf{v. 969}, haplographie vocalique ἐλύραν pour \lecon{εὐλύραν} (Zanetti) ; au \textbf{v. 982}, χαίρειν pour \lecon{χάριν} (Ellebodius) ; au \textbf{v. 995},iotacisme/bêtacisme συῒ κτυπεῖται (« retentit pour un porc » !) au lieu de \lecon{σοὶ κτυπεῖται} (« retentit pour toi [Dionysos] », Zanetti).\par
}
\end{minipage}
\end{adjustbox}
\end{adjustwidth}
\newpage
\enTeteD{Codex Ravennas 429 (Biblioteca Classense) --- Fac-similé du f. 174r (vv. 959–1000)}
\begin{adjustwidth}{-1.5cm}{-3.0cm}
\begin{center}
\includegraphics[width=16.5cm,height=20.5cm,keepaspectratio]{/workspaces/Thesmophories/dossier_ravennas_thesmophories/folios_individuels_ff162v-177r/24_f174r_classense176r_vv0959-1000.jpg}\par
\vspace{2pt}
{\scriptsize\textbf{Codex Ravennas 429, f. 174r (Classense 176r)} --- ΧΟΡΙΚΟΝ, \textbf{vv. 959–1000} ($\text{R}^1, \Rpc, \Rdeux$ (\Mon)).}
\end{center}
\end{adjustwidth}
\newpage
\clearpage\ifodd\value{page}\hbox{}\newpage\fi% --- FOLIO 24 (f. 174r) : SPREAD 2.1 (Texte grec établi + Traduction française + Commentaire 1/2 & 2/2) ---
\clearpage\ifodd\value{page}\hbox{}\newpage\fi
\enTeteG{ΘΕΣΜΟΦΟΡΙΑΖΟΥΣΑΙ --- Texte grec établi (f. 174r, vv. 959–1000)}
\enTeteD{LES THESMOPHORIES --- Traduction française (f. 174r, vv. 959–1000)}
\begin{adjustwidth}{-2cm}{-2cm}
\noindent\textbf{\footnotesize Texte grec établi --- Folio 24 (\texttt{f. 174r}, vv. 959–1000) \hfill ΧΟΡΙΚΟΝ}\\[-2pt]
\noindent\rule{16.5cm}{0.5pt}\\[2pt]
{\fontsize{7.6pt}{8.7pt}\selectfont
\ligneEst{{\tiny\textit{(στροφή α) / cret. $+$ lec.}}}{ἅμα δὲ καὶ γένος Ὀλυμπίων θεῶν}{959}
\ligneEst{{\tiny\textit{$4$ troch\cat}}}{μέλπε καὶ γέραιρε φωνῇ πᾶσα χορομανεῖ τρόπῳ.}{960}
\ligneEst{{\tiny\textit{(ἀντιστρ. α) / cret. $+$ lec.}}}{εἰ δέ τις προσδοκᾷ κακῶς ἐρεῖν}{}
\ligneEst{{\tiny\textit{$4$ troch\cat}}}{ἐν ἱερῷ γυναῖκά μ᾽ οὖσαν ἄνδρας, οὐκ ὀρθῶς φρονεῖ.}{965}
\ligneEst{{\tiny\textit{(μεσῳδός α) / cret.}}}{ἀλλὰ χρή}{}
\ligneEst{{\tiny\textit{dim. troch.}}}{ὥσπερ ἔργον αὖ τι καινὸν}{}
\ligneEst{{\tiny\textit{$4$ troch\cat}}}{{\Taille{11.5}πρῶτον εὐκύκλου χορείας εὐφυᾶ στῆσαι βάσιν.}}{}
\ligneEst{{\tiny\textit{(στροφή β) / dim. iamb.}}}{πρόβαινε ποσὶ τὸν εὐλύραν}{}
\ligneEst{{\tiny\textit{dim. iamb.}}}{μέλπουσα καὶ τὴν τοξοφόρον}{970}
\ligneEst{{\tiny\textit{$2$ iamb\cat}}}{Ἄρτεμιν ἄνασσαν ἁγνήν.}{}
\ligneEst{{\tiny\textit{reiz.}}}{χαῖρ᾽ ὦ ἑκάεργε,}{}
\ligneEst{{\tiny\textit{reiz.}}}{ὄπαζε δὲ νίκην·}{}
\ligneEst{{\tiny\textit{$2$ iamb\cat}}}{Ἥραν τε τὴν τελείαν}{}
\ligneEst{{\tiny\textit{$2$ iamb\cat}}}{μέλψωμεν ὥσπερ εἰκός,}{975}
\ligneEst{{\tiny\textit{trim. iamb.}}}{ἣ πᾶσι τοῖς χοροῖσιν ἐμπαίζει τε καὶ}{}
\ligneEst{{\tiny\textit{$2$ iamb\cat}}}{κλῇδας γάμου φυλάττει.}{}
\ligneEst{{\tiny\textit{(ἀντιστρ. β) / dim. iamb.}}}{Ἑρμῆν τε νόμιον ἄντομαι}{}
\ligneEst{{\tiny\textit{dim. iamb.}}}{καὶ Πᾶνα καὶ Νύμφας φίλας}{}
\ligneEst{{\tiny\textit{$2$ iamb\cat}}}{ἐπιγελάσαι προθύμως}{}
\ligneEst{{\tiny\textit{reiz.}}}{ταῖς ἡμετέραισι}{980}
\ligneEst{{\tiny\textit{reiz.}}}{χαρέντα χορείαις.}{}
\ligneEst{{\tiny\textit{$2$ iamb\cat}}}{ἔξαιρε δὴ προθύμως}{}
\ligneEst{{\tiny\textit{$2$ iamb\cat}}}{διπλῆν χάριν χορείας.}{}
\ligneEst{{\tiny\textit{trim. iamb.}}}{παίσωμεν ὦ γυναῖκες οἷάπερ νόμος,}{}
\ligneEst{{\tiny\textit{$2$ iamb\cat}}}{νηστεύομεν δὲ πάντως.}{}
\ligneEst{{\tiny\textit{(μεσῳδός β) / trim. iamb.}}}{ἀλλ᾽ εἶα πάλλ᾽ ἀνάστρεφ᾽ εὐρύθμῳ ποδί,}{985}
\ligneEst{{\tiny\textit{$2$ iamb\cat}}}{τόρευε πᾶσαν ᾠδήν·}{}
\ligneEst{{\tiny\textit{ia. $+$ mol.}}}{ἡγοῦ δέ γ᾽ ὧδ᾽ αὐτὸς σὺ}{}
\ligneEst{{\tiny\textit{cret. $+$ mol.}}}{κισσοφόρε Βακχεῖε}{}
\ligneEst{{\tiny\textit{aristoph.}}}{δέσποτ᾽· ἐγὼ δὲ κώμοις}{}
\ligneEst{{\tiny\textit{$2$ iamb\cat}}}{σὲ φιλοχόροισι μέλψω.}{}
\ligneEst{{\tiny\textit{(στροφή γ) / aristoph.}}}{†εὔιον ὦ Διόνυσε}{990}
\ligneEst{{\tiny\textit{pherecr.}}}{Βρόμιε καὶ Σεμέλας παῖ,}{}
\ligneEst{{\tiny\textit{hexasyll. aeol.}}}{χοροῖς τερπόμενος}{}
\ligneEst{{\tiny\textit{ia. $+$ arist.}}}{κατ᾽ ὄρεα νυμφᾶν ἐρατοῖς ἐν ὕμνοις,}{}
\ligneEst{{\tiny\textit{$2$ iamb\cat}}}{†εὔιον εὔιον εὐοῖ}{}
\ligneEst{{\tiny\textit{$2$ iamb\cat}}}{< $ \times \ \widearc{\cup \cup} \ \cup $ > ἀναχορεύων.}{}
\ligneEst{{\tiny\textit{(ἀντιστρ. γ) / aristoph.}}}{ἀμφὶ δὲ σοὶ κτυπεῖται}{995}
\ligneEst{{\tiny\textit{pherecr.}}}{Κιθαιρώνιος ἠχώ,}{}
\ligneEst{{\tiny\textit{hexasyll. aeol.}}}{μελάμφυλλά τ᾽ ὄρη}{}
\ligneEst{{\tiny\textit{(?)}}}{δάσκια †καὶ νάπαι πετρώδεις† βρέμονται·}{}
\ligneEst{{\tiny\textit{$2$ iamb\cat}}}{κύκλῳ δὲ περὶ σὲ κισσὸς}{}
\ligneEst{{\tiny\textit{$2$ iamb\cat}}}{εὐπέταλος ἕλικι θάλλει.}{1000}
}
\vspace{2pt}
\noindent\rule{16.5cm}{0.6pt}\\[2pt]
\noindent\textbf{\small Commentaire philologique, métrique et scénique (vv. 959–1000 --- partie 1/2)}\\[2pt]
\begin{adjustbox}{max totalheight=9.8cm, max width=16.5cm, keepaspectratio}%
\parbox[t]{16.5cm}{
\setlength{\columnsep}{9pt}%
\begin{multicols}{2}
\fontsize{8.4pt}{9.7pt}\selectfont
\setlength{\parindent}{0pt}%
\setlength{\parskip}{1.2pt}%
\noindent\textbf{982} \textbf{διπλῆν.} Pollux (IV, 105) compte la διπλῆ parmi les figures de danse pour la tragédie (τραγικῆς ὀρχήσεως σχήματα) :
\citComm{Καὶ μὲν τραγικῆς ὀρχήσεως σχήματα σιμή χείρ, καλαθίσκος, χεὶρ καταπρανής, ξύλου παράληψις, \textbf{διπλῆ}, θερμαυστρίς, κυβίστησις, παραβῆναι τὰ τέτταρα}
\end{multicols}
}%
\end{adjustbox}%
\vfill
\end{adjustwidth}
\newpage
% --- PAGE DE DROITE : Traduction française + Commentaire (2/2) ---
\begin{adjustwidth}{-2cm}{-2cm}
\noindent\textbf{\footnotesize Traduction française --- Folio 24 (\texttt{f. 174r}, vv. 959–1000) \hfill ΧΟΡΙΚΟΝ}\\[-2pt]
\noindent\rule{16.5cm}{0.5pt}\\[2pt]
{\fontsize{7.6pt}{8.7pt}\selectfont
\ligneEst{}{\francais{Et ensemble célèbre et honore de la voix la gent des dieux olympiens,}}{959}
\ligneEst{}{\francais{d'une âme transportée par la danse.}}{960}
\ligneEst{}{\francais{Que si tel croit qu'étant femme je dois médire}}{}
\ligneEst{}{\francais{des hommes pendant nos rites, il s'abuse.}}{965}
\ligneEst{}{\francais{Mais d'abord il convient,}}{}
\ligneEst{}{\francais{c'est la règle,}}{}
\ligneEst{}{\francais{de procéder aux figures élégantes de la danse cyclique. (Danse.)}}{}
\ligneEst{}{\francais{Pars du pied en célébrant le dieu à la belle lyre}}{}
\ligneEst{}{\francais{et Artémis armée de l'arc, la chaste souveraine.}}{970}
\ligneEst{}{\francais{Salut, dieu tutélaire,}}{}
\ligneEst{}{\francais{et donne-nous la victoire.}}{}
\ligneEst{}{\francais{Puis célébrons, comme il sied,}}{}
\ligneEst{}{\francais{Héra la conjugale,}}{}
\ligneEst{}{\francais{qui se récrée}}{975}
\ligneEst{}{\francais{avec tous les chœurs et}}{}
\ligneEst{}{\francais{garde les clefs de l'hymen.}}{}
\ligneEst{}{\francais{Et Hermès pastoral,}}{}
\ligneEst{}{\francais{et Pan et les Nymphes chéries,}}{}
\ligneEst{}{\francais{je les prie de sourire avec bienveillance}}{}
\ligneEst{}{\francais{et de se plaire}}{980}
\ligneEst{}{\francais{à nos danses.}}{}
\ligneEst{}{\francais{Allons, commence avec ardeur}}{}
\ligneEst{}{\francais{la diple, joie de la danse.}}{}
\ligneEst{}{\francais{Ébattons-nous, femmes, suivant l'usage.}}{}
\ligneEst{}{\francais{Aussi bien nous sommes à jeun.}}{}
\ligneEst{}{\francais{Or sus! marque une autre mesure de ton pied cadencé.}}{985}
\ligneEst{}{\francais{Chante à pleine voix.}}{}
\ligneEst{}{\francais{Et guide-nous toi-même,}}{}
\ligneEst{}{\francais{couronné de lierre, seigneur Bacchus.}}{}
\ligneEst{}{\francais{Je te célébrerai}}{}
\ligneEst{}{\francais{en des chants bien rythmés.}}{}
\ligneEst{}{\francais{Évios! fils de Zeus}}{990}
\ligneEst{}{\francais{et de Sémélé, ô Bromios,}}{}
\ligneEst{}{\francais{toi que les chœurs délectent,}}{}
\ligneEst{}{\francais{et qui dans les montagnes, parmi les hymnes aimables des Nymphes,}}{}
\ligneEst{}{\francais{Évios ! Évios ! Évoé !}}{}
\ligneEst{}{\francais{te livres à la danse.}}{}
\ligneEst{}{\francais{Autour de toi}}{995}
\ligneEst{}{\francais{l'écho du Cithéron résonne,}}{}
\ligneEst{}{\francais{et retentissent les montagnes ombreuses}}{}
\ligneEst{}{\francais{et les ravins rocheux,}}{}
\ligneEst{}{\francais{tandis qu'à tes côtés le lierre}}{}
\ligneEst{}{\francais{aux belles feuilles pousse sa vrille tout à l'entour.}}{1000}
}
\vspace{2pt}
\noindent\rule{16.5cm}{0.6pt}\\[2pt]
\noindent\textbf{\small Commentaire philologique, métrique et scénique (vv. 959–1000 --- partie 2/2)}\\[2pt]
\begin{adjustbox}{max totalheight=9.8cm, max width=16.5cm, keepaspectratio}%
\parbox[t]{16.5cm}{
\setlength{\columnsep}{9pt}%
\begin{multicols}{2}
\fontsize{8.4pt}{9.7pt}\selectfont
\setlength{\parindent}{0pt}%
\setlength{\parskip}{1.2pt}%
\noindent Hésychius (δ 1948) ne nous apprend rien davantage :
\citComm{διπλῆ· ὀρχήσεως εἶδος, ἢ κρούματος.}
\noindent Il nous est difficile de reconstituer, même grossièrement, la chorégraphie de ce choral puisqu'Aristophane recourt au vocabulaire technique de la danse de son époque, dont on ne sait plus grand chose.
\end{multicols}
}%
\end{adjustbox}%
\vfill
\end{adjustwidth}
\newpage
\clearpage\ifodd\value{page}\hbox{}\newpage\fi
\enTeteG{ΕΠΕΙΣΟΔΙΟΝ Ε' (Ἀνδρομέδα) --- Codex Ravennas 429, f. 174v (Classense 176v) --- vv. 1001–1043}
\begin{adjustwidth}{-1.5cm}{-3.0cm}
\begin{adjustbox}{max totalheight=21.0cm, max width=16.5cm, keepaspectratio}
\begin{minipage}{16.5cm}
{\fontsize{7.2pt}{8.5pt}\selectfont
\noindent\textbf{Folio 25 --- Codex Ravennas 429, \texttt{f. 174v} (\texttt{Classense 176v}) : ΕΠΕΙΣΟΔΙΟΝ Ε' (Ἀνδρομέδα), vv. 1001–1043}\hfill \textit{Témoins :} $\text{R}^1, \Rpc, \Rdeux$, \PsiQuatre\par\vspace{2pt}
\hrule height 0.4pt\vspace{2.5pt}
}
{\fontsize{8.1pt}{9.5pt}\selectfont
\ligneRav{1001}{τοξ/}{ἐνταῦτα νῦν \difference{οἰμῶξι} πρὸς τὴν αἰτρίαν·}{\textbf{1001} οἰμώξει R : \lecon{οἰμῶξι} Brunck}
\ligneRav{1002}{μ\auDessus{ν}{Η}}{ὦ τοξόθ᾽ ἱκετεύω σε. : μή μ᾽ ἰκετεῦσι σύ·}{\textbf{1002} ἱκετεῦση R : \lecon{ἱκετεῦσι} Brunck}
\ligneRav{1003}{μ\auDessus{ν}{Η}}{χάλασον τὸν ἧλον. : ἀλλὰ ταῦτα \difference{δρᾶσ᾽} ἐγώ·}{\textbf{1003} δρᾷς R : \lecon{δρᾶσ᾽} Voss}
\ligneRav{1004}{μ\auDessus{ν}{Η}}{οἴμοι κακοδαίμων, μᾶλλον \difference{ἐπικρούεις} σύ γε·}{\textbf{1004} ἐπικρούσεις R : \lecon{ἐπικρούεις} \Mon}
\ligneRav{1005}{τοξ/}{ἔτι \difference{μᾶλλο} βοῦλις; : ἀτταταῖ \difference{ἰατταταῖ}·}{\textbf{1005} μᾶλλον R : \lecon{μᾶλλο} Bentley \espApp τατταταί R : \lecon{ἰατταταῖ} Faber}
\ligneRav{1006}{—}{κακοδαίμω\expo{ν} γέρ\inRas{ω}{ο}ν·}{\textbf{1006} κακοδαίμω\expo{ν} γέρ\inRas{ω}{ο}ν \Rpc, M\expo{pc}}
\ligneRav{1007}{}{πέρ᾽ ἐγὼ \difference{'ξινίγκι} πορμός, \difference{ἴνα} πυλάξι σοι·}{\textbf{1007} ξεινίγκι R : \lecon{'ξινίγκι} Enger \espApp ἵνα R : \lecon{ἴνα} Brunck \espApp \lecon{πυλάξι} R : πυλάξῃ Σ}
\ligneRav{}{μ\auDessus{ν}{Η}}{ταυτὶ τὰ βέλτιστ᾽ ἀπολέλαυκ᾽ Εὐριπίδου·}{}
\ligneRav{1009}{}{ἔα· θεοί, Ζεῦ σῶτερ, εἰσὶν ἐλπίδες·}{\textbf{1009} εἰσὶν R : εἴς ἔτ᾽ Naber (Faber ?) \cf 946}
\ligneRav{1010}{}{\difference{ἁνὴρ} ἔοικεν οὐ προδώσειν, ἀλλά μοι}{\textbf{1010} ἀνὴρ R : \lecon{ἁνὴρ} Bothe}
\ligneRav{1011}{}{σημεῖον \difference{ὑπεδήλωσε} Περσεὺς ἐκδραμών,}{\textbf{1011} ὑπεδήλωσεν R : \lecon{ὑπεδήλωσε} Biset}
\ligneRav{}{}{ὅτι δεῖ με γίγνεσθ᾽ Ἀνδρομέδαν· πάντως δέ μοι}{}
\ligneRav{1013}{}{τὰ δέσμ᾽ ὑπάρχει. δῆλον \difference{οὖν τοῦτ᾽} ἔσθ᾽ ὅτι}{\textbf{1013} οὖν R : οὖν νῦν Bizet : \lecon{οὖν τοῦτ᾽} Dobree}
\ligneRav{1014}{}{ἥξει με \difference{σώσων}· οὐ γὰρ ἂν παρέπτετο·}{\textbf{1014} σῶσον R : \lecon{σώσων} \Mon}
\ligneRav{1015}{\auDessus{εὐ}{ρ}}{\hspace{1.1em}φίλαι παρθένοι φίλαι,}{\textbf{1015-21} (Εὐρ.) R : \lecon{(Μνη.)} Tyrwhitt}
\ligneRav{1016}{}{\hspace{1.1em}πῶς ἂν ἀπέλθοιμι καὶ}{\textbf{1016} \lecon{ἀπέλθοιμι} R : ἐπέλθοιμι Ellebodius}
\ligneRav{1017}{}{\hspace{1.1em}τὸν Σκύθην \difference{λάθοιμι};}{\textbf{1017} λάβοιμι R : \lecon{λάθοιμι} Ellebodius}
\ligneRav{1019}{}{\hspace{1.1em}κλύεις; ὦ \difference{προσᾴδουσ'} \difference{ἀϋτὰς} ἐν ἄντροις,}{\textbf{1019} προσᾳδοῦσσαι τὰς R, Σ : \lecon{προσᾴδουσ'} Elmsley, \lecon{ἀϋτὰς} Burges (ἀϋταῖς Sommerstein) : alii alia}
\ligneRav{1020}{}{\hspace{1.1em}κατάνευσον, ἔασον ὡς}{}
\ligneRav{}{}{\hspace{1.1em}τὴν γυναῖκά μ᾽ ἐλθεῖν·}{}
\ligneRav{1022}{μ\auDessus{ν}{Η}}{\hspace{1.1em}ἄνοικτος ὅς μ᾽ ἔδησε τὸν}{\textbf{1022} (Μνη.) praefixit \Rpc, \Mon}
\ligneRav{1023}{}{\hspace{1.1em}πολυπονώτατον βροτῶν·}{\textbf{1023} \lecon{πολυπονώτατον} R, Σ : πολυστονώτατον Burges}
\ligneRav{}{}{\hspace{1.1em}μόλις δὲ γραῖαν ἀποφυγὼν}{}
\ligneRav{1025}{}{\hspace{1.1em}σαπρὰν ἀπωλόμην ὅμως·}{}
\ligneRav{1026}{}{\hspace{1.1em}ὅδε γὰρ ὁ Σκύθης φύλαξ πάλαι ἐφέστηκ᾽· ὀλοὸν}{\textbf{1026} ὁ Σκύθης R : del. Jackson}
\ligneRav{}{}{\hspace{1.1em}ἄφιλον ἐκρέμασε κόραξι δεῖπνον·}{}
\ligneRav{1030}{}{\hspace{1.1em}ὁρᾷς; οὐ χοροῖσιν οὐδ᾽ ὑφ᾽ ἡλίκων νεανίδων}{}
\ligneRav{1031}{}{\hspace{1.1em}ψῆφον \difference{ἕστηκ᾽ ἔχουσα κημόν},}{\textbf{1031} κημὸν ἕστηκ᾽ ἔχουσ᾽ R : \lecon{ἕστηκ᾽ ἔχουσα κημόν} Bibi, cf. Σ ad \italique{Vesp.} 674c}
\ligneRav{}{}{\hspace{1.1em}ἀλλ᾽ ἐν πυκνοῖς δεσμοῖσιν ἐμπεπλεγμένη}{}
\ligneRav{}{}{\hspace{1.1em}κήτει βορὰ Γλαυκέτῃ πρόκειμαι,}{}
\ligneRav{1034}{}{\hspace{1.1em}γαμηλίῳ μὲν οὐ ξὺν}{\textbf{1034-5} ξυμπαιῶνι R : \lecon{ξὺν παιῶνι} Zanetti (\italique{cum carmuine} Divus)}
\ligneRav{1035}{}{\hspace{1.1em}παιῶνι δεσμίῳ δὲ}{}
\ligneRav{}{}{\hspace{1.1em}γοᾶσθέ μ᾽ ὦ γυναῖκες ὡς}{}
\ligneRav{}{}{\hspace{1.1em}μέλεα μὲν πέπονθα μέλεος,}{}
\ligneRav{}{}{\hspace{1.1em}ὢ τάλας ἐγὼ τάλας,}{}
\ligneRav{1039}{}{\hspace{1.1em}ἀπὸ δὲ συγγόνων \difference{ἀλλ᾽} ἄνομα πάθεα,}{\textbf{1039} ἀλλὰν R : \lecon{ἀλλ᾽} Scaliger}
\ligneRav{1040}{}{\hspace{1.1em}†φῶτά \difference{λιτομένα}†,}{\textbf{1040} λιτομέναν R : ἀντομένα (vel -αν) Σ : \lecon{λιτομένα} Enger}
\ligneRav{1041}{}{\hspace{1.1em}πολυδάκρυτον Ἀίδα γόον \difference{ψεφαυγοῦς}·}{\textbf{1041} φεύγουσαν R : χέουσα Rau (\italique{Paratragödia}, p.75) : φλέγουσαν Mugrave (ad Eur. \italique{Or.}1402$=$1396) : \lecon{ψεφαυγοῦς} Bibi, cf Σ ad v.1054, Hesych. ψ 135 \espApp φεύγουσαν R : χέουσα Rau (\italique{Paratragödia}, p.75) : φλέγουσαν Mugrave (ad Eur. \italique{Or.}1402$=$1396) : \lecon{ψεφαυγοῦς} Bibi, cf Σ ad v.1054 e Hesych. ψ 135}
\ligneRav{}{}{\hspace{1.1em}αἰαῖ αἰαῖ ἒ ἒ,}{}
\ligneRav{}{}{\hspace{1.1em}ὃς ἔμ᾽ ἀπεξύρησε πρῶτον,}{}
}
\vspace{3pt}\noindent\rule{4.5cm}{0.4pt}\par\vspace{1.5pt}
{\fontsize{7.0pt}{8.3pt}\selectfont
\noindent\textbf{Loci paralleli (vv. 1001–1043) ---} Tradition directe : Codex Ravennas 429 (R) et scholies anciennes ($\Sigma^{\text{R}}$).\par
}
\vspace{2.5pt}\noindent\rule{4.5cm}{0.4pt}\par\vspace{1.5pt}
{\fontsize{6.9pt}{8.2pt}\selectfont
\noindent\textbf{Suite (en commentaire) de l'apparat critique \& notes paléographiques (\texttt{f. 174v}) ---} \textbf{Témoin papyrologique :} \textbf{\PsiQuatre{} (\italique{P.Oxy.} LXXIII 4935, fr. 1)} débute au \textbf{v. 1043}. --- \textbf{Traitement du grec barbare de l'Archer scythe (vv. 1001–1007) :} $\text{R}^1$ tend parfois à normaliser les finales scythes en grec standard (οἰμώξει, ἱκετεῦση, μᾶλλον, ἵνα) ; au \textbf{v. 1006}, \Rpc{} (suivi de \Mon\expo{p.c.}) ajoute un $\nu$ en exposant et corrige \italique{in rasura} l'oméga de γέρων en omicron : \lecon{κακοδαίμω\expo{ν} γέρ\inRas{ω}{ο}ν}. --- \textbf{Attribution des vv. 1015–1021 :} R attribue le début de la monodie à (Εὐρ.), corrigé en \lecon{(Μνη.)} par Tyrwhitt, tandis qu'au \textbf{v. 1022} \Rpc{} (et \Mon{}) préfixe (Μνη.). --- \textbf{Deux interventions décisives de \textbf{Bibi} (vv. 1031 et 1041) :} --- \textbf{v. 1007 :} ξεινίγκι R : \lecon{'ξινίγκι} Enger \espApp ἵνα R : \lecon{ἴνα} Brunck \espApp \lecon{πυλάξι} R (πυλάξῃ $\Sigma$). --- \textbf{v. 1017 :} λάβοιμι R (confusion minuscule $\theta/\beta$) : \lecon{λάθοιμι} Ellebodius. --- \textbf{v. 1019 :} προσᾳδοῦσσαι τὰς R, $\Sigma$ : \lecon{προσᾴδουσ'} Elmsley, \lecon{ἀϋτὰς} Burges (ἀϋταῖς Sommerstein). --- \textbf{v. 1023 :} \lecon{πολυπονώτατον} R, $\Sigma$ (πολυστονώτατον Burges). --- \textbf{v. 1031 :} κημὸν ἕστηκ' ἔχουσ' R : \lecon{ἕστηκ' ἔχουσα κημόν} \lecon{Bibi}, cf. $\Sigma$ ad \italique{Vesp.} 674c. --- \textbf{v. 1040 :} λιτομέναν R (ἀντομένα vel -αν $\Sigma$, confusion onciale ΛΙ/ΑΝ) : \lecon{λιτομένα} Enger. --- \textbf{v. 1041 :} φεύγουσαν R (χέουσα Rau ; φλέγουσαν Musgrave) : \lecon{ψεφαυγοῦς} \lecon{Bibi}, cf. $\Sigma$ ad v. 1054, Hésychius $\psi\ 135$.\par
}
\end{minipage}
\end{adjustbox}
\end{adjustwidth}
\newpage
\enTeteD{Codex Ravennas 429 (Biblioteca Classense) --- Fac-similé du f. 174v (vv. 1001–1043)}
\begin{adjustwidth}{-1.5cm}{-3.0cm}
\begin{center}
\includegraphics[width=16.5cm,height=20.5cm,keepaspectratio]{/workspaces/Thesmophories/dossier_ravennas_thesmophories/folios_individuels_ff162v-177r/25_f174v_classense176v_vv1001-1043.jpg}\par
\vspace{2pt}
{\scriptsize\textbf{Codex Ravennas 429, f. 174v (Classense 176v)} --- ΕΠΕΙΣΟΔΙΟΝ Ε' (Ἀνδρομέδα), \textbf{vv. 1001–1043} ($\text{R}^1, \Rpc, \Rdeux$, \PsiQuatre).}
\end{center}
\end{adjustwidth}
\newpage
\clearpage\ifodd\value{page}\hbox{}\newpage\fi% --- FOLIO 25 (f. 174v) : SPREAD 2.1 (Texte grec établi + Traduction française + Commentaire 1/8 & 2/8) ---
\clearpage\ifodd\value{page}\hbox{}\newpage\fi
\enTeteG{ΘΕΣΜΟΦΟΡΙΑΖΟΥΣΑΙ --- Texte grec établi (f. 174v, vv. 1001–1043 [1/4])}
\enTeteD{LES THESMOPHORIES --- Traduction française (f. 174v, vv. 1001–1043 [1/4])}
\begin{adjustwidth}{-2cm}{-2cm}
\noindent\textbf{\footnotesize Texte grec établi --- Folio 25 (\texttt{f. 174v}, vv. 1001–1043 [1/4]) \hfill ΕΠΕΙΣΟΔΙΟΝ Ε' (Ἀνδρομέδα)}\\[-2pt]
\noindent\rule{16.5cm}{0.5pt}\\[2pt]
{\fontsize{8.8pt}{10.2pt}\selectfont
\ligneEst{\textbf{Τοξότης} {\tiny\textit{trim. iamb.}}}{ἐνταῦτα νῦν οἰμῶξι πρὸς τὴν αἰτρίαν.}{1001}
\ligneEst{\textbf{Μνησίλοχος}}{ὦ τοξόθ᾽ ἱκετεύω σε.}{}
\ligneEst{\textbf{Τοξότης}}{\decalage{ὦ τοξόθ᾽ ἱκετεύω σε.}μή μ᾽ ἰκετεῦσι σύ.}{}
\ligneEst{\textbf{Μνησίλοχος}}{χάλασον τὸν ἧλον.}{}
\ligneEst{\textbf{Τοξότης}}{\decalage{χάλασον τὸν ἧλον.}ἀλλὰ ταῦτα δρᾶσ᾽ ἐγώ.}{}
\ligneEst{\textbf{Μνησίλοχος}}{οἴμοι κακοδαίμων, μᾶλλον ἐπικρούεις σύ γε.}{}
\ligneEst{\textbf{Τοξότης}}{ἔτι μᾶλλο βοῦλις;}{1005}
\ligneEst{\textbf{Μνησίλοχος}}{\decalage{ἔτι μᾶλλο βοῦλις;}ἀτταταῖ ἰατταταῖ·}{}
\ligneEst{}{κακῶς ἀπόλοιο.}{}
\ligneEst{\textbf{Τοξότης}}{\decalage{κακῶς ἀπόλοιο.}σῖγα κακοδαίμων γέρον.}{}
\ligneEst{}{πέρ᾽ ἐγὼ 'ξινίγκι πορμός, ἴνα πυλάξι σοι.}{}
}
\vspace{2pt}
\noindent\rule{16.5cm}{0.6pt}\\[2pt]
\noindent\textbf{\small Commentaire philologique, métrique et scénique (vv. 1001–1043 --- partie 1/8)}\\[2pt]
\begin{adjustbox}{max totalheight=15.2cm, max width=16.5cm, keepaspectratio}%
\parbox[t]{16.5cm}{
\setlength{\columnsep}{9pt}%
\begin{multicols}{2}
\fontsize{7.0pt}{8.0pt}\selectfont
\setlength{\parindent}{0pt}%
\setlength{\parskip}{1.2pt}%
\noindent\textbf{1001} \textbf{τοξότης.} Aristophane a donné à l'archer un jargon mi-scythe mi-grec, une sorte de parler nègre. On sent qu'il est impossible de rendre cela dans une traduction. (Note de Willems.)\\ À Athènes, comme partout ailleurs, les personnes chargées des missions de police étaient recrutées parmi la lie de l'humanité. En l'occurrence, il y avait en guise de police un contingent de 300 mercenaires Scythes aux ordres des prytanes\marqueur{1}. C'est cette inquiétant contraste, entre l'importance de la mission (les forces de l'«~ordre~», comme on dit...) et la bassesse de ceux qui y sont préposés, qui inspirera à Platon de si longs développements sur le bon choix et l'éducation des gardiens de sa Cité. \latin{Quis custodiet custodes?}
\noindent\textbf{1001} \textbf{ἐνταῦτα νῦν οἰμῶξι πρὸς τὴν αἰτρίαν} $=$ ἐνταῦθα νῦν οἴμωζε πρὸς τὴν αἰθρίαν.
\noindent (Noter le sadisme proprement argousin de l'archer.)
\noindent\textbf{[Notes ad 1001]} \textbf{(1)} Que les archers Scythes fussent directement aux ordres des prytanes, c'est assez bien établi par notre pièce ainsi que le vers 54 des \italique{Acharniens}, cf. aussi \italique{Protagoras} (319c).
\noindent On pourra lire l'article de Jean-Christophe Couvenhes «~L'introduction des archers scythes, esclaves publics, à Athènes~», in \italique{Transferts culturels et droits dans le monde grec et hellénistique} (2012), pp. 99-118, ou l'article plus ancien d'Agnès Plassart «~Les archers d'Athènes~», in \italique{Revue des Études Grecques}, t. 26 (1913), pp. 151-213 dont quelques détails sont discutés dans le premier article.
\noindent\textbf{1007} \textbf{πέρ᾽ ἐγὼ 'ξινίγκι πορμός, ἴνα πυλάξι σοι} $=$ φέρ᾽ ἐγὼ 'ξενέγκω φορμόν, ἵνα φυλάξω σε.
\noindent\textbf{1015-135} \textbf{Parodie de l'\italique{Andromède}.} La pièce perdue d'Euripide \italique{Andromède} fut peut-être représentée la même année que l'\italique{Hélène}\marqueur{1}, en 412, un avant les \italique{Thesmophories}. L'édition des fragments d'Euripide de Jouan \& van Looy\marqueur{2} propose un dossier complet auquel le lecteur pourra se reporter. Le récit le plus circonstancié du mythe d'Andromède se trouve chez Apollodore (II, 4, 3) :
\citComm{Παραγενόμενος δὲ εἰς Αἰθιοπίαν, ἧς ἐβασίλευε Κηφεύς, εὗρε τὴν τούτου θυγατέρα Ἀνδρομέδαν παρακειμένην βορὰν θαλασσίῳ κήτει. Κασσιέπεια γὰρ ἡ Κηφέως γυνὴ Νηρηίσιν ἤρισε περὶ κάλλους, καὶ πασῶν εἶναι κρείσσων ηὔχησεν· ὅθεν αἱ Νηρηίδες ἐμήνισαν, καὶ Ποσειδῶν αὐταῖς συνοργισθεὶς πλήμμυράν τε ἐπὶ τὴν χώραν ἔπεμψε καὶ κῆτος. Ἄμμωνος δὲ χρήσαντος τὴν ἀπαλλαγὴν τῆς συμφορᾶς, ἐὰν ἡ Κασσιεπείας θυγάτηρ Ἀνδρομέδα προτεθῇ τῷ κήτει βορά, τοῦτο ἀναγκασθεὶς ὁ Κηφεὺς ὑπὸ τῶν Αἰθιόπων ἔπραξε, καὶ προσέδησε τὴν θυγατέρα πέτρα. Ταύτην θεασάμενος ὁ Περσεὺς καὶ ἐρασθεὶς ἀναιρήσειν ὑπέσχετο Κηφεῖ τὸ κῆτος, εἰ μέλλει σωθεῖσαν αὐτὴν αὐτῷ δώσειν γυναῖκα. Ἐπὶ τούτοις γενομένων ὅρκων, ὑποστὰς τὸ κῆτος ἔκτεινε καὶ τὴν Ἀνδρομέδαν ἔλυσεν. Ἐπιβουλεύοντος δὲ αὐτῷ Φινέως, ὃς ἦν ἀδελφὸς τοῦ Κηφέως ἐγγεγυημένος πρῶτος τὴν Ἀνδρομέδαν, μαθὼν τὴν ἐπιβουλήν, τὴν Γοργόνα δείξας μετὰ τῶν συνεπιβουλευόντων αὐτὸν ἐλίθωσε παραχρῆμα.\\
Parvenu en Éthiopie, dont Céphée était roi, [Persée] trouva sa fille Andromède exposée pour être dévorée par un monstre marin. Cassiopée, épouse de Céphée, avait osé se comparer aux Néréides pour la beauté, et s'était même vantée de l'emporter sur elles. Les Néréides en furent irritées. Neptune partagea leur indignation, submergea le pays, et y envoya un monstre marin. L'oracle d'Ammon ayant annoncé que ces désastres cesseraient si on exposait Andromède, fille de Cassiopée, pour être dévorée par le monstre, les Éthiopiens forcèrent Céphée à faire ce que l'oracle ordonnait et à attacher sa fille à un rocher. Persée l'ayant vu, il devint amoureux, et promit à Céphée de tuer le monstre s'il voulait la lui donner en mariage. Céphée s'y étant engagé par serment, il attendit le monstre, le tua et délivra Andromède. Phinée, frère de Céphée, à qui Andromède avait été promise avant cet événement, conspira contre lui pour le faire périr ; mais Persée l'ayant découvert, lui montra la tête de la Gorgone, et le changea en pierre, ainsi que tous ceux qui avaient pris part à son complot.}
\end{multicols}
}%
\end{adjustbox}%
\vfill
\end{adjustwidth}
\newpage
% --- PAGE DE DROITE : Traduction française + Commentaire (2/8) ---
\begin{adjustwidth}{-2cm}{-2cm}
\noindent\textbf{\footnotesize Traduction française --- Folio 25 (\texttt{f. 174v}, vv. 1001–1043 [1/4]) \hfill ΕΠΕΙΣΟΔΙΟΝ Ε' (Ἀνδρομέδα)}\\[-2pt]
\noindent\rule{16.5cm}{0.5pt}\\[2pt]
{\fontsize{8.8pt}{10.2pt}\selectfont
\ligneEst{\textbf{Archer Scythe}}{\francais{Ici maintenant, tu vas crier, en plein air.}}{1001}
\ligneEst{\textbf{Mnésiloque}}{\francais{Archer, je t'en supplie...}}{}
\ligneEst{\textbf{Archer}}{\francais{\espace{3.1cm}Ne me supplie pas.}}{}
\ligneEst{\textbf{Mnésiloque}}{\francais{Desserre la cheville.}}{}
\ligneEst{\textbf{Archer}}{\francais{\espace{2.9cm}C'est ce que je vais faire}}{}
\ligneEst{\textbf{Mnésiloque}}{\francais{Aïe, misère! tu la pousses davantage.}}{}
\ligneEst{\textbf{Archer}}{\francais{Encore plus, veux-tu?}}{1005}
\ligneEst{\textbf{Mnésiloque}}{\francais{\espace{2.7cm}Aie, aïe, aïe, aie, aïe, aie !}}{}
\ligneEst{}{\francais{la peste soit de toi.}}{}
\ligneEst{\textbf{Archer}}{\francais{\espace{2.5cm}Silence, misérable vieillard.}}{}
\ligneEst{}{\francais{Allons, je vais, moi, apporter une natte, pour faire ma garde. (Il sort.)}}{}
}
\vspace{2pt}
\noindent\rule{16.5cm}{0.6pt}\\[2pt]
\noindent\textbf{\small Commentaire philologique, métrique et scénique (vv. 1001–1043 --- partie 2/8)}\\[2pt]
\begin{adjustbox}{max totalheight=15.2cm, max width=16.5cm, keepaspectratio}%
\parbox[t]{16.5cm}{
\setlength{\columnsep}{9pt}%
\begin{multicols}{2}
\fontsize{8.4pt}{9.7pt}\selectfont
\setlength{\parindent}{0pt}%
\setlength{\parskip}{1.2pt}%
\noindent\textbf{[Notes ad 1015-135]} \textbf{(1)} C'est ce qu'indiquent les scholies au vers 1012, πιθανῶς· συνδεδίδακται γὰρ τῇ Ἑλένη, ainsi que celle du vers 53 des \italique{Grenouilles}. Le début du prologue d'\italique{Andromède} est repris au vers 1065, «~ὦ νὺξ ἱερὰ~», -- oui : Aristophane a chamboulé l'ordre de la pièce ; scholie à ce vers 1065 : τοῦ προλόγου Ἀνδρομέδας εἰσβολή --, laisse à penser qu'\italique{Andromède} fut jouée au lever du jour, avant donc \italique{Hélène}. Matthew Whright (\italique{Euripides' Escape-Tragedies}, 2005) fait l'hypothèse contraire et avance qu'\italique{Iphigénie en Tauride} concluait la trilogie d'Euripide de 412, agrémenté du drame satyrique du \italique{Cyclope}. Voir la recension que feu Alain Blanchard en fit dans la \italique{Revue des Études Grecques}, tome 119 (2006), pp. 798-800.
\noindent \textbf{(2)} Tome VIII (fragments), première partie, de l'édition d'Euripide aux Belles Lettres (1998).
\noindent\textbf{1015-55} \textbf{Monodie d'Andromède-Mnésiloque.} La structure de ce chant a fait couler beaucoup d'encre... et nous échappe encore\marqueur{1}. Mais bien loin de l'anarchie décrite par les commentateurs, je suis tenté d'y voir une illustration sonore des échos : chaque phrase musicale semble être répétée : l'on peut ainsi regrouper par paires successives les différentes structures métriques (le thème puis son écho, de mètre identique ou très proche). On regroupe ainsi les vers : 1015-1071 \& 1018-1021 ; 1022-1023 \& 1024-1025 ; 1026-1029 \& 1030-1031 ; 1032-1033 \& 1034-1035 ; 1036-1037 \& 1038-1039a ; 1043 \& 1045 ; 1050-1051 \& 1052-1053.
\noindent\textbf{[Notes ad 1015-55]} \textbf{(1)} On pourra aller voir ce qu'en dit Lætitia, pp. 436-445.
\noindent\textbf{1018} \textbf{ὦ προσᾳδοῦσσαι τὰς ἐν ἄντροις.} Comme dans la pièce d'Euripide, la nymphe Écho répond aux plaintes d'Andromède. (Dans la reprise comique, c'est Euripide qui se métamorphose en Écho.) Willems a consacré sa \nieme{25} note (reproduite en annexe) à ce vers et aux diverses retouches qu'il a subies.
\noindent Je reprends la lecture de Burges\marqueur{1}, προσᾴδουσ' ἀϋτάς, satisfaisante quant au sens et qui offre l'avantage de rester très proche du texte du manuscrit.
\noindent\textbf{[Notes ad 1018]} \textbf{(1)} G. Burges, \italique{Classical Journal} 14 (1816) 217. Ou bien voir W. Mitsdörffer, «~Das Mnesichoslied in Aristophanes' \italique{Thesmophoriazusen}~», \italique{Philologus} 98 (1954-55), p. 69.
\noindent\textbf{1026-8} \textbf{ὅδε γὰρ ὁ Σκύθης φύλαξ / πάλαι ἐφέστηκ᾽ ὀλοὸν / ἄφιλον ἐκρέμασε κόραξι δεῖπνον.} Ce passage est sans doute légèrement corrompu. Les tentatives pour le corriger ont toutes été infructueuses\marqueur{1}.
\end{multicols}
}%
\end{adjustbox}%
\vfill
\end{adjustwidth}
\newpage

% --- FOLIO 25 (f. 174v) : SPREAD 2.2 (Texte grec établi + Traduction française + Commentaire 3/8 & 4/8) ---
\clearpage\ifodd\value{page}\hbox{}\newpage\fi
\enTeteG{ΘΕΣΜΟΦΟΡΙΑΖΟΥΣΑΙ --- Texte grec établi (f. 174v, vv. 1001–1043 [2/4])}
\enTeteD{LES THESMOPHORIES --- Traduction française (f. 174v, vv. 1001–1043 [2/4])}
\begin{adjustwidth}{-2cm}{-2cm}
\noindent\textbf{\footnotesize Texte grec établi --- Folio 25 (\texttt{f. 174v}, vv. 1001–1043 [2/4]) \hfill ΕΠΕΙΣΟΔΙΟΝ Ε' (Ἀνδρομέδα)}\\[-2pt]
\noindent\rule{16.5cm}{0.5pt}\\[2pt]
{\fontsize{8.8pt}{10.2pt}\selectfont
\ligneEst{\textbf{Μνησίλοχος}}{ταυτὶ τὰ βέλτιστ᾽ ἀπολέλαυκ᾽ Εὐριπίδου.}{}
\ligneEst{}{ἔα· θεοί, Ζεῦ σῶτερ, εἰσὶν ἐλπίδες.}{}
\ligneEst{}{ἁνὴρ ἔοικεν οὐ προδώσειν, ἀλλά μοι}{1010}
\ligneEst{}{σημεῖον ὑπεδήλωσε Περσεὺς ἐκδραμών,}{}
\ligneEst{}{ὅτι δεῖ με γίγνεσθ᾽ Ἀνδρομέδαν· πάντως δέ μοι}{}
\ligneEst{}{τὰ δέσμ᾽ ὑπάρχει. δῆλον οὖν τοῦτ᾽ ἔσθ᾽ ὅτι}{}
\ligneEst{}{ἥξει με σώσων· οὐ γὰρ ἂν παρέπτετο.}{}
\ligneEst{\textbf{Μνησίλοχος} {\tiny\textit{bacch. $+$ ia.}}}{φίλαι παρθένοι φίλαι,}{1015}
\ligneEst{{\tiny\textit{chor. $+$ cret.}}}{πῶς ἂν ἀπέλθοιμι καὶ}{}
\ligneEst{{\tiny\textit{$+$ ityph.}}}{τὸν Σκύθην λάθοιμι;}{}
\ligneEst{{\tiny\textit{$4$ bacch.}}}{κλύεις; ὦ προσᾴδουσ' ἀϋτὰς ἐν ἄντροις,}{}
}
\vspace{2pt}
\noindent\rule{16.5cm}{0.6pt}\\[2pt]
\noindent\textbf{\small Commentaire philologique, métrique et scénique (vv. 1001–1043 --- partie 3/8)}\\[2pt]
\begin{adjustbox}{max totalheight=15.2cm, max width=16.5cm, keepaspectratio}%
\parbox[t]{16.5cm}{
\setlength{\columnsep}{9pt}%
\begin{multicols}{2}
\fontsize{8.4pt}{9.7pt}\selectfont
\setlength{\parindent}{0pt}%
\setlength{\parskip}{1.2pt}%
\noindent\textbf{[Notes ad 1026-8]} \textbf{(1)} Cf. l'inventaire sans concession qu'en fait Lætitia \italique{The Songs of Aristophanes}, p. 443-444 : son regard infaillible de métricienne aguerrie détecte les invraisemblances métriques des «~corrections~». Sa proposition n'en est pas moins aporétique : elle voit dans ces vers une suite de trois dochmiaques se terminant par un dimètre iambique catalectique tout en concédant que la correption πάλα\lettreBreve{ι} (\latin{vocalis ante vocalem corripitur}) est suspecte : les seuls autres exemples chez Aristophane étant au vers 915, «~κύσ\lettreBreve{ω}. ἄπαγε~», déjà évoqué plus haut, et le peu convaincant «~ἰὼ ἰώ~» au vers 1170 des \italique{Nuées}.
\noindent\textbf{1031} \textbf{ψῆφον κημὸν ἕστηκ᾽ ἔχουσ᾽.} On mettra ce vers en regard d'un fragment d'Hermippe :
\citComm{καὶ πρὸς κύβους ἕστηκ' ἔχων <τὸ> κήθιον.}
\noindent τὸ κήθιον et ὁ κημός sont synonymes : ils désignent dans l'univers ludique le \italique{cornet} dans lequel on secouait  les dés\marqueur{1}. ἡ ψῆφος, comme λίθος, est dans ce contexte un \italique{pion}.
\noindent On constate avec ces deux exemples que `ἕστηκεν ἔχων' devait être l'expression consacrée pour le joueur qui s'apprêtait à lancer les dés.
\noindent On a soupçonné une altération du vers du fait de l'étrange succession initiale de trois longues.
\noindent L'exemple d'Hermippe suscite naturellement l'émendation suivante :
\citComm{ὁρᾷς; οὐ χοροῖσιν οὐδ᾽ ὑφ᾽ ἡλίκων ν\synizese{εα}νίδων\\
ψῆφον ἕστηκ᾽ ἔχουσα κημόν.}
\noindent (κημός placé après ἔχω, les mots sont ainsi agencés dans le même ordre que dans le fragment\marqueur{2}.)\\ L'avantage est qu'on obtient (avec la synizèze de ν\synizese{εα}νίδων) le schéma métrique suivant :
\citComm{$ \cup \ - \ - \ \cup \ - \ | \ \cup \ - \ \cup \ - \ | \ \cup \ - \ - \ \cup \ -$\\
$ - \ \cup \ - \ | \ - \ \cup \ - \ | \ \cup \ - \ \smile $\\
dochmiaque $+$ iambe $+$ dochmiaque \\
crétique $+$ crétique $+$ bacchée,}
\noindent de scansion identique à celle des deux vers précédents, annulant du même coup le surprenant abrègement de πάλα\lettreBreve{ι} en seconde position d'une résolution (qui le faisait scander $ \widearc{\cup \cup}$, cf. commentaire v. 1026-8 n.1).
\end{multicols}
}%
\end{adjustbox}%
\vfill
\end{adjustwidth}
\newpage
% --- PAGE DE DROITE : Traduction française + Commentaire (4/8) ---
\begin{adjustwidth}{-2cm}{-2cm}
\noindent\textbf{\footnotesize Traduction française --- Folio 25 (\texttt{f. 174v}, vv. 1001–1043 [2/4]) \hfill ΕΠΕΙΣΟΔΙΟΝ Ε' (Ἀνδρομέδα)}\\[-2pt]
\noindent\rule{16.5cm}{0.5pt}\\[2pt]
{\fontsize{8.8pt}{10.2pt}\selectfont
\ligneEst{\textbf{Mnésiloque}}{\francais{Ces douceurs-là j'en ai l'obligation à Euripide.}}{}
\ligneEst{}{\francais{(Il aperçoit dans la coulisse Euripide costumé en Persée.) Hein ! dieux ! Zeus sauveur ! il y a espoir.}}{}
\ligneEst{}{\francais{L'homme ne semble pas vouloir m'abandonner. Mais il m'a}}{1010}
\ligneEst{}{\francais{donné à entendre, en apparaissant en Persée,}}{}
\ligneEst{}{\francais{qu'il me faut devenir Andromède. Aussi bien,}}{}
\ligneEst{}{\francais{les chaînes, j'en suis muni. Il est donc clair}}{}
\ligneEst{}{\francais{qu'il viendra me délivrer. Autrement il n'aurait pas dirigé son vol de ce côté.}}{}
\ligneEst{\textbf{Mnésiloque}}{\francais{Vierges, vierges amies, puissé-je}}{1015}
\ligneEst{}{\francais{partir d'ici et}}{}
\ligneEst{}{\francais{tromper le Scythe !}}{}
\ligneEst{}{\francais{M'entends-tu, ô toi qui en des antres répètes mes chants ?}}{}
}
\vspace{2pt}
\noindent\rule{16.5cm}{0.6pt}\\[2pt]
\noindent\textbf{\small Commentaire philologique, métrique et scénique (vv. 1001–1043 --- partie 4/8)}\\[2pt]
\begin{adjustbox}{max totalheight=15.2cm, max width=16.5cm, keepaspectratio}%
\parbox[t]{16.5cm}{
\setlength{\columnsep}{9pt}%
\begin{multicols}{2}
\fontsize{7.6pt}{8.8pt}\selectfont
\setlength{\parindent}{0pt}%
\setlength{\parskip}{1.2pt}%
\noindent Je ne sais ce que vaut cette émendation ; je laisse le soin d'en juger à des métriciens compétents.\\ Il est en tout cas bien évident qu'Andromède-Mnésiloque, se lamentant de son sort présent et regrettant les douces années de son enfance, ne fait aucunement allusion à une quelconque \italique{urne de vote} (κημός ?). Les chœurs de danse et les jeux enfantins, voilà qui suffit à sa mélancolie. Le marigot électoral ? alors que la jeune fille ne pouvait même pas y participer activement ? Mais où donc l'ensemble des commentateurs -- à commencer par le scholiaste -- est-il allé cherché cette abracadabrantesque explication ?
\noindent Qu'Aristophane ait pu utiliser le mot κημός, nom d'une sorte d'entonnoir posé sur l'urne pour en celer le vote, comme synecdoque pour l'\italique{urne} elle-même, cela est vrai mais avec une réserve de taille ayant échappé aux commentateurs : c'est pour permettre un jeu de mots, comme nous l'allons voir. (Ce qui explique qu'il soit le seul auteur recourant à ce vocable avec le sens d'\italique{urne}.)\\
\noindent Cette inadmissible traduction de κημός par «~urne~» [pour le vote] nous oblige à rebondir sur un passage des \italique{Guêpes}.
\noindent Il s'agit du début des \italique{Guêpes} avec la traditionnelle scène des deux esclaves qui présentent la situation : en l'occurrence, le mal dont souffre leur maître Philocléon (et une bonne partie des Athéniens), pris qu'il est de \italique{philéliastite}. Et lorsqu'il voit une dédicace au beau Dêmos, le fils de Pyrilampe, notre malade ne peut s'empêcher d'écrire en marge (παρέγραψε) : κημὸς καλός. Vers 97-99 :
\citComm{καὶ νὴ Δί᾽ ἢν ἴδῃ γέ που γεγραμμένον\\
υἱὸν Πυριλάμπους ἐν θύρᾳ Δῆμον καλόν,\\
ἰὼν παρέγραψε πλησίον `κημὸς καλός.'}
\noindent Qui était ce Dêmos Πυριλάμπους ? On possède quelques autres témoignages à son sujet : Calliclès, l'horrible démagogue du \italique{Gorgias} (481d), en avait fait son éromène.
\noindent Platon souligne malicieusement que, à l'instar de l'autre dêmos que lorgne le démagogue, Calliclès finissait toujours par céder devant son cadet aimé (αἰσθάνομαι οὖν σου ἑκάστοτε, καίπερ ὄντος δεινοῦ, ὅτι ἂν φῇ σου τὰ παιδικὰ καὶ ὅπως ἂν φῇ ἔχειν, οὐ δυναμένου ἀντιλέγειν).
\noindent Eupolis (fr. 227(213) Kassel-Austin) nous le décrit comme \italique{bouché} ; littéralement : il est plein de cérumen (κυψέλη) :
\citComm{καὶ τῷ Πυριλάμπους ἆρα Δήμῳ κυψέλη ἔνεστιν; }
\end{multicols}
}%
\end{adjustbox}%
\vfill
\end{adjustwidth}
\newpage

% --- FOLIO 25 (f. 174v) : SPREAD 2.3 (Texte grec établi + Traduction française + Commentaire 5/8 & 6/8) ---
\clearpage\ifodd\value{page}\hbox{}\newpage\fi
\enTeteG{ΘΕΣΜΟΦΟΡΙΑΖΟΥΣΑΙ --- Texte grec établi (f. 174v, vv. 1001–1043 [3/4])}
\enTeteD{LES THESMOPHORIES --- Traduction française (f. 174v, vv. 1001–1043 [3/4])}
\begin{adjustwidth}{-2cm}{-2cm}
\noindent\textbf{\footnotesize Texte grec établi --- Folio 25 (\texttt{f. 174v}, vv. 1001–1043 [3/4]) \hfill ΕΠΕΙΣΟΔΙΟΝ Ε' (Ἀνδρομέδα)}\\[-2pt]
\noindent\rule{16.5cm}{0.5pt}\\[2pt]
{\fontsize{8.8pt}{10.2pt}\selectfont
\ligneEst{{\tiny\textit{telesill.}}}{κατάνευσον, ἔασον ὡς}{1020}
\ligneEst{{\tiny\textit{$+$ ithyph.}}}{τὴν γυναῖκά μ᾽ ἐλθεῖν.}{}
\ligneEst{{\tiny\textit{dim. iamb.}}}{ἄνοικτος ὅς μ᾽ ἔδησε τὸν}{}
\ligneEst{{\tiny\textit{$+$ lecyth.}}}{πολυπονώτατον βροτῶν·}{}
\ligneEst{{\tiny\textit{dim. iamb.}}}{μόλις δὲ γραῖαν ἀποφυγὼν}{}
\ligneEst{{\tiny\textit{dim. iamb.}}}{σαπρὰν ἀπωλόμην ὅμως.}{1025}
\ligneEst{{\tiny\textit{do. $+$ ia. $+$ do.}}}{ὅδε γὰρ ὁ Σκύθης φύλαξ πάλαι ἐφέστηκ᾽· ὀλοὸν}{}
\ligneEst{{\tiny\textit{$2$ cr. $+$ ba.}}}{ἄφιλον ἐκρέμασε κόραξι δεῖπνον.}{}
\ligneEst{{\tiny\textit{do. $+$ ia. $+$ do.}}}{ὁρᾷς; οὐ χοροῖσιν οὐδ᾽ ὑφ᾽ ἡλίκων νεανίδων}{1030}
\ligneEst{{\tiny\textit{$2$ cr. $+$ ba.}}}{ψῆφον ἕστηκ᾽ ἔχουσα κημόν,}{}
\ligneEst{{\tiny\textit{$+$ trim. iamb.}}}{ἀλλ᾽ ἐν πυκνοῖς δεσμοῖσιν ἐμπεπλεγμένη}{}
}
\vspace{2pt}
\noindent\rule{16.5cm}{0.6pt}\\[2pt]
\noindent\textbf{\small Commentaire philologique, métrique et scénique (vv. 1001–1043 --- partie 5/8)}\\[2pt]
\begin{adjustbox}{max totalheight=15.2cm, max width=16.5cm, keepaspectratio}%
\parbox[t]{16.5cm}{
\setlength{\columnsep}{9pt}%
\begin{multicols}{2}
\fontsize{7.6pt}{8.8pt}\selectfont
\setlength{\parindent}{0pt}%
\setlength{\parskip}{1.2pt}%
\noindent Enfin, Athénée (IX, 397c-d), après avoir décrit le paon comme un animal bruyant qui réveille les dormeurs (τοιόνδε ταὧν, ὃς τοὺς εὕδοντας ἐγείρει) nous apprend que le fils de Pyrilampe en faisait élevage, dont il était d'ailleurs très fier...\\ Que peut-on en conclure ? assurément, que le beau Dêmos était un invétéré babillard, obtus à tout conseil, qui en remontrait à Calliclès au rebours d'une relation maître-élève classique et qui, pas plus que les paons qu'il élevait, ne s'arrêtait jamais de parler.
\noindent Comment lire maintenant le graffiti :
\citComm{κημὸς καλός Δῆμον καλόν}
\noindent (noter l'accusatif de Δῆμον καλόν) ?
\noindent Tout simplement en conformité avec le sens premier de κημός, savoir : une \italique{muselière}.
\citComm{Une bonne muselière [fera taire] le beau Dêmos.}
\noindent Que la muselière (κημός), ait pu désigner aussi une \italique{urne de vote}, voilà qui devait laisser Aristophane bien rêveur !..
\noindent La pique contre la démocratie est évidente :
\citComm{Une bonne muselière [musellera] le bon peuple [pour cinq ans, pourrait-on ajouter à propos de la nôtre, encore plus décadente que celle d'Athènes].}
\noindent Mais Philocléon dans sa naïveté -- qui serait touchante si elle n'était pas criminelle -- ne voit tout ce qu'il y a d'avilissant dans cette mascarade démocratique ; il est satisfait que le bon peuple puisse urner à son aise :
\noindent κημὸς καλός δῆμον καλόν.\\ Autre exemple, un peu plus loin dans la même pièce : Philocléon préside au procès du général Lachès\marqueur{3}, animalisé pour l'occasion en chien \italique{Chapardeur} (Λάβης). L'accusateur public n'est autre que le chien de Dêmos, Cléon. Bdélycléon fait entrer le prévenu et son accusateur ; c'est-à-dire les deux chiens, la canaille et le roquet. Réaction de Philocléon devant ces deux patibulaires (\italique{Guêpes}, vers 852-853):
\citComm{Φιλοκλέων -- ἐς κόρακας, ὡς ἄχθομαι\\
ὁτιὴ 'πελαθόμην τοὺς καδίσκους ἐκφέρειν.\\
Aux corbeaux ! Quel ennui ! j'ai oublié d'apporter les καδίσκοι.}
\noindent On traduit καδίσκοι par \italique{urnes} ; et il n'est pas douteux que dans la parodie de procès, les καδίσκοι soient mis en place d'\italique{urnes}... mais on voit que Philocléon frémit tout simplement devant les deux molosses qu'on lui présente et regrette l'absence de \italique{muselières}, c'est sans doute aussi le sens propre de καδίσκος.\\
\end{multicols}
}%
\end{adjustbox}%
\vfill
\end{adjustwidth}
\newpage
% --- PAGE DE DROITE : Traduction française + Commentaire (6/8) ---
\begin{adjustwidth}{-2cm}{-2cm}
\noindent\textbf{\footnotesize Traduction française --- Folio 25 (\texttt{f. 174v}, vv. 1001–1043 [3/4]) \hfill ΕΠΕΙΣΟΔΙΟΝ Ε' (Ἀνδρομέδα)}\\[-2pt]
\noindent\rule{16.5cm}{0.5pt}\\[2pt]
{\fontsize{8.8pt}{10.2pt}\selectfont
\ligneEst{}{\francais{acquiesce, souffre que}}{1020}
\ligneEst{}{\francais{j'aille retrouver ma femme.}}{}
\ligneEst{}{\francais{Sans merci est celui qui m'enchaîna,}}{}
\ligneEst{}{\francais{des mortels le plus éprouvé.}}{}
\ligneEst{}{\francais{Échappé à grand'peine à une vieille}}{}
\ligneEst{}{\francais{guenipe, je ne laissai pas d'être perdu.}}{1025}
\ligneEst{}{\francais{Car depuis un temps ce Scythe se tient en faction : il m'a suspendu, brisé, sans amis,}}{}
\ligneEst{}{\francais{en pâture aux corbeaux. Vois-tu, ce n'est pas}}{}
\ligneEst{}{\francais{pour des chœurs de danse}}{1030}
\ligneEst{}{\francais{ni parmi des jeunes filles de mon âge}}{}
\ligneEst{}{\francais{que je suis ici debout avec cornet et jeton.}}{}
}
\vspace{2pt}
\noindent\rule{16.5cm}{0.6pt}\\[2pt]
\noindent\textbf{\small Commentaire philologique, métrique et scénique (vv. 1001–1043 --- partie 6/8)}\\[2pt]
\begin{adjustbox}{max totalheight=15.2cm, max width=16.5cm, keepaspectratio}%
\parbox[t]{16.5cm}{
\setlength{\columnsep}{9pt}%
\begin{multicols}{2}
\fontsize{7.6pt}{8.8pt}\selectfont
\setlength{\parindent}{0pt}%
\setlength{\parskip}{1.2pt}%
\noindent\textbf{[Notes ad 1031]} \textbf{(1)} Cf. Pollux (X, 150) sur les ustensiles de jeu :
\citComm{Κυβευτοῦ σκεύη · ἄβαξ ἀβάκιον, κόσκινον, κύβοι, διάσειστοι κύβοι, ἀστράγαλοι, φιμοί, κημοί, τηλία, κήθια κηθίδια, ψῆφοι, πεττοί.}
\noindent φιμός et κημός sont littéralement des \italique{muselières}, avec le sens spécialisé de \italique{cornets} (à dés ou astragales) et sont synonymes de κήθιον. On pourra par exemple se reporter à l'article de Michel Casevitz, «~Sur les termes de jeu : compléments d'histoire des mots~», \italique{Kentron} 36 (2021), p. 209-212. Pour des ouvrages beaucoup plus conséquents, le lecteur pourra lire \italique{Les jeux des Anciens}\expo{2} (1873) de Louis Becq de Fouquières, auteur à la vaste érudition et au style des plus agréables comme le \nieme{XIX} siècle en a été prodigue ; l'édition par Jean Taillardat déjà mentionnée du Περὶ βλασφημιῶν, \italique{Des termes injurieux} et du Περὶ παιδιῶν, \italique{Des jeux grecs}, de Suétone ; la somme de référence qu'est l'article de Hans Lamer «~\latin{Lusoria tabula}~» de la \italique{Paulys Realencyclopädie der classischen Altertumswissenschaft}, \allemand{Band} XIII 2 (1927), col. 1900-2029 ; plus récemment, les \latin{excerpta de ludis} chez Pollux ont été publiés par Salvatore Costanza : \italique{Giulio polluce, onomasticon: excerpta de ludis} accompagnés d'une traduction italienne. Outre l'archéologie, nos témoins principaux quant aux jeux grecs sont Pollux, Hésychius et Eustathe (qui s'appuie sur les traités perdus de Suétone) -- non exempts de contradictions entre eux d'ailleurs.
\noindent \textbf{(2)} Van Leeuwen a recensé d'autres erreurs similaires où l'ordre des mots est inversé dans le \italique{Ravennas} (la comparaison avec les neuf pièces du \italique{Ravennas} qui nous sont connues par d'autres témoins est préférable : la correction est davantage sûre). Cf. l'édition en fac-similé qu'il fit de notre manuscrit : \italique{Comoediae Undecim cum Scholiis : Codex Ravennas 137, 4, A} [$=$ \italique{Rav.} 429 nouvelle cote], \italique{phototypice editus} in \italique{Codices Graeci et Latini photographice depicti}, t. IX (1904), pp. XVIII-XIX.
\noindent \textbf{(3)} La scène suivante est une amusante parodie du procès intenté, ou censé intenté, au stratège athénien Lachès, qui était en effet du dême d'Æxone. Le chien du peuple, c'est Cléon, du dême de Cydathène. Lachès avait commandé la flotte envoyée en Sicile en 427. Rappelé au commencement de 425, il fut accusé d'avoir malversé. Le procès (s'il eut lieu) fut présidé suivant l'usage par un thesmothète (v. 935). Cléon, comme synégore, soutint l'accusation. Lachès fut acquitté, car il remplit encore diverses fonctions publiques. Il est intéressant de constater que si le poète n'épargne pas Cléon, il n'est guère plus indulgent pour Lachès, encore que celui-ci appartînt au parti aristocratique. (Note de Willems.)
\noindent\textbf{1033} \textbf{κήτει βορὰ Γλαυκέτῃ πρόκειμαι.} On sait par le scholiaste des \italique{Oiseaux} (vers 348) que la pièce originale d'Euripide comportait à cet endroit :
\citComm{ἐκθεῖναι κήτεϊ φορβάν,}
\end{multicols}
}%
\end{adjustbox}%
\vfill
\end{adjustwidth}
\newpage

% --- FOLIO 25 (f. 174v) : SPREAD 2.4 (Texte grec établi + Traduction française + Commentaire 7/8 & 8/8) ---
\clearpage\ifodd\value{page}\hbox{}\newpage\fi
\enTeteG{ΘΕΣΜΟΦΟΡΙΑΖΟΥΣΑΙ --- Texte grec établi (f. 174v, vv. 1001–1043 [4/4])}
\enTeteD{LES THESMOPHORIES --- Traduction française (f. 174v, vv. 1001–1043 [4/4])}
\begin{adjustwidth}{-2cm}{-2cm}
\noindent\textbf{\footnotesize Texte grec établi --- Folio 25 (\texttt{f. 174v}, vv. 1001–1043 [4/4]) \hfill ΕΠΕΙΣΟΔΙΟΝ Ε' (Ἀνδρομέδα)}\\[-2pt]
\noindent\rule{16.5cm}{0.5pt}\\[2pt]
{\fontsize{8.8pt}{10.2pt}\selectfont
\ligneEst{{\tiny\textit{$3$ ia. sync\cat}}}{κήτει βορὰ Γλαυκέτῃ πρόκειμαι,}{}
\ligneEst{{\tiny\textit{dim. ia\cat}}}{γαμηλίῳ μὲν οὐ ξὺν}{}
\ligneEst{{\tiny\textit{$+$ 2 ia\cat}}}{παιῶνι δεσμίῳ δὲ}{1035}
\ligneEst{{\tiny\textit{dim. iamb.}}}{γοᾶσθέ μ᾽ ὦ γυναῖκες ὡς}{}
\ligneEst{{\tiny\textit{$2$ troch.}}}{μέλεα μὲν πέπονθα μέλεος,}{}
\ligneEst{{\tiny\textit{lecyth.}}}{ὢ τάλας ἐγὼ τάλας,}{}
\ligneEst{{\tiny\textit{dochm. $+$ ...}}}{ἀπὸ δὲ συγγόνων ἀλλ᾽ ἄνομα πάθεα,}{}
\ligneEst{}{†φῶτά λιτομένα†,}{1040}
\ligneEst{{\tiny\textit{do. $+$ cr. $+$ ba.}}}{πολυδάκρυτον Ἀίδα γόον ψεφαυγοῦς.}{}
\ligneEst{{\tiny\textit{extr. metr.}}}{αἰαῖ αἰαῖ ἒ ἒ,}{}
\ligneEst{{\tiny\textit{dim .troch.}}}{ὃς ἔμ᾽ ἀπεξύρησε πρῶτον,}{1043}
}
\vspace{2pt}
\noindent\rule{16.5cm}{0.6pt}\\[2pt]
\noindent\textbf{\small Commentaire philologique, métrique et scénique (vv. 1001–1043 --- partie 7/8)}\\[2pt]
\begin{adjustbox}{max totalheight=15.2cm, max width=16.5cm, keepaspectratio}%
\parbox[t]{16.5cm}{
\setlength{\columnsep}{9pt}%
\begin{multicols}{2}
\fontsize{7.6pt}{8.8pt}\selectfont
\setlength{\parindent}{0pt}%
\setlength{\parskip}{1.2pt}%
\noindent Γλαυκέτης, l'Athénien glouton (τένθης) est cité au vers 1008 de la \italique{Paix}.
\noindent Trygée y fait une prière pour que la paix revienne à Athènes ; la ville est sens dessus dessous du fait de la guerre : les embargos et autres razzias ont littéralement renversé le marché : au lieu de toutes les bonnes denrées, l'Athénien n'y rencontre que des rappels de ses privations.
\noindent Autre inversion : si la paix revenait, ce serait alors aux citoyens de bousculer les trois parasites du marché aux anguilles. Aristophane joue sur une image qu'il a développée ailleurs\marqueur{1} : comme le chasseur d'anguille trouble l'eau pour capturer sa proie, le politicien véreux trouble les affaires publiques pour mieux s'enrichir. (Remarquer l'utilisation du verbe τυρβάζεσθαι, qu'Aristophane emploie aussi au vers 257 des \italique{Guêpes} : τὸν πηλὸν ὥσπερ ἀτταγᾶς τυρβάσεις βαδίζων.)  On peut mettre notre main au feu que le groupe des trois affreux étaient de francs coquins qui se gobergeaient sur le dos d'Athènes pour mener une vie douce, bonne et grasse\marqueur{2}.
\noindent Γλαυκέτης est mis par paronomase ; \italique{Nomen, omen} : lorsqu'on s'appelle Γλαυκέτης, on est prédestiné à devenir ce poisson vorace γλαῦκος-κῆτος\marqueur{3} !
\noindent (L'image de la femme livrée en pâture (βορά) à la poissonnaille est récurrente chez les Comiques, \italique{confer} Athénée (VII, 315b) citant Platon le Comique ($=$ fr. 57(56) Kassel-Austin)) :
\citComm{σὲ γάρ, γραῦ, συγκατῴκισεν σαπρὰν\\
ὀρφῷσι σελαχίοις τε καὶ φάγροις βοράν}
\noindent (βορά désigne ici la menue fretaille qui sert de pâture aux poissons carnivores comme l'orphe (ὄρφος / ὀρφώς), qu'Athénée dans le même chapitre rapproche du \italique{glauque}, γλαῦκος.)\\
\noindent\textbf{[Notes ad 1033]} \textbf{(1)} Cavaliers, vers 864-867. Voir aussi le vers 361. Érasme a relevé l'image dans ses \italique{Adages} : n° 2576 (\latin{Anguillas captare}).
\noindent \textbf{(2)} Cf. pour Morykhos les \italique{Acharniens}, v. 887 ainsi que les \italique{Guêpes}, v. 506 \& 1142. On notera au passage qu'un Téléas fils de Télénikos apparaît à la même époque dans les \italique{Inscriptiones Graecae} (I\expo{3}, 307 ; 308 ; 331 ; 332 ; 370) à la secrétairie (ἐγραμμάτευε) du Trésor (\textit{h}οι ταμίαι το͂ν \textit{h}ιερο͂ν χρεμάτον τε͂ς Ἀθεναίας). Voir ce qu'en disent les scholies aux vers 167-8 des \italique{Oiseaux}.
\noindent \textbf{(3)} Autre jeu du même acabit sur le nom de Glaucétès : (Platon le Comique, fr. 114(106) K.-A.) :
\citComm{ὦ θεῖε Μόρυχε· πῶς γὰρ οὐ δαίμων ἔφυς;\\
καὶ Γλαυκέτης ἡ ψῆττα, καὶ Λεωγόρας,\\
οἵ ζῆτε τερπνῶς οὐδὲν ἐνθυμούμενοι.}
\end{multicols}
}%
\end{adjustbox}%
\vfill
\end{adjustwidth}
\newpage
% --- PAGE DE DROITE : Traduction française + Commentaire (8/8) ---
\begin{adjustwidth}{-2cm}{-2cm}
\noindent\textbf{\footnotesize Traduction française --- Folio 25 (\texttt{f. 174v}, vv. 1001–1043 [4/4]) \hfill ΕΠΕΙΣΟΔΙΟΝ Ε' (Ἀνδρομέδα)}\\[-2pt]
\noindent\rule{16.5cm}{0.5pt}\\[2pt]
{\fontsize{8.8pt}{10.2pt}\selectfont
\ligneEst{}{\francais{Mais enchaînée en d'étroits liens,}}{}
\ligneEst{}{\francais{je me vois exposée à la voracité du monstre Glaucétès.}}{}
\ligneEst{}{\francais{Lamentez sur moi un péan, non d'hymen, mais de chaînes,}}{1035}
\ligneEst{}{\francais{ô femmes, car,}}{}
\ligneEst{}{\francais{pitoyable, j'ai pitoyablement souffert.}}{}
\ligneEst{}{\francais{Malheureux que je suis, malheureux !}}{}
\ligneEst{}{\francais{Et entre d'autres tortures impies infligées par mes proches,}}{}
\ligneEst{}{\francais{implorant l'homme}}{1040}
\ligneEst{}{\francais{d'une lamentation pleine de larmes de l'obscur Hadès.}}{}
\ligneEst{}{\francais{ah ah! ah ah!}}{}
\ligneEst{}{\francais{l'homme qui me rasa d'abord,}}{1043}
}
\vspace{2pt}
\noindent\rule{16.5cm}{0.6pt}\\[2pt]
\noindent\textbf{\small Commentaire philologique, métrique et scénique (vv. 1001–1043 --- partie 8/8)}\\[2pt]
\begin{adjustbox}{max totalheight=15.2cm, max width=16.5cm, keepaspectratio}%
\parbox[t]{16.5cm}{
\setlength{\columnsep}{9pt}%
\begin{multicols}{2}
\fontsize{7.6pt}{8.8pt}\selectfont
\setlength{\parindent}{0pt}%
\setlength{\parskip}{1.2pt}%
\noindent\textbf{1041} \textbf{Ἀίδα,} Aristophane use du même génitif dorique au vers 1333 des \italique{Grenouilles} (parodie d'Euripide).
\noindent\textbf{1041} \textbf{πολυδάκρυτον Ἀίδα γόον φεύγουσαν.} Le manuscrit porte l'incompréhensible †φεύγουσαν†. Y suppléer par χέουσα(ν), à l'instar de certains éditeurs, c'est faire fi de sa leçon. Je dirais même plus : χέουσα est le mot banal qu'on attendait, il paraît alors très difficile qu'un scribe s'y soit trompé. (σφριγῶσα, \italique{gonflée}, aurait au moins le mérite de se justifier paléographiquement mais le terme paraît impropre en grec à l'image d'yeux gonflés de larmes.) Une autre proposition, `φλέγουσα(ν)', a aussi emporté la faveur de nombreuses éditions. Elle est reprise par Jean Taillardat à propos de l'image de la Muse-flamme (§736 n.2)... mais le grand Taillardat me semble s'abuser : si un poème, un chant, un \italique{thrêne} (qui est avant tout sonore ; la flamme étant celle de la voix selon moi) peuvent être enflammés, cela ne prouve pas pour autant que des pleurs peuvent l'être.\\ Le scholiaste nous fournit une information qui pourrait peut-être bien débrouiller l'écheveau: une glose, σκοτεινήν, \italique{obscur}, se trouve au vers 1054 pour qualifier soi-disant αἰόλαν, \italique{rapide}, \italique{mobile}, \italique{changeant}.
\noindent Sauf qu'αἰόλος n'a jamais eu le sens de σκοτεινός et que Willems a clairement montré qu'il fallait prendre αἰόλος en son sens courant\marqueur{1}.
\noindent Suidas (αι 245) -- ou l'un des lexicographes de la \italique{Souda}, si l'on préfère -- n'a pas manqué néanmoins de relever la glose à l'entrée Αἰόλαν. Mais Suidas est coutumier de ce genre de fausses attributions (voir la note au vers 735).
\noindent Or en cherchant des mots proches de †φεύγουσαν†, je suis retombé sur ce cher Hésychius (ψ 135) :
\citComm{ψεφαυγοῦς· σκοτεινῆς,}
\noindent Génitif dorique (comme Ἀίδα) d'un certain *ψεφαυγής, \italique{au sombre regard}.
\noindent Euripide a une prédilection pour forger des composés en -αυγής. (\italique{Confer} le Chantraine, \italique{sub voce} αὐγή.)  Si l'on fait l'hypothèse de la chute de `ψε', le φαυγους restant aurait bien pu être complété d'un `αν' en regard de l'autre participe †λιτομέναν† juste au-dessus, et donner ainsi notre leçon †φεύγουσαν†.
\noindent\textbf{[Notes ad 1041]} \textbf{(1)} Note n° 27, reproduite ci-après.
\noindent\textbf{1043-51} \textbf{Papyrus d'Oxyrhynque n° 4935.} Un maigre fragment (n° 1) du Papyrus d'Oxyrhynque 4935\marqueur{1} (\nieme{II} siècle ap. J.-C.), reproduit en annexe, nous apporte une leçon nouvelle au vers 1047 : le `μοι' du \italique{Ravennas} n'y figure pas. De plus, il confirme -- comme l'avait déjà fait Willems en avance sur son temps -- le `βάρβαρον' du vers 1051 qui laissait de nombreux éditeurs sceptiques. (Voir sa \nieme{26} note, très instructive sur ce passage et le rôle de l'étoile Sirius.)
\noindent\textbf{[Notes ad 1043-51]} \textbf{(1)} \italique{The Oxyrhynchus Papyri}, LXXIII (2009), ed. A. Benaissa.
\end{multicols}
}%
\end{adjustbox}%
\vfill
\end{adjustwidth}
\newpage
\clearpage\ifodd\value{page}\hbox{}\newpage\fi
\enTeteG{ΕΠΕΙΣΟΔΙΟΝ Ε' (Ἠχώ) --- Codex Ravennas 429, f. 175r (Classense 177r) --- vv. 1044–1081}
\begin{adjustwidth}{-1.5cm}{-3.0cm}
\begin{adjustbox}{max totalheight=21.0cm, max width=16.5cm, keepaspectratio}
\begin{minipage}{16.5cm}
{\fontsize{7.2pt}{8.5pt}\selectfont
\noindent\textbf{Folio 26 --- Codex Ravennas 429, \texttt{f. 175r} (\texttt{Classense 177r}) : ΕΠΕΙΣΟΔΙΟΝ Ε' (Ἠχώ), vv. 1044–1081}\hfill \textit{Témoins :} $\text{R}^1, \Rpc, \Rdeux$, \PsiQuatre\par\vspace{2pt}
\hrule height 0.4pt\vspace{2.5pt}
}
{\fontsize{8.1pt}{9.5pt}\selectfont
\ligneRav{1044}{μ\auDessus{ν}{Η}}{\hspace{1.1em}ὃς ἐμὲ κροκόεντ᾽ ἐνέδυσεν·}{\textbf{1044} \lecon{κροκόεντ᾽} R : κροκόεν τόδ' Bergk \espApp ἐνέδυσεν R : ἀμφέδυσεν Hermann ap Σ}
\ligneRav{1045}{}{\hspace{1.1em}ἐπὶ δὲ τοῖσδε τόδ᾽ ἀνέπεμψεν}{}
\ligneRav{}{}{\hspace{1.1em}ἱερόν, ἔνθα γυναῖκες·}{}
\ligneRav{1047}{}{\hspace{1.1em}ἰώ μοι μοίρας}{\textbf{1047} [ι]\douteux{ω}μοι\douteux{ρ}[ας \PsiQuatre : ἰώ μοι μοίρας R : ἰώ μοίρας Hermann \espApp δαίμων R : δαῖμον Burges \espApp ἀνέτικτε R : \lecon{ἄτεγκτε} Ellebodius : ἅν ἔτικτε Casaubon}
\ligneRav{}{}{\hspace{1.1em}ἄτεγκτε δαίμων·}{}
\ligneRav{}{}{\hspace{1.1em}ὦ κατάρατος ἐγώ·}{}
\ligneRav{}{}{\hspace{1.1em}τίς ἐμὸν οὐκ ἐπόψεται}{}
\ligneRav{}{}{\hspace{1.1em}πάθος ἀμέγαρτον ἐπὶ κακῶν παρουσίᾳ;}{}
\ligneRav{1050}{}{\hspace{1.1em}εἴθε με πυρφόρος αἰθέρος ἀστὴρ—}{}
\ligneRav{1051}{}{\hspace{1.1em}τὸν \difference{βάρβαρον} ἐξολέσειεν·}{\textbf{1051} \lecon{βάρβαρον} \PsiQuatre{}, R, Willems}
\ligneRav{1052}{}{\hspace{1.1em}οὐ γὰρ ἔτ᾽ ἀθανάταν φλόγα \difference{λεύσσειν}}{\textbf{1052} λεύσειν R : \lecon{λεύσσειν} Biset}
\ligneRav{}{}{\hspace{1.1em}ἐστὶν ἐμοὶ φίλον, ὡς ἐκρεμάσθην,}{}
\ligneRav{1054}{}{\hspace{1.1em}λαιμότμητ᾽ ἄχη δία μένων αἰόλαν}{\textbf{1054} δαιμόνων R, Σ : δαιμονῶν Brunck : δαιμόνι᾽ amicus Fritzschii : \lecon{διὰ μένων} Willems}
\ligneRav{1055}{}{\hspace{1.1em}νέκυσιν ἐπὶ πορείαν·}{}
\ligneRav{1056}{ευρ\auDessus{ι}{π} ηχω}{χαῖρ᾽ ὦ φίλη παῖ· τὸν δὲ πατέρα Κηφέα}{\textbf{1056} ευρ\auDessus{ι}{π} ηχω in mg. \Rpc, \Mon}
\ligneRav{}{}{ὅς σ᾽ ἐξέθηκεν ἀπολέσειαν οἱ θεοί·}{}
\ligneRav{1058}{μ\auDessus{ν}{Η}}{σὺ δ᾽ εἶ τίς ἥτις τοὐμὸν \difference{ᾤκτιρας} \difference{πάθος};}{\textbf{1058} ὤκτειρας R : \lecon{ᾤκτιρας} van Herwerden \espApp πάθος R : \lecon{πάθος;} Biset}
\ligneRav{}{\auDessus{εὐ}{ρ}}{Ἠχὼ, λόγων ἀντῳδὸς ἐπικοκκάστρια,}{}
\ligneRav{1060}{}{ἥπερ πέρυσιν ἐν τῷδε \difference{ταὐτῷ} χωρίῳ}{\textbf{1060} ταυτῷ R : \lecon{ταὐτῷ} Dindorf (1830)}
\ligneRav{1061}{}{Εὐριπίδῃ \difference{καὐτὴ} ξυνηγωνιζόμην·}{\textbf{1061} καὶ αὐτὴ R : \lecon{καὐτὴ} Biset}
\ligneRav{1062}{}{ἀλλ᾽ ὦ τέκνον σὲ μὲν \difference{τὸ σαυτῆς} χρὴ ποιεῖν,}{\textbf{1062} τοσαύτης R : \lecon{τὸ σαυτῆς} Dindorf (1839)}
\ligneRav{1063}{—}{\difference{κλαίειν} ἐλεινῶς. : σὲ δ᾽ ἐπικλαίειν ὕστερον·}{\textbf{1063} \lecon{κλαίειν} et \lecon{ἐπικλαίειν} R : κλάειν et ἐπικλάειν Brunck}
\ligneRav{}{\auDessus{εὐ}{ρ}}{ἐμοὶ μελήσει ταῦτά γ᾽. ἀλλ᾽ ἄρχου λόγων·}{}
\ligneRav{1065}{ὡς ἀνδρομ.}{\hspace{1.1em}ὦ νὺξ ἱερὰ}{\textbf{1065} ὁ μνησίλοχος ὡς ἀνδρομέδα in mg. \Rpc, \Mon}
\ligneRav{}{}{\hspace{1.1em}ὡς μακρὸν ἵππευμα διώκεις}{}
\ligneRav{}{}{\hspace{1.1em}ἀστεροειδέα νῶτα διφρεύουσ᾽}{}
\ligneRav{}{}{\hspace{1.1em}αἰθέρος ἱερᾶς}{}
\ligneRav{1069}{—}{\hspace{1.1em}τοῦ σεμνοτάτου δι᾽ Ὀλύμπου; : †δι᾽ Ὀλύμπου;†}{\textbf{1069} \lecon{δι᾽ Ὀλύμπου} R : †δι᾽† ὀλυμπίας Bibi ap. Hesych. (ο 649)}
\ligneRav{1070}{μ\auDessus{ν}{Η}}{\hspace{1.1em}τί ποτ᾽ Ἀνδρομέδα \difference{περίαλλα} κακῶν}{\textbf{1070} περὶ ἄλλα R : \lecon{περίαλλα} Ellebodius}
\ligneRav{1071}{—}{\hspace{1.1em}μέρος ἐξέλαχον; : \difference{μόρος ἔξελ᾽ ἄχος}·}{\textbf{1071} μέρος ἐξέλαχον; R : \lecon{μόρος ἔξελ᾽ ἄχος.} Bibi}
\ligneRav{1072}{μ\auDessus{ν}{Η}}{\hspace{1.1em}θανάτου τλήμων : †θανάτου τλήμων†·}{\textbf{1072} θανάτου τλήμων R : θαλάττης τέμνον Bibi}
\ligneRav{1073}{μ\auDessus{ν}{Η}}{\hspace{1.1em}ἀπολεῖς μ᾽ ὦ γραῦ \difference{στωμυλλομένη}·}{\textbf{1073} στωμυλλωμένη R : \lecon{στωμυλλομένη} Grynaeus}
\ligneRav{1074}{\auDessus{εὐ}{ρ}}{\hspace{1.1em}\difference{στωμυλιομενής}·}{\textbf{1074} στωμυλωμένη R : \lecon{στωμυλιομενής} Bibi}
\ligneRav{1075}{μ\auDessus{ν}{Η}}{\hspace{1.1em}νὴ Δί᾽, ὀχληρά γ᾽ εἰσήρρηκας}{}
\ligneRav{}{—}{\hspace{1.1em}λίαν. : †λίαν†·}{}
\ligneRav{}{μ\auDessus{ν}{Η}}{\hspace{1.1em}ὦγάθ᾽ ἔασόν με μονῳδῆσαι,}{}
\ligneRav{1078}{—}{\hspace{1.1em}καὶ χαριεῖ μοι παῦσαι. : ᾽παΰσαι·}{\textbf{1078} Μνη. χαριεῖ μοι παῦσαι --- Εὐρ. παῦσαι R : \lecon{Μνη. χαριεῖ μοι ᾽παΰσαι --- simul Μνη. et Εὐρ. παῦσαι} Bibi \espApp (Εὐρ.) παῦσαι R : \lecon{(Εὐρ.) ᾽παΰσαι} Bibi}
\ligneRav{1079}{μ\auDessus{ν}{Η}}{\hspace{1.1em}βάλλ᾽ ἐς κόρακας. : \difference{βάλλε σύ κόρακας}·}{\textbf{1079} βάλλ᾽ ἐς κόρακας R : \lecon{βάλλε σύ κόρακας} Bibi}
\ligneRav{1080}{μ\auDessus{ν}{Η}}{\hspace{1.1em}†\difference{τί τὸ κακόν}†; : †τί τὸ κακόν;† : ληρεῖς. : †ληρεῖς†·}{\textbf{1080} \lecon{τί τὸ κακόν} bis R : τί κακόν bis Bentley}
\ligneRav{}{μ\auDessus{ν}{Η}}{\hspace{1.1em}οἴμωζ᾽. : †οἴμωζ᾽†·}{}
}
\vspace{3pt}\noindent\rule{4.5cm}{0.4pt}\par\vspace{1.5pt}
{\fontsize{7.0pt}{8.3pt}\selectfont
\noindent\textbf{Loci paralleli (vv. 1044–1081) ---} \textbf{1049} (παρουσίᾳ) Moeris π 14 \espTest \textbf{1059} Suidas η 686 (ἐπικοκκάστρια codd. GM\expo{a.c.} vel ἐπικοκάστρια cod. I vel ἐπικαγχάστρια cod. A vel ἐπικαχάστρια codd. FV\expo{mg.} vel ἐπικοστράτ textu V) \espTest (ἀντῳδὸς ἐπικοκκάστρια) \Eustathe 1761.26 ἐκτίθεται καινοφώνους λέξεις οἷον [ ... ] καὶ τὸ ἐπικοκκάζειν ὕβρει ἤ κατά τινας ἐπιμηκάζειν, ὃθεν παρὰ Ἀριστοφάνει τὸ `ἀοιδὸς ἐπικοκκάστρια' \espTest \textbf{1066-8} (ὡς μακρὸν -- αἰθέρος) Suidas ι 531 $\sim$ Zonaras 1117 \espTest \textbf{1067} (νῶτα) Moeris ν 6 \espTest \textbf{1067-8} Suidas αι 122 $\sim$ Zonaras 67, 78 (sed ἀστρο- pro ἀστερο-)\par
}
\vspace{2.5pt}\noindent\rule{4.5cm}{0.4pt}\par\vspace{1.5pt}
{\fontsize{6.9pt}{8.2pt}\selectfont
\noindent\textbf{Suite (en commentaire) de l'apparat critique \& notes paléographiques (\texttt{f. 175r}) ---} \textbf{Témoin papyrologique :} \textbf{\PsiQuatre{} (\italique{P.Oxy.} LXXIII 4935, fr. 1)} couvre les \textbf{vv. 1043–1051}. --- \textbf{Deux indications scéniques marginales de \Rpc{} (copiées par \Mon{}) :} --- \textbf{Apport du papyrus \PsiQuatre{} (vv. 1047 et 1051) :} Au v. 1047, \PsiQuatre{} lit [ι]ω μοιρ[ας sans la dittographie μοι de R (ἰώ μοι μοίρας) ; au v. 1051, \PsiQuatre{} confirme βάρβαρον de R (défendu par Alphonse Willems et L. Hauseux). --- \textbf{Cinq restitutions et conjectures de \textbf{Bibi} dans le dialogue avec Écho (vv. 1069, 1071, 1072, 1074, 1078, 1079) :} --- \textbf{v. 1044 :} \lecon{κροκόεντ'} R (κροκόεν τόδ' Bergk) \espApp ἐνέδυσεν R : \lecon{ἀμφέδυσεν} Hermann ap. $\Sigma$. --- \textbf{v. 1047 :} [ι]ω μοιρ[ας \PsiQuatre{} : ἰώ μοι μοίρας R : \lecon{ἰὼ μοίρας} Hermann \espApp δαίμων R : \lecon{δαῖμον} Burges \espApp ἀνέτικτε R (médivision et épenthèse de $\nu$) : \lecon{ἄτεγκτε} Ellebodius (ἅν ἔτικτε Casaubon). --- \textbf{v. 1051 :} \lecon{βάρβαρον} \PsiQuatre{}, R, Willems. --- \textbf{v. 1054 :} δαιμόνων R, $\Sigma$ (δαιμονῶν Brunck ; δαιμόνι' amicus Fritzschii) : \lecon{διὰ μένων} Alphonse Willems (adopté par L. Hauseux). --- \textbf{v. 1056 :} \lecon{ευρ\auDessus{ι}{π} ηχω} in mg. $\text{R}^{\text{pc}}, \text{M}$. --- \textbf{v. 1065 :} \lecon{ὁ μνησίλοχος ὡς ἀνδρομέδα} in mg. $\text{R}^{\text{pc}}, \text{M}$. --- \textbf{v. 1069 :} \lecon{δι' Ὀλύμπου} R (†δι'† ὀλυμπίας \lecon{Bibi} ap. Hésychius $\text{ο}\ 649$). --- \textbf{v. 1071 :} μέρος ἐξέλαχον; R : \lecon{μόρος ἔξελ' ἄχος.} \lecon{Bibi}. --- \textbf{v. 1072 :} \lecon{θανάτου τλήμων} R (θαλάττης τέμνον \lecon{Bibi} conj. in app.). --- \textbf{v. 1074 :} στωμυλωμένη R : \lecon{στωμυλιομενής} \lecon{Bibi}. --- \textbf{v. 1078 :} (Εὐρ.) παῦσαι R : \lecon{(Εὐρ.) 'παΰσαι} \lecon{Bibi}. --- \textbf{v. 1079 :} βάλλ' ἐς κόρακας R : \lecon{βάλλε σύ κόρακας} \lecon{Bibi}.\par
}
\end{minipage}
\end{adjustbox}
\end{adjustwidth}
\newpage
\enTeteD{Codex Ravennas 429 (Biblioteca Classense) --- Fac-similé du f. 175r (vv. 1044–1081)}
\begin{adjustwidth}{-1.5cm}{-3.0cm}
\begin{center}
\includegraphics[width=16.5cm,height=20.5cm,keepaspectratio]{/workspaces/Thesmophories/dossier_ravennas_thesmophories/folios_individuels_ff162v-177r/26_f175r_classense177r_vv1044-1081.jpg}\par
\vspace{2pt}
{\scriptsize\textbf{Codex Ravennas 429, f. 175r (Classense 177r)} --- ΕΠΕΙΣΟΔΙΟΝ Ε' (Ἠχώ), \textbf{vv. 1044–1081} ($\text{R}^1, \Rpc, \Rdeux$, \PsiQuatre).}
\end{center}
\end{adjustwidth}
\newpage
\clearpage\ifodd\value{page}\hbox{}\newpage\fi% --- FOLIO 26 (f. 175r) : SPREAD 2.1 (Texte grec établi + Traduction française + Commentaire 1/8 & 2/8) ---
\clearpage\ifodd\value{page}\hbox{}\newpage\fi
\enTeteG{ΘΕΣΜΟΦΟΡΙΑΖΟΥΣΑΙ --- Texte grec établi (f. 175r, vv. 1044–1081 [1/4])}
\enTeteD{LES THESMOPHORIES --- Traduction française (f. 175r, vv. 1044–1081 [1/4])}
\begin{adjustwidth}{-2cm}{-2cm}
\noindent\textbf{\footnotesize Texte grec établi --- Folio 26 (\texttt{f. 175r}, vv. 1044–1081 [1/4]) \hfill ΕΠΕΙΣΟΔΙΟΝ Ε' (Ἠχώ)}\\[-2pt]
\noindent\rule{16.5cm}{0.5pt}\\[2pt]
{\fontsize{8.8pt}{10.2pt}\selectfont
\ligneEst{}{ὃς ἐμὲ κροκόεντ᾽ ἐνέδυσεν·}{1044}
\ligneEst{{\tiny\textit{dim. troch.}}}{ἐπὶ δὲ τοῖσδε τόδ᾽ ἀνέπεμψεν}{1045}
\ligneEst{{\tiny\textit{pherecr.}}}{ἱερόν, ἔνθα γυναῖκες.}{}
\ligneEst{{\tiny\textit{dochm.}}}{ἰώ μοι μοίρας}{}
\ligneEst{{\tiny\textit{ithyph.}}}{ἄτεγκτε δαίμων·}{}
\ligneEst{{\tiny\textit{hem.}}}{ὦ κατάρατος ἐγώ·}{}
\ligneEst{{\tiny\textit{lecyth.}}}{τίς ἐμὸν οὐκ ἐπόψεται}{}
\ligneEst{{\tiny\textit{trim. iamb.}}}{πάθος ἀμέγαρτον ἐπὶ κακῶν παρουσίᾳ;}{}
\ligneEst{{\tiny\textit{$4$ dactyl.}}}{εἴθε με πυρφόρος αἰθέρος ἀστὴρ—}{1050}
\ligneEst{{\tiny\textit{paroem.}}}{τὸν βάρβαρον ἐξολέσειεν.}{}
\ligneEst{{\tiny\textit{$4$ dactyl.}}}{οὐ γὰρ ἔτ᾽ ἀθανάταν φλόγα λεύσσειν}{}
\ligneEst{{\tiny\textit{$+$ $4$ dact.}}}{ἐστὶν ἐμοὶ φίλον, ὡς ἐκρεμάσθην,}{}
}
\vspace{2pt}
\noindent\rule{16.5cm}{0.6pt}\\[2pt]
\noindent\textbf{\small Commentaire philologique, métrique et scénique (vv. 1044–1081 --- partie 1/8)}\\[2pt]
\begin{adjustbox}{max totalheight=15.2cm, max width=16.5cm, keepaspectratio}%
\parbox[t]{16.5cm}{
\setlength{\columnsep}{9pt}%
\begin{multicols}{2}
\fontsize{8.4pt}{9.7pt}\selectfont
\setlength{\parindent}{0pt}%
\setlength{\parskip}{1.2pt}%
\noindent\textbf{1054-5} \textbf{λαιμότμητ᾽ ἄχη δαιμόνων αἰόλαν / νέκυσιν ἐπὶ πορείαν.} Willems a consacré une note à ce passage\marqueur{1}. Si les vues qu'il y développe sont ingénieuses, elle n'emporte pas pour autant ma conviction.
\noindent Willems propose de construire ἄχη δία μένων là où le \italique{Ravennas} porte ἄχη †δαιμόνων†. À défaut d'autres propositions plus convaincantes, c'est néanmoins la correction que je retiendrai.
\noindent\textbf{[Notes ad 1054-5]} \textbf{(1)} \nieme{27} note aux \italique{Thesmophories} d'Alphonse Willems, reproduite en annexe.
\noindent\textbf{1059} \textbf{ἀντῳδὸς ἐπικοκκάστρια.} Ceci rappelle la jolie épigramme d'Archias (\italique{Anthol. Palat.}, VII, 191) sur une pie (κίσσα\marqueur{1})
\citComm{qui souvent, pareille à l'Écho, faisait entendre, d'un bec prompt à répliquer, une équisonnance bruyante et moqueuse.\\
πολλάκι δὲ κρέξασα πολύθροον, οἷά τις ἀχώ,\\
κέρτομον ἀντῳδοῖς χείλεσιν ἁρμονίαν.}
\noindent (Note de Willems.)
\noindent\textbf{1059} \textbf{Ἠχὼ λόγων ἀντῳδὸς ἐπικοκκάστρια.} Seule compagnie d'Andromède : la nymphe Écho, cachée dans un antre non loin, et qui répète les paroles. L'occasion d'un détournement comique de ces échos était trop belle pour qu'Aristophane n'en profitât pas !
\noindent Euripide – qui joue dans ce passage la nymphe Écho – répond effectivement à l'appel de l'Andromède postiche qu'essaye d'interpréter le peu crédible vieillard Mnésiloque.
\noindent Dans la tragédie éponyme d'Euripide d'antan, l'acteur d'Écho ne se montrait pas sur scène mais renvoyait les répliques depuis un antre à proximité de la falaise où était enchaînée Andromède.
\noindent Ajoutez à cela la μηχανή permettant une entrée aérienne de Persée et vous comprendrez quelque peu l'impression forte et durable que la pièce suscita auprès des spectateurs.\\ Revenons au vers 1059 de notre comédie : Ἠχὼ, λόγων ἀντῳδὸς ἐπικοκκάστρια.\\ On a passé un peu vite sur l'adjectif ἐπικοκκάστρια ; si le sens de λόγων ἀντῳδὸς est clair : « la rediseuse en chantant des paroles », qualificatifs prévisibles de la nymphe Écho (notez tout-de-même que les échos étaient probablement renvoyés en \italique{chantant}, ou au moins en \italique{persiflant}), Aristophane ajoute l'hapax ἐπικοκκάστρια.
\end{multicols}
}%
\end{adjustbox}%
\vfill
\end{adjustwidth}
\newpage
% --- PAGE DE DROITE : Traduction française + Commentaire (2/8) ---
\begin{adjustwidth}{-2cm}{-2cm}
\noindent\textbf{\footnotesize Traduction française --- Folio 26 (\texttt{f. 175r}, vv. 1044–1081 [1/4]) \hfill ΕΠΕΙΣΟΔΙΟΝ Ε' (Ἠχώ)}\\[-2pt]
\noindent\rule{16.5cm}{0.5pt}\\[2pt]
{\fontsize{8.8pt}{10.2pt}\selectfont
\ligneEst{}{\francais{me revêtit de cette crocote,}}{1044}
\ligneEst{}{\francais{puis de surcroît me dépêcha vers ce}}{1045}
\ligneEst{}{\francais{temple, au milieu de femmes.}}{}
\ligneEst{}{\francais{Hélas!}}{}
\ligneEst{}{\francais{démon implacable de mon destin !}}{}
\ligneEst{}{\francais{Ô moi le maudit !}}{}
\ligneEst{}{\francais{Qui, en présence de mes misères, ne tournerait les yeux}}{}
\ligneEst{}{\francais{sur cette inenviable souffrance ?}}{}
\ligneEst{}{\francais{Puisse l'Étoile dévorante de l'éther}}{1050}
\ligneEst{}{\francais{anéantir un infortuné !}}{}
\ligneEst{}{\francais{Car il ne m'est plus doux de voir la lumière immortelle,}}{}
\ligneEst{}{\francais{suspendu que je suis,}}{}
}
\vspace{2pt}
\noindent\rule{16.5cm}{0.6pt}\\[2pt]
\noindent\textbf{\small Commentaire philologique, métrique et scénique (vv. 1044–1081 --- partie 2/8)}\\[2pt]
\begin{adjustbox}{max totalheight=15.2cm, max width=16.5cm, keepaspectratio}%
\parbox[t]{16.5cm}{
\setlength{\columnsep}{9pt}%
\begin{multicols}{2}
\fontsize{8.4pt}{9.7pt}\selectfont
\setlength{\parindent}{0pt}%
\setlength{\parskip}{1.2pt}%
\noindent Scrutons ce mot : ἐπι-κοκκάσ-τρια ; le préfixe ἐπι- apporte selon nous sa nuance d'hostilité, ici de causticité  ; le suffixe -τρια est courant, principalement attesté chez les comiques et en prose\marqueur{1} ; reste κοκκάς, qu'une oreille française ne peut s'empêcher de rapprocher de notre `cocasse', risible mais avec une certaine violence comme le coq dont il tire son origine…  Si le volatile d'Aristophane n'est pas un coq ; κοκκάς renvoie au κόκκυξ ou \italique{coucou}, les Grecs ne prêtaient pas moins à cet oiseau de malheur une méchanceté bien supérieure. \italique{Confer} le magistrale note de Willems sur le \nieme{507} vers des \italique{Oiseaux} et les propos salés (\latin{petulantiæ sales}) qu'on jetait aux laboureurs en imitant le cri du coucou. (Soit dit en passant, on retrouve les mêmes traits dans feu nos traditions orales : le coucou chantait les événements fâcheux (dont en premier lieu, bien-sûr, le cocufiage)\marqueur{2}.\\ Le scholiaste nous est ici précieux\marqueur{3} : il définit ainsi ἐπικοκκάστρια : εἰωθυῖα γελᾶν, γελάστρια, « ayant l'habitude de rire, moqueuse ». Aristophane nous prévient donc que les échos auront quelque chose de piquant et que si on les trouve d'une « \anglais{annoying loquacity} »\marqueur{4} , c'est que, peut-être, le sel attique nous en a échappé.\\ Une autre scholie nous permet de préciser un peu la teneur des échos :
\citComm{ἐπεὶ εἰσήγαγε κακοστένακτον τὴν Ἡχὼ ὁ Εὐριπίδης ἐν τῇ Ανδρομέδᾳ, εἰς τοῦτο παίζει [...]}
\noindent La bonne interprétation de cette scholie réside dans le sens à accorder à l'adjectif κακοστένακτος : \italique{à la méchante lamentation} ?\\ Il est étonnant que les modernes éditions n'aient pas cherché à percer ce que les répliques d'Écho-Euripide pouvaient avoir d'amusantes ou de méchantes… Pourtant, elles disposaient déjà d'un exemple de détournement par la « rediseuse goguenarde » : à la fin de la scène, au vers 1094, alors que l'archer Scythe tente d'attraper Écho-Euripide et que ce dernier s'enfuit sans demander son reste, une scholie nous indiquait clairement ce que disait en son jargon l'archer Scythe :
\citComm{ἀντὶ τοῦ `οὐ χαιρήσεις'\\
À la place de : « tu ne vas pas te réjouir ».}
\end{multicols}
}%
\end{adjustbox}%
\vfill
\end{adjustwidth}
\newpage

% --- FOLIO 26 (f. 175r) : SPREAD 2.2 (Texte grec établi + Traduction française + Commentaire 3/8 & 4/8) ---
\clearpage\ifodd\value{page}\hbox{}\newpage\fi
\enTeteG{ΘΕΣΜΟΦΟΡΙΑΖΟΥΣΑΙ --- Texte grec établi (f. 175r, vv. 1044–1081 [2/4])}
\enTeteD{LES THESMOPHORIES --- Traduction française (f. 175r, vv. 1044–1081 [2/4])}
\begin{adjustwidth}{-2cm}{-2cm}
\noindent\textbf{\footnotesize Texte grec établi --- Folio 26 (\texttt{f. 175r}, vv. 1044–1081 [2/4]) \hfill ΕΠΕΙΣΟΔΙΟΝ Ε' (Ἠχώ)}\\[-2pt]
\noindent\rule{16.5cm}{0.5pt}\\[2pt]
{\fontsize{8.8pt}{10.2pt}\selectfont
\ligneEst{}{λαιμότμητ᾽ ἄχη δία μένων αἰόλαν}{}
\ligneEst{{\tiny\textit{ithyphall.}}}{νέκυσιν ἐπὶ πορείαν.}{1055}
\ligneEst{\textbf{Εὐριπίδης}}{χαῖρ᾽ ὦ φίλη παῖ· τὸν δὲ πατέρα Κηφέα}{}
\ligneEst{}{ὅς σ᾽ ἐξέθηκεν ἀπολέσειαν οἱ θεοί.}{}
\ligneEst{\textbf{Μνησίλοχος}}{σὺ δ᾽ εἶ τίς ἥτις τοὐμὸν ᾤκτιρας πάθος;}{}
\ligneEst{\textbf{Εὐριπίδης}}{Ἠχὼ, λόγων ἀντῳδὸς ἐπικοκκάστρια,}{}
\ligneEst{}{ἥπερ πέρυσιν ἐν τῷδε ταὐτῷ χωρίῳ}{1060}
\ligneEst{}{Εὐριπίδῃ καὐτὴ ξυνηγωνιζόμην.}{}
\ligneEst{}{ἀλλ᾽ ὦ τέκνον σὲ μὲν τὸ σαυτῆς χρὴ ποιεῖν,}{}
\ligneEst{}{κλαίειν ἐλεινῶς.}{}
\ligneEst{\textbf{Μνησίλοχος} {\tiny\textit{trim. iamb.}}}{\decalage{κλαίειν ἐλεινῶς.}σὲ δ᾽ ἐπικλαίειν ὕστερον.}{}
\ligneEst{\textbf{Εὐριπίδης}}{ἐμοὶ μελήσει ταῦτά γ᾽. ἀλλ᾽ ἄρχου λόγων.}{}
\ligneEst{\textbf{Μνησίλοχος} {\tiny\textit{metr. anap.}}}{ὦ νὺξ ἱερὰ}{1065}
}
\vspace{2pt}
\noindent\rule{16.5cm}{0.6pt}\\[2pt]
\noindent\textbf{\small Commentaire philologique, métrique et scénique (vv. 1044–1081 --- partie 3/8)}\\[2pt]
\begin{adjustbox}{max totalheight=15.2cm, max width=16.5cm, keepaspectratio}%
\parbox[t]{16.5cm}{
\setlength{\columnsep}{9pt}%
\begin{multicols}{2}
\fontsize{7.6pt}{8.8pt}\selectfont
\setlength{\parindent}{0pt}%
\setlength{\parskip}{1.2pt}%
\noindent On peut déduire de cette scholie que le Scythe – qui ne tient aucun compte dans son parler des aspirations – avait menacé Écho-Euripide d'un : `οὐ καιρήσεις'. Or, le manuscrit de Ravenne comportait effectivement la réponse ἐπικοκκάστρια d'Euripide-Écho\marqueur{5} : `οὐκ᾽ αἰρήσεις', c'est-à-dire en bon attique : `οὐχ αἱρήσεις', « tu ne m'attraperas pas ! ». Le cas était simple mais il aura fallu tout de même des siècles de philologiques hésitations entre les répartitions : Τοξότης (l'Archer) : « οὐ καιρήσεις » — Εὐριπίδης : « οὐ καιρήσεις » et Τοξότης : « οὐκ αἰρήσεις » — Εὐριπίδης : « οὐκ αἰρήσεις » avant que le grand Wilamowitz\marqueur{6} ne s'avise de la solution conciliatrice, la seule sensée et la seule comique : Τοξότης : « οὐ καιρήσεις » — Εὐριπίδης : « οὐχ αἱρήσεις ».\\ Bref ! les répétitions de la nymphe servent à Aristophane pour son jeu de massacre de la tragédie tant appréciée d'Euripide\marqueur{7} : par de légers détournements phonétiques (facilités par les barbarismes de l'archer Scythe), Écho transforme ce qui devait être une scène d'émotion grave en une foire d'empoigne générale.\\ C'est -- reconnaissons-le -- un exercice des plus périlleux que de s'aventurer dans des tentatives de reconstruction ; laisser ainsi libre cours à l'imagination peut exposer à de rudes retours de bâton ; pour Aristophane, nous devons cependant impérieusement nous y livrer puisqu'il vient de nous prévenir que les échos seraient \italique{moqueurs} ! Avec une réserve cependant : lorsque l'ἐπικοκκάστρια d'une réplique nous est demeurée celée, et que malgré nos efforts les seules propositions que nous avions à présenter n'étaient pas même dignes d'intégrer l'apparat critique -- le cas n'est pas isolé --, nous lui avons apposé la \italique{croix du désespoir} †.. Charge à d'autres de compléter notre ébauche ! je pense, pour ma part, ne pas avoir démérité d'Aristophane.
\noindent\textbf{1060} \textbf{πέρυσιν.} L'\italique{Andromède} avait été produite l'année précédente, cf. la scholie :
\citComm{ἐπεὶ πέρυσιν ἐδιδάχθη <ἡ> Ἀνδρομέδα.}
\noindent (Voir l'introduction sur la question de la date des \italique{Thesmophories}.)
\noindent\textbf{[Notes ad 1059]} \textbf{(1)} Si la pie est un oiseau capable de parler, selon Descartes, l'oiseau dont il question dans cette épigramme serait à en croire le \LSJ \ un corvidé cousin, de la famille des geais : un \italique{geai des chênes}, \latin{Garrulus glandarius}. (NDÉ.)
\noindent \textbf{(1)} Cf. Chantraine pour une longue liste des mots formés sur le même modèle, \italique{La formation des noms en grec ancien}, p. 106-7.
\noindent \textbf{(2)} Pour s'en convaincre, il n'y a qu'à aller feuilleter le \italique{Répertoire} Coirault \italique{des chansons françaises de tradition orale}.
\noindent \textbf{(3)} Les glosateurs avaient sans doute un texte des \italique{Thesmophories} en meilleur état que le nôtre (ce passage est le moins bien conservé de toute la pièce ; et pour cause : le comique a très tôt échappé aux scribes qui le recopiaient). De plus, ils avaient un avantage indéniable sur nous : ils pouvaient aller comparer la parodie avec l'original d'Euripide.
\end{multicols}
}%
\end{adjustbox}%
\vfill
\end{adjustwidth}
\newpage
% --- PAGE DE DROITE : Traduction française + Commentaire (4/8) ---
\begin{adjustwidth}{-2cm}{-2cm}
\noindent\textbf{\footnotesize Traduction française --- Folio 26 (\texttt{f. 175r}, vv. 1044–1081 [2/4]) \hfill ΕΠΕΙΣΟΔΙΟΝ Ε' (Ἠχώ)}\\[-2pt]
\noindent\rule{16.5cm}{0.5pt}\\[2pt]
{\fontsize{8.8pt}{10.2pt}\selectfont
\ligneEst{}{\francais{attendant comme terme d'un supplice jugulant}}{}
\ligneEst{}{\francais{un trajet rapide vers les morts.}}{1055}
\ligneEst{\textbf{Euripide}}{\francais{(Dans la coulisse et sans être vu). Salut, chère enfant. Mais ton père Cepheus}}{}
\ligneEst{}{\francais{qui t'exposa, que les dieux le confondent.}}{}
\ligneEst{\textbf{Mnésiloque}}{\francais{Qui es-tu, toi qui compatis à mon tourment ?}}{}
\ligneEst{\textbf{Euripide}}{\francais{Je suis Écho, la rediseuse goguenarde,}}{}
\ligneEst{}{\francais{celle qui, l'an passé, en ce même lieu,}}{1060}
\ligneEst{}{\francais{prêtai mon aide à Euripide.}}{}
\ligneEst{}{\francais{Songe, mon enfant, à tenir bien ta partie}}{}
\ligneEst{}{\francais{et à gémir lamentablement.}}{}
\ligneEst{\textbf{Mnésiloque}}{\francais{\espace{2.4cm}Et toi, à gémir de plus belle après moi.}}{}
\ligneEst{\textbf{Euripide}}{\francais{Je n'y manquerai pas. Commence.}}{}
\ligneEst{\textbf{Mnésiloque}}{\francais{Ô Nuit sainte,}}{1065}
}
\vspace{2pt}
\noindent\rule{16.5cm}{0.6pt}\\[2pt]
\noindent\textbf{\small Commentaire philologique, métrique et scénique (vv. 1044–1081 --- partie 4/8)}\\[2pt]
\begin{adjustbox}{max totalheight=15.2cm, max width=16.5cm, keepaspectratio}%
\parbox[t]{16.5cm}{
\setlength{\columnsep}{9pt}%
\begin{multicols}{2}
\fontsize{7.6pt}{8.8pt}\selectfont
\setlength{\parindent}{0pt}%
\setlength{\parskip}{1.2pt}%
\noindent \textbf{(4)} Austin-Olson… et pis encore (p. LXX) : « \anglais{much of the effect of the parodies of Helen and Andromeda depends on the fact that the tragic heroines are played by an individual wearing a prominent stagephallus} » (!!!) Mais pour qui prennent-ils Aristophane ?! le plus gras des vaudevillistes serait pris de scrupules devant les procédés que ces brillants professeurs prêtent à Aristophane, le plus fin des Attiques ! Il faudrait leur expliquer qu'il y a un temps pour tout, et qu'un pastiche d'euripidéenne tragédie n'est pas une scène burlesque de déshabillage ou encore une pièce sur la grève du sexe ; c'est calomnier Aristophane que de croire – et plus grave : faire croire à ses lecteurs, circonstance aggravante lorsque l'on prétend à être l'édition de référence des \italique{Thesmophories} ! – qu'il ne soit capable que de jouer d'un seul registre. Au reste, cette opinion est un expédient bien pratique qui dispense de se pencher sur ce que signifie véritablement le texte… ou de reconnaître honnêtement qu'on ne le comprend pas. Enfin, le véritable Aristophane nous dit lui-même ce qu'il pense du recours abusif à ces « phallus de cuir, épais, rouges par le bout, pour donner à rire aux gamins », c'est au vers 538 des \italique{Nuées} (voir aussi la \italique{Paix}, v. 739-753) ; que les futurs exégètes d'Aristophane ruminent bien ce passage avant de prétendre que notre poète concoctait à toutes les sauces des scènes de sexes sémillants. Et qu'ils lisent aussi l'article que Willems y avait consacré, « Le nu dans la Comédie ancienne », t. III, p. 381-400 (reproduit dans les annexes) : Aristophane recourait à l'\italique{olisbos} avec la plus grande parcimonie, et toujours pour de bonnes raisons.
\noindent \textbf{(5)} En attribuant, il est vrai, la réplique à l'archer. Mais les indications de personnages sont de peu de valeur : les anciens parchemins ne comportaient pour la plupart que de légers signes (\italique{paragraphos} ou parfois simple indentation) pour indiquer le changement d'interlocuteur, et ce sans préciser quels étaient ces interlocuteurs ; les indications de personnages furent ajoutées par des scribes très postérieurement… et sont parfois assez folkloriques. (Pour celles du \italique{Ravennas}, quelques indications sont de la main d'Euphrosynus Boninus, celui qui compulsa le manuscrit pour la première édition imprimée des \italique{Thesmophories} -- la \italique{Juntina} (Venise, 1516) des apparats critiques. Von Velsen sut identifier les annotations de la main d'Euphrosynus Boninus sur le \italique{Ravennas} et ne les fit pas apparaître dans son apparat critique. Nous l'avons suivi en cela. Cf. Albert Martin, \italique{Les scolies du manuscrit d'Aristophane à Ravenne} (1882), p. V ; A. von Velsen, \italique{Über den Codex Urbinas der Lysistrata und der Thesmophoriazusen des Aristophanes} (1871).)
\noindent \textbf{(6)} En simple note de ses \italique{Kleine Schriften}, IV (1962), p. 486. Mais ces écrits n'ont été publiés qu'après la mort du grand philologue. Si une partie des travaux modernes rétablissent correctement la menace du Scythe et son écho, il semble que le mérite en revienne plutôt à Sommerstein, qui eut la bonne intuition indépendamment de Wilamowitz.
\noindent \textbf{(7)} Rappelons, par exemple, que c'est en relisant \italique{Andromède} que le Dionysos des \italique{Grenouilles} est pris d'un soudain désir d'Euripide.
\noindent\textbf{1063} \textbf{σὲ δ᾽ ἐπικλάειν ὕστερον.} Puisqu'Écho répète tout, il est logique qu'elle répète \italique{en retour} (ἐπι-) les pleurs versés par Andromède... il y a là cependant une menace de vengeance clairement exprimée par Mnésiloque pour après sa libération (avec l'emploi à double sens du préfixe ἐπι- : hostilité et succession temporelle).
\end{multicols}
}%
\end{adjustbox}%
\vfill
\end{adjustwidth}
\newpage

% --- FOLIO 26 (f. 175r) : SPREAD 2.3 (Texte grec établi + Traduction française + Commentaire 5/8 & 6/8) ---
\clearpage\ifodd\value{page}\hbox{}\newpage\fi
\enTeteG{ΘΕΣΜΟΦΟΡΙΑΖΟΥΣΑΙ --- Texte grec établi (f. 175r, vv. 1044–1081 [3/4])}
\enTeteD{LES THESMOPHORIES --- Traduction française (f. 175r, vv. 1044–1081 [3/4])}
\begin{adjustwidth}{-2cm}{-2cm}
\noindent\textbf{\footnotesize Texte grec établi --- Folio 26 (\texttt{f. 175r}, vv. 1044–1081 [3/4]) \hfill ΕΠΕΙΣΟΔΙΟΝ Ε' (Ἠχώ)}\\[-2pt]
\noindent\rule{16.5cm}{0.5pt}\\[2pt]
{\fontsize{8.8pt}{10.2pt}\selectfont
\ligneEst{}{ὡς μακρὸν ἵππευμα διώκεις}{}
\ligneEst{}{ἀστεροειδέα νῶτα διφρεύουσ᾽}{}
\ligneEst{}{αἰθέρος ἱερᾶς}{}
\ligneEst{}{τοῦ σεμνοτάτου δι᾽ Ὀλύμπου;}{}
\ligneEst{\textbf{Εὐριπίδης}}{\decalage{τοῦ σεμνοτάτου δι᾽ Ὀλύμπου;}†δι᾽ Ὀλύμπου;†}{}
\ligneEst{\textbf{Μνησίλοχος}}{τί ποτ᾽ Ἀνδρομέδα περίαλλα κακῶν}{1070}
\ligneEst{}{μέρος ἐξέλαχον;}{}
\ligneEst{\textbf{Εὐριπίδης}}{\decalage{μέρος ἐξέλαχον;}μόρος ἔξελ᾽ ἄχος.}{}
\ligneEst{\textbf{Μνησίλοχος}}{θανάτου τλήμων}{}
\ligneEst{\textbf{Εὐριπίδης}}{\decalage{θανάτου τλήμων}†θανάτου τλήμων†.}{}
\ligneEst{\textbf{Μνησίλοχος}}{ἀπολεῖς μ᾽ ὦ γραῦ στωμυλλομένη.}{}
\ligneEst{\textbf{Εὐριπίδης}}{στωμυλιομενής.}{}
\ligneEst{\textbf{Μνησίλοχος}}{νὴ Δί᾽, ὀχληρά γ᾽ εἰσήρρηκας}{1075}
}
\vspace{2pt}
\noindent\rule{16.5cm}{0.6pt}\\[2pt]
\noindent\textbf{\small Commentaire philologique, métrique et scénique (vv. 1044–1081 --- partie 5/8)}\\[2pt]
\begin{adjustbox}{max totalheight=15.2cm, max width=16.5cm, keepaspectratio}%
\parbox[t]{16.5cm}{
\setlength{\columnsep}{9pt}%
\begin{multicols}{2}
\fontsize{7.6pt}{8.8pt}\selectfont
\setlength{\parindent}{0pt}%
\setlength{\parskip}{1.2pt}%
\noindent\textbf{1065-72} \textbf{ὦ νὺξ ἱερὰ κτλ.} Nous savons par une scholie que ces vers sont les premiers d'\italique{Andromède} (τοῦ προλόγου Ἀνδρομέδας εἰσβολή). Le mètre est anapestique\marqueur{1} ; les poètes tragiques y recourent effectivement pour certaines scènes de lamentation\marqueur{2}.
\noindent\textbf{1067} \textbf{ἀστεροειδέα νῶτα.} La voûte solide du firmament incorruptible. (Note de Willems.)
\noindent\textbf{[Notes ad 1065-72]} \textbf{(1)} De même que dans le prologue d'\italique{Iphigénie en Aulis}.
\noindent \textbf{(2)} Voir \italique{The Songs of Aristophanes}, p. 57.
\noindent\textbf{1069} \textbf{τοῦ σεμνοτάτου δι᾽ Ὀλύμπου; --- δι᾽ Ὀλύμπου.} En bonne rigueur philologique, nous devons nous demander sur quels éléments probants s'appuyer dans notre aventureuse entreprise de rétablissement de l'\allemand{Urtext} des \italique{Thesmophories}. Car il ne suffit de partir à l'aventure, toutes brides de l'imaginations abattues ; encore faut-il apporter des indices de ce que nous avançons.
\noindent Nous avons déjà évoqué des bribes de scholies ; nous en croiserons d'autres en route.
\noindent Une variante orthographique entre un mot et son écho nous fournira aussi une amorce philologique.
\noindent Cela est bien maigre... est-ce qu'Hésychius, ou les commentateurs qu'il recopie\marqueur{1}, dont nous avons assez montré toute l'importance, avait entre les mains une édition d'Aristophane assez bien conservée pour qu'y apparaisse certains détournements de la nymphe Écho ? Réponse : ce n'est pas impossible !
\noindent voyons cela avec première réplique ἐπικοκκάστρια de la nymphe Écho. \\ Voici la glose, parcellaire, d'Hésychius (ο 649) qui éclaire et justifie peut-être l'interruption d'Écho :
\citComm{ὀλυμπία· δοκός\\
olympia : poutre.}
\noindent La langue d'Aristophane possède la matrice métaphorique : \italique{poutre} $=$ \italique{poutre à laquelle on se pend} $=$ \italique{pendaison}. (Le français dispose d'une matrice semblable avec le mot \italique{corde}.) J'en veux pour preuves : \ \ \ 1° \ ce commentaire\marqueur{2} : [δοκὸς] παρὰ Ἀριστοφάνει ἐν Ταγηνισταῖς σημαίνει τὴν ἀγχόνην, commentaire repris par Hésychius (δ 2129) : δοκόν· τὴν ἀγχόνην.
\noindent \ \ \ 2° \ Dans les \italique{Oiseaux}, après avoir poutré l'importun chresmologue Hiéroclès, on l'envoie se faire pendre εἰς Ἐλύμνιον (vers 1126), ville d'Eubée dont le nom fait entendre à l'oreille grecque sa \italique{poutre} : ἐλύμνιαι· δοκοὶ ὀρόφιναι (Hésychius δ 2227).\\ Rappeler à Andromède sa condition de crucifiée \italique{pendue} à la falaise, voilà une incise qui pourrait être qualifiée de κακοστένακτος.\\ Je m'arrête ici, au bord de la reconstruction tentatrice. Car grande est la tentation de retoucher au texte même d'Andromède en y intégrant le iota d'ὀλυμπία et voir dans `δι᾽ Ὀλύμπου' une mention de Ζεύς Ὀλύμπιος\marqueur{3} ; qu'on se souvienne du vers 272 : Euripide jurait sur l'\italique{Éther, demeure de Zeus} (αἰθέρ᾽, οἴκησιν Διός\marqueur{4}). Le vrai Euripide parle ailleurs du \italique{ciel étoilé de Zeus} (\italique{Ion}, v. 1077 sq.) :
\citComm{καὶ Διὸς ἀστερωπὸς\\
ἀνεχόρευσεν αἰθήρ.}
\end{multicols}
}%
\end{adjustbox}%
\vfill
\end{adjustwidth}
\newpage
% --- PAGE DE DROITE : Traduction française + Commentaire (6/8) ---
\begin{adjustwidth}{-2cm}{-2cm}
\noindent\textbf{\footnotesize Traduction française --- Folio 26 (\texttt{f. 175r}, vv. 1044–1081 [3/4]) \hfill ΕΠΕΙΣΟΔΙΟΝ Ε' (Ἠχώ)}\\[-2pt]
\noindent\rule{16.5cm}{0.5pt}\\[2pt]
{\fontsize{8.8pt}{10.2pt}\selectfont
\ligneEst{}{\francais{quelle longue carrière tu fournis,}}{}
\ligneEst{}{\francais{en guidant ton char sur le dos constellé}}{}
\ligneEst{}{\francais{de l'éther sacré,}}{}
\ligneEst{}{\francais{à travers le très auguste Olympe.}}{}
\ligneEst{\textbf{Euripide}}{\francais{\espace{4cm}Olympe.}}{}
\ligneEst{\textbf{Mnésiloque}}{\francais{Pourquoi seule Andromède reçut-elle tant de maux}}{1070}
\ligneEst{}{\francais{en partage ?}}{}
\ligneEst{\textbf{Euripide}}{\francais{\espace{2.5cm}En partage.}}{}
\ligneEst{\textbf{Mnésiloque}}{\francais{Malheureuse que la mort...}}{}
\ligneEst{\textbf{Euripide}}{\francais{\espace{2.7cm}Malheureuse que la mort.}}{}
\ligneEst{\textbf{Mnésiloque}}{\francais{Tu me tueras, vieille, avec ton babil.}}{}
\ligneEst{\textbf{Euripide}}{\francais{Ton babil.}}{}
\ligneEst{\textbf{Mnésiloque}}{\francais{Par Zeus, ton ingérence m'excède}}{1075}
}
\vspace{2pt}
\noindent\rule{16.5cm}{0.6pt}\\[2pt]
\noindent\textbf{\small Commentaire philologique, métrique et scénique (vv. 1044–1081 --- partie 6/8)}\\[2pt]
\begin{adjustbox}{max totalheight=15.2cm, max width=16.5cm, keepaspectratio}%
\parbox[t]{16.5cm}{
\setlength{\columnsep}{9pt}%
\begin{multicols}{2}
\fontsize{7.6pt}{8.8pt}\selectfont
\setlength{\parindent}{0pt}%
\setlength{\parskip}{1.2pt}%
\noindent La mention de Zeus lorsqu'un poète évoque le ciel constellé est d'ailleurs très naturelle puisque le quatrième astre le plus lumineux, après le Soleil, la lune et Vénus, n'est autre que \italique{Jupiter}, Διὸς ἀστήρ.
\noindent\textbf{[Notes ad 1069]} \textbf{(1)} Cf. le passionnant livre posthume de Pierre Boudreaux, \italique{Le texte d'Aristophane et ses commentateurs}.
\noindent \textbf{(2)} Fragment d'Aristophane 528 K.-A.(515 Kock). Cf. l'édition Kassel-Austin des \italique{Poetae Comici Graeci} pour la juste distinction entre δοκός oxyton et δόκος paroxyton.
\noindent \textbf{(3)} Διός Ὀλυμπίου \italique{contra metrum} cependant.
\noindent \textbf{(4)} Aristophane y revient au vers 100 des \italique{Grenouilles} : αἰθέρα Διὸς δωμάτιον.
\noindent\textbf{1071} \textbf{μέρος ἐξέλαχον;} Sans \italique{tordre} le texte, il me semble que la réplique ἐπικοκκάστρια d'Écho s'entend assez bien -- soit dit avec toute la prudence qu'une telle chirurgie du texte impose :
\citComm{μέρος ἐξέλαχον --- μόρος ἔξελ(ε)\marqueur{1} ἄχος\\
La Mort\marqueur{2} extirpe la douleur.}
\noindent Andromède se plaint des malheurs qui l'accablent ; Écho la rassure à sa manière : la mort l'en délivrera !.. Était-ce ce genre de répliques que le scholiaste qualifiait de `κακοστένακτος' ?\\ On trouve un tour analogue chez Sophocle\marqueur{3} :
\citComm{Ἀλλ᾽ ἔσθ᾽ ὁ θάνατος λοῖσθος ἰατρὸς νόσων.}
\noindent Et dans \italique{la langue des dieux} :
\citComm{\textbf{Le trépas vient tout guérir} ;\\
Mais ne bougeons d'où nous sommes.\\
Plutôt souffrir que mourir,\\
C'est la devise des hommes.}
\noindent\textbf{[Notes ad 1071]} \textbf{(1)} Aoriste \italique{gnomique}.
\noindent \textbf{(2)} Cette acception de μόρος a été relevée par à peu près tous les lexicographes anciens : Hésychius μ 1683 ; Photius μ 538 ; Συναγωγὴ λέξεων χρησίμων μ 267 ; Suidas μ 1257 ; \latin{Etymologicum genuinum} \italique{sub voce} μόρος· ὁ θάνατος ; \&c.
\noindent \textbf{(3)} $=$ fr. 698 Radt $=$ fr. 636 Dindorf. Stobée (IV, 52, 26) a inséré la citation à son \italique{Florilège} et Michel Apostolius (II, 52c) la reprit pour ses \italique{Parœmiæ}.
\noindent\textbf{1072} \textbf{θανάτου τλήμων...} L'on connaît par une scholie la fin du vers original d'\italique{Andromède} :
\citComm{λείπει `μέλλουσα τυχεῖν'.}
\end{multicols}
}%
\end{adjustbox}%
\vfill
\end{adjustwidth}
\newpage

% --- FOLIO 26 (f. 175r) : SPREAD 2.4 (Texte grec établi + Traduction française + Commentaire 7/8 & 8/8) ---
\clearpage\ifodd\value{page}\hbox{}\newpage\fi
\enTeteG{ΘΕΣΜΟΦΟΡΙΑΖΟΥΣΑΙ --- Texte grec établi (f. 175r, vv. 1044–1081 [4/4])}
\enTeteD{LES THESMOPHORIES --- Traduction française (f. 175r, vv. 1044–1081 [4/4])}
\begin{adjustwidth}{-2cm}{-2cm}
\noindent\textbf{\footnotesize Texte grec établi --- Folio 26 (\texttt{f. 175r}, vv. 1044–1081 [4/4]) \hfill ΕΠΕΙΣΟΔΙΟΝ Ε' (Ἠχώ)}\\[-2pt]
\noindent\rule{16.5cm}{0.5pt}\\[2pt]
{\fontsize{8.8pt}{10.2pt}\selectfont
\ligneEst{}{λίαν.}{}
\ligneEst{\textbf{Εὐριπίδης}}{\decalage{λίαν.}†λίαν†.}{}
\ligneEst{\textbf{Μνησίλοχος}}{ὦγάθ᾽ ἔασόν με μονῳδῆσαι,}{}
\ligneEst{}{καὶ χαριεῖ μοι παῦσαι.}{}
\ligneEst{\textbf{Εὐριπίδης}}{\decalage{καὶ χαριεῖ μοι παῦσαι.}᾽παΰσαι.}{}
\ligneEst{\textbf{Μνησίλοχος}}{βάλλ᾽ ἐς κόρακας.}{}
\ligneEst{\textbf{Εὐριπίδης}}{\decalage{βάλλ᾽ ἐς κόρακας.}βάλλε σύ κόρακας.}{}
\ligneEst{\textbf{Μνησίλοχος}}{†τί τὸ κακόν†;}{1080}
\ligneEst{\textbf{Εὐριπίδης}}{\decalage{†τί τὸ κακόν†;}†τί τὸ κακόν;†}{}
\ligneEst{\textbf{Μνησίλοχος}}{\decalage{†τί τὸ κακόν†;†τί τὸ κακόν;†}ληρεῖς.}{}
\ligneEst{\textbf{Εὐριπίδης}}{\decalage{†τί τὸ κακόν†;†τί τὸ κακόν;†ληρεῖς.}†ληρεῖς†.}{}
\ligneEst{\textbf{Μνησίλοχος}}{οἴμωζ᾽.}{1081}
\ligneEst{\textbf{Εὐριπίδης}}{\decalage{οἴμωζ᾽.}†οἴμωζ᾽†.}{}
}
\vspace{2pt}
\noindent\rule{16.5cm}{0.6pt}\\[2pt]
\noindent\textbf{\small Commentaire philologique, métrique et scénique (vv. 1044–1081 --- partie 7/8)}\\[2pt]
\begin{adjustbox}{max totalheight=15.2cm, max width=16.5cm, keepaspectratio}%
\parbox[t]{16.5cm}{
\setlength{\columnsep}{9pt}%
\begin{multicols}{2}
\fontsize{7.6pt}{8.8pt}\selectfont
\setlength{\parindent}{0pt}%
\setlength{\parskip}{1.2pt}%
\noindent\textbf{1073-5} \textbf{ὦ γραῦ στωμυλλωμένη — στωμυλωμένη — νὴ Δί᾽.} Le manuscrit de Ravenne (qui, rappelons-le, est le seul témoin complet de notre pièce) comporte deux orthographes différentes pour στωμυλλωμένη : avec un double lambda -λλ- au vers 1073 et un unique -λ- au vers 1074... Étourderie ou fantaisie de notre scribe jonglant avec une orthographe dont il est peu sûr ? Possible... ou bien preuve q'il recopiait scrupuleusement un texte sans parfois – comme ici – le comprendre ? c'est ce que nous allons tenter de montrer : il se trouve que le verbe στωμύλλω, \italique{babiller}, \italique{déblatérer}, s'écrit effectivement avec deux lambdas en grec classique. Cependant, l'adjectif στωμύλος (\italique{qui parle beaucoup}, et souvent avec nuance méliorative par rapport au verbe) qui en dérive n'en comporte qu'un seul ! Si donc, l'on fait le pari que le scribe ne s'y est pas trompé en conservant l'orthographe, savoir : on est en présence d'un participe moyen-passif au vers 1073 et de l'adjectif au vers 1074, il est alors aisé de retrouver la véritable réponse \italique{ἐπικοκκάστρια} d'Écho : στωμυλιο-μενής, « [je suis] de tempérament prolixe ».
\noindent Le composé est de la même facture que στωμυλιο-συλλεκτάδης, hapax qu'Aristophane mit dans la bouche d'Eschyle au vers 839 des \italique{Grenouilles} pour clore le bec à Euripide, cet intarissable \italique{collectionneur de commérages}. La seconde partie du composé est forgée avec le suffixe -μενής, \italique{au tempérament}, \italique{à la disposition...} (sur le modèle d'εὐ-μενής, δυσ-μενής, πρευ-μενής \&c).\\ Notre conjecture est renforcée par la concessive de Mnésiloque au vers suivant : il abonde, νὴ Δί᾽, ... néanmoins, γ᾽ : « ça oui par Zeus ! du moins... ».
\noindent\textbf{1078} \textbf{καὶ χαριεῖ μοι παῦσαι --- παῦσαι.} Mnésiloque-Andromède, qui ne peut chanter sa monodie correctement, demande à Euripide-Écho de cesser de l'interrompre :
\citComm{καὶ χαριεῖ μοι παῦσαι...}
\noindent «~Tu m'obligeras d'arrêter...~» -- «~De crier~» interrompt facilement Écho. (Et, bien-sûr, en le lui criant à l'oreille :)
\citComm{ (ἐ)παΰσαι.}
\noindent (Le verbe ἐπαΰω, attesté par Eschyle dans ses \italique{Choéphores} (v. 828), est forgé avec le préfixe ἐπι- à l'instar d'ἐπικοκκάστρια et le verbe αὔω, \italique{crier} (cf. le vers 1019).)\\ Une telle provocation explique la réaction d'emportement de Mnésiloque :
\citComm{βάλλ᾽ ἐς κόρακας.}
\noindent\textbf{1079} \textbf{βάλλ᾽ ἐς κόρακας — βάλλ᾽ ἐς κόρακας.} Dans notre entreprise fort périlleuse de dénicher des jeux de mots tout au long de ce passage, en voici un qui est moins spéculatif et qui va nous permettre de reprendre haleine !
\noindent Les Grecs usaient d'une formule d'imprécation : ἐς κόρακας, équivalente de notre « [va] au diable ! ». Aristophane l'emploie fréquemment avec le verbe βάλλειν : \italique{Guêpes} v. 835 ; \italique{Ploutos} v. 782 ; \italique{Nuées} v. 133 – digressons d'ailleurs sur ce passage des \italique{Nuées} puisqu'Aristophane recourt au même jeu de mots qu'ici : Strepsiade arrive devant le pensoir de Socrate et toque à la porte, il se fait recevoir par un peu accueillant :
\citComm{Βάλλ᾽ ἐς κόρακας· / Τίς ἐσθ᾽ ὁ κόψας τήν θύραν;}
\noindent signalons à toutes fins utiles que les κόρακες dont parle le disciple de Socrate sont des marteaux de porte\marqueur{1}.
\end{multicols}
}%
\end{adjustbox}%
\vfill
\end{adjustwidth}
\newpage
% --- PAGE DE DROITE : Traduction française + Commentaire (8/8) ---
\begin{adjustwidth}{-2cm}{-2cm}
\noindent\textbf{\footnotesize Traduction française --- Folio 26 (\texttt{f. 175r}, vv. 1044–1081 [4/4]) \hfill ΕΠΕΙΣΟΔΙΟΝ Ε' (Ἠχώ)}\\[-2pt]
\noindent\rule{16.5cm}{0.5pt}\\[2pt]
{\fontsize{8.8pt}{10.2pt}\selectfont
\ligneEst{}{\francais{au dernier point.}}{}
\ligneEst{\textbf{Euripide}}{\francais{\espace{1cm}Au dernier point.}}{}
\ligneEst{\textbf{Mnésiloque}}{\francais{Mon brave, laisse-moi chanter ma monodie.}}{}
\ligneEst{}{\francais{Tu m'obligeras : cesse !}}{}
\ligneEst{\textbf{Euripide}}{\francais{\espace{3.5cm}(Criard :) C'est ce... que j'allais dire.}}{}
\ligneEst{\textbf{Mnésiloque}}{\francais{Va-t'en aux corbeaux.}}{}
\ligneEst{\textbf{Euripide}}{\francais{\espace{2.9cm}Va-t'en aux corbeaux.}}{}
\ligneEst{\textbf{Mnésiloque}}{\francais{Qu'est-ce qui t'arrive ?}}{1080}
\ligneEst{\textbf{Euripide}}{\francais{\espace{1.5cm}Qu'est-ce qui t'arrive ?}}{}
\ligneEst{\textbf{Mnésiloque}}{\francais{\espace{3cm}Tu veux rire.}}{}
\ligneEst{\textbf{Euripide}}{\francais{\espace{4.2cm}Tu veux rire.}}{}
\ligneEst{\textbf{Mnésiloque}}{\francais{Va te promener..}}{1081}
\ligneEst{\textbf{Euripide}}{\francais{\espace{1.1cm}Va te promener..}}{}
}
\vspace{2pt}
\noindent\rule{16.5cm}{0.6pt}\\[2pt]
\noindent\textbf{\small Commentaire philologique, métrique et scénique (vv. 1044–1081 --- partie 8/8)}\\[2pt]
\begin{adjustbox}{max totalheight=15.2cm, max width=16.5cm, keepaspectratio}%
\parbox[t]{16.5cm}{
\setlength{\columnsep}{9pt}%
\begin{multicols}{2}
\fontsize{7.6pt}{8.8pt}\selectfont
\setlength{\parindent}{0pt}%
\setlength{\parskip}{1.2pt}%
\noindent Le disciple tint donc à peu près ce langage :
\citComm{Βάλλε κόρακας·/ Τίς ἐσθ᾽ ὁ κόψας τὴν θύραν;\\
Bas les marteaux ! Qui est ce tambourineur de porte ?}
\noindent (compatible avec le mètre iambique et surtout avec le bon sens : qui – fût-il disciple de Socrate – s'aviserait de demander l'identité de quelqu'un qu'on vient d'envoyer se faire voir chez les corbeaux grecs ?)\\ Il en va de même dans le passage qui nous intéresse présentement : Écho conseille frivolement à Mnésiloque de retirer les crochets qui le retiennent : « Mais va te faire voir ! — Et toi mets bas le fermoir ! »
\noindent\textbf{[Notes ad 1079]} \textbf{(1)} Sens attesté par Pollux, X, 22.20.
\noindent\textbf{1080} \textbf{τί τὸ κακόν;} La leçon du \italique{Ravennnas}, «~τί τὸ κακόν; --- τί τὸ κακόν;~», ne respecte pas le mètre anapestique. Gardons-nous pour autant de supprimer avec Bentley l'article `τὸ'.
\noindent Si la formule «~τί τὸ κακόν;~» (ou d'autres très approchantes) se rencontre couramment chez Aristophane\marqueur{1}, sa version lapidaire, «~τί κακόν;~», elle, ne se trouve jamais\marqueur{2}. La bonne solution n'est donc pas celle de Bentley, et pour notre part nous nous résignons à \latin{crux}-désespérer †
\noindent\textbf{[Notes ad 1080]} \textbf{(1)} \italique{Thesmophories}, v. 610 ; \italique{Paix}, v. 322 ; \italique{Oiseaux}, v. 1213 ; \italique{Guêpes} v. 1136, τουτὶ τὸ κακὸν τί ἐστι ; \italique{Paix}, v. 181.
\noindent \textbf{(2)} Certains éditeurs prétendent que si : au vers 1085. En vérité, ces faussaires commettent là une double forfaiture en «~corrigeant~» le présent vers \italique{metri gratia} ainsi que le vers 1085 ; le second forfait servant d'alibi au premier : au vers 1085, ils font apparaître «~une seconde fois~» le tour `τί κακόν' alors que le manuscrit porte `σί κακόν', confirmé par le scholiaste (ἀντὶ τοῦ `σοὶ κακόν'). Et quand bien même : c'est dans la bouche d'un Scythe qu'Aristophane la met au vers 1085, l'argument pour la faire dire aussi à Mnésiloque est irrecevable. En dernière analyse, ces éditeurs manquent cruellement de bon sens : le Scythe en son jargon ne tient déjà aucun compte des aspirations, de la grammaire, il ignore sa conjugaison et ses déclinaisons, il iotacise à l'envi... et l'on voudrait encore qu'il sigmatise ses `τ' ?! un tel massacre du grec aurait achevé d'être tout-à-fait inintelligible aux spectateurs.
\noindent\textbf{1080-2} \textbf{τί τὸ κακόν; [...] ὀτότυζ᾽.} Quel est le jeu de scène de tout ce passage, dont les répliques d'Écho demeurent obscures ? Il me semble qu'on peut peut y répondre avec vraisemblance : Euripide-Écho s'est approchée de Mnésiloque (cf. v. 1090, γυνὴ πλησίον αὕτη) ; la raison en est évidente : Euripide tente de libérer Mnésiloque.
\noindent Écho parvient cependant à se cacher dans un premier temps du Scythe (οὖτος τί λαλῖς; --- πῶτε τὸ πωνή;) ; c'est l'iconographie qui nous fournit la réponse : l'acteur jouant Écho se cachait derrière un grand voile\marqueur{1}.
\noindent Et lorsque l'archer Scythe veut mettre la main sur Euripide, ce dernier \italique{s'envole} (ποῖ ποῖ πεύγει -- certainement à l'aide de la μηχανή), lançant un narquois `οὐχ αἱρήσεις'.
\noindent\textbf{[Notes ad 1080-2]} \textbf{(1)} Sur un vase attique à figure rouge du \nieme{IV} siècle av. J.-C. (n° 7864 au catalogue de Beazley) représentant Dionysos et sa suite, on voit la nymphe Écho, ailée, se voiler le visage derrière son manteau. Ce détail est reproduit en annexe.
\end{multicols}
}%
\end{adjustbox}%
\vfill
\end{adjustwidth}
\newpage
\clearpage\ifodd\value{page}\hbox{}\newpage\fi
\enTeteG{ΕΠΕΙΣΟΔΙΟΝ Ε' (Περσεύς) --- Codex Ravennas 429, f. 175v (Classense 177v) --- vv. 1082–1127}
\begin{adjustwidth}{-1.5cm}{-3.0cm}
\begin{adjustbox}{max totalheight=21.0cm, max width=16.5cm, keepaspectratio}
\begin{minipage}{16.5cm}
{\fontsize{7.2pt}{8.5pt}\selectfont
\noindent\textbf{Folio 27 --- Codex Ravennas 429, \texttt{f. 175v} (\texttt{Classense 177v}) : ΕΠΕΙΣΟΔΙΟΝ Ε' (Περσεύς), vv. 1082–1127}\hfill \textit{Témoins :} $\text{R}^1, \Rpc, \Rdeux$ (\Mon)\par\vspace{2pt}
\hrule height 0.4pt\vspace{2.5pt}
}
{\fontsize{8.1pt}{9.5pt}\selectfont
\ligneRav{1082}{μ\auDessus{ν}{Η}}{\hspace{1.1em}ὀτότυζ᾽. : †ὀτότυζ᾽†·}{\textbf{1082} \lecon{ὀτότυζ᾽} R : ὀλόφυγ᾽ Bibi}
\ligneRav{1083}{τοξ/}{\hspace{1.1em}οὖτος τί λαλ\inRas{εῖ}{ῖ}ς; \hfill \textbf{[Ἠχώ :]} οὖτος τί λαλ\inRas{εῖ}{ῖ}ς;}{\textbf{1083} οῦτος οὗτος R : \lecon{οὖτος} bis Brunck \espApp λαλεῖς bis \Rac : \lecon{λαλῖς} bis \Rdeux, Σ}
\ligneRav{1084}{τοξ/}{\hspace{1.1em}πρυτάνεις καλέσω· \hfill \textbf{[Ἠχώ :]} πρυτάνεις καλέσω·}{}
\ligneRav{1085}{τοξ/}{\hspace{1.1em}σί κακόν; \hfill \textbf{[Ἠχώ :]} σί κακόν;}{\textbf{1085} σί bis R, sch. : τί Coulon}
\ligneRav{1086}{τοξ/}{\hspace{1.1em}\difference{†}π\inRas{ρ}{}ῶτε\difference{†} τὸ φωνή; \hfill \textbf{[Ἠχώ :]} \difference{†}π\inRas{ρ}{}ῶτε\difference{†} τὸ φωνή;}{\textbf{1086} π\inRas{ρ}{}ωτε bis \Rpc : πωτε bis \Mon}
\ligneRav{1087}{τοξ/}{\hspace{1.1em}σὺ λαλ\auDessus{ῖ}{εῖ}ς; \hfill \textbf{[Ἠχώ :]} σὺ λαλ\auDessus{ῖ}{εῖ}ς;}{\textbf{1087} \lecon{λαλῖς} bis R : λαλεῖς bis \Rpc}
\ligneRav{1088}{τοξ/}{\hspace{1.1em}\difference{κλαύσαι}· \hfill \textbf{[Ἠχώ :]} \difference{κλαύσαι}·}{\textbf{1088} κλαύσαιμι bis R : κλαύσει bis Brunck : \lecon{κλαύσαι} bis Rogers : κλαῦσ' ἔτι bis Austin}
\ligneRav{1089}{τοξ/}{\hspace{1.1em}κακκάσκι μοι; \hfill \textbf{[Ἠχώ :]} κακκάσκι μοι;}{\textbf{1089} κακκασκι bis R : κἀγκάσκις bis van Leeuwen (<ς> add. Fritzsche) : κἀγκάσκι bis Coulon : alii alia}
\ligneRav{1090}{χορ/}{\hspace{1.1em}μὰ Δί᾽ ἀλλὰ γυνὴ πλησίον αὕτη·}{\textbf{1090} (Μνη.) edd. plerique : \lecon{(Χορ.)} Bibi, cf. Σ}
\ligneRav{1091}{\auDessus{εὐ}{ρ}}{\hspace{1.1em}\difference{πλήξασ᾿ αὕτη}·}{\textbf{1091} πλησίον αὕτη R : \lecon{πλήξασ᾿ αὕτη} Bibi}
\ligneRav{1092}{τοξ/}{\hspace{1.1em}ποῦ ᾿στ᾽ ἠ μιαρά; : καὶ δὴ φεύγει·}{\textbf{1092} 'σθ' ἡ R : ᾿στ᾽ ἠ Brunck \espApp (Μνη.) edd. plerique : \lecon{(Χορ.)} Bibi}
\ligneRav{1093}{τοξ/}{\hspace{1.1em}ποῖ ποῖ \difference{πεύγεις}; : ποῖ ποῖ πέτει·}{\textbf{1093} φεύγεις R : \lecon{πεύγεις} Enger \espApp (Εὐρ.) ποῖ ποῖ πέτει om. R : \lecon{suppl.} Bibi}
\ligneRav{1094}{τοξ/}{\hspace{1.1em}οὐ καιρήσεις. : \difference{οὐχ αἱρήσεις}·}{\textbf{1094} (Τοξ.) οὐ καιρήσεις om. R : \lecon{suppl.} Brunck ap. Σ \espApp οὐκ᾽ αἰρήσεις; R : \lecon{οὐχ αἱρήσεις.} Wilamowitz (\italique{Kleine Schriften} IV, p.486n.1)}
\ligneRav{1095}{τοξ/}{\hspace{1.1em}ἔτι γὰρ γρύζεις; : †ἔτι γὰρ γρύζεις;†}{}
\ligneRav{1096}{τοξ/}{\hspace{1.1em}λαβὲ \difference{τὴ} μιαρά. : †λαβὲ τὴ μιαρά.†}{\textbf{1096} τῇ et τῆ R : \lecon{τὴ} bis Bothe}
\ligneRav{}{τοξ/}{\hspace{1.1em}λάλο καὶ κατάρατο γύναικο·}{}
\ligneRav{1098}{ὡς περσευς}{ὦ θεοὶ τίν᾽ εἰς γῆν βαρβάρων ἀφίγμεθα}{\textbf{1098} ευρ\auDessus{ι}{π} ως περσευς in mg. add. \Rpc, \Mon}
\ligneRav{}{}{ταχεῖ πεδίλῳ; διὰ μέσου γὰρ αἰθέρος}{}
\ligneRav{1100}{}{τέμνων κέλευθον πόδα τίθημ᾽ ὑπόπτερον}{}
\ligneRav{1101}{}{Περσεὺς πρὸς Ἄργος \difference{ναυστολῶν} τὸ Γοργόνος}{\textbf{1101} ναυτολῶν R : \lecon{ναυστολῶν} Canini (\italique{navigans} iam Divus)}
\ligneRav{1102}{—}{$[\![\text{in marg. inf. \Rdeux{} : }$κάρα κομίζων : τί λέγι; τὴ \difference{Γόργος} πέρι$]\!]$}{\textbf{1102} Γοργόνος R : \lecon{Γόργος} Σ, Bothe}
\ligneRav{1103}{—}{$[\![\text{in marg. inf. \Rdeux{} : }$τὸ γραμματέο σὺ τὴ \difference{κεπαλή}; : τὴν Γοργόνος$]\!]$}{\textbf{1103} κεφαλή R : \lecon{κεπαλή} Bothe}
\ligneRav{}{—}{ἔγωγε φημί. : Γόργο τοι κἀγὼ λέγι·}{}
\ligneRav{1105}{\auDessus{εὐ}{ρ}}{ἔα· τίν᾽ ὄχθον τόνδ᾽ ὁρῶ καὶ παρθένον}{}
\ligneRav{}{}{θεαῖς ὁμοίαν ναῦν ὅπως ὡρμισμένην;}{}
\ligneRav{1107}{κ\auDessus{η}{δ} ἀνδρ.}{ὦ ξένε \difference{κατοίκτιρόν} με τὴν παναθλίαν,}{\textbf{1107} κ\auDessus{η}{δ} ως ανδρομ\auDessus{ε}{δ} in mg. \Rpc, \Mon \espApp κατοικτειρόν R : \lecon{κατοίκτιρόν} van Herwerden}
\ligneRav{1108}{—}{λῦσόν με δεσμῶν. : οὐκὶ μὶ λαλῆσι σύ;}{\textbf{1108} \lecon{μὶ λαλῆσι} R : μὴ λαλῆσι Brunck}
\ligneRav{1109}{}{κατάρατο τολμᾷς ἀποτανουμένη \difference{λαλῖς};}{\textbf{1109} λαλᾶς R : \lecon{λαλῖς;} Bergk (cf. vv. 1082 et 1087) : λαλᾷς; Hall-Geldart}
\ligneRav{1110}{\auDessus{εὐ}{ρ}}{ὦ παρθέν᾽ \difference{οἰκτίρω} σὲ κρεμαμένην ὁρῶν·}{\textbf{1110} οἰκτείρω R : \lecon{οἰκτίρω} van Herwerden}
\ligneRav{}{τοξ/}{οὐ παρτέν᾽ ἐστίν, ἀλλ᾽ ἀμαρτωλὴ γέρων}{}
\ligneRav{}{—}{καὶ κλέπτο καὶ πανοῦργο. : ληρεῖς ὦ Σκύθα·}{}
\ligneRav{}{}{αὕτη γάρ ἐστιν Ἀνδρομέδα παῖς Κηφέως·}{}
\ligneRav{1114}{τοξ/}{σκέψαι τὸ \difference{κύστο}· \difference{μή τι} \difference{μικκὸν} παίνεται;}{\textbf{1114} σκύτο R : \lecon{κύστο} Scaliger grat. metr. : σῦκο Sommerstein \espApp μῆτι R : \lecon{μή τι} \Mon (\italique{numquid} Divus) \espApp μικτὸν R : \lecon{μικκὸν} Bentley : μικρὸ Enger}
\ligneRav{1115}{\auDessus{εὐ}{ρ}}{φέρε δεῦρό μοι τὴν χεῖρ᾽, ἵν᾽ ἅψωμαι κόρης·}{\textbf{1115} δεῦρο δεῦρο R : \lecon{δεῦρο} Biset}
\ligneRav{}{}{φέρε Σκύθ᾽· ἀνθρώποισι γὰρ νοσήματα}{}
\ligneRav{}{}{ἅπασίν ἐστιν· ἐμὲ δὲ καὐτὸν τῆς κόρης}{}
\ligneRav{}{—}{ταύτης ἔρως εἴληφεν. : οὐ ζηλῶσί σε·}{}
\ligneRav{1119}{}{ἀτὰρ εἰ \difference{τὸ πρωκτὸ} δεῦρο περιεστραμμένον,}{\textbf{1119} τῶ πρωκτῶ R : \lecon{τὸ πρωκτὸ} Biset}
\ligneRav{1120}{}{οὐκ \difference{ἐπτόνησά σ᾽} αὐτὸ πυγίζεις ἄγων·}{\textbf{1120} ἐπτόνησας R : \lecon{ἐπτόνησά σ᾽} Bentley}
\ligneRav{}{\auDessus{εὐ}{ρ}}{τί δ᾽ οὐκ ἐᾷς λύσαντά μ᾽ αὐτὴν ὦ Σκύθα}{}
\ligneRav{1122}{}{πεσεῖν \difference{ἐς} εὐνὴν καὶ γαμήλιον λέχος;}{\textbf{1122} ἐστ᾽ R : ἐς τ᾽ \Mon : \lecon{ἐς} Grynaeus}
\ligneRav{}{τοξ/}{εἰ σπόδρ᾽ ἐπιτυμεῖς τὴ γέροντο πύγισο,}{}
\ligneRav{1124}{}{τὴ σανίδο τρήσας \difference{ἐξόπιστο} πρώκτισον·}{\textbf{1124} ἐξόπισθο R : \lecon{ἐξόπιστο} Daubuz}
\ligneRav{1125}{\auDessus{εὐ}{ρ}}{μὰ Δί᾽ ἀλλὰ λύσω \difference{δεσμά}. : μαστιγῶ σ᾽ ἄρα·}{\textbf{1125} δέμας R : \lecon{δεσμά} Grynaeus}
\ligneRav{1126}{\auDessus{εὐ}{ρ}}{καὶ μὴν ποιήσω τοῦτο : \difference{τὸ κεπαλή σ᾽} (τὸ κ\inRas{αὶ }{ε}παλῆς) ἄρα}{\textbf{1126} τὸ κεπαλῆς (ex καὶ παλῆς) \Rpc, \Mon : \lecon{τὸ κεπαλή σ᾽} Brunck}
\ligneRav{1127}{}{τῇ ξιπομάκαιρᾳ \difference{ἀποκεκόψι} (ἀποκεκόψ\inRas{᾽}{οι}) τουτογί·}{\textbf{1127} ἀποκέκοιψ᾽ \Rac : ἀποκεκόψοι \Rpc : ἀποκεκόψο \Mon : \lecon{ἀποκεκόψι} Brunck \espApp πέπυσται R : \lecon{πέπαισται} Σ}
}
\vspace{3pt}\noindent\rule{4.5cm}{0.4pt}\par\vspace{1.5pt}
{\fontsize{7.0pt}{8.3pt}\selectfont
\noindent\textbf{Loci paralleli (vv. 1082–1127) ---} \textbf{1125} (δεσμά) Moeris δ 28\par
}
\vspace{2.5pt}\noindent\rule{4.5cm}{0.4pt}\par\vspace{1.5pt}
{\fontsize{6.9pt}{8.2pt}\selectfont
\noindent\textbf{Suite (en commentaire) de l'apparat critique \& notes paléographiques (\texttt{f. 175v}) ---} \textbf{Omission des vv. 1102–1103 par $\text{R}^1$ et ajout par \Rdeux{} dans la marge inférieure :} Le copiste $\text{R}^1$ avait sauté les \textbf{vv. 1102–1103} ($\text{Γόργος πέρι λέγις τὸ γραμματέο / σὺ τὴν κεπαλή;}$ — omis également par \Mon{} qui n'a pas regardé le bas de page) ; le réviseur \textbf{\Rdeux{} les a rajoutés tout en bas du folio dans la marge inférieure}, avec un signe de renvoi. --- \textbf{Analyse critique des vv. 1085–1086 par Louis Hauseux (\texttt{vv 642-fin-bis.tex}, ll. 4520, 4588–4612) :} --- \textbf{Deux notes scéniques marginales de \Rpc{} :} Au \textbf{v. 1098}, \lecon{ευρ\auDessus{ι}{π} ως περσευς} ; au \textbf{v. 1107}, \lecon{κ\auDessus{η}{δ} ως ανδρομ\auDessus{ε}{δ}}. --- \textbf{Prothèse de $\sigma$ au v. 1114 :} R écrit \lecon{σκύτο} pour \lecon{κύστο} (Scaliger, = $\text{κύσθον}$ en prononciation scythe), parallèle exact invoqué par Louis Hauseux pour sa conjecture τὸν κύσθον au v. 557. --- \textbf{v. 1082 :} \lecon{ὀτότυζ'} R (ὀλόφυγ' \lecon{Bibi} conj. in app.). --- \textbf{v. 1083 :} οῦτος οὗτος R : \lecon{οὖτος} bis Brunck \espApp λαλεῖς bis \Rac{} : \lecon{λαλῖς} bis $\text{R}^2, \Sigma$. --- \textbf{v. 1085 :} \textbf{σί κακόν; bis} R, $\Sigma^{\text{R}}$ (ἀντὶ τοῦ 'σοὶ κακόν', défendu par \lecon{Bibi} contre τί κακόν; Coulon). --- \textbf{v. 1086 :} π\inRas{ρ}{}ωτε bis \Rpc{} (πωτε bis \Mon{}) : \lecon{†πῶτε†} bis \lecon{Bibi} (rejetant l'interprétation πόθεν de $\Sigma^{\text{R}}$). --- \textbf{vv. 1090 et 1092 :} (Μνη.) edd. plerique : \lecon{(Χορ.)} \lecon{Bibi}, cf. $\Sigma^{\text{R}}$. --- \textbf{v. 1091 :} πλησίον αὕτη R : \lecon{πλήξασ' αὕτη} \lecon{Bibi}. --- \textbf{v. 1093 :} φεύγεις R : \lecon{πεύγεις} Enger \espApp \lecon{<(Εὐρ.) ποῖ ποῖ πέτει>} om. R : \textbf{suppl. \textbf{Bibi}}. --- \textbf{v. 1094 :} <(Τοξ.) οὐχ αἱρήσεις> om. R : \textbf{suppl. Brunck} ap. $\Sigma$ \espApp οὐκ' αἰρήσεις; R : \lecon{οὐχ αἱρήσεις.} Wilamowitz (\italique{Kleine Schriften} IV, p. 486 n. 1). --- \textbf{v. 1098 :} \lecon{ευρ\auDessus{ι}{π} ως περσευς} in mg. add. $\text{R}^{\text{pc}}, \text{M}$. --- \textbf{vv. 1102–1103 :} om. $\text{R}^1, \text{M}$ : \textbf{in mg. inferiore add. \Rdeux{}} \espApp \textbf{v. 1102 :} Γοργόνος R : \lecon{Γόργος} $\Sigma$, Bothe \espApp \textbf{v. 1103 :} κεφαλή R : \lecon{κεπαλή} Bothe. --- \textbf{v. 1107 :} \lecon{κ\auDessus{η}{δ} ως ανδρομ\auDessus{ε}{δ}} in mg. $\text{R}^{\text{pc}}, \text{M}$ \espApp κατοικτειρόν R : \lecon{κατοίκτιρόν} van Herwerden. --- \textbf{v. 1115 :} δεῦρο δεῦρο R (dittographie) : \lecon{δεῦρο} Biset. --- \textbf{v. 1125 :} δέμας R (métathèse de $\sigma$) : \lecon{δεσμά} Grynaeus. --- \textbf{v. 1127 :} ἀποκέκοιψ' \Rac{} (ἀποκεκόψοι \Rpc{}, ἀποκεκόψο \Mon{}) : \lecon{ἀποκεκόψι} Brunck \espApp πέπυσται R (faute récurrente de R, cf. v. 1227) : \lecon{πέπαισται} $\Sigma$.\par
}
\end{minipage}
\end{adjustbox}
\end{adjustwidth}
\newpage
\enTeteD{Codex Ravennas 429 (Biblioteca Classense) --- Fac-similé du f. 175v (vv. 1082–1127)}
\begin{adjustwidth}{-1.5cm}{-3.0cm}
\begin{center}
\includegraphics[width=16.5cm,height=20.5cm,keepaspectratio]{/workspaces/Thesmophories/dossier_ravennas_thesmophories/folios_individuels_ff162v-177r/27_f175v_classense177v_vv1082-1127.jpg}\par
\vspace{2pt}
{\scriptsize\textbf{Codex Ravennas 429, f. 175v (Classense 177v)} --- ΕΠΕΙΣΟΔΙΟΝ Ε' (Περσεύς), \textbf{vv. 1082–1127} ($\text{R}^1, \Rpc, \Rdeux$ (\Mon)).}
\end{center}
\end{adjustwidth}
\newpage
\clearpage\ifodd\value{page}\hbox{}\newpage\fi% --- FOLIO 27 (f. 175v) : SPREAD 2.1 (Texte grec établi + Traduction française + Commentaire 1/4 & 2/4) ---
\clearpage\ifodd\value{page}\hbox{}\newpage\fi
\enTeteG{ΘΕΣΜΟΦΟΡΙΑΖΟΥΣΑΙ --- Texte grec établi (f. 175v, vv. 1082–1127 [1/2])}
\enTeteD{LES THESMOPHORIES --- Traduction française (f. 175v, vv. 1082–1127 [1/2])}
\begin{adjustwidth}{-2cm}{-2cm}
\noindent\textbf{\footnotesize Texte grec établi --- Folio 27 (\texttt{f. 175v}, vv. 1082–1127 [1/2]) \hfill ΕΠΕΙΣΟΔΙΟΝ Ε' (Περσεύς)}\\[-2pt]
\noindent\rule{16.5cm}{0.5pt}\\[2pt]
{\fontsize{7.6pt}{8.7pt}\selectfont
\ligneEst{\textbf{Μνησίλοχος}}{\decalage{οἴμωζ᾽.†οἴμωζ᾽†.}ὀτότυζ᾽.}{}
\ligneEst{\textbf{Εὐριπίδης}}{\decalage{οἴμωζ᾽.†οἴμωζ᾽†.ὀτότυζ᾽.}†ὀτότυζ᾽†.}{}
\ligneEst{\textbf{Τοξότης}}{οὖτος τί λαλῖς;}{}
\ligneEst{\textbf{Εὐριπίδης}}{\decalage{οὖτος τί λαλῖς;}†οὖτος τί λαλῖς;†}{}
\ligneEst{\textbf{Τοξότης}}{πρυτάνεις καλέσω;}{}
\ligneEst{\textbf{Εὐριπίδης}}{\decalage{πρυτάνεις καλέσω;}†πρυτάνεις καλέσω;†}{}
\ligneEst{\textbf{Τοξότης}}{σί κακόν;}{1085}
\ligneEst{\textbf{Εὐριπίδης}}{\decalage{σί κακόν;}†σί κακόν;†}{}
\ligneEst{\textbf{Τοξότης}}{†πῶτε† τὸ πωνή;}{}
\ligneEst{\textbf{Εὐριπίδης}}{\decalage{†πῶτε† τὸ πωνή;}†πῶτε τὸ πωνή;†}{}
\ligneEst{\textbf{Τοξότης}}{σὺ λαλῖς;}{}
\ligneEst{\textbf{Εὐριπίδης}}{\decalage{σὺ λαλῖς;}†σὺ λαλῖς;†}{}
\ligneEst{\textbf{Τοξότης}}{κλαύσαι.}{}
\ligneEst{\textbf{Εὐριπίδης}}{\decalage{κλαύσαι.}†κλαύσαι.†}{}
\ligneEst{\textbf{Τοξότης}}{κἀκκάσκι μοι;}{}
\ligneEst{\textbf{Εὐριπίδης}}{\decalage{κἀκκάσκι μοι;}†κἀκκάσκι μοι;†}{}
\ligneEst{\textbf{Χορός}}{μὰ Δί᾽ ἀλλὰ γυνὴ πλησίον αὕτη.}{1090}
\ligneEst{\textbf{Εὐριπίδης}}{πλήξασ᾿ αὕτη.}{}
\ligneEst{\textbf{Τοξότης}}{ποῦ ᾿στ᾽ ἠ μιαρά;}{}
\ligneEst{\textbf{Χορός}}{\decalage{ποῦ ᾿στ᾽ ἠ μιαρά;}καὶ δὴ φεύγει.}{}
\ligneEst{\textbf{Τοξότης}}{ποῖ ποῖ πεύγεις;}{}
\ligneEst{\textbf{Εὐριπίδης}}{\decalage{ποῖ ποῖ πεύγεις;}ποῖ ποῖ πέτει.}{}
\ligneEst{\textbf{Τοξότης}}{οὐ καιρήσεις.}{}
\ligneEst{\textbf{Εὐριπίδης}}{\decalage{οὐ καιρήσεις.}οὐχ αἱρήσεις.}{}
\ligneEst{\textbf{Τοξότης}}{ἔτι γὰρ γρύζεις;}{1095}
\ligneEst{\textbf{Εὐριπίδης}}{\decalage{ἔτι γὰρ γρύζεις;}†ἔτι γὰρ γρύζεις;†}{}
\ligneEst{\textbf{Τοξότης}}{λαβὲ τὴ μιαρά.}{}
\ligneEst{\textbf{Εὐριπίδης}}{\decalage{λαβὲ τὴ μιαρά.}†λαβὲ τὴ μιαρά.†}{}
\ligneEst{\textbf{Τοξότης}}{λάλο καὶ κατάρατο γύναικο.}{}
\ligneEst{\textbf{Εὐριπίδης} {\tiny\textit{trim. iamb.}}}{ὦ θεοὶ τίν᾽ εἰς γῆν βαρβάρων ἀφίγμεθα}{}
\ligneEst{}{ταχεῖ πεδίλῳ; διὰ μέσου γὰρ αἰθέρος}{}
\ligneEst{}{τέμνων κέλευθον πόδα τίθημ᾽ ὑπόπτερον}{1100}
\ligneEst{}{Περσεὺς πρὸς Ἄργος ναυστολῶν τὸ Γοργόνος}{}
}
\vspace{2pt}
\noindent\rule{16.5cm}{0.6pt}\\[2pt]
\noindent\textbf{\small Commentaire philologique, métrique et scénique (vv. 1082–1127 --- partie 1/4)}\\[2pt]
\begin{adjustbox}{max totalheight=9.8cm, max width=16.5cm, keepaspectratio}%
\parbox[t]{16.5cm}{
\setlength{\columnsep}{9pt}%
\begin{multicols}{2}
\fontsize{7.6pt}{8.8pt}\selectfont
\setlength{\parindent}{0pt}%
\setlength{\parskip}{1.2pt}%
\noindent\textbf{1082} \textbf{ὀτότυζ᾽ --- ὀτότυζ᾽.} Signalons à l'attention du lecteur -- que nous espérons exégète à ses heures -- qu'Aristophane invente le gentilé Ὀτοτύξιοι au vers 1043 des \italique{Oiseaux} par jeu de consonance avec Ὀλοφύξιοι (v. 1042), \italique{habitants d'Olophyxos}. Y a-t-il ici un jeu analogue ?
\noindent\textbf{1083} \textbf{λαλῖς.} Scholie à propos de la graphie :
\citComm{χωρὶς τοῦ ε γράφεται· ὁ γὰρ Σκύθης βαρβαρίζει.}
\noindent\textbf{1085} \textbf{σί κακόν;} Commentaire du scholiaste :
\citComm{ἀντὶ τοῦ `σοὶ κακόν'.}
\noindent\textbf{1086} \textbf{πρωτε το πωνη;} Le scholiaste interprète ainsi le jargon scythe :
\citComm{πόθεν ἡ φωνή; θαυμάζει τὴν ἠχώ.}
\noindent On lui objectera que le \italique{Ravennas} comporte `πρωτε'. Et quand bien même on adopterait la correction `πωτε' (qu'une main postérieure a corrigé sur le manuscrit -- peut-être d'après la scholie elle-même), qu'il n'en resterait pas moins qu'à aucun autre endroit le Scythe ne confond un omicron ο avec un oméga ω\marqueur{1}.
\noindent\textbf{[Notes ad 1086]} \textbf{(1)} Voir la seconde note du commentaire au vers 1080 à propos du nombre déjà bien suffisant de fautes commises par le Scythe.
\noindent\textbf{1089} \textbf{κακκασκι μοι;} $=$ καὶ ἐγχάσκει μοι. Cf. \italique{Guêpes}, vers 721 et le scholiaste : καταγελᾷς μοι.
\noindent\textbf{1090} \textbf{μὰ Δί᾽ ἀλλὰ γυνὴ πλησίον αὕτη.} C'est le chœur qui donne cette indication à l'archer désemparé, comme le suggérait timidement le scholiaste :
\citComm{τοῦτο ὁ χορὸς ἢ ὁ γέρων.}
\noindent Aucun éditeur cependant n'a suivi la première hypothèse du scholiaste, quoique la seule vraisemblable -- aussi agacé soit-il, Mnésiloque ne peut en aucun cas dénoncer son complice, son dernier espoir de survie.
\end{multicols}
}%
\end{adjustbox}%
\vfill
\end{adjustwidth}
\newpage
% --- PAGE DE DROITE : Traduction française + Commentaire (2/4) ---
\begin{adjustwidth}{-2cm}{-2cm}
\noindent\textbf{\footnotesize Traduction française --- Folio 27 (\texttt{f. 175v}, vv. 1082–1127 [1/2]) \hfill ΕΠΕΙΣΟΔΙΟΝ Ε' (Περσεύς)}\\[-2pt]
\noindent\rule{16.5cm}{0.5pt}\\[2pt]
{\fontsize{7.6pt}{8.7pt}\selectfont
\ligneEst{\textbf{Mnésiloque}}{\francais{\espace{2.3cm}Foin de toi.}}{}
\ligneEst{\textbf{Euripide}}{\francais{\espace{3.5cm}Foin de toi.}}{}
\ligneEst{\textbf{Archer}}{\francais{Holà ! qu'as-tu à bavarder ?}}{}
\ligneEst{\textbf{Euripide}}{\francais{\espace{2.6cm}Holà ! qu'as-tu à bavarder ?}}{}
\ligneEst{\textbf{Archer}}{\francais{J'appellerai les prytanes.}}{}
\ligneEst{\textbf{Euripide}}{\francais{\espace{2.7cm}J'appellerai les prytanes.}}{}
\ligneEst{\textbf{Archer}}{\francais{Qu'est-ce qui arrive ?}}{1085}
\ligneEst{\textbf{Euripide}}{\francais{\espace{1.6cm}Qu'est-ce qui arrive ?}}{}
\ligneEst{\textbf{Archer}}{\francais{D'où part cette voix ?}}{}
\ligneEst{\textbf{Euripide}}{\francais{\espace{2.4cm}D'où part cette voix ?}}{}
\ligneEst{\textbf{Archer}}{\francais{Est-ce toi qui bavardes?}}{}
\ligneEst{\textbf{Euripide}}{\francais{\espace{1.8cm}Est-ce toi qui bavardes?}}{}
\ligneEst{\textbf{Archer}}{\francais{Gare à toi.}}{}
\ligneEst{\textbf{Euripide}}{\francais{\espace{1.4cm}Gare à toi.}}{}
\ligneEst{\textbf{Archer}}{\francais{Tu te gausses de moi.}}{}
\ligneEst{\textbf{Euripide}}{\francais{\espace{2.2cm}Tu te gausses de moi.}}{}
\ligneEst{\textbf{Chœur}}{\francais{Non, par Zeus : c'est cette femme, là tout près.}}{1090}
\ligneEst{\textbf{Euripide}}{\francais{Là tout près.}}{}
\ligneEst{\textbf{Archer}}{\francais{Où est la canaille ?}}{}
\ligneEst{\textbf{Chœur}}{\francais{\espace{2.2cm}La voilà qui se sauve.}}{}
\ligneEst{\textbf{Archer}}{\francais{Où, où te sauves-tu ?}}{}
\ligneEst{\textbf{Euripide}}{\francais{\decalage{Où, où te sauves-tu ?}Où, où voles-tu ?}}{}
\ligneEst{\textbf{Archer}}{\francais{Il t'en cuira.}}{}
\ligneEst{\textbf{Euripide}}{\francais{\espace{2.1cm}Tu ne m'auras.}}{}
\ligneEst{\textbf{Archer}}{\francais{Tu maronnes encore ?}}{1095}
\ligneEst{\textbf{Euripide}}{\francais{\espace{2.5cm}Tu maronnes encore ?}}{}
\ligneEst{\textbf{Archer}}{\francais{Arrête, la canaille.}}{}
\ligneEst{\textbf{Euripide}}{\francais{\espace{2.3cm}Arrête, la canaille.}}{}
\ligneEst{\textbf{Archer}}{\francais{La bavarde, la maudite femme.}}{}
\ligneEst{\textbf{Euripide}}{\francais{(En Persée :) « En quel pays barbare, ô dieux, ces talonnières}}{}
\ligneEst{}{\francais{ont-elles amené Persée ? Au milieu de l'éther}}{}
\ligneEst{}{\francais{me frayant un chemin, prestement je fends l'air d'un pied ailé,}}{1100}
\ligneEst{}{\francais{portant à la cité argienne la tête de}}{}
}
\vspace{2pt}
\noindent\rule{16.5cm}{0.6pt}\\[2pt]
\noindent\textbf{\small Commentaire philologique, métrique et scénique (vv. 1082–1127 --- partie 2/4)}\\[2pt]
\begin{adjustbox}{max totalheight=9.8cm, max width=16.5cm, keepaspectratio}%
\parbox[t]{16.5cm}{
\setlength{\columnsep}{9pt}%
\begin{multicols}{2}
\fontsize{7.6pt}{8.8pt}\selectfont
\setlength{\parindent}{0pt}%
\setlength{\parskip}{1.2pt}%
\noindent\textbf{1091} \textbf{πλησίον αὕτη.} Après un long calvaire parsemé de \latin{cruces desperationis}, il nous semble enfin saisir le sel ἐπικοκκάστρια de la réplique d'Écho, fournissant par la même occasion un indice sur l'action.\\ Sur l'indication du chœur, l'archer Scythe est en train d'inspecter \italique{à proximité} (πλησίον) de Mnésiloque. Il se rapproche effectivement d'Euripide (ce dernier est toujours camouflé derrière son manteau).
\noindent Quand l'archer arrive à portée de bras, Euripide lui assène alors un coup : πλήξασ᾿ αὕτη.\\ Ce coup explique la fureur du Scythe, qui change de langage et passe à l'offensive : ποῦ ᾿στ᾽ ἠ μιαρά;
\noindent\textbf{1093} \textbf{ποῖ ποῖ φεύγεις;} Il y a longtemps que les scribes chargés de recopier cette scène n'en saisissaient plus le comique ; ne nous étonnons donc pas qu'elle soit la plus altérée de toute la comédie : il manque plusieurs répliques (cf. le vers 1094 ; où la lacune est certaine et peut-être comblée grâce à une ancienne scholie). C'est le cas aussi ici.\\ Euripide s'est approché de Mnésiloque pour le délivrer. Sur les indications du chœur (μὰ Δί᾽ ἀλλὰ γυνὴ πλησίον αὕτη), le Scythe a fini par repérer Euripide, et ce malgré son camouflage. C'est ce coup qu'il est bon de partir ; et Euripide de déloger sans trompette (καὶ δὴ φεύγει) mais non sans abandonner sa verve moqueuse. \\ Remettons-nous bien la scène : Euripide incarne la nymphe ailée Écho -- sa fuite est donc aérienne ; et sur scène, l'acteur est tiré dans les airs à l'aide de la μηχανή. La réplique ἐπικοκκάστρια d'Écho fait vraisemblablement allusion à cette \italique{envolée}. Or pour la circonstance, nous ne sommes pas tout-à-fait démunis puisque nous possédons une comédie entière, les \italique{Oiseaux}, avec des jeux théâtraux similaires.\\ En \italique{Oiseaux} 226, lorsque la Huppe se dispose à chanter, elle entonne : ἐποποποῖ ποποῖ, ποποποποῖ ποποῖ.\\ De même, lorsque le chœur au vers 310 répond à l'appel de la Huppe, c'est au cri de : ποποποποποπο ποῦ μ' ὃς ἐκάλεσε;\\ C'est ainsi qu'Aristophane transcrivait les onomatopées d'oiseaux\marqueur{1} ; on comprend mieux pourquoi lorqu'Iris descend du ciel au moyen de la μηχανή, Pisthétaire lui lance ces paroles ailées (vers 1199) : αὕτη σύ, ποῖ ποῖ ποῖ πέτει;\\ Il en va de même dans notre passage, alors que le Scythe poursuit Écho en train de s'envoler, il lui jette à la figure des cris d'oiseaux : ποῖ ποῖ πεύγεις; -- qui fait entendre ποῖ ποῖ πέτει -- qu'on rétablit donc comme réplique ἐπικοκκάστρια manquante d'Écho.
\end{multicols}
}%
\end{adjustbox}%
\vfill
\end{adjustwidth}
\newpage

% --- FOLIO 27 (f. 175v) : SPREAD 2.2 (Texte grec établi + Traduction française + Commentaire 3/4 & 4/4) ---
\clearpage\ifodd\value{page}\hbox{}\newpage\fi
\enTeteG{ΘΕΣΜΟΦΟΡΙΑΖΟΥΣΑΙ --- Texte grec établi (f. 175v, vv. 1082–1127 [2/2])}
\enTeteD{LES THESMOPHORIES --- Traduction française (f. 175v, vv. 1082–1127 [2/2])}
\begin{adjustwidth}{-2cm}{-2cm}
\noindent\textbf{\footnotesize Texte grec établi --- Folio 27 (\texttt{f. 175v}, vv. 1082–1127 [2/2]) \hfill ΕΠΕΙΣΟΔΙΟΝ Ε' (Περσεύς)}\\[-2pt]
\noindent\rule{16.5cm}{0.5pt}\\[2pt]
{\fontsize{7.6pt}{8.7pt}\selectfont
\ligneEst{}{κάρα κομίζων.}{}
\ligneEst{\textbf{Τοξότης}}{\decalage{κάρα κομίζων.}τί λέγι; τὴ Γόργος πέρι}{}
\ligneEst{}{τὸ γραμματέο σὺ τὴ κεπαλή;}{}
\ligneEst{\textbf{Εὐριπίδης}}{\decalage{τὸ γραμματέο σὺ τὴ κεπαλή;}τὴν Γοργόνος}{}
\ligneEst{}{ἔγωγε φημί.}{}
\ligneEst{\textbf{Τοξότης}}{\decalage{ἔγωγε φημί.}Γόργο τοι κἀγὼ λέγι.}{}
\ligneEst{\textbf{Εὐριπίδης}}{ἔα· τίν᾽ ὄχθον τόνδ᾽ ὁρῶ καὶ παρθένον}{1105}
\ligneEst{}{θεαῖς ὁμοίαν ναῦν ὅπως ὡρμισμένην;}{}
\ligneEst{\textbf{Μνησίλοχος}}{ὦ ξένε κατοίκτιρόν με τὴν παναθλίαν,}{}
\ligneEst{}{λῦσόν με δεσμῶν.}{}
\ligneEst{\textbf{Τοξότης}}{\decalage{λῦσόν με δεσμῶν.}οὐκὶ μὶ λαλῆσι σύ;}{}
\ligneEst{}{κατάρατο τολμᾷς ἀποτανουμένη λαλῖς;}{}
\ligneEst{\textbf{Εὐριπίδης}}{ὦ παρθέν᾽ οἰκτίρω σὲ κρεμαμένην ὁρῶν.}{1110}
\ligneEst{\textbf{Τοξότης}}{οὐ παρτέν᾽ ἐστίν, ἀλλ᾽ ἀμαρτωλὴ γέρων}{}
\ligneEst{}{καὶ κλέπτο καὶ πανοῦργο.}{}
\ligneEst{\textbf{Εὐριπίδης}}{\decalage{καὶ κλέπτο καὶ πανοῦργο.}ληρεῖς ὦ Σκύθα.}{}
\ligneEst{}{αὕτη γάρ ἐστιν Ἀνδρομέδα παῖς Κηφέως.}{}
\ligneEst{\textbf{Τοξότης}}{σκέψαι τὸ κύστο· μή τι μικκὸν παίνεται;}{}
\ligneEst{\textbf{Εὐριπίδης}}{φέρε δεῦρό μοι τὴν χεῖρ᾽, ἵν᾽ ἅψωμαι κόρης·}{1115}
\ligneEst{}{φέρε Σκύθ᾽· ἀνθρώποισι γὰρ νοσήματα}{}
\ligneEst{}{ἅπασίν ἐστιν· ἐμὲ δὲ καὐτὸν τῆς κόρης}{}
\ligneEst{}{ταύτης ἔρως εἴληφεν.}{}
\ligneEst{\textbf{Τοξότης}}{\decalage{ταύτης ἔρως εἴληφεν.}οὐ ζηλῶσί σε·}{}
\ligneEst{}{ἀτὰρ εἰ τὸ πρωκτὸ δεῦρο περιεστραμμένον,}{}
\ligneEst{}{οὐκ ἐπτόνησά σ᾽ αὐτὸ πυγίζεις ἄγων.}{1120}
\ligneEst{\textbf{Εὐριπίδης}}{τί δ᾽ οὐκ ἐᾷς λύσαντά μ᾽ αὐτὴν ὦ Σκύθα}{}
\ligneEst{}{πεσεῖν ἐς εὐνὴν καὶ γαμήλιον λέχος;}{}
\ligneEst{\textbf{Τοξότης}}{εἰ σπόδρ᾽ ἐπιτυμεῖς τὴ γέροντο πύγισο,}{}
\ligneEst{}{τὴ σανίδο τρήσας ἐξόπιστο πρώκτισον.}{}
\ligneEst{\textbf{Εὐριπίδης}}{μὰ Δί᾽ ἀλλὰ λύσω δεσμά.}{1125}
\ligneEst{\textbf{Τοξότης}}{\decalage{μὰ Δί᾽ ἀλλὰ λύσω δεσμά.}μαστιγῶ σ᾽ ἄρα.}{}
\ligneEst{\textbf{Εὐριπίδης}}{καὶ μὴν ποιήσω τοῦτο.}{}
\ligneEst{\textbf{Τοξότης}}{\decalage{καὶ μὴν ποιήσω τοῦτο.}τὸ κεπαλή σ᾽ ἄρα}{}
\ligneEst{}{τὸ ξιπομάκαιραν ἀποκεκόψι τουτοΐ.}{1127}
}
\vspace{2pt}
\noindent\rule{16.5cm}{0.6pt}\\[2pt]
\noindent\textbf{\small Commentaire philologique, métrique et scénique (vv. 1082–1127 --- partie 3/4)}\\[2pt]
\begin{adjustbox}{max totalheight=9.8cm, max width=16.5cm, keepaspectratio}%
\parbox[t]{16.5cm}{
\setlength{\columnsep}{9pt}%
\begin{multicols}{2}
\fontsize{7.6pt}{8.8pt}\selectfont
\setlength{\parindent}{0pt}%
\setlength{\parskip}{1.2pt}%
\noindent\textbf{[Notes ad 1093]} \textbf{(1)} Pour d'autres exemples de telles onomatopées chez Aristophane, voir l'article de Loredana Di Virgilio «~Some thoughts on humor and music in Aristophanes' plays~», in \italique{Mediterranean World}, 24 (2019), pp. 36 sqq.
\noindent\textbf{1094} \textbf{<οὐ καιρήσεις.> --- οὐκ᾽ αἰρήσεις.} Le texte manquant a été rétabli à partir de la scholie :
\citComm{ἀντὶ τοῦ οὐ χαιρήσεις.}
\noindent \italique{Confer} la note \italique{supra} au vers 1056.
\noindent\textbf{1098-102} \textbf{ὦ θεοὶ ... κάρα κομίζων.} La scholie à propos de la tirade de Persée\marqueur{1} :
\citComm{εἰς Περσέα ἐξ Ἀνδρομέδας τρία τὰ πρῶτα. καὶ λοιπὸν (i.e. vv. 1101-2) ἐπέζευξεν τὰ ἑξῆς.}
\noindent\textbf{[Notes ad 1098-102]} \textbf{(1)} Fr. 9 de l'\italique{Andromède} dans l'édition Jouan-van Looy $=$ fr. 124 Kannicht.
\noindent\textbf{1102-3} \textbf{τί λέγι; τὴ Γόργος πέρι / τὸ γραμματέο σὺ τὴ κεπαλή;} $=$ τί λέγεις; τοῦ Γόργου φέρεις / τοῦ γραμματέως σὺ τὴν κεφαλήν;\\ Et la scholie (plus que douteuse) :
\citComm{ὁ δὲ Γόργος γραμματεὺς, ἀλλὰ καὶ βάρβαρος.}
\noindent Il n'y a en grec qu'une différence d'accent entre Γοργών (la Gorgone) et Γόργων (Gorgon).
\noindent\textbf{1105} \textbf{ἔα· τίν᾽ ὄχθον τόνδ᾽ ὁρῶ καὶ παρθένον.} On connaît le texte original d'Euripide grâce à une scholie\marqueur{1} :
\citComm{πάλιν ἐξ Ἀνδρομέδας·\\
ἔα, τίν᾽ ὄχθον τόνδ᾽ ὁρῶ περίρρυτον\\
ἀφρῷ θαλάσσης παρθένου τ᾽ εἰκώ τίνα\\
$[$ ἐξ αὐτομόρφων λαΐνων τυκισμάτων,
σοφῆς ἄγαλμα χειρός; $]$ }
\end{multicols}
}%
\end{adjustbox}%
\vfill
\end{adjustwidth}
\newpage
% --- PAGE DE DROITE : Traduction française + Commentaire (4/4) ---
\begin{adjustwidth}{-2cm}{-2cm}
\noindent\textbf{\footnotesize Traduction française --- Folio 27 (\texttt{f. 175v}, vv. 1082–1127 [2/2]) \hfill ΕΠΕΙΣΟΔΙΟΝ Ε' (Περσεύς)}\\[-2pt]
\noindent\rule{16.5cm}{0.5pt}\\[2pt]
{\fontsize{7.6pt}{8.7pt}\selectfont
\ligneEst{}{\francais{Gorgone. »}}{}
\ligneEst{\textbf{Archer}}{\francais{\espace{2.2cm}Que dis-tu? Tu portes la tête de Gorgon}}{}
\ligneEst{}{\francais{le scribe ?}}{}
\ligneEst{\textbf{Euripide}}{\francais{\espace{4.4cm}De Gorgone, ai-je dit.}}{}
\ligneEst{}{\francais{ἔγωγε φημί.}}{}
\ligneEst{\textbf{Archer}}{\francais{\espace{2cm}Gorgon, parfaitement, je dis de même.}}{}
\ligneEst{\textbf{Euripide}}{\francais{« Ciel! Quelle est cette vierge aux déesses pareille amarrée à ce roc,}}{1105}
\ligneEst{}{\francais{comme un navire au port ? »}}{}
\ligneEst{\textbf{Mnésiloque}}{\francais{« Étranger, prends pitié de moi, la désolée.}}{}
\ligneEst{}{\francais{Détache mes liens. »}}{}
\ligneEst{\textbf{Archer}}{\francais{\espace{2.8cm}Veux-tu ne pas bavarder ?}}{}
\ligneEst{}{\francais{Maudit, oses-tu bien ? Tu vas mourir et tu bavardes.}}{}
\ligneEst{\textbf{Euripide}}{\francais{« Ô vierge, une pitié me prend quand je te vois suspendue à ce roc. »}}{1110}
\ligneEst{\textbf{Archer}}{\francais{Ce n'est pas une vierge, c'est un vieux coquin,}}{}
\ligneEst{}{\francais{et un voleur, et un fripon.}}{}
\ligneEst{\textbf{Euripide}}{\francais{\espace{3.9cm}Tu plaisantes, ô Scythe.}}{}
\ligneEst{}{\francais{C'est Andromède, fille de Cépheus.}}{}
\ligneEst{\textbf{Archer}}{\francais{Regarde son cas : te paraît-il petit ?}}{}
\ligneEst{\textbf{Euripide}}{\francais{Donne-moi la main, que j'atteigne à la jeune fille.}}{1115}
\ligneEst{}{\francais{Voyons, Scythe « chaque homme a ses faiblesses » :}}{}
\ligneEst{}{\francais{la mienne est d'être féru}}{}
\ligneEst{}{\francais{d'amour pour cette jeune fille.}}{}
\ligneEst{\textbf{Archer}}{\francais{\espace{3.3cm}Je ne suis pas jaloux.}}{}
\ligneEst{}{\francais{S'il tournait les fesses par ici,}}{}
\ligneEst{}{\francais{je ne te refuserais pas d'en passer ta fantaisie.}}{1120}
\ligneEst{\textbf{Euripide}}{\francais{Que ne me permets-tu de la délier, ô Scythe}}{}
\ligneEst{}{\francais{« et de me laisser choir dans le lit nuptial »}}{}
\ligneEst{\textbf{Archer}}{\francais{Si tu as tant envie de jouir de ce vieillard,}}{}
\ligneEst{}{\francais{perce-moi l'ais par derrière et enfourne.}}{}
\ligneEst{\textbf{Euripide}}{\francais{Non, par Zeus, je briserai ses liens.}}{1125}
\ligneEst{\textbf{Archer}}{\francais{\espace{3.7cm}Alors je te fouette.}}{}
\ligneEst{\textbf{Euripide}}{\francais{N'importe, je vais le faire.}}{}
\ligneEst{\textbf{Archer}}{\francais{\espace{3.2cm}Alors je te coupe la tête}}{}
\ligneEst{}{\francais{avec ce sabre.}}{1127}
}
\vspace{2pt}
\noindent\rule{16.5cm}{0.6pt}\\[2pt]
\noindent\textbf{\small Commentaire philologique, métrique et scénique (vv. 1082–1127 --- partie 4/4)}\\[2pt]
\begin{adjustbox}{max totalheight=9.8cm, max width=16.5cm, keepaspectratio}%
\parbox[t]{16.5cm}{
\setlength{\columnsep}{9pt}%
\begin{multicols}{2}
\fontsize{8.4pt}{9.7pt}\selectfont
\setlength{\parindent}{0pt}%
\setlength{\parskip}{1.2pt}%
\noindent\textbf{1106} \textbf{ναῦν ὅπως ὡρμισμένην;} Quoique le scholiaste ne l'indique pas, le vers 1106 est probablement lui aussi repris de l'\italique{Andromède} d'Euripide. On trouve en tout cas la même image chez l'auteur d'\italique{Hercule Furieux} (vers 1094-1097), Hercule dont la situation était analogue à celle d'Andromède :
\citComm{τί δεσμοῖς ναῦς ὅπως ὡρμισμένος\\
νεανίαν θώρακα καὶ βραχίονα,\\
πρὸς ἡμιθραύστῳ λαΐνῳ τυκίσματι\\
ἧμαι, νεκροῖσι γείτονας θάκους ἔχων;}
\noindent\textbf{[Notes ad 1105]} \textbf{(1)} Dixième fragment de l'\italique{Andromède} dans l'édition Jouan-van Looy $=$ fr. 125, 1-2 Kannicht. La scholie rapporte les seuls deux premiers vers ; le fragment a pu être complété en le recoupant avec d'autres témoignages.
\noindent\textbf{1114} \textbf{σκέψαι τὸ σκύτο·} Fait remarquable : l'erreur du scribe, copiant τὸ σκύτο en place de τὸ κύστο ($=$τὸν κύσθον), est la même qu'au vers 557 : métathèse du `s' (cf. le commentaire à ce vers).
\noindent Le scholiaste devait avoir devant les yeux la bonne leçon qui précise :
\citComm{δείκνυσιν αὐτῷ τὸ αἰδοῖον.}
\noindent Second fait remarquable : en première rédaction, le \italique{Ravennas} faisait apparaître un `σ' supplémentaire au début d'un mot du vers 1187 : τὸ \textbf{σ}ποστίον à la de τὸ πόστιον ($=$ τὸ πόσθιον)...  on constate qu'il s'agit là-encore d'un organe sexuel -- mâle cette fois.
\noindent\textbf{1127} \textbf{τὸ ξιπομάκαιραν.} La harpé de Persée. (Note de Willems.)
\end{multicols}
}%
\end{adjustbox}%
\vfill
\end{adjustwidth}
\newpage
\clearpage\ifodd\value{page}\hbox{}\newpage\fi
\enTeteG{ΧΟΡΙΚΟΝ / ΕΞΟΔΟΣ --- Codex Ravennas 429, f. 176r (Classense 178r) --- vv. 1128–1168}
\begin{adjustwidth}{-1.5cm}{-3.0cm}
\begin{adjustbox}{max totalheight=21.0cm, max width=16.5cm, keepaspectratio}
\begin{minipage}{16.5cm}
{\fontsize{7.2pt}{8.5pt}\selectfont
\noindent\textbf{Folio 28 --- Codex Ravennas 429, \texttt{f. 176r} (\texttt{Classense 178r}) : ΧΟΡΙΚΟΝ / ΕΞΟΔΟΣ, vv. 1128–1168}\hfill \textit{Témoins :} $\text{R}^1, \Rpc, \Rdeux$ (\Mon)\par\vspace{2pt}
\hrule height 0.4pt\vspace{2.5pt}
}
{\fontsize{8.1pt}{9.5pt}\selectfont
\ligneRav{1128}{\auDessus{εὐ}{ρ}}{αἰαῖ· τί δράσω; πρὸς τίνας στρεφθῶ λόγους;}{}
\ligneRav{1129}{}{ἀλλ᾽ οὐκ \difference{ἂν ἐνδέξαιτο} βάρβαρος φύσις·}{\textbf{1129} ἂν δέξαιτο R : ἂν ἀνδέξαιτο Ellebodius : \lecon{ἂν ἐνδέξαιτο} Küster}
\ligneRav{1130}{}{σκαιοῖσι γάρ τοι καινὰ προσφέρων σοφὰ}{}
\ligneRav{}{}{μάτην ἀναλίσκοις ἄν, ἀλλ᾽ ἄλλην τινὰ}{}
\ligneRav{1132}{}{\difference{τούτῳ} πρέπουσαν μηχανὴν προσοιστέον·}{\textbf{1132} τοῦτο R : \lecon{τούτῳ} Zanetti}
\ligneRav{1133}{τοξ/}{μιαρὸς ἀλώπηξ, \difference{οἶον} \difference{ἐπιτήκιζί} μοι·}{\textbf{1133} οἷον R : \lecon{οἶον} Brunck \espApp ἐπιτηκίζει R : \lecon{ἐπιτήκιζί} anonymus Parisinus}
\ligneRav{}{μ\auDessus{ν}{Η}}{μέμνησο Περσεῦ μ᾽ ὡς καταλείπεις ἀθλίαν·}{}
\ligneRav{1135}{τοξ/}{ἔτι γὰρ σὺ \difference{τῇ μάστιγ᾽ ἂν} (τῆ μαστ\inRas{η}{ι}γᾶν) ἐπεθύμεις λαβεῖν;}{\textbf{1135} τῆ μαστιγᾶν (ex μαστηγαν) \Rpc \espApp λαβεῖν· R : \lecon{λαβεῖν;} Scaliger}
\ligneRav{}{χορ/}{\hspace{1.1em}Παλλάδα τὴν φιλόχορον ἐμοὶ}{}
\ligneRav{1138}{}{\hspace{1.1em}δεῦρο καλεῖν νόμος \difference{εἰς} χορόν,}{\textbf{1138} ἐς R : \lecon{εἰς} von Velsen}
\ligneRav{1139}{}{\hspace{1.1em}παρθένον ἄζυγα \difference{κούρην},}{\textbf{1139} κόρην R : \lecon{κούρην} Bothe}
\ligneRav{1140}{}{\hspace{1.1em}ἣ πόλιν ἡμετέραν ἔχει}{}
\ligneRav{}{}{\hspace{1.1em}καὶ κράτος φανερὸν μόνη}{}
\ligneRav{}{}{\hspace{1.1em}κλῃδοῦχός τε καλεῖται·}{}
\ligneRav{}{}{\hspace{1.1em}φάνηθ᾽ ὦ τυράννους στυγοῦσ᾽ ὥσπερ εἰκός·}{}
\ligneRav{1145}{}{\hspace{1.1em}δῆμός τοί σε καλεῖ γυναι-}{}
\ligneRav{}{}{\hspace{1.1em} κῶν· ἔχουσα δέ μοι μόλοις}{}
\ligneRav{}{}{\hspace{1.1em}εἰρήνην φιλέορτον·}{}
\ligneRav{1148}{}{\hspace{1.1em}\difference{ἥκετέ τ᾽} εὔφρονες ἵλαοι,}{\textbf{1148} ἥκετ᾽ R : \lecon{ἥκετέ τ᾽} Enger : ἥκετε δ᾽ Hermann}
\ligneRav{}{}{\hspace{1.1em}πότνιαι, ἄλσος ἐς ὑμέτερον,}{}
\ligneRav{1150}{}{\hspace{1.1em}οὗ δὴ ἀνδράσιν οὐ \difference{θεμιτεύετον} εἰσορᾶν}{\textbf{1150-1} \lecon{οὗ δὴ ἀνδράσιν} R : οὗ δῆτ' ἀνδράσιν Burges : ἄνδρας ἵν(α) Hermann \espApp θεμιτὸν R : θέμις Burges : θεμιτ' Bothe : \lecon{θεμιτεύετον} Bibi}
\ligneRav{}{}{\hspace{1.1em}ὄργια σεμνὰ θεαῖν ἵνα λαμπάσι·}{}
\ligneRav{}{}{\hspace{1.1em}φαίνετον ἄμβροτον ὄψιν·}{}
\ligneRav{1155}{}{\hspace{1.1em}μόλετον, ἔλθετον, ἀντόμεθ᾽ ὦ}{}
\ligneRav{1156}{}{\hspace{1.1em}Θεσμοφόρω \difference{πολυποτνία},}{\textbf{1156} πολὺ πότνια R : \lecon{πολυποτνία} Hermann}
\ligneRav{}{}{\hspace{1.1em}εἰ καὶ πρότερόν ποτ᾽ ἐπηκόω ἤλθετον}{}
\ligneRav{1159}{}{\hspace{1.1em}νῦν ἀφίκεσθον †ἱκετεύομεν ἐνθάδ᾽ ἡμῖν†·}{\textbf{1159} \lecon{ἐνθάδ᾽ ἡμῖν} R : ἐνθάδε χἠμῖν Hermann}
\ligneRav{1160}{\auDessus{εὐ}{ρ}}{γυναῖκες εἰ βούλεσθε τὸν λοιπόν χρόνον}{}
\ligneRav{}{}{σπονδὰς ποιήσασθαι πρὸς ἐμέ, νυνὶ πάρα,}{}
\ligneRav{}{}{ἐφ᾽ ᾧτ᾽ ἀκοῦσαι μηδὲν ὑπ᾽ ἐμοῦ μηδαμὰ}{}
\ligneRav{}{}{κακὸν τὸ λοιπόν. ταῦτ᾽ ἐπικηρυκεύομαι·}{}
\ligneRav{}{χορ/}{χρείᾳ δὲ ποίᾳ τόνδ᾽ ἐπεσφέρεις λόγον;}{}
\ligneRav{1165}{\auDessus{εὐ}{ρ}}{ὅδ᾽ ἐστὶν \difference{οὑν} τῇ σανίδι κηδεστὴς ἐμός·}{\textbf{1165} οὖν R : \lecon{οὑν} Biset}
\ligneRav{1166}{}{ἢν οὖν κομίσωμαι τοῦτον, οὐδὲν μή ποτε}{\textbf{1166} \lecon{μή} R : μου Bentley}
\ligneRav{1167}{}{κακῶς \difference{ἀκούσητ᾽}· ἢν δὲ μὴ \difference{πίθησθέ} μοι,}{\textbf{1167} ἀκούσαιτ' R : \lecon{ἀκούσητ᾽} Elmsley (ad \italique{Ach.} v.295) \espApp πείθησθέ R : \lecon{πίθησθέ} Hirschig (ad \italique{Vesp} v.155)}
\ligneRav{}{}{ἃ νῦν ὑποικουρεῖτε τοῖσιν ἀνδράσιν}{}
}
\vspace{3pt}\noindent\rule{4.5cm}{0.4pt}\par\vspace{1.5pt}
{\fontsize{7.0pt}{8.3pt}\selectfont
\noindent\textbf{Loci paralleli (vv. 1128–1168) ---} \textbf{1139} (ἄζυγα) \Phryn 100.3\par
}
\vspace{2.5pt}\noindent\rule{4.5cm}{0.4pt}\par\vspace{1.5pt}
{\fontsize{6.9pt}{8.2pt}\selectfont
\noindent\textbf{Suite (en commentaire) de l'apparat critique \& notes paléographiques (\texttt{f. 176r}) ---} \textbf{Dittographie anticipatrice exponctuée après le v. 1143 :} Entre les vv. 1143 et 1144, \Rac{} avait écrit par anticipation déformée du v. 1144 (στυγοῦσ' ὥσπερ τὸν τύραννον) les mots \lecon{στυγνὰς ὡσσε} ; \Rpc{} les a supprimés (\italique{del.}), mais le copiste de \Mon{} les a néanmoins transcrits sous la forme \lecon{στυγνὰς ὡς ἐ}. --- \textbf{Glose interne et conjecture métrique de \textbf{Bibi} au v. 1151 :} R comporte la glose interne ὅπου (expliquant un ἵνα ou οὗ). Là où R porte \lecon{οὗ δὴ ἀνδράσιν οὐ θεμιτὸν} (déficient de deux syllabes dans la responsion lyrique), \textbf{Louis Hauseux (\textbf{Bibi}) conserve intact οὗ δὴ ἀνδράσιν de R et restitue le duel verbal οὐ θεμιτεύετον} (« où toutes deux [les Déesses] interdisent par leur loi sacrée aux hommes de voir leurs orgies... », haplographie de la finale -εύετον en -όν). --- \textbf{vv. 1150–1151 :} \lecon{οὗ δὴ ἀνδράσιν} R (avec glose interne ὅπου ; οὗ δῆτ' ἀνδράσιν Burges ; ἄνδρας ἵν(α) Hermann) \espApp θεμιτὸν R (θέμις Burges ; θεμιτ' Bothe) : \lecon{θεμιτεύετον} \lecon{Bibi}. --- \textbf{v. 1153 :} λαμπάσι R : \lecon{λαμπάσι.} \lecon{Bibi}.\par
}
\end{minipage}
\end{adjustbox}
\end{adjustwidth}
\newpage
\enTeteD{Codex Ravennas 429 (Biblioteca Classense) --- Fac-similé du f. 176r (vv. 1128–1168)}
\begin{adjustwidth}{-1.5cm}{-3.0cm}
\begin{center}
\includegraphics[width=16.5cm,height=20.5cm,keepaspectratio]{/workspaces/Thesmophories/dossier_ravennas_thesmophories/folios_individuels_ff162v-177r/28_f176r_classense178r_vv1128-1168.jpg}\par
\vspace{2pt}
{\scriptsize\textbf{Codex Ravennas 429, f. 176r (Classense 178r)} --- ΧΟΡΙΚΟΝ / ΕΞΟΔΟΣ, \textbf{vv. 1128–1168} ($\text{R}^1, \Rpc, \Rdeux$ (\Mon)).}
\end{center}
\end{adjustwidth}
\newpage
\clearpage\ifodd\value{page}\hbox{}\newpage\fi% --- FOLIO 28 (f. 176r) : SPREAD 2.1 (Texte grec établi + Traduction française + Commentaire 1/4 & 2/4) ---
\clearpage\ifodd\value{page}\hbox{}\newpage\fi
\enTeteG{ΘΕΣΜΟΦΟΡΙΑΖΟΥΣΑΙ --- Texte grec établi (f. 176r, vv. 1128–1168 [1/2])}
\enTeteD{LES THESMOPHORIES --- Traduction française (f. 176r, vv. 1128–1168 [1/2])}
\begin{adjustwidth}{-2cm}{-2cm}
\noindent\textbf{\footnotesize Texte grec établi --- Folio 28 (\texttt{f. 176r}, vv. 1128–1168 [1/2]) \hfill ΧΟΡΙΚΟΝ / ΕΞΟΔΟΣ}\\[-2pt]
\noindent\rule{16.5cm}{0.5pt}\\[2pt]
{\fontsize{8.4pt}{9.7pt}\selectfont
\ligneEst{\textbf{Εὐριπίδης}}{αἰαῖ· τί δράσω; πρὸς τίνας στρεφθῶ λόγους;}{1128}
\ligneEst{}{ἀλλ᾽ οὐκ ἂν ἐνδέξαιτο βάρβαρος φύσις.}{}
\ligneEst{}{σκαιοῖσι γάρ τοι καινὰ προσφέρων σοφὰ}{1130}
\ligneEst{}{μάτην ἀναλίσκοις ἄν, ἀλλ᾽ ἄλλην τινὰ}{}
\ligneEst{}{τούτῳ πρέπουσαν μηχανὴν προσοιστέον.}{}
\ligneEst{\textbf{Τοξότης}}{μιαρὸς ἀλώπηξ, οἶον ἐπιτήκιζί μοι.}{}
\ligneEst{\textbf{Μνησίλοχος}}{μέμνησο Περσεῦ μ᾽ ὡς καταλείπεις ἀθλίαν.}{}
\ligneEst{\textbf{Τοξότης}}{ἔτι γὰρ σὺ τὴ μάστιγαν ἐπιτυμεῖς λαβεῖν;}{1135}
\ligneEst{\textbf{Χορός} {\tiny\textit{ibyc.}}}{Παλλάδα τὴν φιλόχορον ἐμοὶ}{}
\ligneEst{{\tiny\textit{ibyc.}}}{δεῦρο καλεῖν νόμος εἰς χορόν,}{}
\ligneEst{{\tiny\textit{ibyc\cat}}}{παρθένον ἄζυγα κούρην,}{}
\ligneEst{{\tiny\textit{(στροφή α) / glyc.}}}{ἣ πόλιν ἡμετέραν ἔχει}{1140}
\ligneEst{{\tiny\textit{glyc.}}}{καὶ κράτος φανερὸν μόνη}{}
\ligneEst{{\tiny\textit{pherecr.}}}{κλῃδοῦχός τε καλεῖται.}{}
\ligneEst{{\tiny\textit{(μεσῳδός α) / $4$ bacch.}}}{φάνηθ᾽ ὦ τυράννους στυγοῦσ᾽ ὥσπερ εἰκός.}{}
\ligneEst{{\tiny\textit{(ἀντιστρ. α) / glyc.}}}{δῆμός τοί σε καλεῖ γυναι-}{1145}
\ligneEst{{\tiny\textit{$+$ glyc.}}}{\alinea κῶν· ἔχουσα δέ μοι μόλοις}{}
\ligneEst{{\tiny\textit{pherecr.}}}{εἰρήνην φιλέορτον.}{}
}
\vspace{2pt}
\noindent\rule{16.5cm}{0.6pt}\\[2pt]
\noindent\textbf{\small Commentaire philologique, métrique et scénique (vv. 1128–1168 --- partie 1/4)}\\[2pt]
\begin{adjustbox}{max totalheight=13.0cm, max width=16.5cm, keepaspectratio}%
\parbox[t]{16.5cm}{
\setlength{\columnsep}{9pt}%
\begin{multicols}{2}
\fontsize{8.4pt}{9.7pt}\selectfont
\setlength{\parindent}{0pt}%
\setlength{\parskip}{1.2pt}%
\noindent\textbf{1130} \textbf{σκαιοῖσι γάρ τοι καινὰ προσφέρων σοφὰ.} Cf. le reproche qu'Euripide mit dans la bouche de Médée dans la tragédie de ce nom (vers 298) :
\citComm{σκαιοῖσι μὲν γὰρ καινὰ προσφέρων σοφὰ.}
\noindent\textbf{1133} \textbf{οἷον ἐπιτηκίζει μοι} $=$ οἷον ἐπιθήκιζέ μοι. À propos des images animales du trompeur : du renard, ἡ ἀλώπηξ, et du singe, ὁ πίθηκος, vus comme les représentant de la ruse et de la fourbe, le grand Taillardat a consacré deux paragraphes auxquels je renvoie le lecteur : §§405-406. Notons que sur le même modèle que πιθηκίζειν (qui se rencontre aussi au vers 1290 des \italique{Guêpes}), il existe chez Aristophane le verbe ἀλωπεκίζειν (\italique{Guêpes}, vers 1241).
\noindent\textbf{1136-59} \textbf{Chant ibycéen du chœur.} Il a été beaucoup discuté du caractère strophique ou a-strophique de ce chant. Si la balance des commentateurs penche plutôt en faveur d'un chant astrophique\marqueur{1}, il est tentant, à l'instar d'Hermann, von Velsen ou Prato, de décomposer ce chant en deux parties avec pour chacune la division classique du chœur en deux demi-chœurs (strophe -- antistrophe). Voici la structure que je propose :
\par\vspace{1pt}\noindent\adjustbox{max width=\columnwidth}{$\displaystyle\begin{alignedat}{4}
\text{1\expo{ère} partie } & \text{ : invocation d'Athéna, } && \text{ vv. 1140-1142 } &&& \text{ \ } = \text{ vv. 1145-1147 ; }\\
\text{2\expo{nde} partie } & \text{ : invocation des Thesmophores, } && \text{ vv. 1148-1153 } &&& \simeq \text{ vv. 1155-1159.}
\end{alignedat}$}\par\vspace{1pt}
\end{multicols}
}%
\end{adjustbox}%
\vfill
\end{adjustwidth}
\newpage
% --- PAGE DE DROITE : Traduction française + Commentaire (2/4) ---
\begin{adjustwidth}{-2cm}{-2cm}
\noindent\textbf{\footnotesize Traduction française --- Folio 28 (\texttt{f. 176r}, vv. 1128–1168 [1/2]) \hfill ΧΟΡΙΚΟΝ / ΕΞΟΔΟΣ}\\[-2pt]
\noindent\rule{16.5cm}{0.5pt}\\[2pt]
{\fontsize{8.4pt}{9.7pt}\selectfont
\ligneEst{\textbf{Euripide}}{\francais{Aïe, aïe! que faire ? à quelles raisons recourir ?}}{1128}
\ligneEst{}{\francais{Mais cette nature barbare ne les admettrait pas.}}{}
\ligneEst{}{\francais{« Proposer à des sots des doctrines nouvelles,}}{1130}
\ligneEst{}{\francais{c'est perdre sa peine. » Eh bien, dressons quelqu'autre}}{}
\ligneEst{}{\francais{machine, qui soit de sa portée. (Il s'en va.)}}{}
\ligneEst{\textbf{Archer}}{\francais{La canaille de renard, comme il pensait m'embabouiner !}}{}
\ligneEst{\textbf{Mnésiloque}}{\francais{Souviens-toi, Persée, que tu me laisses en vilaine posture.}}{}
\ligneEst{\textbf{Archer}}{\francais{As-tu encore envie d'attraper le fouet? (Il se recouche et s'endort.)}}{1135}
\ligneEst{\textbf{Chœur}}{\francais{C'est ma coutume d'appeler parmi nous Pallas,}}{}
\ligneEst{}{\francais{amie du chœur, la jeune fille,}}{}
\ligneEst{}{\francais{la vierge libre du joug,}}{}
\ligneEst{}{\francais{gardienne de notre acropole,}}{1140}
\ligneEst{}{\francais{seule en possession d'un visible pouvoir,}}{}
\ligneEst{}{\francais{et qui a nom Clédouque.}}{}
\ligneEst{}{\francais{Apparais, toi qui détestes les tyrans, comme faire se doit.}}{}
\ligneEst{}{\francais{Le peuple des femmes t'appelle.}}{1145}
\ligneEst{}{\francais{Viens à moi avec}}{}
\ligneEst{}{\francais{la Paix amie des fêtes.}}{}
}
\vspace{2pt}
\noindent\rule{16.5cm}{0.6pt}\\[2pt]
\noindent\textbf{\small Commentaire philologique, métrique et scénique (vv. 1128–1168 --- partie 2/4)}\\[2pt]
\begin{adjustbox}{max totalheight=13.0cm, max width=16.5cm, keepaspectratio}%
\parbox[t]{16.5cm}{
\setlength{\columnsep}{9pt}%
\begin{multicols}{2}
\fontsize{8.4pt}{9.7pt}\selectfont
\setlength{\parindent}{0pt}%
\setlength{\parskip}{1.2pt}%
\noindent (Nous n'avons pas cependant forcé la concordance métrique au point de retoucher le texte lorsqu'il ne s'y prêtait que de mauvaise grâce ; si les vers 1140 ou 1148 peuvent se voir comme des glyconéens à base trissyllabique\marqueur{2}, ainsi, à la rigueur, que le vers 1155\marqueur{3}, il subsiste un problème de correspondance entre le vers 1153, ὄργια σεμνὰ θεαῖν ἵνα λαμπάσι, et le vers 1159, νῦν ἀφίκεσθον ἱκετεύομεν ἐνθάδ᾽ ἡμῖν.)\\ Mis de cette façon en ordre, le chant est entrecoupé entre chaque paire strophique d'un intermède, μεσῳδός, commençant par un impératif de φαίνω : aux vers 1143-4, φάνηθ᾽ ὦ τυράννους στυγοῦσ᾽ ὥσπερ εἰκός, et au vers 1154, φαίνετον ἄμβροτον ὄψιν.
\noindent\textbf{[Notes ad 1136-59]} \textbf{(1)} L'analyse la plus poussée de ce chant est celle de Lætitia, \italique{The Songs of Aristophanes}, pp. 446-453 ; on pourra aussi consulter les non moins magistraux ouvrages : \italique{The Lyric Metres of Greek Drama}\expo{2} (1968) pp. 166-167 d'Amy M. Dale ainsi que \italique{The Verse of Greek Comedy} (1912), §387, de White. White ne perçoit de répons antistrophique que pour les $6 = 3 + 3$ premiers vers.
\noindent \textbf{(2)} Cf. Euripide, \italique{Iphigénie en Aulis}, vers 759.
\noindent \textbf{(3)} Voir ce qu'en dit Amy, \italique{The Lyric Metres of Greek Drama}\expo{2}, p. 166 n.2.
\end{multicols}
}%
\end{adjustbox}%
\vfill
\end{adjustwidth}
\newpage

% --- FOLIO 28 (f. 176r) : SPREAD 2.2 (Texte grec établi + Traduction française + Commentaire 3/4 & 4/4) ---
\clearpage\ifodd\value{page}\hbox{}\newpage\fi
\enTeteG{ΘΕΣΜΟΦΟΡΙΑΖΟΥΣΑΙ --- Texte grec établi (f. 176r, vv. 1128–1168 [2/2])}
\enTeteD{LES THESMOPHORIES --- Traduction française (f. 176r, vv. 1128–1168 [2/2])}
\begin{adjustwidth}{-2cm}{-2cm}
\noindent\textbf{\footnotesize Texte grec établi --- Folio 28 (\texttt{f. 176r}, vv. 1128–1168 [2/2]) \hfill ΧΟΡΙΚΟΝ / ΕΞΟΔΟΣ}\\[-2pt]
\noindent\rule{16.5cm}{0.5pt}\\[2pt]
{\fontsize{8.4pt}{9.7pt}\selectfont
\ligneEst{{\tiny\textit{(στροφή β) / glyc.}}}{ἥκετέ τ᾽ εὔφρονες ἵλαοι,}{}
\ligneEst{{\tiny\textit{ibyc.}}}{πότνιαι, ἄλσος ἐς ὑμέτερον,}{}
\ligneEst{{\tiny\textit{dact. aeol.}}}{οὗ δὴ ἀνδράσιν οὐ θεμιτεύετον εἰσορᾶν}{1150}
\ligneEst{{\tiny\textit{tetram. dact.}}}{ὄργια σεμνὰ θεαῖν ἵνα λαμπάσι.}{}
\ligneEst{{\tiny\textit{(μεσῳδός β) / trim. dact.}}}{φαίνετον ἄμβροτον ὄψιν.}{}
\ligneEst{{\tiny\textit{(ἀντιστρ. β) / glyc.}}}{μόλετον, ἔλθετον, ἀντόμεθ᾽ ὦ}{1155}
\ligneEst{{\tiny\textit{ibyc.}}}{Θεσμοφόρω πολυποτνία,}{}
\ligneEst{{\tiny\textit{dact. aeol.}}}{εἰ καὶ πρότερόν ποτ᾽ ἐπηκόω ἤλθετον}{}
\ligneEst{}{νῦν ἀφίκεσθον †ἱκετεύομεν ἐνθάδ᾽ ἡμῖν†.}{}
\ligneEst{\textbf{Εὐριπίδης} {\tiny\textit{trim. iamb.}}}{γυναῖκες εἰ βούλεσθε τὸν λοιπόν χρόνον}{1160}
\ligneEst{}{σπονδὰς ποιήσασθαι πρὸς ἐμέ, νυνὶ πάρα,}{}
\ligneEst{}{ἐφ᾽ ᾧτ᾽ ἀκοῦσαι μηδὲν ὑπ᾽ ἐμοῦ μηδαμὰ}{}
\ligneEst{}{κακὸν τὸ λοιπόν. ταῦτ᾽ ἐπικηρυκεύομαι.}{}
\ligneEst{\textbf{Κορυφαῖος}}{χρείᾳ δὲ ποίᾳ τόνδ᾽ ἐπεσφέρεις λόγον;}{}
\ligneEst{\textbf{Εὐριπίδης}}{ὅδ᾽ ἐστὶν οὑν τῇ σανίδι κηδεστὴς ἐμός.}{1165}
\ligneEst{}{ἢν οὖν κομίσωμαι τοῦτον, οὐδὲν μή ποτε}{}
\ligneEst{}{κακῶς ἀκούσητ᾽· ἢν δὲ μὴ πίθησθέ μοι,}{}
\ligneEst{}{ἃ νῦν ὑποικουρεῖτε τοῖσιν ἀνδράσιν}{1168}
}
\vspace{2pt}
\noindent\rule{16.5cm}{0.6pt}\\[2pt]
\noindent\textbf{\small Commentaire philologique, métrique et scénique (vv. 1128–1168 --- partie 3/4)}\\[2pt]
\begin{adjustbox}{max totalheight=13.0cm, max width=16.5cm, keepaspectratio}%
\parbox[t]{16.5cm}{
\setlength{\columnsep}{9pt}%
\begin{multicols}{2}
\fontsize{8.4pt}{9.7pt}\selectfont
\setlength{\parindent}{0pt}%
\setlength{\parskip}{1.2pt}%
\noindent\textbf{1150-3} \textbf{οὗ δὴ ἀνδράσιν οὐ / θεμιτὸν εἰσορᾶν / ὄργια σεμνὰ θεαῖν / ἵνα λαμπάσι φαί-/} Les mots †οὐ θεμιτὸν† ont été diversement retouchés, jamais de façon satisfaisante. Pour ma part, j'y verrais le verbe θεμιτεύω au duel.
\noindent La première partie du chant est une prière adressée à Athéna ; un scribe -- point bien réveillé -- n'a pas tout-de-suite vu que la seconde partie s'adressait aux deux déesses thesmophores.
\noindent Au premier vers de cette seconde partie (1148), il aura omis de mettre le verbe au pluriel : ἥκετ<έ τ᾽> εὔφρονες (simple \italique{haplographie}, comme on dit).
\noindent \italique{Bis repetita} : non encore averti par la longue énumération des duels aux vers 1154 et suivants, il persévère dans son erreur et le premier duel lui échappe tout aussi haplographiquement :
\citComm{θεμι<τεύε>τον.}
\noindent Le verbe θεμιτεύω, \italique{permettre selon la loi ou le rite}, est utilisé par Euripide avec le même complément, ὄργια, au vers 79 des \italique{Bacchantes} :
\citComm{τά τε ματρὸς μεγάλας ὄρ-\\
\alinea για Κυβέλας θεμιτεύων\\
Vénérant les orgies de la mère Cybèle la grande}
\end{multicols}
}%
\end{adjustbox}%
\vfill
\end{adjustwidth}
\newpage
% --- PAGE DE DROITE : Traduction française + Commentaire (4/4) ---
\begin{adjustwidth}{-2cm}{-2cm}
\noindent\textbf{\footnotesize Traduction française --- Folio 28 (\texttt{f. 176r}, vv. 1128–1168 [2/2]) \hfill ΧΟΡΙΚΟΝ / ΕΞΟΔΟΣ}\\[-2pt]
\noindent\rule{16.5cm}{0.5pt}\\[2pt]
{\fontsize{8.4pt}{9.7pt}\selectfont
\ligneEst{}{\francais{Venez, augustes déesses, bienveillantes et propices,}}{}
\ligneEst{}{\francais{dans votre enceinte sacrée,}}{}
\ligneEst{}{\francais{où vous ne permettez pas aux hommes de contempler}}{1150}
\ligneEst{}{\francais{vos saintes orgies à la lueur des torches.}}{}
\ligneEst{}{\francais{Découvrez vos faces immortelles.}}{}
\ligneEst{}{\francais{Venez, approchez, nous vous en conjurons,}}{1155}
\ligneEst{}{\francais{Thesmophores très vénérées.}}{}
\ligneEst{}{\francais{Si jamais auparavant vous répondîtes à notre appel,}}{}
\ligneEst{}{\francais{venez aujourd'hui, sur nos instantes prières, ici au milieu de nous.}}{}
\ligneEst{\textbf{Euripide}}{\francais{(Déguisé en vieille procureuse et portant une pectis ou harpe lydienne. Il est accompagné d'une danseuse et d'un jeune flûtiste). Femmes, si vous voulez pour le reste du temps}}{1160}
\ligneEst{}{\francais{faire la paix avec moi, c'est le moment,}}{}
\ligneEst{}{\francais{à charge que jamais plus à l'avenir je ne dirai de vous}}{}
\ligneEst{}{\francais{le moindre mal. Telles sont mes propositions.}}{}
\ligneEst{\textbf{Coryphée}}{\francais{Qu'est-ce qui te décide à faire ces ouvertures ?}}{}
\ligneEst{\textbf{Euripide}}{\francais{Cet homme attaché à l'ais est mon beau-père.}}{1165}
\ligneEst{}{\francais{Que je le recouvre, et jamais}}{}
\ligneEst{}{\francais{je ne médirai de vous. Si vous refusez,}}{}
\ligneEst{}{\francais{toutes vos pratiques à huis-clos, je les révélerai}}{1168}
}
\vspace{2pt}
\noindent\rule{16.5cm}{0.6pt}\\[2pt]
\noindent\textbf{\small Commentaire philologique, métrique et scénique (vv. 1128–1168 --- partie 4/4)}\\[2pt]
\begin{adjustbox}{max totalheight=13.0cm, max width=16.5cm, keepaspectratio}%
\parbox[t]{16.5cm}{
\setlength{\columnsep}{9pt}%
\begin{multicols}{2}
\fontsize{8.4pt}{9.7pt}\selectfont
\setlength{\parindent}{0pt}%
\setlength{\parskip}{1.2pt}%
\noindent Cerise sur le gâteau : cette émendation permet de rétablir la correspondance métrique avec le vers 1158 :
\citComm{οὗ δ\synizese{ὴ ἀ}νδράσιν\marqueur{1} οὐ θεμιτεύετον εἰσορᾶν\\
εἰ καὶ πρότερόν ποτ᾽ ἐπηκόω ἤλθετον\\
$ - \ - \ \cup \ \cup \ - \ \cup \ \cup \ - \ \cup \ \cup \ - \ \cup \ -$\marqueur{2}}
\noindent\textbf{[Notes ad 1150-53]} \textbf{(1)} Signalons la synalèphe au début du vers : δ\synizese{ὴ ἀ}νδράσιν -- je reprends le vocable de Koster dans son \italique{Traité de métrique grecque}. Il la distingue de la synizèse, terme qu'il applique lorsque la contraction entre deux voyelles a lieu au sein d'un même mot. Subtilité de grammairien ? pas que : la synizèse semble obéir à des lois différentes quant à la longueur des voyelles. Prato, lui, évoque à ce vers la présence d'une \italique{crase}.
\noindent \textbf{(2)} Séquence interprétée par Lætitia comme un \italique{dactyle éolien}, c'est-à-dire dans son jargon une base (ici unissyllabique $ \times $, permise par l'éolisme du vers) sur laquelle viennent se fixer un nombre arbitraire de dactyles purs ($ - \cup \cup $) et finissant par un \italique{dodrans} ($ - \cup \cup - \cup - $), incontournable en contexte choriambique.
\noindent\textbf{1151} \textbf{ἵνα,} Une glose à l'intérieur du texte du manuscrit\marqueur{} précise quel sens donner à (ἵνα : ) ὅπου\marqueur{1}.
\noindent\textbf{[Notes ad 1151]} \textbf{(1)} On retrouve une glose identique dans de nombreux lexiques : Συναγωγή ι 82 ; Hésychius ι 648 ; Photius ι 133 ; Suidas ι 360 ; \&c.
\end{multicols}
}%
\end{adjustbox}%
\vfill
\end{adjustwidth}
\newpage
\clearpage\ifodd\value{page}\hbox{}\newpage\fi
\enTeteG{ΕΞΟΔΟΣ --- Codex Ravennas 429, f. 176v (Classense 178v) --- vv. 1169–1210}
\begin{adjustwidth}{-1.5cm}{-3.0cm}
\begin{adjustbox}{max totalheight=21.0cm, max width=16.5cm, keepaspectratio}
\begin{minipage}{16.5cm}
{\fontsize{7.2pt}{8.5pt}\selectfont
\noindent\textbf{Folio 29 --- Codex Ravennas 429, \texttt{f. 176v} (\texttt{Classense 178v}) : ΕΞΟΔΟΣ, vv. 1169–1210}\hfill \textit{Témoins :} $\text{R}^1, \Rpc, \Rdeux$, \PsiTrois, \PsiQuatre\par\vspace{2pt}
\hrule height 0.4pt\vspace{2.5pt}
}
{\fontsize{8.1pt}{9.5pt}\selectfont
\ligneRav{1169}{\auDessus{εὐ}{ρ}}{ἀπὸ τῆς στρατιᾶς παροῦσιν ὑμῶν διαβαλῶ·}{}
\ligneRav{1170}{χορ/}{τὰ μὲν παρ᾽ ἡμῶν ἴσθι σοι πεπεισμένα·}{}
\ligneRav{1171}{}{τὸν βάρβαρον δὲ τοῦτον αὐτὸς \difference{πεῖθε} σύ·}{\textbf{1171} πεῖσαι R : \lecon{πεῖθε} Biset}
\ligneRav{1172}{\auDessus{εὐ}{ρ}}{οὐδὲν μὰ Δί᾽ ἐμὸν \inRas{γ᾽}{} ὤν, ὥσπερ ἐχθὲς \difference{ὤμνυες}·}{\textbf{1172} Εὐριπίδης ἐν σχήματι προαγωγοῦ γραός Σ \espApp \lecon{ἐμὸν} \inRas{γ᾽}{} \Rdeux}
\ligneRav{1173}{}{τί δ᾽; οὐκ \expo{ἔφραζον} (ἔφ\expo{ραζ}ον \Rpc) ταῦτ᾽ ἐγώ σοι;}{\textbf{1173} \lecon{ἔφραζον} (ex ἔφρον) \Rpc}
\ligneRav{1174}{}{πρῶτον μὲν οὖν δίελθε \difference{κἀνακάλπασον}·}{\textbf{1174} κἀνακόλπασον R : \lecon{κἀνακάλπασον} Hermann, cf. Σ, Suid. κ 623, Zonar. p.1171}
\ligneRav{1175}{}{σὺ δ᾽ ὦ \difference{Τερηδὸν} ἐπαναφύσα Περσικόν·}{\textbf{1175} Τερηδὼν R, Suid. τ 340 AGFM\expo{ac} : \lecon{Τερηδὸν} Suid. VM\expo{pc}}
\ligneRav{1176}{τοξ/}{τί τὸ βόμβο τοῦτο; κῶμό τις ἀνεγεῖρί μοι;}{\textbf{1176} κωμοτίς ἀν- R : κωμοδρίαν Σ : \lecon{κῶμό τις ἀν-} van Velsen}
\ligneRav{}{\auDessus{εὐ}{ρ}}{ἡ παῖς ἔμελλε προμελετᾶν ὦ τοξότα·}{}
\ligneRav{1178}{}{ὀρχησομένη γὰρ \difference{ἔρχεθ᾽} ὡς ἄνδρας τινάς·}{\textbf{1178} ἔρχετ᾽ R : \lecon{ἔρχεθ᾽} Brunck}
\ligneRav{1179}{τοξ/}{ὀρκῆσι καὶ μελετῆσι, οὐ \difference{κωλῦσ᾽} ἐγώ·}{\textbf{1179} κωλύσ᾽ R : \lecon{κωλῦσ᾽} Blaydes (κωλῦσι $=$ κωλύσω), cf. ad 1003, 1196, 1201, 1225}
\ligneRav{1180}{}{\difference{ὠς} ἐλαπρός, \difference{ὤσπερ} ψύλλο κατὰ τὸ κῴδιο·}{\textbf{1180} ὡς et ὥσπερ R : \lecon{ὠς} et \lecon{ὤσπερ} Brunck}
\ligneRav{}{\auDessus{εὐ}{ρ}}{φέρε θοἰμάτιον ἄνωθεν ὦ τέκνον τοδί·}{}
\ligneRav{1182}{}{καθιζομένη δ᾽ ἐπὶ \difference{τοῖσι} γόνασι τοῦ Σκύθου}{\textbf{1182} τοῖς R : \lecon{τοῖσι} Bentley}
\ligneRav{}{—}{τὼ πόδε πρότεινον, ἵν᾽ ὑπολύσω. : ναῖκι ναὶ}{}
\ligneRav{1184}{}{κάτησο κάτησο, ναίκι ναί, τυγάτριον·}{\textbf{1184} \lecon{ναίκι ναὶ} R : ναίκι ναίκι Bothe}
\ligneRav{1185}{}{ἀνίστασ᾽ ὀρθός· \difference{ὡς στεριπὸ} (ὡ σ\inRas{π}{τ}εριπὸ) τὸ τιττί᾽, ὥσπερ γογγύλη·}{\textbf{1185} ὡ στεριπὸ (ex σπεριπο) \Rpc : \lecon{ὡς στεριπὸ} (sic) \Mon \espApp ὥσπερ R : \lecon{ὤσπερ} Brunck \espApp \lecon{γογγυλί} \PsiTrois, Enger : γογγύλῃ R}
\ligneRav{}{\auDessus{εὐ}{ρ}}{αὔλει σὺ θᾶττον· ἔτι δέδοικας τὸν Σκύθην;}{}
\ligneRav{1187}{τοξ/}{καλόν γε τὸ πυγή· \difference{κλαυσῖ γ᾽} (κλαυσ\inRas{᾽ ε}{}ϊ γ᾽) ἂν μὴ ᾽νδον μένῃ·}{\textbf{1187} κλαυσ\inRas{᾽ ε}{}ϊ γ᾽ \Rdeux : κλαυσεῖ γ᾽ \Mon \espApp ἂν R : \lecon{ἢν} Dindorf (1839) \ligneApp{post 1187} μ$]$\douteux{ε}νο\douteux{ }$[$ \PsiTrois : ἀνακύπτη· καὶ παρακύπτι· ἀπεψωλημένος· R : \lecon{del.} Ellebodius}
\ligneRav{}{}{\hspace{2.2em}\textbf{[παρεπιγραφή R\expo{1}, \PsiTrois{} :]} \lecon{ἀνακύπτη· καὶ παρακύπτι· ἀπεψωλημένος·}}{\textbf{post 1187} \lecon{ἀνακύπτη· καὶ παρακύπτι· ἀπεψωλημένος·} R, μ]ενο[ \PsiTrois}
\ligneRav{1188}{}{εἶἑν· καλὴ τὸ σχῆμα περὶ τὸ \inRas{σ}{}ποστίον·}{\textbf{1188} τὸ \inRas{σ}{}ποστίον \Rpc : τὸ ποστίον \Mon}
\ligneRav{1189}{\auDessus{εὐ}{ρ}}{καλῶς ἔχει. λαβὲ θοἰμάτιον· ὥρα \difference{'στὶ} νῷν}{\textbf{1189} 'στιν R : \lecon{'στὶ} \Mon}
\ligneRav{1190}{—}{ἤδη βαδίζειν. : \difference{οὐκὶ} \difference{πιλῆσι} πρῶτά με;}{\textbf{1190} τί οὐκ ἐπιλήσει R : \lecon{οὐκὶ} (Biset) \lecon{πιλῆσι} Porson (ad \italique{Aristophanica} p.222), cf. Σ `οὐχὶ φιλήσειν'}
\ligneRav{1191}{\auDessus{εὐ}{ρ}}{πάνυ γε· φίλησον αὐτόν. : \difference{ὂ ὂ ὂ} παπαπαπαῖ,}{\textbf{1191} ὂ R : \lecon{ὂ ὂ ὂ} \Mon}
\ligneRav{1192}{}{\difference{ὠς} γλυκερὸ τὸ γλῶσσ᾽, \difference{ὤσπερ} Ἀττικὸς \difference{μέλι}·}{\textbf{1192} ὡς et ὥσπερ R : \lecon{ὠς} et \lecon{ὤσπερ} Brunck \espApp \lecon{μέλι} \PsiTrois, \Mon : μέλις R}
\ligneRav{}{—}{τί οὐ κατεύδει παρ᾽ ἐμέ; : χαῖρε τοξότα,}{}
\ligneRav{1194}{—}{οὐ γὰρ γένοιτ᾽ ἂν τοῦτο. : \difference{ναὶ ναὶ} \difference{γρᾴδιο}·}{\textbf{1194} ναὶ R : ναίκι Bentley : \lecon{ναὶ ναὶ} Brunck \espApp γράδιον R : \lecon{γρᾴδιο} Bothe}
\ligneRav{1195}{—}{ἐμοὶ \inRas{σὺ}{} \difference{κάρισο σὺ} τοῦτο : δώσεις οὖν δραχμή;}{\textbf{1195} \inRas{σὺ}{} κάρισοσ οὐ \Rpc : κάρισο σοῦ \Mon : \lecon{κάρισο σὺ} Brunck}
\ligneRav{}{τοξ/}{ναὶ ναῖκι δῶσι. : τἀργύριον τοίνυν φέρε·}{}
\ligneRav{1197}{τοξ/}{ἀλλ᾽ οὐκ \difference{ἔκὠδέν}· ἀλλὰ τὸ συμβήνην λαβέ·}{\textbf{1197} έκ' ὠδέν R : ἐχω οὐδέν Suid. σ 1273, Photius σ 660 : \lecon{ἔκὠδέν} Bergk \espApp \lecon{συμβήνην} R : συβήνην Σ, Suid., Photius}
\ligneRav{1198}{}{\difference{κομίζι σ᾽ αὐτίκ᾽} (κ\inRas{ω}{ο}μίζεις αὐτοῖς) : ὦ τάλαιν᾽ ἐγώ·}{\textbf{1198} κ\inRas{ω}{ο}μίζεις αὐτοῖς \Rpc : κομιεῖς αὐτό (αὖθις Brunck) Biset : (Τοξότ.) κομίσι σ' αὐτός (αὖτις anonym. Parisinus) Enger : (Τοξότ.) \lecon{κομίζι σ᾽ αὐτίκ᾽} Bibi \espApp ἀκουλούτι R : \lecon{ἀκολούτι} \Mon : ἀκουλούτει Brunck}
\ligneRav{}{}{σὺ δὲ τοῦτο τήρει τὴ γέροντο, γρᾴδιο·}{}
\ligneRav{1200}{—}{ὄνομα δέ σοι τί ἔστιν; : Ἀρτεμισία·}{}
\ligneRav{1201}{τοξ/}{μεμνῆσι τοίνυν τοὔνομ᾽· \difference{Ἀρταμουξία}·}{\textbf{1201} ἀρτομουξία R : \lecon{Ἀρταμουξία} ed. Lipsiensis (1814)}
\ligneRav{}{\auDessus{εὐ}{ρ}}{Ἑρμῆ δόλιε ταυτὶ μὲν ἔτι καλῶς ποιεῖς·}{}
\ligneRav{1203}{}{σὺ μὲν οὖν ἀπότρεχε παιδάριον τουτὶ λαβών·}{\textbf{1203} \lecon{τουτὶ} R : ταυτὶ Dobree, cf. Σ}
\ligneRav{}{}{ἐγὼ δὲ λύσω τόνδε. σὺ δ᾽ ὅπως ἀνδρικῶς}{}
\ligneRav{1205}{}{ὅταν λυθῇς τάχιστα φεύξει καὶ τενεῖς}{}
\ligneRav{}{}{ὡς τὴν γυναῖκα καὶ τὰ παιδί᾽ οἴκαδε·}{}
\ligneRav{}{μ\auDessus{ν}{Η}}{ἐμοὶ μελήσει ταῦτά γ᾽ ἢν ἅπαξ λυθῶ·}{}
\ligneRav{1208}{\auDessus{εὐ}{ρ}}{ὠς τάκιστα σὺ λέλυσο\inRas{ν}{} (cf. \PsiQuatre) καὶ φεῦγ᾽·}{\textbf{1208} ca. 6 litt. $]$· σ\douteux{ο}$[$νεργον \PsiQuatre : \lecon{λέλυσο}\inRas{ν}{} \Rpc : λέλυσαι Bentley}
\ligneRav{}{—}{ἥκοντα καταλαβεῖν. : ἐγὼ δὴ τοῦτο δρῶ·}{}
\ligneRav{1210}{τοξ/}{σὺ δὲ τυγά\inRas{θ}{τ}ριον (\Rpc) κόσμιε καὶ φύλαττε·}{\textbf{1210} τυγάτριον ex τυγάθριον ut videtur \Rpc}
}
\vspace{3pt}\noindent\rule{4.5cm}{0.4pt}\par\vspace{1.5pt}
{\fontsize{7.0pt}{8.3pt}\selectfont
\noindent\textbf{Loci paralleli (vv. 1169–1210) ---} \textbf{1169} (στρατιᾶς) Photius σ 608 στρατιάν \espTest \textbf{1174} (κἀνακάλπασον) Photius α 1514 ἀνακαλπάζει ; cf. Zonaras 1171 καλπάζειν· τὸ ἄκρως βαδίζειν $\simeq$ Scholium $=$ Suidas κ 623 καλπάζειν \espTest \textbf{1175} Suidas τ \espTest (Περσικόν) Pollux IV 100 ; cf. Hesychius π 2005 (?) \espTest \textbf{1183} (τὼ πόδε) Hesychius τ 1758 τὼ πόδε· τοὺς πόδας δυϊκῶς \espTest \textbf{1187} (πυγή) Hesychius π 4279 πυγή· τὸ κάθισμα \espTest \textbf{1197} Photius σ 660 συβήνη $=$ Suidas σ 1273 (sed ἐχω οὐδέν pro έκ' ὠδέν, τὸ omisso et συβήνην pro συμβήνην) \espTest (συμβήνην) Hesychius σ 2137 συμβίνης· καπροβόλον. ἐμβόλιον\par
}
\vspace{2.5pt}\noindent\rule{4.5cm}{0.4pt}\par\vspace{1.5pt}
{\fontsize{6.9pt}{8.2pt}\selectfont
\noindent\textbf{Suite (en commentaire) de l'apparat critique \& notes paléographiques (\texttt{f. 176v}) ---} \textbf{v. 1172 :} $\Sigma^{\text{R}}$ Εὐριπίδης ἐν σχήματι προαγωγοῦ γραός \espApp \textbf{ἐμὸν \inRas{γ'}{}} \Rdeux{}. --- \textbf{v. 1174 :} κἀνακόλπασον R (confusion minuscule $\alpha \to \text{ο}$) : \lecon{κἀνακάλπασον} Hermann, cf. $\Sigma$, Souda $\kappa\ 623$, Zonaras p. 1171 (dans $\Sigma^{\text{R}}$, ἁβρῶς R : \lecon{ἄκρως} \lecon{Bibi} ex Zonaras). --- \textbf{v. 1176 :} κωμοτίς ἀν- R (κωμοδρίαν $\Sigma$) : \lecon{κῶμό τις ἀν-} von Velsen. --- \textbf{v. 1190 :} τί οὐκ ἐπιλήσει R : \lecon{οὐκὶ} (Biset) \lecon{πιλῆσι} Porson (\italique{Aristophanica}, p. 222, cf. $\Sigma$ οὐχὶ φιλήσειν). --- \textbf{v. 1197 :} έκ' ὠδέν R (ἐχω οὐδέν Souda $\sigma\ 1273$, Photius $\sigma\ 660$) : \lecon{ἔκὠδέν} Bergk \espApp \lecon{συμβήνην} R, défendu par \lecon{Bibi} (ap. Hésychius $\sigma\ 2137$) : συβήνην $\Sigma$, Souda, Photius. --- \textbf{v. 1198 :} γϋ. ante v. 1198 (Boninus) et σκ'. ante ἀκολ. in R \espApp κ\inRas{ω}{ο}μίζεις αὐτοῖς \Rpc{} (κομιεῖς αὐτό Biset ; αὖθις Brunck ; (Τοξότ.) κομίσι σ' αὐτός Enger) : \lecon{(Τοξότ.) κομίζι σ' αὐτίκ'} \lecon{Bibi} \espApp ἀκουλούτι R : \lecon{ἀκολούτι} \Mon{} (ἀκουλούτει Brunck). --- \textbf{v. 1203 :} \lecon{τουτὶ} R (ταυτὶ Dobree, cf. $\Sigma$).\par
}
\end{minipage}
\end{adjustbox}
\end{adjustwidth}
\newpage
\enTeteD{Codex Ravennas 429 (Biblioteca Classense) --- Fac-similé du f. 176v (vv. 1169–1210)}
\begin{adjustwidth}{-1.5cm}{-3.0cm}
\begin{center}
\includegraphics[width=16.5cm,height=20.5cm,keepaspectratio]{/workspaces/Thesmophories/dossier_ravennas_thesmophories/folios_individuels_ff162v-177r/29_f176v_classense178v_vv1169-1210.jpg}\par
\vspace{2pt}
{\scriptsize\textbf{Codex Ravennas 429, f. 176v (Classense 178v)} --- ΕΞΟΔΟΣ, \textbf{vv. 1169–1210} ($\text{R}^1, \Rpc, \Rdeux$, \PsiTrois, \PsiQuatre).}
\end{center}
\end{adjustwidth}
\newpage
\clearpage\ifodd\value{page}\hbox{}\newpage\fi% --- FOLIO 29 (f. 176v) : SPREAD 2.1 (Texte grec établi + Traduction française + Commentaire 1/8 & 2/8) ---
\clearpage\ifodd\value{page}\hbox{}\newpage\fi
\enTeteG{ΘΕΣΜΟΦΟΡΙΑΖΟΥΣΑΙ --- Texte grec établi (f. 176v, vv. 1169–1210 [1/4])}
\enTeteD{LES THESMOPHORIES --- Traduction française (f. 176v, vv. 1169–1210 [1/4])}
\begin{adjustwidth}{-2cm}{-2cm}
\noindent\textbf{\footnotesize Texte grec établi --- Folio 29 (\texttt{f. 176v}, vv. 1169–1210 [1/4]) \hfill ΕΞΟΔΟΣ}\\[-2pt]
\noindent\rule{16.5cm}{0.5pt}\\[2pt]
{\fontsize{8.8pt}{10.2pt}\selectfont
\ligneEst{}{ἀπὸ τῆς στρατιᾶς παροῦσιν ὑμῶν διαβαλῶ.}{1169}
\ligneEst{\textbf{Κορυφαῖος}}{τὰ μὲν παρ᾽ ἡμῶν ἴσθι σοι πεπεισμένα·}{1170}
\ligneEst{}{τὸν βάρβαρον δὲ τοῦτον αὐτὸς πεῖθε σύ.}{}
\ligneEst{\textbf{Εὐριπίδης}}{ἐμὸν ἔργον ἐστίν· καὶ σὸν ὦλάφιον ἅ σοι}{}
\ligneEst{}{καθ᾽ ὁδὸν ἔφραζον ταῦτα μεμνῆσθαι ποιεῖν.}{}
\ligneEst{}{πρῶτον μὲν οὖν δίελθε κἀνακάλπασον.}{}
\ligneEst{}{σὺ δ᾽ ὦ Τερηδὸν ἐπαναφύσα Περσικόν.}{1175}
\ligneEst{\textbf{Τοξότης}}{τί τὸ βόμβο τοῦτο; κῶμό τις ἀνεγεῖρί μοι;}{}
\ligneEst{\textbf{Εὐριπίδης}}{ἡ παῖς ἔμελλε προμελετᾶν ὦ τοξότα.}{}
\ligneEst{}{ὀρχησομένη γὰρ ἔρχεθ᾽ ὡς ἄνδρας τινάς.}{}
\ligneEst{\textbf{Τοξότης}}{ὀρκῆσι καὶ μελετῆσι, οὐ κωλῦσ᾽ ἐγώ.}{}
\ligneEst{}{ὠς ἐλαπρός, ὤσπερ ψύλλο κατὰ τὸ κῴδιο.}{1180}
}
\vspace{2pt}
\noindent\rule{16.5cm}{0.6pt}\\[2pt]
\noindent\textbf{\small Commentaire philologique, métrique et scénique (vv. 1169–1210 --- partie 1/8)}\\[2pt]
\begin{adjustbox}{max totalheight=15.2cm, max width=16.5cm, keepaspectratio}%
\parbox[t]{16.5cm}{
\setlength{\columnsep}{9pt}%
\begin{multicols}{2}
\fontsize{7.0pt}{8.0pt}\selectfont
\setlength{\parindent}{0pt}%
\setlength{\parskip}{1.2pt}%
\noindent\textbf{1168-9} \textbf{ἃ νῦν ὑποικουρεῖτε τοῖσιν ἀνδράσιν / ἀπὸ τῆς στρατιᾶς παροῦσιν.} À propos de la légèreté des Athéniennes pendant le service de leur mari, cf. le vers 495 et le commentaire \italique{ad locum}.
\noindent\textbf{1170} \textbf{τὰ μὲν παρ᾽ ἡμῶν ἴσθι σοι πεπεισμένα·} Une scholie intérieure au texte nous précise :
\citComm{Ἀττικῶς ἀντὶ τοῦ `πεπείσμεθα'.}
\noindent\textbf{1174} \textbf{πρῶτον μὲν οὖν δίελθε κἀνακόλπασον.} (Que le lecteur curieux aille faire son miel de la note qu'a consacré Jacques Moreau à ce vers\marqueur{1}.) Le \italique{Ravennas} porte au vers 1174 la leçon : πρῶτον μὲν οὖν δίελθε κἀνακ\textbf{ό}λπασον.
\noindent Le verbe ἀνακολπάζειν est un hapax, certes, mais dont le sens se laisserait entendre -- s'il n'y avait lieu comme nous l'allons voir de le soupçonner : ὁ κόλπος désigne au propre un \italique{pli}, \italique{creux} ; aussi bien le pli d'un vêtement -- notamment au niveau de la ceinture -- que, chez la mère ou la nourrice, le \italique{genou} ou le \italique{sein}. On trouve chez Homère l'expression κόλπον ἀνιεμένη\marqueur{2} lorsqu'Hécube fait tomber le haut de son vêtement et découvre son sein.
\noindent Il n'y aurait rien d'étonnant à ce qu'ἀνακολπάζειν ait pu signifier \italique{retrousser un vêtement}\marqueur{3}.
\noindent Pour définir les ζειραί, ces vêtements qu'on \italique{resserre au niveau de la ceinture}, le rédacteur anonyme de l'\latin{Etymologicum Magnum} use du verbe ἀνακολπόω :
\citComm{ζειραί: χιτῶνες ἀνακεκολπωμένοι\marqueur{4}, ἢ ἀνάκωλοι.}
\noindent Si le \italique{Ravennas} ne comportait en sus du texte au vers 1174 une scholie marginale venant renverser tout notre fragile édifice, ἀνακολπάζειν ferait partie de ces nombreux hapax dont le sens n'a été inféré qu'à grand renfort d'imagination sur de maigres indices étymologiques.\\ Cependant un ancien commentateur eut l'heureuse idée de gloser la leçon que son manuscrit donnait :
\citComm{σ\douteux{η}[ * * * ] τοῦτο. ση[μαίνει] δὲ τὸ `ἁβρῶς βαδίζειν'.}
\noindent Exit donc la leçon ἀνακολπάζειν, sans rapport avec le \italique{déplacement langoureux}\marqueur{5} dont parle le scholiaste, nullement étayée et d'ailleurs inepte : la danseuse court-vêtue n'a aucun vêtement pesant à retrousser ; au moment d'enflammer le Scythe, Euripide lui fera sauter les préliminaires en lui ordonnant à brûle-pourpoint de retirer tout bonnement son léger ἱμάτιον.\\ Hermann\marqueur{6} proposa de lire le verbe ἀνακαλπάζειν\marqueur{7}.
\noindent Son hypothèse a été reprise depuis lors par toutes les éditions ; elle s'appuyait effectivement sur un indice de taille : à l'adresse κ 623 de la \italique{Souda}, l'on retrouve notre glose précédée du lemme καλπάζειν :
\citComm{καλπάζειν· τὸ ἁβρῶς βαδίζειν.}
\noindent Le verbe est formé à partir de ἡ κάλπη, \italique{le trot} (ou le \italique{galop} ?), qui désigne une sorte de course aux jeux olympiques, où le cavalier, en approchant du but, devait prendre le trot, sauter à terre et suivre son cheval.
\noindent Quant au composé, ἀνακαλπάζειν, Photius (α 1514) nous en a conservé une trace en recopiant la glose d'un grammairien versé dans le \italique{Bon usage} attique :
\citComm{ἀνακαλπάζει· τινὲς μὲν ὡς οὐ δόκιμον ἐφυλάξαντο τὴν φωνήν, Αἰσχύλος δὲ ἐχρήσατο Μυσοῖς \{ὡς δόκιμον\}· λέγει γάρ `εἶδον καλπάζοντας ἐν αἰχμαῖς'. ὁμοίως Σοφοκλῆς καὶ Ἀριστοφάνης καὶ Πλάτων καὶ ἕτεροι.\\
ἀνακαλπάζει· certains défendirent [l'usage de] de ce mot, comme non attesté. Eschyle pourtant l'utilisa dans ses \italique{Mysiens}. Il dit en effet : «~Je les vis trotter au milieu des lances~». De même Sophocle, Aristophane, Platon et d'autres.}
\noindent C'est bien le sens de \italique{trotter}, \italique{galoper}, c'est-à-dire \italique{courir} qu'il faut donner au verbe <ἀνα>καλπάζειν, témoin le fragment d'Eschyle cité par Photius\marqueur{8} : il y fait allusion à une locution proverbiale\marqueur{9} :
\citComm{`δι᾽ ὀξείας δραμεῖν'· ἐπὶ τῶν διακινδυνευόντων. ὀξεῖαν γὰρ λέγουσι τὴν λόγχην\\
«~Courir en traversant les pointes~» : à propos de ceux qui prennent des risques ; on appelle pointe la lance.}
\end{multicols}
}%
\end{adjustbox}%
\vfill
\end{adjustwidth}
\newpage
% --- PAGE DE DROITE : Traduction française + Commentaire (2/8) ---
\begin{adjustwidth}{-2cm}{-2cm}
\noindent\textbf{\footnotesize Traduction française --- Folio 29 (\texttt{f. 176v}, vv. 1169–1210 [1/4]) \hfill ΕΞΟΔΟΣ}\\[-2pt]
\noindent\rule{16.5cm}{0.5pt}\\[2pt]
{\fontsize{8.8pt}{10.2pt}\selectfont
\ligneEst{}{\francais{à vos maris, à leur retour de l'armée.}}{1169}
\ligneEst{\textbf{Coryphée}}{\francais{Pour ce qui est de nous, sache que nous acquiesçons ;}}{1170}
\ligneEst{}{\francais{mais ce barbare, c'est à toi de le convaincre.}}{}
\ligneEst{\textbf{Euripide}}{\francais{Ça, c'est mon affaire. (À la danseuse ) La tienne, Élaphion,}}{}
\ligneEst{}{\francais{est de songer à t'y prendre comme je t'ai dit chemin faisant.}}{}
\ligneEst{}{\francais{Ainsi donc, passe d'abord à l'autre bout, et reviens en courant sur tes pointes.}}{}
\ligneEst{}{\francais{(Au flûtiste :) Toi, Térédon, joue-nous sur la flûte une Persique. (Térédon joue de la flûte, pendant que la jeune fille danse.)}}{1175}
\ligneEst{\textbf{Archer}}{\francais{(Se réveillant.) Qu'est-ce que ce ronflement? c'est un comos qui me réveille ?}}{}
\ligneEst{\textbf{Euripide}}{\francais{Cette enfant, archer, pensait s'essayer.}}{}
\ligneEst{}{\francais{Elle s'en va danser chez des particuliers.}}{}
\ligneEst{\textbf{Archer}}{\francais{Qu'elle danse, qu'elle s'essaie, je ne m'y oppose pas.}}{}
\ligneEst{}{\francais{Qu'elle est vive! On dirait une puce sur une peau de bique.}}{1180}
}
\vspace{2pt}
\noindent\rule{16.5cm}{0.6pt}\\[2pt]
\noindent\textbf{\small Commentaire philologique, métrique et scénique (vv. 1169–1210 --- partie 2/8)}\\[2pt]
\begin{adjustbox}{max totalheight=15.2cm, max width=16.5cm, keepaspectratio}%
\parbox[t]{16.5cm}{
\setlength{\columnsep}{9pt}%
\begin{multicols}{2}
\fontsize{7.0pt}{8.0pt}\selectfont
\setlength{\parindent}{0pt}%
\setlength{\parskip}{1.2pt}%
\noindent Dans le passage d'Eschyle, καλπάζοντας est un permutant métaphorique -- donc synonyme -- de δραμόντας.\\ Willems \& nombre d'exégètes voulurent rattacher <ἀνα>καλπάζειν à un terme technique de danse : \italique{se déplacer} ou \italique{danser sur les pointes}. Le sens leur a-t-il été suggéré par l'image du cheval au trot ? Dans le célèbre roman de Longus (II, 37), lorsque Daphnis imite le dieu Pan avec ses sabots, il se met à courir
\citComm{ἐπ' ἄκρων τῶν ὀνύχων.}
\noindent Là est le point : car l'on sait le type d'entraînement auquel se soumet la danseuse : au vers 1180, le Scythe qui la voit si \italique{légère} se mettre à \italique{sauter} («~comme une puce~», dixit le Scythe...), il s'exclame : ὡς ἐλαπρός, ὥσπερ ψύλλο κατὰ τὸ κῴδιο.
\noindent Moreau avait donc raison de souligner que la glose présentant notre danseuse aguichant l'Archer \italique{mollement} (ἁβρῶς βαδίζειν) ne tenait pas ; mais l'erreur venait de la glose et non de la correction d'Hermann : on eut le tort de vouloir conserver à tout prix cette glose telle que rapportée par le \italique{Ravennas} et Suidas ; elle n'était pourtant en conformité ni avec l'action ni avec le type d'entraînement -- hyperactif -- de la danseuse.
\noindent Le lexique de Jean Zonaras\marqueur{10} nous en livre la bonne version (simple substitution de ἄ\textbf{κ}ρως à ἁ\textbf{β}ρῶς) :
\citComm{καλπάζειν· τὸ ἄκρως βαδίζειν.}
\noindent Willems avait donc vu juste : l'exercice de la danseuse devait évoquer la course de la κάλπη : commencer par un \italique{trot}, c'est-à-dire par un déplacement en \italique{pointes} (τὸ ἄκρως βαδίζειν) qui s'accélérait pour terminer par un saut ; ce qu'Euripide désigne par le verbe ἀνακαλπάζειν.\\ Cette interprétation nous est confirmée par Eustathe (\italique{Commentaire. à l'}Iliade, XX, v. 228 p. 1206, à propos du mot `σκιρτῷεν') :
\citComm{
ἔστι δὲ σκιρτᾶν οὐ τὸ ἁπλῶς θέειν, ἀλλὰ τὸ ἀνέτως καὶ δίχα ἐπιταγῆς καὶ πηδητικῶς καὶ οἷον ὀρχηστικῶς \textbf{κατὰ τοὺς καλπάζοντας}.\\

«~Bondir~» ((σκιρτᾶν)) n'est pas le simple fait de courir, mais de le faire librement, sans ordre, fougueusement, ainsi que le font les \textbf{καλπάζοντες} en dansant.}
\noindent Terminons par Hésychius (κ 519)\marqueur{11}, qui ne le comprenait pas autrement :
\citComm{καλπάζει· ὀξυποδεῖ.}
\noindent ὀξυποδέω signifie \italique{accélérer le pas}\marqueur{12}. Notons aussi l'autre glose d'Hésychius (σ 822) :
\citComm{σκαλ\{α\}πάζειν· ῥεμβωδῶς βαδίζειν.\\
σκαλπάζειν : se déplacer en tournoyant.}
\noindent\textbf{1175} \textbf{Τερηδὼν,} ἡ αὐλητρίς. (Note du scholiaste ; pourquoi pense-t-il que Térédon soit \italique{une} joueuse de flûte ?)
\noindent\textbf{1175} \textbf{Περσικόν.} Euripide demande au flûtiste d'accompagner la danseuse avec un air \italique{persan} -- les oreilles ? Or il existe une danse du même nom\marqueur{1} que Xénophon nous décrit dans son \italique{Anabase} (VI, 1, 10) :
\citComm{μετὰ τοῦτο Μυσὸς\marqueur{2} εἰσῆλθεν ἐν ἑκατέρᾳ τῇ χειρὶ ἔχων πέλτην, καὶ τοτὲ μὲν ὡς δύο ἀντιταττομένων μιμούμενος ὠρχεῖτο, τοτὲ δὲ ὡς πρὸς ἕνα ἐχρῆτο ταῖς πέλταις, τοτὲ δ᾽ ἐδινεῖτο καὶ ἐξεκυβίστα ἔχων τὰς πέλτας, ὥστε ὄψιν καλὴν φαίνεσθαι. τέλος δὲ \textbf{τὸ περσικὸν ὠρχεῖτο} κρούων τὰς πέλτας καὶ ὤκλαζε καὶ ἐξανίστατο· καὶ ταῦτα πάντα ἐν ῥυθμῷ ἐποίει πρὸς τὸν αὐλόν.\\
Un Mysien\marqueur{2} entra ensuite ; il tenait, dans chacune de ses mains, un bouclier léger ; quelquefois il s'en servait en dansant, comme s'il eût eu à se défendre contre deux adversaires ; quelquefois, comme s'il n'eût eu affaire qu'à un seul. Souvent il tournait et faisait le saut périlleux sans lâcher ses boucliers : c'était un spectacle agréable qu'il donnait ; il finit par \textbf{danser la Persique} ; frappant d'un bouclier sur l'autre ; il s'accroupissait, se relevait, et exécutait tous ces mouvements en mesure et au son de la flûte.}
\end{multicols}
}%
\end{adjustbox}%
\vfill
\end{adjustwidth}
\newpage

% --- FOLIO 29 (f. 176v) : SPREAD 2.2 (Texte grec établi + Traduction française + Commentaire 3/8 & 4/8) ---
\clearpage\ifodd\value{page}\hbox{}\newpage\fi
\enTeteG{ΘΕΣΜΟΦΟΡΙΑΖΟΥΣΑΙ --- Texte grec établi (f. 176v, vv. 1169–1210 [2/4])}
\enTeteD{LES THESMOPHORIES --- Traduction française (f. 176v, vv. 1169–1210 [2/4])}
\begin{adjustwidth}{-2cm}{-2cm}
\noindent\textbf{\footnotesize Texte grec établi --- Folio 29 (\texttt{f. 176v}, vv. 1169–1210 [2/4]) \hfill ΕΞΟΔΟΣ}\\[-2pt]
\noindent\rule{16.5cm}{0.5pt}\\[2pt]
{\fontsize{8.8pt}{10.2pt}\selectfont
\ligneEst{\textbf{Εὐριπίδης}}{φέρε θοἰμάτιον ἄνωθεν ὦ τέκνον τοδί·}{}
\ligneEst{}{καθιζομένη δ᾽ ἐπὶ τοῖσι γόνασι τοῦ Σκύθου}{}
\ligneEst{}{τὼ πόδε πρότεινον, ἵν᾽ ὑπολύσω.}{}
\ligneEst{\textbf{Τοξότης}}{\decalage{τὼ πόδε πρότεινον, ἵν᾽ ὑπολύσω.}ναῖκι ναὶ}{}
\ligneEst{}{κάτησο κάτησο, ναίκι ναί, τυγάτριον.}{}
\ligneEst{}{οἴμ᾽ ὠς στέριπο τὸ τιττί᾽, ὤσπερ γογγυλί.}{1185}
\ligneEst{\textbf{Εὐριπίδης}}{αὔλει σὺ θᾶττον· ἔτι δέδοικας τὸν Σκύθην;}{}
\ligneEst{\textbf{Τοξότης}}{καλό γε τὸ πυγή. κλαῦσί γ᾽ ἢν μὴ 'νδον μένῃς.}{}
\ligneEst{}{εἶεν· καλὴ τὸ σκῆμα περὶ τὸ πόστιον.}{}
\ligneEst{\textbf{Εὐριπίδης}}{καλῶς ἔχει. λαβὲ θοἰμάτιον· ὥρα 'στὶ νῷν}{}
\ligneEst{}{ἤδη βαδίζειν.}{1190}
\ligneEst{\textbf{Τοξότης}}{\decalage{ἤδη βαδίζειν.}οὐκὶ πιλῆσι πρῶτά με;}{}
\ligneEst{\textbf{Εὐριπίδης}}{πάνυ γε· φίλησον αὐτόν.}{}
}
\vspace{2pt}
\noindent\rule{16.5cm}{0.6pt}\\[2pt]
\noindent\textbf{\small Commentaire philologique, métrique et scénique (vv. 1169–1210 --- partie 3/8)}\\[2pt]
\begin{adjustbox}{max totalheight=15.2cm, max width=16.5cm, keepaspectratio}%
\parbox[t]{16.5cm}{
\setlength{\columnsep}{9pt}%
\begin{multicols}{2}
\fontsize{7.0pt}{8.0pt}\selectfont
\setlength{\parindent}{0pt}%
\setlength{\parskip}{1.2pt}%
\noindent L'accroupissement (ὀκλάζω) est le mouvement principal de cette danse et lui a donné son nom ; c'est le scholiaste aux \italique{Thesmophories} qui nous l'apprend :
\citComm{Περσικὸν τὸ ὄκλασμα καλεῖται.}
\noindent Pollux (IV, 100), compulsant un peu vite les \italique{Thesmophories} et ses commentateurs, croit que le terme ὄκλασμα apparaît dans le texte-même d'Aristophane : il range l'ὄκλασμα auprès du κνισμός parmi les danses guerrières et frénétique :
\citComm{καὶ ὄκλασμα· οὕτω γὰρ ἐν Θεσμοφοριαζούσαις ὀνομάζεται τὸ ὄρχημα τὸ Περσικὸν καὶ σύντονον. τὴνδ' αὐτὴν καὶ ὑγρὰν ὠνόμαζον.}
\noindent L'on trouve d'ailleurs l'une des scholies des \italique{Thesmophories} les plus savantes à notre vers 1175\marqueur{3} :
\citComm{Περσικὸν ὄρχημα, περὶ οὗ Ἰόβας μακρὸν πεποίηται λόγον ἐν τοῖς περὶ τῆς θεατρικῆς ἱστορίας. ὥστε λελύσθαι τὴν Σελεύκοὺ πρότασιν. προτείνει γὰρ ἐν τῷ πρὸς Ζήνωνα προτατικῷ τὸ Περσικὸν ὄρχημα.}
\noindent Pour ma part, je doute que l'exercice auquel se livre notre danseuse Ἐλάφιον soit cette danse guerrière illustrée par le Mysien de Xénophon ; l'adjectif Περσικός était assez commun pour qu'il pût désigner à la fois un air de flûte et une forme de danse sans qu'il y eut nécessairement de rapport entre les deux.
\noindent\textbf{[Notes ad 1174]} \textbf{(1)} «~Note sur le vers 1174 des \italique{Thesmophories} d'Aristophane~», \italique{L'antiquité classique}, tome 7, fasc. 2 (1938) pp. 227-231. À peine âgé de vingt ans, le jeune étudiant d'alors fait montre d'une érudition et d'un esprit critique peu communs ; il faut dire qu'il avait pour professeur à l'Université de Bruxelles Henri Grégoire -- le successeur d'Alphonse Willems -- et dont l'assistant n'était autre que Roger Goossens ! bouillon intellectuel dont on imagine sans peine l'extraordinaire émulation qui devait en résulter parmi les élèves.
\noindent \textbf{(2)} \italique{Iliade}, XX, 80. Et l'explication d'Hésychius (κ 3393) : κόλπον ἀνιεμένη· τὸ κατὰ τὸ στῆθος μέρος ἀνέλκουσα, ἵνα τὸν μαστὸν δείξῃ.
\noindent \textbf{(3)} Ou, à la rigueur, «~tenir haut les seins~» comme le traduit Germaine Prudhommeau.
\noindent \textbf{(4)} †ἀνακεκολλαμένοι† dans divers lexiques, notamment le lexique platonicien de Timée le Sophiste ou celui de Photius. La correction est-elle du fait du rédacteur de l'\latin{Etymologicum Magnum} ?
\noindent \textbf{(5)} Ou \italique{lascif}, dira-t-on, au sens que le \nieme{XXI} siècle prête à ce mot ; c'est-à-dire avec \italique{langueur}, un \italique{manque de vigueur} qui trahit un fort penchant pour la sensualité. Est-ce l'influence de \italique{las} qui aura changé radicalement le sens originel de \italique{lascif} (\italique{lubrique}, \italique{folâtre} ; sens que rapportent encore exclusivement nos dictionnaires) ?
\noindent \textbf{(6)} Dans sa recension des \italique{Thesmophories} éditées par Enger en 1844 (\italique{Aristophanis Comoediae cum scholiis}), in \italique{Zeitschrift für altertumswissenschaft} (1845), coll. 905-920. Article réédité par Fritzsche dans les \italique{Godofredi Hermanni Opuscula}, vol. 8 ; Hermann y discute de notre vers à la page 307.
\noindent \textbf{(7)} Au vers 1201, un `α' est lu fautivement de la même manière en `ο'.
\noindent \textbf{(8)} $=$ fr. 144a Radt.
\noindent \textbf{(9)} Transmis par les parémiographes Zénobius (III, 13), Suidas (δ 1186 ; ο 414) et Érasme (2014). Il n'est pas impossible que le proverbe soit lui-même tiré de la danse, et d'une éventuelle \italique{danse des épées}. Sur l'utilisation d'épées pour certaines acrobaties périlleuses, voir \italique{La danse grecque antique} de Maurice Emmanuel (1896), fig. 556 et 557 ou l'article «~Histoire anthropologique des danses acrobatiques~» de	Myriam Peignist, \italique{Corps} (2009/2), n° 7, pp. 29-38. La danseuse se livre-t-elle à ce genre d'acrobaties guerrières pour attirer l'attention du Scythe ? cela paraît peu vraisemblable ; il est préférable d'interpréter le fragment d'Eschyle de façon imagée et ne pas rattacher directement le verbe καλπάζειν à une quelconque danse des épées.
\noindent Un proverbe analogue et que l'on connaît en français, \italique{être sur le fil du rasoir}, apparaît déjà chez Homère (\italique{Iliade} X, 173).
\end{multicols}
}%
\end{adjustbox}%
\vfill
\end{adjustwidth}
\newpage
% --- PAGE DE DROITE : Traduction française + Commentaire (4/8) ---
\begin{adjustwidth}{-2cm}{-2cm}
\noindent\textbf{\footnotesize Traduction française --- Folio 29 (\texttt{f. 176v}, vv. 1169–1210 [2/4]) \hfill ΕΞΟΔΟΣ}\\[-2pt]
\noindent\rule{16.5cm}{0.5pt}\\[2pt]
{\fontsize{8.8pt}{10.2pt}\selectfont
\ligneEst{\textbf{Euripide}}{\francais{Tends-moi ce vêtement, ma fille, en le tirant par en haut.}}{}
\ligneEst{}{\francais{Assieds-toi sur les genoux du Scythe,}}{}
\ligneEst{}{\francais{et avance les pieds, que je les déchausse.}}{}
\ligneEst{\textbf{Archer}}{\francais{\espace{4.5cm}Si fait, si,}}{}
\ligneEst{}{\francais{mets-toi, si fait, ma fillette.}}{}
\ligneEst{}{\francais{(Pendant qu'Euripide ôte les chaussures à la danseuse.) Ouais! est-il ferme, ce téton ! une vraie rave. (La jeune fille déchaussée et nue danse au son de la flûte.)}}{1185}
\ligneEst{\textbf{Euripide}}{\francais{Joue donc plus vite. As-tu encore peur du Scythe ?}}{}
\ligneEst{\textbf{Archer}}{\francais{Ah! que voilà de belles fesses ! — Gare à toi si tu ne restes dans ton coin.}}{}
\ligneEst{}{\francais{(Il s'arrange.) — À la bonne heure ! Elle a repris de la tenue, la broquette.}}{}
\ligneEst{\textbf{Euripide}}{\francais{(À la danseuse). C'est assez. Reprends ton vêtement. Il est temps que}}{}
\ligneEst{}{\francais{nous partions.}}{1190}
\ligneEst{\textbf{Archer}}{\francais{\espace{2.1cm}Sans me donner un baiser?}}{}
\ligneEst{\textbf{Euripide}}{\francais{Très volontiers. (À la danseuse.) Baise-le.}}{}
}
\vspace{2pt}
\noindent\rule{16.5cm}{0.6pt}\\[2pt]
\noindent\textbf{\small Commentaire philologique, métrique et scénique (vv. 1169–1210 --- partie 4/8)}\\[2pt]
\begin{adjustbox}{max totalheight=15.2cm, max width=16.5cm, keepaspectratio}%
\parbox[t]{16.5cm}{
\setlength{\columnsep}{9pt}%
\begin{multicols}{2}
\fontsize{7.0pt}{8.0pt}\selectfont
\setlength{\parindent}{0pt}%
\setlength{\parskip}{1.2pt}%
\noindent \textbf{(10)} Édité par Tittmann sous le titre de \italique{Iohannis Zonarae Lexicon} (1808), p. 1171. Ce lexique de la fin du \nieme{XII} ou début du \nieme{XIII} siècle, s'il faut en croire Klaus Alpers, serait plutôt à attribuer à un certain Nicéphore. Le Pseudo-Zonaras a puisé abondamment dans la \italique{Souda} (ou des sources très voisines), notamment parmi les adresses de la lettre Α à Κ et les gloses touchant à Aristophane. Notre cas est intéressant, puisqu'il est peu vraisemblable que la transposition de β et κ soit une erreur ou correction du copiste ; il faut donc en conclure que si le Pseudo-Zonaras est allé piocher sa glose dans la \italique{Souda}, il disposait d'un meilleur manuscrit que ceux qui nous sont parvenus. Klaus Alpers a consacré un article très riche au lexique du Pseudo-Zonaras dans la monumentale \italique{Realencyclopädie der classischen Altertumswissenschaft}, 10/1 (1972), à l'entrée \allemand{'Zonarae' Lexikon}, coll. 732-763. Il le résume dans \italique{Der Neue Pauly}, 12/2 (2003), col. 831. Pour une étude plus générale de l'ensemble du corpus des anciens lexiques en langue grecque, l'on consultera la synthèse du même auteur dans l'\italique{Historisches Wörterbuch der Rhetorik}, 5 (2001) à l'entrée \allemand{Lexikographie} B, coll. 194-210.
\noindent \textbf{(11)} Contrairement à l'usage de notre commentaire, nous n'avons réservé Hésychius que pour la fin de cette note. La raison en est que la glose de l'unicodex -- un \italique{Marcianus} du \nieme{XV} siècle -- est légèrement corrompue et donc sujette à caution :
\citComm{καλπάγει ὀξυπόδης}
\noindent (Le lexique de Cyrille ajoute : σακκιζει (sic, immo βαδίζει).)
\noindent \textbf{(12)} Cf. le passage de l'\italique{Oreste} d'Euripide, vers 1550, où Ménélas furieux accourt ὀξύπους ($=$ ταχύπους). Ce sens est celui rapporté avec raison par le \LSJ. Le mot latin \latin{addensator}, \italique{celui qui rend plus serré} (en parlant du foulon et de son pas), est glosé dans certains manuscrits gréco-latin ὀξυποδητής.
\noindent \textbf{(1)} Hésychius (π 2005) : περσικά· [ ... ] καὶ ὄρχησις Περσική.
\noindent \textbf{(2)} Je cite longuement le passage de Xénophon pour faire apparaître la race du danseur. Dans le cas où la danse persique aurait effectivement un rapport avec le verbe ἀνακαλπάζειν du vers 1174 -- ce que je ne crois pas --, on notera qu'il faisait partie du peuple d'Asie qui servit de titre à la pièce d'Eschyle, citée par Photius à l'adresse ἀνακαλπάζει et dont on a parlé au vers précédent.
\noindent \textbf{(3)} Elle est discutée par Jacques Moreau dans l'article déjà mentionné : «~Note sur le vers 1174 des \italique{Thesmophories} d'Aristophane~», \italique{L'antiquité classique}, tome 7, fasc. 2 (1938) pp. 227-231.
\noindent\textbf{1176} \textbf{κωμοτίς ἀνεγεῖρί μοι;} Variante de lecture du scholiaste : διχῶς κωμοδρίαν.
\noindent\textbf{1172 et 1180} \textbf{ὦλάφιον / ὡς ἐλαπρός, ὥσπερ ψύλλο κατὰ τὸ κῴδιο.} Les verbes signifiant \italique{sauter}, \italique{danser} (πηδᾶν, ἄλλεσθαι) attirent presque inévitablement une comparaison, et le plus souvent avec un animal.
\noindent Les jeunes filles de Sparte dansent en bondissant comme des pouliches, ἇτε πῶλοι ταὶ κόραι / πὰρ τὸν Εὐρώταν / ἀμπαδίοντι (ἀμπαδίοντι Coulon : ἀμπάλλοντι codd.) (\italique{Lysistr.} 1308)\marqueur{1}.
\noindent La Muse laconienne danse et saute comme une biche, ποδοῖν τε πάδη / ἇ τις ἔλαφος (\italique{ib}. 1317)\marqueur{2}.
\noindent La danseuse des \italique{Thesmophories} porte d'ailleurs le sobriquet de \italique{Bichette} (Ἐλάφιον, v. 1172).
\noindent Bichette danse avec tant de grâce devant l'archer Scythe qu'il s'écrie : « qu'elle est légère ! une vraie puce sur ma peau de mouton » (\italique{Thesm}. 1180). Cette comparaison dénuée de poésie sent le corps de garde et ses peaux de mouton où grouille la vermine (on dormait communément sur des \italique{peaux de mouton}, κῴδια). (Note de Taillardat, §791.)
\noindent\textbf{[Notes ad 1172 et 1180]} \textbf{(1)} Comparaison banale, cf. Eur. \italique{Or}. 45 (à propos d'Oreste furieux): πηδᾷ δρομαῖος, πῶλος ὣς ὑπὸ ζυγοῦ ; Eubul. 75, 6 $[=$ Athénée, III, 108ab$]$. Le πωλίον symbolise la vivacité, Ar. \italique{Guêp}. 188. En outre, πῶλος désigne couramment les jeunes femmes : Anacréon 88 \italique{Diehl} [$=$ 417 \italique{Poetae Melici Graeci} Page$]$ ; Eur. \italique{Héc.} 142, \italique{Hipp}. 516, etc. ; Cratin. 87 $[=$ 94 Kassel-Austin$]$ ; Épicrat. 9, 4; Eubul. 84 $[=$ 82 Kassel-Austin$]$.
\end{multicols}
}%
\end{adjustbox}%
\vfill
\end{adjustwidth}
\newpage

% --- FOLIO 29 (f. 176v) : SPREAD 2.3 (Texte grec établi + Traduction française + Commentaire 5/8 & 6/8) ---
\clearpage\ifodd\value{page}\hbox{}\newpage\fi
\enTeteG{ΘΕΣΜΟΦΟΡΙΑΖΟΥΣΑΙ --- Texte grec établi (f. 176v, vv. 1169–1210 [3/4])}
\enTeteD{LES THESMOPHORIES --- Traduction française (f. 176v, vv. 1169–1210 [3/4])}
\begin{adjustwidth}{-2cm}{-2cm}
\noindent\textbf{\footnotesize Texte grec établi --- Folio 29 (\texttt{f. 176v}, vv. 1169–1210 [3/4]) \hfill ΕΞΟΔΟΣ}\\[-2pt]
\noindent\rule{16.5cm}{0.5pt}\\[2pt]
{\fontsize{8.8pt}{10.2pt}\selectfont
\ligneEst{\textbf{Τοξότης}}{\decalage{πάνυ γε· φίλησον αὐτόν.}ὂ ὂ ὂ παπαπαπαῖ,}{}
\ligneEst{}{ὠς γλυκερὸ τὸ γλῶσσ᾽, ὤσπερ Ἀττικὸς μέλι.}{}
\ligneEst{}{τί οὐ κατεύδει παρ᾽ ἐμέ;}{}
\ligneEst{\textbf{Εὐριπίδης}}{\decalage{τί οὐ κατεύδει παρ᾽ ἐμέ;}χαῖρε τοξότα,}{}
\ligneEst{}{οὐ γὰρ γένοιτ᾽ ἂν τοῦτο.}{}
\ligneEst{\textbf{Τοξότης}}{\decalage{οὐ γὰρ γένοιτ᾽ ἂν τοῦτο.}ναὶ ναὶ γρᾴδιο.}{}
\ligneEst{}{ἐμοὶ κάρισο σὺ τοῦτο.}{1195}
\ligneEst{\textbf{Εὐριπίδης}}{\decalage{ἐμοὶ κάρισο σὺ τοῦτο.}δώσεις οὖν δραχμήν;}{}
\ligneEst{\textbf{Τοξότης}}{ναὶ ναῖκι δῶσι.}{}
\ligneEst{\textbf{Εὐριπίδης}}{\decalage{ναὶ ναῖκι δῶσι.}τἀργύριον τοίνυν φέρε.}{}
\ligneEst{\textbf{Τοξότης}}{ἀλλ᾽ οὐκ ἔκὠδέν· ἀλλὰ τὸ συμβήνην λαβέ.}{}
\ligneEst{}{ἔπειτα κομίζι σ᾽ αὐτίκ᾽. ἀκολούτι τέκνον.}{}
\ligneEst{}{σὺ δὲ τοῦτο τήρει τὴ γέροντο, γρᾴδιο.}{}
}
\vspace{2pt}
\noindent\rule{16.5cm}{0.6pt}\\[2pt]
\noindent\textbf{\small Commentaire philologique, métrique et scénique (vv. 1169–1210 --- partie 5/8)}\\[2pt]
\begin{adjustbox}{max totalheight=15.2cm, max width=16.5cm, keepaspectratio}%
\parbox[t]{16.5cm}{
\setlength{\columnsep}{9pt}%
\begin{multicols}{2}
\fontsize{7.0pt}{8.0pt}\selectfont
\setlength{\parindent}{0pt}%
\setlength{\parskip}{1.2pt}%
\noindent \textbf{(2)} Le cerf (ἔλαφος) est le symbole de l'agilité, \italique{Nuées} 351. (Et aussi, comme chez Homère, de la lâcheté.)
\noindent\textbf{1181} \textbf{φέρε θοἰμάτιον ἄνωθεν.} Sur le texte et le sens de ce vers 1181, voir la note 28 d'Alphonse Willems, reproduite ci-après.
\noindent\textbf{1185} \textbf{ὡς στεριπὸ τὸ τιττί᾽} $=$ ὡς στέριφον τὸ τιτθίον.
\noindent\textbf{1185-93} \textbf{Papyrus d'Oxyrhynque n° 3840.} Un mince débris contenant quelques lettres de notre passage fut déniché parmi les papyrus d'Oxyrhynque\marqueur{1}. Il y a néanmoins quelques informations à en tirer : il confirme la conjecture d'Enger au vers 1185 (γογγυλί à la place de γογγύλῃ), ainsi que μέλι au vers 1192. Parsons note la présence d'une ligne supplémentaire après le vers 1187, là où le \italique{Ravennas} possède l'une de ses rares didascalies\marqueur{2} : ἀνακύπτη· καὶ παρακύπτι· ἀπεψωλημένος· (même jeu de scène avec l'olisbos qu'aux vers 643-648).
\noindent\textbf{[Notes ad 1185-93]} \textbf{(1)} \italique{The Oxyrhynchus Papyri}, vol. LVI (1989), n°3840. Étudié par Peter J. Parsons, aidé de Colin Austin -- déjà en train de préparer son édition des \italique{Thesmophories}. La transcription qu'en fit Parsons et son analyse sont recopiées en annexe. Voir aussi le papyrus d'Oxyrhynque n° 5132 (\italique{The Oxyrhynchus Papyri} LXXVIII) pour un semblant d'\italique{addendum} : la présence de doubles points `$ : $', comme au vers 1190, pour indiquer un changement de locuteur est confirmée par un autre fragment du même papyrus au vers 1209.
\noindent \textbf{(2)} Pour notre pièce où les scholies très anciennes et érudites se font rares, on ne relève d'autres didascalies qu'aux vers 100, 129 et 276.
\noindent\textbf{1187} \textbf{κλαῦσί γ᾽ ἂν μὴ 'νδον μένῃς.} πρὸς τὸ πέος λέγει ἐπιπλήττων ἐντεινομένῳ. (Note du scholiaste.)\\ L'archer suit d'un œil émérillonné tous les mouvements de la danseuse. Son agitation se trahit d'une manière par trop visible. Le Scythe, qui s'en aperçoit, reprend « gare à toi, etc .. »
\noindent\textbf{1192} \textbf{ὠς γλυκερὸ τὸ γλῶσσ᾽, ὤσπερ Ἀττικὸς μέλις.} Si l'image des baisers plus doux que le miel\marqueur{1} ou le nectar nous paraît banale, c'est à une longue tradition héritée de ces culs-terreux de latins que nous le devons\marqueur{2} ; pour ce qui est des auteurs de langue grecque, on ne peut guère citer que Alciphron (I, 38, 7), postérieur d'environ six cents ans à Aristophane.
\noindent Les spectateurs d'Aristophane devaient ressentir avec cette comparaison une confusion burlesque : l'épithète poétique μελίγλωσσος existe en grec, mais l'image d'une \italique{langue pleine de miel} signifie tout autre chose.
\noindent Ainsi parlait Taillardat (§§739-743) :\\ «~Le miel étant le symbole de la douceur\marqueur{3}, les poètes se plaisent à dire que leurs chants sont de miel\marqueur{4}. Cette image est chez Aristophane : μελιγλώσσων ἐπέων ... ἀοιδάν (\italique{Ois}. 908) donne un exemple de plus de l'épithète banale μελίγλωσσος\marqueur{5} $[$ ... $]$ On dira donc tout naturellement que la lèvre d'un poète est enduite de miel : c'est le cas de la bouche de Sophocle, Σοφοκλέους τοῦ μέλιτι κεχριµένου... τὸ στόμα (\italique{fr}. Aristophane 581 $[=$ 598 Kassel-Austin$]$). Composant des poèmes doux comme le miel, les poètes aiment à se comparer aux abeilles, ainsi qu'en témoigne Platon (\italique{Ion}, 534b).~»\\ Aristophane fait rire son public aux dépens de ces images trop usées : le Scythe rustique vante la \italique{langue douce comme le miel} de la danseuse, mais c'est une langue toute matérielle -- et qui appartient d'ailleurs à une actrice muette. Il y a des impropriétés de termes qui sont des malpropretés du cœur.
\noindent\textbf{[Notes ad 1192]} \textbf{(1)} Et l'Attique est souvent citée pour son miel. Cf. la \italique{Paix}, vers 252-254 ; Antiphane fr. 177(179) Kassel-Austin $=$ Athénée II, 43b \& III, 74d ; le corpus hippocratique mentionne maintes fois du miel Attique.
\noindent \textbf{(2)} Cf. le \italique{Dictionnaire des sentences latines et grecques} de Renzo Tosi, §1385 \latin{Nectar dulcius} : Catulle 99, 1-2 ; Horace, \italique{Carmina} 1, 13, 14-16 ; \&c. Honoré de Balzac a fameusement repris cette longue tradition : « Les baisers d'une femme sincère ont un miel divin qui semble mettre dans cette caresse une âme, un feu subtil par lequel le cœur est pénétré. Les baisers dénués de cette onction savoureuse sont âpres et secs. »
\end{multicols}
}%
\end{adjustbox}%
\vfill
\end{adjustwidth}
\newpage
% --- PAGE DE DROITE : Traduction française + Commentaire (6/8) ---
\begin{adjustwidth}{-2cm}{-2cm}
\noindent\textbf{\footnotesize Traduction française --- Folio 29 (\texttt{f. 176v}, vv. 1169–1210 [3/4]) \hfill ΕΞΟΔΟΣ}\\[-2pt]
\noindent\rule{16.5cm}{0.5pt}\\[2pt]
{\fontsize{8.8pt}{10.2pt}\selectfont
\ligneEst{\textbf{Archer}}{\francais{\espace{3.9cm}Oh! oh! oh! papapapai!}}{}
\ligneEst{}{\francais{Que sa langue est douce. C'est comme du miel attique.}}{}
\ligneEst{}{\francais{Pourquoi ne couche-t-elle pas avec moi ?}}{}
\ligneEst{\textbf{Euripide}}{\francais{\espace{3.7cm}Adieu, archer,}}{}
\ligneEst{}{\francais{cela ne se peut.}}{}
\ligneEst{\textbf{Archer}}{\francais{\espace{3.7cm}Si, si, petite vieille,}}{}
\ligneEst{}{\francais{fais-moi ce plaisir.}}{1195}
\ligneEst{\textbf{Euripide}}{\francais{\espace{3.3cm}Me donneras-tu une drachme ?}}{}
\ligneEst{\textbf{Archer}}{\francais{Si fait, si, je te la donnerai.}}{}
\ligneEst{\textbf{Euripide}}{\francais{\espace{2.4cm}Voyons ton argent.}}{}
\ligneEst{\textbf{Archer}}{\francais{Je n'en ai pas. Mais prends ma lance. (Il la lui confie.)}}{}
\ligneEst{}{\francais{Je te l'apporte sur-le-champ. Suis-moi, mon enfant.}}{}
\ligneEst{}{\francais{(À Euripide :) Toi, petite vieille, garde ce bonhomme.}}{}
}
\vspace{2pt}
\noindent\rule{16.5cm}{0.6pt}\\[2pt]
\noindent\textbf{\small Commentaire philologique, métrique et scénique (vv. 1169–1210 --- partie 6/8)}\\[2pt]
\begin{adjustbox}{max totalheight=15.2cm, max width=16.5cm, keepaspectratio}%
\parbox[t]{16.5cm}{
\setlength{\columnsep}{9pt}%
\begin{multicols}{2}
\fontsize{7.0pt}{8.0pt}\selectfont
\setlength{\parindent}{0pt}%
\setlength{\parskip}{1.2pt}%
\noindent \textbf{(3)} \italique{Guêpes} 878 (voir §361), \italique{Thesm.} 1192, fr. 587, 622 $[=$ 681, 633 K.-A.$]$ ; Alexis 146, 6 $[=$ 150.6 K.-A.$]$ ; Ménand. 708 (?), etc.
\noindent \textbf{(4)} Pind. \italique{Ol.} 11, 4 : μελιγάρυες ὕμνοι ; id. \italique{Ném.} 11,18 : μελιγδούποισι ... μελιζέμεν ἀοιδαῖς ; id. \italique{Isthm.} 2, 32 : μελικόμπων ἀοιδᾶν ; Alcman 41 $[=$ Athénée XIV, 633a $=$ fr. adesp. 954b \italique{Poetae Melici Graeci}, ed. Page$]$ : μέλεα μελιπτέρωτα ; Simonid. 40 D $[=$ fr 595.3 \italique{Poetae Melici Graeci} Page$]$ : μελιαδέα γᾶρυν, etc. On notera que, déjà chez Homère (\italique{Il.} I, 249), la voix de Nestor est dite $[$ἀπὸ γλώσσης$]$ μέλιτος γλυκίων. D'autre part Pindare, \italique{Ol.} 7, 8, appelle encore la poésie un nectar, νέκταρ.
\noindent \textbf{(5)} Eschyl. \italique{Prom.} 172 ; Bacchyl. 3, 97 ; fr. 4, 25 Snell (μελιγλώσσων ἀοιδᾶν ἄνθεα).
\noindent\textbf{1195} \textbf{ἐμοὶ κάρισο σὺ τοῦτο} $=$ ἐμοὶ χάρισαι σὺ τοῦτο.
\noindent\textbf{1197-216} \textbf{ἀλλὰ τὸ συμβήνην λαβέ $[$ ... $]$ ὀρτῶς δὲ συβήνη 'στί· καταβηνῆσι γάρ.} Ce passage a été si maltraité par les commentateurs -- anciens comme modernes -- que nous nous sommes cru obligé à commettre cette note rectificatrice. \italique{Recte}, descendons-donc dans l'arène (ὀρτῶς δὲ συμβήνη 'στί· καταβηνῆσι γάρ.)\\ L'origine de tous les malentendus est la glose d'un commentateur ancien qui sera allé un peu vite : le \italique{Ravennas} comporte deux scholies pour définir la συμβήνην (la première scholie se trouve à l'intérieur même du texte) :
\citComm{τὴν τοξοθήκην\\
συβήνη· αὐλοθήκη. λέγουσι δὲ καὶ τὸν φαρετρεῶνα συβήνην.}
\noindent On retrouve cette scholie chez Hésychius (σ 2136) :
\citComm{συβήνη· αὐλοθήκη, ἢ τοξοθήκη. ἢ ὁ ναυτικὸς χιτών.}
\noindent Suidas cite notre passage dans une glose qu'on croise à l'identique, ou peu s'en faut, dans d'autres lexiques\marqueur{1}:
\citComm{Συβήνη: ἡ δερματίνη αὐλοθήκη\marqueur{2}. ἢ ἡ φαρέτρα. Θεσμοφοριαζούσαις Ἀριστοφάνης· ἀλλ' οὐκ ἔχω οὐδέν. ἀλλὰ συβήνην λάβε.}
\noindent Cette unanimité ne doit pas nous tromper : elle prouve seulement l'ancienneté de la première rédaction qui aura diversement infusé dans de nombreux lexiques. Je dirais même plus : ces gloses concordent si bien qu'il est certain qu'elles remontent toutes à une même source -- certes ancienne -- mais unique, et donc faillible.\\ Que le terme συβήνη -- qui n'est d'ailleurs pas celui transmis par notre manuscrit comportant συ\textbf{μ}βήνην --, que συβήνη ait pu désigner un \italique{étui à flûte}, je n'en disconviens pas. À la rigueur, que par dérivation le mot ait aussi pu être appliqué à un \italique{carquois}, pourquoi pas\marqueur{3} ? Mais il serait absurde de lui donner ce sens dans notre passage : croit-on que l'Archer, tout Scythe qu'il fût, pensait un seul instant qu'une vieille maquerelle pût accepter son carquois dont elle n'aurait que faire au lieu d'une drachme sonnante et trébuchante ? ce serait tout bonnement ridicule\marqueur{4} ! Le Scythe, Euripide, Aristophane, le savent bien ; en fait n'importe quel homme, pourvu qu'il ne fût pas versé dans l'exégèse d'Aristophane :
\citComm{Chez la belle Suzon, pas d'argent, pas de cuisse...\marqueur{5}}
\noindent En réalité, le Scythe demande à Euripide de garder son prisonnier quelques instants, le temps qu'il aille chercher l'argent nécessaire.
\noindent Encore une fois, c'est Hésychius (σ 2137) qui détenait la solution :
\citComm{συμβίνης· καπροβόλον. ἐμβόλιον.}
\noindent Le mot est apparenté à σιβύνη, une \italique{lance} (Hésychius nous précise pour la chasse au \italique{sanglier}, κάπρος), dont il est une légère déformation qui prépare le jeu de mots scabreux du vers 1215.
\noindent Le texte est aussi clair qu'un ciel d'azur au milieu de l'été : le Scythe confie sa lance à Euripide pour la garde de Mnésiloque.
\end{multicols}
}%
\end{adjustbox}%
\vfill
\end{adjustwidth}
\newpage

% --- FOLIO 29 (f. 176v) : SPREAD 2.4 (Texte grec établi + Traduction française + Commentaire 7/8 & 8/8) ---
\clearpage\ifodd\value{page}\hbox{}\newpage\fi
\enTeteG{ΘΕΣΜΟΦΟΡΙΑΖΟΥΣΑΙ --- Texte grec établi (f. 176v, vv. 1169–1210 [4/4])}
\enTeteD{LES THESMOPHORIES --- Traduction française (f. 176v, vv. 1169–1210 [4/4])}
\begin{adjustwidth}{-2cm}{-2cm}
\noindent\textbf{\footnotesize Texte grec établi --- Folio 29 (\texttt{f. 176v}, vv. 1169–1210 [4/4]) \hfill ΕΞΟΔΟΣ}\\[-2pt]
\noindent\rule{16.5cm}{0.5pt}\\[2pt]
{\fontsize{8.8pt}{10.2pt}\selectfont
\ligneEst{}{ὄνομα δέ σοι τί ἔστιν;}{1200}
\ligneEst{\textbf{Εὐριπίδης}}{\decalage{ὄνομα δέ σοι τί ἔστιν;}Ἀρτεμισία.}{}
\ligneEst{\textbf{Τοξότης}}{μεμνῆσι τοίνυν τοὔνομ᾽· Ἀρταμουξία.}{}
\ligneEst{\textbf{Εὐριπίδης}}{Ἑρμῆ δόλιε ταυτὶ μὲν ἔτι καλῶς ποιεῖς.}{}
\ligneEst{}{σὺ μὲν οὖν ἀπότρεχε παιδάριον τουτὶ λαβών·}{}
\ligneEst{}{ἐγὼ δὲ λύσω τόνδε. σὺ δ᾽ ὅπως ἀνδρικῶς}{}
\ligneEst{}{ὅταν λυθῇς τάχιστα φεύξει καὶ τενεῖς}{1205}
\ligneEst{}{ὡς τὴν γυναῖκα καὶ τὰ παιδί᾽ οἴκαδε.}{}
\ligneEst{\textbf{Μνησίλοχος}}{ἐμοὶ μελήσει ταῦτά γ᾽ ἢν ἅπαξ λυθῶ.}{}
\ligneEst{\textbf{Εὐριπίδης}}{λέλυσο. σὸν ἔργον, φεῦγε πρὶν τὸν τοξότην}{}
\ligneEst{}{ἥκοντα καταλαβεῖν.}{}
\ligneEst{\textbf{Μνησίλοχος}}{\decalage{ἥκοντα καταλαβεῖν.}ἐγὼ δὴ τοῦτο δρῶ.}{}
\ligneEst{\textbf{Τοξότης}}{ὦ γρᾴδι᾽ ὠς καρίεντό σοι τὸ τυγάτριο,}{1210}
}
\vspace{2pt}
\noindent\rule{16.5cm}{0.6pt}\\[2pt]
\noindent\textbf{\small Commentaire philologique, métrique et scénique (vv. 1169–1210 --- partie 7/8)}\\[2pt]
\begin{adjustbox}{max totalheight=15.2cm, max width=16.5cm, keepaspectratio}%
\parbox[t]{16.5cm}{
\setlength{\columnsep}{9pt}%
\begin{multicols}{2}
\fontsize{7.0pt}{8.0pt}\selectfont
\setlength{\parindent}{0pt}%
\setlength{\parskip}{1.2pt}%
\noindent Et nous nous garderons bien de biffer le μ de la leçon du manuscrit, συ\textbf{μ}βήνην, qui aura survécu à plus de mille ans de tradition manuscrite mais pas à la griffe implacable des éditeurs d'Aristophane.\\ Le début du vers suivant, altéré, ἔπειτα †κ\inRas{ω}{ο}μίζεις αὐτοῖς†, se laisse lire sans difficulté :
\citComm{ἔπειτα κομίζι ($=$ κομίζω) σ(οὶ) αὐτίκ(α)\\
Ensuite je te l'apporte $[$la drachme$]$ sur-le-champ.}
\noindent La dernière version du manuscrit attribue, il est vrai, cette réplique à Euripide. Ces indications marginales de personnages ne doivent pas nous arrêter : elles sont récentes et même pour certaines du \nieme{XVI} siècle\marqueur{6} (c'est le cas de cette annotation, «~γϋ.~», que porte le manuscrit (la même qu'au vers 1189) et mis pour γυνή -- c'est-à-dire la \italique{vieille}, γραῦς ; c'est-à-dire Euripide. Von Velsen ne l'a pas intégrée à son apparat critique et avait donc reconnu l'écriture d'Euphrosynus Boninus).\\ Notre émendation appelle quelque justification : l'hypothétique confusion αυτικ / αυτοις peut aisément s'expliquer en constatant la ressemblance entre le digramme `ις' et le `κ'\marqueur{7}. Continuons.\\ Le Scythe part chercher l'argent, accompagné de la jeune beauté, qui se montre si caressante selon les instructions d'Euripide (vers 1172 ; l'attitude de la danseuse est la version muette de celle de la Myrrhine dans \italique{Lysistrata}) qu'il en revient tout émoustillé et en proie à une promesse de félicité qui point\marqueur{8}. Il déchante très vite : Euripide et Mnésiloque ont pris la poudre d'escampette. Élaphion décampe à ce moment sous les imprécations du Scythe (ἀπότρεκ᾽ ὠς τάκιστα σύ).
\noindent Le Scythe s'arme donc de sa lance, laissée sur place, et s'apprête à descendre (καταβαίνειν) à leur poursuite :
\citComm{ὀρτῶς δὲ συμβήνη ᾽στί· καταβηνῆσι γάρ.}
\noindent Je n'explique pas le trait d'esprit -- d'esprit \italique{Scythe} ; si vous ne le comprenez pas, laissez-là Aristophane et allez faire des mathématiques !\\ Nul besoin de retailler le vers : καταβαίνω est le terme consacré lorsqu'un acteur descend de la σκηνή vers l'ὀρχήστρα, d'autant plus adapté quand la scène se situe sur un lieu en hauteur comme la colline de la Pnyx\marqueur{9} ; cf. les \italique{Guêpes} (vers 1514) ou l'\italique{Assemblée des femmes} (vers 1152)\marqueur{10}.
\noindent\textbf{[Notes ad 1197-216]} \textbf{(1)} Photius (σ 660), \latin{Etymologicum magnum}, \latin{genuinum}, Zonaras, \&c. Pollux (VII, 153 ; X, 153) : καὶ τὸ μὲν ἀγγεῖον τὸ τὰς γλώττας ἔχον γλωττοκομεῖον, ἣ δὲ τῶν αὐλῶν θήκη συβήνη. Pour un dossier plus complet, je renvoie le lecteur à la \allemand{Realencyclopädie} : \italique{Paulys Realencyclopädie der classischen Altertumswissenschaft}, IV 1 (1931), coll. 1011-1012.
\noindent \textbf{(2)} Dans certains manuscrits du \italique{Zonarae Lexicon} (p. 1686), on trouve βελοθήκη en place d'αὐλοθήκη.
\noindent \textbf{(3)} Encore qu'il me semble plus vraisemblable que ce sens ait été inféré par notre scholiaste, puis scrupuleusement intégré à tous les lexiques.
\noindent \textbf{(4)} À tous les vulgarisateurs d'Aristophane, tournebouleurs et avilisseurs de ses onze chefs-d'œuvre, rappelons la maxime de bon goût qu'Aristophane s'imposait :
\citComm{ὡς ἐγὼ φοβοῦμαι περὶ τῶν μελλόντων ῥηθήσεσθαι, οὔ τι μὴ γελοῖα εἴπω — τοῦτο μὲν γὰρ ἂν κέρδος εἴη καὶ τῆς ἡμετέρας μούσης ἐπιχώριον — ἀλλὰ μὴ καταγέλαστα.\\
Ce n'est pas de provoquer le rire que je redoute -- cela serait tout gain, et conforme à ma muse -- mais de dire des choses ridicules.\\
(Trad. P. Boutang.)}
\noindent (Tradition reçue en droite ligne d'Aristodème, témoin des événements, qui en fit le récit à Apollodore, qui lui-même le transmit à Glaucon, frère de Platon, qui -- on peut l'imaginer ainsi -- confia bientôt, après peut-être encore quelques méandres, à son frère Platon le précieux compte-rendu, lequel Platon finit par coucher par écrit cette fameuse nuit de \italique{Banquet}, 189b, quelques quarante années après les faits.)
\end{multicols}
}%
\end{adjustbox}%
\vfill
\end{adjustwidth}
\newpage
% --- PAGE DE DROITE : Traduction française + Commentaire (8/8) ---
\begin{adjustwidth}{-2cm}{-2cm}
\noindent\textbf{\footnotesize Traduction française --- Folio 29 (\texttt{f. 176v}, vv. 1169–1210 [4/4]) \hfill ΕΞΟΔΟΣ}\\[-2pt]
\noindent\rule{16.5cm}{0.5pt}\\[2pt]
{\fontsize{8.8pt}{10.2pt}\selectfont
\ligneEst{}{\francais{Mais comment te nommes-tu ?}}{1200}
\ligneEst{\textbf{Euripide}}{\francais{\espace{3.5cm}Artémisia.}}{}
\ligneEst{\textbf{Archer}}{\francais{Je le retiendrai, ce nom : Artamouxia. (Il s'en va avec la danseuse.)}}{}
\ligneEst{\textbf{Euripide}}{\francais{Hermès, dieu de la ruse, jusqu'ici tu régis tout à souhait.}}{}
\ligneEst{}{\francais{Toi, mon petit, enfuis-toi, avec ces instruments. (Il lui donne ses pectis.)}}{}
\ligneEst{}{\francais{Pendant ce temps, je le détacherai.}}{}
\ligneEst{}{\francais{(À Mnésiloque :) Quant à toi, dès que tu seras libre, vois à fuir résolument}}{1205}
\ligneEst{}{\francais{et cours retrouver chez toi, ta femme et tes enfants.}}{}
\ligneEst{\textbf{Mnésiloque}}{\francais{J'en fais mon affaire, une fois délivré.}}{}
\ligneEst{\textbf{Euripide}}{\francais{Sois libre. Le reste te regarde : fuis avant que l'archer}}{}
\ligneEst{}{\francais{ne revienne et ne t'empoigne.}}{}
\ligneEst{\textbf{Mnésiloque}}{\francais{\espace{2.9cm}C'est ce que je fais. (Euripide et Mnésiloque s'enfuient.)}}{}
\ligneEst{\textbf{Archer}}{\francais{(Revenant en scène). Ah! petite vieille, la charmante fillette que tu as là,}}{1210}
}
\vspace{2pt}
\noindent\rule{16.5cm}{0.6pt}\\[2pt]
\noindent\textbf{\small Commentaire philologique, métrique et scénique (vv. 1169–1210 --- partie 8/8)}\\[2pt]
\begin{adjustbox}{max totalheight=15.2cm, max width=16.5cm, keepaspectratio}%
\parbox[t]{16.5cm}{
\setlength{\columnsep}{9pt}%
\begin{multicols}{2}
\fontsize{7.0pt}{8.0pt}\selectfont
\setlength{\parindent}{0pt}%
\setlength{\parskip}{1.2pt}%
\noindent \textbf{(5)} Le proverbe d'origine détourné par Brassens (\italique{Grand-père}) étant «~point d'argent, point de Suisse~». Quel sera l'homme d'assez bonne volonté pour nous offrir le colossal équivalent brassinien du maître-œuvre de Taillardat, quel est le preux qui prendra à bras-le-corps les quelques deux cents chansons de Brassens que ce dernier a bien voulu nous conserver ? Espérons qu'il ne soit pas besoin d'attendre encore vingt-cinq siècles comme pour Aristophane. (C'est quand le vers est troussé qu'il faut l'analyser.) Je tiens à la disposition de cette bonne âme éventuelle un matériau qui poserait de solides fondations à ces \italique{Images de Brassens} en fomentation.
\noindent \textbf{(6)} Cf. \italique{supra}, v. 1059 n. 5.
\noindent \textbf{(7)} Il faudrait alors supposer l'existence d'une forme intermédiaire αυτιις ou αυτις, bien vite corrigée en αυτοις. Une forme intermédiaire, ou un \italique{chaînon manquant}, pour reprendre le jargon des darwiniens. Ne poussons donc pas plus loin la «~preuve~» afin d'éviter l'écueil de ces demi-savants, capables de tout prouver et même son contraire !
\noindent \textbf{(8)} Savez-vous ce qu'on peut lire dans le commentaire le plus universitairement reçu, le plus académiquement académisant -- celui d'Austin \& Olson, puisqu'il faut le nommer ? Je vous le donne en mille : le Scythe a consommé la belle Élaphion... il revient donc, précisent leurs auteurs, avec l'\italique{olisbos} abaissé (!!!) Quand on sait -- ou qu'on ne sait pas, comme Austin \& Olson -- combien Aristophane adorait ses chères actrices, qu'il ne manquait jamais de rendre grâce à leur beauté, toujours plein d'attention et de prévenance vis-à-vis de la voracité de la lie athénienne... un archer Scythe... la lie de la lie... suggérer une telle monstruosité ?! «~Paraît qu'elle s'vend même à des flics. / Quelle décadence ! / Y'a plus d'moralité publique / Dans notre France...~» En France peut-être, mais au théâtre d'Aristophane, certainement pas ! (Lisez la note de Willems, t. III pp. 394-395, in « Le nu dans la Comédie ancienne des Grecs », article reproduit en annexe ; purgatif si nécessaire quand on est souillé par de telles salauderies !) A\&O ne nous indiquent pas si le public pouvait se délecter en sus de bruitages émanant des coulisses, quand ça y fornique... Mais pour qui prennent-ils Aristophane ?.. voici donc la littérature qu'on inflige à nos pauvres étudiants... je tiens le cas pour pendable !
\citComm{et si, par un malheur, j'en avois fait autant,\\
je m'irois, de regret, pendre tout à l'instant.}
\noindent \textbf{(9)} Le Scythe s'interrompt quelques instants de parler, pause soulignée par l'interjection \italique{extra metrum} οἴμοι qu'il lance une fois descendu : il se demande vers laquelle des deux entrées de l'orchestre se tourner (τί δρᾶσι; ποῖ τὸ γρᾴδι᾽;)
\noindent \textbf{(10)} \italique{Confer} surtout la note d'Alphonse Willems au vers 1341 des \italique{Guêpes}. Lisez Willems ! aucune excuse : la note est reproduite en annexe.
\noindent\textbf{1201} \textbf{Ἀρταμουξία.} L'archer scythifie le nom d'Ἀρτεμισία : il l'entend comme un composé de ἀρτα-, multiplie les `α' et ajoute un `ξ'\marqueur{1}. Cf. le parler de Ψευδ\textbf{αρτά}βας dans les \italique{Acharniens} (v. 100) :
\citComm{ἰαρταμὰν ἐξάρξαν ἀπισσόνα σάτρα.}
\noindent\textbf{[Notes ad 1201]} \textbf{(1)} Voir le commentaire vv. 91-2 d'Olson aux \italique{Acharniens} ainsi que l'article «~The Archer Scene in Aristophanes' \italique{Thesmophoriazusae}~» d'Édith M. Hall, in \italique{Philologus} 133 (1989), p. 40. On peut citer Ἀρταβάτας, Ξέρξης, Ἀρταξέρξης, Ἀρταφρένης, Ἀρτάβαζος \&c.
\noindent\textbf{1203-11} \textbf{Papyrus d'Oxyrhynque n° 4935 fr. 2 \& n° 5132.} Deux fragments d'un même papyrus d'Oxyrhynque\marqueur{1} contiennent quelques lettres de ce passage. (Voir en annexe.) Ils n'apportent aucune leçon originale. Au vers 1208, l'espace manque pour être compatible avec la leçon \italique{ante correctionem} λέλυσον du \italique{Ravennas} ; l'état du fragment ne permet pas pour autant de trancher entre les corrections λέλυσο (\Rpc) et λέλυσαι (Bentley).
\noindent\textbf{[Notes ad 1203-11]} \textbf{(1)} N° 4935 fr. 2, in \italique{The Oxyrhynchus Papyri}, LXXIII (2009), ed. A. Benaissa. N° 5132, \italique{The Oxyrhynchus Papyri}, LXXVIII (2012), ed. W. B. Henry.
\end{multicols}
}%
\end{adjustbox}%
\vfill
\end{adjustwidth}
\newpage
\clearpage\ifodd\value{page}\hbox{}\newpage\fi
\enTeteG{ΕΞΟΔΟΣ (Τέλος) --- Codex Ravennas 429, f. 177r (Classense 179r) --- vv. 1211–1231}
\begin{adjustwidth}{-1.5cm}{-3.0cm}
\begin{adjustbox}{max totalheight=21.0cm, max width=16.5cm, keepaspectratio}
\begin{minipage}{16.5cm}
{\fontsize{7.2pt}{8.5pt}\selectfont
\noindent\textbf{Folio 30 --- Codex Ravennas 429, \texttt{f. 177r} (\texttt{Classense 179r}) : ΕΞΟΔΟΣ (Τέλος), vv. 1211–1231}\hfill \textit{Témoins :} $\text{R}^1, \Rpc, \Rdeux$ (\Mon)\par\vspace{2pt}
\hrule height 0.4pt\vspace{2.5pt}
}
{\fontsize{8.1pt}{9.5pt}\selectfont
\ligneRav{1211}{τοξ/}{κοὐ \difference{δύσκολ᾽} ἀλλὰ πρᾶο. ποῦ τὸ γρᾴδιο;}{\textbf{1211} δύσκολλ᾽ R : \lecon{δύσκολ᾽} Grynaeus}
\ligneRav{}{}{οἴμ᾽ ὠς ἀπόλωλο· ποῦ τὸ γέροντ᾽ ἐντευτενί;}{}
\ligneRav{1213}{}{ὦ γρᾴδι᾽, ὦ \difference{γρᾶ᾽}. οὐκ ἐπαινῶ γρᾴδιο·}{\textbf{1213} γρᾶίο R : \lecon{γρᾶ᾽} Bentley}
\ligneRav{}{}{Ἀρταμουξία·}{}
\ligneRav{1214}{}{διέβαλλέ μ᾽ \difference{ἡ γραῦς}. ἀπότρεκ᾽ ὠς τάκιστα σύ·}{\textbf{1214} \lecon{διέβαλλε} R : διέβαλέ Σ, Suid. δ 892, Eust. p. 1406.19 \espApp ὁ γραῦς R : \lecon{ἡ γραῦς} Suid., Eust. \espApp ὡς τάχιστα R : ὠς τάκιστα Brunck}
\ligneRav{1215}{}{ὀρτῶς δὲ \difference{συμβήνη} 'στί· καταβηνῆσι γάρ·}{\textbf{1215} συβήνη R : \lecon{συμβήνη} Bibi, cf. v.1197, Hesych. σ 2137}
\ligneRav{1216}{}{οἴμοι, τί \difference{δρᾶσι}; ποῖ τὸ \difference{γρᾴδιο};}{\textbf{1216} δράσει R : \lecon{δρᾶσι} Blaydes (1842) \espApp γράδιο R : \lecon{γρᾴδιο} \Mon}
\ligneRav{}{}{Ἀρταμουξία·}{}
\ligneRav{1217}{χορ/}{τὴν \difference{γραῦν} ἐρωτᾷς, ἣ 'φερεν τὰς πηκτίδας;}{\textbf{1217} γραῦ R : \lecon{γραῦν} Suid. π 1502}
\ligneRav{}{τοξ/}{ναὶ ναῖκι. εἶδες αὐτό; : ταύτῃ γ᾽ οἴχεται}{}
\ligneRav{1219}{}{αὐτή τ᾽ ἐκείνη καὶ \difference{γέρων} τις εἵπετο·}{\textbf{1219} γέρον R : \lecon{γέρων} \Mon}
\ligneRav{1220}{τοξ/}{κροκῶτ᾽ ἔκοντο τὴ γέροντο; : φήμ᾽ ἐγώ·}{}
\ligneRav{}{}{ἔτ᾽ ἂν καταλάβοις, εἰ διώκοις ταυτῃί·}{}
\ligneRav{1222}{τοξ/}{ὦ μιαρὸ \difference{γρᾷο}· πότερα τρέξι τὴν \difference{ὀδό};}{\textbf{1222} γραῦ et ὁδό R : \lecon{γρᾷο} et \lecon{ὀδό} Brunck}
\ligneRav{}{}{Ἀρταμουξία·}{}
\ligneRav{}{χορ/}{ὀρθὴν ἄνω δίωκε. ποῖ θεῖς; οὐ πάλιν}{}
\ligneRav{1224}{}{τῃδὶ \difference{διώξει}; \difference{᾽ς} τοὔμπαλιν τρέχεις σύ γε·}{\textbf{1224} δ᾽ ὤξεις R : \lecon{διώξει} Elmsley (ad \italique{Ach.} 278) \lecon{᾽ς} Cobet (\italique{Mnemos.} 2 (1853), p.213)}
\ligneRav{1225}{τοξ/}{κακόδαιμον, ἀλλὰ τρέξι. Ἀρταμουξία·}{}
\ligneRav{1226}{χορ/}{\hspace{1.1em}τρέχε νυν \difference{κατ᾽ αὐτοὺς} \difference{ἐς} κόρακας ἐπουρίσας·}{\textbf{1226} κατὰ τοὺς R : \lecon{κατ᾽ αὐτοὺς} Jackson (p.82) \lecon{ἐς} O. Bachmann}
\ligneRav{1227}{}{\hspace{1.1em}ἀλλὰ \difference{πέπαισται} μετρίως ἡμῖν·}{\textbf{1227} πέπυσται R : \lecon{πέπαισται} Σ}
\ligneRav{}{}{\hspace{1.1em}ὥσθ᾽ ὥρα δή 'στι βαδίζειν}{}
\ligneRav{}{}{\hspace{1.1em}οἴκαδ᾽ ἑκάστῃ. τὼ Θεσμοφόρω δ᾽}{}
\ligneRav{1230}{}{\hspace{1.1em}ἡμῖν ἀγαθὴν}{}
\ligneRav{1231}{}{\hspace{1.1em}τούτων χάριν ἀνταποδοίτην·}{\textbf{1231} ἀνταδοῖτον R : \lecon{ἀνταποδοῖτον} Grynaeus}
\vspace{3pt}\hrule height 0.3pt\vspace{2pt}
\centerline{\footnotesize\textbf{$\sim\;\cdot\;\sim\;\cdot\;\sim\;\cdot\;\sim\;\cdot\;\sim$ (\italique{korônis} ornementale de clôture)}}\vspace{2pt}
\hrule height 0.3pt\vspace{2pt}
{\scriptsize\noindent\textbf{[Moitié inférieure du \texttt{f. 177r} — Intitulé et \italique{Hypothesis} des \italique{Ecclesiazusae} (sans liste des personnages) :]}\\
\centerline{\textbf{$\ast\;\cdot\;\ast$ ἐκκλησιαζουσῶν ἀριστοφάν(ους) :$\sim$}}
\noindent αἱ γυναῖκες ἐν σκίροις περὶ τῶν κατὰ τὴν πόλιν ἐκκλησιάσασαι, ἔδοξαν ἀνδρῶν εἶναι, καὶ σχηματίσασαι ἑαυτὰς καὶ παραδοῦναι σφίσι τὴν πόλιν, δημηγορήσασαί τε μία ἐξ αὐτῶν $\cdot$ οἱ δὲ μηχαναῖς τοῦ δόξαι αὐτὰς ἄνδρας εἶναι, τοιαῦται $\cdot$ πώγωνας περιθέτοις $\cdot$ καὶ ἀνδρείαν ἀναλαμβάνουσι στολήν $\cdot$ προασκήσασαι καὶ προνοήσασαι $\cdot$ τὸ σῶμα αὑτῶν ὡς ὅτι μάλιστα ἀνδρικὸν εἶναι δόξαι $\cdot$ μία δὲ ἐξ αὐτῶν, πραξαγόρα μὲν ὀνόματι οὖσα προσέρχεται $\cdot$ κατὰ τὰς συνθήκας καὶ φησίν $\cdot$ ὦ λαμπρὸν ὄμμα :$\sim$}
}
\vspace{3pt}\noindent\rule{4.5cm}{0.4pt}\par\vspace{1.5pt}
{\fontsize{7.0pt}{8.3pt}\selectfont
\noindent\textbf{Loci paralleli (vv. 1211–1231) ---} \textbf{1214} Suidas δ 892 διέβαλέ μ' ἡ γραῦς ; \Eustathe 1406.19 \espTest \textbf{1217} Suidas π 1502 \espTest \textbf{1231} Subscriptio \PsiUn \ fr. 6 ΑΡΙΣΤΟΦΑΝ$]$ΟΥΣ ΘΕΣΜΟΦΟΡΙ$]$ΑΖΥΣΑΙ $]$ΧΗΗΙ ($=$ 1201, versuum numerus)\par
}
\vspace{2.5pt}\noindent\rule{4.5cm}{0.4pt}\par\vspace{1.5pt}
{\fontsize{6.9pt}{8.2pt}\selectfont
\noindent\textbf{Suite (en commentaire) de l'apparat critique \& notes paléographiques (\texttt{f. 177r}) ---} \textbf{Témoin papyrologique :} \textbf{\PsiUn{} (\italique{PSI} XI 1194, fr. 6)} conserve après le \textbf{v. 1231} la \italique{subscriptio} stichométrique antique en chiffres acrophoniques attiques : ΑΡΙΣΤΟΦΑΝ]ΟΥΣ ΘΕΣΜΟΦΟΡΙ]ΑΖΥΣΑΙ ]ΧΗΗΙ (= \textbf{1201 vers}). --- \textbf{Harmonisation de συμβήνη au v. 1215 par \textbf{Bibi} :} Là où R avait conservé συμβήνην avec $-\mu-$ au v. 1197 mais écrit συβήνη sans $-\mu-$ au v. 1215, \textbf{\textbf{Bibi} rétablit συμβήνη} au v. 1215 pour maintenir le jeu de mots avec $\text{συμβινεῖν / καταβηνῆσι}$ (cf. Hésychius $\sigma\ 2137$). --- \textbf{Notes marginales de locuteur au v. 1219 :} R porte Σκ. en marge devant le v. 1219 et Χορ. devant καὶ γέρον à l'intérieur du vers. --- \textbf{Faute récurrente au v. 1227 :} Comme au v. 1127, $\text{R}^1$ écrit πέπυσται au lieu de \lecon{πέπαισται} (« nous avons assez joué », conservé par $\Sigma^{\text{R}}$). --- \textbf{Haplographie syllabique au dernier vers (v. 1231) :} $\text{R}^1$ écrit ἀνταδοῖτον (saut de la syllabe -πο-) au lieu de ἀνταποδοῖτον (Grynaeus) / \lecon{ἀνταποδοίτην} (optatif duel adopté dans le texte par L. Hauseux). --- \textbf{v. 1214 :} \lecon{διέβαλλε} R (διέβαλέ $\Sigma$, Souda $\delta\ 892$, Eustathe p. 1406.19) \espApp ὁ γραῦς R : \lecon{ἡ γραῦς} Souda, Eustathe \espApp ὡς τάχιστα R : \lecon{ὠς τάκιστα} Brunck (adopté par L. Hauseux). --- \textbf{v. 1215 :} συβήνη R : \lecon{συμβήνη} \lecon{Bibi} (cf. v. 1197, Hésychius $\sigma\ 2137$). --- \textbf{v. 1219 :} Σκ. ante versum in mg. et Χορ. ante καὶ γέρον in R \espApp γέρον R : \lecon{γέρων} \Mon{}. --- \textbf{v. 1227 :} πέπυσται R : \lecon{πέπαισται} $\Sigma$.\par
}
\end{minipage}
\end{adjustbox}
\end{adjustwidth}
\newpage
\enTeteD{Codex Ravennas 429 (Biblioteca Classense) --- Fac-similé du f. 177r (vv. 1211–1231)}
\begin{adjustwidth}{-1.5cm}{-3.0cm}
\begin{center}
\includegraphics[width=16.5cm,height=20.5cm,keepaspectratio]{/workspaces/Thesmophories/dossier_ravennas_thesmophories/folios_individuels_ff162v-177r/30_f177r_classense179r_vv1211-1231_fin.jpg}\par
\vspace{2pt}
{\scriptsize\textbf{Codex Ravennas 429, f. 177r (Classense 179r)} --- ΕΞΟΔΟΣ (Τέλος), \textbf{vv. 1211–1231} ($\text{R}^1, \Rpc, \Rdeux$ (\Mon)).}
\end{center}
\end{adjustwidth}
\newpage
\clearpage\ifodd\value{page}\hbox{}\newpage\fi% --- FOLIO 30 (f. 177r) : SPREAD 2.1 (Texte grec établi + Traduction française + Commentaire 1/2 & 2/2) ---
\clearpage\ifodd\value{page}\hbox{}\newpage\fi
\enTeteG{ΘΕΣΜΟΦΟΡΙΑΖΟΥΣΑΙ --- Texte grec établi (f. 177r, vv. 1211–1231)}
\enTeteD{LES THESMOPHORIES --- Traduction française (f. 177r, vv. 1211–1231)}
\begin{adjustwidth}{-2cm}{-2cm}
\noindent\textbf{\footnotesize Texte grec établi --- Folio 30 (\texttt{f. 177r}, vv. 1211–1231) \hfill ΕΞΟΔΟΣ (Τέλος)}\\[-2pt]
\noindent\rule{16.5cm}{0.5pt}\\[2pt]
{\fontsize{7.6pt}{8.7pt}\selectfont
\ligneEst{}{κοὐ δύσκολ᾽ ἀλλὰ πρᾶο. ποῦ τὸ γρᾴδιο;}{1211}
\ligneEst{}{οἴμ᾽ ὠς ἀπόλωλο· ποῦ τὸ γέροντ᾽ ἐντευτενί;}{}
\ligneEst{}{ὦ γρᾴδι᾽, ὦ γρᾶ᾽. οὐκ ἐπαινῶ γρᾴδιο.}{}
\ligneEst{}{Ἀρταμουξία.}{}
\ligneEst{}{διέβαλλέ μ᾽ ἡ γραῦς. ἀπότρεκ᾽ ὠς τάκιστα σύ·}{}
\ligneEst{}{ὀρτῶς δὲ συμβήνη 'στί· καταβηνῆσι γάρ.}{1215}
\ligneEst{}{οἴμοι, τί δρᾶσι; ποῖ τὸ γρᾴδιο;}{}
\ligneEst{}{Ἀρταμουξία.}{}
\ligneEst{\textbf{Κορυφαῖος}}{τὴν γραῦν ἐρωτᾷς, ἣ 'φερεν τὰς πηκτίδας;}{}
\ligneEst{\textbf{Τοξότης}}{ναὶ ναῖκι. εἶδες αὐτό;}{}
\ligneEst{\textbf{Κορυφαῖος}}{\decalage{ναὶ ναῖκι. εἶδες αὐτό;}ταύτῃ γ᾽ οἴχεται}{}
\ligneEst{}{αὐτή τ᾽ ἐκείνη καὶ γέρων τις εἵπετο.}{}
\ligneEst{\textbf{Τοξότης}}{κροκῶτ᾽ ἔκοντο τὴ γέροντο;}{1220}
\ligneEst{\textbf{Κορυφαῖος}}{\decalage{κροκῶτ᾽ ἔκοντο τὴ γέροντο;}φήμ᾽ ἐγώ.}{}
\ligneEst{}{ἔτ᾽ ἂν καταλάβοις, εἰ διώκοις ταυτῃί.}{}
\ligneEst{\textbf{Τοξότης}}{ὦ μιαρὸ γρᾷο· πότερα τρέξι τὴν ὀδό;}{}
\ligneEst{}{Ἀρταμουξία.}{}
\ligneEst{\textbf{Κορυφαῖος}}{ὀρθὴν ἄνω δίωκε. ποῖ θεῖς; οὐ πάλιν}{}
\ligneEst{}{τῃδὶ διώξει; ᾽ς τοὔμπαλιν τρέχεις σύ γε.}{}
\ligneEst{\textbf{Τοξότης}}{κακόδαιμον, ἀλλὰ τρέξι. Ἀρταμουξία.}{1225}
\ligneEst{\textbf{Κορυφαῖος} {\tiny\textit{anap.}}}{τρέχε νυν κατ᾽ αὐτοὺς ἐς κόρακας ἐπουρίσας.}{}
\ligneEst{}{ἀλλὰ πέπαισται μετρίως ἡμῖν·}{}
\ligneEst{}{ὥσθ᾽ ὥρα δή 'στι βαδίζειν}{}
\ligneEst{}{οἴκαδ᾽ ἑκάστῃ. τὼ Θεσμοφόρω δ᾽}{}
\ligneEst{}{ἡμῖν ἀγαθὴν}{1230}
\ligneEst{}{τούτων χάριν ἀνταποδοίτην.}{1231}
}
\vspace{2pt}
\noindent\rule{16.5cm}{0.6pt}\\[2pt]
\noindent\textbf{\small Commentaire philologique, métrique et scénique (vv. 1211–1231 --- partie 1/2)}\\[2pt]
\begin{adjustbox}{max totalheight=9.8cm, max width=16.5cm, keepaspectratio}%
\parbox[t]{16.5cm}{
\setlength{\columnsep}{9pt}%
\begin{multicols}{2}
\fontsize{8.4pt}{9.7pt}\selectfont
\setlength{\parindent}{0pt}%
\setlength{\parskip}{1.2pt}%
\noindent\textbf{1216} \textbf{οἴμοι, τί δράσει; ποῖ τὸ γράδιο·} Il manque un pied à ce vers.
\end{multicols}
}%
\end{adjustbox}%
\vfill
\end{adjustwidth}
\newpage
% --- PAGE DE DROITE : Traduction française + Commentaire (2/2) ---
\begin{adjustwidth}{-2cm}{-2cm}
\noindent\textbf{\footnotesize Traduction française --- Folio 30 (\texttt{f. 177r}, vv. 1211–1231) \hfill ΕΞΟΔΟΣ (Τέλος)}\\[-2pt]
\noindent\rule{16.5cm}{0.5pt}\\[2pt]
{\fontsize{7.6pt}{8.7pt}\selectfont
\ligneEst{}{\francais{et pas méchante... gentille! Où est la vieille ?}}{1211}
\ligneEst{}{\francais{Ah ! je suis perdu! Où s'est éclipsé le vieillard ?}}{}
\ligneEst{}{\francais{Ah! petite vieille, petite vieille ! Ce n'est pas bien, petite vieille.}}{}
\ligneEst{}{\francais{(Appelant.) Artamouxia !}}{}
\ligneEst{}{\francais{Elle m'a trompé, la vieille. (La danseuse s'esquive.) Va-t'en au plus vite, toi.}}{}
\ligneEst{}{\francais{(Reprenant sa lance :) La lance est droite, je descends.}}{1215}
\ligneEst{}{\francais{Ah ! que vais-je faire ? Où est la petite vieille ?}}{}
\ligneEst{}{\francais{Artamouxia !}}{}
\ligneEst{\textbf{Coryphée}}{\francais{Tu demandes après la vieille qui portait les pectis ?}}{}
\ligneEst{\textbf{Archer}}{\francais{Si fait, si. L'as-tu vue ?}}{}
\ligneEst{\textbf{Coryphée}}{\francais{\espace{3.54cm}C'est par là qu'elle a disparu,}}{}
\ligneEst{}{\francais{elle et un vieux qui la suivait.}}{}
\ligneEst{\textbf{Archer}}{\francais{Il avait une crocote, le vieux ?}}{1220}
\ligneEst{\textbf{Coryphée}}{\francais{\espace{3.9cm}Tout juste.}}{}
\ligneEst{}{\francais{Tu peux encore les rattraper en tirant de ce côté. (Indiquant une autre direction.)}}{}
\ligneEst{\textbf{Archer}}{\francais{Ô la vieille canaille! Par quel chemin prendre ?}}{}
\ligneEst{}{\francais{Artamouxia ! (ll court vers l'entrée de l'orchestre.)}}{}
\ligneEst{\textbf{Coryphée}}{\francais{Monte tout droit. Où cours-tu ?}}{}
\ligneEst{}{\francais{Veux-tu bien revenir ici ? Tu cours à contresens.}}{}
\ligneEst{\textbf{Archer}}{\francais{Malheureux !... Mais je me sauve. Artamouxia !}}{1225}
\ligneEst{\textbf{Coryphée}}{\francais{Cours maintenant, cours aux corbeaux, et qu'un bon vent te seconde.}}{}
\ligneEst{}{\francais{Mais nous nous sommes assez ébattues.}}{}
\ligneEst{}{\francais{Il est temps pour chacune de rentrer au logis.}}{}
\ligneEst{}{\francais{Puissent, en retour, les Thesmophores}}{}
\ligneEst{}{\francais{nous accorder}}{1230}
\ligneEst{}{\francais{leur précieuse faveur !}}{1231}
}
\vspace{2pt}
\noindent\rule{16.5cm}{0.6pt}\\[2pt]
\noindent\textbf{\small Commentaire philologique, métrique et scénique (vv. 1211–1231 --- partie 2/2)}\\[2pt]
\begin{adjustbox}{max totalheight=9.8cm, max width=16.5cm, keepaspectratio}%
\parbox[t]{16.5cm}{
\setlength{\columnsep}{9pt}%
\begin{multicols}{2}
\fontsize{8.4pt}{9.7pt}\selectfont
\setlength{\parindent}{0pt}%
\setlength{\parskip}{1.2pt}%
\noindent\textbf{1129-31} \textbf{τὼ Θεσμοφόρω δ᾽ / ἡμῖν ἀγαθὴν / τούτων χάριν ἀνταποδοίτην.} Selon son habitude -- qui était peut-être d'usage dans l'Ancienne comédie --, Artistophane conclue sa pièce par des souhaits de succès (il vise bien-sûr la victoire au concours comique).
\end{multicols}
}%
\end{adjustbox}%
\vfill
\end{adjustwidth}
\newpage

\end{document}
